% Auto-generated from data/segments.json + layout.json (+ transforms) — do not edit by hand.
% volume: 18Khu01
% mode: reading
% segments: 6563
% transform_rules: 0
% toc_marks: 610

% Volume layout defaults (layout.json ``layout``).
\csromanlayoutapply{1}{1.5}{21.6pt}{5pt}{6.3pt}{65pt}{2.5em}%

% Reading physical-page layout (layout.json ``page_layout_reading_mode``).
\csromanreadingpagelayoutvolume{1}{1.5}{21.6pt}{5pt}{6.3pt}{65pt}{2.5em}%
\csromanreadingpagelayoutenable

\csromanheader{2}{\textbf{ขุทฺทกนิกาย}}
\csromanmatikaanchor{mtk.0}
\csromantocmarknopage{chapter}{ขุทฺทกปาฐปาฬิ}{mtk.0}
\bookmark[dest={mtk.0},level=0]{ขุทฺทกปาฐปาฬิ}
\gambhira{\textbf{ขุทฺทกปาฐปาฬิ}}
\namakkaram{นโม ตสฺส ภควโต อรหโต สมฺมาสมฺพุทฺธสฺส.}
\csromanmatikaanchor{mtk.1}
\csromantocmarkref{section}{1. สรณตฺตย}{mtk.1}
\bookmark[dest={mtk.1},level=1]{1. สรณตฺตย}
\csromanheader{1}{1. สรณตฺตย}
\prose{\csromanfolio{1}พุทฺธํ สรณํ คจฺฉามิ.}

\prose{ธมฺมํ สรณํ คจฺฉามิ.}

\prose{สํฆํ สรณํ คจฺฉามิ.}

\vspace{\tipitakaparskip}
\csromancenter{ทุติยมฺปิ พุทฺธํ สรณํ คจฺฉามิ.}
\csromancenter{ทุติยมฺปิ ธมฺมํ สรณํ คจฺฉามิ.}
\csromancenter{ทุติยมฺปิ สํฆํ สรณํ คจฺฉามิ.}
\csromancenter{ตติยมฺปิ พุทฺธํ สรณํ คจฺฉามิ.}
\csromancenter{ตติยมฺปิ ธมฺมํ สรณํ คจฺฉามิ.}
\csromancenter{ตติยมฺปิ สํฆํ สรณํ คจฺฉามิ.\csromanspacer{} สรณตฺตยํ\footnote{สรณคมนํ นิฏฺฐิตํ. (สฺยา) 2. เวรมณีสิกฺขาปทํ. (สี, สฺยา)}.}
\csromanmatikaanchor{mtk.2}
\csromantocmarkref{section}{2. ทสสิกฺขาปท}{mtk.2}
\bookmark[dest={mtk.2},level=1]{2. ทสสิกฺขาปท}
\csromanheader{1}{2. ทสสิกฺขาปท}
\proseitem{1}{ปาณาติปาตา เวรมณิสิกฺขาปทํ2 สมาทิยามิ.}

\proseitem{2}{อทินฺนาทานา เวรมณิสิกฺขาปทํ สมาทิยามิ.}

\proseitem{3}{\csromanfolio{2}อพฺรหฺมจริยา เวรมณิสิกฺขาปทํ สมาทิยามิ.}

\proseitem{4}{มุสาวาทา เวรมณิสิกฺขาปทํ สมาทิยามิ.}

\proseitem{5}{สุราเมรยมชฺชปมาทฏฺฐานา เวรมณิสิกฺขาปทํ สมาทิยามิ.}

\proseitem{6}{วิกาลโภชนา เวรมณิสิกฺขาปทํ สมาทิยามิ.}

\proseitem{7}{นจฺจคีต วาทิต วิสูก ทสฺสนา เวรมณิสิกฺขาปทํ สมาทิยามิ.}

\csromanheader{2}{8. มาลา คนฺธ วิเลปน ธารณ มณฺฑน วิภูสนฏฺฐานา}
\prose{เวรมณิสิกฺขาปทํ สมาทิยามิ.}

\proseitem{9}{อุจฺจาสยน มหาสยนา เวรมณิสิกฺขาปทํ สมาทิยามิ.}

\proseitem{10}{ชาตรูป รชตปฏิคฺคหณา เวรมณิสิกฺขาปทํ สมาทิยามิ.\csromanspacer{} ทสสิกฺขาปทํ\footnote{ทสสิกฺขาปทํ นิฏฺฐิตํ. (สฺยา)}.}

\csromanmatikaanchor{mtk.3}
\csromantocmarkref{section}{3. ทฺวตฺติํสาการ}{mtk.3}
\bookmark[dest={mtk.3},level=1]{3. ทฺวตฺติํสาการ}
\csromanheader{1}{3. ทฺวตฺติํสาการ}
\csromanheader{2}{อตฺถิ อิมสฺมิํ กาเย\csromandash{}}
\prose{เกสา โลมา นขา ทนฺตา ตโจ.}

\prose{มํสํ นฺหารุ\footnote{นหารุ (สี, อิ), นหารู (สฺยา, กํ) 3. อฏฺฐี (สฺยา, กํ) 4. ( ) สพฺพตฺถ นตฺถิ, อฏฺฐกถายํ ววตฺถานวิธิ จ ทฺวตฺติํสสงฺขฺยา จ มนสิ กาตพฺพา.} อฏฺฐิ3 อฏฺฐิมิญฺชํ วกฺกํ.}

\prose{หทยํ ยกนํ กิโลมกํ ปิหกํ ปปฺผาสํ.}

\csromanheader{2}{อนฺตํ อนฺตคุณํ อุทริยํ กรีสํ (มตฺถลุงฺคํ).4}
\prose{ปิตฺตํ เสมฺหํ ปุพฺโพ โลหิตํ เสโท เมโท.}

\prose{อสฺสุ วสา เขโฬ สิงฺฆาณิกา ลสิกา มุตฺตนฺติ.\footnote{มุตฺตํ, มตฺถเก มตฺถลุงฺคนฺติ (สฺยา)} ทฺวตฺติํสาการํ.}

\csromanmatikaanchor{mtk.4}
\csromanmatikaanchor{mtk.5}
\csromantocmarkref{section}{4. กุมารปญฺหา}{mtk.4}
\bookmark[dest={mtk.4},level=1]{4. กุมารปญฺหา}
\csromantocmarkref{section}{5. มงฺคลสุตฺต}{mtk.5}
\bookmark[dest={mtk.5},level=1]{5. มงฺคลสุตฺต}
\csromanheader{1}{5. มงฺคลสุตฺต \textbf{4. กุมารปญฺหา}}
\proseitem{1}{\csromanfolio{3}เอกํ นาม กิํ? สพฺเพ สตฺตา อาหารฏฺฐิติกา.}

\proseitem{2}{เทฺว นาม กิํ? นามญฺจ รูปญฺจ.}

\proseitem{3}{ตีณิ นาม กิํ? ติสฺโส เวทนา.}

\proseitem{4}{จตฺตาริ นาม กิํ? จตฺตาริ อริยสจฺจานิ.}

\proseitem{5}{ปญฺจ นาม กิํ? ปญฺจุปาทานกฺขนฺธา.}

\proseitem{6}{ฉ นาม กิํ? ฉ อชฺฌตฺติกานิ อายตนานิ.}

\proseitem{7}{สตฺต นาม กิํ? สตฺต โพชฺฌงฺคา.}

\proseitem{8}{อฏฺฐ นาม กิํ? อริโย อฏฺฐงฺคิโก มคฺโค.}

\proseitem{9}{นว นาม กิํ? นว สตฺตาวาสา.}

\proseitem{10}{ทส นาม กิํ? ทสหงฺเคหิ สมนฺนาคโต “อรหา”ติ วุจฺจตีติ.\csromanspacer{} กุมารปญฺหา.}

\csromanheader{2}{5. มงฺคลสุตฺต}
\proseitem{1}{\csromansymbolfootnote{*}{อุปริ 318 ปิฏฺเฐปิ.}เอวํ เม สุตํ\csromandash{}เอกํ สมยํ ภควา สาวตฺถิยํ วิหรติ เชตวเน อนาถปิณฺฑิกสฺส อาราเม.\csromanspacer{} อถ โข อญฺญตรา เทวตา อภิกฺกนฺตาย รตฺติยา อภิกฺกนฺตวณฺณา เกวลกปฺปํ เชตวนํ โอภาเสตฺวา เยน ภควา เตนุปสงฺกมิ, อุปสงฺกมิตฺวา ภควนฺตํ อภิวาเทตฺวา เอกมนฺตํ อฏฺฐาสิ, เอกมนฺตํ ฐิตา โข สา เทวตา ภควนฺตํ คาถาย อชฺฌภาสิ\csromandash{}}

\proseitem{2}{พหู เทวา มนุสฺสา จ, มงฺคลานิ อจินฺตยุํ.}

\csromangathameasure{อากงฺขมานา โสตฺถานํ, พฺรูหิ มงฺคลมุตฺตมํ.}
\csromangathagroup{\csromangathastanza{อากงฺขมานา โสตฺถานํ, พฺรูหิ มงฺคลมุตฺตมํ.}}

\proseitem{3}{อเสวนา จ พาลานํ, ปณฺฑิตานญฺจ เสวนา.}

\csromangathameasure{ปูชา จ ปูชเนยฺยานํ, เอตํ มงฺคลมุตฺตมํ.}
\csromangathagroup{\csromangathastanza{ปูชา จ ปูชเนยฺยานํ\footnote{ปูชนียานํ (สี, สฺยา, กํ, อิ)}, เอตํ มงฺคลมุตฺตมํ.}}

\csromanmatikaanchor{mtk.6}
\csromantocmarkref{section}{6. รตนสุตฺต}{mtk.6}
\bookmark[dest={mtk.6},level=1]{6. รตนสุตฺต}
\proseitem{4}{\csromanfolio{4}ปติรูปเทสวาโส จ, ปุพฺเพ จ กตปุญฺญตา.}

\csromangathameasure{อตฺตสมฺมาปณิธิ จ, เอตํ มงฺคลมุตฺตมํ.}
\csromangathagroup{\csromangathastanza{อตฺตสมฺมาปณิธิ\footnote{อตฺตสมฺมาปณีธี (กตฺถจิ)} จ, เอตํ มงฺคลมุตฺตมํ.}}

\proseitem{5}{พาหุสจฺจญฺจ สิปฺปญฺจ, วินโย จ สุสิกฺขิโต.}

\csromangathameasure{สุภาสิตา จ ยา วาจา, เอตํ มงฺคลมุตฺตมํ.}
\csromangathagroup{\csromangathastanza{สุภาสิตา จ ยา วาจา, เอตํ มงฺคลมุตฺตมํ.}}

\proseitem{6}{มาตาปิตุ อุปฏฺฐานํ, ปุตฺตทารสฺส สงฺคโห.}

\csromangathameasure{อนากุลา จ กมฺมนฺตา, เอตํ มงฺคลมุตฺตมํ.}
\csromangathagroup{\csromangathastanza{อนากุลา จ กมฺมนฺตา, เอตํ มงฺคลมุตฺตมํ.}}

\proseitem{7}{ทานญฺจ ธมฺมจริยา จ, ญาตกานญฺจ สงฺคโห.}

\csromangathameasure{อนวชฺชานิ กมฺมานิ, เอตํ มงฺคลมุตฺตมํ.}
\csromangathagroup{\csromangathastanza{อนวชฺชานิ กมฺมานิ, เอตํ มงฺคลมุตฺตมํ.}}

\proseitem{8}{อารตี วิรตี ปาปา, มชฺชปานา จ สํยโม.}

\csromangathameasure{อปฺปมาโท จ ธมฺเมสุ, เอตํ มงฺคลมุตฺตมํ.}
\csromangathagroup{\csromangathastanza{อปฺปมาโท จ ธมฺเมสุ, เอตํ มงฺคลมุตฺตมํ.}}

\proseitem{9}{คารโว จ นิวาโต จ, สนฺตุฏฺฐิ จ กตญฺญุตา.}

\prose{กาเลน ธมฺมสฺสวนํ\footnote{ธมฺมสฺสวณํ (ก, สี), ธมฺมสวนํ (ก, สี)}, เอตํ มงฺคลมุตฺตมํ.}

\proseitem{10}{ขนฺตี จ โสวจสฺสตา, สมณานญฺจ ทสฺสนํ.}

\csromangathameasure{กาเลน ธมฺมสากจฺฉา, เอตํ มงฺคลมุตฺตมํ.}
\csromangathagroup{\csromangathastanza{กาเลน ธมฺมสากจฺฉา, เอตํ มงฺคลมุตฺตมํ.}}

\proseitem{11}{ตโป จ พฺรหฺมจริยญฺจ, อริยสจฺจาน ทสฺสนํ.}

\csromangathameasure{นิพฺพานสจฺฉิกิริยา จ, เอตํ มงฺคลมุตฺตมํ.}
\csromangathagroup{\csromangathastanza{นิพฺพานสจฺฉิกิริยา จ, เอตํ มงฺคลมุตฺตมํ.}}

\proseitem{12}{ผุฏฺฐสฺส โลกธมฺเมหิ, จิตฺตํ ยสฺส น กมฺปติ.}

\csromangathameasure{อโสกํ วิรชํ เขมํ, เอตํ มงฺคลมุตฺตมํ.}
\csromangathagroup{\csromangathastanza{อโสกํ วิรชํ เขมํ, เอตํ มงฺคลมุตฺตมํ.}}

\proseitem{13}{เอตาทิสานิ กตฺวาน, สพฺพตฺถ มปราชิตา.}

\csromangathameasure{สพฺพตฺถ โสตฺถิํ คจฺฉนฺติ, ตํ เตสํ มงฺคลมุตฺตมนฺติ. มงฺคลสุตฺตํ.}
\csromangathagroup{\csromangathastanza{สพฺพตฺถ โสตฺถิํ คจฺฉนฺติ, ตํ เตสํ มงฺคลมุตฺตมนฺติ.\csromanspacer{} มงฺคลสุตฺตํ.}}

\proseitem{1}{\csromansymbolfootnote{*}{อุปริ 312 ปิฏฺเฐปิ.}ยานีธ ภูตานิ สมาคตานิ,}

\prose{ภุมฺมานิ\footnote{ภูมานิ (ก)} วา ยานิ ว อนฺตลิกฺเข.}

\csromangathameasure{สพฺเพว ภูตา สุมนา ภวนฺตุ, \\ อโถปิ สกฺกจฺจ สุณนฺตุ ภาสิตํ. \\ ตสฺมา หิ ภูตา นิสาเมถ สพฺเพ, \\ เมตฺตํ กโรถ มานุสิยา ปชาย. \\ ทิวา จ รตฺโต จ หรนฺติ เย พลิํ, \\ ตสฺมา หิ เน รกฺขถ อปฺปมตฺตา.}
\csromangathagroup{\csromangathastanza{สพฺเพว ภูตา สุมนา ภวนฺตุ, \\ อโถปิ สกฺกจฺจ สุณนฺตุ ภาสิตํ. \\ ตสฺมา หิ ภูตา นิสาเมถ สพฺเพ, \\ เมตฺตํ กโรถ มานุสิยา ปชาย. \\ ทิวา จ รตฺโต จ หรนฺติ เย พลิํ, \\ ตสฺมา หิ เน รกฺขถ อปฺปมตฺตา.}}

\proseitem{3}{\csromanfolio{5}ยํกิญฺจิ วิตฺตํ อิธ วา หุรํ วา,}

\prose{สคฺเคสุ วา ยํ รตนํ ปณีตํ.}

\csromangathameasure{น โน สมํ อตฺถิ ตถาคเตน, \\ อิทมฺปิ พุทฺเธ รตนํ ปณีตํ.}
\csromangathagroup{\csromangathastanza{น โน สมํ อตฺถิ ตถาคเตน, \\ อิทมฺปิ พุทฺเธ รตนํ ปณีตํ.}}

\csromanhangingbegin
\hangingheaditem{4}{ขยํ วิราคํ อมตํ ปณีตํ,}
\hangingbody{ยทชฺฌคา สกฺยมุนี สมาหิโต.}
\hangingbody{น เตน ธมฺเมน สมตฺถิ กิญฺจิ,}
\hangingbody{อิทมฺปิ ธมฺเม รตนํ ปณีตํ.}
\hangingbody{เอเตน สจฺเจน สุวตฺติ โหตุ.}
\csromanhangingend

\csromangathameasure{ยํ พุทฺธเสฏฺโฐ ปริวณฺณยี สุจิํ, \\ สมาธิมานนฺตริกญฺญมาหุ. \\ สมาธินา เตน สโม น วิชฺชติ, \\ อิทมฺปิ ธมฺเม รตนํ ปณีตํ.}
\csromangathagroup{\csromangathastanza{ยํ พุทฺธเสฏฺโฐ ปริวณฺณยี สุจิํ, \\ สมาธิมานนฺตริกญฺญมาหุ. \\ สมาธินา เตน สโม น วิชฺชติ, \\ อิทมฺปิ ธมฺเม รตนํ ปณีตํ.}}

\proseitem{6}{\csromansymbolfootnote{*}{ขุ 2. 62 ปิฏฺเฐปิ.}เย ปุคฺคลา อฏฺฐ สตํ ปสตฺถา,}

\prose{จตฺตาริ เอตานิ ยุคานิ โหนฺติ.}

\csromangathameasure{เต ทกฺขิเณยฺยา สุคตสฺส สาวกา, \\ เอเตสุ ทินฺนานิ มหปฺผลานิ. \\ เย สุปฺปยุตฺตา มนสา ทเฬฺหน, \\ นิกฺกามิโน โคตมสาสนมฺหิ. \\ เต ปตฺติปตฺตา อมตํ วิคยฺห, \\ ลทฺธา มุธา นิพฺพุติํ ภุญฺชมานา. \\ เอเตน สจฺเจน สุวตฺติ โหตุ.}
\csromangathagroup{\csromangathastanza{เต ทกฺขิเณยฺยา สุคตสฺส สาวกา, \\ เอเตสุ ทินฺนานิ มหปฺผลานิ. \\ เย สุปฺปยุตฺตา มนสา ทเฬฺหน, \\ นิกฺกามิโน โคตมสาสนมฺหิ.}\csromangathastanza{เต ปตฺติปตฺตา อมตํ วิคยฺห, \\ ลทฺธา มุธา นิพฺพุติํ\footnote{นิพฺพุติ (ก)} ภุญฺชมานา. \\ เอเตน สจฺเจน สุวตฺติ โหตุ.}}

\proseitem{8}{\csromanfolio{6}\csromansymbolfootnote{*}{ขุ 10. 144-5 ปิฏฺเฐสุปิ.}ยถินฺทขีโล ปถวิสฺสิโต\footnote{ปฐวิสฺสิโต (ก, สี), ปถวิํสิโต (ก, สี, สฺยา, กํ, อิ)} สิยา,}

\prose{จตุพฺภิ วาเตหิ อสมฺปกมฺปิโย.}

\csromangathameasure{ตถูปมํ สปฺปุริสํ วทามิ, \\ โย อริยสจฺจานิ อเวจฺจ ปสฺสติ.}
\csromangathagroup{\csromangathastanza{ตถูปมํ สปฺปุริสํ วทามิ, \\ โย อริยสจฺจานิ อเวจฺจ ปสฺสติ.}}

\proseitem{9}{\csromansymbolmark{*}เย อริยสจฺจานิ วิภาวยนฺติ,}

\prose{คมฺภีรปญฺเญน สุเทสิตานิ.}

\csromangathameasure{กิญฺจาปิ เต โหนฺติ ภุสํ ปมตฺตา, \\ น เต ภวํ อฏฺฐมมาทิยนฺติ.}
\csromangathagroup{\csromangathastanza{กิญฺจาปิ เต โหนฺติ ภุสํ ปมตฺตา, \\ น เต ภวํ อฏฺฐมมาทิยนฺติ.}}

\proseitem{10}{สหาว’สฺส ทสฺสนสมฺปทาย\footnote{สหาวสทฺทสฺสนสมฺปทาย (ก)},}

\prose{ตย’สฺสุ ธมฺมา ชหิตา ภวนฺติ.}

\csromangathameasure{สกฺกายทิฏฺฐี วิจิกิจฺฉิตญฺจ, \\ สีลพฺพตํ วาปิ ยทตฺถิ กิญฺจิ.}
\csromangathagroup{\csromangathastanza{สกฺกายทิฏฺฐี วิจิกิจฺฉิตญฺจ, \\ สีลพฺพตํ วาปิ ยทตฺถิ กิญฺจิ.}}

\proseitem{11}{จตูหปาเยหิ จ วิปฺปมุตฺโต,}

\prose{ฉจฺจาภิฐานานิ\footnote{ฉ จาภิฐานานิ (สี, สฺยา)} อภพฺพ กาตุํ\footnote{อภพฺโพ กาตุํ (สี)}.}

\proseitem{12}{กิญฺจาปิ โส กมฺม\footnote{กมฺมํ (สี, สฺยา, กํ, อิ)} กโรติ ปาปกํ,}

\prose{กาเยน วาจา อุท เจตสา วา.}

\csromangathameasure{อภพฺพ โส ตสฺส ปฏิจฺฉทาย, \\ อภพฺพตา ทิฏฺฐปทสฺส วุตฺตา. \\ วนปฺปคุมฺเพ ยถ ผุสฺสิตคฺเค, \\ คิมฺหานมาเส ปฐมสฺมิํ คิเมฺห. \\ ตถูปมํ ธมฺมวรํ อเทสยิ, \\ นิพฺพานคามิํ ปรมํ หิตาย. \\ อิทมฺปิ พุทฺเธ รตนํ ปณีตํ, \\ วโร วรญฺญู วรโท วราหโร, \\ อนุตฺตโร ธมฺมวรํ อเทสยิ. \\ อิทมฺปิ พุทฺเธ รตนํ ปณีตํ, \\ ขีณํ ปุราณํ นว นตฺถิสมฺภวํ, \\ วิรตฺตจิตฺตายติเก ภวสฺมิํ. \\ เต ขีณพีชา อวิรูฬฺหิฉนฺทา, \\ นิพฺพนฺติ ธีรา ยถายํ ปทีโป. \\ ยานีธ ภูตานิ สมาคตานิ, \\ พุทฺธํ นมสฺสาม สุวตฺถิ โหตุ. \\ ยานีธ ภูตานิ สมาคตานิ, \\ ธมฺมํ นมสฺสาม สุวตฺถิ โหตุ.}
\csromangathagroup{\csromangathastanza{อภพฺพ\footnote{อภพฺโพ (พหูสุ)} โส ตสฺส ปฏิจฺฉทาย\footnote{ปฏิจฺฉาทาย (สี)}, \\ อภพฺพตา ทิฏฺฐปทสฺส วุตฺตา. \\ วนปฺปคุมฺเพ ยถ\footnote{อภพฺโพ (พหูสุ)} ผุสฺสิตคฺเค, \\ คิมฺหานมาเส ปฐมสฺมิํ\footnote{อภพฺโพ (พหูสุ)} คิเมฺห.}\csromangathastanza{\csromanfolio{7}ตถูปมํ ธมฺมวรํ อเทสยิ\footnote{อเทสยี (สี)}, \\ นิพฺพานคามิํ ปรมํ หิตาย. \\ อิทมฺปิ พุทฺเธ รตนํ ปณีตํ, \\ วโร วรญฺญู วรโท วราหโร,}\csromangathastanza{อนุตฺตโร ธมฺมวรํ อเทสยิ. \\ อิทมฺปิ พุทฺเธ รตนํ ปณีตํ, \\ ขีณํ ปุราณํ นว นตฺถิสมฺภวํ, \\ วิรตฺตจิตฺตายติเก ภวสฺมิํ.}\csromangathastanza{เต ขีณพีชา อวิรูฬฺหิฉนฺทา, \\ นิพฺพนฺติ ธีรา ยถายํ\footnote{ยถยํ (ก)} ปทีโป. \\ ยานีธ ภูตานิ สมาคตานิ, \\ พุทฺธํ นมสฺสาม สุวตฺถิ โหตุ. \\ ยานีธ ภูตานิ สมาคตานิ, \\ ธมฺมํ นมสฺสาม สุวตฺถิ โหตุ.}}

\proseitem{18}{\csromanfolio{8}ยานีธ ภูตานิ สมาคตานิ,}

\prose{สํฆํ นมสฺสาม สุวตฺถิ โหตูติ.\csromanspacer{} รตนสุตฺตํ.}

\csromanheader{1}{7. ติโรกุฏฺฏสุตฺต}
\proseitem{1}{\csromansymbolfootnote{*}{ขุ 2. 129 ปิฏฺเฐ เปตวตฺถุมฺหิปิ.}ติโรกุฏฺเฏสุ ติฏฺฐนฺติ, สนฺธิสิงฺฆาฏเกสุ จ.}

\csromangathameasure{ทฺวารพาหาสุ ติฏฺฐนฺติ, อาคนฺตฺวาน สกํ ฆรํ.}
\csromangathagroup{\csromangathastanza{ทฺวารพาหาสุ ติฏฺฐนฺติ, อาคนฺตฺวาน สกํ ฆรํ.}}

\csromanhangingbegin
\hangingheaditem{2}{ปหูเต อนฺนปานมฺหิ, ขชฺชโภชฺเช อุปฏฺฐิเต.}
\hangingbody{น เตสํ โกจิ สรติ, สตฺตานํ กมฺมปจฺจยา.}
\csromanhangingend

\csromanhangingbegin
\hangingheaditem{3}{เอวํ ททนฺติ ญาตีนํ, เย โหนฺติ อนุกมฺปกา.}
\hangingbody{สุจิํ ปณีตํ กาเลน, กปฺปิยํ ปานโภชนํ.}
\hangingbody{อิทํ โว ญาตีนํ โหตุ, สุขิตา โหนฺตุ ญาตโย.}
\csromanhangingend

\csromanhangingbegin
\hangingheaditem{4}{เต จ ตตฺถ สมาคนฺตฺวา, ญาติเปตา สมาคตา.}
\hangingbody{ปหูเต อนฺนปานมฺหิ, สกฺกจฺจํ อนุโมทเร.}
\csromanhangingend

\csromanhangingbegin
\hangingheaditem{5}{จิรํ ชีวนฺตุ โน ญาตี, เยสํ เหตุ ลภามเส.}
\hangingbody{อมฺหากญฺจ กตา ปูชา, ทายกา จ อนิปฺผลา.}
\csromanhangingend

\csromanhangingbegin
\hangingheaditem{6}{น หิ ตตฺถ กสี\footnote{กสิ (ก)} อตฺถิ, โครกฺเขตฺถ น วิชฺชติ.}
\hangingbody{วณิชฺชา ตาทิสี นตฺถิ, หิรญฺเญน กยากยํ2.}
\hangingbody{อิโต ทินฺเนน ยาเปนฺติ, เปตา กาลคตา3 ตหิํ.}
\csromanhangingend

\csromanhangingbegin
\hangingheaditem{7}{อุนฺนเม อุทกํ วุฏฺฐํ, ยถา นินฺนํ ปวตฺตติ.}
\hangingbody{เอวเมวํ อิโต ทินฺนํ, เปตานํ อุปกปฺปติ.}
\csromanhangingend

\csromanhangingbegin
\hangingheaditem{8}{ยถา วาริวหา ปูรา, ปริปูเรนฺติ สาครํ.}
\hangingbody{เอวเมวํ อิโต ทินฺนํ, เปตานํ อุปกปฺปติ.}
\csromanhangingend

\csromanheader{1}{8. นิธิกณฺฑสุตฺต}
\csromanhangingbegin
\hangingheaditem{9}{\csromanfolio{9}อทาสิ เม อกาสิ เม, ญาติ มิตฺตา\footnote{ญาติ มิตฺโต (?)} สขา จ เม.}
\hangingbody{เปตานํ ทกฺขิณํ ทชฺชา, ปุพฺเพ กตมนุสฺสรํ}
\csromanhangingend

\csromangathameasure{น หิ รุณฺณํ วา โสโก วา, ยา จญฺญา ปริเทวนา. \\ น ตํ เปตานมตฺถาย, เอวํ ติฏฺฐนฺติ ญาตโย. \\ อยญฺจ โข ทกฺขิณา ทินฺนา, สํฆมฺหิ สุปฺปติฏฺฐิตา. \\ ทีฆรตฺตํ หิตายสฺส, ฐานโส อุปกปฺปติ. \\ โส ญาติ ธมฺโม จ อยํ นิทสฺสิโต, \\ เปตาน ปูชา จ กตา อุฬารา. \\ พลญฺจ ภิกฺขูนมนุปฺปทินฺนํ, \\ ตุเมฺหหิ ปุญฺญํ ปสุตํ อนปฺปกนฺติ. ติโรกุฏฺฏสุตฺตํ.}
\csromangathagroup{\csromangathastanza{น หิ รุณฺณํ วา โสโก วา, ยา จญฺญา ปริเทวนา. \\ น ตํ เปตานมตฺถาย, เอวํ ติฏฺฐนฺติ ญาตโย.}\csromangathastanza{อยญฺจ โข ทกฺขิณา ทินฺนา, สํฆมฺหิ สุปฺปติฏฺฐิตา. \\ ทีฆรตฺตํ หิตายสฺส, ฐานโส อุปกปฺปติ.}\csromangathastanza{โส ญาติ ธมฺโม จ อยํ นิทสฺสิโต, \\ เปตาน ปูชา จ กตา อุฬารา. \\ พลญฺจ ภิกฺขูนมนุปฺปทินฺนํ\footnote{...มนุปฺปทินฺนวา (ก)}, \\ ตุเมฺหหิ ปุญฺญํ ปสุตํ อนปฺปกนฺติ.\csromanspacer{} ติโรกุฏฺฏสุตฺตํ.}}

\csromanheader{2}{8. นิธิกณฺฑสุตฺต}
\csromanhangingbegin
\hangingheaditem{1}{นิธิํ นิเธติ ปุริโส, คมฺภีเร โอทกนฺติเก.}
\hangingbody{อตฺเถ กิจฺเจ สมุปฺปนฺเน, อตฺถาย เม ภวิสฺสติ.}
\csromanhangingend

\csromanhangingbegin
\hangingheaditem{2}{ราชโต วา ทุรุตฺตสฺส, โจรโต ปีฬิตสฺส วา.}
\hangingbody{อิณสฺส วา ปโมกฺขาย, ทุพฺภิกฺเข อาปทาสุ วา.}
\hangingbody{เอตทตฺถาย โลกสฺมิํ, นิธิ นาม นิธียติ.}
\csromanhangingend

\csromanhangingbegin
\hangingheaditem{3}{ตาว’สฺสุนิหิโต\footnote{ตาว สุนิหิโต (สี)} สนฺโต, คมฺภีเร โอทกนฺติเก.}
\hangingbody{น สพฺโพ สพฺพทา เอว, ตสฺส ตํ อุปกปฺปติ.}
\csromanhangingend

\csromanhangingbegin
\hangingheaditem{4}{นิธิ วา ฐานา จวติ, สญฺญา วาสฺส วิมุยฺหติ.}
\hangingbody{นาคา วา อปนาเมนฺติ, ยกฺขา วาปิ หรนฺติ นํ.}
\csromanhangingend

\csromanhangingbegin
\hangingheaditem{5}{อปฺปิยา วาปิ ทายาทา, อุทฺธรนฺติ อปสฺสโต.}
\hangingbody{ยทา ปุญฺญกฺขโย โหติ, สพฺพเมตํ วินสฺสติ.}
\csromanhangingend

\csromanhangingbegin
\hangingheaditem{6}{ยสฺส ทาเนน สีเลน, สํยเมน ทเมน จ.}
\hangingbody{นิธี สุนิหิโต โหติ, ปุริสสฺส วา.}
\csromanhangingend

\csromanhangingbegin
\hangingheaditem{7}{\csromanfolio{10}เจติยมฺหิ จ สํเฆ วา, ปุคฺคเล อติถีสุ วา.}
\hangingbody{มาตริ ปิตริ จาปิ1, อโถ เชฏฺฐมฺหิ ภาตริ.}
\csromanhangingend

\csromanhangingbegin
\hangingheaditem{8}{เอโส นิธิ สุนิหิโต, อเชยฺโย อนุคามิโก.}
\hangingbody{ปหาย คมนีเยสุ, เอตํ อาทาย คจฺฉติ.}
\csromanhangingend

\csromanhangingbegin
\hangingheaditem{9}{อสาธารณมญฺเญสํ, อโจราหรโณ นิธิ.}
\hangingbody{กยิราถ ธีโร ปุญฺญานิ, โย นิธิ อนุคามิโก.}
\csromanhangingend

\proseitem{10}{เอส เทวมนุสฺสานํ, สพฺพกามทโท นิธิ.}

\csromangathameasure{ยํ ยเทวาภิปตฺเถนฺติ, สพฺพเมเตน ลพฺภติ.}
\csromangathagroup{\csromangathastanza{ยํ ยเทวาภิปตฺเถนฺติ, สพฺพเมเตน ลพฺภติ.}}

\proseitem{11}{สุวณฺณตา สุสรตา, สุสณฺฐานา สุรูปตา\footnote{สุสณฺฐานสุรูปตา (สี), สุสณฺฐานํ สุรูปตา (สฺยา, กํ)}.}

\csromangathameasure{อาธิปจฺจปริวารา, สพฺพเมเตน ลพฺภติ.}
\csromangathagroup{\csromangathastanza{อาธิปจฺจปริวารา, สพฺพเมเตน ลพฺภติ.}}

\proseitem{12}{ปเทสรชฺชํ อิสฺสริยํ, จกฺกวตฺติสุขํ ปิยํ.}

\csromangathameasure{เทวรชฺชมฺปิ ทิพฺเพสุ, สพฺพเมเตน ลพฺภติ.}
\csromangathagroup{\csromangathastanza{เทวรชฺชมฺปิ ทิพฺเพสุ, สพฺพเมเตน ลพฺภติ.}}

\csromanmatikaanchor{mtk.7}
\csromantocmarkref{section}{7. ติโรกุฏฺฏสุตฺต}{mtk.7}
\bookmark[dest={mtk.7},level=1]{7. ติโรกุฏฺฏสุตฺต}
\proseitem{13}{มานุสฺสิกา จ สมฺปตฺติ, เทวโลเก จ ยา รติ.}

\csromangathameasure{ยา จ นิพฺพานสมฺปตฺติ, สพฺพเมเตน ลพฺภติ.}
\csromangathagroup{\csromangathastanza{ยา จ นิพฺพานสมฺปตฺติ, สพฺพเมเตน ลพฺภติ.}}

\proseitem{14}{มิตฺตสมฺปทมาคมฺม, โยนิโสว\footnote{โยนิโส เว (สี), โยนิโส เจ (สฺยา), โยนิโส จ (?)} ปยุญฺชโต.}

\csromangathameasure{วิชฺชาวิมุตฺติวสีภาโว, สพฺพเมเตน ลพฺภติ.}
\csromangathagroup{\csromangathastanza{วิชฺชาวิมุตฺติวสีภาโว, สพฺพเมเตน ลพฺภติ.}}

\proseitem{15}{ปฏิสมฺภิทา วิโมกฺขา จ, ยา จ สาวกปารมี.}

\csromangathameasure{ปจฺเจกโพธิ พุทฺธภูมิ, สพฺพเมเตน ลพฺภติ.}
\csromangathagroup{\csromangathastanza{ปจฺเจกโพธิ พุทฺธภูมิ, สพฺพเมเตน ลพฺภติ.}}

\proseitem{16}{เอวํ มหตฺถิกา เอสา, ยทิทํ ปุญฺญสมฺปทา.}

\csromangathameasure{ตสฺมา ธีรา ปสํสนฺติ, ปณฺฑิตา กตปุญฺญตนฺติ. นิธิกณฺฑสุตฺตํ.}
\csromangathagroup{\csromangathastanza{ตสฺมา ธีรา ปสํสนฺติ, ปณฺฑิตา กตปุญฺญตนฺติ.\csromanspacer{} นิธิกณฺฑสุตฺตํ.}}

\csromanheader{1}{9. เมตฺตสุตฺต}
\proseitem{1}{\csromansymbolfootnote{*}{อุปริ 300 ปิฏฺเฐปิ.}กรณียมตฺถกุสเลน, ยนฺตสนฺตํ ปทํ อภิสเมจฺจ.}

\csromanmatikaanchor{mtk.8}
\csromantocmarkref{section}{8. นิธิกณฺฑสุตฺต}{mtk.8}
\bookmark[dest={mtk.8},level=1]{8. นิธิกณฺฑสุตฺต}
\csromangathameasure{สกฺโก อุชู จ สุหุชู จ, สูวโจ จสฺส มุทุ อนติมานี.}
\csromangathagroup{\csromangathastanza{สกฺโก อุชู จ สุหุชู\footnote{สูชู (สี)} จ, สูวโจ จสฺส มุทุ อนติมานี.}}

\csromanheader{2}{9. เมตฺตสุตฺต}
\csromangathameasure{สนฺตุสฺสโก จ สุภโร จ, อปฺปกิจฺโจ จ สลฺลหุกวุตฺติ. \\ สนฺตินฺทฺริโย จ นิปโก จ, อปฺปคพฺโภ กุเลสฺวนนุคิทฺโธ. \\ น จ ขุทฺทมาจเร กิญฺจิ, เยน วิญฺญู ปเร อุปวเทยฺยุํ. \\ สุขิโน ว เขมิโน โหนฺตุ, สพฺพสตฺตา ภวนฺตุ สุขิตตฺตา. \\ เย เกจิ ปาณภูตตฺถิ, ตสา วา ถาวรา ว’นวเสสา. \\ ทีฆา วา เย ว มหนฺตา, มชฺฌิมา รสฺสกา อณุกถูลา. \\ ทิฏฺฐา วา เย ว อทิฏฺฐา, เย ว ทูเร วสนฺติ อวิทูเร. \\ ภูตา ว สมฺภเวสี ว, สพฺพสตฺตา ภวนฺตุ สุขิตตฺตา. \\ น ปโร ปรํ นิกุพฺเพถ, นาติมญฺเญถ กตฺถจิ น กญฺจิ. \\ พฺยาโรสนา ปฏิฆสญฺญ, นาญฺญมญฺญสฺส ทุกฺขมิจฺเฉยฺย. \\ มาตา ยถา นิยํปุตฺต, มายุสา เอกปุตฺตมนุรกฺเข. \\ เอวมฺปิ สพฺพภูเตสุ, มานสํ ภาวเย อปริมาณํ. \\ เมตฺตญฺจ สพฺพโลกสฺมิ, มานสํ ภาวเย อปริมาณํ. \\ อุทฺธํ อโธ จ ติริยญฺจ, อสมฺพาธํ อเวร’มสปตฺตํ.}
\csromangathagroup{\csromangathastanza{\csromanfolio{11}สนฺตุสฺสโก จ สุภโร จ, อปฺปกิจฺโจ จ สลฺลหุกวุตฺติ. \\ สนฺตินฺทฺริโย จ นิปโก จ, อปฺปคพฺโภ กุเลสฺวนนุคิทฺโธ.}\csromangathastanza{น จ ขุทฺทมาจเร กิญฺจิ, เยน วิญฺญู ปเร อุปวเทยฺยุํ. \\ สุขิโน ว เขมิโน โหนฺตุ, สพฺพสตฺตา\footnote{สพฺเพ สตฺตา (สี, สฺยา)} ภวนฺตุ สุขิตตฺตา.}\csromangathastanza{เย เกจิ ปาณภูตตฺถิ, ตสา วา ถาวรา ว’นวเสสา. \\ ทีฆา วา เย ว มหนฺตา\footnote{มหนฺต (?)}, มชฺฌิมา รสฺสกา อณุกถูลา.}\csromangathastanza{ทิฏฺฐา วา เย ว อทิฏฺฐา\footnote{อทิฏฺฐ (?)}, เย ว\footnote{เย จ (สี, สฺยา, กํ, อิ)} ทูเร วสนฺติ อวิทูเร. \\ ภูตา ว\footnote{อทิฏฺฐ (?)} สมฺภเวสี ว\footnote{เย จ (สี, สฺยา, กํ, อิ)}, สพฺพสตฺตา ภวนฺตุ สุขิตตฺตา.}\csromangathastanza{น ปโร ปรํ นิกุพฺเพถ, นาติมญฺเญถ กตฺถจิ น กญฺจิ\footnote{นํ กญฺจิ (สี, อิ), นํ กิญฺจิ (สฺยา), น กิญฺจิ (ก)}. \\ พฺยาโรสนา ปฏิฆสญฺญ, นาญฺญมญฺญสฺส ทุกฺขมิจฺเฉยฺย.}\csromangathastanza{มาตา ยถา นิยํปุตฺต, มายุสา เอกปุตฺตมนุรกฺเข. \\ เอวมฺปิ สพฺพภูเตสุ, มานสํ ภาวเย อปริมาณํ.}\csromangathastanza{เมตฺตญฺจ สพฺพโลกสฺมิ, มานสํ ภาวเย อปริมาณํ. \\ อุทฺธํ อโธ จ ติริยญฺจ, อสมฺพาธํ อเวร’มสปตฺตํ.}}

\proseitem{9}{ติฏฺฐํ จรํ นิสินฺโน ว6, สยาโน ยาวตา’สฺส วิตมิทฺโธ\footnote{วิคตมิทฺโธ (พหูสุ)},}

\csromangathameasure{เอตํ สติํ อธิฏฺเฐยฺย, พฺรหฺมเมตํ วิหารมิธมาหุ.}
\csromangathagroup{\csromangathastanza{เอตํ สติํ อธิฏฺเฐยฺย, พฺรหฺมเมตํ วิหารมิธมาหุ.}}

\proseitem{10}{ทิฏฺฐิญฺจ อนุปคฺคมฺม,สีลวา ทสฺสเนน สมฺปนฺโน.}

\csromangathameasure{กาเมสุ วินย เคธํ, น หิ ชาตุคฺคพฺภเสยฺย ปุน เรตีติ. เมตฺตสุตฺตํ.}
\csromangathagroup{\csromangathastanza{กาเมสุ วินย\footnote{วิเนยฺย (สี)} เคธํ, น หิ ชาตุคฺคพฺภเสยฺย ปุน เรตีติ.\csromanspacer{} เมตฺตสุตฺตํ.}}

\nitthitam{\textbf{ขุทฺทกปาฐปาฬิ นิฏฺฐิตา.}}
\csromanheader{2}{\textbf{ขุทฺทกนิกาย}}
\gambhira{\textbf{ธมฺมปทปาฬิ}}
\namakkaram{นโม ตสฺส ภควโต อรหโต สมฺมาสมฺพุทฺธสฺส.}
\csromanheader{1}{1. ยมกวคฺค}
\csromanheader{2}{1. จกฺขุปาลตฺเถรวตฺถุ}
\csromangathameasure{มโนปุพฺพงฺคมา ธมฺมา, มโนเสฏฺฐา มโนมยา. \\ มนสา เจ ปทุฏฺเฐน, ภาสติ วา กโรติ วา. \\ ตโต นํ ทุกฺขมเนฺวติ, จกฺกํว วหโต ปทํ.}
\csromangathagroup{\csromangathastanza{\csromanfolio{13}มโนปุพฺพงฺคมา ธมฺมา, มโนเสฏฺฐา มโนมยา. \\ มนสา เจ ปทุฏฺเฐน, ภาสติ วา กโรติ วา. \\ ตโต นํ ทุกฺขมเนฺวติ, จกฺกํว วหโต ปทํ.}}

\csromanheader{2}{2. มฏฺฐกุณฺฑลีวตฺถุ}
\csromangathameasure{*มโนปุพฺพงฺคมา ธมฺมา, มโนเสฏฺฐา มโนมยา. \\ มนสา เจ ปสนฺเนน, ภาสติ วา กโรติ วา. \\ ตโต นํ สุขมเนฺวติ, ฉายาว อนปายินี.}
\csromangathagroup{\csromangathastanza{\csromansymbolfootnote{*}{ขุ 10. 112 ปิฏฺเฐปิ.}มโนปุพฺพงฺคมา ธมฺมา, มโนเสฏฺฐา มโนมยา. \\ มนสา เจ ปสนฺเนน, ภาสติ วา กโรติ วา. \\ ตโต นํ สุขมเนฺวติ, ฉายาว อนปายินี\footnote{อนุปายินี (ก)}.}}

\csromanheader{2}{3. ติสฺสตฺเถรวตฺถุ}
\csromanhangingbegin
\hangingheaditem{3}{อกฺโกจฺฉิ มํ อวธิ มํ, อชินิ\footnote{อชินี (?)} มํ อหาสิ เม.}
\hangingbody{เย จ ตํ อุปนยฺหนฺติ, เวรํ เตสํ น สมฺมติ.}
\csromanhangingend

\csromanhangingbegin
\hangingheaditem{4}{อกฺโกจฺฉิ มํ อวธิ มํ, อชินิ มํ อหาสิ เม.}
\hangingbody{เย จ ตํ นุปนยฺหนฺติ, เวรํ เตสูปสมฺมติ.}
\csromanhangingend

\gambhira{ธมฺมปทปาฬิ \textbf{4. กาฬยกฺขินีวตฺถุ}}
\csromangathameasure{น หิ เวเรน เวรานิ, สมฺมนฺตีธ กุทาจนํ. \\ อเวเรน จ สมฺมนฺติ, เอส ธมฺโม สนนฺตโน.}
\csromangathagroup{\csromangathastanza{\csromanfolio{14}น หิ เวเรน เวรานิ, สมฺมนฺตีธ กุทาจนํ. \\ อเวเรน จ สมฺมนฺติ, เอส ธมฺโม สนนฺตโน.}}

\csromanheader{2}{5. โกสมฺพกวตฺถุ}
\proseitem{6}{\csromansymbolfootnote{*}{ขุ 2. 269 ปิฏฺเฐปิ.}ปเร จ น วิชานนฺติ, มยเมตฺถ ยมามเส.}

\csromangathameasure{เย จ ตตฺถ วิชานนฺติ, ตโต สมฺมนฺติ เมธคา.}
\csromangathagroup{\csromangathastanza{เย จ ตตฺถ วิชานนฺติ, ตโต สมฺมนฺติ เมธคา.}}

\csromanheader{2}{6. มหากาฬตฺเถรวตฺถุ}
\csromanmatikaanchor{mtk.9}
\csromantocmarkref{section}{9. เมตฺตสุตฺต}{mtk.9}
\bookmark[dest={mtk.9},level=1]{9. เมตฺตสุตฺต}
\csromangathameasure{สุภานุปสฺสิํ วิหรนฺตํ, อินฺทฺริเยสุ อสํวุตํ. \\ โภชนมฺหิ จามตฺตญฺญุํ, กุสีตํ หีนวีริยํ. \\ ตํ เว ปสหติ มาโร, วาโต รุกฺขํว ทุพฺพลํ. \\ อสุภานุปสฺสิํ วิหรนฺตํ, อินฺทฺริเยสุ สุสํวุตํ. \\ โภชนมฺหิ จ มตฺตญฺญุํ, สทฺธํ อารทฺธวีริยํ. \\ ตํ เว นปฺปสหติ มาโร, วาโต เสลํว ปพฺพตํ.}
\csromangathagroup{\csromangathastanza{สุภานุปสฺสิํ วิหรนฺตํ, อินฺทฺริเยสุ อสํวุตํ. \\ โภชนมฺหิ จามตฺตญฺญุํ, กุสีตํ หีนวีริยํ.}\csromangathastanza{ตํ เว ปสหติ มาโร, วาโต รุกฺขํว ทุพฺพลํ. \\ อสุภานุปสฺสิํ วิหรนฺตํ, อินฺทฺริเยสุ สุสํวุตํ.}\csromangathastanza{โภชนมฺหิ จ มตฺตญฺญุํ, สทฺธํ อารทฺธวีริยํ. \\ ตํ เว นปฺปสหติ มาโร, วาโต เสลํว ปพฺพตํ.}}

\csromanheader{2}{7. เทวทตฺตวตฺถุ}
\proseitem{9}{\csromansymbolfootnote{+}{ขุ 2. 342; ขุ 5. 60, 374 ปิฏฺเฐสุปิ. ++ ขุ 2. 248 ปิฏฺเฐปิ}อนิกฺกสาโว กาสาวํ, โย วตฺถํ ปริทหิสฺสติ.}

\csromangathameasure{อเปโต ทมสจฺเจน, น โส กาสาวมรหติ. \\ โย จ วนฺตกสาว’สฺส, สีเลสุ สุสมาหิโต. \\ อุเปโต ทมสจฺเจน, ส เว กาสาวมรหติ.}
\csromangathagroup{\csromangathastanza{อเปโต ทมสจฺเจน, น โส กาสาวมรหติ. \\ โย จ วนฺตกสาว’สฺส, สีเลสุ สุสมาหิโต. \\ อุเปโต ทมสจฺเจน, ส เว กาสาวมรหติ.}}

\csromanheader{2}{8. สาริปุตฺตตฺเถรวตฺถุ}
\csromangathameasure{อสาเร สารมติโน, สาเร จาสารทสฺสิโน. \\ เต สารํ นาธิคจฺฉนฺติ, มิจฺฉาสงฺกปฺปโคจรา. \\ สารญฺจ สารโต ญตฺวา, อสารญฺจ อสารโต. \\ เต สารํ อธิคจฺฉนฺติ, สมฺมาสงฺกปฺปโคจรา.}
\csromangathagroup{\csromangathastanza{อสาเร สารมติโน, สาเร จาสารทสฺสิโน. \\ เต สารํ นาธิคจฺฉนฺติ, มิจฺฉาสงฺกปฺปโคจรา.}\csromangathastanza{สารญฺจ สารโต ญตฺวา, อสารญฺจ อสารโต. \\ เต สารํ อธิคจฺฉนฺติ, สมฺมาสงฺกปฺปโคจรา.}}

\csromanheader{2}{9. นนฺทตฺเถรวตฺถุ}
\proseitem{13}{\csromansymbolmark{+}+ ยถา อคารํ ทุจฺฉนฺนํ, วุฏฺฐี สมติวิชฺฌติ.}

\csromangathameasure{เอวํ อภาวิตํ จิตฺตํ, ราโค สมติวิชฺฌติ.}
\csromangathagroup{\csromangathastanza{เอวํ อภาวิตํ จิตฺตํ, ราโค สมติวิชฺฌติ.}}

\chapterheadpageread{1. ยมกวคฺค 14. ยถา อคารํ สุฉนฺนํ, วุฏฺฐี น สมติวิชฺฌติ.}
\csromangathameasure{เอวํ สุภาวิตํ จิตฺตํ, ราโค น สมติวิชฺฌติ.}
\csromangathagroup{\csromangathastanza{\csromanfolio{15}เอวํ สุภาวิตํ จิตฺตํ, ราโค น สมติวิชฺฌติ.}}

\csromanheader{2}{10. จุนฺทสูกริกวตฺถุ}
\proseitem{15}{\csromansymbolfootnote{*}{ขุ 10. 172 ปิฏฺเฐปิ.}อิธ โสจติ เปจฺจ โสจติ.}

\prose{ปาปการี อุภยตฺถ โสจติ.}

\csromangathameasure{โส โสจติ โส วิหญฺญติ, \\ ทิสฺวา กมฺมกิลิฏฺฐมตฺตโน.}
\csromangathagroup{\csromangathastanza{โส โสจติ โส วิหญฺญติ, \\ ทิสฺวา กมฺมกิลิฏฺฐมตฺตโน.}}

\csromanmatikaanchor{mtk.10}
\csromantocmarknopage{chapter}{ธมฺมปทปาฬิ}{mtk.10}
\bookmark[dest={mtk.10},level=0]{ธมฺมปทปาฬิ}
\csromanheader{2}{11. ธมฺมิก-อุปาสกวตฺถุ}
\csromangathameasure{อิธ โมทติ เปจฺจ โมทติ, \\ กตปุญฺโญ อุภยตฺถ โมทติ. \\ โส โมทติ โส ปโมทติ, \\ ทิสฺวา กมฺมวิสุทฺธิมตฺตโน.}
\csromangathagroup{\csromangathastanza{อิธ โมทติ เปจฺจ โมทติ, \\ กตปุญฺโญ อุภยตฺถ โมทติ. \\ โส โมทติ โส ปโมทติ, \\ ทิสฺวา กมฺมวิสุทฺธิมตฺตโน.}}

\csromanmatikaanchor{mtk.11}
\csromantocmarkref{section}{1. ยมกวคฺค}{mtk.11}
\bookmark[dest={mtk.11},level=1]{1. ยมกวคฺค}
\csromanheader{2}{12. เทวทตฺตวตฺถุ}
\csromangathameasure{อิธ ตปฺปติ เปจฺจ ตปฺปติ, \\ ปาปการี อุภยตฺถ ตปฺปติ. \\ “ปาปํ เม กตนฺ”ติ ตปฺปติ, \\ ภิยฺโย ตปฺปติ ทุคฺคติํ คโต.}
\csromangathagroup{\csromangathastanza{อิธ ตปฺปติ เปจฺจ ตปฺปติ, \\ ปาปการี\footnote{ปาปการิ (?)} อุภยตฺถ ตปฺปติ. \\ “ปาปํ เม กตนฺ”ติ ตปฺปติ, \\ ภิยฺโย\footnote{ปาปการิ (?)} ตปฺปติ ทุคฺคติํ คโต.}}

\csromanheader{2}{13. สุมนเทวีวตฺถุ}
\proseitem{18}{อิธ นนฺทติ เปจฺจ นนฺทติ,}

\prose{กตปุญฺโญ อุภยตฺถ นนฺทติ.}

\csromangathameasure{“ปุญฺญํ เม กตนฺ”ติ นนฺทติ, \\ ภิยฺโย นนฺทติ สุคฺคติํ คโต.}
\csromangathagroup{\csromangathastanza{“ปุญฺญํ เม กตนฺ”ติ นนฺทติ, \\ ภิยฺโย นนฺทติ สุคฺคติํ คโต.}}

\csromanheader{2}{14. เทฺวสหายกภิกฺขุวตฺถุ}
\proseitem{19}{พหุมฺปิ เจ สํหิต\footnote{สหิตํ (สี, สฺยา, กํ, อิ)} ภาสมาโน,}

\prose{น ตกฺกโร โหติ นโร ปมตฺโต.}

\gambhira{ธมฺมปทปาฬิ}
\csromangathameasure{โคโปว คาโว คณยํ ปเรสํ, \\ น ภาควา สามญฺญสฺส โหติ. \\ อปฺปมฺปิ เจ สํหิต ภาสมาโน, \\ ธมฺมสฺส โหติ อนุธมฺมจารี. \\ ราคญฺจ โทสญฺจ ปหาย โมหํ, \\ สมฺมปฺปชาโน สุวิมุตฺตจิตฺโต. \\ อนุปาทิยาโน อิธ วา หุรํ วา, \\ ส ภาควา สามญฺญสฺส โหติ. ยมกวคฺโค ปฐโม.}
\csromangathagroup{\csromangathastanza{\csromanfolio{16}โคโปว คาโว คณยํ ปเรสํ, \\ น ภาควา สามญฺญสฺส โหติ. \\ อปฺปมฺปิ เจ สํหิต ภาสมาโน, \\ ธมฺมสฺส โหติ\footnote{โหตี (สี, อิ)} อนุธมฺมจารี.}\csromangathastanza{ราคญฺจ โทสญฺจ ปหาย โมหํ, \\ สมฺมปฺปชาโน สุวิมุตฺตจิตฺโต. \\ อนุปาทิยาโน อิธ วา หุรํ วา, \\ ส ภาควา สามญฺญสฺส โหติ.\csromanspacer{} ยมกวคฺโค ปฐโม.}}

\csromanheader{1}{2. อปฺปมาทวคฺค}
\csromanheader{2}{1. สามาวตีวตฺถุ}
\proseitem{21}{\csromansymbolfootnote{*}{ขุ 10. 30 ปิฏฺเฐ เนตฺติยมฺปิ.}อปฺปมาโท อมตปทํ\footnote{อมตํ ปทํ (ก)}, ปมาโท มจฺจุโน ปทํ.}

\csromangathameasure{อปฺปมตฺตา น มียนฺติ, เย ปมตฺตา ยถา มตา.}
\csromangathagroup{\csromangathastanza{อปฺปมตฺตา น มียนฺติ, เย ปมตฺตา ยถา มตา.}}

\proseitem{22}{เอวํ\footnote{เอตํ (สี, สฺยา, กํ, อิ)} วิเสสโต ญตฺวา, อปฺปมาทมฺหิ ปณฺฑิตา.}

\csromangathameasure{อปฺปมาเท ปโมทนฺติ, อริยานํ โคจเร รตา.}
\csromangathagroup{\csromangathastanza{อปฺปมาเท ปโมทนฺติ, อริยานํ โคจเร รตา.}}

\proseitem{23}{เต ฌายิโน สาตติกา, นิจฺจํ ทฬฺหปรกฺกมา,}

\csromangathameasure{ผุสนฺติ ธีรา นิพฺพานํ, โยคกฺเขมํ อนุตฺตรํ.}
\csromangathagroup{\csromangathastanza{ผุสนฺติ ธีรา นิพฺพานํ, โยคกฺเขมํ อนุตฺตรํ.}}

\csromanheader{2}{2. กุมฺภโฆสกเสฏฺฐิวตฺถุ}
\csromangathameasure{อุฏฺฐานวโต สตีมโต, \\ สุจิกมฺมสฺส นิสมฺมการิโน. \\ สญฺญตสฺส ธมฺมชีวิโน, \\ อปฺปมตฺตสฺส ยโสภิวฑฺฒติ.}
\csromangathagroup{\csromangathastanza{อุฏฺฐานวโต สตีมโต\footnote{สติมโต (สี, สฺยา, ก)}, \\ สุจิกมฺมสฺส นิสมฺมการิโน. \\ สญฺญตสฺส ธมฺมชีวิโน, \\ อปฺปมตฺตสฺส\footnote{สติมโต (สี, สฺยา, ก)} ยโสภิวฑฺฒติ.}}

\vspace{\tipitakaparskip}
\csromancenter{อปฺปมาทวคฺค \textbf{3.}\csromanspacer{}\textbf{ จูฬปนฺถกวตฺถุ} 25.\csromanspacer{} อุฏฺฐาเนน’ปฺปมาเทน, สํยเมน ทเมน จ.}
\vspace{0.5\baselineskip}
\csromangathameasure{ทีปํ กยิราถ เมธาวี, ยํ โอโฆ นาภิกีรติ.}
\csromangathagroup{\csromangathastanza{\csromanfolio{17}ทีปํ กยิราถ เมธาวี, ยํ โอโฆ นาภิกีรติ.}}

\csromanheader{2}{4. พาลนกฺขตฺตสงฺฆุฏฺฐวตฺถุ}
\proseitem{26}{ปมาทมนุยุญฺชนฺติ, พาลา ทุมฺเมธิโน ชนา.}

\csromangathameasure{อปฺปมาทญฺจ เมธาวี, ธนํ เสฏฺฐํว รกฺขติ.}
\csromangathagroup{\csromangathastanza{อปฺปมาทญฺจ เมธาวี, ธนํ เสฏฺฐํว รกฺขติ.}}

\proseitem{27}{มา ปมาทมนุยุญฺเชถ, มา กามรติสนฺถวํ\footnote{สนฺธวํ (ก)}.}

\csromangathameasure{อปฺปมตฺโต หิ ฌายนฺโต, ปปฺโปติ วิปุลํ สุขํ.}
\csromangathagroup{\csromangathastanza{อปฺปมตฺโต หิ ฌายนฺโต, ปปฺโปติ วิปุลํ สุขํ.}}

\csromanheader{2}{5. มหากสฺสปตฺเถรวตฺถุ}
\proseitem{28}{ปมาทํ อปฺปมาเทน, ยทา นุทติ ปณฺฑิโต.}

\csromangathameasure{ปญฺญาปาสาทมารุยฺห, อโสโก โสกินิํ ปชํ. \\ ปพฺพตฏฺโฐว ภูมฏฺเฐ, ธีโร พาเล อเวกฺขติ.}
\csromangathagroup{\csromangathastanza{ปญฺญาปาสาทมารุยฺห, อโสโก โสกินิํ ปชํ. \\ ปพฺพตฏฺโฐว ภูมฏฺเฐ\footnote{ภุมฺมฏฺเฐ (สี, สฺยา)}, ธีโร พาเล อเวกฺขติ.}}

\csromanheader{2}{6. ปมตฺตาปมตฺตเทฺวสหายกภิกฺขุวตฺถุ}
\proseitem{29}{อปฺปมตฺโต ปมตฺเตสุ, สตฺเตสุ พหุชาคโร.}

\csromangathameasure{อพลสฺสํว สีฆสฺโส, หิตฺวา ยาติ สุเมธโส.}
\csromangathagroup{\csromangathastanza{อพลสฺสํว สีฆสฺโส, หิตฺวา ยาติ สุเมธโส.}}

\csromanheader{2}{7. มฆวตฺถุ}
\proseitem{30}{อปฺปมาเทน มฆวา, เทวานํ เสฏฺฐตํ คโต.}

\csromangathameasure{อปฺปมาทํ ปสํสนฺติ, ปมาโท ครหิโต สทา.}
\csromangathagroup{\csromangathastanza{อปฺปมาทํ ปสํสนฺติ, ปมาโท ครหิโต สทา.}}

\csromanheader{2}{8. อญฺญตรภิกฺขุวตฺถุ}
\proseitem{31}{อปฺปมาทรโต ภิกฺขุ, ปมาเท ภยทสฺสิ วา.}

\csromangathameasure{สํโยชนํ อณุํ ถูลํ, ฑหํ อคฺคีว คจฺฉติ.}
\csromangathagroup{\csromangathastanza{สํโยชนํ อณุํ ถูลํ, ฑหํ อคฺคีว คจฺฉติ.}}

\csromanheader{2}{9. นิคมวาสิติสฺสตฺเถรวตฺถุ}
\proseitem{32}{อปฺปมาทรโต ภิกฺขุ, ปมาเท ภยทสฺสิ วา.}

\csromangathameasure{อภพฺโพ ปริหานาย, นิพฺพานสฺเสว สนฺติเก. อปฺปมาทวคฺโค ทุติโย.}
\csromangathagroup{\csromangathastanza{อภพฺโพ ปริหานาย, นิพฺพานสฺเสว สนฺติเก.\csromanspacer{} อปฺปมาทวคฺโค ทุติโย.}}

\csromanheader{1}{ธมฺมปทปาฬิ \textbf{3. จิตฺตวคฺค}}
\csromanheader{2}{1. เมฆิยตฺเถรวตฺถุ}
\csromangathameasure{ผนฺทนํ จปลํ จิตฺตํ, ทูรกฺขํ ทุนฺนิวารยํ. \\ อุชุํ กโรติ เมธาวี, อุสุกาโรว เตชนํ. \\ วาริโชว ถเล ขิตฺโต, โอกโมกต-อุพฺภโต. \\ ปริผนฺทติทํ จิตฺตํ, มารเธยฺยํ ปหาตเว.}
\csromangathagroup{\csromangathastanza{\csromanfolio{18}ผนฺทนํ จปลํ จิตฺตํ, ทูรกฺขํ\footnote{ทุรกฺขํ (สพฺพตฺถ)} ทุนฺนิวารยํ. \\ อุชุํ กโรติ เมธาวี, อุสุกาโรว เตชนํ.}\csromangathastanza{วาริโชว ถเล ขิตฺโต, โอกโมกต-อุพฺภโต. \\ ปริผนฺทติทํ จิตฺตํ, มารเธยฺยํ ปหาตเว.}}

\csromanheader{2}{2. อญฺญตรภิกฺขุวตฺถุ}
\csromangathameasure{ทุนฺนิคฺคหสฺส ลหุโน, ยตฺถกามนิปาติโน. \\ จิตฺตสฺส ทมโถ สาธุ, จิตฺตํ ทนฺตํ สุขาวหํ.}
\csromangathagroup{\csromangathastanza{ทุนฺนิคฺคหสฺส ลหุโน, ยตฺถกามนิปาติโน. \\ จิตฺตสฺส ทมโถ สาธุ, จิตฺตํ ทนฺตํ สุขาวหํ.}}

\csromanheader{2}{3. อญฺญตร-อุกฺกณฺฐิตภิกฺขุวตฺถุ}
\csromangathameasure{สุทุทฺทสํ สุนิปุณํ, ยตฺถกามนิปาตินํ. \\ จิตฺตํ รกฺเขถ เมธาวี, จิตฺตํ คุตฺตํ สุขาวหํ.}
\csromangathagroup{\csromangathastanza{สุทุทฺทสํ สุนิปุณํ, ยตฺถกามนิปาตินํ. \\ จิตฺตํ รกฺเขถ เมธาวี, จิตฺตํ คุตฺตํ สุขาวหํ.}}

\csromanheader{2}{4. สํฆรกฺขิตภาคิเนยฺยตฺเถรวตฺถุ}
\csromangathameasure{ทูรงฺคมํ เอกจรํ, อสรีรํ คุหาสยํ. \\ เย จิตฺตํ สํยมิสฺสนฺติ, โมกฺขนฺติ มารพนฺธนา.}
\csromangathagroup{\csromangathastanza{ทูรงฺคมํ เอกจรํ\footnote{เอกจารํ (ก)}, อสรีรํ คุหาสยํ. \\ เย จิตฺตํ สํยมิสฺสนฺติ, โมกฺขนฺติ มารพนฺธนา.}}

\csromanheader{2}{5. จิตฺถหตฺถตฺเถรวตฺถุ}
\csromangathameasure{อนวฏฺฐิตจิตฺตสฺส, สทฺธมฺมํ อวิชานโต. \\ ปริปฺลวปสาทสฺส, ปญฺญา น ปริปูรติ. \\ อนวสฺสุตจิตฺตสฺส, อนนฺวาหตเจตโส. \\ ปุญฺญปาปปหีนสฺส, นตฺถิ ชาครโต ภยํ.}
\csromangathagroup{\csromangathastanza{อนวฏฺฐิตจิตฺตสฺส, สทฺธมฺมํ อวิชานโต. \\ ปริปฺลวปสาทสฺส, ปญฺญา น ปริปูรติ.}\csromangathastanza{อนวสฺสุตจิตฺตสฺส, อนนฺวาหตเจตโส. \\ ปุญฺญปาปปหีนสฺส, นตฺถิ ชาครโต ภยํ.}}

\csromanheader{2}{6. ปญฺจสตภิกฺขุวตฺถุ}
\proseitem{40}{\csromansymbolfootnote{*}{ขุ 10. 177 ปิฏฺเฐปิ.}กุมฺภูปมํ กายมิมํ วิทิตฺวา,}

\prose{นครูปมํ จิตฺตมิทํ ฐเปตฺวา.}

\csromanmatikaanchor{mtk.12}
\csromantocmarkref{section}{2. อปฺปมาทวคฺค}{mtk.12}
\bookmark[dest={mtk.12},level=1]{2. อปฺปมาทวคฺค}
\csromangathameasure{โยเธถ มารํ ปญฺญาวุเธน, \\ ชิตญฺจ รกฺเข อนิเวสโน สิยา.}
\csromangathagroup{\csromangathastanza{โยเธถ มารํ ปญฺญาวุเธน, \\ ชิตญฺจ รกฺเข อนิเวสโน สิยา.}}

\vspace{\tipitakaparskip}
\csromancenter{ปุปฺผวคฺค \textbf{7.}\csromanspacer{}\textbf{ ปูติคตฺตติสฺสตฺเถรวตฺถุ} 41.\csromanspacer{} อจิรํ วต’ยํ กาโย, ปถวิํ อธิเสสฺสติ.}
\vspace{0.5\baselineskip}
\csromangathameasure{ฉุทฺโธ อเปตวิญฺญาโณ, นิรตฺถํว กลิงฺครํ.}
\csromangathagroup{\csromangathastanza{\csromanfolio{19}ฉุทฺโธ อเปตวิญฺญาโณ, นิรตฺถํว กลิงฺครํ.}}

\csromanheader{2}{8. นนฺทโคปาลกวตฺถุ}
\csromangathameasure{ทิโส ทิสํ ยํ ตํ กยิรา, เวรี วา ปน เวรินํ. \\ มิจฺฉาปณิหิตํ จิตฺตํ, ปาปิโย นํ ตโต กเร.}
\csromangathagroup{\csromangathastanza{ทิโส ทิสํ ยํ ตํ กยิรา, เวรี วา ปน เวรินํ. \\ มิจฺฉาปณิหิตํ จิตฺตํ, ปาปิโย\footnote{ปาปิยํ (?)} นํ ตโต กเร.}}

\csromanheader{2}{9. โสเรยฺยวตฺถุ}
\csromangathameasure{น ตํ มาตา ปิตา กยิรา, อญฺเญ วาปิ จ ญาตกา. \\ สมฺมาปณิหิตํ จิตฺตํ, เสยฺยโส นํ ตโต กเร. จิตฺตวคฺโค ตติโย.}
\csromangathagroup{\csromangathastanza{น ตํ มาตา ปิตา กยิรา, อญฺเญ วาปิ จ ญาตกา. \\ สมฺมาปณิหิตํ จิตฺตํ, เสยฺยโส นํ ตโต กเร.\csromanspacer{} จิตฺตวคฺโค ตติโย.}}

\csromanheader{1}{4. ปุปฺผวคฺค}
\csromanheader{2}{1. ปถวิกถาปสุตปญฺจสตภิกฺขุวตฺถุ}
\csromangathameasure{โก อิมํ ปถวิํ วิเจสฺสติ, \\ ยมโลกญฺจ อิมํ สเทวกํ. \\ โก ธมฺมปทํ สุเทสิตํ, \\ กุสโล ปุปฺผมิว ปเจสฺสติ.}
\csromangathagroup{\csromangathastanza{โก อิมํ\footnote{โกมํ (ก)} ปถวิํ วิเจสฺสติ\footnote{วิเชสฺสติ (สี, สฺยา, อิ)}, \\ ยมโลกญฺจ อิมํ สเทวกํ. \\ โก ธมฺมปทํ สุเทสิตํ, \\ กุสโล ปุปฺผมิว ปเจสฺสติ\footnote{โกมํ (ก)}.}}

\proseitem{45}{เสโข ปถวิํ วิเจสฺสติ,}

\prose{ยมโลกญฺจ อิมํ สเทวกํ.}

\csromangathameasure{เสโข ธมฺมปทํ สุเทสิตํ, \\ กุสโล ปุปฺผมิว ปเจสฺสติ.}
\csromangathagroup{\csromangathastanza{เสโข ธมฺมปทํ สุเทสิตํ, \\ กุสโล ปุปฺผมิว ปเจสฺสติ.}}

\vspace{\tipitakaparskip}
\csromancenter{ธมฺมปทปาฬิ \textbf{2.}\csromanspacer{}\textbf{ มรีจิกมฺมฏฺฐานิกภิกฺขุวตฺถุ} 46.\csromanspacer{} เผณูปมํ กายมิมํ วิทิตฺวา,}
\prose{\csromanfolio{20}มรีจิธมฺมํ อภิสมฺพุธาโน.}

\csromangathameasure{เฉตฺวาน มารสฺส ปปุปฺผกานิ, \\ อทสฺสนํ มจฺจุราชสฺส คจฺเฉ.}
\csromangathagroup{\csromangathastanza{เฉตฺวาน มารสฺส ปปุปฺผกานิ\footnote{สปุปฺผกานิ (ฏีกา)}, \\ อทสฺสนํ มจฺจุราชสฺส คจฺเฉ.}}

\vspace{\tipitakaparskip}
\csromancenter{วิฏฏูภวตฺถุ 47.\csromanspacer{} ปุปฺผานิ เหว ปจินนฺตํ, พฺยาสตฺตมนสํ\footnote{พฺยาสตฺตมานสํ (ก)} นรํ.}
\vspace{0.5\baselineskip}
\csromangathameasure{สุตฺตํ คามํ มโหโฆว, มจฺจุ อาทาย คจฺฉติ.}
\csromangathagroup{\csromangathastanza{สุตฺตํ คามํ มโหโฆว, มจฺจุ อาทาย คจฺฉติ.}}

\csromanheader{2}{4. ปติปูชิกกุมาริวตฺถุ}
\csromangathameasure{ปุปฺผานิ เหว ปจินนฺตํ, พฺยาสตฺตมนสํ นรํ. \\ อติตฺตญฺเญว กาเมสุ, อนฺตโก กุรุเต วสํ.}
\csromangathagroup{\csromangathastanza{ปุปฺผานิ เหว ปจินนฺตํ, พฺยาสตฺตมนสํ นรํ. \\ อติตฺตญฺเญว กาเมสุ, อนฺตโก กุรุเต วสํ.}}

\csromanheader{2}{5. มจฺฉริยโกสิยเสฏฺฐิวตฺถุ}
\proseitem{49}{\csromansymbolfootnote{+}{ขุ 10. 159 ปิฏฺเฐปิ.}ยถาปิ ภมโร ปุปฺผํ, วณฺณคนฺธมเหฐยํ\footnote{วณฺณคนฺธมโปฐยํ (ก)}.}

\csromangathameasure{ปเลติ รสมาทาย, เอวํ คาเม มุนี จเร.}
\csromangathagroup{\csromangathastanza{ปเลติ รสมาทาย, เอวํ คาเม มุนี จเร.}}

\csromanheader{2}{6. ปาเวยฺย-อาชีวกวตฺถุ}
\csromangathameasure{น ปเรสํ วิโลมานิ, น ปเรสํ กตากตํ. \\ อตฺตโนว อเวกฺเขยฺย, กตานิ อกตานิ จ.}
\csromangathagroup{\csromangathastanza{น ปเรสํ วิโลมานิ, น ปเรสํ กตากตํ. \\ อตฺตโนว อเวกฺเขยฺย, กตานิ อกตานิ จ.}}

\vspace{\tipitakaparskip}
\csromancenter{ฉตฺตปาณิ-อุปาสกวตฺถุ 51. * ยถาปิ รุจิรํ ปุปฺผํ, วณฺณวนฺตํ อคนฺธกํ.}
\vspace{0.5\baselineskip}
\csromangathameasure{เอวํ สุภาสิตา วาจา, อผลา โหติ อกุพฺพโต.}
\csromangathagroup{\csromangathastanza{เอวํ สุภาสิตา วาจา, อผลา โหติ อกุพฺพโต.}}

\proseitem{52}{ยถาปิ รุจิรํ ปุปฺผํ, วณฺณวนฺตํ สคนฺธกํ\footnote{สุคนฺธกํ (ก)}.}

\csromangathameasure{เอวํ สุภาสิตา วาจา, สผลา โหติ กุพฺพโต.}
\csromangathagroup{\csromangathastanza{เอวํ สุภาสิตา วาจา, สผลา โหติ กุพฺพโต\footnote{สกุพฺพโต (สี, อิ), ปกุพฺพโต (สี-ฏฺฐ), สุกุพฺพโต (สฺยา, กํ)}.}}

\chapterheadpageread{4. ปุปฺผวคฺค \textbf{8. วิสาขาวตฺถุ}}
\csromangathameasure{ยถาปิ ปุปฺผราสิมฺหา, กยิรา มาลาคุเณ พหู. \\ เอวํ ชาเตน มจฺเจน, กตฺตพฺพํ กุสลํ พหุํ.}
\csromangathagroup{\csromangathastanza{\csromanfolio{21}ยถาปิ ปุปฺผราสิมฺหา, กยิรา มาลาคุเณ พหู. \\ เอวํ ชาเตน มจฺเจน, กตฺตพฺพํ กุสลํ พหุํ.}}

\csromanheader{2}{9. อานนฺทตฺเถรปญฺหาวตฺถุ}
\proseitem{54}{\csromansymbolfootnote{*}{อํ 1. 227 ปิฏฺเฐปิ.}น ปุปฺผคนฺโธ ปฏิวาตเมติ,}

\prose{น จนฺทนํ ตครมลฺลิกา\footnote{ตคฺครมลฺลิกา (ก)} วา.}

\csromangathameasure{สตญฺจ คนฺโธ ปฏิวาตเมติ, \\ สพฺพา ทิสา สปฺปุริโส ปวายติ.}
\csromangathagroup{\csromangathastanza{สตญฺจ คนฺโธ ปฏิวาตเมติ, \\ สพฺพา ทิสา สปฺปุริโส ปวายติ.}}

\csromanmatikaanchor{mtk.13}
\csromantocmarkref{section}{3. จิตฺตวคฺค}{mtk.13}
\bookmark[dest={mtk.13},level=1]{3. จิตฺตวคฺค}
\proseitem{55}{จนฺทนํ ตครํ วาปิ, อุปฺปลํ อถ วสฺสิกี.}

\csromangathameasure{เอเตสํ คนฺธชาตานํ, สีลคนฺโธ อนุตฺตโร.}
\csromangathagroup{\csromangathastanza{เอเตสํ คนฺธชาตานํ, สีลคนฺโธ อนุตฺตโร.}}

\csromanheader{2}{10. มหากสฺสปตฺเถรปิณฺฑปาตทินฺนวตฺถุ}
\proseitem{56}{อปฺปมตฺโต อยํ คนฺโธ, ยายํ ตครจนฺทนี\footnote{ยฺวายํ ตครจนฺทนํ (ก)}.}

\csromangathameasure{โย จ สีลวตํ คนฺโธ, วาติ เทเวสุ อุตฺตโม.}
\csromangathagroup{\csromangathastanza{โย จ สีลวตํ คนฺโธ, วาติ เทเวสุ อุตฺตโม.}}

\csromanheader{2}{11. โคธิกตฺเถรปรินิพฺพานวตฺถุ}
\proseitem{57}{เตสํ สมฺปนฺนสีลานํ, อปฺปมาทวิหารินํ.}

\csromangathameasure{สมฺมทญฺญา วิมุตฺตานํ, มาโร มคฺคํ น วินฺทติ.}
\csromangathagroup{\csromangathastanza{สมฺมทญฺญา วิมุตฺตานํ, มาโร มคฺคํ น วินฺทติ.}}

\vspace{\tipitakaparskip}
\csromancenter{ครหทินฺนวตฺถุ 58.\csromanspacer{} ยถา สงฺการธานสฺมิํ\footnote{สงฺการฐานสฺมิํ (ก)}, อุชฺฌิตสฺมิํ มหาปเถ.}
\vspace{0.5\baselineskip}
\csromangathameasure{ปทุมํ ตตฺถ ชาเยถ, สุจิคนฺธํ มโนรมํ.}
\csromangathagroup{\csromangathastanza{ปทุมํ ตตฺถ ชาเยถ, สุจิคนฺธํ มโนรมํ.}}

\proseitem{59}{เอวํ สงฺการภูเตสุ, อนฺธภูเต\footnote{อนฺธีภูเต (ก)} ปุถุชฺชเน.}

\csromangathameasure{อติโรจติ ปญฺญาย, สมฺมาสมฺพุทฺธสาวโก. ปุปฺผวคฺโค จตุตฺโถ.}
\csromangathagroup{\csromangathastanza{อติโรจติ ปญฺญาย, สมฺมาสมฺพุทฺธสาวโก.\csromanspacer{} ปุปฺผวคฺโค จตุตฺโถ.}}

\csromanheader{1}{ธมฺมปทปาฬิ \textbf{5. พาลวคฺค}}
\csromanheader{2}{1. อญฺญตรปุริสวตฺถุ}
\csromangathameasure{ทีฆา ชาครโต รตฺติ, ทีฆํ สนฺตสฺส โยชนํ. \\ ทีโฆ พาลาน สํสาโร, สทฺธมฺมํ อวิชานตํ.}
\csromangathagroup{\csromangathastanza{\csromanfolio{22}ทีฆา ชาครโต รตฺติ, ทีฆํ สนฺตสฺส โยชนํ. \\ ทีโฆ พาลาน สํสาโร, สทฺธมฺมํ อวิชานตํ.}}

\csromanheader{2}{2. มหากสฺสปสทฺธิวิหาริกวตฺถุ}
\proseitem{61}{จรญฺเจ นาธิคจฺเฉยฺย, เสยฺยํ สทิสมตฺตโน.}

\csromangathameasure{เอกจริยํ ทฬฺหํ กยิรา, นตฺถิ พาเล สหายตา.}
\csromangathagroup{\csromangathastanza{เอกจริยํ\footnote{เอกจฺจริยํ (ก)} ทฬฺหํ กยิรา, นตฺถิ พาเล สหายตา.}}

\csromanheader{2}{3. อานนฺทเสฏฺฐิวตฺถุ}
\csromangathameasure{ปุตฺตา ม’ตฺถิ ธนํ ม’ตฺถิ, อิติ พาโล วิหญฺญติ. \\ อตฺตา หิ อตฺตโน นตฺถิ, กุโต ปุตฺตา กุโต ธนํ.}
\csromangathagroup{\csromangathastanza{ปุตฺตา ม’ตฺถิ ธนํ ม’ตฺถิ\footnote{ปุตฺตมตฺถิ ธนมตฺถิ (ก)}, อิติ พาโล วิหญฺญติ. \\ อตฺตา หิ\footnote{ปุตฺตมตฺถิ ธนมตฺถิ (ก)} อตฺตโน นตฺถิ, กุโต ปุตฺตา กุโต ธนํ.}}

\csromanheader{2}{4. คณฺฐิเภทกโจรวตฺถุ}
\csromangathameasure{โย พาโล มญฺญติ พาลฺยํ, ปณฺฑิโต วาปิ เตน โส. \\ พาโล จ ปณฺฑิตมานี, ส เว “พาโล”ติ วุจฺจติ.}
\csromangathagroup{\csromangathastanza{โย พาโล มญฺญติ พาลฺยํ, ปณฺฑิโต วาปิ เตน โส. \\ พาโล จ ปณฺฑิตมานี, ส เว “พาโล”ติ วุจฺจติ.}}

\csromanheader{2}{5. อุทายิตฺเถรวตฺถุ}
\proseitem{64}{ยาวชีวมฺปิ เจ พาโล, ปณฺฑิตํ ปยิรุปาสติ.}

\csromanmatikaanchor{mtk.14}
\csromantocmarkref{section}{4. ปุปฺผวคฺค}{mtk.14}
\bookmark[dest={mtk.14},level=1]{4. ปุปฺผวคฺค}
\csromangathameasure{น โส ธมฺมํ วิชานาติ, ทพฺพี สูปรสํ ยถา.}
\csromangathagroup{\csromangathastanza{น โส ธมฺมํ วิชานาติ, ทพฺพี สูปรสํ ยถา.}}

\csromanheader{2}{6. ติํสปาเวยฺยกภิกฺขุวตฺถุ}
\proseitem{65}{มุหุตฺตมปิ เจ วิญฺญู, ปณฺฑิตํ ปยิรุปาสติ.}

\csromangathameasure{ขิปฺปํ ธมฺมํ วิชานาติ, ชิวฺหา สูปรสํ ยถา.}
\csromangathagroup{\csromangathastanza{ขิปฺปํ ธมฺมํ วิชานาติ, ชิวฺหา สูปรสํ ยถา.}}

\csromanheader{2}{7. สุปฺปพุทฺธกุฏฺฐิวตฺถุ}
\proseitem{66}{\csromansymbolfootnote{*}{ขุ 10. 110 ปิฏฺเฐปิ.}จรนฺติ พาลา ทุมฺเมธา, อมิตฺเตเนว อตฺตนา.}

\csromangathameasure{กโรนฺตา ปาปกํ กมฺมํ, ยํ โหติ กฏุกปฺผลํ.}
\csromangathagroup{\csromangathastanza{กโรนฺตา ปาปกํ กมฺมํ, ยํ โหติ กฏุกปฺผลํ.}}

\chapterheadpageread{5. พาลวคฺค \textbf{8. กสฺสกวตฺถุ}}
\csromangathameasure{น ตํ กมฺมํ กตํ สาธุ, ยํ กตฺวา อนุตปฺปติ. \\ ยสฺส อสฺสุมุโข โรทํ, วิปากํ ปฏิเสวติ.}
\csromangathagroup{\csromangathastanza{\csromanfolio{23}น ตํ กมฺมํ กตํ สาธุ, ยํ กตฺวา อนุตปฺปติ. \\ ยสฺส อสฺสุมุโข โรทํ, วิปากํ ปฏิเสวติ.}}

\csromanheader{2}{9. สุมนมาลาการวตฺถุ}
\csromangathameasure{ตญฺจ กมฺมํ กตํ สาธุ, ยํ กตฺวา นานุตปฺปติ. \\ ยสฺส ปตีโต สุมโน, วิปากํ ปฏิเสวติ.}
\csromangathagroup{\csromangathastanza{ตญฺจ กมฺมํ กตํ สาธุ, ยํ กตฺวา นานุตปฺปติ. \\ ยสฺส ปตีโต สุมโน, วิปากํ ปฏิเสวติ.}}

\csromanheader{2}{10. อุปฺปลวณฺณตฺเถรีวตฺถุ}
\csromangathameasure{มธุํวา มญฺญติ พาโล, ยาว ปาปํ น ปจฺจติ. \\ ยทา จ ปจฺจติ ปาปํ, อถ ทุกฺขํ นิคจฺฉติ.}
\csromangathagroup{\csromangathastanza{มธุํวา\footnote{มธุวา (สพฺพตฺถ) มธุํ วา (ที 1 ฏีกา ปสฺสิตพฺพา)} มญฺญติ พาโล, ยาว ปาปํ น ปจฺจติ. \\ ยทา จ ปจฺจติ ปาปํ, อถ\footnote{มธุวา (สพฺพตฺถ) มธุํ วา (ที 1 ฏีกา ปสฺสิตพฺพา)} ทุกฺขํ นิคจฺฉติ.}}

\vspace{\tipitakaparskip}
\csromancenter{ชมฺพุกตฺเถรวตฺถุ 70.\csromanspacer{} มาเส มาเส กุสคฺเคน, พาโล ภุญฺเชยฺย โภชนํ.}
\vspace{0.5\baselineskip}
\csromangathameasure{น โส สงฺขาตธมฺมานํ, กลํ อคฺฆติ โสฬสิํ.}
\csromangathagroup{\csromangathastanza{น โส สงฺขาตธมฺมานํ\footnote{สงฺขตธมฺมานํ (สี, อิ, ก)}, กลํ อคฺฆติ โสฬสิํ.}}

\vspace{\tipitakaparskip}
\csromancenter{อหิเปตวตฺถุ 71. * น หิ ปาปํ กตํ กมฺมํ, สชฺชุขีรํว มุจฺจติ.}
\vspace{0.5\baselineskip}
\csromangathameasure{qอหนฺตํ พาลมเนฺวติ, ภสฺมจฺฉนฺโนว ปาวโก.}
\csromangathagroup{\csromangathastanza{qอหนฺตํ พาลมเนฺวติ, ภสฺมจฺฉนฺโนว\footnote{ภสฺมาฉนฺโนว (สฺยา, อิ, ก)} ปาวโก.}}

\csromanheader{2}{13. สฏฺฐิกูฏเปตวตฺถุ}
\csromangathameasure{ยาวเทว อนตฺถาย, ญตฺตํ พาลสฺส ชายติ. \\ หนฺติ พาลสฺส สุกฺกํสํ, มุทฺธมสฺส วิปาตยํ.}
\csromangathagroup{\csromangathastanza{ยาวเทว อนตฺถาย, ญตฺตํ\footnote{ญาตฺตํ (?)} พาลสฺส ชายติ. \\ หนฺติ พาลสฺส สุกฺกํสํ, มุทฺธมสฺส วิปาตยํ.}}

\csromanheader{2}{14. จิตฺตคหปติวตฺถุ}
\csromangathameasure{อสนฺตํ ภาวนมิจฺเฉยฺย, ปุเรกฺขารญฺจ ภิกฺขุสุ. \\ อาวาเสสุ จ อิสฺสริยํ, ปูชํ ปรกุเลสุ จ.}
\csromangathagroup{\csromangathastanza{อสนฺตํ ภาวนมิจฺเฉยฺย\footnote{อสนฺตํ ภาวมิจฺเฉยฺย (สฺยา), อสนฺตภาวนมิจฺเฉยฺย (ก)}, ปุเรกฺขารญฺจ ภิกฺขุสุ. \\ อาวาเสสุ จ อิสฺสริยํ, ปูชํ ปรกุเลสุ จ.}}

\csromanheader{2}{ธมฺมปทปาฬิ 74. มเมว กต มญฺญนฺตุ, คิหี ปพฺพชิตา อุโภ.}
\csromangathameasure{มเมวาติวสา อสฺสุ, กิจฺจากิจฺเจสุ กิสฺมิจิ. \\ อิติ พาลสฺส สงฺกปฺโป, อิจฺฉา มาโน จ วฑฺฒติ.}
\csromangathagroup{\csromangathastanza{\csromanfolio{24}มเมวาติวสา อสฺสุ, กิจฺจากิจฺเจสุ กิสฺมิจิ. \\ อิติ พาลสฺส สงฺกปฺโป, อิจฺฉา มาโน จ วฑฺฒติ.}}

\csromanheader{2}{15. วนวาสิติสฺสสามเณรวตฺถุ}
\csromangathameasure{อญฺญา หิ ลาภูปนิสา, อญฺญา นิพฺพานคามินี. \\ เอวเมตํ อภิญฺญาย, ภิกฺขุ พุทฺธสฺส สาวโก. \\ สกฺการํ นาภินนฺเทยฺย, วิเวกมนุพฺรูหเย. พาลวคฺโค ปญฺจโม.}
\csromangathagroup{\csromangathastanza{อญฺญา หิ ลาภูปนิสา, อญฺญา นิพฺพานคามินี. \\ เอวเมตํ อภิญฺญาย, ภิกฺขุ พุทฺธสฺส สาวโก. \\ สกฺการํ นาภินนฺเทยฺย, วิเวกมนุพฺรูหเย.\csromanspacer{} พาลวคฺโค ปญฺจโม.}}

\csromanheader{1}{6. ปณฺฑิตวคฺค}
\csromanheader{2}{1. ราธตฺเถรวตฺถุ}
\csromangathameasure{นิธีนํว ปวตฺตารํ, ยํ ปสฺเส วชฺชทสฺสินํ. \\ นิคฺคยฺหวาทิํ เมธาวิํ, ตาทิสํ ปณฺฑิตํ ภเช. \\ ตาทิสํ ภชมานสฺส, เสยฺโย โหติ น ปาปิโย.}
\csromangathagroup{\csromangathastanza{นิธีนํว ปวตฺตารํ, ยํ ปสฺเส วชฺชทสฺสินํ. \\ นิคฺคยฺหวาทิํ เมธาวิํ, ตาทิสํ ปณฺฑิตํ ภเช. \\ ตาทิสํ ภชมานสฺส, เสยฺโย โหติ น ปาปิโย.}}

\csromanheader{2}{2. อสฺสชิปุนพฺพสุกวตฺถุ}
\csromangathameasure{โอวเทยฺยา’นุสาเสยฺย, อสพฺภา จ นิวารเย. \\ สตํ หิ โส ปิโย โหติ, อสตํ โหติ อปฺปิโย.}
\csromangathagroup{\csromangathastanza{โอวเทยฺยา’นุสาเสยฺย, อสพฺภา จ นิวารเย. \\ สตํ หิ โส ปิโย โหติ, อสตํ โหติ อปฺปิโย.}}

\csromanheader{2}{3. ฉนฺนตฺเถรวตฺถุ}
\csromangathameasure{น ภเช ปาปเก มิตฺเต, น ภเช ปุริสาธเม. \\ ภเชถ มิตฺเต กลฺยาเณ, ภเชถ ปุริสุตฺตเม.}
\csromangathagroup{\csromangathastanza{น ภเช ปาปเก มิตฺเต, น ภเช ปุริสาธเม. \\ ภเชถ มิตฺเต กลฺยาเณ, ภเชถ ปุริสุตฺตเม.}}

\csromanheader{2}{4. มหากปฺปินตฺเถรวตฺถุ}
\csromangathameasure{ธมฺมปีติ สุขํ เสติ, วิปฺปสนฺเนน เจตสา. \\ อริยปฺปเวทิเต ธมฺเม, สทา รมติ ปณฺฑิโต.}
\csromangathagroup{\csromangathastanza{ธมฺมปีติ สุขํ เสติ, วิปฺปสนฺเนน เจตสา. \\ อริยปฺปเวทิเต ธมฺเม, สทา รมติ ปณฺฑิโต.}}

\chapterheadpageread{6. ปณฺฑิตวคฺค \textbf{5. ปณฺฑิตสามเณรวตฺถุ}}
\proseitem{80}{\csromanfolio{25}\csromansymbolfootnote{*}{ขุ 2. 224, 334.}อุทกํ หิ นยนฺติ เนตฺติกา, อุสุการา นมยนฺติ\footnote{ทมยนฺติ (ก)} เตชนํ.}

\csromangathameasure{ทารุํ นมยนฺติ1 ตจฺฉกา, อตฺตานํ ทมยนฺติ ปณฺฑิตา.}
\csromangathagroup{\csromangathastanza{ทารุํ นมยนฺติ1 ตจฺฉกา, อตฺตานํ ทมยนฺติ ปณฺฑิตา.}}

\vspace{\tipitakaparskip}
\csromancenter{ลกุณฺฑกภทฺทิยตฺเถรวตฺถุ 81.\csromanspacer{} เสโล ยถา เอกฆโน\footnote{เอกคฺฆโน (ก)}, วาเตน น สมีรติ.}
\vspace{0.5\baselineskip}
\csromangathameasure{เอวํ นินฺทาปสํสาสุ, น สมิญฺชนฺติ ปณฺฑิตา.}
\csromangathagroup{\csromangathastanza{เอวํ นินฺทาปสํสาสุ, น สมิญฺชนฺติ ปณฺฑิตา.}}

\vspace{\tipitakaparskip}
\csromancenter{กาณมาตาวตฺถุ 82.\csromanspacer{} ยถาปิ รหโท คมฺภีโร, วิปฺปสนฺโน อนาวิโล.}
\vspace{0.5\baselineskip}
\csromangathameasure{เอวํ ธมฺมานิ สุตฺวาน, วิปฺปสีทนฺติ ปณฺฑิตา.}
\csromangathagroup{\csromangathastanza{เอวํ ธมฺมานิ สุตฺวาน, วิปฺปสีทนฺติ ปณฺฑิตา.}}

\csromanheader{2}{8. ปญฺจสตภิกฺขุวตฺถุ}
\csromangathameasure{สพฺพตฺถ เว สปฺปุริสา จชนฺติ, \\ น กามกามา ลปยนฺติ สนฺโต. \\ สุเขน ผุฏฺฐา อถ วา ทุเขน, \\ น อุจฺจาวจํ ปณฺฑิตา ทสฺสยนฺติ.}
\csromangathagroup{\csromangathastanza{สพฺพตฺถ เว สปฺปุริสา จชนฺติ, \\ น กามกามา ลปยนฺติ สนฺโต. \\ สุเขน ผุฏฺฐา อถ วา ทุเขน, \\ น อุจฺจาวจํ\footnote{โนจฺจาวจํ (สี-ฏฺฐ)} ปณฺฑิตา ทสฺสยนฺติ.}}

\csromanheader{2}{9. ธมฺมิกตฺเถรวตฺถุ}
\csromangathameasure{น อตฺตเหตุ น ปรสฺส เหตุ, \\ น ปุตฺตมิจฺเฉ น ธนํ น รฏฺฐํ. \\ น อิจฺเฉยฺย อธมฺเมน สมิทฺธิมตฺตโน, \\ ส สีลวา ปญฺญวา ธมฺมิโก สิยา.}
\csromangathagroup{\csromangathastanza{น อตฺตเหตุ น ปรสฺส เหตุ, \\ น ปุตฺตมิจฺเฉ น ธนํ น รฏฺฐํ. \\ น อิจฺเฉยฺย\footnote{นยิจฺเฉ (อิ), นิจฺเฉ (?)} อธมฺเมน สมิทฺธิมตฺตโน, \\ ส สีลวา ปญฺญวา ธมฺมิโก สิยา.}}

\csromanheader{2}{10. ธมฺมสฺสวนตฺเถรวตฺถุ}
\csromanmatikaanchor{mtk.15}
\csromantocmarkref{section}{5. พาลวคฺค}{mtk.15}
\bookmark[dest={mtk.15},level=1]{5. พาลวคฺค}
\csromangathameasure{อปฺปกา เต มนุสฺเสสุ, เย ชนา ปารคามิโน. \\ อถายํ อิตรา ปชา, ตีรเมวานุธาวติ. \\ เย จ โข สมฺมทกฺขาเต, ธมฺเม ธมฺมานุวตฺติโน. \\ เต ชนา ปารเมสฺสนฺติ, มจฺจุเธยฺยํ สุทุตฺตรํ.}
\csromangathagroup{\csromangathastanza{อปฺปกา เต มนุสฺเสสุ, เย ชนา ปารคามิโน. \\ อถายํ อิตรา ปชา, ตีรเมวานุธาวติ.}\csromangathastanza{เย จ โข สมฺมทกฺขาเต, ธมฺเม ธมฺมานุวตฺติโน. \\ เต ชนา ปารเมสฺสนฺติ, มจฺจุเธยฺยํ สุทุตฺตรํ.}}

\csromanheader{2}{ธมฺมปทปาฬิ \textbf{11. ปญฺจสต-อาคนฺตุกภิกฺขุวตฺถุ}}
\proseitem{87}{\csromanfolio{26}\csromansymbolfootnote{*}{อํ 3. 446, 463-4 ปิฏฺเฐสุปิ.}กณฺหํ ธมฺมํ วิปฺปหาย, สุกฺกํ ภาเวถ ปณฺฑิโต.}

\csromangathameasure{โอกา อโนกมาคมฺม, วิเวเก ยตฺถ ทูรมํ.}
\csromangathagroup{\csromangathastanza{โอกา อโนกมาคมฺม, วิเวเก ยตฺถ ทูรมํ.}}

\proseitem{88}{\csromansymbolmark{*}ตตฺราภิรติมิจฺเฉยฺย, หิตฺวา กาเม อกิญฺจโน.}

\csromangathameasure{ปริโยทเปยฺย อตฺตานํ, จิตฺตเกฺลเสหิ ปณฺฑิโต.}
\csromangathagroup{\csromangathastanza{ปริโยทเปยฺย\footnote{ปริโยทาเปยฺย (?)} อตฺตานํ, จิตฺตเกฺลเสหิ ปณฺฑิโต.}}

\proseitem{89}{\csromansymbolmark{*}เยสํ สมฺโพธิยงฺเคสุ, สมฺมา จิตฺตํ สุภาวิตํ.}

\csromangathameasure{อาทานปฏินิสฺสคฺเค, อนุปาทาย เย รตา. \\ ขีณาสวา ชุติมนฺโต, เต โลเก ปรินิพฺพุตา. ปณฺฑิตวคฺโค ฉฏฺโฐ.}
\csromangathagroup{\csromangathastanza{อาทานปฏินิสฺสคฺเค, อนุปาทาย เย รตา. \\ ขีณาสวา ชุติมนฺโต, เต โลเก ปรินิพฺพุตา.\csromanspacer{} ปณฺฑิตวคฺโค ฉฏฺโฐ.}}

\csromanheader{1}{7. อรหนฺตวคฺค}
\csromanheader{2}{1. ชีวกปญฺหวตฺถุ}
\proseitem{90}{คตทฺธิโน วิโสกสฺส, วิปฺปมุตฺตสฺส สพฺพธิ.}

\csromangathameasure{สพฺพคนฺถปฺปหีนสฺส, ปริฬาโห น วิชฺชติ.}
\csromangathagroup{\csromangathastanza{สพฺพคนฺถปฺปหีนสฺส, ปริฬาโห น วิชฺชติ.}}

\csromanheader{2}{2. มหากสฺสปตฺเถรวตฺถุ}
\csromangathameasure{อุยฺยุญฺชนฺติ สตีมนฺโต, น นิเกเต รมนฺติ เต. \\ หํสาว ปลฺลลํ หิตฺวา, โอกโมกํ ชหนฺติ เต.}
\csromangathagroup{\csromangathastanza{อุยฺยุญฺชนฺติ สตีมนฺโต, น นิเกเต รมนฺติ เต. \\ หํสาว ปลฺลลํ หิตฺวา, โอกโมกํ ชหนฺติ เต.}}

\csromanheader{2}{3. เพลฏฺฐสีสตฺเถรวตฺถุ}
\csromangathameasure{เยสํ สนฺนิจโย นตฺถิ, เย ปริญฺญาตโภชนา. \\ สุญฺญโต อนิมิตฺโต จ, วิโมกฺโข เยสํ โคจโร. \\ อากาเสว สกุนฺตานํ, คติ เตสํ ทุรนฺนยา.}
\csromangathagroup{\csromangathastanza{เยสํ สนฺนิจโย นตฺถิ, เย ปริญฺญาตโภชนา. \\ สุญฺญโต อนิมิตฺโต จ, วิโมกฺโข เยสํ โคจโร. \\ อากาเสว สกุนฺตานํ\footnote{สกุณานํ (ก)}, คติ เตสํ ทุรนฺนยา.}}

\chapterheadpageread{7. อรหนฺตวคฺค \textbf{4. อนุรุทฺธตฺเถรวตฺถุ}}
\csromangathameasure{ยสฺสาสวา ปริกฺขีณา, อาหาเร จ อนิสฺสิโต. \\ สุญฺญโต อนิมิตฺโต จ, วิโมกฺโข ยสฺส โคจโร. \\ อากาเสว สกุนฺตานํ, ปทํ ตสฺส ทุรนฺนยํ.}
\csromangathagroup{\csromangathastanza{\csromanfolio{27}ยสฺสาสวา ปริกฺขีณา, อาหาเร จ อนิสฺสิโต. \\ สุญฺญโต อนิมิตฺโต จ, วิโมกฺโข ยสฺส โคจโร. \\ อากาเสว สกุนฺตานํ, ปทํ ตสฺส ทุรนฺนยํ.}}

\vspace{\tipitakaparskip}
\csromancenter{มหากจฺจายนตฺเถรวตฺถุ 94. * ยสฺสินฺทฺริยานิ สมถงฺคตานิ\footnote{สมถํ คตานิ (สี, อิ)},}
\prose{อสฺสา ยถา สารถินา สุทนฺตา.}

\csromangathameasure{ปหีนมานสฺส อนาสวสฺส, \\ เทวาปิ ตสฺส ปิหยนฺติ ตาทิโน.}
\csromangathagroup{\csromangathastanza{ปหีนมานสฺส อนาสวสฺส, \\ เทวาปิ ตสฺส ปิหยนฺติ ตาทิโน.}}

\csromanheader{2}{6. สาริปุตฺตตฺเถรวตฺถุ}
\csromangathameasure{ปถวิสโม โน วิรุชฺฌติ, \\ อินฺทขิลุปโม ตาทิ สุพฺพโต. \\ รหโทว อเปตกทฺทโม, \\ สํสารา น ภวนฺติ ตาทิโน.}
\csromangathagroup{\csromangathastanza{ปถวิสโม โน วิรุชฺฌติ, \\ อินฺทขิลุปโม\footnote{อินฺทขีลูปโม (สี, สฺยา, อิ)} ตาทิ สุพฺพโต. \\ รหโทว อเปตกทฺทโม, \\ สํสารา น ภวนฺติ ตาทิโน.}}

\csromanheader{2}{7. โกสมฺพิวาสิติสฺสตฺเถรสามเณรวตฺถุ}
\csromangathameasure{สนฺตํ ตสฺส มนํ โหติ, สนฺตา วาจา จ กมฺม จ. \\ สมฺมทญฺญา วิมุตฺตสฺส, อุปสนฺตสฺส ตาทิโน.}
\csromangathagroup{\csromangathastanza{สนฺตํ ตสฺส มนํ โหติ, สนฺตา วาจา จ กมฺม จ. \\ สมฺมทญฺญา วิมุตฺตสฺส, อุปสนฺตสฺส ตาทิโน.}}

\csromanheader{2}{8. สาริปุตฺตตฺเถรวตฺถุ}
\csromangathameasure{อสฺสทฺโธ อกตญฺญู จ, สนฺธิจฺเฉโท จ โย นโร. \\ หตาวกาโส วนฺตาโส, ส เว อุตฺตมโปริโส.}
\csromangathagroup{\csromangathastanza{อสฺสทฺโธ อกตญฺญู จ, สนฺธิจฺเฉโท จ โย นโร. \\ หตาวกาโส วนฺตาโส, ส เว อุตฺตมโปริโส.}}

\csromanheader{2}{9. ขทิรวนิยเรวตตฺเถรวตฺถุ}
\proseitem{98}{\csromansymbolfootnote{+}{ขุ 2. 344 ปิฏฺเฐปิ.}คาเม วา ยทิ วารญฺเญ, นินฺเน วา ยทิ วา ถเล.}

\csromangathameasure{ยตฺถ อรหนฺโต วิหรนฺติ, ตํ ภูมิรามเณยฺยกํ.}
\csromangathagroup{\csromangathastanza{ยตฺถ อรหนฺโต วิหรนฺติ, ตํ ภูมิรามเณยฺยกํ.}}

\gambhira{ธมฺมปทปาฬิ \textbf{10. อญฺญตร-อิตฺถิวตฺถุ}}
\csromangathameasure{รมณียานิ อรญฺญานิ, ยตฺถ น รมตี ชโน. \\ วีตราคา รมิสฺสนฺติ, น เต กามคเวสิโน. อรหนฺตวคฺโค สตฺตโม.}
\csromangathagroup{\csromangathastanza{\csromanfolio{28}รมณียานิ อรญฺญานิ, ยตฺถ น รมตี ชโน. \\ วีตราคา รมิสฺสนฺติ, น เต กามคเวสิโน.\csromanspacer{} อรหนฺตวคฺโค สตฺตโม.}}

\csromanheader{1}{8. สหสฺสวคฺค}
\vspace{\tipitakaparskip}
\csromancenter{ตมฺพทาฐิกโจรฆาตกวตฺถุ 100.\csromanspacer{} สหสฺสมปิ เจ วาจา, อนตฺถปทสํหิตา.}
\vspace{0.5\baselineskip}
\csromangathameasure{เอกํ อตฺถปทํ เสยฺโย, ยํ สุตฺวา อุปสมฺมติ.}
\csromangathagroup{\csromangathastanza{เอกํ อตฺถปทํ เสยฺโย, ยํ สุตฺวา อุปสมฺมติ.}}

\csromanheader{2}{2. พาหิยทารุจีริยตฺเถรวตฺถุ 101. สหสฺสมปิ เจ คาถา, อนตฺถปทสํหิตา.}
\csromanmatikaanchor{mtk.16}
\csromantocmarkref{section}{6. ปณฺฑิตวคฺค}{mtk.16}
\bookmark[dest={mtk.16},level=1]{6. ปณฺฑิตวคฺค}
\csromangathameasure{เอกํ คาถาปทํ เสยฺโย, ยํ สุตฺวา อุปสมฺมติ.}
\csromangathagroup{\csromangathastanza{เอกํ คาถาปทํ เสยฺโย, ยํ สุตฺวา อุปสมฺมติ.}}

\csromanheader{2}{3. กุณฺฑลเกสิเถรีวตฺถุ 102. โย จ คาถา สตํ ภาเส, อนตฺถปทสํหิตา\footnote{อนตฺถปทสญฺหิตํ (ก) วิเสสนํ เหตํ คาถาติปทสฺส.}.}
\csromangathameasure{เอกํ ธมฺมปทํ เสยฺโย, ยํ สุตฺวา อุปสมฺมติ.}
\csromangathagroup{\csromangathastanza{เอกํ ธมฺมปทํ เสยฺโย, ยํ สุตฺวา อุปสมฺมติ.}}

\proseitem{103}{โย สหสฺสํ สหสฺเสน, สงฺคาเม มานุเส ชิเน.}

\csromangathameasure{เอกญฺจ เชยฺยมตฺตานํ, ส เว สงฺคามชุตฺตโม.}
\csromangathagroup{\csromangathastanza{เอกญฺจ เชยฺยมตฺตานํ\footnote{อตฺตานํ (สี, อิ)}, ส เว สงฺคามชุตฺตโม.}}

\vspace{\tipitakaparskip}
\csromancenter{อนตฺถปุจฺฉกพฺราหฺมณวตฺถุ 104.\csromanspacer{} อตฺตา หเว ชิตํ เสยฺโย, ยา จายํ อิตรา ปชา.}
\vspace{0.5\baselineskip}
\csromangathameasure{อตฺตทนฺตสฺส โปสสฺส, นิจฺจํ สญฺญตจาริโน.}
\csromangathagroup{\csromangathastanza{อตฺตทนฺตสฺส โปสสฺส, นิจฺจํ สญฺญตจาริโน.}}

\proseitem{105}{เนว เทโว น คนฺธพฺโพ, น มาโร สห พฺรหฺมุนา.}

\csromangathameasure{ชิตํ อปชิตํ กยิรา, ตถารูปสฺส ชนฺตุโน.}
\csromangathagroup{\csromangathastanza{ชิตํ อปชิตํ กยิรา, ตถารูปสฺส ชนฺตุโน.}}

\vspace{\tipitakaparskip}
\csromancenter{สหสฺสวคฺค \textbf{5.}\csromanspacer{}\textbf{ สาริปุตฺตตฺเถรสฺส มาตุลพฺราหฺมณวตฺถุ} 106.\csromanspacer{} มาเส มาเส สหสฺเสน, โย ยเชถ สตํ สมํ.}
\vspace{0.5\baselineskip}
\csromangathameasure{เอกญฺจ ภาวิตตฺตานํ, มุหุตฺตมปิ ปูชเย. \\ สาเยว ปูชนา เสยฺโย, ยญฺเจ วสฺสสตํ หุตํ.}
\csromangathagroup{\csromangathastanza{\csromanfolio{29}เอกญฺจ ภาวิตตฺตานํ, มุหุตฺตมปิ ปูชเย. \\ สาเยว ปูชนา เสยฺโย, ยญฺเจ วสฺสสตํ หุตํ.}}

\vspace{\tipitakaparskip}
\csromancenter{สาริปุตฺตตฺเถรสฺส ภาคิเนยฺยวตฺถุ 107.\csromanspacer{} โย จ วสฺสสตํ ชนฺตุ, อคฺคิํ ปริจเร วเน.}
\vspace{0.5\baselineskip}
\csromangathameasure{เอกญฺจ ภาวิตตฺตานํ, มุหุตฺตมปิ ปูชเย. \\ สาเยว ปูชนา เสยฺโย, ยญฺเจ วสฺสสตํ หุตํ. \\ สาริปุตฺตตฺเถรสฺส สหายกพฺราหฺมณวตฺถุ 108. ยํกิญฺจิ ยิฏฺฐํ ว หุตํ ว โลเก, \\ สํวจฺฉรํ ยเชถ ปุญฺญเปกฺโข. \\ สพฺพมฺปิ ตํ น จตุภาคเมติ, \\ อภิวาทนา อุชฺชุคเตสุ เสยฺโย.}
\csromangathagroup{\csromangathastanza{เอกญฺจ ภาวิตตฺตานํ, มุหุตฺตมปิ ปูชเย. \\ สาเยว ปูชนา เสยฺโย, ยญฺเจ วสฺสสตํ หุตํ.}\csromangathastanza{สาริปุตฺตตฺเถรสฺส สหายกพฺราหฺมณวตฺถุ 108.\csromanspacer{} ยํกิญฺจิ ยิฏฺฐํ ว หุตํ ว\footnote{ยิฏฺฐญฺจ หุตญฺจ (ก)} โลเก, \\ สํวจฺฉรํ ยเชถ ปุญฺญเปกฺโข. \\ สพฺพมฺปิ ตํ น จตุภาคเมติ, \\ อภิวาทนา อุชฺชุคเตสุ เสยฺโย.}}

\vspace{\tipitakaparskip}
\csromancenter{อายุวฑฺฒนกุมารวตฺถุ 109.\csromanspacer{} อภิวาทนสีลิสฺส, นิจฺจํ วุฑฺฒาปจายิโน\footnote{วทฺธาปจายิโน (สี, อิ)}.}
\vspace{0.5\baselineskip}
\csromangathameasure{จตฺตาโร ธมฺมา วฑฺฒนฺติ, อายุ วณฺโณ สุขํ พลํ.}
\csromangathagroup{\csromangathastanza{จตฺตาโร ธมฺมา วฑฺฒนฺติ, อายุ วณฺโณ สุขํ พลํ.}}

\vspace{\tipitakaparskip}
\csromancenter{สํกิจฺจสามเณรวตฺถุ 110.\csromanspacer{} โย จ วสฺสสตํ ชีเว, ทุสฺสีโล อสมาหิโต.}
\vspace{0.5\baselineskip}
\csromangathameasure{เอกาหํ ชีวิตํ เสยฺโย, สีลวนฺตสฺส ฌายิโน.}
\csromangathagroup{\csromangathastanza{เอกาหํ ชีวิตํ เสยฺโย, สีลวนฺตสฺส ฌายิโน.}}

\vspace{\tipitakaparskip}
\csromancenter{ขาณุโกณฺฑญฺญตฺเถรวตฺถุ 111.\csromanspacer{} โย จ วสฺสสตํ ชีเว, ทุปฺปญฺโญ อสมาหิโต.}
\vspace{0.5\baselineskip}
\csromangathameasure{เอกาหํ ชีวิตํ เสยฺโย, ปญฺญวนฺตสฺส ฌายิโน.}
\csromangathagroup{\csromangathastanza{เอกาหํ ชีวิตํ เสยฺโย, ปญฺญวนฺตสฺส ฌายิโน.}}

\vspace{\tipitakaparskip}
\csromancenter{สพฺพทาสตฺเถรวตฺถุ 112.\csromanspacer{} โย จ วสฺสสตํ ชีเว, กุสีโต หีนวีริโย.}
\vspace{0.5\baselineskip}
\csromangathameasure{เอกาหํ ชีวิตํ เสยฺโย, วีริยมารภโต ทฬฺหํ.}
\csromangathagroup{\csromangathastanza{เอกาหํ ชีวิตํ เสยฺโย, วีริยมารภโต ทฬฺหํ.}}

\vspace{\tipitakaparskip}
\csromancenter{ธมฺมปทปาฬิ \textbf{12.}\csromanspacer{}\textbf{ ปฏาจาราเถรีวตฺถุ} 113.\csromanspacer{} โย จ วสฺสสตํ ชีเว, อปสฺสํ อุทยพฺพยํ.}
\vspace{0.5\baselineskip}
\csromangathameasure{เอกาหํ ชีวิตํ เสยฺโย, ปสฺสโต อุทยพฺพยํ.}
\csromangathagroup{\csromangathastanza{\csromanfolio{30}เอกาหํ ชีวิตํ เสยฺโย, ปสฺสโต อุทยพฺพยํ.}}

\vspace{\tipitakaparskip}
\csromancenter{กิสาโคตมีวตฺถุ 114.\csromanspacer{} โย จ วสฺสสตํ ชีเว, อปสฺสํ อมตํ ปทํ.}
\vspace{0.5\baselineskip}
\csromangathameasure{เอกาหํ ชีวิตํ เสยฺโย, ปสฺสโต อมตํ ปทํ.}
\csromangathagroup{\csromangathastanza{เอกาหํ ชีวิตํ เสยฺโย, ปสฺสโต อมตํ ปทํ.}}

\vspace{\tipitakaparskip}
\csromancenter{พหุปุตฺติกตฺเถรีวตฺถุ 115.\csromanspacer{} โย จ วสฺสสตํ ชีเว, อปสฺสํ ธมฺมมุตฺตมํ.}
\vspace{0.5\baselineskip}
\csromangathameasure{เอกาหํ ชีวิตํ เสยฺโย, ปสฺสโต ธมฺมมุตฺตมํ. สหสฺสวคฺโค อฏฺฐโม.}
\csromangathagroup{\csromangathastanza{เอกาหํ ชีวิตํ เสยฺโย, ปสฺสโต ธมฺมมุตฺตมํ.\csromanspacer{} สหสฺสวคฺโค อฏฺฐโม.}}

\csromanheader{1}{9. ปาปวคฺค}
\vspace{\tipitakaparskip}
\csromancenter{จูเฬกสาฏกพฺราหฺมณวตฺถุ 116.\csromanspacer{} อภิตฺถเรถ กลฺยาเณ, ปาปา จิตฺตํ นิวารเย.}
\vspace{0.5\baselineskip}
\csromangathameasure{ทนฺธํ หิ กโรโต ปุญฺญํ, ปาปสฺมิํ รมตี มโน.}
\csromangathagroup{\csromangathastanza{ทนฺธํ หิ กโรโต ปุญฺญํ, ปาปสฺมิํ รมตี มโน.}}

\vspace{\tipitakaparskip}
\csromancenter{เสยฺยสกตฺเถรวตฺถุ 117.\csromanspacer{} ปาปญฺเจ ปุริโส กยิรา, น นํ\footnote{น ตํ (สี, อิ)} กยิรา ปุนปฺปุนํ.}
\vspace{0.5\baselineskip}
\csromangathameasure{น ตมฺหิ ฉนฺทํ กยิราถ, ทุกฺโข ปาปสฺส อุจฺจโย.}
\csromangathagroup{\csromangathastanza{น ตมฺหิ ฉนฺทํ กยิราถ, ทุกฺโข ปาปสฺส อุจฺจโย.}}

\vspace{\tipitakaparskip}
\csromancenter{ลาชเทวธีตาวตฺถุ 118.\csromanspacer{} ปุญฺญญฺเจ ปุริโส กยิรา, กยิรา นํ\footnote{กยิราเถตํ (สี, สฺยา), กยิราเถนํ (อิ)} ปุนปฺปุนํ.}
\vspace{0.5\baselineskip}
\csromangathameasure{ตมฺหิ ฉนฺทํ กยิราถ, สุโข ปุญฺญสฺส อุจฺจโย.}
\csromangathagroup{\csromangathastanza{ตมฺหิ ฉนฺทํ กยิราถ, สุโข ปุญฺญสฺส อุจฺจโย.}}

\vspace{\tipitakaparskip}
\csromancenter{ปาปวคฺค \textbf{4.}\csromanspacer{}\textbf{ อนาถปิณฺฑิกเสฏฺฐิวตฺถุ} 119.\csromanspacer{} ปาโปปิ ปสฺสติ ภทฺรํ, ยาว ปาปํ น ปจฺจติ.}
\vspace{0.5\baselineskip}
\csromanmatikaanchor{mtk.17}
\csromantocmarkref{section}{7. อรหนฺตวคฺค}{mtk.17}
\bookmark[dest={mtk.17},level=1]{7. อรหนฺตวคฺค}
\csromangathameasure{ยทา จ ปจฺจติ ปาปํ, อถ ปาโป ปาปานิ ปสฺสติ. \\ ภโทฺรปิ ปสฺสติ ปาปํ, ยาว ภทฺรํ น ปจฺจติ. \\ ยทา จ ปจฺจติ ภทฺรํ, อถ ภโทฺร ภทฺรานิ ปสฺสติ.}
\csromangathagroup{\csromangathastanza{\csromanfolio{31}ยทา จ ปจฺจติ ปาปํ, อถ ปาโป ปาปานิ\footnote{อถ ปาปานิ (?)} ปสฺสติ. \\ ภโทฺรปิ ปสฺสติ ปาปํ, ยาว ภทฺรํ น ปจฺจติ. \\ ยทา จ ปจฺจติ ภทฺรํ, อถ ภโทฺร ภทฺรานิ\footnote{อถ ปาปานิ (?)} ปสฺสติ.}}

\vspace{\tipitakaparskip}
\csromancenter{อสญฺญตปริกฺขารวตฺถุ 121.\csromanspacer{} มาวมญฺเญถ\footnote{มาปฺปมญฺเญถ (สี, สฺยา, อิ)} ปาปสฺส, น มนฺทํ\footnote{น มํ ตํ (สี, อิ), น มตฺตํ (สฺยา)} อาคมิสฺสติ.}
\vspace{0.5\baselineskip}
\csromangathameasure{อุทพินฺทุนิปาเตน, อุทกุมฺโภปิ ปูรติ. \\ พาโล ปูรติ ปาปสฺส, โถกํ โถกมฺปิ อาจินํ.}
\csromangathagroup{\csromangathastanza{อุทพินฺทุนิปาเตน, อุทกุมฺโภปิ ปูรติ. \\ พาโล ปูรติ\footnote{ปูรติ พาโล (สี, ก), อาปูรติ พาโล (สฺยา)} ปาปสฺส, โถกํ โถกมฺปิ\footnote{โถก โถกมฺปิ (สี, อิ)} อาจินํ.}}

\vspace{\tipitakaparskip}
\csromancenter{พิฬาลปาทกเสฏฺฐิวตฺถุ 122.\csromanspacer{} มาวมญฺเญถ ปุญฺญสฺส, น มนฺทํ อาคมิสฺสติ.}
\vspace{0.5\baselineskip}
\csromangathameasure{อุทพินฺทุนิปาเตน, อุทกุมฺโภปิ ปูรติ. \\ ธีโร ปูรติ ปุญฺญสฺส, โถกํ โถกมฺปิ อาจินํ.}
\csromangathagroup{\csromangathastanza{อุทพินฺทุนิปาเตน, อุทกุมฺโภปิ ปูรติ. \\ ธีโร ปูรติ ปุญฺญสฺส, โถกํ โถกมฺปิ อาจินํ.}}

\vspace{\tipitakaparskip}
\csromancenter{มหาธนวาณิชวตฺถุ 123.\csromanspacer{} วาณิโชว ภยํ มคฺคํ, อปฺปสตฺโถ มหทฺธโน.}
\vspace{0.5\baselineskip}
\csromangathameasure{วิสํ ชีวิตุกาโมว, ปาปานิ ปริวชฺชเย.}
\csromangathagroup{\csromangathastanza{วิสํ ชีวิตุกาโมว, ปาปานิ ปริวชฺชเย.}}

\vspace{\tipitakaparskip}
\csromancenter{กุกฺกุฏมิตฺตเนสาทวตฺถุ 124.\csromanspacer{} ปาณิมฺหิ เจ วโณ นาสฺส, หเรยฺย ปาณินา วิสํ.}
\vspace{0.5\baselineskip}
\csromangathameasure{นาพฺพณํ วิสมเนฺวติ, นตฺถิ ปาปํ อกุพฺพโต.}
\csromangathagroup{\csromangathastanza{นาพฺพณํ วิสมเนฺวติ, นตฺถิ ปาปํ อกุพฺพโต.}}

\vspace{\tipitakaparskip}
\csromancenter{ธมฺมปทปาฬิ \textbf{9.}\csromanspacer{}\textbf{ โกกสุนขลุทฺทกวตฺถุ} 125. * โย อปฺปทุฏฺฐสฺส นรสฺส ทุสฺสติ,}
\prose{\csromanfolio{32}สุทฺธสฺส โปสสฺส อนงฺคณสฺส.}

\csromangathameasure{ตเมว พาลํ ปจฺเจติ ปาปํ, \\ สุขุโม รโช ปฏิวาตํว ขิตฺโต.}
\csromangathagroup{\csromangathastanza{ตเมว พาลํ ปจฺเจติ ปาปํ, \\ สุขุโม รโช ปฏิวาตํว ขิตฺโต.}}

\vspace{\tipitakaparskip}
\csromancenter{มณิการกุลูปกติสฺสตฺเถรวตฺถุ 126.\csromanspacer{} คพฺภเมเก อุปฺปชฺชนฺติ, นิรยํ ปาปกมฺมิโน.}
\vspace{0.5\baselineskip}
\csromangathameasure{สคฺคํ สุคติโน ยนฺติ, ปรินิพฺพนฺติ อนาสวา.}
\csromangathagroup{\csromangathastanza{สคฺคํ สุคติโน ยนฺติ, ปรินิพฺพนฺติ อนาสวา.}}

\vspace{\tipitakaparskip}
\csromancenter{ตโยชนวตฺถุ 127.\csromanspacer{} น อนฺตลิกฺเข น สมุทฺทมชฺเฌ,}
\prose{น ปพฺพตานํ วิวรํ ปวิสฺส\footnote{ปวิสํ (สฺยา)}.}

\csromangathameasure{น วิชฺชตี โส ชคติปฺปเทโส, \\ ยตฺถฏฺฐิโต มุจฺเจยฺย ปาปกมฺมา.}
\csromangathagroup{\csromangathastanza{น วิชฺชตี\footnote{น วิชฺชติ (ก-สี, อิ, ก)} โส ชคติปฺปเทโส, \\ ยตฺถฏฺฐิโต\footnote{น วิชฺชติ (ก-สี, อิ, ก)} มุจฺเจยฺย ปาปกมฺมา.}}

\vspace{\tipitakaparskip}
\csromancenter{สุปฺปพุทฺธสกฺยวตฺถุ 128.\csromanspacer{} น อนฺตลิกฺเข น สมุทฺทมชฺเฌ,}
\prose{น ปพฺพตานํ วิวรํ ปวิสฺส.}

\csromangathameasure{น วิชฺชตี โส ชคติปฺปเทโส, \\ ยตฺถฏฺฐิตํ นปฺปสเหยฺย มจฺจุ. ปาปวคฺโค นวโม.}
\csromangathagroup{\csromangathastanza{น วิชฺชตี โส ชคติปฺปเทโส, \\ ยตฺถฏฺฐิตํ\footnote{ยตฺรฏฺฐิตํ (สฺยา)} นปฺปสเหยฺย มจฺจุ.\csromanspacer{} ปาปวคฺโค นวโม.}}

\csromanheader{1}{10. ทณฺฑวคฺค}
\vspace{\tipitakaparskip}
\csromancenter{ฉพฺพคฺคิยภิกฺขุวตฺถุ 129.\csromanspacer{} สพฺเพ ตสนฺติ ทณฺฑสฺส, สพฺเพ ภายนฺติ มจฺจุโน.}
\vspace{0.5\baselineskip}
\csromangathameasure{อตฺตานํ อุปมํ กตฺวา, น หเนยฺย น ฆาตเย.}
\csromangathagroup{\csromangathastanza{อตฺตานํ อุปมํ กตฺวา, น หเนยฺย น ฆาตเย.}}

\vspace{\tipitakaparskip}
\csromancenter{ทณฺฑวคฺค \textbf{2.}\csromanspacer{}\textbf{ ฉพฺพคฺคิยวตฺถุ} 130.\csromanspacer{} สพฺเพ ตสนฺติ ทณฺฑสฺส, สพฺเพสํ ชีวิตํ ปิยํ.}
\vspace{0.5\baselineskip}
\csromangathameasure{อตฺตานํ อุปมํ กตฺวา, น หเนยฺย น ฆาตเย.}
\csromangathagroup{\csromangathastanza{\csromanfolio{33}อตฺตานํ อุปมํ กตฺวา, น หเนยฺย น ฆาตเย.}}

\vspace{\tipitakaparskip}
\csromancenter{สมฺพหุลกุมารกวตฺถุ 131. * สุขกามานิ ภูตานิ, โย ทณฺเฑน วิหิํสติ.}
\vspace{0.5\baselineskip}
\csromangathameasure{อตฺตโน สุขเมสาโน, เปจฺจ โส น ลภเต สุขํ.}
\csromangathagroup{\csromangathastanza{อตฺตโน สุขเมสาโน, เปจฺจ โส น ลภเต สุขํ.}}

\proseitem{132}{สุขกามานิ ภูตานิ, โย ทณฺเฑน น หิํสติ.}

\csromanmatikaanchor{mtk.18}
\csromantocmarkref{section}{8. สหสฺสวคฺค}{mtk.18}
\bookmark[dest={mtk.18},level=1]{8. สหสฺสวคฺค}
\csromangathameasure{อตฺตโน สุขเมสาโน, เปจฺจ โส ลภเต สุขํ.}
\csromangathagroup{\csromangathastanza{อตฺตโน สุขเมสาโน, เปจฺจ โส ลภเต สุขํ.}}

\vspace{\tipitakaparskip}
\csromancenter{โกณฺฑธานตฺเถรวตฺถุ 133.\csromanspacer{} มาโวจ ผรุสํ กญฺจิ, วุตฺตา ปฏิวเทยฺยุ ตํ\footnote{ปฏิวเทยฺยุํ ตํ (ก)}.}
\vspace{0.5\baselineskip}
\csromangathameasure{ทุกฺขา หิ สารมฺภกถา, ปฏิทณฺฑา ผุเสยฺยุ ตํ. \\ สเจ เนเรสิ อตฺตานํ, กํโส อุปหโต ยถา. \\ เอส ปตฺโตสิ นิพฺพานํ, สารมฺโภ เต น วิชฺชติ.}
\csromangathagroup{\csromangathastanza{ทุกฺขา หิ สารมฺภกถา, ปฏิทณฺฑา ผุเสยฺยุ ตํ\footnote{ผุเสยฺยุํ ตํ (ก)}. \\ สเจ เนเรสิ อตฺตานํ, กํโส อุปหโต ยถา. \\ เอส ปตฺโตสิ นิพฺพานํ, สารมฺโภ เต น วิชฺชติ.}}

\vspace{\tipitakaparskip}
\csromancenter{อุโปสถิก-อิตฺถีนํวตฺถุ 135.\csromanspacer{} ยถา ทณฺเฑน โคปาโล, คาโว ปาเชติ โคจรํ.}
\vspace{0.5\baselineskip}
\csromangathameasure{เอวํ ชรา จ มจฺจุ จ, อายุํ ปาเชนฺติ ปาณินํ.}
\csromangathagroup{\csromangathastanza{เอวํ ชรา จ มจฺจุ จ, อายุํ ปาเชนฺติ ปาณินํ.}}

\vspace{\tipitakaparskip}
\csromancenter{อชครเปตวตฺถุ 136.\csromanspacer{} อถ ปาปานิ กมฺมานิ, กรํ พาโล น พุชฺฌติ.}
\vspace{0.5\baselineskip}
\csromangathameasure{เสหิ กมฺเมหิ ทุมฺเมโธ, อคฺคิฑฑฺโฒว ตปฺปติ.}
\csromangathagroup{\csromangathastanza{เสหิ กมฺเมหิ ทุมฺเมโธ, อคฺคิฑฑฺโฒว ตปฺปติ.}}

\vspace{\tipitakaparskip}
\csromancenter{มหาโมคฺคลฺลานตฺเถรวตฺถุ 137.\csromanspacer{} โย ทณฺเฑน อทณฺเฑสุ, อปฺปทุฏฺเฐสุ ทุสฺสติ.}
\vspace{0.5\baselineskip}
\csromangathameasure{ทสนฺนมญฺญตรํ ฐานํ, ขิปฺปเมว นิคจฺฉติ.}
\csromangathagroup{\csromangathastanza{ทสนฺนมญฺญตรํ ฐานํ, ขิปฺปเมว นิคจฺฉติ.}}

\csromanheader{2}{ธมฺมปทปาฬิ 138. เวทนํ ผรุสํ ชานิํ, สรีรสฺส ว เภทนํ\footnote{สรีรสฺส ปเภทนํ (สฺยา)}.}
\csromangathameasure{ครุกํ วาปิ อาพาธํ, จิตฺตกฺเขปํ ว ปาปุเณ. \\ ราชโต วา อุปสคฺคํ, อพฺภกฺขานํ ว ทารุณํ. \\ ปริกฺขยํ ว ญาตีนํ, โภคานํ ว ปภงฺคุรํ. \\ อถ วาสฺส อคารานิ, อคฺคิ ฑหติ ปาวโก. \\ กายสฺส เภทา ทุปฺปญฺโญ, นิรยํ โสปปชฺชติ.}
\csromangathagroup{\csromangathastanza{\csromanfolio{34}ครุกํ วาปิ อาพาธํ, จิตฺตกฺเขปํ ว\footnote{จิตฺตกฺเขปญฺจ (ก)} ปาปุเณ. \\ ราชโต วา อุปสคฺคํ\footnote{จิตฺตกฺเขปญฺจ (ก)}, อพฺภกฺขานํ ว\footnote{จิตฺตกฺเขปญฺจ (ก)} ทารุณํ.}\csromangathastanza{ปริกฺขยํ ว\footnote{ปริกฺขยญฺจ (ก)} ญาตีนํ, โภคานํ ว\footnote{โภคานญฺจ (ก)} ปภงฺคุรํ\footnote{ปภงฺคุนํ (ก)}. \\ อถ วาสฺส อคารานิ, อคฺคิ ฑหติ\footnote{ปริกฺขยญฺจ (ก)} ปาวโก. \\ กายสฺส เภทา ทุปฺปญฺโญ, นิรยํ โสปปชฺชติ\footnote{ปริกฺขยญฺจ (ก)}.}}

\vspace{\tipitakaparskip}
\csromancenter{พหุภณฺฑิกภิกฺขุวตฺถุ 141.\csromanspacer{} น นคฺคจริยา น ชฏา น ปงฺกา,}
\prose{นานาสกา ถณฺฑิลสายิกา วา.}

\csromangathameasure{รโชชลฺลํ อุกฺกุฏิกปฺปธานํ, \\ โสเธนฺติ มจฺจํ อวิติณฺณกงฺขํ.}
\csromangathagroup{\csromangathastanza{รโชชลฺลํ อุกฺกุฏิกปฺปธานํ, \\ โสเธนฺติ มจฺจํ อวิติณฺณกงฺขํ.}}

\vspace{\tipitakaparskip}
\csromancenter{สนฺตติมหามตฺตวตฺถุ 142.\csromanspacer{} อลงฺกโต เจปิ สมํ จเรยฺย,}
\prose{สนฺโต ทนฺโต นิยโต พฺรหฺมจารี.}

\csromangathameasure{สพฺเพสุ ภูเตสุ นิธาย ทณฺฑํ, \\ โส พฺราหฺมโณ โส สมโณ ส ภิกฺขุ.}
\csromangathagroup{\csromangathastanza{สพฺเพสุ ภูเตสุ นิธาย ทณฺฑํ, \\ โส พฺราหฺมโณ โส สมโณ ส ภิกฺขุ.}}

\vspace{\tipitakaparskip}
\csromancenter{ปิโลติกติสฺสตฺเถรวตฺถุ 143.\csromanspacer{} หิรีนิเสโธ ปุริโส, โกจิ โลกสฺมิ วิชฺชติ.}
\vspace{0.5\baselineskip}
\csromangathameasure{โย นิทฺทํ อปโพเธติ, อสฺโส ภโทฺร กสามิว.}
\csromangathagroup{\csromangathastanza{โย นิทฺทํ\footnote{นินฺทํ (สี, อิ) สํ 1. 8 ปิฏฺเฐปิ.} อปโพเธติ\footnote{อปโพธติ (สี, สฺยา, อิ)}, อสฺโส ภโทฺร กสามิว.}}

\csromanheader{1}{11. ชราวคฺค 144. อสฺโส ยถา ภโทฺร กสานิวิฏฺโฐ,}
\prose{\csromanfolio{35}อาตาปิโน สํเวคิโน ภวาถ.}

\csromangathameasure{สทฺธาย สีเลน จ วีริเยน จ, \\ สมาธินา ธมฺมวินิจฺฉเยน จ. \\ สมฺปนฺนวิชฺชาจรณา ปติสฺสตา, \\ ชหิสฺสถ ทุกฺขมิทํ อนปฺปกํ.}
\csromangathagroup{\csromangathastanza{สทฺธาย สีเลน จ วีริเยน จ, \\ สมาธินา ธมฺมวินิจฺฉเยน จ. \\ สมฺปนฺนวิชฺชาจรณา ปติสฺสตา, \\ ชหิสฺสถ\footnote{ปหสฺสถ (สี, สฺยา, อิ)} ทุกฺขมิทํ อนปฺปกํ.}}

\vspace{\tipitakaparskip}
\csromancenter{สุขสามเณรวตฺถุ 145. * อุทกํ หิ นยนฺติ เนตฺติกา, อุสุการา นมยนฺติ เตชนํ.}
\vspace{0.5\baselineskip}
\csromangathameasure{ทารุํ นมยนฺติ ตจฺฉกา, อตฺตานํ ทมยนฺติ สุพฺพตา. ทณฺฑวคฺโค ทสโม.}
\csromangathagroup{\csromangathastanza{ทารุํ นมยนฺติ ตจฺฉกา, อตฺตานํ ทมยนฺติ สุพฺพตา.\csromanspacer{} ทณฺฑวคฺโค ทสโม.}}

\chapterhead{11. ชราวคฺค}
\vspace{\tipitakaparskip}
\csromancenter{วิสาขาย สหายิกานํ วตฺถุ 146.\csromanspacer{} โก นุ หาโส\footnote{กินฺนุหาโส (ก)} กิมานนฺโท, นิจฺจํ ปชฺชลิเต สติ.}
\csromancenter{อนฺธกาเรน โอนทฺธา, ปทีปํ น คเวสถ,}
\csromancenter{สิริมาวตฺถุ 147.\csromanspacer{} ปสฺส จิตฺตกตํ พิมฺพํ, อรุกายํ สมุสฺสิตํ.}
\vspace{0.5\baselineskip}
\csromangathameasure{อาตุรํ พหุสงฺกปฺปํ, ยสฺส นตฺถิ ธุวํ ฐิติ.}
\csromangathagroup{\csromangathastanza{อาตุรํ พหุสงฺกปฺปํ, ยสฺส นตฺถิ ธุวํ ฐิติ.}}

\vspace{\tipitakaparskip}
\csromancenter{อุตฺตราเถรีวตฺถุ 148.\csromanspacer{} ปริชิณฺณมิทํ รูปํ, โรคนีฬํ\footnote{โรคนิฑฺฒํ (สี, อิ), โรคนิทฺธํ (สฺยา)} ปภงฺคุรํ.}
\vspace{0.5\baselineskip}
\csromangathameasure{ภิชฺชติ ปูติสนฺเทโห, มรณนฺตํ หิ ชีวิตํ.}
\csromangathagroup{\csromangathastanza{ภิชฺชติ ปูติสนฺเทโห, มรณนฺตํ หิ ชีวิตํ.}}

\csromanmatikaanchor{mtk.19}
\csromantocmarkref{section}{9. ปาปวคฺค}{mtk.19}
\bookmark[dest={mtk.19},level=1]{9. ปาปวคฺค}
\vspace{\tipitakaparskip}
\csromancenter{ธมฺมปทปาฬิ \textbf{4.}\csromanspacer{}\textbf{ สมฺพหุล-อธิมานิกภิกฺขุวตฺถุ} 149.\csromanspacer{} ยานิมานิ อปตฺถานิ\footnote{อปตฺตานิ (ก), ยานิมานิ’ปวิทฺธานิ (?)}, อลาพูเนว\footnote{อลาปูเนว (สี, สฺยา, อิ)} สารเท.}
\vspace{0.5\baselineskip}
\csromangathameasure{กาโปตกานิ อฏฺฐีนิ, ตานิ ทิสฺวาน กา รติ. \\ ชนปทกลฺยาณีรูปนนฺทาเถรีวตฺถุ 150. อฏฺฐีนํ นครํ กตํ, มํสโลหิตเลปนํ. \\ ยตฺถ ชรา จ มจฺจุ จ, มาโน มกฺโข จ โอหิโต.}
\csromangathagroup{\csromangathastanza{\csromanfolio{36}กาโปตกานิ อฏฺฐีนิ, ตานิ ทิสฺวาน กา รติ. \\ ชนปทกลฺยาณีรูปนนฺทาเถรีวตฺถุ 150.\csromanspacer{} อฏฺฐีนํ นครํ กตํ, มํสโลหิตเลปนํ. \\ ยตฺถ ชรา จ มจฺจุ จ, มาโน มกฺโข จ โอหิโต.}}

\vspace{\tipitakaparskip}
\csromancenter{มลฺลิกาเทวีวตฺถุ 151.\csromanspacer{} ชีรนฺติ เว ราชรถา สุจิตฺตา,}
\prose{อโถ สรีรมฺปิ ชรํ อุเปติ.}

\csromangathameasure{สตญฺจ ธมฺโม น ชรํ อุเปติ, \\ สนฺโต หเว สพฺภิ ปเวทยนฺติ.}
\csromangathagroup{\csromangathastanza{สตญฺจ ธมฺโม น ชรํ อุเปติ, \\ สนฺโต หเว สพฺภิ ปเวทยนฺติ.}}

\vspace{\tipitakaparskip}
\csromancenter{ลาฬุทายีเถรวตฺถุ 152.\csromanspacer{} อปฺปสฺสุตายํ ปุริโส, พลีพทฺโทว\footnote{พลิวทฺโทว (สี, สฺยา, อิ) พลิพทฺโธว (ก)} ชีรติ.}
\vspace{0.5\baselineskip}
\csromangathameasure{มํสานิ ตสฺส วฑฺฒนฺติ, ปญฺญา ตสฺส น วฑฺฒติ.}
\csromangathagroup{\csromangathastanza{มํสานิ ตสฺส วฑฺฒนฺติ, ปญฺญา ตสฺส น วฑฺฒติ.}}

\vspace{\tipitakaparskip}
\csromancenter{อุทานวตฺถุ 153.\csromanspacer{} อเนกชาติสํสารํ, สนฺธาวิสฺสํ อนิพฺพิสํ.}
\vspace{0.5\baselineskip}
\csromangathameasure{คหการํ คเวสนฺโต, ทุกฺขา ชาติ ปุนปฺปุนํ. \\ คหการก ทิฏฺโฐสิ, ปุน เคหํ น กาหสิ. \\ สพฺพา เต ผาสุกา ภคฺคา, คหกูฏํ วิสงฺขตํ. \\ วิสงฺขารคตํ จิตฺตํ, ตณฺหานํ ขยมชฺฌคา.}
\csromangathagroup{\csromangathastanza{คหการํ\footnote{คหการกํ (สี, สฺยา, อิ)} คเวสนฺโต, ทุกฺขา ชาติ ปุนปฺปุนํ. \\ คหการก ทิฏฺโฐสิ, ปุน เคหํ น กาหสิ.}\csromangathastanza{สพฺพา เต ผาสุกา ภคฺคา, คหกูฏํ วิสงฺขตํ. \\ วิสงฺขารคตํ จิตฺตํ, ตณฺหานํ ขยมชฺฌคา.}}

\vspace{\tipitakaparskip}
\csromancenter{มหาธนเสฏฺฐิปุตฺตวตฺถุ 155. * อจริตฺวา พฺรหฺมจริยํ, อลทฺธา โยพฺพเน ธนํ.}
\vspace{0.5\baselineskip}
\csromangathameasure{ชิณฺณโกญฺจาว ฌายนฺติ, ขีณมจฺเฉว ปลฺลเล.}
\csromangathagroup{\csromangathastanza{ชิณฺณโกญฺจาว ฌายนฺติ, ขีณมจฺเฉว ปลฺลเล.}}

\csromanheader{1}{12. อตฺตวคฺค 156. อจริตฺวา พฺรหฺมจริยํ, อลทฺธา โยพฺพเน ธนํ.}
\csromangathameasure{เสนฺติ จาปาติขีณาว, ปุราณานิ อนุตฺถุนํ. ชราวคฺโค เอกาทสโม.}
\csromangathagroup{\csromangathastanza{\csromanfolio{37}เสนฺติ จาปาติขีณาว, ปุราณานิ อนุตฺถุนํ.\csromanspacer{} ชราวคฺโค เอกาทสโม.}}

\chapterhead{12. อตฺตวคฺค}
\vspace{\tipitakaparskip}
\csromancenter{โพธิราชกุมารวตฺถุ 157.\csromanspacer{} อตฺตานญฺเจ ปิยํ ชญฺญา, รกฺเขยฺย นํ สุรกฺขิตํ.}
\vspace{0.5\baselineskip}
\csromangathameasure{ติณฺณํ อญฺญตรํ ยามํ, ปฏิชคฺเคยฺย ปณฺฑิโต.}
\csromangathagroup{\csromangathastanza{ติณฺณํ อญฺญตรํ ยามํ, ปฏิชคฺเคยฺย ปณฺฑิโต.}}

\vspace{\tipitakaparskip}
\csromancenter{อุปนนฺทสกฺยปุตฺตตฺเถรวตฺถุ 158.\csromanspacer{} อตฺตานเมว ปฐมํ, ปติรูเป นิเวสเย.}
\vspace{0.5\baselineskip}
\csromangathameasure{อถญฺญมนุสาเสยฺย, น กิลิสฺเสยฺย ปณฺฑิโต.}
\csromangathagroup{\csromangathastanza{อถญฺญมนุสาเสยฺย, น กิลิสฺเสยฺย ปณฺฑิโต.}}

\vspace{\tipitakaparskip}
\csromancenter{ปธานิกติสฺสตฺเถรวตฺถุ 159.\csromanspacer{} อตฺตานญฺเจ ตถา กยิรา, ยถาญฺญมนุสาสติ.}
\vspace{0.5\baselineskip}
\csromangathameasure{สุทนฺโต วต ทเมถ, อตฺตา หิ กิร ทุทฺทโม.}
\csromangathagroup{\csromangathastanza{สุทนฺโต วต ทเมถ, อตฺตา หิ กิร ทุทฺทโม.}}

\vspace{\tipitakaparskip}
\csromancenter{กุมารกสฺสปมาตุตฺเถริวตฺถุ 160.\csromanspacer{} อตฺตา หิ อตฺตโน นาโถ, โก หิ นาโถ ปโร สิยา.}
\vspace{0.5\baselineskip}
\csromangathameasure{อตฺตนา หิ สุทนฺเตน, นาถํ ลภติ ทุลฺลภํ.}
\csromangathagroup{\csromangathastanza{อตฺตนา หิ สุทนฺเตน, นาถํ ลภติ ทุลฺลภํ.}}

\vspace{\tipitakaparskip}
\csromancenter{มหากาล-อุปาสกวตฺถุ 161.\csromanspacer{} อตฺตนา หิ กตํ ปาปํ, อตฺตชํ อตฺตสมฺภวํ.}
\vspace{0.5\baselineskip}
\csromangathameasure{อภิมตฺถติ ทุมฺเมธํ, วชิรํว’สฺมมยํ มณิํ.}
\csromangathagroup{\csromangathastanza{อภิมตฺถติ\footnote{อภิมนฺถติ (สี, อิ) ขุ 10. 158 ปิฏฺเฐปิ.} ทุมฺเมธํ, วชิรํว’สฺมมยํ\footnote{วชิรํว’มฺห มยํ (สฺยา, ก)} มณิํ.}}

\vspace{\tipitakaparskip}
\csromancenter{เทวทตฺตวตฺถุ 162.\csromanspacer{} ยสฺส อจฺจนฺตทุสฺสิลฺยํ, มาลุวา สาลมิโวตฺถตํ.}
\vspace{0.5\baselineskip}
\csromangathameasure{กโรติ โส ตถ’ตฺตานํ, ยถา นํ อิจฺฉตี ทิโส.}
\csromangathagroup{\csromangathastanza{กโรติ โส ตถ’ตฺตานํ, ยถา นํ อิจฺฉตี ทิโส.}}

\vspace{\tipitakaparskip}
\csromancenter{ธมฺมปทปาฬิ \textbf{7.}\csromanspacer{}\textbf{ สํฆเภทปริสกฺกนวตฺถุ} 163.\csromanspacer{} สุกรานิ อสาธูนิ, อตฺตโน อหิตานิ จ.}
\vspace{0.5\baselineskip}
\csromangathameasure{ยํ เว หิตญฺจ สาธุญฺจ, ตํ เว ปรมทุกฺกรํ.}
\csromangathagroup{\csromangathastanza{\csromanfolio{38}ยํ เว หิตญฺจ สาธุญฺจ, ตํ เว ปรมทุกฺกรํ.}}

\vspace{\tipitakaparskip}
\csromancenter{กาลตฺเถรวตฺถุ 164.\csromanspacer{} โย สาสนํ อรหตํ, อริยานํ ธมฺมชีวินํ.}
\vspace{0.5\baselineskip}
\csromangathameasure{ปฏิกฺโกสติ ทุมฺเมโธ, ทิฏฺฐิํ นิสฺสาย ปาปิกํ. \\ ผลานิ กฏฺฐกสฺเสว, อตฺตฆาตาย ผลฺลติ.}
\csromangathagroup{\csromangathastanza{ปฏิกฺโกสติ ทุมฺเมโธ, ทิฏฺฐิํ นิสฺสาย ปาปิกํ. \\ ผลานิ กฏฺฐกสฺเสว, อตฺตฆาตาย\footnote{อตฺตฆญฺญาย (สี, สฺยา, อิ)} ผลฺลติ.}}

\csromanmatikaanchor{mtk.20}
\csromantocmarkref{section}{10. ทณฺฑวคฺค}{mtk.20}
\bookmark[dest={mtk.20},level=1]{10. ทณฺฑวคฺค}
\vspace{\tipitakaparskip}
\csromancenter{จูฬกาล-อุปาสกวตฺถุ 165.\csromanspacer{} อตฺตนา หิ\footnote{อตฺตนาว (สี, สฺยา, อิ) ขุ 8. 95 ปิฏฺเฐปิ.} กตํ ปาปํ, อตฺตนา สํกิลิสฺสติ.}
\vspace{0.5\baselineskip}
\csromangathameasure{อตฺตนา อกตํ ปาปํ, อตฺตนาว วิสุชฺฌติ. \\ สุทฺธี อสุทฺธิ ปจฺจตฺตํ, นาญฺโญ อญฺญํ วิโสธเย.}
\csromangathagroup{\csromangathastanza{อตฺตนา อกตํ ปาปํ, อตฺตนาว วิสุชฺฌติ. \\ สุทฺธี อสุทฺธิ ปจฺจตฺตํ, นาญฺโญ อญฺญํ\footnote{นาญฺญมญฺโญ (สี)} วิโสธเย.}}

\vspace{\tipitakaparskip}
\csromancenter{อตฺตทตฺถตฺเถรวตฺถุ 166.\csromanspacer{} อตฺตทตฺถํ ปรตฺเถน, พหุนาปิ น หาปเย.}
\vspace{0.5\baselineskip}
\csromangathameasure{อตฺตทตฺถมภิญฺญาย, สทตฺถปสุโต สิยา. อตฺตวคฺโค ทฺวาทสโม.}
\csromangathagroup{\csromangathastanza{อตฺตทตฺถมภิญฺญาย, สทตฺถปสุโต สิยา.\csromanspacer{} อตฺตวคฺโค ทฺวาทสโม.}}

\csromanheader{1}{13. โลกวคฺค}
\vspace{\tipitakaparskip}
\csromancenter{ทหรภิกฺขุวตฺถุ 167.\csromanspacer{} หีนํ ธมฺมํ น เสเวยฺย, ปมาเทน น สํวเส.}
\vspace{0.5\baselineskip}
\csromangathameasure{มิจฺฉาทิฏฺฐิํ น เสเวยฺย, น สิยา โลกวฑฺฒโน.}
\csromangathagroup{\csromangathastanza{มิจฺฉาทิฏฺฐิํ น เสเวยฺย, น สิยา โลกวฑฺฒโน.}}

\vspace{\tipitakaparskip}
\csromancenter{สุทฺโธทนวตฺถุ 168.\csromanspacer{} อุตฺติฏฺเฐ นปฺปมชฺเชยฺย, ธมฺมํ สุจริตํ จเร.}
\vspace{0.5\baselineskip}
\csromangathameasure{ธมฺมจารี สุขํ เสติ, อสฺมิํ โลเก ปรมฺหิ จ.}
\csromangathagroup{\csromangathastanza{ธมฺมจารี สุขํ เสติ, อสฺมิํ โลเก ปรมฺหิ จ.}}

\chapterheadpageread{13. โลกวคฺค 169. ธมฺมํ จเร สุจริตํ, น นํ ทุจฺจริตํ จเร.}
\csromangathameasure{ธมฺมจารี สุขํ เสติ, อสฺมิํ โลเก ปรมฺหิ จ.}
\csromangathagroup{\csromangathastanza{\csromanfolio{39}ธมฺมจารี สุขํ เสติ, อสฺมิํ โลเก ปรมฺหิ จ.}}

\vspace{\tipitakaparskip}
\csromancenter{ปญฺจสตวิปสฺสกภิกฺขุวตฺถุ 170.\csromanspacer{} ยถา ปุพฺพุฬกํ\footnote{พุพฺพุฬกํ (สี, อิ)} ปสฺเส, ยถา ปสฺเส มรีจิกํ.}
\vspace{0.5\baselineskip}
\csromangathameasure{เอวํ โลกํ อเวกฺขนฺตํ, มจฺจุราชา น ปสฺสติ.}
\csromangathagroup{\csromangathastanza{เอวํ โลกํ อเวกฺขนฺตํ, มจฺจุราชา น ปสฺสติ.}}

\vspace{\tipitakaparskip}
\csromancenter{อภยราชกุมารวตฺถุ 171.\csromanspacer{} เอถ ปสฺสถิมํ โลกํ, จิตฺตํ ราชรถูปมํ.}
\vspace{0.5\baselineskip}
\csromangathameasure{ยตฺถ พาลา วิสีทนฺติ, นตฺถิ สงฺโค วิชานตํ.}
\csromangathagroup{\csromangathastanza{ยตฺถ พาลา วิสีทนฺติ, นตฺถิ สงฺโค วิชานตํ.}}

\vspace{\tipitakaparskip}
\csromancenter{สมฺมชฺชนตฺเถรวตฺถุ 172.\csromanspacer{} โย จ ปุพฺเพ ปมชฺชิตฺวา, ปจฺฉา โส นปฺปมชฺชติ.}
\vspace{0.5\baselineskip}
\csromangathameasure{โส’มํ โลกํ ปภาเสติ, อพฺภา มุตฺโตว จนฺทิมา.}
\csromangathagroup{\csromangathastanza{โส’มํ โลกํ ปภาเสติ, อพฺภา มุตฺโตว จนฺทิมา.}}

\vspace{\tipitakaparskip}
\csromancenter{องฺคุลิมาลตฺเถรวตฺถุ 173.\csromanspacer{} ยสฺส ปาปํ กตํ กมฺมํ, กุสเลน ปิธียติ\footnote{ปิถียติ (สี, สฺยา, อิ)}.}
\vspace{0.5\baselineskip}
\csromangathameasure{โส’มํ โลกํ ปภาเสติ, อพฺภา มุตฺโตว จนฺทิมา.}
\csromangathagroup{\csromangathastanza{โส’มํ โลกํ ปภาเสติ, อพฺภา มุตฺโตว จนฺทิมา.}}

\vspace{\tipitakaparskip}
\csromancenter{เปสการธีตาวตฺถุ 174.\csromanspacer{} อนฺธภูโต\footnote{อนฺธีภูโต (ก)} อยํ โลโก, ตนุเก’ตฺถ วิปสฺสติ.}
\vspace{0.5\baselineskip}
\csromangathameasure{สกุโณ ชาลมุตฺโตว, อปฺโป สคฺคาย คจฺฉติ.}
\csromangathagroup{\csromangathastanza{สกุโณ ชาลมุตฺโตว, อปฺโป สคฺคาย คจฺฉติ.}}

\vspace{\tipitakaparskip}
\csromancenter{ติํสภิกฺขุวตฺถุ 175.\csromanspacer{} หํสา’ทิจฺจปเถ ยนฺติ, อากาเส ยนฺติ อิทฺธิยา.}
\vspace{0.5\baselineskip}
\csromangathameasure{นียนฺติ ธีรา โลกมฺหา, เชตฺวา มารํ สวาหินิํ.}
\csromangathagroup{\csromangathastanza{นียนฺติ ธีรา โลกมฺหา, เชตฺวา มารํ สวาหินิํ\footnote{สวาหนํ (สฺยา, ก)}.}}

\vspace{\tipitakaparskip}
\csromancenter{จิญฺจมาณวิกาวตฺถุ 176.\csromanspacer{} เอกํ ธมฺมํ อตีตสฺส, มุสาวาทิสฺส ชนฺตุโน.}
\vspace{0.5\baselineskip}
\csromangathameasure{วิติณฺณปรโลกสฺส, นตฺถิ ปาปํ อการิยํ.}
\csromangathagroup{\csromangathastanza{วิติณฺณปรโลกสฺส, นตฺถิ ปาปํ อการิยํ.}}

\csromanheader{2}{ธมฺมปทปาฬิ \textbf{10. อสทิสทานวตฺถุ} 177. น เว กทริยา เทวโลกํ วชนฺติ,}
\prose{\csromanfolio{40}พาลา หเว นปฺปสํสนฺติ ทานํ.}

\csromangathameasure{ธีโร จ ทานํ อนุโมทมาโน, \\ เตเนว โส โหติ สุขี ปรตฺถ.}
\csromangathagroup{\csromangathastanza{ธีโร จ ทานํ อนุโมทมาโน, \\ เตเนว โส โหติ สุขี ปรตฺถ.}}

\vspace{\tipitakaparskip}
\csromancenter{อนาถปิณฺฑิกปุตฺตกาลวตฺถุ 178.\csromanspacer{} ปถพฺยา เอกรชฺเชน, สคฺคสฺส คมเนน วา.}
\vspace{0.5\baselineskip}
\csromangathameasure{สพฺพโลกาธิปจฺเจน, โสตาปตฺติผลํ วรํ. โลกวคฺโค เตรสโม.}
\csromangathagroup{\csromangathastanza{สพฺพโลกาธิปจฺเจน, โสตาปตฺติผลํ วรํ.\csromanspacer{} โลกวคฺโค เตรสโม.}}

\csromanmatikaanchor{mtk.21}
\csromanmatikaanchor{mtk.22}
\csromantocmarknopage{chapter}{อุทานปาฬิ}{mtk.21}
\bookmark[dest={mtk.21},level=0]{อุทานปาฬิ}
\csromantocmarkref{section}{11. ชราวคฺค}{mtk.22}
\bookmark[dest={mtk.22},level=1]{11. ชราวคฺค}
\csromanheader{1}{14. พุทฺธวคฺค}
\vspace{\tipitakaparskip}
\csromancenter{มารธีตรวตฺถุ 179.\csromanspacer{} ยสฺส ชิตํ นาวชียติ,}
\prose{ชิตํ ยสฺส\footnote{ชิตมสฺส (สี, สฺยา, อิ), ชิตํ มสฺส (ก)} โน’ยาติ โกจิ โลเก.}

\csromangathameasure{ตํ พุทฺธมนนฺตโคจรํ, \\ อปทํ เกน ปเทน เนสฺสถ. \\ ยสฺส ชาลินี วิสตฺติกา, \\ ตณฺหา นตฺถิ กุหิญฺจิ เนตเว. \\ ตํ พุทฺธมนนฺตโคจรํ, \\ อปทํ เกน ปเทน เนสฺสถ.}
\csromangathagroup{\csromangathastanza{ตํ พุทฺธมนนฺตโคจรํ, \\ อปทํ เกน ปเทน เนสฺสถ. \\ ยสฺส ชาลินี วิสตฺติกา, \\ ตณฺหา นตฺถิ กุหิญฺจิ เนตเว. \\ ตํ พุทฺธมนนฺตโคจรํ, \\ อปทํ เกน ปเทน เนสฺสถ.}}

\vspace{\tipitakaparskip}
\csromancenter{เทโวโรหณวตฺถุ 181.\csromanspacer{} เย ฌานปสุตา ธีรา, เนกฺขมฺมูปสเม รตา.}
\vspace{0.5\baselineskip}
\csromangathameasure{เทวาปิ เตสํ ปิหยนฺติ, สมฺพุทฺธานํ สตีมตํ.}
\csromangathagroup{\csromangathastanza{เทวาปิ เตสํ ปิหยนฺติ, สมฺพุทฺธานํ สตีมตํ.}}

\vspace{\tipitakaparskip}
\csromancenter{พุทฺธวคฺค \textbf{3.}\csromanspacer{}\textbf{ เอรกปตฺตนาคราชวตฺถุ} 182.\csromanspacer{} กิจฺโฉ มนุสฺสปฏิลาโภ, กิจฺฉํ มจฺจาน ชีวิตํ.}
\vspace{0.5\baselineskip}
\csromangathameasure{กิจฺฉํ สทฺธมฺมสฺสวนํ, กิจฺโฉ พุทฺธานมุปฺปาโท.}
\csromangathagroup{\csromangathastanza{\csromanfolio{41}กิจฺฉํ สทฺธมฺมสฺสวนํ, กิจฺโฉ พุทฺธานมุปฺปาโท.}}

\vspace{\tipitakaparskip}
\csromancenter{อานนฺทตฺเถรปญฺหวตฺถุ 183. * สพฺพปาปสฺส อกรณํ, กุสลสฺส อุปสมฺปทา\footnote{กุสลสฺสูปสมฺปทา (สฺยา)}.}
\vspace{0.5\baselineskip}
\csromangathameasure{สจิตฺตปริโยทปนํ, เอตํ พุทฺธาน สาสนํ.}
\csromangathagroup{\csromangathastanza{สจิตฺตปริโยทปนํ\footnote{สจิตฺตปริโยทาปนํ (?)}, เอตํ พุทฺธาน สาสนํ.}}

\proseitem{184}{\csromansymbolfootnote{*}{ที 2. 42; อุปริ 127; ขุ 10. 37 ปิฏฺเฐสุปิ.}ขนฺตี ปรมํ ตโป ติติกฺขา,}

\prose{นิพฺพานํ\footnote{นิพฺพาณํ (ก-สี, อิ)} ปรมํ วทนฺติ พุทฺธา.}

\csromangathameasure{น หิ ปพฺพชิโต ปรูปฆาตี, \\ น สมโณ โหติ ปรํ วิเหฐยนฺโต.}
\csromangathagroup{\csromangathastanza{น หิ ปพฺพชิโต ปรูปฆาตี, \\ น\footnote{อยํ นกาโร สี-สฺยา-อิ-โปตฺถเกสุ น ทิสฺสติ.} สมโณ โหติ ปรํ วิเหฐยนฺโต.}}

\proseitem{185}{\csromansymbolmark{*}อนูปวาโท อนูปฆาโต\footnote{อนุปวาโท อนุปฆาโต (สฺยา, ก)}, ปาติโมกฺเข จ สํวโร.}

\csromangathameasure{มตฺตญฺญุตา จ ภตฺตสฺมิํ, ปนฺตญฺจ สยนาสนํ. \\ อธิจิตฺเต จ อาโยโค, เอตํ พุทฺธาน สาสนํ.}
\csromangathagroup{\csromangathastanza{มตฺตญฺญุตา จ ภตฺตสฺมิํ, ปนฺตญฺจ สยนาสนํ. \\ อธิจิตฺเต จ อาโยโค, เอตํ พุทฺธาน สาสนํ.}}

\vspace{\tipitakaparskip}
\csromancenter{อนภิรตภิกฺขุวตฺถุ 186.\csromanspacer{} น กหาปณวสฺเสน, ติตฺติ กาเมสุ วิชฺชติ.}
\vspace{0.5\baselineskip}
\csromangathameasure{อปฺปสฺสาทา ทุขา กามา, อิติ วิญฺญาย ปณฺฑิโต. \\ อปิ ทิพฺเพสุ กาเมสุ, รติํ โส นาธิคจฺฉติ. \\ ตณฺหกฺขยรโต โหติ, สมฺมาสมฺพุทฺธสาวโก.}
\csromangathagroup{\csromangathastanza{อปฺปสฺสาทา ทุขา กามา, อิติ วิญฺญาย ปณฺฑิโต. \\ อปิ ทิพฺเพสุ กาเมสุ, รติํ โส นาธิคจฺฉติ. \\ ตณฺหกฺขยรโต โหติ, สมฺมาสมฺพุทฺธสาวโก.}}

\vspace{\tipitakaparskip}
\csromancenter{อคฺคิทตฺตพฺราหฺมณวตฺถุ 188.\csromanspacer{} พหุํ เว สรณํ ยนฺติ, ปพฺพตานิ วนานิ จ.}
\vspace{0.5\baselineskip}
\csromangathameasure{อารามรุกฺขเจตฺยานิ, มนุสฺสา ภยตชฺชิตา. \\ เนตํ โข สรณํ เขมํ, เนตํ สรณมุตฺตมํ. \\ เนตํ สรณมาคมฺม, สพฺพทุกฺขา ปมุจฺจติ.}
\csromangathagroup{\csromangathastanza{อารามรุกฺขเจตฺยานิ, มนุสฺสา ภยตชฺชิตา. \\ เนตํ โข สรณํ เขมํ, เนตํ สรณมุตฺตมํ. \\ เนตํ สรณมาคมฺม, สพฺพทุกฺขา ปมุจฺจติ.}}

\csromanheader{2}{ธมฺมปทปาฬิ 190. โย จ พุทฺธญฺจ ธมฺมญฺจ, สํฆญฺจ สรณํ คโต.}
\csromangathameasure{จตฺตาริ อริยสจฺจานิ, สมฺมปฺปญฺญาย ปสฺสติ. \\ ทุกฺขํ ทุกฺขสมุปฺปาทํ, ทุกฺขสฺส จ อติกฺกมํ. \\ อริยํ จฏฺฐงฺคิกํ มคฺคํ, ทุกฺขูปสมคามินํ. \\ เอตํ โข สรณํ เขมํ, เอตํ สรณมุตฺตมํ. \\ เอตํ สรณมาคมฺม, สพฺพทุกฺขา ปมุจฺจติ.}
\csromangathagroup{\csromangathastanza{\csromanfolio{42}จตฺตาริ อริยสจฺจานิ, สมฺมปฺปญฺญาย ปสฺสติ. \\ ทุกฺขํ ทุกฺขสมุปฺปาทํ, ทุกฺขสฺส จ อติกฺกมํ.}\csromangathastanza{อริยํ จฏฺฐงฺคิกํ มคฺคํ, ทุกฺขูปสมคามินํ. \\ เอตํ โข สรณํ เขมํ, เอตํ สรณมุตฺตมํ. \\ เอตํ สรณมาคมฺม, สพฺพทุกฺขา ปมุจฺจติ.}}

\vspace{\tipitakaparskip}
\csromancenter{อานนฺทตฺเถรปญฺหวตฺถุ 193.\csromanspacer{} ทุลฺลโภ ปุริสาชญฺโญ, น โส สพฺพตฺถ ชายติ.}
\vspace{0.5\baselineskip}
\csromangathameasure{ยตฺถ โส ชายติ ธีโร, ตํ กุลํ สุขเมธติ.}
\csromangathagroup{\csromangathastanza{ยตฺถ โส ชายติ ธีโร, ตํ กุลํ สุขเมธติ.}}

\vspace{\tipitakaparskip}
\csromancenter{สมฺพหุลภิกฺขุวตฺถุ 194.\csromanspacer{} สุโข พุทฺธานมุปฺปาโท, สุขา สทฺธมฺมเทสนา.}
\vspace{0.5\baselineskip}
\csromangathameasure{สุขา สํฆสฺส สามคฺคี, สมคฺคานํ ตโป สุโข. \\ กสฺสปทสพลสฺส สุวณฺณเจติยวตฺถุ 195. ปูชารเห ปูชยโต, พุทฺเธ ยทิ ว สาวเก. \\ ปปญฺจสมติกฺกนฺเต, ติณฺณโสกปริทฺทเว. \\ เต ตาทิเส ปูชยโต, นิพฺพุเต อกุโตภเย. \\ น สกฺกา ปุญฺญํ สงฺขาตุํ, อิเมตฺตมปิ เกนจิ. พุทฺธวคฺโค จุทฺทสโม. ปฐมภาณวารํ.}
\csromangathagroup{\csromangathastanza{สุขา สํฆสฺส สามคฺคี, สมคฺคานํ ตโป สุโข. \\ กสฺสปทสพลสฺส สุวณฺณเจติยวตฺถุ 195.\csromanspacer{} ปูชารเห ปูชยโต, พุทฺเธ ยทิ ว สาวเก.}\csromangathastanza{ปปญฺจสมติกฺกนฺเต, ติณฺณโสกปริทฺทเว. \\ เต ตาทิเส ปูชยโต, นิพฺพุเต อกุโตภเย. \\ น สกฺกา ปุญฺญํ สงฺขาตุํ, อิเมตฺตมปิ เกนจิ.\csromanspacer{} พุทฺธวคฺโค จุทฺทสโม.\csromanspacer{} ปฐมภาณวารํ.}}

\csromanheader{1}{15. สุขวคฺค}
\csromanmatikaanchor{mtk.23}
\csromantocmarkref{section}{12. อตฺตวคฺค}{mtk.23}
\bookmark[dest={mtk.23},level=1]{12. อตฺตวคฺค}
\vspace{\tipitakaparskip}
\csromancenter{ญาติกลหวูปสมนวตฺถุ 197.\csromanspacer{} สุสุขํ วต ชีวาม, เวริเนสุ อเวริโน.}
\vspace{0.5\baselineskip}
\csromangathameasure{เวริเนสุ มนุสฺเสสุ, วิหราม อเวริโน.}
\csromangathagroup{\csromangathastanza{เวริเนสุ มนุสฺเสสุ, วิหราม อเวริโน.}}

\chapterheadpageread{15. สุขวคฺค 198. สุสุขํ วต ชีวาม, อาตุเรสุ อนาตุรา.}
\csromangathameasure{อาตุเรสุ มนุสฺเสสุ, วิหราม อนาตุรา. \\ สุสุขํ วต ชีวาม, อุสฺสุเกสุ อนุสฺสุกา. \\ อุสฺสุเกสุ มนุสฺเสสุ, วิหราม อนุสฺสุกา.}
\csromangathagroup{\csromangathastanza{\csromanfolio{43}อาตุเรสุ มนุสฺเสสุ, วิหราม อนาตุรา. \\ สุสุขํ วต ชีวาม, อุสฺสุเกสุ อนุสฺสุกา. \\ อุสฺสุเกสุ มนุสฺเสสุ, วิหราม อนุสฺสุกา.}}

\vspace{\tipitakaparskip}
\csromancenter{มารวตฺถุ 200.\csromanspacer{} สุสุขํ วต ชีวาม, เยสํ โน นตฺถิ กิญฺจนํ.}
\vspace{0.5\baselineskip}
\csromangathameasure{ปีติภกฺขา ภวิสฺสาม, เทวา อาภสฺสรา ยถา.}
\csromangathagroup{\csromangathastanza{ปีติภกฺขา ภวิสฺสาม, เทวา อาภสฺสรา ยถา.}}

\vspace{\tipitakaparskip}
\csromancenter{โกสลรญฺโญ ปราชยวตฺถุ 201.\csromanspacer{} ชยํ เวรํ ปสวติ, ทุกฺขํ เสติ ปราชิโต.}
\vspace{0.5\baselineskip}
\csromangathameasure{อุปสนฺโต สุขํ เสติ, หิตฺวา ชยปราชยํ.}
\csromangathagroup{\csromangathastanza{อุปสนฺโต สุขํ เสติ, หิตฺวา ชยปราชยํ.}}

\vspace{\tipitakaparskip}
\csromancenter{อญฺญตรกุลทาริกาวตฺถุ 202.\csromanspacer{} นตฺถิ ราคสโม อคฺคิ, นตฺถิ โทสสโม กลิ.}
\vspace{0.5\baselineskip}
\csromangathameasure{นตฺถิ ขนฺธสมา ทุกฺขา, นตฺถิ สนฺติปรํ สุขํ.}
\csromangathagroup{\csromangathastanza{นตฺถิ ขนฺธสมา\footnote{ขนฺธาทิสา (สี, สฺยา, อิ, รูปสิทฺธิยา สเมติ)} ทุกฺขา, นตฺถิ สนฺติปรํ สุขํ.}}

\vspace{\tipitakaparskip}
\csromancenter{เอก-อุปาสกวตฺถุ 203.\csromanspacer{} ชิฆจฺฉาปรมา โรคา, สงฺขารปรมา\footnote{สงฺขารา ปรมา (พหูสุ)} ทุขา.}
\vspace{0.5\baselineskip}
\csromangathameasure{เอตํ ญตฺวา ยถาภูตํ, นิพฺพานํ ปรมํ สุขํ.}
\csromangathagroup{\csromangathastanza{เอตํ ญตฺวา ยถาภูตํ, นิพฺพานํ ปรมํ สุขํ.}}

\vspace{\tipitakaparskip}
\csromancenter{ปเสนทิโกสลวตฺถุ 204.\csromanspacer{} อาโรคฺยปรมา ลาภา, สนฺตุฏฺฐิปรมํ ธนํ.}
\vspace{0.5\baselineskip}
\csromangathameasure{วิสฺสาสปรมา ญาติ, นิพฺพานํ ปรมํ สุขํ.}
\csromangathagroup{\csromangathastanza{วิสฺสาสปรมา ญาติ\footnote{วิสฺสาสปรโม ญาติ (ก-สี), วิสฺสาสปรมา ญาตี (สี-ฏฺฐ), วิสฺสาสา ปรมา ญาติ (ก)}, นิพฺพานํ ปรมํ\footnote{นิพฺพาณปรมํ (ก-สี)} สุขํ.}}

\vspace{\tipitakaparskip}
\csromancenter{ติสฺสตฺเถรวตฺถุ 205.\csromanspacer{} ปวิเวกรสํ ปิตฺวา\footnote{ปีตฺวา (สี, สฺยา, กํ, อิ)}, รสํ อุปสมสฺส จ.}
\vspace{0.5\baselineskip}
\csromangathameasure{นิทฺทโร โหติ นิปฺปาโป, ธมฺมปีติรสํ ปิวํ.}
\csromangathagroup{\csromangathastanza{นิทฺทโร โหติ นิปฺปาโป, ธมฺมปีติรสํ ปิวํ.}}

\csromanheader{2}{ธมฺมปทปาฬิ \textbf{8. สกฺกวตฺถุ} 206. สาหุ ทสฺสนมริยานํ, สนฺนิวาโส สทา สุโข.}
\csromangathameasure{อทสฺสเนน พาลานํ, นิจฺจเมว สุขี สิยา. \\ พาลสงฺคตจารี หิ, ทีฆมทฺธาน โสจติ. \\ ทุกฺโข พาเลหิ สํวาโส, อมิตฺเตเนว สพฺพทา. \\ ธีโร จ สุขสํวาโส, ญาตีนํว สมาคโม.}
\csromangathagroup{\csromangathastanza{\csromanfolio{44}อทสฺสเนน พาลานํ, นิจฺจเมว สุขี สิยา. \\ พาลสงฺคตจารี\footnote{พาลสงฺคติจารี (ก)} หิ, ทีฆมทฺธาน โสจติ.}\csromangathastanza{ทุกฺโข พาเลหิ สํวาโส, อมิตฺเตเนว สพฺพทา. \\ ธีโร จ สุขสํวาโส, ญาตีนํว สมาคโม.}}

\proseitem{208}{2ตสฺมา หิ\csromandash{}}

\csromangathameasure{ธีรญฺจ ปญฺญญฺจ พหุสฺสุตญฺจ, \\ โธรยฺหสีลํ วตวนฺตมริยํ. \\ ตํ ตาทิสํ สปฺปุริสํ สุเมธํ, \\ ภเชถ นกฺขตฺตปถํว จนฺทิมา. สุขวคฺโค ปนฺนรสโม.}
\csromangathagroup{\csromangathastanza{ธีรญฺจ ปญฺญญฺจ พหุสฺสุตญฺจ, \\ โธรยฺหสีลํ วตวนฺตมริยํ. \\ ตํ ตาทิสํ สปฺปุริสํ สุเมธํ, \\ ภเชถ นกฺขตฺตปถํว จนฺทิมา\footnote{ตสฺมา หิ ธีรํ ปญฺญญฺจ, พหุสฺสุตญฺจ โธรยฺหํ. สีลํ ธุตวตมริยํ, ตํ ตาทิสํ สปฺปุริสํ. สุเมธํ ภเชถ นกฺขตฺตปถํว จนฺทิมา. (ก)}.\csromanspacer{} สุขวคฺโค ปนฺนรสโม.}}

\csromanheader{1}{16. ปิยวคฺค}
\vspace{\tipitakaparskip}
\csromancenter{ตโยชนปพฺพชิตวตฺถุ 209.\csromanspacer{} อโยเค ยุญฺช’มตฺตานํ, โยคสฺมิญฺจ อโยชยํ.}
\vspace{0.5\baselineskip}
\csromanmatikaanchor{mtk.24}
\csromantocmarkref{section}{13. โลกวคฺค}{mtk.24}
\bookmark[dest={mtk.24},level=1]{13. โลกวคฺค}
\csromangathameasure{อตฺถํ หิตฺวา ปิยคฺคาหี, ปิเหต’ตฺตานุโยคินํ. \\ มา ปิเยหิ สมาคญฺฉิ, อปฺปิเยหิ กุทาจนํ. \\ ปิยานํ อทสฺสนํ ทุกฺขํ, อปฺปิยานญฺจ ทสฺสนํ.}
\csromangathagroup{\csromangathastanza{อตฺถํ หิตฺวา ปิยคฺคาหี, ปิเหต’ตฺตานุโยคินํ. \\ มา ปิเยหิ สมาคญฺฉิ, อปฺปิเยหิ กุทาจนํ. \\ ปิยานํ อทสฺสนํ ทุกฺขํ, อปฺปิยานญฺจ ทสฺสนํ.}}

\proseitem{211}{ตสฺมา ปิยํ น กยิราถ, ปิยาปาโย หิ ปาปโก.}

\csromangathameasure{คนฺถา เตสํ น วิชฺชนฺติ, เยสํ นตฺถิ ปิยาปฺปิยํ.}
\csromangathagroup{\csromangathastanza{คนฺถา เตสํ น วิชฺชนฺติ, เยสํ นตฺถิ ปิยาปฺปิยํ.}}

\vspace{\tipitakaparskip}
\csromancenter{อญฺญตรกุฏุมฺพิกวตฺถุ 212.\csromanspacer{} ปิยโต ชายตี โสโก, ปิยโต ชายตี\footnote{ชายเต (ก)} ภยํ.}
\vspace{0.5\baselineskip}
\csromangathameasure{ปิยโต วิปฺปมุตฺตสฺส, นตฺถิ โสโก กุโต ภยํ.}
\csromangathagroup{\csromangathastanza{ปิยโต วิปฺปมุตฺตสฺส, นตฺถิ โสโก กุโต ภยํ.}}

\chapterheadpageread{16. ปิยวคฺค \textbf{3. วิสาขาวตฺถุ} 213. เปมโต ชายตี โสโก, เปมโต ชายตี ภยํ.}
\csromangathameasure{เปมโต วิปฺปมุตฺตสฺส, นตฺถิ โสโก กุโต ภยํ.}
\csromangathagroup{\csromangathastanza{\csromanfolio{45}เปมโต วิปฺปมุตฺตสฺส, นตฺถิ โสโก กุโต ภยํ.}}

\vspace{\tipitakaparskip}
\csromancenter{ลิจฺฉวีวตฺถุ 214.\csromanspacer{} รติยา ชายตี โสโก, รติยา ชายตี ภยํ.}
\vspace{0.5\baselineskip}
\csromangathameasure{รติยา วิปฺปมุตฺตสฺส, นตฺถิ โสโก กุโต ภยํ.}
\csromangathagroup{\csromangathastanza{รติยา วิปฺปมุตฺตสฺส, นตฺถิ โสโก กุโต ภยํ.}}

\vspace{\tipitakaparskip}
\csromancenter{อนิตฺถิคนฺธกุมารวตฺถุ 215.\csromanspacer{} กามโต ชายตี โสโก, กามโต ชายตี ภยํ.}
\vspace{0.5\baselineskip}
\csromangathameasure{กามโต วิปฺปมุตฺตสฺส, นตฺถิ โสโก กุโต ภยํ.}
\csromangathagroup{\csromangathastanza{กามโต วิปฺปมุตฺตสฺส, นตฺถิ โสโก กุโต ภยํ.}}

\vspace{\tipitakaparskip}
\csromancenter{อญฺญตรพฺราหฺมณวตฺถุ 216.\csromanspacer{} ตณฺหาย ชายตี\footnote{ชายเต (ก)} โสโก, ตณฺหาย ชายตี1 ภยํ.}
\vspace{0.5\baselineskip}
\csromangathameasure{ตณฺหาย วิปฺปมุตฺตสฺส, นตฺถิ โสโก กุโต ภยํ.}
\csromangathagroup{\csromangathastanza{ตณฺหาย วิปฺปมุตฺตสฺส, นตฺถิ โสโก กุโต ภยํ.}}

\vspace{\tipitakaparskip}
\csromancenter{ปญฺจสตทารกวตฺถุ 217.\csromanspacer{} สีลทสฺสนสมฺปนฺนํ, ธมฺมฏฺฐํ สจฺจเวทินํ.}
\vspace{0.5\baselineskip}
\csromangathameasure{อตฺตโน กมฺม กุพฺพานํ, ตํ ชโน กุรุเต ปิยํ.}
\csromangathagroup{\csromangathastanza{อตฺตโน กมฺม กุพฺพานํ, ตํ ชโน กุรุเต ปิยํ.}}

\vspace{\tipitakaparskip}
\csromancenter{เอก-อนาคามิตฺเถรวตฺถุ 218.\csromanspacer{} ฉนฺทชาโต อนกฺขาเต, มนสา จ ผุโฏ สิยา.}
\vspace{0.5\baselineskip}
\csromangathameasure{กาเมสุ จ อปฺปฏิพทฺธจิตฺโต, “อุทฺธํโสโต”ติ วุจฺจติ.}
\csromangathagroup{\csromangathastanza{กาเมสุ จ อปฺปฏิพทฺธจิตฺโต\footnote{อปฺปฏิพนฺธจิตฺโต (ก)}, “อุทฺธํโสโต”ติ วุจฺจติ.}}

\vspace{\tipitakaparskip}
\csromancenter{นนฺทิยวตฺถุ 219. * จิรปฺปวาสิํ ปุริสํ, ทูรโต โสตฺถิมาคตํ.}
\vspace{0.5\baselineskip}
\csromangathameasure{ญาติมิตฺตา สุหชฺชา จ, อภินนฺทนฺติ อาคตํ.}
\csromangathagroup{\csromangathastanza{ญาติมิตฺตา สุหชฺชา จ, อภินนฺทนฺติ อาคตํ.}}

\proseitem{220}{\csromansymbolfootnote{*}{ขุ 2. 73 ปิฏฺเฐปิ.}ตเถว กตปุญฺญมฺปิ, อสฺมา โลกา ปรํ คตํ.}

\csromangathameasure{ปุญฺญานิ ปฏิคณฺหนฺติ, ปิยํ ญาตีว อาคตํ. ปิยวคฺโค โสฬสโม.}
\csromangathagroup{\csromangathastanza{ปุญฺญานิ ปฏิคณฺหนฺติ, ปิยํ ญาตีว อาคตํ.\csromanspacer{} ปิยวคฺโค โสฬสโม.}}

\csromanheader{1}{ธมฺมปทปาฬิ \textbf{17. โกธวคฺค}}
\vspace{\tipitakaparskip}
\csromancenter{โรหินีขตฺติยกญฺญาวตฺถุ 221.\csromanspacer{} โกธํ ชเห วิปฺปชเหยฺย มานํ,}
\prose{\csromanfolio{46}สํโยชนํ สพฺพมติกฺกเมยฺย.}

\csromangathameasure{ตํ นามรูปสฺมิมสชฺชมานํ, \\ อกิญฺจนํ นานุปตนฺติ ทุกฺขา.}
\csromangathagroup{\csromangathastanza{ตํ นามรูปสฺมิมสชฺชมานํ, \\ อกิญฺจนํ นานุปตนฺติ ทุกฺขา.}}

\vspace{\tipitakaparskip}
\csromancenter{อญฺญตรภิกฺขุวตฺถุ 222.\csromanspacer{} โย เว อุปฺปติตํ โกธํ, รถํ ภนฺตํว วารเย\footnote{ธารเย (สี, สฺยา, อิ)}.}
\vspace{0.5\baselineskip}
\csromangathameasure{ตมหํ สารถิํ พฺรูมิ, รสฺมิคฺคาโห อิตโร ชโน.}
\csromangathagroup{\csromangathastanza{ตมหํ สารถิํ พฺรูมิ, รสฺมิคฺคาโห อิตโร ชโน.}}

\csromanmatikaanchor{mtk.25}
\csromantocmarkref{section}{14. พุทฺธวคฺค}{mtk.25}
\bookmark[dest={mtk.25},level=1]{14. พุทฺธวคฺค}
\vspace{\tipitakaparskip}
\csromancenter{อุตฺตรา-อุปาสิกวตฺถุ 223.\csromanspacer{} อกฺโกเธน ชิเน โกธํ, อสาธุํ สาธุนา ชิเน.}
\vspace{0.5\baselineskip}
\csromangathameasure{ชิเน กทริยํ ทาเนน, สจฺเจนา’ลิกวาทินํ.}
\csromangathagroup{\csromangathastanza{ชิเน กทริยํ ทาเนน, สจฺเจนา’ลิกวาทินํ.}}

\vspace{\tipitakaparskip}
\csromancenter{มหาโมคฺคลฺลานปญฺหวตฺถุ 224.\csromanspacer{} สจฺจํ ภเณ น กุชฺเฌยฺย, ทชฺชา อปฺปมฺปิ\footnote{ทชฺชา’ปฺปสฺมิมฺปิ (สี, อิ), ทชฺชา อปฺปสฺมิ (สฺยา, ก)} ยาจิโต.}
\vspace{0.5\baselineskip}
\csromangathameasure{เอเตหิ ตีหิ ฐาเนหิ, คจฺเฉ เทวาน สนฺติเก.}
\csromangathagroup{\csromangathastanza{เอเตหิ ตีหิ ฐาเนหิ, คจฺเฉ เทวาน สนฺติเก.}}

\vspace{\tipitakaparskip}
\csromancenter{พุทฺธปิตุพฺราหฺมณวตฺถุ 225.\csromanspacer{} อหิํสกา เย มุนโย\footnote{อหิํสกายา มุนโย (ก)}, นิจฺจํ กาเยน สํวุตา.}
\vspace{0.5\baselineskip}
\csromangathameasure{เต ยนฺติ อจฺจุตํ ฐานํ, ยตฺถ คนฺตฺวา น โสจเร.}
\csromangathagroup{\csromangathastanza{เต ยนฺติ อจฺจุตํ ฐานํ, ยตฺถ คนฺตฺวา น โสจเร.}}

\vspace{\tipitakaparskip}
\csromancenter{ปุณฺณทาสีวตฺถุ 226.\csromanspacer{} สทา ชาครมานานํ, อโหรตฺตานุสิกฺขินํ.}
\vspace{0.5\baselineskip}
\csromangathameasure{นิพฺพานํ อธิมุตฺตานํ, อตฺถํ คจฺฉนฺติ อาสวา.}
\csromangathagroup{\csromangathastanza{นิพฺพานํ อธิมุตฺตานํ, อตฺถํ คจฺฉนฺติ อาสวา.}}

\csromanheader{1}{18. มลวคฺค \textbf{7. อตุล-อุปาสกวตฺถุ} 227. โปราณเมตํ อตุล, เนตํ อชฺชตนามิว.}
\csromangathameasure{นินฺทนฺติ ตุณฺหิมาสีนํ, นินฺทนฺติ พหุภาณินํ. \\ มิตภาณิมฺปิ นินฺทนฺติ, นตฺถิ โลเก อนินฺทิโต. \\ น จาหุ น จ ภวิสฺสติ, น เจตรหิ วิชฺชติ. \\ เอกนฺตํ นินฺทิโต โปโส, เอกนฺตํ วา ปสํสิโต. \\ ยํ เจ วิญฺญู ปสํสนฺติ, อนุวิจฺจ สุเว สุเว. \\ อจฺฉิทฺทวุตฺติํ เมธาวิํ, ปญฺญาสีลสมาหิตํ. \\ นิกฺขํ ชมฺโพนทสฺเสว, โก ตํ นินฺทิตุมรหติ. \\ เทวาปิ นํ ปสํสนฺติ, พฺรหฺมุนาปิ ปสํสิโต.}
\csromangathagroup{\csromangathastanza{\csromanfolio{47}นินฺทนฺติ ตุณฺหิมาสีนํ, นินฺทนฺติ พหุภาณินํ. \\ มิตภาณิมฺปิ นินฺทนฺติ, นตฺถิ โลเก อนินฺทิโต.}\csromangathastanza{น จาหุ น จ ภวิสฺสติ, น เจตรหิ วิชฺชติ. \\ เอกนฺตํ นินฺทิโต โปโส, เอกนฺตํ วา ปสํสิโต.}\csromangathastanza{ยํ เจ วิญฺญู ปสํสนฺติ, อนุวิจฺจ สุเว สุเว. \\ อจฺฉิทฺทวุตฺติํ\footnote{อจฺฉินฺนวุตฺติํ (ก)} เมธาวิํ, ปญฺญาสีลสมาหิตํ.}\csromangathastanza{นิกฺขํ\footnote{เนกฺขํ (สี, สฺยา, อิ)} ชมฺโพนทสฺเสว, โก ตํ นินฺทิตุมรหติ. \\ เทวาปิ นํ ปสํสนฺติ, พฺรหฺมุนาปิ ปสํสิโต.}}

\vspace{\tipitakaparskip}
\csromancenter{ฉพฺพคฺคิยวตฺถุ 231.\csromanspacer{} กายปฺปโกปํ รกฺเขยฺย, กาเยน สํวุโต สิยา.}
\vspace{0.5\baselineskip}
\csromangathameasure{กายทุจฺจริตํ หิตฺวา, กาเยน สุจริตํ จเร. \\ วจีปโกปํ รกฺเขยฺย, วาจาย สํวุโต สิยา. \\ วจีทุจฺจริตํ หิตฺวา, วาจาย สุจริตํ จเร. \\ มโนปโกปํ รกฺเขยฺย, มนสา สํวุโต สิยา. \\ มโนทุจฺจริตํ หิตฺวา, มนสา สุจริตํ จเร.}
\csromangathagroup{\csromangathastanza{กายทุจฺจริตํ หิตฺวา, กาเยน สุจริตํ จเร. \\ วจีปโกปํ รกฺเขยฺย, วาจาย สํวุโต สิยา.}\csromangathastanza{วจีทุจฺจริตํ หิตฺวา, วาจาย สุจริตํ จเร. \\ มโนปโกปํ รกฺเขยฺย, มนสา สํวุโต สิยา. \\ มโนทุจฺจริตํ หิตฺวา, มนสา สุจริตํ จเร.}}

\proseitem{234}{กาเยน สํวุตา ธีรา, อโถ วาจาย สํวุตา.}

\csromangathameasure{มนสา สํวุตา ธีรา, เต เว สุปริสํวุตา. โกธวคฺโค สตฺตรสโม.}
\csromangathagroup{\csromangathastanza{มนสา สํวุตา ธีรา, เต เว สุปริสํวุตา.\csromanspacer{} โกธวคฺโค สตฺตรสโม.}}

\chapterhead{18. มลวคฺค}
\vspace{\tipitakaparskip}
\csromancenter{โคฆาตกปุตฺตวตฺถุ 235.\csromanspacer{} ปณฺฑุปลาโสว ทานิสิ, ยมปุริสาปิ จ เต\footnote{ตํ (สี, สฺยา, กํ, อิ)} อุปฏฺฐิตา.}
\vspace{0.5\baselineskip}
\csromangathameasure{อุยฺโยคมุเข จ ติฏฺฐสิ, ปาเถยฺยมฺปิ จ เต น วิชฺชติ.}
\csromangathagroup{\csromangathastanza{อุยฺโยคมุเข จ ติฏฺฐสิ, ปาเถยฺยมฺปิ จ เต น วิชฺชติ.}}

\csromanheader{2}{ธมฺมปทปาฬิ 236. โส กโรหิ ทีปมตฺตโน, ขิปฺปํ วายม ปณฺฑิโต ภว.}
\csromangathameasure{นิทฺธนฺตมโล อนงฺคโณ, ทิพฺพํ อริยภูมิํ อุเปหิสิ. \\ อุปนีตวโย จ ทานิสิ, สมฺปยาโตสิ ยมสฺส สนฺติกํ. \\ วาโส เต นตฺถิ อนฺตรา, ปาเถยฺยมฺปิ จ เต น วิชฺชติ. \\ โส กโรหิ ทีปมตฺตโน, ขิปฺปํ วายม ปณฺฑิโต ภว. \\ นิทฺธนฺตมโล อนงฺคโณ, น ปุนํ ชาติชรํ อุเปหิสิ.}
\csromangathagroup{\csromangathastanza{\csromanfolio{48}นิทฺธนฺตมโล อนงฺคโณ, ทิพฺพํ อริยภูมิํ อุเปหิสิ\footnote{ทิพฺพํ อริยภูมิเมหิสิ (สี, สฺยา, อิ), ทิพฺพมริยภูมิํ อุเปหิสิ (?)}. \\ อุปนีตวโย จ ทานิสิ, สมฺปยาโตสิ ยมสฺส สนฺติกํ.}\csromangathastanza{วาโส\footnote{วาโสปิ จ (พหูสุ)} เต นตฺถิ อนฺตรา, ปาเถยฺยมฺปิ จ เต น วิชฺชติ. \\ โส กโรหิ ทีปมตฺตโน, ขิปฺปํ วายม ปณฺฑิโต ภว. \\ นิทฺธนฺตมโล อนงฺคโณ, น ปุนํ ชาติชรํ\footnote{วาโสปิ จ (พหูสุ)} อุเปหิสิ.}}

\vspace{\tipitakaparskip}
\csromancenter{อญฺญตรพฺราหฺมณวตฺถุ 239.\csromanspacer{} อนุปุพฺเพน เมธาวี, โถกํ โถกํ ขเณ ขเณ.}
\vspace{0.5\baselineskip}
\csromangathameasure{กมฺมาโร รชตสฺเสว, นิทฺธเม มลมตฺตโน.}
\csromangathagroup{\csromangathastanza{กมฺมาโร รชตสฺเสว, นิทฺธเม มลมตฺตโน.}}

\vspace{\tipitakaparskip}
\csromancenter{ติสฺสตฺเถรวตฺถุ 240. * อยสาว มลํ สมุฏฺฐิตํ\footnote{สมุฏฺฐาย (ก)}, ตตุฏฺฐาย\footnote{ตทุฏฺฐาย (สี, สฺยา, อิ)} ตเมว ขาทติ.}
\vspace{0.5\baselineskip}
\csromangathameasure{เอวํ อติโธนจารินํ, สานิ กมฺมานิ นยนฺติ ทุคฺคติํ.}
\csromangathagroup{\csromangathastanza{เอวํ อติโธนจารินํ, สานิ กมฺมานิ\footnote{สกกมฺมานิ (สี, อิ)} นยนฺติ ทุคฺคติํ.}}

\vspace{\tipitakaparskip}
\csromancenter{ลาฬุทายีวตฺถุ 241.\csromanspacer{} อสชฺฌายมลา มนฺตา, อนุฏฺฐานมลา ฆรา.}
\vspace{0.5\baselineskip}
\csromangathameasure{มลํ วณฺณสฺส โกสชฺชํ, ปมาโท รกฺขโต มลํ.}
\csromangathagroup{\csromangathastanza{มลํ วณฺณสฺส โกสชฺชํ, ปมาโท รกฺขโต มลํ.}}

\vspace{\tipitakaparskip}
\csromancenter{อญฺญตรกุลปุตฺตวตฺถุ 242.\csromanspacer{} มลิตฺถิยา ทุจฺจริตํ, มจฺเฉรํ ททโต มลํ.}
\vspace{0.5\baselineskip}
\csromangathameasure{มลา เว ปาปกา ธมฺมา, อสฺมิํ โลเก ปรมฺหิ จ. \\ ตโต มลา มลตรํ, อวิชฺชา ปรมํ มลํ. \\ เอตํ มลํ ปหนฺตฺวาน, นิมฺมลา โหถ ภิกฺขโว.}
\csromangathagroup{\csromangathastanza{มลา เว ปาปกา ธมฺมา, อสฺมิํ โลเก ปรมฺหิ จ. \\ ตโต มลา มลตรํ, อวิชฺชา ปรมํ มลํ. \\ เอตํ มลํ ปหนฺตฺวาน, นิมฺมลา โหถ ภิกฺขโว.}}

\vspace{\tipitakaparskip}
\csromancenter{จูฬสาริภิกฺขุวตฺถุ 244.\csromanspacer{} สุชีวํ อหิริเกน, กากสูเรน ธํสินา.}
\vspace{0.5\baselineskip}
\csromanmatikaanchor{mtk.26}
\csromantocmarkref{section}{15. สุขวคฺค}{mtk.26}
\bookmark[dest={mtk.26},level=1]{15. สุขวคฺค}
\csromangathameasure{ปกฺขนฺทินา ปคพฺเภน, สํกิลิฏฺเฐน ชีวิตํ.}
\csromangathagroup{\csromangathastanza{ปกฺขนฺทินา ปคพฺเภน, สํกิลิฏฺเฐน ชีวิตํ.}}

\chapterheadpageread{18. มลวคฺค 245. หิรีมตา จ ทุชฺชีวํ, นิจฺจํ สุจิคเวสินา.}
\csromangathameasure{อลีเนนา’ปฺปคพฺเภน, สุทฺธาชีเวน ปสฺสตา.}
\csromangathagroup{\csromangathastanza{\csromanfolio{49}อลีเนนา’ปฺปคพฺเภน, สุทฺธาชีเวน ปสฺสตา.}}

\vspace{\tipitakaparskip}
\csromancenter{ปญฺจ-อุปาสกวตฺถุ 246.\csromanspacer{} โย ปาณมติปาเตติ, มุสาวาทญฺจ ภาสติ.}
\vspace{0.5\baselineskip}
\csromangathameasure{โลเก อทินฺนมาทิยติ, ปรทารญฺจ คจฺฉติ. \\ สุราเมรยปานญฺจ, โย นโร อนุยุญฺชติ. \\ อิเธว เมโส โลกสฺมิํ, มูลํ ขณติ อตฺตโน. \\ เอวํ โภ ปุริส ชานาหิ, ปาปธมฺมา อสญฺญตา. \\ มา ตํ โลโภ อธมฺโม จ, จิรํ ทุกฺขาย รนฺธยุํ.}
\csromangathagroup{\csromangathastanza{โลเก อทินฺนมาทิยติ, ปรทารญฺจ คจฺฉติ. \\ สุราเมรยปานญฺจ, โย นโร อนุยุญฺชติ.}\csromangathastanza{อิเธว เมโส โลกสฺมิํ, มูลํ ขณติ อตฺตโน. \\ เอวํ โภ ปุริส ชานาหิ, ปาปธมฺมา อสญฺญตา. \\ มา ตํ โลโภ อธมฺโม จ, จิรํ ทุกฺขาย รนฺธยุํ.}}

\vspace{\tipitakaparskip}
\csromancenter{ติสฺสทหรวตฺถุ 249.\csromanspacer{} ททาติ เว ยถาสทฺธํ, ยถาปสาทนํ\footnote{ยตฺถ ปสาทนํ (กตฺถจิ)} ชโน.}
\vspace{0.5\baselineskip}
\csromangathameasure{ตตฺถ โย จ มงฺกุ ภวติ, ปเรสํ ปานโภชเน. \\ น โส ทิวา วา รตฺติํ วา, สมาธิมธิคจฺฉติ. \\ ยสฺส เจตํ สมุจฺฉินฺนํ, มูลฆจฺจํ สมูหตํ. \\ ส เว ทิวา วา รตฺติํ วา, สมาธิมธิคจฺฉติ.}
\csromangathagroup{\csromangathastanza{ตตฺถ โย จ มงฺกุ ภวติ\footnote{ตตฺถ เจ มํกุ โย โหติ (สี), ตตฺถ โย มงฺกุโต โหติ (สฺยา)}, ปเรสํ ปานโภชเน. \\ น โส ทิวา วา รตฺติํ วา, สมาธิมธิคจฺฉติ.}\csromangathastanza{ยสฺส เจตํ สมุจฺฉินฺนํ, มูลฆจฺจํ\footnote{มูลฆจฺฉํ (ก)} สมูหตํ. \\ ส เว ทิวา วา รตฺติํ วา, สมาธิมธิคจฺฉติ.}}

\vspace{\tipitakaparskip}
\csromancenter{ปญฺจ-อุปาสกวตฺถุ 251.\csromanspacer{} นตฺถิ ราคสโม อคฺคิ, นตฺถิ โทสสโม คโห.}
\vspace{0.5\baselineskip}
\csromangathameasure{นตฺถิ โมหสมํ ชาลํ, นตฺถิ ตณฺหาสมา นที.}
\csromangathagroup{\csromangathastanza{นตฺถิ โมหสมํ ชาลํ, นตฺถิ ตณฺหาสมา นที.}}

\vspace{\tipitakaparskip}
\csromancenter{เมณฺฑกเสฏฺฐิวตฺถุ 252.\csromanspacer{} สุทสฺสํ วชฺชมญฺเญสํ, อตฺตโน ปน ทุทฺทสํ.}
\vspace{0.5\baselineskip}
\csromangathameasure{ปเรสํ หิ โส วชฺชานิ, โอปุนาติ ยถา ภุสํ. \\ อตฺตโน ปน ฉาเทติ, กลิํว กิตวา สโฐ.}
\csromangathagroup{\csromangathastanza{ปเรสํ หิ โส วชฺชานิ, โอปุนาติ\footnote{โอผุนาติ (ก)} ยถา ภุสํ. \\ อตฺตโน ปน ฉาเทติ, กลิํว กิตวา สโฐ.}}

\vspace{\tipitakaparskip}
\csromancenter{ธมฺมปทปาฬิ \textbf{11.}\csromanspacer{}\textbf{ อุชฺฌานสญฺญิตฺเถรวตฺถุ} 253.\csromanspacer{} ปรวชฺชานุปสฺสิสฺส, นิจฺจํ อุชฺฌานสญฺญิโน.}
\vspace{0.5\baselineskip}
\csromangathameasure{อาสวา ตสฺส วฑฺฒนฺติ, อารา โส อาสวกฺขยา.}
\csromangathagroup{\csromangathastanza{\csromanfolio{50}อาสวา ตสฺส วฑฺฒนฺติ, อารา โส อาสวกฺขยา.}}

\vspace{\tipitakaparskip}
\csromancenter{สุภทฺทปริพฺพาชกวตฺถุ 254.\csromanspacer{} อากาเสว ปทํ นตฺถิ, สมโณ นตฺถิ พาหิเร.}
\vspace{0.5\baselineskip}
\csromangathameasure{ปปญฺจาภิรตา ปชา, นิปฺปปญฺจา ตถาคตา. \\ อากาเสว ปทํ นตฺถิ, สมโณ นตฺถิ พาหิเร. \\ สงฺขารา สสฺสตา นตฺถิ, นตฺถิ พุทฺธานมิญฺชิตํ. มลวคฺโค อฏฺฐารสโม.}
\csromangathagroup{\csromangathastanza{ปปญฺจาภิรตา ปชา, นิปฺปปญฺจา ตถาคตา. \\ อากาเสว ปทํ นตฺถิ, สมโณ นตฺถิ พาหิเร. \\ สงฺขารา สสฺสตา นตฺถิ, นตฺถิ พุทฺธานมิญฺชิตํ.\csromanspacer{} มลวคฺโค อฏฺฐารสโม.}}

\csromanheader{1}{19. ธมฺมฏฺฐวคฺค}
\vspace{\tipitakaparskip}
\csromancenter{วินิจฺฉยมหามตฺตวตฺถุ 256.\csromanspacer{} น เตน โหติ ธมฺมฏฺโฐ, เยนตฺถํ สาหสา\footnote{สหสา (สี, สฺยา, อิ)} นเย.}
\vspace{0.5\baselineskip}
\csromangathameasure{โย จ อตฺถํ อนตฺถญฺจ, อุโภ นิจฺเฉยฺย ปณฺฑิโต. \\ อสาหเสน ธมฺเมน, สเมน นยตี ปเร. \\ ธมฺมสฺส คุตฺโต เมธาวี, “ธมฺมฏฺโฐ”ติ ปวุจฺจติ.}
\csromangathagroup{\csromangathastanza{โย จ อตฺถํ อนตฺถญฺจ, อุโภ นิจฺเฉยฺย ปณฺฑิโต. \\ อสาหเสน ธมฺเมน, สเมน นยตี ปเร. \\ ธมฺมสฺส คุตฺโต เมธาวี, “ธมฺมฏฺโฐ”ติ ปวุจฺจติ.}}

\vspace{\tipitakaparskip}
\csromancenter{ฉพฺพคฺคิยวตฺถุ 258.\csromanspacer{} น เตน ปณฺฑิโต โหติ, ยาวตา พหุ ภาสติ.}
\vspace{0.5\baselineskip}
\csromangathameasure{เขมี อเวรี อภโย, “ปณฺฑิโต”ติ ปวุจฺจติ.}
\csromangathagroup{\csromangathastanza{เขมี อเวรี อภโย, “ปณฺฑิโต”ติ ปวุจฺจติ.}}

\vspace{\tipitakaparskip}
\csromancenter{เอกุทานขีณาสววตฺถุ 259.\csromanspacer{} น ตาวตา ธมฺมธโร, ยาวตา พหุ ภาสติ.}
\vspace{0.5\baselineskip}
\csromangathameasure{โย จ อปฺปมฺปิ สุตฺวาน, ธมฺมํ กาเยน ปสฺสติ. \\ ส เว ธมฺมธโร โหติ, โย ธมฺมํ นปฺปมชฺชติ.}
\csromangathagroup{\csromangathastanza{โย จ อปฺปมฺปิ สุตฺวาน, ธมฺมํ กาเยน ปสฺสติ. \\ ส เว ธมฺมธโร โหติ, โย ธมฺมํ นปฺปมชฺชติ.}}

\vspace{\tipitakaparskip}
\csromancenter{ธมฺมฏฺฐวคฺค \textbf{4.}\csromanspacer{}\textbf{ ลกุณฺฑกภทฺทิยตฺเถรวตฺถุ} 260.\csromanspacer{} น เตน เถโร โส โหติ\footnote{เถโร โหติ (สี, สฺยา)}, เยนสฺส ปลิตํ สิโร.}
\vspace{0.5\baselineskip}
\csromanmatikaanchor{mtk.27}
\csromantocmarkref{section}{16. ปิยวคฺค}{mtk.27}
\bookmark[dest={mtk.27},level=1]{16. ปิยวคฺค}
\csromangathameasure{ปริปกฺโก วโย ตสฺส, “โมฆชิณฺโณ”ติ วุจฺจติ. \\ ยมฺหิ สจฺจญฺจ ธมฺโม จ, อหิํสา สํยโม ทโม. \\ ส เว วนฺตมโล ธีโร, “เถโร” อิติ ปวุจฺจติ.}
\csromangathagroup{\csromangathastanza{\csromanfolio{51}ปริปกฺโก วโย ตสฺส, “โมฆชิณฺโณ”ติ วุจฺจติ. \\ ยมฺหิ สจฺจญฺจ ธมฺโม จ, อหิํสา สํยโม ทโม. \\ ส เว วนฺตมโล ธีโร, “เถโร” อิติ\footnote{โส เถโรติ (สฺยา, ก)} ปวุจฺจติ.}}

\vspace{\tipitakaparskip}
\csromancenter{สมฺพหุลภิกฺขุวตฺถุ 262.\csromanspacer{} น วากฺกรณมตฺเตน, วณฺณโปกฺขรตาย วา.}
\vspace{0.5\baselineskip}
\csromangathameasure{สาธุรูโป นโร โหติ, อิสฺสุกี มจฺฉรี สโฐ. \\ ยสฺส เจตํ สมุจฺฉินฺนํ, มูลฆจฺจํ สมูหตํ. \\ ส วนฺตโทโส เมธาวี, “สาธุรูโป”ติ วุจฺจติ.}
\csromangathagroup{\csromangathastanza{สาธุรูโป นโร โหติ, อิสฺสุกี มจฺฉรี สโฐ. \\ ยสฺส เจตํ สมุจฺฉินฺนํ, มูลฆจฺจํ สมูหตํ. \\ ส วนฺตโทโส เมธาวี, “สาธุรูโป”ติ วุจฺจติ.}}

\vspace{\tipitakaparskip}
\csromancenter{หตฺถกวตฺถุ 264.\csromanspacer{} น มุณฺฑเกน สมโณ, อพฺพโต อลิกํ ภณํ.}
\vspace{0.5\baselineskip}
\csromangathameasure{อิจฺฉาโลภสมาปนฺโน, สมโณ กิํ ภวิสฺสติ. \\ โย จ สเมติ ปาปานิ, อณุํถูลานิ สพฺพโส. \\ สมิตตฺตา หิ ปาปานํ, “สมโณ”ติ ปวุจฺจติ.}
\csromangathagroup{\csromangathastanza{อิจฺฉาโลภสมาปนฺโน, สมโณ กิํ ภวิสฺสติ. \\ โย จ สเมติ ปาปานิ, อณุํถูลานิ สพฺพโส. \\ สมิตตฺตา หิ ปาปานํ, “สมโณ”ติ ปวุจฺจติ.}}

\vspace{\tipitakaparskip}
\csromancenter{อญฺญตรพฺราหฺมณวตฺถุ 266.\csromanspacer{} น เตน ภิกฺขุ โส โหติ, ยาวตา ภิกฺขเต ปเร.}
\vspace{0.5\baselineskip}
\csromangathameasure{วิสฺสํ ธมฺมํ สมาทาย, ภิกฺขุ โหติ น ตาวตา. \\ โย’ธ ปุญฺญญฺจ ปาปญฺจ, พาเหตฺวา พฺรหฺมจริยวา. \\ สงฺขาย โลเก จรติ, ส เว “ภิกฺขู”ติ วุจฺจติ.}
\csromangathagroup{\csromangathastanza{วิสฺสํ ธมฺมํ สมาทาย, ภิกฺขุ โหติ น ตาวตา. \\ โย’ธ ปุญฺญญฺจ ปาปญฺจ, พาเหตฺวา พฺรหฺมจริยวา\footnote{พฺรหฺมจริยํ (ก)}. \\ สงฺขาย โลเก จรติ, ส เว “ภิกฺขู”ติ วุจฺจติ.}}

\vspace{\tipitakaparskip}
\csromancenter{ติตฺถิยวตฺถุ 268. * น โมเนน มุนี โหติ, มูฬฺหรูโป อวิทฺทสุ.}
\vspace{0.5\baselineskip}
\csromangathameasure{โย จ ตุลํว ปคฺคยฺห, วรมาทาย ปณฺฑิโต.}
\csromangathagroup{\csromangathastanza{โย จ ตุลํว ปคฺคยฺห, วรมาทาย ปณฺฑิโต.}}

\csromanheader{2}{ธมฺมปทปาฬิ 269. ปาปานิ ปริวชฺเชติ, ส มุนี เตน โส มุนิ.}
\csromangathameasure{โย มุนาติ อุโภ โลเก, “มุนิ” เตน ปวุจฺจติ.}
\csromangathagroup{\csromangathastanza{\csromanfolio{52}โย มุนาติ อุโภ โลเก, “มุนิ” เตน ปวุจฺจติ.}}

\vspace{\tipitakaparskip}
\csromancenter{พาลิสิกวตฺถุ 270.\csromanspacer{} น เตน อริโย โหติ, เยน ปาณานิ หิํสติ.}
\vspace{0.5\baselineskip}
\csromangathameasure{อหิํสา สพฺพปาณานํ, “อริโย”ติ ปวุจฺจติ. \\ สมฺพหุลสีลาทิสมฺปนฺนภิกฺขุวตฺถุ 271. น สีลพฺพตมตฺเตน, พาหุสจฺเจน วา ปน. \\ อถ วา สมาธิลาเภน, วิวิตฺตสยเนน วา. \\ ผุสามิ เนกฺขมฺมสุขํ, อปุถุชฺชนเสวิตํ. \\ ภิกฺขุ วิสฺสาสมาปาทิ, อปฺปตฺโต อาสวกฺขยํ. ธมฺมฏฺฐวคฺโค เอกูนวีสติโม.}
\csromangathagroup{\csromangathastanza{อหิํสา สพฺพปาณานํ, “อริโย”ติ ปวุจฺจติ. \\ สมฺพหุลสีลาทิสมฺปนฺนภิกฺขุวตฺถุ 271.\csromanspacer{} น สีลพฺพตมตฺเตน, พาหุสจฺเจน วา ปน.}\csromangathastanza{อถ วา สมาธิลาเภน, วิวิตฺตสยเนน วา. \\ ผุสามิ เนกฺขมฺมสุขํ, อปุถุชฺชนเสวิตํ. \\ ภิกฺขุ วิสฺสาสมาปาทิ, อปฺปตฺโต อาสวกฺขยํ.\csromanspacer{} ธมฺมฏฺฐวคฺโค เอกูนวีสติโม.}}

\csromanheader{1}{20. มคฺควคฺค}
\vspace{\tipitakaparskip}
\csromancenter{ปญฺจสตภิกฺขุวตฺถุ 273. * มคฺคานฏฺฐงฺคิโก เสฏฺโฐ, สจฺจานํ จตุโร ปทา.}
\vspace{0.5\baselineskip}
\csromangathameasure{วิราโค เสฏฺโฐ ธมฺมานํ, ทฺวิปทานญฺจ จกฺขุมา. \\ เอเสว มคฺโค นตฺถญฺโญ, ทสฺสนสฺส วิสุทฺธิยา. \\ เอตญฺหิ ตุเมฺห ปฏิปชฺชถ, มารสฺเสตํ ปโมหนํ. \\ เอตญฺหิ ตุเมฺห ปฏิปนฺนา, ทุกฺขสฺสนฺตํ กริสฺสถ. \\ อกฺขาโต โว มยา มคฺโค, อญฺญาย สลฺลกนฺตนํ3. \\ ตุเมฺหหิ กิจฺจมาตปฺปํ, อกฺขาตาโร ตถาคตา. \\ ปฏิปนฺนา ปโมกฺขนฺติ, ฌายิโน มารพนฺธนา.}
\csromangathagroup{\csromangathastanza{วิราโค เสฏฺโฐ ธมฺมานํ, ทฺวิปทานญฺจ จกฺขุมา. \\ เอเสว\footnote{เอโสว (สี, อิ)} มคฺโค นตฺถญฺโญ, ทสฺสนสฺส วิสุทฺธิยา.}\csromangathastanza{เอตญฺหิ ตุเมฺห ปฏิปชฺชถ, มารสฺเสตํ ปโมหนํ. \\ เอตญฺหิ ตุเมฺห ปฏิปนฺนา, ทุกฺขสฺสนฺตํ กริสฺสถ.\csromanspacer{}\csromanspacer{}}\csromangathastanza{อกฺขาโต โว\footnote{อกฺขาโต เว (สี, อิ) 3. สลฺลสนฺถนํ (สี, อิ), สลฺลสตฺถนํ (สฺยา)} มยา มคฺโค, อญฺญาย สลฺลกนฺตนํ3. \\ ตุเมฺหหิ กิจฺจมาตปฺปํ, อกฺขาตาโร ตถาคตา. \\ ปฏิปนฺนา ปโมกฺขนฺติ, ฌายิโน มารพนฺธนา.}}

\vspace{\tipitakaparskip}
\csromancenter{มคฺควคฺค \textbf{1.}\csromanspacer{}\textbf{ อนิจฺจลกฺขณวตฺถุ} 277. * “สพฺเพ สงฺขารา อนิจฺจา”ติ, ยทา ปญฺญาย ปสฺสติ.}
\vspace{0.5\baselineskip}
\csromangathameasure{อถ นิพฺพินฺทติ ทุกฺเข, เอส มคฺโค วิสุทฺธิยา.}
\csromangathagroup{\csromangathastanza{\csromanfolio{53}อถ นิพฺพินฺทติ ทุกฺเข, เอส มคฺโค วิสุทฺธิยา.}}

\vspace{\tipitakaparskip}
\csromancenter{ทุกฺขลกฺขณวตฺถุ 278. * “สพฺเพ สงฺขารา ทุกฺขา”ติ, ยทา ปญฺญาย ปสฺสติ.}
\vspace{0.5\baselineskip}
\csromangathameasure{อถ นิพฺพินฺทติ ทุกฺเข, เอส มคฺโค วิสุทฺธิยา.}
\csromangathagroup{\csromangathastanza{อถ นิพฺพินฺทติ ทุกฺเข, เอส มคฺโค วิสุทฺธิยา.}}

\csromanmatikaanchor{mtk.28}
\csromantocmarkref{section}{17. โกธวคฺค}{mtk.28}
\bookmark[dest={mtk.28},level=1]{17. โกธวคฺค}
\vspace{\tipitakaparskip}
\csromancenter{อนตฺตลกฺขณวตฺถุ 279. * “สพฺเพ ธมฺมา อนตฺตา”ติ, ยทา ปญฺญาย ปสฺสติ.}
\vspace{0.5\baselineskip}
\csromangathameasure{อถ นิพฺพินฺทติ ทุกฺเข, เอส มคฺโค วิสุทฺธิยา.}
\csromangathagroup{\csromangathastanza{อถ นิพฺพินฺทติ ทุกฺเข, เอส มคฺโค วิสุทฺธิยา.}}

\vspace{\tipitakaparskip}
\csromancenter{ปธานกมฺมิกติสฺสตฺเถรวตฺถุ 280.\csromanspacer{} อุฏฺฐานกาลมฺหิ อนุฏฺฐหาโน,}
\prose{ยุวา พลี อาลสิยํ อุเปโต.}

\csromangathameasure{สํสนฺนสงฺกปฺปมโน กุสีโต, \\ ปญฺญาย มคฺคํ อลโส น วินฺทติ.}
\csromangathagroup{\csromangathastanza{สํสนฺนสงฺกปฺปมโน\footnote{อสมฺปนฺนสงฺกปฺปมโน (ก)} กุสีโต, \\ ปญฺญาย มคฺคํ อลโส น วินฺทติ.}}

\vspace{\tipitakaparskip}
\csromancenter{สูกรเปตวตฺถุ 281. + วาจานุรกฺขี มนสา สุสํวุโต,}
\prose{กาเยน จ นากุสลํ กยิรา\footnote{อกุสลํ น กยิรา (สี, สฺยา, กํ, อิ)}.}

\csromangathameasure{เอเต ตโย กมฺมปเถ วิโสธเย, \\ อาราธเย มคฺค’มิสิปฺปเวทิตํ.}
\csromangathagroup{\csromangathastanza{เอเต ตโย กมฺมปเถ วิโสธเย, \\ อาราธเย มคฺค’มิสิปฺปเวทิตํ.}}

\vspace{\tipitakaparskip}
\csromancenter{โปฏฺฐิลตฺเถรวตฺถุ 282.\csromanspacer{} โยคา เว ชายตี\footnote{ชายเต (กตฺถจิ)} ภูริ, อโยคา ภูริสงฺขโย.}
\vspace{0.5\baselineskip}
\csromangathameasure{เอตํ เทฺวธาปถํ ญตฺวา, ภวาย วิภวาย จ. \\ ตถา’ตฺตานํ นิเวเสยฺย, ยถา ภูริ ปวฑฺฒติ.}
\csromangathagroup{\csromangathastanza{เอตํ เทฺวธาปถํ ญตฺวา, ภวาย วิภวาย จ. \\ ตถา’ตฺตานํ นิเวเสยฺย, ยถา ภูริ ปวฑฺฒติ.}}

\vspace{\tipitakaparskip}
\csromancenter{ธมฺมปทปาฬิ \textbf{8.}\csromanspacer{}\textbf{ ปญฺจมหลฺลกภิกฺขุวตฺถุ} 283.\csromanspacer{} วนํ ฉินฺทถ มา รุกฺขํ, วนโต ชายเต ภยํ.}
\vspace{0.5\baselineskip}
\csromangathameasure{เฉตฺวา วนญฺจ วนถญฺจ, นิพฺพนา โหถ ภิกฺขโว. \\ ยาว หิ วนโถ น ฉิชฺชติ, \\ อณุมตฺโตปิ นรสฺส นาริสุ. \\ ปฏิพทฺธมโนว ตาว โส, \\ วจฺโฉ ขีรปโกว มาตริ.}
\csromangathagroup{\csromangathastanza{\csromanfolio{54}เฉตฺวา วนญฺจ วนถญฺจ, นิพฺพนา โหถ ภิกฺขโว.}\csromangathastanza{ยาว หิ วนโถ น ฉิชฺชติ, \\ อณุมตฺโตปิ นรสฺส นาริสุ. \\ ปฏิพทฺธมโนว\footnote{ปฏิพนฺธมโนว (ก)} ตาว โส, \\ วจฺโฉ ขีรปโกว\footnote{ปฏิพนฺธมโนว (ก)} มาตริ.}}

\vspace{\tipitakaparskip}
\csromancenter{สุวณฺณการตฺเถรวตฺถุ 285. * อุจฺฉินฺท สิเนหมตฺตโน, กุมุทํ สารทิกํว ปาณินา.}
\vspace{0.5\baselineskip}
\csromangathameasure{สนฺติมคฺคเมว พฺรูหย, นิพฺพานํ สุคเตน เทสิตํ.}
\csromangathagroup{\csromangathastanza{สนฺติมคฺคเมว พฺรูหย, นิพฺพานํ สุคเตน เทสิตํ.}}

\csromanmatikaanchor{mtk.29}
\csromantocmarkref{section}{18. มลวคฺค}{mtk.29}
\bookmark[dest={mtk.29},level=1]{18. มลวคฺค}
\vspace{\tipitakaparskip}
\csromancenter{มหาธนวาณิชวตฺถุ 286.\csromanspacer{} อิธ วสฺสํ วสิสฺสามิ, อิธ เหมนฺตคิมฺหิสุ.}
\vspace{0.5\baselineskip}
\csromangathameasure{อิติ พาโล วิจินฺเตติ, อนฺตรายํ น พุชฺฌติ.}
\csromangathagroup{\csromangathastanza{อิติ พาโล วิจินฺเตติ, อนฺตรายํ น พุชฺฌติ.}}

\vspace{\tipitakaparskip}
\csromancenter{กิสาโคตมีวตฺถุ 287.\csromanspacer{} ตํ ปุตฺตปสุสมฺมตฺตํ, พฺยาสตฺตมนสํ นรํ.}
\vspace{0.5\baselineskip}
\csromangathameasure{สุตฺตํ คามํ มโหโฆว, มจฺจุ อาทาย คจฺฉติ.}
\csromangathagroup{\csromangathastanza{สุตฺตํ คามํ มโหโฆว, มจฺจุ อาทาย คจฺฉติ.}}

\vspace{\tipitakaparskip}
\csromancenter{ปฏาจาราวตฺถุ 288.\csromanspacer{} น สนฺติ ปุตฺตา ตาณาย, น ปิตา นาปิ พนฺธวา.}
\vspace{0.5\baselineskip}
\csromangathameasure{อนฺตเกนา’ธิปนฺนสฺส, นตฺถิ ญาตีสุ ตาณตา. \\ เอตมตฺถวสํ ญตฺวา, ปณฺฑิโต สีลสํวุโต. \\ นิพฺพานคมนํ มคฺคํ, ขิปฺปเมว วิโสธเย. มคฺควคฺโค วีสติโม.}
\csromangathagroup{\csromangathastanza{อนฺตเกนา’ธิปนฺนสฺส, นตฺถิ ญาตีสุ ตาณตา. \\ เอตมตฺถวสํ ญตฺวา, ปณฺฑิโต สีลสํวุโต. \\ นิพฺพานคมนํ มคฺคํ, ขิปฺปเมว วิโสธเย.\csromanspacer{} มคฺควคฺโค วีสติโม.}}

\csromanheader{1}{21. ปกิณฺณกวคฺค \textbf{21. ปกิณฺณกวคฺค}}
\vspace{\tipitakaparskip}
\csromancenter{อตฺตโน ปุพฺพกมฺมวตฺถุ 290.\csromanspacer{} มตฺตาสุขปริจฺจาคา, ปสฺเส เจ วิปุลํ สุขํ.}
\vspace{0.5\baselineskip}
\csromangathameasure{จเช มตฺตาสุขํ ธีโร, สมฺปสฺสํ วิปุลํ สุขํ.}
\csromangathagroup{\csromangathastanza{\csromanfolio{55}จเช มตฺตาสุขํ ธีโร, สมฺปสฺสํ วิปุลํ สุขํ.}}

\vspace{\tipitakaparskip}
\csromancenter{กุกฺกุฏณฺฑขาทิกาวตฺถุ 291.\csromanspacer{} ปรทุกฺขูปธาเนน, อตฺตโน\footnote{โย อตฺตโน (สฺยา, อิ, ก)} สุขมิจฺฉติ.}
\vspace{0.5\baselineskip}
\csromangathameasure{เวรสํสคฺคสํสฏฺโฐ, เวรา โส น ปริมุจฺจติ.}
\csromangathagroup{\csromangathastanza{เวรสํสคฺคสํสฏฺโฐ, เวรา โส น ปริมุจฺจติ.}}

\vspace{\tipitakaparskip}
\csromancenter{ภทฺทิยานํ ภิกฺขูนํ วตฺถุ 292.\csromanspacer{} ยํ หิ กิจฺจํ อปวิทฺธํ\footnote{ตทปวิทฺธํ (สี, สฺยา)}, อกิจฺจํ ปน กรียติ.}
\vspace{0.5\baselineskip}
\csromangathameasure{อุนฺนฬานํ ปมตฺตานํ, เตสํ วฑฺฒนฺติ อาสวา.}
\csromangathagroup{\csromangathastanza{อุนฺนฬานํ ปมตฺตานํ, เตสํ วฑฺฒนฺติ อาสวา.}}

\proseitem{293}{\csromansymbolfootnote{*}{ขุ 10. 27 ปิฏฺเฐ เนตฺติยมฺปิ.}เยสญฺจ สุสมารทฺธา, นิจฺจํ กายคตา สติ.}

\csromangathameasure{อกิจฺจํ เต น เสวนฺติ, กิจฺเจ สาตจฺจการิโน. \\ สตานํ สมฺปชานานํ, อตฺถํ คจฺฉนฺติ อาสวา.}
\csromangathagroup{\csromangathastanza{อกิจฺจํ เต น เสวนฺติ, กิจฺเจ สาตจฺจการิโน. \\ สตานํ สมฺปชานานํ, อตฺถํ คจฺฉนฺติ อาสวา.}}

\vspace{\tipitakaparskip}
\csromancenter{ลกุณฺฑกภทฺทิยวตฺถุ 294.\csromanspacer{} มาตรํ ปิตรํ หนฺตฺวา, ราชาโน เทฺว จ ขตฺติเย.}
\vspace{0.5\baselineskip}
\csromangathameasure{รฏฺฐํ สานุจรํ หนฺตฺวา, อนีโฆ ยาติ พฺราหฺมโณ.}
\csromangathagroup{\csromangathastanza{รฏฺฐํ สานุจรํ หนฺตฺวา, อนีโฆ ยาติ พฺราหฺมโณ.}}

\proseitem{295}{มาตรํ ปิตรํ หนฺตฺวา, ราชาโน เทฺว จ โสตฺถิเย.}

\csromangathameasure{เวยคฺฆปญฺจมํ หนฺตฺวา, อนีโฆ ยาติ พฺราหฺมโณ.}
\csromangathagroup{\csromangathastanza{เวยคฺฆปญฺจมํ หนฺตฺวา, อนีโฆ ยาติ พฺราหฺมโณ.}}

\vspace{\tipitakaparskip}
\csromancenter{ทารุสากฏิกปุตฺตวตฺถุ 296.\csromanspacer{} สุปฺปพุทฺธํ ปพุชฺฌนฺติ, สทา โคตมสาวกา.}
\vspace{0.5\baselineskip}
\csromangathameasure{เยสํ ทิวา จ รตฺโต จ, นิจฺจํ พุทฺธคตา สติ. \\ สุปฺปพุทฺธํ ปพุชฺฌนฺติ, สทา โคตมสาวกา. \\ เยสํ ทิวา จ รตฺโต จ, นิจฺจํ ธมฺมคตา สติ.}
\csromangathagroup{\csromangathastanza{เยสํ ทิวา จ รตฺโต จ, นิจฺจํ พุทฺธคตา สติ. \\ สุปฺปพุทฺธํ ปพุชฺฌนฺติ, สทา โคตมสาวกา. \\ เยสํ ทิวา จ รตฺโต จ, นิจฺจํ ธมฺมคตา สติ.}}

\csromanheader{2}{ธมฺมปทปาฬิ 298. สุปฺปพุทฺธํ ปพุชฺฌนฺติ, สทา โคตมสาวกา.}
\csromangathameasure{เยสํ ทิวา จ รตฺโต จ, นิจฺจํ สํฆคตา สติ. \\ สุปฺปพุทฺธํ ปพุชฺฌนฺติ, สทา โคตมสาวกา. \\ เยสํ ทิวา จ รตฺโต จ, นิจฺจํ กายคตา สติ. \\ สุปฺปพุทฺธํ ปพุชฺฌนฺติ, สทา โคตมสาวกา. \\ เยสํ ทิวา จ รตฺโต จ, อหิํสาย รโต มโน. \\ สุปฺปพุทฺธํ ปพุชฺฌนฺติ, สทา โคตมสาวกา. \\ เยสํ ทิวา จ รตฺโต จ, ภาวนาย รโต มโน.}
\csromangathagroup{\csromangathastanza{\csromanfolio{56}เยสํ ทิวา จ รตฺโต จ, นิจฺจํ สํฆคตา สติ. \\ สุปฺปพุทฺธํ ปพุชฺฌนฺติ, สทา โคตมสาวกา.}\csromangathastanza{เยสํ ทิวา จ รตฺโต จ, นิจฺจํ กายคตา สติ. \\ สุปฺปพุทฺธํ ปพุชฺฌนฺติ, สทา โคตมสาวกา.}\csromangathastanza{เยสํ ทิวา จ รตฺโต จ, อหิํสาย รโต มโน. \\ สุปฺปพุทฺธํ ปพุชฺฌนฺติ, สทา โคตมสาวกา. \\ เยสํ ทิวา จ รตฺโต จ, ภาวนาย รโต มโน.}}

\vspace{\tipitakaparskip}
\csromancenter{วชฺชิปุตฺตกภิกฺขุวตฺถุ 302.\csromanspacer{} ทุปฺปพฺพชฺชํ ทุรภิรมํ, ทุราวาสา ฆรา ทุขา.}
\vspace{0.5\baselineskip}
\csromangathameasure{ทุกฺโข’สมานสํวาโส, ทุกฺขานุปติตทฺธคู. \\ ตสฺมา น จทฺธคู สิยา, น จ ทุกฺขานุปติโต สิยา.}
\csromangathagroup{\csromangathastanza{ทุกฺโข’สมานสํวาโส, ทุกฺขานุปติตทฺธคู. \\ ตสฺมา น จทฺธคู สิยา, น จ\footnote{ตสฺมา น จทฺธคู น จ, (ก)} ทุกฺขานุปติโต สิยา\footnote{ทุกฺขานุปาติโต (?)}.}}

\vspace{\tipitakaparskip}
\csromancenter{จิตฺตคหปติวตฺถุ 303.\csromanspacer{} สทฺโธ สีเลน สมฺปนฺโน, ยโสโภคสมปฺปิโต.}
\vspace{0.5\baselineskip}
\csromangathameasure{ยํ ยํ ปเทสํ ภชติ, ตตฺถ ตตฺเถว ปูชิโต.}
\csromangathagroup{\csromangathastanza{ยํ ยํ ปเทสํ ภชติ, ตตฺถ ตตฺเถว ปูชิโต.}}

\vspace{\tipitakaparskip}
\csromancenter{จูฬสุภทฺทาวตฺถุ 304. * ทูเร สนฺโต ปกาเสนฺติ, หิมวนฺโตว ปพฺพโต.}
\vspace{0.5\baselineskip}
\csromangathameasure{อสนฺเตตฺถ น ทิสฺสนฺติ, รตฺติํ ขิตฺตา ยถา สรา.}
\csromangathagroup{\csromangathastanza{อสนฺเตตฺถ น ทิสฺสนฺติ, รตฺติํ ขิตฺตา ยถา สรา.}}

\vspace{\tipitakaparskip}
\csromancenter{เอกวิหาริตฺเถรวตฺถุ 305.\csromanspacer{} เอกาสนํ เอกเสยฺยํ, เอโก จรมตนฺทิโต.}
\vspace{0.5\baselineskip}
\csromangathameasure{เอโก ทมยมตฺตานํ, วนนฺเต รมิโต สิยา. ปกิณฺณกวคฺโค เอกวีสติโม.}
\csromangathagroup{\csromangathastanza{เอโก ทมยมตฺตานํ, วนนฺเต รมิโต สิยา.\csromanspacer{} ปกิณฺณกวคฺโค เอกวีสติโม.}}

\csromanheader{1}{22. นิรยวคฺค \textbf{22. นิรยวคฺค}}
\vspace{\tipitakaparskip}
\csromancenter{สุนฺทรี ปริพฺพาชิกาวตฺถุ 306. + อภูตวาที นิรยํ อุเปติ,}
\prose{\csromanfolio{57}โย วาปิ\footnote{โย จาปิ (สี, อิก)} กตฺวา น กโรมิ’จาห\footnote{น กโรมีติ จาห (สฺยา)}.}

\prose{อุโภปิ เต เปจฺจ สมา ภวนฺติ.}

\prose{นิหีนกมฺมา มนุชา ปรตฺถ.}

\vspace{\tipitakaparskip}
\csromancenter{ทุจฺจริตผลปีฬิตวตฺถุ 307. * กาสาวกณฺฐา พหโว, ปาปธมฺมา อสญฺญตา.}
\vspace{0.5\baselineskip}
\csromangathameasure{ปาปา ปาเปหิ กมฺเมหิ, นิรยํ เต อุปปชฺชเร.}
\csromangathagroup{\csromangathastanza{ปาปา ปาเปหิ กมฺเมหิ, นิรยํ เต อุปปชฺชเร.}}

\vspace{\tipitakaparskip}
\csromancenter{วคฺคุมุทาตีริยภิกฺขุวตฺถุ 308. * เสยฺโย อโยคุโฬ ภุตฺโต, ตตฺโต อคฺคิสิขูปโม.}
\vspace{0.5\baselineskip}
\csromangathameasure{ยญฺเจ ภุญฺเชยฺย ทุสฺสีโล, รฏฺฐปิณฺฑมสญฺญโต.}
\csromangathagroup{\csromangathastanza{ยญฺเจ ภุญฺเชยฺย ทุสฺสีโล, รฏฺฐปิณฺฑมสญฺญโต.}}

\csromanmatikaanchor{mtk.30}
\csromantocmarkref{section}{19. ธมฺมฏฺฐวคฺค}{mtk.30}
\bookmark[dest={mtk.30},level=1]{19. ธมฺมฏฺฐวคฺค}
\vspace{\tipitakaparskip}
\csromancenter{เขมกเสฏฺฐิปุตฺตวตฺถุ 309.\csromanspacer{} จตฺตาริ ฐานานิ นโร ปมตฺโต,}
\prose{อาปชฺชติ ปรทารูปเสวี.}

\csromangathameasure{อปุญฺญลาภํ น นิกามเสยฺยํ, \\ นินฺทํ ตตียํ นิรยํ จตุตฺถํ. \\ อปุญฺญลาโภ จ คตี จ ปาปิกา, \\ ภีตสฺส ภีตาย รตี จ โถกิกา. \\ ราชา จ ทณฺฑํ ครุกํ ปเณติ, \\ ตสฺมา นโร ปรทารํ น เสเว.}
\csromangathagroup{\csromangathastanza{อปุญฺญลาภํ น นิกามเสยฺยํ, \\ นินฺทํ ตตียํ นิรยํ จตุตฺถํ. \\ อปุญฺญลาโภ จ คตี จ ปาปิกา, \\ ภีตสฺส ภีตาย รตี จ โถกิกา. \\ ราชา จ ทณฺฑํ ครุกํ ปเณติ, \\ ตสฺมา นโร ปรทารํ น เสเว.}}

\vspace{\tipitakaparskip}
\csromancenter{ทุพฺพจภิกฺขุวตฺถุ 311.\csromanspacer{} กุโส ยถา ทุคฺคหิโต, หตฺถเมวา’นุกนฺตติ.}
\vspace{0.5\baselineskip}
\csromangathameasure{สามญฺญํ ทุปฺปรามฏฺฐํ, นิรยายุ’ปกฑฺฒติ.}
\csromangathagroup{\csromangathastanza{สามญฺญํ ทุปฺปรามฏฺฐํ, นิรยายุ’ปกฑฺฒติ.}}

\csromanheader{2}{ธมฺมปทปาฬิ 312. ยํกิญฺจิ สิถิลํ กมฺมํ, สํกิลิฏฺฐญฺจ ยํ วตํ.}
\csromangathameasure{สงฺกสฺสรํ พฺรหฺมจริยํ, น ตํ โหติ มหปฺผลํ. \\ กยิรา เจ กยิราเถนํ, ทฬฺหเมนํ ปรกฺกเม. \\ สิถิโล หิ ปริพฺพาโช, ภิยฺโย อากิรเต รชํ.}
\csromangathagroup{\csromangathastanza{\csromanfolio{58}สงฺกสฺสรํ พฺรหฺมจริยํ, น ตํ โหติ มหปฺผลํ. \\ กยิรา เจ กยิราเถนํ\footnote{กยิรา นํ (ก)}, ทฬฺหเมนํ ปรกฺกเม. \\ สิถิโล หิ ปริพฺพาโช, ภิยฺโย อากิรเต รชํ.}}

\vspace{\tipitakaparskip}
\csromancenter{อิสฺสาปกต-อิตฺถิวตฺถุ 314.\csromanspacer{} อกตํ ทุกฺกฏํ เสยฺโย, ปจฺฉา ตปฺปติ ทุกฺกฏํ.}
\vspace{0.5\baselineskip}
\csromangathameasure{กตญฺจ สุกตํ เสยฺโย, ยํ กตฺวา นานุตปฺปติ.}
\csromangathagroup{\csromangathastanza{กตญฺจ สุกตํ เสยฺโย, ยํ กตฺวา นานุตปฺปติ.}}

\vspace{\tipitakaparskip}
\csromancenter{สมฺพหุลภิกฺขุวตฺถุ 315.\csromanspacer{} นครํ ยถา ปจฺจนฺตํ, คุตฺตํ สนฺตรพาหิรํ.}
\vspace{0.5\baselineskip}
\csromangathameasure{เอวํ โคเปถ อตฺตานํ, ขโณ โว มา อุปจฺจคา. \\ ขณาตีตา หิ โสจนฺติ, นิรยมฺหิ สมปฺปิตา.}
\csromangathagroup{\csromangathastanza{เอวํ โคเปถ อตฺตานํ, ขโณ โว\footnote{ขโณ เว (สี, อิ, ก)} มา อุปจฺจคา. \\ ขณาตีตา หิ โสจนฺติ, นิรยมฺหิ สมปฺปิตา.}}

\vspace{\tipitakaparskip}
\csromancenter{นิคณฺฐวตฺถุ 316.\csromanspacer{} อลชฺชิตาเย ลชฺชนฺติ, ลชฺชิตาเย น ลชฺชเร.}
\vspace{0.5\baselineskip}
\csromangathameasure{มิจฺฉาทิฏฺฐิสมาทานา, สตฺตา คจฺฉนฺติ ทุคฺคติํ. \\ อภเย ภยทสฺสิโน, ภเย จาภยทสฺสิโน. \\ มิจฺฉาทิฏฺฐิสมาทานา, สตฺตา คจฺฉนฺติ ทุคฺคติํ.}
\csromangathagroup{\csromangathastanza{มิจฺฉาทิฏฺฐิสมาทานา, สตฺตา คจฺฉนฺติ ทุคฺคติํ. \\ อภเย ภยทสฺสิโน, ภเย จาภยทสฺสิโน. \\ มิจฺฉาทิฏฺฐิสมาทานา, สตฺตา คจฺฉนฺติ ทุคฺคติํ.}}

\vspace{\tipitakaparskip}
\csromancenter{ติตฺถิยสาวกวตฺถุ 318.\csromanspacer{} อวชฺเช วชฺชมติโน, วชฺเช จาวชฺชทสฺสิโน.}
\vspace{0.5\baselineskip}
\csromangathameasure{มิจฺฉาทิฏฺฐิสมาทานา, สตฺตา คจฺฉนฺติ ทุคฺคติํ. \\ วชฺชญฺจ วชฺชโต ญตฺวา, อวชฺชญฺจ อวชฺชโต. \\ สมฺมาทิฏฺฐิสมาทานา, สตฺตา คจฺฉนฺติ สุคฺคติํ. นิรยวคฺโค ทฺวาวีสติโม.}
\csromangathagroup{\csromangathastanza{มิจฺฉาทิฏฺฐิสมาทานา, สตฺตา คจฺฉนฺติ ทุคฺคติํ. \\ วชฺชญฺจ วชฺชโต ญตฺวา, อวชฺชญฺจ อวชฺชโต. \\ สมฺมาทิฏฺฐิสมาทานา, สตฺตา คจฺฉนฺติ สุคฺคติํ.\csromanspacer{} นิรยวคฺโค ทฺวาวีสติโม.}}

\csromanheader{1}{23. นาควคฺค \textbf{23. นาควคฺค}}
\vspace{\tipitakaparskip}
\csromancenter{อตฺตทนฺตวตฺถุ 320.\csromanspacer{} อหํ นาโคว สงฺคาเม, จาปโต ปติตํ สรํ.}
\vspace{0.5\baselineskip}
\csromangathameasure{อติวากฺยํ ติติกฺขิสฺสํ, ทุสฺสีโล หิ พหุชฺชโน.}
\csromangathagroup{\csromangathastanza{\csromanfolio{59}อติวากฺยํ ติติกฺขิสฺสํ, ทุสฺสีโล หิ พหุชฺชโน.}}

\proseitem{321}{\csromansymbolfootnote{*}{ขุ 7.187; ขุ 8. 66 ปิฏฺเฐสุปิ.}ทนฺตํ นยนฺติ สมิติํ, ทนฺตํ ราชา’ภิรูหติ.}

\csromangathameasure{ทนฺโต เสฏฺโฐ มนุสฺเสสุ, โย’ติวากฺยํ ติติกฺขติ.}
\csromangathagroup{\csromangathastanza{ทนฺโต เสฏฺโฐ มนุสฺเสสุ, โย’ติวากฺยํ ติติกฺขติ.}}

\proseitem{322}{\csromansymbolmark{*}วรมสฺสตรา ทนฺตา, อาชานียา จ\footnote{อาชานียาว (สฺยา)} สินฺธวา.}

\csromangathameasure{กุญฺชรา จ มหานาคา, อตฺตทนฺโต ตโต วรํ.}
\csromangathagroup{\csromangathastanza{กุญฺชรา จ\footnote{กุญฺชราว (สฺยา)} มหานาคา, อตฺตทนฺโต ตโต วรํ.}}

\csromanmatikaanchor{mtk.31}
\csromantocmarkref{section}{20. มคฺควคฺค}{mtk.31}
\bookmark[dest={mtk.31},level=1]{20. มคฺควคฺค}
\vspace{\tipitakaparskip}
\csromancenter{หตฺถาจริยปุพฺพกภิกฺขุวตฺถุ 323. * น หิ เอเตหิ ยาเนหิ, คจฺเฉยฺย อคตํ ทิสํ.}
\vspace{0.5\baselineskip}
\csromangathameasure{ยถา’ตฺตนา สุทนฺเตน, ทนฺโต ทนฺเตน คจฺฉติ.}
\csromangathagroup{\csromangathastanza{ยถา’ตฺตนา สุทนฺเตน, ทนฺโต ทนฺเตน คจฺฉติ.}}

\vspace{\tipitakaparskip}
\csromancenter{ปริชิณฺณพฺราหฺมณปุตฺตวตฺถุ 324.\csromanspacer{} ธนปาโล\footnote{ธนปาลโก (สี, สฺยา, กํ, อิ)} นาม กุญฺชโร,}
\prose{กฏุกเภทโน\footnote{กฏุกปฺปเภทโน (สี, สฺยา, อิ)} ทุนฺนิวารโย.}

\csromangathameasure{พทฺโธ กพฬํ น ภุญฺชติ, \\ สุมรติ นาควนสฺส กุญฺชโร.}
\csromangathagroup{\csromangathastanza{พทฺโธ กพฬํ น ภุญฺชติ, \\ สุมรติ\footnote{สุสรติ (ก)} นาควนสฺส กุญฺชโร.}}

\vspace{\tipitakaparskip}
\csromancenter{ปเสนทิโกสลวตฺถุ 325. + มิทฺธี ยทา โหติ มหคฺฆโส จ,}
\prose{นิทฺทายิตา สมฺปริวตฺตสายี.}

\csromangathameasure{มหาวราโหว นิวาปปุฏฺโฐ, \\ ปุนปฺปุนํ คพฺภมุเปติ มนฺโท.}
\csromangathagroup{\csromangathastanza{มหาวราโหว นิวาปปุฏฺโฐ, \\ ปุนปฺปุนํ คพฺภมุเปติ มนฺโท.}}

\csromanheader{2}{ธมฺมปทปาฬิ \textbf{5. สานุสามเณรวตฺถุ} 326. อิทํ ปุเร จิตฺตมจาริ จาริกํ,}
\prose{\csromanfolio{60}เยนิจฺฉกํ ยตฺถกามํ ยถาสุขํ.}

\csromangathameasure{ตทชฺชหํ นิคฺคเหสฺสามิ โยนิโส, \\ หตฺถิปฺปภินฺนํ วิย องฺกุสคฺคโห.}
\csromangathagroup{\csromangathastanza{ตทชฺชหํ นิคฺคเหสฺสามิ โยนิโส, \\ หตฺถิปฺปภินฺนํ วิย องฺกุสคฺคโห.}}

\vspace{\tipitakaparskip}
\csromancenter{ปาเวยฺยกหตฺถิวตฺถุ 327.\csromanspacer{} อปฺปมาทรตา โหถ, สจิตฺตมนุรกฺขถ.}
\vspace{0.5\baselineskip}
\csromangathameasure{ทุคฺคา อุทฺธรถ’ตฺตานํ, ปงฺเก สนฺโนว กุญฺชโร.}
\csromangathagroup{\csromangathastanza{ทุคฺคา อุทฺธรถ’ตฺตานํ, ปงฺเก สนฺโนว\footnote{สตฺโตว (สี, อิ)} กุญฺชโร.}}

\vspace{\tipitakaparskip}
\csromancenter{สมฺพหุลภิกฺขุวตฺถุ 328. * สเจ ลเภถ นิปกํ สหายํ,}
\prose{สทฺธิํจรํ สาธุวิหาริธีรํ.}

\csromangathameasure{อภิภุยฺย สพฺพานิ ปริสฺสยานิ, \\ จเรยฺย เตน’ตฺตมโน สตีมา.}
\csromangathagroup{\csromangathastanza{อภิภุยฺย สพฺพานิ ปริสฺสยานิ, \\ จเรยฺย เตน’ตฺตมโน สตีมา.}}

\proseitem{329}{\csromansymbolfootnote{*}{วิ 3. 496; ม 3. 192; อุปริ 286 ปิฏฺเฐสุปิ.}โน เจ ลเภถ นิปกํ สหายํ,}

\prose{สทฺธิํจรํ สาธุวิหาริธีรํ.}

\csromangathameasure{ราชาว รฏฺฐํ วิชิตํ ปหาย, \\ เอโก จเร มาตงฺค’รญฺเญว นาโค.}
\csromangathagroup{\csromangathastanza{ราชาว รฏฺฐํ วิชิตํ ปหาย, \\ เอโก จเร มาตงฺค’รญฺเญว นาโค.}}

\proseitem{330}{\csromansymbolmark{*}เอกสฺส จริตํ เสยฺโย,}

\prose{นตฺถิ พาเล สหายตา.}

\csromangathameasure{เอโก จเร น จ ปาปานิ กยิรา, \\ อปฺโปสฺสุกฺโก มาตงฺค’รญฺเญว นาโค.}
\csromangathagroup{\csromangathastanza{เอโก จเร น จ ปาปานิ กยิรา, \\ อปฺโปสฺสุกฺโก มาตงฺค’รญฺเญว นาโค.}}

\vspace{\tipitakaparskip}
\csromancenter{มารวตฺถุ 331.\csromanspacer{} อตฺถมฺหิ ชาตมฺหิ สุขา สหายา,}
\prose{ตุฏฺฐี สุขา ยา อิตรีตเรน.}

\csromangathameasure{ปุญฺญํ สุขํ ชีวิตสงฺขยมฺหิ, \\ สพฺพสฺส ทุกฺขสฺส สุขํ ปหานํ.}
\csromangathagroup{\csromangathastanza{ปุญฺญํ สุขํ ชีวิตสงฺขยมฺหิ, \\ สพฺพสฺส ทุกฺขสฺส สุขํ ปหานํ.}}

\csromanheader{1}{24. ตณฺหาวคฺค 332. สุขา มตฺเตยฺยตา โลเก,}
\prose{\csromanfolio{61}อโถ เปตฺเตยฺยตา สุขา.}

\csromangathameasure{สุขา สามญฺญตา โลเก, \\ อโถ พฺรหฺมญฺญตา สุขา. \\ สุขํ ยาว ชรา สีลํ, สุขา สทฺธา ปติฏฺฐิตา. \\ สุโข ปญฺญาย ปฏิลาโภ, ปาปานํ อกรณํ สุขํ. นาควคฺโค เตวีสติโม.}
\csromangathagroup{\csromangathastanza{สุขา สามญฺญตา โลเก, \\ อโถ พฺรหฺมญฺญตา สุขา.}\csromangathastanza{สุขํ ยาว ชรา สีลํ, สุขา สทฺธา ปติฏฺฐิตา. \\ สุโข ปญฺญาย ปฏิลาโภ, ปาปานํ อกรณํ สุขํ.\csromanspacer{} นาควคฺโค เตวีสติโม.}}

\chapterhead{24. ตณฺหาวคฺค}
\vspace{\tipitakaparskip}
\csromancenter{กปิลมจฺฉวตฺถุ 334.\csromanspacer{} มนุชสฺส ปมตฺตจาริโน,}
\prose{ตณฺหา วฑฺฒติ มาลุวา วิย.}

\csromanmatikaanchor{mtk.32}
\csromantocmarkref{section}{21. ปกิณฺณกวคฺค}{mtk.32}
\bookmark[dest={mtk.32},level=1]{21. ปกิณฺณกวคฺค}
\csromangathameasure{โส ปฺลวตี หุรา หุรํ, \\ ผลมิจฺฉํว วนสฺมิ วานโร. \\ ยํ เอสา สหเต ชมฺมี, ตณฺหา โลเก วิสตฺติกา. \\ โสกา ตสฺส ปวฑฺฒนฺติ, อภิวฏฺฐํว พีรณํ. \\ โย เจตํ สหเต ชมฺมิํ, ตณฺหํ โลเก ทุรจฺจยํ. \\ โสกา ตมฺหา ปปตนฺติ, อุทพินฺทุว โปกฺขรา. \\ ตํ โว วทามิ ภทฺทํ โว, ยาวนฺเต'ตฺถ สมาคตา. \\ ตณฺหาย มูลํ ขณถ, อุสีรตฺโถว พีรณํ. \\ มา โว นฬํว โสโตว, มาโร ภญฺชิ ปุนปฺปุนํ.}
\csromangathagroup{\csromangathastanza{โส ปฺลวตี\footnote{ปฺลวติ (สี, อิ), ปลเวตี (ก), อุปฺลวตี (?)} หุรา หุรํ, \\ ผลมิจฺฉํว วนสฺมิ วานโร.}\csromangathastanza{ยํ เอสา สหเต ชมฺมี, ตณฺหา โลเก วิสตฺติกา. \\ โสกา ตสฺส ปวฑฺฒนฺติ, อภิวฏฺฐํว\footnote{อภิวฑฺฒํว (สฺยา), อภิวฏฺฏํว (อิ), อภิวุฑฺฒํว (ก)} พีรณํ.}\csromangathastanza{โย เจตํ สหเต ชมฺมิํ, ตณฺหํ โลเก ทุรจฺจยํ. \\ โสกา ตมฺหา ปปตนฺติ, อุทพินฺทุว โปกฺขรา.}\csromangathastanza{ตํ โว วทามิ ภทฺทํ โว, ยาวนฺเต'ตฺถ สมาคตา. \\ ตณฺหาย มูลํ ขณถ, อุสีรตฺโถว พีรณํ. \\ มา โว นฬํว โสโตว, มาโร ภญฺชิ ปุนปฺปุนํ.}}

\csromanheader{2}{ธมฺมปทปาฬิ \textbf{2. สูกรโปติกาวตฺถุ} 338. * ยถาปิ มูเล อนุปทฺทเว ทเฬฺห,}
\prose{\csromanfolio{62}ฉินฺโนปิ รุกฺโข ปุนเรว รูหติ.}

\csromangathameasure{เอวมฺปิ ตณฺหานุสเย อนูหเต, \\ นิพฺพตฺตตี ทุกฺขมิทํ ปุนปฺปุนํ. \\ ยสฺส ฉตฺติํสติ โสตา, มนาปสวนา ภุสา. \\ มหา วหนฺติ ทุทฺทิฏฺฐิํ, สงฺกปฺปา ราคนิสฺสิตา. \\ สวนฺติ สพฺพธิ โสตา, ลตา อุปฺปชฺช ติฏฺฐติ. \\ ตญฺจ ทิสฺวา ลตํ ชาตํ, มูลํ ปญฺญาย ฉินฺทถ. \\ สริตานิ สิเนหิตานิ จ, \\ โสมนสฺสานิ ภวนฺติ ชนฺตุโน. \\ เต สาตสิตา สุเขสิโน, \\ เต เว ชาติชรูปคา นรา. \\ ตสิณาย ปุรกฺขตา ปชา, \\ ปริสปฺปนฺติ สโสว พนฺธิโต. \\ สํโยชนสงฺคสตฺตกา, \\ ทุกฺขมุเปนฺติ ปุนปฺปุนํ จิราย. \\ ตสิณาย ปุรกฺขตา ปชา. \\ ปริสปฺปนฺติ สโสว พนฺธิโต. \\ ตสฺมา ตสิณํ วิโนทเย, \\ อากงฺขนฺต วิราคมตฺตโน.}
\csromangathagroup{\csromangathastanza{เอวมฺปิ ตณฺหานุสเย อนูหเต, \\ นิพฺพตฺตตี ทุกฺขมิทํ ปุนปฺปุนํ.}\csromangathastanza{ยสฺส ฉตฺติํสติ โสตา, มนาปสวนา ภุสา. \\ มหา\footnote{วาหา (สี, สฺยา, อิ)} วหนฺติ ทุทฺทิฏฺฐิํ, สงฺกปฺปา ราคนิสฺสิตา.}\csromangathastanza{สวนฺติ สพฺพธิ โสตา, ลตา อุปฺปชฺช\footnote{อุพฺภิชฺช (สี, สฺยา, กํ, อิ)} ติฏฺฐติ. \\ ตญฺจ ทิสฺวา ลตํ ชาตํ, มูลํ ปญฺญาย ฉินฺทถ.}\csromangathastanza{สริตานิ สิเนหิตานิ จ, \\ โสมนสฺสานิ ภวนฺติ ชนฺตุโน. \\ เต สาตสิตา สุเขสิโน, \\ เต เว ชาติชรูปคา นรา.}\csromangathastanza{ตสิณาย ปุรกฺขตา ปชา, \\ ปริสปฺปนฺติ สโสว พนฺธิโต\footnote{พาธิโต (พหูสุ)}. \\ สํโยชนสงฺคสตฺตกา, \\ ทุกฺขมุเปนฺติ ปุนปฺปุนํ จิราย.}\csromangathastanza{ตสิณาย ปุรกฺขตา ปชา. \\ ปริสปฺปนฺติ สโสว พนฺธิโต. \\ ตสฺมา ตสิณํ วิโนทเย, \\ อากงฺขนฺต\footnote{ภิกฺขุ อากงฺขี (สี), ภิกฺขุ อากงฺขํ (สฺยา)} วิราคมตฺตโน.}}

\vspace{\tipitakaparskip}
\csromancenter{วิพฺภนฺตภิกฺขุวตฺถุ 344.\csromanspacer{} โย นิพฺพนโถ วนาธิมุตฺโต,}
\prose{วนมุตฺโต วนเมว ธาวติ.}

\csromangathameasure{ตํ ปุคฺคลเมถ ปสฺสถ, \\ มุตฺโต พนฺธนเมว ธาวติ.}
\csromangathagroup{\csromangathastanza{ตํ ปุคฺคลเมถ ปสฺสถ, \\ มุตฺโต พนฺธนเมว ธาวติ.}}

\chapterheadpageread{24. ตณฺหาวคฺค \textbf{4. พนฺธนาคารวตฺถุ} 345. * น ตํ ทฬฺหํ พนฺธนมาหุ ธีรา,}
\prose{\csromanfolio{63}ยทายสํ ทารุชปพฺพชญฺจ\footnote{ทารุชํ พพฺพชญฺจ (สี, อิ)}.}

\csromangathameasure{สารตฺตรตฺตา มณิกุณฺฑเลสุ, \\ ปุตฺเตสุ ทาเรสุ จ ยา อเปกฺขา. \\ เอตํ ทฬฺหํ พนฺธนมาหุ ธีรา, \\ โอหารินํ สิถิลํ ทุปฺปมุญฺจํ. \\ เอตมฺปิ เฉตฺวาน ปริพฺพชนฺติ, \\ อนเปกฺขิโน กามสุขํ ปหาย.}
\csromangathagroup{\csromangathastanza{สารตฺตรตฺตา มณิกุณฺฑเลสุ, \\ ปุตฺเตสุ ทาเรสุ จ ยา อเปกฺขา. \\ เอตํ ทฬฺหํ พนฺธนมาหุ ธีรา, \\ โอหารินํ สิถิลํ ทุปฺปมุญฺจํ. \\ เอตมฺปิ เฉตฺวาน ปริพฺพชนฺติ, \\ อนเปกฺขิโน กามสุขํ ปหาย.}}

\vspace{\tipitakaparskip}
\csromancenter{เขมาเถรีวตฺถุ 347.\csromanspacer{} เย ราครตฺตา’นุปตนฺติ โสตํ,}
\prose{สยํกตํ มกฺกฏโกว ชาลํ.}

\csromangathameasure{เอตมฺปิ เฉตฺวาน วชนฺติ ธีรา, \\ อนเปกฺขิโน สพฺพทุกฺขํ ปหาย.}
\csromangathagroup{\csromangathastanza{เอตมฺปิ เฉตฺวาน วชนฺติ ธีรา, \\ อนเปกฺขิโน สพฺพทุกฺขํ ปหาย.}}

\vspace{\tipitakaparskip}
\csromancenter{อุคฺคเสนวตฺถุ 348.\csromanspacer{} มุญฺจ ปุเร มุญฺจ ปจฺฉโต,}
\prose{มชฺเฌ มุญฺจ ภวสฺส ปารคู.}

\csromangathameasure{สพฺพตฺถ วิมุตฺตมานโส, \\ น ปุนํ ชาติชรํ อุเปหิสิ.}
\csromangathagroup{\csromangathastanza{สพฺพตฺถ วิมุตฺตมานโส, \\ น ปุนํ ชาติชรํ อุเปหิสิ.}}

\vspace{\tipitakaparskip}
\csromancenter{จูฬธนุคฺคหปณฺฑิตวตฺถุ 349.\csromanspacer{} วิตกฺกมถิตสฺส ชนฺตุโน,}
\prose{ติพฺพราคสฺส สุภานุปสฺสิโน.}

\csromangathameasure{ภิยฺโย ตณฺหา ปวฑฺฒติ, \\ เอส โข ทฬฺหํ กโรติ พนฺธนํ.}
\csromangathagroup{\csromangathastanza{ภิยฺโย ตณฺหา ปวฑฺฒติ, \\ เอส โข ทฬฺหํ\footnote{เอส คาฬํ (ก)} กโรติ พนฺธนํ.}}

\csromanheader{2}{ธมฺมปทปาฬิ 350. วิตกฺกูปสเม จ\footnote{วิตกฺกูปสเมว (ก)} โย รโต,}
\prose{\csromanfolio{64}อสุภํ ภาวยเต สทา สโต.}

\csromangathameasure{เอส โข พฺยนฺติ กาหิติ, \\ เอส2 เฉจฺฉติ มารพนฺธนํ.}
\csromangathagroup{\csromangathastanza{เอส\footnote{เอโส (?)} โข พฺยนฺติ กาหิติ, \\ เอส2 เฉจฺฉติ มารพนฺธนํ.}}

\csromanmatikaanchor{mtk.33}
\csromantocmarkref{section}{22. นิรยวคฺค}{mtk.33}
\bookmark[dest={mtk.33},level=1]{22. นิรยวคฺค}
\vspace{\tipitakaparskip}
\csromancenter{มารวตฺถุ 351.\csromanspacer{} นิฏฺฐงฺคโต อสนฺตาสี, วีตตโณฺห อนงฺคโณ.}
\vspace{0.5\baselineskip}
\csromangathameasure{อจฺฉินฺทิ ภวสลฺลานิ, อนฺติโมยํ สมุสฺสโย. \\ วีตตโณฺห อนาทาโน, นิรุตฺติปทโกวิโท. \\ อกฺขรานํ สนฺนิปาตํ, ชญฺญา ปุพฺพาปรานิ จ. \\ ส เว “อนฺติมสารีโร, มหาปญฺโญ มหาปุริโส”ติ วุจฺจติ.}
\csromangathagroup{\csromangathastanza{อจฺฉินฺทิ ภวสลฺลานิ, อนฺติโมยํ สมุสฺสโย. \\ วีตตโณฺห อนาทาโน, นิรุตฺติปทโกวิโท.}\csromangathastanza{อกฺขรานํ สนฺนิปาตํ, ชญฺญา ปุพฺพาปรานิ จ. \\ ส เว “อนฺติมสารีโร, มหาปญฺโญ มหาปุริโส”ติ วุจฺจติ.}}

\vspace{\tipitakaparskip}
\csromancenter{อุปกาชีวกวตฺถุ 353. * สพฺพาภิภู สพฺพวิทูหมสฺมิ,}
\prose{สพฺเพสุ ธมฺเมสุ อนูปลิตฺโต.}

\csromangathameasure{สพฺพญฺชโห ตณฺหกฺขเย วิมุตฺโต, \\ สยํ อภิญฺญาย กมุทฺทิเสยํ.}
\csromangathagroup{\csromangathastanza{สพฺพญฺชโห ตณฺหกฺขเย วิมุตฺโต, \\ สยํ อภิญฺญาย กมุทฺทิเสยํ.}}

\vspace{\tipitakaparskip}
\csromancenter{สกฺกปญฺหวตฺถุ 354.\csromanspacer{} สพฺพทานํ ธมฺมทานํ ชินาติ, สพฺพรสํ ธมฺมรโส ชินาติ.}
\vspace{0.5\baselineskip}
\csromangathameasure{สพฺพรติํ ธมฺมรติ ชินาติ, ตณฺหกฺขโย สพฺพทุกฺขํ ชินาติ.}
\csromangathagroup{\csromangathastanza{สพฺพรติํ ธมฺมรติ ชินาติ, ตณฺหกฺขโย สพฺพทุกฺขํ ชินาติ.}}

\vspace{\tipitakaparskip}
\csromancenter{อปุตฺตกเสฏฺฐิวตฺถุ 355.\csromanspacer{} หนนฺติ โภคา ทุมฺเมธํ, โน จ ปารคเวสิโน.}
\vspace{0.5\baselineskip}
\csromangathameasure{โภคตณฺหาย ทุมฺเมโธ, หนฺติ อญฺเญว อตฺตนํ.}
\csromangathagroup{\csromangathastanza{โภคตณฺหาย ทุมฺเมโธ, หนฺติ อญฺเญว อตฺตนํ.}}

\vspace{\tipitakaparskip}
\csromancenter{องฺกุรวตฺถุ 356.\csromanspacer{} ติณโทสานิ เขตฺตานิ, ราคโทสา อยํ ปชา.}
\vspace{0.5\baselineskip}
\csromangathameasure{ตสฺมา หิ วีตราเคสุ, ทินฺนํ โหติ มหปฺผลํ.}
\csromangathagroup{\csromangathastanza{ตสฺมา หิ วีตราเคสุ, ทินฺนํ โหติ มหปฺผลํ.}}

\csromanheader{1}{25. ภิกฺขุวคฺค 357. ติณโทสานิ เขตฺตานิ, โทสโทสา อยํ ปชา.}
\csromangathameasure{ตสฺมา หิ วีตโทเสสุ, ทินฺนํ โหติ มหปฺผลํ. \\ ติณโทสานิ เขตฺตานิ, โมหโทสา อยํ ปชา. \\ ตสฺมา หิ วีตโมเหสุ, ทินฺนํ โหติ มหปฺผลํ. \\ ติณโทสานิ เขตฺตานิ, อิจฺฉาโทสา อยํ ปชา. \\ ตสฺมา หิ วิคติจฺเฉสุ, ทินฺนํ โหติ มหปฺผลํ. ( ) ตณฺหาวคฺโค จตุวีสติโม.}
\csromangathagroup{\csromangathastanza{\csromanfolio{65}ตสฺมา หิ วีตโทเสสุ, ทินฺนํ โหติ มหปฺผลํ. \\ ติณโทสานิ เขตฺตานิ, โมหโทสา อยํ ปชา.}\csromangathastanza{ตสฺมา หิ วีตโมเหสุ, ทินฺนํ โหติ มหปฺผลํ. \\ ติณโทสานิ เขตฺตานิ, อิจฺฉาโทสา อยํ ปชา. \\ ตสฺมา หิ วิคติจฺเฉสุ, ทินฺนํ โหติ มหปฺผลํ. ( )\footnote{(ติณโทสานิ เขตฺตานิ, ตณฺหาโทสา อยํ ปชา. ตสฺมา หิ วีตตเณฺหสุ, ทินฺนํ โหติ มหปฺผลํ.) กตฺถจิ โปตฺถเก ทิสฺสติ.} ตณฺหาวคฺโค จตุวีสติโม.}}

\chapterhead{25. ภิกฺขุวคฺค}
\vspace{\tipitakaparskip}
\csromancenter{ปญฺจภิกฺขุวตฺถุ 360.\csromanspacer{} จกฺขุนา สํวโร สาธุ, สาธุ โสเตน สํวโร.}
\vspace{0.5\baselineskip}
\csromangathameasure{ฆาเนน สํวโร สาธุ, สาธุ ชิวฺหาย สํวโร. \\ กาเยน สํวโร สาธุ, สาธุ วาจาย สํวโร. \\ มนสา สํวโร สาธุ, สาธุ สพฺพตฺถ สํวโร. \\ สพฺพตฺถ สํวุโต ภิกฺขุ, สพฺพทุกฺขา ปมุจฺจติ.}
\csromangathagroup{\csromangathastanza{ฆาเนน สํวโร สาธุ, สาธุ ชิวฺหาย สํวโร. \\ กาเยน สํวโร สาธุ, สาธุ วาจาย สํวโร.}\csromangathastanza{มนสา สํวโร สาธุ, สาธุ สพฺพตฺถ สํวโร. \\ สพฺพตฺถ สํวุโต ภิกฺขุ, สพฺพทุกฺขา ปมุจฺจติ.}}

\vspace{\tipitakaparskip}
\csromancenter{หํสฆาตกภิกฺขุวตฺถุ 362.\csromanspacer{} หตฺถสํยโต ปาทสํยโต,}
\prose{วาจาสํยโต สํยตุตฺตโม.}

\csromangathameasure{อชฺฌตฺตรโต สมาหิโต, \\ เอโก สนฺตุสิโต ตมาหุ ภิกฺขุํ.}
\csromangathagroup{\csromangathastanza{อชฺฌตฺตรโต สมาหิโต, \\ เอโก สนฺตุสิโต ตมาหุ ภิกฺขุํ.}}

\vspace{\tipitakaparskip}
\csromancenter{โกกาลิกวตฺถุ 363.\csromanspacer{} โย มุขสํยโต ภิกฺขุ, มนฺตภาณี อนุทฺธโต.}
\vspace{0.5\baselineskip}
\csromangathameasure{อตฺถํ ธมฺมญฺจ ทีเปติ, มธุรํ ตสฺส ภาสิตํ.}
\csromangathagroup{\csromangathastanza{อตฺถํ ธมฺมญฺจ ทีเปติ, มธุรํ ตสฺส ภาสิตํ.}}

\vspace{\tipitakaparskip}
\csromancenter{ธมฺมปทปาฬิ \textbf{4.}\csromanspacer{}\textbf{ ธมฺมารามตฺเถรวตฺถุ} 364.\csromanspacer{} ธมฺมาราโม ธมฺมรโต, ธมฺมํ อนุวิจินฺตยํ.}
\vspace{0.5\baselineskip}
\csromangathameasure{ธมฺมํ อนุสฺสรํ ภิกฺขุ, สทฺธมฺมา น ปริหายติ.}
\csromangathagroup{\csromangathastanza{\csromanfolio{66}ธมฺมํ อนุสฺสรํ ภิกฺขุ, สทฺธมฺมา น ปริหายติ.}}

\vspace{\tipitakaparskip}
\csromancenter{วิปกฺขเสวกภิกฺขุวตฺถุ 365.\csromanspacer{} สลาภํ นาติมญฺเญยฺย, นา’ญฺเญสํ ปิหยํ จเร.}
\vspace{0.5\baselineskip}
\csromangathameasure{อญฺเญสํ ปิหยํ ภิกฺขุ, สมาธิํ นาธิคจฺฉติ. \\ อปฺปลาโภปิ เจ ภิกฺขุ, สลาภํ นาติมญฺญติ. \\ ตํ เว เทวา ปสํสนฺติ, สุทฺธาชีวิํ อตนฺทิตํ.}
\csromangathagroup{\csromangathastanza{อญฺเญสํ ปิหยํ ภิกฺขุ, สมาธิํ นาธิคจฺฉติ. \\ อปฺปลาโภปิ เจ ภิกฺขุ, สลาภํ นาติมญฺญติ. \\ ตํ เว เทวา ปสํสนฺติ, สุทฺธาชีวิํ อตนฺทิตํ.}}

\vspace{\tipitakaparskip}
\csromancenter{ปญฺจคฺคทายกพฺราหฺมณวตฺถุ 367.\csromanspacer{} สพฺพโส นามรูปสฺมิํ, ยสฺส นตฺถิ มมายิตํ.}
\vspace{0.5\baselineskip}
\csromanmatikaanchor{mtk.34}
\csromantocmarkref{section}{23. นาควคฺค}{mtk.34}
\bookmark[dest={mtk.34},level=1]{23. นาควคฺค}
\csromangathameasure{อสตา จ น โสจติ, ส เว “ภิกฺขู”ติ วุจฺจติ.}
\csromangathagroup{\csromangathastanza{อสตา จ น โสจติ, ส เว “ภิกฺขู”ติ วุจฺจติ.}}

\vspace{\tipitakaparskip}
\csromancenter{สมฺพหุลภิกฺขุวตฺถุ 368.\csromanspacer{} เมตฺตาวิหารี โย ภิกฺขุ, ปสนฺโน พุทฺธสาสเน.}
\vspace{0.5\baselineskip}
\csromangathameasure{อธิคจฺเฉ ปทํ สนฺตํ, สงฺขารูปสมํ สุขํ. \\ สิญฺจ ภิกฺขุ อิมํ นาวํ, สิตฺตา เต ลหุเมสฺสติ. \\ เฉตฺวา ราคญฺจ โทสญฺจ, ตโต นิพฺพานเมหิสิ. \\ ปญฺจ ฉินฺเท ปญฺจ ชเห, ปญฺจ จุตฺตริ ภาวเย. \\ ปญฺจสงฺคาติโค ภิกฺขุ, “โอฆติณฺโณ”ติ วุจฺจติ. \\ ฌาย ภิกฺขุ มา ปมาโท, \\ มา เต กามคุเณ รเมสฺสุ จิตฺตํ. \\ มา โลหคุฬํ คิลี ปมตฺโต, \\ มา กนฺทิ “ทุกฺขมิทนฺ”ติ ทยฺหมาโน. \\ นตฺถิ ฌานํ อปญฺญสฺส, ปญฺญา นตฺถิ อฌายโต. \\ ยมฺหิ ฌานญฺจ ปญฺญา จ, ส เว นิพฺพานสนฺติเก.}
\csromangathagroup{\csromangathastanza{อธิคจฺเฉ ปทํ สนฺตํ, สงฺขารูปสมํ สุขํ. \\ สิญฺจ ภิกฺขุ อิมํ นาวํ, สิตฺตา เต ลหุเมสฺสติ.}\csromangathastanza{เฉตฺวา ราคญฺจ โทสญฺจ, ตโต นิพฺพานเมหิสิ. \\ ปญฺจ ฉินฺเท ปญฺจ ชเห, ปญฺจ จุตฺตริ ภาวเย. \\ ปญฺจสงฺคาติโค ภิกฺขุ, “โอฆติณฺโณ”ติ วุจฺจติ.}\csromangathastanza{ฌาย ภิกฺขุ\footnote{ฌาย ตุวํ ภิกฺขุ (?)} มา ปมาโท\footnote{มา จ ปมาโท (สี, สฺยา, อิ)}, \\ มา เต กามคุเณ รเมสฺสุ\footnote{ฌาย ตุวํ ภิกฺขุ (?)} จิตฺตํ. \\ มา โลหคุฬํ คิลี ปมตฺโต, \\ มา กนฺทิ “ทุกฺขมิทนฺ”ติ ทยฺหมาโน. \\ นตฺถิ ฌานํ อปญฺญสฺส, ปญฺญา นตฺถิ อฌายโต\footnote{อฌายิโน (ก)}. \\ ยมฺหิ ฌานญฺจ ปญฺญา จ, ส เว นิพฺพานสนฺติเก.}}

\chapterheadpageread{25. ภิกฺขุวคฺค 373. สุญฺญาคารํ ปวิฏฺฐสฺส, สนฺตจิตฺตสฺส ภิกฺขุโน.}
\csromangathameasure{อมานุสี รติ โหติ, สมฺมา ธมฺมํ วิปสฺสโต. \\ ยโต ยโต สมฺมสติ, ขนฺธานํ อุทยพฺพยํ. \\ ลภตี ปีติปาโมชฺชํ, อมตํ ตํ วิชานตํ. \\ ตตฺรายมาทิ ภวติ, อิธ ปญฺญสฺส ภิกฺขุโน. \\ อินฺทฺริยคุตฺติ สนฺตุฏฺฐิ, ปาติโมกฺเข จ สํวโร. \\ มิตฺเต ภชสฺสุ กลฺยาเณ, สุทฺธาชีเว อตนฺทิเต. \\ ปฏิสนฺถารวุตฺยสฺส, อาจารกุสโล สิยา. \\ ตโต ปาโมชฺชพหุโล, ทุกฺขสฺสนฺตํ กริสฺสติ.}
\csromangathagroup{\csromangathastanza{\csromanfolio{67}อมานุสี รติ โหติ, สมฺมา ธมฺมํ วิปสฺสโต. \\ ยโต ยโต สมฺมสติ, ขนฺธานํ อุทยพฺพยํ.}\csromangathastanza{ลภตี\footnote{ลภติ (อิ), ลภเต (ก)} ปีติปาโมชฺชํ, อมตํ ตํ วิชานตํ. \\ ตตฺรายมาทิ ภวติ, อิธ ปญฺญสฺส ภิกฺขุโน.}\csromangathastanza{อินฺทฺริยคุตฺติ สนฺตุฏฺฐิ, ปาติโมกฺเข จ สํวโร. \\ มิตฺเต ภชสฺสุ กลฺยาเณ, สุทฺธาชีเว อตนฺทิเต.}\csromangathastanza{ปฏิสนฺถารวุตฺยสฺส\footnote{ปฏิสนฺธารวุตฺยสฺส (ก)}, อาจารกุสโล สิยา. \\ ตโต ปาโมชฺชพหุโล, ทุกฺขสฺสนฺตํ กริสฺสติ.}}

\vspace{\tipitakaparskip}
\csromancenter{ปญฺจสตภิกฺขุวตฺถุ 377.\csromanspacer{} วสฺสิกา วิย ปุปฺผานิ, มทฺทวานิ\footnote{มชฺชวานิ (ก, ฏีกา), ปจฺจวานิ (ก-ฏฺฐ)} ปมุญฺจติ.}
\vspace{0.5\baselineskip}
\csromangathameasure{เอวํ ราคญฺจ โทสญฺจ, วิปฺปมุญฺเจถ ภิกฺขโว.}
\csromangathagroup{\csromangathastanza{เอวํ ราคญฺจ โทสญฺจ, วิปฺปมุญฺเจถ ภิกฺขโว.}}

\vspace{\tipitakaparskip}
\csromancenter{สนฺตกายตฺเถรวตฺถุ 378.\csromanspacer{} สนฺตกาโย สนฺตวาโจ, สนฺตวา สุสมาหิโต\footnote{สนฺตมโน สุสมาหิโต (สฺยา, อิ), สนฺตมโน สมาหิโต (ก)}.}
\vspace{0.5\baselineskip}
\csromangathameasure{วนฺตโลกามิโส ภิกฺขุ, “อุปสนฺโต”ติ วุจฺจติ.}
\csromangathagroup{\csromangathastanza{วนฺตโลกามิโส ภิกฺขุ, “อุปสนฺโต”ติ วุจฺจติ.}}

\vspace{\tipitakaparskip}
\csromancenter{นงฺคลกุลตฺเถรวตฺถุ 379.\csromanspacer{} อตฺตนา โจทยตฺตานํ, ปฏิมํเสถ อตฺตนา\footnote{ปฏิมาเส อตฺตมตฺตนา (สี, อิ), ปฏิมํเส ตมตฺตนา (สฺยา)}.}
\vspace{0.5\baselineskip}
\csromangathameasure{โส อตฺตคุตฺโต สติมา, สุขํ ภิกฺขุ วิหาหิสิ.}
\csromangathagroup{\csromangathastanza{โส อตฺตคุตฺโต สติมา, สุขํ ภิกฺขุ วิหาหิสิ.}}

\proseitem{380}{อตฺตา หิ อตฺตโน นาโถ,}

\prose{(โก หิ นาโถ ปโร สิยา)\footnote{( ) วิเทสโปตฺถเกสุ นตฺถิ.}}

\prose{อตฺตา หิ อตฺตโน คติ.}

\csromangathameasure{ตสฺมา สํยมมตฺตานํ, \\ อสฺสํ ภทฺรํว วาณิโช.}
\csromangathagroup{\csromangathastanza{ตสฺมา สํยมมตฺตานํ\footnote{สํยมย’ตฺตานํ (สี, อิ)}, \\ อสฺสํ ภทฺรํว วาณิโช.}}

\vspace{\tipitakaparskip}
\csromancenter{ธมฺมปทปาฬิ \textbf{11.}\csromanspacer{}\textbf{ วกฺกลิตฺเถรวตฺถุ} 381.\csromanspacer{} ปาโมชฺชพหุโล ภิกฺขุ, ปสนฺโน พุทฺธสาสเน.}
\vspace{0.5\baselineskip}
\csromangathameasure{อธิคจฺเฉ ปทํ สนฺตํ, สงฺขารูปสมํ สุขํ.}
\csromangathagroup{\csromangathastanza{\csromanfolio{68}อธิคจฺเฉ ปทํ สนฺตํ, สงฺขารูปสมํ สุขํ.}}

\vspace{\tipitakaparskip}
\csromancenter{สุมนสามเณรวตฺถุ 382.\csromanspacer{} โย หเว ทหโร ภิกฺขุ, ยุญฺชติ พุทฺธสาสเน.}
\vspace{0.5\baselineskip}
\csromangathameasure{โส’มํ โลกํ ปภาเสติ, อพฺภา มุตฺโตว จนฺทิมา. ภิกฺขุวคฺโค ปญฺจวีสติโม.}
\csromangathagroup{\csromangathastanza{โส’มํ\footnote{โส อิมํ (สี, สฺยา, กํ, อิ)} โลกํ ปภาเสติ, อพฺภา มุตฺโตว จนฺทิมา.\csromanspacer{} ภิกฺขุวคฺโค ปญฺจวีสติโม.}}

\csromanheader{1}{26. พฺราหฺมณวคฺค}
\vspace{\tipitakaparskip}
\csromancenter{ปสาทพหุลพฺราหฺมณวตฺถุ 383.\csromanspacer{} ฉินฺท โสตํ ปรกฺกมฺม, กาเม ปนุท พฺราหฺมณ.}
\vspace{0.5\baselineskip}
\csromangathameasure{สงฺขารานํ ขยํ ญตฺวา, อกตญฺญูสิ พฺราหฺมณ.}
\csromangathagroup{\csromangathastanza{สงฺขารานํ ขยํ ญตฺวา, อกตญฺญูสิ พฺราหฺมณ.}}

\vspace{\tipitakaparskip}
\csromancenter{สมฺพหุลภิกฺขุวตฺถุ 384.\csromanspacer{} ยทา ทฺวเยสุ ธมฺเมสุ, ปารคู โหติ พฺราหฺมโณ.}
\vspace{0.5\baselineskip}
\csromangathameasure{อถสฺส สพฺเพ สํโยคา, อตฺถํ คจฺฉนฺติ ชานโต.}
\csromangathagroup{\csromangathastanza{อถสฺส สพฺเพ สํโยคา, อตฺถํ คจฺฉนฺติ ชานโต.}}

\vspace{\tipitakaparskip}
\csromancenter{มารวตฺถุ 385.\csromanspacer{} ยสฺส ปารํ อปารํ วา, ปาราปารํ น วิชฺชติ.}
\vspace{0.5\baselineskip}
\csromangathameasure{วีตทฺทรํ วิสํยุตฺตํ, ตมหํ พฺรูมิ พฺราหฺมณํ.}
\csromangathagroup{\csromangathastanza{วีตทฺทรํ วิสํยุตฺตํ, ตมหํ พฺรูมิ พฺราหฺมณํ.}}

\vspace{\tipitakaparskip}
\csromancenter{อญฺญตร พฺราหฺมณวตฺถุ 386.\csromanspacer{} ฌายิํ วิรช’มาสีนํ, กตกิจฺจ’มนาสวํ.}
\vspace{0.5\baselineskip}
\csromangathameasure{อุตฺตมตฺถ’มนุปฺปตฺตํ, ตมหํ พฺรูมิ พฺราหฺมณํ.}
\csromangathagroup{\csromangathastanza{อุตฺตมตฺถ’มนุปฺปตฺตํ, ตมหํ พฺรูมิ พฺราหฺมณํ.}}

\vspace{\tipitakaparskip}
\csromancenter{พฺราหฺมณวคฺค \textbf{5.}\csromanspacer{}\textbf{ อานนฺทตฺเถรวตฺถุ} 387.\csromanspacer{} ทิวา ตปติ อาทิจฺโจ, รตฺติ’มาภาติ จนฺทิมา.}
\vspace{0.5\baselineskip}
\csromangathameasure{สนฺนทฺโธ ขตฺติโย ตปติ, ฌายี ตปติ พฺราหฺมโณ. \\ อถ สพฺพมโหรตฺติํ, พุทฺโธ ตปติ เตชสา.}
\csromangathagroup{\csromangathastanza{\csromanfolio{69}สนฺนทฺโธ ขตฺติโย ตปติ, ฌายี ตปติ พฺราหฺมโณ. \\ อถ สพฺพมโหรตฺติํ\footnote{สพฺพมโหรตฺตํ (?)}, พุทฺโธ ตปติ เตชสา.}}

\vspace{\tipitakaparskip}
\csromancenter{อญฺญตรพฺราหฺมณปพฺพชิตวตฺถุ 38\textbf{8.}\csromanspacer{} พาหิตปาโปติ พฺราหฺมโณ,}
\prose{สมจริยา “สมโณ”ติ วุจฺจติ.}

\csromangathameasure{ปพฺพาชยมตฺตโน มลํ, \\ ตสฺมา “ปพฺพชิโต”ติ วุจฺจติ.}
\csromangathagroup{\csromangathastanza{ปพฺพาชยมตฺตโน มลํ, \\ ตสฺมา “ปพฺพชิโต”ติ วุจฺจติ.}}

\csromanmatikaanchor{mtk.35}
\csromantocmarkref{section}{24. ตณฺหาวคฺค}{mtk.35}
\bookmark[dest={mtk.35},level=1]{24. ตณฺหาวคฺค}
\vspace{\tipitakaparskip}
\csromancenter{สาริปุตฺตตฺเถรวตฺถุ 389.\csromanspacer{} น พฺราหฺมณสฺส ปหเรยฺย, นาสฺส มุญฺเจถ พฺราหฺมโณ.}
\vspace{0.5\baselineskip}
\csromangathameasure{ธี พฺราหฺมณสฺส หนฺตารํ, ตโต ธี ยสฺส มุญฺจติ. \\ น พฺราหฺมณสฺเส’ตทกิญฺจิ เสยฺโย, \\ ยทา นิเสโธ มนโส ปิเยหิ. \\ ยโต ยโต หิํสมโน นิวตฺตติ, \\ ตโต ตโต สมฺมติเมว ทุกฺขํ.}
\csromangathagroup{\csromangathastanza{ธี\footnote{ธิ (สฺยา, พฺยากรเณสุ)} พฺราหฺมณสฺส หนฺตารํ, ตโต ธี ยสฺส\footnote{โย + อสฺส = ยสฺส} มุญฺจติ.}\csromangathastanza{น พฺราหฺมณสฺเส’ตทกิญฺจิ เสยฺโย, \\ ยทา นิเสโธ มนโส ปิเยหิ. \\ ยโต ยโต หิํสมโน นิวตฺตติ, \\ ตโต ตโต สมฺมติเมว ทุกฺขํ.}}

\vspace{\tipitakaparskip}
\csromancenter{มหาปชาปติโคตมีวตฺถุ 391.\csromanspacer{} ยสฺส กาเยน วาจาย, มนสา นตฺถิ ทุกฺกฏํ.}
\vspace{0.5\baselineskip}
\csromangathameasure{สํวุตํ ตีหิ ฐาเนหิ, ตมหํ พฺรูมิ พฺราหฺมณํ.}
\csromangathagroup{\csromangathastanza{สํวุตํ ตีหิ ฐาเนหิ, ตมหํ พฺรูมิ พฺราหฺมณํ.}}

\vspace{\tipitakaparskip}
\csromancenter{สาริปุตฺตตฺเถรวตฺถุ 392.\csromanspacer{} ยมฺหา ธมฺมํ วิชาเนยฺย, สมฺมาสมฺพุทฺธเทสิตํ.}
\vspace{0.5\baselineskip}
\csromangathameasure{สกฺกจฺจํ ตํ นมสฺเสยฺย, อคฺคิหุตฺตํว พฺราหฺมโณ.}
\csromangathagroup{\csromangathastanza{สกฺกจฺจํ ตํ นมสฺเสยฺย, อคฺคิหุตฺตํว พฺราหฺมโณ.}}

\vspace{\tipitakaparskip}
\csromancenter{ชฏิลพฺราหฺมณวตฺถุ 393.\csromanspacer{} น ชฏาหิ น โคตฺเตน, น ชจฺจา โหติ พฺราหฺมโณ.}
\vspace{0.5\baselineskip}
\csromangathameasure{ยมฺหิ สจฺจญฺจ ธมฺโม จ, โส สุจี โส จ พฺราหฺมโณ.}
\csromangathagroup{\csromangathastanza{ยมฺหิ สจฺจญฺจ ธมฺโม จ, โส สุจี โส จ พฺราหฺมโณ.}}

\vspace{\tipitakaparskip}
\csromancenter{ธมฺมปทปาฬิ \textbf{11.}\csromanspacer{}\textbf{ กุหกพฺราหฺมณวตฺถุ} 394.\csromanspacer{} กิํ เต ชฏาหิ ทุมฺเมธ, กิํ เต อชินสาฏิยา.}
\vspace{0.5\baselineskip}
\csromangathameasure{อพฺภนฺตรํ เต คหนํ, พาหิรํ ปริมชฺชสิ.}
\csromangathagroup{\csromangathastanza{\csromanfolio{70}อพฺภนฺตรํ เต คหนํ, พาหิรํ ปริมชฺชสิ.}}

\vspace{\tipitakaparskip}
\csromancenter{กิสาโคตมีวตฺถุ 395.\csromanspacer{} ปํสุกูลธรํ ชนฺตุํ, กิสํ ธมนิสนฺถ ตํ.}
\vspace{0.5\baselineskip}
\csromangathameasure{เอกํ วนสฺมิํ ฌายนฺตํ, ตมหํ พฺรูมิ พฺราหฺมณํ.}
\csromangathagroup{\csromangathastanza{เอกํ วนสฺมิํ ฌายนฺตํ, ตมหํ พฺรูมิ พฺราหฺมณํ.}}

\vspace{\tipitakaparskip}
\csromancenter{เอกพฺราหฺมณวตฺถุ 396. * น จาหํ พฺราหฺมณํ พฺรูมิ, โยนิชํ มตฺติสมฺภวํ.}
\vspace{0.5\baselineskip}
\csromangathameasure{โภวาทิ นาม โส โหติ, สเจ โหติ สกิญฺจโน. \\ อกิญฺจนํ อนาทานํ, ตมหํ พฺรูมิ พฺราหฺมณํ.}
\csromangathagroup{\csromangathastanza{โภวาทิ นาม โส โหติ, สเจ โหติ สกิญฺจโน. \\ อกิญฺจนํ อนาทานํ, ตมหํ พฺรูมิ พฺราหฺมณํ.}}

\vspace{\tipitakaparskip}
\csromancenter{อุคฺคเสนเสฏฺฐิปุตฺตวตฺถุ 397. * สพฺพสํโยชนํ เฉตฺวา, โย เว น ปริตสฺสติ.}
\vspace{0.5\baselineskip}
\csromangathameasure{สงฺคาติคํ วิสํยุตฺตํ, ตมหํ พฺรูมิ พฺราหฺมณํ.}
\csromangathagroup{\csromangathastanza{สงฺคาติคํ วิสํยุตฺตํ, ตมหํ พฺรูมิ พฺราหฺมณํ.}}

\vspace{\tipitakaparskip}
\csromancenter{เทฺวพฺราหฺมณวตฺถุ 398. * เฉตฺวา นทฺธิํ\footnote{นนฺธิํ (ก-สี), นนฺทิํ (อิ)} วรตฺตญฺจ, สนฺทานํ\footnote{สนฺทามํ (สี)} สหนุกฺกมํ.}
\vspace{0.5\baselineskip}
\csromangathameasure{อุกฺขิตฺตปลิฆํ พุทฺธํ, ตมหํ พฺรูมิ พฺราหฺมณํ.}
\csromangathagroup{\csromangathastanza{อุกฺขิตฺตปลิฆํ พุทฺธํ, ตมหํ พฺรูมิ พฺราหฺมณํ.}}

\vspace{\tipitakaparskip}
\csromancenter{อกฺโกสกภารทฺวาชวตฺถุ 399. * อกฺโกสํ วธพนฺธญฺจ, อทุฏฺโฐ โย ติติกฺขติ.}
\vspace{0.5\baselineskip}
\csromangathameasure{ขนฺตีพลํ พลานีกํ, ตมหํ พฺรูมิ พฺราหฺมณํ.}
\csromangathagroup{\csromangathastanza{ขนฺตีพลํ พลานีกํ, ตมหํ พฺรูมิ พฺราหฺมณํ.}}

\vspace{\tipitakaparskip}
\csromancenter{สาริปุตฺตตฺเถรวตฺถุ 400. * อกฺโกธนํ วตวนฺตํ, สีลวนฺตํ อนุสฺสทํ.}
\vspace{0.5\baselineskip}
\csromangathameasure{ทนฺตํ อนฺติมสารีรํ, ตมหํ พฺรูมิ พฺราหฺมณํ.}
\csromangathagroup{\csromangathastanza{ทนฺตํ อนฺติมสารีรํ, ตมหํ พฺรูมิ พฺราหฺมณํ.}}

\vspace{\tipitakaparskip}
\csromancenter{พฺราหฺมณวคฺค \textbf{18.}\csromanspacer{}\textbf{ อุปฺปลวณฺณาตฺเถรีวตฺถุ} 401. * วาริ โปกฺขรปตฺเตว, อารคฺเคริว สาสโป.}
\vspace{0.5\baselineskip}
\csromangathameasure{โย น ลิมฺปติ กาเมสุ, ตมหํ พฺรูมิ พฺราหฺมณํ.}
\csromangathagroup{\csromangathastanza{\csromanfolio{71}โย น ลิมฺปติ\footnote{ลิปฺปติ (สี, อิ)} กาเมสุ, ตมหํ พฺรูมิ พฺราหฺมณํ.}}

\vspace{\tipitakaparskip}
\csromancenter{อญฺญตร พฺราหฺมณวตฺถุ 402. * โย ทุกฺขสฺส ปชานาติ, อิเธว ขยมตฺตโน.}
\vspace{0.5\baselineskip}
\csromangathameasure{ปนฺนภารํ วิสํยุตฺตํ, ตมหํ พฺรูมิ พฺราหฺมณํ.}
\csromangathagroup{\csromangathastanza{ปนฺนภารํ วิสํยุตฺตํ, ตมหํ พฺรูมิ พฺราหฺมณํ.}}

\vspace{\tipitakaparskip}
\csromancenter{เขมาภิกฺขุนีวตฺถุ 403. * คมฺภีรปญฺญํ เมธาวิํ, มคฺคามคฺคสฺส โกวิทํ.}
\vspace{0.5\baselineskip}
\csromangathameasure{อุตฺตมตฺถ’มนุปฺปตฺตํ, ตมหํ พฺรูมิ พฺราหฺมณํ.}
\csromangathagroup{\csromangathastanza{อุตฺตมตฺถ’มนุปฺปตฺตํ, ตมหํ พฺรูมิ พฺราหฺมณํ.}}

\vspace{\tipitakaparskip}
\csromancenter{ปพฺภารวาสิติสฺสตฺเถรวตฺถุ 404. * อสํสฏฺฐํ คหฏฺเฐหิ, อนาคาเรหิ จูภยํ.}
\vspace{0.5\baselineskip}
\csromangathameasure{อโนกสาริ’มปฺปิจฺฉํ, ตมหํ พฺรูมิ พฺราหฺมณํ.}
\csromangathagroup{\csromangathastanza{อโนกสาริ’มปฺปิจฺฉํ, ตมหํ พฺรูมิ พฺราหฺมณํ.}}

\vspace{\tipitakaparskip}
\csromancenter{อญฺญตรภิกฺขุวตฺถุ 405. * นิธาย ทณฺฑํ ภูเตสุ, ตเสสุ ถาวเรสุ จ.}
\vspace{0.5\baselineskip}
\csromangathameasure{โย น หนฺติ น ฆาเตติ, ตมหํ พฺรูมิ พฺราหฺมณํ.}
\csromangathagroup{\csromangathastanza{โย น หนฺติ น ฆาเตติ, ตมหํ พฺรูมิ พฺราหฺมณํ.}}

\vspace{\tipitakaparskip}
\csromancenter{สามเณรานํ วตฺถุ 406. * อวิรุทฺธํ วิรุทฺเธสุ, อตฺตทณฺเฑสุ นิพฺพุตํ.}
\vspace{0.5\baselineskip}
\csromangathameasure{สาทาเนสุ อนาทานํ, ตมหํ พฺรูมิ พฺราหฺมณํ.}
\csromangathagroup{\csromangathastanza{สาทาเนสุ อนาทานํ, ตมหํ พฺรูมิ พฺราหฺมณํ.}}

\vspace{\tipitakaparskip}
\csromancenter{มหาปนฺถกตฺเถรวตฺถุ 407. * ยสฺส ราโค จ โทโส จ, มาโน มกฺโข จ ปาติโต.}
\vspace{0.5\baselineskip}
\csromangathameasure{สาสโปริว อารคฺคา, ตมหํ พฺรูมิ พฺราหฺมณํ.}
\csromangathagroup{\csromangathastanza{สาสโปริว อารคฺคา\footnote{อารคฺเค (ก)}, ตมหํ พฺรูมิ พฺราหฺมณํ.}}

\vspace{\tipitakaparskip}
\csromancenter{ธมฺมปทปาฬิ \textbf{25.}\csromanspacer{}\textbf{ ปิลินฺทวจฺฉตฺเถรวตฺถุ} 408. * อกกฺกสํ วิญฺญาปนิํ, คิรํ สจฺจ’มุทีรเย.}
\vspace{0.5\baselineskip}
\csromangathameasure{ยาย นาภิสเช กญฺจิ, ตมหํ พฺรูมิ พฺราหฺมณํ.}
\csromangathagroup{\csromangathastanza{\csromanfolio{72}ยาย นาภิสเช กญฺจิ\footnote{กิญฺจิ (ก)}, ตมหํ พฺรูมิ พฺราหฺมณํ.}}

\vspace{\tipitakaparskip}
\csromancenter{อญฺญตรตฺเถรวตฺถุ 409. * โย’ธ ทีฆํ ว รสฺสํ วา, อณุํ ถูลํ สุภาสุภํ.}
\vspace{0.5\baselineskip}
\csromangathameasure{โลเก อทินฺนํ นาทิยติ, ตมหํ พฺรูมิ พฺราหฺมณํ.}
\csromangathagroup{\csromangathastanza{โลเก อทินฺนํ นาทิยติ\footnote{นาเทติ (ม 2. 411)}, ตมหํ พฺรูมิ พฺราหฺมณํ.}}

\vspace{\tipitakaparskip}
\csromancenter{สาริปุตฺตตฺเถรวตฺถุ 410. * อาสา ยสฺส น วิชฺชนฺติ, อสฺมิํ โลเก ปรมฺหิ จ.}
\vspace{0.5\baselineskip}
\csromangathameasure{นิราสาสํ วิสํยุตฺตํ, ตมหํ พฺรูมิ พฺราหฺมณํ.}
\csromangathagroup{\csromangathastanza{นิราสาสํ\footnote{นิราสยํ (สี, สฺยา, อิ), นิราสกํ (?)} วิสํยุตฺตํ, ตมหํ พฺรูมิ พฺราหฺมณํ.}}

\vspace{\tipitakaparskip}
\csromancenter{มหาโมคฺคลฺลานตฺเถรวตฺถุ 411. * ยสฺสาลยา น วิชฺชนฺติ, อญฺญาย อกถํกถี.}
\vspace{0.5\baselineskip}
\csromangathameasure{อมโตคธ’มนุปฺปตฺตํ, ตมหํ พฺรูมิ พฺราหฺมณํ.}
\csromangathagroup{\csromangathastanza{อมโตคธ’มนุปฺปตฺตํ, ตมหํ พฺรูมิ พฺราหฺมณํ.}}

\vspace{\tipitakaparskip}
\csromancenter{เรวตตฺเถรวตฺถุ 412. * โย’ธ ปุญฺญญฺจ ปาปญฺจ, อุโภ สงฺค’มุปจฺจคา.}
\vspace{0.5\baselineskip}
\csromangathameasure{อโสกํ วิรชํ สุทฺธํ, ตมหํ พฺรูมิ พฺราหฺมณํ.}
\csromangathagroup{\csromangathastanza{อโสกํ วิรชํ สุทฺธํ, ตมหํ พฺรูมิ พฺราหฺมณํ.}}

\vspace{\tipitakaparskip}
\csromancenter{จนฺทาภตฺเถรวตฺถุ 413. * จนฺทํว วิมลํ สุทฺธํ, วิปฺปสนฺนมนาวิลํ.}
\vspace{0.5\baselineskip}
\csromangathameasure{นนฺทีภวปริกฺขีณํ, ตมหํ พฺรูมิ พฺราหฺมณํ.}
\csromangathagroup{\csromangathastanza{นนฺทีภวปริกฺขีณํ, ตมหํ พฺรูมิ พฺราหฺมณํ.}}

\vspace{\tipitakaparskip}
\csromancenter{สีวลิตฺเถรวตฺถุ 414. * โย’มํ\footnote{โย อิมํ (สี, สฺยา, กํ, อิ)} ปลิปถํ ทุคฺคํ, สํสารํ โมหมจฺจคา.}
\vspace{0.5\baselineskip}
\csromangathameasure{ติณฺโณ ปารงฺคโต ฌายี, อเนโช อกถํกถี. \\ อนุปาทาย นิพฺพุโต, ตมหํ พฺรูมิ พฺราหฺมณํ.}
\csromangathagroup{\csromangathastanza{ติณฺโณ ปารงฺคโต\footnote{ปารคโต (สี, สฺยา, กํ, อิ)} ฌายี, อเนโช อกถํกถี. \\ อนุปาทาย นิพฺพุโต, ตมหํ พฺรูมิ พฺราหฺมณํ.}}

\vspace{\tipitakaparskip}
\csromancenter{พฺราหฺมณวคฺค \textbf{32.}\csromanspacer{}\textbf{ สุนฺทรสมุทฺทตฺเถรวตฺถุ} 415. * โย’ธ กาเม ปหนฺตฺวาน\footnote{ปหตฺวาน (สี, อิ)}, อนาคาโร ปริพฺพเช.}
\vspace{0.5\baselineskip}
\csromangathameasure{กามภวปริกฺขีณํ, ตมหํ พฺรูมิ พฺราหฺมณํ.}
\csromangathagroup{\csromangathastanza{\csromanfolio{73}กามภวปริกฺขีณํ, ตมหํ พฺรูมิ พฺราหฺมณํ.}}

\vspace{\tipitakaparskip}
\csromancenter{ชฏิลตฺเถรวตฺถุ 416. * โย’ธ ตณฺหํ ปหนฺตฺวาน, อนาคาโร ปริพฺพเช.}
\vspace{0.5\baselineskip}
\csromangathameasure{ตณฺหาภวปริกฺขีณํ, ตมหํ พฺรูมิ พฺราหฺมณํ.}
\csromangathagroup{\csromangathastanza{ตณฺหาภวปริกฺขีณํ, ตมหํ พฺรูมิ พฺราหฺมณํ\footnote{อิทํ คาถาทฺวยํ วิเทสโปตฺถเกสุ สกิเทว ทสฺสิตํ.}.}}

\vspace{\tipitakaparskip}
\csromancenter{โชติกตฺเถรวตฺถุ 416. * โย’ธ ตณฺหํ ปหนฺตฺวาน, อนาคาโร ปริพฺพเช.}
\vspace{0.5\baselineskip}
\csromangathameasure{ตณฺหาภวปริกฺขีณํ, ตมหํ พฺรูมิ พฺราหฺมณํ2.}
\csromangathagroup{\csromangathastanza{ตณฺหาภวปริกฺขีณํ, ตมหํ พฺรูมิ พฺราหฺมณํ2.}}

\vspace{\tipitakaparskip}
\csromancenter{นฏปุตฺตกตฺเถรวตฺถุ 417. * หิตฺวา มานุสกํ โยคํ, ทิพฺพํ โยคํ อุปจฺจคา.}
\vspace{0.5\baselineskip}
\csromangathameasure{สพฺพโยควิสํยุตฺตํ, ตมหํ พฺรูมิ พฺราหฺมณํ.}
\csromangathagroup{\csromangathastanza{สพฺพโยควิสํยุตฺตํ, ตมหํ พฺรูมิ พฺราหฺมณํ.}}

\vspace{\tipitakaparskip}
\csromancenter{นฏปุตฺตกตฺเถรวตฺถุ 418. * หิตฺวา รติํ จ อรติํ จ, สีติภูตํ นิรูปธิํ.}
\vspace{0.5\baselineskip}
\csromangathameasure{สพฺพโลกาภิภุํ วีรํ, ตมหํ พฺรูมิ พฺราหฺมณํ.}
\csromangathagroup{\csromangathastanza{สพฺพโลกาภิภุํ วีรํ, ตมหํ พฺรูมิ พฺราหฺมณํ.}}

\vspace{\tipitakaparskip}
\csromancenter{วงฺคีสตฺเถรวตฺถุ 419. * จุติํ โย เวทิ สตฺตานํ, อุปปตฺติญฺจ สพฺพโส.}
\vspace{0.5\baselineskip}
\csromangathameasure{อสตฺตํ สุคตํ พุทฺธํ, ตมหํ พฺรูมิ พฺราหฺมณํ.}
\csromangathagroup{\csromangathastanza{อสตฺตํ สุคตํ พุทฺธํ, ตมหํ พฺรูมิ พฺราหฺมณํ.}}

\proseitem{420}{\csromansymbolfootnote{*}{อุปริ 377; ม 2. 411 ปิฏฺเฐสุปิ.}ยสฺส คติํ น ชานนฺติ, เทวา คนฺธพฺพมานุสา.}

\csromangathameasure{ขีณาสวํ อรหนฺตํ, ตมหํ พฺรูมิ พฺราหฺมณํ.}
\csromangathagroup{\csromangathastanza{ขีณาสวํ อรหนฺตํ, ตมหํ พฺรูมิ พฺราหฺมณํ.}}

\vspace{\tipitakaparskip}
\csromancenter{ธมฺมปทปาฬิ \textbf{38.}\csromanspacer{}\textbf{ ธมฺมทินฺนาตฺเถรีวตฺถุ} 421.\csromanspacer{} ยสฺส ปุเร จ ปจฺฉา จ, มชฺเฌ จ นตฺถิ กิญฺจนํ.}
\vspace{0.5\baselineskip}
\csromangathameasure{อกิญฺจนํ อนาทานํ, ตมหํ พฺรูมิ พฺราหฺมณํ.}
\csromangathagroup{\csromangathastanza{\csromanfolio{74}อกิญฺจนํ อนาทานํ, ตมหํ พฺรูมิ พฺราหฺมณํ.}}

\vspace{\tipitakaparskip}
\csromancenter{องฺคุลิมาลตฺเถรวตฺถุ 422.\csromanspacer{} อุสภํ ปวรํ วีรํ, มเหสิํ วิชิตาวินํ.}
\vspace{0.5\baselineskip}
\csromangathameasure{อเนชํ นฺหาตกํ พุทฺธํ, ตมหํ พฺรูมิ พฺราหฺมณํ.}
\csromangathagroup{\csromangathastanza{อเนชํ นฺหาตกํ\footnote{นหาตกํ (สี, สฺยา, กํ, อิ)} พุทฺธํ, ตมหํ พฺรูมิ พฺราหฺมณํ.}}

\vspace{\tipitakaparskip}
\csromancenter{เทวหิตพฺราหฺมณวตฺถุ 423.\csromanspacer{} ปุพฺเพนิวาสํ โย เวทิ, สคฺคาปายญฺจ ปสฺสติ.}
\vspace{0.5\baselineskip}
\csromangathameasure{อโถ ชาติกฺขยํ ปตฺโต, อภิญฺญาโวสิโต มุนิ. \\ สพฺพโวสิตโวสานํ, ตมหํ พฺรูมิ พฺราหฺมณํ. พฺราหฺมณวคฺโค ฉพฺพีสติโม.}
\csromangathagroup{\csromangathastanza{อโถ ชาติกฺขยํ ปตฺโต, อภิญฺญาโวสิโต มุนิ. \\ สพฺพโวสิตโวสานํ, ตมหํ พฺรูมิ พฺราหฺมณํ.\csromanspacer{} พฺราหฺมณวคฺโค ฉพฺพีสติโม.}}

\prose{(เอตฺตาวตา สพฺพปฐเม ยมกวคฺเค จุทฺทส วตฺถูนิ, อปฺปมาทวคฺเค นว, จิตฺตวคฺเค นว, ปุปฺผวคฺเค ทฺวาทส, พาลวคฺเค ปนฺนรส, ปณฺฑิตวคฺเค เอกาทส, อรหนฺตวคฺเค ทส, สหสฺสวคฺเค จุทฺทส, ปาปวคฺเค ทฺวาทส, ทณฺฑวคฺเค เอกาทส, ชราวคฺเค นว, อตฺตวคฺเค ทส, โลกวคฺเค เอกาทส, พุทฺธวคฺเค นว\footnote{อฏฺฐ (ก)}, สุขวคฺเค อฏฺฐ, ปิยวคฺเค นว, โกธวคฺเค อฏฺฐ, มลวคฺเค ทฺวาทส, ธมฺมฏฺฐวคฺเค ทส, มคฺควคฺเค ทฺวาทส, ปกิณฺณกวคฺเค นว, นิรยวคฺเค นว, นาควคฺเค อฏฺฐ, ตณฺหาวคฺเค ทฺวาทส, ภิกฺขุวคฺเค ทฺวาทส, พฺราหฺมณวคฺเค จตฺตาลีสาติ ปญฺจาธิกานิ ตีณิ วตฺถุสตานิ.}

\csromanmatikaanchor{mtk.36}
\csromantocmarkref{section}{25. ภิกฺขุวคฺค}{mtk.36}
\bookmark[dest={mtk.36},level=1]{25. ภิกฺขุวคฺค}
\csromangathameasure{สเตวีสจตุสฺสตา, จตุสจฺจวิภาวินา. \\ สตตฺตยญฺจ วตฺถูนํ, ปญฺจาธิกํ สมุฏฺฐิตาติ).}
\csromangathagroup{\csromangathastanza{สเตวีสจตุสฺสตา, จตุสจฺจวิภาวินา. \\ สตตฺตยญฺจ วตฺถูนํ, ปญฺจาธิกํ สมุฏฺฐิตาติ)\footnote{( ) เอตฺถนฺตเร ปาโฐ วิเทสโปตฺถเกสุ นตฺถิ, อฏฺฐกถาสุเยว ทิสฺสติ.}.}}

\csromanheader{1}{อุทฺทานคาถาโย}
\csromangathameasure{1 ธมฺมปเท วคฺคานมุทฺทานํ ยมกปฺปมาโท จิตฺตํ, \\ ปุปฺผํ พาเลน ปณฺฑิโต. อรหนฺโต สหสฺสญฺจ, \\ ปาปํ ทณฺเฑน เต ทส. ชรา อตฺตา จ โลโก จ, \\ พุทฺโธ สุขํ ปิเยน จ. โกโธ มลญฺจ ธมฺมฏฺโฐ, \\ มคฺควคฺเคน วีสติ. ปกิณฺณํ นิรโย นาโค, \\ ตณฺหา ภิกฺขุ จ พฺราหฺมโณ. เอเต ฉพฺพีสติ วคฺคา, \\ เทสิตาทิจฺจพนฺธุนา.}
\csromangathagroup{\csromangathastanza{\csromanfolio{75}1 ธมฺมปเท วคฺคานมุทฺทานํ ยมกปฺปมาโท จิตฺตํ, \\ ปุปฺผํ พาเลน ปณฺฑิโต. อรหนฺโต สหสฺสญฺจ,}\csromangathastanza{ปาปํ ทณฺเฑน เต ทส. ชรา อตฺตา จ โลโก จ, \\ พุทฺโธ สุขํ ปิเยน จ. โกโธ มลญฺจ ธมฺมฏฺโฐ,}\csromangathastanza{มคฺควคฺเคน วีสติ. ปกิณฺณํ นิรโย นาโค, \\ ตณฺหา ภิกฺขุ จ พฺราหฺมโณ. เอเต ฉพฺพีสติ วคฺคา, \\ เทสิตาทิจฺจพนฺธุนา.}}

\csromanheader{2}{\textbf{คาถานมุทฺทานํ}}
\csromanheader{2}{ยมเก วีสติ คาถา, อปฺปมาทมฺหิ ทฺวาทส.}
\csromangathameasure{เอกาทส จิตฺตวคฺเค, ปุปฺผวคฺคมฺหิ โสฬส.}
\csromangathagroup{\csromangathastanza{เอกาทส จิตฺตวคฺเค, ปุปฺผวคฺคมฺหิ โสฬส.}}

\csromanheader{2}{พาเล จ โสฬส คาถา, ปณฺฑิตมฺหิ จตุทฺทส.}
\csromanheader{2}{อรหนฺเต ทส คาถา, สหสฺเส โหนฺติ โสฬส.}
\csromangathameasure{เตรส ปาปวคฺคมฺหิ, ทณฺฑมฺหิ ทส สตฺต จ. \\ เอกาทส ชรา วคฺเค, อตฺตวคฺคมฺหิ ตา ทส.}
\csromangathagroup{\csromangathastanza{เตรส ปาปวคฺคมฺหิ, ทณฺฑมฺหิ ทส สตฺต จ. \\ เอกาทส ชรา วคฺเค, อตฺตวคฺคมฺหิ ตา ทส.}}

\gambhira{ธมฺมปทปาฬิ}
\csromangathameasure{ทฺวาทส โลกวคฺคมฺหิ, พุทฺธวคฺคมฺหิ ฐารส. \\ สุเข จ ปิยวคฺเค จ, คาถาโย โหนฺติ ทฺวาทส. \\ จุทฺทส โกธวคฺคมฺหิ, มลวคฺเคกวีสติ. \\ สตฺตรส จ ธมฺมฏฺเฐ, มคฺควคฺเค สตฺตรส.}
\csromangathagroup{\csromangathastanza{\csromanfolio{76}ทฺวาทส โลกวคฺคมฺหิ, พุทฺธวคฺคมฺหิ ฐารส\footnote{โสฬส (สพฺพตฺถ)}. \\ สุเข จ ปิยวคฺเค จ, คาถาโย โหนฺติ ทฺวาทส.}\csromangathastanza{จุทฺทส โกธวคฺคมฺหิ, มลวคฺเคกวีสติ. \\ สตฺตรส จ ธมฺมฏฺเฐ, มคฺควคฺเค สตฺตรส.}}

\csromanheader{2}{ปกิณฺเณ โสฬส คาถา, นิรเย นาเค จ จุทฺทส.}
\csromangathameasure{ฉพฺพีส ตณฺหาวคฺคมฺหิ, เตวีส ภิกฺขุวคฺคิกา. \\ เอกตาลีสคาถาโย, พฺราหฺมเณ วคฺคมุตฺตเม. \\ คาถาสตานิ จตฺตาริ, เตวีส จ ปุนาปเร. \\ ธมฺมปเท นิปาตมฺหิ, เทสิตาทิจฺจพนฺธุนาติ1.}
\csromangathagroup{\csromangathastanza{ฉพฺพีส ตณฺหาวคฺคมฺหิ, เตวีส ภิกฺขุวคฺคิกา. \\ เอกตาลีสคาถาโย, พฺราหฺมเณ วคฺคมุตฺตเม.}\csromangathastanza{คาถาสตานิ จตฺตาริ, เตวีส จ ปุนาปเร. \\ ธมฺมปเท นิปาตมฺหิ, เทสิตาทิจฺจพนฺธุนาติ1.}}

\gambhira{\textbf{ธมฺมปทปาฬิ สมตฺตา.}}
\csromanheader{2}{\textbf{ขุทฺทกนิกาย}}
\gambhira{\textbf{อุทานปาฬิ}}
\namakkaram{นโม ตสฺส ภควโต อรหโต สมฺมาสมฺพุทฺธสฺส.}
\chapterhead{1. โพธิวคฺค}
\csromanheader{2}{1. ปฐมโพธิสุตฺต}
\proseitem{1}{\csromanfolio{77}\csromansymbolfootnote{*}{วิ 3. 1 ปิฏฺเฐปิ.}เอวํ เม สุตํ\csromandash{}เอกํ สมยํ ภควา อุรุเวลายํ วิหรติ นชฺชา เนรญฺชราย ตีเร โพธิรุกฺขมูเล ปฐมาภิสมฺพุทฺโธ.\csromanspacer{} เตน โข ปน สมเยน ภควา สตฺตาหํ เอกปลฺลงฺเกน นิสินฺโน โหติ วิมุตฺติสุขปฏิสํเวที\footnote{วิมุตฺติสุขํ ปฏิสํเวที (สฺยา, อิ, ก)}.\csromanspacer{} อถ โข ภควา ตสฺส สตฺตาหสฺส อจฺจเยน ตมฺหา สมาธิมฺหา วุฏฺฐหิตฺวา รตฺติยา ปฐมํ ยามํ ปฏิจฺจสมุปฺปาทํ อนุโลมํ สาธุกํ มนสากาสิ\csromandash{}}

\prose{อิติ อิมสฺมิํ สติ อิทํ โหติ, อิมสฺสุปฺปาทา อิทํ อุปฺปชฺชติ, ยทิทํ อวิชฺชาปจฺจยา สงฺขารา, สงฺขารปจฺจยา วิญฺญาณํ, วิญฺญาณปจฺจยา นามรูปํ, นามรูปปจฺจยา สฬายตนํ, สฬายตนปจฺจยา ผสฺโส, ผสฺสปจฺจยา เวทนา, เวทนาปจฺจยา ตณฺหา, ตณฺหาปจฺจยา อุปาทานํ, อุปาทานปจฺจยา ภโว, ภวปจฺจยา ชาติ, ชาติปจฺจยา ชรามรณํ โสกปริเทวทุกฺขโทมนสฺสุปายาสา สมฺภวนฺติ, เอวเมตสฺส เกวลสฺส ทุกฺขกฺขนฺธสฺส สมุทโย โหตีติ.}

\prose{อถ โข ภควา เอตมตฺถํ วิทิตฺวา ตายํ เวลายํ อิมํ อุทานํ อุทาเนสิ\csromandash{}}

\gambhira{อุทานปาฬิ}
\csromangathameasure{“ยทา หเว ปาตุภวนฺติ ธมฺมา, \\ อาตาปิโน ฌายโต พฺราหฺมณสฺส. \\ อถสฺส กงฺขา วปยนฺติ สพฺพา, \\ ยโต ปชานาติ สเหตุธมฺมนฺ”ติ.ปฐมํ.}
\csromangathagroup{\csromangathastanza{\csromanfolio{78}“ยทา หเว ปาตุภวนฺติ ธมฺมา, \\ อาตาปิโน ฌายโต พฺราหฺมณสฺส. \\ อถสฺส กงฺขา วปยนฺติ สพฺพา, \\ ยโต ปชานาติ สเหตุธมฺมนฺ”ติ.\csromanspacer{}\csromanspacer{}ปฐมํ.}}

\csromanheader{2}{2. ทุติยโพธิสุตฺต}
\proseitem{2}{\csromansymbolfootnote{*}{วิ 3. 2 ปิฏฺเฐปิ.}เอวํ เม สุตํ\csromandash{}เอกํ สมยํ ภควา อุรุเวลายํ วิหรติ นชฺชา เนรญฺชราย ตีเร โพธิรุกฺขมูเล ปฐมาภิสมฺพุทฺโธ.\csromanspacer{} เตน โข ปน สมเยน ภควา สตฺตาหํ เอกปลฺลงฺเกน นิสินฺโน โหติ วิมุตฺติสุขปฏิสํเวที.\csromanspacer{} อถ โข ภควา ตสฺส สตฺตาหสฺส อจฺจเยน ตมฺหา สมาธิมฺหา วุฏฺฐหิตฺวา รตฺติยา มชฺฌิมํ ยามํ ปฏิจฺจสมุปฺปาทํ ปฏิโลมํ สาธุกํ มนสากาสิ\csromandash{}}

\prose{อิติ อิมสฺมิํ อสติ อิทํ น โหติ, อิมสฺส นิโรธา อิทํ นิรุชฺฌติ, ยทิทํ อวิชฺชานิโรธา สงฺขารนิโรโธ, สงฺขารนิโรธา วิญฺญาณนิโรโธ, วิญฺญาณนิโรธา นามรูปนิโรโธ, นามรูปนิโรธา สฬายตนนิโรโธ, สฬายตนนิโรธา ผสฺสนิโรโธ, ผสฺสนิโรธา เวทนานิโรโธ, เวทนานิโรธา ตณฺหานิโรโธ, ตณฺหานิโรธา อุปาทานนิโรโธ, อุปาทานนิโรธา ภวนิโรโธ, ภวนิโรธา ชาตินิโรโธ, ชาตินิโรธา ชรามรณํ โสกปริเทวทุกฺขโทมนุสฺสุปายาสา นิรุชฺฌนฺติ, เอวเมตสฺส เกวลสฺส ทุกฺขกฺขนฺธสฺส นิโรโธ โหตีติ.}

\prose{อถ โข ภควา เอตมตฺถํ วิทิตฺวา ตายํ เวลายํ อิมํ อุทานํ อุทาเนสิ\csromandash{}}

\csromangathameasure{“ยทา หเว ปาตุภวนฺติ ธมฺมา, \\ อาตาปิโน ฌายโต พฺราหฺมณสฺส. \\ อถสฺส กงฺขา วปยนฺติ สพฺพา, \\ ยโต ขยํ ปจฺจยานํ อเวที”ติ.ทุติยํ.}
\csromangathagroup{\csromangathastanza{“ยทา หเว ปาตุภวนฺติ ธมฺมา, \\ อาตาปิโน ฌายโต พฺราหฺมณสฺส. \\ อถสฺส กงฺขา วปยนฺติ สพฺพา, \\ ยโต ขยํ ปจฺจยานํ อเวที”ติ.\csromanspacer{}\csromanspacer{}ทุติยํ.}}

\csromanheader{2}{3. ตติยโพธิสุตฺต}
\proseitem{3}{\csromansymbolmark{*}เอวํ เม สุตํ\csromandash{}เอกํ สมยํ ภควา อุรุเวลายํ วิหรติ นชฺชา เนรญฺชราย ตีเร โพธิรุกฺขมูเล ปฐมาภิสมฺพุทฺโธ.\csromanspacer{} เตน โข ปน}

\vspace{\tipitakaparskip}
\csromancenter{โพธิวคฺค สมเยน ภควา สตฺตาหํ เอกปลฺลงฺเกน นิสินฺโน โหติ วิมุตฺติสุขปฏิสํเวที.\csromanspacer{} อถ โข ภควา ตสฺส สตฺตาหสฺส อจฺจเยน ตมฺหา สมาธิมฺหา วุฏฺฐหิตฺวา รตฺติยา ปจฺฉิมํ ยามํ ปฏิจฺจสมุปฺปาทํ อนุโลมปฏิโลมํ สาธุกํ มนสากาสิ\csromandash{}}
\prose{\csromanfolio{79}อิติ อิมสฺมิํ สติ อิทํ โหติ, อิมสฺสุปฺปาทา อิทํ อุปฺปชฺชติ, อิมสฺมิํ อสติ อิทํ น โหติ, อิมสฺส นิโรธา อิทํ นิรุชฺฌติ, ยทิทํ อวิชฺชาปจฺจยา สงฺขารา, สงฺขารปจฺจยา วิญฺญาณํ, วิญฺญาณปจฺจยา นามรูปํ, นามรูปปจฺจยา สฬายตนํ, สฬายตนปจฺจยา ผสฺโส, ผสฺสปจฺจยา เวทนา, เวทนาปจฺจยา ตณฺหา, ตณฺหาปจฺจยา อุปาทานํ, อุปาทานปจฺจยา ภโว, ภวปจฺจยา ชาติ, ชาติปจฺจยา ชรามรณํ โสกปริเทวทุกฺขโทมนสฺสุปายาสา สมฺภวนฺติ, เอวเมตสฺส เกวลสฺส ทุกฺขกฺขนฺธสฺส สมุทโย โหติ.}

\prose{อวิชฺชาย เตฺวว อเสสวิราคนิโรธา สงฺขารนิโรโธ, สงฺขารนิโรธา วิญฺญาณนิโรโธ, วิญฺญาณนิโรธา นามรูปนิโรโธ, นามรูปนิโรธา สฬายตนนิโรโธ, สฬายตนนิโรธา ผสฺสนิโรโธ, ผสฺสนิโรธา เวทนานิโรโธ, เวทนานิโรธา ตณฺหานิโรโธ, ตณฺหานิโรธา อุปาทานนิโรโธ, อุปาทานนิโรธา ภวนิโรโธ, ภวนิโรธา ชาตินิโรโธ, ชาตินิโรธา ชรามรณํ โสกปริเทวทุกฺขโทมนสฺสุปายาสา นิรุชฺฌนฺติ, เอวเมตสฺส เกวลสฺส ทุกฺขกฺขนฺธสฺส นิโรโธ โหตีติ.}

\prose{อถ โข ภควา เอตมตฺถํ วิทิตฺวา ตายํ เวลายํ อิมํ อุทานํ อุทาเนสิ\csromandash{}}

\csromangathameasure{“ยทา หเว ปาตุภวนฺติ ธมฺมา, \\ อาตาปิโน ฌายโต พฺราหฺมณสฺส. \\ วิธูปยํ ติฏฺฐติ มารเสนํ, \\ สูริโยว โอภาสยมนฺตลิกฺขนฺ”ติ.ตติยํ.}
\csromangathagroup{\csromangathastanza{“ยทา หเว ปาตุภวนฺติ ธมฺมา, \\ อาตาปิโน ฌายโต พฺราหฺมณสฺส. \\ วิธูปยํ ติฏฺฐติ มารเสนํ, \\ สูริโยว\footnote{สุริโยว (สี, สฺยา, กํ, อิ)} โอภาสยมนฺตลิกฺขนฺ”ติ.\csromanspacer{}\csromanspacer{}ตติยํ.}}

\vspace{\tipitakaparskip}
\csromancenter{อุทานปาฬิ \textbf{4.}\csromanspacer{}\textbf{ หุํหุงฺกสุตฺต} 4. * เอวํ เม สุตํ\csromandash{}เอกํ สมยํ ภควา อุรุเวลายํ วิหรติ นชฺชา เนรญฺชราย ตีเร อชปาลนิโคฺรเธ ปฐมาภิสมฺพุทฺโธ.\csromanspacer{} เตน โข ปน สมเยน ภควา สตฺตาหํ เอกปลฺลงฺเกน นิสินฺโน โหติ วิมุตฺติสุขปฏิสํเวที.\csromanspacer{} อถ โข ภควา ตสฺส สตฺตาหสฺส อจฺจเยน ตมฺหา สมาธิมฺหา วุฏฺฐาสิ.}
\csromanmatikaanchor{mtk.37}
\csromanmatikaanchor{mtk.43}
\csromanmatikaanchor{mtk.44}
\csromantocmarkref{section}{26. พฺราหฺมณวคฺค}{mtk.37}
\bookmark[dest={mtk.37},level=1]{26. พฺราหฺมณวคฺค}
\prose{\csromanfolio{80}อถ โข อญฺญตโร หุํหุงฺกชาติโก\footnote{หุหุงฺกชาติโก (สี, สฺยา, กํ, อิ)} พฺราหฺมโณ เยน ภควา เตนุปสงฺกมิ, อุปสงฺกมิตฺวา ภควตา สทฺธิํ สมฺโมทิ, สมฺโมทนียํ กถํ สารณียํ วีติสาเรตฺวา เอกมนฺตํ อฏฺฐาสิ, เอกมนฺตํ ฐิโต โข โส พฺราหฺมโณ ภควนฺตํ เอตทโวจ “กิตฺตาวตา นุ โข โภ โคตม พฺราหฺมโณ โหติ, กตเม จ ปน พฺราหฺมณกรณา\footnote{พฺราหฺมณการกา (ก)} ธมฺมา”ติ.}

\prose{อถ โข ภควา เอตมตฺถํ วิทิตฺวา ตายํ เวลายํ อิมํ อุทานํ อุทาเนสิ\csromandash{}}

\csromangathameasure{+“โย พฺราหฺมโณ พาหิตปาปธมฺโม, \\ นิหุํหุงฺโก นิกฺกสาโว ยตตฺโต. \\ เวทนฺตคู วูสิตพฺรหฺมจริโย, \\ ธมฺเมน โส พฺรหฺมวาทํ วเทยฺย. \\ ยสฺสุ’สฺสทา นตฺถิ กุหิญฺจิ โลเก”ติ.จตุตฺถํ.}
\csromangathagroup{\csromangathastanza{\csromansymbolfootnote{+}{ขุ 10. 129 ปิฏฺเฐปิ.}“โย พฺราหฺมโณ พาหิตปาปธมฺโม, \\ นิหุํหุงฺโก\footnote{นิหุหุงฺโก (สี, สฺยา, กํ, อิ)} นิกฺกสาโว ยตตฺโต. \\ เวทนฺตคู วูสิตพฺรหฺมจริโย, \\ ธมฺเมน โส พฺรหฺมวาทํ วเทยฺย.}\csromangathastanza{ยสฺสุ’สฺสทา นตฺถิ กุหิญฺจิ โลเก”ติ.\csromanspacer{}\csromanspacer{}จตุตฺถํ.}}

\vspace{\tipitakaparskip}
\csromancenter{พฺราหฺมณสุตฺต 5.\csromanspacer{} เอวํ เม สุตํ\csromandash{}เอกํ สมยํ ภควา สาวตฺถิยํ วิหรติ เชตวเน อนาถปิณฺฑิกสฺส อาราเม.\csromanspacer{} เตน โข ปน สมเยน อายสฺมา จ สาริปุตฺโต อายสฺมา จ มหาโมคฺคลฺลาโน อายสฺมา จ มหากสฺสโป อายสฺมา จ มหากจฺจาโน\footnote{มหากจฺจายโน (สี, อิ, ก)} อายสฺมา จ มหาโกฏฺฐิโก อายสฺมา จ มหากปฺปิโน อายสฺมา จ มหาจุนฺโท อายสฺมา จ อนุรุทฺโธ อายสฺมา จ เรวโต อายสฺมา จ นนฺโท\footnote{อานนฺโท (สี, อิ)} เยน ภควา เตนุปสงฺกมิํสุ.}
\chapterheadpageread{1. โพธิวคฺค}
\prose{\csromanfolio{81}อทฺทสา โข ภควา เต อายสฺมนฺเต ทูรโตว อาคจฺฉนฺเต, ทิสฺวาน ภิกฺขู อามนฺเตสิ “เอเต ภิกฺขเว พฺราหฺมณา อาคจฺฉนฺติ, เอเต ภิกฺขเว พฺราหฺมณา อาคจฺฉนฺตี”ติ.\csromanspacer{} เอวํ วุตฺเต อญฺญตโร พฺราหฺมณชาติโก ภิกฺขุ ภควนฺตํ เอตทโวจ “กิตฺตาวตา นุ โข ภนฺเต พฺราหฺมโณ โหติ, กตเม จ ปน พฺราหฺมณกรณา ธมฺมา”ติ.}

\prose{อถ โข ภควา เอตมตฺถํ วิทิตฺวา ตายํ เวลายํ อิมํ อุทานํ อุทาเนสิ\csromandash{}}

\prose{\csromansymbolfootnote{*}{ขุ 10. 129 ปิฏฺเฐปิ.}“พาหิตฺวา ปาปเก ธมฺเม, เย จรนฺติ สทา สตา.}

\prose{ขีณสํโยชนา พุทฺธา, เต เว\footnote{เตว (สี)} โลกสฺมิ พฺราหฺมณา”ติ.\csromanspacer{} ปญฺจมํ.}

\csromanheader{2}{6. มหากสฺสปสุตฺต}
\proseitem{6}{เอวํ เม สุตํ\csromandash{}เอกํ สมยํ ภควา ราชคเห วิหรติ เวฬุวเน กลนฺทกนิวาเป.\csromanspacer{} เตน โข ปน สมเยน อายสฺมา มหากสฺสโป ปิปฺปลิคุหายํ\footnote{วิปฺผลิคุหายํ (สฺยา), สิมฺพลิคุหายํ (ก)} วิหรติ อาพาธิโก\footnote{อาพาธิโก โหติ (สฺยา, อิ) สํ 3. 71 ปิฏฺเฐปิ ปสฺสิตพฺพํ.} ทุกฺขิโต พาฬฺหคิลาโน.\csromanspacer{} อถ โข อายสฺมา มหากสฺสโป อปเรน สมเยน ตมฺหา อาพาธา วุฏฺฐาสิ, อถ โข อายสฺมโต มหากสฺสปสฺส ตมฺหา อาพาธา วุฏฺฐิตสฺส เอตทโหสิ “ยํนูนาหํ ราชคหํ ปิณฺฑาย ปวิเสยฺยนฺ”ติ.}

\prose{เตน โข ปน สมเยน ปญฺจมตฺตานิ เทวตาสตานิ อุสฺสุกฺกํ อาปนฺนานิ โหนฺติ อายสฺมโต มหากสฺสปสฺส ปิณฺฑปาตปฏิลาภาย.\csromanspacer{} อถ โข อายสฺมา มหากสฺสโป ตานิ ปญฺจมตฺตานิ เทวตาสตานิ ปฏิกฺขิปิตฺวา ปุพฺพณฺหสมยํ นิวาเสตฺวา ปตฺตจีวรมาทาย ราชคหํ ปิณฺฑาย ปาวิสิ เยน ทลิทฺทวิสิขา กปณวิสิขา เปสการวิสิขา, อทฺทสา โข ภควา อายสฺมนฺตํ มหากสฺสปํ ราชคเห ปิณฺฑาย จรนฺตํ เยน ทลิทฺทวิสิขา กปณวิสิขา เปสการวิสิขา.}

\prose{อถ โข ภควา เอตมตฺถํ วิทิตฺวา ตายํ เวลายํ อิมํ อุทานํ อุทาเนสิ\csromandash{}}

\prose{“อนญฺญโปสิ’มญฺญาตํ, ทนฺตํ สาเร ปติฏฺฐิตํ.}

\prose{ขีณาสวํ วนฺตโทสํ, ตมหํ พฺรูมิ พฺราหฺมณนฺ”ติ.\csromanspacer{}\csromanspacer{}ฉฏฺฐํ.}

\csromanheader{2}{อุทานปาฬิ \textbf{7. อชกลาปกสุตฺต}}
\proseitem{7}{\csromanfolio{82}เอวํ เม สุตํ\csromandash{}เอกํ สมยํ ภควา ปาวายํ\footnote{ปาฏลิยํ (อิ)} วิหรติ อชกลาปเก เจติเย อชกลาปกสฺส ยกฺขสฺส ภวเน.\csromanspacer{} เตน โข ปน สมเยน ภควา รตฺตนฺธการติมิสายํ อพฺโภกาเส นิสินฺโน โหติ, เทโว จ เอกเมกํ ผุสายติ.\csromanspacer{} อถ โข อชกลาปโก ยกฺโข ภควโต ภยํ ฉมฺภิตตฺตํ โลมหํสํ อุปฺปาเทตุกาโม เยน ภควา เตนุปสงฺกมิ, อุปสงฺกมิตฺวา ภควโต อวิทูเร ติกฺขตฺตุํ “อกฺกุโล ปกฺกุโล”ติ อกฺกุลปกฺกุลิกํ อกาสิ “เอโส เต สมณ ปิสาโจ”ติ.}

\prose{อถ โข ภควา เอตมตฺถํ วิทิตฺวา ตายํ เวลายํ อิมํ อุทานํ อุทาเนสิ\csromandash{}}

\csromangathameasure{“ยทา สเกสุ ธมฺเมสุ, ปารคู โหติ พฺราหฺมโณ. \\ อถ เอตํ ปิสาจญฺจ, ปกฺกุลญฺจา’ติวตฺตตี”ติ.สตฺตมํ.}
\csromangathagroup{\csromangathastanza{“ยทา สเกสุ ธมฺเมสุ, ปารคู โหติ พฺราหฺมโณ. \\ อถ เอตํ ปิสาจญฺจ, ปกฺกุลญฺจา’ติวตฺตตี”ติ.\csromanspacer{}\csromanspacer{}สตฺตมํ.}}

\csromanheader{2}{8. สงฺคามชิสุตฺต}
\proseitem{8}{เอวํ เม สุตํ\csromandash{}เอกํ สมยํ ภควา สาวตฺถิยํ วิหรติ เชตวเน อนาถปิณฺฑิกสฺส อาราเม.\csromanspacer{} เตน โข ปน สมเยน อายสฺมา สงฺคามชิ สาวตฺถิํ อนุปฺปตฺโต โหติ ภควนฺตํ ทสฺสนาย.\csromanspacer{} อสฺโสสิ โข อายสฺมโต สงฺคามชิสฺส ปุราณทุติยิกา “อยฺโย กิร สงฺคามชิ สาวตฺถิํ อนุปฺปตฺโต”ติ, สา ทารกํ อาทาย เชตวนํ อคมาสิ.}

\prose{เตน โข ปน สมเยน อายสฺมา สงฺคามชิ อญฺญตรสฺมิํ รุกฺขมูเล ทิวาวิหารํ นิสินฺโน โหติ.\csromanspacer{} อถ โข อายสฺมโต สงฺคามชิสฺส ปุราณทุติยิกา เยนายสฺมา สงฺคามชิ เตนุปสงฺกมิ, อุปสงฺกมิตฺวา อายสฺมนฺตํ สงฺคามชิํ เอตทโวจ “ขุทฺทปุตฺตํ หิ\footnote{ขุทฺทปุตฺตามฺหิ (สี)} สมณ โปส มนฺ”ติ.\csromanspacer{} เอวํ วุตฺเต อายสฺมา สงฺคามชิ ตุณฺหี อโหสิ.}

\prose{ทุติยมฺปิ โข อายสฺมโต สงฺคามชิสฺส ปุราณทุติยิกา อายสฺมนฺตํ สงฺคามชิํ เอตทโวจ “ขุทฺทปุตฺตํ หิ สมณ โปส มนฺ”ติ.\csromanspacer{} ทุติยมฺปิ โข อายสฺมา สงฺคามชิ ตุณฺหี อโหสิ.}

\chapterheadpageread{1. โพธิวคฺค}
\prose{\csromanfolio{83}ตติยมฺปิ โข อายสฺมโต สงฺคามชิสฺส ปุราณทุติยิกา อายสฺมนฺตํ สงฺคามชิํ เอตทโวจ “ขุทฺทปุตฺตํ หิ สมณ โปส มนฺ”ติ.\csromanspacer{} ตติยมฺปิ โข อายสฺมา สงฺคามชิ ตุณฺหี อโหสิ.}

\prose{อถ โข อายสฺมโต สงฺคามชิสฺส ปุรณทุติยิกา ตํ ทารกํ อายสฺมโต สงฺคามชิสฺส ปุรโต นิกฺขิปิตฺวา ปกฺกามิ\footnote{ปกฺกมิ (ก) เอวมุปริปิ.} “เอโส\footnote{เอส (สี, ก)} เต สมณ ปุตฺโต, โปส นนฺ”ติ.}

\prose{อถ โข อายสฺมา สงฺคามชิ ตํ ทารกํ เนว โอโลเกสิ นาปิ อาลปิ.\csromanspacer{} อถ โข อายสฺมโต สงฺคามชิสฺส ปุราณทุติยิกา อวิทูรํ\footnote{อวิทูเร (สฺยา, อิ)} คนฺตฺวา อปโลเกนฺตี อทฺทส อายสฺมนฺตํ สงฺคามชิํ ตํ ทารกํ เนว โอโลเกนฺตํ นาปิ อาลปนฺตํ, ทิสฺวานสฺสา เอตทโหสิ “น จายํ สมโณ ปุตฺเตนปิ อตฺถิโก”ติ ตโต ปฏินิวตฺติตฺวา ทารกํ อาทาย ปกฺกามิ.\csromanspacer{} อทฺทสา โข ภควา ทิพฺเพน จกฺขุนา วิสุทฺเธน อติกฺกนฺตมานุสเกน อายสฺมโต สงฺคามชิสฺส ปุราณทุติยิกาย เอวรูปํ วิปฺปการํ.}

\prose{อถ โข ภควา เอตมตฺถํ วิทิตฺวา ตายํ เวลายํ อิมํ อุทานํ อุทาเนสิ\csromandash{} * “อายนฺติํ นาภินนฺทติ, ปกฺกมนฺติํ น โสจติ.}

\prose{สงฺคา สงฺคามชิํ มุตฺตํ, ตมหํ พฺรูมิ พฺราหฺมณนฺ”ติ.\csromanspacer{}\csromanspacer{}}

\prose{อฏฺฐมํ.}

\csromanheader{2}{9. ชฏิลสุตฺต}
\proseitem{9}{เอวํ เม สุตํ\csromandash{}เอกํ สมยํ ภควา คยายํ วิหรติ คยาสีเส.\csromanspacer{} เตน โข ปน สมเยน สมฺพหุลา ชฏิลา สีตาสุ เหมนฺติกาสุ รตฺตีสุ อนฺตรฏฺฐเก หิมปาตสมเย คยายํ อุมฺมุชฺชนฺติปิ นิมุชฺชนฺติปิ, อุมฺมุชฺชนิมุชฺชมฺปิ กโรนฺติ โอสิญฺจนฺติปิ, อคฺคิมฺปิ ชุหนฺติ “อิมินา สุทฺธี”ติ.}

\prose{อทฺทสา โข ภควา เต สมฺพหุเล ชฏิเล สีตาสุ เหมนฺติกาสุ รตฺตีสุ อนฺตรฏฺฐเก หิมปาตสมเย คยายํ อุมฺมุชฺชนฺเตปิ นิมุชฺชนฺเตปิ อุมฺมุชฺชนิมุชฺชมฺปิ กโรนฺเต\footnote{อุมฺมุชฺชนิมุชฺชํ กโรนฺเตปิ (สี, อิ, ก)} โอสิญฺจนฺเตปิ อคฺคิมฺปิ ชุหนฺเต “อิมินา สุทฺธี”ติ.}

\gambhira{อุทานปาฬิ}
\prose{\csromanfolio{84}อถ โข ภควา เอตมตฺถํ วิทิตฺวา ตายํ เวลายํ อิมํ อุทานํ อุทาเนสิ\csromandash{}}

\csromangathameasure{*“น อุทเกน สุจี โหติ, พเหฺวตฺถ นฺหายตี ชโน. \\ ยมฺหิ สจฺจญฺจ ธมฺโม จ, โส สุจี โส จ พฺราหฺมโณ”ติ. นวมํ.}
\csromangathagroup{\csromangathastanza{\csromansymbolfootnote{*}{ขุ 10-130-ปิฏฺเฐปิ.}“น อุทเกน สุจี โหติ, พเหฺวตฺถ นฺหายตี\footnote{นหายตี (สี)} ชโน. \\ ยมฺหิ สจฺจญฺจ ธมฺโม จ, โส สุจี โส จ พฺราหฺมโณ”ติ.\csromanspacer{} นวมํ.}}

\csromanheader{2}{10. พาหิยสุตฺต}
\proseitem{10}{เอวํ เม สุตํ\csromandash{}เอกํ สมยํ ภควา สาวตฺถิยํ วิหรติ เชตวเน อนาถปิณฺฑิกสฺส อาราเม.\csromanspacer{} เตน โข ปน สมเยน พาหิโย ทารุจีริโย สุปฺปารเก ปฏิวสติ สมุทฺทตีเร สกฺกโต ครุกโต มานิโต ปูชิโต อปจิโต ลาภี จีวร ปิณฺฑปาตเสนาสน คิลานปจฺจย เภสชฺช ปริกฺขารานํ.\csromanspacer{} อถ โข พาหิยสฺส ทารุจีริยสฺส รโหคตสฺส ปฏิสลฺลีนสฺส เอวํ เจตโส ปริวิตกฺโก อุทปาทิ “เย โข เกจิ โลเก อรหนฺโต วา อรหตฺตมคฺคํ วา สมาปนฺนา, อหํ เตสํ อญฺญตโร”ติ. อถ โข พาหิยสฺส ทารุจีริยสฺส ปุราณสาโลหิตา เทวตา อนุกมฺปิกา อตฺถกามา พาหิยสฺส ทารุจีริยสฺส เจตสา เจโตปริวิตกฺกมญฺญาย เยน พาหิโย ทารุจีริโย เตนุปสงฺกมิ, อุปสงฺกมิตฺวา พาหิยํ ทารุจีริยํ เอตทโวจ “เนว โข ตฺวํ พาหิย อรหา นาปิ อรหตฺตมคฺคํ วา สมาปนฺโน, สาปิ เต ปฏิปทา นตฺถิ, ยาย ตฺวํ อรหา วา อสฺส\footnote{อสฺสสิ (สฺยา, ก)} อรหตฺตมคฺคํ วา สมาปนฺโน”ติ.}

\prose{อถ เก จรหิ เทวเต โลเก อรหนฺโต วา อรหตฺตมคฺคํ วา สมาปนฺนาติ.\csromanspacer{} อตฺถิ พาหิย อุตฺตเรสุ ชนปเทสุ\footnote{ชนปเท (สี)} สาวตฺถี นาม นครํ, ตตฺถ โส ภควา เอตรหิ วิหรติ อรหํ สมฺมาสมฺพุทฺโธ, โส หิ พาหิย ภควา อรหา เจว อรหตฺตาย จ ธมฺมํ เทเสตีติ.}

\prose{อถ โข พาหิโย ทารุจีริโย ตาย เทวตาย สํเวชิโต ตาวเทว สุปฺปารกมฺหา ปกฺกามิ สพฺพตฺถ เอกรตฺติปริวาเสน เยน สวตฺถี เชตวนํ อนาถปิณฺฑิกสฺส อาราโม เตนุปสงฺกมิ.\csromanspacer{} เตน โข ปน}

\vspace{\tipitakaparskip}
\csromancenter{โพธิวคฺค สมเยน สมฺพหุลา ภิกฺขู อพฺโภกาเส จงฺกมนฺติ.\csromanspacer{} อถ โข พาหิโย ทารุจีริโย เยน เต ภิกฺขู เตนุปสงฺกมิ, อุปสงฺกมิตฺวา เต ภิกฺขู เอตทโวจ “กหํ นุ โข ภนฺเต เอตรหิ ภควา วิหรติ อรหํ สมฺมาสมฺพุทฺโธ, ทสฺสนกามมฺหา มยํ ตํ ภควนฺตํ อรหนฺตํ สมฺมาสมฺพุทฺธนฺ”ติ.\csromanspacer{} อนฺตรฆรํ ปวิฏฺโฐ โข พาหิย ภควา ปิณฺฑายาติ.}
\prose{\csromanfolio{85}อถ โข พาหิโย ทารุจีริโย ตรมานรูโป เชตวนา นิกฺขมิตฺวา สาวตฺถิํ ปวิสิตฺวา อทฺทส ภควนฺตํ สาวตฺถิยํ ปิณฺฑาย จรนฺตํ ปาสาทิกํ ปสาทนียํ สนฺตินฺทฺริยํ สนฺตมานสํ อุตฺตมทมถสมถมนุปฺปตฺตํ ทนฺตํ คุตฺตํ ยตินฺทฺริยํ นาคํ, ทิสฺวาน เยน ภควา เตนุปสงฺกมิ, อุปสงฺกมิตฺวา ภควโต ปาเท สิรสา นิปติตฺวา ภควนฺตํ เอตทโวจ “เทเสตุ เม ภนฺเต ภควา ธมฺมํ, เทเสตุ สุคโต ธมฺมํ, ยํ มมสฺส ทีฆรตฺตํ หิตาย สุขายา”ติ.\csromanspacer{} เอวํ วุตฺเต ภควา พาหิยํ ทารุจีริยํ เอตทโวจ “อกาโล โข ตาว พาหิย, อนฺตรฆรํ ปวิฏฺฐมฺหา ปิณฺฑายา”ติ.}

\prose{ทุติยมฺปิ โข พาหิโย ทารุจีริโย ภควนฺตํ เอตทโวจ “ทุชฺชานํ โข ปเนตํ ภนฺเต ภควโต วา ชีวิตนฺตรายานํ, มยฺหํ วา ชีวิตนฺตรายานํ.\csromanspacer{} เทเสตุ เม ภนฺเต ภควา ธมฺมํ, เทเสตุ สุคโต ธมฺมํ, ยํ มมสฺส ทีฆรตฺตํ หิตาย สุขายา”ติ.\csromanspacer{} ทุติยมฺปิ โข ภควา พาหิยํ ทารุจีริยํ เอตทโวจ “อกาโล โข ตาว พาหิย, อนฺตรฆรํ ปวิฏฺฐมฺหา ปิณฺฑายา”ติ.}

\prose{ตติยมฺปิ โข พาหิโย ทารุจีริโย ภควนฺตํ เอตทโวจ “ทุชฺชานํ โข ปเนตํ ภนฺเต ภควโต วา ชีวิตนฺตรายานํ, มยฺหํ วา ชีวิตนฺตรายานํ.\csromanspacer{} เทเสตุ เม ภนฺเต ภควา ธมฺมํ, เทเสตุ สุคโต ธมฺมํ, ยํ มมสฺส ทีฆรตฺตํ หิตาย สุขายา”ติ.}

\prose{ตสฺมาติห เต พาหิย เอวํ สิกฺขิตพฺพํ “ทิฏฺเฐ ทิฏฺฐมตฺตํ ภวิสฺสติ, สุเต สุตมตฺตํ ภวิสฺสติ, มุเต มุตมตฺตํ ภวิสฺสติ, วิญฺญาเต วิญฺญาตมตฺตํ ภวิสฺสตี”ติ, เอวญฺหิ เต พาหิย สิกฺขิตพฺพํ.\csromanspacer{} ยโต โข เต พาหิย ทิฏฺเฐ ทิฏฺฐมตฺตํ ภวิสฺสติ, สุเต สุตมตฺตํ ภวิสฺสติ, มุเต มุตมตฺตํ ภวิสฺสติ, วิญฺญาเต วิญฺญาตมตฺตํ ภวิสฺสติ, ตโต ตฺวํ พาหิย น เตน.\csromanspacer{} ยโต ตฺวํ พาหิย น เตน, ตโต ตฺวํ พาหิย น}

\vspace{\tipitakaparskip}
\csromancenter{อุทานปาฬิ ตตฺถ.\csromanspacer{} ยโต ตฺวํ พาหิย น ตตฺถ, ตโต ตฺวํ พาหิย เนวิธ น หุรํ น อุภยมนฺตเรน.\csromanspacer{} เอเสวนฺโต ทุกฺขสฺสาติ.}
\prose{\csromanfolio{86}อถ โข พาหิยสฺส ทารุจีริยสฺส ภควโต อิมาย สํขิตฺตาย ธมฺมเทสนาย ตาวเทว อนุปาทาย อาสเวหิ จิตฺตํ วิมุจฺจิ.}

\prose{อถ โข ภควา พาหิยํ ทารุจีริยํ อิมินา สํขิตฺเตน โอวาเทน โอวทิตฺวา ปกฺกามิ.\csromanspacer{} อถ โข อจิรปกฺกนฺตสฺส ภควโต พาหิยํ ทารุจีริยํ คาวี ตรุณวจฺฉา อธิปติตฺวา\footnote{อธิปาเตตฺวา (สี, สฺยา, อิ), อธิปาติตฺวา (ก)} ชีวิตา โวโรเปสิ.}

\prose{อถ โข ภควา สาวตฺถิยํ ปิณฺฑาย จริตฺวา ปจฺฉาภตฺตํ ปิณฺฑปาตปฏิกฺกนฺโต สมฺพหุเลหิ ภิกฺขูหิ สทฺธิํ นครมฺหา นิกฺขมิตฺวา อทฺทส พาหิยํ ทารุจีริยํ กาลงฺกตํ\footnote{กาลกต (สี, สฺยา, กํ)}, ทิสฺวาน ภิกฺขู อามนฺเตสิ “คณฺหถ ภิกฺขเว พาหิยสฺส ทารุจีริยสฺส สรีรกํ, มญฺจกํ อาโรเปตฺวา นีหริตฺวา ฌาเปถ, ถูปญฺจสฺส กโรถ, สพฺรหฺมจารี โว ภิกฺขเว กาลงฺกโต”ติ.}

\prose{“เอวํ ภนฺเต”ติ โข เต ภิกฺขู ภควโต ปฏิสฺสุตฺวา พาหิยสฺส ทารุจีริยสฺส สรีรกํ มญฺจกํ อาโรเปตฺวา นีหริตฺวา ฌาเปตฺวา ถูปญฺจสฺส กตฺวา เยน ภควา เตนุปสงฺกมิํสุ, อุปสงฺกมิตฺวา ภควนฺตํ อภิวาเทตฺวา เอกมนฺตํ นิสีทิํสุ, เอกมนฺตํ นิสินฺนา โข เต ภิกฺขู ภควนฺตํ เอตทโวจุํ “ทฑฺฒํ ภนฺเต พาหิยสฺส ทารุจีริยสฺส สรีรํ, ถูโป จสฺส กโต, ตสฺส กา คติ, โก อภิสมฺปราโย”ติ.\csromanspacer{} ปณฺฑิโต ภิกฺขเว พาหิโย ทารุจีริโย ปจฺจปาทิ ธมฺมสฺสานุธมฺมํ, น จ มํ ธมฺมาธิกรณํ วิเหเสสิ, ปรินิพฺพุโต ภิกฺขเว พาหิโย ทารุจีริโยติ.}

\prose{อถ โข ภควา เอตมตฺถํ วิทิตฺวา ตายํ เวลายํ อิมํ อุทานํ อุทาเนสิ\csromandash{}}

\csromangathameasure{*“ยตฺถ อาโป จ ปถวี, เตโช วาโย น คาธติ. \\ น ตตฺถ สุกฺกา โชตนฺติ, อาทิจฺโจ นปฺปกาสติ. \\ น ตตฺถ จนฺทิมา ภาติ, ตโม ตตฺถ น วิชฺชติ.}
\csromangathagroup{\csromangathastanza{\csromansymbolfootnote{*}{ขุ 10. 129 ปิฏฺเฐปิ.}“ยตฺถ อาโป จ ปถวี, เตโช วาโย น คาธติ. \\ น ตตฺถ สุกฺกา โชตนฺติ, อาทิจฺโจ นปฺปกาสติ. \\ น ตตฺถ จนฺทิมา ภาติ, ตโม ตตฺถ น วิชฺชติ.}}

\chapterheadpageread{2. มุจลินฺทวคฺค}
\csromangathameasure{ยทา จ อตฺตนา’เวทิ, มุนิ โมเนน พฺราหฺมโณ. \\ อถ รูปา อรูปา จ, สุขทุกฺขา ปมุจฺจตี”ติ.ทสมํ.}
\csromangathagroup{\csromangathastanza{\csromanfolio{87}ยทา จ อตฺตนา’เวทิ\footnote{เวมี (ก) ขุ 10. 129 ปิฏฺเฐปิ.}, มุนิ โมเนน พฺราหฺมโณ. \\ อถ รูปา อรูปา จ, สุขทุกฺขา ปมุจฺจตี”ติ.\csromanspacer{}\csromanspacer{}ทสมํ.}}

\vspace{\tipitakaparskip}
\csromancenter{(อยมฺปิ อุทาโน วุตฺโต ภควตา อิติ เม สุตนฺติ.)\footnote{( ) สฺยามโปตฺถเก นตฺติ.}.\csromanspacer{} โพธิวคฺโค ปฐโม.}
\titlehead{\textbf{ตสฺสุทฺทานํ}}
\csromangathameasure{ตโย โพธี จ หุํหุงฺโก พฺราหฺมณา กสฺสเปน จ. อช สงฺคาม ชฏิลา, \\ พาหิเยนาติ เต ทสาติ.}
\csromangathagroup{\csromangathastanza{ตโย โพธี จ หุํหุงฺโก\footnote{ตโย จ โพธิ นิโคฺรโธ (สพฺพตฺถ)} พฺราหฺมณา\footnote{เต เถรา (สี, สฺยา, อิ), เถโร (ก)} กสฺสเปน จ. อช\footnote{ตโย จ โพธิ นิโคฺรโธ (สพฺพตฺถ)} สงฺคาม ชฏิลา, \\ พาหิเยนาติ เต ทสาติ.}}

\chapterhead{2. มุจลินฺทวคฺค}
\csromanheader{2}{1. มุจลินฺทสุตฺต}
\proseitem{11}{\csromansymbolfootnote{*}{วิ 3. 3 ปิฏฺเฐปิ.}เอวํ เม สุตํ\csromandash{}เอกํ สมยํ ภควา อุรุเวลายํ วิหรติ นชฺชา เนรญฺชราย ตีเร มุจลินฺทมูเล ปฐมาภิสมฺพุทฺโธ.\csromanspacer{} เตน โข ปน สมเยน ภควา สตฺตาหํ เอกปลฺลงฺเกน นิสินฺโน โหติ วิมุตฺติสุขปฏิสํเวที.}

\prose{เตน โข ปน สมเยน มหา อกาลเมโฆ อุทปาทิ สตฺตาหวทฺทลิกา สีตวาตทุทฺทินี.\csromanspacer{} อถ โข มุจลินฺโท นาคราชา สกภวนา นิกฺขมิตฺวา ภควโต กายํ สตฺตกฺขตฺตุํ โภเคหิ ปริกฺขิปิตฺวา อุปริมุทฺธนิ มหนฺตํ ผณํ วิหจฺจ อฏฺฐาสิ “มา ภควนฺตํ สีตํ, มา ภควนฺตํ อุณฺหํ, มา ภควนฺตํ ฑํส มกส วาตาตป สรีสป\footnote{สิริํสป (สี, สฺยา, กํ, อิ)} สมฺผสฺโส”ติ.}

\prose{อถ โข ภควา ตสฺส สตฺตาหสฺส อจฺจเยน ตมฺหา สมาธิมฺหา วุฏฺฐาสิ.\csromanspacer{} อถ โข มุจลินฺโท นาคราชา วิทฺธํ วิคตวลาหกํ เทวํ วิทิตฺวา ภควโต กายา โภเค วินิเวเฐตฺวา สกวณฺณํ ปฏิสํหริตฺวา มาณวกวณฺณํ อภินิมฺมินิตฺวา ภควโต ปุรโต อฏฺฐาสิ ปญฺชลิโก ภควนฺตํ นมสฺสมาโน.}

\gambhira{อุทานปาฬิ}
\csromanmatikaanchor{mtk.38}
\csromanmatikaanchor{mtk.52}
\csromantocmarkref{section}{อุทฺทานคาถา}{mtk.38}
\bookmark[dest={mtk.38},level=1]{อุทฺทานคาถา}
\prose{\csromanfolio{88}อถ โข ภควา เอตมตฺถํ วิทิตฺวา ตายํ เวลายํ อิมํ อุทานํ อุทาเนสิ\csromandash{} * “สุโข วิเวโก ตุฏฺฐสฺส, สุตธมฺมสฺส ปสฺสโต.}

\prose{อพฺยาปชฺชํ สุขํ โลเก, ปาณภูเตสุ สํยโม.}

\prose{สุขา วิราคตา โลเก, กามานํ สมติกฺกโม.}

\prose{อสฺมิมานสฺส โย วินโย, เอตํ เว ปรมํ สุขนฺ”ติ.\csromanspacer{}\csromanspacer{}ปฐมํ.}

\vspace{\tipitakaparskip}
\csromancenter{ราชสุตฺต 12.\csromanspacer{} เอวํ เม สุตํ\csromandash{}เอกํ สมยํ ภควา สาวตฺถิยํ วิหรติ เชตวเน อนาถปิณฺฑิกสฺส อาราเม.\csromanspacer{} เตน โข ปน สมเยน สมฺพหุลานํ ภิกฺขูนํ ปจฺฉาภตฺตํ ปิณฺฑปาตปฏิกฺกนฺตานํ อุปฏฺฐานสาลายํ สนฺนิสินฺนานํ สนฺนิปติตานํ อยมนฺตรากถา อุทปาทิ “โก นุ โข อาวุโส อิเมสํ ทฺวินฺนํ ราชูนํ มหทฺธนตโร วา มหาโภคตโร วา มหาโกสตโร วา มหาวิชิตตโร วา มหาวาหนตโร วา มหพฺพลตโร วา มหิทฺธิกตโร วา มหานุภาวตโร วา ราชา วา มาคโธ เสนิโย พิมฺพิสาโร, ราชา วา ปเสนทิ โกสโล”ติ.\csromanspacer{} อยญฺจรหิ เตสํ ภิกฺขูนํ อนฺตรากถา โหติ วิปฺปกตา.}
\prose{อถ โข ภควา สายนฺหสมยํ ปฏิสลฺลานา วุฏฺฐิโต เยนุปฏฺฐานสาลา เตนุปสงฺกมิ, อุปสงฺกมิตฺวา ปญฺญตฺเต อาสเน นิสีทิ, นิสชฺช โข ภควา ภิกฺขู อามนฺเตสิ “กาย นุตฺถ ภิกฺขเว เอตรหิ กถาย สนฺนิสินฺนา สนฺนิปติตา, กา จ ปน โว อนฺตรากถา วิปฺปกตา”ติ.}

\prose{อิธ ภนฺเต อมฺหากํ ปจฺฉาภตฺตํ ปิณฺฑปาตปฏิกฺกนฺตานํ อุปฏฺฐานสาลายํ สนฺนิสินฺนานํ สนฺนิปติตานํ อยมนฺตรากถา อุทปาทิ “โก นุ โข อาวุโส อิเมสํ ทฺวินฺนํ ราชูนํ มหทฺธนตโร วา มหาโภคตโร วา มหาโกสตโร วา มหาวิชิตตโร วา มหาวาหนตโร วา มหพฺพลตโร วา มหิทฺธิกตโร วา มหานุภาวตโร วา ราชา วา มาคโธ เสนิโย พิมฺพิสาโร, ราชา วา ปเสนทิ โกสโล”ติ.\csromanspacer{} อยํ โข โน ภนฺเต อนฺตรากถา วิปฺปกตา, อถ ภควา อนุปฺปตฺโตติ.}

\chapterheadpageread{2. มุจลินฺทวคฺค}
\csromanmatikaanchor{mtk.39}
\csromanmatikaanchor{mtk.53}
\csromantocmarknopage{chapter}{1. โพธิวคฺค}{mtk.39}
\bookmark[dest={mtk.39},level=0]{1. โพธิวคฺค}
\prose{\csromanfolio{89}น เขฺวตํ ภิกฺขเว ตุมฺหากํ ปติรูปํ กุลปุตฺตานํ สทฺธา อคารสฺมา อนคาริยํ ปพฺพชิตานํ, ยํ ตุเมฺห เอวรูปิํ กถํ กเถยฺยาถ, สนฺนิปติตานํ โว ภิกฺขเว ทฺวยํ กรณียํ ธมฺมี วา กถา อริโย วา ตุณฺหีภาโวติ.}

\prose{อถ โข ภควา เอตมตฺถํ วิทิตฺวา ตายํ เวลายํ อิมํ อุทานํ อุทาเนสิ\csromandash{}}

\csromangathameasure{“ยญฺจ กามสุขํ โลเก, ยญฺจิทํ ทิวิยํ สุขํ. \\ ตณฺหกฺขยสุขสฺเสเต, กลํ นาคฺฆนฺติ โสฬสินฺ”ติ.ทุติยํ.}
\csromangathagroup{\csromangathastanza{“ยญฺจ กามสุขํ โลเก, ยญฺจิทํ ทิวิยํ สุขํ. \\ ตณฺหกฺขยสุขสฺเสเต, กลํ นาคฺฆนฺติ โสฬสินฺ”ติ.\csromanspacer{}\csromanspacer{}ทุติยํ.}}

\vspace{\tipitakaparskip}
\csromancenter{ทณฺฑสุตฺต 13.\csromanspacer{} เอวํ เม สุตํ\csromandash{}เอกํ สมยํ ภควา สาวตฺถิยํ วิหรติ เชตวเน อนาถปิณฺฑิกสฺส อาราเม.\csromanspacer{} เตน โข ปน สมเยน สมฺพหุลา กุมารกา อนฺตรา จ สาวตฺถิํ อนฺตรา จ เชตวนํ อหิํ ทณฺเฑน หนนฺติ.\csromanspacer{} อถ โข ภควา ปุพฺพณฺหสมยํ นิวาเสตฺวา ปตฺตจีวรมาทาย สาวตฺถิํ ปิณฺฑาย ปาวิสิ, อทฺทสา โข ภควา สมฺพหุเล กุมารเก อนฺตรา จ สาวตฺถิํ อนฺตรา จ เชตวนํ อหิํ ทณฺเฑน หนนฺเต.}
\prose{อถ โข ภควา เอตมตฺถํ วิทิตฺวา ตายํ เวลายํ อิมํ อุทานํ อุทาเนสิ\csromandash{}}

\prose{\csromansymbolfootnote{*}{เหฏฺฐา 33 ปิฏฺเฐปิ.}“สุขกามานิ ภูตานิ, โย ทณฺเฑน วิหิํสติ.}

\csromangathameasure{อตฺตโน สุขเมสาโน, เปจฺจ โส น ลภเต สุขํ.}
\csromangathagroup{\csromangathastanza{อตฺตโน สุขเมสาโน, เปจฺจ โส น ลภเต สุขํ.}}

\prose{สุขกามานิ ภูตานิ, โย ทณฺเฑน น หิํสติ.}

\csromangathameasure{อตฺตโน สุขเมสาโน, เปจฺจ โส ลภเต สุขนฺ”ติ. ตติยํ.}
\csromangathagroup{\csromangathastanza{อตฺตโน สุขเมสาโน, เปจฺจ โส ลภเต สุขนฺ”ติ.\csromanspacer{} ตติยํ.}}

\csromanheader{2}{4. สกฺการสุตฺต}
\proseitem{14}{\csromansymbolfootnote{+}{อุปริ 128 ปิฏฺเฐปิ.}เอวํ เม สุตํ\csromandash{}เอกํ สมยํ ภควา สาวตฺถิยํ วิหรติ เชตวเน อนาถปิณฺฑิกสฺส อาราเม.\csromanspacer{} เตน โข ปน สมเยน ภควา สกฺกโต โหติ ครุกโต มานิโต ปูชิโต อปจิโต ลาภี จีวร ปิณฺฑปาตเสนาสน คิลานปจฺจย เภสชฺช ปริกฺขารานํ, ภิกฺขุสํโฆปิ สกฺกโต}

\vspace{\tipitakaparskip}
\csromancenter{อุทานปาฬิ โหติ ครุกโต มานิโต ปูชิโต อปจิโต ลาภี จีวรปิณฺฑปาตเสนาสน คิลานปจฺจย เภสชฺช ปริกฺขารานํ.\csromanspacer{} อญฺญติตฺถิยา ปน ปริพฺพาชกา อสกฺกตา โหนฺติ อครุกตา อมานิตา อปูชิตา อนปจิตา\footnote{น อปจิตา (สฺยา, อิ)} น ลาภิโน จีวร ปิณฺฑปาต เสนาสน คิลานปจฺจย เภสชฺช ปริกฺขารานํ.\csromanspacer{} อถ โข เต อญฺญติตฺถิยา ปริพฺพาชกา ภควโต สกฺการํ อสหมานา ภิกฺขุสํฆสฺส จ, คาเม จ อรญฺเญ จ ภิกฺขู ทิสฺวา อสพฺภาหิ ผรุสาหิ วาจาหิ อกฺโกสนฺติ ปริภาสนฺติ โรเสนฺติ วิเหเสนฺติ.}
\prose{\csromanfolio{90}อถ โข สมฺพหุลา ภิกฺขู เยน ภควา เตนุปสงฺกมิํสุ, อุปสงฺกมิตฺวา ภควนฺตํ อภิวาเทตฺวา เอกมนฺตํ นิสีทิํสุ, เอกมนฺตํ นิสินฺนา โข เต ภิกฺขู ภควนฺตํ เอตทโวจุํ “เอตรหิ ภนฺเต ภควา สกฺกโต ครุกโต มานิโต ปูชิโต อปจิโต ลาภี จีวร ปิณฺฑปาตเสนาสน คิลานปจฺจย เภสชฺช ปริกฺขารานํ, ภิกฺขุสํโฆปิ สกฺกโต ครุกโต มานิโต ปูชิโต อปจิโต ลาภี จีวร ปิณฺฑปาตเสนาสน คเลนปจฺจย เภสชฺช ปริกฺขารานํ.\csromanspacer{} อญฺญติตฺถิยา ปน ปริพฺพาชกา อสกฺกตา อครุกตา อมานิตา อปูชิตา อนปจิตา น ลาภิโน จีวรปิณฺฑปาต เสนาสน คิลานปจฺจย เภสชฺช ปริกฺขารานํ.\csromanspacer{} อถ โข เต ภนฺเต อญฺญติตฺถิยา ปริพฺพาชกา ภควโต สกฺการํ อสหมานา ภิกฺขุสํฆสฺส จ, คาเม จ อรญฺเญ จ ภิกฺขู ทิสฺวา อสพฺภาหิ ผรุสาหิ วาจาหิ อกฺโกสนฺติ ปริภาสนฺติ โรเสนฺติ วิเหเสนฺตี”ติ.}

\prose{อถ โข ภควา เอตมตฺถํ วิทิตฺวา ตายํ เวลายํ อิมํ อุทานํ อุทาเนสิ\csromandash{}}

\csromangathameasure{“คาเม อรญฺเญ สุขทุกฺขผุฏฺโฐ, \\ เนวตฺตโต โน ปรโต ทเหถ. \\ ผุสนฺติ ผสฺสา อุปธิํ ปฏิจฺจ, \\ นิรูปธิํ เกน ผุเสยฺยุ ผสฺสา”ติ.จตุตฺถํ.}
\csromangathagroup{\csromangathastanza{“คาเม อรญฺเญ สุขทุกฺขผุฏฺโฐ, \\ เนวตฺตโต โน ปรโต ทเหถ. \\ ผุสนฺติ ผสฺสา อุปธิํ ปฏิจฺจ, \\ นิรูปธิํ เกน ผุเสยฺยุ ผสฺสา”ติ.\csromanspacer{}\csromanspacer{}จตุตฺถํ.}}

\csromanheader{2}{5. อุปาสกสุตฺต}
\proseitem{15}{เอวํ เม สุตํ\csromandash{}เอกํ สมยํ ภควา สาวตฺถิยํ วิหรติ เชตวเน อนาถปิณฺฑิกสฺส อาราเม.\csromanspacer{} เตน โข ปน สมเยน อญฺญตโร}

\vspace{\tipitakaparskip}
\csromancenter{มุจลินฺทวคฺค อิจฺฉานงฺคลโก อุปาสโก สาวตฺถิํ อนุปฺปตฺโต โหติ เกนจิเทว กรณีเยน.\csromanspacer{} อถ โข โส อุปาสโก สาวตฺถิยํ ตํ กรณียํ ตีเรตฺวา เยน ภควา เตนุปสงฺกมิ, อุปสงฺกมิตฺวา ภควนฺตํ อภิวาเทตฺวา เอกมนฺตํ นิสีทิ.\csromanspacer{} เอกมนฺตํ นิสินฺนํ โข ตํ อุปาสกํ ภควา เอตทโวจ “จิรสฺสํ โข ตฺวํ อุปาสก อิมํ ปริยายมกาสิ, ยทิทํ อิธาคมนายา”ติ.}
\prose{\csromanfolio{91}จิรปฏิกาหํ ภนฺเต ภควนฺตํ ทสฺสนาย อุปสงฺกมิตุกาโม, อปิ จาหํ เกหิจิ เกหิจิ กิจฺจกรณีเยหิ พฺยาวโฏ, เอวาหํ นาสกฺขิํ ภควนฺตํ ทสฺสนาย อุปสงฺกมิตุนฺติ.}

\prose{อถ โข ภควา เอตมตฺถํ วิทิตฺวา ตายํ เวลายํ อิมํ อุทานํ อุทาเนสิ\csromandash{}}

\csromangathameasure{“สุขํ วต ตสฺส น โหติ กิญฺจิ, \\ สงฺขาตธมฺมสฺส พหุสฺสุตสฺส. \\ สกิญฺจนํ ปสฺส วิหญฺญมานํ, \\ ชโน ชนสฺมิํ ปฏิพนฺธรูโป”ติ.ปญฺจมํ.}
\csromangathagroup{\csromangathastanza{“สุขํ วต ตสฺส น โหติ กิญฺจิ, \\ สงฺขาตธมฺมสฺส พหุสฺสุตสฺส. \\ สกิญฺจนํ ปสฺส วิหญฺญมานํ, \\ ชโน ชนสฺมิํ ปฏิพนฺธรูโป”ติ.\csromanspacer{}\csromanspacer{}ปญฺจมํ.}}

\csromanheader{2}{6. คพฺภินีสุตฺต}
\proseitem{16}{เอวํ เม สุตํ\csromandash{}เอกํ สมยํ ภควา สาวตฺถิยํ วิหรติ เชตวเน อนาถปิณฺฑิกสฺส อาราเม.\csromanspacer{} เตน โข ปน สมเยน อญฺญตรสฺส ปริพฺพาชกสฺส ทหรมาณวิกา ปชาปติ โหติ คพฺภินี อุปวิชญฺญา.\csromanspacer{} อถ โข สา ปริพฺพาชิกา ตํ ปริพฺพาชกํ เอตทโวจ “คจฺฉ ตฺวํ พฺราหฺมณ เตลํ อาหร, ยํ เม วิชาตาย ภวิสฺสตี”ติ.}

\prose{เอวํ วุตฺเต โส ปริพฺพาชโก ตํ ปริพฺพาชิกํ เอตทโวจ “กุโต ปนาหํ โภติ\footnote{โภติยา (สฺยา, อิ, ก)} เตลํ อาหรามี”ติ.\csromanspacer{} ทุติยมฺปิ โข สา ปริพฺพาชิกา ตํ ปริพฺพาชกํ เอตทโวจ “คจฺฉ ตฺวํ พฺราหฺมณ เตลํ อาหร, ยํ เม วิชาตาย ภวิสฺสตี”ติ.\csromanspacer{} ทุติยมฺปิ โข โส ปริพฺพาชโก ตํ ปริพฺพาชิกํ เอตทโวจ “กุโต ปนาหํ โภติ เตลํ อาหรามี”ติ.\csromanspacer{} ตติยมฺปิ โข สา ปริพฺพาชิกา ตํ ปริพฺพาชกํ เอตทโวจ “คจฺฉ ตฺวํ พฺราหฺมณ เตลํ อาหร, ยํ เม วิชาตาย ภวิสฺสตี”ติ.}

\gambhira{อุทานปาฬิ}
\prose{\csromanfolio{92}เตน โข ปน สมเยน รญฺโญ ปเสนทิสฺส โกสลสฺส โกฏฺฐาคาเร สมณสฺส วา พฺราหฺมณสฺส วา สปฺปิสฺส วา เตลสฺส วา ยาวทตฺถํ ปาตุํ ทียติ\footnote{ทิยฺยติ (สี, ก)} โน นีหริตุํ.}

\prose{อถ โข ตสฺส ปริพฺพาชกสฺส เอตทโหสิ “รญฺโญ โข ปน ปเสนทิสฺส โกสลสฺส โกฏฺฐาคาเร สมณสฺส วา พฺราหฺมณสฺส วา สปฺปิสฺส วา เตลสฺส วา ยาวทตฺถํ ปาตุํ ทียติ โน นีหริตุํ, ยํนูนาหํ รญฺโญ ปเสนทิสฺส โกสลสฺส โกฏฺฐาคารํ คนฺตฺวา เตลสฺส ยาวทตฺถํ ปิวิตฺวา ฆรํ อาคนฺตฺวา อุจฺฉทฺทิตฺวาน\footnote{อุคฺคิริตฺวาน (สี, สฺยา, อิ), อุจฺฉทิตฺวา (สี-สฺยา-ฏฺฐ), อุจฺฉฑฺฑิตฺวาน (ก)} ทเทยฺยํ ยํ อิมิสฺสา วิชาตาย ภวิสฺสตี”ติ.}

\prose{อถ โข โส ปริพฺพาชโก รญฺโญ ปเสนทิสฺส โกสลสฺส โกฏฺฐาคารํ คนฺตฺวา เตลสฺส ยาวทตฺถํ ปิวิตฺวา ฆรํ อาคนฺตฺวา เนว สกฺโกติ อุทฺธํ กาตุํ น ปน อโธ, โส ทุกฺขาหิ ติพฺพาหิ\footnote{ติปฺปาหิ (สฺยา)} ขราหิ กฏุกาหิ เวทนาหิ ผุฏฺโฐ อาวฏฺฏติ ปริวฏฺฏติ.}

\prose{อถ โข ภควา ปุพฺพณฺหสมยํ นิวาเสตฺวา ปตฺตจีวรมาทาย สาวตฺถิํ ปิณฺฑาย ปาวิสิ, อทฺทสา โข ภควา ตํ ปริพฺพาชกํ ทุกฺขาหิ ติพฺพาหิ ขราหิ กฏุกาหิ เวทนาหิ ผุฏฺฐํ อาวฏฺฏมานํ ปริวฏฺฏมานํ.}

\prose{อถ โข ภควา เอตมตฺถํ วิทิตฺวา ตายํ เวลายํ อิมํ อุทานํ อุทาเนสิ\csromandash{}}

\csromangathameasure{“สุขิโน วต เย อกิญฺจนา, \\ เวทคุโน หิ ชนา อกิญฺจนา. \\ สกิญฺจนํ ปสฺส วิหญฺญมานํ, \\ ชโน ชนสฺมิํ ปฏิพนฺธจิตฺโต”4ติ.ฉฏฺฐํ.}
\csromangathagroup{\csromangathastanza{“สุขิโน วต เย อกิญฺจนา, \\ เวทคุโน หิ ชนา อกิญฺจนา. \\ สกิญฺจนํ ปสฺส วิหญฺญมานํ, \\ ชโน ชนสฺมิํ ปฏิพนฺธจิตฺโต”4ติ.\csromanspacer{}\csromanspacer{}ฉฏฺฐํ.}}

\csromanheader{2}{7. เอกปุตฺตกสุตฺต}
\proseitem{17}{เอวํ เม สุตํ\csromandash{}เอกํ สมยํ ภควา สาวตฺถิยํ วิหรติ เชตวเน อนาถปิณฺฑิกสฺส อาราเม.\csromanspacer{} เตน โข ปน สมเยน อญฺญตรสฺส อุปาสกสฺส เอกปุตฺตโก ปิโย มนาโป กาลงฺกโต โหติ.}

\chapterheadpageread{2. มุจลินฺทวคฺค}
\prose{\csromanfolio{93}อถ โข สมฺพหุลา อุปาสกา อลฺลวตฺถา อลฺลเกสา ทิวา ทิวสฺส เยน ภควา เตนุปสงฺกมิํสุ, อุปสงฺกมิตฺวา ภควนฺตํ อภิวาเทตฺวา เอกมนฺตํ นิสีทิํสุ, เอกมนฺตํ นิสินฺเน โข เต อุปาสเก ภควา เอตทโวจ “กิํ นุ โข ตุเมฺห อุปาสกา อลฺลวตฺถา อลฺลเกสา อิธูปสงฺกมนฺตา ทิวา ทิวสฺสา”ติ.}

\prose{เอวํ วุตฺเต โส อุปาสโก ภควนฺตํ เอตทโวจ “มยฺหํ โข ภนฺเต เอกปุตฺตโก ปิโย มนาโป กาลงฺกโต, เตน มยํ อลฺลวตฺถา อลฺลเกสา อิธูปสงฺกมนฺตา ทิวา ทิวสฺสา”ติ.}

\prose{อถ โข ภควา เอตมตฺถํ วิทิตฺวา ตายํ เวลายํ อิมํ อุทานํ อุทาเนสิ\csromandash{}}

\csromangathameasure{“ปิยรูปสฺสาทคธิตาเส, \\ เทวกายา ปุถุ มนุสฺสา จ. \\ อฆาวิโน ปริชุนฺนา, \\ มจฺจุราชสฺส วสํ คจฺฉนฺติ. \\ เย เว ทิวา จ รตฺโต จ, \\ อปฺปมตฺตา ชหนฺติ ปิยรูปํ. \\ เต เว ขณนฺติ อฆมูลํ, \\ มจฺจุโน อามิสํ ทุรติวตฺตนฺ”ติ.สตฺตมํ.}
\csromangathagroup{\csromangathastanza{“ปิยรูปสฺสาทคธิตาเส\footnote{ปิยรูปสฺสาตคธิตาเส (สี, อิ)}, \\ เทวกายา ปุถุ มนุสฺสา จ. \\ อฆาวิโน ปริชุนฺนา, \\ มจฺจุราชสฺส วสํ คจฺฉนฺติ.}\csromangathastanza{เย เว ทิวา จ รตฺโต จ, \\ อปฺปมตฺตา ชหนฺติ ปิยรูปํ. \\ เต เว ขณนฺติ อฆมูลํ, \\ มจฺจุโน อามิสํ ทุรติวตฺตนฺ”ติ.\csromanspacer{}\csromanspacer{}สตฺตมํ.}}

\csromanheader{2}{8. สุปฺปวาสาสุตฺต}
\proseitem{18}{เอวํ เม สุตํ\csromandash{}เอกํ สมยํ ภควา กุณฺฑิกายํ\footnote{กุณฺฑิยายํ (สี, สฺยา, อิ)} วิหรติ กุณฺฑธานวเน\footnote{กุณฺฑิฏฺฐานวเน (สฺย, อิ)}.\csromanspacer{} เตน โข ปน สมเยน สุปฺปวาสา โกลิยธีตา สตฺต วสฺสานิ คพฺภํ ธเรติ สตฺตาหํ มูฬฺหคพฺภา, สา ทุกฺขาหิ ติพฺพาหิ ขราหิ กฏุกาหิ เวทนาหิ ผุฏฺฐา ตีหิ วิตกฺเกหิ อธิวาเสติ\csromandash{} “สมฺมาสมฺพุทฺโธ วต โส ภควา, โย อิมสฺส เอวรูปสฺส ทุกฺขสฺส ปหานาย ธมฺมํ เทเสติ.\csromanspacer{} สุปฺปฏิปนฺโน วต ตสฺส ภควโต สาวกสํโฆ, โย อิมสฺส เอวรูปสฺส ทุกฺขสฺส ปหานาย ปฏิปนฺโน.\csromanspacer{} สุสุขํ วต ตํ นิพฺพานํ, ยตฺถิทํ เอวรูปํ ทุกฺขํ น สํวิชฺชตี”ติ.}

\gambhira{อุทานปาฬิ}
\prose{\csromanfolio{94}อถ โข สุปฺปวาสา โกลิยธีตา สามิกํ อามนฺเตสิ “เอหิ ตฺวํ อยฺยปุตฺถ เยน ภควา เตนุปสงฺกม, อุปสงฺกมิตฺวา มม วจเนน ภควโต ปาเท สิรสา วนฺทาหิ, อปฺปาพาธํ อปฺปาตงฺกํ ลหุฏฺฐานํ พลํ ผาสุวิหารํ ปุจฺฉ ‘สุปฺปวาสา ภนฺเต โกลิยธีตา ภควโต ปาเท สิรสา วนฺทติ, อปฺปาพาธํ อปฺปาตงฺกํ ลหุฏฺฐานํ พลํ ผาสุวิหารํ ปุจฺฉตี’ติ.\csromanspacer{} เอวญฺจ วเทหิ ‘สุปฺปวาสา ภนฺเต โกลิยธีตา สตฺต วสฺสานิ คพฺภํ ธาเรติ สตฺตาหํ มูฬฺหคพฺภา, สา ทุกฺขาหิ ติพฺพาหิ ขราหิ กฏุกาหิ เวทนาหิ ผุฏฺฐา ตีหิ วิตกฺเกหิ อธิวาเสติ\csromandash{}สมฺมาสมฺพุทฺโธ วต โส ภควา, โย อิมสฺส เอวรูปสฺส ทุกฺขสฺส ปหานาย ธมฺมํ เทเสติ.\csromanspacer{} สุปฺปฏิปนฺโน วต ตสฺส ภควโต สาวกสํโฆ, โย อิมสฺส เอวรูปสฺส ทุกฺขสฺส ปหานาย ปฏิปนฺโน.\csromanspacer{} สุสุขํ วต ตํ นิพฺพานํ, ยตฺถิทํ เอวรูปํ ทุกฺขํ น สํวิชฺชตี’ติ”.}

\prose{“ปรมนฺ”ติ โข โส โกลิยปุตฺโต สุปฺปวาสาย โกลิยธีตาย ปฏิสฺสุตฺวา เยน ภควา เตนุปสงฺกมิ, อุปสงฺกมิตฺวา ภควนฺตํ อภิวาเทตฺวา เอกมนฺตํ นิสีทิ, เอกมนฺตํ นิสินฺโน โข โกลิยปุตฺโต ภควนฺตํ เอตทโวจ “สุปฺปวาสา ภนฺเต โกลิยธีตา ภควโต ปาเท สิรสา วนฺทติ, อปฺปาพาธํ อปฺปาตงฺกํ ลหุฏฺฐานํ พลํ ผาสุวิหารํ ปุจฺฉติ, เอวญฺจ วเทติ ‘สุปฺปวาสา ภนฺเต โกลิยธีตา สตฺต วสฺสานิ คพฺภํ ธาเรติ สตฺตาหํ มูฬฺหคพฺภา, สา ทุกฺขาหิ ติพฺพาหิ ขราหิ กฏุกาหิ เวทนาหิ ผุฏฺฐา ตีหี วิตกฺเกหิ อธิวาเสติ\csromandash{}สมฺมาสมฺพุทฺโธ วต โส ภควา, โย อิมสฺส เอวรูปสฺส ทุกฺขสฺส ปหานาย ธมฺมํ เทเสติ.\csromanspacer{} สุปฺปฏิปนฺโน วต ตสฺส ภควโต สาวกสํโฆ, โย อิมสฺส เอวรูปสฺส ทุกฺขสฺส ปหานาย ปฏิปนฺโน.\csromanspacer{} สุสุขํ วต นิพฺพานํ, ยตฺถิทํ เอวรูปํ ทุกฺขํ น สํวิชฺชตี’ติ”.}

\prose{“สุขินี โหตุ สุปฺปวาสา โกลิยธีตา อโรคา, อโรคํ ปุตฺตํ วิชายตู”ติ สห วจนา จ ปน ภควโต สุปฺปวาสา โกลิยธีตา สุขินี อโรคา อโรคํ ปุตฺตํ วิชายิ.}

\csromanmatikaanchor{mtk.40}
\csromantocmarkref{subsection}{1. ปฐมโพธิสุตฺต}{mtk.40}
\bookmark[dest={mtk.40},level=2]{1. ปฐมโพธิสุตฺต}
\prose{“เอวํ ภนฺเต”ติ โข โส โกลิยปุตฺโต ภควโต ภาสิตํ อภินนฺทิตฺวา อนุโมทิตฺวา อุฏฺฐายาสนา ภควนฺตํ อภิวาเทตฺวา ปทกฺขิณํ กตฺวา เยน สกํ ฆรํ เตน ปจฺจายาสิ.\csromanspacer{} อทฺทสา โข โส โกลิยปุตฺโต สุปฺปวาสํ โกลิยธีตรํ สุขินิํ อโรคํ,}

\vspace{\tipitakaparskip}
\csromancenter{มุจลินฺทวคฺค อโรคํ ปุตฺตํ วิชาตํ, ทิสฺวานสฺส เอตทโหสิ “อจฺฉริยํ วต โภ, อพฺภุตํ วต โภ, ตถาคตสฺส มหิทฺธิกตา มหานุภาวตา, ยตฺร หิ นามายํ สุปฺปวาสา โกลิยธิตา สห วจนา จ ปน\footnote{สห วจนา ปน (อิ), สห วจนา (?)} ภควโต สุขินี อโรคา อโรคํ ปุตฺตํ วิชายิสฺสตี”ติ อตฺตมโน ปมุทิโต ปีติโสมนสฺสชาโต อโหสิ.}
\prose{\csromanfolio{95}อถ โข สุปฺปวาสา โกลิยธีตา สามิกํ อามนฺเตสิ “เอหิ ตฺวํ อยฺยปุตฺต เยน ภควา เตนุปสงฺกม, อุปสงฺกมิตฺวา มม วจเนน ภควโต ปาเท สิรสา วนฺทาหิ ‘สุปฺปวาสา ภนฺเต โกลิยธีตา ภควโต ปาเท สิรสา วนฺทตี’ติ, เอวญฺจ วเทหิ ‘สุปฺปวาสา ภนฺเต โกลิยธีตา สตฺต วสฺสานิ คพฺภํ ธาเรติ สตฺตาหํ มูฬฺหคพฺภา, สา เอตรหิ สุขินี อโรคา อโรคํ ปุตฺตํ วิชาตา, สา สตฺตาหํ พุทฺธปฺปมุขํ ภิกฺขุสํฆํ ภตฺเตน นิมนฺเตติ ‘อธิวาเสตุ กิร ภนฺเต ภควา สุปฺปวาสาย โกลิยธีตาย สตฺต ภตฺตานิ สทฺธิํ ภิกฺขุสํเฆนา’ติ”.}

\prose{“ปรมนฺ”ติ โข โส โกลิยปุตฺโต สุปฺปวาสาย โกลิยธีตาย ปฏิสฺสุตฺวา เยน ภควา เตนุปสงฺกมิ, อุปสงฺกมิตฺวา ภควนฺตํ อภิวาเทตฺวา เอกมนฺตํ นิสีทิ, เอกมนฺตํ นิสินฺโน โข โส โกลิยปุตฺโต ภควนฺตํ เอตทโวจ\csromandash{}}

\prose{สุปฺปวาสา ภนฺเต โกลิยธีตา ภควโต ปาเท สิรสา วนฺทติ, เอวญฺจ วเทติ “สุปฺปวาสา ภนฺเต โกลิยธีตา สตฺต วสฺสานิ คพฺภํ ธาเรติ สตฺตาหํ มูฬฺหคพฺภา, สา เอตรหิ สุขินี อโรคา อโรคํ ปุตฺตํ วิชาตา, สา สตฺตาหํ พุทฺธปฺปมุธํ ภิกฺขุสํฆํ ภตฺเตน นิมนฺเตติ ‘อธิวาเสตุ กิร ภนฺเต ภควา สุปฺปวาสาย โกลิยธีตาย สตฺต ภตฺตานิ สทฺธิํ ภิกฺขุสํเฆนา’ติ”.}

\prose{เตน โข ปน สมเยน อญฺญตเรน อุปาสเกน พุทฺธปฺปมุโข ภิกฺขุสํโฆ สฺวาตนาย ภตฺเตน นิมนฺติโต โหติ, โส จ อุปาสโก อายสฺมโต มหาโมคฺคลฺลานสฺส\footnote{มหาโมคฺคลานสฺส (ก)} อุปฏฺฐาโก โหติ.\csromanspacer{} อถ โข ภควา อายสฺมนฺตํ มหาโมคฺคลฺลานํ อามนฺเตสิ “เอหิ ตฺวํ โมคฺคลฺลาน เยน โส อุปาสโก เตนุปสงฺกม, อุปสงฺกมิตฺวา ตํ อุปาสกํ เอวํ วเทหิ ‘สุปฺปวาสา อาวุโส โกลิยธีตา}

\csromanmatikaanchor{mtk.41}
\csromantocmarkref{subsection}{2. ทุติยโพธิสุตฺต}{mtk.41}
\bookmark[dest={mtk.41},level=2]{2. ทุติยโพธิสุตฺต}
\vspace{\tipitakaparskip}
\csromancenter{อุทานปาฬิ สตฺต วสฺสานิ คพฺภํ ธาเรสิ สตฺตาหํ มูฬฺหคพฺภา, สา เอตรหิ สุขินี อโรคา อโรคํ ปุตฺตํ วิชาตา, สา สตฺตาหํ พุทฺธปฺปมุขํ ภิกฺขุสํฆํ ภตฺเถน นิมนฺเตติ.\csromanspacer{} กโรตุ สุปฺปวาสา โกลิยธีตา สตฺต ภตฺตานิ, ปจฺฉา ตฺวํ กริสฺสสี’ติ\footnote{กริสฺสสีติ สญฺญาเปหิ (ก)}, ตุเยฺหโส อุปฏฺฐาโก”ติ.}
\prose{\csromanfolio{96}“เอวํ ภนฺเต”ติ โข อายสฺมา มหาโมคฺคลฺลาโน ภควโต ปฏิสฺสุตฺวา เยน โส อุปาสโก เตนุปสงฺกมิ, อุปสงฺกมิตฺวา ตํ อุปาสกํ เอตทโวจ “สุปฺปวาสา อาวุโส โกลิยธีตา สตฺต วสฺสานิ คพฺภํ ธาเรติ สตฺตาหํ มูฬฺหคพฺภา, สา เอตรหิ สุขินี อโรคา อโรคํ ปุตฺตํ วิชาตา, สา สตฺตาหํ พุทฺธปฺปมุขํ ภิกฺขุสํฆํ ภตฺเตน นิมนฺเตติ, กโรตุ สุปฺปวาสา โกลิยธีตา สตฺต ภตฺตานิ, ปจฺฉา ตฺวํ กริสฺสสี”ติ.}

\prose{สเจ เม ภนฺเต อยฺโย มหาโมคฺคลฺลาโน ติณฺณํ ธมฺมานํ ปาฏิโภโค\csromandash{}โภคานญฺจ ชีวิตสฺส จ สทฺธาย จ, กโรตุ สุปฺปวาสา โกลิยธีตา สตฺต ภตฺตานิ, ปจฺฉาหํ กริสฺสามีติ.\csromanspacer{} ทฺวินฺนํ โข เต อหํ\footnote{ทฺวินฺนํ โข เตสํ (อิ), ทฺวินฺนํ โข เนสํ (ก)} อาวุโส ธมฺมานํ ปาฏิโภโค\csromandash{}โภคานญฺจ ชีวิตสฺส จ, สทฺธาย ปน ตฺวํเยว ปาฏิโภโคติ.}

\prose{สเจ เม ภนฺเต อยฺโย มหาโมคฺคลฺลาโน ทฺวินฺนํ ธมฺมานํ ปาฏิโภโค\csromandash{}โภคานญฺจ ชีวิตสฺส จ, กโรตุ สุปฺปวาสา โกลิยธีตา สตฺต ภตฺตานิ, ปจฺฉาหํ กริสฺสามีติ.}

\prose{อถ โข อายสฺมา มหาโมคฺคลฺลาโน ตํ อุปาสกํ สญฺญาเปตฺวา เยน ภควา เตนุปสงฺกมิ, อุปสงฺกมิตฺวา ภควนฺตํ เอตทโวจ “สญฺญตฺโต\footnote{สญฺญาโต (สฺยา)} ภนฺเต โส อุปาสโก มยา, กโรตุ สุปฺปวาสา โกลิยธีตา สตฺต ภตฺตานิ, ปจฺฉา โส กริสฺสตี”ติ.}

\csromanmatikaanchor{mtk.42}
\csromantocmarkref{subsection}{3. ตติยโพธิสุตฺต}{mtk.42}
\bookmark[dest={mtk.42},level=2]{3. ตติยโพธิสุตฺต}
\csromantocmarkref{subsection}{4. หุํหุงฺกสุตฺต}{mtk.43}
\bookmark[dest={mtk.43},level=2]{4. หุํหุงฺกสุตฺต}
\csromantocmarkref{subsection}{5. พฺราหฺมณสุตฺต}{mtk.44}
\bookmark[dest={mtk.44},level=2]{5. พฺราหฺมณสุตฺต}
\prose{อถ โข สุปฺปวาสา โกลิยธีตา สตฺตาหํ พุทฺธปฺปมุขํ ภิกฺขุสํฆํ ปณีเตน ขาทนีเยน โภชนีเยน สหตฺถา สนฺตปฺเปสิ สมฺปวาเรสิ, ตญฺจ ทารกํ ภควนฺตํ วนฺทาเปสิ สพฺพญฺจ ภิกฺขุสํฆํ.}

\chapterheadpageread{2. มุจลินฺทวคฺค}
\prose{\csromanfolio{97}อถ โข อายสฺมา สาริปุตฺโต ตํ ทารกํ เอตทโวจ “กจฺจิ เต ทารก ขมนียํ กจฺจิ ยาปนียํ กจฺจิ น กิญฺจิ ทุกฺขนฺ”ติ.\csromanspacer{} กุโต เม ภนฺเต สาริปุตฺต ขมนียํ กุโต ยาปนียํ, สตฺต เม วสฺสานิ โลหิตกุมฺภิยํ วุตฺถานีติ.}

\prose{อถ โข สุปฺปวาสา โกลิยธีตา “ปุตฺโต เม ธมฺมเสนาปตินา สทฺธิํ มนฺเตตี”ติ อตฺตมนา ปมุทิตา ปีติโสมนสฺสชาตา อโหสิ.\csromanspacer{} อถ โข ภควา (สุปฺปวาสํ โกลิยธีตรํ อตฺตมนํ ปมุทิตํ ปีติโสมนสฺสชาตํ วิทิตฺวา1)\footnote{( ) นตฺถิ อิงฺคลิสโปตฺถเก.} สุปฺปวาสํ โกลิยธีตรํ เอตทโวจ “อิจฺเฉยฺยาสิ ตฺวํ สุปฺปวาเส อญฺญมฺปิ เอวรูปํ ปุตฺตนฺ”ติ.\csromanspacer{} อิจฺเฉยฺยามหํ ภควา อญฺญานิปิ เอวรูปานิ สตฺต ปุตฺตานีติ.}

\prose{อถ โข ภควา เอตมตฺถํ วิทิตฺวา ตายํ เวลายํ อิมํ อุทานํ อุทาเนสิ\csromandash{}}

\csromangathameasure{“อสาตํ สาตรูเปน, ปิยรูเปน อปฺปิยํ. \\ ทุกฺขํ สุขสฺส รูเปน, ปมตฺตมติวตฺตตี”ติ.อฏฺฐมํ.}
\csromangathagroup{\csromangathastanza{“อสาตํ สาตรูเปน, ปิยรูเปน อปฺปิยํ. \\ ทุกฺขํ สุขสฺส รูเปน, ปมตฺตมติวตฺตตี”ติ.\csromanspacer{}\csromanspacer{}อฏฺฐมํ.}}

\csromanheader{2}{9. วิสาขาสุตฺต}
\proseitem{19}{เอวํ เม สุตํ\csromandash{}เอกํ สมยํ ภควา สาวตฺถิยํ วิหรติ ปุพฺพาราเม มิคารมาตุปาสาเท.\csromanspacer{} เตน โข ปน สมเยน วิสาขาย มิคารมาตุยา โกจิเทว อตฺโถ รญฺเญ ปเสนทิมฺหิ โกสเล ปฏิพทฺโธ\footnote{ปฏิพนฺโธ (อิ, ก)} โหติ.\csromanspacer{} ตํ ราชา ปเสนทิ โกสโล น ยถาธิปฺปายํ ตีเรติ.}

\prose{อถ โข วิสาขา มิคารมาตา ทิวา ทิวสฺส\footnote{ทิวาทิวสฺเสว (สฺยา), ทิวาทิวสฺเสเยว (อิ), ทิวา ทิวสฺสเยว (ก)} เยน ภควา เตนุปสงฺกมิ, อุปสงฺกมิตฺวา ภควนฺตํ อภิวาเทตฺวา เอกมนฺตํ นิสีทิ, เอกมนฺตํ นิสินฺนํ โข วิสาขํ มิคารมาตรํ ภควา เอตทโวจ “หนฺท กุโต นุ ตฺวํ วิสาเข อาคจฺฉสิ ทิวา ทิวสฺสา”ติ.\csromanspacer{} อิธ เม ภนฺเต โกจิเทว อตฺโถ รญฺเญ ปเสนทิมฺหิ โกสเล ปฏิพทฺโธ, ตํ ราชา ปเสนทิ โกสโล น ยถาธิปฺปายํ ตีเรตีติ.}

\gambhira{อุทานปาฬิ}
\prose{\csromanfolio{98}อถ โข ภควา เอตมตฺถํ วิทิตฺวา ตายํ เวลายํ อิมํ อุทานํ อุทาเนสิ\csromandash{}}

\csromangathameasure{“สพฺพํ ปรวสํ ทุกฺขํ, สพฺพํ อิสฺสริยํ สุขํ. \\ สาธารเณ วิหญฺญนฺติ, โยคา หิ ทุรติกฺกมา”ติ.นวมํ.}
\csromangathagroup{\csromangathastanza{“สพฺพํ ปรวสํ ทุกฺขํ, สพฺพํ อิสฺสริยํ สุขํ. \\ สาธารเณ วิหญฺญนฺติ, โยคา หิ ทุรติกฺกมา”ติ.\csromanspacer{}\csromanspacer{}นวมํ.}}

\csromanheader{2}{10. ภทฺทิยสุตฺต}
\proseitem{20}{\csromansymbolfootnote{*}{วิ 4. 399 ปิฏฺเฐปิ.}เอวํ เม สุตํ\csromandash{}เอกํ สมยํ ภควา อนุปิยายํ วิหรติ อมฺพวเน.\csromanspacer{} เตน โข ปน สมเยน อายสฺมา ภทฺทิโย กาฬีโคธาย ปุตฺโต อรญฺญคโตปิ รุกฺขมูลคโตปิ สุญฺญาคารคโตปิ อภิกฺขณํ อุทานํ อุทาเนสิ “อโห สุขํ อโห สุขนฺ”ติ.}

\prose{อสฺโสสุํ โข สมฺพหุลา ภิกฺขู อายสฺมโต ภทฺทิยสฺส กาฬีโคธาย ปุตฺตสฺส อรญฺญคตสฺสปิ รุกฺขมูลคตสฺสปิ สุญฺญาคารคตสฺสปิ อภิกฺขณํ อุทานํ อุทาเนนฺตสฺส “อโห สุขํ อโห สุขนฺ”ติ, สุตฺวาน เนสํ เอตทโหสิ “นิสฺสํสยํ โข อาวุโส อายสฺมา ภทฺทิโย กาฬีโคธาย ปุตฺโต อนภิรโต พฺรหฺมจริยํ จรติ, ยํ’ส ปุพฺเพ อคาริยภูตสฺส\footnote{อคาริกภูตสฺส (สฺยา)} รชฺชสุขํ, โส ตมนุสฺสรมาโน อรญฺญคโตปิ รุกฺขมูลคโตปิ สุญฺญาคารคโตปิ อภิกฺขณํ อุทานํ อุทาเนสิ “อโห สุขํ อโห สุขนฺ”ติ.}

\prose{อถ โข สมฺพหุลา ภิกฺขู เยน ภควา เตนุปสงฺกมิํสุ, อุปสงฺกมิตฺวา ภควนฺตํ อภิวาเทตฺวา เอกมนฺตํ นิสีทิํสุ, เอกมนฺตํ นิสินฺนา โข เต ภิกฺขู ภควนฺตํ เอตทโวจุํ “อายสฺมา ภนฺเต ภทฺทิโย กาฬีโคธาย ปุตฺโต อรญฺญคโตปิ รุกฺขมูลคโตปิ สุญฺญาคารคโตปิ อภิกฺขณํ อุทานํ อุทาเนสิ ‘อโห สุขํ อโห สุขนฺ’ติ, นิสฺสํสยํ โข ภนฺเต อายสฺมา ภทฺทิโย กาฬีโคธาย ปุตฺโต อนภิรโต พฺรหฺมจริยํ จรติ, ยํ’ส ปุพฺเพ อคาริยภูตสฺส รชฺชสุขํ, โส ตมนุสฺสรมาโน อรญฺญคโตปิ รุกฺขมูลคโตปิ สุญฺญาคารคโตปิ อภิกฺขณํ อุทานํ อุทาเนสิ ‘อโห สุขํ อโห สุขนฺ’ติ”.}

\chapterheadpageread{2. มุจลินฺทวคฺค}
\prose{\csromanfolio{99}อถ โข ภควา อญฺญตรํ ภิกฺขุํ อามนฺเตสิ “เอหิ ตฺวํ ภิกฺขุ มม วจเนน ภทฺทิยํ ภิกฺขุํ อามนฺเตหิ “สตฺถา ตํ อาวุโส ภทฺทิย อามนฺเตตี’ติ”.}

\csromanmatikaanchor{mtk.45}
\csromantocmarkref{subsection}{6. มหากสฺสปสุตฺต}{mtk.45}
\bookmark[dest={mtk.45},level=2]{6. มหากสฺสปสุตฺต}
\prose{“เอวํ ภนฺเต”ติ โข โส ภิกฺขุ ภควโต ปฏิสฺสุตฺวา เยนายสฺมา ภทฺทิโย กาฬีโคธาย ปุตฺโต เตนุปสงฺกมิ, อุปสงฺกมิตฺวา ภทฺทิยํ กาฬีโคธาย ปุตฺตํ เอตทโวจ “สตฺถา ตํ อาวุโส ภทฺทิย อามนฺเตตี”ติ. “เอวมาวุโส”ติ โข อายสฺมา ภทฺทิโย กาฬีโคธาย ปุตฺโต ตสฺส ภิกฺขุโน ปฏิสฺสุตฺวา เยน ภควา เตนุปสงฺกมิ, อุปสงฺกมิตฺวา ภควนฺตํ อภิวาเทตฺวา เอกมนฺตํ นิสีทิ, เอกมนฺตํ นิสินฺนํ โข อายสฺมนฺตํ ภทฺทิยํ กาฬีโคธาย ปุตฺตํ ภควา เอตทโวจ\csromandash{}}

\prose{สจฺจํ กิร ตฺวํ ภทฺทิย อรญฺญคโตปิ รุกฺขมูลคโตปิ สุญฺญาคารคโตปิ อภิกฺขณํ อุทานํ อุทาเนสิ “อโห สุขํ อโห สุขนฺ”ติ, เอวํ ภนฺเตติ.}

\prose{กิํ ปน\footnote{กํ ปน (สฺยา, อิ)} ตฺวํ ภทฺทิย อตฺถวสํ สมฺปสฺสมาโน อรญฺญคโตปิ รุกฺขมูลคโตปิ สุญฺญาคารคโตปิ อภิกฺขณํ อุทานํ อุทาเนสิ “อโห สุขํ อโห สุขนฺ”ติ.\csromanspacer{} ปุพฺเพ เม ภนฺเต อคาริยภูตสฺส รชฺชํ กาเรนฺตสฺส อนฺโตปิ อนฺเตปุเร รกฺขา สุสํวิหิตา อโหสิ, พหิปิ อนฺเตปุเร รกฺขา สุสํวิหิตา อโหสิ, อนฺโตปิ นคเร รกฺขา สุสํวิหิตา อโหสิ, พหิปิ นคเร รกฺขา สุสํวิหิตา อโหสิ, อนฺโตปิ ชนปเท รกฺขา สุสํวิหิตา อโหสิ, พหิปิ ชนปเท รกฺขา สุสํวิหิตา อโหสิ.\csromanspacer{} โส โข อหํ ภนฺเต เอวํ รกฺขิโต โคปิโต สนฺโต ภีโต อุพฺพิคฺโค อุสฺสงฺกี อุตฺราสี วิหาสิํ, เอตรหิ โข ปนาหํ ภนฺเต อรญฺญคโตปิ รุกฺขมูลคโตปิ สุญฺญาคารคโตปิ เอโก\footnote{เอกโก (สฺยา, อิ)} อภีโต อนุพฺพิคฺโค อนุสฺสงฺกี อนุตฺราสี อปฺโปสฺสุกฺโก ปนฺนโลโม ปรทตฺตวุตฺโต\footnote{ปรทวุตฺโต (ก-สี, สฺยา, อิ)} มิคภูเตน เจตสา วิหรามิ, อิมํ\footnote{อิทํ (สี, ก)} โข อหํ ภนฺเต อตฺถวสํ สมฺปสฺสมาโน อรญฺญคโตปิ รุกฺขมูลคโตปิ สุญฺญาคารคโตปิ อภิกฺขณํ อุทานํ อุทาเนสิํ\footnote{อุทาเนมิ (ก)} “อโห สุขํ อโห สุขนฺ”ติ.}

\prose{อถ โข ภควา เอตมตฺถํ วิทิตฺวา ตายํ เวลายํ อิมํ อุทานํ อุทาเนสิ\csromandash{}}

\gambhira{อุทานปาฬิ}
\csromangathameasure{“ยสฺสนฺตรโต น สนฺติ โกปา, \\ อิติภวาภวตญฺจ วีติวตฺโต. \\ ตํ วิคตภยํ สุขิํ อโสกํ, \\ เทวา นานุภวนฺติ ทสฺสนายา”ติ.ทสมํ. มุจลินฺทวคฺโค ทุติโย.}
\csromangathagroup{\csromangathastanza{\csromanfolio{100}“ยสฺสนฺตรโต น สนฺติ โกปา, \\ อิติภวาภวตญฺจ วีติวตฺโต. \\ ตํ วิคตภยํ สุขิํ อโสกํ, \\ เทวา นานุภวนฺติ ทสฺสนายา”ติ.\csromanspacer{}\csromanspacer{}ทสมํ.\csromanspacer{} มุจลินฺทวคฺโค ทุติโย.}}

\csromanmatikaanchor{mtk.46}
\csromantocmarkref{subsection}{7. อชกลาปกสุตฺต}{mtk.46}
\bookmark[dest={mtk.46},level=2]{7. อชกลาปกสุตฺต}
\titlehead{\textbf{ตสฺสุทฺทานํ}}
\csromangathameasure{มุจลินฺโท ราชา ทณฺเฑน, สกฺกาโร อุปาสเกน จ. \\ คพฺภินี เอกปุตฺโต จ, สุปฺปวาสา วิสาขา จ.}
\csromangathagroup{\csromangathastanza{มุจลินฺโท ราชา ทณฺเฑน, สกฺกาโร อุปาสเกน จ. \\ คพฺภินี เอกปุตฺโต จ, สุปฺปวาสา วิสาขา จ.}}

\csromanheader{2}{กาฬีโคธาย ภทฺทิโยติ.}
\csromansectionrule
\chapterhead{3. นนฺทวคฺค}
\csromanmatikaanchor{mtk.47}
\csromantocmarkref{subsection}{8. สงฺคามชิสุตฺต}{mtk.47}
\bookmark[dest={mtk.47},level=2]{8. สงฺคามชิสุตฺต}
\csromanheader{2}{1. กมฺมวิปากชสุตฺต}
\proseitem{21}{เอวํ เม สุตํ\csromandash{}เอกํ สมยํ ภควา สาวตฺถิยํ วิหรติ เชตวเน อนาถปิณฺฑิกสฺส อาราเม.\csromanspacer{} เตน โข ปน สมเยน อญฺญตโร ภิกฺขุ ภควโต อวิทูเร นิสินฺโน โหติ ปลฺลงฺกํ อาภุชิตฺวา อุชุํ กายํ ปณิธาย ปุราณกมฺมวิปากชํ ทุกฺขํ ติพฺพํ ขรํ กฏุกํ เวทนํ อธิวาเสนฺโต สโต สมฺปชาโน อวิหญฺญมาโน.}

\prose{อทฺทสา โข ภควา ตํ ภิกฺขุํ อวิทูเร นิสินฺนํ ปลฺลงฺกํ อาภุชิตฺวา อุชุํ กายํ ปณิธาย ปุราณกมฺมวิปากชํ ทุกฺขํ ติพฺพํ ขรํ กฏุกํ เวทนํ อธิวาเสนฺตํ สตํ สมฺปชานํ อวิหญฺญมานํ.}

\prose{อถ โข ภควา เอตมตฺถํ วิทิตฺวา ตายํ เวลายํ อิมํ อุทานํ อุทาเนสิ\csromandash{}}

\csromangathameasure{“สพฺพกมฺมชหสฺส ภิกฺขุโน, \\ ธุนมานสฺส ปุเร กตํ รชํ. \\ อมมสฺส ฐิตสฺส ตาทิโน, \\ อตฺโถ นตฺถิ ชนํ ลเปตเว”ติ.ปฐมํ.}
\csromangathagroup{\csromangathastanza{“สพฺพกมฺมชหสฺส ภิกฺขุโน, \\ ธุนมานสฺส ปุเร กตํ รชํ. \\ อมมสฺส ฐิตสฺส ตาทิโน, \\ อตฺโถ นตฺถิ ชนํ ลเปตเว”ติ.\csromanspacer{}\csromanspacer{}ปฐมํ.}}

\csromanheader{2}{3. นนฺทวคฺค \textbf{2. นนฺทสุตฺต}}
\proseitem{22}{\csromanfolio{101}เอวํ เม สุตํ\csromandash{}เอกํ สมยํ ภควา สาวตฺถิยํ วิหรติ เชตวเน อนาถปิณฺฑิกสฺส อาราเม.\csromanspacer{} เตน โข ปน สมเยน อายสฺมา นนฺโท ภควโต ภาตา มาตุจฺฉาปุตฺโต สมฺพหุลานํ ภิกฺขูนํ เอวมาโรเจติ “อนภิรโต อหํ อาวุโส พฺรหฺมจริยํ จรามิ, น สกฺโกมิ พฺรหฺมจริยํ สนฺธาเรตุํ, สิกฺขํ ปจฺจกฺขาย หีนายาวตฺติสฺสามี”ติ.}

\prose{อถ โข อญฺญตโร ภิกฺขุ เยน ภควา เตนุปสงฺกมิ, อุปสงฺกมิตฺวา ภควนฺตํ อภิวาเทตฺวา เอกมนฺตํ นิสีทิ, เอกมนฺตํ นิสินฺโน โข โส ภิกฺขุ ภควนฺตํ เอตทโวจ “อายสฺมา ภนฺเต นนฺโท ภควโต ภาตา มาตุจฺฉาปุตฺโต สมฺพหุลานํ ภิกฺขูนํ เอวมาโรเจติ ‘อนภิรโต อหํ อาวุโส พฺรหฺมจริยํ จรามิ, น สกฺโกมิ พฺรหฺมจริยํ สนฺธาเรตุํ, สิกฺขํ ปจฺจกฺขาย หีนายาวตฺติสฺสามี’ติ”.}

\prose{อถ โข ภควา อญฺญตรํ ภิกฺขุํ อามนฺเต สิ “เอหิ ตฺวํ ภิกฺขุ มม วจเนน นนฺทํ ภิกฺขุํ อามนฺเตหิ ‘สตฺถา ตํ อาวุโส นนฺท อามนฺเตตี’ติ”. “เอวํ ภนฺเต”ติ โข โส ภิกฺขุ ภควโต ปฏิสฺสุตฺวา เยนายสฺมา นนฺโท เตนุปสงฺกมิ, อุปสงฺกมิตฺวา อายสฺมนฺตํ นนฺทํ เอตทโวจ “สตฺถา ตํ อาวุโส นนฺท อามนฺเตตี”ติ.}

\prose{“เอวมาวุโส”ติ โข อายสฺมา นนฺโท ตสฺส ภิกฺขุโน ปฏิสฺสุตฺวา เยน ภควา เตนุปสงฺกมิ, อุปสงฺกมิตฺวา ภควนฺตํ อภิวาเทตฺวา เอกมนฺตํ นิสีทิ, เอกมนฺตํ นิสินฺนํ โข อายสฺมนฺตํ นนฺทํ ภควา เอตทโวจ\csromandash{}}

\prose{สจฺจํ กิร ตฺวํ นนฺท สมฺพหุลานํ ภิกฺขูนํ เอวมาโรเจสิ “อนภิรโต อหํ อาวุโส พฺรหฺมจริยํ จรามิ, น สกฺโกมิ พฺรหฺมจริยํ สนฺธาเรตุํ, สิกฺขํ ปจฺจกฺขาย หีนายาวตฺติสฺสามี”ติ.\csromanspacer{} เอวํ ภนฺเตติ.}

\csromanmatikaanchor{mtk.48}
\csromantocmarkref{subsection}{9. ชติลสุตฺต}{mtk.48}
\bookmark[dest={mtk.48},level=2]{9. ชติลสุตฺต}
\prose{กิสฺส ปน ตฺวํ นนฺท อนภิรโต พฺรหฺมจริยํ จรสิ, น สกฺโกสิ พฺรหฺมจริยํ สนฺธาเรตุํ, สิกฺขํ ปจฺจกฺขาย หีนายาวตฺติสฺสสีติ.\csromanspacer{} สากิยานี มํ\footnote{มม (สฺยา, อฏฺฐกถา โอโลเกตพฺพา)} ภนฺเต ชนปทกลฺยาณี ฆรา นิกฺขมนฺตสฺส\footnote{นิกฺขมนฺตํ (อฏฺฐกถายํ ปาฐนฺตรํ)} อุปฑฺฒุลฺลิขิเตหิ เกเสหิ อปโลเกตฺวา มํ เอตทโวจ “ตุวฏํ โข อยฺยปุตฺต อาคจฺเฉยฺยาสี”ติ.}

\vspace{\tipitakaparskip}
\csromancenter{อุทานปาฬิ โส โข อหํ ภนฺเต ตมนุสฺสรมาโน อนภิรโต พฺรหฺมจริยํ จรามิ, น สกฺโกมิ พฺรหฺมจริยํ สนฺธาเรตุํ, สิกฺขํ ปจฺจกฺขาย หีนายาวตฺติสฺสามีติ.}
\prose{\csromanfolio{102}อถ โข ภควา อายสฺมนฺตํ นนฺทํ พาหายํ คเหตฺวา เสยฺยถาปิ นาม พลวา ปุริโส สมิญฺชิตํ\footnote{สมฺมิญฺชิตํ (สี, สฺยา, กํ, อิ)} วา พาหํ ปสาเรยฺย, ปสาริตํ วา พาหํ สมิญฺเชยฺย\footnote{สมฺมิญฺเชยฺย (สี, สฺยา, กํ, อิ)}.\csromanspacer{} เอวเมวํ เชตวเน อนฺตรหิโต เทเวสุ ตาวติํเสสุ ปาตุรโหสิ.}

\prose{เตน โข ปน สมเยน ปญฺจมตฺตานิ อจฺฉราสตานิ สกฺกสฺส เทวานมินฺทสฺส อุปฏฺฐานํ อาคตานิ โหนฺติ กกุฏปาทานิ.\csromanspacer{} อถ โข ภควา อายสฺมนฺตํ นนฺทํ อามนฺเตสิ “ปสฺสสิ โน ตฺวํ นนฺท อิมานิ ปญฺจ อจฺฉราสตานิ กกุฏปาทานี”ติ.\csromanspacer{} เอวํ ภนฺเตติ.}

\prose{ตํ กิํ มญฺญสิ นนฺท, กตมา นุ โข อภิรูปตรา วา ทสฺสนียตรา วา ปาสาทิกตรา วา สากิยานี วา ชนปทกลฺยาณี, อิมานิ วา ปญฺจ อจฺฉราสตานิ กกุฏปาทานีติ.\csromanspacer{} เสยฺยถาปิ ภนฺเต ปลุฏฺฐมกฺกฏี กณฺณนาสจฺฉินฺนา.\csromanspacer{} เอวเมวํ โข ภนฺเต สากิยานี ชนปทกลฺยาณี อิเมสํ ปญฺจนฺนํ อจฺฉราสตานํ อุปนิธาย สงฺขฺยมฺปิ\footnote{สงฺขมฺปิ (สี)} โนเปติ, กลภาคมฺปิ โนเปติ, อุปนิธิมฺปิ โนเปติ.\csromanspacer{} อถ โข อิมานิ ปญฺจ อจฺฉราสตานิ อภิรูปตรานิ เจว ทสฺสนียตรานิ จ ปาสาทิกตรานิ จาติ.}

\prose{อภิรม นนฺท อภิรม นนฺท, อหํ เต ปาฏิโภโค ปญฺจนฺนํ อจฺฉราสตานํ ปฏิลาภาย กกุฏปาทานนฺติ.\csromanspacer{} สเจ เม ภนฺเต ภควา ปาฏิโภโค ปญฺจนฺนํ อจฺฉราสตานํ ปฏิลาภาย กกุฏปาทานํ, อภิรมิสฺสามหํ ภนฺเต ภควติ พฺรหฺมจริเยติ\footnote{ภควา พฺรหฺมจริเยติ (สฺยา, อิ), ภควา พฺรหฺมจริยนฺติ (ก)}.}

\csromanmatikaanchor{mtk.49}
\csromantocmarkref{subsection}{10. พาหิยสุตฺต}{mtk.49}
\bookmark[dest={mtk.49},level=2]{10. พาหิยสุตฺต}
\prose{อถ โข ภควา อายสฺมนฺตํ นนฺทํ พาหายํ คเหตฺวา เสยฺยถาปิ นาม พลวา ปุริโส สมิญฺชิตํ วา พาหํ ปสาเรยฺย, ปสาริตํ วา พาหํ สมิญฺเชยฺย.\csromanspacer{} เอวเมวํ เทเวสุ ตาวติํเสสุ อนฺตรหิโต เชตวเน ปาตุรโหสิ.}

\chapterheadpageread{3. นนฺทวคฺค}
\prose{\csromanfolio{103}อสฺโสสุํ โข ภิกฺขู “อายสฺมา กิร นนฺโท ภควโต ภาตา มาตุจฺฉาปุตฺโต อจฺฉรานํ เหตุ พฺรหฺมจริยํ จรติ, ภควา กิรสฺส ปาฏิโภโค ปญฺจนฺนํ อจฺฉราสตานํ ปฏิลาภาย กกุฏปาทานนฺ”ติ.}

\prose{อถ โข อายสฺมโต นนฺทสฺส สหายกา ภิกฺขู อายสฺมนฺตํ นนฺทํ ภตกวาเทน จ อุปกฺกิตกวาเทน จ สมุทาจรนฺติ “ภตโก กิรายสฺมา นนฺโท, อุปกฺกิตโก กิรายสฺมา นนฺโท อจฺฉรานํ เหตุ พฺรหฺมจริยํ จรติ, ภควา กิรสฺส ปาฏิโภโค ปญฺจนฺนํ อจฺฉราสตานํ ปฏิลาภาย กกุฏปาทานนฺ”ติ.}

\prose{อถ โข อายสฺมา นนฺโท สหายกานํ ภิกฺขูนํ ภตกวาเทน จ อุปกฺกิตกวาเทน จ อฏฺฏียมาโน หรายมาโน ชิคุจฺฉมาโน เอโก วูปกฏฺโฐ อปฺปมตฺโต อาตาปี ปหิตตฺโต วิหรนฺโต นจิรสฺเสว ยสฺสตฺถาย กุลปุตฺตา สมฺมเทว อคารสฺมา อนคาริยํ ปพฺพชนฺติ, ตทนุตฺตรํ พฺรหฺมจริยปริโยสานํ ทิฏฺเฐว ธมฺเม สยํ อภิญฺญา สจฺฉิกตฺวา อุปสมฺปชฺช วิหาสิ, “ขีณา ชาติ, วุสิตํ พฺรหฺมจริยํ, กตํ กรณียํ, นาปรํ อิตฺถตฺตายา”ติ อพฺภญฺญาสิ.\csromanspacer{} อญฺญตโร โข ปนายสฺมา นนฺโท อรหตํ อโหสิ.}

\prose{อถ โข อญฺญตรา เทวตา อภิกฺกนฺตาย รตฺติยา อภิกฺกนฺตวณฺณา เกวลกปฺปํ เชตวนํ โอภาเสตฺวา เยน ภควา เตนุปสงฺกมิ, อุปสงฺกมิตฺวา ภควนฺตํ อภิวาเทตฺวา เอกมนฺตํ อฏฺฐาสิ, เอกมนฺตํ ฐิตา โข สา เทวตา ภควนฺตํ เอตทโวจ “อายสฺมา ภนฺเต นนฺโท ภควโต ภาตา มาตุจฺฉาปุตฺโต อาสวานํ ขยา อนาสวํ เจโตวิมุตฺติํ ปญฺญาวิมุตฺติํ ทิฏฺเฐว ธมฺเม สยํ อภิญฺญา สจฺฉิกตฺวา อุปสมฺปชฺช วิหรตี”ติ.\csromanspacer{} ภควโตปิ โข ญาณํ อุทปาทิ “นนฺโท อาสวานํ ขยา อนาสวํ เจโตวิมุตฺติํ ปญฺญาวิมุตฺติํ ทิฏฺเฐว ธมฺเม สยํ อภิญฺญา สจฺฉิกตฺวา อุปสมฺปชฺช วิหรตี”ติ.}

\prose{อถ โข อายสฺมา นนฺโท ตสฺสา รตฺติยา อจฺจเยน เยน ภควา เตนุปสงฺกมิ, อุปสงฺกมิตฺวา ภควนฺตํ อภิวาเทตฺวา เอกมนฺตํ นิสีทิ, เอกมนฺตํ นิสินฺโน โข อายสฺมา นนฺโท ภควนฺตํ เอตทโวจ “ยํ เม ภนฺเต ภควา ปาฏิโภโค ปญฺจนฺนํ อจฺฉราสตานํ ปฏิลาภาย กกุฏปาทานํ, มุญฺจามหํ ภนฺเต ภควนฺตํ เอตสฺมา ปฏิสฺสวา”ติ.\csromanspacer{} มยาปิ}

\vspace{\tipitakaparskip}
\csromancenter{อุทานปาฬิ โข ตฺวํ นนฺท\footnote{โข เต นนฺท (สี, สฺยา, อิ), โข นนฺท (ก)} เจตสา เจโต ปริจฺจ วิทิโต “นนฺโท อาสวานํ ขยา อนาสวํ เจโตวิมุตฺติํ ปญฺญาวิมุตฺติํ ทิฏฺเฐว ธมฺเม สยํ อภิญฺญา สจฺฉิกตฺวา อุปสมฺปชฺช วิหรตี”ติ.\csromanspacer{} เทวตาปิ เม เอตมตฺถํ อาโรเจสิ “อายสฺมา ภนฺเต นนฺโท ภควโต ภาตา มาตุจฺฉาปุตฺโต อาสวานํ ขยา อนาสวํ เจโตวิมุตฺติํ ปญฺญาวิมุตฺติํ ทิฏฺเฐว ธมฺเม สยํ อภิญฺญา สจฺฉิกตฺวา อุปสมฺปชฺช วิหรตี”ติ.\csromanspacer{} ยเทว โข เต นนฺท อนุปาทาย อาสเวหิ จิตฺตํ วิมุตฺตํ, อถาหํ มุตฺโต เอตสฺมา ปฏิสฺสวาติ.}
\prose{\csromanfolio{104}อถ โข ภควา เอตมตฺถํ วิทิตฺวา ตายํ เวลายํ อิมํ อุทานํ อุทาเนสิ\csromandash{}}

\csromangathameasure{“ยสฺส นิตฺติณฺโณ ปงฺโก, \\ มทฺทิโต กามกณฺฑโก. \\ โมหกฺขยํ อนุปฺปตฺโต, \\ สุขทุกฺเขสุ น เวธตี ส ภิกฺขู”ติ.ทุติยํ.}
\csromangathagroup{\csromangathastanza{“ยสฺส นิตฺติณฺโณ ปงฺโก, \\ มทฺทิโต กามกณฺฑโก. \\ โมหกฺขยํ อนุปฺปตฺโต, \\ สุขทุกฺเขสุ น เวธตี ส ภิกฺขู”ติ.\csromanspacer{}\csromanspacer{}ทุติยํ.}}

\csromanheader{2}{3. ยโสชสุตฺต}
\proseitem{23}{เอวํ เม สุตํ\csromandash{}เอกํ สมยํ ภควา สาวตฺถิยํ วิหรติ เชตวเน อนาถปิณฺฑิกสฺส อาราเม.\csromanspacer{} เตน โข ปน สมเยน ยโสชปฺปมุขานิ ปญฺจมตฺตานิ ภิกฺขุสตานิ สาวตฺถิํ อนุปฺปตฺตานิ โหนฺติ ภควนฺตํ ทสฺสนาย, เต’ธ โข อาคนฺตุกา ภิกฺขู เนวาสิเกหิ ภิกฺขูหิ สทฺธิํ ปฏิสมฺโมทมานา เสนาสนานิ ปญฺญาปยมานา ปตฺตจีวรานิ ปฏิสามยมานา อุจฺจาสทฺทา มหาสทฺทา\footnote{อุจฺจาสทฺทมหาสทฺทา (ก)} อเหสุํ.}

\prose{อถ โข ภควา อายสฺมนฺตํ อานนฺทํ อามนฺเตสิ “เก ปเนเต อานนฺท อุจฺจาสทฺทา มหาสทฺทา เกวฏฺฏา มญฺเญ มจฺฉวิโลเป”ติ.\csromanspacer{} เอตานิ ภนฺเต ยโสชปฺปมุขานิ ปญฺจมตฺตานิ ภิกฺขุสตานิ สาวตฺถิํ อนุปฺปตฺตานิ ภควนฺตํ ทสฺสนาย, เต’เต อาคนฺตุกา ภิกฺขู เนวาสิเกหิ ภิกฺขูหิ สทฺธิํ ปฏิสมฺโมทมานา เสนาสนานิ ปญฺญาปยมานา ปตฺตจีวรานิ ปฏิสามยมานา อุจฺจาสทฺทา มหาสทฺทาติ.\csromanspacer{} เตนหานนฺท มม วจเนน เต ภิกฺขู อามนฺเตหิ “สตฺถา อายสฺมนฺเต อามนฺเตตี”ติ.}

\chapterheadpageread{3. นนฺทวคฺค}
\prose{\csromanfolio{105}“เอวํ ภนฺเต”ติ โข อายสฺมา อานนฺโท ภควโต ปฏิสฺสุตฺวา เยน เต ภิกฺขู เตนุปสงฺกมิ, อุปสงฺกมิตฺวา เต ภิกฺขู เอตทโวจ “สตฺถา อายสฺมนฺเต อามนฺเตตี”ติ. “เอวมาวุโส”ติ โข เต ภิกฺขู อายสฺมโต อานนฺทสฺส ปฏิสฺสุตฺวา เยน ภควา เตนุปสงฺกมิํสุ, อุปสงฺกมิตฺวา ภควนฺตํ อภิวาเทตฺวา เอกมนฺตํ นิสีทิํสุ, เอกมนฺตํ นิสินฺเน โข เต ภิกฺขู ภควา เอตทโวจ\csromandash{}}

\prose{กิํ นุ ตุเมฺห ภิกฺขเว อุจฺจาสทฺทา มหาสทฺทา เกวฏฺฏา มญฺเญ มจฺฉวิโลเปติ.\csromanspacer{} เอวํ วุตฺเต อายสฺมา ยโสโช ภควนฺตํ เอตทโวจ “อิมานิ ภนฺเต ปญฺจมตฺตานิ ภิกฺขุสตานิ สาวตฺถิํ อนุปฺปตฺตานิ ภควนฺตํ ทสฺสนาย, เต’เม อาคนฺตุกา ภิกฺขู เนวาสิเกหิ ภิกฺขูหิ สทฺธิํ ปฏิสมฺโมทมานา เสนาสนานิ ปญฺญาปยมานา ปตฺตจีวรานิ ปฏิสามยมานา อุจฺจาสทฺทา มหาสทฺทา”ติ.\csromanspacer{} คจฺฉถ ภิกฺขเว, ปณาเมมิ โว\footnote{โว ปณาเมมิ (สพฺพตฺถ) ม 2. 120 ปิฏฺเฐ ปสฺสิตพฺพํ.} “น โว มม สนฺติเก วตฺถพฺพนฺ”ติ.}

\prose{“เอวํ ภนฺเต”ติ โข เต ภิกฺขู ภควโต ปฏิสฺสุตฺวา อุฏฺฐายาสนา ภควนฺตํ อภิวาเทตฺวา ปทกฺขิณํ กตฺวา เสนาสนํ สํสาเมตฺวา\footnote{ปฏิสํสาเมตฺวา (สฺยา)} ปตฺตจีวรมาทาย เยน วชฺชี เตน จาริกํ ปกฺกมิํสุ, วชฺชีสุ อนุปุพฺเพน จาริกํ จรมานา เยน วคฺคุมุทา นที เตนุปสงฺกมิํสุ, อุปสงฺกมิตฺวา วคฺคุมุทาย นทิยา ตีเร ปณฺณกุฏิโย กริตฺวา วสฺสํ อุปคจฺฉิํสุ.}

\csromanmatikaanchor{mtk.50}
\csromantocmarknopage{chapter}{2. มุจลินฺทวคฺค}{mtk.50}
\bookmark[dest={mtk.50},level=0]{2. มุจลินฺทวคฺค}
\prose{อถ โข อายสฺมา ยโสโช วสฺสูปคโต\footnote{วสฺสูปคเต (ก)} ภิกฺขู อามนฺเตสิ “ภควตา มยํ อาวุโส ปณามิตา อตฺถกาเมน หิเตสินา อนุกมฺปเกน อนุกมฺปํ อุปาทาย, หนฺท มยํ อาวุโส ตถา วิหารํ กปฺเปม, ยถา โน วิหรตํ ภควา อตฺตมโน อสฺสา”ติ. “เอวมาวุโส”ติ โข เต ภิกฺขู อายสฺมโต ยโสชสฺส ปจฺจสฺโสสุํ.\csromanspacer{} อถ โข เต ภิกฺขู วูปกฏฺฐา อปฺปมตฺตา อาตาปิโน ปหิตตฺตา วิหรนฺตา เตเนวนฺตรวสฺเสน สพฺเพว ติสฺโส วิชฺชา สจฺฉากํ สุ.}

\prose{อถ โข ภควา สาวตฺถิยํ ยถาภิรนฺตํ วิหริตฺวา เยน เวสาลี เตน จาริกํ ปกฺกามิ, อนุปุพฺเพน จาริกํ จรมาโน เยน เวสาลี ตทวสริ, ตตฺร สุทํ ภควา เวสาลิยํ วิหรติ มหาวเน กูฏาคารสาลายํ.}

\gambhira{อุทานปาฬิ}
\prose{\csromanfolio{106}อถ โข ภควา วคฺคุมุทาตีริยานํ ภิกฺขูนํ เจตสา เจโต ปริจฺจ มนสิ กริตฺวา อายสฺมนฺตํ อานนฺทํ อามนฺเตสิ “อาโลกชาตา วิย เม อานนฺท เอสา ทิสา, โอภาสชาตา วิย เม อานนฺท เอสา ทิสา, ยสฺสํ ทิสายํ\footnote{ยายํ (ก)} วคฺคุมุทาตีริยา ภิกฺขู วิหรนฺติ, คนฺตุํ อปฺปฏิกูลาสิ เม มนสิ กาตุํ, ปหิเณยฺยาสิ ตฺวํ อานนฺท วคฺคุมุทาตีริยานํ ภิกฺขูนํ สนฺติเก ทูตํ, ‘สตฺถา อายสฺมนฺเต อามนฺเตติ, สตฺถา อายสฺมนฺตานํ ทสฺสนกาโม’ติ.}

\prose{“เอวํ ภนฺเต”ติ โข อายสฺมา อานนฺโท ภควโต ปฏิสฺสุตฺวา เยน อญฺญตโร ภิกฺขุ เตนุปสงฺกมิ, อุปสงฺกมิตฺวา ตํ ภิกฺขุํ เอตทโวจ “เอหิ ตฺวํ อาวุโส เยน วคฺคุมุทาตีริยา ภิกฺขู เตนุปสงฺกม, อุปสงฺกมิตฺวา วคฺคุมุทาตีริเย ภิกฺขู เอวํ วเทหิ ‘สตฺถา อายสฺมนฺเต อามนฺเตติ, สตฺถา อายสฺมนฺตานํ ทสฺสนกาโม’ติ”.}

\prose{“เอวมาวุโส”ติ โข โส ภิกฺขุ อายสฺมโต อานนฺทสฺส ปฏิสฺสุตฺวา เสยฺยถาปิ นาม พลวา ปุริโส สมิญฺชิตํ วา พาหํ ปสาเรยฺย, ปสาริตํ วา พาหํ สมิญฺเชยฺย, เอวเมวํ มหาวเน กูฏาคารสาลายํ อนฺตรหิโต วคฺคุมุทาย นทิยา ตีเร เตสํ ภิกฺขูนํ ปุรโต ปาตุรโหสิ.\csromanspacer{} อถ โข โส ภิกฺขุ วคฺคุมุทาตีริเย ภิกฺขู เอตทโวจ “สตฺถา อายสฺมนฺเต อามนฺเตติ, สตฺถา อายสฺมนฺตานํ ทสฺสนกาโม”ติ.}

\csromanmatikaanchor{mtk.51}
\csromantocmarkref{subsection}{1. มุจลินฺทสุตฺต}{mtk.51}
\bookmark[dest={mtk.51},level=2]{1. มุจลินฺทสุตฺต}
\csromantocmarkref{subsection}{2. ราชสุตฺต}{mtk.52}
\bookmark[dest={mtk.52},level=2]{2. ราชสุตฺต}
\csromantocmarkref{subsection}{3. ทณฺฑสุตฺต}{mtk.53}
\bookmark[dest={mtk.53},level=2]{3. ทณฺฑสุตฺต}
\prose{“เอวมาวุโส”ติ โข เต ภิกฺขู ตสฺส ภิกฺขุโน ปฏิสฺสุตฺวา เสนาสนํ สํสาเมตฺวา ปตฺตจีวรมาทาย เสยฺยถาปิ นาม พลวา ปุริโส สมิญฺชิตํ วา พาหํ ปสาเรยฺย, ปสาริตํ วา พาหํ สมิญฺเชยฺย, เอวเมวํ วคฺคุมุทาย นทิยา ตีเร อนฺตรหิตา มหาวเน กูฏาคารสาลายํ ภควโต สมฺมุเข ปาตุรเหสุํ.\csromanspacer{} เตน โข ปน สมเยน ภควา อาเนญฺเชน สมาธินา นิสินฺโน โหติ.\csromanspacer{} อถ โข เตสํ ภิกฺขูนํ เอตทโหสิ “กตเมน นุ โข ภควา วิหาเรน เอตรหิ วิหรตี”ติ.\csromanspacer{} อถ โข เตสํ ภิกฺขูนํ เอตทโหสิ “อาเนญฺเชน โข ภควา วิหาเรน เอตรหิ วิหรตี”ติ สพฺเพว อาเนญฺชสมาธินา นิสีทิํสุ.}

\chapterheadpageread{3. นนฺทวคฺค}
\prose{\csromanfolio{107}อถ โข อายสฺมา อานนฺโท อภิกฺกนฺตาย รตฺติยา นิกฺขนฺเต ปฐเม ยาเม อุฏฺฐายาสนา เอกํสํ อุตฺตราสงฺคํ\footnote{จีวรํ (สพฺพตฺถ) วิ 4. 418; อํ 3. 44 ปิฏฺเฐสุปิ.} กริตฺวา เยน ภควา เตนญฺชลิํ ปณาเมตฺวา ภควนฺตํ เอตทโวจ “อภิกฺกนฺตา ภนฺเต รตฺติ, นิกฺขนฺโต ปฐโม ยาโม, จิรนิสินฺนา อาคนฺตุกา ภิกฺขู, ปฏิสมฺโมทตุ ภนฺเต ภควา อาคนฺตุเกหิ ภิกฺขูหี”ติ.\csromanspacer{} เอวํ วุตฺเต ภควา ตุณฺหี อโหสิ.}

\prose{ทุติยมฺปิ โข อายสฺมา อานนฺโท อภิกฺกนฺตาย รตฺติยา นิกฺขนฺเต มชฺฌิเม ยาเม อุฏฺฐายาสนา เอกํสํ อุตฺตราสงฺคํ กริตฺวา เยน ภควา เตนญฺชลิํ ปณาเมตฺวา ภควนฺตํ เอตทโวจ “อภิกฺกนฺตา ภนฺเต รตฺติ, นิกฺขนฺโต มชฺฌิโม ยาโม, จิรนิสินฺนา อาคนฺตุกา ภิกฺขู, ปฏิสมฺโมทตุ ภนฺเต ภควา อาคนฺตุเกหิ ภิกฺขูหี”ติ.\csromanspacer{} ทุติยมฺปิ โข ภควา ตุณฺหี อโหสิ.}

\prose{ตติยมฺปิ โข อายสฺมา อานนฺโท อภิกฺกนฺตาย รตฺติยา นิกฺขนฺเต ปจฺฉิเม ยาเม อุทฺธสฺเต อรุเณ นนฺทิมุขิยา รตฺติยา อุฏฺฐายาสนา เอกํสํ อุตฺตราสงฺคํ กริตฺวา เยน ภควา เตนญฺชลิํ ปณาเมตฺวา ภควนฺตํ เอตทโวจ “อภิกฺกนฺตา ภนฺเต รตฺติ, นิกฺขนฺโต ปจฺฉิโม ยาโม, อุทฺธสฺโต อรุโณ, นนฺทิมุขี รตฺติ, จิรนิสินฺนา อาคนฺตุกา ภิกฺขู, ปฏิสมฺโมทตุ ภนฺเต ภควา อาคนฺตุเกหิ ภิกฺขูหี”ติ.}

\prose{อถ โข ภควา ตมฺหา สมาธิมฺหา วุฏฺฐหิตฺวา อายสฺมนฺตํ อานนฺทํ อามนฺเตสิ “สเจ โข ตฺวํ อานนฺท ชาเนยฺยาสิ, เอตฺตกมฺปิ เต นปฺปฏิภาเยยฺย\footnote{นปฺปฏิเภยฺย (?) อํ 3. 417 ปิฏฺเฐปิ.}, อหญฺจ อานนฺท อิมานิ จ ปญฺจ ภิกฺขุสตานิ สพฺเพว อาเนญฺชสมาธินา นิสีทิมฺหา”ติ.}

\prose{อถ โข ภควา เอตมตฺถํ วิทิตฺวา ตายํ เวลายํ อิมํ อุทานํ อุทาเนสิ\csromandash{}}

\csromangathameasure{“ยสฺส ชิโต กามกณฺฑโก, \\ อกฺโกโส จ วโธ จ พนฺธนญฺจ. \\ ปพฺพโตว โส ฐิโต อเนโช, \\ สุขทุกฺเขสุ น เวธตี ส ภิกฺขู”ติ.ตติยํ.}
\csromangathagroup{\csromangathastanza{“ยสฺส ชิโต กามกณฺฑโก, \\ อกฺโกโส จ วโธ จ พนฺธนญฺจ. \\ ปพฺพโตว\footnote{ปพฺพโต วิย (สี, สฺยา, อิ)} โส ฐิโต อเนโช, \\ สุขทุกฺเขสุ น เวธตี ส ภิกฺขู”ติ.\csromanspacer{}\csromanspacer{}ตติยํ.}}

\csromanheader{2}{อุทานปาฬิ \textbf{4. สาริปุตฺตสุตฺต}}
\csromanmatikaanchor{mtk.54}
\csromanmatikaanchor{mtk.66}
\csromantocmarkref{subsection}{4. สกฺการสุตฺต}{mtk.54}
\bookmark[dest={mtk.54},level=2]{4. สกฺการสุตฺต}
\proseitem{24}{\csromanfolio{108}เอวํ เม สุตํ\csromandash{}เอกํ สมยํ ภควา สาวตฺถิยํ วิหรติ เชตวเน อนาถปิณฺฑิกสฺส อาราเม.\csromanspacer{} เตน โข ปน สมเยน อายสฺมา สาริปุตฺโต ภควโต อวิทูเร นิสินฺโน โหติ ปลฺลงฺกํ อาภุชิตฺวา อุชุํ กายํ ปณิธาย ปริมุขํ สติํ อุปฏฺฐเปตฺวา.\csromanspacer{} อทฺทสา โข ภควา อายสฺมนฺตํ สาริปุตฺตํ อวิทูเร นิสินฺนํ ปลฺลงฺกํ อาภุชิตฺวา อุชุํ กายํ ปณิธาย ปริมุขํ สติํ อุปฏฺฐเปตฺวา.}

\prose{อถ โข ภควา เอตมตฺถํ วิทิตฺวา ตายํ เวลายํ อิมํ อุทานํ อุทาเนสิ\csromandash{}}

\csromangathameasure{*“ยถาปิ ปพฺพโต เสโล, อจโล สุปฺปติฏฺฐิโต. \\ เอวํ โมหกฺขยา ภิกฺขุ, ปพฺพโตว น เวธตี”ติ.จตุตฺถํ.}
\csromangathagroup{\csromangathastanza{\csromansymbolfootnote{*}{ขุ 2. 345 ปิฏฺเฐปิ.}“ยถาปิ ปพฺพโต เสโล, อจโล สุปฺปติฏฺฐิโต. \\ เอวํ โมหกฺขยา ภิกฺขุ, ปพฺพโตว น เวธตี”ติ.\csromanspacer{}\csromanspacer{}จตุตฺถํ.}}

\csromanheader{2}{5. มหาโมคฺคลฺลานสุตฺต}
\proseitem{25}{เอวํ เม สุตํ\csromandash{}เอกํ สมยํ ภควา สาวตฺถิยํ วิหรติ เชตวเน อนาถปิณฺฑิกสฺส อาราเม.\csromanspacer{} เตน โข ปน สมเยน อายสฺมา มหาโมคฺคลฺลาโน ภควโต อวิทูเร นิสินฺโน โหติ ปลฺลงฺกํ อาภุชิตฺวา อุชุํ กายํ ปณิธาย กายคตาย สติยา อชฺฌตฺตํ สูปฏฺฐิตาย.\csromanspacer{} อทฺทสา โข ภควา อายสฺมนฺตํ มหาโมคฺคลฺลานํ อวิทูเร นิสินฺนํ ปลฺลงฺกํ อาภุชิตฺวา อุชุํ กายํ ปณิธาย กายคตาย สติยา อชฺฌตฺตํ สูปฏฺฐิตาย.}

\prose{อถ โข ภควา เอตมตฺถํ วิทิตฺวา ตายํ เวลายํ อิมํ อุทานํ อุทาเนสิ\csromandash{}}

\prose{\csromansymbolfootnote{+}{ขุ 10. 179 ปิฏฺเฐปิ.}“สติ กายคตา อุปฏฺฐิตา,}

\prose{ฉสุ ผสฺสายตเนสุ สํวุโต.}

\prose{สตตํ ภิกฺขุ สมาหิโต,}

\prose{ชญฺญา นิพฺพานมตฺตโน”ติ.\csromanspacer{}\csromanspacer{}ปญฺจมํ.}

\csromanheader{2}{3. นนฺทวคฺค \textbf{6. ปิลินฺทวจฺฉสุตฺต}}
\proseitem{26}{\csromanfolio{109}เอวํ เม สุตํ\csromandash{}เอกํ สมยํ ภควา ราชคเห วิหรติ เวฬุวเน กลนฺทกนิวาเป.\csromanspacer{} เตน โข ปน สมเยน อายสฺมา ปิลินฺทวจฺโฉ\footnote{ปิลินฺทิวจฺโฉ (สี)} ภิกฺขู วสลวาเทน สมุทาจรติ.\csromanspacer{} อถ โข สมฺพหุลา ภิกฺขู เยน ภควา เตนุปสงฺกมิํสุ, อุปสงฺกมิตฺวา ภควนฺตํ อภิวาเทตฺวา เอกมนฺตํ นิสีทิํสุ, เอกมนฺตํ นิสินฺนา โข เต ภิกฺขู ภควนฺตํ เอตทโวจุํ “อายสฺมา ภนฺเต ปิลินฺทวจฺโฉ ภิกฺขู วสลวาเทน สมุทาจรตี”ติ.}

\prose{อถ โข ภควา อญฺญตรํ ภิกฺขุํ อามนฺเตสิ “เอหิ ตฺวํ ภิกฺขุ มม วจเนน ปิลินฺทวจฺฉํ ภิกฺขุํ อามนฺเตหิ ‘สตฺถา ตํ อาวุโส ปิลินฺทวจฺฉ\footnote{วจฺฉ (สฺยา)} อามนฺเตตี’ติ”. “เอวํ ภนฺเต”ติ โข โส ภิกฺขุ ภควโต ปฏิสฺสุตฺวา เยนายสฺมา ปิลินฺทวจฺโฉ เตนุปสงฺกมิ, อุปสงฺกมิตฺวา อายสฺมนฺตํ ปิลินฺทวจฺฉํ เอตทโวจ “สตฺถา ตํ อาวุโส ปิลินฺทวจฺฉ อามนฺเตตี”ติ.}

\prose{“เอวมาวุโส”ติ โข อายสฺมา ปิลินฺทวจฺโฉ ตสฺส ภิกฺขุโน ปฏิสฺสุตฺวา เยน ภควา เตนุปสงฺกมิ, อุปสงฺกมิตฺวา ภควนฺตํ อภิวาเทตฺวา เอกมนฺตํ นิสีทิ, เอกมนฺตํ นิสินฺนํ โข อายสฺมนฺตํ ปิลินฺทวจฺฉํ ภควา เอตทโวจ “สจฺจํ กิร ตฺวํ วจฺฉ ภิกฺขู วสลวาเทน สมุทาจรสี”ติ.\csromanspacer{} เอวํ ภนฺเตติ.}

\prose{อถ โข ภควา อายสฺมโต ปิลินฺทวจฺฉสฺส ปุพฺเพนิวาสํ มนสิ กริตฺวา ภิกฺขู อามนฺเตสิ\csromandash{}มา โข ตุเมฺห ภิกฺขเว วจฺฉสฺส ภิกฺขุโน อุชฺฌายิตฺถ, น ภิกฺขเว วจฺโฉ โทสนฺตโร ภิกฺขู วสลวาเทน สมุทาจรติ, วจฺฉสฺส ภิกฺขเว ภิกฺขุโน ปญฺจ ชาติสตานิ อพฺโพกิณฺณานิ พฺราหฺมณกุเล ปจฺจาชาตานิ, โส ตสฺส วสลวาโท ทีฆรตฺตํ สมุทาจิณฺโณ\footnote{อชฺฌาจิณฺโณ (สฺยา, อิ, ก, อฏฺฐกถายํ ปาฐนฺตรํ)}, เตนายํ วจฺโฉ ภิกฺขู วสลวาเทน สมุทาจรตีติ.}

\prose{อถ โข ภควา เอตมตฺถํ วิทิตฺวา ตายํ เวลายํ อิมํ อุทานํ อุทาเนสิ\csromandash{}}

\gambhira{อุทานปาฬิ}
\csromangathameasure{“ยมฺหี น มายา วสตี นมาโน, \\ โย วีตโลโภ อมโม นิราโส. \\ ปณุนฺนโกโธ อภินิพฺพุตตฺโต, \\ โส พฺราหฺมโณ โส สมโณ ส ภิกฺขู”ติ.ฉฏฺฐํ.}
\csromangathagroup{\csromangathastanza{\csromanfolio{110}“ยมฺหี น มายา วสตี นมาโน, \\ โย วีตโลโภ อมโม นิราโส. \\ ปณุนฺนโกโธ\footnote{ปนุณฺณโกโธ (ก) ขุ 10. 175 ปิฏฺเฐปิ.} อภินิพฺพุตตฺโต, \\ โส พฺราหฺมโณ โส สมโณ ส ภิกฺขู”ติ.\csromanspacer{}\csromanspacer{}ฉฏฺฐํ.}}

\csromanheader{2}{7. สกฺกุทานสุตฺต}
\csromanmatikaanchor{mtk.55}
\csromantocmarkref{subsection}{5. อุปาสกสุตฺต}{mtk.55}
\bookmark[dest={mtk.55},level=2]{5. อุปาสกสุตฺต}
\proseitem{27}{เอวํ เม สุตํ\csromandash{}เอกํ สมยํ ภควา ราชคเห วิหรติ เวฬุวเน กลนฺทกนิวาเป.\csromanspacer{} เตน โข ปน สมเยน อายสฺมา มหากสฺสโป ปิปฺปลิคุหายํ วิหรติ, สตฺตาหํ เอกปลฺลงฺเกน นิสินฺโน โหติ อญฺญตรํ\footnote{นิสินฺโน อญฺญตรํ (สฺยา, ก)} สมาธิํ สมาปชฺชิตฺวา.\csromanspacer{} อถ โข อายสฺมา มหากสฺสโป ตสฺส สตฺตาหสฺส อจฺจเยน ตมฺหา สมาธิมฺหา วุฏฺฐาสิ.\csromanspacer{} อถ โข อายสฺมโต มหากสฺสปสฺส ตมฺหา สมาธิมฺหา วุฏฺฐิตสฺส เอตทโหสิ “ยํนูนาหํ ราชคหํ ปิณฺฑาย ปวิเสยฺยนฺ”ติ.}

\prose{เตน โข ปน สมเยน ปญฺจมตฺตานิ เทวตาสตานิ อุสฺสุกฺกํ อาปนฺนานิ โหนฺติ อายสฺมโต มหากสฺสปสฺส ปิณฺฑปาตปฏิลาภาย.\csromanspacer{} อถ โข อายสฺมา มหากสฺสโป ตานิ ปญฺจมตฺตานิ เทวตาสตานิ ปฏิกฺขิปิตฺวา ปุพฺพณฺหสมยํ นิวาเสตฺวา ปตฺตจีวรมาทาย ราชคหํ ปิณฺฑาย ปาวิสิ.}

\prose{เตน โข ปน สมเยน สกฺโก เทวานมินฺโท อายสฺมโต มหากสฺสปสฺส ปิณฺฑปาตํ ทาตุกาโม โหติ.\csromanspacer{} เปสการวณฺณํ อภินิมฺมินิตฺวา ตนฺตํ วินาติ สุชา\footnote{สุชาตา (สฺยา, อิ, ก)} อสุรกญฺญา ตสรํ ปูเรติ.\csromanspacer{} อถ โข อายสฺมา มหากสฺสโป ราชคเห สปทานํ ปิณฺฑาย จรมาโน เยน สกฺกสฺส เทวานมินฺทสฺส นิเวสนํ เตนุปสงฺกมิ.\csromanspacer{} อทฺทสา โข สกฺโก เทวานมินฺโท อายสฺมนฺตํ มหากสฺสปํ ทูรโตว อาคจฺฉนฺตํ, ทิสฺวาน ฆรา นิกฺขมิตฺวา ปจฺจุคฺคนฺตฺวา หตฺถโต ปตฺตํ คเหตฺวา ฆรํ ปวิสิตฺวา\footnote{ปวิเสตฺวา (ก)} ฆฏิยา โอทนํ อุทฺธริตฺวา ปตฺตํ ปูเรตฺวา อายสฺมโต มหากสฺสปสฺส อทาสิ.\csromanspacer{} โส อโหสิ ปิณฺฑปาโต อเนกสูโป อเนกพฺยญฺชโน อเนกรสพฺยญฺชโน\footnote{อเนกสูปรสพฺยญฺชโน (สี, อิ)}.\csromanspacer{} อถ โข อายสฺมโต มหากสฺสปสฺส เอตทโหสิ “โก นุ โข อยํ สตฺโต, ยสฺสายํ เอวรูโป}

\vspace{\tipitakaparskip}
\csromancenter{นนฺทวคฺค อิทฺธานุภาโว”ติ.\csromanspacer{} อถ โข อายสฺมโต มหากสฺสปสฺส เอตทโหสิ “สกฺโก โข อยํ เทวานมินฺโท”ติ อิติ วิทิตฺวา สกฺกํ เทวานมินฺทํ เอตทโวจ “กตํ โข เต อิทํ โกสิย, มา\footnote{มาสฺสุ (สี, สฺยา)} ปุนปิ เอวรูปมกาสี”ติ.\csromanspacer{} อมฺหากมฺปิ ภนฺเต กสฺสป ปุญฺเญน อตฺโถ, อมฺหากมฺปิ ปุญฺเญน กรณียนฺติ.}
\prose{\csromanfolio{111}อถ โข สกฺโก เทวานมินฺโท อายสฺมนฺตํ มหากสฺสปํ อภิวาเทตฺวา ปทกฺขิณํ กตฺวา เวหาสํ อพฺภุคฺคนฺตฺวา อากาเส อนฺตลิกฺเข ติกฺขตฺตุํ อุทานํ อุทาเนสิ “อโห ทานํ ปรมทานํ\footnote{ปรมํ ทานํ (อิ, ก)} กสฺสเป สุปฺปติฏฺฐิตํ, อโห ทานํ ปรมทานํ กสฺสเป สุปฺปติฏฺฐิตํ, อโห ทานํ ปรมทานํ กสฺสเป สุปฺปติฏฺฐิตนฺ”ติ.\csromanspacer{} อสฺโสสิ โข ภควา ทิพฺพาย โสตธาตุยา วิสุทฺธาย อติกฺกนฺตมานุสิกาย สกฺกสฺส เทวนมินฺทสฺส เวหาสํ อพฺภุคฺคนฺตฺวา อากาเส อนฺตลิกฺเข ติกฺขตฺตุํ อุทานํ อุทาเนนฺตสฺส “อโห ทานํ ปรมทานํ กสฺสเป สุปฺปติฏฺฐิตํ, อโห ทานํ ปรมทานํ กสฺสเป สุปฺปติฏฺฐิตํ, อโห ทานํ ปรมทานํ กสฺสเป สุปฺปติฏฺฐิตนฺ”ติ.}

\prose{อถ โข ภควา เอตมตฺถํ วิทิตฺวา ตายํ เวลายํ อิมํ อุทานํ อุทาเนสิ\csromandash{}}

\csromanmatikaanchor{mtk.56}
\csromantocmarkref{subsection}{6. คพฺภินีสุตฺต}{mtk.56}
\bookmark[dest={mtk.56},level=2]{6. คพฺภินีสุตฺต}
\csromangathameasure{“ปิณฺฑปาติกสฺส ภิกฺขุโน, \\ อตฺตภรสฺส อนญฺญโปสิโน. \\ เทวา ปิหยนฺติ ตาทิโน, \\ อุปสนฺตสฺส สทา สตีมโต”ติ.สตฺตมํ.}
\csromangathagroup{\csromangathastanza{“ปิณฺฑปาติกสฺส ภิกฺขุโน, \\ อตฺตภรสฺส อนญฺญโปสิโน. \\ เทวา ปิหยนฺติ ตาทิโน, \\ อุปสนฺตสฺส สทา สตีมโต”ติ.\csromanspacer{}\csromanspacer{}สตฺตมํ.}}

\csromanheader{2}{8. ปิณฺฑปาติกสุตฺต}
\proseitem{28}{เอวํ เม สุตํ\csromandash{}เอกํ สมยํ ภควา สาวตฺถิยํ วิหรติ เชตวเน อนาถปิณฺฑิกสฺส อาราเม.\csromanspacer{} เตน โข ปน สมเยน สมฺพหุลานํ ภิกฺขูนํ ปจฺฉาภตฺตํ ปิณฺฑปาตปฏิกฺกนฺตานํ กเรริมณฺฑลมาเฬ สนฺนิสินฺนานํ สนฺนิปติตานํ อยมนฺตรากถา อุทปาทิ\csromandash{}}

\prose{ปิณฺฑปาติโก อาวุโส ภิกฺขุ ปิณฺฑาย จรนฺโต ลภติ กาเลน กาลํ มนาปิเก จกฺขุนา รูเป ปสฺสิตุํ, ลภติ กาเลน กาลํ มนาปิเก โสเตน สทฺเท โสตุํ, ลภติ กาเลน กาลํ มนาปิเก}

\vspace{\tipitakaparskip}
\csromancenter{อุทานปาฬิ ฆาเนน คนฺเธ ฆายิตุํ, ลภติ กาเลน กาลํ มนาปิเก ชิวฺหาย รเส สายิตุํ, ลภติ กาเลน กาลํ มนาปิเก กาเยน โผฏฺฐพฺเพ ผุสิตุํ, ปิณฺฑปาติโก อาวุโส ภิกฺขุ สกฺกโต ครุกโต มานิโต ปูชิโต อปจิโต ปิณฺฑาย จรติ.\csromanspacer{} หนฺทาวุโส มยมฺปิ ปิณฺฑปาติกา โหม, มยมฺปิ ลจฺฉาม กาเลน กาลํ มนาปิเก จกฺขุนา รูเป ปสฺสิตุํ, มยมฺปิ ลจฺฉาม กาเลน กาลํ มนาปิเก โสเตน สทฺเท โสตุํ, มยมฺปิ ลจฺฉาม กาเลน กาลํ มนาปิเก ฆาเนน คนฺเธ ฆายิตุํ, มยมฺปิ ลจฺฉาม กาเลน กาลํ มนาปิเก ชิวฺหาย รเส สายิตุํ, มยมฺปิ ลจฺฉาม กาเลน กาลํ มนาปิเก กาเยน โผฏฺฐพฺเพ ผุสิตุํ, มยมฺปิ สกฺกตา ครุกตา มานิตา ปูชิตา อปจิตา ปิณฺฑาย จริสฺสามาติ.\csromanspacer{} อยญฺจรหิ เตสํ ภิกฺขูนํ อนฺตรากถา โหติ วิปฺปกตา.}
\prose{\csromanfolio{112}อถ โข ภควา สายนฺหสมยํ ปฏิสลฺลานา วุฏฺฐิโต เยน กเรริมณฺฑลมาโฬ เตนุปสงฺกมิ, อุปสงฺกมิตฺวา ปญฺญตฺเต อาสเน นิสีทิ, นิสชฺช โข ภควา ภิกฺขู อามนฺเตสิ “กาย นุตฺถ ภิกฺขเว เอตรหิ กถาย สนฺนิสินฺนา, กา จ ปน โว อนฺตรากถา วิปฺปกตา”ติ.}

\prose{อิธ ภนฺเต อมฺหากํ ปจฺฉาภตฺตํ ปิณฺฑปาตปฏิกฺกนฺตานํ กเรริมณฺฑลมาเฬ สนฺนิสินฺนานํ สนฺนิปติตานํ อยมนฺตรากถา อุทปาทิ\csromandash{}}

\prose{ปิณฺฑปาติโก อาวุโส ภิกฺขุ ปิณฺฑาย จรนฺโต ลภติ กาเลน กาลํ มนาปิเก จกฺขุนา รูเป ปสฺสิตุํ, ลภติ กาเลน กาลํ มนาปิเก โสเตน สทฺเท โสตุํ, ลภติ กาเลน กาลํ มนาปิเก ฆาเนน คนฺเธ ฆายิตุํ, ลภติ กาเลน กาลํ มนาปิเก ชิวฺหาย รเส สายิตุํ, ลภติ กาเลน กาลํ มนาปิเก กาเยน โผฏฺฐพฺเพ ผุสิตุํ, ปิณฺฑปาติโก อาวุโส ภิกฺขุ สกฺกโต ครุกโต มานิโต ปูชิโต อปจิโต ปิณฺฑาย จรติ.\csromanspacer{} หนฺทาวุโส มยมฺปิ ปิณฺฑปาติกา โหม, มยมฺปิ ลจฺฉาม กาเลน กาลํ มนาปิเก จกฺขุนา รูเป ปสฺสิตุํ ฯเปฯ กาเยน โผฏฺฐพฺเพ ผุสิตุํ, มยมฺปิ สกฺกตา ครุกตา มานิตา ปูชิตา อปจิตา ปิณฺฑาย จริสฺสามาติ.\csromanspacer{} อยํ โข โน ภนฺเต อนฺตรากถา วิปฺปกตา, อถ ภควา อนุปฺปตฺโตติ.}

\chapterheadpageread{3. นนฺทวคฺค}
\prose{\csromanfolio{113}น เขฺวตํ ภิกฺขเว ตุมฺหากํ ปติรูปํ กุลปุตฺตานํ สทฺธา อคารสฺมา อนคาริยํ ปพฺพชิตานํ, ยํ ตุเมฺห เอวรูปิํ กถํ กเถยฺยาถ, สนฺนิปติตานํ โว ภิกฺขเว ทฺวยํ กรณียํ ธมฺมี วา กถา อริโย วา ตุณฺหีภาโวติ.}

\csromanmatikaanchor{mtk.57}
\csromantocmarkref{subsection}{7. เอกปุตฺตกสุตฺต}{mtk.57}
\bookmark[dest={mtk.57},level=2]{7. เอกปุตฺตกสุตฺต}
\prose{อถ โข ภควา เอตมตฺถํ วิทิตฺวา ตายํ เวลายํ อิมํ อุทานํ อุทาเนสิ\csromandash{}}

\csromangathameasure{“ปิณฺฑปาติกสฺส ภิกฺขุโน, \\ อตฺตภรสฺส อนญฺญโปสิโน. \\ เทวา ปิหยนฺติ ตาทิโน, \\ โน เจ สทฺทสิโลกนิสฺสิโต”ติ.อฏฺฐมํ.}
\csromangathagroup{\csromangathastanza{“ปิณฺฑปาติกสฺส ภิกฺขุโน, \\ อตฺตภรสฺส อนญฺญโปสิโน. \\ เทวา ปิหยนฺติ ตาทิโน, \\ โน เจ สทฺทสิโลกนิสฺสิโต”ติ.\csromanspacer{}\csromanspacer{}อฏฺฐมํ.}}

\csromanheader{2}{9. สิปฺปสุตฺต}
\proseitem{29}{เอวํ เม สุตํ\csromandash{}เอกํ สมยํ ภควา สาวตฺถิยํ วิหรติ เชตวเน อนาถปิณฺฑิกสฺส อาราเม.\csromanspacer{} เตน โข ปน สมเยน สมฺพหุลานํ ภิกฺขูนํ ปจฺฉาภตฺตํ ปิณฺฑปาตปฏิกฺกนฺตานํ มณฺฑลมาเฬ สนฺนิสินฺนานํ สนฺนิปติตนํ อยมนฺตรากถา อุทปาทิ “โก นุ โข อาวุโส สิปฺปํ ชานาติ, โก กิํ สิปฺปํ สิกฺขิ, กตรํ สิปฺปํ สิปฺปานํ อคฺคนฺ”ติ.}

\prose{ตตฺเถกจฺเจ เอวมาหํสุ “หตฺถิสิปฺปํ สิปฺปานํ อคฺคนฺ”ติ, เอกจฺเจ เอวมาหํสุ “อสฺสสิปฺปํ สิปฺปานํ อคฺคนฺ”ติ, เอกจฺเจ เอวมาหํสุ “รถสิปฺปํ สิปฺปานํ อคฺคนฺ”ติ, เอกจฺเจ เอวมาหํสุ “ธนุสิปฺปํ สิปฺปานํ อคฺคนฺ”ติ, เอกจฺเจ เอวมาหํสุ “ถรุสิปฺปํ สิปฺปานํ อคฺคนฺ”ติ, เอกจฺเจ เอวมาหํสุ “มุทฺทาสิปฺปํ สิปฺปานํ อคฺคนฺ”ติ, เอกจฺเจ เอวมาหํสุ “คณนาสิปฺปํ สิปฺปานํ อคฺคนฺ”ติ, เอกจฺเจ เอวมาหํสุ “สงฺขานสิปฺปํ สิปฺปานํ อคฺคนฺ”ติ, เอกจฺเจ เอวมาหํสุ “เลขาสิปฺปํ สิปฺปานํ อคฺคนฺ”ติ, เอกจฺเจ เอวมาหํสุ “กาเวยฺยสิปฺปํ\footnote{กาพฺยสิปฺปํ (สฺยา)} สิปฺปานํ อคฺคนฺ”ติ, เอกจฺเจ เอวมาหํสุ “โลกายตสิปฺปํ สิปฺปานํ อคฺคนฺ”ติ, เอกจฺเจ เอวมาหํสุ “ขตฺตวิชฺชาสิปฺปํ สิปฺปานํ อคฺคนฺ”ติ.\csromanspacer{} อยญฺจรหิ เตสํ ภิกฺขูนํ อนฺตรากถา โหติ วิปฺปกตา.}

\prose{อถ โข ภควา สายนฺหสมยํ ปฏิสลฺลานา วุฏฺฐิโต เยน มณฺฑลมาโฬ เตนุปสงฺกมิ, อุปสงฺกมิตฺวา ปญฺญตฺเต อาสเน นิสีทิ,}

\vspace{\tipitakaparskip}
\csromancenter{อุทานปาฬิ นิสชฺช โข ภควา ภิกฺขู อามนฺเตสิ “กาย นุตฺถ ภิกฺขเว เอตรหิ กถาย สนฺนิสินฺนา, กา จ ปน โว อนฺตรากถา วิปฺปกตา”ติ.}
\prose{\csromanfolio{114}อิธ ภนฺเต อมฺหากํ ปจฺฉาภตฺตํ ปิณฺฑปาตปฏิกฺกนฺตานํ มณฺฑลมาเฬ สนฺนิสินฺนานํ อยมนฺตรากถา อุทปาทิ “โก นุ โข อาวุโส สิปฺปํ ชานาติ, โก กิํ สิปฺปํ สิกฺขิ, กตรํ สิปฺปํ สิปฺปานํ อคฺคนฺ”ติ.}

\csromanmatikaanchor{mtk.58}
\csromantocmarkref{subsection}{8. สุปฺปวาสาสุตฺต}{mtk.58}
\bookmark[dest={mtk.58},level=2]{8. สุปฺปวาสาสุตฺต}
\prose{ตตฺเถกจฺเจ เอวมาหํสุ “หตฺถิสิปฺปํ สิปฺปานํ อคฺคนฺ”ติ, เอกจฺเจ เอวมาหํสุ “อสฺสสิปฺปํ สิปฺปานํ อคฺคนฺ”ติ, เอกจฺเจ เอวมาหํสุ “รถสิปฺปํ สิปฺปานํ อคฺคนฺ”ติ, เอกจฺเจ เอวมาหํสุ “ธนุสิปฺปํ สิปฺปานํ อคฺคนฺ”ติ, เอกจฺเจ เอวมาหํสุ “ถรุสิปฺปํ สิปฺปานํ อคฺคนฺ”ติ, เอกจฺเจ เอวมาหํสุ “มุทฺทาสิปฺปํ สิปฺปานํ อคฺคนฺ”ติ, เอกจฺเจ เอวมาหํสุ “คณนาสิปฺปํ สิปฺปานํ อคฺคนฺ”ติ, เอกจฺเจ เอวมาหํสุ “สงฺขานสิปฺปํ สิปฺปานํ อคฺคนฺ”ติ, เอกจฺเจ เอวมาหํสุ “เลขาสิปฺปํ สิปฺปานํ อคฺคนฺ”ติ, เอกจฺเจ เอวมาหํสุ “กาเวยฺยสิปฺปํ สิปฺปานํ อคฺคนฺ”ติ, เอกจฺเจ เอวมาหํสุ “โลกายตสิปฺปํ สิปฺปานํ อคฺคนฺ”ติ, เอกจฺเจ เอวมาหํสุ “ขตฺตวิชฺชาสิปฺปํ สิปฺปานํ อคฺคนฺ”ติ.\csromanspacer{} อยํ โข โน ภนฺเต อนฺตรากถา โหติ วิปฺปกตา, อถ ภควา อนุปฺปตฺโตติ.}

\prose{น เขฺวตํ ภิกฺขเว ตุมฺหากํ ปติรูปํ กุลปุตฺตานํ สทฺธา อคารสฺมา อนคาริยํ ปพฺพชิตานํ, ยํ ตุเมฺห เอวรูปิํ กถํ กเถยฺยาถ, สนฺนิปติตานํ โว ภิกฺขเว ทฺวยํ กรณียํ ธมฺมี วา กถา อริโย วา ตุณฺหีภาโวติ.}

\prose{อถ โข ภควา เอตมตฺถํ วิทิตฺวา ตายํ เวลายํ อิมํ อุทานํ อุทาเนสิ\csromandash{}}

\csromangathameasure{“อสิปฺปชีวี ลหุ อตฺถกาโม, \\ ยตินฺทฺริโย สพฺพธิ วิปฺปมุตฺโต. \\ อโนกสารี อมโม นิราโส, \\ หิตฺวา มานํ เอกจโร ส ภิกฺขู”ติ.นวมํ.}
\csromangathagroup{\csromangathastanza{“อสิปฺปชีวี ลหุ อตฺถกาโม, \\ ยตินฺทฺริโย สพฺพธิ วิปฺปมุตฺโต. \\ อโนกสารี อมโม นิราโส, \\ หิตฺวา มานํ เอกจโร ส ภิกฺขู”ติ.\csromanspacer{}\csromanspacer{}นวมํ.}}

\csromanheader{2}{10. โลกสุตฺต}
\proseitem{30}{เอวํ เม สุตํ\csromandash{}เอกํ สมยํ ภควา อุรุเวลายํ วิหรติ นชฺชา เนรญฺชราย ตีเร โพธิรุกฺขมูเล ปฐมาภิสมฺพุทฺโธ.\csromanspacer{} เตน โข ปน}

\vspace{\tipitakaparskip}
\csromancenter{นนฺทวคฺค สมเยน ภควา สตฺตาหํ เอกปลฺลงฺเกน นิสินฺโน โหติ วิมุตฺติสุขปฏิสํเวที.}
\prose{\csromanfolio{115}อถ โข ภควา ตสฺส สตฺตาหสฺส อจฺจเยน ตมฺหา สมาธิมฺหา วุฏฺฐหิตฺวา พุทฺธจกฺขุนา โลกํ โวโลเกสิ.\csromanspacer{} อทฺทสา โข ภควา พุทฺธจกฺขุนา โวโลเกนฺโต สตฺเต อเนเกหิ สนฺตาเปหิ สนฺตปฺปมาเน อเนเกหิ จ ปริฬาเหหิ ปริฑยฺหมาเน ราคเชหิปิ โทสเชหิปิ โมหเชหิปิ\footnote{โมหเชหิปีติ (สพฺพตฺถ)}.}

\prose{อถ โข ภควา เอตมตฺถํ วิทิตฺวา ตายํ เวลายํ อิมํ อุทานํ อุทาเนสิ\csromandash{}}

\csromangathameasure{*“อยํ โลโก สนฺตาปชาโต, \\ ผสฺสปเรโต โรคํ วทติ อตฺตโต. \\ เยน เยน หิ มญฺญติ, \\ ตโต ตํ โหติ อญฺญถา. \\ อญฺญถาภาวี ภวสตฺโต โลโก, \\ ภวปเรโต ภวเมวาภินนฺทติ. \\ ยทภินนฺทติ ตํ ภยํ, \\ ยสฺส ภายติ ตํ ทุกฺขํ. \\ ภววิปฺปหานาย โข ปนิทํ พฺรหฺมจริยํ วุสฺสติ.}
\csromangathagroup{\csromangathastanza{\csromansymbolfootnote{*}{ขุ 10. 135, 186 ปิฏฺเฐสุปิ.}“อยํ โลโก สนฺตาปชาโต, \\ ผสฺสปเรโต โรคํ วทติ อตฺตโต. \\ เยน เยน หิ มญฺญติ\footnote{เยน หิ มญฺญติ (สฺยา, อิ)}, \\ ตโต ตํ โหติ อญฺญถา.}\csromangathastanza{อญฺญถาภาวี ภวสตฺโต โลโก, \\ ภวปเรโต ภวเมวาภินนฺทติ. \\ ยทภินนฺทติ ตํ ภยํ, \\ ยสฺส ภายติ ตํ ทุกฺขํ.}\csromangathastanza{ภววิปฺปหานาย โข ปนิทํ พฺรหฺมจริยํ วุสฺสติ.}}

\prose{เย หิ เกจิ สมณา วา พฺราหฺมณา วา ภเวน ภวสฺส วิปฺปโมกฺขมาหํสุ, สพฺเพ เต ‘อวิปฺปมุตฺตา ภวสฺมา’ติ วทามิ.\csromanspacer{} เย วา ปน เกจิ สมณา วา พฺราหฺมณา วา วิภเวน ภวสฺส นิสฺสรณมาหํสุ, สพฺเพ เต อนิสฺสฏา ภวสฺมา’ติ วทามิ.}

\prose{อุปธิํ หิ ปฏิจฺจ ทุกฺขมิทํ สมฺโภติ, สพฺพุปาทานกฺขยา นตฺถิ ทุกฺขสฺส สมฺภโว, โลกมิมํ ปสฺส, ปุถู อวิชฺชาย ปเรตา ภูตา ภูตรตา อปริมุตฺตา, เย หิ เกจิ ภวา สพฺพธิ สพฺพตฺถตาย, สพฺเพ เต ภวา อนิจฺจา ทุกฺขา วิปริณามธมฺมาติ.}

\gambhira{อุทานปาฬิ}
\csromangathameasure{เอวเมตํ ยถาภูตํ, สมฺมปฺปญฺญาย ปสฺสโต. \\ ภวตณฺหา ปหียติ, วิภวํ นาภินนฺทติ. \\ สพฺพโส ตณฺหานํ ขยา, \\ อเสสวิราคนิโรโธ นิพฺพานํ. \\ ตสฺส นิพฺพุตสฺส ภิกฺขุโน, \\ อนุปาทา ปุนพฺภโว น โหติ. \\ อภิภูโต มาโร วิชิตสงฺคาโม, \\ อุปจฺจคา สพฺพภวานิ ตาที”ติ.ทสมํ. นนฺทวคฺโค ตติโย.}
\csromangathagroup{\csromangathastanza{\csromanfolio{116}เอวเมตํ ยถาภูตํ, สมฺมปฺปญฺญาย ปสฺสโต. \\ ภวตณฺหา ปหียติ, วิภวํ นาภินนฺทติ.}\csromangathastanza{สพฺพโส ตณฺหานํ ขยา, \\ อเสสวิราคนิโรโธ นิพฺพานํ. \\ ตสฺส นิพฺพุตสฺส ภิกฺขุโน, \\ อนุปาทา\footnote{อนุปาทานา (สี)} ปุนพฺภโว น โหติ. \\ อภิภูโต มาโร วิชิตสงฺคาโม, \\ อุปจฺจคา สพฺพภวานิ ตาที”ติ.\csromanspacer{}\csromanspacer{}ทสมํ.\csromanspacer{} นนฺทวคฺโค ตติโย.}}

\titlehead{\textbf{ตสฺสุทฺทานํ}}
\csromangathameasure{กมฺมํ นนฺโท ยโสโช จ, สาริปุตฺโต จ โกลิโต. \\ ปิลินฺโท กสฺสโป ปิณฺโฑ, สิปฺปํ โลเกน เต ทสาติ.}
\csromangathagroup{\csromangathastanza{กมฺมํ นนฺโท ยโสโช จ, สาริปุตฺโต จ โกลิโต. \\ ปิลินฺโท\footnote{ปิลินฺทิ (สี)} กสฺสโป ปิณฺโฑ, สิปฺปํ โลเกน เต ทสาติ.}}

\chapterhead{4. เมฆิยวคฺค}
\csromanheader{2}{1. เมฆิยสุตฺต}
\proseitem{31}{เอวํ เม สุตํ\csromandash{}เอกํ สมยํ ภควา จาลิกายํ วิหรติ จาลิเก ปพฺพเต.\csromanspacer{} เตน โข ปน สมเยน อายสฺมา เมฆิโย ภควโต อุปฏฺฐาโก โหติ.\csromanspacer{} อถ โข อายสฺมา เมฆิโย เยน ภควา เตนุปสงฺกมิ, อุปสงฺกมิตฺวา ภควนฺตํ อภิวาเทตฺวา เอกมนฺตํ อฏฺฐาสิ, เอกมนฺตํ ฐิโต โข อายสฺมา เมฆิโย ภควนฺตํ เอตทโวจ “อิจฺฉามหํ ภนฺเต ชนฺตุคามํ ปิณฺฑาย ปวิสิตุนฺ”ติ.\csromanspacer{} ยสฺสทานิ ตฺวํ เมฆิย กาลํ มญฺญสีติ.}

\prose{อถ โข อายสฺมา เมฆิโย ปุพฺพณฺหสมยํ นิวาเสตฺวา ปตฺตจีวรมาทาย ชนฺตุคามํ ปิณฺฑาย ปาวิสิ, ชนฺตุคาเม ปิณฺฑาย จริตฺวา ปจฺฉาภตฺตํ ปิณฺฑปาตปฏิกฺกนฺโต เยน กิมิกาฬาย นทิยา ตีรํ เตนุปสงฺกมิ,}

\csromanmatikaanchor{mtk.59}
\csromantocmarkref{subsection}{9. วิสาขาสุตฺต}{mtk.59}
\bookmark[dest={mtk.59},level=2]{9. วิสาขาสุตฺต}
\vspace{\tipitakaparskip}
\csromancenter{เมฆิยวคฺค 1อทฺทสา โข อายสฺมา เมฆิโย1 กิมิกาฬาย นทิยา ตีเร ชงฺฆาวิหารํ\footnote{ชงฺฆวิหารํ (ก)} อนุจงฺกมมาโน อนุวิจรมาโน\footnote{อนุวิจรมาโน อทฺทสา โข (สี, สฺยา, อิ), อนุวิจรมาโน อทฺทส (ก)} อมฺพวนํ ปาสาทิกํ มนุญฺญํ\footnote{อิทํ ปทํ วิเทสโปตฺถเกสุ นตฺถิ, องฺคุตฺตเรปิ. 5-5. อุปสงฺกมิตฺวา (สพฺพตฺถ)} รมณียํ, ทิสฺวานสฺส เอตทโหสิ “ปาสาทิกํ วติทํ อมฺพวนํ มนุญฺญํ รมณียํ, อลํ วติทํ กุลปุตฺตสฺส ปธานตฺถิกสฺส ปธานาย.\csromanspacer{} สเจ มํ ภควา อนุชาเนยฺย, อาคจฺเฉยฺยาหํ อิมํ อมฺพวนํ ปธานายา”ติ.}
\prose{\csromanfolio{117}อถ โข อายสฺมา เมฆิโย เยน ภควา เตนุปสงฺกมิ, อุปสงฺกมิตฺวา ภควนฺตํ อภิวาเทตฺวา เอกมนฺตํ นิสีทิ, เอกมนฺตํ นิสินฺโน โข อายสฺมา เมฆิโย ภควนฺตํ เอตทโวจ\csromandash{}}

\prose{อิธาหํ ภนฺเต ปุพฺพณฺหสมยํ นิวาเสตฺวา ปตฺตจีวรมาทาย ชนฺตุคามํ ปิณฺฑาย ปาวิสิํ, ชนฺตุคาเม ปิณฺฑาย จริตฺวา ปจฺฉาภตฺตํ ปิณฺฑปาตปฏิกฺกนฺโต เยน กิมิกาฬาย นทิยา ตีรํ เตนุปสงฺกมิํ, 5 อทฺทสํ โข อหํ ภนฺเต5 กิมิกาฬาย นทิยา ตีเร ชงฺฆาวิหารํ อนุจงฺกมมาโน อนุวิจรมาโน\footnote{อนุวิจรมาโน อทฺทสํ (สพฺพตฺถ)} อมฺพวนํ ปาสาทิกํ มนุญฺญํ รมณียํ, ทิสฺวาน เม เอตทโหสิ “ปาสาทิกํ วติทํ อมฺพวนํ มนุญฺญํ รมณียํ, อลํ วติทํ กุลปุตฺตสฺส ปธานตฺถิกสฺส ปธานาย.\csromanspacer{} สเจ มํ ภควา อนุชาเนยฺย, อาคจฺเฉยฺยาหํ อิมํ อมฺพวนํ ปธานายา”ติ.\csromanspacer{} สเจ มํ ภนฺเต ภควา อนุชานาติ\footnote{อนุชาเนยฺย (อํ 3. 167)}, คจฺเฉยฺยาหํ ตํ อมฺพวนํ ปธานายาติ.}

\prose{เอวํ วุตฺเต ภควา อายสฺมนฺตํ เมฆิยํ เอตทโวจ “อาคเมหิ ตาว เมฆิย, เอกกมฺหิ\footnote{เอกกมฺหา (สี, อิ), เอกโกมฺหิ (สฺยา)} ตาว, ยาว อญฺโญปิ โกจิ ภิกฺขุ อาคจฺฉตี”ติ.}

\prose{ทุติยมฺปิ โข อายสฺมา เมฆิโย ภควนฺตํ เอตทโวจ “ภควโต ภนฺเต นตฺถิ กิญฺจิ อุตฺตริ\footnote{อุตฺตริํ (สี, สฺยา, กํ, อิ)} กรณียํ, นตฺถิ กตสฺส วา ปติจโย.\csromanspacer{} มยฺหํ โข ปน ภนฺเต อตฺถิ อุตฺตริ กรณียํ, อตฺถิ กตสฺส ปติจโย.\csromanspacer{} สเจ มํ ภควา อนุชานาติ, คจฺเฉยฺยาหํ ตํ อมฺพวนํ ปธานายา”ติ.\csromanspacer{} ทุติยมฺปิ}

\vspace{\tipitakaparskip}
\csromancenter{อุทานปาฬิ โข ภควา อายสฺมนฺตํ เมฆิยํ เอตทโวจ “อาคเมหิ ตาว เมฆิย, เอกกมฺหิ ตาว, ยาว อญฺโญปิ โกจิ ภิกฺขุ อาคจฺฉตี”ติ.}
\csromanmatikaanchor{mtk.60}
\csromantocmarkref{subsection}{10. ภทฺทิยสุตฺต}{mtk.60}
\bookmark[dest={mtk.60},level=2]{10. ภทฺทิยสุตฺต}
\prose{\csromanfolio{118}ตติยมฺปิ โข อายสฺมา เมฆิโย ภควนฺตํ เอตทโวจ “ภควโต ภนฺเต นตฺถิ กิญฺจิ อุตฺตริ กรณียํ, นตฺถิ กตสฺส วา ปติจโย.\csromanspacer{} มยฺหํ โข ปน ภนฺเต อตฺถิ อุตฺตริ กรณียํ, อตฺถิ กตสฺส ปติจโย.\csromanspacer{} สเจ มํ ภควา อนุชานาติ, คจฺเฉยฺยาหํ ตํ อมฺพวนํ ปธานายา”ติ. “ปธานนฺ”ติ โข เมฆิย วทมานํ กินฺติ วเทยฺยาม, ยสฺสทานิ ตฺวํ เมฆิย กาลํ มญฺญสีติ.}

\prose{อถ โข อายสฺมา เมฆิโย อุฏฺฐายาสนา ภควนฺตํ อภิวาเทตฺวา ปทกฺขิณํ กตฺวา เยน ตํ อมฺพวนํ เตนุปสงฺกมิ, อุปสงฺกมิตฺวา ตํ อมฺพวนํ อชฺโฌคาเหตฺวา\footnote{อชฺโฌคเหตฺวา (สี, สฺยา, อิ)} อญฺญตรสฺมิํ รุกฺขมูเล ทิวาวิหารํ นิสีทิ.\csromanspacer{} อถ โข อายสฺมโต เมฆิยสฺส ตสฺมิํ อมฺพวเน วิหรนฺตสฺส เยภุยฺเยน ตโย ปาปกา อกุสลา วิตกฺกา สมุทาจรนฺติ.\csromanspacer{} เสยฺยถิทํ, กามวิตกฺโก พฺยาปาทวิตกฺโก วิหิํสาวิตกฺโก\footnote{วิตกฺโกติ (สี, อิ, ก)}. อถ โข อายสฺมโต เมฆิยสฺส เอตทโหสิ “อจฺฉริยํ วต โภ, อพฺภุตํ วต โภ, สทฺธาย จ วตมฺหา อคารสฺมา อนคาริยํ ปพฺพชิตา, อถ จ ปนิเมหิ ตีหิ ปาปเกหิ อกุสเลหิ วิตกฺเกหิ อนฺวาสตฺตา.\csromanspacer{} เสยฺยถิทํ, กามวิตกฺเกน พฺยาปาทวิตกฺเกน วิหิํสาวิตกฺเกน.”}

\prose{อถ โข อายสฺมา เมฆิโย สายนฺหสมยํ ปฏิสลฺลานา วุฏฺฐิโต เยน ภควา เตนุปสงฺกมิ, อุปสงฺกมิตฺวา ภควนฺตํ อภิวาเทตฺวา เอกมนฺตํ นิสีทิ, เอกมนฺตํ นิสินฺโน โข อายสฺมา เมฆิโย ภควนฺตํ เอตทโวจ “อิธ มยฺหํ ภนฺเต ตสฺมิํ อมฺพวเน วิหรนฺตสฺส เยภุยฺเยน ตโย ปาปกา อกุสลา วิตกฺกา สมุทาจรนฺติ.\csromanspacer{} เสยฺยถิทํ, กามวิตกฺโก พฺยาปาทวิตกฺโก วิหิํสาวิตกฺโก.\csromanspacer{} ตสฺส มยฺหํ ภนฺเต เอตทโหสิ “อจฺฉริยํ วต โภ อพฺภุตํ วต โภ, สทฺธาย จ วตมฺหา อคารสฺมา อนคาริยํ ปพฺพชิตา, อถ จ ปนิเมหิ ตีหิ ปาปเกหิ}

\vspace{\tipitakaparskip}
\csromancenter{เมฆิยวคฺค อกุสเลหิ วิตกฺเกหิ อนฺวาสตฺตา.\csromanspacer{} เสยฺยถิทํ, กามวิตกฺเกน พฺยาปาทวิตกฺเกน วิหิํสาวิตกฺเกน.”}
\prose{\csromanfolio{119}อปริปกฺกาย เมฆิย เจโตวิมุตฺติยา ปญฺจ ธมฺมา ปริปากาย สํวตฺตนฺติ.\csromanspacer{} กตเม ปญฺจ?}

\prose{อิธ เมฆิย ภิกฺขุ กลฺยาณมิตฺโต โหติ กลฺยาณสหาโย กลฺยาณสมฺปวงฺโก, อปริปกฺกาย เมฆิย เจโตวิมุตฺติยา อยํ ปฐโม ธมฺโม ปริปากาย สํวตฺตติ.}

\prose{ปุน จปรํ เมฆิย ภิกฺขุ สีลวา โหติ ปาติโมกฺขสํวรสํวุโต วิหรติ อาจารโคจรสมฺปนฺโน, อณุมตฺเตสุ วชฺเชสุ ภยทสฺสาวี สมาทาย สิกฺขติ สิกฺขาปเทสุ, อปริปกฺกาย เมฆิย เจโตวิมุตฺติยา อยํ ทุติโย ธมฺโม ปริปากาย สํวตฺตติ.}

\prose{ปุน จปรํ เมฆิย ภิกฺขุ ยายํ กถา อภิสลฺเลขิกา เจโตวิวรณสปฺปายา เอกนฺตนิพฺพิทาย วิราคาย นิโรธาย อุปสมาย อภิญฺญาย สมฺโพธาย นิพฺพานาย สํวตฺตติ, เสยฺยถิทํ, อปฺปิจฺฉกถา สนฺตุฏฺฐิกถา ปวิเวกกถา อสํสคฺคกถา วีริยารมฺภกถา สีลกถา สมาธิกถา ปญฺญากถา วิมุตฺติกถา วิมุตฺติญาณทสฺสนกถา, เอวรูปาย กถาย นิกามลาภี โหติ อกิจฺฉลาภี อกสิรลาภี, อปริปากาย เมฆิย เจโตวิมุตฺติยา อยํ ตติโย ธมฺโม ปริปากาย สํวตฺตติ.}

\prose{ปุน จปรํ เมฆิย ภิกฺขุ อารทฺธวีริโย วิหรติ อกุสลานํ ธมฺมานํ ปหานาย กุสลานํ ธมฺมานํ อุปสมฺปทาย\footnote{อุปฺปาทาย (สฺยา)} ถามวา ทฬฺหปรกฺกโม อนิกฺขิตฺตธุโร กุสเลสุ ธมฺเมสุ, อปริปกฺกาย เมฆิย เจโตวิมุตฺติยา อยํ จตุตฺโถ ธมฺโม ปริปากาย สํวตฺตติ.}

\prose{ปุน จปรํ เมฆิย ภิกฺขุ ปญฺญวา โหติ อุทยตฺถคามินิยา ปญฺญาย สมนฺนาคโต อริยาย นิพฺเพธิกาย สมฺมา ทุกฺขกฺขยคามินิยา, อปริปกฺกาย เมฆิย เจโตวิมุตฺติยา อยํ ปญฺจโม ธมฺโม ปริปากาย สํวตฺตติ.\csromanspacer{} อปริปกฺกาย เมฆิย เจโตวิมุตฺติยา อิเม ปญฺจ ธมฺมา ปริปากาย สํวตฺตนฺติ.}

\gambhira{อุทานปาฬิ}
\prose{\csromanfolio{120}กลฺยาณมิตฺตสฺเสตํ เมฆิย ภิกฺขุโน ปาฏิกงฺขํ กลฺยาณสหายสฺส กลฺยาณสมฺปวงฺกสฺส “ยํ สีลวา ภวิสฺสติ ปาติโมกฺขสํวรสํวุโต วิหริสฺสติ อาจารโคจรสมฺปนฺโน, อณุมตฺเตสุ วชฺเชสุ ภยทสฺสาวี สมาทาย สิกฺขิสฺสติ สิกฺขาปเทสุ”.}

\prose{กลฺยาณมิตฺตสฺเสตํ เมฆิย ภิกฺขุโน ปาฏิกงฺขํ กลฺยาณสหายสฺส กลฺยาณสมฺปวงฺกสฺส “ยํ ยายํ กถา อภิสลฺเลขิกา เจโตวิวรณสปฺปายา เอกนฺตนิพฺพิทาย วิราคาย นิโรธาย อุปสมาย อภิญฺญาย สมฺโพธาย นิพฺพานาย สํวตฺตติ, เสยฺยถิทํ, อปฺปิจฺฉกถา สนฺตุฏฺฐิกถา ปวิเวกกถา อสํสคฺคกถา วีริยารมฺภกถา สีลกถา สมาธิกถา ปญฺญากถา วิมุตฺติกถา วิมุตฺติญาณทสฺสนกถา, เอวรูปาย กถาย นิกามลาภี ภวิสฺสติ อกิจฺฉลาภี อกสิรลาภี”.}

\prose{กลฺยาณมิตฺตสฺเสตํ เมฆิย ภิกฺขุโน ปาฏิกงฺขํ กลฺยาณสหายสฺส กลฺยาณสมฺปวงฺกสฺส “ยํ อารทฺธวีริโย วิหริสฺสติ อกุสลานํ ธมฺมานํ ปหานาย กุสลานํ ธมฺมานํ อุปสมฺปทาย ถามวา ทฬฺหปรกฺกโม อนิกฺขิตฺตธุโร กุสเลสุ ธมฺเมสุ”.}

\prose{กลฺยาณมิตฺตสฺเสตํ เมฆิย ภิกฺขุโน ปาฏิกงฺขํ กลฺยาณสหายสฺส กลฺยาณสมฺปวงฺกสฺส “ยํ ปญฺญวา ภวิสฺสติ อุทยตฺถคามินิยา ปญฺญาย สมนฺนาคโต อริยาย นิพฺเพธิกาย สมฺมา ทุกฺขกฺขยคามินิยา”.}

\csromanmatikaanchor{mtk.61}
\csromanmatikaanchor{mtk.62}
\csromantocmarknopage{chapter}{3. นนฺทวคฺค}{mtk.61}
\bookmark[dest={mtk.61},level=0]{3. นนฺทวคฺค}
\csromantocmarkref{subsection}{1. กมฺมวิปากชสุตฺต}{mtk.62}
\bookmark[dest={mtk.62},level=2]{1. กมฺมวิปากชสุตฺต}
\prose{เตน จ ปน เมฆิย ภิกฺขุนา อิเมสุ ปญฺจสุ ธมฺเมสุ ปติฏฺฐาย จตฺตาโร ธมฺมา อุตฺตริ ภาเวตพฺพา\csromandash{}อสุภา ภาเวตพฺพา ราคสฺส ปหานาย, เมตฺตา ภาเวตพฺพา พฺยาปาทสฺส ปหานาย, อานาปานสฺสติ ภาเวตพฺพา วิตกฺกุปจฺเฉทาย, อนิจฺจสญฺญา ภาเวตพฺพา อสฺมิมานสมุคฺฆาตาย.\csromanspacer{} อนิจฺจสญฺญิโน หิ เมฆิย อนตฺตสญฺญา สณฺฐาติ, อนตฺตสญฺญี อสฺมิมานสมุคฺฆาตํ ปาปุณาติ, ทิฏฺเฐว ธมฺเม นิพฺพานนฺติ.}

\prose{อถ โข ภควา เอตมตฺถํ วิทิตฺวา ตายํ เวลายํ อิมํ อุทานํ อุทาเนสิ\csromandash{}}

\chapterheadpageread{4. เมฆิยวคฺค}
\csromangathameasure{“ขุทฺทา วิตกฺกา สุขุมา วิตกฺกา, \\ อนุคตา มนโส อุปฺปิลาวา2. \\ เอเต อวิทฺวา มนโส วิตกฺเก, \\ หุรา หุรํ ธาวติ ภนฺตจิตฺโต. \\ เอเต จ วิทฺวา มนโส วิตกฺเก, \\ อาตาปิโย สํวรตี สตีมา. \\ อนุคเต มนโส อุปฺปิลาเว, \\ อเสสเมเต ปชหาสิ พุทฺโธ”ติ.ปฐมํ.}
\csromangathagroup{\csromangathastanza{\csromanfolio{121}“ขุทฺทา วิตกฺกา สุขุมา วิตกฺกา, \\ อนุคตา\footnote{อนุคฺคตา (สี, ก, อฏฺฐกถายํ ปาฐนฺตรํ.) 2. อุพฺพิลาปา (สี, สฺยา, อิ)} มนโส อุปฺปิลาวา2. \\ เอเต อวิทฺวา มนโส วิตกฺเก, \\ หุรา หุรํ ธาวติ ภนฺตจิตฺโต.}\csromangathastanza{เอเต จ วิทฺวา มนโส วิตกฺเก, \\ อาตาปิโย สํวรตี สตีมา. \\ อนุคเต มนโส อุปฺปิลาเว, \\ อเสสเมเต ปชหาสิ พุทฺโธ”ติ.\csromanspacer{}\csromanspacer{}ปฐมํ.}}

\csromanmatikaanchor{mtk.63}
\csromantocmarkref{subsection}{2. นนฺทสุตฺต}{mtk.63}
\bookmark[dest={mtk.63},level=2]{2. นนฺทสุตฺต}
\csromanheader{2}{2. อุทฺธตสุตฺต}
\proseitem{32}{เอวํ เม สุตํ\csromandash{}เอกํ สมยํ ภควา กุสินารายํ วิหรติ อุปวตฺตเน มลฺลานํ สาลวเน.\csromanspacer{} เตน โข ปน สมเยน สมฺพหุลา ภิกฺขู ภควโต อวิทูเร อรญฺญกุฏิกายํ วิหรนฺติ อุทฺธตา อุนฺนฬา จปลา มุขรา วิกิณฺณวาจา มุฏฺฐสฺสติโน อสมฺปชานา อสมาหิตา วิพฺภนฺตจิตฺตา ปากตินฺทฺริยา.}

\prose{อทฺทสา โข ภควา เต สมฺพหุเล ภิกฺขู อวิทูเร อรญฺญกุฏิกายํ วิหรนฺเต อุทฺธเต อุนฺนเฬ จปเล มุขเร วิกิณฺณวาเจ มุฏฺฐสฺสติโน อสมฺปชาเน อสมาหิเต วิพฺภนฺตจิตฺเต ปากตินฺทฺริเย.}

\prose{อถ โข ภควา เอตมตฺถํ วิทิตฺวา ตายํ เวลายํ อิมํ อุทานํ อุทาเนสิ\csromandash{}}

\prose{\csromansymbolfootnote{+}{ขุ 10. 71 ปิฏฺเฐปิ.}“อรกฺขิเตน กาเยน\footnote{จิตฺเตน (เนตฺติยํ)}, มิจฺฉาทิฏฺฐิหเตน\footnote{มิจฺฉาทิฏฺฐิคเตน (พหุสุ)} จ.}

\csromangathameasure{ถินมิทฺธา ภิภูเตน, วสํ มารสฺส คจฺฉติ.}
\csromangathagroup{\csromangathastanza{ถินมิทฺธา\footnote{ถีนมิทฺธา (สี, สฺยา, กํ, อิ)} ภิภูเตน, วสํ มารสฺส คจฺฉติ.}}

\prose{\csromansymbolfootnote{*}{ขุ 10. 41 ปิฏฺเฐปิ.}ตสฺมา รกฺขิตจิตฺต’สฺส, สมฺมาสงฺกปฺปโคจโร.}

\csromangathameasure{สมฺมาทิฏฺฐิปุเรกฺขาโร, ญตฺวาน อุทยพฺพยํ. \\ ถินมิทฺธาภิภู ภิกฺขุ, สพฺพา ทุคฺคติโย ชเห”ติ.ทุติยํ.}
\csromangathagroup{\csromangathastanza{สมฺมาทิฏฺฐิปุเรกฺขาโร, ญตฺวาน อุทยพฺพยํ. \\ ถินมิทฺธาภิภู ภิกฺขุ, สพฺพา ทุคฺคติโย ชเห”ติ.\csromanspacer{}\csromanspacer{}ทุติยํ.}}

\csromanheader{2}{อุทานปาฬิ \textbf{3. โคปาลกสุตฺต}}
\proseitem{33}{\csromanfolio{122}เอวํ เม สุตํ\csromandash{}เอกํ สมยํ ภควา โกสเลสุ จาริกํ จรติ มหตา ภิกฺขุสํเฆน สทฺธิํ.\csromanspacer{} อถ โข ภควา มคฺคา โอกฺกมฺม เยน อญฺญตรํ รุกฺขมูลํ เตนุปสงฺกมิ, อุปสงฺกมิตฺวา ปญฺญตฺเต อาสเน นิสีทิ.}

\prose{อถ โข อญฺญตโร โคปาลโก เยน ภควา เตนุปสงฺกมิ, อุปสงฺกมิตฺวา ภควนฺตํ อภิวาเทตฺวา เอกมนฺตํ นิสีทิ.\csromanspacer{} เอกมนฺตํ นิสินฺนํ โข ตํ โคปาลกํ ภควา ธมฺมิยา กถาย สนฺทสฺเสสิ สมาทเปสิ\footnote{สมาทาเปสิ (?)} สมุตฺเตเชสิ สมฺปหํเสสิ.}

\prose{อถ โข โส โคปาลโก ภควตา ธมฺมิยา กถาย สนฺทสฺสิโต สมาทปิโต\footnote{สมาทาปิโต (?)} สมุตฺเตชิโต สมฺปหํสิโต ภควนฺตํ เอตทโวจ “อธิวาเสตุ เม ภนฺเต ภควา สฺวาตนาย ภตฺตํ สทฺธิํ ภิกฺขุสํเฆนา”ติ.\csromanspacer{} อธิวาเสสิ ภควา ตุณฺหีภาเวน.\csromanspacer{} อถ โข โส โคปาลโก ภควโต อธิวาสนํ วิทิตฺวา อุฏฺฐายาสนา ภควนฺตํ อภิวาเทตฺวา ปทกฺขิณํ กตฺวา ปกฺกามิ.}

\prose{อถ โข โส โคปาลโก ตสฺสา รตฺติยา อจฺจเยน สเก นิเวสเน ปหูตํ อปฺโปทกปายาสํ ปฏิยาทาเปตฺวา นวญฺจ สปฺปิํ ภควโต กาลํ อาโรเจสิ “กาโล ภนฺเต นิฏฺฐิตํ ภตฺตนฺ”ติ.\csromanspacer{} อถ โข ภควา ปุพฺพณฺหสมยํ นิวาเสตฺวา ปตฺตจีวรมาทาย สทฺธิํ ภิกฺขุสํเฆน เยน ตสฺส โคปาลกสฺส นิเวสนํ เตนุปสงฺกมิ, อุปสงฺกมิตฺวา ปญฺญตฺเต อาสเน นิสีทิ.\csromanspacer{} อถ โข โส โคปาลโก พุทฺธปฺปมุขํ ภิกฺขุสํฆํ อปฺโปทกปายาเสน\footnote{อปฺโปทกปายเสน จ (ก)} นเวน จ สปฺปินา สหตฺถา สนฺตปฺเปสิ สมฺปวาเรสิ.\csromanspacer{} อถ โข โส โคปาลโก ภควนฺตํ ภุตฺตาวิํ โอนีตปตฺตปาณิํ อญฺญตรํ นีจํ อาสนํ คเหตฺวา เอกมนฺตํ นิสีทิ, เอกมนฺตํ นิสินฺนํ โข ตํ โคปาลกํ ภควา ธมฺมิยา กถาย สนฺทสฺเสตฺวา สมาทเปตฺวา สมุตฺเตเชตฺวา สมฺปหํเสตฺวา อุฏฺฐายาสนา ปกฺกามิ.\csromanspacer{} อถ โข อจิรปกฺกนฺตสฺส ภควโต ตํ โคปาลกํ อญฺญตโร ปุริโส สีมนฺตริกาย ชีวิตา โวโรเปสิ.}

\chapterheadpageread{4. เมฆิยวคฺค}
\prose{\csromanfolio{123}อถ โข สมฺพหุลา ภิกฺขู เยน ภควา เตนุปสงฺกมิํสุ, อุปสงฺกมิตฺวา ภควนฺตํ อภิวาเทตฺวา เอกมนฺตํ นิสีทิํสุ, เอกมนฺตํ นิสินฺนา โข เต ภิกฺขู ภควนฺตํ เอตทโวจุํ “เยน ภนฺเต โคปาลเกน อชฺช พุทฺธปฺปมุโข ภิกฺขุสํโฆ อปฺโปทกปายาเสน นเวน จ สปฺปินา สหตฺถา สนฺตปฺปิโต สมฺปวาริโต, โส กิร ภนฺเต โคปาลโก อญฺญตเรน ปุริเสน สีมนฺตริกาย ชีวิตา โวโรปิโต”ติ.}

\prose{อถ โข ภควา เอตมตฺถํ วิทิตฺวา ตายํ เวลายํ อิมํ อุทานํ อุทาเนสิ\csromandash{}}

\csromangathameasure{“ทิโส ทิสํ ยํ ตํ กยิรา, เวรี วา ปน เวรินํ. \\ มิจฺฉาปณิหิตํ จิตฺตํ, ปาปิโย นํ ตโต กเร”ติ.ตติยํ.}
\csromangathagroup{\csromangathastanza{“ทิโส ทิสํ ยํ ตํ กยิรา, เวรี วา ปน เวรินํ. \\ มิจฺฉาปณิหิตํ จิตฺตํ, ปาปิโย นํ ตโต กเร”ติ.\csromanspacer{}\csromanspacer{}ตติยํ.}}

\csromanheader{2}{4. ยกฺขปหารสุตฺต}
\proseitem{34}{เอวํ เม สุตํ\csromandash{}เอกํ สมยํ ภควา ราชคเห วิหรติ เวฬุวเน กลนฺทกนิวาเป.\csromanspacer{} เตน โข ปน สมเยน อายสฺมา จ สาริปุตฺโต, อายสฺมา จ มหาโมคฺคลฺลาโน กโปตกนฺทรายํ วิหรนฺติ.\csromanspacer{} เตน โข ปน สมเยน อายสฺมา สาริปุตฺโต ชุณฺหาย รตฺติยา นโวโรปิเตหิ เกเสหิ อพฺโภกาเส นิสินฺโน โหติ อญฺญตรํ สมาธิํ สเมปชฺชิตฺวา.}

\prose{เตน โข ปน สมเยน เทฺว ยกฺขา สหายกา อุตฺตราย ทิสาย ทกฺขิณํ ทิสํ คจฺฉนฺติ เกนจิเทว กรณีเยน.\csromanspacer{} อทฺทสํสุ โข เต ยกฺขา อายสฺมนฺตํ สาริปุตฺตํ ชุณฺหาย รตฺติยา นโวโรปิเตหิ เกเสหิ อพฺโภกาเส นิสินฺนํ, ทิสฺวาน เอโก ยกฺโข ทุติยํ ยกฺขํ เอตทโวจ “ปฏิภาติ มํ สมฺม อิมสฺส สมณสฺส สีเส ปหารํ ทาตุนฺ”ติ.\csromanspacer{} เอวํ วุตฺเต โส ยกฺโข ตํ ยกฺขํ เอตทโวจ “อลํ สมฺม มา สมณํ อาสาเทสิ, อุฬาโร โส สมฺม สมโณ มหิทฺธิโก มหานุภาโว”ติ.}

\prose{ทุติยมฺปิ โข โส ยกฺโข ตํ ยกฺขํ เอตทโวจ “ปฏิภาติ มํ สมฺม อิมสฺส สมณสฺส สีเส ปหารํ ทาตุนฺ”ติ.\csromanspacer{} ทุติยมฺปิ โข โส ยกฺโข ตํ ยกฺขํ เอตทโวจ “อลํ สมฺม มา สมณํ อาสาเทสิ, อุฬาโร โส สมโณ มหิทฺธิโก มหานุภาโว”ติ.\csromanspacer{} ตติยมฺปิ โข โส ยกฺโข ตํ ยกฺขํ เอตทโวจ “ปฏิภาติ มํ สมฺม อิมสฺส สมณสฺส สีเส}

\vspace{\tipitakaparskip}
\csromancenter{อุทานปาฬิ ปหารํ ทาตุนฺ”ติ.\csromanspacer{} ตติยมฺปิ โข โส ยกฺโข ตํ ยกฺขํ เอตทโวจ “อลํ สมฺม มา สมณํ อาสาเทสิ, อุฬาโร โส สมฺม สมโณ มหิทฺธิโก มหานุภาโว”ติ.}
\csromanmatikaanchor{mtk.64}
\csromantocmarkref{subsection}{3. ยโสชสุตฺต}{mtk.64}
\bookmark[dest={mtk.64},level=2]{3. ยโสชสุตฺต}
\prose{\csromanfolio{124}อถ โข โส ยกฺโข ตํ ยกฺขํ อนาทิยิตฺวา อายสฺมโต สาริปุตฺตตฺเถรสฺส สีเส ปหารํ อทาสิ.\csromanspacer{} ตาว มหา ปหาโร อโหสิ, อปิ เตน ปหาเรน สตฺตรตนํ วา อฑฺฒฏฺฐมรตนํ วา นาคํ โอสาเทยฺย, มหนฺตํ วา ปพฺพตกูฏํ ปทาเลยฺย.\csromanspacer{} อถ จ ปน โส ยกฺโข “ฑยฺหามิ ฑยฺหามี”ติ วตฺวา ตตฺเถว มหานิรยํ อปตาสิ\footnote{อวตฺถาสิ (ก-สี)}.}

\prose{อทฺทสา โข อายสฺมา มหาโมคฺคลฺลาโน ทิพฺเพน จกฺขุนา วิสุทฺเธน อติกฺกนฺตมานุสเกน เตน ยกฺเขน อายสฺมโต สาริปุตฺตตฺเถรสฺส สีเส ปหารํ ทียมานํ, ทิสฺวา เยน อายสฺมา สาริปุตฺโต เตนุปสงฺกมิ, อุปสงฺกมิตฺวา อายสฺมนฺตํ สาริปุตฺตํ เอตทโวจ “กจฺจิ เต อาวุโส ขมนียํ, กจฺจิ ยาปนียํ, กจฺจิ น กิญฺจิ ทุกฺขนฺ”ติ.\csromanspacer{} ขมนียํ เม อาวุโส โมคฺคลฺลาน, ยาปนียํ เม อาวุโส โมคฺคลฺลาน, อปิ จ เม สีสํ โถกํ ทุกฺขนฺติ.}

\prose{อจฺฉริยํ อาวุโส สาริปุตฺต, อพฺภุตํ อาวุโส สาริปุตฺต, ยาว\footnote{ยํ ตฺวํ (สี, ก), ยํ (สฺยา)} มหิทฺธิโก อายสฺมา สาริปุตฺโต มหานุภาโว, อิธ เต อาวุโส สาริปุตฺต อญฺญตโร ยกฺโข สีเส ปหารํ อทาสิ.\csromanspacer{} ตาว มหา ปหาโร อโหสิ, อปิ เตน ปหาเรน สตฺตรตนํ วา อฑฺฒฏฺฐมรตนํ วา นาคํ โอสาเทยฺย, มหนฺตํ วา ปพฺพตกูฏํ ปทาเลยฺย.\csromanspacer{} อถ จ ปนายสฺมา สาริปุตฺโต เอวมาห “ขมนียํ เม อาวุโส โมคฺคลฺลาน, ยาปนียํ เม อาวุโส โมคฺคลฺลาน, อปิ จ เม สีสํ โถกํ ทุกฺขนฺ”ติ.}

\prose{อจฺฉริยํ อาวุโส โมคฺคลฺลาน, อพฺภุตํ อาวุโส โมคฺคลฺลาน, ยาว\footnote{ยํ (สฺยา)} มหิทฺธิโก อายสฺมา มหาโมคฺคลฺลาโน มหานุภาโว, ยตฺร หิ นาม ยกฺขมฺปิ ปสฺสิสฺสติ, มยํ ปเนตรหิ ปํสุปิสาจกมฺปิ น ปสฺสามาติ.}

\prose{อสฺโสสิ โข ภควา ทิพฺพาย โสตธาตุยา วิสุทฺธาย อติกฺกนฺตมานุสิกาย เตสํ อุภินฺนํ มหานาคานํ อิมํ เอวรูปํ กถาสลฺลาปํ.}

\chapterheadpageread{4. เมฆิยวคฺค}
\prose{\csromanfolio{125}อถ โข ภควา เอตมตฺถํ วิทิตฺวา ตายํ เวลายํ อิมํ อุทานํ อุทาเนสิ\csromandash{}}

\prose{\csromansymbolfootnote{*}{ขุ 2. 257; ขุ 10. 129 ปิฏฺเฐสุปิ.}“ยสฺส เสลูปมํ จิตฺตํ, ฐิตํ นานุปกมฺปติ.}

\prose{วิรตฺตํ รชนีเยสุ, โกปเนยฺเย น กุปฺปติ.}

\prose{ยสฺเสวํ ภาวิตํ จิตฺตํ, กุโต ตํ ทุกฺขเมสฺสตี”ติ.\csromanspacer{}\csromanspacer{}จตุตฺถํ.}

\csromanheader{2}{5. นาคสุตฺต}
\proseitem{35}{เอวํ เม สุตํ\csromandash{}เอกํ สมยํ ภควา โกสมฺพิยํ วิหรติ โฆสิตาราเม.\csromanspacer{} เตน โข ปน สมเยน ภควา อากิณฺโณ วิหรติ ภิกฺขูหิ ภิกฺขุนีหิ อุปาสเกหิ อุปาสิกาหิ ราชูหิ ราชมหามตฺเตหิ ติตฺถิเยหิ ติตฺถิยสาวเกหิ, อากิณฺโณ ทุกฺขํ น ผาสุ วิหรติ.\csromanspacer{} อถ โข ภควโต เอตทโหสิ “อหํ โข เอตรหิ อากิณฺโณ วิหรามิ ภิกฺขูหิ ภิกฺขุนีหิ อุปาสเกหิ อุปาสิกาหิ ราชูหิ ราชมหามตฺเตหิ ติตฺถิเยหิ ติตฺถิยสาวเกหิ, อากิณฺโณ ทุกฺขํ น ผาสุ วิหรามิ, ยํนูนาหํ เอโก คณสฺมา วูปกฏฺโฐ วิหเรยฺยนฺ”ติ.}

\prose{อถ โข ภควา ปุพฺพณฺหสมยํ นิวาเสตฺวา ปตฺตจีวรมาทาย โกสมฺพิํ ปิณฺฑาย ปาวิสิ, โกสมฺพิยํ ปิณฺฑาย จริตฺวา ปจฺฉาภตฺตํ ปิณฺฑปาตปฏิกฺกนฺโต สามํ เสนาสนํ สํสาเมตฺวา ปตฺตจีวรมาทาย อนามนฺเตตฺวา อุปฏฺฐากํ อนปโลเกตฺวา ภิกฺขุสํฆํ เอโก อทุติโย เยน ปาลิเลยฺยกํ เตน จาริกํ ปกฺกามิ.\csromanspacer{} อนุปุพฺเพน จาริกํ จรมาโน เยน ปาลิเลยฺยกํ ตทวสริ, ตตฺร สุทํ ภควา ปาลิเลยฺยเก วิหรติ รกฺขิตวนสณฺเฑ ภทฺทสาลมูเล.}

\prose{อญฺญตโรปิ โข หตฺถินาโค อากิณฺโณ วิหรติ หตฺถีหิ หตฺถินีหิ หตฺถิกลเภหิ หตฺถิจฺฉาเปหิ, ฉินฺนคฺคานิ เจว ติณานิ ขาทติ, โอภคฺโคภคฺคญฺจสฺส สาขาภงฺคํ ขาทนฺติ, อาวิลานิ จ ปานียานิ ปิวติ, โอคาหา จสฺส อุตฺติณฺณสฺส หตฺถินิโย กายํ อุปนิฆํสนฺติโย คจฺฉนฺติ, อากิณฺโณ ทุกฺขํ น ผาสุ วิหรติ.\csromanspacer{} อถ โข ตสฺส หตฺถินาคสฺส เอตทโหสิ “อหํ โข เอตรหิ อากิณฺโณ วิหรามิ หตฺถีหิ หตฺถินีหิ หตฺถิกลเภหิ หตฺถิจฺฉาเปหิ, ฉินฺนคฺคานิ เจว ติณานิ ขาทามิ, โอภคฺโคภคฺคญฺจ เม สาขาภงฺคํ ขาทนฺติ, อาวิลานิ จ ปานียานิ ปิวามิ, โอคาหา จ เม}

\vspace{\tipitakaparskip}
\csromancenter{อุทานปาฬิ อุตฺติณฺณสฺส หตฺถินิโย กายํ อุปนิฆํสนฺติโย คจฺฉนฺติ, อากิณฺโณ ทุกฺขํ น ผาสุ วิหรามิ, ยํนูนาหํ เอโก คณสฺมา วูปกฏฺโฐ วิหเรยฺยนฺ”ติ.}
\prose{\csromanfolio{126}อถ โข โส หตฺถินาโค ยูถา อปกฺกมฺม เยน ปาลิเลยฺยกํ รกฺขิตวนสณฺโฑ ภทฺทสาลมูลํ, เยน ภควา เตนุปสงฺกมิ, ตตฺร สุทํ\footnote{อุปสงฺกมิตฺวา ตตฺร สุทํ (สฺยา, อิ, ก) 2-2. อปฺปหริตญฺจ กโรติ, โสณฺฑาย (พหุสุ)} โส หตฺถินาโค ยสฺมิํ ปเทเส ภควา วิหรติ, ตํ ปเทสํ 2 อปฺปหริตํ กโรติ, โสณฺฑาย จ2 ภควโต ปานียํ ปริโภชนียํ อุปฏฺฐาเปติ\footnote{อุปฏฺฐเปติ (สี, สฺยา, กํ, อิ)}.}

\prose{อถ โข ภควโต รโหคตสฺส ปฏิสลฺลีนสฺส เอวํ เจตโส ปริวิตกฺโก อุทปาทิ “อหํ โข ปุพฺเพ อากิณฺโณ วิหาสิํ ภิกฺขูหิ ภิกฺขุนีหิ อุปาสเกหิ อุปาสิกาหิ ราชูหิ ราชมหามตฺเตหิ ติตฺถิเยหิ ติตฺถิยสาวเกหิ, อากิณฺโณ ทุกฺขํ น ผาสุ วิหาสิํ.\csromanspacer{} โสมฺหิ เอตรหิ อนากิณฺโณ วิหรามิ ภิกฺขูหิ ภิกฺขุนีหิ อุปาสเกหิ อุปาสิกาหิ ราชูหิ ราชมหามตฺเตหิ ติตฺถิเยหิ ติตฺถิยสาวเกหิ, อนากิณฺโณ สุขํ ผาสุ วิหรามี”ติ.}

\prose{ตสฺสปิ โข หตฺถินาคสฺส เอวํ เจตโส ปริวิตกฺโก อุทปาทิ “อหํ โข ปุพฺเพ อากิณฺโณ วิหาสิํ หตฺถีหิ หตฺถินีหิ หตฺถิกลเภหิ หตฺถิจฺฉาเปหิ, ฉินฺนคฺคานิ เจว ติณานิ ขาทิํ, โอภคฺโคภคฺคญฺจ เม สาขาภงฺคํ ขาทิํสุ, อาวิลานิ จ ปานียานิ อปายิํ, โอคาหา จ เม อุตฺติณฺณสฺส หตฺถินิโย กายํ อุปนิฆํสนฺติโย อคมํสุ, อากิณฺโณ ทุกฺขํ น ผาสุ วิหาสิํ.\csromanspacer{} โสมฺหิ เอตรหิ อนากิณฺโณ วิหรามิ หตฺถีหิ หตฺถินีหิ หตฺถิกลเภหิ หตฺถิจฺฉาเปหิ, อจฺฉินฺนคฺคานิ เจว ติณานิ ขาทามิ, โอภคฺโคภคฺคญฺจ เม สาขาภงฺคํ น ขาทนฺติ, อนาวิลานิ จ ปานียานิ ปิวามิ, โอคาหา จ เม อุตฺติณฺณสฺส หตฺถินิโย น กายํ อุปนิฆํสนฺติโย คจฺฉนฺติ, อนากิณฺโณ สุขํ ผาสุ วิหรามี”ติ.}

\prose{อถ โข ภควา อตฺตโน จ ปวิเวกํ วิทิตฺวา ตสฺส จ หตฺถินาคสฺส เจตสา เจโตปริวิตกฺกมญฺญาย ตายํ เวลายํ อิมํ อุทานํ อุทาเนสิ\csromandash{}}

\chapterheadpageread{4. เมฆิยวคฺค}
\csromangathameasure{“เอตํ นาคสฺส นาเคน, อีสาทนฺตสฺส หตฺถิโน. \\ สเมติ จิตฺตํ จิตฺเตน, ยเทโก รมตี มโน”ติ.ปญฺจมํ.}
\csromangathagroup{\csromangathastanza{\csromanfolio{127}“เอตํ\footnote{เอวํ (ก) วิ 3. 500 ปิฏฺเฐปิ.} นาคสฺส นาเคน, อีสาทนฺตสฺส หตฺถิโน. \\ สเมติ จิตฺตํ จิตฺเตน, ยเทโก รมตี มโน”ติ.\csromanspacer{}\csromanspacer{}ปญฺจมํ.}}

\csromanmatikaanchor{mtk.65}
\csromantocmarkref{subsection}{4. สาริปุตฺตสุตฺต}{mtk.65}
\bookmark[dest={mtk.65},level=2]{4. สาริปุตฺตสุตฺต}
\csromantocmarkref{subsection}{5. มหาโมคฺคลานสุตฺต}{mtk.66}
\bookmark[dest={mtk.66},level=2]{5. มหาโมคฺคลานสุตฺต}
\csromanheader{2}{6. ปิณฺโฑลสุตฺต}
\proseitem{36}{เอวํ เม สุตํ\csromandash{}เอกํ สมยํ ภควา สาวตฺถิยํ วิหรติ เชตวเน อนาถปิณฺฑิกสฺส อาราเม.\csromanspacer{} เตน โข ปน สมเยน อายสฺมา ปิณฺโฑลภารทฺวาโช ภควโต อวิทูเร นิสินฺโน โหติ ปลฺลงฺกํ อาภุชิตฺวา อุชุํ กายํ ปณิธาย อารญฺญิโก, ปิณฺฑปาติโก, ปํสุกูลิโก, เตจีวริโก อปฺปิจฺโฉ สนฺตุฏฺโฐ ปวิวิตฺโต อสํสฏฺโฐ อารทฺธวีริโย\footnote{อารทฺธวิริโย (สี, สฺยา, กํ, อิ)} ธุตวาโท อธิจิตฺตมนุยุตฺโต.}

\prose{อทฺทสา โข ภควา อายสฺมนฺตํ ปิณฺโฑลภารทฺวาชํ อวิทูเร นิสินฺนํ ปลฺลงฺกํ อาภุชิตฺวา อุชุํ กายํ ปณิธาย อารญฺญิกํ, ปิณฺฑปาติกํ, ปํสุกูลิกํ, เตจีวริกํ อปฺปิจฺฉํ สนฺตุฏฺฐํ ปวิวิตฺตํ อสํสฏฺฐํ อารทฺธวีริยํ ธุตวาทํ อธิจิตฺตมนุยุตฺตํ.}

\prose{อถ โข ภควา เอตมตฺถํ วิทิตฺวา ตายํ เวลายํ อิมํ อุทานํ อุทาเนสิ\csromandash{}}

\csromangathameasure{*“อนูปวาโท อนูปฆาโต3, ปาติโมกฺเข จ สํวโร. \\ มตฺตญฺญุตา จ ภตฺตสฺมิํ, ปนฺตญฺจ สยนาสนํ. \\ อธิจิตฺเต จ อาโยโค, เอตํ พุทฺธาน สาสนนฺ”ติ.ฉฏฺฐํ.}
\csromangathagroup{\csromangathastanza{\csromansymbolfootnote{*}{ที 2. 42; เหฏฺฐา 41. 3 อนุปวาโท อนุปฆาโต (สฺยา, อิ, ก)}“อนูปวาโท อนูปฆาโต3, ปาติโมกฺเข จ สํวโร. \\ มตฺตญฺญุตา จ ภตฺตสฺมิํ, ปนฺตญฺจ สยนาสนํ. \\ อธิจิตฺเต จ อาโยโค, เอตํ พุทฺธาน สาสนนฺ”ติ.\csromanspacer{}\csromanspacer{}ฉฏฺฐํ.}}

\csromanheader{2}{7. สาริปุตฺตสุตฺต}
\proseitem{37}{เอวํ เม สุตํ\csromandash{}เอกํ สมยํ ภควา สาวตฺถิยํ วิหรติ เชตวเน อนาถปิณฺฑิกสฺส อาราเม.\csromanspacer{} เตน โข ปน สมเยน อายสฺมา สาริปุตฺโต ภควโต อวิทูเร นิสินฺโน โหติ ปลฺลงฺกํ อาภุชิตฺวา อุชุํ กายํ ปณิธาย อปฺปิจฺโฉ สนฺตุฏฺโฐ ปวิวิตฺโต อสํสฏฺโฐ อารทฺธวีริโย อธิจิตฺตมนุยุตฺโต.}

\gambhira{อุทานปาฬิ}
\csromanmatikaanchor{mtk.67}
\csromanmatikaanchor{mtk.80}
\csromantocmarkref{subsection}{6. ปิลินฺทวจฺฉสุตฺต}{mtk.67}
\bookmark[dest={mtk.67},level=2]{6. ปิลินฺทวจฺฉสุตฺต}
\prose{\csromanfolio{128}อทฺทสา โข ภควา อายสฺมนฺตํ สาริปุตฺตํ อวิทูเร นิสินฺนํ ปลฺลงฺกํ อาภุชิตฺวา อุชุํ กายํ ปณิธาย อปฺปิจฺฉํ สนฺตุฏฺฐํ ปวิวิตฺตํ อสํสฏฺฐํ อารทฺธวีริยํ อธิจิตฺตมนุยุตฺตํ.}

\prose{อถ โข ภควา เอตมตฺถํ วิทิตฺวา ตายํ เวลายํ อิมํ อุทานํ อุทาเนสิ\csromandash{}}

\csromangathameasure{*“อธิเจตโส อปฺปมชฺชโต, \\ มุนิโน โมนปเถสุ สิกฺขโต. \\ โสกา น ภวนฺติ ตาทิโน, \\ อุปสนฺตสฺส สทา สตีมโต”ติ.สตฺตมํ.}
\csromangathagroup{\csromangathastanza{\csromansymbolfootnote{*}{วิ 2. 77; ขุ 2. 234 ปิฏฺเฐสุปิ.}“อธิเจตโส อปฺปมชฺชโต, \\ มุนิโน โมนปเถสุ สิกฺขโต. \\ โสกา น ภวนฺติ ตาทิโน, \\ อุปสนฺตสฺส สทา สตีมโต”ติ.\csromanspacer{}\csromanspacer{}สตฺตมํ.}}

\csromanheader{2}{8. สุนฺทรี สุตฺต}
\proseitem{38}{\csromansymbolfootnote{+}{เหฏฺฐา 89 ปิฏฺเฐปิ.}เอวํ เม สุตํ\csromandash{}เอกํ สมยํ ภควา สาวตฺถิยํ วิหรติ เชตวเน อนาถปิณฺฑิกสฺส อาราเม.\csromanspacer{} เตน โข ปน สมเยน ภควา สกฺกโต โหติ ครุกโต มานิโต ปูชิโต อปจิโต ลาภี จีวรปิณฺฑปาตเสนาสนคิลานปจฺจยเภสชฺชปริกฺขารานํ, ภิกฺขุสํโฆปิ สกฺกโต โหติ ครุกโต มานิโต ปูชิโต อปจิโต ลาภี จีวรปิณฺฑปาตเสนาสนคิลานปจฺจยเภสชฺชปริกฺขารานํ.\csromanspacer{} อญฺญติตฺถิยา ปน ปริพฺพาชกา อสกฺกตา โหนฺติ อครุกตา อมานิตา อปูชิตา อนปจิตา น ลาภิโน จีวรปิณฺฑปาตเสนาสนคิลานปจฺจยเภสชฺชปริกฺขารานํ.}

\prose{อถ โข เต อญฺญติตฺถิยา ปริพฺพาชกา ภควโต สกฺการํ อสหมานา ภิกฺขุสํฆสฺส จ เยน สุนฺทรี ปริพฺพาชิกา เตนุปสงฺกมิํสุ, อุปสงฺกมิตฺวา สุนฺทริํ ปริพฺพาชิกํ เอตทโวจุํ “อุสฺสหสิ ตฺวํ ภคินิ ญาตีนํ อตฺถํ กาตุนฺ”ติ.\csromanspacer{} กฺยาหํ อยฺยา กโรมิ, กิํ มยา น สกฺกา\footnote{กิํ มยา สกฺกา (สฺยา, อิ)} กาตุํ, ชีวิตมฺปิ เม ปริจฺจตฺตํ ญาตีนํ อตฺถายาติ.}

\prose{เตน หิ ภคินิ อภิกฺขณํ เชตวนํ คจฺฉาหีติ. “เอวํ อยฺยา”ติ โข สุนฺทรี ปริพฺพาชิกา เตสํ อญฺญติตฺถิยานํ ปริพฺพาชกานํ ปฏิสฺสุตฺวา อภิกฺขณํ เชตวนํ อคมาสิ.}

\chapterheadpageread{4. เมฆิยวคฺค}
\prose{\csromanfolio{129}ยทา เต อญฺญิํสุ อญฺญติตฺถิยา ปริพฺพาชกา “โวทิฏฺฐา โข สุนฺทรี ปริพฺพาชิกา พหุชเนน อภิกฺขณํ เชตวนํ คจฺฉนฺตี”ติ\footnote{คจฺฉตีติ (สี, สฺยา, กํ, อิ)}, อถ นํ ชีวิตา โวโรเปตฺวา ตตฺเถว เชตวนสฺส ปริขากูเป นิกฺขิปิตฺวา\footnote{นิขนิตฺวา (สี, สฺยา, อิ)} เยน ราชา ปเสนทิ โกสโล เตนุปสงฺกมิํสุ, อุปสงฺกมิตฺวา ราชานํ ปเสนทิํ โกสลํ เอตทโวจุํ “ยา สา มหาราช สุนฺทรี ปริพฺพาชิกา, สา โน น ทิสฺสตี”ติ.\csromanspacer{} กตฺถ ปน ตุเมฺห อาสงฺกถาติ.\csromanspacer{} เชตวเน มหาราชาติ.\csromanspacer{} เตน หิ เชตวนํ วิจินถาติ.}

\prose{อถ โข เต อญฺญติตฺถิยา ปริพฺพาชกา เชตวนํ วิจินิตฺวา ยถานิกฺขิตฺตํ ปริขากูปา อุทฺธริตฺวา มญฺจกํ อาโรเปตฺวา สาวตฺถิํ ปเวเสตฺวา รถิยาย รถิยํ สิงฺฆาฏเกน สิงฺฆาฏกํ อุปสงฺกมิตฺวา มนุสฺเส อุชฺฌาเปสุํ\csromandash{}}

\prose{ปสฺสถายฺยา สมณานํ สกฺยปุตฺติยานํ กมฺมํ, อลชฺชิโน อิเม สมณา สกฺยปุตฺติยา ทุสฺสีลา ปาปธมฺมา มุสาวาทิโน อพฺรหฺมจาริโน, อิเม หิ นาม ธมฺมจาริโน สมจาริโน พฺรหฺมจาริโน สจฺจวาทิโน สีลวนฺโต กลฺยาณธมฺมา ปฏิชานิสฺสนฺติ, นตฺถิ อิเมสํ สามญฺญํ, นตฺถิ อิเมสํ พฺรหฺมญฺญํ, นฏฺฐํ อิเมสํ สามญฺญํ, นฏฺฐํ อิเมสํ พฺรหฺมญฺญํ, กุโต อิเมสํ สามญฺญํ, กุโต อิเมสํ พฺรหฺมญฺญํ, อปคตา อิเม สามญฺญา, อปคตา อิเมพฺรหฺมญฺญา, กถํ หิ นาม ปุริโส ปุริสกิจฺจํ กริตฺวา อิตฺถิํ ชีวิตา โวโรเปสฺสตีติ.}

\csromanmatikaanchor{mtk.68}
\csromantocmarkref{subsection}{7. สกฺกุทานสุตฺต}{mtk.68}
\bookmark[dest={mtk.68},level=2]{7. สกฺกุทานสุตฺต}
\prose{เตน โข ปน สมเยน สาวตฺถิยํ มนุสฺสา ภิกฺขู ทิสฺวา อสพฺภาหิ ผรุสาหิ วาจาหิ อกฺโกสนฺติ ปริภาสนฺติ โรสนฺติ วิเหสนฺติ\csromandash{}}

\prose{อลชฺชิโน อิเม สมณา สกฺยปุตฺติยา ทุสฺสีลา ปาปธมฺมา มุสาวาทิโน อพฺรหฺมจาริโน, อิเม หิ นาม ธมฺมจาริโน สมจาริโน พฺรหฺมจาริโน สจฺจวาทิโน สีลวนฺโต กลฺยาณธมฺมา ปฏิชานิสฺสนฺติ, นตฺถิ อิเมสํ สามญฺญํ, นตฺถิ อิเมสํ พฺรหฺมญฺญํ, นฏฺฐํ อิเมสํ สามญฺญํ, นฏฺฐํ อิเมสํ พฺรหฺมญฺญํ, กุโต อิเมสํ สามญฺญํ, กุโต อิเมสํ พฺรหฺมญฺญํ, อปคตา อิเม สามญฺญา, อปคตา อิเม พฺรหฺมญฺญา, กถํ หิ นาม ปุริโส ปุริสกิจฺจํ กริตฺวา อิตฺถิํ ชีวิตา โวโรเปสฺสตีติ.}

\prose{อถ โข สมฺพหุลา ภิกฺขู ปุพฺพณฺหสมยํ นิวาเสตฺวา ปตฺตจีวรมาทาย สาวตฺถิํ ปิณฺฑาย ปาวิสิํสุ, สาวตฺถิยํ ปิณฺฑาย จริตฺวา ปจฺฉาภตฺตํ}

\vspace{\tipitakaparskip}
\csromancenter{อุทานปาฬิ ปิณฺฑปาตปฏิกฺกนฺตา เยน ภควา เตนุปสงฺกมิํสุ, อุปสงฺกมิตฺวา ภควนฺตํ อภิวาเทตฺวา เอกมนฺตํ นิสีทิํสุ, เอกมนฺตํ นิสินฺนา โข เต ภิกฺขู ภควนฺตํ เอตทโวจุํ\csromandash{}}
\prose{\csromanfolio{130}เอตรหิ ภนฺเต สาวตฺถิยํ มนุสฺสา ภิกฺขู ทิสฺวา อสพฺภาหิ ผรุสาหิ วาจาหิ อกฺโกสนฺติ ปริภาสนฺติ โรสนฺติ วิเหสนฺติ “อลชฺชิโน อิเม สมณา สกฺยปุตฺติยา ทุสฺสีลา ปาปธมฺมา มุสาวาทิโน อพฺรหฺมจาริโน, อิเม หิ นาม ธมฺมจาริโน สมจาริโน พฺรหฺมจาริโน สจฺจวาทิโน สีลวนฺโต กลฺยาณธมฺมา ปฏิชานิสฺสนฺติ, นตฺถิ อิเมสํ สามญฺญํ, นตฺถิ อิเมสํ พฺรหฺมญฺญํ, นฏฺฐํ อิเมสํ สามญฺญํ, นฏฺฐํ อิเมสํ พฺรหฺมญฺญํ, กุโต อิเมสํ สามญฺญํ, กุโต อิเมสํ พฺรหฺมญฺญํ, อปคตา อิเม สามญฺญา, อปคตา อิเม พฺรหฺมญฺญา, กถํ หิ นาม ปุริโส ปุริสกิจฺจํ กริตฺวา อิตฺถิํ ชีวิตา โวโรเปสฺสตี”ติ.}

\prose{เนโส ภิกฺขเว สทฺโท จิรํ ภวิสฺสติ, สตฺตาหเมว ภวิสฺสติ, สตฺตาหสฺส อจฺจเยน อนฺตรธายิสฺสติ.\csromanspacer{} เตน หิ ภิกฺขเว เย มนุสฺสา ภิกฺขู ทิสฺวา อสพฺภาหิ ผรุสาหิ วาจาหิ อกฺโกสนฺติ ปริภาสนฺติ โรสนฺติ วิเหสนฺติ, เต ตุเมฺห อิมาย คาถาย ปฏิโจเทถ\csromandash{}}

\csromangathameasure{“อภูตวาที นิรยํ อุเปติ, \\ โย วาปิ กตฺวา น กโรมิ’จาห. \\ อุโภปิ เต เปจฺจ สมา ภวนฺติ, \\ นิหีนกมฺมา มนุชา ปรตฺถา”ติ.}
\csromangathagroup{\csromangathastanza{“อภูตวาที นิรยํ อุเปติ, \\ โย วาปิ\footnote{โย จาปิ (สี, อิ, ก)} กตฺวา น กโรมิ’จาห. \\ อุโภปิ เต เปจฺจ สมา ภวนฺติ, \\ นิหีนกมฺมา มนุชา ปรตฺถา”ติ.}}

\prose{อถ โข เต ภิกฺขู ภควโต สนฺติเก อิมํ คาถํ ปริยาปุณิตฺวา เย มนุสฺสา ภิกฺขู ทิสฺวา อสพฺภาหิ ผรุสาหิ วาจาหิ อกฺโกสนฺติ ปริภาสนฺติ โรสนฺติ วิเหสนฺติ, เต อิมาย คาถาย ปฏิโจเทนฺติ\csromandash{}}

\csromanmatikaanchor{mtk.69}
\csromantocmarkref{subsection}{8. ปิณฺฑปาติกสุตฺต}{mtk.69}
\bookmark[dest={mtk.69},level=2]{8. ปิณฺฑปาติกสุตฺต}
\csromangathameasure{“อภูตวาที นิรยํ อุเปติ, \\ โย วาปิ กตฺวา น กโรมิ’จาห. \\ อุโภปิ เต เปจฺจ สมา ภวนฺติ, \\ นิหีนกมฺมา มนุชา ปรตฺถา”ติ.}
\csromangathagroup{\csromangathastanza{“อภูตวาที นิรยํ อุเปติ, \\ โย วาปิ กตฺวา น กโรมิ’จาห. \\ อุโภปิ เต เปจฺจ สมา ภวนฺติ, \\ นิหีนกมฺมา มนุชา ปรตฺถา”ติ.}}

\chapterheadpageread{4. เมฆิยวคฺค}
\prose{\csromanfolio{131}มนุสฺสานํ เอตทโหสิ “อการกา อิเม สมณา สกฺยปุตฺติยา, นยิเมหิ กตํ, สปนฺติเม สมณา สกฺยปุตฺติยา”ติ, เนว โส สทฺโท จิรํ อโหสิ, สตฺตาหเมว อโหสิ, สตฺตาหสฺส อจฺจเยน อนฺตรธายิ.}

\prose{อถ โข สมฺพหุลา ภิกฺขู เยน ภควา เตนุปสงฺกมิํสุ, อุปสงฺกมิตฺวา ภควนฺตํ อภิวาเทตฺวา เอกมนฺตํ นิสีทิํสุ, เอกมนฺตํ นิสินฺนา โข เต ภิกฺขู ภควโต เอตทโวจุํ\csromandash{}}

\prose{อจฺฉริยํ ภนฺเต อพฺภุตํ ภนฺเต, ยาว สุภาสิตํ จิทํ ภนฺเต ภควตา “เนโส ภิกฺขเว สทฺโท จิรํ ภวิสฺสติ, สตฺตาหเมว ภวิสฺสติ, สตฺตาหสฺส อจฺจเยน อนฺตรธายิสฺสตี”ติ, อนฺตรหิโต โส ภนฺเต สทฺโทติ.}

\prose{อถ โข ภควา เอตมตฺถํ วิทิตฺวา ตายํ เวลายํ อิมํ อุทานํ อุทาเนสิ\csromandash{}}

\csromangathameasure{“ตุทนฺติ วาจาย ชนา อสญฺญตา, \\ สเรหิ สงฺคามคตํว กุญฺชรํ. \\ สุตฺวาน วากฺยํ ผรุสํ อุทีริตํ, \\ อธิวาสเย ภิกฺขุ อทุฏฺฐจิตฺโต”ติ.อฏฺฐมํ.}
\csromangathagroup{\csromangathastanza{“ตุทนฺติ วาจาย ชนา อสญฺญตา, \\ สเรหิ สงฺคามคตํว กุญฺชรํ. \\ สุตฺวาน วากฺยํ ผรุสํ อุทีริตํ, \\ อธิวาสเย ภิกฺขุ อทุฏฺฐจิตฺโต”ติ.\csromanspacer{}\csromanspacer{}อฏฺฐมํ.}}

\csromanheader{2}{9. อุปเสนสุตฺต}
\proseitem{39}{เอวํ เม สุตํ\csromandash{}เอกํ สมยํ ภควา ราชคเห วิหรติ เวฬุวเน กลนฺทกนิวาเป.\csromanspacer{} อถ โข อายสฺมโต อุปเสนสฺส วงฺคนฺตปุตฺตสฺส รโหคตสฺส ปฏิสลฺลีนสฺส เอวํ เจตโส ปริวิตกฺโก อุทปาทิ “ลาภา วต เม, สุลทฺธํ วต เม, สตฺถา จ เม ภควา อรหํ สมฺมาสมฺพุทฺโธ, สฺวากฺขาเต จมฺหิ ธมฺมวินเย อคารสฺมา อนคาริยํ ปพฺพชิโต, สพฺรหฺมจาริโน จ เม สีลวนฺโต กลฺยาณธมฺมา, สีเลสุ จมฺหิ ปริปูรการี, สุสมาหิโต จมฺหิ เอกคฺคจิตฺโต, อรหา จมฺหิ ขีณาสโว, มหิทฺธิโก จมฺหิ มหานุภาโว, ภทฺทกํ เม ชีวิตํ ภทฺทกํ มรณนฺ”ติ.}

\prose{อถ โข ภควา อายสฺมโต อุปเสนสฺส วงฺคนฺตปุตฺตสฺส เจตสา เจโตปริวิตกฺกมญฺญาย ตายํ เวลายํ อิมํ อุทานํ อุทาเนสิ\csromandash{}}

\gambhira{อุทานปาฬิ}
\csromanmatikaanchor{mtk.70}
\csromantocmarkref{subsection}{9. สิปฺปสุตฺต}{mtk.70}
\bookmark[dest={mtk.70},level=2]{9. สิปฺปสุตฺต}
\csromangathameasure{“ยํ ชีวิตํ น ตปติ, มรณนฺเต น โสจติ. \\ ส เว ทิฏฺฐปโท ธีโร, โสกมชฺเฌ น โสจติ. \\ อุจฺฉินฺนภวตณฺหสฺส, สนฺตจิตฺตสฺส ภิกฺขุโน. \\ วิกฺขีโณ ชาติสํสาโร, นตฺถิ ตสฺส ปุนพฺภโว”ติ.นวมํ.}
\csromangathagroup{\csromangathastanza{\csromanfolio{132}“ยํ ชีวิตํ น ตปติ, มรณนฺเต น โสจติ. \\ ส เว ทิฏฺฐปโท ธีโร, โสกมชฺเฌ น โสจติ.}\csromangathastanza{อุจฺฉินฺนภวตณฺหสฺส, สนฺตจิตฺตสฺส ภิกฺขุโน. \\ วิกฺขีโณ ชาติสํสาโร, นตฺถิ ตสฺส ปุนพฺภโว”ติ.\csromanspacer{}\csromanspacer{}นวมํ.}}

\csromanheader{2}{10. สาริปุตฺต-อุปสมสุตฺต}
\proseitem{40}{เอวํ เม สุตํ\csromandash{}เอกํ สมยํ ภควา สาวตฺถิยํ วิหรติ เชตวเน อนาถปิณฺฑิกสฺส อาราเม.\csromanspacer{} เตน โข ปน สมเยน อายสฺมา สาริปุตฺโต ภควโต อวิทูเร นิสินฺโน โหติ ปลฺลงฺกํ อาภุชิตฺวา อุชุํ กายํ ปณิธาย อตฺตโน อุปสมํ ปจฺจเวกฺขมาโน.}

\prose{อทฺทสา โข ภควา อายสฺมนฺตํ สาริปุตฺตํ อวิทูเร นิสินฺนํ ปลฺลงฺกํ อาภุชิตฺวา อุชุํ กายํ ปณิธาย อตฺตโน อุปสมํ ปจฺจเวกฺขมานํ.}

\prose{อถ โข ภควา เอตมตฺถํ วิทิตฺวา ตายํ เวลายํ อิมํ อุทานํ อุทาเนสิ\csromandash{}}

\csromangathameasure{“อุปสนฺตสนฺตจิตฺตสฺส, เนตฺติจฺฉินฺนสฺส ภิกฺขุโน. \\ วิกฺขีโณ ชาติสํสาโร, มุตฺโต โส มารพนฺธนา”ติ. ทสมํ. เมฆิยวคฺโค จตุตฺโถ.}
\csromangathagroup{\csromangathastanza{“อุปสนฺตสนฺตจิตฺตสฺส, เนตฺติจฺฉินฺนสฺส ภิกฺขุโน. \\ วิกฺขีโณ ชาติสํสาโร, มุตฺโต โส มารพนฺธนา”ติ.\csromanspacer{} ทสมํ.\csromanspacer{} เมฆิยวคฺโค จตุตฺโถ.}}

\titlehead{\textbf{ตสฺสุทฺทานํ}}
\csromangathameasure{เมฆิโย อุทฺธตา โคปาโล, ยกฺโข นาเคน ปญฺจมํ. \\ ปิณฺโฑโล สาริปุตฺโต จ, สุนฺทรี ภวติ อฏฺฐมํ. \\ อุปเสโน วงฺคนฺตปุตฺโต, สาริปุตฺโต จ เต ทสาติ.}
\csromangathagroup{\csromangathastanza{เมฆิโย อุทฺธตา โคปาโล, ยกฺโข\footnote{ชุณฺหา (สี, สฺยา, อิ), ชุณฺหํ (ก)} นาเคน ปญฺจมํ. \\ ปิณฺโฑโล สาริปุตฺโต จ, สุนฺทรี ภวติ อฏฺฐมํ. \\ อุปเสโน วงฺคนฺตปุตฺโต, สาริปุตฺโต จ เต ทสาติ.}}

\chapterheadpageread{5. โสณวคฺค \textbf{5. โสณวคฺค}\footnote{มหาวคฺค (อฏฺฐกถาย สเมติ)}}
\csromanmatikaanchor{mtk.71}
\csromantocmarkref{subsection}{10. โลกสุตฺต}{mtk.71}
\bookmark[dest={mtk.71},level=2]{10. โลกสุตฺต}
\csromanheader{2}{1. ปิยตรสุตฺต}
\proseitem{41}{\csromanfolio{133}เอวํ เม สุตํ\csromandash{}เอกํ สมยํ ภควา สาวตฺถิยํ วิหรติ เชตวเน อนาถปิณฺฑิกสฺส อาราเม.\csromanspacer{} เตน โข ปน สมเยน ราชา ปเสนทิ โกสโล มลฺลิกาย เทวิยา สทฺธิํ อุปริปาสาทวรคโต โหติ.\csromanspacer{} อถ โข ราชา ปเสนทิ โกสโล มลฺลิกํ เทวิํ เอตทโวจ “อตฺถิ นุ โข เต มลฺลิเก โกจญฺโญ อตฺตนา ปิยตโร”ติ.}

\prose{นตฺถิ โข เม มหาราช โกจญฺโญ อตฺตนา ปิยตโร, ตุยฺหํ ปน มหาราช อตฺถญฺโญ โกจิ อตฺตนา ปิยตโรติ.\csromanspacer{} มยฺหมฺปิ โข มลฺลิเก นตฺถญฺโญ โกจิ อตฺตนา ปิยตโรติ.}

\prose{อถ โข ราชา ปเสนทิ โกสโล ปาสาทา โอโรหิตฺวา เยน ภควา เตนุปสงฺกมิ, อุปสงฺกมิตฺวา ภควนฺตํ อภิวาเทตฺวา เอกมนฺตํ นิสีทิ, เอกมนฺตํ นิสินฺโน โข ราชา ปเสนทิ โกสโล ภควนฺตํ เอตทโวจ\csromandash{}}

\prose{อิธาหํ ภนฺเต มลฺลิกาย เทวิยา สทฺธิํ อุปริปาสาทวรคโต มลฺลิกํ เทวิํ เอตทโวจํ “อตฺถิ นุ โข เต มลฺลิเก โกจญฺโญ อตฺตนา ปิยตโร”ติ.\csromanspacer{} เอวํ วุตฺเต มลฺลิกา เทวี มํ เอตทโวจ “นตฺถิ โข เม มหาราช โกจญฺโญ อตฺตนา ปิยตโร.\csromanspacer{} ตุยฺหํ ปน มหาราช อตฺถญฺโญ โกจิ อตฺตนา ปิยตโร”ติ.\csromanspacer{} เอวํ วุตฺเต อหํ ภนฺเต มลฺลิกํ เทวิํ เอตทโวจํ “มยฺหมฺปิ โข มลฺลิเก นตฺถญฺโญ โกจิ อตฺตนา ปิยตโร”ติ.}

\prose{อถ โข ภควา เอตมตฺถํ วิทิตฺวา ตายํ เวลายํ อิมํ อุทานํ อุทาเนสิ\csromandash{}}

\csromangathameasure{*“สพฺพา ทิสา อนุปริคมฺม เจตสา, \\ เนวชฺฌคา ปิยตรมตฺตนา กฺวจิ. \\ เอวํ ปิโย ปุถุ อตฺตา ปเรสํ, \\ ตสฺมา น หิํเส ปร’มตฺตกาโม”ติ.ปฐมํ.}
\csromangathagroup{\csromangathastanza{\csromansymbolfootnote{*}{ขุ 10. 142 ปิฏฺเฐปิ.}“สพฺพา ทิสา อนุปริคมฺม เจตสา, \\ เนวชฺฌคา ปิยตรมตฺตนา กฺวจิ. \\ เอวํ ปิโย ปุถุ อตฺตา ปเรสํ, \\ ตสฺมา น หิํเส ปร’มตฺตกาโม”ติ.\csromanspacer{}\csromanspacer{}ปฐมํ.}}

\csromanheader{2}{อุทานปาฬิ \textbf{2. อปฺปายุกสุตฺต}}
\proseitem{42}{\csromanfolio{134}เอวํ เม สุตํ\csromandash{}เอกํ สมยํ ภควา สาวตฺถิยํ วิหรติ เชตวเน อนาถปิณฺฑิกสฺส อาราเม.\csromanspacer{} อถ โข อายสฺมา อานนฺโท สายนฺหสมยํ ปฏิสลฺลานา\footnote{ปฏิสลฺลาณา (สี)} วุฏฺฐิโต เยน ภควา เตนุปสงฺกมิ, อุปสงฺกมิตฺวา ภควนฺตํ อภิวาเทตฺวา เอกมนฺตํ นิสีทิ, เอกมนฺตํ นิสินฺโน โข อายสฺมา อานนฺโท ภควนฺตํ เอตทโวจ “อจฺฉริยํ ภนฺเต อพฺภุตํ ภนฺเต, ยาว อปฺปายุกา หิ ภนฺเต ภควโต มาตา อโหสิ, สตฺตาหชาเต ภควติ ภควโต มาตา กาลมกาสิ, ตุสิตํ กายํ อุปปชฺชี”ติ.}

\prose{เอวเมตํ อานนฺท\footnote{เอวเมตํ อานนฺท เอวเมตํ อานนฺท (สฺยา)}, อปฺปายุกา หิ อานนฺท โพธิสตฺตมาตโร โหนฺติ, สตฺตาหชาเตสุ โพธิสตฺเตสุ โพธิสตฺตมาตโร กาลํ กโรนฺติ, ตุสิตํ กายํ อุปปชฺชนฺตีติ.}

\prose{อถ โข ภควา เอตมตฺถํ วิทิตฺวา ตายํ เวลายํ อิมํ อุทานํ อุทาเนสิ\csromandash{}}

\csromangathameasure{*“เย เกจิ ภูตา ภวิสฺสนฺติ เย วาปิ, \\ สพฺเพ คมิสฺสนฺติ ปหาย เทหํ. \\ ตํ สพฺพชานิํ กุสโล วิทิตฺวา, \\ อาตาปิโย พฺรหฺมจริยํ จเรยฺยา”ติ.ทุติยํ.}
\csromangathagroup{\csromangathastanza{\csromansymbolfootnote{*}{ขุ 10. 142 ปิฏฺเฐปิ.}“เย เกจิ ภูตา ภวิสฺสนฺติ เย วาปิ, \\ สพฺเพ คมิสฺสนฺติ ปหาย เทหํ. \\ ตํ สพฺพชานิํ กุสโล วิทิตฺวา, \\ อาตาปิโย พฺรหฺมจริยํ จเรยฺยา”ติ.\csromanspacer{}\csromanspacer{}ทุติยํ.}}

\csromanheader{2}{3. สุปฺปพุทฺธกุฏฺฐิสุตฺต}
\proseitem{43}{เอวํ เม สุตํ\csromandash{}เอกํ สมยํ ภควา ราชคเห วิหรติ เวฬุวเน กลนฺทกนิวาเป.\csromanspacer{} เตน โข ปน สมเยน ราชคเห สุปฺปพุทฺโธ นาม กุฏฺฐี อโหสิ มนุสฺสทลิทฺโท มนุสฺสกปโณ มนุสฺสวราโก.\csromanspacer{} เตน โข ปน สมเยน ภควา มหติยา ปริสาย ปริวุโต ธมฺมํ เทเสนฺโต นิสินฺโน โหติ.}

\prose{อทฺทสา โข สุปฺปพุทฺโธ กุฏฺฐี ตํ มหาชนกายํ ทูรโตว สนฺนิปติตํ, ทิสฺวานสฺส เอตทโหสิ “นิสฺสํสยํ โข เอตฺถ กิญฺจิ ขาทนียํ วา โภชนียํ วา ภาชียติ\footnote{ภาชียิสฺสติ (สี)}, ยํนูนาหํ เยน โส มหาชนกาโย เตนุปสงฺกเมยฺยํ, อปฺเปว นาเมตฺถ กิญฺจิ ขาทนียํ วา โภชนียํ วา ลเภยฺยนฺ”ติ.}

\csromanmatikaanchor{mtk.72}
\csromantocmarknopage{chapter}{4. เมฆิยวคฺค}{mtk.72}
\bookmark[dest={mtk.72},level=0]{4. เมฆิยวคฺค}
\chapterheadpageread{5. โสณวคฺค}
\csromanmatikaanchor{mtk.73}
\csromantocmarkref{subsection}{1. เมฆิยสุตฺต}{mtk.73}
\bookmark[dest={mtk.73},level=2]{1. เมฆิยสุตฺต}
\prose{\csromanfolio{135}อถ โข สุปฺปพุทฺโธ กุฏฺฐี เยน โส มหาชนกาโย เตนุปสงฺกมิ.\csromanspacer{} อทฺทสา โข สุปฺปพุทฺโธ กุฏฺฐี ภควนฺตํ มหติยา ปริสาย ปริวุตํ ธมฺมํ เทเสนฺตํ นิสินฺนํ, ทิสฺวานสฺส เอตทโหสิ “น โข เอตฺถ กิญฺจิ ขาทนียํ วา โภชนียํ วา ภาชียติ, สมโณ อยํ โคตโม ปริสติ ธมฺมํ เทเสติ, ยํนูนาหมฺปิ ธมฺมํ สุเณยฺยนฺ”ติ ตตฺเถว เอกมนฺตํ นิสีทิ “อหมฺปิ ธมฺมํ โสสฺสามี”ติ.}

\prose{อถ โข ภควา สพฺพาวนฺตํ ปริสํ เจตสา เจโต ปริจฺจ มนสากาสิ “โก นุ โข อิธ ภพฺโพ ธมฺมํ วิญฺญาตุนฺ”ติ.\csromanspacer{} อทฺทสา โข ภควา สุปฺปพุทฺธํ กุฏฺฐิํ ตสฺสํ ปริสายํ นิสินฺนํ, ทิสฺวานสฺส เอตทโหสิ “อยํ โข อิธ ภพฺโพ ธมฺมํ วิญฺญาตุนฺ”ติ สุปฺปพุทฺธํ กุฏฺฐิํ อารพฺภ อนุปุพฺพิํ กถํ\footnote{อานุปุพฺพิกถํ (สี), อนุปุพฺพิกถํ (สฺยา, อิ, ก)} กเถสิ.\csromanspacer{} เสยฺยถิทํ, ทานกถํ สีลกถํ สคฺคกถํ กามานํ อาทีนวํ โอการํ สํกิเลสํ เนกฺขมฺเม\footnote{เนกฺขมฺเมจ (สี, สฺยา, อิ)} อานิสํสํ ปกาเสสิ.\csromanspacer{} ยทา ภควา อญฺญาสิ สุปฺปพุทฺธํ กุฏฺฐิํ กลฺลจิตฺตํ มุทุจิตฺตํ วินีวรณจิตฺตํ อุทคฺคจิตฺตํ ปสนฺนจิตฺตํ.\csromanspacer{} อถ ยา พุทฺธานํ สามุกฺกํสิกา ธมฺมเทสนา, ตํ ปกาเสสิ ทุกฺขํ สมุทยํ นิโรธํ มคฺคํ.\csromanspacer{} เสยฺยถาปิ นาม สุทฺธํ วตฺถํ อปคตกาฬกํ สมฺมเทว รชนํ ปฏิคฺคเณฺหยฺย, เอวเมวํ สุปฺปพุทฺธสฺส กุฏฺฐิสฺส ตสฺมิํ เยว อาสเน วิรชํ วีตมลํ ธมฺมจกฺขุํ อุทปาทิ “ยํ กิญฺจิ สมุทยธมฺมํ, สพฺพํ ตํ นิโรธธมฺมนฺ”ติ.}

\prose{อถ โข สุปฺปพุทฺโธ กุฏฺฐี ทิฏฺฐธมฺโม ปตฺตธมฺโม วิทิตธมฺโม ปริโยคาฬฺหธมฺโม ติณฺณวิจิกิจฺโฉ วิคตกถํกโถ เวสารชฺชปฺปตฺโต อปรปฺปจฺจโย สตฺถุ สาสเน อุฏฺฐายาสนา เยน ภควา เตนุปสงฺกมิ, อุปสงฺกมิตฺวา ภควนฺตํ อภิวาเทตฺวา เอกมนฺตํ นิสีทิ, เอกมนฺตํ นิสินฺโน โข สุปฺปพุทฺโธ กุฏฺฐี ภควนฺตํ เอตทโวจ\csromandash{}}

\prose{อภิกฺกนฺตํ ภนฺเต อภิกฺกนฺตํ ภนฺเต, เสยฺยถาปิ ภนฺเต นิกฺกุชฺชิตํ วา อุกฺกุชฺเชยฺย, ปฏิจฺฉนฺนํ วา วิวเรยฺย, มูฬฺหสฺส วา มคฺคํ อาจิกฺเขยฺย, อนฺธกาเร วา เตลปชฺโชตํ ธาเรยฺย “จกฺขุมนฺโต รูปานิ ทกฺขนฺตี”ติ.\csromanspacer{} เอวเมวํ ภควตา อเนกปริยาเยน ธมฺโม ปกาสิโต, เอสาหํ ภนฺเต ภควนฺตํ สรณํ คจฺฉามิ ธมฺมญฺจ ภิกฺขุสํฆญฺจ, อุปาสกํ มํ ภควา ธาเรตุ “อชฺชตคฺเค ปาณุเปตํ สรณํ คตนฺ”ติ.}

\gambhira{อุทานปาฬิ}
\prose{\csromanfolio{136}อถ โข สุปฺปพุทฺโธ กุฏฺฐี ภควตา ธมฺมิยา กถาย สนฺทสฺสิโต สเมทปิโต สมุตฺเตชิโต สมฺปหํสิโต ภควตา ภาสิตํ อภินนฺทิตฺวา อนุโมทิตฺวา อุฏฺฐายาสนา ภควนฺตํ อภิวาเทตฺวา ปทกฺขิณํ กตฺวา ปกฺกามิ.\csromanspacer{} อถ โข อจิรปกฺกนฺตํ สุปฺปพุทฺธํ กุฏฺฐิํ คาวี ตรุณวจฺฉา อธิปติตฺวา ชีวิตา โวโรเปสิ.}

\prose{อถ โข สมฺพหุลา ภิกฺขู เยน ภควา เตนุปสงฺกมิํสุ, อุปสงฺกมิตฺวา ภควนฺตํ อภิวาเทตฺวา เอกมนฺตํ นิสีทิํสุ, เอกมนฺตํ นิสินฺนา โข เต ภิกฺขู ภควนฺตํ เอตทโวจุํ “โย โส ภนฺเต สุปฺปพุทฺโธ นาม กุฏฺฐี ภควตา ธมฺมิยา กถาย สนฺทสฺสิโต สมาทปิโต สมุตฺเตชิโต สมฺปหํสิโต, โส กาลงฺกโต, ตสฺส กา คติ, โก อภิสมฺปราโย”ติ.}

\prose{ปณฺฑิโต ภิกฺขเว สุปฺปพุทฺโธ กุฏฺฐี ปจฺจปาทิ ธมฺมสฺสานุธมฺมํ, น จ มํ ธมฺมาธิกรณํ วิเหเสสิ, สุปฺปพุทฺโธ ภิกฺขเว กุฏฺฐี ติณฺณํ สํโยชนานํ ปริกฺขยา โสตาปนฺโน อวินิปาตธมฺโม นิยโต สมฺโพธิปรายโณติ.}

\prose{เอวํ วุตฺเต อญฺญตโร ภิกฺขุ ภควนฺตํ เอตทโวจ “โก นุ โข ภนฺเต เหตุ โก ปจฺจโย, เยน สุปฺปพุทฺโธ กุฏฺฐี อโหสิ มนุสฺสทลิทฺโท มนุสฺสกปโณ มนุสฺสวราโก”ติ.}

\prose{ภูตปุพฺพํ ภิกฺขเว สุปฺปพุทฺโธ กุฏฺฐี อิมสฺมิํเยว ราชคเห เสฏฺฐิปุตฺโต อโหสิ, โส อุยฺยานภูมิํ นิยฺยนฺโต อทฺทส ตครสิขิํ\footnote{ตคฺครสิขิํ (ก)} ปจฺเจกพุทฺธํ นครํ ปิณฺฑาย ปวิสนฺตํ, ทิสฺวานสฺส เอตทโหสิ “กฺวายํ กุฏฺฐี กุฏฺฐิจีวเรน วิจรตี”ติ นิฏฺฐุภิตฺวา อปสพฺยโต\footnote{อปพฺยามโต (สฺยา, สํ 1. 228 ปิฏฺเฐปิ)} กริตฺวา ปกฺกามิ.\csromanspacer{} โส ตสฺส กมฺมสฺส วิปาเกน พหูนิ วสฺสสตานิ พหูนิ วสฺสสหสฺสานิ พหูนิ วสฺสสตสหสฺสานิ นิรเย ปจฺจิตฺถ, ตสฺเสว กมฺมสฺส วิปากาวเสเสน อิมสฺมิํเยว ราชคเห กุฏฺฐี อโหสิ มนุสฺสทลิทฺโท มนุสฺสกปโณ มนุสฺสวราโก.\csromanspacer{} โส ตถาคตปฺปเวทิตํ ธมฺมวินยํ อาคมฺม สทฺธํ สมาทิยิ, สีลํ สมาทิยิ, สุตํ สมาทิยิ, จาคํ สมาทิยิ, ปญฺญํ สมาทิยิ.\csromanspacer{} โส ตถาคตปฺปเวทิตํ ธมฺมวินยํ อาคมฺม สทฺธํ สมาทิยิตฺวา สีลํ สมาทิยิตฺวา สุตํ สมาทิยิตฺวา จาคํ สมาทิยิตฺวา ปญฺญํ สมาทิยิตฺวา}

\vspace{\tipitakaparskip}
\csromancenter{โสณวคฺค กายสฺส เภทา ปรํ มรณา สุคติํ สคฺคํ โลกํ อุปปนฺโน เทวานํ ตาวติํสานํ สหพฺยตํ, โส ตตฺถ อญฺเญ เทเว อติโรจติ วณฺเณน เจว ยสสา จาติ.}
\prose{\csromanfolio{137}อถ โข ภควา เอตมตฺถํ วิทิตฺวา ตายํ เวลายํ อิมํ อุทานํ อุทาเนสิ\csromandash{}}

\csromangathameasure{“จกฺขุมา วิสมานีว, วิชฺชมาเน ปรกฺกเม*. \\ ปณฺฑิโต ชีวโลกสฺมิํ, ปาปานิ ปริวชฺชเย”ติ.ตติยํ.}
\csromangathagroup{\csromangathastanza{“จกฺขุมา วิสมานีว, วิชฺชมาเน ปรกฺกเม*. \\ ปณฺฑิโต ชีวโลกสฺมิํ, ปาปานิ ปริวชฺชเย”ติ.\csromanspacer{}\csromanspacer{}ตติยํ.}}

\csromanheader{2}{4. กุมารกสุตฺต}
\proseitem{44}{เอวํ เม สุตํ\csromandash{}เอกํ สมยํ ภควา สาวตฺถิยํ วิหรติ เชตวเน อนาถปิณฺฑิกสฺส อาราเม.\csromanspacer{} เตน โข ปน สมเยน สมฺพหุลา กุมารกา อนฺตรา จ สาวตฺถิํ อนฺตรา จ เชตวนํ มจฺฉเก พาเธนฺติ.}

\prose{อถ โข ภควา ปุพฺพณฺหสมยํ นิวาเสตฺวา ปตฺตจีวรมาทาย สาวตฺถิํ ปิณฺฑาย ปาวิสิ, อทฺทสา โข ภควา เต สมฺพหุเล กุมารเก อนฺตรา จ สาวตฺถิํ อนฺตรา จ เชตวนํ มจฺฉเก พาเธนฺเต, ทิสฺวาน เยน เต กุมารกา เตนุปสงฺกมิ, อุปสงฺกมิตฺวา เต กุมารเก เอตทโวจ “ภายถ โว ตุเมฺห กุมารกา ทุกฺขสฺส, อปฺปิยํ โว ทุกฺขนฺ”ติ.\csromanspacer{} เอวํ ภนฺเต, ภายาม มยํ ภนฺเต ทุกฺขสฺส, อปฺปิยํ โน ทุกฺขนฺติ.}

\prose{อถ โข ภควา เอตมตฺถํ วิทิตฺวา ตายํ เวลายํ อิมํ อุทานํ อุทาเนสิ\csromandash{}}

\prose{\csromansymbolfootnote{+}{ขุ 10. 109 ปิฏฺเฐปิ.}“สเจ ภายถ ทุกฺขสฺส, สเจ โว ทุกฺขมปฺปิยํ.}

\csromangathameasure{มากตฺถ ปาปกํ กมฺมํ, อาวิ วา ยทิ วา รโห. \\ สเจ จ ปาปกํ กมฺมํ, กริสฺสถ กโรถ วา. \\ น โว ทุกฺขา ปมุตฺยตฺถิ, อุเปจฺจปิ ปลายตนฺ”ติ.จตุตฺถํ.}
\csromangathagroup{\csromangathastanza{มากตฺถ ปาปกํ กมฺมํ, อาวิ วา ยทิ วา รโห. \\ สเจ จ ปาปกํ กมฺมํ, กริสฺสถ กโรถ วา. \\ น โว ทุกฺขา ปมุตฺยตฺถิ, อุเปจฺจปิ\footnote{อุปจฺจปิ (ก), อุปฺปจฺจปิ (?), อุปฺปติตฺวาปิ อิติ อตฺโถ.} ปลายตนฺ”ติ.\csromanspacer{}\csromanspacer{}จตุตฺถํ.}}

\csromanheader{2}{อุทานปาฬิ \textbf{5. อุโปสถสุตฺต}}
\proseitem{45}{\csromanfolio{138}เอวํ เม สุตํ\csromandash{}เอกํ สมยํ ภควา สาวตฺถิยํ วิหรติ ปุพฺพาราเม มิคารมาตุปาสาเท.\csromanspacer{} เตน โข ปน สมเยน ภควา ตทหุโปสเถ ภิกฺขุสํฆปริวุโต นิสินฺโน โหติ.}

\prose{อถ โข อายสฺมา อานนฺโท อภิกฺกนฺตาย รตฺติยา นิกฺขนฺเต ปฐเม ยาเม อุฏฺฐายาสนา เอกํสํ อุตฺตราสงฺคํ\footnote{จีวรํ (สพฺพตฺถ)} กริตฺวา เยน ภควา เตนญฺชลิํ ปณาเมตฺวา ภควนฺตํ เอตทโวจ “อภิกฺกนฺตา ภนฺเต รตฺติ, นิกฺขนฺโต ปฐโม ยาโม, จิรนิสินฺโน ภิกฺขุสํโฆ, อุทฺทิสตุ ภนฺเต ภควา ภิกฺขูนํ ปาติโมกฺขนฺ”ติ.\csromanspacer{} เอวํ วุตฺเต ภควา ตุณฺหี อโหสิ.}

\prose{ทุติยมฺปิ โข อายสฺมา อานนฺโท อภิกฺกนฺตาย รตฺติยา นิกฺขนฺเต มชฺฌิเม ยาเม อุฏฺฐายาสนา เอกํสํ อุตฺตราสงฺคํ กริตฺวา เยน ภควา เตนญฺชลิํ ปณาเมตฺวา ภควนฺตํ เอตทโวจ “อภิกฺกนฺตา ภนฺเต รตฺติ, นิกฺขนฺโต มชฺฌิโม ยาโม, จิรนิสินฺโน ภิกฺขุสํโฆ, อุทฺทิสตุ ภนฺเต ภควา ภิกฺขูนํ ปาติโมกฺขนฺ”ติ.\csromanspacer{} ทุติยมฺปิ โข ภควา ตุณฺหี อโหสิ.}

\prose{ตติยมฺปิ โข อายสฺมา อานนฺโท อภิกฺกนฺตาย รตฺติยา นิกฺขนฺเต ปจฺฉิเม ยาเม อุทฺธสฺเต อรุเณ นนฺทิมุขิยา รตฺติยา อุฏฺฐายาสนา เอกํสํ อุตฺตราสงฺคํ กริตฺวา เยน ภควา เตนญฺชลิํ ปณาเมตฺวา ภควนฺตํ เอตทโวจ “อภิกฺกนฺตา ภนฺเต รตฺติ, นิกฺขนฺโต ปจฺฉิโม ยาโม, อุทฺธสฺโต อรุโณ, นนฺทิมุขี รตฺติ, จิรนิสินฺโน ภิกฺขุสํโฆ, อุทฺทิสตุ ภนฺเต ภควา ภิกฺขูนํ ปาติโมกฺขนฺ”ติ.\csromanspacer{} อปริสุทฺธา อานนฺท ปริสาติ.}

\prose{อถ โข อายสฺมโต มหาโมคฺคลฺลานสฺส เอตทโหสิ “กํ นุ โข ภควา ปุคฺคลํ สนฺธาย เอวมาห ‘อปริสุทฺธา อานนฺท ปริสา’ติ”.\csromanspacer{} อถ โข อายสฺมา มหาโมคฺคลฺลาโน สพฺพาวนฺตํ ภิกฺขุสํฆํ เจตสา เจโต ปริจฺจ มนสากาสิ, อทฺทสา โข อายสฺมา มหาโมคฺคลฺลาโน ตํ ปุคฺคลํ ทุสฺสีลํ ปาปธมฺมํ อสุจิํ สงฺกสฺสรสมาจารํ ปฏิจฺฉนฺนกมฺมนฺตํ อสฺสมณํ สมณปฏิญฺญํ อพฺรหฺมจาริํ พฺรหฺมจาริปฏิญฺญํ อนฺโตปูติํ อวสฺสุตํ กสมฺพุชาตํ มชฺเฌ ภิกฺขุสํฆสฺส นิสินฺนํ, ทิสฺวาน อุฏฺฐายาสนา เยน โส ปุคฺคโล เตนุปสงฺกมิ, อุปสงฺกมิตฺวา ตํ ปุคฺคลํ}

\vspace{\tipitakaparskip}
\csromancenter{โสณวคฺค เอตทโวจ “อุฏฺเฐหิ อาวุโส, ทิฏฺโฐสิ ภควตา, นตฺถิ เต ภิกฺขูหิ สทฺธิํ สํวาโส”ติ.\csromanspacer{} เอวํ วุตฺเต\footnote{อถ โข (สพฺพตฺถ), วิ 4. 419; อํ 3. 45 ปิฏฺเฐ ปสฺสิตพฺพํ.} โส ปุคฺคโล ตุณฺหี อโหสิ.}
\prose{\csromanfolio{139}ทุติยมฺปิ โข อายสฺมา มหาโมคฺคลฺลาโน ตํ ปุคฺคลํ เอตทโวจ “อุฏฺเฐหิ อาวุโส, ทิฏฺโฐสิ ภควตา, นตฺถิ เต ภิกฺขูหิ สทฺธิํ สํวาโส”ติ.\csromanspacer{} ทุติยมฺปิ โข ฯเปฯ.\csromanspacer{} ตติยมฺปิ โข โส ปุคฺคโล ตุณฺหี อโหสิ.}

\prose{อถ โข อายสฺมา มหาโมคฺคลฺลาโน ตํ ปุคฺคลํ พาหาย คเหตฺวา พหิทฺวารโกฏฺฐกา นิกฺขาเมตฺวา สูจิฆฏิกํ ทตฺวา เยน ภควา เตนุปสงฺกมิ, อุปสงฺกมิตฺวา ภควนฺตํ เอตทโวจ “นิกฺขามิโต ภนฺเต โส ปุคฺคโล มยา, ปริสุทฺธา ปริสา, อุทฺทิสตุ ภนฺเต ภควา ภิกฺขูนํ ปาติโมกฺขนฺ”ติ.\csromanspacer{} อจฺฉริยํ โมคฺคลฺลาน อพฺภุตํ โมคฺคลฺลาน, ยาว พาหาคหณาปิ นาม โส โมฆปุริโส อาคเมสฺสตีติ.}

\prose{อถ โข ภควา ภิกฺขู อามนฺเตสิ\csromandash{}น ทานาหํ ภิกฺขเว อิโต ปรํ\footnote{น ทานาหํ ภิกฺขเว อชฺชตคฺเค (อํ 3. 45 ปิฏฺเฐ)} อุโปสถํ กริสฺสามิ, ปาติโมกฺขํ อุทฺทิสิสฺสามิ.\csromanspacer{} ตุเมฺหว ทานิ ภิกฺขเว อิโต ปรํ อุโปสถํ กเรยฺยาถ, ปาติโมกฺขํ อุทฺทิเสยฺยาถ.\csromanspacer{} อฏฺฐานเมตํ ภิกฺขเว อนวกาโส ยํ ตถาคโต อปริสุทฺธาย ปริสาย อุโปสถํ กเรยฺย, ปาติโมกฺขํ อุทฺทิเสยฺย.}

\prose{อฏฺฐิเม ภิกฺขเว มหาสมุทฺเท อจฺฉริยา อพฺภุตา ธมฺมา, เย ทิสฺวา ทิสฺวา อสุรา มหาสมุทฺเท อภิรมนฺติ.\csromanspacer{} กตเม อฏฺฐ?}

\prose{มหาสมุทฺโท ภิกฺขเว อนุปุพฺพนินฺโน อนุปุพฺพโปโณ อนุปุพฺพปพฺภาโร น อายตเกเนว ปปาโต.\csromanspacer{} ยมฺปิ\footnote{ยํ (สี, สฺยา, ก)} ภิกฺขเว มหาสมุทฺโท อนุปุพฺพนินฺโน อนุปุพฺพโปโณ อนุปุพฺพปพฺภาโร น อายตเกเนว ปปาโต.\csromanspacer{} อยํ ภิกฺขเว มหาสมุทฺเท ปฐโม อจฺฉริโย อพฺภุโต ธมฺโม, ยํ ทิสฺวา ทิสฺวา อสุรา มหาสมุทฺเท อภิรมนฺติ.}

\gambhira{อุทานปาฬิ}
\prose{\csromanfolio{140}ปุน จปรํ ภิกฺขเว มหาสมุทฺโท ฐิตธมฺโม เวลํ นาติวตฺตติ.\csromanspacer{} ยมฺปิ ภิกฺขเว มหาสมุทฺโท ฐิตธมฺโม เวลํ นาติวตฺตติ.\csromanspacer{} อยํ\footnote{อยมฺปิ (สพฺพตฺถ)} ภิกฺขเว มหาสมุทฺเท ทุติโย อจฺฉริโย อพฺภุโต ธมฺโม, ยํ ทิสฺวา ทิสฺวา อสุรา มหาสมุทฺเท อภิรมนฺติ.}

\prose{ปุน จปรํ ภิกฺขเว มหาสมุทฺโท น มเตน กุณเปน สํวสติ.\csromanspacer{} ยํ โหติ มหาสมุทฺเท มตํ กุณปํ, ตํ ขิปฺปเมว\footnote{ขิปฺปญฺเญว (สี), ขิปฺปํเยว (ก)} ตีรํ วาเหติ, ถลํ อุสฺสาเรติ.\csromanspacer{} ยมฺปิ ภิกฺขเว มหาสมุทฺโท น มเตน กุณเปน สํวสติ.\csromanspacer{} ยํ โหติ มหาสมุทฺเท มตํ กุณปํ, ตํ ขิปฺปเมว ตีรํ วาเหติ, ถลํ อุสฺสาเรติ.\csromanspacer{} อยํ ภิกฺขเว มหาสมุทฺเท ตติโย อจฺฉริโย อพฺภุโต ธมฺโม, ยํ ทิสฺวา ทิสฺวา อสุรา มหาสมุทฺเท อภิรมนฺติ.}

\prose{ปุน จปรํ ภิกฺขเว ยา กาจิ มหานทิโย.\csromanspacer{} เสยฺยถิทํ, คงฺคา ยมุนา อจิรวตี สรภู มหี.\csromanspacer{} ตา มหาสมุทฺทํ ปตฺวา\footnote{ปตฺตา (สฺยา, อิ, ก)} ชหนฺติ ปุริมานิ นามโคตฺตานิ, “มหาสมุทฺโท”เตฺวว สงฺขํ คจฺฉนฺติ.\csromanspacer{} ยมฺปิ ภิกฺขเว ยา กาจิ มหานทิโย.\csromanspacer{} เสยฺยถิทํ, คงฺคา ยมุนา อจิรวตี สรภู มหี.\csromanspacer{} ตา มหาสมุทฺทํ ปตฺวา ชหนฺติ ปุริมานิ นามโคตฺตานิ, “มหาสมุทฺโท”เตฺวว สงฺขํ คจฺฉนฺติ.\csromanspacer{} อยํ ภิกฺขเว มหาสมุทฺเท จตุตฺโถ อจฺฉริโย อพฺภุโต ธมฺโม, ยํ ทิสฺวา ทิสฺวา อสุรา มหาสมุทฺเท อภิรมนฺติ.}

\prose{ปุน จปรํ ภิกฺขเว ยา จ โลเก สวนฺติโย มหาสมุทฺทํ อปฺเปนฺติ, ยา จ อนฺตลิกฺขา ธารา ปปตนฺติ, น เตน มหาสมุทฺทสฺส อูนตฺตํ วา ปูรตฺตํ วา ปญฺญายติ.\csromanspacer{} ยมฺปิ ภิกฺขเว ยา จ โลเก สวนฺติโย มหาสมุทฺทํ อปฺเปนฺติ, ยา จ อนฺตลิกฺขา ธารา ปปตนฺติ, น เตน มหาสมุทฺทสฺส อูนตฺตํ วา ปูรตฺตํ วา ปญฺญายติ.\csromanspacer{} อยํ ภิกฺขเว มหาสมุทฺเท ปญฺจโม อจฺฉริโย อพฺภุโต ธมฺโม, ยํ ทิสฺวา ทิสฺวา อสุรา มหาสมุทฺเท อภิรมนฺติ.}

\csromanmatikaanchor{mtk.74}
\csromantocmarkref{subsection}{2. อุทฺธตสุตฺต}{mtk.74}
\bookmark[dest={mtk.74},level=2]{2. อุทฺธตสุตฺต}
\prose{ปุน จปรํ ภิกฺขเว มหาสมุทฺโท เอกรโส โลณรโส.\csromanspacer{} ยมฺปิ ภิกฺขเว มหาสมุทฺโท เอกรโส โลณรโส.\csromanspacer{} อยํ ภิกฺขเว มหาสมุทฺเท ฉฏฺโฐ อจฺฉริโย อพฺภุโต ธมฺโม, ยํ ทิสฺวา ทิสฺวา อสุรา มหาสมุทฺเท อภิรมนฺติ.}

\prose{ปุน จปรํ ภิกฺขเว มหาสมุทฺโท พหุรตโน อเนกรตโน.\csromanspacer{} ตตฺริมานิ รตนานิ, เสยฺยถิทํ, มุตฺตา มณิ เวฬุริโย สงฺโข สิลา}

\vspace{\tipitakaparskip}
\csromancenter{โสณวคฺค ปวาฬํ รชตํ ชาตรูปํ โลหิตงฺโค\footnote{โลหิตงฺโก (สี, อิ), โลหิตโก (?)} มสารคลฺลํ.\csromanspacer{} ยมฺปิ ภิกฺขเว มหาสมุทฺโท พหุรตโน อเนกรตโน.\csromanspacer{} ตตฺริมานิ รตนานิ, เสยฺยถิทํ, มุตฺตา มณิ เวฬุริโย สงฺโข สิลา ปวาฬํ รชตํ ชาตรูปํ โลหิตงฺโค มสารคลฺลํ.\csromanspacer{} อยํ ภิกฺขเว มหาสมุทฺเท สตฺตโม อจฺฉริโย อพฺภุโต ธมฺโม, ยํ ทิสฺวา ทิสฺวา อสุรา มหาสมุทฺเท อภิรมนฺติ.}
\prose{\csromanfolio{141}ปุน จปรํ ภิกฺขเว มหาสมุทฺโท มหตํ ภูตานํ อาวาโส.\csromanspacer{} ตตฺริเม ภูตา, ติมิ ติมิงฺคโล ติมิติมิงฺคโล\footnote{ติมิ ติมิงฺคโล ติมิรปิงฺคโล (สี, อิ, อํ 3. 40)} อสุรา นาคา คนฺธพฺพา.\csromanspacer{} สนฺติ มหาสมุทฺเท โยชนสติกาปิ อตฺตภาวา ทฺวิโยชนสติกาปิ อตฺตภาวา ติโยชนสติกาปิ อตฺตภาวา จตุโยชนสติกาปิ อตฺตภาวา ปญฺจโยชนสติกาปิ อตฺตภาวา.\csromanspacer{} ยมฺปิ ภิกฺขเว มหาสมุทฺโท มหตํ ภูตานํ อาวาโส.\csromanspacer{} ตตฺริเม ภูตา, ติมิ ติมิงฺคโล ติมิติมิงฺคโล อสุรา นาคา คนฺธพฺพา.\csromanspacer{} สนฺติ มหาสมุทฺเท โยชนสติกาปิ อตฺตภาวา ทฺวิโยชนสติกาปิ อตฺตภาวา ฯเปฯ ปญฺจโยชนสติกาปิ อตฺตภาวา.\csromanspacer{} อยํ ภิกฺขเว มหาสมุทฺเท อฏฺฐโม อจฺฉริโย อพฺภุโต ธมฺโม, ยํ ทิสฺวา ทิสฺวา อสุรา มหาสมุทฺเท อภิรมนฺติ.\csromanspacer{} อิเม โข ภิกฺขเว อฏฺฐ มหาสมุทฺเท อจฺฉริยา อพฺภุตา ธมฺมา, เย ทิสฺวา ทิสฺวา อสุรา มหาสมุทฺเท อภิรมนฺติ.}

\prose{เอวเมวํ โข ภิกฺขเว อิมสฺมิํ ธมฺมวินเย อฏฺฐ อจฺฉริยา อพฺภุตา ธมฺมา, เย ทิสฺวา ทิสฺวา ภิกฺขู อิมสฺมิํ ธมฺมวินเย อภิรมนฺติ.\csromanspacer{} กตเม อฏฺฐ?}

\prose{เสยฺยถาปิ ภิกฺขเว มหาสมุทฺโท อนุปุพฺพนินฺโน อนุปุพฺพโปโณ อนุปุพฺพปพฺภาโร น อายตเกเนว ปปาโต.\csromanspacer{} เอวเมวํ โข ภิกฺขเว อิมสฺมิํ ธมฺมวินเย อนุปุพฺพสิกฺขา อนุปุพฺพกิริยา อนุปุพฺพปฏิปทา น อายตเกเนว อญฺญาปฏิเวโธ.\csromanspacer{} ยมฺปิ ภิกฺขเว อิมสฺมิํ ธมฺมวินเย อนุปุพฺพสิกฺขา อนุปุพฺพกิริยา อนุปุพฺพปฏิปทา น อายตเกเนว อญฺญาปฏิเวโธ.\csromanspacer{} อยํ ภิกฺขเว อิมสฺมิํ ธมฺมวินเย ปฐโม อจฺฉริโย อพฺภุโต ธมฺโม, ยํ ทิสฺวา ทิสฺวา ภิกฺขู อิมสฺมิํ ธมฺมวินเย อภิรมนฺติ.}

\gambhira{อุทานปาฬิ}
\prose{\csromanfolio{142}เสยฺยถาปิ ภิกฺขเว มหาสมุทฺโท ฐิตธมฺโม เวลํ นาติวตฺตติ.\csromanspacer{} เอวเมวํ โข ภิกฺขเว ยํ มยา สาวกานํ สิกฺขาปทํ ปญฺญตฺตํ, ตํ มม สาวกา ชีวิตเหตุปิ นาติกฺกมนฺติ.\csromanspacer{} ยมฺปิ ภิกฺขเว มยา สาวกานํ สิกฺขาปทํ ปญฺญตฺตํ, ตํ มม สาวกา ชีวิตเหตุปิ นาติกฺกมนฺติ.\csromanspacer{} อยํ ภิกฺขเว อิมสฺมิํ ธมฺมวินเย ทุติโย อจฺฉริโย อพฺภุโต ธมฺโม, ยํ ทิสฺวา ทิสฺวา ภิกฺขู อิมสฺมิํ ธมฺมวินเย อภิรมนฺติ.}

\csromanmatikaanchor{mtk.75}
\csromantocmarkref{subsection}{3. โคปาลกสุตฺต}{mtk.75}
\bookmark[dest={mtk.75},level=2]{3. โคปาลกสุตฺต}
\prose{เสยฺยถาปิ ภิกฺขเว มหาสมุทฺโท น มเตน กุณเปน สํวสติ.\csromanspacer{} ยํ โหติ มหาสมุทฺเท มตํ กุณปํ, ตํ ขิปฺปเมว ตีรํ วาเหติ, ถลํ อุสฺสาเรติ.\csromanspacer{} เอวเมวํ โข ภิกฺขเว โย โส ปุคฺคโล ทุสฺสีโล ปาปธมฺโม อสุจิ สงฺกสฺสรสมาจาโร ปฏิจฺฉนฺนกมฺมนฺโต อสฺสมโณ สมณปฏิญฺโญ อพฺรหฺมจารี พฺรหฺมจาริปฏิญฺโญ อนฺโตปูติ อวสฺสุโต กสมฺพุชาโต, น เตน สํโฆ สํวสติ, อถ โข นํ ขิปฺปเมว สนฺนิปติตฺวา อุกฺขิปติ.\csromanspacer{} กิญฺจาปิ โส โหติ มชฺเฌ ภิกฺขุสํฆสฺส นิสินฺโน, อถ โข โส อารกาว สํฆมฺหา, สํโฆ จ เตน.\csromanspacer{} ยมฺปิ ภิกฺขเว โย โส ปุคฺคโล ทุสฺสีโล ปาปธมฺโม อสุจิ สงฺกสฺสรสมาจาโร ปฏิจฺฉนฺนกมฺมนฺโต อสฺสมโณ สมณปฏิญฺโญ อพฺรหฺมจารี พฺรหฺมจาริปฏิญฺโญ อนฺโตปูติ อวสฺสุโต กสมฺพุชาโต, น เตน สํโฆ สํวสติ.\csromanspacer{} ขิปฺปเมว นํ สนฺนิปติตฺวา อุกฺขิปติ.\csromanspacer{} กิญฺจาปิ โส โหติ มชฺเฌ ภิกฺขุสํฆสฺส นิสินฺโน, อถ โข โส อารกาว สํฆมฺหา, สํโฆ จ เตน.\csromanspacer{} อยํ ภิกฺขเว อิมสฺมิํ ธมฺมวินเย ตติโย อจฺฉริโย อพฺภุโต ธมฺโม, ยํ ทิสฺวา ทิสฺวา ภิกฺขู อิมสฺมิํ ธมฺมวินเย อภิรมนฺติ.}

\prose{เสยฺยถาปิ ภิกฺขเว ยา กาจิ มหานทิโย.\csromanspacer{} เสยฺยถิทํ, คงฺคา ยมุนา อจิรวตี สรภู มหี.\csromanspacer{} ตา มหาสมุทฺทํ ปตฺวา ชหนฺติ ปุริมานิ นามโคตฺตานิ, “มหาสมุทฺโท”เตฺวว สงฺขํ คจฺฉนฺติ.\csromanspacer{} เอวเมวํ โข ภิกฺขเว จตฺตาโร วณฺณา ขตฺติยา พฺราหฺมณา เวสฺสา สุทฺทา, เต ตถาคตปฺปเวทิเต ธมฺมวินเย อคารสฺมา อนคาริยํ ปพฺพชิตฺวา\footnote{ปพฺพชิตา (ก-สี)} ชหนฺติ ปุริมานิ นามโคตฺตานิ, “สมณา สกฺยปุตฺติยา”เตฺวว สงฺขํ คจฺฉนฺติ.\csromanspacer{} ยมฺปิ ภิกฺขเว จตฺตาโร วณฺณา ขตฺติยา พฺราหฺมณา เวสฺสา สุทฺทา, เต ตถาคตปฺปเวทิเต ธมฺมวินเย อคารสฺมา อนคาริยํ ปพฺพชิตฺวา ชหนฺติ ปุริมานิ}

\vspace{\tipitakaparskip}
\csromancenter{โสณวคฺค นามโคตฺตานิ, “สมณา สกฺยปุตฺติยา”เตฺวว สงฺขํ คจฺฉนฺติ.\csromanspacer{} อยํ ภิกฺขเว อิมสฺมิํ ธมฺมวินเย จตุตฺโถ อจฺฉริโย อพฺภุโต ธมฺโม, ยํ ทิสฺวา ทิสฺวา ภิกฺขู อิมสฺมิํ ธมฺมวินเย อภิรมนฺติ.}
\prose{\csromanfolio{143}เสยฺยถาปิ ภิกฺขเว ยา จ โลเก สวนฺติโย มหาสมุทฺทํ อปฺเปนฺติ, ยา จ อนฺตลิกฺขา ธารา ปปตนฺติ, น เตน มหาสมุทฺทสฺส อูนตฺตํ วา ปูรตฺตํ วา ปญฺญายติ.\csromanspacer{} เอวเมวํ โข ภิกฺขเว พหู เจปิ ภิกฺขู อนุปาทิเสสาย นิพฺพานธาตุยา ปรินิพฺพายนฺติ, น เตน นิพฺพานธาตุยา อูนตฺตํ วา ปูรตฺตํ วา ปญฺญายติ.\csromanspacer{} ยมฺปิ ภิกฺขเว พหู เจปิ ภิกฺขู อนุปาทิเสสาย นิพฺพานธาตุยา ปรินิพฺพายนฺติ, น เตน นิพฺพานธาตุยา อูนตฺตํ วา ปูรตฺตํ วา ปญฺญายติ.\csromanspacer{} อยํ ภิกฺขเว อิมสฺมิํ ธมฺมวินเย ปญฺจโม อจฺฉริโย อพฺภุโต ธมฺโม, ยํ ทิสฺวา ทิสฺวา ภิกฺขู อิมสฺมิํ ธมฺมวินเย อภิรมนฺติ.}

\prose{เสยฺยถาปิ ภิกฺขเว มหาสมุทฺโท เอกรโส โลณรโส.\csromanspacer{} เอวเมวํ โข ภิกฺขเว อยํ ธมฺมวินโย เอกรโส วิมุตฺติรโส.\csromanspacer{} ยมฺปิ ภิกฺขเว อยํ ธมฺมวินโย เอกรโส วิมุตฺติรโส.\csromanspacer{} อยํ ภิกฺขเว อิมสฺมิํ ธมฺมวินเย ฉฏฺโฐ อจฺฉริโย อพฺภุโต ธมฺโม, ยํ ทิสฺวา ทิสฺวา ภิกฺขู อิมสฺมิํ ธมฺมวินเย อภิรมนฺติ.}

\prose{เสยฺยถาปิ ภิกฺขเว มหาสมุทฺโท พหุรตโน อเนกรตโน.\csromanspacer{} ตตฺริมานิ รตนานิ, เสยฺยถิทํ, มุตฺตา มณิ เวฬุริโย สงฺโข สิลา ปวาฬํ รชตํ ชาตรูปํ โลหิตงฺโค มสารคลฺลํ.\csromanspacer{} เอวเมวํ โข ภิกฺขเว อยํ ธมฺมวินโย พหุรตโน อเนกรตโน.\csromanspacer{} ตตฺริมานิ รตนานิ, เสยฺยถิทํ, จตฺตาโร สติปฏฺฐานา จตฺตาโร สมฺมปฺปธานา จตฺตาโร อิทฺธิปาทา ปญฺจินฺทฺริยานิ ปญฺจ พลานิ สตฺต โพชฺฌงฺคา อริโย อฏฺฐงฺคิโก มคฺโค.\csromanspacer{} ยมฺปิ ภิกฺขเว อยํ ธมฺมวินโย พหุรตโน อเนกรตโน.\csromanspacer{} ตตฺริมานิ รตนานิ, เสยฺยถิทํ, จตฺตาโร สติปฏฺฐานา จตฺตาโร สมฺมปฺปธานา จตฺตาโร อิทฺธิปาทา ปญฺจินฺทฺริยานิ ปญฺจ พลานิ สตฺต โพชฺฌงฺคา อริโย อฏฺฐงฺคิโก มคฺโค.\csromanspacer{} อยํ ภิกฺขเว อิมสฺมิํ ธมฺมวินเย สตฺตโม อจฺฉริโย อพฺภุโต ธมฺโม, ยํ ทิสฺวา ทิสฺวา ภิกฺขู อิมสฺมิํ ธมฺมวินเย อภิรมนฺติ.}

\prose{เสยฺยถาปิ ภิกฺขเว มหาสมุทฺโท มหตํ ภูตานํ อาวาโส.\csromanspacer{} ตตฺริเม ภูตา, ติมิ ติมิงฺคโล ติมิติมิงฺคโล อสุรา นาคา คนฺธพฺพา.}

\vspace{\tipitakaparskip}
\csromancenter{อุทานปาฬิ สนฺติ มหาสมุทฺเท โยชนสติกาปิ อตฺตภาวา ทฺวิโยชนสติกาปิ อตฺตภาวา ติโยชนสติกาปิ อตฺตภาวา จตุโยชนสติกาปิ อตฺตภาวา ปญฺจโยชนสติกาปิ อตฺตภาวา.\csromanspacer{} เอวเมวํ โข ภิกฺขเว อยํ ธมฺมวินโย มหตํ ภูตานํ อาวาโส.\csromanspacer{} ตตฺริเม ภูตา, โสตาปนฺโน โสตาปตฺติผลสจฺฉิกิริยาย ปฏิปนฺโน, สกทาคามี สกทาคามิผลสจฺฉิกิริยาย ปฏิปนฺโน, อนาคามี อนาคามิผลสจฺฉิกิริยาย ปฏิปนฺโน, อรหา อรหตฺตาย\footnote{อรหตฺตผลสจฺฉิกิริยาย (สี)} ปฏิปนฺโน.\csromanspacer{} ยมฺปิ ภิกฺขเว อยํ ธมฺมวินโย มหตํ ภูตานํ อาวาโส.\csromanspacer{} ตตฺริเม ภูตา, โสตาปนฺโน โสตาปตฺติผลสจฺฉิกิริยาย ปฏิปนฺโน, สกทาคามี สกทาคามิผลสจฺฉิกิริยาย ปฏิปนฺโน, อนาคามี อนาคามิผลสจฺฉิกิริยาย ปฏิปนฺโน, อรหา อรหตฺตาย ปฏิปนฺโน.\csromanspacer{} อยํ ภิกฺขเว อิมสฺมิํ ธมฺมวินเย อฏฺฐโม อจฺฉริโย อพฺภุโต ธมฺโม, ยํ ทิสฺวา ทิสฺวา ภิกฺขู อิมสฺมิํ ธมฺมวินเย อภิรมนฺติ.\csromanspacer{} อิเม โข ภิกฺขเว อิมสฺมิํ ธมฺมวินเย อฏฺฐ อจฺฉริยา อพฺภุตา ธมฺมา, เย ทิสฺวา ทิสฺวา ภิกฺขู อิมสฺมิํ ธมฺมวินเย อภิรมนฺตีติ.}
\prose{\csromanfolio{144}อถ โข ภควา เอตมตฺถํ วิทิตฺวา ตายํ เวลายํ อิมํ อุทานํ อุทาเนสิ\csromandash{}}

\csromanmatikaanchor{mtk.76}
\csromantocmarkref{subsection}{4. ยกฺขปหารสุตฺต}{mtk.76}
\bookmark[dest={mtk.76},level=2]{4. ยกฺขปหารสุตฺต}
\prose{\csromansymbolfootnote{*}{ขุ 2. 289; วิ 4. 424; วิ 5. 263; ขุ 10. 132, 144 ปิฏฺเฐสุปิ.}“ฉนฺนมติวสฺสติ, วิวฏํ นาติวสฺสติ.}

\csromangathameasure{ตสฺมา ฉนฺนํ วิวเรถ, เอวํ ตํ นาติวสฺสตี”ติ.ปญฺจมํ.}
\csromangathagroup{\csromangathastanza{ตสฺมา ฉนฺนํ วิวเรถ, เอวํ ตํ นาติวสฺสตี”ติ.\csromanspacer{}\csromanspacer{}ปญฺจมํ.}}

\csromanheader{2}{6. โสณสุตฺต}
\proseitem{46}{เอวํ เม สุตํ\csromandash{}เอกํ สมยํ ภควา สาวตฺถิยํ วิหรติ เชตวเน อนาถปิณฺฑิกสฺส อาราเม.\csromanspacer{} เตน โข ปน สมเยน อายสฺมา มหากจฺจาโน อวนฺตีสุ วิหรติ กุรรฆเร\footnote{กุรุรฆเร (สฺยา, วิ 3. 283), กุลฆเร (ก)} ปวตฺเต ปพฺพเต.\csromanspacer{} เตน โข ปน สมเยน โสโณ อุปาสโก กุฏิกณฺโณ อายสฺมโต มหากจฺจานสฺส อุปฏฺฐาโก โหติ.}

\prose{อถ โข โสณสฺส อุปาสกสฺส กุฏิกณฺณสฺส รโหคตสฺส ปฏิสลฺลีนสฺส เอวํ เจตโส ปริวิตกฺโก อุทปาทิ “ยถา ยถา โข อยฺโย มหากจฺจาโน ธมฺมํ เทเสติ, นยิทํ สุกรํ อคารํ อชฺฌาวสตา เอกนฺตปริปุณฺณํ เอกนฺตปริสุทฺธํ สงฺขลิขิตํ พฺรหฺมจริยํ จริตุํ,}

\vspace{\tipitakaparskip}
\csromancenter{โสณวคฺค ยํนูนาหํ เกสมสฺสุํ โอหาเรตฺวา กาสายานิ วตฺถานิ อจฺฉาเทตฺวา อคารสฺมา อนคาริยํ ปพฺพเชยฺยนฺ”ติ.}
\prose{\csromanfolio{145}อถ โข โสโณ อุปาสโก กุฏิกณฺโณ เยนายสฺมา มหากจฺจาโน เตนุปสงฺกมิ, อุปสงฺกมิตฺวา อายสฺมนฺตํ มหากจฺจานํ อภิวาเทตฺวา เอกมนฺตํ นิสีทิ, เอกมนฺตํ นิสินฺโน โข โสโณ อุปาสโก กุฏิกณฺโณ อายสฺมนฺตํ มหากจฺจานํ เอตทโวจ\csromandash{}}

\prose{อิธ มยฺหํ ภนฺเต รโหคตสฺส ปฏิสลฺลีนสฺส เอวํ เจตโส ปริวิตกฺโก อุทปาทิ “ยถา ยถา โข อยฺโย มหากจฺจาโน ธมฺมํ เทเสติ, นยิทํ สุกรํ อคารํ อชฺฌาวสตา เอกนฺตปริปุณฺณํ เอกนฺตปริสุทฺธํ สงฺขลิขิตํ พฺรหฺมจริยํ จริตุํ, ยํนูนาหํ เกสมสฺสุํ โอหาเรตฺวา กาสายานิ วตฺถานิ อจฺฉาเทตฺวา อคารสฺมา อนคาริยํ ปพฺพเชยฺยนฺ”ติ, ปพฺพาเชตุ มํ ภนฺเต อยฺโย มหากจฺจาโนติ.}

\prose{เอวํ วุตฺเต อายสฺมา มหากจฺจาโน โสณํ อุปาสกํ กุฏิกณฺณํ เอตทโวจ “ทุกฺกรํ โข โสณ ยาวชีวํ เอกภตฺตํ เอกเสยฺยํ พฺรหฺมจริยํ, อิงฺฆ ตฺวํ โสณ ตตฺเถว อาคาริกภูโต สมาโน พุทฺธานํ สาสนํ อนุยุญฺช กาลยุตฺตํ เอกภตฺตํ เอกเสยฺยํ พฺรหฺมจริยนฺ”ติ.\csromanspacer{} อถ โข โสณสฺส อุปาสกสฺส กุฏิกณฺณสฺส โย อโหสิ ปพฺพชฺชาภิสงฺขาโร, โส ปฏิปสฺสมฺภิ.}

\prose{ทุติยมฺปิ โข ฯเปฯ.\csromanspacer{} ทุติยมฺปิ โข อายสฺมา มหากจฺจาโน โสณํ อุปาสกํ กุฏิกณฺณํ เอตทโวจ “ทุกฺกรํ โข โสณ ยาวชีวํ เอกภตฺตํ เอกเสยฺยํ พฺรหฺมจริยํ, อิงฺฆ ตฺวํ โสณ ตตฺเถว อาคาริกภูโต สมาโน พุทฺธานํ สาสนํ อนุยุญฺช กาลยุตฺตํ เอกภตฺตํ เอกเสยฺยํ พฺรหฺมจริยนฺ”ติ.\csromanspacer{} ทุติยมฺปิ โข โสณสฺส อุปาสกสฺส กุฏิกณฺณสฺส โย อโหสิ ปพฺพชฺชาภิสงฺขาโร, โส ปฏิปสฺสมฺภิ.}

\prose{ตติยมฺปิ โข โสณสฺส อุปาสกสฺส กุฏิกณฺณสฺส รโหคตสฺส ปฏิสลฺลีนสฺส เอวํ เจตโส ปริวิตกฺโก อุทปาทิ “ยถา ยถา โข อยฺโย มหากจฺจาโน ธมฺมํ เทเสติ, นยิทํ สุกรํ อคารํ อชฺฌาวสตา เอกนฺตปริปุณฺณํ เอกนฺตปริสุทฺธํ สงฺขลิขิตํ พฺรหฺมจริยํ จริตุํ, ยํนูนาหํ เกสมสฺสุํ โอหาเรตฺวา กาสายานิ วตฺถานิ อจฺฉาเทตฺวา อคารสฺมา อนคาริยํ ปพฺพเชยฺยนฺ”ติ.\csromanspacer{} ตติยมฺปิ โข โสโณ อุปาสโก}

\vspace{\tipitakaparskip}
\csromancenter{อุทานปาฬิ กุฏิกณฺโณ เยนายสฺมา มหากจฺจาโน เตนุปสงฺกมิ, อุปสงฺกมิตฺวา อายสฺมนฺตํ มหากจฺจานํ อภิวาเทตฺวา เอกมนฺตํ นิสีทิ, เอกมนฺตํ นิสินฺโน โข โสโณ อุปสโก กุฏิกณฺโณ อายสฺมนฺตํ มหากจฺจานํ เอตทโวจ\csromandash{}}
\prose{\csromanfolio{146}อิธ มยฺหํ ภนฺเต รโหคตสฺส ปฏิสลฺลีนสฺส เอวํ เจตโส ปริวิตกฺโก อุทปาทิ “ยถา ยถา โข อยฺโย มหากจฺจาโน ธมฺมํ เทเสติ, นยิทํ สุกรํ อคารํ อชฺฌาวสตา เอกนฺตปริปุณฺณํ เอกนฺตปริสุทฺธํ สงฺขลิขิตํ พฺรหฺมจริยํ จริตุํ, ยํนูนาหํ เกสมสฺสุํ โอหาเรตฺวา กาสายานิ วตฺถานิ อจฺฉาเทตฺวา อคารสฺมา อนคาริยํ ปพฺพเชยฺยนฺ”ติ, ปพฺพาเชตุ มํ ภนฺเต อยฺโย มหากจฺจาโนติ.}

\prose{อถ โข อายสฺมา มหากจฺจาโน โสณํ อุปาสกํ กุฏิกณฺณํ ปพฺพาเชสิ.\csromanspacer{} เตน โข ปน สมเยน อวนฺติ ทกฺขิณาปโถ\footnote{อวนฺติ ทกฺขิณปโถ (สี)} อปฺปภิกฺขุโก โหติ.\csromanspacer{} อถ โข อายสฺมา มหากจฺจาโน ติณฺณํ วสฺสานํ อจฺจเยน กิจฺเฉน กสิเรน ตโต ตโต ทสวคฺคํ ภิกฺขุสํฆํ สนฺนิปาเตตฺวา อายสฺมนฺตํ โสณํ อุปสมฺปาเทสิ.}

\prose{อถ โข อายสฺมโต โสณสฺส วสฺสํวุฏฺฐสฺส\footnote{วสฺสํวุตฺถสฺส (สี, สฺยา, กํ, อิ)} รโหคตสฺส ปฏิสลฺลีนสฺส เอวํ เจตโส ปริวิตกฺโก อุทปาทิ “น โข เม โส ภควา สมฺมุขา ทิฏฺโฐ, อปิ จ สุโตเยว เม ‘โส ภควา อีทิโส จ อีทิโส จา’ติ.\csromanspacer{} สเจ มํ อุปชฺฌาโต อนุชาเนยฺย, คจฺเฉยฺยาหํ ตํ ภควนฺตํ ทสฺสนาย อรหนฺตํ สมฺมาสมฺพุทฺธนฺ”ติ.}

\csromanmatikaanchor{mtk.77}
\csromantocmarkref{subsection}{5. นาคสุตฺต}{mtk.77}
\bookmark[dest={mtk.77},level=2]{5. นาคสุตฺต}
\prose{อถ โข อายสฺมา โสโณ สายนฺหสมยํ ปฏิสลฺลานา วุฏฺฐิโต เยนายสฺมา มหากจฺจาโน เตนุปสงฺกมิ, อุปสงฺกมิตฺวา อายสฺมนฺตํ มหากจฺจานํ อภิวาเทตฺวา เอกมนฺตํ นิสีทิ, เอกมนฺตํ นิสินฺโน โข อายสฺมา โสโณ อายสฺมนฺตํ มหากจฺจานํ เอตทโวจ\csromandash{}}

\prose{อิธ มยฺหํ ภนฺเต รโหคตสฺส ปฏิสลฺลีนสฺส เอวํ เจตโส ปริวิตกฺโก อุทปาทิ “น โข เม โส ภควา สมฺมุขา ทิฏฺโฐ, อปิ จ สุโตเยว เม ‘โส ภควา อีทิโส จ อีทิโส จา’ติ.\csromanspacer{} สเจ มํ}

\vspace{\tipitakaparskip}
\csromancenter{โสณวคฺค อุปชฺฌาโย อนุชาเนยฺย, คจฺเฉยฺยาหํ ตํ ภควนฺตํ ทสฺสนาย อรหนฺตํ สมฺมาสมฺพุทฺธนฺ”ติ. ( )\footnote{(คจฺเฉยฺยาหํ ภนฺเต ตํ ภควนฺตํ ทสฺสนาย อรหนฺตํ สมฺมาสมฺพุทฺธํ, สเจ มํ อุปชฺฌาโย อนุชานาตีติ) (วิ 3. 285)}}
\prose{\csromanfolio{147}สาธุ สาธุ โสณ, คจฺฉ ตฺวํ โสณ ตํ ภควนฺตํ ทสฺสนาย อรหนฺตํ สมฺมาสมฺพุทฺธํ\footnote{สมฺมาสมฺพุทฺธนฺติ (สพฺพตฺถ)}, ทกฺขิสฺสสิ ตฺวํ โสณ ตํ ภควนฺตํ ปาสาทิกํ ปสาทนียํ สนฺตินฺทฺริยํ สนฺตมานสํ อุตฺตมทมถสมถมนุปฺปตฺตํ ทนฺตํ คุตฺตํ ยตินฺทฺริยํ นาคํ, ทิสฺวาน มม วจเนน ภควโต ปาเท สิรสา วนฺทาหิ, อปฺปาพาธํ อปฺปาตงฺกํ ลหุฏฺฐานํ พลํ ผาสุวิหารํ\footnote{ผาสุวิหารญฺจ (สี)} ปุจฺฉ “อุปชฺฌาโย เม ภนฺเต อายสฺมา มหากจฺจาโน ภควโต ปาเท สิรสา วนฺทติ, อปฺปาพาธํ อปฺปาตงฺกํ ลหุฏฺฐานํ พลํ ผาสุวิหารํ3 ปุจฺฉตี”ติ.}

\prose{“เอวํ ภนฺเต”ติ โข อายสฺมา โสโณ อายสฺมโต มหากจฺจานสฺส ภาสิตํ อภินนฺทิตฺวา อนุโมทิตฺวา อุฏฺฐายาสนา อายสฺมนฺตํ มหากจฺจานํ อภิวาเทตฺวา ปทกฺขิณํ กตฺวา เสนาสนํ สํสาเมตฺวา ปตฺตจีวรมาทาย เยน สาวตฺถิ เตน จาริกํ ปกฺกามิ.\csromanspacer{} อนุปุพฺเพน จาริกํ จรมาโน เยน สาวตฺถิ เชตวนํ อนาถปิณฺฑิกสฺส อาราโม, เยน ภควา เตนุปสงฺกมิ, อุปสงฺกมิตฺวา ภควนฺตํ อภิวาเทตฺวา เอกมนฺตํ นิสีทิ, เอกมนฺตํ นิสินฺโน โข อายสฺมา โสโณ ภควนฺตํ เอตทโวจ “อุปชฺฌาโย เม ภนฺเต อายสฺมา มหากจฺจาโน ภควโต ปาเท สิรสา วนฺทติ, อปฺปาพาธํ อปฺปาตงฺกํ ลหุฏฺฐานํ พลํ ผาสุวิหารํ3 ปุจฺฉตี”ติ.}

\prose{กจฺจิ ภิกฺขุ ขมนียํ, กจฺจิ ยาปนียํ, กจฺจิสิ อปฺปกิลมเถน อทฺธานํ อาคโต, น จ ปิณฺฑเกน กิลนฺโตสีติ.\csromanspacer{} ขมนียํ ภควา, ยาปนียํ ภควา, อปฺปกิลมเถน จาหํ ภนฺเต อทฺธานํ อาคโต, น จ ปิณฺฑเกน กิลนฺโตมฺหีติ.}

\prose{อถ โข ภควา อายสฺมนฺตํ อานนฺทํ อามนฺเตสิ “อิมสฺสานนฺท อาคนฺตุกสฺส ภิกฺขุโน เสนาสนํ ปญฺญาเปหี”ติ.\csromanspacer{} อถ โข อายสฺมโต อานนฺทสฺส เอตทโหสิ “ยสฺส โข มํ ภควา}

\vspace{\tipitakaparskip}
\csromancenter{อุทานปาฬิ อาณาเปติ ‘อิมสฺสานนฺท อาคนฺตุกสฺส ภิกฺขุโน เสนาสนํ ปญฺญาเปหี’ติ, อิจฺฉติ ภควา เตน ภิกฺขุนา สทฺธิํ เอกวิหาเร วตฺถุํ, อิจฺฉติ ภควา อายสฺมตา โสเณน สทฺธิํ เอกวิหาเร วตฺถุนฺ”ติ, ยสฺมิํ วิหาเร ภควา วิหรติ, ตสฺมิํ วิหาเร อายสฺมโต โสณสฺส เสนาสนํ ปญฺญาเปสิ.}
\prose{\csromanfolio{148}อถ โข ภควา พหุเทว รตฺติํ อพฺโภกาเส นิสชฺชาย วีตินาเมตฺวา ปาเท ปกฺขาเลตฺวา วิหารํ ปาวิสิ.\csromanspacer{} อายสฺมาปิ โข โสโณ พหุเทว รตฺติํ อพฺโภกาเส นิสชฺชาย วีตินาเมตฺวา ปาเท ปกฺขาเลตฺวา วิหารํ ปาวิสิ.\csromanspacer{} อถ โข ภควา รตฺติยา ปจฺจูสสมยํ ปจฺจุฏฺฐาย อายสฺมนฺตํ โสณํ อชฺเฌสิ “ปฏิภาตุ ตํ ภิกฺขุ ธมฺโม ภาสิตุนฺ”ติ.}

\prose{“เอวํ ภนฺเต”ติ โข อายสฺมา โสโณ ภควโต ปฏิสฺสุตฺวา * โสฬส อฏฺฐกวคฺคิกานิ สพฺพาเนว สเรน อภณิ.\csromanspacer{} อถ โข ภควา อายสฺมโต โสณสฺส สรภญฺญปริโยสาเน อพฺภนุโมทิ “สาธุ สาธุ ภิกฺขุ, สุคฺคหิตานิ เต ภิกฺขุ โสฬส อฏฺฐกวคฺคิกานิ สุมนสิกตานิ สูปธาริตานิ, กลฺยาณิยาสิ\footnote{กลฺยาณิยา จ (ก), กลฺยาณิยา จาสิ (?)} วาจาย สมนฺนาคโต วิสฺสฏฺฐาย อเนลคฬาย อตฺถสฺส วิญฺญาปนิยา.\csromanspacer{} กติวสฺโสสิ ตฺวํ ภิกฺขู”ติ.\csromanspacer{} เอกวสฺโส อหํ ภควาติ.\csromanspacer{} กิสฺส ปน ตฺวํ ภิกฺขุ เอวํ จิรํ อกาสีติ.\csromanspacer{} จิรํ ทิฏฺโฐ\footnote{จิรทิฏฺโฐ (สี)} เม ภนฺเต กาเมสุ อาทีนโว, อปิ จ สมฺพาโธ ฆราวาโส พหุกิจฺโจ พหุกรณีโยติ.}

\prose{อถ โข ภควา เอตมตฺถํ วิทิตฺวา ตายํ เวลยํ อิมํ อุทานํ อุทาเนสิ\csromandash{}}

\csromanmatikaanchor{mtk.78}
\csromantocmarkref{subsection}{6. ปิณฺโฑลสุตฺต}{mtk.78}
\bookmark[dest={mtk.78},level=2]{6. ปิณฺโฑลสุตฺต}
\csromangathameasure{“ทิสฺวา อาทีนวํ โลเก, ญตฺวา ธมฺมํ นิรูปธิํ. \\ อริโย น รมตี ปาเป, ปาเป น รมตี สุจี”ติ.ฉฏฺฐํ.}
\csromangathagroup{\csromangathastanza{“ทิสฺวา อาทีนวํ โลเก, ญตฺวา ธมฺมํ นิรูปธิํ. \\ อริโย น รมตี ปาเป, ปาเป น รมตี สุจี”ติ.\csromanspacer{}\csromanspacer{}ฉฏฺฐํ.}}

\csromanheader{2}{7. กงฺขาเรวตสุตฺต}
\proseitem{47}{เอวํ เม สุตํ\csromandash{}เอกํ สมยํ ภควา สาวตฺถิยํ วิหรติ เชตวเน อนาถปิณฺฑิกสฺส อาราเม.\csromanspacer{} เตน โข ปน สมเยน อายสฺมา กงฺขาเรวโต ภควโต อวิทูเร นิสินฺโน โหติ ปลฺลงฺกํ อาภุชิตฺวา อุชุํ กายํ ปณิธาย อตฺตโน กงฺขาวิตรณวิสุทฺธิํ ปจฺจเวกฺขมาโน.}

\chapterheadpageread{5. โสณวคฺค}
\prose{\csromanfolio{149}อทฺทสา โข ภควา อายสฺมนฺตํ กงฺขาเรวตํ อวิทูเร นิสินฺนํ ปลฺลงฺกํ อาภุชิตฺวา อุชุํ กายํ ปณิธาย อตฺตโน กงฺขาวิตรณวิสุทฺธิํ ปจฺจเวกฺขมานํ.}

\csromanmatikaanchor{mtk.79}
\csromantocmarkref{subsection}{7. สาริปุตฺตสุตฺต}{mtk.79}
\bookmark[dest={mtk.79},level=2]{7. สาริปุตฺตสุตฺต}
\csromantocmarkref{subsection}{8. สุนฺทรีสุตฺต}{mtk.80}
\bookmark[dest={mtk.80},level=2]{8. สุนฺทรีสุตฺต}
\prose{อถ โข ภควา เอตมตฺถํ วิทิตฺวา ตายํ เวลายํ อิมํ อุทานํ อุทาเนสิ\csromandash{}}

\csromangathameasure{“ยา กาจิ กงฺขา อิธ วา หุรํ วา, \\ สกเวทิยา วา ปรเวทิยา วา. \\ เย ฌายิโน ตา ปชหนฺติ สพฺพา, \\ อาตาปิโน พฺรหฺมจริยํ จรนฺตา”ติ.สตฺตมํ.}
\csromangathagroup{\csromangathastanza{“ยา กาจิ กงฺขา อิธ วา หุรํ วา, \\ สกเวทิยา วา ปรเวทิยา วา. \\ เย ฌายิโน ตา ปชหนฺติ สพฺพา, \\ อาตาปิโน พฺรหฺมจริยํ จรนฺตา”ติ.\csromanspacer{}\csromanspacer{}สตฺตมํ.}}

\csromanheader{2}{8. สํฆเภทสุตฺต}
\proseitem{48}{เอวํ เม สุตํ\csromandash{}เอกํ สมยํ ภควา ราชคเห วิหรติ เวฬุวเน กลนฺทกนิวาเป.\csromanspacer{} เตน โข ปน สมเยน อายสฺมา อานนฺโท ตทหุโปสเถ ปุพฺพณฺหสมยํ นิวาเสตฺวา ปตฺตจีวรมาทาย ราชคหํ ปิณฺฑาย ปาวิสิ.}

\prose{อทฺทสา โข เทวทตฺโต อายสฺมนฺตํ อานนฺทํ ราชคเห ปิณฺฑาย จรนฺตํ, ทิสฺวาน เยนายสฺมา อานนฺโท เตนุปสงฺกมิ, อุปสงฺกมิตฺวา อายสฺมนฺตํ อานนฺทํ เอตทโวจ “อชฺชตคฺเค ทานาหํ อาวุโส อานนฺท อญฺญเตฺรว ภควตา, อญฺญตฺร ภิกฺขุสํฆา อุโปสถํ กริสฺสามิ สํฆกมฺมานิ จา”ติ.}

\prose{อถ โข อายสฺมา อานนฺโท ราชคเห ปิณฺฑาย จริตฺวา ปจฺฉาภตฺตํ ปิณฺฑปาตปฏิกฺกนฺโต เยน ภควา เตนุปสงฺกมิ, อุปสงฺกมิตฺวา ภควนฺตํ อภิวาเทตฺวา เอกมนฺตํ นิสีทิ, เอกมนฺตํ นิสินฺโน โข อายสฺมา อานนฺโท ภควนฺตํ เอตทโวจ\csromandash{}}

\prose{“อิธาหํ ภนฺเต ปุพฺพณฺหสมยํ นิวาเสตฺวา ปตฺตจีวรมาทาย ราชคหํ ปิณฺฑาย ปาวิสิํ, อทฺทสา โข มํ ภนฺเต เทวทตฺโต ราชคเห ปิณฺฑาย จรนฺตํ, ทิสฺวาน เยนาหํ เตนุปสงฺกมิ, อุปสงฺกมิตฺวา มํ เอตทโวจ “อชฺชตคฺเค ทานาหํ อาวุโส อานนฺท อญฺญเตฺรว ภควตา, อญฺญตฺร ภิกฺขุสํฆา อุโปสถํ กริสฺสามิ สํฆกมฺมานิ จา”ติ.\csromanspacer{} อชฺช ภนฺเต เทวทตฺโต สํฆํ ภินฺทิสฺสติ, อุโปสถญฺจ กริสฺสติ สํฆกมฺมานิ จาติ.}

\gambhira{อุทานปาฬิ}
\prose{\csromanfolio{150}อถ โข ภควา เอตมตฺถํ วิทิตฺวา ตายํ เวลายํ อิมํ อุทานํ อุทาเนสิ\csromandash{}}

\csromangathameasure{“สุกรํ สาธุนา สาธุ, สาธุ ปาเปน ทุกฺกรํ. \\ ปาปํ ปาเปน สุกรํ, ปาปมริเยหิ ทุกฺกรนฺ”ติ.อฏฺฐมํ.}
\csromangathagroup{\csromangathastanza{“สุกรํ สาธุนา สาธุ, สาธุ ปาเปน ทุกฺกรํ\footnote{สุกรํ สาธุนา สาธุํ, สาธุํ ปาเปน ทุกฺกรํ (ก) วิ 4. 361 ปิฏฺเฐปิ ปสฺสิตพฺพํ.}. \\ ปาปํ ปาเปน สุกรํ, ปาปมริเยหิ ทุกฺกรนฺ”ติ.\csromanspacer{}\csromanspacer{}อฏฺฐมํ.}}

\csromanheader{2}{9. สธายมานสุตฺต}
\proseitem{49}{เอวํ เม สุตํ\csromandash{}เอกํ สมยํ ภควา โกสเลสุ จาริกํ จรติ มหตา ภิกฺขุสํเฆน สทฺธิํ.\csromanspacer{} เตน โข ปน สมเยน สมฺพหุลา มาณวกา ภควโต อวิทูเร สธายมานรูปา\footnote{สทฺทายมานรูปา (สฺยา, อิ, อฏฺฐกถายํ ปาฐนฺตรํ.), ปถายมานรูปา (ก), วธายมานรูปา (ก-สี, ก-ฏฺฐ), สทฺธายมานรูปา (?), สทฺธุธาตุยา สธุธาตุยา วา สิทฺธมิทนฺติ เวทิตพฺพํ.} อติกฺกมนฺติ.\csromanspacer{} อทฺทสา โข ภควา สมฺพหุเล มาณวเก อวิทูเร สธายมานรูเป อติกฺกนฺเต.}

\prose{อถ โข ภควา เอตมตฺถํ วิทิตฺวา ตายํ เวลายํ อิมํ อุทานํ อุทาเนสิ\csromandash{}}

\csromangathameasure{“ปริมุฏฺฐา ปณฺฑิตาภาสา, วาจาโคจรภาณิโน. \\ ยาวิจฺฉนฺติ มุขายามํ, เยน นีตา น ตํ วิทู”ติ.นวมํ.}
\csromangathagroup{\csromangathastanza{“ปริมุฏฺฐา ปณฺฑิตาภาสา, วาจาโคจรภาณิโน. \\ ยาวิจฺฉนฺติ มุขายามํ, เยน นีตา น ตํ วิทู”ติ.\csromanspacer{}\csromanspacer{}นวมํ.}}

\csromanheader{2}{10. จูฬปนฺถกสุตฺต}
\proseitem{50}{เอวํ เม สุตํ\csromandash{}เอกํ สมยํ ภควา สาวตฺถิยํ วิหรติ เชตวเน อนาถปิณฺฑิกสฺส อาราเม.\csromanspacer{} เตน โข ปน สมเยน อายสฺมา จูฬปนฺถโก\footnote{จุลฺลปนฺถโก (สี), จูลปนฺถโก (อิ)} ภควโต อวิทูเร นิสินฺโน โหติ ปลฺลงฺกํ อาภุชิตฺวา อุชุํ กายํ ปณิธาย ปริมุขํ สติํ อุปฏฺฐเปตฺวา.}

\prose{อทฺทสา โข ภควา อายสฺมนฺตํ จูฬปนฺถกํ อวิทูเร นิสินฺนํ ปลฺลงฺกํ อาภุชิตฺวา อุชุํ กายํ ปณิธาย ปริมุขํ สติํ อุปฏฺฐเปตฺวา.}

\prose{อถ โข ภควา เอตมตฺถํ วิทิตฺวา ตายํ เวลายํ อิมํ อุทานํ อุทานํ อุทาเนสิ\csromandash{}}

\chapterheadpageread{6. ชจฺจนฺธวคฺค}
\csromangathameasure{“ฐิเตน กาเยน ฐิเตน เจตสา, \\ ติฏฺฐํ นิสินฺโน อุท วา สยาโน. \\ เอตํ สติํ ภิกฺขุ อธิฏฺฐหาโน, \\ ลเภถ ปุพฺพาปริยํ วิเสสํ. \\ ลทฺธาน ปุพฺพาปริยํ วิเสสํ, \\ อทสฺสนํ มจฺจุราชสฺส คจฺเฉ”ติ.ทสมํ. โสณวคฺโค ปญฺจโม.}
\csromangathagroup{\csromangathastanza{\csromanfolio{151}“ฐิเตน กาเยน ฐิเตน เจตสา, \\ ติฏฺฐํ นิสินฺโน อุท วา สยาโน. \\ เอตํ\footnote{เอวํ (ก)} สติํ ภิกฺขุ อธิฏฺฐหาโน, \\ ลเภถ ปุพฺพาปริยํ วิเสสํ. \\ ลทฺธาน ปุพฺพาปริยํ วิเสสํ, \\ อทสฺสนํ มจฺจุราชสฺส คจฺเฉ”ติ.\csromanspacer{}\csromanspacer{}ทสมํ.\csromanspacer{} โสณวคฺโค\footnote{เอวํ (ก)} ปญฺจโม.}}

\titlehead{\textbf{ตสฺสุทฺทานํ}}
\csromangathameasure{ปิโย อปฺปายุกา กุฏฺฐี, กุมารกา อุโปสโถ. \\ โสโณ จ เรวโต เภโท, สธาย ปนฺถเกน จาติ.}
\csromangathagroup{\csromangathastanza{ปิโย อปฺปายุกา กุฏฺฐี, กุมารกา อุโปสโถ. \\ โสโณ จ เรวโต เภโท, สธาย ปนฺถเกน จาติ.}}

\chapterhead{6. ชจฺจนฺธวคฺค}
\csromanheader{2}{1. อายุสงฺขาโรสฺสชฺชนสุตฺต}
\proseitem{51}{\csromansymbolfootnote{*}{ที 2. 90; สํ 3. 230; อํ 3. 129 ปิฏฺเฐสุ.}เอวํ เม สุตํ\csromandash{}เอกํ สมยํ ภควา เวสาลิยํ วิหรติ มหาวเน กูฏาคารสาลายํ.\csromanspacer{} อถ โข ภควา ปุพฺพณฺหสมยํ นิวาเสตฺวา ปตฺตจีวรมาทาย เวสาลิํ ปิณฺฑาย ปาวิสิ, เวสาลิยํ ปิณฺฑาย จริตฺวา ปจฺฉาภตฺตํ ปิณฺฑปาตปฏิกฺกนฺโต อายสฺมนฺตํ อานนฺทํ อามนฺเตสิ “คณฺหาหิ อานนฺท นิสีทนํ, เยน จาปาลํ\footnote{ปาวาลํ (สฺยา)} เจติยํ เตนุปสงฺกมิสฺสาม ทิวาวิหารายา”ติ.}

\prose{“เอวํ ภนฺเต”ติ โข อายสฺมา อานนฺโท ภควโต ปฏิสฺสุตฺวา นิสีทนํ อาทาย ภควนฺตํ ปิฏฺฐิโต ปิฏฺฐิโต อนุพนฺธิ.\csromanspacer{} อถ โข ภควา เยน จาปาลํ เจติยํ เตนุปสงฺกมิ, อุปสงฺกมิตฺวา ปญฺญตฺเต อาสเน นิสีทิ, นิสชฺช โข ภควา อายสฺมนฺตํ อานนฺทํ อามนฺเตสิ\csromandash{}}

\prose{รมณียา อานนฺท เวสาลี, รมณียํ อุเทนํ เจติยํ, รมณียํ โคตมกํ เจติยํ, รมณียํ สตฺตมฺพํ เจติยํ, รมณียํ พหุปุตฺตํ เจติยํ, รมณียํ สารนฺททํ เจติยํ, รมณียํ จาปาลํ เจติยํ.\csromanspacer{} ยสฺส กสฺสจิ}

\vspace{\tipitakaparskip}
\csromancenter{อุทานปาฬิ อานนฺท จตฺตาโร อิทฺธิปาทา ภาวิตา พหุลีกตา ยานีกตา วตฺถุกตา อนุฏฺฐิตา ปริจิตา สุสมารทฺธา, โส อากงฺขมาโน ( )\footnote{(อานนฺท) (ก)} กปฺปํ วา ติฏฺเฐยฺย กปฺปาวเสสํ วา.\csromanspacer{} ตถาคตสฺส โข อานนฺท จตฺตาโร อิทฺธิปาทา ภาวิตา พหุลีกตา ยานีกตา วตฺถุกตา อนุฏฺฐิตา ปริจิตา สุสมารทฺธา, อากงฺขมาโน อานนฺท ตถาคโต กปฺปํ วา ติฏฺเฐยฺย กปฺปาวเสสํ วาติ.}
\prose{\csromanfolio{152}เอวมฺปิ โข อายสฺมา อานนฺโท ภควตา โอฬาริเก นิมิตฺเต กยิรมาเน โอฬาริเก โอภาเส กยิรมาเน นาสกฺขิ ปฏิวิชฺฌิตุํ.\csromanspacer{} น ภควนฺตํ ยาจิ “ติฏฺฐตุ ภนฺเต ภควา กปฺปํ, ติฏฺฐตุ สุคโต กปฺปํ พหุชนหิตาย พหุชนสุขาย โลกานุกมฺปาย อตฺถาย หิตาย สุขาย เทวมนุสฺสานนฺ”ติ, ยถา ตํ มาเรน ปริยุฏฺฐิตจิตฺโต.\csromanspacer{} ทุติยมฺปิ โข ฯเปฯ.\csromanspacer{} ตติยมฺปิ โข ภควา อายสฺมนฺตํ อานนฺทํ อามนฺเตสิ\csromandash{}}

\csromanmatikaanchor{mtk.81}
\csromantocmarkref{subsection}{9. อุปเสนสุตฺต}{mtk.81}
\bookmark[dest={mtk.81},level=2]{9. อุปเสนสุตฺต}
\prose{รมณียา อานนฺท เวสาลี, รมณียํ อุเทนํ เจติยํ, รมณียํ โคตมกํ เจติยํ, รมณียํ สตฺตมฺพํ เจติยํ, รมณียํ พหุปุตฺตํ เจติยํ, รมณียํ สารนฺททํ เจติยํ, รมณียํ จาปาลํ เจติยํ.\csromanspacer{} ยสฺส กสฺสจิ อานนฺท จตฺตาโร อิทฺธิปาทา ภาวิตา พหุลีกตา ยานีกตา วตฺถุกตา อนุฏฺฐิตา ปริจิตา สุสมารทฺธา, โส อากงฺขมาโน กปฺปํ วา ติฏฺเฐยฺย กปฺปาวเสสํ วา.\csromanspacer{} ตถาคตสฺส โข อานนฺท จตฺตาโร อิทฺธิปาทา ภาวิตา พหุลีกตา ยานีกตา วตฺถุกตา อนุฏฺฐิตา ปริจิตา สุสมารทฺธา, อากงฺขมาโน อานนฺท ตถาคโต กปฺปํ วา ติฏฺเฐยฺย กปฺปาวเสสํ วาติ.}

\prose{เอวมฺปิ โข อายสฺมา อานนฺโท ภควตา โอฬาริเก นิมิตฺเต กยิรมาเน โอฬาริเก โอภาเส กยิรมาเน นาสกฺขิ ปฏิวิชฺฌิตุํ.\csromanspacer{} น ภควนฺตํ ยาจิ “ติฏฺฐตุ ภนฺเต ภควา กปฺปํ, ติฏฺฐตุ สุคโต กปฺปํ พหุชนหิตาย พหุชนสุขาย โลกานุกมฺปาย อตฺถาย หิตาย สุขาย เทวมนุสฺสานนฺ”ติ, ยถา ตํ มาเรน ปริยุฏฺฐิตจิตฺโต.}

\prose{อถ โข ภควา อายสฺมนฺตํ อานนฺทํ อามนฺเตสิ “คจฺฉ ตฺวํ อานนฺท, ยสฺสทานิ กาลํ มญฺญสี”ติ. “เอวํ ภนฺเต”ติ โข อายสฺมา อานนฺโท}

\vspace{\tipitakaparskip}
\csromancenter{ชจฺจนฺธวคฺค ภควโต ปฏิสฺสุตฺวา อุฏฺฐายาสนา ภควนฺตํ อภิวาเทตฺวา ปทกฺขิณํ กตฺวา อวิทูเร อญฺญตรสฺมิํ รุกฺขมูเล นิสีทิ.}
\prose{\csromanfolio{153}อถ โข มาโร ปาปิมา อจิรปกฺกนฺเต อายสฺมนฺเต อานนฺเท เยน ภควา เตนุปสงฺกมิ, อุปสงฺกมิตฺวา เอกมนฺตํ อฏฺฐาสิ, เอกมนฺตํ ฐิโต โข มาโร ปาปิมา ภควนฺตํ เอตทโวจ\csromandash{}}

\prose{“ปรินิพฺพาตุ ทานิ ภนฺเต ภควา, ปรินิพฺพาตุ สุคโต, ปรินิพฺพานกาโล ทานิ ภนฺเต ภควโต.\csromanspacer{} ภาสิตา โข ปเนสา ภนฺเต ภควตา วาจา ‘น ตาวาหํ ปาปิม ปรินิพฺพายิสฺสามิ, ยาว เม ภิกฺขู น สาวกา ภวิสฺสนฺติ วิยตฺตา วินีตา วิสารทา\footnote{วิสารทา ปตฺตโยคกฺเขมา (อํ 3. 130), วิสารทปฺปตฺตา โยคกฺเขมา (สี, อิ, ก), วิสารทปฺปตฺตา โยคกฺเขมกามา (สฺยา) ที 2, 88 ปิฏฺเฐปิ ปสฺสิตพฺพํ.} พหุสฺสุตา ธมฺมธรา ธมฺมานุธมฺมปฺปฏิปนฺนา สามีจิปฺปฏิปนฺนา อนุธมฺมจาริโน, สกํ อาจริยกํ อุคฺคเหตฺวา อาจิกฺขิสฺสนฺติ เทเสสฺสนฺติ ปญฺญเปสฺสนฺติ ปฏฺฐเปสฺสนฺติ วิวริสฺสนฺติ วิภชิสฺสนฺติ อุตฺตานีกริสฺสนฺติ, อุปฺปนฺนํ ปรปฺปวาทํ สหธมฺเมน สุนิคฺคหิตํ นิคฺคเหตฺวา สปฺปาฏิหาริยํ ธมฺมํ เทเสสฺสนฺตี’ติ.\csromanspacer{} เอตรหิ โข ปน ภนฺเต\footnote{สนฺติ โข ปน ภนฺเต เอตรหิ (สี, อิ, สํ 3. 229)} ภิกฺขู ภควโต สาวกา วิยตฺตา วินีตา วิสารทา พหุสฺสุตา ธมฺมธรา ธมฺมานุธมฺมปฺปฏิปนฺนา สามีจิปฺปฏิปนฺนา อนุธมฺมจาริโน, สกํ อาจริยกํ อุคฺคเหตฺวา อาจิกฺขนฺติ เทเสนฺติ ปญฺญเปนฺติ ปฏฺฐเปนฺติ วิวรนฺติ วิภชนฺติ อุตฺตานีกโรนฺติ, อุปฺปนฺนํ ปรปฺปวาทํ สหธมฺเมน สุนิคฺคหิตํ นิคฺคเหตฺวา สปฺปาฏิหาริยํ ธมฺมํ เทเสนฺติ.\csromanspacer{} ปรินิพฺพาตุ ทานิ ภนฺเต ภควา, ปรินิพฺพาตุ สุคโต, ปรินิพฺพานกาโล ทานิ ภนฺเต ภควโต.}

\csromanmatikaanchor{mtk.82}
\csromantocmarkref{subsection}{10. สาริปุตฺต-อุปสมสุตฺต}{mtk.82}
\bookmark[dest={mtk.82},level=2]{10. สาริปุตฺต-อุปสมสุตฺต}
\prose{ภาสิตา โข ปเนสา ภนฺเต ภควตา วาจา ‘น ตาวาหํ ปาปิม ปรินิพฺพายิสฺสามิ, ยาว เม ภิกฺขุนิโย น สาวิกา ภวิสฺสนฺติ วิยตฺตา วินีตา วิสารทา พหุสฺสุตา ธมฺมธรา ธมฺมานุธมฺมปฺปฏิปนฺนา สามีจิปฺปฏิปนฺนา อนุธมฺมจารินิโย, สกํ อาจริยกํ อุคฺคเหตฺวา อาจิกฺขิสฺสนฺติ เทเสสฺสนฺติ ปญฺญเปสฺสนฺติ ปฏฺฐเปสฺสนฺติ วิวริสฺสนฺติ วิภชิสฺสนฺติ อุตฺตานีกริสฺสนฺติ, อุปฺปนฺนํ ปรปฺปวาทํ สหธมฺเมน สุนิคฺคหิตํ นิคฺคเหตฺวา สปฺปาฏิหาริยํ ธมฺมํ เทเสสฺสนฺตี’ติ.\csromanspacer{} เอตรหิ โข ปน ภนฺเต ภิกฺขุนิโย ภควโต สาวิกา วิยตฺตา วินีตา วิสารทา พหุสฺสุตา ธมฺมธรา ธมฺมานุธมฺมปฺปฏิปนฺนา}

\vspace{\tipitakaparskip}
\csromancenter{อุทานปาฬิ สามีจิปฺปฏิปนฺนา อนุธมฺมจารินิโย, สกํ อาจริยกํ อุคฺคเหตฺวา อาจิกฺขนฺติ เทเสนฺติ ปญฺญเปนฺติ ปฏฺฐเปนฺติ วิวรนฺติ วิภชนฺติ อุตฺตานีกโรนฺติ, อุปฺปนฺนํ ปรปฺปวาทํ สหธมฺเมน สุนิคฺคหิตํ นิคฺคเหตฺวา สปฺปาฏิหาริยํ ธมฺมํ เทเสนฺติ.\csromanspacer{} ปรินิพฺพาตุ ทานิ ภนฺเต ภควา, ปรินิพฺพาตุ สุคโต, ปรินิพฺพานกาโล ทานิ ภนฺเต ภควโต.}
\prose{\csromanfolio{154}ภาสิตา โข ปเนสา ภนฺเต ภควตา วาจา ‘น ตาวาหํ ปาปิม ปรินิพฺพายิสฺสามิ, ยาว เม อุปาสกา น สาวกา ภวิสฺสนฺติ วิยตฺตา วินีตา วิสารทา พหุสฺสุตา ธมฺมธรา ธมฺมานุธมฺมปฺปฏิปนฺนา สามีจิปฺปฏิปนฺนา อนุธมฺมจาริโน, สกํ อาจริยกํ อุคฺคเหตฺวา อาจิกฺขิสฺสนฺติ เทเสสฺสนฺติ ปญฺญเปสฺสนฺติ ปฏฺฐเปสฺสนฺติ วิวริสฺสนฺติ วิภชิสฺสนฺติ อุตฺตานีกริสฺสนฺติ, อุปฺปนฺนํ ปรปฺปวาทํ สหธมฺเมน สุนิคฺคหิตํ นิคฺคเหตฺวา สปฺปาฏิหาริยํ ธมฺมํ เทเสสฺสนฺตี’ติ.\csromanspacer{} เอตรติ โข ปน ภนฺเต อุปาสกา ภควโต สาวกา วิยตฺตา วินีตา วิสารทา พหุสฺสุตา ธมฺมธรา ธมฺมานุธมฺมปฺปฏิปนฺนา สามีจิปฺปฏิปนฺนา อนุธมฺมจาริโน, สกํ อาจริยกํ อุคฺคเหตฺวา อาจิกฺขนฺติ เทเสนฺติ ปญฺญเปนฺติ ปฏฺฐเปนฺติ วิวรนฺติ วิภชนฺติ อุตฺตานีกโรนฺติ, อุปฺปนฺนํ ปรปฺปวาทํ สหธมฺเมน สุนิคฺคหิตํ นิคฺคเหตฺวา สปฺปาฏิหาริยํ ธมฺมํ เทเสนฺติ.\csromanspacer{} ปรินิพฺพาตุ ทานิ ภนฺเต ภควา, ปรินิพฺพาตุ สุคโต, ปรินิพฺพานกาโล ทานิ ภนฺเต ภควโต.}

\prose{ภสิตา โข ปเนสา ภนฺเต ภควตา วาจา ‘น ตาวาหํ ปาปิม ปรินิพฺพายิสฺสามิ, ยาว เม อุปาสิกา น สาวิกา ภวิสฺสนฺติ วิยตฺตา วินีตา วิสารทา พหุสฺสุตา ธมฺมธรา ธมฺมานุธมฺมปฺปฏิปนฺนา สามีจิปฺปฏิปนฺนา อนุธมฺมจารินิโย, สกํ อาจริยกํ อุคฺคเหตฺวา อาจิกฺขิสฺสนฺติ เทเสสฺสนฺติ ปญฺญเปสฺสนฺติ ปฏฺฐเปสฺสนฺติ วิวริสฺสนฺติ วิภชิสฺสนฺติ อุตฺตานีกริสฺสนฺติ, อุปฺปนฺนํ ปรปฺปวาทํ สหธมฺเมน สุนิคฺคหิตํ นิคฺคเหตฺวา สปฺปาฏิหาริยํ ธมฺมํ เทเสสฺสนฺตี’ติ.\csromanspacer{} เอตรหิ โข ปน ภนฺเต อุปาสิกา ภควโต สาวิกา วิยตฺตา วินีตา วิสารทา พหุสฺสุตา ธมฺมธรา ธมฺมานุธมฺมปฺปฏิปนฺนา สามีจิปฺปฏิปนฺนา อนุธมฺมจารินิโย, สกํ อาจริยกํ อุคฺคเหตฺวา อาจิกฺขนฺติ เทเสนฺติ ปญฺญเปนฺติ ปฏฺฐเปนฺติ วิวรนฺติ วิภชนฺติ อุตฺตานีกโรนฺติ, อุปฺปนฺนํ ปรปฺปวาทํ สหธมฺเมน สุนิคฺคหิตํ นิคฺคเหตฺวา สปฺปาฏิหาริยํ ธมฺมํ เทเสนฺติ.\csromanspacer{} ปรินิพฺพาตุ ทานิ ภนฺเต ภควา, ปรินิพฺพาตุ สุคโต, ปรินิพฺพานกาโล ทานิ ภนฺเต ภควโต.}

\chapterheadpageread{6. ชจฺจนฺธวคฺค}
\prose{\csromanfolio{155}ภาสิตา โข ปเนสา ภนฺเต ภควตา วาจา ‘น ตาวาหํ ปาปิม ปรินิพฺพายิสฺสามิ, ยาว เม อิทํ พฺรหฺมจริยํ น อิทฺธญฺเจว ภวิสฺสติ ผีตญฺจ วิตฺถาริกํ พาหุชญฺญํ ปุถุภูตํ, ยาว เทวมนุสฺเสหิ สุปฺปกาสิตนฺ’ติ.\csromanspacer{} เอตรหิ โข ปน ภนฺเต\footnote{ตยิทํ ภนฺเต (สํ 3. 229)} ภควโต พฺรหฺมจริยํ อิทฺธญฺเจว ผีตญฺจ วิตฺถาริกํ พาหุชญฺญํ ปุถุภูตํ, ยาว เทวมนุสฺเสหิ สุปฺปกาสิตํ.\csromanspacer{} ปรินิพฺพาตุ ทานิ ภนฺเต ภควา, ปรินิพฺพาตุ สุคโต, ปรินิพฺพานกาโล ทานิ ภนฺเต ภควโต”ติ.}

\prose{เอวํ วุตฺเต ภควา มารํ ปาปิมนฺตํ เอตทโวจ “อปฺโปสฺสุกฺโก ตฺวํ ปาปิม โหหิ, น จิรํ ตถาคตสฺส ปรินิพฺพานํ ภวิสฺสติ, อิโต ติณฺณํ มาสานํ อจฺจเยน ตถาคโต ปรินิพฺพายิสฺสตี”ติ.}

\csromanmatikaanchor{mtk.83}
\csromantocmarknopage{chapter}{5. โสณวคฺค}{mtk.83}
\bookmark[dest={mtk.83},level=0]{5. โสณวคฺค}
\prose{อถ โข ภควา จาปาเล เจติเย สโต สมฺปชาโน อายุสงฺขารํ โอสฺสชิ, โอสฺสฏฺเฐ จ ภควตา อายุสงฺขาเร มหาภูมิจาโล อโหสิ ภิํสนโก โลมหํโส, เทวทุนฺทุภิโย\footnote{เทวทุทฺรภิโย (ก)} จ ผลิํสุ.}

\csromanmatikaanchor{mtk.84}
\csromantocmarkref{subsection}{1. ปิยตรสุตฺต}{mtk.84}
\bookmark[dest={mtk.84},level=2]{1. ปิยตรสุตฺต}
\prose{อถ โข ภควา เอตมตฺถํ วิทิตฺวา ตายํ เวลายํ อิมํ อุทานํ อุทาเนสิ\csromandash{}}

\csromangathameasure{“ตุลมตุลญฺจ สมฺภวํ, \\ ภวสงฺขารมวสฺสชิ มุนิ. \\ อชฺฌตฺตรโต สมาหิโต, \\ อภินฺทิ กวจมิวตฺตสมฺภวนฺ”ติ.ปฐมํ.}
\csromangathagroup{\csromangathastanza{“ตุลมตุลญฺจ สมฺภวํ, \\ ภวสงฺขารมวสฺสชิ มุนิ. \\ อชฺฌตฺตรโต สมาหิโต, \\ อภินฺทิ กวจมิวตฺตสมฺภวนฺ”ติ.\csromanspacer{}\csromanspacer{}ปฐมํ.}}

\csromanheader{2}{2. สตฺตชฏิลสุตฺต}
\proseitem{52}{เอวํ เม สุตํ\csromandash{}เอกํ สมยํ ภควา สาวตฺถิยํ วิหรติ ปุพฺพาราเม มิคารมาตุปาสาเท.\csromanspacer{} เตน โข ปน สมเยน ภควา สายนฺหสมยํ ปฏิสลฺลานา วุฏฺฐิโต พหิทฺวารโกฏฺฐเก นิสินฺโน โหติ.\csromanspacer{} อถ โข ราชา ปเสนทิ โกสโล เยน ภควา เตนุปสงฺกมิ, อุปสงฺกมิตฺวา ภควนฺตํ อภิวาเทตฺวา เอกมนฺตํ นิสีทิ.}

\gambhira{อุทานปาฬิ}
\prose{\csromanfolio{156}เตน โข ปน สมเยน สตฺต จ ชฏิลา สตฺต จ นิคณฺฐา สตฺต จ อเจลกา สตฺต จ เอกสาฏกา สตฺต จ ปริพฺพาชกา ปรูฬฺหกจฺฉนขโลมา ขาริวิวิธมาทาย\footnote{ขารีวิธมาทาย (ก, สํ 1. 77; ที 1. 95 ปิฏฺเฐสุปิ)} ภควโต อวิทูเร อติกฺกมนฺติ.}

\prose{อทฺทสา โข ราชา ปเสนทิ โกสโล เต สตฺต จ ชฏิเล สตฺต จ นิคณฺเฐ สตฺต จ อเจลเก สตฺต จ เอกสาฏเก สตฺต จ ปริพฺพาชเก ปรูฬฺหกจฺฉนขโลเม ขาริวิวิธมาทาย ภควโต อวิทูเร อติกฺกมนฺเต, ทิสฺวาน อุฏฺฐายาสนา เอกํสํ อุตฺตราสงฺคํ กริตฺวา ทกฺขิณชาณุมณฺฑลํ ปถวิยํ\footnote{ปฐวิยํ (สี, สา, อิ)} นิหนฺตฺวา เยน เต สตฺต จ ชฏิลา สตฺต จ นิคณฺฐา สตฺต จ อเจลกา สตฺต จ เอกสาฏกา สตฺต จ ปริพฺพาชกา เตนญฺชลิํ ปณาเมตฺวา ติกฺขตฺตุํ นามํ สาเวสิ “ราชาหํ ภนฺเต ปเสนทิ โกสโล, ราชาหํ ภนฺเต ปเสนทิ โกสโล, ราชาหํ ภนฺเต ปเสนทิ โกสโล”ติ.}

\csromanmatikaanchor{mtk.85}
\csromantocmarkref{subsection}{2. อปฺปายุกสุตฺต}{mtk.85}
\bookmark[dest={mtk.85},level=2]{2. อปฺปายุกสุตฺต}
\prose{อถ โข ราชา ปเสนทิ โกสโล อจิรปกฺกนฺเตสุ เตสุ สตฺตสุ จ ชฏิเลสุ สตฺตสุ จ นิคณฺเฐสุ สตฺตสุ จ อเจลเกสุ สตฺตสุ จ เอกสาฏเกสุ สตฺตสุ จ ปริพฺพาชเกสุ เยน ภควา เตนุปสงฺกมิ, อุปสงฺกมิตฺวา ภควนฺตํ อภิวาเทตฺวา เอกมนฺตํ นิสีทิ, เอกมนฺตํ นิสินฺโน โข ราชา ปเสนทิ โกสโล ภควนฺตํ เอตทโวจ “เย โข\footnote{เย จ โข (สี), เย จ เต (สฺยา), เย นุ เกจิ โข (อิ), เย เต (สํ 1.78), เย นุ โข เกจิ (?)} ภนฺเต โลเก อรหนฺโต วา อรหตฺตมคฺคํ วา สมาปนฺนา, เอเต เตสํ อญฺญตเร”ติ\footnote{อญฺญตราติ (สี, ก), อญฺญตโรติ (สฺยา, อิ)}.}

\prose{ทุชฺชานํ โข เอตํ มหาราช ตยา คิหินา กามโภคินา ปุตฺตสมฺพาธสยนํ อชฺฌาวสนฺเตน กาสิกจนฺทนํ ปจฺจนุโภนฺเตน มาลาคนฺธวิเลปนํ ธารยนฺเตน ชาตรูปรชตํ สาทิยนฺเตน “อิเม วา อรหนฺโต อิเม วา อรหตฺตมคฺคํ สมาปนฺนา”ติ.}

\prose{สํวาเสน โข มหาราช สีลํ เวทิตพฺพํ, ตญฺจ โข ทีเฆน อทฺธุนา น อิตฺตรํ\footnote{น อิตฺตเรน (สฺยา, สี-สฺยา-ฏฺฐ)}, มนสิกโรตา โน อมนสิกโรตา, ปญฺญวตา โน ทุปฺปญฺเญน.\csromanspacer{} สํโวหาเรน โข มหาราช โสเจยฺยํ เวทิตพฺพํ, ตญฺจ โข ทีเฆน}

\vspace{\tipitakaparskip}
\csromancenter{ชจฺจนฺธวคฺค อทฺธุนา น อิตฺตรํ, มนสิกโรตา โน อมนสิกโรตา, ปญฺญวตา โน ทุปฺปญฺเญน.\csromanspacer{} อาปทาสุ โข มหาราช ถาโม เวทิตพฺโพ, โส จ โข ทีเฆน อทฺธุนา น อิตฺตรํ, มนสิกโรตา โน อมนสิกโรตา, ปญฺญวตา โน ทุปฺปญฺเญน.\csromanspacer{} สากจฺฉาย โข มหาราช ปญฺญา เวทิตพฺพา, สา จ โข ทีเฆน อทฺธุนา น อิตฺตรํ, มนสิกโรตา โน อมนสิกโรตา, ปญฺญวตา โน ทุปฺปญฺเญนาติ.}
\prose{\csromanfolio{157}อจฺฉริยํ ภนฺเต, อพฺภุตํ ภนฺเต, ยาว สุภาสิตํ จิทํ\footnote{สุภาสิตมิทํ (สํ 1. 78)} ภนฺเต ภควตา\csromandash{}“ทุชฺชานํ โข เอตํ มหาราช ตยา คิหินา ปุตฺตสมฺพาธสยนํ อชฺฌาวสนฺเตน กาสิกจนฺทนํ ปจฺจนุโภนฺเตน มาลาคนฺธวิเลปนํ ธารยนฺเตน ชาตรูปรชตํ สาทิยนฺเตน ‘อิเม วา อรหนฺโต อิเม วา อรหตฺตมคฺคํ สมาปนฺนา’ติ.\csromanspacer{} สํวาเสน โข มหาราช สีลํ เวทิตพฺพํ ฯเปฯ.\csromanspacer{} สากจฺฉาย โข มหาราช ปญฺญา เวทิตพฺพา, สา จ โข ทีเฆน อทฺธุนา น อิตฺตรํ, มนสิกโรตา โน อมนสิกโรตา, ปญฺญวตา โน ทุปฺปญฺเญนา”ติ.}

\csromanmatikaanchor{mtk.86}
\csromantocmarkref{subsection}{3. สุปฺปพุทฺธกุฏฺฐิสุตฺต}{mtk.86}
\bookmark[dest={mtk.86},level=2]{3. สุปฺปพุทฺธกุฏฺฐิสุตฺต}
\prose{เอเต ภนฺเต มม ปุริสา โจรา\footnote{จรา (สํ 1. 79)} โอจรกา ชนปทํ โอจริตฺวา คจฺฉนฺติ, เตหิ ปฐมํ โอจิณฺณํ อหํ ปจฺฉา โอสาริสฺสามิ\footnote{โอตริสฺสามิ (สี, สฺยา, อิ), โอยายิสฺสามิ (สี-สฺยา-ฏฺฐ), โอสาปยิสฺสามิ (สํ 1. 79)}.\csromanspacer{} อิทานิ เต ภนฺเต ตํ รโชชลฺลํ ปวาเหตฺวา สุนฺหาตา สุวิลิตฺตา กปฺปิตเกสมสฺสู โอทาตวตฺถวสนา ปญฺจหิ กามคุเณหิ สมปฺปิตา สมงฺคิภูตา ปริจาเรสฺสนฺตี4ติ.}

\prose{อถ โข ภควา เอตมตฺถํ วิทิตฺวา ตายํ เวลายํ อิมํ อุทานํ อุทาเนสิ\csromandash{}}

\prose{“น วายเมยฺย สพฺพตฺถ, นาญฺญสฺส ปุริโส สิยา.}

\prose{นาญฺญํ นิสฺสาย ชีเวยฺย, ธมฺเมน น วณิํ\footnote{วาณิํ (สี), วณี (สฺยา, อิ), วาณิชํ (ก)} จเร”ติ.\csromanspacer{}\csromanspacer{}ทุติยํ.}

\csromanheader{2}{3. ปจฺจเวกฺขณสุตฺต}
\proseitem{53}{เอวํ เม สุตํ\csromandash{}เอกํ สมยํ ภควา สาวตฺถิยํ วิหรติ เชตวเน อนาถปิณฺฑิกสฺส อาราเม.\csromanspacer{} เตน โข ปน สมเยน ภควา อตฺตโน}

\vspace{\tipitakaparskip}
\csromancenter{อุทานปาฬิ อเนเก ปาปเก อกุสเล ธมฺเม ปหีเน ปจฺจเวกฺขมาโน นิสินฺโน โหติ อเนเก จ กุสเล ธมฺเม ภาวนาปาริปูริํ คเต.}
\prose{\csromanfolio{158}อถ โข ภควา 1 อตฺตโน อเนเก ปาปเก อกุสเล ธมฺเม ปหีเน วิทิตฺวา อเนเก จ กุสเล ธมฺเม ภาวนาปาริปูริํ คเต1 ตายํ เวลายํ อิมํ อุทานํ อุทาเนสิ\csromandash{}}

\csromangathameasure{“อหุ ปุพฺเพ ตทา นาหุ, นาหุ ปุพฺเพ ตทา อหุ. \\ น จาหุ น จ ภวิสฺสติ, น เจตรหิ วิชฺชตี”ติ.ตติยํ.}
\csromangathagroup{\csromangathastanza{“อหุ ปุพฺเพ ตทา นาหุ, นาหุ ปุพฺเพ ตทา อหุ. \\ น จาหุ น จ ภวิสฺสติ, น เจตรหิ วิชฺชตี”ติ.\csromanspacer{}\csromanspacer{}ตติยํ.}}

\csromanheader{2}{4. ปฐมนานาติตฺถิยสุตฺต}
\proseitem{54}{เอวํ เม สุตํ\csromandash{}เอกํ สมยํ ภควา สาวตฺถิยํ วิหรติ เชตวเน อนาถปิณฺฑิกสฺส อาราเม.\csromanspacer{} เตน โข ปน สมเยน สมฺพหุลา นานาติตฺถิยสมณพฺราหฺมณปริพฺพาชกา สาวตฺถิยํ ปฏิวสนฺติ นานาทิฏฺฐิกา นานาขนฺติกา นานารุจิกา นานาทิฏฺฐินิสฺสยนิสฺสิตา\csromandash{}}

\prose{สนฺเตเก สมณพฺราหฺมณา เอวํวาทิโน เอวํทิฏฺฐิโน “สสฺสโต โลโก, อิทเมว สจฺจํ โมฆมญฺญนฺ”ติ.\csromanspacer{} สนฺติ ปเนเก สมณพฺราหฺมณา เอวํวาทิโน เอวํทิฏฺฐิโน “อสสฺสโต โลโก, อิทเมว สจฺจํ โมฆมญฺญนฺ”ติ.\csromanspacer{} สนฺเตเก สมณพฺราหฺมณา เอวํวาทิโน เอวํทิฏฺฐิโน “อนฺตวา โลโก, อิทเมว สจฺจํ โมฆมญฺญนฺ”ติ.\csromanspacer{} สนฺติ ปเนเก สมณพฺราหฺมณา เอวํวาทิโน เอวํทิฏฺฐิโน “อนนฺตวา โลโก, อิทเมว สจฺจํ โมฆมญฺญนฺ”ติ.\csromanspacer{} สนฺเตเก สมณพฺราหฺมณา เอวํวาทิโน เอวํทิฏฺฐิโน “ตํ ชีวํ ตํ สรีรํ, อิทเมว สจฺจํ โมฆมญฺญนฺ”ติ.\csromanspacer{} สนฺติ ปเนเก สมณพฺราหฺมณา เอวํวาทิโน เอวํทิฏฺฐิโน “อญฺญํ ชีวํ อญฺญํ สรีรํ, อิทเมว สจฺจํ โมฆมญฺญนฺ”ติ.\csromanspacer{} สนฺเตเก สมณพฺราหฺมณา เอวํวาทิโน เอวํทิฏฺฐิโน “โหติ ตถาคโต ปรํ มรณา, อิทเมว สจฺจํ โมฆมญฺญนฺ”ติ.\csromanspacer{} สนฺติ ปเนเก สมณพฺราหฺมณา เอวํวาทิโน เอวํทิฏฺฐิโน “น โหติ ตถาคโต ปรํ มรณา, อิทเมว สจฺจํ โมฆมญฺญนฺ”ติ.\csromanspacer{} สนฺเตเก สมณพฺราหฺมณา เอวํวาทิโน เอวํทิฏฺฐิโน “โหติ จ น จ โหติ ตถาคโต ปรํ มรณา, อิทเมว สจฺจํ โมฆมญฺญนฺ”ติ.\csromanspacer{} สนฺติ}

\vspace{\tipitakaparskip}
\csromancenter{ชจฺจนฺธวคฺค ปเนเก สมณพฺราหฺมณา เอวํวาทิโน เอวํทิฏฺฐิโน “เนว โหติ น น โหติ ตถาคโต ปรํ มรณา, อิทเมว สจฺจํ โมฆมญฺญนฺ”ติ.}
\prose{\csromanfolio{159}เต ภณฺฑนชาตา กลหชาตา วิวาทาปนฺนา อญฺญมญฺญํ มุขสตฺตีหิ วิตุทนฺตา วิหรนฺติ “เอทิโส ธมฺโม, เนทิโส ธมฺโม.\csromanspacer{} เนทิโส ธมฺโม, เอทิโส ธมฺโม”ติ.}

\prose{อถ โข สมฺพหุลา ภิกฺขู ปุพฺพณฺหสมยํ นิวาเสตฺวา ปตฺตจีวรมาทาย สาวตฺถิํ ปิณฺฑาย ปาวิสิํสุ, สาวตฺถิยํ ปิณฺฑาย จริตฺวา ปจฺฉาภตฺตํ ปิณฺฑปาตปฏิกฺกนฺตา เยน ภควา เตนุปสงฺกมิํสุ, อุปสงฺกมิตฺวา ภควนฺตํ อภิวาเทตฺวา เอกมนฺตํ นิสีทิํสุ, เอกมนฺตํ นิสินฺนา โข เต ภิกฺขู ภควนฺตํ เอตทโวจุํ\csromandash{}}

\prose{“อิธ ภนฺเต สมฺพหุลา นานาติตฺถิยสมณพฺราหฺมณปริพฺพาชกา สาวตฺถิยํ ปฏิวสนฺติ นานาทิฏฺฐิกา นานาขนฺติกา นานารุจิกา นานาทิฏฺฐินิสฺสยนิสฺสิตา\csromandash{}}

\prose{สนฺเตเก สมณพฺราหฺมณา เอวํวาทิโน เอวํทิฏฺฐิโน ‘สสฺสโต โลโก, อิทเมว สจฺจํ โมฆมญฺญนฺ’ติ ฯเปฯ.\csromanspacer{} เต ภณฺฑนชาตา กลหชาตา วิวาทาปนฺนา อญฺญมญฺญํ มุขสตฺตีหิ วิตุทนฺตา วิหรนฺติ ‘เอทิโส ธมฺโม, เนทิโส ธมฺโม.\csromanspacer{} เนทิโส ธมฺโม, เอทิโส ธมฺโม’ติ.”}

\csromanmatikaanchor{mtk.87}
\csromantocmarkref{subsection}{4. กุมารกสุตฺต}{mtk.87}
\bookmark[dest={mtk.87},level=2]{4. กุมารกสุตฺต}
\prose{อญฺญติตฺถิยา ภิกฺขเว ปริพฺพาชกา อนฺธา อจกฺขุกา อตฺถํ น ชานนฺติ อนตฺถํ น ชานนฺติ, ธมฺมํ น ชานนฺติ อธมฺมํ น ชานนฺติ.\csromanspacer{} เต อตฺถํ อชานนฺตา อนตฺถํ อชานนฺตา ธมฺมํ อชานนฺตา อธมฺมํ อชานนฺตา ภณฺฑนชาตา กลหชาตา วิวาทาปนฺนา อญฺญมญฺญํ มุขสตฺตีหิ วิตุทนฺตา วิหรนฺติ “เอทิโส ธมฺโม, เนทิโส ธมฺโม.\csromanspacer{} เนทิโส ธมฺโม, เอทิโส ธมฺโม”ติ.}

\prose{ภูตปุพฺพํ ภิกฺขเว อิมิสฺสาเยว สาวตฺถิยา อญฺญตโร ราชา อโหสิ.\csromanspacer{} อถ โข ภิกฺขเว โส ราชา อญฺญตรํ ปุริสํ อามนฺเตสิ “เอหิ ตฺวํ อมฺโภ ปุริส ยาวตกา สาวตฺถิยา ชจฺจนฺธา, เต สพฺเพ เอกชฺฌํ สนฺนิปาเตหี”ติ. “เอวํ เทวา”ติ โข ภิกฺขเว โส ปุริโส ตสฺส รญฺโญ ปฏิสฺสุตฺวา ยาวตกา สาวตฺถิยา ชจฺจนฺธา, เต สพฺเพ คเหตฺวา เยน โส ราชา เตนุปสงฺกมิ, อุปสงฺกมิตฺวา ตํ ราชานํ เอตทโวจ “สนฺนิปาติตา โข เต เทว ยาวตกา สาวตฺถิยา}

\vspace{\tipitakaparskip}
\csromancenter{อุทานปาฬิ ชจฺจนฺธา”ติ.\csromanspacer{} เตน หิ ภเณ ชจฺจนฺธานํ หตฺถิํ ทสฺเสหีติ. “เอวํ เทวา”ติ โข ภิกฺขเว โส ปุริโส ตสฺส รญฺโญ ปฏิสฺสุตฺวา ชจฺจนฺธานํ หตฺถิํ ทสฺเสสิ\csromandash{}}
\prose{\csromanfolio{160}เอกจฺจานํ ชจฺจนฺธานํ หตฺถิสฺส สีสํ ทสฺเสสิ “เอทิโส ชจฺจนฺธา หตฺถี”ติ, เอกจฺจานํ ชจฺจนฺธานํ หตฺถิสฺส กณฺณํ ทสฺเสสิ “เอทิโส ชจฺจนฺธา หตฺถี”ติ, เอกจฺจานํ ชจฺจนฺธานํ หตฺถิสฺส ทนฺตํ ทสฺเสสิ “เอทิโส ชจฺจนฺธา หตฺถี”ติ, เอกจฺจานํ ชจฺจนฺธานํ หตฺถิสฺส โสณฺฑํ ทสฺเสสิ “เอทิโส ชจฺจนฺธา หตฺถี”ติ, เอกจฺจานํ ชจฺจนฺธานํ หตฺถิสฺส กายํ ทสฺเสสิ “เอทิโส ชจฺจนฺธา หตฺถี”ติ, เอกจฺจานํ ชจฺจนฺธานํ หตฺถิสฺส ปาทํ ทสฺเสสิ “เอทิโส ชจฺจนฺธา หตฺถี”ติ, เอกจฺจานํ ชจฺจนฺธานํ หตฺถิสฺส สตฺถิํ\footnote{ปิฏฺฐิํ (สฺยา)} ทสฺเสสิ “เอทิโส ชจฺจนฺธา หตฺถี”ติ, เอกจฺจานํ ชจฺจนฺธานํ หตฺถิสฺส นงฺคุฏฺฐํ ทสฺเสสิ “เอทิโส ชจฺจนฺธา หตฺถี”ติ, เอกจฺจานํ ชจฺจนฺธานํ หตฺถิสฺส วาลธิํ ทสฺเสสิ “เอทิโส ชจฺจนฺธา หตฺถี”ติ,}

\prose{อถ โข ภิกฺขเว โส ปุริโส ชจฺจนฺธานํ หตฺถิํ ทสฺเสตฺวา เยน โส ราชา เตนุปสงฺกมิ, อุปสงฺกมิตฺวา ตํ ราชานํ เอตทโวจ “ทิฏฺโฐ โข เตหิ เทว ชจฺจนฺเธหิ หตฺถี, ยสฺสทานิ กาลํ มญฺญสี”ติ.}

\prose{อถ โข ภิกฺขเว โส ราชา เยน เต ชจฺจนฺธา เตนุปสงฺกมิ, อุปสงฺกมิตฺวา เต ชจฺจนฺเธ เอตทโวจ “ทิฏฺโฐ โว ชจฺจนฺธา หตฺถี”ติ.\csromanspacer{} เอวํ เทว, ทิฏฺโฐ โน หตฺถีติ.\csromanspacer{} วเทถ ชจฺจนฺธา กีทิโส หตฺถีติ.}

\csromanmatikaanchor{mtk.88}
\csromantocmarkref{subsection}{5. อุโปสถสุตฺต}{mtk.88}
\bookmark[dest={mtk.88},level=2]{5. อุโปสถสุตฺต}
\prose{เยหิ ภิกฺขเว ชจฺจนฺเธหิ หตฺถิสฺส สีสํ ทิฏฺฐํ อโหสิ, เต เอวมาหํสุ “เอทิโส เทว หตฺถี เสยฺยถาปิ กุมฺโภ”ติ.}

\prose{เยหิ ภิกฺขเว ชจฺจนฺเธหิ หตฺถิสฺส กณฺโณ ทิฏฺโฐ อโหสิ, เต เอวมาหํสุ “เอทิโส เทว หตฺถี เสยฺยถาปิ สุปฺโป”ติ.}

\prose{เยหิ ภิกฺขเว ชจฺจนฺเธหิ หตฺถิสฺส ทนฺโต ทิฏฺโฐ อโหสิ, เต เอวมาหํสุ “เอทิโส เทว หตฺถี เสยฺยถาปิ ขีโล”ติ.}

\prose{เยหิ ภิกฺขเว ชจฺจนฺเธหิ หตฺถิสฺส โสณฺโฑ ทิฏฺโฐ อโหสิ, เต เอวมาหํสุ “เอทิโส เทว หตฺถี เสยฺยถาปิ นงฺคลีสา”ติ.}

\prose{เยหิ ภิกฺขเว ชจฺจนฺเธหิ หตฺถิสฺส กาโย ทิฏฺโฐ อโหสิ, เต เอวมาหํสุ “เอทิโส เทว หตฺถี เสยฺยถาปิ โกฏฺโฐ”ติ.}

\chapterheadpageread{6. ชจฺจนฺธวคฺค}
\prose{\csromanfolio{161}เยหิ ภิกฺขเว ชจฺจนฺเธหิ หตฺถิสฺส ปาโท ทิฏฺโฐ อโหสิ, เต เอวมาหํสุ “เอทิโสเทว หตฺถี เสยฺยถาปิ ถูโณ”ติ.}

\prose{เยหิ ภิกฺขเว ชจฺจนฺเธหิ หตฺถิสฺส สตฺถิ ทิฏฺโฐ\footnote{ปิฏฺฐิ ทิฏฺฐา (ก-สี, สฺยา, อิ), สตฺถิ ทิฏฺฐา (ก-สี)} อโหสิ, เต เอวมาหํสุ “เอทิโส เทว หตฺถี เสยฺยถาปิ อุทุกฺขโล”ติ.}

\prose{เยหิ ภิกฺขเว ชจฺจนฺเธหิ หตฺถิสฺส นงฺคุฏฺฐํ ทิฏฺฐํ อโหสิ, เต เอวมาหํสุ “เอทิโส เทว หตฺถี เสยฺยถาปิ มุสโล”ติ.}

\prose{เยหิ ภิกฺขเว ชจฺจนฺเธหิ หตฺถิสฺส วาลธิ ทิฏฺโฐ อโหสิ, เต เอวมาหํสุ “เอทิโส เทว หตฺถี เสยฺยถาปิ สมฺมชฺชนี”ติ.}

\prose{เต “เอทิโส หตฺถี, เนทิโส หตฺถี.\csromanspacer{} เนทิโส หตฺถี, เอทิโส หตฺถี”ติ อญฺญมญฺญํ มุฏฺฐีหิ สํสุมฺภิํสุ\footnote{สํยุชฺฌิํสุ (ก-สี, สฺยา, อิ)}.\csromanspacer{} เตน จ ปน ภิกฺขเว โส ราชา อตฺตมโน อโหสิ.}

\prose{เอวเมวํ โข ภิกฺขเว อญฺญติตฺถิยา ปริพฺพาชกา อนฺธา อจกฺขุกา เต อตฺถํ น ชานนฺติ อนตฺถํ นชานนฺติ, ธมฺมํ น ชานนฺติ อธมฺมํ น ชานนฺติ.\csromanspacer{} เต อตฺถํ อชานนฺตา อนตฺถํ อชานนฺตา ธมฺมํ อชานนฺตา อธมฺมํ อชานนฺตา ภณฺฑนชาตา กลหชาตา วิวาทาปนฺนา อญฺญมญฺญํ มุขสตฺตีหิ วิตุทนฺตา วิหรนฺติ “เอทิโส ธมฺโม, เนทิโส ธมฺโม.\csromanspacer{} เนทิโส ธมฺโม, เอทิโส ธมฺโม”ติ.}

\prose{อถ โข ภควา เอตมตฺถํ วิทิตฺวา ตายํ เวลายํ อิมํ อุทานํ อุทาเนสิ\csromandash{}}

\csromangathameasure{“อิเมสุ กิร สชฺชนฺติ, เอเก สมณพฺราหฺมณา. \\ วิคฺคยฺห นํ วิวทนฺติ, ชนา เอกงฺคทสฺสิโน”ติ.จตุตฺถํ.}
\csromangathagroup{\csromangathastanza{“อิเมสุ กิร สชฺชนฺติ, เอเก สมณพฺราหฺมณา. \\ วิคฺคยฺห นํ วิวทนฺติ, ชนา เอกงฺคทสฺสิโน”ติ.\csromanspacer{}\csromanspacer{}จตุตฺถํ.}}

\csromanheader{2}{5. ทุติยนานาติตฺถิยสุตฺต}
\proseitem{55}{เอวํ เม สุตํ\csromandash{}เอกํ สมยํ ภควา สาวตฺถิยํ วิหรติ เชตวเน อนาถปิณฺฑิกสฺส อาราเม.\csromanspacer{} เตน โข ปน สมเยน สมฺพหุลา นานาติตฺถิยสมณพฺราหฺมณปริพฺพาชกา สาวตฺถิยํ ปฏิวสนฺติ นานาทิฏฺฐิกา นานาขนฺติกา นานารุจิกา นานาทิฏฺฐินิสฺสยนิสฺสิตา\csromandash{}}

\gambhira{อุทานปาฬิ}
\prose{\csromanfolio{162}สนฺเตเก สมณพฺราหฺมณา เอวํวาทิโน เอวํทิฏฺฐิโน “สสฺสโต อตฺตา จ โลโก จ, อิทเมว สจฺจํ โมฆมญฺญนฺ”ติ.\csromanspacer{} สนฺติ ปเนเก สมณพฺราหฺมณา เอวํวาทิโน เอวํทิฏฺฐิโน “อสสฺสโต อตฺตา จ โลโก จ, อิทเมว สจฺจํ โมฆมญฺญนฺ”ติ.\csromanspacer{} สนฺเตเก สมณพฺราหฺมณา เอวํวาทิโน เอวํทิฏฺฐิโน “สสฺสโต จ อสสฺสโต จ\footnote{สสฺสโต อสสฺสโต (สี)} อตฺตา จ โลโก จ, อิทเมว สจฺจํ โมฆมญฺญนฺ”ติ.\csromanspacer{} สนฺติ ปเนเก สมณพฺราหฺมณา เอวํวาทิโน เอวํทิฏฺฐิโน “เนว สสฺสโต นาสสฺสโต อตฺตา จ โลโก จ, อิทเมว สจฺจํ โมฆมญฺญนฺ”ติ.\csromanspacer{} สนฺเตเก สมณพฺราหฺมณา เอวํวาทิโน เอวํทิฏฺฐิโน “สยํกโต อตฺตา จ โลโก จ, อิทเมว สจฺจํ โมฆมญฺญนฺ”ติ.\csromanspacer{} สนฺติ ปเนเก สมณพฺราหฺมณา เอวํวาทิโน เอวํทิฏฺฐิโน “ปรํกโต อตฺตา จ โลโก จ, อิทเมว สจฺจํ โมฆมญฺญนฺ”ติ.\csromanspacer{} สนฺเตเก สมณพฺราหฺมณา เอวํวาทิโน เอวํทิฏฺฐิโน “สยํกโต จ ปรํกโต จ\footnote{สยํกโต ปรํกโต (สี)} อตฺตา จ โลโก จ, อิทเมว สจฺจํ โมฆมญฺญนฺ”ติ.\csromanspacer{} สนฺติ ปเนเก สมณพฺราหฺมณา เอวํวาทิโน เอวํทิฏฺฐิโน “อสยํกาโร อปรํกาโร\footnote{อสยํกาโร จ อปรํกาโร จ (สฺยา, อิ)} อธิจฺจสมุปฺปนฺโน อตฺตา จ โลโก จ, อิทเมว สจฺจํ โมฆมญฺญนฺ”ติ.\csromanspacer{} สนฺเตเก สมณพฺราหฺมณา เอวํวาทิโน เอวํทิฏฺฐิโน “สสฺสตํ สุขทุกฺขํ อตฺตา จ โลโก จ, อิทเมว สจฺจํ โมฆมญฺญนฺ”ติ.\csromanspacer{} สนฺติ ปเนเก สมณพฺราหฺมณา เอวํวาทิโน เอวํทิฏฺฐิโน “อสสฺสตํ สุขทุกฺขํ อตฺตา จ โลโก จ, อิทเมว สจฺจํ โมฆมญฺญนฺ”ติ.\csromanspacer{} สนฺเตเก สมณพฺราหฺมณา เอวํวาทิโน เอวํทิฏฺฐิโน “สสฺสตญฺจ อสสฺสตญฺจ\footnote{สสฺสตํ อสสฺสตํ (สี)} สุขทุกฺขํ อตฺตา จ โลโก จ, อิทเมว สจฺจํ โมฆมญฺญนฺ”ติ.\csromanspacer{} สนฺติ ปเนเก สมณพฺราหฺมณา เอวํวาทิโน เอวํทิฏฺฐิโน “เนว สสฺสตํ นาสสฺสตํ สุขทุกฺขํ อตฺตา จ โลโก จ, อิทเมว สจฺจํ โมฆมญฺญนฺ”ติ.\csromanspacer{} สนฺเตเก สมณพฺราหฺมณา เอวํวาทิโน เอวํทิฏฺฐิโน “สยํกตํ สุขทุกฺขํ อตฺตา จ โลโก จ, อิทเมว สจฺจํ โมฆมญฺญนฺ”ติ.\csromanspacer{} สนฺติ ปเนเก สมณพฺราหฺมณา เอวํวาทิโน เอวํทิฏฺฐิโน “ปรํกตํ สุขทุกฺขํ อตฺตา จ โลโก จ, อิทเมว สจฺจํ โมฆมญฺญนฺ”ติ.\csromanspacer{} สนฺเตเก สมณพฺราหฺมณา เอวํวาทิโน เอวํทิฏฺฐิโน “สยํกตญฺจ ปรํกตญฺจ\footnote{สยํกตํ ปรํกตํ (สี)} สุขทุกฺขํ อตฺตา จ}

\vspace{\tipitakaparskip}
\csromancenter{ชจฺจนฺธวคฺค โลโก จ, อิทเมว สจฺจํ โมฆมญฺญนฺ”ติ.\csromanspacer{} สนฺติ ปเนเก สมณพฺราหฺมณา เอวํวาทิโน เอวํทิฏฺฐิโน “อสยํการํ อปรํการํ อธิจฺจสมุปฺปนฺนํ สุขทุกฺขํ อตฺตา จ โลโก จ, อิทเมว สจฺจํ โมฆมญฺญนฺ”ติ.}
\prose{\csromanfolio{163}เต ภณฺฑนชาตา กลหชาตา วิวาทาปนฺนา อญฺญมญฺญํ มุขสตฺตีหิ วิตุทนฺตา วิหรนฺติ “เอทิโส ธมฺโม, เนทิโส ธมฺโม.\csromanspacer{} เนทิโส ธมฺโม, เอทิโส ธมฺโม”ติ.}

\prose{อถ โข สมฺพหุลา ภิกฺขู ปุพฺพณฺหสมยํ นิวาเสตฺวา ปตฺตจีวรมาทาย สาวตฺถิํ ปิณฺฑาย ปาวิสิํสุ, สาวตฺถิยํ ปิณฺฑาย จริตฺวา ปจฺฉาภตฺตํ ปิณฺฑปาตปฏิกฺกนฺตา เยน ภควา เตนุปสงฺกมิํสุ, อุปสงฺกมิตฺวา ภควนฺตํ อภิวาเทตฺวา เอกมนฺตํ นิสีทิํสุ, เอกมนฺตํ นิสินฺนา โข เต ภิกฺขู ภควนฺตํ เอตทโวจุํ\csromandash{}}

\prose{“อิธ ภนฺเต สมฺพหุลา นานาติตฺถิยสมณพฺราหฺมณปริพฺพาชกา สาวตฺถิยํ ปฏิวสนฺติ นานาทิฏฺฐิกา นานาขนฺติกา นานารุจิกา นานาทิฏฺฐินิสฺสยนิสฺสิตา\csromandash{}}

\prose{สนฺเตเก สมณพฺราหฺมณา เอวํวาทิโน เอวํทิฏฺฐิโน ‘สสฺสโต อตฺตา จ โลโก จ, อิทเมว สจฺจํ โมฆมญฺญนฺ’ติ ฯเปฯ.\csromanspacer{} เต ภณฺฑนชาตา กลหชาตา วิวาทาปนฺนา อญฺญมญฺญํ มุขสตฺตีหิ วิตุทนฺตา วิหรนฺติ ‘เอทิโส ธมฺโม, เนทิโส ธมฺโม.\csromanspacer{} เนทิโส ธมฺโม, เอทิโส ธมฺโม’ติ”.}

\prose{อญฺญติตฺถิยา ภิกฺขเว ปริพฺพาชกา อนฺธา อจกฺขุกา อตฺถํ น ชานนฺติ อนตฺถํ น ชานนฺติ, ธมฺมํ น ชานนฺติ อธมฺมํ น ชานนฺติ.\csromanspacer{} เต อตฺถํ อชานนฺตา อนตฺถํ อชานนฺตา ธมฺมํ อชานนฺตา อธมฺมํ อชานนฺตา ภณฺฑนชาตา กลหชาตา วิวาทาปนฺนา อญฺญมญฺญํ มุขสตฺตีหิ วิตุทนฺตา วิหรนฺติ “เอทิโส ธมฺโม, เนทิโส ธมฺโม.\csromanspacer{} เนทิโส ธมฺโม, เอทิโส ธมฺโม”ติ.}

\prose{อถ โข ภควา เอตมตฺถํ วิทิตฺวา ตายํ เวลายํ อิมํ อุทานํ อุทาเนสิ\csromandash{}}

\csromangathameasure{“อิเมสุ กิร สชฺชนฺติ, เอเก สมณพฺราหฺมณา. \\ อนฺตราว วิสีทนฺติ, อปฺปตฺวาว ตโมคธนฺ”ติ.ปญฺจมํ.}
\csromangathagroup{\csromangathastanza{“อิเมสุ กิร สชฺชนฺติ, เอเก สมณพฺราหฺมณา. \\ อนฺตราว วิสีทนฺติ, อปฺปตฺวาว ตโมคธนฺ”ติ.\csromanspacer{}\csromanspacer{}ปญฺจมํ.}}

\csromanheader{2}{อุทานปาฬิ \textbf{6. ตติยนานาติตฺถิยสุตฺต}}
\proseitem{56}{\csromanfolio{164}เอวํ เม สุตํ\csromandash{}เอกํ สมยํ ภควา สาวตฺถิยํ วิหรติ เชตวเน อนาถปิณฺฑิกสฺส อาราเม.\csromanspacer{} เตน โข ปน สมเยน สมฺพหุลา นานาติตฺถิยสมณพฺราหฺมณปริพฺพาชกา สาวตฺถิยํ ปฏิวสนฺติ นานาทิฏฺฐิกา นานาขนฺติกา นานารุจิกา นานาทิฏฺฐินิสฺสยนิสฺสิตา\csromandash{}}

\prose{สนฺเตเก สมณพฺราหฺมณา เอวํวาทิโน เอวํทิฏฺฐิโน “สสฺสโต อตฺตา จ โลโก จ, อิทเมว สจฺจํ โมฆมญฺญนฺ”ติ.\csromanspacer{} สนฺติ ปเนเก สมณพฺราหฺมณา เอวํวาทิโน เอวํทิฏฺฐิโน “อสสฺสโต อตฺตา จ โลโก จ, อิทเมว สจฺจํ โมฆมญฺญนฺ”ติ.\csromanspacer{} สนฺเตเก สมณพฺราหฺมณา เอวํวาทิโน เอวํทิฏฺฐิโน “สสฺสโต จ อสสฺสโต จ อตฺตา จ โลโก จ, อิทเมว สจฺจํ โมฆมญฺญนฺ”ติ.\csromanspacer{} สนฺติ ปเนเก สมณพฺราหฺมณา เอวํวาทิโน เอวํทิฏฺฐิโน “เนว สสฺสโต นาสสฺสโต อตฺตา จ โลโก จ, อิทเมว สจฺจํ โมฆมญฺญนฺ”ติ.\csromanspacer{} สนฺเตเก สมณพฺราหฺมณา เอวํวาทิโน เอวํทิฏฺฐิโน “สยํกโต อตฺตา จ โลโก จ, อิทเมว สจฺจํ โมฆมญฺญนฺ”ติ.\csromanspacer{} สนฺติ ปเนเก สมณพฺราหฺมณา เอวํวาทิโน เอวํทิฏฺฐิโน “ปรํกโต อตฺตา จ โลโก จ, อิทเมว สจฺจํ โมฆมญฺญนฺ”ติ.\csromanspacer{} สนฺเตเก สมณพฺราหฺมณา เอวํวาทิโน เอวํทิฏฺฐิโน “สยํกโต จ ปรํกโต จ อตฺตา จ โลโก จ, อิทเมว สจฺจํ โมฆมญฺญนฺ”ติ.\csromanspacer{} สนฺติ ปเนเก สมณพฺราหฺมณา เอวํวาทิโน เอวํทิฏฺฐิโน “อสยํกาโร อปรํกาโร อธิจฺจสมุปฺปนฺโน อตฺตา จ โลโก จ, อิทเมว สจฺจํ โมฆมญฺญนฺ”ติ.\csromanspacer{} สนฺเตเก สมณพฺราหฺมณา เอวํวาทิโน เอวํทิฏฺฐิโน “สสฺสตํ สุขทุกฺขํ อตฺตา จ โลโก จ, อิทเมว สจฺจํ โมฆมญฺญนฺ”ติ.\csromanspacer{} สนฺติ ปเนเก สมณพฺราหฺมณา เอวํวาทิโน เอวํทิฏฺฐิโน “อสสฺสตํ สุขทุกฺขํ อตฺตา จ โลโก จ, อิทเมว สจฺจํ โมฆมญฺญนฺ”ติ.\csromanspacer{} สนฺเตเก สมณพฺราหฺมณา เอวํวาทิโน เอวํทิฏฺฐิโน “สสฺสตญฺจ อสสฺสตญฺจ สุขทุกฺขํ อตฺตา จ โลโก จ, อิทเมว สจฺจํ โมฆมญฺญนฺ”ติ.\csromanspacer{} สนฺติ ปเนเก สมณพฺราหฺมณา เอวํวาทิโน เอวํทิฏฺฐิโน “เนว สสฺสตํ นาสสฺสตํ สุขทุกฺขํ อตฺตา จ โลโก จ, อิทเมว สจฺจํ โมฆมญฺญนฺ”ติ.\csromanspacer{} สนฺเตเก สมณพฺราหฺมณา เอวํวาทิโน เอวํทิฏฺฐิโน “สยํกตํ สุขทุกฺขํ อตฺตา จ โลโก จ, อิทเมว สจฺจํ โมฆมญฺญนฺ”ติ.\csromanspacer{} สนฺติ ปเนเก สมณพฺราหฺมณา เอวํวาทิโน เอวํทิฏฺฐิโน “ปรํกตํ สุขทุกฺขํ อตฺตา จ โลโก จ, อิทเมว สจฺจํ}

\vspace{\tipitakaparskip}
\csromancenter{ชจฺจนฺธวคฺค โมฆมญฺญนฺ”ติ.\csromanspacer{} สนฺเตเก สมณพฺราหฺมณา เอวํวาทิโน เอวํทิฏฺฐิโน “สยํกตญฺจ ปรํกตญฺจ สุขทุกฺขํ อตฺตา จ โลโก จ, อิทเมว สจฺจํ โมฆมญฺญนฺ”ติ.\csromanspacer{} สนฺติ ปเนเก สมณพฺราหฺมณา เอวํวาทิโน เอวํทิฏฺฐิโน “อสยํการํ อปรํการํ อธิจฺจสมุปฺปนฺนํ สุขทุกฺขํ อตฺตา จ โลโก จ, อิทเมว สจฺจํ โมฆมญฺญนฺ”ติ.}
\prose{\csromanfolio{165}เต ภณฺฑนชาตา กลหชาตา วิวาทาปนฺนา อญฺญมญฺญํ มุขสตฺตีหิ วิตุทนฺตา วิหรนฺติ “เอทิโส ธมฺโม, เนทิโส ธมฺโม.\csromanspacer{} เนทิโส ธมฺโม, เอทิโส ธมฺโม”ติ.}

\prose{อถ โข สมฺพหุลา ภิกฺขู ปุพฺพณฺหสมยํ นิวาเสตฺวา ปตฺตจีวรมาทาย สาวตฺถิํ ปิณฺฑาย ปาวิสิํสุ, สาวตฺถิยํ ปิณฺฑาย จริตฺวา ปจฺฉาภตฺตํ ปิณฺฑปาตปฏิกฺกนฺตา เยน ภควา เตนุปสงฺกมิํสุ, อุปสงฺกมิตฺวา ภควนฺตํ อภิวาเทตฺวา เอกมนฺตํ นิสีทิํสุ, เอกมนฺตํ นิสินฺนา โข เต ภิกฺขู ภควนฺตํ เอตทโวจุํ\csromandash{}}

\prose{“อิธ ภนฺเต สมฺพหุลา นานาติตฺถิยสมณพฺราหฺมณปริพฺพาชกา สาวตฺถิยํ ปฏิวสนฺติ นานาทิฏฺฐิกา นานาขนฺติกา นานารุจิกา นานาทิฏฺฐินิสฺสยนิสฺสิตา\csromandash{} สนฺเตเก สมณพฺราหฺมณา เอวํวาทิโน เอวํทิฏฺฐิโน ‘สสฺสโต อตฺตา จ โลโก จ, อิทเมว สจฺจํ โมฆมญฺญนฺ’ติ ฯเปฯ.\csromanspacer{} เต ภณฺฑนชาตา กลหชาตา วิวาทาปนฺนา อญฺญมญฺญํ มุขสตฺตีหิ วิตุทนฺตา วิหรนฺติ ‘เอทิโส ธมฺโม, เนทิโส ธมฺโม.\csromanspacer{} เนทิโส ธมฺโม, เอทิโส ธมฺโม’ติ”.}

\prose{อญฺญติตฺถิยา ภิกฺขเว ปริพฺพาชกา อนฺธา อจกฺขุกา, เต อตฺถํ น ชานนฺติ อนตฺถํ น ชานนฺติ, ธมฺมํ น ชานนฺติ อธมฺมํ น ชานนฺติ.\csromanspacer{} เต อตฺถํ อชานนฺตา อนตฺถํ อชานนฺตา ธมฺมํ อชานนฺตา อธมฺมํ อชานนฺตา ภณฺฑนชาตา กลหชาตา วิวาทาปนฺนา อญฺญมญฺญํ มุขสตฺตีหิ วิตุทนฺตา วิหรนฺติ “เอทิโส ธมฺโม, เนทิโส ธมฺโม.\csromanspacer{} เนทิโส ธมฺโม, เอทิโส ธมฺโม”ติ.}

\prose{อถ โข ภควา เอตมตฺถํ วิทิตฺวา ตายํ เวลายํ อิมํ อุทานํ อุทาเนสิ\csromandash{}}

\csromanmatikaanchor{mtk.89}
\csromantocmarkref{subsection}{6. โสณสุตฺต}{mtk.89}
\bookmark[dest={mtk.89},level=2]{6. โสณสุตฺต}
\gambhira{อุทานปาฬิ}
\csromanmatikaanchor{mtk.90}
\csromanmatikaanchor{mtk.101}
\csromantocmarkref{subsection}{7. กงฺขาเรวตสุตฺต}{mtk.90}
\bookmark[dest={mtk.90},level=2]{7. กงฺขาเรวตสุตฺต}
\csromangathameasure{“อหงฺการปสุตายํ ปชา, ปรํการูปสํหิตา. \\ เอตเทเก นาพฺภญฺญํสุ, น นํ ‘สลฺลนฺ’ติ อทฺทสุํ. \\ เอตญฺจ สลฺลํ ปฏิกจฺจ ปสฺสโต, \\ ‘อหํ กโรมี’ติ น ตสฺส โหติ, \\ ‘ปโร กโรตี’ติ น ตสฺส โหติ. \\ มานุเปตา อยํ ปชา, \\ มานคนฺถา มานวินิพทฺธา.}
\csromangathagroup{\csromangathastanza{\csromanfolio{166}“อหงฺการปสุตายํ ปชา, ปรํการูปสํหิตา. \\ เอตเทเก นาพฺภญฺญํสุ, น นํ ‘สลฺลนฺ’ติ อทฺทสุํ.}\csromangathastanza{เอตญฺจ สลฺลํ ปฏิกจฺจ\footnote{ปฏิคจฺจ (สี, สฺยา, กํ, อิ)} ปสฺสโต, \\ ‘อหํ กโรมี’ติ น ตสฺส โหติ, \\ ‘ปโร กโรตี’ติ น ตสฺส โหติ. \\ มานุเปตา อยํ ปชา,}\csromangathastanza{มานคนฺถา มานวินิพทฺธา\footnote{มานวินิพนฺธา (สี)}.}}

\prose{ทิฏฺฐีสุ สารมฺภกถา, สํสารํ นาติวตฺตตี”ติ.\csromanspacer{}\csromanspacer{}ฉฏฺฐํ. \textbf{7.}\csromanspacer{}\textbf{ สุภูติสุตฺต}}

\proseitem{57}{เอวํ เม สุตํ\csromandash{}เอกํ สมยํ ภควา สาวตฺถิยํ วิหรติ เชตวเน อนาถปิณฺฑิกสฺส อาราเม.\csromanspacer{} เตน โข ปน สมเยน อายสฺมา สุภูติ ภควโต อวิทูเร นิสินฺโน โหติ ปลฺลงฺกํ อาภุชิตฺวา อุชุํ กายํ ปณิธาย อวิตกฺกํ สมาธิํ สเมปชฺชิตฺวา.}

\prose{อทฺทสา โข ภควา อายสฺมนฺตํ สุภูติํ อวิทูเร นิสินฺนํ ปลฺลงฺกํ อาภุชิตฺวา อุชุํ กายํ ปณิธาย อวิตกฺกํ สมาธิํ สมาปนฺนํ.}

\prose{อถ โข ภควา เอตมตฺถํ วิทิตฺวา ตายํ เวลายํ อิมํ อุทานํ อุทาเนสิ\csromandash{}}

\csromangathameasure{“ยสฺส วิตกฺกา วิธูปิตา, \\ อชฺฌตฺตํ สุวิกปฺปิตา อเสสา. \\ ตํ สงฺคมติจฺจ อรูปสญฺญี, \\ จตุโยคาติคโต น ชาตุ เมตี”ติ.สตฺตมํ.}
\csromangathagroup{\csromangathastanza{“ยสฺส วิตกฺกา วิธูปิตา, \\ อชฺฌตฺตํ สุวิกปฺปิตา อเสสา. \\ ตํ สงฺคมติจฺจ อรูปสญฺญี, \\ จตุโยคาติคโต น ชาตุ เมตี”ติ\footnote{น ชาติเมตีติ (สฺยา, อิ, ฏฺฐ-ปาฐนฺตรํ)}.\csromanspacer{}\csromanspacer{}สตฺตมํ.}}

\csromanheader{2}{8. คณิกาสุตฺต}
\proseitem{58}{เอวํ เม สุตํ\csromandash{}เอกํ สมยํ ภควา ราชคเห วิหรติ เวฬุวเน กลนฺทกนิวาเป.\csromanspacer{} เตน โข ปน สมเยน ราชคเห เทฺว ปูคา อญฺญตริสฺสา คณิกาย สารตฺตา โหนฺติ ปฏิพทฺธจิตฺตา ภณฺฑนชาตา กลหชาตา วิวาทาปนฺนา อญฺญมญฺญํ ปาณีหิปิ อุปกฺกมนฺติ, เลฑฺฑูหิปิ}

\vspace{\tipitakaparskip}
\csromancenter{ชจฺจนฺธวคฺค อุปกฺกมนฺติ, ทณฺเฑหิปิ อุปกฺกมนฺติ, สตฺเถหิปิ อุปกฺกมนฺติ.\csromanspacer{} เต ตตฺถ มรณมฺปิ นิคจฺฉนฺติ มรณมตฺตมฺปิ ทุกฺขํ.}
\prose{\csromanfolio{167}อถ โข สมฺพหุลา ภิกฺขู ปุพฺพณฺหสมยํ นิวาเสตฺวา ปตฺตจีวรมาทาย ราชคหํ ปิณฺฑาย ปาวิสิํสุ, ราชคเห ปิณฺฑาย จริตฺวา ปจฺฉาภตฺตํ ปิณฺฑปาตปฏิกฺกนฺตา เยน ภควา เตนุปสงฺกมิํสุ, อุปสงฺกมิตฺวา ภควนฺตํ อภิวาเทตฺวา เอกมนฺตํ นิสีทิํสุ, เอกมนฺตํ นิสินฺนา โข เต ภิกฺขู ภควนฺตํ เอตทโวจุํ\csromandash{}}

\prose{อิธ ภนฺเต ราชคเห เทฺว ปูคา อญฺญตริสฺสา คณิกาย สารตฺตา ปฏิพทฺธจิตฺตา ภณฺฑนชาตา กลหชาตา วิวาทาปนฺนา อญฺญมญฺญํ ปาณีหิปิ อุปกฺกมนฺติ, เลฑฺฑูหิปิ อุปกฺกมนฺติ, ทณฺเฑหิปิ อุปกฺกมนฺติ, สตฺเถหิปิ อุปกฺกมนฺติ.\csromanspacer{} เต ตตฺถ มรณมฺปิ นิคจฺฉนฺติ มรณมตฺตมฺปิ ทุกฺขนฺติ.}

\prose{อถ โข ภควา เอตมตฺถํ วิทิตฺวา ตายํ เวลายํ อิมํ อุทานํ อุทาเนสิ\csromandash{} “ยญฺจ ปตฺตํ ยญฺจ ปตฺตพฺพํ, อุภยเมตํ รชานุกิณฺณํ, อาตุรสฺสานุสิกฺขโต.\csromanspacer{} เย จ สิกฺขาสารา สีลพฺพตํ ชีวิตํ พฺรหฺมจริยํ อุปฏฺฐานสารา, อยเมโก อนฺโต.\csromanspacer{} เย จ เอวํวาทิโน ‘นตฺถิ กาเมสุ โทโส’ติ, อยํ ทุติโย อนฺโต.\csromanspacer{} อิจฺเจเต อุโภ อนฺตา กฏสิวฑฺฒนา, กฏสิโย ทิฏฺฐิํ วฑฺเฒนฺติ.\csromanspacer{} เอเต เต อุโภ อนฺเต อนภิญฺญาย โอลียนฺติ เอเก, อติธาวนฺติ เอเก.\csromanspacer{} เย จ โข เต อภิญฺญาย ตตฺร จ นาเหสุํ, เตน จ นามญฺญิํสุ, วฏฺฏํ เตสํ นตฺถิ ปญฺญาปนายา”ติ.\csromanspacer{}\csromanspacer{}อฏฺฐมํ.}

\csromanheader{2}{9. อุปาติธาวนฺติสุตฺต}
\proseitem{59}{เอวํ เม สุตํ\csromandash{}เอกํ สมยํ ภควา สาวตฺถิยํ วิหรติ เชตวเน อนาถปิณฺฑิกสฺส อาราเม.\csromanspacer{} เตน โข ปน สมเยน ภควา รตฺตนฺธการติมิสายํ อพฺโภกาเส นิสินฺโน โหติ เตลปฺปทีเปสุ ฌายมาเนสุ.}

\prose{เตน โข ปน สมเยน สมฺพหุลา อธิปาตกา เตสุ เตลปฺปทีเปสุ อาปาตปริปาตํ อนยํ อาปชฺชนฺติ, พฺยสนํ อาปชฺชนฺติ, 1 อนยพฺยสนํ}

\vspace{\tipitakaparskip}
\csromancenter{อุทานปาฬิ อาปชฺชนฺติ1.\csromanspacer{} อทฺทสา โข ภควา เต สมฺพหุเล อธิปาตเก เตสุ เตลปฺปทีเปสุ อาปาตปริปาตํ อนยํ อาปชฺชนฺเต พฺยสนํ อาปชฺชนฺเต อนยพฺยสนํ อาปชฺชนฺเต.}
\prose{\csromanfolio{168}อถ โข ภควา เอตมตฺถํ วิทิตฺวา ตายํ เวลายํ อิมํ อุทานํ อุทาเนสิ\csromandash{}}

\csromangathameasure{“อุปาติธาวนฺติ น สารเมนฺติ, \\ นวํ นวํ พนฺธนํ พฺรูหยนฺติ. \\ ปตนฺติ ปชฺโชตมิวาธิปาตกา, \\ ทิฏฺเฐ สุเต อิติเหเก นิวิฏฺฐา”ติ.นวมํ.}
\csromangathagroup{\csromangathastanza{“อุปาติธาวนฺติ น สารเมนฺติ, \\ นวํ นวํ พนฺธนํ พฺรูหยนฺติ. \\ ปตนฺติ ปชฺโชตมิวาธิปาตกา\footnote{...ธิปาตา (สี, สฺยา)}, \\ ทิฏฺเฐ สุเต อิติเหเก นิวิฏฺฐา”ติ.\csromanspacer{}\csromanspacer{}นวมํ.}}

\csromanheader{2}{10. อุปฺปชฺชนฺติสุตฺต}
\proseitem{60}{เอวํ เม สุตํ\csromandash{}เอกํ สมยํ ภควา สาวตฺถิยํ วิหรติ เชตวเน อนาถปิณฺฑิกสฺส อาราเม.\csromanspacer{} อถ โข อายสฺมา อานนฺโท เยน ภควา เตนุปสงฺกมิ, อุปสงฺกมิตฺวา ภควนฺตํ อภิวาเทตฺวา เอกมนฺตํ นิสีทิ, เอกมนฺตํ นิสินฺโน โข อายสฺมา อานนฺโท ภควนฺตํ เอตทโวจ\csromandash{}}

\prose{ยาวกีวญฺจ ภนฺเต ตถาคตา โลเก นุปฺปชฺชนฺติ อรหนฺโต สมฺมาสมฺพุทฺธา, ตาว อญฺญติตฺถิยา ปริพฺพาชกา สกฺกตา โหนฺติ ครุกตา มานิตา ปูชิตา อปจิตา, ลาภิโน จีวร ปิณฺฑปาต เสนาสน คิลานปจฺจย เภสชฺช ปริกฺขารานํ.\csromanspacer{} ยโต จ โข ภนฺเต ตถาคตา โลเก อุปฺปชฺชนฺติ อรหนฺโต สมฺมาสมฺพุทฺธา, อถ อญฺญติตฺถิยา ปริพฺพาชกา อสกฺกตา โหนฺติ อครุกตา อมานิตา อปูชิตา อนปจิตา, น ลาภิโน จีวร ปิณฺฑปาต เสนาสน คิลานปจฺจย เภสชฺช ปริกฺขารานํ.\csromanspacer{} ภควาเยว\footnote{ภควา เจว (สฺยา)} ทานิ ภนฺเต สกฺกโต โหติ ครุกโต มานิโต ปูชิโต อปจิโต, ลาภี จีวร ปิณฺฑปาต เสนาสน คิลานปจฺจยเภสชฺช ปริกฺขารานํ ภิกฺขุสํโฆ จาติ.}

\chapterheadpageread{6. ชจฺจนฺธวคฺค}
\prose{\csromanfolio{169}เอวเมตํ อานนฺท, ยาวกีวญฺจ อานนฺท ตถาคตา โลเก นุปฺปชฺชนฺติ อรหนฺโต สมฺมาสมฺพุทฺธา.\csromanspacer{} ตาว อญฺญติตฺถิยา ปริพฺพาชกา สกฺกตา โหนฺติ ครุกตา มานิตา ปูชิตา อปจิตา, ลาภิโน จีวร ปิณฺฑปาตเสนาสน คิลานปจฺจยเภสชฺช ปริกฺขารานํ.\csromanspacer{} ยโต จ โข อานนฺท ตถาคตา โลเก อุปฺปชฺชนฺติ อรหนฺโต สมฺมาสมฺพุทฺธา, อถ อญฺญติตฺถิยา ปริพฺพาชกา อสกฺกตา โหนฺติ อครุกตา อมานิตา อปูชิตา อนปจิตา, น ลาภิโน จีวรปิณฺฑปาตเสนาสนคิลานปจฺจยเภสชฺช ปริกฺขารานํ.\csromanspacer{} ตถาคโตว\footnote{ตถาคโต เจว (สฺยา)} ทานิ สกฺกโต โหติ ครุกโต มานิโต ปูชิโต อปจิโต, ลาภี จีวร ปิณฺฑปาต เสนาสน คิลานปจฺจย เภสชฺช ปริกฺขารานํ ภิกฺขุสํโฆ จาติ.}

\prose{อถ โข ภควา เอตมตฺถํ วิทิตฺวา ตายํ เวลายํ อิมํ อุทานํ อุทาเนสิ\csromandash{}}

\csromangathameasure{“โอภาสติ ตาว โส กิมิ, \\ ยาว น อุนฺนมเต ปภงฺกโร.}
\csromangathagroup{\csromangathastanza{“โอภาสติ ตาว โส กิมิ, \\ ยาว น อุนฺนมเต\footnote{อุคฺคมติ (สี), อุนฺนมติ (สฺยา)} ปภงฺกโร.}}

\prose{(ส)\footnote{( ) นตฺถิ สี-สฺยา-โปตฺถเกสุ.} เวโรจนมฺหิ อุคฺคเต,}

\prose{หตปฺปโภ โหติ น จาปิ ภาสติ.}

\csromangathameasure{เอวํ โอภาสิตเมว ตกฺกิกานํ, \\ ยาว สมฺมาสมฺพุทฺธา โลเก นุปฺปชฺชนฺติ. \\ น ตกฺกิกา สุชฺฌนฺติ น จาปิ สาวกา, \\ ทุทฺทิฏฺฐี น ทุกฺขา ปมุจฺจเร”ติ.ทสมํ.}
\csromangathagroup{\csromangathastanza{เอวํ โอภาสิตเมว ตกฺกิกานํ\footnote{ติตฺถิยานํ (สี, สฺยา, อิ)}, \\ ยาว สมฺมาสมฺพุทฺธา โลเก นุปฺปชฺชนฺติ. \\ น ตกฺกิกา สุชฺฌนฺติ น จาปิ สาวกา, \\ ทุทฺทิฏฺฐี น ทุกฺขา ปมุจฺจเร”ติ.\csromanspacer{}\csromanspacer{}ทสมํ.}}

\csromanmatikaanchor{mtk.91}
\csromantocmarkref{subsection}{8. สํฆเภทสุตฺต}{mtk.91}
\bookmark[dest={mtk.91},level=2]{8. สํฆเภทสุตฺต}
\titlehead{\textbf{ตสฺสุทฺทานํ}}
\csromangathameasure{อายุชฏิลเวกฺขณา, ตโย ติตฺถิยา สุภูติ. \\ คณิกา อุปาติ นวโม, อุปฺปชฺชนฺติ จ เต ทสาติ. ชจฺจนฺธวคฺโค ฉฏฺโฐ.}
\csromangathagroup{\csromangathastanza{อายุชฏิลเวกฺขณา, ตโย ติตฺถิยา สุภูติ. \\ คณิกา อุปาติ นวโม, อุปฺปชฺชนฺติ จ เต ทสาติ.\csromanspacer{} ชจฺจนฺธวคฺโค ฉฏฺโฐ.}}

\chapterheadpageread{อุทานปาฬิ \textbf{7. จูฬวคฺค}}
\csromanheader{2}{1. ปฐมลกุณฺฑกภทฺทิยสุตฺต}
\proseitem{61}{\csromanfolio{170}เอวํ เม สุตํ\csromandash{}เอกํ สมยํ ภควา สาวตฺถิยํ วิหรติ เชตวเน อนาถปิณฺฑิกสฺส อาราเม.\csromanspacer{} เตน โข ปน สมเยน อายสฺมา สาริปุตฺโต อายสฺมนฺตํ ลกุณฺฑกภทฺทิยํ อเนกปริยาเยน ธมฺมิยา กถาย สนฺทสฺเสติ สมาทเปติ\footnote{สมาทาเปติ (?)} สมุตฺเตเชติ สมฺปหํเสติ.}

\prose{อถ โข อายสฺมโต ลกุณฺฑกภทฺทิยสฺส อายสฺมตา สาริปุตฺเตน อเนกปริยาเยน ธมฺมิยา กถาย สนฺทสฺสิยมานสฺส สมาทปิยมานสฺส สมุตฺเตชิยมานสฺส สมฺปหํสิยมานสฺส อนุปาทาย อาสเวหิ จิตฺตํ วิมุจฺจิ.}

\prose{อทฺทสา โข ภควา อายสฺมนฺตํ ลกุณฺฑกภทฺทิยํ อายสฺมตา สาริปุตฺเตน อเนกปริยาเยน ธมฺมิยา กถาย สนฺทสฺสิยมานํ สมาทปิยมานํ สมุตฺเตชิยมานํ สมฺปหํสิยมานํ อนุปาทาย อาสเวหิ จิตฺตํ วิมุตฺตํ\footnote{วิมุตฺตจิตฺตํ (?)}.}

\prose{อถ โข ภควา เอตมตฺถํ วิทิตฺวา ตายํ เวลายํ อิมํ อุทานํ อุทาเนสิ\csromandash{} * “อุทฺธํ อโธ สพฺพธิ วิปฺปมุตฺโต, อยํหมสฺมีติ\footnote{อยมหมสฺมีติ (สี, สฺยา, อิ)} อนานุปสฺสี.}

\csromanmatikaanchor{mtk.92}
\csromantocmarkref{subsection}{9. สธายมานสุตฺต}{mtk.92}
\bookmark[dest={mtk.92},level=2]{9. สธายมานสุตฺต}
\prose{เอวํ วิมุตฺโต อุทตาริ โอฆํ, อติณฺณปุพฺพํ อปุนพฺภวายา”ติ.\csromanspacer{} ปฐมํ.}

\csromanheader{2}{2. ทุติยลกุณฺฑกภทฺทิยสุตฺต}
\proseitem{62}{เอวํ เม สุตํ\csromandash{}เอกํ สมยํ ภควา สาวตฺถิยํ วิหรติ เชตวเน อนาถปิณฺฑิกสฺส อาราเม.\csromanspacer{} เตน โข ปน สมเยน อายสฺมา สาริปุตฺโต อายสฺมนฺตํ ลกุณฺฑกภทฺทิยํ เสขํ\footnote{เสกฺโขติ (สฺยา), เสโขติ (อิ)} มญฺญมาโน ภิยฺโยโส มตฺตาย อเนกปริยาเยน ธมฺมิยา กถาย สนฺทสฺเสติ สมาทเปติ สมุตฺเตเชติ สมฺปหํเสติ.}

\chapterheadpageread{7. จูฬวคฺค}
\csromanmatikaanchor{mtk.93}
\csromantocmarkref{subsection}{10. จูฬปนฺถกสุตฺต}{mtk.93}
\bookmark[dest={mtk.93},level=2]{10. จูฬปนฺถกสุตฺต}
\prose{\csromanfolio{171}อทฺทสา โข ภควา อายสฺมนฺตํ สาริปุตฺตํ อายสฺมนฺตํ ลกุณฺฑกภทฺทิยํ เสขํ มญฺญมานํ ภิยฺโยโส มตฺตาย อเนกปริยาเยน ธมฺมิยา กถาย สนฺทสฺเสนฺตํ สมาทเปนฺตํ สมุตฺเตเชนฺตํ สมฺปหํเสนฺตํ.}

\prose{อถ โข ภควา เอตมตฺถํ วิทิตฺวา ตายํ เวลายํ อิมํ อุทานํ อุทาเนสิ\csromandash{}}

\csromangathameasure{“อจฺเฉจฺฉิ วฏฺฏํ พฺยคา นิราสํ, วิสุกฺขา สริตา น สนฺทติ. \\ ฉินฺนํ วฏฺฏํ น วตฺตติ, เอเสวนฺโต ทุกฺขสฺสา”ติ.ทุติยํ.}
\csromangathagroup{\csromangathastanza{“อจฺเฉจฺฉิ\footnote{อจฺเฉชฺชิ (ก-สี), อจฺฉิชฺชิ (ก, สี, สฺยา), อฉิชฺชิ (ก)} วฏฺฏํ พฺยคา นิราสํ, วิสุกฺขา สริตา น สนฺทติ. \\ ฉินฺนํ วฏฺฏํ น วตฺตติ, เอเสวนฺโต ทุกฺขสฺสา”ติ.\csromanspacer{}\csromanspacer{}ทุติยํ.}}

\csromanheader{2}{3. ปฐมสตฺตสุตฺต}
\csromanmatikaanchor{mtk.94}
\csromantocmarknopage{chapter}{6. ชจฺจนฺธวคฺค}{mtk.94}
\bookmark[dest={mtk.94},level=0]{6. ชจฺจนฺธวคฺค}
\proseitem{63}{เอวํ เม สุตํ\csromandash{}เอกํ สมยํ ภควา สาวตฺถิยํ วิหรติ เชตวเน อนาถปิณฺฑิกสฺส อาราเม.\csromanspacer{} เตน โข ปน สมเยน สาวตฺถิยา มนุสฺสา เยภุยฺเยน กาเมสุ อติเวลํ สตฺตา ( )\footnote{(โหนฺติ) (พหูสุ) อฏฺฐกถาย สํสนฺเทตพฺพํ.} รตฺตา คิทฺธา คถิตา\footnote{คธิตา (ก)} มุจฺฉิตา อชฺโฌสนฺนา\footnote{อชฺโฌปนฺนา (ก)} สมฺมตฺตกชาตา กาเมสุ วิหรนฺติ.}

\prose{อถ โข สมฺพหุลา ภิกฺขู ปุพฺพณฺหสมยํ นิวาเสตฺวา ปตฺตจีวรมาทาย สาวตฺถิยํ ปิณฺฑาย ปาวิสิํสุ, สาวตฺถิยํ ปิณฺฑาย จริตฺวา ปจฺฉาภตฺตํ ปิณฺฑปาตปฏิกฺกนฺตา เยน ภควา เตนุปสงฺกมิํสุ, อุปสงฺกมิตฺวา ภควนฺตํ อภิวาเทตฺวา เอกมนฺตํ นิสีทิํสุ, เอกมนฺตํ นิสินฺนา โข เต ภิกฺขู ภควนฺตํ เอตทโวจุํ “อิธ ภนฺเต สาวตฺถิยา มนุสฺสา เยภุยฺเยน กาเมสุ อติเวลํ สตฺตา รตฺตา คิทฺธา คถิตา มุจฺฉิตา อชฺโฌสนฺนา สมฺมตฺตกชาตา กาเมสุ วิหรนฺตี”ติ.}

\prose{อถ โข ภควา เอตมตฺถํ วิทิตฺวา ตายํ เวลายํ อิมํ อุทานํ อุทาเนสิ\csromandash{}}

\csromangathameasure{*“กาเมสุ สตฺตา กามสงฺคสตฺตา, \\ สํโยชเน วชฺชมปสฺสมานา. \\ น หิ ชาตุ สํโยชนสงฺคสตฺตา, \\ โอฆํ ตเรยฺยุํ วิปุลํ มหนฺตนฺ”ติ.ตติยํ.}
\csromangathagroup{\csromangathastanza{\csromansymbolfootnote{*}{ขุ 10. 174 ปิฏฺเฐปิ.}“กาเมสุ สตฺตา กามสงฺคสตฺตา, \\ สํโยชเน วชฺชมปสฺสมานา. \\ น หิ ชาตุ สํโยชนสงฺคสตฺตา, \\ โอฆํ ตเรยฺยุํ วิปุลํ มหนฺตนฺ”ติ.\csromanspacer{}\csromanspacer{}ตติยํ.}}

\csromanheader{2}{อุทานปาฬิ \textbf{4. ทุติยสตฺตสุตฺต}}
\csromanmatikaanchor{mtk.95}
\csromantocmarkref{subsection}{1. อายุสงฺขาโรสฺสชฺชนสุตฺต}{mtk.95}
\bookmark[dest={mtk.95},level=2]{1. อายุสงฺขาโรสฺสชฺชนสุตฺต}
\proseitem{64}{\csromanfolio{172}เอวํ เม สุภํ\csromandash{}เอกํ สมยํ ภควา สาวตฺถิยํ วิหรติ เชตวเน อนาถปิณฺฑิกสฺส อาราเม.\csromanspacer{} เตน โข ปน สมเยน สาวตฺถิยา มนุสฺสา เยภุยฺเยน กาเมสุ สตฺตา ( )\footnote{(โหนฺติ) (พหูสุ) อฏฺฐกถาย สํสนฺเทตพฺพํ.} รตฺตา คิทฺธา คถิตา มุจฺฉิตา อชฺโฌสนฺนา อนฺธีกตา สมฺมตฺตกชาตา กาเมสุ วิหรนฺติ.}

\prose{อถ โข ภควา ปุพฺพณฺหสมยํ นิวาเสตฺวา ปตฺตจีวรมาทาย สาวตฺถิํ ปิณฺฑาย ปาวิสิ.\csromanspacer{} อทฺทสา โข ภควา สาวตฺถิยา เต มนุสฺเส เยภุยฺเยน กาเมสุ สตฺเต รตฺเต คิทฺเธ คถิเต มุจฺฉิเต อชฺโฌสนฺเน อนฺธีกเต สมฺมตฺตกชาเต กาเมสุ วิหรนฺเต.}

\prose{อถ โข ภควา เอตมตฺถํ วิทิตฺวา ตายํ เวลายํ อิมํ อุทานํ อุทาเนสิ\csromandash{} * “กามนฺธา ชาลสญฺฉนฺนา, ตณฺหาฉทนฉาทิตา.}

\prose{ปมตฺตพนฺธุนา พทฺธา, มจฺฉาว กุมินามุเข.}

\prose{ชรามรณมเนฺวนฺติ\footnote{ชรามรณํ คจฺฉนฺติ (สี, สฺยา)}, วจฺโฉ ขีรปโกว มาตรนฺ”ติ.\csromanspacer{}\csromanspacer{}จตุตฺถํ.}

\csromanheader{2}{5. อปรลกุณฺฑกภทฺทิยสุตฺต}
\proseitem{65}{เอวํ เม สุตํ\csromandash{}เอกํ สมยํ ภควา สาวตฺถิยํ วิหรติ เชตวเน อนาถปิณฺฑิกสฺส อาราเม.\csromanspacer{} เตน โข ปน สมเยน อายสฺมา ลกุณฺฑกภทฺทิโย สมฺพหุลานํ ภิกฺขูนํ ปิฏฺฐิโต ปิฏฺฐิโต เยน ภควา เตนุปสงฺกมิ.}

\prose{อทฺทสา โข ภควา อายสฺมนฺตํ ลกุณฺฑกภทฺทิยํ ทูรโตว สมฺพหุลานํ ภิกฺขูนํ ปิฏฺฐิโต ปิฏฺฐิโต อาคจฺฉนฺตํ ทุพฺพณฺณํ ทุทฺทสิกํ โอโกฏิมกํ เยภุยฺเยน ภิกฺขูนํ ปริภูตรูปํ, ทิสฺวาน ภิกฺขู อามนฺเตสิ\csromandash{}}

\prose{ปสฺสถ โน ตุเมฺห ภิกฺขเว เอตํ ภิกฺขุํ ทูรโตว สมฺพหุลานํ ภิกฺขูนํ ปิฏฺฐิโต ปิฏฺฐิโต อาคจฺฉนฺตํ ทุพฺพณฺณํ ทุทฺทสิกํ โอโกฏิมกํ เยภุยฺเยน ภิกฺขูนํ ปริภูตรูปนฺติ.\csromanspacer{} เอวํ ภนฺเตติ.}

\chapterheadpageread{7. จูฬวคฺค}
\csromanmatikaanchor{mtk.96}
\csromanmatikaanchor{mtk.112}
\csromantocmarkref{subsection}{2. สตฺตชฏิลสุตฺต}{mtk.96}
\bookmark[dest={mtk.96},level=2]{2. สตฺตชฏิลสุตฺต}
\prose{\csromanfolio{173}เอโส ภิกฺขเว ภิกฺขุ มหิทฺธิโก มหานุภาโว, น จ สา สมาปตฺติ สุลภรูปา, ยา เตน ภิกฺขุนา อสมาปนฺนปุพฺพา.\csromanspacer{} ยสฺส จตฺถาย\footnote{ยสฺสตฺถาย (สี, ก)} กุลปุตฺตา สมฺมเทว อคารสฺมา อนคาริยํ ปพฺพชนฺติ, ตทนุตฺตรํ พฺรหฺมจริยปริโยสานํ ทิฏฺเฐว ธมฺเม สยํ อภิญฺญา สจฺฉิกตฺวา อุปสมฺปชฺช วิหรตีติ.}

\prose{อถ โข ภควา เอตมตฺถํ วิทิตฺวา ตายํ เวลายํ อิมํ อุทานํ อุทาเนสิ\csromandash{}}

\prose{“เนลงฺโค เสตปจฺฉาโท, เอกาโร วตฺตตี รโถ.}

\prose{อนีฆํ ปสฺส อายนฺตํ, ฉินฺนโสตํ อพนฺธนนฺ”ติ.\csromanspacer{}\csromanspacer{}ปญฺจมํ.}

\csromanheader{2}{6. ตณฺหาสงฺขยสุตฺต}
\proseitem{66}{เอวํ เม สุตํ\csromandash{}เอกํ สมยํ ภควา สาวตฺถิยํ วิหรติ เชตวเน อนาถปิณฺฑิกสฺส อาราเม.\csromanspacer{} เตน โข ปน สมเยน อายสฺมา อญฺญาสิโกณฺฑญฺโญ\footnote{อญฺญาตโกณฺฑญฺโญ (สพฺพตฺถ)} ภควโต อวิทูเร นิสินฺโน โหติ ปลฺลงฺกํ อาภุชิตฺวา อุชุํ กายํ ปณิธาย ตณฺหาสงฺขยวิมุตฺติํ ปจฺจเวกฺขมาโน.}

\prose{อทฺทสา โข ภควา อายสฺมนฺตํ อญฺญาสิโกณฺฑญฺญํ อวิทูเร นิสินฺนํ ปลฺลงฺกํ อาภุชิตฺวา อุชุํ กายํ ปณิธาย ตณฺหาสงฺขยวิมุตฺติํ ปจฺจเวกฺขมานํ.}

\prose{อถ โข ภควา เอตมตฺถํ วิทิตฺวา ตายํ เวลายํ อิมํ อุทานํ อุทาเนสิ\csromandash{}}

\prose{“ยสฺส มูลํ ฉมา นตฺถิ, ปณฺณา นตฺถิ กุโต ลตา.}

\prose{ตํ ธีรํ พนฺธนา มุตฺตํ, โก ตํ นินฺทิตุมรหติ.}

\prose{เทวาปิ นํ ปสํสนฺติ, พฺรหฺมุนาปิ ปสํสิโต”ติ.\csromanspacer{}\csromanspacer{}ฉฏฺฐํ.}

\csromanheader{2}{7. ปปญฺจกฺขยสุตฺต}
\proseitem{67}{เอวํ เม สุตํ\csromandash{}เอกํ สมยํ ภควา สาวตฺถิยํ วิหรติ เชตวเน อนาถปิณฺฑิกสฺส อาราเม.\csromanspacer{} เตน โข ปน สมเยน ภควา อตฺตโน ปปญฺจสญฺญาสงฺขาปหานํ ปจฺจเวกฺขมาโน นิสินฺโน โหติ.}

\gambhira{อุทานปาฬิ}
\prose{\csromanfolio{174}อถ โข ภควา อตฺตโน ปปญฺจสญฺญาสงฺขาปหานํ วิทิตฺวา ตายํ เวลายํ อิมํ อุทานํ อุทาเนสิ\csromandash{}}

\csromangathameasure{*“ยสฺส ปปญฺจา ฐิติ จ นตฺถิ, \\ สนฺทานํ ปลิฆญฺจ วีติวตฺโต. \\ ตํ นิตฺตณฺหํ มุนิํ จรนฺตํ, \\ นาวชานาติ สเทวโกปิ โลโก”ติ.สตฺตมํ.}
\csromangathagroup{\csromangathastanza{\csromansymbolfootnote{*}{ขุ 10. 32 ปิฏฺเฐปิ.}“ยสฺส ปปญฺจา ฐิติ จ นตฺถิ, \\ สนฺทานํ ปลิฆญฺจ วีติวตฺโต. \\ ตํ นิตฺตณฺหํ มุนิํ จรนฺตํ, \\ นาวชานาติ สเทวโกปิ โลโก”ติ.\csromanspacer{}\csromanspacer{}สตฺตมํ.}}

\csromanheader{2}{8. กจฺจานสุตฺต}
\proseitem{68}{เอวํ เม สุตํ\csromandash{}เอกํ สมยํ ภควา สาวตฺถิยํ วิหรติ เชตวเน อนาถปิณฺฑิกสฺส อาราเม.\csromanspacer{} เตน โข ปน สมเยน อายสฺมา มหากจฺจาโน ภควโต อวิทูเร นิสินฺโน โหติ ปลฺลงฺกํ อาภุชิตฺวา อุชุํ กายํ ปณิธาย กายคตาย สติยา อชฺฌตฺตํ ปริมุขํ สูปฏฺฐิตาย.}

\prose{อทฺทสา โข ภควา อายสฺมนฺตํ มหากจฺจานํ อวิทูเร นิสินฺนํ ปลฺลงฺกํ อาภุชิตฺวา อุชุํ กายํ ปณิธาย กายคตาย สติยา อชฺฌตฺตํ ปริมุขํ สูปฏฺฐิตาย.}

\prose{อถ โข ภควา เอตมตฺถํ วิทิตฺวา ตายํ เวลายํ อิมํ อุทานํ อุทาเนสิ\csromandash{}}

\csromangathameasure{“ยสฺส สิยา สพฺพทา สติ, \\ สตตํ กายคตา อุปฏฺฐิตา. \\ โน จสฺส โน จ เม สิยา, \\ น ภวิสฺสติ น จ เม ภวิสฺสติ. \\ อนุปุพฺพวิหาริ ตตฺถ โส, \\ กาเลเนว ตเร วิสตฺติกนฺ”ติ.อฏฺฐมํ.}
\csromangathagroup{\csromangathastanza{“ยสฺส สิยา สพฺพทา สติ, \\ สตตํ กายคตา อุปฏฺฐิตา. \\ โน จสฺส โน จ เม สิยา, \\ น ภวิสฺสติ น จ เม ภวิสฺสติ. \\ อนุปุพฺพวิหาริ ตตฺถ โส, \\ กาเลเนว ตเร วิสตฺติกนฺ”ติ.\csromanspacer{}\csromanspacer{}อฏฺฐมํ.}}

\csromanheader{2}{9. อุทปานสุตฺต}
\proseitem{69}{เอวํ เม สุตํ\csromandash{}เอกํ สมยํ ภควา มลฺเลสุ จาริกํ จรมาโน มหตา ภิกฺขุสํเฆน สทฺธิํ เยน ถูณํ\footnote{ถูนํ (สี, สฺยา, อิ)} นาม มลฺลานํ พฺราหฺมณคาโม ตทวสริ.\csromanspacer{} อสฺโสสุํ โข ถูเณยฺยกา พฺราหฺมณคหปติกา “สมโณ ขลุ โภ โคตโม สกฺยปุตฺโต สกฺยกุลา ปพฺพชิโต มลฺเลสุ จาริกํ}

\vspace{\tipitakaparskip}
\csromancenter{จูฬวคฺค จรมาโน มหตา ภิกฺขุสํเฆน สทฺธิํ ถูณํ อนุปฺปตฺโต”ติ. ( )\footnote{(อถ โข เต ถูเณยฺยกา พฺราหฺมณคหปติกา) (?)} อุทปานํ ติณสฺส จ ภุสสฺส จ ยาว มุขโต ปูเรสุํ “มา เต มุณฺฑกา สมณกา ปานียํ อปํสู”ติ.}
\prose{\csromanfolio{175}อถ โข ภควา มคฺคา โอกฺกมฺม เยน รุกฺขมูลํ เตนุปสงฺกมิ, อุปสงฺกมิตฺวา ปญฺญตฺเต อาสเน นิสีทิ, นิสชฺช โข ภควา อายสฺมนฺตํ อานนฺทํ อามนฺเตสิ “อิงฺฆ เม ตฺวํ อานนฺท เอตมฺหา อุทปานา ปานียํ อาหรา”ติ.}

\prose{เอวํ วุตฺเต อายสฺมา อานนฺโท ภควนฺตํ เอตทโวจ “อิทานิ โส ภนฺเต อุทปาโน ถูเณยฺยเกหิ พฺราหฺมณคหปติเกหิ ติณสฺส จ ภุสสฺส จ ยาว มุขโต ปูริโต ‘มา เต มุณฺฑกา สมณกา ปานียํ อปํสู’ติ”.}

\csromanmatikaanchor{mtk.97}
\csromantocmarkref{subsection}{3. ปจฺจเวกฺขณสุตฺต}{mtk.97}
\bookmark[dest={mtk.97},level=2]{3. ปจฺจเวกฺขณสุตฺต}
\prose{ทุติยมฺปิ โข ฯเปฯ.\csromanspacer{} ตติยมฺปิ โข ภควา อายสฺมนฺตํ อานนฺทํ อามนฺเตสิ “อิงฺฆ เม ตฺวํ อานนฺท เอตมฺหา อุทปานา ปานียํ อาหรา”ติ. “เอวํ ภนฺเต”ติ โข อายสฺมา อานนฺโท ภควโต ปฏิสฺสุตฺวา ปตฺตํ คเหตฺวา เยน โส อุทปาโน เตนุปสงฺกมิ.\csromanspacer{} อถ โข โส อุทปาโน อายสฺมนฺเต อานนฺเท อุปสงฺกมนฺเต สพฺพํ ตํ ติณญฺจ ภุสญฺจ มุขโต โอวมิตฺวา อจฺฉสฺส อุทกสฺส อนาวิลสฺส วิปฺปสนฺนสฺส ยาว มุขโต ปูริโต วิสฺสนฺทนฺโต\footnote{วิสฺสนฺโท (ก)} มญฺเญ อฏฺฐาสิ.}

\prose{อถ โข อายสฺมโต อานนฺทสฺส เอตทโหสิ “อจฺฉริยํ วต โภ, อพฺภุตํ วต โภ, ตถาคตสฺส มหิทฺธิกตา มหานุภาวตา, อยํ หิ โส อุทปาโน มยิ อุปสงฺกมนฺเต สพฺพํ ตํ ติณญฺจ ภุสญฺจ มุขโต โอวมิตฺวา อจฺฉสฺส อุทกสฺส อนาวิลสฺส วิปฺปสนฺนสฺส ยาว มุขโต ปูริโต วิสฺสนฺทนฺโต มญฺเญ ฐิโต”ติ ปตฺเตน ปานียํ อาทาย เยน ภควา เตนุปสงฺกมิ, อุปสงฺกมิตฺวา ภควนฺตํ เอตทโวจ “อจฺฉริยํ ภนฺเต, อพฺภุตํ ภนฺเต, ตถาคตสฺส มหิทฺธิกตา มหานุภาวตา, อยํ หิ โส ภนฺเต อุทปาโน มยิ อุปสงฺกมนฺเต สพฺพํ ตํ ติณญฺจ ภุสญฺจ มุขโต โอวมิตฺวา อจฺฉสฺส อุทกสฺส อนาวิลสฺส วิปฺปสนฺนสฺส ยาว มุขโต ปูริโต วิสฺสนฺทนฺโต มญฺเญ อฏฺฐาสิ, ปิวตุ ภควา ปานียํ, ปิวตุ สุคโต ปานียนฺ”ติ.}

\gambhira{อุทานปาฬิ}
\prose{\csromanfolio{176}อถ โข ภควา เอตมตฺถํ วิทิตฺวา ตายํ เวลายํ อิมํ อุทานํ อุทาเนสิ\csromandash{}}

\csromangathameasure{“กิํ กยิรา อุทปาเนน, \\ อาปา เจ สพฺพทา สิยุํ. \\ ตณฺหาย มูลโต เฉตฺวา, \\ กิสฺส ปริเยสนํ จเร”ติ.นวมํ.}
\csromangathagroup{\csromangathastanza{“กิํ กยิรา อุทปาเนน, \\ อาปา เจ สพฺพทา สิยุํ. \\ ตณฺหาย มูลโต เฉตฺวา, \\ กิสฺส ปริเยสนํ จเร”ติ.\csromanspacer{}\csromanspacer{}นวมํ.}}

\csromanmatikaanchor{mtk.98}
\csromantocmarkref{subsection}{4. ปฐมนานาติตฺถิยสุตฺต}{mtk.98}
\bookmark[dest={mtk.98},level=2]{4. ปฐมนานาติตฺถิยสุตฺต}
\csromanheader{2}{10. อุเตนสุตฺต}
\proseitem{70}{เอวํ เม สุตํ\csromandash{}เอกํ สมยํ ภควา โกสมฺพิยํ วิหรติ โฆสิตาราเม.\csromanspacer{} เตน โข ปน สมเยน รญฺโญ อุเตนสฺส\footnote{อุเทนสฺส (สี, สฺยา, อิ) 2. ปญฺจ อิตฺถิสตานิ (สี, สฺยา, อิ)} อุยฺยานคตสฺส อนฺเตปุรํ ทฑฺฒํ โหติ, ปญฺจ จ อิตฺถิสตานิ2 กาลงฺกตานิ โหนฺติ สามาวตีปมุขานิ.}

\prose{อถ โข สมฺพหุลา ภิกฺขู ปุพฺพณฺหสมยํ นิวาเสตฺวา ปตฺตจีวรมาทาย โกสมฺพิํ ปิณฺฑาย ปาวิสิํสุ, โกสมฺพิยํ ปิณฺฑาย จริตฺวา ปจฺฉาภตฺตํ ปิณฺฑปาตปฏิกฺกนฺตา เยน ภควา เตนุปสงฺกมิํสุ, อุปสงฺกมิตฺวา ภควนฺตํ อภิวาเทตฺวา เอกมนฺตํ นิสีทิํสุ, เอกมนฺตํ นิสินฺนา โข เต ภิกฺขู ภควนฺตํ เอตทโวจุํ “อิธ ภนฺเต รญฺโญ อุเตนสฺส อุยฺยานคตสฺส อนฺเตปุรํ ทฑฺฒํ, ปญฺจ จ อิตฺถิสตานิ กาลงฺกตานิ สามาวตีปมุขานิ, ตาสํ ภนฺเต อุปาสิกานํ กา คติ, โก อภิสมฺปราโย”ติ.}

\prose{สนฺเตตฺถ ภิกฺขเว อุปาสิกาโย โสตาปนฺนา, สนฺติ สกทาคามินิโย, สนฺติ อนาคามินิโย, สพฺพา ตา ภิกฺขเว อุปาสิกาโย อนิปฺผลา กาลงฺกตาติ.}

\prose{อถ โข ภควา เอตมตฺถํ วิทิตฺวา ตายํ เวลายํ อิมํ อุทานํ อุทาเนสิ\csromandash{}}

\csromangathameasure{*“โมหสมฺพนฺธโน โลโก, ภพฺพรูโปว ทิสฺสติ. \\ อุปธิพนฺธโน พาโล, ตมสา ปริวาริโต. \\ สสฺสโตริว ขายติ, ปสฺสโต นตฺถิ กิญฺจนนฺ”ติ.ทสมํ.}
\csromangathagroup{\csromangathastanza{\csromansymbolfootnote{*}{ขุ 10. 53 ปิฏฺเฐปิ.}“โมหสมฺพนฺธโน โลโก, ภพฺพรูโปว ทิสฺสติ. \\ อุปธิพนฺธโน\footnote{อุปธิสมฺพนฺธโน (ก-สี)} พาโล, ตมสา ปริวาริโต. \\ สสฺสโตริว\footnote{อุปธิสมฺพนฺธโน (ก-สี)} ขายติ, ปสฺสโต นตฺถิ กิญฺจนนฺ”ติ.\csromanspacer{}\csromanspacer{}ทสมํ.}}

\chapterheadpageread{8. ปาฏลิคามิยวคฺค \textbf{ตสฺสุทฺทานํ}}
\csromangathameasure{เทฺว ภทฺทิยา เทฺว จ สตฺตา, ลกุณฺฑโก ตณฺหากฺขโย. \\ ปปญฺจกฺขโย จ กจฺจาโน, อุทปานญฺจ อุเตโนติ. จูฬวคฺโค สตฺตโม.}
\csromangathagroup{\csromangathastanza{\csromanfolio{177}เทฺว ภทฺทิยา เทฺว จ สตฺตา, ลกุณฺฑโก ตณฺหากฺขโย. \\ ปปญฺจกฺขโย จ กจฺจาโน, อุทปานญฺจ อุเตโนติ.\csromanspacer{} จูฬวคฺโค\footnote{จุลฺลวคฺโค (สี), จูลวคฺโค (อิ)} สตฺตโม.}}

\chapterhead{8. ปาฏลิคามิยวคฺค}
\csromanheader{2}{1. ปฐมนิพฺพานปฏิสํยุตฺตสุตฺต}
\proseitem{71}{เอวํ เม สุตํ\csromandash{}เอกํ สมยํ ภควา สาวตฺถิยํ วิหรติ เชตวเน อนาถปิณฺฑิกสฺส อาราเม.\csromanspacer{} เตน โข ปน สมเยน ภควา ภิกฺขู นิพฺพานปฏิสํยุตฺตาย ธมฺมิยา กถาย สนฺทสฺเสติ สมาทเปติ สมุตฺเตเชติ สมฺปหํเสติ.\csromanspacer{} เต’ธ ภิกฺขู\footnote{เต จ ภิกฺขู (สี, สฺยา, อิ, ตทฏฺฐกถาปิ โอโลเกตพฺพา.)} อฏฺฐิํ กตฺวา\footnote{อฏฺฐีกตฺวา (สี, สฺยา), อฏฺฐิกตฺวา (อิ)} มนสิ กตฺวา สพฺพํ เจตโส\footnote{สพฺพํ เจตสา (อิติปิ อญฺญสุตฺเตสุ)} สมนฺนาหริตฺวา โอหิตโสตา ธมฺมํ สุณนฺติ.}

\prose{อถ โข ภควา เอตมตฺถํ วิทิตฺวา ตายํ เวลายํ อิมํ อุทานํ อุทาเนสิ\csromandash{} “อตฺถิ ภิกฺขเว ตทายตนํ, ยตฺถ เนว ปถวี น อาโป น เตโช น วาโย, น อากาสานญฺจายตนํ น วิญฺญาณญฺจายตนํ น อากิญฺจญฺญายตนํ น เนวสญฺญานาสญฺญายตนํ, นายํ โลโก น ปรโลโก, น อุโภ จนฺทิมสูริยา.\csromanspacer{} ตตฺราปาหํ ภิกฺขเว เนว อาคติํ วทามิ น คติํ น ฐิติํ น จุติํ น อุปปตฺติํ, อปฺปติฏฺฐํ อปฺปวตฺตํ อนารมฺมณเมเวตํ.\csromanspacer{} เอเสวนฺโต ทุกฺขสฺสา”ติ.\csromanspacer{}\csromanspacer{}ปฐมํ.}

\csromanheader{2}{2. ทุติยนิพฺพานปฏิสํยุตฺตสุตฺต}
\proseitem{72}{เอวํ เม สุตํ\csromandash{}เอกํ สมยํ ภควา สาวตฺถิยํ วิหรติ เชตวเน อนาถปิณฺฑิกสฺส อาราเม.\csromanspacer{} เตน โข ปน สมเยน ภควา ภิกฺขู}

\vspace{\tipitakaparskip}
\csromancenter{อุทานปาฬิ นิพฺพานปฏิสํยุตฺตาย ธมฺมิยา กถาย สนฺทสฺเสติ สมาทเปติ สมุตฺเตเชติ สมฺปหํเสติ.\csromanspacer{} เต’ธ ภิกฺขู อฏฺฐิํ กตฺวา มนสิ กตฺวา สพฺพํ เจตโส สมนฺนาหริตฺวา โอหิตโสตา ธมฺมํ สุณนฺติ.}
\prose{\csromanfolio{178}อถ โข ภควา เอตมตฺถํ วิทิตฺวา ตายํ เวลายํ อิมํ อุทานํ อุทาเนสิ\csromandash{}}

\csromangathameasure{“ทุทฺทสํ อนตํ นาม, น หิ สจฺจํ สุทสฺสนํ. \\ ปฏิวิทฺธา ตณฺหา ชานโต, ปสฺสโต นตฺถิ กิญฺจนนฺ”ติ.ทุติยํ.}
\csromangathagroup{\csromangathastanza{“ทุทฺทสํ อนตํ นาม, น หิ สจฺจํ สุทสฺสนํ. \\ ปฏิวิทฺธา ตณฺหา ชานโต, ปสฺสโต นตฺถิ กิญฺจนนฺ”ติ.\csromanspacer{}\csromanspacer{}ทุติยํ.}}

\csromanheader{2}{3. ตติยนิพฺพานปฏิสํยุตฺตสุตฺต}
\proseitem{73}{เอวํ เม สุตํ\csromandash{}เอกํ สมยํ ภควา สาวตฺถิยํ วิหรติ เชตวเน อนาถปิณฺฑิกสฺส อาราเม.\csromanspacer{} เตน โข ปน สมเยน ภควา ภิกฺขู นิพฺพานปฏิสํยุตฺตาย ธมฺมิยา กถาย สนฺทสฺเสติ สมาทเปติ สมุตฺเตเชติ สมฺปหํเสติ.\csromanspacer{} เต’ธ ภิกฺขู อฏฺฐิํ กตฺวา มนสิ กตฺวา สพฺพํ เจตโส สมนฺนาหริตฺวา โอหิตโสตา ธมฺมํ สุณนฺติ.}

\prose{อถ โข ภควา เอตมตฺถํ วิทิตฺวา ตายํ เวลายํ อิมํ อุทานํ อุทาเนสิ\csromandash{} “อตฺถิ ภิกฺขเว อชาตํ อภูตํ อกตํ อสงฺขตํ, โน เจตํ ภิกฺขเว อภวิสฺส อชาตํ อภูตํ อกตํ อสงฺขตํ, นยิธ ชาตสฺส ภูตสฺส กตสฺส สงฺขตสฺส นิสฺสรณํ ปญฺญาเยถ.\csromanspacer{} ยสฺมา จ โข ภิกฺขเว อตฺถิ อชาตํ อภูตํ อกตํ อสงฺขตํ.\csromanspacer{} ตสฺมา ชาตสฺส ภูตสฺส กตสฺส สงฺขตสฺส นิสฺสรณํ ปญฺญายตี”ติ.\csromanspacer{}\csromanspacer{}ตติยํ.}

\csromanheader{2}{4. จตุตฺถนิพฺพานปฏิสํยุตฺตสุตฺต}
\proseitem{74}{เอวํ เม สุตํ\csromandash{}เอกํ สมยํ ภควา สาวตฺถิยํ วิหรติ เชตวเน อนาถปิณฺฑิกสฺส อาราเม.\csromanspacer{} เตน โข ปน สมเยน ภควา ภิกฺขู นิพฺพานปฏิสํยุตฺตาย ธมฺมิยา กถาย สนฺทสฺเสติ สมาทเปติ สมุตฺเตเชติ สมฺปหํเสติ.\csromanspacer{} เต’ธ ภิกฺขู อฏฺฐิํ กตฺวา มนสิ กตฺวา สพฺพํ เจตโส สมนฺนาหริตฺวา โอหิตโสตา ธมฺมํ สุณนฺติ.}

\prose{อถ โข ภควา เอตมตฺถํ วิทิตฺวา ตายํ เวลายํ อิมํ อุทานํ อุทาเนสิ\csromandash{}}

\vspace{\tipitakaparskip}
\csromancenter{ปาฏลิคามิยวคฺค + “นิสฺสิตสฺส จลิตํ, อนิสฺสิตสฺส จลิตํ นตฺถิ, จลิเต อสติ ปสฺสทฺธิ, ปสฺสทฺธิยา สติ นติ น โหติ, นติยา อสติ อาคติคติ น โหติ, อาคติคติยา อสติ จุตูปปาโต น โหติ, จุตูปปาเต อสติ เนวิธ น หุรํ น อุภยมนฺตเรน\footnote{น อุภยมนฺตเร (สพฺพตฺถ) ม 3. 310 ปิฏฺเฐ; สํ 2. 284 ปิฏฺเฐ จ ปสฺสิตพฺพํ.}.\csromanspacer{} เอเสวนฺโต ทุกฺขสฺสา”ติ.\csromanspacer{}\csromanspacer{}จตุตฺถํ.}
\csromanheader{2}{5. จุนฺทสุตฺต}
\proseitem{75}{\csromanfolio{179}\csromansymbolfootnote{*}{ที 2. 105 ปิฏฺเฐปิ.}เอวํ เม สุตํ\csromandash{}เอกํ สมยํ ภควา มลฺเลสุ จาริกํ จรมาโน มหตา ภิกฺขุสํเฆน สทฺธิํ เยน ปาวา ตทวสริ.\csromanspacer{} ตตฺร สุทํ ภควา ปาวายํ วิหรติ จุนฺทสฺส กมฺมารปุตฺตสฺส อมฺพวเน.}

\prose{อสฺโสสิ โข จุนฺโท กมฺมารปุตฺโต “ภควา กิร มลฺเลสุ จาริกํ จรมาโน มหตา ภิกฺขุสํเฆน สทฺธิํ ปาวํ อนุปฺปตฺโต ปาวายํ วิหรติ มยฺหํ อมฺพวเน”ติ.\csromanspacer{} อถ โข จุนฺโท กมฺมารปุตฺโต เยน ภควา เตนุปสงฺกมิ, อุปสงฺกมิตฺวา ภควนฺตํ อภิวาเทตฺวา เอกมนฺตํ นิสีทิ, เอกมนฺตํ นิสินฺนํ โข จุนฺทํ กมฺมารปุตฺตํ ภควา ธมฺมิยา กถาย สนฺทสฺเสสิ สมาทเปสิ สมุตฺเตเชสิ สมฺปหํเสสิ.\csromanspacer{} อถ โข จุนฺโท กมฺมารปุตฺโต ภควตา ธมฺมิยา กถาย สนฺทสฺสิโต สมาทปิโต สมุตฺเตชิโต สมฺปหํสิโต ภควนฺตํ เอตทโวจ “อธิวาเสตุ เม ภนฺเต ภควา สฺวาตนาย ภตฺตํ สทฺธิํ ภิกฺขุสํเฆนา”ติ.\csromanspacer{} อธิวาเสสิ ภควา ตุณฺหีภาเวน.}

\prose{อถ โข จุนฺโท กมฺมารปุตฺโต ภควโต อธิวาสนํ วิทิตฺวา อุฏฺฐายาสนา ภควนฺตํ อภิวาเทตฺวา ปทกฺขิณํ กตฺวา ปกฺกามิ.\csromanspacer{} อถ โข จุนฺโท กมฺมารปุตฺโต ตสฺสา รตฺติยา อจฺจเยน สเก นิเวสเน ปณีตํ ขาทนียํ โภชนียํ ปฏิยาทาเปตฺวา ปหูตญฺจ สูกรมทฺทวํ, ภควโต กาลํ อาโรจาเปสิ “กาโล ภนฺเต นิฏฺฐิตํ ภตฺตนฺ”ติ.}

\csromanmatikaanchor{mtk.99}
\csromantocmarkref{subsection}{5. ทุติยนานาติตฺถิยสุตฺต}{mtk.99}
\bookmark[dest={mtk.99},level=2]{5. ทุติยนานาติตฺถิยสุตฺต}
\prose{อถ โข ภควา ปุพฺพณฺหสมยํ นิวาเสตฺวา ปตฺตจีวรมาทาย สทฺธิํ ภิกฺขุสํเฆน เยน จุนฺทสฺส กมฺมารปุตฺตสฺส นิเวสนํ เตนุปสงฺกมิ, อุปสงฺกมิตฺวา ปญฺญตฺเต อาสเน นิสีทิ, นิสชฺช โข ภควา จุนฺทํ กมฺมารปุตฺตํ อามนฺเตสิ “ยํ เต จุนฺท สูกรมทฺทวํ ปฏิยตฺตํ เตน มํ ปริวิส, ยํ ปนญฺญํ}

\vspace{\tipitakaparskip}
\csromancenter{อุทานปาฬิ ขาทนียํ โภชนียํ ปฏิยตฺตํ เตน ภิกฺขุสํฆํ ปริวิสา”ติ. “เอวํ ภนฺเต”ติ โข จุนฺโท กมฺมารปุตฺโต ภควโต ปฏิสฺสุตฺวา ยํ อโหสิ สูกรมทฺทวํ ปฏิยตฺตํ, เตน ภควนฺตํ ปริวิสิ, ยํ ปนญฺญํ ขาทนียํ โภชนียํ ปฏิยตฺตํ, เตน ภิกฺขุสํฆํ ปริวิสิ.}
\prose{\csromanfolio{180}อถ โข ภควา จุนฺทํ กมฺมารปุตฺตํ อามนฺเตสิ “ยํ เต จุนฺท สูกรมทฺทวํ อวสิฏฺฐํ, ตํ โสพฺเภ นิขณาหิ, นาหํ ตํ จุนฺท ปสฺสามิ สเทวเก โลเก สมารเก สพฺรหฺมเก สสฺสมณพฺราหฺมณิยา ปชาย สเทวมนุสฺสาย, ยสฺส ตํ ปริภุตฺตํ สมฺมา ปริณามํ คจฺเฉยฺย อญฺญตฺร ตถาคตสฺสา”ติ\footnote{อญฺญตฺร ตถาคเตนาติ (ก-สี)}. “เอวํ ภนฺเต”ติ โข จุนฺโท กมฺมารปุตฺโต ภควโต ปฏิสฺสุตฺวา ยํ อโหสิ สูกรมทฺทวํ อวสิฏฺฐํ, ตํ โสพฺเภ นิขณิตฺวา เยน ภควา เตนุปสงฺกมิ, อุปสงฺกมิตฺวา ภควนฺตํ อภิวาเทตฺวา เอกมนฺตํ นิสีทิ, เอกมนฺตํ นิสินฺนํ โข จุนฺทํ กมฺมารปุตฺตํ ภควา ธมฺมิยา กถาย สนฺทสฺเสตฺวา สมาทเปตฺวา สมุตฺเตเชตฺวา สมฺปหํเสตฺวา อุฏฺฐายาสนา ปกฺกามิ.}

\prose{อถ โข ภควโต จุนฺทสฺส กมฺมารปุตฺตสฺส ภตฺตํ ภุตฺตาวิสฺส ขโร อาพาโธ อุปฺปชฺชิ, โลหิตปกฺขนฺทิกา ปพาฬฺหา\footnote{พาฬฺหา (สี, สฺยา, อิ)} เวทนา วตฺตนฺติ มารณนฺติกา.\csromanspacer{} ตตฺร สุทํ ภควา สโต สมฺปชาโน อธิวาเสสิ อวิหญฺญมาโน.\csromanspacer{} อถ โข ภควา อายสฺมนฺตํ อานนฺทํ อามนฺเตสิ “อายามานนฺท เยน กุสินารา เตนุปสงฺกมิสฺสามา”ติ. “เอวํ ภนฺเต”ติ โข อายสฺมา อานนฺโท ภควโต ปจฺจสฺโสสิ.}

\csromangathameasure{จุนฺทสฺส ภตฺตํ ภุญฺชิตฺวา, กมฺมารสฺสาติ เม สุตํ. \\ อาพาธํ สมฺผุสี ธีโร, ปพาฬฺหํ มารณนฺติกํ. \\ ภุตฺตสฺส จ สูกรมทฺทเวน, พฺยาธิปฺปพาโฬฺห อุทปาทิ สตฺถุโน. \\ วิริจฺจมาโน ภควา อโวจ, คจฺฉามหํ กุสินารํ นครนฺติ.}
\csromangathagroup{\csromangathastanza{จุนฺทสฺส ภตฺตํ ภุญฺชิตฺวา, กมฺมารสฺสาติ เม สุตํ. \\ อาพาธํ สมฺผุสี ธีโร, ปพาฬฺหํ มารณนฺติกํ.}\csromangathastanza{ภุตฺตสฺส จ สูกรมทฺทเวน, พฺยาธิปฺปพาโฬฺห อุทปาทิ สตฺถุโน. \\ วิริจฺจมาโน\footnote{วิริญฺจมาโน (?) วิเรจมาโน (ที 2. 106)} ภควา อโวจ, คจฺฉามหํ กุสินารํ นครนฺติ.}}

\prose{อถ โข ภควา มคฺคา โอกฺกมฺม เยน อญฺญตรํ รุกฺขมูลํ เตนุปสงฺกมิ, อุปสงฺกมิตฺวา อายสฺมนฺตํ อานนฺทํ อามนฺเตสิ “อิงฺฆ เม ตฺวํ อานนฺท จตุคฺคุณํ สงฺฆาฏิํ ปญฺญาเปหิ, กิลนฺโตสฺมิ อานนฺท, นิสีทิสฺสามี”ติ. “เอวํ ภนฺเต”ติ โข อายสฺมา อานนฺโท ภควโต ปฏิสฺสุตฺวา จตุคฺคุณํ}

\vspace{\tipitakaparskip}
\csromancenter{ปาฏลิคามิยวคฺค สงฺฆาฏิํ ปญฺญาเปสิ.\csromanspacer{} นิสีทิ ภควา ปญฺญตฺเต อาสเน, นิสชฺช โข ภควา อายสฺมนฺตํ อานนฺทํ อามนฺเตสิ “อิงฺฆ เม ตฺวํ อานนฺท, ปานียํ อาหร, ปิปาสิโตสฺมิ อานนฺท, ปิวิสฺสามี”ติ.}
\prose{\csromanfolio{181}เอวํ วุตฺเต อายสฺมา อานนฺโท ภควนฺตํ เอตทโวจ “อิทานิ ภนฺเต ปญฺจมตฺตานิ สกฏสตานิ อติกฺกนฺตานิ, ตํ จกฺกจฺฉินฺนํ อุทกํ ปริตฺตํ ลุฬิตํ อาวิลํ สนฺทติ.\csromanspacer{} อยํ ภนฺเต กุกุธา\footnote{กกุตฺถา (สี), กุกุฏา (สฺยา), กกุฏฺฐา (ก)} นที อวิทูเร อจฺโฉทกา สาโตทกา สีโตทกา เสโตทกา สุปติตฺถา รมณียา, เอตฺถ ภควา ปานียญฺจ ปิวิสฺสติ คตฺตานิ จ สีตีกริสฺสตี”ติ\footnote{สีติํ กริสฺสตีติ (สี), สีตํ กริสฺสตีติ (สฺยา, อิ, ก)}.}

\prose{ทุติยมฺปิ โข ฯเปฯ.\csromanspacer{} ตติยมฺปิ โข ภควา อายสฺมนฺตํ อานนฺทํ อามนฺเตสิ “อิงฺฆ เม ตฺวํ อานนฺท, ปานียํ อาหร, ปิปาสิโตสฺมิ อานนฺท, ปิวิสฺสามี”ติ. “เอวํ ภนฺเต”ติ โข อายสฺมา อานนฺโท ภควโต ปฏิสฺสุตฺวา ปตฺตํ คเหตฺวา เยน สา นที เตนุปสงฺกมิ.\csromanspacer{} อถ โข สา นที จกฺกจฺฉินฺนา ปริตฺตา ลุฬิตา อาวิลา สนฺทมานา อายสฺมนฺเต อานนฺเท อุปสงฺกมนฺเต อจฺฉา วิปฺปสนฺนา อนาวิลา สนฺทติ.}

\prose{อถ โข อายสฺมโต อานนฺทสฺส เอตทโหสิ “อจฺฉริยํ วต โภ, อพฺภุตํ วต โภ, ตถาคตสฺส มหิทฺธิกตา มหานุภาวตา, อยํ หิ สา นที จกฺกจฺฉินฺนา ปริตฺตา ลุฬิตา อาวิลา สนฺทมานา มยิ อุปสงฺกมนฺเต อจฺฉา วิปฺปสนฺนา อนาวิลา สนฺทตี”ติ, ปตฺเตน ปานียํ อาทาย เยน ภควา เตนุปสงฺกมิ, อุปสงฺกมิตฺวา ภควนฺตํ เอตทโวจ “อจฺฉริยํ ภนฺเต, อพฺภุตํ ภนฺเต, ตถาคตสฺส มหิทฺธิกตา มหานุภาวตา, อยํ หิ สา ภนฺเต นที จกฺกจฺฉินฺนา ปริตฺตา ลุฬิตา อาวิลา สนฺทมานา มยิ อุปสงฺกมนฺเต อจฺฉา วิปฺปสนฺนา อนาวิลา สนฺทติ, ปิวตุ ภควา ปานียํ, ปิวตุ สุคโต ปานียนฺ”ติ.}

\prose{อถ โข ภควา ปานียํ อปายิ\footnote{อปาสิ (สี)}.\csromanspacer{} อถ โข ภควา มหตา ภิกฺขุสํเฆน สทฺธิํ เยน กุกุธา นที เตนุปสงฺกมิ, อุปสงฺกมิตฺวา กุกุธํ นทิํ อชฺโฌคาเหตฺวา นฺหตฺวา จ ปิวิตฺวา จ ปจฺจุตฺตริตฺวา เยน อมฺพวนํ เตนุปสงฺกมิ, อุปสงฺกมิตฺวา อายสฺมนฺตํ จุนฺทกํ อามนฺเตสิ “อิงฺฆ เม ตฺวํ จุนฺทก จตุคฺคุณํ สงฺฆาฏิํ ปญฺญาเปหิ, กิลนฺโตสฺมิ จุนฺทก,}

\csromanmatikaanchor{mtk.100}
\csromantocmarkref{subsection}{6. ตติยนานาติตฺถิยสุตฺต}{mtk.100}
\bookmark[dest={mtk.100},level=2]{6. ตติยนานาติตฺถิยสุตฺต}
\csromantocmarkref{subsection}{7. สุภูติสุตฺต}{mtk.101}
\bookmark[dest={mtk.101},level=2]{7. สุภูติสุตฺต}
\vspace{\tipitakaparskip}
\csromancenter{อุทานปาฬิ นิปชฺชิสฺสามี”ติ. “เอวํ ภนฺเต”ติ โข อายสฺมา จุนฺทโก ภควโต ปฏิสฺสุตฺวา จตุคฺคุณํ สงฺฆาฏิํ ปญฺญาเปสิ.\csromanspacer{} อถ โข ภควา ทกฺขิเณน ปสฺเสน สีหเสยฺยํ กปฺเปสิ ปาเท ปาทํ อจฺจาธาย สโต สมฺปชาโน อุฏฺฐานสญฺญํ มนสิ กริตฺวา, อายสฺมา ปน จุนฺทโก ตตฺเถว ภควโต ปุรโต นิสีทิ.}
\csromancenter{คนฺตฺวาน พุทฺโธ นทิกํ กุกุธํ, (ที 2. 112)}
\csromancenter{อจฺโฉทกํ สาตุทกํ\footnote{สาโตทกํ (สพฺพตฺถ)} วิปฺปสนฺนํ.}
\vspace{0.5\baselineskip}
\csromangathameasure{โอคาหิ สตฺถา สุกิลนฺตรูโป, \\ ตถาคโต อปฺปฏิโม’ธ โลเก. \\ นฺหตฺวา จ ปิวิตฺวา จุทตาริ สตฺถา, \\ ปุรกฺขโต ภิกฺขุคณสฺส มชฺเฌ. \\ สตฺถา ปวตฺตา ภควา อิธ ธมฺเม, \\ อุปาคมิ อมฺพวนํ มเหสิ. \\ อามนฺตยิ จุนฺทกํ นาม ภิกฺขุํ, \\ จตุคฺคุณํ สนฺถร เม นิปชฺชํ. \\ โส โจทิโต ภาวิตตฺเตน จุนฺโท, \\ จตุคฺคุณํ สนฺถริ ขิปฺปเมว. \\ นิปชฺชิ สตฺถา สุกิลนฺตรูโป, \\ จุนฺโทปิ ตตฺถ ปมุเข นิสีทีติ.}
\csromangathagroup{\csromangathastanza{\csromanfolio{182}โอคาหิ สตฺถา สุกิลนฺตรูโป, \\ ตถาคโต อปฺปฏิโม’ธ โลเก. \\ นฺหตฺวา จ ปิวิตฺวา จุทตาริ\footnote{นฺหตฺวา จ อุตฺตริ (ก)} สตฺถา, \\ ปุรกฺขโต ภิกฺขุคณสฺส มชฺเฌ.}\csromangathastanza{สตฺถา ปวตฺตา ภควา อิธ ธมฺเม, \\ อุปาคมิ อมฺพวนํ มเหสิ. \\ อามนฺตยิ จุนฺทกํ นาม ภิกฺขุํ, \\ จตุคฺคุณํ สนฺถร\footnote{ปตฺถร (สี, อิ)} เม นิปชฺชํ.}\csromangathastanza{โส โจทิโต ภาวิตตฺเตน จุนฺโท, \\ จตุคฺคุณํ สนฺถริ\footnote{ปตฺถริ (สี, อิ)} ขิปฺปเมว. \\ นิปชฺชิ สตฺถา สุกิลนฺตรูโป, \\ จุนฺโทปิ ตตฺถ ปมุเข นิสีทีติ.}}

\prose{อถ โข ภควา อายสฺมนฺตํ อานนฺทํ อามนฺเตสิ “สิยา โข ปนานนฺท จุนฺทสฺส กมฺมารปุตฺตสฺส โกจิ วิปฺปฏิสารํ อุปทเหยฺย ‘ตสฺส เต อาวุโส จุนฺท อลาภา, ตสฺส เต ทุลฺลทฺธํ, ยสฺส เต ตถาคโต ปจฺฉิมํ ปิณฺฑปาตํ ภุญฺชิตฺวา ปรินิพฺพุโต’ติ”.\csromanspacer{} จุนฺทสฺสานนฺท กมฺมารปุตฺตสฺส เอวํ วิปฺปฏิสาโร ปฏิวิโนเทตพฺโพ\csromandash{}}

\prose{“ตสฺส เต อาวุโส จุนฺท ลาภา, ตสฺส เต สุลทฺธํ, ยสฺส เต ตถาคโต ปจฺฉิมํ ปิณฺฑปาตํ ปริภุญฺชิตฺวา ปรินิพฺพุโต.\csromanspacer{} สมฺมุขา เมตํ อาวุโส จุนฺท ภควโต สุตํ, สมฺมุขา ปฏิคฺคหิตํ\csromandash{}‘เทฺวเม ปิณฺฑปาตา}

\vspace{\tipitakaparskip}
\csromancenter{ปาฏลิคามิยวคฺค สมสมผลา สมสมวิปากา อติวิย อญฺเญหิ ปิณฺฑปาเตหิ มหปฺผลตรา จ มหานิสํสตรา จ.\csromanspacer{} กตเม เทฺว? ยญฺจ ปิณฺฑปาตํ ปริภุญฺชิตฺวา ตถาคโต อนุตฺตรํ สมฺมาสมฺโมธิํ อภิสมฺพุชฺฌติ, ยญฺจ ปิณฺฑปาตํ ปริภุญฺชิตฺวา อนุปาทิเสสาย นิพฺพานธาตุยา ปรินิพฺพายติ, อิเม เทฺว ปิณฺฑปาตา สมสมผลา สมสมวิปากา อติวิย อญฺเญหิ ปิณฺฑปาเตหิ มหปฺผลตรา จ มหานิสํสตรา จ’.}
\prose{\csromanfolio{183}อายุสํวตฺตนิกํ อายสฺมตา จุนฺเทน กมฺมารปุตฺเตน กมฺมํ อุปจิตํ, วณฺณสํวตฺตนิกํ อายสฺมตา จุนฺเทน กมฺมารปุตฺเตน กมฺมํ อุปจิตํ, สุขสํวตฺตนิกํ อายสฺมตา จุนฺเทน กมฺมารปุตฺเตน กมฺมํ อุปจิตํ, สคฺคสํวตฺตนิกํ อายสฺมตา จุนฺเทน กมฺมารปุตฺเตน กมฺมํ อุปจิตํ, ยสสํวตฺตนิกํ อายสฺมตา จุนฺเทน กมฺมารปุตฺเตน กมฺมํ อุปจิตํ, อาธิปเตยฺยสํวตฺตนิกํ อายสฺมตา จุนฺเทน กมฺมารปุตฺเตน กมฺมํ อุปจิตนฺ”ติ.\csromanspacer{} จุนฺทสฺสานนฺท กมฺมารปุตฺตสฺส เอวํ วิปฺปฏิสาโร ปฏิวิโนเทตพฺโพติ.}

\prose{อถ โข ภควา เอตมตฺถํ วิทิตฺวา ตายํ เวลายํ อิมํ อุทานํ อุทาเนสิ\csromandash{}}

\csromangathameasure{+“ททโต ปุญฺญํ ปวฑฺฒติ, \\ สํยมโต เวรํ น จียติ. \\ กุสโล จ ชหาติ ปาปกํ, \\ ราคโทสโมหกฺขยา สนิพฺพุโต”ติ.ปญฺจมํ.}
\csromangathagroup{\csromangathastanza{\csromansymbolfootnote{+}{ที 2. 113; ขุ 10. 187 ปิฏฺเฐสุปิ.}“ททโต ปุญฺญํ ปวฑฺฒติ, \\ สํยมโต เวรํ น จียติ. \\ กุสโล จ ชหาติ ปาปกํ, \\ ราคโทสโมหกฺขยา สนิพฺพุโต”ติ\footnote{ปรินิพฺพุโตติ (สี, สฺยา, อิ)}.\csromanspacer{}\csromanspacer{}ปญฺจมํ.}}

\csromanheader{2}{6. ปาฏลิคามิยสุตฺต}
\proseitem{76}{\csromansymbolfootnote{*}{วิ 3. 321; ม 2. 16; สํ 2. 389 ปิฏฺเฐสุปิ.}เอวํ เม สุตํ\csromandash{}เอกํ สมยํ ภควา มคเธสุ จาริกํ จรมาโน มหตา ภิกฺขุสํเฆน สทฺธิํ เยน ปาฏลิคาโม ตทวสริ.\csromanspacer{} อสฺโสสุํ โข ปาฏลิคามิยา\footnote{ปาฏลิคามิกา (ที 2. 71)} อุปาสกา “ภควา กิร มคเธสุ จาริกํ จรมาโน มหตา ภิกฺขุสํเฆน สทฺธิํ ปาฏลิคามํ อนุปฺปตฺโต”ติ.\csromanspacer{} อถ โข ปาฏลิคามิยา อุปาสกา เยน ภควา เตนุปสงฺกมิํสุ, อุปสงฺกมิตฺวา ภควนฺตํ อภิวาเทตฺวา เอกมนฺตํ นิสีทิํสุ.\csromanspacer{} เอกมนฺตํ นิสินฺนา โข ปาฏลิคามิยา อุปาสกา ภควนฺตํ เอตทโวจุํ “อธิวาเสตุ โน ภนฺเต ภควา อาวสถาคารนฺ”ติ.\csromanspacer{} อธิวาเสสิ ภควา ตุณฺหีภาเวน.}

\gambhira{อุทานปาฬิ}
\prose{\csromanfolio{184}อถ โข ปาฏลิคามิยา อุปาสกา ภควโต อธิวาสนํ วิทิตฺวา อุฏฺฐายาสนา ภควนฺตํ อภิวาเทตฺวา ปทกฺขิณํ กตฺวา เยนาวสถาคารํ เตนุปสงฺกมิํสุ, อุปสงฺกมิตฺวา สพฺพสนฺถริํ อาวสถาคารํ สนฺถริตฺวา อาสนานิ ปญฺญาเปตฺวา อุทกมณิกํ ปติฏฺฐาเปตฺวา เตลปฺปทีปํ อาโรเปตฺวา เยน ภควา เตนุปสงฺกมิํสุ, อุปสงฺกมิตฺวา ภควนฺตํ อภิวาเทตฺวา เอกมนฺตํ อฏฺฐํสุ, เอกมนฺตํ ฐิตา โข ปาฏลิคามิยา อุปาสกา ภควนฺตํ เอตทโวจุํ “สพฺพสนฺถริสนฺถตํ\footnote{สพฺพสนฺถริํ สนฺถตํ (สี, สฺยา, อิ)} ภนฺเต อาวสถาคารํ, อาสนานิ ปญฺญตฺตานิ, อุทกมณิโก ปติฏฺฐาปิโต\footnote{อุทกมณิกํ ปติฏฺฐาปิตํ (สฺยา)}, เตลปฺปทีโป อาโรปิโต.\csromanspacer{} ยสฺส ทานิ ภนฺเต ภควา กาลํ มญฺญตี”ติ.}

\prose{อถ โข ภควา นิวาเสตฺวา ปตฺตจีวรมาทาย สทฺธิํ ภิกฺขุสํเฆน เยน อาวสถาคารํ เตนุปสงฺกมิ, อุปสงฺกมิตฺวา ปาเท ปกฺขาเลตฺวา อาวสถาคารํ ปวิสิตฺวา มชฺฌิมํ ถมฺภํ นิสฺสาย ปุรตฺถาภิมุโข นิสีทิ, ภิกฺขุสํโฆปิ โข ปาเท ปกฺขาเลตฺวา อาวสถาคารํ ปวิสิตฺวา ปจฺฉิมํ ภิตฺติํ นิสฺสาย ปุรตฺถาภิมุโข นิสีทิ ภควนฺตํเยว ปุรกฺขตฺวา, ปาฏลิคามิยาปิ โข อุปาสกา ปาเท ปกฺขาเลตฺวา อาวสถาคารํ ปวิสิตฺวา ปุรตฺถิมํ ภิตฺติํ นิสฺสาย ปจฺฉิมาภิมุขา นิสีทิํสุ ภควนฺตํเยว ปุรกฺขตฺวา.\csromanspacer{} อถ โข ภควา ปาฏลิคามิเย อุปาสเก อามนฺเตสิ\csromandash{}}

\prose{\csromansymbolfootnote{*}{อํ 2. 221; ที 3.196 ปิฏฺเฐปิ.}ปญฺจิเม คหปตโย อาทีนวา ทุสฺสีลสฺส สีลวิปตฺติยา.\csromanspacer{} กตเม ปญฺจ? อิธ คหปตโย ทุสฺสีโล สีลวิปนฺโน ปมาทาธิกรณํ มหติํ โภคชานิํ นิคจฺฉติ, อยํ ปฐโม อาทีนโว ทุสฺสีลสฺส สีลวิปตฺติยา.}

\prose{ปุน จปรํ คหปตโย ทุสฺสีลสฺส สีลวิปนฺนสฺส ปาปโก กิตฺติสทฺโท อพฺภุคฺคจฺฉติ, อยํ ทุติโย อาทีนโว ทุสฺสีลสฺส สีลวิปตฺติยา.}

\csromanmatikaanchor{mtk.102}
\csromantocmarkref{subsection}{8. คณิกาสุตฺต}{mtk.102}
\bookmark[dest={mtk.102},level=2]{8. คณิกาสุตฺต}
\prose{ปุน จปรํ คหปตโย ทุสฺสีโล สีลวิปนฺโน ยญฺญเทว ปริสํ อุปสงฺกมติ ยทิ ขตฺติยปริสํ ยทิ พฺราหฺมณปริสํ ยทิ คหปติปริสํ ยทิ สมณปริสํ, อวิสารโท อุปสงฺกมติ มงฺกุภูโต, อยํ ตติโย อาทีนโว ทุสฺสีลสฺส สีลวิปตฺติยา.}

\prose{ปุน จปรํ คหปตโย ทุสฺสีโล สีลวิปนฺโน สมฺมูโฬฺห กาลํ กโรติ, อยํ จตุตฺโถ อาทีนโว ทุสฺสีลสฺส สีลวิปตฺติยา.}

\chapterheadpageread{8. ปาฏลิคามิยวคฺค}
\prose{\csromanfolio{185}ปุน จปรํ คหปตโย ทุสฺสีโล สีลวิปนฺโน กายสฺส เภทา ปรํ มรณา อปายํ ทุคฺคติํ วินิปาตํ นิรยํ อุปปชฺชติ, อยํ ปญฺจโม อาทีนโว ทุสฺสีลสฺส สีลวิปตฺติยา.\csromanspacer{} อิเม โข คหปตโย ปญฺจ อาทีนวา ทุสฺสีลสฺส สีลวิปตฺติยา.}

\prose{ปญฺจิเม คหปตโย อานิสํสา สีลวโต สีลสมฺปทาย.\csromanspacer{} กตเม ปญฺจ? อิธ คหปตโย สีลวา สีลสมฺปนฺโน อปฺปมาทาธิกรณํ มหนฺตํ โภคกฺขนฺธํ อธิคจฺฉติ, อยํ ปฐโม อานิสํโส สีลวโต สีลสมฺปทาย.}

\prose{ปุน จปรํ คหปตโย สีลวโต สีลสมฺปนฺนสฺส กลฺยาโณ กิตฺติสทฺโท อพฺภุคฺคจฺฉติ, อยํ ทุติโย อานิสํโส สีลวโต สีลสมฺปทาย.}

\csromanmatikaanchor{mtk.103}
\csromantocmarkref{subsection}{9. อุปาติธาวนฺติสุตฺต}{mtk.103}
\bookmark[dest={mtk.103},level=2]{9. อุปาติธาวนฺติสุตฺต}
\prose{ปุน จปรํ คหปตโย สีลวา สีลสมฺปนฺโน ยญฺญเทว ปริสํ อุปสงฺกมติ ยทิ ขตฺติยปริสํ ยทิ พฺราหฺมณปริสํ ยทิ คหปติปริสํ ยทิ สมณปริสํ วิสารโท อุปสงฺกมติ อมงฺกุภูโต, อยํ ตติโย อานิสํโส สีลวโต สีลสมฺปทาย.}

\prose{ปุน จปรํ คหปตโย สีลวา สีลสมฺปนฺโน อสมฺมูโฬฺห กาลํ กโรติ, อยํ จตุตฺโถ อานิสํโส สีลวโต สีลสมฺปทาย.}

\prose{ปุน จปรํ คหปตโย สีลวา สีลสมฺปนฺโน กายสฺส เภทา ปรํ มรณา สุคติํ สคฺคํ โลกํ อุปปชฺชติ, อยํ ปญฺจโม อานิสํโส สีลวโต สีลสมฺปทาย.\csromanspacer{} อิเม โข คหปตโย ปญฺจ อานิสํสา สีลวโต สีลสมฺปทายาติ.}

\prose{อถ โข ภควา ปาฏลิคามิเย อุปาสเก พหุเทว รตฺติํ ธมฺมิยา กถาย สนฺทสฺเสตฺวา สมาทเปตฺวา สมุตฺเตเชตฺวา สมฺปหํเสตฺวา อุยฺโยเชสิ “อภิกฺกนฺตา โข คหปตโย รตฺติ, ยสฺสทานิ ตุเมฺห กาลํ มญฺญถา”ติ. 1อถ โข ปาฏลิคามิยา อุปาสกา ภควโต ภาสิตํ อภินนฺทิตฺวา อนุโมทิตฺวา1 อุฏฺฐายาสนา ภควนฺตํ อภิวาเทตฺวา ปทกฺขิณํ กตฺวา ปกฺกมิํสุ.\csromanspacer{} อถ โข ภควา อจิรปกฺกนฺเตสุ ปาฏลิคามิเยสุ อุปาสเกสุ สุญฺญาคารํ ปาวิสิ.}

\gambhira{อุทานปาฬิ}
\prose{\csromanfolio{186}เตน โข ปน สมเยน สุนิธวสฺสการา\footnote{สุนีธวสฺสการา (สี, สฺยา, อิ)} มคธมหามตฺตา ปาฏลิคาเม นครํ มาเปนฺติ วชฺชีนํ ปฏิพาหาย.\csromanspacer{} เตน โข ปน สมเยน สมฺพหุลา เทวตาโย สหสฺสสหสฺเสว\footnote{สหสฺเสว (สฺยา, ก), สหสฺสสฺเสว (อิ)} ปาฏลิคาเม วตฺถูนิ ปริคฺคณฺหนฺติ.\csromanspacer{} ยสฺมิํ ปเทเส มเหสกฺขา เทวตา วตฺถูนิ ปริคฺคณฺหนฺติ, มเหสกฺขานํ ตตฺถ รญฺญํ ราชมหามตฺตานํ จิตฺตานิ นมนฺติ นิเวสนานิ มาเปตุํ.\csromanspacer{} ยสฺมิํ ปเทเส มชฺฌิมา เทวตา วตฺถูนิ ปริคฺคณฺหนฺติ, มชฺฌิมานํ ตตฺถ รญฺญํ ราชมหามตฺตานํ จิตฺตานิ นมนฺติ นิเวสนานิ มาเปตุํ.\csromanspacer{} ยสฺมิํ ปเทเส นีจา เทวตา วตฺถูนิ ปริคฺคณฺหนฺติ, นีจานํ ตตฺถ รญฺญํ ราชมหามตฺตานํ จิตฺตานิ นมนฺติ นิเวสนานิ มาเปตุํ.}

\csromanmatikaanchor{mtk.104}
\csromantocmarkref{subsection}{10. อุปฺปชฺชนฺติสุตฺต}{mtk.104}
\bookmark[dest={mtk.104},level=2]{10. อุปฺปชฺชนฺติสุตฺต}
\prose{อทฺทสา โข ภควา ทิพฺเพน จกฺขุนา วิสุทฺเธน อติกฺกนฺตมานุสเกน ตา เทวตาโย สหสฺสสหสฺเสว ปาฏลิคาเม วตฺถูนิ ปริคฺคณฺหนฺติโย, ยสฺมิํ ปเทเส มเหสกฺขา เทวตา วตฺถูนิ ปริคฺคณฺหนฺติ, มเหสกฺขานํ ตตฺถ รญฺญํ ราชมหามตฺตานํ จิตฺตานิ นมนฺติ นิเวสนานิ มาเปตุํ.\csromanspacer{} ยสฺมิํ ปเทเส มชฺฌิมา เทวตา วตฺถูนิ ปริคฺคณฺหนฺติ, มชฺฌิมานํ ตตฺถ รญฺญํ ราชมหามตฺตานํ จิตฺตานิ นมนฺติ นิเวสนานิ มาเปตุํ.\csromanspacer{} ยสฺมิํ ปเทเส นีจา เทวตา วตฺถูนิ ปริคฺคณฺหนฺติ, นีจานํ ตตฺถ รญฺญํ ราชมหามตฺตานํ จิตฺตานิ นมนฺติ นิเวสนานิ มาเปตุํ.\csromanspacer{} อถ โข ภควา ตสฺสา รตฺติยา ปจฺจูสสมเย ปจฺจุฏฺฐาย อายสฺมนฺตํ อานนฺทํ อามนฺเตสิ\csromandash{}}

\prose{เก นุ โข\footnote{โก นุ โข (สพฺพตฺถ)} อานนฺท ปาฏลิคาเม นครํ มาเปนฺตีติ\footnote{มาเปตีติ (สพฺพตฺถ)}.\csromanspacer{} สุนิธวสฺสการา ภนฺเต มคธมหามตฺตา ปาฏลิคาเม นครํ มาเปนฺติ วชฺชีนํ ปฏิพาหายาติ.\csromanspacer{} เสยฺยถาปิ อานนฺท เทเวหิ ตาวติํเสหิ สทฺธิํ มนฺเตตฺวา, เอวเมวํ โข อานนฺท สุนิธวสฺสการา มคธมหามตฺตา ปาฏลิคาเม นครํ มาเปนฺติ วชฺชีนํ ปฏิพาหายา.\csromanspacer{} อิธาหํ อานนฺท อทฺทสํ ทิพฺเพน จกฺขุนา วิสุทฺเธน อติกฺกนฺตมานุสเกน สมฺพหุลา เทวตาโย สหสฺสสหสฺเสว ปาฏลิคาเม วตฺถูนิ ปริคฺคณฺหนฺติโย.\csromanspacer{} ยสฺมิํ ปเทเส มเหสกฺขา เทวตา วตฺถูนิ ปริคฺคณฺหนฺติ, มเหสกฺขานํ ตตฺถ รญฺญํ ราชมหามตฺตานํ จิตฺตานิ นมนฺติ นิเวสนานิ มาเปตุํ.\csromanspacer{} ยสฺมิํ ปเทเส มชฺฌิมา เทวตา วตฺถูนิ ปริคฺคณฺหนฺติ, มชฺฌิมานํ ตตฺถ รญฺญํ ราชมหามตฺตานํ จิตฺตานิ นมนฺติ นิเวสนานิ มาเปตุํ.}

\vspace{\tipitakaparskip}
\csromancenter{ปาฏลิคามิยวคฺค ยสฺมิํ ปเทเส นีจา เทวตา วตฺถูนิ ปริคฺคณฺหนฺติ, นีจานํ ตตฺถ รญฺญํ ราชมหามตฺตานํ จิตฺตานิ นมนฺติ นิเวสนานิ มาเปตุํ.\csromanspacer{} ยาวตา อานนฺท อริยํ อายตนํ ยาวตา วณิปฺปโถ, อิทํ อคฺคนครํ ภวิสฺสติ ปาฏลิปุตฺตํ ปุฏเภทนํ.\csromanspacer{} ปาฏลิปุตฺตสฺส โข อานนฺท ตโย อนฺตรายา ภวิสฺสนฺติ อคฺคิโต วา อุทกโต วา มิถุเภทโต วาติ.}
\prose{\csromanfolio{187}อถ โข สุนิธวสฺสการา มคธมหามตฺตา เยน ภควา เตนุปสงฺกมิํสุ, อุปสงฺกมิตฺวา ภควตา สทฺธิํ สมฺโมทิํสุ, สมฺโมทนียํ กถํ สารณียํ\footnote{สาราณียํ (สี, สฺยา, กํ, อิ)} วีติสาเรตฺวา เอกมนฺตํ อฏฺฐํสุ, เอกมนฺตํ ฐิตา โข สุนิธวสฺสการา มคธมหามตฺตา ภควนฺตํ เอตทโวจุํ “อธิวาเสตุ โน ภวํ โคตโม อชฺชตนาย ภตฺตํ สทฺธิํ ภิกฺขุสํเฆนา”ติ.\csromanspacer{} อธิวาเสสิ ภควา ตุณฺหีภาเวน.}

\prose{อถ โข สุนิธวสฺสการา มคธมหามตฺตา ภควโต อธิวาสนํ วิทิตฺวา เยน สโก อาวสโถ เตนุปสงฺกมิํสุ, อุปสงฺกมิตฺวา สเก อาวสเถ ปณีตํ ขาทนียํ โภชนียํ ปฏิยาทาเปตฺวา ภควโต กาลํ อาโรเจสุํ “กาโล โภ โคตม, นิฏฺฐิตํ ภตฺตนฺ”ติ.}

\prose{อถ โข ภควา ปุพฺพณฺหสมยํ นิวาเสตฺวา ปตฺตจีวรมาทาย สทฺธิํ ภิกฺขุสํเฆน เยน สุนิธวสฺสการานํ มคธมหามตฺตานํ อาวสโถ เตนุปสงฺกมิ, อุปสงฺกมิตฺวา ปญฺญตฺเต อาสเน นิสีทิ.\csromanspacer{} อถ โข สุนิธวสฺสการา มคธมหามตฺตา พุทฺธปฺปมุขํ ภิกฺขุสํฆํ ปณีเตน ขาทนีเยน โภชนีเยน สหตฺถา สนฺตปฺเปสุํ สมฺปวาเรสุํ.}

\prose{อถ โข สุนิธวสฺสการา มคธมหามตฺตา ภควนฺตํ ภุตฺตาวิํ โอนีตปตฺตปาณิํ อญฺญตรํ นีจํ อาสนํ คเหตฺวา เอกมนฺตํ นิสีทิํสุ, เอกมนฺตํ นิสินฺเน โข สุนิธวสฺสกาเร มคธมหามตฺเต ภควา อิมาหิ คาถาหิ อนุโมทิ\csromandash{}}

\csromangathameasure{“ยสฺมิํ ปเทเส กปฺเปติ, วาสํ ปณฺฑิตชาติโย. \\ สีลวนฺเตตฺถ โภเชตฺวา, สญฺญเต พฺรหฺมจารโย.}
\csromangathagroup{\csromangathastanza{“ยสฺมิํ ปเทเส กปฺเปติ, วาสํ ปณฺฑิตชาติโย. \\ สีลวนฺเตตฺถ โภเชตฺวา, สญฺญเต พฺรหฺมจารโย\footnote{พฺรหฺมจาริโน (สฺยา), พฺรหฺมจริเย (อิ, ก)}.}}

\gambhira{อุทานปาฬิ}
\csromangathameasure{ยา ตตฺถ เทวตา อาสุํ, ตาสํ ทกฺขิณมาทิเส. \\ ตา ปูชิตา ปูชยนฺติ, มานิตา มานยนฺติ นํ. \\ ตโต นํ อนุกมฺปนฺติ, มาตา ปุตฺตํว โอรสํ. \\ เทวตานุกมฺปิโต โปโส, สทา ภทฺรานิ ปสฺสตี”ติ.}
\csromangathagroup{\csromangathastanza{\csromanfolio{188}ยา ตตฺถ เทวตา อาสุํ, ตาสํ ทกฺขิณมาทิเส. \\ ตา ปูชิตา ปูชยนฺติ, มานิตา มานยนฺติ นํ.}\csromangathastanza{ตโต นํ อนุกมฺปนฺติ, มาตา ปุตฺตํว โอรสํ. \\ เทวตานุกมฺปิโต โปโส, สทา ภทฺรานิ ปสฺสตี”ติ.}}

\prose{อถ โข ภควา สุนิธวสฺสการานํ มคธมหามตฺตานํ อิมาหิ คาถาหิ อนุโมทิตฺวา อุฏฺฐายาสนา ปกฺกามิ.}

\csromanmatikaanchor{mtk.105}
\csromantocmarknopage{chapter}{7. จูฬวคฺค}{mtk.105}
\bookmark[dest={mtk.105},level=0]{7. จูฬวคฺค}
\prose{เตน โข ปน สมเยน สุนิธวสฺสการา มคธมหามตฺตา ภควนฺตํ ปิฏฺฐิโต ปิฏฺฐิโต อนุพนฺธา โหนฺติ “เยนชฺช สมโณ โคตโม ทฺวาเรน นิกฺขมิสฺสติ, ตํ โคตมทฺวารํ นาม ภวิสฺสติ, เยน ติตฺเถน คงฺคํ นทิํ ตริสฺสติ, ตํ โคตมติตฺถํ นาม ภวิสฺสตี”ติ.}

\csromanmatikaanchor{mtk.106}
\csromantocmarkref{subsection}{1. ปฐมลกุณฺฑกภทฺทิยสุตฺต}{mtk.106}
\bookmark[dest={mtk.106},level=2]{1. ปฐมลกุณฺฑกภทฺทิยสุตฺต}
\prose{อถ โข ภควา เยน ทฺวาเรน นิกฺขมิ, ตํ โคตมทฺวารํ นาม อโหสิ.\csromanspacer{} อถ โข ภควา เยน คงฺคา นที เตนุปสงฺกมิ.\csromanspacer{} เตน โข ปน สมเยน คงฺคา นที ปูรา โหติ สมติตฺติกา กากเปยฺยา, อปฺเปกจฺเจ มนุสฺสา นาวํ ปริเยสนฺติ, อปฺเปกจฺเจ อุฬุมฺปํ ปริเยสนฺติ, อปฺเปกจฺเจ กุลฺลํ พนฺธนฺติ อปารา ปารํ คนฺตุกามา.\csromanspacer{} อถ โข ภควา เสยฺยถาปิ นาม พลวา ปุริโส สมิญฺชิตํ วา พาหํ ปสาเรยฺย, ปสาริตํ วา พาหํ สมิญฺเชยฺย, เอวเมวํ คงฺคาย นทิยา โอริมตีเร\footnote{โอริมตีรา (พหูสุ) วิ 3. 325; ที 2. 76 ปิฏฺเฐสุ ปสฺสิตพฺพํ.} อนฺตรหิโต ปาริมตีเร ปจฺจุฏฺฐาสิ สทฺธิํ ภิกฺขุสํเฆน.}

\prose{อทฺทสา โข ภควา เต มนุสฺเส อปฺเปกจฺเจ นาวํ ปริเยสนฺเต อปฺเปกจฺเจ อุฬุมฺปํ ปริเยสนฺเต อปฺเปกจฺเจ กุลฺลํ พนฺธนฺเต อปารา ปารํ คนฺตุกาเม.}

\prose{อถ โข ภควา เอตมตฺถํ วิทิตฺวา ตายํ เวลายํ อิมํ อุทานํ อุทาเนสิ\csromandash{}}

\csromangathameasure{“เย ตรนฺติ อณฺณวํ สรํ, \\ เสตุํ กตฺวาน วิสชฺช ปลฺลลานิ. \\ กุลฺลํ หิ ชโน ปพนฺธติ, \\ ติณฺณา เมธาวิโน ชนา”ติ.ฉฏฺฐํ.}
\csromangathagroup{\csromangathastanza{“เย ตรนฺติ อณฺณวํ สรํ, \\ เสตุํ กตฺวาน วิสชฺช ปลฺลลานิ. \\ กุลฺลํ หิ ชโน ปพนฺธติ\footnote{พนฺธติ (สฺยา, อิ)}, \\ ติณฺณา\footnote{พนฺธติ (สฺยา, อิ)} เมธาวิโน ชนา”ติ.\csromanspacer{}\csromanspacer{}ฉฏฺฐํ.}}

\csromanheader{2}{8. ปาฏลิคามิยวคฺค \textbf{7. ทฺวิธาปถสุตฺต}}
\proseitem{77}{\csromanfolio{189}เอวํ เม สุตํ\csromandash{}เอกํ สมยํ ภควา โกสเลสุ อทฺธานมคฺคปฏิปนฺโน โหติ อายสฺมตา นาคสมาเลน ปจฺฉาสมเณน.\csromanspacer{} อทฺทสา โข อายสฺมา นาคสมาโล อนฺตรามคฺเค ทฺวิธาปถํ\footnote{เทฺวธาปถํ (สี)}, ทิสฺวาน ภควนฺตํ เอตทโวจ “อยํ ภนฺเต ภควา ปนฺโถ, อิมินา คจฺฉามา”ติ.\csromanspacer{} เอวํ วุตฺเต ภควา อายสฺมนฺตํ นาคสมาลํ เอตทโวจ “อยํ นาคสมาล ปนฺโถ, อิมินา คจฺฉามา”ติ.}

\csromanmatikaanchor{mtk.107}
\csromantocmarkref{subsection}{2. ทุติยลกุณฺฑกภทฺทิยสุตฺต}{mtk.107}
\bookmark[dest={mtk.107},level=2]{2. ทุติยลกุณฺฑกภทฺทิยสุตฺต}
\prose{ทุติยมฺปิ โข ฯเปฯ.\csromanspacer{} ตติยมฺปิ โข อายสฺมา นาคสมาโล ภควนฺตํ เอตทโวจ “อยํ ภนฺเต ภควา ปนฺโถ, อิมินา คจฺฉามา”ติ.\csromanspacer{} ตติยมฺปิ โข ภควา อายสฺมนฺตํ นาคสมาลํ เอตทโวจ “อยํ นาคสมาล ปนฺโถ, อิมินา คจฺฉามา”ติ.\csromanspacer{} อถ โข อายสฺมา นาคสมาโล ภควโต ปตฺตจีวรํ ตตฺเถว ฉมายํ นิกฺขิปิตฺวา ปกฺกามิ “อิทํ ภนฺเต ภควโต ปตฺตจีวรนฺ”ติ.}

\prose{อถ โข อายสฺมโต นาคสมาลสฺส เตน ปนฺเถน คจฺฉนฺตสฺส อนฺตรามคฺเค โจรา นิกฺขมิตฺวา หตฺเถหิ จ ปาเทหิ จ อาโกเฏสุํ, ปตฺตญฺจ ภินฺทิํสุ, สงฺฆาฏิญฺจ วิปฺผาเลสุํ.\csromanspacer{} อถ โข อายสฺมา นาคสมาโล ภินฺเนน ปตฺเตน วิปฺผาลิตาย สงฺฆาฏิยา เยน ภควา เตนุปสงฺกมิ, อุปสงฺกมิตฺวา ภควนฺตํ อภิวาเทตฺวา เอกมนฺตํ นิสีทิ, เอกมนฺตํ นิสินฺโน โข อายสฺมา นาคสมาโล ภควนฺตํ เอตทโวจ “อิธ มยฺหํ ภนฺเต เตน ปนฺเถน คจฺฉนฺตสฺส อนฺตรามคฺเค โจรา นิกฺขมิตฺวา หตฺเถหิ จ ปาเทหิ จ อาโกเฏสุํ, ปตฺตญฺจ ภินฺทิํสุ, สงฺฆาฏิญฺจ วิปฺผาเลสุนฺ”ติ.}

\prose{อถ โข ภควา เอตมตฺถํ วิทิตฺวา ตายํ เวลายํ อิมํ อุทานํ อุทาเนสิ\csromandash{}}

\csromangathameasure{“สทฺธิํ จรเมกโต วสํ, \\ มิสฺโส อญฺญชเนน เวทคู. \\ วิทฺวา ปชหาติ ปาปกํ, \\ โกญฺโจ ขีรปโกว นินฺนคนฺ”ติ.สตฺตมํ.}
\csromangathagroup{\csromangathastanza{“สทฺธิํ จรเมกโต วสํ, \\ มิสฺโส อญฺญชเนน เวทคู. \\ วิทฺวา ปชหาติ ปาปกํ, \\ โกญฺโจ ขีรปโกว นินฺนคนฺ”ติ.\csromanspacer{}\csromanspacer{}สตฺตมํ.}}

\csromanheader{2}{อุทานปาฬิ \textbf{8. วิสาขาสุตฺต}}
\proseitem{78}{\csromanfolio{190}เอวํ เม สุตํ\csromandash{}เอกํ สมยํ ภควา สาวตฺถิยํ วิหรติ ปุพฺพาราเม มิคารมาตุปาสาเท.\csromanspacer{} เตน โข ปน สมเยน วิสาขาย มิคารมาตุยา นตฺตา กาลงฺกตา โหติ ปิยา มนาปา.\csromanspacer{} อถ โข วิสาขา มิคารมาตา อลฺลวตฺถา อลฺลเกสา ทิวา ทิวสฺส เยน ภควา เตนุปสงฺกมิ, อุปสงฺกมิตฺวา ภควนฺตํ อภิวาเทตฺวา เอกมนฺตํ นิสีทิ, เอกมนฺตํ นิสินฺนํ โข วิสาขํ มิคารมาตรํ ภควา เอตทโวจ\csromandash{}}

\csromanmatikaanchor{mtk.108}
\csromantocmarkref{subsection}{3. ปฐมสตฺตสุตฺต}{mtk.108}
\bookmark[dest={mtk.108},level=2]{3. ปฐมสตฺตสุตฺต}
\prose{หนฺท กุโต นุ ตฺวํ วิสาเข อาคจฺฉสิ อลฺลวตฺถา อลฺลเกสา อิธูปสงฺกนฺตา ทิวา ทิวสฺสาติ.\csromanspacer{} นตฺตา เม ภนฺเต ปิยา มนาปา กาลงฺกตา, เตนาหํ อลฺลวตฺถา อลฺลเกสา อิธูปสงฺกนฺตา ทิวา ทิวสฺสาติ.\csromanspacer{} อิจฺเฉ ยฺยาสิ ตฺวํ วิสาเข ยาวติกา\footnote{ยาวตกา (?)} สาวตฺถิยา มนุสฺสา, ตาวติเก\footnote{ตาวตเก (?)} ปุตฺเต จ นตฺตาโร จาติ.\csromanspacer{} อิจฺเฉยฺยาหํ ภควา\footnote{อิจฺเฉยฺยาหํ ภนฺเต ภควา (สฺยา)} ยาวติกา สาวตฺถิยา มนุสฺสา, ตาวติเก ปุตฺเต จ นตฺตาโร จาติ.}

\prose{กีวพหุกา ปน วิสาเข สาวตฺถิยา มนุสฺสา เทวสิกํ กาลํ กโรนฺตีติ.\csromanspacer{} ทสปิ ภนฺเต สาวตฺถิยา มนุสฺสา เทวสิกํ กาลํ กโรนฺติ, นวปิ ภนฺเต.\csromanspacer{} อฏฺฐปิ ภนฺเต.\csromanspacer{} สตฺตปิ ภนฺเต.\csromanspacer{} ฉปิ ภนฺเต.\csromanspacer{} ปญฺจปิ ภนฺเต.\csromanspacer{} จตฺตาโรปิ ภนฺเต.\csromanspacer{} ตีณิปิ ภนฺเต.\csromanspacer{} เทฺวปิ ภนฺเต สาวตฺถิยา มนุสฺสา เทวสิกํ กาลํ กโรนฺติ, เอโกปิ ภนฺเต สาวตฺถิยา มนุสฺโส เทวสิกํ กาลํ กโรติ, อวิวิตฺตา ภนฺเต สาวตฺถิ มนุสฺเสหิ กาลํ กโรนฺเตหีติ.}

\prose{ตํ กิํ มญฺญสิ วิสาเข, อปิ นุ ตฺวํ กทาจิ กรหจิ อนลฺลวตฺถา วา ภเวยฺยาสิ อนลฺลเกสา วาติ.\csromanspacer{} โน เหตํ ภนฺเต, อลํ เม ภนฺเต ตาว พหุเกหิ ปุตฺเตหิ จ นตฺตาเรหิ จาติ.}

\prose{เยสํ โข วิสาเข สตํ ปิยานิ, สตํ เตสํ ทุกฺขานิ, เยสํ นวุติ ปิยานิ, นวุติ เตสํ ทุกฺขานิ, เยสํ อสีติ ปิยานิ, อสีติ เตสํ ทุกฺขานิ, เยสํ สตฺตติ ปิยานิ, สตฺตติ เตสํ ทุกฺขานิ, เยสํ สฏฺฐิ ปิยานิ, สฏฺฐิเตสํ ทุกฺขานิ, เยสํ ปญฺญาสํ ปิยานิ, ปญฺญาสํ เตสํ ทุกฺขานิ, เยสํ จตฺตารีสํ ปิยานิ, จตฺตารีสํ เตสํ ทุกฺขานิ, เยสํ ติํสํ ปิยานิ, ติํสํ เตสํ ทุกฺขานิ, เยสํ วีสติ ปิยานิ, วีสติ เตสํ ทุกฺขานิ, เยสํ}

\vspace{\tipitakaparskip}
\csromancenter{ปาฏลิคามิยวคฺค ทส ปิยานิ, ทส เตสํ ทุกฺขานิ, เยสํ นว ปิยานิ, นว เตสํ ทุกฺขานิ, เยสํ อฏฺฐ ปิยานิ, อฏฺฐ เตสํ ทุกฺขานิ, เยสํ สตฺต ปิยานิ, สตฺต เตสํ ทุกฺขานิ, เยสํ ฉ ปิยานิ, ฉ เตสํ ทุกฺขานิ, เยสํ ปญฺจ ปิยานิ, ปญฺจ เตสํ ทุกฺขานิ, เยสํ จตฺตาริ ปิยานิ, จตฺตาริ เตสํ ทุกฺขานิ, เยสํ ตีณิ ปิยานิ, ตีณิ เตสํ ทุกฺขานิ, เยสํ เทฺว ปิยานิ, เทฺว เตสํ ทุกฺขานิ, เยสํ เอกํ ปิยํ, เอกํ เตสํ ทุกฺขํ, เยสํ นตฺถิ ปิยํ, นตฺถิ เตสํ ทุกฺขํ, “อโสกา เต วิรชา อนุปายาสา”ติ วทามีติ.}
\csromanmatikaanchor{mtk.109}
\csromantocmarkref{subsection}{4. ทุติยสตฺตสุตฺต}{mtk.109}
\bookmark[dest={mtk.109},level=2]{4. ทุติยสตฺตสุตฺต}
\prose{\csromanfolio{191}อถ โข ภควา เอตมตฺถํ วิทิตฺวา ตายํ เวลายํ อิมํ อุทานํ อุทาเนสิ\csromandash{}}

\csromangathameasure{*“เย เกจิ โสกา ปริเทวิตา วา, \\ ทุกฺขา ว โลกสฺมิมเนกรูปา. \\ ปิยํ ปฏิจฺจปฺปภวนฺติ เอเต, \\ ปิเย อสนฺเต น ภวนฺติ เอเต. \\ ตสฺมา หิ เต สุขิโน วีตโสกา, \\ เยสํ ปิยํ นตฺถิ กุหิญฺจิ โลเก. \\ ตสฺมา อโสกํ วิรชํ ปตฺถยาโน, \\ ปิยํ น กยิราถ กุหิญฺจิ โลเก”ติ.อฏฺฐมํ.}
\csromangathagroup{\csromangathastanza{\csromansymbolfootnote{*}{ขุ 10. 57, 178 ปิฏฺเฐสุปิ.}“เย เกจิ โสกา ปริเทวิตา วา, \\ ทุกฺขา ว\footnote{ทุกฺขา จ (สพฺพตฺถ)} โลกสฺมิมเนกรูปา. \\ ปิยํ ปฏิจฺจปฺปภวนฺติ เอเต, \\ ปิเย อสนฺเต น ภวนฺติ เอเต.}\csromangathastanza{ตสฺมา หิ เต สุขิโน วีตโสกา, \\ เยสํ ปิยํ นตฺถิ กุหิญฺจิ โลเก. \\ ตสฺมา อโสกํ วิรชํ ปตฺถยาโน, \\ ปิยํ น กยิราถ กุหิญฺจิ โลเก”ติ.\csromanspacer{}\csromanspacer{}อฏฺฐมํ.}}

\csromanheader{2}{9. ปฐมทพฺพสุตฺต}
\proseitem{79}{เอวํ เม สุตํ\csromandash{}เอกํ สมยํ ภควา ราชคเห วิหรติ เวฬุวเน กลนฺทกนิวาเป.\csromanspacer{} อถ โข อายสฺมา ทพฺโพ มลฺลปุตฺโต เยน ภควา เตนุปสงฺกมิ, อุปสงฺกมิตฺวา ภควนฺตํ อภิวาเทตฺวา เอกมนฺตํ นิสีทิ, เอกมนฺตํ นิสินฺโน โข อายสฺมา ทพฺโพ มลฺลปุตฺโต ภควนฺตํ เอตทโวจ “ปรินิพฺพาน กาโล เม ทานิ สุคตา”ติ.\csromanspacer{} ยสฺสทานิ ตฺวํ ทพฺพ กาลํ มญฺญสีติ.}

\prose{อถ โข อายสฺมา ทพฺโพ มลฺลปุตฺโต อุฏฺฐายาสนา ภควนฺตํ อภิวาเทตฺวา ปทกฺขิณํ กตฺวา เวหาสํ อพฺภุคฺคนฺตฺวา อากาเส อนฺตลิกฺเข ปลฺลงฺเกน นิสีทิตฺวา เตโชธาตุํ สมาปชฺชิตฺวา วุฏฺฐหิตฺวา ปรินิพฺพายิ.}

\csromanmatikaanchor{mtk.110}
\csromantocmarkref{subsection}{5. อปรลกุณฺฑกภทฺทิยสุตฺต}{mtk.110}
\bookmark[dest={mtk.110},level=2]{5. อปรลกุณฺฑกภทฺทิยสุตฺต}
\gambhira{อุทานปาฬิ}
\prose{\csromanfolio{192}อถ โข อายสฺมโต ทพฺพสฺส มลฺลปุตฺตสฺส เวหาสํ อพฺภุคฺคนฺตฺวา อากาเส อนฺตลิกฺเข ปลฺลงฺเกน นิสีทิตฺวา เตโชธาตุํ สมาปชฺชิตฺวา วุฏฺฐหิตฺวา ปรินิพฺพุตสฺส สรีรสฺส ฌายมานสฺส ฑยฺหมานสฺส เนว ฉาริกา ปญฺญายิตฺถ น มสิ.\csromanspacer{} เสยฺยถาปิ นาม สปฺปิสฺส วา เตลสฺส วา ฌายมานสฺส ฑยฺหมานสฺส เนว ฉาริกา ปญฺญายติ น มสิ, เอวเมวํ อายสฺมโต ทพฺพสฺส มลฺลปุตฺตสฺส เวหาสํ อพฺภุคฺคนฺตฺวา อากาเส อนฺตลิกฺเข ปลฺลงฺเกน นิสีทิตฺวา เตโชธาตุํ สมาปชฺชิตฺวา วุฏฺฐหิตฺวา ปรินิพฺพุตสฺส สรีรสฺส ฌายมานสฺส ฑยฺหมานสฺส เนว ฉาริกา ปญฺญายิตฺถ น มสีติ.}

\prose{อถ โข ภควา เอตมตฺถํ วิทิตฺวา ตายํ เวลายํ อิมํ อุทานํ อุทาเนสิ\csromandash{}}

\csromangathameasure{“อเภทิ กาโย นิโรธิ สญฺญา, \\ เวทนา สีติภวิํสุ สพฺพา. \\ วูปสมิํสุ สงฺขารา, \\ วิญฺญาณํ อตฺถมาคมา”ติ.นวมํ.}
\csromangathagroup{\csromangathastanza{“อเภทิ กาโย นิโรธิ สญฺญา, \\ เวทนา สีติภวิํสุ\footnote{ปีติทหํสุ (สฺยา, อิ), สีติทหิํสุ (ก)} สพฺพา. \\ วูปสมิํสุ สงฺขารา, \\ วิญฺญาณํ อตฺถมาคมา”ติ.\csromanspacer{}\csromanspacer{}นวมํ.}}

\csromanheader{2}{10. ทุติยทพฺพสุตฺต}
\proseitem{80}{เอวํ เม สุตํ\csromandash{}เอกํ สมยํ ภควา สาวตฺถิยํ วิหรติ เชตวเน อนาถปิณฺฑิกสฺส อาราเม.\csromanspacer{} ตตฺร โข ภควา ภิกฺขู อามนฺเตสิ “ภิกฺขโว”ติ. “ภทนฺเต”ติ เต ภิกฺขู ภควโต ปจฺจสฺโสสุํ.\csromanspacer{} ภควา เอตทโวจ\csromandash{}}

\prose{ทพฺพสฺส ภิกฺขเว มลฺลปุตฺตสฺส เวหาสํ อพฺภุคฺคนฺตฺวา อากาเส อนฺตลิกฺเข ปลฺลงฺเกน นิสีทิตฺวา เตโชธาตุํ สมาปชฺชิตฺวา วุฏฺฐหิตฺวา ปรินิพฺพุตสฺส สรีรสฺส ฌายมานสฺส ฑยฺหมานสฺส เนว ฉาริกา ปญฺญายิตฺถ น มสิ.\csromanspacer{} เสยฺยถาปิ นาม สปฺปิสฺส วา เตลสฺส วา ฌายมานสฺส ฑยฺหมานสฺส เนว ฉาริกา ปญฺญายติ น มสิ, เอวเมวํ โข ภิกฺขเว ทพฺพสฺส มลฺลปุตฺตสฺส เวหาสํ อพฺภุคฺคนฺตฺวา อากาเส อนฺตลิกฺเข ปลฺลงฺเกน นิสีทิตฺวา เตโชธาตุํ สมาปชฺชิตฺวา วุฏฺฐหิตฺวา ปรินิพฺพุตสฺส สรีรสฺส ฌายมานสฺส ฑยฺหมานสฺส เนว ฉาริกา ปญฺญายิตฺถ น มสีติ.}

\chapterheadpageread{8. ปาฏลิคามิยวคฺค}
\prose{\csromanfolio{193}อถ โข ภควา เอตมตฺถํ วิทิตฺวา ตายํ เวลายํ อิมํ อุทานํ อุทาเนสิ\csromandash{}}

\csromanmatikaanchor{mtk.111}
\csromantocmarkref{subsection}{6. ตณฺหาสงฺขยสุตฺต}{mtk.111}
\bookmark[dest={mtk.111},level=2]{6. ตณฺหาสงฺขยสุตฺต}
\csromantocmarkref{subsection}{7. ปปญฺจขยสุตฺต}{mtk.112}
\bookmark[dest={mtk.112},level=2]{7. ปปญฺจขยสุตฺต}
\prose{“อโยฆนหตสฺเสว, ชลโต ชาตเวทโส\footnote{ชาตเวทสฺส (สฺยา)}.}

\prose{อนุปุพฺพูปสนฺตสฺส, ยถา น ญายเต คติ.}

\prose{เอวํ สมฺมาวิมุตฺตานํ, กามพนฺโธฆตารินํ.}

\prose{ปญฺญาเปตุํ คติ นตฺถิ, ปตฺตานํ อจลํ สุขนฺ”ติ.\csromanspacer{}\csromanspacer{}ทสมํ.\csromanspacer{} ปาฏลิคามิยวคฺโค\footnote{ปาฏลิคามวคฺโค (ก)} อฏฺฐโม.}

\titlehead{\textbf{ตสฺสุทฺทานํ}}
\csromangathameasure{นิพฺพานา จตุโร วุตฺตา, จุนฺโท ปาฏลิคามิยา. \\ ทฺวิธาปโถ วิสาขา จ, ทพฺเพน สห เต ทสาติ.}
\csromangathagroup{\csromangathastanza{นิพฺพานา จตุโร วุตฺตา, จุนฺโท ปาฏลิคามิยา. \\ ทฺวิธาปโถ วิสาขา จ, ทพฺเพน สห เต ทสาติ.}}

\csromanheader{2}{\textbf{อุทาเน วคฺคานมุทฺทานํ}}
\csromangathameasure{วคฺคมิทํ ปฐมํ วรโพธิ, วคฺคมิทํ ทุติยํ มุจลินฺโท. \\ นนฺทกวคฺควโร ตติโย ตุ, เมฆิยวคฺควโร จ จตุตฺโถ. \\ ปญฺจมวคฺควรนฺติธ โสโณ, ฉฏฺฐมวคฺควรนฺติ ชจฺจนฺโธ. \\ สตฺตมวคฺควรนฺติ จ จูโฬ, ปาฏลิคามิยมฏฺฐมวคฺโค. \\ อสีติมนูนกสุตฺตวรํ, วคฺคมิทฏฺฐกํ สุวิภตฺตํ. \\ ทสฺสิตํ จกฺขุมตา วิมเลน, อทฺธา หิ ตํ อุทานมิตีทมาหุ.}
\csromangathagroup{\csromangathastanza{วคฺคมิทํ ปฐมํ วรโพธิ, วคฺคมิทํ ทุติยํ มุจลินฺโท. \\ นนฺทกวคฺควโร ตติโย ตุ, เมฆิยวคฺควโร จ จตุตฺโถ.}\csromangathastanza{ปญฺจมวคฺควรนฺติธ โสโณ, ฉฏฺฐมวคฺควรนฺติ ชจฺจนฺโธ\footnote{ฉฏฺฐมวคฺควรํ ตุ ตมนฺโธ (สี, ก)}. \\ สตฺตมวคฺควรนฺติ จ จูโฬ, ปาฏลิคามิยมฏฺฐมวคฺโค\footnote{ฉฏฺฐมวคฺควรํ ตุ ตมนฺโธ (สี, ก)}.}\csromangathastanza{อสีติมนูนกสุตฺตวรํ, วคฺคมิทฏฺฐกํ สุวิภตฺตํ. \\ ทสฺสิตํ จกฺขุมตา วิมเลน, อทฺธา หิ ตํ อุทานมิตีทมาหุ\footnote{อตฺถาเยตํ อุทานมิติมาหุ (ก), สทฺธา หิ ตํ อุทานนฺติทมาหุ (สฺยา, กํ, อิ)}.}}

\nitthitam{\textbf{อุทานปาฬิ นิฏฺฐิตา.}}
\csromanheader{2}{\textbf{ขุทฺทกนิกาย}}
\csromanmatikaanchor{mtk.113}
\csromantocmarkref{subsection}{8. กจฺจานสุตฺต}{mtk.113}
\bookmark[dest={mtk.113},level=2]{8. กจฺจานสุตฺต}
\gambhira{\textbf{อิติวุตฺตกปาฬิ}}
\namakkaram{นโม ตสฺส ภควโต อรหโต สมฺมาสมฺพุทฺธสฺส.}
\chapterhead{1. เอกกนิปาต}
\csromanheader{1}{1. ปฐมวคฺค}
\csromanheader{2}{1. โลภสุตฺต}
\csromanmatikaanchor{mtk.114}
\csromantocmarkref{subsection}{9. อุทปานสุตฺต}{mtk.114}
\bookmark[dest={mtk.114},level=2]{9. อุทปานสุตฺต}
\csromanheader{2}{1. วุตฺตํ เหตํ ภควตา วุตฺตมรหตาติ เม สุตํ\csromandash{}}
\csromanmatikaanchor{mtk.115}
\csromanmatikaanchor{mtk.130}
\csromantocmarkref{subsection}{10. อุเตนสุตฺต}{mtk.115}
\bookmark[dest={mtk.115},level=2]{10. อุเตนสุตฺต}
\prose{\csromanfolio{195}เอกธมฺมํ ภิกฺขเว ปชหถ, อหํ โว ปาฏิโภโค อนาคามิตาย.\csromanspacer{} กตมํ เอกธมฺมํ? โลภํ ภิกฺขเว เอกธมฺมํ ปชหถ, อหํ โว ปาฏิโภโค อนาคามิตายาติ.\csromanspacer{} เอตมตฺถํ ภควา อโวจ, ตตฺเถตํ อิติ วุจฺจติ\csromandash{}}

\csromangathameasure{“เยน โลเภน ลุทฺธาเส, สตฺตา คจฺฉนฺติ ทุคฺคติํ. \\ ตํ โลภํ สมฺมทญฺญาย, ปชหนฺติ วิปสฺสิโน. \\ ปหาย น ปุนายนฺติ, อิมํ โลกํ กุทาจนนฺ”ติ.}
\csromangathagroup{\csromangathastanza{“เยน โลเภน ลุทฺธาเส, สตฺตา คจฺฉนฺติ ทุคฺคติํ. \\ ตํ โลภํ สมฺมทญฺญาย, ปชหนฺติ วิปสฺสิโน. \\ ปหาย น ปุนายนฺติ, อิมํ โลกํ กุทาจนนฺ”ติ.}}

\prose{อยมฺปิ อตฺโถ วุตฺโต ภควตา อิติ เม สุตนฺติ.\csromanspacer{}\csromanspacer{}ปฐมํ.}

\csromanheader{2}{2. โทสสุตฺต}
\csromanheader{2}{2. วุตฺตํ เหตํ ภควตา วุตฺตมรหตาติ เม สุตํ\csromandash{}}
\prose{เอกธมฺมํ ภิกฺขเว ปชหถ, อหํ โว ปาฏิโภโค อนาคามิตาย.\csromanspacer{} กตมํ เอกธมฺมํ? โทสํ ภิกฺขเว เอกธมฺมํ ปชหถ, อหํ โว ปาฏิโภโค อนาคามิตายาติ.\csromanspacer{} เอตมตฺถํ ภควา อโวจ, ตตฺเถตํ อิติ วุจฺจติ\csromandash{}}

\gambhira{อิติวุตฺตกปาฬิ}
\csromangathameasure{“เยน โทเสน ทุฏฺฐาเส, สตฺตา คจฺฉนฺติ ทุคฺคติํ. \\ ตํ โทสํ สมฺมทญฺญาย, ปชหนฺติ วิปสฺสิโน. \\ ปหาย น ปุนายนฺติ, อิมํ โลกํ กุทาจนนฺ”ติ.}
\csromangathagroup{\csromangathastanza{\csromanfolio{196}“เยน โทเสน ทุฏฺฐาเส, สตฺตา คจฺฉนฺติ ทุคฺคติํ. \\ ตํ โทสํ สมฺมทญฺญาย, ปชหนฺติ วิปสฺสิโน. \\ ปหาย น ปุนายนฺติ, อิมํ โลกํ กุทาจนนฺ”ติ.}}

\prose{อยมฺปิ อตฺโถ วุตฺโต ภควตา อิติ เม สุตนฺติ.\csromanspacer{}\csromanspacer{}ทุติยํ.}

\csromanheader{2}{3. โมหสุตฺต}
\csromanheader{2}{3. วุตฺตํ เหตํ ภควตา วุตฺตมรหตาติ เม สุตํ\csromandash{}}
\prose{เอกธมฺมํ ภิกฺขเว ปชหถ, อหํ โว ปาฏิโภโค อนาคามิตาย.\csromanspacer{} กตมํ เอกธมฺมํ? โมหํ ภิกฺขเว เอกธมฺมํ ปชหถ, อหํ โว ปาฏิโภโค อนาคามิตายาติ.\csromanspacer{} เอตมตฺถํ ภควา อโวจ, ตตฺเถตํ อิติ วุจฺจติ\csromandash{}}

\csromangathameasure{“เยน โมเหน มูฬฺหาเส, สตฺตา คจฺฉนฺติ ทุคฺคติํ. \\ ตํ โมหํ สมฺมทญฺญาย, ปชหนฺติ วิปสฺสิโน. \\ ปหาย น ปุนายนฺติ, อิมํ โลกํ กุทาจนนฺ”ติ.}
\csromangathagroup{\csromangathastanza{“เยน โมเหน มูฬฺหาเส, สตฺตา คจฺฉนฺติ ทุคฺคติํ. \\ ตํ โมหํ สมฺมทญฺญาย, ปชหนฺติ วิปสฺสิโน. \\ ปหาย น ปุนายนฺติ, อิมํ โลกํ กุทาจนนฺ”ติ.}}

\prose{อยมฺปิ อตฺโถ วุตฺโต ภควตา อิติ เม สุตนฺติ.\csromanspacer{}\csromanspacer{}ตติยํ.}

\csromanheader{2}{4. โกธสุตฺต}
\csromanheader{2}{4. วุตฺตํ เหตํ ภควตา วุตฺตมรหตาติ เม สุตํ\csromandash{}}
\prose{เอกธมฺมํ ภิกฺขเว ปชหถ, อหํ โว ปาฏิโภโค อนาคามิตาย.\csromanspacer{} กตมํ เอกธมฺมํ? โกธํ ภิกฺขเว เอกธมฺมํ ปชหถ, อหํ โว ปาฏิโภโค อนาคามิตายาติ.\csromanspacer{} เอตมตฺถํ ภควา อโวจ, ตตฺเถตํ อิติ วุจฺจติ\csromandash{}}

\csromangathameasure{“เยน โกเธน กุทฺธาเส, สตฺตา คจฺฉนฺติ ทุคฺคติํ. \\ ตํ โกธํ สมฺมทญฺญาย, ปชหนฺติ วิปสฺสิโน. \\ ปหาย น ปุนายนฺติ, อิมํ โลกํ กุทาจนนฺ”ติ.}
\csromangathagroup{\csromangathastanza{“เยน โกเธน กุทฺธาเส, สตฺตา คจฺฉนฺติ ทุคฺคติํ. \\ ตํ โกธํ สมฺมทญฺญาย, ปชหนฺติ วิปสฺสิโน. \\ ปหาย น ปุนายนฺติ, อิมํ โลกํ กุทาจนนฺ”ติ.}}

\csromanmatikaanchor{mtk.116}
\csromanmatikaanchor{mtk.117}
\csromantocmarknopage{chapter}{8. ปาฏลิคามิยวคฺค}{mtk.116}
\bookmark[dest={mtk.116},level=0]{8. ปาฏลิคามิยวคฺค}
\csromantocmarkref{subsection}{1. ปฐมนิพฺพานปฏิสํยุตฺตสุตฺต}{mtk.117}
\bookmark[dest={mtk.117},level=2]{1. ปฐมนิพฺพานปฏิสํยุตฺตสุตฺต}
\prose{อยมฺปิ อตฺโถ วุตฺโต ภควตา อิติ เม สุตนฺติ.\csromanspacer{}\csromanspacer{}จตุตฺถํ.}

\csromanheader{2}{1. เอกกปิปาต \textbf{5. มกฺขสุตฺต}}
\csromanheader{2}{5. วุตฺตํ เหตํ ภควตา วุตฺตมรหตาติ เม สุตํ\csromandash{}}
\csromanmatikaanchor{mtk.118}
\csromanmatikaanchor{mtk.131}
\csromantocmarkref{subsection}{2. ทุติยนิพฺพานปฏิสํยุตฺตสุตฺต}{mtk.118}
\bookmark[dest={mtk.118},level=2]{2. ทุติยนิพฺพานปฏิสํยุตฺตสุตฺต}
\prose{\csromanfolio{197}เอกธมฺมํ ภิกฺขเว ปชหถ, อหํ โว ปาฏิโภโค อนาคามิตาย.\csromanspacer{} กตมํ เอกธมฺมํ? มกฺขํ ภิกฺขเว เอกธมฺมํ ปชหถ, อหํ โว ปาฏิโภโค อนาคามิตายาติ.\csromanspacer{} เอตมตฺถํ ภควา อโวจ, ตตฺเถตํ อิติ วุจฺจติ\csromandash{}}

\csromangathameasure{“เยน มกฺเขน มกฺขาเส, สตฺตา คจฺฉนฺติ ทุคฺคติํ. \\ ตํ มกฺขํ สมฺมทญฺญาย, ปชหนฺติ วิปสฺสิโน. \\ ปหาย น ปุนายนฺติ, อิมํ โลกํ กุทาจนนฺ”ติ.}
\csromangathagroup{\csromangathastanza{“เยน มกฺเขน มกฺขาเส\footnote{มกฺขิตาเส (สฺยา)}, สตฺตา คจฺฉนฺติ ทุคฺคติํ. \\ ตํ มกฺขํ สมฺมทญฺญาย, ปชหนฺติ วิปสฺสิโน. \\ ปหาย น ปุนายนฺติ, อิมํ โลกํ กุทาจนนฺ”ติ.}}

\prose{อยมฺปิ อตฺโถ วุตฺโต ภควตา อิติ เม สุตนฺติ.\csromanspacer{}\csromanspacer{}ปญฺจมํ.}

\csromanheader{2}{6. มานสุตฺต}
\csromanheader{2}{6. วุตฺตํ เหตํ ภควตา วุตฺตมรหตาติ เม สุตํ\csromandash{}}
\csromanmatikaanchor{mtk.119}
\csromantocmarkref{subsection}{3. ตติยนิพฺพานปฏิสํยุตฺตสุตฺต}{mtk.119}
\bookmark[dest={mtk.119},level=2]{3. ตติยนิพฺพานปฏิสํยุตฺตสุตฺต}
\prose{เอกธมฺมํ ภิกฺขเว ปชหถ, อหํ โว ปาฏิโภโค อนาคามิตาย.\csromanspacer{} กตมํ เอกธมฺมํ? มานํ ภิกฺขเว เอกธมฺมํ ปชหถ, อหํ โว ปาฏิโภโค อนาคามิตายาติ.\csromanspacer{} เอตมตฺถํ ภควา อโวจ, ตตฺเถตํ อิติ วุจฺจติ\csromandash{}}

\csromangathameasure{“เยน มาเนน มตฺตาเส, สตฺตา คจฺฉนฺติ ทุคฺคติํ. \\ ตํ มานํ สมฺมทญฺญาย, ปชหนฺติ วิปสฺสิโน. \\ ปหาย น ปุนายนฺติ, อิมํ โลกํ กุทาจนนฺ”ติ.}
\csromangathagroup{\csromangathastanza{“เยน มาเนน มตฺตาเส, สตฺตา คจฺฉนฺติ ทุคฺคติํ. \\ ตํ มานํ สมฺมทญฺญาย, ปชหนฺติ วิปสฺสิโน. \\ ปหาย น ปุนายนฺติ, อิมํ โลกํ กุทาจนนฺ”ติ.}}

\prose{อยมฺปิ อตฺโถ วุตฺโต ภควตา อิติ เม สุตนฺติ.\csromanspacer{}\csromanspacer{}ฉฏฺฐํ.}

\csromanmatikaanchor{mtk.120}
\csromantocmarkref{subsection}{4. จตุตฺถนิพฺพานปฏิสํยุตฺตสุตฺต}{mtk.120}
\bookmark[dest={mtk.120},level=2]{4. จตุตฺถนิพฺพานปฏิสํยุตฺตสุตฺต}
\csromanheader{2}{7. สพฺพปริญฺญาสุตฺต}
\csromanheader{2}{7. วุตฺตํ เหตํ ภควตา วุตฺตมรหตาติ เม สุตํ\csromandash{}}
\prose{สพฺพํ ภิกฺขเว อนภิชานํ อปริชานํ ตตฺถ จิตฺตํ อวิราชยํ อปฺปชหํ อภพฺโพ ทุกฺขกฺขยาย.\csromanspacer{} สพฺพญฺจ โข ภิกฺขเว อภิชานํ ปริชานํ ตตฺถ จิตฺตํ วิราชยํ ปชหํ ภพฺโพ ทุกฺขกฺขยายาติ.\csromanspacer{} เอตมตฺถํ ภควา อโวจ, ตตฺเถตํ อิติ วุจฺจติ\csromandash{}}

\gambhira{อิติวุตฺตกปาฬิ}
\csromanmatikaanchor{mtk.121}
\csromantocmarkref{subsection}{5. จุนฺทสุตฺต}{mtk.121}
\bookmark[dest={mtk.121},level=2]{5. จุนฺทสุตฺต}
\csromangathameasure{“โย สพฺพํ สพฺพโต ญตฺวา, สพฺพตฺเถสุ น รชฺชติ. \\ ส เว สพฺพปริญฺญา โส, สพฺพทุกฺขมุปจฺจคา”ติ.}
\csromangathagroup{\csromangathastanza{\csromanfolio{198}“โย สพฺพํ สพฺพโต ญตฺวา, สพฺพตฺเถสุ น รชฺชติ. \\ ส เว สพฺพปริญฺญา\footnote{สพฺพํ ปริญฺญา (สฺยา, อิ)} โส, สพฺพทุกฺขมุปจฺจคา”ติ\footnote{สพฺพํ ทุกฺขํ อุปจฺจคาติ (สฺยา), สพฺพทุกฺขํ อุปจฺจคาติ (อิ-ฏฺฐ)}.}}

\prose{อยมฺปิ อตฺโถ วุตฺโต ภควตา อิติ เม สุตนฺติ.\csromanspacer{}\csromanspacer{}สตฺตมํ.}

\csromanheader{2}{8. มานปริญฺญาสุตฺต}
\csromanheader{2}{8. วุตฺตํ เหตํ ภควตา วุตฺตมรหตาติ เม สุตํ\csromandash{}}
\prose{มานํ ภิกฺขเว อนภิชานํ อปริชานํ ตตฺถ จิตฺตํ อวิราชยํ อปฺปชหํ อภพฺโพ ทุกฺขกฺขยาย.\csromanspacer{} มานญฺจ โข ภิกฺขเว อภิชานํ ปริชานํ ตตฺถ จิตฺตํ วิราชยํ ปชหํ ภพฺโพ ทุกฺขกฺขยายาติ.\csromanspacer{} เอตมตฺถํ ภควา อโวจ, ตตฺเถตํ อิติ วุจฺจติ\csromandash{}}

\csromangathameasure{“มานุเปตา อยํ ปชา, มานคนฺถา ภเว รตา. \\ มานํ อปริชานนฺตา, อาคนฺตาโร ปุนพฺภวํ. \\ เย จ มานํ ปหนฺตฺวาน, วิมุตฺตา มานสงฺขเย. \\ เต มานคนฺถาภิภุโน, สพฺพทุกฺขมุปจฺจคุนฺ”ติ.}
\csromangathagroup{\csromangathastanza{“มานุเปตา อยํ ปชา, มานคนฺถา ภเว รตา. \\ มานํ อปริชานนฺตา, อาคนฺตาโร ปุนพฺภวํ.}\csromangathastanza{เย จ มานํ ปหนฺตฺวาน, วิมุตฺตา มานสงฺขเย. \\ เต มานคนฺถาภิภุโน, สพฺพทุกฺขมุปจฺจคุนฺ”ติ\footnote{สพฺพทุกฺขํ อุปจฺจคุนฺติ (อิ)}.}}

\prose{อยมฺปิ อตฺโถ วุตฺโต ภควตา อิติ เม สุตนฺติ.\csromanspacer{}\csromanspacer{}อฏฺฐมํ.}

\csromanheader{2}{9. โลภปริญฺญาสุตฺต}
\csromanheader{2}{9. วุตฺตํ เหตํ ภควตา วุตฺตมรหตาติ เม สุตํ\csromandash{}}
\prose{โลภํ ภิกฺขเว อนภิชานํ อปริชานํ ตตฺถ จิตฺตํ อวิราชยํ อปฺปชหํ อภพฺโพ ทุกฺขกฺขยาย.\csromanspacer{} โลภญฺจ โข ภิกฺขเว อภิชานํ ปริชานํ ตตฺถ จิตฺตํ วิราชยํ ปชหํ ภพฺโพ ทุกฺขกฺขยายาติ.\csromanspacer{} เอตมตฺถํ ภควา อโวจ, ตตฺเถตํ อิติ วุจฺจติ\csromandash{}}

\csromangathameasure{“เยน โลเภน ลุทฺธาเส, สตฺตา คจฺฉนฺติ ทุคฺคติํ. \\ ตํ โลภํ สมฺมทญฺญาย, ปชหนฺติ วิปสฺสิโน. \\ ปหาย น ปุนายนฺติ, อิมํ โลกํ กุทาจนนฺ”ติ.}
\csromangathagroup{\csromangathastanza{“เยน โลเภน ลุทฺธาเส, สตฺตา คจฺฉนฺติ ทุคฺคติํ. \\ ตํ โลภํ สมฺมทญฺญาย, ปชหนฺติ วิปสฺสิโน. \\ ปหาย น ปุนายนฺติ, อิมํ โลกํ กุทาจนนฺ”ติ.}}

\prose{อยมฺปิ อตฺโถ วุตฺโต ภควตา อิติ เม สุตนฺติ.\csromanspacer{}\csromanspacer{}นวมํ.}

\csromanheader{2}{1. เอกกปิปาต \textbf{10. โทสปริญฺญาสุตฺต}}
\csromanheader{2}{10. วุตฺตํ เหตํ ภควตา วุตฺตมรหตาติ เม สุตํ\csromandash{}}
\csromanmatikaanchor{mtk.122}
\csromanmatikaanchor{mtk.133}
\csromantocmarkref{subsection}{6. ปาฏลิคามิยสุตฺต}{mtk.122}
\bookmark[dest={mtk.122},level=2]{6. ปาฏลิคามิยสุตฺต}
\prose{\csromanfolio{199}โทสํ ภิกฺขเว อนภิชานํ อปริชานํ ตตฺถ จิตฺตํ อวิราชยํ อปฺปชหํ อภพฺโพ ทุกฺขกฺขยาย.\csromanspacer{} โทสญฺจ โข ภิกฺขเว อภิชานํ ปริชานํ ตตฺถ จิตฺตํ วิราชยํ ปชหํ ภพฺโพ ทุกฺขกฺขยายาติ.\csromanspacer{} เอตมตฺถํ ภควา อโวจ, ตตฺเถตํ อิติ วุจฺจติ\csromandash{}}

\csromangathameasure{“เยน โทเสน ทุฏฺฐาเส, สตฺตา คจฺฉนฺติ ทุคฺคติํ. \\ ตํ โทสํ สมฺมทญฺญาย, ปชหนฺติ วิปสฺสิโน. \\ ปหาย น ปุนายนฺติ, อิมํ โลกํ กุทาจนนฺ”ติ.}
\csromangathagroup{\csromangathastanza{“เยน โทเสน ทุฏฺฐาเส, สตฺตา คจฺฉนฺติ ทุคฺคติํ. \\ ตํ โทสํ สมฺมทญฺญาย, ปชหนฺติ วิปสฺสิโน. \\ ปหาย น ปุนายนฺติ, อิมํ โลกํ กุทาจนนฺ”ติ.}}

\prose{อยมฺปิ อตฺโถ วุตฺโต ภควตา อิติ เม สุตนฺติ.\csromanspacer{}\csromanspacer{}ทสมํ.\csromanspacer{} ปาฏิโภควคฺโค ปฐโม.}

\titlehead{\textbf{ตสฺสุทฺทานํ}}
\csromangathameasure{ราคโทสา อถ โมโห, โกธมกฺขา มานํ สพฺพํ. \\ มานโต ราคโทสา ปุน เทฺว, ปกาสิตา วคฺคมาหุ ปฐมนฺติ.}
\csromangathagroup{\csromangathastanza{ราคโทสา อถ โมโห, โกธมกฺขา มานํ สพฺพํ. \\ มานโต ราคโทสา ปุน เทฺว, ปกาสิตา วคฺคมาหุ ปฐมนฺติ.}}

\csromanheader{1}{2. ทุติยวคฺค}
\csromanheader{2}{1. โมหปริญฺญาสุตฺต}
\csromanheader{2}{11. วุตฺตํ เหตํ ภควตา วุตฺตมรหตาติ เม สุตํ\csromandash{}}
\prose{โมหํ ภิกฺขเว อนภิชานํ อปริชานํ ตตฺถ จิตฺตํ อวิราชยํ อปฺปชหํ อภพฺโพ ทุกฺขกฺขยาย.\csromanspacer{} โมหญฺจ โข ภิกฺขเว อภิชานํ ปริชานํ ตตฺถ จิตฺตํ วิราชยํ ปชหํ ภพฺโพ ทุกฺขกฺขยายาติ.\csromanspacer{} เอตมตฺถํ ภควา อโวจ, ตตฺเถตํ อิติ วุจฺจติ\csromandash{}}

\csromangathameasure{“เยน โมเหน มูฬฺหาเส, สตฺตา คจฺฉนฺติ ทุคฺคติํ. \\ ตํ โมหํ สมฺมทญฺญาย, ปชหนฺติ วิปสฺสิโน. \\ ปหาย น ปุนายนฺติ, อิมํ โลกํ กุทาจนนฺ”ติ.}
\csromangathagroup{\csromangathastanza{“เยน โมเหน มูฬฺหาเส, สตฺตา คจฺฉนฺติ ทุคฺคติํ. \\ ตํ โมหํ สมฺมทญฺญาย, ปชหนฺติ วิปสฺสิโน. \\ ปหาย น ปุนายนฺติ, อิมํ โลกํ กุทาจนนฺ”ติ.}}

\prose{อยมฺปิ อตฺโถ วุตฺโต ภควตา อิติ เม สุตนฺติ.\csromanspacer{}\csromanspacer{}ปฐมํ.}

\gambhira{อิติวุตฺตกปาฬิ \textbf{2. โกธปริญฺญาสุตฺต}}
\csromanheader{2}{12. วุตฺตํ เหตํ ภควตา วุตฺตมรหตาติ เม สุตํ\csromandash{}}
\csromanmatikaanchor{mtk.123}
\csromanmatikaanchor{mtk.134}
\csromantocmarkref{subsection}{7. ทฺวิธาปถสุตฺต}{mtk.123}
\bookmark[dest={mtk.123},level=2]{7. ทฺวิธาปถสุตฺต}
\prose{\csromanfolio{200}โกธํ ภิกฺขเว อนภิชานํ อปริชานํ ตตฺถ จิตฺตํ อวิราชยํ อปฺปชหํ อภพฺโพ ทุกฺขกฺขยาย.\csromanspacer{} โกธญฺจ โข ภิกฺขเว อภิชานํ ปริชานํ ตตฺถ จิตฺตํ วิราชยํ ปชหํ ภพฺโพ ทุกฺขกฺขยายาติ.\csromanspacer{} เอตมตฺถํ ภควา อโวจ, ตตฺเถตํ อิติ วุจฺจติ\csromandash{}}

\csromangathameasure{“เยน โกเธน กุทฺธาเส, สตฺตา คจฺฉนฺติ ทุคฺคติํ. \\ ตํ โกธํ สมฺมทญฺญาย, ปชหนฺติ วิปสฺสิโน. \\ ปหาย น ปุนายนฺติ, อิมํ โลกํ กุทาจนนฺ”ติ.}
\csromangathagroup{\csromangathastanza{“เยน โกเธน กุทฺธาเส, สตฺตา คจฺฉนฺติ ทุคฺคติํ. \\ ตํ โกธํ สมฺมทญฺญาย, ปชหนฺติ วิปสฺสิโน. \\ ปหาย น ปุนายนฺติ, อิมํ โลกํ กุทาจนนฺ”ติ.}}

\prose{อยมฺปิ อตฺโถ วุตฺโต ภควตา อิติ เม สุตนฺติ.\csromanspacer{}\csromanspacer{}ทุติยํ.}

\csromanheader{2}{3. มกฺขปริญฺญาสุตฺต}
\csromanheader{2}{13. วุตฺตํ เหตํ ภควตา วุตฺตมรหตาติ เม สุตํ\csromandash{}}
\prose{มกฺขํ ภิกฺขเว อนภิชานํ อปริชานํ ตตฺถ จิตฺตํ อวิราชยํ อปฺปชหํ อภพฺโพ ทุกฺขกฺขยาย.\csromanspacer{} มกฺขญฺจ โข ภิกฺขเว อภิชานํ ปริชานํ ตตฺถ จิตฺตํ วิราชยํ ปชหํ ภพฺโพ ทุกฺขกฺขยายาติ.\csromanspacer{} เอตมตฺถํ ภควา อโวจ, ตตฺเถตํ อิติ วุจฺจติ\csromandash{}}

\csromangathameasure{“เยน มกฺเขน มกฺขาเส, สตฺตา คจฺฉนฺติ ทุคฺคติํ. \\ ตํ มกฺขํ สมฺมทญฺญาย, ปชหนฺติ วิปสฺสิโน. \\ ปหาย น ปุนายนฺติ, อิมํ โลกํ กุทาจนนฺ”ติ.}
\csromangathagroup{\csromangathastanza{“เยน มกฺเขน มกฺขาเส, สตฺตา คจฺฉนฺติ ทุคฺคติํ. \\ ตํ มกฺขํ สมฺมทญฺญาย, ปชหนฺติ วิปสฺสิโน. \\ ปหาย น ปุนายนฺติ, อิมํ โลกํ กุทาจนนฺ”ติ.}}

\prose{อยมฺปิ อตฺโถ วุตฺโต ภควตา อิติ เม สุตนฺติ.\csromanspacer{}\csromanspacer{}ตติยํ.}

\csromanheader{2}{4. อวิชฺชานีวรณสุตฺต}
\csromanheader{2}{14. วุตฺตํ เหตํ ภควตา วุตฺตมรหตาติ เม สุตํ\csromandash{}}
\prose{นาหํ ภิกฺขเว อญฺญํ เอกนีวรณมฺปิ สมนุปสฺสามิ, เยน\footnote{เยเนวํ (?)} นีวรเณน นิวุตา ปชา ทีฆรตฺตํ สนฺธาวนฺติ สํสรนฺติ, ยถยิทํ ภิกฺขเว อวิชฺชานีวรณํ\footnote{อวิชฺชานีวรเณน (?)}.\csromanspacer{} อวิชฺชานีวรเณน หิ ภิกฺขเว นิวุตา ปชา ทีฆรตฺตํ สนฺธาวนฺติ สํสรนฺตีติ.\csromanspacer{} เอตมตฺถํ ภควา อโวจ, ตตฺเถตํ อิติ วุจฺจติ\csromandash{}}

\csromanheader{2}{1. เอกกปิปาต}
\csromanmatikaanchor{mtk.124}
\csromanmatikaanchor{mtk.135}
\csromantocmarkref{subsection}{8. วิสาขาสุตฺต}{mtk.124}
\bookmark[dest={mtk.124},level=2]{8. วิสาขาสุตฺต}
\csromangathameasure{“นตฺถญฺโญ เอกธมฺโมปิ, เยเนวํ นิวุตา ปชา. \\ สํสรนฺติ อโหรตฺตํ, ยถา โมเหน อาวุตา. \\ เย จ โมหํ ปหนฺตฺวาน, ตโมขนฺธํ ปทาลยุํ. \\ น เต ปุน สํสรนฺติ, เหตุ เตสํ น วิชฺชตี”ติ.}
\csromangathagroup{\csromangathastanza{\csromanfolio{201}“นตฺถญฺโญ เอกธมฺโมปิ, เยเนวํ\footnote{เยเนว (สี, อิ, ก)} นิวุตา ปชา. \\ สํสรนฺติ อโหรตฺตํ, ยถา โมเหน อาวุตา.}\csromangathastanza{เย จ โมหํ ปหนฺตฺวาน, ตโมขนฺธํ\footnote{ตโมกฺขนฺธํ (สี, สฺยา, อิ)} ปทาลยุํ. \\ น เต ปุน สํสรนฺติ, เหตุ เตสํ น วิชฺชตี”ติ.}}

\prose{อยมฺปิ อตฺโถ วุตฺโต ภควตา อิติ เม สุตนฺติ.\csromanspacer{}\csromanspacer{}จตุตฺถํ.}

\csromanheader{2}{5. ตณฺหาสํโยชนสุตฺต}
\csromanheader{2}{15. วุตฺตํ เหตํ ภควตา วุตฺตมรหตาติ เม สุตํ\csromandash{}}
\prose{นาหํ ภิกฺขเว อญฺญํ เอกสํโยชนมฺปิ สมนุปสฺสามิ, เยน\footnote{เยเนวํ (สฺยา)} สํโยชเนน สํยุตฺตา สตฺตา ทีฆรตฺตํ สนฺธาวนฺติ สํสรนฺติ, ยถยิทํ ภิกฺขเว ตณฺหาสํโยชนํ\footnote{ตณฺหาสํโยชเนน (?)}.\csromanspacer{} ตณฺหาสํโยชเนน หิ ภิกฺขเว สํยุตฺตา สตฺตา ทีฆรตฺตํ สนฺธาวนฺติ สํสรนฺตีติ.\csromanspacer{} เอตมตฺถํ ภควา อโวจ, ตตฺเถตํ อิติ วุจฺจติ\csromandash{}}

\csromangathameasure{*“ตณฺหาทุติโย ปุริโส, ทีฆมทฺธาน สํสรํ. \\ อิตฺถภาวญฺญถาภาวํ, สํสารํ นาติวตฺตติ. \\ เอตมาทีนวํ ญตฺวา, ตณฺหํ ทุกฺขสฺส สมฺภวํ. \\ วีตตโณฺห อนาทาโน, สโต ภิกฺขุ ปริพฺพเช”ติ.}
\csromangathagroup{\csromangathastanza{\csromansymbolfootnote{*}{อุปริ 268; ขุ 7. 362; ขุ 8. 211 ปิฏฺเฐสุปิ.}“ตณฺหาทุติโย ปุริโส, ทีฆมทฺธาน สํสรํ. \\ อิตฺถภาวญฺญถาภาวํ\footnote{อิตฺถมฺภาวญฺญถาภาวํ (สฺยา)}, สํสารํ นาติวตฺตติ.}\csromangathastanza{เอตมาทีนวํ\footnote{เอวมาทีนวํ (สี, อิ, ก)} ญตฺวา, ตณฺหํ\footnote{ตณฺหา (สี, ก)} ทุกฺขสฺส สมฺภวํ. \\ วีตตโณฺห อนาทาโน, สโต ภิกฺขุ ปริพฺพเช”ติ.}}

\prose{อยมฺปิ อตฺโถ วุตฺโต ภควตา อิติ เม สุตนฺติ.\csromanspacer{}\csromanspacer{}ปญฺจมํ.}

\csromanheader{2}{6. ปฐมเสขสุตฺต}
\csromanheader{2}{16. วุตฺตํ เหตํ ภควตา วุตฺตมรหตาติ เม สุตํ\csromandash{}}
\prose{เสขสฺส ภิกฺขเว ภิกฺขุโน อปฺปตฺตมานสสฺส อนุตฺตรํ โยคกฺเขมํ ปตฺถยมานสฺส วิหรโต “อชฺฌตฺติกํ องฺคนฺ”ติ กริตฺวา นาญฺญํ เอกงฺคมฺปิ สมนุปสฺสามิ, ยํ เอวํ พหูปการํ, ยถยิทํ ภิกฺขเว โยนิโสมนสิกาโร.\csromanspacer{} โยนิโส ภิกฺขเว ภิกฺขุ มนสิ กโรนฺโต อกุสลํ}

\vspace{\tipitakaparskip}
\csromancenter{อิติวุตฺตกปาฬิ ปชหติ, กุสลํ ภาเวตีติ.\csromanspacer{} เอตมตฺถํ ภควา อโวจ, ตตฺเถตํ อิติ วุจฺจติ\csromandash{}}
\vspace{0.5\baselineskip}
\csromanmatikaanchor{mtk.125}
\csromanmatikaanchor{mtk.136}
\csromantocmarkref{subsection}{9. ปฐมทพฺพสุตฺต}{mtk.125}
\bookmark[dest={mtk.125},level=2]{9. ปฐมทพฺพสุตฺต}
\csromangathameasure{“โยนิโสมนสิกาโร, ธมฺโม เสขสฺส ภิกฺขุโน. \\ นตฺถญฺโญ เอวํ พหุกาโร, อุตฺตมตฺถสฺส ปตฺติยา. \\ โยนิโส ปทหํ ภิกฺขุ, ขยํ ทุกฺขสฺส ปาปุเณ”ติ.}
\csromangathagroup{\csromangathastanza{\csromanfolio{202}“โยนิโสมนสิกาโร, ธมฺโม เสขสฺส ภิกฺขุโน. \\ นตฺถญฺโญ เอวํ พหุกาโร, อุตฺตมตฺถสฺส ปตฺติยา. \\ โยนิโส ปทหํ ภิกฺขุ, ขยํ ทุกฺขสฺส ปาปุเณ”ติ.}}

\prose{อยมฺปิ อตฺโถ วุตฺโต ภควตา อิติ เม สุตนฺติ.\csromanspacer{}\csromanspacer{}ฉฏฺฐํ.}

\csromanheader{2}{7. ทุติยเสขสุตฺต}
\csromanheader{2}{17. วุตฺตํ เหตํ ภควตา วุตฺตมรหตาติ เม สุตํ\csromandash{}}
\prose{เสขสฺส ภิกฺขเว ภิกฺขุโน อปฺปตฺตมานสสฺส อนุตฺตรํ โยคกฺเขมํ ปตฺถยมานสฺส วิหรโต “พาหิรํ องฺคนฺ”ติ กริตฺวา นาญฺญํ เอกงฺคมฺปิ สมนุปสฺสามิ, ยํ เอวํ พหูปการํ, ยถยิทํ ภิกฺขเว กลฺยาณมิตฺตตา.\csromanspacer{} กลฺยาณมิตฺโต ภิกฺขเว ภิกฺขุ อกุสลํ ปชหติ, กุสลํ ภาเวตีติ.\csromanspacer{} เอตมตฺถํ ภควา อโวจ, ตตฺเถตํ อิติ วุจฺจติ\csromandash{}}

\csromangathameasure{“กลฺยาณมิตฺโต โย ภิกฺขุ, สปฺปติสฺโส สคารโว. \\ กรํ มิตฺตานํ วจนํ, สมฺปชาโน ปติสฺสโต. \\ ปาปุเณ อนุปุพฺเพน, สพฺพสํโยชนกฺขยนฺ”ติ.}
\csromangathagroup{\csromangathastanza{“กลฺยาณมิตฺโต โย ภิกฺขุ, สปฺปติสฺโส สคารโว. \\ กรํ มิตฺตานํ วจนํ, สมฺปชาโน ปติสฺสโต. \\ ปาปุเณ อนุปุพฺเพน, สพฺพสํโยชนกฺขยนฺ”ติ.}}

\prose{อยมฺปิ อตฺโถ วุตฺโต ภควตา อิติ เม สุตนฺติ.\csromanspacer{}\csromanspacer{}สตฺตมํ.}

\csromanheader{2}{8. สํฆเภทสุตฺต}
\csromanheader{2}{18. วุตฺตํ เหตํ ภควตา วุตฺตมรหตาติ เม สุตํ\csromandash{}}
\prose{เอกธมฺโม ภิกฺขเว โลเก อุปฺปชฺชมาโน อุปฺปชฺชติ พหุชนาหิตาย พหุชนาสุขาย พหุโน ชนสฺส อนตฺถาย อหิตาย ทุกฺขาย เทวมนุสฺสานํ.\csromanspacer{} กตโม เอกธมฺโม? สํฆเภโท.\csromanspacer{} สํเฆ โข ปน ภิกฺขเว ภินฺเน อญฺญมญฺญํ ภณฺฑนานิ เจว โหนฺติ, อญฺญมญฺญํ ปริภาสา จ โหนฺติ, อญฺญมญฺญํ ปริกฺเขปา จ โหนฺติ, อญฺญมญฺญํ ปริจฺจชนา จ โหนฺติ.\csromanspacer{} ตตฺถ อปฺปสนฺนา เจว นปฺปสีทนฺติ, ปสนฺนานญฺจ เอกจฺจานํ อญฺญถตฺตํ โหตีติ.\csromanspacer{} เอตมตฺถํ ภควา อโวจ, ตตฺเถตํ อิติ วุจฺจติ\csromandash{}}

\csromanheader{2}{1. เอกกปิปาต}
\csromangathameasure{“อาปายิโก เนรยิโก, กปฺปฏฺโฐ สํฆเภทโก. \\ วคฺคาราโม อธมฺมฏฺโฐ, โยคกฺเขมา ปธํสติ. \\ สํฆํ สมคฺคํ เภตฺวาน, กปฺปํ นิรยมฺหิ ปจฺจตี”ติ.}
\csromangathagroup{\csromangathastanza{\csromanfolio{203}“อาปายิโก เนรยิโก, กปฺปฏฺโฐ สํฆเภทโก. \\ วคฺคาราโม อธมฺมฏฺโฐ, โยคกฺเขมา ปธํสติ\footnote{โยคกฺเขมโต ธํสติ (สฺยา, อิ), โยคกฺเขมา วิธํสติ (สี, ก)}. \\ สํฆํ สมคฺคํ เภตฺวาน\footnote{โยคกฺเขมโต ธํสติ (สฺยา, อิ), โยคกฺเขมา วิธํสติ (สี, ก)}, กปฺปํ นิรยมฺหิ ปจฺจตี”ติ.}}

\prose{อยมฺปิ อตฺโถ วุตฺโต ภควตา อิติ เม สุตนฺติ.\csromanspacer{}\csromanspacer{}อฏฺฐมํ.}

\csromanheader{2}{9. สํฆสามคฺคีสุตฺต}
\csromanheader{2}{19. วุตฺตํ เหตํ ภควตา วุตฺตมรหตาติ เม สุตํ\csromandash{}}
\prose{เอกธมฺโม ภิกฺขเว โลเก อุปฺปชฺชมาโน อุปฺปชฺชติ พหุชนหิตาย พหุชนสุขาย พหุโน ชนสฺส อตฺถาย หิตาย สุขาย เทวมนุสฺสานํ.\csromanspacer{} กตโม เอกธมฺโม? สํฆสามคฺคี.\csromanspacer{} สํเฆ โข ปน ภิกฺขเว สมคฺเค น เจว อญฺญมญฺญํ ภณฺฑนานิ โหนฺติ, น จ อญฺญมญฺญํ ปริภาสา โหนฺติ, น จ อญฺญมญฺญํ ปริกฺเขปา โหนฺติ, น จ อญฺญมญฺญํ ปริจฺจชนา โหนฺติ.\csromanspacer{} ตตฺถ อปฺปสนฺนา เจว ปสีทนฺติ, ปสนฺนานญฺจ ภิยฺโยภาโว โหตีติ.\csromanspacer{} เอตมตฺถํ ภควา อโวจ, ตตฺเถตํ อิติ วุจฺจติ\csromandash{}}

\csromangathameasure{*“สุขา สํฆสฺส สามคฺคี, สมคฺคานญฺจนุคฺคโห. \\ สมคฺครโต ธมฺมฏฺโฐ, โยคกฺเขมา น ธํสติ. \\ สํฆํ สมคฺคํ กตฺวาน, กปฺปํ สคฺคมฺหิ โมทตี”ติ.}
\csromangathagroup{\csromangathastanza{\csromansymbolfootnote{*}{อํ 1. 8; ขุ 10.109 ปิฏฺเฐสุปิ.}“สุขา สํฆสฺส สามคฺคี, สมคฺคานญฺจนุคฺคโห. \\ สมคฺครโต ธมฺมฏฺโฐ, โยคกฺเขมา น ธํสติ. \\ สํฆํ สมคฺคํ กตฺวาน, กปฺปํ สคฺคมฺหิ โมทตี”ติ.}}

\prose{อยมฺปิ อตฺโถ วุตฺโต ภควตา อิติ เม สุตนฺติ.\csromanspacer{}\csromanspacer{}นวมํ.}

\csromanheader{2}{10. ปทุฏฺฐจิตฺตสุตฺต}
\csromanheader{2}{20. วุตฺตํ เหตํ ภควตา วุตฺตมรหตาติ เม สุตํ\csromandash{}}
\prose{\csromansymbolmark{*}อิธาหํ ภิกฺขเว เอกจฺจํ ปุคฺคลํ ปทุฏฺฐจิตฺตํ เอวํ เจตสา เจโต ปริจฺจ ปชานามิ “อิมมฺหิ จายํ สมเย ปุคฺคโล กาลํ กเรยฺย, ยถาภตํ นิกฺขิตฺโต เอวํ นิรเย”.\csromanspacer{} ตํ กิสฺส เหตุ? จิตฺตํ หิสฺส ภิกฺขเว ปทุฏฺฐํ.\csromanspacer{} เจโตปโทสเหตุ โข ปน ภิกฺขเว เอวมิเธกจฺเจ สตฺตา กายสฺส เภทา ปรํ มรณา อปายํ ทุคฺคติํ วินิปาตํ นิรยํ อุปปชฺชนฺตีติ.\csromanspacer{} เอตมตฺถํ ภควา อโวจ, ตตฺเถตํ อิติ วุจฺจติ\csromandash{}}

\gambhira{อิติวุตฺตกปาฬิ}
\csromangathameasure{“ปทุฏฺฐจิตฺตํ ญตฺวาน, เอกจฺจํ อิธ ปุคฺคลํ. \\ เอตมตฺถญฺจ พฺยากาสิ, พุทฺโธ ภิกฺขูน สนฺติเก. \\ อิมมฺหิ จายํ สมเย, กาลํ กยิราถ ปุคฺคโล. \\ นิรยํ อุปปชฺเชยฺย, จิตฺตํ หิสฺส ปทูสิตํ. \\ ยถา หริตฺวา นิกฺขิเปยฺย, เอวเมวํ ตถาวิโธ. \\ เจโตปโทสเหตุ หิ, สตฺตา คจฺฉนฺติ ทุคฺคตินฺ”ติ.}
\csromangathagroup{\csromangathastanza{\csromanfolio{204}“ปทุฏฺฐจิตฺตํ ญตฺวาน, เอกจฺจํ อิธ ปุคฺคลํ. \\ เอตมตฺถญฺจ พฺยากาสิ, พุทฺโธ ภิกฺขูน สนฺติเก.}\csromangathastanza{อิมมฺหิ จายํ สมเย, กาลํ กยิราถ ปุคฺคโล. \\ นิรยํ อุปปชฺเชยฺย, จิตฺตํ หิสฺส ปทูสิตํ.}\csromangathastanza{ยถา หริตฺวา นิกฺขิเปยฺย, เอวเมวํ ตถาวิโธ. \\ เจโตปโทสเหตุ หิ, สตฺตา คจฺฉนฺติ ทุคฺคตินฺ”ติ.}}

\prose{อยมฺปิ อตฺโถ วุตฺโต ภควตา อิติ เม สุตนฺติ.\csromanspacer{}\csromanspacer{}ทสมํ.\csromanspacer{} ทุติโย วคฺโค.\csromanspacer{} \textbf{ตสฺสุทฺทานํ}}

\csromangathameasure{โมโห โกโธ อถ มกฺโข, วิชฺชา ตณฺหา เสขทุเว จ. \\ เภโท สามคฺคิปุคฺคโล, วคฺคมาหุ ทุติยนฺติ วุจฺจตีติ.}
\csromangathagroup{\csromangathastanza{โมโห โกโธ อถ มกฺโข, วิชฺชา ตณฺหา เสขทุเว จ. \\ เภโท สามคฺคิปุคฺคโล\footnote{โมหโกธ อถ มกฺขาคโต, มูหา กามเสกฺขทุเว. เภทสามคฺคปุคฺคโล จ (สี, ก) โมหโกธา อถ มกฺโข, โมหกามา เสกฺขา ทุเว. เภทโมทา ปุคฺคโลจ (สฺยา, อิ)}, วคฺคมาหุ ทุติยนฺติ วุจฺจตีติ.}}

\csromanheader{1}{3. ตติยวคฺค}
\csromanheader{2}{1. ปสนฺนจิตฺตสุตฺต}
\csromanheader{2}{21. วุตฺตํ เหตํ ภควตา วุตฺตมรหตาติ เม สุตํ\csromandash{}}
\prose{\csromansymbolfootnote{*}{อํ 1. 8; ขุ 10. 118 ปิฏฺเฐสุปิ.}อิธาหํ ภิกฺขเว เอกจฺจํ ปุคฺคลํ ปสนฺนจิตฺตํ เอวํ เจตสา เจโต ปริจฺจ ปชานามิ “อิมมฺหิ จายํ สมเย ปุคฺคโล กาลํ กเรยฺย, ยถาภตํ นิกฺขิตฺโต เอวํ สคฺเค”.\csromanspacer{} ตํ กิสฺส เหตุ? จิตฺตํ หิสฺส ภิกฺขเว ปสนฺนํ, เจโตปสาทเหตุ โข ปน ภิกฺขเว เอวมิเธกจฺเจ สตฺตา กายสฺส เภทา ปรํ มรณา สุคติํ สคฺคํ โลกํ อุปปชฺชนฺตีติ.\csromanspacer{} เอตมตฺถํ ภควา อโวจ, ตตฺเถตํ อิติ วุจฺจติ\csromandash{}}

\csromangathameasure{“ปสนฺนจิตฺตํ ญตฺวาน, เอกจฺจํ อิธปุคฺคลํ. \\ เอตมตฺถญฺจ พฺยากาสิ, พุทฺโธ ภิกฺขูน สนฺติเก.}
\csromangathagroup{\csromangathastanza{“ปสนฺนจิตฺตํ ญตฺวาน, เอกจฺจํ อิธปุคฺคลํ. \\ เอตมตฺถญฺจ พฺยากาสิ, พุทฺโธ ภิกฺขูน สนฺติเก.}}

\csromanheader{2}{1. เอกกปิปาต}
\csromanmatikaanchor{mtk.126}
\csromantocmarkref{subsection}{10. ทุติยทพฺพสุตฺต}{mtk.126}
\bookmark[dest={mtk.126},level=2]{10. ทุติยทพฺพสุตฺต}
\csromangathameasure{อิมมฺหิ จายํ สมเย, กาลํ กยิราถ ปุคฺคโล. \\ สุคติํ อุปปชฺเชยฺย, จิตฺตํ หิสฺส ปสาทิตํ. \\ ยถา หริตฺวา นิกฺขิเปยฺย, เอวเมวํ ตถาวิโธ. \\ เจโตปสาทเหตุ หิ, สตฺตา คจฺฉนฺติ สุคฺคตินฺ”ติ.}
\csromangathagroup{\csromangathastanza{\csromanfolio{205}อิมมฺหิ จายํ สมเย, กาลํ กยิราถ ปุคฺคโล. \\ สุคติํ อุปปชฺเชยฺย, จิตฺตํ หิสฺส ปสาทิตํ.}\csromangathastanza{ยถา หริตฺวา นิกฺขิเปยฺย, เอวเมวํ ตถาวิโธ. \\ เจโตปสาทเหตุ หิ, สตฺตา คจฺฉนฺติ สุคฺคตินฺ”ติ.}}

\prose{อยมฺปิ อตฺโถ วุตฺโต ภควตา อิติ เม สุตนฺติ.\csromanspacer{}\csromanspacer{}ปฐมํ.}

\csromanheader{2}{2. เมตฺตสุตฺต}
\csromanheader{2}{22. วุตฺตํ เหตํ ภควตา วุตฺตมรหตาติ เม สุตํ\csromandash{}}
\prose{มา ภิกฺขเว ปุญฺญานํ ภายิตฺถ, สุขสฺเสตํ ภิกฺขเว อธิวจนํ อิฏฺฐสฺส กนฺตสฺส ปิยสฺส มนาปสฺส ยทิทํ ปุญฺญานิ\footnote{ปุญฺญานนฺติ, (อํ 2. 465)}.\csromanspacer{} อภิชานามิ โข ปนาหํ ภิกฺขเว ทีฆรตฺตํ กตานํ ปุญฺญานํ อิฏฺฐํ กนฺตํ ปิยํ มนาปํ วิปากํ ปจฺจนุภูตํ, สตฺต วสฺสานิ เมตฺตจิตฺตํ ภาเวตฺวา สตฺต สํวฏฺฏวิวฏฺฏกปฺเป นยิมํ โลกํ ปุนราคมาสิํ.\csromanspacer{} สํวฏฺฏมาเน สุทํ ภิกฺขเว กปฺเป อาภสฺสรูปโค โหมิ, วิวฏฺฏมาเน กปฺเป สุญฺญํ พฺรหฺมวิมานํ อุปปชฺชามิ.}

\prose{ตตฺร สุทํ ภิกฺขเว พฺรหฺมา โหมิ มหาพฺรหฺมา อภิภู อนภิภูโต อญฺญทตฺถุทโส วสวตฺตี.\csromanspacer{} ฉตฺติํสกฺขตฺตุํ โข ปนาหํ ภิกฺขเว สกฺโก อโหสิํ เทวานมินฺโท, อเนกสตกฺขตฺตุํ ราชา อโหสิํ จกฺกวตฺตี ธมฺมิโก ธมฺมราชา จาตุรนฺโต วิชิตาวี ชนปทตฺถาวริยปฺปตฺโต สตฺตรตนสมนฺนาคโต, โก ปน วาโท ปเทสรชฺชสฺส.}

\prose{ตสฺส มยฺหํ ภิกฺขเว เอตทโหสิ “กิสฺส นุ โข เม อิทํ กมฺมสฺส ผลํ กิสฺส กมฺมสฺส วิปาโก, เยนาหํ เอตรหิ เอวํมหิทฺธิโก เอวํมหานุภาโว”ติ.\csromanspacer{} ตสฺส มยฺหํ ภิกฺขเว เอตทโหสิ “ติณฺณํ โข เม อิทํ กมฺมานํ ผลํ ติณฺณํ กมฺมานํ วิปาโก, เยนาหํ เอตรหิ เอวํมหิทฺธิโก เอวํมหานุภาโวติ.\csromanspacer{} เสยฺยถิทํ\footnote{เสยฺยถีทํ (สี, สฺยา, กํ, อิ)}, ทานสฺส ทมสฺส สํยมสฺสาติ.\csromanspacer{} เอตมตฺถํ ภควา อโวจ, ตตฺเถตํ อิติ วุจฺจติ\csromandash{}}

\csromangathameasure{“ปุญฺญเมว โส สิกฺเขยฺย, อายตคฺคํ สุขุทฺรยํ. \\ ทานญฺจ สมจริยญฺจ, เมตฺตจิตฺตญฺจ ภาวเย.}
\csromangathagroup{\csromangathastanza{“ปุญฺญเมว โส สิกฺเขยฺย, อายตคฺคํ สุขุทฺรยํ. \\ ทานญฺจ สมจริยญฺจ, เมตฺตจิตฺตญฺจ ภาวเย.}}

\gambhira{อิติวุตฺตกปาฬิ}
\csromangathameasure{เอเต ธมฺเม ภาวยิตฺวา, ตโย สุขสมุทฺทเย. \\ อพฺยาปชฺชํ สุขํ โลกํ, ปณฺฑิโต อุปปชฺชตี”ติ.}
\csromangathagroup{\csromangathastanza{\csromanfolio{206}เอเต ธมฺเม ภาวยิตฺวา, ตโย สุขสมุทฺทเย\footnote{สุขสมุทฺรเย (สี-ฏฺฐ)}. \\ อพฺยาปชฺชํ\footnote{สุขสมุทฺรเย (สี-ฏฺฐ)} สุขํ โลกํ, ปณฺฑิโต อุปปชฺชตี”ติ.}}

\prose{อยมฺปิ อตฺโต วุตฺโถ ภควตา อิติ เม สุตนฺติ.\csromanspacer{}\csromanspacer{}ทุติยํ.}

\csromanheader{2}{3. อุภยตฺถสุตฺต}
\csromanheader{2}{23. วุตฺตํ เหตํ ภควตา วุตฺตมรหตาติ เม สุตํ\csromandash{}}
\prose{เอกธมฺโม ภิกฺขเว ภาวิโต พหุลีกโต อุโภ อตฺเถ สมธิคยฺห ติฏฺฐติ ทิฏฺฐธมฺมิกญฺเจว อตฺถํ สมฺปรายิกญฺจ.\csromanspacer{} กตโม เอกธมฺโม? อปฺปมาโท กุสเลสุ ธมฺเมสุ.\csromanspacer{} อยํ โข ภิกฺขเว เอกธมฺโม ภาวิโต พหุลีกโต อุโภ อตฺเถ สมธิคยฺห ติฏฺฐติ ทิฏฺฐธมฺมิกญฺเจว อตฺถํ สมฺปรายิกญฺจาติ.\csromanspacer{} เอตมตฺถํ ภควา อโวจ, ตตฺเถตํ อิติ วุจฺจติ\csromandash{}}

\csromangathameasure{*“อปฺปมาทํ ปสํสนฺติ, ปุญฺญกิริยาสุ ปณฺฑิตา. \\ อปฺปมตฺโต อุโภ อตฺเถ, อธิคณฺหาติ ปณฺฑิโต. \\ *ทิฏฺเฐ ธมฺเม จ โย อตฺโถ, โย จตฺโถ สมฺปรายิโก. \\ อตฺถาภิสมยา ธีโร, ‘ปณฺฑิโต’ติ ปวุจฺจตี”ติ.}
\csromangathagroup{\csromangathastanza{\csromansymbolfootnote{*}{สํ 1. 87, 90; อํ 2. 42 ปิฏฺเฐสุปิ.}“อปฺปมาทํ ปสํสนฺติ, ปุญฺญกิริยาสุ ปณฺฑิตา. \\ อปฺปมตฺโต อุโภ อตฺเถ, อธิคณฺหาติ ปณฺฑิโต.}\csromangathastanza{\csromansymbolmark{*}ทิฏฺเฐ ธมฺเม จ โย อตฺโถ, โย จตฺโถ สมฺปรายิโก. \\ อตฺถาภิสมยา ธีโร, ‘ปณฺฑิโต’ติ ปวุจฺจตี”ติ.}}

\prose{อยมฺปิ อตฺโถ วุตฺโต ภควตา อิติ เม สุตนฺติ.\csromanspacer{} ตติยํ.}

\csromanmatikaanchor{mtk.127}
\csromantocmarknopage{chapter}{อิติวุตฺตกปาฬิ}{mtk.127}
\bookmark[dest={mtk.127},level=0]{อิติวุตฺตกปาฬิ}
\csromanheader{2}{4. อฏฺฐิปุญฺชสุตฺต}
\csromanheader{2}{24. วุตฺตํ เหตํ ภควตา วุตฺตมรหตาติ เม สุตํ\csromandash{}}
\csromanmatikaanchor{mtk.128}
\csromantocmarknopage{chapter}{1. เอกกนิปาต}{mtk.128}
\bookmark[dest={mtk.128},level=0]{1. เอกกนิปาต}
\prose{เอกปุคฺคลสฺส ภิกฺขเว กปฺปํ สนฺธาวโต สํสรโต สิยา เอวํ มหา อฏฺฐิกงฺกโล อฏฺฐิปุญฺโช อฏฺฐิราสิ, ยถายํ เวปุลฺโล ปพฺพโต, สเจ สํหารโก อสฺส, สมฺภตญฺจ น วินสฺเสยฺยาติ.\csromanspacer{} เอตมตฺถํ ภควา อโวจ ตตฺเถตํ อิติ วุจฺจติ\csromandash{}}

\csromanmatikaanchor{mtk.129}
\csromantocmarknopage{section}{1. ปฐมวคฺค}{mtk.129}
\bookmark[dest={mtk.129},level=1]{1. ปฐมวคฺค}
\csromantocmarkref{subsection}{1-6. โลภสุตฺตาทิ}{mtk.130}
\bookmark[dest={mtk.130},level=2]{1-6. โลภสุตฺตาทิ}
\csromantocmarkref{subsection}{7-10. สพฺพปริญฺญาสุตฺตาทิ}{mtk.131}
\bookmark[dest={mtk.131},level=2]{7-10. สพฺพปริญฺญาสุตฺตาทิ}
\csromangathameasure{“เอกสฺเสเกน กปฺเปน, ปุคฺคลสฏฺฐิสญฺชโย. \\ สิยา สพฺพตสโม ราสิ, อิติ วุตฺตํ มเหสินา. \\ โส โข ปนายํ อกฺขาโต, เวปุลฺโล ปพฺพโต มหา. \\ อุตฺตโร คิชฺฌกูฏสฺส, มคธานํ คิริพฺพเช.}
\csromangathagroup{\csromangathastanza{“เอกสฺเสเกน กปฺเปน, ปุคฺคลสฏฺฐิสญฺชโย. \\ สิยา สพฺพตสโม ราสิ, อิติ วุตฺตํ มเหสินา.}\csromangathastanza{โส โข ปนายํ อกฺขาโต, เวปุลฺโล ปพฺพโต มหา. \\ อุตฺตโร คิชฺฌกูฏสฺส, มคธานํ คิริพฺพเช.}}

\csromanheader{2}{1. เอกกปิปาต}
\csromangathameasure{ยโต จ อริยสจฺจานิ, สมฺมปฺปญฺญาย ปสฺสติ. \\ ทุกฺขํ ทุกฺขสมุปฺปาทํ, ทุกฺขสฺส จ อติกฺกมํ. \\ อริยญฺจฏฺฐงฺคิกํ มคฺคํ, ทุกฺขูปสมคามินํ. \\ ส สตฺตกฺขตฺตุํปรมํ, สนฺธาวิตฺวาน ปุคฺคโล. \\ ทุกฺขสฺสนฺตกโร โหติ, สพฺพสํโยชนกฺขยา”ติ.}
\csromangathagroup{\csromangathastanza{\csromanfolio{207}ยโต จ อริยสจฺจานิ, สมฺมปฺปญฺญาย ปสฺสติ. \\ ทุกฺขํ ทุกฺขสมุปฺปาทํ, ทุกฺขสฺส จ อติกฺกมํ.}\csromangathastanza{อริยญฺจฏฺฐงฺคิกํ มคฺคํ, ทุกฺขูปสมคามินํ. \\ ส สตฺตกฺขตฺตุํปรมํ, สนฺธาวิตฺวาน ปุคฺคโล. \\ ทุกฺขสฺสนฺตกโร โหติ, สพฺพสํโยชนกฺขยา”ติ.}}

\prose{อยมฺปิ อตฺโถ วุตฺโต ภควตา อิติ เม สุตนฺติ.\csromanspacer{}\csromanspacer{}จตุตฺถํ.}

\csromanheader{2}{5. มุสาวาทสุตฺต}
\csromanheader{2}{25. วุตฺตํ เหตํ ภควตา วุตฺตมรหตาติ เม สุตํ\csromandash{}}
\prose{เอกธมฺมํ อตีตสฺส ภิกฺขเว ปุริสปุคฺคลสฺส นาหํ ตสฺส กิญฺจิ ปาปกมฺมํ อกรณียนฺติ วทามิ.\csromanspacer{} กตมํ เอกธมฺมํ? ยทิทํ\footnote{ยถยิทํ (สี, สฺยา, ก), ยถายิทํ (อิ)} ภิกฺขเว สมฺปชานมุสาวาโทติ.\csromanspacer{} เอตมตฺถํ ภควา อโวจ, ตตฺเถตํ อิติ วุจฺจติ\csromandash{}}

\csromangathameasure{“เอกธมฺมํ อตีตสฺส, มุสาวาทิสฺส ชนฺตุโน. \\ วิติณฺณปรโลกสฺส, นตฺถิ ปาปํ อการิยนฺ”ติ.}
\csromangathagroup{\csromangathastanza{“เอกธมฺมํ อตีตสฺส, มุสาวาทิสฺส ชนฺตุโน. \\ วิติณฺณปรโลกสฺส, นตฺถิ ปาปํ อการิยนฺ”ติ.}}

\prose{อยมฺปิ อตฺโถ วุตฺโต ภควตา อิติ เม สุตนฺติ.\csromanspacer{}\csromanspacer{}ปญฺจมํ.}

\csromanheader{2}{6. ทานสุตฺต}
\csromanheader{2}{26. วุตฺตํ เหตํ ภควตา วุตฺตมรหตาติ เม สุตํ\csromandash{}}
\prose{เอวญฺเจ ภิกฺขเว สตฺตา ชาเนยฺยุํ ทานสํวิภาคสฺส วิปากํ, ยถาหํ ชานามิ, น อทตฺวา ภุญฺเชยฺยุํ, น จ เนสํ มจฺเฉรมลํ จิตฺตํ ปริยาทาย ติฏฺเฐยฺย.\csromanspacer{} โยปิ เนสํ อสฺส จริโม อาโลโป จริมํ กพฬํ, ตโตปิ น อสํวิภชิตฺวา ภุญฺเชยฺยุํ, สเจ เนสํ ปฏิคฺคาหกา อสฺสุ.\csromanspacer{} ยสฺมา จ โข ภิกฺขเว สตฺตา น เอวํ ชานนฺติ ทานสํวิภาคสฺส วิปากํ ยถาหํ ชานามิ, ตสฺมา อทตฺวา ภุญฺชนฺติ, มจฺเฉรมลญฺจ เนสํ จิตฺตํ ปริยาทาย ติฏฺฐตีติ.\csromanspacer{} เอตมตฺถํ ภควา อโวจ, ตตฺเถตํ อิติ วุจฺจติ\csromandash{}}

\csromangathameasure{“เอวญฺเจ สตฺตา ชาเนยฺยุํ, ยถาวุตฺตํ มเหสินา. \\ วิปากํ สํวิภาคสฺส, ยถา โหติ มหปฺผลํ.}
\csromangathagroup{\csromangathastanza{“เอวญฺเจ สตฺตา ชาเนยฺยุํ, ยถาวุตฺตํ มเหสินา. \\ วิปากํ สํวิภาคสฺส, ยถา โหติ มหปฺผลํ.}}

\gambhira{อิติวุตฺตกปาฬิ}
\csromangathameasure{วิเนยฺย มจฺเฉรมลํ, วิปฺปสนฺเนน เจตสา. \\ ทชฺชุํ กาเลน อริเยสุ, ยตฺถ ทินฺนํ มหปฺผลํ. \\ อนฺนญฺจ ทตฺวา พหุโน, ทกฺขิเณยฺเยสุ ทกฺขิณํ. \\ อิโต จุตา มนุสฺสตฺตา, สคฺคํ คจฺฉนฺติ ทายกา. \\ เต จ สคฺคคตา ตตฺถ, โมทนฺติ กามกามิโน. \\ วิปากํ สํวิภาคสฺส, อนุโภนฺติ อมจฺฉรา”ติ.}
\csromangathagroup{\csromangathastanza{\csromanfolio{208}วิเนยฺย มจฺเฉรมลํ, วิปฺปสนฺเนน เจตสา. \\ ทชฺชุํ กาเลน อริเยสุ, ยตฺถ ทินฺนํ มหปฺผลํ.}\csromangathastanza{อนฺนญฺจ ทตฺวา\footnote{ทตฺวาน (สฺยา)} พหุโน, ทกฺขิเณยฺเยสุ ทกฺขิณํ. \\ อิโต จุตา มนุสฺสตฺตา, สคฺคํ คจฺฉนฺติ ทายกา.}\csromangathastanza{เต จ สคฺคคตา\footnote{สคฺคํ คตา (สี, อิ, ก)} ตตฺถ, โมทนฺติ กามกามิโน. \\ วิปากํ สํวิภาคสฺส, อนุโภนฺติ อมจฺฉรา”ติ.}}

\prose{อยมฺปิ อตฺโถ วุตฺโต ภควตา อิติ เม สุตนฺติ.\csromanspacer{}\csromanspacer{}ฉฏฺฐํ.}

\csromanheader{2}{7. เมตฺตาภาวนาสุตฺต}
\csromanheader{2}{27. วุตฺตํ เหตํ ภควตา วุตฺตมรหตาติ เม สุตํ\csromandash{}}
\prose{ยานิ กานิจิ ภิกฺขเว โอปธิกานิ ปุญฺญกิริยวตฺถูนิ, สพฺพานิ ตานิ เมตฺตาย เจโตวิมุตฺติยา กลํ นาคฺฆนฺติ โสฬสิํ, เมตฺตาเยว ตานิ เจโตวิมุตฺติ อธิคฺคเหตฺวา ภาสเต จ ตปเต จ วิโรจติ จ.}

\prose{เสยฺยถาปิ ภิกฺขเว ยา กาจิ ตารกรูปานํ ปภา, สพฺพา ตา จนฺทิยา ปภาย กลํ นาคฺฆนฺติ โสฬสิํ, จนฺทปภาเยว ตา อธิคฺคเหตฺวา ภาสเต จ ตปเต จ วิโรจติ จ.\csromanspacer{} เอวเมวํ โข ภิกฺขเว ยานิ กานิจิ โอปธิกานิ ปุญฺญกิริยวตฺถูนิ, สพฺพานิ ตานิ เมตฺตาย เจโตวิมุตฺติยา กลํ นาคฺฆนฺติ โสฬสิํ, เมตฺตาเยว ตานิ เจโตวิมุตฺติ อธิคฺคเหตฺวา ภาสเต จ ตปเต จ วิโรจติ จ.}

\prose{เสยฺยถาปิ ภิกฺขเว วสฺสานํ ปจฺฉิเม มาเส สรทสมเย วิทฺเธ วิคตวลาหเก เทเว\footnote{นเภ (สี)} อาทิจฺโจ นภํ อพฺภุสฺสกฺกมาโน\footnote{อพฺภุคฺคมมาโน (ก-ฏฺฐ)} สพฺพํ อากาสคตํ\footnote{อากาสํ (สฺยา)} ตมคตํ อภิวิหจฺจ\footnote{อภิหจฺจ (สฺยา)} ภาสเต จ ตปเต จ วิโรจติ จ.\csromanspacer{} เอวเมวํ โข ภิกฺขเว ยานิ กานิจิ โอปธิกานิ ปุญฺญกิริยวตฺถูนิ, สพฺพานิ ตานิ เมตฺตาย เจโตวิมุตฺติยา กลํ นาคฺฆนฺติ โสฬสิํ, เมตฺตาเยว ตานิ เจโตวิมุตฺติ อธิคฺคเหตฺวา ภาสเต จ ตปเต จ วิโรจติ จ.}

\csromanheader{2}{1. เอกกปิปาต}
\prose{\csromanfolio{209}เสยฺยถาปิ ภิกฺขเว รตฺติยา ปจฺจูสสมยํ โอสธิตารกา ภาสเต จ ตปเต จ วิโรจติ จ.\csromanspacer{} เอวเมวํ โข ภิกฺขเว ยานิ กานิจิ โอปธิกานิ ปุญฺญกิริยวตฺถูนิ, สพฺพานิ ตานิ เมตฺตาย เจโตวิมุตฺติยา กลํ นาคฺฆนฺติ โสฬสิํ, เมตฺตาเยว ตานิ เจโตวิมุตฺติ อธิคฺคเหตฺวา ภาสเต จ ตปเต จ วิโรจติ จาติ.\csromanspacer{} เอตมตฺถํ ภควา อโวจ, ตตฺเถตํ อิติ วุจฺจติ\csromandash{}}

\csromangathameasure{“โย จ เมตฺตํ ภาวยติ, อปฺปมาณํ ปฏิสฺสโต. \\ ตนู สํโยชนา โหนฺติ, ปสฺสโต อุปธิกฺขยํ. \\ เอกมฺปิ เจ ปาณมทุฏฺฐจิตฺโต, \\ เมตฺตายติ กุสโล เตน โหติ. \\ สพฺเพ จ ปาเณ มนสานุกมฺปํ, \\ ปหูตมริโย ปกโรติ ปุญฺญํ. \\ เย สตฺตสณฺฑํ ปถวิํ วิชิตฺวา, \\ ราชิสโย ยชมานา’นุปริยคา. \\ อสฺสเมธํ ปุริสเมธํ, \\ สมฺมาปาสํ วาชเปยฺยํ นิรคฺคฬํ. \\ เมตฺตสฺส จิตฺตสฺส สุภาวิตสฺส, \\ กลมฺปิ เต นานุภวนฺติ โสฬสิํ, \\ จนฺทปฺปภา ตารคณาว สพฺเพ. \\ โย น หนฺติ น ฆาเตติ, \\ น ชินาติ น ชาปเย. เมตฺตํโส สพฺพภูเตสุ, \\ เวรํ ตสฺส น เกนจี”ติ.}
\csromangathagroup{\csromangathastanza{“โย จ เมตฺตํ ภาวยติ, อปฺปมาณํ ปฏิสฺสโต. \\ ตนู\footnote{ตนุ (สี)} สํโยชนา โหนฺติ, ปสฺสโต อุปธิกฺขยํ.}\csromangathastanza{เอกมฺปิ เจ ปาณมทุฏฺฐจิตฺโต, \\ เมตฺตายติ กุสโล เตน โหติ. \\ สพฺเพ จ ปาเณ มนสานุกมฺปํ, \\ ปหูตมริโย ปกโรติ ปุญฺญํ.}\csromangathastanza{เย\footnote{โย (สี)} สตฺตสณฺฑํ ปถวิํ วิชิตฺวา, \\ ราชิสโย\footnote{โย (สี)} ยชมานา’นุปริยคา. \\ อสฺสเมธํ ปุริสเมธํ, \\ สมฺมาปาสํ วาชเปยฺยํ นิรคฺคฬํ.}\csromangathastanza{เมตฺตสฺส จิตฺตสฺส สุภาวิตสฺส, \\ กลมฺปิ เต นานุภวนฺติ โสฬสิํ, \\ จนฺทปฺปภา ตารคณาว สพฺเพ. \\ โย น หนฺติ น ฆาเตติ, \\ น ชินาติ น ชาปเย. เมตฺตํโส สพฺพภูเตสุ, \\ เวรํ ตสฺส น เกนจี”ติ.}}

\prose{อยมฺปิ อตฺโถ วุตฺโต ภควตา อิติ เม สุตนฺติ.\csromanspacer{}\csromanspacer{}สตฺตมํ.\csromanspacer{} ตติโย วคฺโค.}

\gambhira{อิติวุตฺตกปาฬิ \textbf{ตสฺสุทฺทานํ}}
\csromangathameasure{จิตฺตํ เมตฺตํ อุโภ อตฺเถ, ปุญฺชํ เวปุลฺลปพฺพตํ. \\ สมฺปชานมุสาวาโท, ทานญฺจ เมตฺตภาวนา. \\ สตฺติมานิ จ3 สุตฺตานิ, ปุริมานิ จ วีสติ. \\ เอกธมฺเมสุ สุตฺตนฺตา, สตฺตวีสติสงฺคหาติ.}
\csromangathagroup{\csromangathastanza{\csromanfolio{210}จิตฺตํ เมตฺตํ\footnote{ฌายี (สี, สฺยา), ฌายิ (อิ, ก)} อุโภ อตฺเถ, ปุญฺชํ เวปุลฺลปพฺพตํ. \\ สมฺปชานมุสาวาโท, ทานญฺจ เมตฺตภาวนา\footnote{ฌายี (สี, สฺยา), ฌายิ (อิ, ก)}.}\csromangathastanza{สตฺติมานิ จ3 สุตฺตานิ, ปุริมานิ จ วีสติ. \\ เอกธมฺเมสุ สุตฺตนฺตา, สตฺตวีสติสงฺคหาติ.}}

\nitthitam{\textbf{เอกกนิปาโต นิฏฺฐิโต.}}
\chapterheadpageread{2. ทุกนิปาต}
\chapterhead{1. ปฐมวคฺค}
\csromanheader{2}{1. ทุกฺขวิหารสุตฺต}
\csromanmatikaanchor{mtk.132}
\csromanmatikaanchor{mtk.147}
\csromanmatikaanchor{mtk.148}
\csromantocmarknopage{section}{2. ทุติยวคฺค}{mtk.132}
\bookmark[dest={mtk.132},level=1]{2. ทุติยวคฺค}
\csromantocmarkref{subsection}{1-3. โมหปริญฺญาสุตฺตาทิ}{mtk.133}
\bookmark[dest={mtk.133},level=2]{1-3. โมหปริญฺญาสุตฺตาทิ}
\csromantocmarkref{subsection}{4-5. อวิชฺชานีวรณสุตฺตาทิ}{mtk.134}
\bookmark[dest={mtk.134},level=2]{4-5. อวิชฺชานีวรณสุตฺตาทิ}
\csromantocmarkref{subsection}{6-7. เสขสุตฺตาทิ}{mtk.135}
\bookmark[dest={mtk.135},level=2]{6-7. เสขสุตฺตาทิ}
\csromantocmarkref{subsection}{8-9. สํฆเภทสุตฺตาทิ}{mtk.136}
\bookmark[dest={mtk.136},level=2]{8-9. สํฆเภทสุตฺตาทิ}
\proseitem{28}{\csromanfolio{211}( )\footnote{(เทฺว ธมฺเม อนุกฺกฏิ) กตฺถจิ.} วุตฺตํ เหตํ ภควตา วุตฺตมรหตาติ เม สุตํ\csromandash{}}

\prose{ทฺวีหิ ภิกฺขเว ธมฺเมหิ สมนฺนาคโต ภิกฺขุ ทิฏฺเฐว ธมฺเม ทุกฺขํ วิหรติ สวิฆาตํ ส-อุปายาสํ สปริฬาหํ, กายสฺส เภทา ปรํ มรณา ทุคฺคติ ปาฏิกงฺขา.\csromanspacer{} กตเมหิ ทฺวีหิ? อินฺทฺริเยสุ อคุตฺตทฺวารตาย\footnote{อคุตฺตทฺวาโร (ฏฺฐ)} จ, โภชเน อมตฺตญฺญุตาย\footnote{อมตฺตญฺญู (ฏฺฐ)} จ.\csromanspacer{} อิเมหิ โข ภิกฺขเว ทฺวีหิ ธมฺเมหิ สมนฺนาคโต ภิกฺขุ ทิฏฺเฐว ธมฺเม ทุกฺขํ วิหรติ สวิฆาตํ ส-อุปายาสํ สปริฬาหํ กายสฺส เภทา ปรํ มรณา ทุคฺคติ ปาฏิกงฺขาติ.\csromanspacer{} เอตมตฺถํ ภควา อโวจ, ตตฺเถตํ อิติ วุจฺจติ\csromandash{}}

\csromangathameasure{“จกฺขุ โสตญฺจ ฆานญฺจ, ชิวฺหา กาโย ตถา มโน. \\ เอตานิ ยสฺส ทฺวารานิ, อคุตฺตานิ จ ภิกฺขุโน. \\ โภชนมฺหิ อมตฺตญฺญู, อินฺทฺริเยสุ อสํวุโต. \\ กายทุกฺขํ เจโตทุกฺขํ, ทุกฺขํ โส อธิคจฺฉติ. \\ qอยฺหมาเนน กาเยน, ฑยฺหมาเนน เจตสา. \\ ทิวา วา ยทิ วา รตฺติํ, ทุกฺขํ วิหรติ ตาทิโส”ติ.}
\csromangathagroup{\csromangathastanza{“จกฺขุ โสตญฺจ ฆานญฺจ, ชิวฺหา กาโย ตถา มโน. \\ เอตานิ ยสฺส ทฺวารานิ, อคุตฺตานิ จ\footnote{อคุตฺตานิธ (ก)} ภิกฺขุโน.}\csromangathastanza{โภชนมฺหิ อมตฺตญฺญู, อินฺทฺริเยสุ อสํวุโต. \\ กายทุกฺขํ เจโตทุกฺขํ, ทุกฺขํ โส อธิคจฺฉติ.}\csromangathastanza{qอยฺหมาเนน กาเยน, ฑยฺหมาเนน เจตสา. \\ ทิวา วา ยทิ วา รตฺติํ, ทุกฺขํ วิหรติ ตาทิโส”ติ.}}

\prose{อยมฺปิ อตฺโถ วุตฺโต ภควตา อิติ เม สุตนฺติ.\csromanspacer{}\csromanspacer{}ปฐมํ.}

\csromanheader{2}{2. สุขวิหารสุตฺต}
\csromanheader{2}{29. วุตฺตํ เหตํ ภควตา วุตฺตมรหตาติ เม สุตํ\csromandash{}}
\prose{ทฺวีหิ ภิกฺขเว ธมฺเมหิ สมนฺนาคโต ภิกฺขุ ทิฏฺเฐว ธมฺเม สุขํ วิหรติ อวิฆาตํ อนุปายาสํ อปริฬาหํ, กายสฺส เภทา ปรํ มรณา สุคติ}

\vspace{\tipitakaparskip}
\csromancenter{อิติวุตฺตกปาฬิ ปาฏิกงฺขา.\csromanspacer{} กตเมหิ ทฺวีหิ? อินฺทฺริเยสุ คุตฺตทฺวารตาย จ, โภชเน มตฺตญฺญุตาย จ.\csromanspacer{} อิเมหิ โข ภิกฺขเว ทฺวีหิ ธมฺเมหิ สมนฺนาคโต ภิกฺขุ ทิฏฺเฐว ธมฺเม สุขํ วิหรติ อวิฆาตํ อนุปายาสํ อปริฬาหํ, กายสฺส เภทา ปรํ มรณา สุคติ ปาฏิกงฺขาติ.\csromanspacer{} เอตมตฺถํ ภควา อโวจ, ตตฺเถตํ อิติ วุจฺจติ\csromandash{}}
\vspace{0.5\baselineskip}
\csromanmatikaanchor{mtk.137}
\csromanmatikaanchor{mtk.149}
\csromantocmarkref{subsection}{10. ปทุฏฺฐจิตฺตสุตฺต}{mtk.137}
\bookmark[dest={mtk.137},level=2]{10. ปทุฏฺฐจิตฺตสุตฺต}
\csromangathameasure{“จกฺขุ โสตญฺจ ฆานญฺจ, ชิวฺหา กาโย ตถา มโน. \\ เอตานิ ยสฺส ทฺวารานิ, สุคุตฺตานิ จ ภิกฺขุโน. \\ โภชนมฺหิ จ มตฺตญฺญู, อินฺทฺริเยสุ จ สํวุโต. \\ กายสุขํ เจโตสุขํ, สุขํ โส อธิคจฺฉติ. \\ อฑยฺหมาเนน กาเยน, อฑยฺหมาเนน เจตสา. \\ ทิวา วา ยทิ วา รตฺติํ, สุขํ วิหรติ ตาทิโส”ติ.}
\csromangathagroup{\csromangathastanza{\csromanfolio{212}“จกฺขุ โสตญฺจ ฆานญฺจ, ชิวฺหา กาโย ตถา\footnote{อโถ (สี, สฺยา, ก)} มโน. \\ เอตานิ ยสฺส ทฺวารานิ, สุคุตฺตานิ จ ภิกฺขุโน.}\csromangathastanza{โภชนมฺหิ จ มตฺตญฺญู, อินฺทฺริเยสุ จ สํวุโต. \\ กายสุขํ เจโตสุขํ, สุขํ โส อธิคจฺฉติ.}\csromangathastanza{อฑยฺหมาเนน กาเยน, อฑยฺหมาเนน เจตสา. \\ ทิวา วา ยทิ วา รตฺติํ, สุขํ วิหรติ ตาทิโส”ติ.}}

\prose{อยมฺปิ อตฺโถ วุตฺโต ภควตา อิติ เม สุตนฺติ.\csromanspacer{}\csromanspacer{}ทุติยํ.}

\csromanheader{2}{3. ตปนียสุตฺต}
\csromanheader{2}{30. วุตฺตํ เหตํ ภควตา วุตฺตมรหตาติ เม สุตํ\csromandash{}}
\prose{เทฺวเม ภิกฺขเว ธมฺมา ตปนียา.\csromanspacer{} กตเม เทฺว? อิธ ภิกฺขเว เอกจฺโจ อกตกลฺยาโณ โหติ อกตกุสโล อกตภีรุตฺตาโณ กตปาโป กตลุทฺโท กตกิพฺพิโส, โส “อกตํ เม กลฺยาณนฺ”ติปิ ตปฺปติ, “กตํ เม ปาปนฺ”ติปิ ตปฺปติ.\csromanspacer{} อิเม โข ภิกฺขเว เทฺว ธมฺมา ตปนียาติ.\csromanspacer{} เอตมตฺถํ ภควา อโวจ, ตตฺเถตํ อิติ วุจฺจติ\csromandash{}}

\csromangathameasure{“กายทุจฺจริตํ กตฺวา, วจีทุจฺจริตานิ จ. \\ มโนทุจฺจริตํ กตฺวา, ยญฺจญฺญํ โทสสญฺหิตํ. \\ อกตฺวา กุสลํ กมฺมํ, กตฺวานา’กุสลํ พหุํ. \\ กายสฺส เภทา ทุปฺปญฺโญ, นิรยํ โสปปชฺชตี”ติ.}
\csromangathagroup{\csromangathastanza{“กายทุจฺจริตํ กตฺวา, วจีทุจฺจริตานิ จ. \\ มโนทุจฺจริตํ กตฺวา, ยญฺจญฺญํ โทสสญฺหิตํ.}\csromangathastanza{อกตฺวา กุสลํ กมฺมํ, กตฺวานา’กุสลํ พหุํ. \\ กายสฺส เภทา ทุปฺปญฺโญ, นิรยํ โสปปชฺชตี”ติ\footnote{นิรยํ โส อุปปชฺชตีติ (สี, สฺยา, กํ, อิ)}.}}

\prose{อยมฺปิ อตฺโถ วุตฺโต ภควตา อิติ เม สุตนฺติ.\csromanspacer{}\csromanspacer{}ตติยํ.}

\csromanheader{2}{2. ทุกนิปาต \textbf{4. อตปนียสุตฺต}}
\csromanheader{2}{31. วุตฺตํ เหตํ ภควตา วุตฺตมรหตาติ เม สุตํ\csromandash{}}
\csromanmatikaanchor{mtk.138}
\csromanmatikaanchor{mtk.150}
\csromantocmarknopage{section}{3. ตติยวคฺค}{mtk.138}
\bookmark[dest={mtk.138},level=1]{3. ตติยวคฺค}
\prose{\csromanfolio{213}เทฺวเม ภิกฺขเว ธมฺมา อตปนียา.\csromanspacer{} กตเม เทฺว? อิธ ภิกฺขเว เอกจฺโจ กตกลฺยาโณ โหติ กตกุสโล กตภีรุตฺตาโณ อกตปาโป อกตลุทฺโท อกตกิพฺพิโส, โส “กตํ เม กลฺยาณนฺ”ติปิ น ตปฺปติ, “อกตํ เม ปาปนฺ”ติปิ น ตปฺปติ.\csromanspacer{} อิเม โข ภิกฺขเว เทฺว ธมฺมา อตปนียาติ.\csromanspacer{} เอตมตฺถํ ภควา อโวจ, ตตฺเถตํ อิติ วุจฺจติ\csromandash{}}

\csromangathameasure{“กายทุจฺจริตํ หิตฺวา, วจีทุจฺจริตานิ จ. \\ มโนทุจฺจริตํ หิตฺวา, ยญฺจญฺญํ โทสสญฺหิตํ. \\ อกตฺวา’กุสลํ กมฺมํ, กตฺวาน กุสลํ พหุํ. \\ กายสฺส เภทา สปฺปญฺโญ, สคฺคํ โส อุปปชฺชตี”ติ.}
\csromangathagroup{\csromangathastanza{“กายทุจฺจริตํ หิตฺวา, วจีทุจฺจริตานิ จ. \\ มโนทุจฺจริตํ หิตฺวา, ยญฺจญฺญํ โทสสญฺหิตํ.}\csromangathastanza{อกตฺวา’กุสลํ กมฺมํ, กตฺวาน กุสลํ พหุํ. \\ กายสฺส เภทา สปฺปญฺโญ, สคฺคํ โส อุปปชฺชตี”ติ.}}

\prose{อยมฺปิ อตฺโถ วุตฺโต ภควตา อิติ เม สุตนฺติ.\csromanspacer{}\csromanspacer{}จตุตฺถํ.}

\csromanheader{2}{5. ปฐมสีลสุตฺต}
\csromanheader{2}{32. วุตฺตํ เหตํ ภควตา วุตฺตมรหตาติ เม สุตํ\csromandash{}}
\prose{ทฺวีหิ ภิกฺขเว ธมฺเมหิ สมนฺนาคโต ปุคฺคโล ยถาภตํ นิกฺขิตฺโต เอวํ นิรเย.\csromanspacer{} กตเมหิ ทฺวีหิ? ปาปเกน จ สีเลน, ปาปิกาย จ ทิฏฺฐิยา.\csromanspacer{} อิเมหิ โข ภิกฺขเว ทฺวีหิ ธมฺเมหิ สมนฺนาคโต ปุคฺคโล ยถาภตํ นิกฺขิตฺโต เอวํ นิรเยติ.\csromanspacer{} เอตมตฺถํ ภควา อโวจ, ตตฺเถตํ อิติ วุจฺจติ\csromandash{}}

\csromangathameasure{“ปาปเกน จ สีเลน, ปาปิกาย จ ทิฏฺฐิยา. \\ เอเตหิ ทฺวีหิ ธมฺเมหิ, โย สมนฺนาคโต นโร. \\ กายสฺส เภทา ทุปฺปญฺโญ, นิรยํ โสปปชฺชตี”ติ.}
\csromangathagroup{\csromangathastanza{“ปาปเกน จ สีเลน, ปาปิกาย จ ทิฏฺฐิยา. \\ เอเตหิ ทฺวีหิ ธมฺเมหิ, โย สมนฺนาคโต นโร. \\ กายสฺส เภทา ทุปฺปญฺโญ, นิรยํ โสปปชฺชตี”ติ.}}

\prose{อยมฺปิ อตฺโถ วุตฺโต ภควตา อิติ เม สุตนฺติ.\csromanspacer{}\csromanspacer{}ปญฺจมํ.}

\csromanheader{2}{6. ทุติยสีลสุตฺต}
\csromanheader{2}{33. วุตฺตํ เหตํ ภควตา วุตฺตมรหตาติ เม สุตํ\csromandash{}}
\prose{ทฺวีหิ ภิกฺขเว ธมฺเมหิ สมนฺนาคโต ปุคฺคโล ยถาภตํ นิกฺขิตฺโต เอวํ สคฺเค.\csromanspacer{} กตเมหิ ทฺวีหิ? ภทฺทเกน จ สีเลน, ภทฺทิกาย จ ทิฏฺฐิยา.}

\vspace{\tipitakaparskip}
\csromancenter{อิติวุตฺตกปาฬิ อิเมหิ โข ภิกฺขเว ทฺวีหิ ธมฺเมหิ สมนฺนาคโต ปุคฺคโล ยถาภตํ นิกฺขิตฺโต เอวํ สคฺเคติ.\csromanspacer{} เอตมตฺถํ ภควา อโวจ, ตตฺเถตํ อิติ วุจฺจติ\csromandash{}}
\vspace{0.5\baselineskip}
\csromanmatikaanchor{mtk.139}
\csromanmatikaanchor{mtk.152}
\csromantocmarkref{subsection}{1. ปสนฺนจิตฺตสุตฺต}{mtk.139}
\bookmark[dest={mtk.139},level=2]{1. ปสนฺนจิตฺตสุตฺต}
\csromangathameasure{“ภทฺทเกน จ สีเลน, ภทฺทิกาย จ ทิฏฺฐิยา. \\ เอเตหิ ทฺวีหิ ธมฺเมหิ, โย สมนฺนาคโต นโร. \\ กายสฺส เภทา สปฺปญฺโญ, สคฺคํ โส อุปปชฺชตี”ติ.}
\csromangathagroup{\csromangathastanza{\csromanfolio{214}“ภทฺทเกน จ สีเลน, ภทฺทิกาย จ ทิฏฺฐิยา. \\ เอเตหิ ทฺวีหิ ธมฺเมหิ, โย สมนฺนาคโต นโร. \\ กายสฺส เภทา สปฺปญฺโญ, สคฺคํ โส อุปปชฺชตี”ติ.}}

\prose{อยมฺปิ อตฺโถ วุตฺโต ภควตา อิติ เม สุตนฺติ.\csromanspacer{}\csromanspacer{}ฉฏฺฐํ.}

\csromanheader{2}{7. อาตาปีสุตฺต}
\csromanheader{2}{34. วุตฺตํ เหตํ ภควตา วุตฺตมรหตาติ เม สุตํ\csromandash{}}
\prose{อนาตาปี ภิกฺขเว ภิกฺขุ อโนตฺตาปี\footnote{อโนตฺตปฺปี (พหูสุ) อฏฺฐกถา ปสฺสิตพฺพา.} อภพฺโพ สมฺโพธาย, อภพฺโพ นิพฺพานาย, อภพฺโพ อนุตฺตรสฺส โยคกฺเขมสฺส อธิคมาย.\csromanspacer{} อาตาปี จ โข ภิกฺขเว ภิกฺขุ โอตฺตาปี\footnote{โอตฺตปฺปี (พหูสุ)} ภพฺโพ สมฺโพธาย, ภพฺโพ นิพฺพานาย, ภพฺโพ อนุตฺตรสฺส โยคกฺเขมสฺส อธิคมายาติ.\csromanspacer{} เอตมตฺถํ ภควา อโวจ, ตตฺเถตํ อิติ วุจฺจติ\csromandash{}}

\csromangathameasure{“อนาตาปี อโนตฺตาปี, กุสีโต หีนวีริโย. \\ โย ถีนมิทฺธพหุโล, อหิรีโก อนาทโร. \\ อภพฺโพ ตาทิโส ภิกฺขุ, \\ ผุฏฺฐุํ สมฺโพธิมุตฺตมํ. \\ โย จ สติมา นิปโก ฌายี, \\ อาตาปี โอตฺตาปี จ อปฺปมตฺโต. \\ สํโยชนํ ชาติชราย เฉตฺวา, \\ อิเธว สมฺโพธิมนุตฺตรํ ผุเส”ติ.}
\csromangathagroup{\csromangathastanza{“อนาตาปี อโนตฺตาปี, กุสีโต หีนวีริโย. \\ โย ถีนมิทฺธพหุโล, อหิรีโก อนาทโร.}\csromangathastanza{อภพฺโพ ตาทิโส ภิกฺขุ, \\ ผุฏฺฐุํ สมฺโพธิมุตฺตมํ. \\ โย จ สติมา นิปโก ฌายี, \\ อาตาปี โอตฺตาปี จ อปฺปมตฺโต. \\ สํโยชนํ ชาติชราย เฉตฺวา, \\ อิเธว สมฺโพธิมนุตฺตรํ ผุเส”ติ.}}

\prose{อยมฺปิ อตฺโถ วุตฺโต ภควตา อิติ เม สุตนฺติ.\csromanspacer{}\csromanspacer{}สตฺตมํ.}

\csromanheader{2}{8. ปฐมนกุหนสุตฺต}
\csromanheader{2}{35. วุตฺตํ เหตํ ภควตา วุตฺตมรหตาติ เม สุตํ\csromandash{}}
\prose{นยิทํ ภิกฺขเว พฺรหฺมจริยํ วุสฺสติ ชนกุหนตฺถํ, น ชนลปนตฺถํ, น ลาภสกฺการสิโลกานิสํสตฺถํ, น “อิติ มํ ชโน ชานาตู”ติ, อถ โข อิทํ}

\vspace{\tipitakaparskip}
\csromancenter{ทุกนิปาต ภิกฺขเว พฺรหฺมจริยํ วุสฺสติ สํวรตฺถญฺเจว ปหานตฺถญฺจาติ.\csromanspacer{} เอตมตฺถํ ภควา อโวจ, ตตฺเถตํ อิติ วุจฺจติ\csromandash{}}
\vspace{0.5\baselineskip}
\csromangathameasure{*“สํวรตฺถํ ปหานตฺถํ, พฺรหฺมจริยํ อนีติหํ. \\ อเทสยิ โส ภควา, นิพฺพาโนคธคามินํ. \\ เอส มคฺโค มหตฺเตหิ, อนุยาโต มเหสิภิ. \\ เย เย ตํ ปฏิปชฺชนฺติ, ยถา พุทฺเธน เทสิตํ. \\ ทุกฺขสฺสนฺตํ กริสฺสนฺติ, สตฺถุสาสนการิโน”ติ.}
\csromangathagroup{\csromangathastanza{\csromanfolio{215}\csromansymbolfootnote{*}{อํ 1. 334 ปิฏฺเฐปิ.}“สํวรตฺถํ ปหานตฺถํ, พฺรหฺมจริยํ อนีติหํ. \\ อเทสยิ โส ภควา, นิพฺพาโนคธคามินํ.}\csromangathastanza{เอส มคฺโค มหตฺเตหิ\footnote{มหนฺเตหิ (สี, ก), มหตฺเถหิ (สฺยา)}, อนุยาโต มเหสิภิ\footnote{มเหสิโน (สี, ก)}. \\ เย เย ตํ ปฏิปชฺชนฺติ, ยถา พุทฺเธน เทสิตํ. \\ ทุกฺขสฺสนฺตํ กริสฺสนฺติ, สตฺถุสาสนการิโน”ติ.}}

\prose{อยมฺปิ อตฺโถ วุตฺโต ภควตา อิติ เม สุตนฺติ.\csromanspacer{}\csromanspacer{}อฏฺฐมํ.}

\csromanheader{2}{9. ทุติยนกุหนสุตฺต}
\csromanheader{2}{36. วุตฺตํ เหตํ ภควตา วุตฺตมรหตาติ เม สุตํ\csromandash{}}
\prose{นยิทํ ภิกฺขเว พฺรหฺมจริยํ วุสฺสติ ชนกุหนตฺถํ, น ชนลปนตฺถํ, น ลาภสกฺการสิโลกานิสํสตฺถํ, น “อิติ มํ ชโน ชานาตู”ติ, อถ โข อิทํ ภิกฺขเว พฺรหฺมจริยํ วุสฺสติ อภิญฺญตฺถญฺเจว ปริญฺญตฺถญฺจาติ.\csromanspacer{} เอตมตฺถํ ภควา อโวจ, ตตฺเถตํ อิติ วุจฺจติ\csromandash{}}

\csromangathameasure{“อภิญฺญตฺถํ ปริญฺญตฺถํ, พฺรหฺมจริยํ อนีติหํ. \\ อเทสยิ โส ภควา, นิพฺพาโนคธคามินํ. \\ เอส มคฺโค มหตฺเตหิ, อนุยาโต มเหสิภิ. \\ เย เย ตํ ปฏิปชฺชนฺติ, ยถา พุทฺเธน เทสิตํ. \\ ทุกฺขสฺสนฺตํ กริสฺสนฺติ, สตฺถุสาสนการิโน”ติ.}
\csromangathagroup{\csromangathastanza{“อภิญฺญตฺถํ ปริญฺญตฺถํ, พฺรหฺมจริยํ อนีติหํ. \\ อเทสยิ โส ภควา, นิพฺพาโนคธคามินํ.}\csromangathastanza{เอส มคฺโค มหตฺเตหิ, อนุยาโต มเหสิภิ. \\ เย เย ตํ ปฏิปชฺชนฺติ, ยถา พุทฺเธน เทสิตํ. \\ ทุกฺขสฺสนฺตํ กริสฺสนฺติ, สตฺถุสาสนการิโน”ติ.}}

\prose{อยมฺปิ อตฺโถ วุตฺโต ภควตา อิติ เม สุตนฺติ.\csromanspacer{}\csromanspacer{}นวมํ.}

\csromanheader{2}{10. โสมนสฺสสุตฺต}
\csromanheader{2}{37. วุตฺตํ เหตํ ภควตา วุตฺตมรหตาติ เม สุตํ\csromandash{}}
\prose{ทฺวีหิ ภิกฺขเว ธมฺเมหิ สมนฺนาคโต ภิกฺขุ ทิฏฺเฐว ธมฺเม สุขโสมนสฺสพหุโล วิหรติ, โยนิ จสฺส\footnote{โยนิโส (สี, สฺยา, อิ), โยนิสฺส (ก)} อารทฺธา โหติ อาสวานํ ขยาย.}

\vspace{\tipitakaparskip}
\csromancenter{อิติวุตฺตกปาฬิ กตเมหิ ทฺวีหิ? สํเวชนีเยสุ ฐาเนสุ สํเวชเนน, สํวิคฺคสฺส จ โยนิโส ปธาเนน.\csromanspacer{} อิเมหิ โข ภิกฺขเว ทฺวีหิ ธมฺเมหิ สมนฺนาคโต ภิกฺขุ ทิฏฺเฐว ธมฺเม สุขโสมนสฺสพหุโล วิหรติ, โยนิ จสฺส อารทฺธา โหติ อาสวานํ ขยายาติ.\csromanspacer{} เอตมตฺถํ ภควา อโวจ, ตตฺเถตํ อิติ วุจฺจติ\csromandash{}}
\vspace{0.5\baselineskip}
\csromangathameasure{“สํเวชนียฏฺฐาเนสุ, สํวิชฺเชเถว ปณฺฑิโต. \\ อาตาปี นิปโก ภิกฺขุ, ปญฺญาย สมเวกฺขิย. \\ เอวํ วิหารี อาตาปี, สนฺตวุตฺติ อนุทฺธโต. \\ เจโตสมถมนุยุตฺโต, ขยํ ทุกฺขสฺส ปาปุเณ”ติ.}
\csromangathagroup{\csromangathastanza{\csromanfolio{216}“สํเวชนียฏฺฐาเนสุ\footnote{สํเวชนีเยสุ ฐาเนสุ (สฺยา, อิ)}, สํวิชฺเชเถว ปณฺฑิโต. \\ อาตาปี นิปโก ภิกฺขุ, ปญฺญาย สมเวกฺขิย.}\csromangathastanza{เอวํ วิหารี อาตาปี, สนฺตวุตฺติ อนุทฺธโต. \\ เจโตสมถมนุยุตฺโต, ขยํ ทุกฺขสฺส ปาปุเณ”ติ.}}

\prose{อยมฺปิ อตฺโถ วุตฺโต ภควตา อิติ เม สุตนฺติ.\csromanspacer{}\csromanspacer{}ทสมํ.\csromanspacer{} ปฐโม วคฺโค.}

\titlehead{\textbf{ตสฺสุทฺทานํ}}
\csromangathameasure{เทฺว จ ภิกฺขู ตปนียา, ตปนียา ปรตฺเถหิ. \\ อาตาปี นกุหนา เทฺว, โสมนสฺเสน เต ทสาติ.}
\csromangathagroup{\csromangathastanza{เทฺว จ ภิกฺขู ตปนียา, ตปนียา ปรตฺเถหิ. \\ อาตาปี\footnote{เทฺว ปาทา (ก), เทฺว อาตาปี (สี)} นกุหนา เทฺว\footnote{น กุหนา จ (สพฺพตฺถ)}, โสมนสฺเสน เต ทสาติ.}}

\chapterhead{2. ทุติยวคฺค}
\csromanheader{2}{1. วิตกฺกสุตฺต}
\csromanheader{2}{38. วุตฺตํ เหตํ ภควตา วุตฺตมรหตาติ เม สุตํ\csromandash{}}
\prose{ตถาคตํ ภิกฺขเว อรหนฺตํ สมฺมาสมฺพุทฺธํ เทฺว วิตกฺกา พหุลํ สมุทาจรนฺติ, เขโม จ วิตกฺโก, ปวิเวโก จ\footnote{วิเวโก จ (สฺยา)}.\csromanspacer{} อพฺยาปชฺชาราโม\footnote{อพฺยาปชฺฌาราโม (ก), อพฺยาพชฺฌาราโม (?)} ภิกฺขเว ตถาคโต อพฺยาปชฺชรโต, ตเมนํ ภิกฺขเว ตถาคตํ อพฺยาปชฺชารามํ อพฺยาปชฺชรตํ เอเสว วิตกฺโก พหุลํ สมุทาจรติ “อิมายาหํ อิริยาย น กญฺจิ พฺยาพาเธมิ ตสํ วา ถาวรํ วา”ติ.}

\csromanheader{2}{2. ทุกนิปาต}
\prose{\csromanfolio{217}ปวิเวการาโม ภิกฺขเว ตถาคโต ปวิเวกรโต, ตเมนํ ภิกฺขเว ตถาคตํ ปวิเวการามํ ปวิเวกรตํ เอเสว วิตกฺโก พหุลํ สมุทาจรติ “ยํ อกุสลํ ตํ ปหีนนฺ”ติ.}

\prose{ตสฺมาติห ภิกฺขเว ตุเมฺหปิ อพฺยาปชฺชารามา วิหรถ อพฺยาปชฺชรตา, เตสํ โว ภิกฺขเว ตุมฺหากํ อพฺยาปชฺชารามานํ วิหรตํ อพฺยาปชฺชรตานํ เอเสว วิตกฺโก พหุลํ สมุทาจริสฺสติ “อิมาย มยํ อิริยาย น กญฺจิ พฺยาพาเธม ตสํ วา ถาวรํ วา”ติ.}

\prose{ปวิเวการามา ภิกฺขเว วิหรถ ปวิเวกรตา, เตสํ โว ภิกฺขเว ตุมฺหากํ ปวิเวการามานํ วิหรตํ ปวิเวกรตานํ เอเสว วิตกฺโก พหุลํ สมุทาจริสฺสติ “กิํ อกุสลํ กิํ อปฺปหีนํ กิํ ปชหามา”ติ.\csromanspacer{} เอตมตฺถํ ภควา อโวจ, ตตฺเถตํ อิติ วุจฺจติ\csromandash{}}

\csromangathameasure{“ตถาคตํ พุทฺธมสยฺหสาหินํ, \\ ทุเว วิตกฺกา สมุทาจรนฺติ นํ. \\ เขโม วิตกฺโก ปฐโม อุทีริโต, \\ ตโต วิเวโก ทุติโย ปกาสิโต. \\ ตโมนุทํ ปารคตํ มเหสิํ, \\ ตํ ปตฺติปตฺตํ วสิมํ อนาสวํ. \\ วิสนฺตรํ ตณฺหกฺขเย วิมุตฺตํ, \\ ตํ เว มุนิํ อนฺติมเทหธาริํ, \\ มารญฺชหํ พฺรูมิ ชราย ปารคุํ. \\ *เสเล ยถา ปพฺพตมุทฺธนิฏฺฐิโต, \\ ยถาปิ ปสฺเส ชนตํ สมนฺตโต. \\ ตถูปมํ ธมฺมมยํ สุเมโธ, \\ ปาสาทมารุยฺห สมนฺตจกฺขุ. \\ โสกาวติณฺณํ ชนตมเปตโสโก, \\ อเวกฺขติ ชาติชราภิภูตนฺ”ติ.}
\csromangathagroup{\csromangathastanza{“ตถาคตํ พุทฺธมสยฺหสาหินํ, \\ ทุเว วิตกฺกา สมุทาจรนฺติ นํ. \\ เขโม วิตกฺโก ปฐโม อุทีริโต, \\ ตโต วิเวโก ทุติโย ปกาสิโต.}\csromangathastanza{ตโมนุทํ ปารคตํ มเหสิํ, \\ ตํ ปตฺติปตฺตํ วสิมํ อนาสวํ. \\ วิสนฺตรํ\footnote{เวสฺสนฺตรํ (สี, ก), วิสฺสนฺตรํ (อิ)} ตณฺหกฺขเย วิมุตฺตํ, \\ ตํ เว มุนิํ อนฺติมเทหธาริํ,}\csromangathastanza{มารญฺชหํ\footnote{มารชหํ (สฺยา), มานชหํ (สี, ก), มานํ ชหํ (อิ)} พฺรูมิ ชราย ปารคุํ. \\ \csromansymbolfootnote{*}{วิ 3. 6; ที 2. 33; ม 1. 225; ม 2. 292; สํ 1. 139 ปิฏฺเฐสุปิ.}เสเล ยถา ปพฺพตมุทฺธนิฏฺฐิโต, \\ ยถาปิ ปสฺเส ชนตํ สมนฺตโต. \\ ตถูปมํ ธมฺมมยํ สุเมโธ,}\csromangathastanza{ปาสาทมารุยฺห สมนฺตจกฺขุ. \\ โสกาวติณฺณํ ชนตมเปตโสโก, \\ อเวกฺขติ ชาติชราภิภูตนฺ”ติ.}}

\prose{อยมฺปิ อตฺโถ วุตฺโต ภควตา อิติ เม สุตนฺติ.\csromanspacer{}\csromanspacer{}ปฐมํ.}

\csromanheader{2}{อิติวุตฺตกปาฬิ \textbf{2. เทสนาสุตฺต}}
\proseitem{39}{\csromanfolio{218}วุตฺตํ เหตํ ภควตา วุตฺตมรหตาติ เม สุตํ\csromandash{}}

\prose{ตถาคตสฺส ภิกฺขเว อรหโต สมฺมาสมฺพุทฺธสฺส เทฺว ธมฺมเทสนา ปริยาเยน ภวนฺติ.\csromanspacer{} กตมา เทฺว? ปาปํ ปาปกโต ปสฺสถาติ อยํ ปฐมา ธมฺมเทสนา, ปาปํ ปาปกโต ทิสฺวา ตตฺถ นิพฺพินฺทถ วิรชฺชถ วิมุจฺจถาติ อยํ ทุติยา ธมฺมเทสนา.\csromanspacer{} ตถาคตสฺส ภิกฺขเว อรหโต สมฺมาสมฺพุทฺธสฺส อิมา เทฺว ธมฺมเทสนา ปริยาเยน ภวนฺตีติ.\csromanspacer{} เอตมตฺถํ ภควา อโวจ, ตตฺเถตํ อิติ วุจฺจติ\csromandash{}}

\csromangathameasure{“ตถาคตสฺส พุทฺธสฺส, สพฺพภูตานุกมฺปิโน, \\ ปริยายวจนํ ปสฺส, \\ เทฺว จ ธมฺมา ปกาสิตา. \\ ปาปกํ ปสฺสถ เจตํ, \\ ตตฺถ จาปิ วิรชฺชถ. ตโต วิรตฺตจิตฺตาเส, \\ ทุกฺขสฺสนฺตํ กริสฺสถา”ติ.}
\csromangathagroup{\csromangathastanza{“ตถาคตสฺส พุทฺธสฺส, สพฺพภูตานุกมฺปิโน, \\ ปริยายวจนํ ปสฺส, \\ เทฺว จ ธมฺมา ปกาสิตา. \\ ปาปกํ ปสฺสถ เจตํ\footnote{เจกํ (สี, อิ), เฉกา (สฺยา)}, \\ ตตฺถ จาปิ วิรชฺชถ. ตโต วิรตฺตจิตฺตาเส, \\ ทุกฺขสฺสนฺตํ กริสฺสถา”ติ.}}

\prose{อยมฺปิ อตฺโถ วุตฺโต ภควตา อิติ เม สุตนฺติ.\csromanspacer{}\csromanspacer{}ทุติยํ.}

\csromanheader{2}{3. วิชฺชาสุตฺต}
\csromanmatikaanchor{mtk.140}
\csromantocmarkref{subsection}{2. เมตฺตสุตฺต}{mtk.140}
\bookmark[dest={mtk.140},level=2]{2. เมตฺตสุตฺต}
\csromanheader{2}{40. วุตฺตํ เหตํ ภควตา วุตฺตมรหตาติ เม สุตํ\csromandash{}}
\prose{อวิชฺชา ภิกฺขเว ปุพฺพงฺคมา อกุสลานํ ธมฺมานํ สมาปตฺติยา, อนฺวเทว อหิริกํ อโนตฺตปฺปํ, วิชฺชา จ โข ภิกฺขเว ปุพฺพงฺคมา กุสลานํ ธมฺมานํ สมาปตฺติยา, อนฺวเทว หิโรตฺตปฺปนฺติ.\csromanspacer{} เอตมตฺถํ ภควา อโวจ, ตตฺเถตํ อิติ วุจฺจติ\csromandash{}}

\csromangathameasure{“ยา กาจิมา ทุคฺคติโย, อสฺมิํ โลเก ปรมฺหิ จ. \\ อวิชฺชามูลิกา สพฺพา, อิจฺฉาโลภสมุสฺสยา. \\ ยโต จ โหติ ปาปิจฺโฉ, อหิรีโก อนาทโร. \\ ตโต ปาปํ ปสวติ, อปายํ เตน คจฺฉติ. \\ ตสฺมา ฉนฺทญฺจ โลภญฺจ, อวิชฺชญฺจ วิราชยํ. \\ วิชฺชํ อุปฺปาทยํ ภิกฺขุ, สพฺพา ทุคฺคติโย ชเห”ติ.}
\csromangathagroup{\csromangathastanza{“ยา กาจิมา ทุคฺคติโย, อสฺมิํ โลเก ปรมฺหิ จ. \\ อวิชฺชามูลิกา สพฺพา, อิจฺฉาโลภสมุสฺสยา.}\csromangathastanza{ยโต จ โหติ ปาปิจฺโฉ, อหิรีโก อนาทโร. \\ ตโต ปาปํ ปสวติ, อปายํ เตน คจฺฉติ.}\csromangathastanza{ตสฺมา ฉนฺทญฺจ โลภญฺจ, อวิชฺชญฺจ วิราชยํ. \\ วิชฺชํ อุปฺปาทยํ ภิกฺขุ, สพฺพา ทุคฺคติโย ชเห”ติ.}}

\prose{อยมฺปิ อตฺโถ วุตฺโต ภควตา อิติ เม สุตนฺติ.\csromanspacer{}\csromanspacer{}ตติยํ.\csromanspacer{} ปฐมภาณวาโร.}

\csromanheader{2}{2. ทุกนิปาต \textbf{4. ปญฺญาปริหีนสุตฺต}}
\csromanheader{2}{41. วุตฺตํ เหตํ ภควตา วุตฺตมรหตาติ เม สุตํ\csromandash{}}
\prose{\csromanfolio{219}เต ภิกฺขเว สตฺตา สุปริหีนา, เย อริยาย ปญฺญาย ปริหีนา, เต ทิฏฺเฐว ธมฺเม ทุกฺขํ วิหรนฺติ สวิฆาตํ ส-อุปายาสํ สปริฬาหํ, กายสฺส เภทา ปรํ มรณา ทุคฺคติ ปาฏิกงฺขา.\csromanspacer{} เต\footnote{เต จ โข (?)} ภิกฺขเว สตฺตา อปริหีนา เย อริยาย ปญฺญาย อปริหีนา, เต ทิฏฺเฐว ธมฺเม สุขํ วิหรนฺติ อวิฆาตํ อนุปายาสํ อปริฬาหํ, กายสฺส เภทา ปรํ มรณา สุคติ ปาฏิกงฺขาติ.\csromanspacer{} เอตมตฺถํ ภควา อโวจ, ตตฺเถตํ อิติ วุจฺจติ\csromandash{}}

\csromanmatikaanchor{mtk.141}
\csromantocmarkref{subsection}{3. อุภยตฺถสุตฺต}{mtk.141}
\bookmark[dest={mtk.141},level=2]{3. อุภยตฺถสุตฺต}
\csromangathameasure{“ปญฺญาย ปริหาเนน, ปสฺส โลกํ สเทวกํ. \\ นิวิฏฺฐํ นามรูปสฺมิํ, อิทํ สจฺจนฺติ มญฺญติ. \\ +ปญฺญา หิ เสฏฺฐา โลกสฺมิํ, ยายํ นิพฺเพธคามินี. \\ ยาย สมฺมา ปชานาติ, ชาติภวปริกฺขยํ. \\ เตสํ เทวา มนุสฺสา จ, สมฺพุทฺธานํ สตีมตํ. \\ ปิหยนฺติ หาสปญฺญานํ, สรีรนฺติมธารินนฺ”ติ.}
\csromangathagroup{\csromangathastanza{“ปญฺญาย ปริหาเนน, ปสฺส โลกํ สเทวกํ. \\ นิวิฏฺฐํ นามรูปสฺมิํ, อิทํ สจฺจนฺติ มญฺญติ.}\csromangathastanza{\csromansymbolfootnote{+}{ขุ 10. 143 ปิฏฺเฐปิ.}ปญฺญา หิ เสฏฺฐา โลกสฺมิํ, ยายํ นิพฺเพธคามินี. \\ ยาย สมฺมา ปชานาติ, ชาติภวปริกฺขยํ.}\csromangathastanza{เตสํ เทวา มนุสฺสา จ, สมฺพุทฺธานํ สตีมตํ. \\ ปิหยนฺติ หาสปญฺญานํ\footnote{หาสุปญฺญานํ (สี-ฏฺฐ)}, สรีรนฺติมธารินนฺ”ติ.}}

\prose{อยมฺปิ อตฺโถ วุตฺโต ภควตา อิติ เม สุตนฺติ.\csromanspacer{}\csromanspacer{}จตุตฺถํ.}

\csromanheader{2}{5. สุกฺกธมฺมสุตฺต}
\csromanheader{2}{42. วุตฺตํ เหตํ ภควตา วุตฺตมรหตาติ เม สุตํ\csromandash{}}
\csromanmatikaanchor{mtk.142}
\csromantocmarkref{subsection}{4. อฏฺฐิปุญฺชสุตฺต}{mtk.142}
\bookmark[dest={mtk.142},level=2]{4. อฏฺฐิปุญฺชสุตฺต}
\prose{\csromansymbolfootnote{*}{อํ 1. 53 ปิฏฺเฐปิ.}เทฺวเม ภิกฺขเว สุกฺกา ธมฺมา โลกํ ปาเลนฺติ.\csromanspacer{} กตเม เทฺว? หิรี\footnote{หิริ (สี, สฺยา, กํ, อิ)} จ โอตฺตปฺปญฺจ.\csromanspacer{} อิเม เจ ภิกฺขเว เทฺว สุกฺกา ธมฺมา โลกํ น ปาเลยฺยุํ, นยิธ ปญฺญาเยถ “มาตา”ติ วา “มาตุจฺฉา”ติ วา “มาตุลานี”ติ วา “อาจริยภริยา”ติ วา “ครูนํ ทารา”ติ วา, สมฺเภทํ โลโก อคมิสฺส, ยถา อเชฬกา กุกฺกุฏสูกรา โสณสิงฺคาลา\footnote{โสณสิคาลา (สี, สฺยา, กํ, อิ)}.\csromanspacer{} ยสฺมา จ โข ภิกฺขเว อิเม เทฺว สุกฺกา ธมฺมา โลกํ ปาเลนฺติ, ตสฺมา ปญฺญายติ “มาตา”ติ วา “มาตุจฺฉา”ติ วา “มาตุลานี”ติ วา “อาจริยภริยา”ติ วา “ครูนํ ทารา”ติ วาติ.\csromanspacer{} เอตมตฺถํ ภควา อโวจ, ตตฺเถตํ อิติ วุจฺจติ\csromandash{}}

\gambhira{อิติวุตฺตกปาฬิ}
\csromangathameasure{“เยสํ เจ หิริ-โอตฺตปฺปํ, สพฺพทา จ น วิชฺชติ. \\ โวกฺกนฺตา สุกฺกมูลา เต, ชาติมรณคามิโน. \\ เยสญฺจ หิริ-โอตฺตปฺปํ, สทา สมฺมา อุปฏฺฐิตา. \\ วิรูฬฺหพฺรหฺมจริยา เต, สนฺโต ขีณปุนพฺภวา”ติ.}
\csromangathagroup{\csromangathastanza{\csromanfolio{220}“เยสํ เจ หิริ-โอตฺตปฺปํ, สพฺพทา จ น วิชฺชติ. \\ โวกฺกนฺตา สุกฺกมูลา เต, ชาติมรณคามิโน.}\csromangathastanza{เยสญฺจ หิริ-โอตฺตปฺปํ, สทา สมฺมา อุปฏฺฐิตา. \\ วิรูฬฺหพฺรหฺมจริยา เต, สนฺโต ขีณปุนพฺภวา”ติ.}}

\prose{อยมฺปิ อตฺโถ วุตฺโต ภควตา อิติ เม สุตนฺติ.\csromanspacer{}\csromanspacer{}ปญฺจมํ.}

\csromanheader{2}{6. อชาตสุตฺต}
\csromanheader{2}{43. วุตฺตํ เหตํ ภควตา วุตฺตมรหตาติ เม สุตํ\csromandash{}}
\prose{อตฺถิ ภิกฺขเว อชาตํ อภูตํ อกตํ อสงฺขตํ, โน เจตํ ภิกฺขเว อภวิสฺส อชาตํ อภูตํ อกตํ อสงฺขตํ, นยิธ ชาตสฺส ภูตสฺส กตสฺส สงฺขตสฺส นิสฺสรณํ ปญฺญาเยถ.\csromanspacer{} ยสฺมา จ โข ภิกฺขเว อตฺถิ อชาตํ อภูตํ อกตํ อสงฺขตํ, ตสฺมา ชาตสฺส ภูตสฺส กตสฺส สงฺขตสฺส นิสฺสรณํ ปญฺญายตีติ.\csromanspacer{} เอตมตฺถํ ภควา อโวจ, ตตฺเถตํ อิติ วุจฺจติ\csromandash{}}

\csromanmatikaanchor{mtk.143}
\csromantocmarkref{subsection}{5. มุสาวาทสุตฺต}{mtk.143}
\bookmark[dest={mtk.143},level=2]{5. มุสาวาทสุตฺต}
\csromangathameasure{“ชาตํ ภูตํ สมุปฺปนฺนํ, กตํ สงฺขตมทฺธุวํ. \\ ชรามรณสงฺฆาฏํ, โรคนีฬํ ปภงฺคุรํ. \\ อาหารเนตฺติปฺปภวํ, นาลํ ตทภินนฺทิตุํ. \\ ตสฺส นิสฺสรณํ สนฺตํ, อตกฺกาวจรํ ธุวํ. \\ อชาตํ อสมุปฺปนฺนํ, อโสกํ วิรชํ ปทํ. \\ นิโรโธ ทุกฺขธมฺมานํ, สงฺขารูปสโม สุโข”ติ.}
\csromangathagroup{\csromangathastanza{“ชาตํ ภูตํ สมุปฺปนฺนํ, กตํ สงฺขตมทฺธุวํ. \\ ชรามรณสงฺฆาฏํ, โรคนีฬํ\footnote{โรคนิฑฺฒํ (สี)} ปภงฺคุรํ\footnote{ปภงฺคุนํ (ก-สี, ก), ปภงฺคุณํ (สฺยา)}.}\csromangathastanza{อาหารเนตฺติปฺปภวํ, นาลํ ตทภินนฺทิตุํ. \\ ตสฺส นิสฺสรณํ สนฺตํ, อตกฺกาวจรํ ธุวํ.}\csromangathastanza{อชาตํ อสมุปฺปนฺนํ, อโสกํ วิรชํ ปทํ. \\ นิโรโธ ทุกฺขธมฺมานํ, สงฺขารูปสโม สุโข”ติ.}}

\prose{อยมฺปิ อตฺโถ วุตฺโต ภควตา อิติ เม สุตนฺติ.\csromanspacer{}\csromanspacer{}ฉฏฺฐํ.}

\csromanheader{2}{7. นิพฺพานธาตุสุตฺต}
\csromanheader{2}{44. วุตฺตํ เหตํ ภควตา วุตฺตมรหตาติ เม สุตํ\csromandash{}}
\csromanmatikaanchor{mtk.144}
\csromantocmarkref{subsection}{6. ทานสุตฺต}{mtk.144}
\bookmark[dest={mtk.144},level=2]{6. ทานสุตฺต}
\prose{เทฺวมา ภิกฺขเว นิพฺพานธาตุโย.\csromanspacer{} กตมา เทฺว? ส-อุปาทิเสสา จ นิพฺพานธาตุ, อนุปาทิเสสา จ นิพฺพานธาตุ.}

\csromanheader{2}{2. ทุกนิปาต}
\prose{\csromanfolio{221}กตมา จ ภิกฺขเว ส-อุปาทิเสสา นิพฺพานธาตุ? อิธ ภิกฺขเว ภิกฺขุ อรหํ โหติ ขีณาสโว วุสิตวา กตกรณีโย โอหิตภาโร อนุปฺปตฺตสทตฺโถ ปริกฺขีณภวสํโยชโน สมฺมทญฺญา วิมุตฺโต, ตสฺส ติฏฺฐนฺเตว ปญฺจินฺทฺริยานิ เยสํ อวิฆาตตฺตา\footnote{อวิคตตฺตา (สี-ฏฺฐ)} มนาปามนาปํ ปจฺจนุโภติ, สุขทุกฺขํ ปฏิสํเวเทติ.\csromanspacer{} ตสฺส โย ราคกฺขโย โทสกฺขโย โมหกฺขโย.\csromanspacer{} อยํ วุจฺจติ ภิกฺขเว ส-อุปาทิเสสา นิพฺพานธาตุ.}

\prose{กตมา จ ภิกฺขเว อนุปาทิเสสา นิพฺพานธาตุ? อิธ ภิกฺขเว ภิกฺขุ อรหํ โหติ ขีณาสโว วุสิตวา กตกรณีโย โอหิตภาโร อนุปฺปตฺตสทตฺโถ ปริกฺขีณภวสํโยชโน สมฺมทญฺญา วิมุตฺโต, ตสฺส อิเธว ภิกฺขเว สพฺพเวทยิตานิ อนภินนฺทิตานิ สีติ ภวิสฺสนฺติ\footnote{สีตีภวิสฺสนฺติ (?)}.\csromanspacer{} อยํ วุจฺจติ ภิกฺขเว อนุปาทิเสสา นิพฺพานธาตุ.\csromanspacer{} อิมา โข ภิกฺขเว เทฺว นิพฺพานธาตุโยติ.\csromanspacer{} เอตมตฺถํ ภควา อโวจ, ตตฺเถกํ อิติ วุจฺจติ\csromandash{}}

\csromangathameasure{“ทุเว อิมา จกฺขุมตา ปกาสิตา, \\ นิพฺพานธาตู อนิสฺสิเตน ตาทินา. \\ เอกา หิ ธาตุ อิธ ทิฏฺฐธมฺมิกา, \\ ส-อุปาทิเสสา ภวเนตฺติสงฺขยา. \\ อนุปาทิเสสา ปน สมฺปรายิกา, \\ ยมฺหิ นิรุชฺฌนฺติ ภวานิ สพฺพโส. \\ เย เอตทญฺญาย ปทํ อสงฺขตํ, \\ วิมุตฺตจิตฺตา ภวเนตฺติสงฺขยา. \\ เต ธมฺมสาราธิคมา ขเย รตา, \\ ปหํสุ เต สพฺพภวานิ ตาทิโน”ติ.}
\csromangathagroup{\csromangathastanza{“ทุเว อิมา จกฺขุมตา ปกาสิตา, \\ นิพฺพานธาตู อนิสฺสิเตน ตาทินา. \\ เอกา หิ ธาตุ อิธ ทิฏฺฐธมฺมิกา, \\ ส-อุปาทิเสสา ภวเนตฺติสงฺขยา.}\csromangathastanza{อนุปาทิเสสา ปน สมฺปรายิกา, \\ ยมฺหิ นิรุชฺฌนฺติ ภวานิ สพฺพโส. \\ เย เอตทญฺญาย ปทํ อสงฺขตํ, \\ วิมุตฺตจิตฺตา ภวเนตฺติสงฺขยา. \\ เต ธมฺมสาราธิคมา ขเย รตา, \\ ปหํสุ เต สพฺพภวานิ ตาทิโน”ติ.}}

\prose{อยมฺปิ อตฺโถ วุตฺโต ภควตา อิติ เม สุตนฺติ.\csromanspacer{}\csromanspacer{}สตฺตมํ.}

\csromanheader{2}{8. ปฏิสลฺลานสุตฺต}
\csromanheader{2}{45. วุตฺตํ เหตํ ภควตา วุตฺตมรหตาติ เม สุตํ\csromandash{}}
\csromanmatikaanchor{mtk.145}
\csromantocmarkref{subsection}{7. เมตฺตาภาวนาสุตฺต}{mtk.145}
\bookmark[dest={mtk.145},level=2]{7. เมตฺตาภาวนาสุตฺต}
\prose{ปฏิสลฺลานารามา\footnote{ปฏิสลฺลาณารามา (ก)} ภิกฺขเว วิหรถ ปฏิสลฺลานรตา อชฺฌตฺตํ เจโตสมถมนุยุตฺตา อนิรากตชฺฌานา วิปสฺสนาย สมนฺนาคตา พฺรูเหตาโร}

\vspace{\tipitakaparskip}
\csromancenter{อิติวุตฺตกปาฬิ สุญฺญาคารานํ, ปฏิสลฺลานารามานํ ภิกฺขเว วิหรตํ ปฏิสลฺลานรตานํ อชฺฌตฺตํ เจโตสมถมนุยุตฺตานํ อนิรากตชฺฌานานํ วิปสฺสนาย สมนฺนาคตานํ พฺรูเหตานํ สุญฺญาคารานํ ทฺวินฺนํ ผลานํ อญฺญตรํ ผลํ ปาฏิกงฺขํ ทิฏฺเฐว ธมฺเม อญฺญา, สติ วา อุปาทิเสเส อนาคามิตาติ.\csromanspacer{} เอตมตฺถํ ภควา อโวจ, ตตฺเถตํ อิติ วุจฺจติ\csromandash{}}
\vspace{0.5\baselineskip}
\csromangathameasure{“เย สนฺตจิตฺตา นิปกา, สติมนฺโต จ ฌายิโน. \\ สมฺมา ธมฺมํ วิปสฺสนฺติ, กาเมสุ อนเปกฺขิโน. \\ อปฺปมาทรตา สนฺตา, ปมาเท ภยทสฺสิโน. \\ อภพฺพา ปริหานาย, นิพฺพานสฺเสว สนฺติเก”ติ.}
\csromangathagroup{\csromangathastanza{\csromanfolio{222}“เย สนฺตจิตฺตา นิปกา, สติมนฺโต จ\footnote{สติมนฺโตว (สี, ก)} ฌายิโน. \\ สมฺมา ธมฺมํ วิปสฺสนฺติ, กาเมสุ อนเปกฺขิโน.}\csromangathastanza{อปฺปมาทรตา สนฺตา, ปมาเท ภยทสฺสิโน. \\ อภพฺพา ปริหานาย, นิพฺพานสฺเสว สนฺติเก”ติ.}}

\prose{อยมฺปิ อตฺโถ วุตฺโต ภควตา อิติ เม สุตนฺติ.\csromanspacer{}\csromanspacer{}อฏฺฐมํ.}

\csromanheader{2}{9. สิกฺขานิสํสสุตฺต}
\csromanheader{2}{46. วุตฺตํ เหตํ ภควตา วุตฺตมรหตาติ เม สุตํ\csromandash{}}
\prose{สิกฺขานิสํสา ภิกฺขเว วิหรถ ปญฺญุตฺตรา วิมุตฺติสารา สตาธิปเตยฺยา, สิกฺขานิสํสานํ ภิกฺขเว วิหรตํ ปญฺญุตฺตรานํ วิมุตฺติสารานํ สตาธิปเตยฺยานํ ทฺวินฺนํ ผลานํ อญฺญตรํ ผลํ ปาฏิกงฺขํ ทิฏฺเฐว ธมฺเม อญฺญา, สติ วา อุปาทิเสเส อนาคามิตาติ.\csromanspacer{} เอตมตฺถํ ภควา อโวจ, ตตฺเถตํ อิติ วุจฺจติ\csromandash{}}

\csromangathameasure{“ปริปุณฺณสิกฺขํ อปหานธมฺมํ, \\ ปญฺญุตฺตรํ ชาติขยนฺตทสฺสิํ. \\ ตํ เว มุนิํ อนฺติมเทหธาริํ, \\ มารญฺชหํ พฺรูมิ ชราย ปารคุํ. \\ ตสฺมา สทา ฌานรตา สมาหิตา, \\ อาตาปิโน ชาติขยนฺตทสฺสิโน. \\ มารํ สเสนํ อภิภุยฺย ภิกฺขโว, \\ ภวถ ชาติมรณสฺส ปารคา”ติ.}
\csromangathagroup{\csromangathastanza{“ปริปุณฺณสิกฺขํ\footnote{ปริปุณฺณเสขํ (สี), ปริปุณฺณเสกฺขํ (สฺยา)} อปหานธมฺมํ, \\ ปญฺญุตฺตรํ ชาติขยนฺตทสฺสิํ. \\ ตํ เว มุนิํ อนฺติมเทหธาริํ, \\ มารญฺชหํ พฺรูมิ ชราย ปารคุํ.}\csromangathastanza{ตสฺมา สทา ฌานรตา สมาหิตา, \\ อาตาปิโน ชาติขยนฺตทสฺสิโน. \\ มารํ สเสนํ อภิภุยฺย ภิกฺขโว, \\ ภวถ ชาติมรณสฺส ปารคา”ติ.}}

\prose{อยมฺปิ อตฺโถ วุตฺโต ภควตา อิติ เม สุตนฺติ.\csromanspacer{}\csromanspacer{}นวมํ.}

\csromanheader{2}{2. ทุกนิปาต \textbf{10. ชาคริยสุตฺต}}
\csromanheader{2}{47. วุตฺตํ เหตํ ภควตา วุตฺตมรหตาติ เม สุตํ\csromandash{}}
\prose{\csromanfolio{223}ชาคโร จสฺส ภิกฺขเว ภิกฺขุ วิหเรยฺย สโต สมฺปชาโน สมาหิโต ปมุทิโต วิปฺปสนฺโน จ, ตตฺถ กาลวิปสฺสี จ กุสเลสุ ธมฺเมสุ.\csromanspacer{} ชาครสฺส ภิกฺขเว ภิกฺขุโน วิหรโต สตสฺส สมฺปชานสฺส สมาหิตสฺส ปมุทิตสฺส วิปฺปสนฺนสฺส ตตฺถ กาลวิปสฺสิโน กุสเลสุ ธมฺเมสุ ทฺวินฺนํ ผลานํ อญฺญตรํ ผลํ ปาฏิกงฺขํ ทิฏฺเฐว ธมฺเม อญฺญา, สติ วา อุปาทิเสเส อนาคามิตาติ.\csromanspacer{} เอตมตฺถํ ภควา อโวจ, ตตฺเถตํ อิติ วุจฺจติ\csromandash{}}

\csromangathameasure{“ชาครนฺตา สุณาเถตํ, เย สุตฺตา เต ปพุชฺฌถ. \\ สุตฺตา ชาคริตํ เสยฺโย, นตฺถิ ชาครโต ภยํ. \\ โย ชาคโร จ สติมา สมฺปชาโน, \\ สมาหิโต มุทิโต วิปฺปสนฺโน จ. \\ กาเลน โส สมฺมา ธมฺมํ ปริวีมํสมาโน, \\ เอโกทิภูโต วิหเน ตมํ โส. \\ ตสฺมา หเว ชาคริยํ ภเชถ, \\ อาตาปี ภิกฺขุ นิปโก ฌานลาภี. \\ สํโยชนํ ชาติชราย เฉตฺวา, \\ อิเธว สมฺโพธิมนุตฺตรํ ผุเส”ติ.}
\csromangathagroup{\csromangathastanza{“ชาครนฺตา สุณาเถตํ, เย สุตฺตา เต ปพุชฺฌถ. \\ สุตฺตา ชาคริตํ เสยฺโย, นตฺถิ ชาครโต ภยํ.}\csromangathastanza{โย ชาคโร จ สติมา สมฺปชาโน, \\ สมาหิโต มุทิโต วิปฺปสนฺโน จ. \\ กาเลน โส สมฺมา ธมฺมํ ปริวีมํสมาโน, \\ เอโกทิภูโต วิหเน ตมํ โส.}\csromangathastanza{ตสฺมา หเว ชาคริยํ ภเชถ, \\ อาตาปี ภิกฺขุ นิปโก ฌานลาภี. \\ สํโยชนํ ชาติชราย เฉตฺวา, \\ อิเธว สมฺโพธิมนุตฺตรํ ผุเส”ติ.}}

\csromanmatikaanchor{mtk.146}
\csromantocmarknopage{chapter}{2. ทุกนิปาต}{mtk.146}
\bookmark[dest={mtk.146},level=0]{2. ทุกนิปาต}
\csromantocmarknopage{section}{1. ปฐมวคฺค}{mtk.147}
\bookmark[dest={mtk.147},level=1]{1. ปฐมวคฺค}
\csromantocmarkref{subsection}{1-2. ทุกฺขวิหารสุตฺตาทิ}{mtk.148}
\bookmark[dest={mtk.148},level=2]{1-2. ทุกฺขวิหารสุตฺตาทิ}
\csromantocmarkref{subsection}{3-4. ตปนียสุตฺตาทิ}{mtk.149}
\bookmark[dest={mtk.149},level=2]{3-4. ตปนียสุตฺตาทิ}
\csromantocmarkref{subsection}{5-6. สีลสุตฺตานิ}{mtk.150}
\bookmark[dest={mtk.150},level=2]{5-6. สีลสุตฺตานิ}
\prose{อยมฺปิ อตฺโถ วุตฺโต ภควตา อิติ เม สุตนฺติ.\csromanspacer{}\csromanspacer{}ทสมํ.}

\csromanheader{2}{11. อาปายิกสุตฺต}
\csromanheader{2}{48. วุตฺตํ เหตํ ภควตา วุตฺตมรหตาติ เม สุตํ\csromandash{}}
\prose{เทฺวเม ภิกฺขเว อาปายิกา เนรยิกา อิทมปฺปหาย.\csromanspacer{} กตเม เทฺว? โย จ อพฺรหฺมจารี พฺรหฺมจาริปฏิญฺโญ, โย จ ปริปุณฺณํ ปริสุทฺธํ พฺรหฺมจริยํ จรนฺตํ อมูลเกน อพฺรหฺมจริเยน อนุทฺธํเสติ.\csromanspacer{} อิเม โข ภิกฺขเว เทฺว อาปายิกา เนรยิกา อิทมปฺปหายาติ.\csromanspacer{} เอตมตฺถํ ภควา อโวจ, ตตฺเถตํ อิติ วุจฺจติ\csromandash{}}

\gambhira{อิติวุตฺตกปาฬิ}
\csromangathameasure{+“อภูตวาที นิรยํ อุเปติ, \\ โย วาปิ กตฺวา น กโรมิ’จาห. \\ อุโภปิ เต เปจฺจ สมา ภวนฺติ, \\ นิหีนกมฺมา มนุชา ปรตฺถ. \\ *กาสาวกณฺฐา พหโว, ปาปธมฺมา อสญฺญตา. \\ ปาปา ปาเปหิ กมฺเมหิ, นิรยํ เต อุปปชฺชเร. \\ *เสยฺโย อโยคุโฬ ภุตฺโต, ตตฺโต อคฺคิสิขูปโม. \\ ยญฺเจ ภุญฺเชยฺย ทุสฺสีโล, รฏฺฐปิณฺฑมสญฺญโต”ติ.}
\csromangathagroup{\csromangathastanza{\csromanfolio{224}\csromansymbolfootnote{+}{เหฏฺฐา 57 ปิฏฺเฐปิ.}“อภูตวาที นิรยํ อุเปติ, \\ โย วาปิ กตฺวา น กโรมิ’จาห. \\ อุโภปิ เต เปจฺจ สมา ภวนฺติ, \\ นิหีนกมฺมา มนุชา ปรตฺถ. \\ \csromansymbolfootnote{*}{เหฏฺฐา 57, อุปริ 256; วิ 1. 116 ปิฏฺเฐสุปิ.}กาสาวกณฺฐา พหโว, ปาปธมฺมา อสญฺญตา. \\ ปาปา ปาเปหิ กมฺเมหิ, นิรยํ เต อุปปชฺชเร.}\csromangathastanza{\csromansymbolmark{*}เสยฺโย อโยคุโฬ ภุตฺโต, ตตฺโต อคฺคิสิขูปโม. \\ ยญฺเจ ภุญฺเชยฺย ทุสฺสีโล, รฏฺฐปิณฺฑมสญฺญโต”ติ.}}

\prose{อยมฺปิ อตฺโถ วุตฺโต ภควตา อิติ เม สุตนฺติ.\csromanspacer{}\csromanspacer{}เอกาทสมํ.}

\csromanheader{2}{12. ทิฏฺฐิคตสุตฺต}
\csromanheader{2}{49. วุตฺตํ เหตํ ภควตา วุตฺตมรหตาติ เม สุตํ\csromandash{}}
\prose{ทฺวีหิ ภิกฺขเว ทิฏฺฐิคเตหิ ปริยุฏฺฐิตา เทวมนุสฺสา โอลียนฺติ เอเก, อติธาวนฺติ เอเก, จกฺขุมนฺโต จ ปสฺสนฺติ.}

\prose{กถญฺจ ภิกฺขเว โอลียนฺติ เอเก? ภวารามา ภิกฺขเว เทวมนุสฺสา ภวรตา ภวสมฺมุทิตา, เตสํ ภวนิโรธาย ธมฺเม เทสิยมาเน จิตฺตํ น ปกฺขนฺทติ น ปสีทติ น สนฺติฏฺฐติ นาธิมุจฺจติ.\csromanspacer{} เอวํ โข ภิกฺขเว โอลียนฺติ เอเก.}

\prose{กถญฺจ ภิกฺขเว อติธาวนฺติ เอเก? ภเวเนว โข ปเนเก อฏฺฏียมานา หรายมานา ชิคุจฺฉมานา วิภวํ อภินนฺทนฺติ, ยโต กิร โภ อยํ อตฺตา\footnote{สตฺโต (สี, ก)} กายสฺส เภทา ปรํ มรณา อุจฺฉิชฺชติ วินสฺสติ น โหติ ปรํ มรณา, เอตํ สนฺตํ เอตํ ปณีตํ เอตํ ยาถาวนฺติ.\csromanspacer{} เอวํ โข ภิกฺขเว อติธาวนฺติ เอเก.}

\prose{กถญฺจ ภิกฺขเว จกฺขุมนฺโต ปสฺสนฺติ? อิธ ภิกฺขุ ภูตํ ภูตโต ปสฺสติ, ภูตํ ภูตโต ทิสฺวา ภูตสฺส นิพฺพิทาย วิราคาย นิโรธาย ปฏิปนฺโน โหติ.\csromanspacer{} เอวํ โข ภิกฺขเว จกฺขุมนฺโต ปสฺสนฺตีติ.\csromanspacer{} เอตมตฺถํ ภควา อโวจ, ตตฺเถตํ อิติ วุจฺจติ\csromandash{}}

\csromanheader{2}{2. ทุกนิปาต}
\csromangathameasure{“เย ภูตํ ภูตโต ทิสฺวา, ภูตสฺส จ อติกฺกมํ. \\ ยถาภูเต วิมุจฺจนฺติ, ภวตณฺหา ปริกฺขยา. \\ ส เว ภูตปริญฺโญ โส, วีตตโณฺห ภวาภเว. \\ ภูตสฺส วิภวา ภิกฺขุ, นาคจฺฉติ ปุนพฺภวนฺ”ติ.}
\csromangathagroup{\csromangathastanza{\csromanfolio{225}“เย\footnote{โย (สฺยา, ก)} ภูตํ ภูตโต ทิสฺวา, ภูตสฺส จ อติกฺกมํ. \\ ยถาภูเต วิมุจฺจนฺติ, ภวตณฺหา ปริกฺขยา.}\csromangathastanza{ส เว\footnote{สเจ (ก-สี, สฺยา, อิ)} ภูตปริญฺโญ โส, วีตตโณฺห ภวาภเว. \\ ภูตสฺส วิภวา ภิกฺขุ, นาคจฺฉติ ปุนพฺภวนฺ”ติ.}}

\prose{อยมฺปิ อตฺโถ วุตฺโต ภควตา อิติ เม สุตนฺติ.\csromanspacer{}\csromanspacer{}ทฺวาทสมํ.\csromanspacer{} ทุติโย วคฺโค.}

\titlehead{\textbf{ตสฺสุทฺทานํ}}
\csromangathameasure{เทฺว อินฺทฺริยา เทฺว ตปนียา, สีเลน อปเร ทุเว. \\ อโนตฺตาปี กุหนา เทฺว จ, สํเวชนีเยน เต ทส. \\ วิตกฺกา เทสนา วิชฺชา, ปญฺญา ธมฺเมน ปญฺจมํ. \\ อชาตํ ธาตุสลฺลานํ, สิกฺขา ชาคริเยน จ. \\ อปายทิฏฺฐิยา เจว, พาวีสติ ปกาสิตาติ.}
\csromangathagroup{\csromangathastanza{เทฺว อินฺทฺริยา เทฺว ตปนียา, สีเลน อปเร ทุเว. \\ อโนตฺตาปี กุหนา เทฺว จ, สํเวชนีเยน เต ทส.}\csromangathastanza{วิตกฺกา เทสนา วิชฺชา, ปญฺญา ธมฺเมน ปญฺจมํ. \\ อชาตํ ธาตุสลฺลานํ, สิกฺขา ชาคริเยน จ. \\ อปายทิฏฺฐิยา เจว\footnote{เยว (สี, สฺยา)}, พาวีสติ ปกาสิตาติ.}}

\nitthitam{\textbf{ทุกนิปาโต นิฏฺฐิโต.}}
\csromanheader{1}{3. ติกนิปาต}
\chapterhead{1. ปฐมวคฺค}
\csromanheader{3}{1. มูลสุตฺต}
\csromanheader{2}{50. วุตฺตํ เหตํ ภควตา วุตฺตมรหตาติ เม สุตํ\csromandash{}}
\prose{\csromanfolio{226}ตีณิมานิ ภิกฺขเว อกุสลมูลานิ.\csromanspacer{} กตมานิ ตีณิ? โลโภ อกุสลมูลํ, โทโส อกุสลมูลํ, โมโห อกุสลมูลํ.\csromanspacer{} อิมานิ โข ภิกฺขเว ตีณิ อกุสลมูลานีติ.\csromanspacer{} เอตมตฺถํ ภควา อโวจ, ตตฺเถตํ อิติ วุจฺจติ\csromandash{}}

\csromangathameasure{*“โลโภ โทโส จ โมโห จ, ปุริสํ ปาปเจตสํ. \\ หิํสนฺติ อตฺตสมฺภูตา, ตจสารํว สมฺผลนฺ”ติ.}
\csromangathagroup{\csromangathastanza{\csromansymbolfootnote{*}{ขุ 8. 252 ปิฏฺเฐปิ.}“โลโภ โทโส จ โมโห จ, ปุริสํ ปาปเจตสํ. \\ หิํสนฺติ อตฺตสมฺภูตา, ตจสารํว สมฺผลนฺ”ติ.}}

\prose{อยมฺปิ อตฺโถ วุตฺโต ภควตา อิติ เม สุตนฺติ.\csromanspacer{}\csromanspacer{}ปฐมํ.}

\csromanheader{3}{2. ธาตุสุตฺต}
\csromanheader{2}{51. วุตฺตํ เหตํ ภควตา วุตฺตมรหตาติ เม สุตํ\csromandash{}}
\prose{ติสฺโส อิมา ภิกฺขเว ธาตุโย.\csromanspacer{} กตมา ติสฺโส? รูปธาตุ, อรูปธาตุ, นิโรธธาตุ.\csromanspacer{} อิมา โข ภิกฺขเว ติสฺโส ธาตุโยติ.\csromanspacer{} เอตมตฺถํ ภควา อโวจ, ตตฺเถตํ อิติ วุจฺจติ\csromandash{}}

\csromangathameasure{“รูปธาตุํ ปริญฺญาย, อารุปฺเปสุ อสณฺฐิตา. \\ นิโรเธ เย วิมุจฺจนฺติ, เต ชนา มจฺจุหายิโน. \\ กาเยน อมตํ ธาตุํ, ผุสยิตฺวา นิรูปธิํ. \\ อุปธิปฺปฏินิสฺสคฺคํ, สจฺฉิกตฺวา อนาสโว. \\ เทเสติ สมฺมาสมฺพุทฺโธ, อโสกํ วิรชํ ปทนฺ”ติ.}
\csromangathagroup{\csromangathastanza{“รูปธาตุํ\footnote{รูปธาตุ (สพฺพตฺถ)} ปริญฺญาย, อารุปฺเปสุ อสณฺฐิตา. \\ นิโรเธ เย วิมุจฺจนฺติ, เต ชนา มจฺจุหายิโน.}\csromangathastanza{กาเยน อมตํ ธาตุํ, ผุสยิตฺวา\footnote{ผุสฺสยิตฺวา (สฺยา), ผสฺสยิตฺวา (อิ)} นิรูปธิํ. \\ อุปธิปฺปฏินิสฺสคฺคํ, สจฺฉิกตฺวา อนาสโว. \\ เทเสติ สมฺมาสมฺพุทฺโธ, อโสกํ วิรชํ ปทนฺ”ติ.}}

\prose{อยมฺปิ อตฺโถ วุตฺโต ภควตา อิติ เม สุตนฺติ.\csromanspacer{}\csromanspacer{}ทุติยํ.}

\csromanheader{2}{3. ติกนิปาต \textbf{3. ปฐมเวทนาสุตฺต}}
\csromanheader{2}{52. วุตฺตํ เหตํ ภควตา วุตฺตมรหตาติ เม สุตํ\csromandash{}}
\csromanmatikaanchor{mtk.151}
\csromanmatikaanchor{mtk.171}
\csromantocmarkref{subsection}{7. อาตาปีสุตฺต}{mtk.151}
\bookmark[dest={mtk.151},level=2]{7. อาตาปีสุตฺต}
\csromantocmarkref{subsection}{8-9. นกุหนสุตฺตาทิ}{mtk.152}
\bookmark[dest={mtk.152},level=2]{8-9. นกุหนสุตฺตาทิ}
\prose{\csromanfolio{227}ติสฺโส อิมา ภิกฺขเว เวทนา.\csromanspacer{} กตมา ติสฺโส? สุขา เวทนา, ทุกฺขา เวทนา, อทุกฺขมสุขา เวทนา.\csromanspacer{} อิมา โข ภิกฺขเว ติสฺโส เวทนาติ.\csromanspacer{} เอตมตฺถํ ภควา อโวจ, ตตฺเถตํ อิติ วุจฺจติ\csromandash{}}

\csromangathameasure{“สมาหิโต สมฺปชาโน, สโต พุทฺธสฺส สาวโก. \\ เวทนา จ ปชานาติ, เวทนานญฺจ สมฺภวํ. \\ ยตฺถ เจตา นิรุชฺฌนฺติ, มคฺคญฺจ ขยคามินํ. \\ เวทนานํ ขยา ภิกฺขุ, นิจฺฉาโต ปรินิพฺพุโต”ติ.}
\csromangathagroup{\csromangathastanza{“สมาหิโต สมฺปชาโน, สโต พุทฺธสฺส สาวโก. \\ เวทนา จ ปชานาติ, เวทนานญฺจ สมฺภวํ.}\csromangathastanza{ยตฺถ เจตา นิรุชฺฌนฺติ, มคฺคญฺจ ขยคามินํ. \\ เวทนานํ ขยา ภิกฺขุ, นิจฺฉาโต ปรินิพฺพุโต”ติ.}}

\prose{อยมฺปิ อตฺโถ วุตฺโต ภควตา อิติ เม สุตนฺติ.\csromanspacer{}\csromanspacer{}ตติยํ.}

\csromanheader{2}{4. ทุติยเวทนาสุตฺต}
\csromanheader{2}{53. วุตฺตํ เหตํ ภควตา วุตฺตมรหตาติ เม สุตํ\csromandash{}}
\prose{ติสฺโส อิมา ภิกฺขเว เวทนา.\csromanspacer{} กตมา ติสฺโส? สุขา เวทนา, ทุกฺขา เวทนา, อทุกฺขมสุขา เวทนา.\csromanspacer{} สุขา ภิกฺขเว เวทนา ทุกฺขโต ทฏฺฐพฺพา, ทุกฺขา เวทนา สลฺลโต ทฏฺฐพฺพา, อทุกฺขมสุขา เวทนา อนิจฺจโต ทฏฺฐพฺพา.\csromanspacer{} ยโต โข ภิกฺขเว ภิกฺขุโน สุขา เวทนา ทุกฺขโต ทิฏฺฐา โหติ, ทุกฺขา เวทนา สลฺลโต ทิฏฺฐา โหติ, อทุกฺขมสุขา เวทนา อนิจฺจโต ทิฏฺฐา โหติ.\csromanspacer{} อยํ วุจฺจติ ภิกฺขเว ภิกฺขุ อริโย สมฺมทฺทโส อจฺเฉจฺฉิ\footnote{อจฺเฉชฺชิ (สี, อิ), อจฺฉิชฺชิ (ก)} ตณฺหํ, วิวตฺตยิ\footnote{วาวตฺตยิ (สี-ฏฺฐ)} สํโยชนํ, สมฺมา มานาภิสมยา อนฺตมกาสิ ทุกฺขสฺสาติ.\csromanspacer{} เอตมตฺถํ ภควา อโวจ, ตตฺเถตํ อิติ วุจฺจติ\csromandash{}}

\csromangathameasure{*“โย สุขํ ทุกฺขโต อทฺท, ทุกฺขมทฺทกฺขิ สลฺลโต. \\ อทุกฺขมสุขํ สนฺตํ, อทฺทกฺขิ นํ อนิจฺจโต. \\ ส เว สมฺมทฺทโส ภิกฺขุ, ยโต ตตฺถ วิมุจฺจติ. \\ อภิญฺญาโวสิโต สนฺโต, ส เว โยคาติโค มุนี”ติ.}
\csromangathagroup{\csromangathastanza{\csromansymbolfootnote{*}{สํ 2. 409; ขุ 2. 344 ปิฏฺเฐสุปิ.}“โย สุขํ ทุกฺขโต อทฺท\footnote{ทกฺขิ (สี, อิ, ก), อทกฺขิ (สฺยา)}, ทุกฺขมทฺทกฺขิ สลฺลโต. \\ อทุกฺขมสุขํ สนฺตํ, อทฺทกฺขิ นํ อนิจฺจโต.}\csromangathastanza{ส เว สมฺมทฺทโส ภิกฺขุ, ยโต ตตฺถ วิมุจฺจติ. \\ อภิญฺญาโวสิโต สนฺโต, ส เว โยคาติโค มุนี”ติ.}}

\prose{อยมฺปิ อตฺโถ วุตฺโต ภควตา อิติ เม สุตนฺติ.\csromanspacer{}\csromanspacer{}จตุตฺถํ.}

\gambhira{อิติวุตฺตกปาฬิ}
\csromanheader{2}{5. ปฐม-เอสนาสุตฺต}
\csromanheader{2}{54. วุตฺตํ เหตํ ภควตา วุตฺตมรหตาติ เม สุตํ\csromandash{}}
\csromanmatikaanchor{mtk.153}
\csromanmatikaanchor{mtk.172}
\csromanmatikaanchor{mtk.173}
\csromantocmarkref{subsection}{10. โสมนสฺสสุตฺต}{mtk.153}
\bookmark[dest={mtk.153},level=2]{10. โสมนสฺสสุตฺต}
\prose{\csromanfolio{228}ติสฺโส อิมา ภิกฺขเว เอสนา.\csromanspacer{} กตมา ติสฺโส? กาเมสนา, ภเวสนา, พฺรหฺมจริเยสนา.\csromanspacer{} อิมา โข ภิกฺขเว ติสฺโส เอสนาติ.\csromanspacer{} เอตมตฺถํ ภควา อโวจ, ตตฺเถตํ อิติ วุจฺจติ\csromandash{}}

\csromangathameasure{“สมาหิโต สมฺปชาโน, สโต พุทฺธสฺส สาวโก. \\ เอสนา จ ปชานาติ, เอสนานญฺจ สมฺภวํ. \\ ยตฺถ เจตา นิรุชฺฌนฺติ, มคฺคญฺจ ขยคามินํ. \\ เอสนานํ ขยา ภิกฺขุ, นิจฺฉาโต ปรินิพฺพุโต”ติ.}
\csromangathagroup{\csromangathastanza{“สมาหิโต สมฺปชาโน, สโต พุทฺธสฺส สาวโก. \\ เอสนา จ ปชานาติ, เอสนานญฺจ สมฺภวํ.}\csromangathastanza{ยตฺถ เจตา นิรุชฺฌนฺติ, มคฺคญฺจ ขยคามินํ. \\ เอสนานํ ขยา ภิกฺขุ, นิจฺฉาโต ปรินิพฺพุโต”ติ.}}

\prose{อยมฺปิ อตฺโถ วุตฺโต ภควตา อิติ เม สุตนฺติ.\csromanspacer{}\csromanspacer{}ปญฺจมํ.}

\csromanheader{2}{6. ทุติย-เอสนาสุตฺต}
\csromanheader{2}{55. วุตฺตํ เหตํ ภควตา วุตฺตมรหตาติ เม สุตํ\csromandash{}}
\prose{ติสฺโส อิมา ภิกฺขเว เอสนา.\csromanspacer{} กตมา ติสฺโส? กาเมสนา, ภเวสนา, พฺรหฺมจริเยสนา.\csromanspacer{} อิมา โข ภิกฺขเว ติสฺโส เอสนาติ.\csromanspacer{} เอตมตฺถํ ภควา อโวจ, ตตฺเถตํ อิติ วุจฺจติ\csromandash{}}

\csromangathameasure{“กาเมสนา ภเวสนา, พฺรหฺมจริเยสนา สห. \\ อิติ สจฺจปรามาโส, ทิฏฺฐิฏฺฐานา สมุสฺสยา. \\ สพฺพราควิรตฺตสฺส, ตณฺหกฺขยวิมุตฺติโน. \\ เอสนา ปฏินิสฺสฏฺฐา, ทิฏฺฐิฏฺฐานา สมูหตา. \\ เอสนานํ ขยา ภิกฺขุ, นิราโส อกถํกถี”ติ.}
\csromangathagroup{\csromangathastanza{“กาเมสนา ภเวสนา, พฺรหฺมจริเยสนา สห. \\ อิติ สจฺจปรามาโส, ทิฏฺฐิฏฺฐานา สมุสฺสยา.}\csromangathastanza{สพฺพราควิรตฺตสฺส, ตณฺหกฺขยวิมุตฺติโน. \\ เอสนา ปฏินิสฺสฏฺฐา, ทิฏฺฐิฏฺฐานา สมูหตา. \\ เอสนานํ ขยา ภิกฺขุ, นิราโส อกถํกถี”ติ.}}

\prose{อยมฺปิ อตฺโถ วุตฺโต ภควตา อิติ เม สุตนฺติ.\csromanspacer{}\csromanspacer{}ฉฏฺฐํ.}

\csromanheader{2}{7. ปฐม-อาสวสุตฺต}
\csromanheader{2}{56. วุตฺตํ เหตํ ภควตา วุตฺตมรหตาติ เม สุตํ\csromandash{}}
\prose{ตโยเม ภิกฺขเว อาสวา.\csromanspacer{} กตเม ตโย? กามาสโว, ภวาสโว, อวิชฺชาสโว.\csromanspacer{} อิเม โข ภิกฺขเว ตโย อาสวาติ.\csromanspacer{} เอตมตฺถํ ภควา อโวจ, ตตฺเถตํ อิติ วุจฺจติ\csromandash{}}

\csromanheader{2}{3. ติกนิปาต}
\csromangathameasure{“สมาหิโต สมฺปชาโน, สโต พุทฺธสฺส สาวโก. \\ อาสเว จ ปชานาติ, อาสวานญฺจ สมฺภวํ. \\ ยตฺถ เจตา นิรุชฺฌนฺติ, มคฺคญฺจ ขยคามินํ. \\ อาสวานํ ขยา ภิกฺขุ, นิจฺฉาโต ปรินิพฺพุโต”ติ.}
\csromangathagroup{\csromangathastanza{\csromanfolio{229}“สมาหิโต สมฺปชาโน, สโต พุทฺธสฺส สาวโก. \\ อาสเว จ ปชานาติ, อาสวานญฺจ สมฺภวํ.}\csromangathastanza{ยตฺถ เจตา นิรุชฺฌนฺติ, มคฺคญฺจ ขยคามินํ. \\ อาสวานํ ขยา ภิกฺขุ, นิจฺฉาโต ปรินิพฺพุโต”ติ.}}

\prose{อยมฺปิ อตฺโถ วุตฺโต ภควตา อิติ เม สุตนฺติ.\csromanspacer{}\csromanspacer{}สตฺตมํ.}

\csromanheader{2}{8. ทุติย-อาสวสุตฺต}
\csromanmatikaanchor{mtk.154}
\csromanmatikaanchor{mtk.155}
\csromantocmarknopage{section}{2. ทุติยวคฺค}{mtk.154}
\bookmark[dest={mtk.154},level=1]{2. ทุติยวคฺค}
\csromantocmarkref{subsection}{1. วิตกฺกสุตฺต}{mtk.155}
\bookmark[dest={mtk.155},level=2]{1. วิตกฺกสุตฺต}
\csromanheader{2}{57. วุตฺตํ เหตํ ภควตา วุตฺตมรหตาติ เม สุตํ\csromandash{}}
\prose{ตโยเม ภิกฺขเว อาสวา.\csromanspacer{} กตเม ตโย? กามาสโว, ภวาสโว, อวิชฺชาสโว.\csromanspacer{} อิเม โข ภิกฺขเว ตโย อาสวาติ.\csromanspacer{} เอตมตฺถํ ภควา อโวจ, ตตฺเถตํ อิติ วุจฺจติ\csromandash{}}

\csromangathameasure{“ยสฺส กามาสโว ขีโณ, อวิชฺชา จ วิราชิตา. \\ ภวาสโว ปริกฺขีโณ, วิปฺปมุตฺโต นิรูปธิ. \\ ธาเรติ อนฺติมํ เทหํ, เชตฺวา มารํ สวาหินินฺ”ติ.}
\csromangathagroup{\csromangathastanza{“ยสฺส กามาสโว ขีโณ, อวิชฺชา จ วิราชิตา. \\ ภวาสโว ปริกฺขีโณ, วิปฺปมุตฺโต นิรูปธิ. \\ ธาเรติ อนฺติมํ เทหํ, เชตฺวา มารํ สวาหินินฺ”ติ\footnote{สวาหนนฺติ (พหูสุ)}.}}

\prose{อยมฺปิ อตฺโถ วุตฺโต ภควตา อิติ เม สุตนฺติ.\csromanspacer{}\csromanspacer{}อฏฺฐมํ.}

\csromanheader{3}{9. ตณฺหาสุตฺต}
\csromanheader{2}{58. วุตฺตํ เหตํ ภควตา วุตฺตมรหตาติ เม สุตํ\csromandash{}}
\prose{ติสฺโส อิมา ภิกฺขเว ตณฺหา.\csromanspacer{} กตมา ติสฺโส? กามกณฺหา, ภวตณฺหา, วิภวตณฺหา.\csromanspacer{} อิมา โข ภิกฺขเว ติสฺโส ตณฺหาติ.\csromanspacer{} เอตมตฺถํ ภควา อโวจ, ตตฺเถตํ อิติ วุจฺจติ\csromandash{}}

\csromangathameasure{“ตณฺหาโยเคน สํยุตฺตา, รตฺตจิตฺตา ภวาภเว. \\ เต โยคยุตฺตา มารสฺส, อโยคกฺเขมิโน ชนา. \\ สตฺตา คจฺฉนฺติ สํสารํ, ชาตีมรณคามิโน. \\ เย จ ตณฺหํ ปหนฺตฺวาน, วีตตณฺหา ภวาภเว. \\ เต เว ปารงฺคตา โลเก, เย ปตฺตา อาสวกฺขยนฺ”ติ.}
\csromangathagroup{\csromangathastanza{“ตณฺหาโยเคน สํยุตฺตา, รตฺตจิตฺตา ภวาภเว. \\ เต โยคยุตฺตา มารสฺส, อโยคกฺเขมิโน ชนา.}\csromangathastanza{สตฺตา คจฺฉนฺติ สํสารํ, ชาตีมรณคามิโน. \\ เย จ ตณฺหํ ปหนฺตฺวาน, วีตตณฺหา\footnote{นิตฺตณฺหา จ (สี, ก)} ภวาภเว. \\ เต เว\footnote{นิตฺตณฺหา จ (สี, ก)} ปารงฺคตา\footnote{นิตฺตณฺหา จ (สี, ก)} โลเก, เย ปตฺตา อาสวกฺขยนฺ”ติ.}}

\prose{อยมฺปิ อตฺโถ วุตฺโต ภควตา อิติ เม สุตนฺติ.\csromanspacer{}\csromanspacer{}นวมํ.}

\gambhira{อิติวุตฺตกปาฬิ}
\csromanheader{3}{10. มารเธยฺยสุตฺต}
\csromanmatikaanchor{mtk.156}
\csromantocmarkref{subsection}{2. เทสนาสุตฺต}{mtk.156}
\bookmark[dest={mtk.156},level=2]{2. เทสนาสุตฺต}
\csromanheader{2}{59. วุตฺตํ เหตํ ภควตา วุตฺตมรหตาติ เม สุตํ\csromandash{}}
\prose{\csromanfolio{230}ตีหิ ภิกฺขเว ธมฺเมหิ สมนฺนาคโต ภิกฺขุ อติกฺกมฺม มารเธยฺยํ อาทิจฺโจว วิโรจติ.\csromanspacer{} กตเมหิ ตีหิ? อิธ ภิกฺขเว ภิกฺขุ อเสเขน สีลกฺขนฺเธน สมนฺนาคโต โหติ, อเสเขน สมาธิกฺขนฺเธน สมนฺนาคโต โหติ, อเสเขน ปญฺญากฺขนฺเธน สมนฺนาคโต โหติ.\csromanspacer{} อิเมหิ โข ภิกฺขเว ตีหิ ธมฺเมหิ สมนฺนาคโต ภิกฺขุ อติกฺกมฺม มารเธยฺยํ อาทิจฺโจว วิโรจตีติ.\csromanspacer{} เอตมตฺถํ ภควา อโวจ, ตตฺเถตํ อิติ วุจฺจติ\csromandash{}}

\csromangathameasure{“สีลํ สมาธิ ปญฺญา จ, ยสฺส เอเต สุภาวิตา. \\ อติกฺกมฺม มารเธยฺยํ, อาทิจฺโจว วิโรจตี”ติ.}
\csromangathagroup{\csromangathastanza{“สีลํ สมาธิ ปญฺญา จ, ยสฺส เอเต สุภาวิตา. \\ อติกฺกมฺม มารเธยฺยํ, อาทิจฺโจว วิโรจตี”ติ.}}

\prose{อยมฺปิ อตฺโถ วุตฺโต ภควตา อิติ เม สุตนฺติ.\csromanspacer{}\csromanspacer{}ทสมํ.\csromanspacer{} ปฐโม วคฺโค.}

\titlehead{\textbf{ตสฺสุทฺทานํ}}
\csromangathameasure{มูลธาตุ อถ เวทนา ทุเว, \\ เอสนา จ ทุเว อาสวา ทุเว. \\ ตณฺหาโต จ อถ มารเธยฺยโต, \\ วคฺคมาหุ ปฐมนฺติ มุตฺตมนฺติ.}
\csromangathagroup{\csromangathastanza{มูลธาตุ อถ เวทนา ทุเว, \\ เอสนา จ ทุเว อาสวา ทุเว. \\ ตณฺหาโต จ อถ\footnote{ตณฺหาโต อถ (สฺยา)} มารเธยฺยโต, \\ วคฺคมาหุ ปฐมนฺติ มุตฺตมนฺติ.}}

\csromanmatikaanchor{mtk.157}
\csromantocmarkref{subsection}{3. วิชฺชาสุตฺต}{mtk.157}
\bookmark[dest={mtk.157},level=2]{3. วิชฺชาสุตฺต}
\chapterhead{2. ทุติยวคฺค}
\csromanheader{3}{1. ปุญฺญกิริยวตฺถุสุตฺต}
\csromanheader{2}{60. วุตฺตํ เหตํ ภควตา วุตฺตมรหตาติ เม สุตํ\csromandash{}}
\prose{\csromansymbolfootnote{*}{อํ 3. 72 ปิฏฺเฐปิ.}ตีณิมานิ ภิกฺขเว ปุญฺญกิริยวตฺถูนิ.\csromanspacer{} กตมานิ ตีณิ? ทานมยํ ปุญฺญกิริยวตฺถุ, สีลมยํ ปุญฺญกิริยวตฺถุ, ภาวนามยํ ปุญฺญกิริยวตฺถุ.}

\vspace{\tipitakaparskip}
\csromancenter{ติกนิปาต อิมานิ โข ภิกฺขเว ตีณิ ปุญฺญกิริยวตฺถูนีติ.\csromanspacer{} เอตมตฺถํ ภควา อโวจ, ตตฺเถตํ อิติ วุจฺจติ\csromandash{}}
\vspace{0.5\baselineskip}
\csromangathameasure{“ปุญฺญเมว โส สิกฺเขยฺย, อายตคฺคํ สุขุทฺรยํ. \\ ทานญฺจ สมจริยญฺจ, เมตฺตจิตฺตญฺจ ภาวเย. \\ เอเต ธมฺเม ภาวยิตฺวา, ตโย สุขสมุทฺทเย. \\ อพฺยาพชฺฌํ สุขํ โลกํ, ปณฺฑิโต อุปปชฺชตี”ติ.}
\csromangathagroup{\csromangathastanza{\csromanfolio{231}“ปุญฺญเมว โส สิกฺเขยฺย, อายตคฺคํ สุขุทฺรยํ. \\ ทานญฺจ สมจริยญฺจ, เมตฺตจิตฺตญฺจ ภาวเย.}\csromangathastanza{เอเต ธมฺเม ภาวยิตฺวา, ตโย สุขสมุทฺทเย. \\ อพฺยาพชฺฌํ\footnote{อํ 1. 353; อํ 2. 296 ปิฏฺเฐสุปิ.} สุขํ โลกํ, ปณฺฑิโต อุปปชฺชตี”ติ.}}

\csromanmatikaanchor{mtk.158}
\csromantocmarkref{subsection}{4. ปญฺญาปริหีนสุตฺต}{mtk.158}
\bookmark[dest={mtk.158},level=2]{4. ปญฺญาปริหีนสุตฺต}
\prose{อยมฺปิ อตฺโถ วุตฺโต ภควตา อิติ เม สุตนฺติ.\csromanspacer{}\csromanspacer{}ปฐมํ.}

\csromanheader{3}{2. จกฺขุสุตฺต}
\csromanheader{2}{61. วุตฺตํ เหตํ ภควตา วุตฺตมรหตาติ เม สุตํ\csromandash{}}
\prose{ตีณิมานิ ภิกฺขเว จกฺขูนิ.\csromanspacer{} กตมานิ ตีณิ? มํสจกฺขุ ทิพฺพจกฺขุ, ปญฺญาจกฺขุ.\csromanspacer{} อิมานิ โข ภิกฺขเว ตีณิ จกฺขูนีติ.\csromanspacer{} เอตมตฺถํ ภควา อโวจ, ตตฺเถตํ อิติ วุจฺจติ\csromandash{}}

\csromangathameasure{“มํสจกฺขุ ทิพฺพจกฺขุ, ปญฺญาจกฺขุ อนุตฺตรํ. \\ เอตานิ ตีณิ จกฺขูนิ, อกฺขาสิ ปุริสุตฺตโม. \\ มํสจกฺขุสฺส อุปฺปาโท, มคฺโค ทิพฺพสฺส จกฺขุโน. \\ ยโต ญาณํ อุทปาทิ, ปญฺญาจกฺขุ อนุตฺตรํ. \\ ยสฺส จกฺขุสฺส ปฏิลาภา, สพฺพทุกฺขา ปมุจฺจตี”ติ.}
\csromangathagroup{\csromangathastanza{“มํสจกฺขุ ทิพฺพจกฺขุ, ปญฺญาจกฺขุ อนุตฺตรํ. \\ เอตานิ ตีณิ จกฺขูนิ, อกฺขาสิ ปุริสุตฺตโม.}\csromangathastanza{มํสจกฺขุสฺส อุปฺปาโท, มคฺโค ทิพฺพสฺส จกฺขุโน. \\ ยโต ญาณํ อุทปาทิ, ปญฺญาจกฺขุ อนุตฺตรํ. \\ ยสฺส จกฺขุสฺส ปฏิลาภา, สพฺพทุกฺขา ปมุจฺจตี”ติ.}}

\prose{อยมฺปิ อตฺโถ วุตฺโต ภควตา อิติ เม สุตนฺติ.\csromanspacer{}\csromanspacer{}ทุติยํ.}

\csromanmatikaanchor{mtk.159}
\csromantocmarkref{subsection}{5. สุกฺกธมฺมสุตฺต}{mtk.159}
\bookmark[dest={mtk.159},level=2]{5. สุกฺกธมฺมสุตฺต}
\csromanheader{3}{3. อินฺทฺริยสุตฺต}
\csromanheader{2}{62. วุตฺตํ เหตํ ภควตา วุตฺตมรหตาติ เม สุตํ\csromandash{}}
\prose{ตีณิมานิ ภิกฺขเว อินฺทฺริยานิ.\csromanspacer{} กตมานิ ตีณิ? อนญฺญาตญฺญสฺสามีตินฺทฺริยํ, อญฺญินฺทฺริยํ, อญฺญาตาวินฺทฺริยํ.\csromanspacer{} อิมานิ โข ภิกฺขเว ตีณิ อินฺทฺริยานีติ.\csromanspacer{} เอตมตฺถํ ภควา อโวจ, ตตฺเถตํ อิติ วุจฺจติ\csromandash{}}

\csromangathameasure{“เสขสฺส สิกฺขมานสฺส, อุชุมคฺคานุสาริโน. \\ ขยสฺมิํ ปฐมํ ญาณํ, ตโต อญฺญา อนนฺตรา.}
\csromangathagroup{\csromangathastanza{“เสขสฺส สิกฺขมานสฺส, อุชุมคฺคานุสาริโน. \\ ขยสฺมิํ ปฐมํ ญาณํ, ตโต อญฺญา อนนฺตรา.}}

\gambhira{อิติวุตฺตกปาฬิ}
\csromangathameasure{ตโต อญฺญา วิมุตฺตสฺส, ญาณํ เว โหติ ตาทิโน.}
\csromangathagroup{\csromangathastanza{\csromanfolio{232}ตโต อญฺญา วิมุตฺตสฺส, ญาณํ เว โหติ ตาทิโน.}}

\vspace{\tipitakaparskip}
\csromancenter{“อกุปฺปา เม วิมุตฺตี”ติ, ภวสํโยชนกฺขยา.*}
\vspace{0.5\baselineskip}
\csromanmatikaanchor{mtk.160}
\csromantocmarkref{subsection}{6. อชาตสุตฺต}{mtk.160}
\bookmark[dest={mtk.160},level=2]{6. อชาตสุตฺต}
\csromangathameasure{ส เว อินฺทฺริยสมฺปนฺโน, สนฺโต สนฺติปเท รโต. \\ ธาเรติ อนฺติมํ เทหํ, เชตฺวา มารํ สวาหินินฺ”ติ.}
\csromangathagroup{\csromangathastanza{ส เว\footnote{สเจ (สี, สฺยา)} อินฺทฺริยสมฺปนฺโน, สนฺโต สนฺติปเท รโต. \\ ธาเรติ อนฺติมํ เทหํ, เชตฺวา มารํ สวาหินินฺ”ติ.}}

\prose{อยมฺปิ อตฺโถ วุตฺโต ภควตา อิติ เม สุตนฺติ.\csromanspacer{}\csromanspacer{}ตติยํ.}

\csromanheader{3}{4. อทฺธาสุตฺต}
\csromanheader{2}{63. วุตฺตํ เหตํ ภควตา วุตฺตมรหตาติ เม สุตํ\csromandash{}}
\prose{ตโยเม ภิกฺขเว อทฺธา.\csromanspacer{} กตเม ตโย? อตีโต อทฺธา, อนาคโต อทฺธา, ปจฺจุปฺปนฺโน อทฺธา.\csromanspacer{} อิเม โข ภิกฺขเว ตโย อทฺธาติ.\csromanspacer{} เอตมตฺถํ ภควา อโวจ, ตตฺเถตํ อิติ วุจฺจติ\csromandash{}}

\csromanmatikaanchor{mtk.161}
\csromantocmarkref{subsection}{7. นิพฺพานธาตุสุตฺต}{mtk.161}
\bookmark[dest={mtk.161},level=2]{7. นิพฺพานธาตุสุตฺต}
\csromangathameasure{“อกฺเขยฺยสญฺญิโน สตฺตา, อกฺเขยฺยสฺมิํ ปติฏฺฐิตา. \\ อกฺเขยฺยํ อปริญฺญาย, โยคมายนฺติ มจฺจุโน. \\ อกฺเขยฺยญฺจ ปริญฺญาย, อกฺขาตารํ น มญฺญติ. \\ ผุฏฺโฐ วิโมกฺโข มนสา, สนฺติปทมนุตฺตรํ. \\ ส เว อกฺเขยฺยสมฺปนฺโน, สนฺโต สนฺติปเท รโต. \\ สงฺขาย เสวี ธมฺมฏฺโฐ, สงฺขฺยํ โนเปติ เวทคู”ติ.}
\csromangathagroup{\csromangathastanza{“อกฺเขยฺยสญฺญิโน สตฺตา, อกฺเขยฺยสฺมิํ ปติฏฺฐิตา. \\ อกฺเขยฺยํ อปริญฺญาย, โยคมายนฺติ มจฺจุโน.}\csromangathastanza{อกฺเขยฺยญฺจ ปริญฺญาย, อกฺขาตารํ น มญฺญติ. \\ ผุฏฺโฐ วิโมกฺโข มนสา, สนฺติปทมนุตฺตรํ.}\csromangathastanza{ส เว\footnote{สเจ (ก)} อกฺเขยฺยสมฺปนฺโน, สนฺโต สนฺติปเท รโต. \\ สงฺขาย เสวี ธมฺมฏฺโฐ, สงฺขฺยํ โนเปติ เวทคู”ติ.}}

\prose{อยมฺปิ อตฺโถ วุตฺโต ภควตา อิติ เม สุตนฺติ.\csromanspacer{}\csromanspacer{}จตุตฺถํ.}

\csromanheader{3}{5. ทุจฺจริตสุตฺต}
\csromanheader{2}{64. วุตฺตํ เหตํ ภควตา วุตฺตมรหตาติ เม สุตํ\csromandash{}}
\prose{ตีณิมานิ ภิกฺขเว ทุจฺจริตานิ, กตมานิ ตีณิ? กายทุจฺจริตํ, วจีทุจฺจริตํ, มโนทุจฺจริตํ.\csromanspacer{} อิมานิ โข ภิกฺขเว ตีณิ ทุจฺจริตานีติ.\csromanspacer{} เอตมตฺถํ ภควา อโวจ, ตตฺเถตํ อิติ วุจฺจติ\csromandash{}}

\csromangathameasure{“กายทุจฺจริตํ กตฺวา, วจีทุจฺจริตานิ จ. \\ มโนทุจฺจริตํ กตฺวา, ยญฺจญฺญํ โทสสํหิตํ.}
\csromangathagroup{\csromangathastanza{“กายทุจฺจริตํ กตฺวา, วจีทุจฺจริตานิ จ. \\ มโนทุจฺจริตํ กตฺวา, ยญฺจญฺญํ โทสสํหิตํ.}}

\csromanheader{2}{3. ติกนิปาต}
\csromangathameasure{อกตฺวา กุสลํ กมฺมํ, กตฺวานา’กุสลํ พหุํ. \\ กายสฺส เภทา ทุปฺปญฺโญ, นิรยํ โสปปชฺชตี”ติ.}
\csromangathagroup{\csromangathastanza{\csromanfolio{233}อกตฺวา กุสลํ กมฺมํ, กตฺวานา’กุสลํ พหุํ. \\ กายสฺส เภทา ทุปฺปญฺโญ, นิรยํ โสปปชฺชตี”ติ.}}

\prose{อยมฺปิ อตฺโถ วุตฺโต ภควตา อิติ เม สุตนฺติ.\csromanspacer{}\csromanspacer{}ปญฺจมํ.}

\csromanmatikaanchor{mtk.162}
\csromantocmarkref{subsection}{8. ปฏิสลฺลานสุตฺต}{mtk.162}
\bookmark[dest={mtk.162},level=2]{8. ปฏิสลฺลานสุตฺต}
\csromanheader{3}{6. สุจริตสุตฺต}
\csromanheader{2}{65. วุตฺตํ เหตํ ภควตา วุตฺตมรหตาติ เม สุตํ\csromandash{}}
\prose{ตีณิมานิ ภิกฺขเว สุจริตานิ.\csromanspacer{} กตมานิ ตีณิ? กายสุจริตํ, วจีสุจริตํ, มโนสุจริตํ.\csromanspacer{} อิมานิ โข ภิกฺขเว ตีณิ สุจริตานีติ.\csromanspacer{} เอตมตฺถํ ภควา อโวจ, ตตฺเถตํ อิติ วุจฺจติ\csromandash{}}

\csromangathameasure{“กายทุจฺจริตํ หิตฺวา, วจีทุจฺจริตานิ จ. \\ มโนทุจฺจริตํ หิตฺวา, ยญฺจญฺญํ โทสสํหิตํ. \\ อกตฺวา’กุสลํ กมฺมํ, กตฺวาน กุสลํ พหุํ. \\ กายสฺส เภทา สปฺปญฺโญ, สคฺคํ โส อุปปชฺชตี”ติ.}
\csromangathagroup{\csromangathastanza{“กายทุจฺจริตํ หิตฺวา, วจีทุจฺจริตานิ จ. \\ มโนทุจฺจริตํ หิตฺวา, ยญฺจญฺญํ โทสสํหิตํ.}\csromangathastanza{อกตฺวา’กุสลํ กมฺมํ, กตฺวาน กุสลํ พหุํ. \\ กายสฺส เภทา สปฺปญฺโญ, สคฺคํ โส อุปปชฺชตี”ติ.}}

\prose{อยมฺปิ อตฺโถ วุตฺโต ภควตา อิติ เม สุตนฺติ.\csromanspacer{}\csromanspacer{}ฉฏฺฐํ.}

\csromanheader{3}{7. โสเจยฺยสุตฺต}
\csromanmatikaanchor{mtk.163}
\csromantocmarkref{subsection}{9. สิกฺขานิสํสสุตฺต}{mtk.163}
\bookmark[dest={mtk.163},level=2]{9. สิกฺขานิสํสสุตฺต}
\csromanheader{2}{66. วุตฺตํ เหตํ ภควตา วุตฺตมรหตาติ เม สุตํ\csromandash{}}
\prose{ตีณิมานิ ภิกฺขเว โสเจยฺยานิ.\csromanspacer{} กตมานิ ตีณิ? กายโสเจยฺยํ, วจีโสเจยฺยํ, มโนโสเจยฺยํ.\csromanspacer{} อิมานิ โข ภิกฺขเว ตีณิ โสเจยฺยานีติ.\csromanspacer{} เอตมตฺถํ ภควา อโวจ, ตตฺเถตํ อิติ วุจฺจติ\csromandash{}}

\csromangathameasure{“กายสุจิํ วจีสุจิํ, เจโตสุจิมนาสวํ. \\ สุจิํ โสเจยฺยสมฺปนฺนํ, อาหุ สพฺพปฺปหายินนฺ”ติ.}
\csromangathagroup{\csromangathastanza{“กายสุจิํ วจีสุจิํ\footnote{วาจาสุจิํ (ก)}, เจโตสุจิมนาสวํ. \\ สุจิํ โสเจยฺยสมฺปนฺนํ, อาหุ สพฺพปฺปหายินนฺ”ติ\footnote{วาจาสุจิํ (ก)}.}}

\prose{อยมฺปิ อตฺโถ วุตฺโต ภควตา อิติ เม สุตนฺติ.\csromanspacer{}\csromanspacer{}สตฺตมํ.}

\gambhira{อิติวุตฺตกปาฬิ}
\csromanheader{3}{8. โมเนยฺยสุตฺต}
\csromanmatikaanchor{mtk.164}
\csromantocmarkref{subsection}{10. ชาคริยสุตฺต}{mtk.164}
\bookmark[dest={mtk.164},level=2]{10. ชาคริยสุตฺต}
\csromanheader{2}{67. วุตฺตํ เหตํ ภควตา วุตฺตมรหตาติ เม สุตํ\csromandash{}}
\csromanmatikaanchor{mtk.165}
\csromanmatikaanchor{mtk.185}
\csromantocmarkref{subsection}{11. อาปายิกสุตฺต}{mtk.165}
\bookmark[dest={mtk.165},level=2]{11. อาปายิกสุตฺต}
\prose{\csromanfolio{234}ตีณิมานิ ภิกฺขเว โมเนยฺยานิ.\csromanspacer{} กตมานิ ตีณิ? กายโมเนยฺยํ, วจีโมเนยฺยํ, มโนโมเนยฺยํ.\csromanspacer{} อิมานิ โข ภิกฺขเว ตีณิ โมเนยฺยานีติ.\csromanspacer{} เอตมตฺถํ ภควา อโวจ, ตตฺเถตํ อิติ วุจฺจติ\csromandash{}}

\csromangathameasure{“กายมุนิํ วจีมุนิํ, มโนมุนิมนาสวํ. \\ มุนิํ โมเนยฺยสมฺปนฺนํ, อาหุ นินฺหาตปาปกนฺ”ติ.}
\csromangathagroup{\csromangathastanza{“กายมุนิํ วจีมุนิํ, มโนมุนิมนาสวํ. \\ มุนิํ โมเนยฺยสมฺปนฺนํ, อาหุ นินฺหาตปาปกนฺ”ติ\footnote{อาหุ สพฺพปฺปหายินนฺติ (อํ 1. 277; ขุ 8. 73)}.}}

\prose{อยมฺปิ อตฺโถ วุตฺโต ภควตา อิติ เม สุตนฺติ.\csromanspacer{}\csromanspacer{}อฏฺฐมํ.}

\csromanheader{2}{9. ปฐมราคสุตฺต}
\csromanheader{2}{68. วุตฺตํ เหตํ ภควตา วุตฺตมรหตาติ เม สุตํ\csromandash{}}
\prose{ยสฺส กสฺสจิ ภิกฺขเว ราโค อปฺปหีโน โทโส อปฺปหีโน โมโห อปฺปหีโน, อยํ วุจฺจติ ภิกฺขเว พทฺโธ\footnote{พนฺโธ (พหูสุ)} มารสฺส, ปฏิมุกฺก’สฺส มารปาโส ยถากามกรณีโย\footnote{ยถา กามกรณีโย จ (สี, สฺยา, อิ, ก)} ปาปิมโต.\csromanspacer{} ยสฺส กสฺสจิ ภิกฺขเว ราโค ปหีโน โทโส ปหีโน โมโห ปหีโน, อยํ วุจฺจติ ภิกฺขเว อพทฺโธ มารสฺส, โอมุกฺก’สฺส มารปาโส น ยถา กามกรณีโย\footnote{น ยถากามกรณีโย จ (สฺยา)} ปาปิมโตติ.\csromanspacer{} เอตมตฺถํ ภควา อโวจ, ตตฺเถตํ อิติ วุจฺจติ\csromandash{}}

\csromangathameasure{“ยสฺส ราโค จ โทโส จ, อวิชฺชา จ วิราชิตา. \\ ตํ ภาวิตตฺตญฺญตรํ, พฺรหฺมภูตํ ตถาคตํ. \\ พุทฺธํ เวรภยาตีตํ, อาหุ สพฺพปฺปหายินนฺ”ติ.}
\csromangathagroup{\csromangathastanza{“ยสฺส ราโค จ โทโส จ, อวิชฺชา จ วิราชิตา. \\ ตํ ภาวิตตฺตญฺญตรํ, พฺรหฺมภูตํ ตถาคตํ. \\ พุทฺธํ เวรภยาตีตํ, อาหุ สพฺพปฺปหายินนฺ”ติ.}}

\prose{อยมฺปิ อตฺโถ วุตฺโต ภควตา อิติ เม สุตนฺติ.\csromanspacer{}\csromanspacer{}นวมํ.}

\csromanheader{2}{10. ทุติยราคสุตฺต}
\csromanheader{2}{69. วุตฺตํ เหตํ ภควตา วุตฺตมรหตาติ เม สุตํ\csromandash{}}
\prose{ยสฺส กสฺสจิ ภิกฺขเว ภิกฺขุสฺส วา ภิกฺขุนิยา วา ราโค อปฺปหีโน โทโส อปฺปหีโน โมโห อปฺปหีโน, อยํ วุจฺจติ ภิกฺขเว}

\vspace{\tipitakaparskip}
\csromancenter{ติกนิปาต อติณฺโณ\footnote{อตริ (ก)} สมุทฺทํ ส-อูมิํ สวีจิํ สาวฏฺฏํ สคหํ สรกฺขสํ.\csromanspacer{} ยสฺส กสฺสจิ ภิกฺขเว ภิกฺขุสฺส วา ภิกฺขุนิยา วา ราโค ปหีโน โทโส ปหีโน โมโห ปหีโน, อยํ วุจฺจติ ภิกฺขเว อตริ สมุทฺทํ ส-อูมิํ สวีจิํ สาวฏฺฏํ สคหํ สรกฺขสํ, ติณฺโณ ปารงฺคโต\footnote{ปารคโต (สี-ฏฺฐ, สฺยา)} ถเล ติฏฺฐติ พฺราหฺมโณติ.\csromanspacer{} เอตมตฺถํ ภควา อโวจ, ตตฺเถตํ อิติ วุจฺจติ\csromandash{}}
\vspace{0.5\baselineskip}
\csromangathameasure{“ยสฺส ราโค จ โทโส จ, อวิชฺชา จ วิราชิตา. \\ โสมํ สมุทฺทํ สคหํ สรกฺขสํ, \\ ส-อูมิภยํ ทุตฺตรํ อจฺจตาริ. \\ สงฺคาติโค มจฺจุชโห นิรูปธิ, \\ ปหาสิ ทุกฺขํ อปุนพฺภวาย. \\ อตฺถงฺคโต โส น ปมาณเมติ, \\ อโมหยิ มจฺจุราชนฺติ พฺรูมี”ติ.}
\csromangathagroup{\csromangathastanza{\csromanfolio{235}“ยสฺส ราโค จ โทโส จ, อวิชฺชา จ วิราชิตา.}\csromangathastanza{โสมํ สมุทฺทํ สคหํ สรกฺขสํ, \\ ส-อูมิภยํ ทุตฺตรํ อจฺจตาริ. \\ สงฺคาติโค มจฺจุชโห นิรูปธิ, \\ ปหาสิ ทุกฺขํ อปุนพฺภวาย. \\ อตฺถงฺคโต โส น ปมาณเมติ, \\ อโมหยิ มจฺจุราชนฺติ พฺรูมี”ติ.}}

\csromanmatikaanchor{mtk.166}
\csromantocmarkref{subsection}{12. ทิฏฺฐิคตสุตฺต}{mtk.166}
\bookmark[dest={mtk.166},level=2]{12. ทิฏฺฐิคตสุตฺต}
\prose{อยมฺปิ อตฺโถ วุตฺโต ภควตา อิติ เม สุตนฺติ.\csromanspacer{}\csromanspacer{}ทสมํ.\csromanspacer{} ทุติโย วคฺโค.}

\titlehead{\textbf{ตสฺสุทฺทานํ}}
\csromangathameasure{ปุญฺญํ จกฺขุ อถ อินฺทฺริยานิ, อทฺธา จ จริตํ ทุเว โสจิ. \\ มุโน อถ ราคทุเว, ปุน วคฺคมาหุ ทุติยมุตฺตมนฺ”ติ.}
\csromangathagroup{\csromangathastanza{ปุญฺญํ จกฺขุ อถ อินฺทฺริยานิ\footnote{อตฺถินฺทฺริยา (สฺยา)}, อทฺธา จ จริตํ ทุเว โสจิ\footnote{สุจิ (สฺยา)}. \\ มุโน\footnote{อตฺถินฺทฺริยา (สฺยา)} อถ ราคทุเว, ปุน วคฺคมาหุ ทุติยมุตฺตมนฺ”ติ.}}

\chapterhead{3. ตติยวคฺค}
\csromanheader{3}{1. มิจฺฉาทิฏฺฐิกสุตฺต}
\csromanheader{2}{70. วุตฺตํ เหตํ ภควตา วุตฺตมรหตาติ เม สุตํ\csromandash{}}
\prose{ทิฏฺฐา มยา ภิกฺขเว สตฺตา กายทุจฺจริเตน สมนฺนาคตา วจีทุจฺจริเตน สมนฺนาคตา มโนทุจฺจริเตน สมนฺนาคตา อริยานํ อุปวาทกา มิจฺฉาทิฏฺฐิกา มิจฺฉาทิฏฺฐิกมฺมสมาทานา, เต กายสฺส เภทา ปรํ มรณา อปายํ ทุคฺคติํ วินิปาตํ นิรยํ อุปปนฺนา.}

\gambhira{อิติวุตฺตกปาฬิ}
\prose{\csromanfolio{236}ตํ โข ปนาหํ ภิกฺขเว นาญฺญสฺส สมณสฺส วา พฺราหฺมณสฺส วา สุตฺวา วทามิ.\csromanspacer{} ทิฏฺฐา มยา ภิกฺขเว สตฺตา กายทุจฺจริเตน สมนฺนาคตา วจีทุจฺจริเตน สมนฺนาคตา มโนทุจฺจริเตน สมนฺนาคตา อริยานํ อุปวาทกา มิจฺฉาทิฏฺฐิกา มิจฺฉาทิฏฺฐิกมฺมสมาทานา, เต กายสฺส เภทา ปรํ มรณา อปายํ ทุคฺคติํ วินิปาตํ นิรยํ อุปปนฺนา.\csromanspacer{} อปิ จ ภิกฺขเว ยเทว สามํ ญาตํ สามํ ทิฏฺฐํ สามํ วิทิตํ, ตเทวาหํ วทามิ.}

\prose{ทิฏฺฐา มยา ภิกฺขเว สตฺตา กายทุจฺจริเตน สมนฺนาคตา วจีทุจฺจริเตน สมนฺนาคตา มโนทุจฺจริเตน สมนฺนาคตา อริยานํ อุปวาทกา มิจฺฉาทิฏฺฐิกา มิจฺฉาทิฏฺฐิกมฺมสมาทานา, เต กายสฺส เภทา ปรํ มรณา อปายํ ทุคฺคติํ วินิปาตํ นิรยํ อุปปนฺนาติ.\csromanspacer{} เอตมตฺถํ ภควา อโวจ, ตตฺเถตํ อิติ วุจฺจติ\csromandash{}}

\csromangathameasure{“มิจฺฉา มนํ ปณิธาย, มิจฺฉา วาจญฺจ ภาสิย. \\ มิจฺฉา กมฺมานิ กตฺวาน, กาเยน อิธ ปุคฺคโล. \\ อปฺปสฺสุตา’ปุญฺญกโร, อปฺปสฺมิํ อิธ ชีวิเต. \\ กายสฺส เภทา ทุปฺปญฺโญ, นิรยํ โสปปชฺชตี”ติ.}
\csromangathagroup{\csromangathastanza{“มิจฺฉา มนํ ปณิธาย, มิจฺฉา วาจญฺจ ภาสิย\footnote{มิจฺฉา วาจํ อภาสิย (สพฺพตฺถ)}. \\ มิจฺฉา กมฺมานิ กตฺวาน, กาเยน อิธ ปุคฺคโล.}\csromangathastanza{อปฺปสฺสุตา’ปุญฺญกโร\footnote{อปฺปสฺสุโตปุญฺญกโร (สี), อปฺปสฺสุโต อปุญฺญกโร (สฺยา, อิ)}, อปฺปสฺมิํ อิธ ชีวิเต. \\ กายสฺส เภทา ทุปฺปญฺโญ, นิรยํ โสปปชฺชตี”ติ.}}

\prose{อยมฺปิ อตฺโถ วุตฺโต ภควตา อิติ เม สุตนฺติ.\csromanspacer{}\csromanspacer{}ปฐมํ.}

\csromanheader{3}{2. สมฺมาทิฏฺฐิกสุตฺต}
\csromanmatikaanchor{mtk.167}
\csromantocmarknopage{section}{3. ติกนิปาต}{mtk.167}
\bookmark[dest={mtk.167},level=1]{3. ติกนิปาต}
\csromanheader{2}{71. วุตฺตํ เหตํ ภควตา วุตฺตมรหตาติ เม สุตํ\csromandash{}}
\prose{ทิฏฺฐา มยา ภิกฺขเว สตฺตา กายสุจริเตน สมนฺนาคตา วจีสุจริเตน สมนฺนาคตา มโนสุจริเตน สมนฺนาคตา อริยานํ อนุปวาทกา สมฺมาทิฏฺฐิกา สเมทิฏฺฐิกมฺมสมาทานา, เต กายสฺส เภทา ปรํ มรณา สุคติํ สคฺคํ โลกํ อุปปนฺนา.}

\csromanmatikaanchor{mtk.168}
\csromanmatikaanchor{mtk.169}
\csromantocmarknopage{subsection}{1. ปฐมวคฺค}{mtk.168}
\bookmark[dest={mtk.168},level=2]{1. ปฐมวคฺค}
\csromantocmarkref{subsubsection}{1. มูลสุตฺต}{mtk.169}
\bookmark[dest={mtk.169},level=3]{1. มูลสุตฺต}
\prose{ตํ โข ปนาหํ ภิกฺขเว นาญฺญสฺส สมณสฺส วา พฺราหฺมณสฺส วา สุตฺวา วทามิ.\csromanspacer{} ทิฏฺฐา มยา ภิกฺขเว สตฺตา กายสุจริเตน สมนฺนาคตา วจีสุจริเตน สมนฺนาคตา มโนสุจริเตน สมนฺนาคตา อริยานํ อนุปวาทกา สมฺมาทิฏฺฐิกา สเมทิฏฺฐิกมฺมสมาทานา, เต กายสฺส เภทา}

\vspace{\tipitakaparskip}
\csromancenter{ติกนิปาต ปรํ มรณา สุคติํ สคฺคํ โลกํ อุปปนฺนา.\csromanspacer{} อปิ จ ภิกฺขเว ยเทว สามํ ญาตํ สามํ ทิฏฺฐํ สามํ วิทิตํ, ตเทวาหํ วทามิ.}
\prose{\csromanfolio{237}ทิฏฺฐา มยา ภิกฺขเว สตฺตา กายสุจริเตน สมนฺนาคตา วจีสุจริเตน สมนฺนาคตา มโนสุจริเตน สมนฺนาคตา อริยานํ อนุปวาทกา สมฺมาทิฏฺฐิกา สมฺมาทิฏฺฐิกมฺมสมาทานา, เต กายสฺส เภทา ปรํ มรณา สุคติํ สคฺคํ โลกํ อุปปนฺนาติ.\csromanspacer{} เอตมตฺถํ ภควา อโวจ, ตตฺเถตํ อิติ วุจฺจติ\csromandash{}}

\csromangathameasure{“สมฺมา มนํ ปณิธาย, สมฺมา วาจญฺจ ภาสิย. \\ สมฺมา กมฺมานิ กตฺวาน, กาเยน อิธ ปุคฺคโล. \\ พหุสฺสุโต ปุญฺญกโร, อปฺปสฺมิํ อิธ ชีวิเต. \\ กายสฺส เภทา สปฺปญฺโญ, สคฺคํ โส อุปปชฺชตี”ติ.}
\csromangathagroup{\csromangathastanza{“สมฺมา มนํ ปณิธาย, สมฺมา วาจญฺจ ภาสิย\footnote{สมฺมา วาจํ อภาสิย (สพฺพตฺถ)}. \\ สมฺมา กมฺมานิ กตฺวาน, กาเยน อิธ ปุคฺคโล.}\csromangathastanza{พหุสฺสุโต ปุญฺญกโร, อปฺปสฺมิํ อิธ ชีวิเต. \\ กายสฺส เภทา สปฺปญฺโญ, สคฺคํ โส อุปปชฺชตี”ติ.}}

\csromanmatikaanchor{mtk.170}
\csromantocmarkref{subsubsection}{2. ธาตุสุตฺต}{mtk.170}
\bookmark[dest={mtk.170},level=3]{2. ธาตุสุตฺต}
\csromantocmarkref{subsubsection}{3-4. เวทนาสุตฺตาทิ}{mtk.171}
\bookmark[dest={mtk.171},level=3]{3-4. เวทนาสุตฺตาทิ}
\csromantocmarkref{subsubsection}{5-6. เอสนาสุตฺตาทิ}{mtk.172}
\bookmark[dest={mtk.172},level=3]{5-6. เอสนาสุตฺตาทิ}
\csromantocmarkref{subsubsection}{7-8. อาสวสุตฺตาทิ}{mtk.173}
\bookmark[dest={mtk.173},level=3]{7-8. อาสวสุตฺตาทิ}
\prose{อยมฺปิ อตฺโถ วุตฺโต ภควตา อิติ เม สุตนฺติ.\csromanspacer{}\csromanspacer{}ทุติยํ.}

\csromanheader{3}{3. นิสฺสรณิยสุตฺต}
\csromanheader{2}{72. วุตฺตํ เหตํ ภควตา วุตฺตมรหตาติ เม สุตํ\csromandash{}}
\prose{ติสฺโส อิมา ภิกฺขเว นิสฺสรณิยา\footnote{นิสฺสารณียา (อํ 2. 214)} ธาตุโย.\csromanspacer{} กตมา ติสฺโส? กามานเมตํ นิสฺสรณํ ยทิทํ เนกฺขมฺมํ, รูปานเมตํ นิสฺสรณํ ยทิทํ อารุปฺปํ, ยํ โข ปน กิญฺจิ ภูตํ สงฺขตํ ปฏิจฺจสมุปฺปนฺนํ, นิโรโธ ตสฺส นิสฺสรณํ.\csromanspacer{} อิมา โข ภิกฺขเว ติสฺโส นิสฺสรณิยา ธาตุโยติ.\csromanspacer{} เอตมตฺถํ ภควา อโวจ, ตตฺเถตํ อิติ วุจฺจติ\csromandash{}}

\csromangathameasure{“กามนิสฺสรณํ ญตฺวา, รูปานญฺจ อติกฺกมํ. \\ สพฺพสงฺขารสมถํ, ผุสํ อาตาปิ สพฺพทา. \\ ส เว สมฺมทฺทโส ภิกฺขุ, ยโต ตตฺถ วิมุจฺจติ. \\ อภิญฺญาโวสิโต สนฺโต, ส เว โยคาติโค มุนี”ติ.}
\csromangathagroup{\csromangathastanza{“กามนิสฺสรณํ ญตฺวา, รูปานญฺจ อติกฺกมํ. \\ สพฺพสงฺขารสมถํ, ผุสํ อาตาปิ สพฺพทา.}\csromangathastanza{ส เว สมฺมทฺทโส ภิกฺขุ, ยโต ตตฺถ วิมุจฺจติ. \\ อภิญฺญาโวสิโต สนฺโต, ส เว โยคาติโค มุนี”ติ.}}

\prose{อยมฺปิ อตฺโถ วุตฺโต ภควตา อิติ เม สุตนฺติ.\csromanspacer{}\csromanspacer{}ตติยํ.}

\gambhira{อิติวุตฺตกปาฬิ}
\csromanheader{3}{4. สนฺตตรสุตฺต}
\csromanheader{2}{73. วุตฺตํ เหตํ ภควตา วุตฺตมรหตาติ เม สุตํ\csromandash{}}
\prose{\csromanfolio{238}รูเปหิ ภิกฺขเว อรูปา\footnote{อารุปฺปา (สี)} สนฺตตรา, อรูเปหิ นิโรโธ สนฺตตโรติ.\csromanspacer{} เอตมตฺถํ ภควา อโวจ, ตตฺเถตํ อิติ วุจฺจติ\csromandash{}}

\csromangathameasure{“เย จ รูปูปคา สตฺตา, เย จ อรูปฏฺฐายิโน. \\ นิโรธํ อปฺปชานนฺตา, อาคนฺตาโร ปุนพฺภวํ. \\ เย จ รูเป ปริญฺญาย, อรูเปสุ อสณฺฐิตา. \\ นิโรเธ เย วิมุจฺจนฺติ, เต ชนา มจฺจุหายิโน. \\ กาเยน อมตธาตุํ, ผุสยิตฺวา นิรูปธิํ. \\ อุปธิปฺปฏินิสฺสคฺคํ, สจฺฉิกตฺวา อนาสโว. \\ เทเสติ สมฺมาสมฺพุทฺโธ, อโสกํ วิรชํ ปทนฺ”ติ.}
\csromangathagroup{\csromangathastanza{“เย จ รูปูปคา สตฺตา, เย จ อรูปฏฺฐายิโน\footnote{อารุปฺปฏฺฐายิโน (สี)}. \\ นิโรธํ อปฺปชานนฺตา, อาคนฺตาโร ปุนพฺภวํ.}\csromangathastanza{เย จ รูเป ปริญฺญาย, อรูเปสุ อสณฺฐิตา. \\ นิโรเธ เย วิมุจฺจนฺติ, เต ชนา มจฺจุหายิโน.}\csromangathastanza{กาเยน อมตธาตุํ, ผุสยิตฺวา นิรูปธิํ. \\ อุปธิปฺปฏินิสฺสคฺคํ, สจฺฉิกตฺวา อนาสโว. \\ เทเสติ สมฺมาสมฺพุทฺโธ, อโสกํ วิรชํ ปทนฺ”ติ.}}

\prose{อยมฺปิ อตฺโถ วุตฺโต ภควตา อิติ เม สุตนฺติ.\csromanspacer{}\csromanspacer{}จตุตฺถํ.}

\csromanheader{3}{5. ปุตฺตสุตฺต}
\csromanheader{2}{74. วุตฺตํ เหตํ ภควตา วุตฺตมรหตาติ เม สุตํ\csromandash{}}
\prose{ตโยเม ภิกฺขเว ปุตฺตา สนฺโต สํวิชฺชมานา โลกสฺมิํ.\csromanspacer{} กตเม ตโย? อติชาโต อนุชาโต อวชาโตติ.}

\prose{กถญฺจ ภิกฺขเว ปุตฺโต อติชาโต โหติ? อิธ ภิกฺขเว ปุตฺตสฺส มาตาปิตโร โหนฺติ\csromandash{}น พุทฺธํ สรณํ คตา, น ธมฺมํ สรณํ คตา, น สํฆํ สรณํ คตา, ปาณาติปาตา อปฺปฏิวิรตา, อทินฺนาทานา อปฺปฏิวิรตา, กาเมสุมิจฺฉาจารา อปฺปฏิวิรตา, มุสาวาทา อปฺปฏิวิรตา, สุราเมรยมชฺชปมาทฏฺฐานา อปฺปฏิวิรตา, ทุสฺสีลา ปาปธมฺมา.\csromanspacer{} ปุตฺโต จ เนสํ โหติ\csromandash{}พุทฺธํ สรณํ คโต, ธมฺมํ สรณํ คโต, สํฆํ สรณํ คโต, ปาณาติปาตา ปฏิวิรโต, อทินฺนาทานา ปฏิวิรโต, กาเมสุมิจฺฉาจารา ปฏิวิรโต, มุสาวาทา ปฏิวิรโต, สุราเมรยมชฺชปมาทฏฺฐานา ปฏิวิรโต, สีลวา กลฺยาณธมฺโม.\csromanspacer{} เอวํ โข ภิกฺขเว ปุตฺโต อติชาโต โหติ.}

\csromanheader{2}{3. ติกนิปาต}
\prose{\csromanfolio{239}กถญฺจ ภิกฺขเว ปุตฺโต อนุชาโต โหติ? อิธ ภิกฺขเว ปุตฺตสฺส มาตาปิตโร โหนฺติ\csromandash{}พุทฺธํ สรณํ คตา, ธมฺมํ สรณํ คตา, สํฆํ สรณํ คตา, ปาณาติปาตา ปฏิวิรตา, อทินฺนาทานา ปฏิวิรตา, กาเมสุมิจฺฉาจารา ปฏิวิรตา, มุสาวาทา ปฏิวิรตา, สุราเมรยมชฺชปมาทฏฺฐานา ปฏิวิรตา, สีลวนฺโต กลฺยาณธมฺมา.\csromanspacer{} ปุตฺโตปิ เนสํ โหติ\csromandash{}พุทฺธํ สรณํ คโต, ธมฺมํ สรณํ คโต, สํฆํ สรณํ คโต, ปาณาติปาตา ปฏิวิรโต, อทินฺนาทานา ปฏิวิรโต, กาเมสุมิจฺฉาจารา ปฏิวิรโต, มุเสวาทา ปฏิวิรโต, สุราเมรยมชฺชปมาทฏฺฐานา ปฏิวิรโต, สีลวา กลฺยาณธมฺโม.\csromanspacer{} เอวํ โข ภิกฺขเว ปุตฺโต อนุชาโต โหติ.}

\prose{กถญฺจ ภิกฺขเว ปุตฺโต อวชาโต โหติ? อิธ ภิกฺขเว ปุตฺตสฺส มาตาปิตโร โหนฺติ\csromandash{}พุทฺธํ สรณํ คตา, ธมฺมํ สรณํ คตา, สํฆํ สรณํ คตา, ปาณาติปาตา ปฏิวิรตา, อทินฺนาทานา ปฏิวิรตา, กาเมสุมิจฺฉาจารา ปฏิวิรตา, มุสาวาทา ปฏิวิรตา, สุราเมรยมชฺชปมาทฏฺฐานา ปฏิวิรตา, สีลวนฺโต กลฺยาณธมฺมา.\csromanspacer{} ปุตฺโต จ เนสํ โหติ\csromandash{}น พุทฺธํ สรณํ คโต, น ธมฺมํ สรณํ คโต, น สํฆํ สรณํ คโต, ปาณาติปาตา อปฺปฏิวิรโต, อทินฺนาทานา อปฺปฏิวิรโต, กาเมสุมิจฺฉาจารา อปฺปฏิวิรโต, มุสาวาทา อปฺปฏิวิรโต, สุราเมรยมชฺชปมาทฏฺฐานา อปฺปฏิวิรโต, ทุสฺสีโล ปาปธมฺโม.\csromanspacer{} เอวํ โข ภิกฺขเว ปุตฺโต อวชาโต โหติ.\csromanspacer{} อิเม โข ภิกฺขเว ตโย ปุตฺตา สนฺโต สํวิชฺชมานา โลกสฺมินฺติ.\csromanspacer{} เอตมตฺถํ ภควา อโวจ, ตตฺเถตํ อิติ วุจฺจติ\csromandash{}}

\csromangathameasure{“อติชาตํ อนุชาตํ, ปุตฺตมิจฺฉนฺติ ปณฺฑิตา. \\ อวชาตํ น อิจฺฉนฺติ, โย โหติ กุลคนฺธโน. \\ เอเต โข ปุตฺตา โลกสฺมิํ, เย ภวนฺติ อุปาสกา. \\ สทฺธา สีเลน สมฺปนฺนา, วทญฺญู วีตมจฺฉรา. \\ จนฺโท อพฺภฆนา มุตฺโต, ปริสาสุ วิโรจเร”ติ.}
\csromangathagroup{\csromangathastanza{“อติชาตํ อนุชาตํ, ปุตฺตมิจฺฉนฺติ ปณฺฑิตา. \\ อวชาตํ น อิจฺฉนฺติ, โย โหติ กุลคนฺธโน.}\csromangathastanza{เอเต โข ปุตฺตา โลกสฺมิํ, เย ภวนฺติ อุปาสกา. \\ สทฺธา สีเลน สมฺปนฺนา, วทญฺญู วีตมจฺฉรา. \\ จนฺโท อพฺภฆนา มุตฺโต, ปริสาสุ วิโรจเร”ติ.}}

\prose{อยมฺปิ อตฺโถ วุตฺโต ภควตา อิติ เม สุตนฺติ.\csromanspacer{}\csromanspacer{}ปญฺจมํ.}

\csromanheader{3}{6. อวุฏฺฐิกสุตฺต}
\csromanheader{2}{75. วุตฺตํ เหตํ ภควตา วุตฺตมรหตาติ เม สุตํ\csromandash{}}
\prose{ตโยเม ภิกฺขเว ปุคฺคลา สนฺโต สํวิชฺชมานา โลกสฺมิํ.\csromanspacer{} กตเม ตโย? อวุฏฺฐิกสโม, ปเทสวสฺสี, สพฺพตฺถาภิวสฺสี.}

\gambhira{อิติวุตฺตกปาฬิ}
\prose{\csromanfolio{240}กถญฺจ ภิกฺขเว ปุคฺคโล อวุฏฺฐิกสโม โหติ? อิธ ภิกฺขเว เอกจฺโจ ปุคฺคโล สพฺเพสญฺเญว น ทาตา โหติ สมณพฺราหฺมณกปณทฺธิกวนิพฺพกยาจกานํ\footnote{...วณิพฺพกยาจกานํ (สี)} อนฺนํ ปานํ วตฺถํ ยานํ มาลาคนฺธวิเลปนํ เสยฺยาวสถปทีเปยฺยํ.\csromanspacer{} เอวํ โข ภิกฺขเว ปุคฺคโล อวุฏฺฐิกสโม โหติ.}

\prose{กถญฺจ ภิกฺขเว ปุคฺคโล ปเทสวสฺสี โหติ? อิธ ภิกฺขเว เอกจฺโจ ปุคฺคโล เอกจฺจานํ ทาตา โยติ, เอกจฺจานํ น ทาตา (โหติ)\footnote{( ) นตฺถิ สฺยามโปตฺถเก.} สมณพฺราหฺมณกปณทฺธิกวนิพฺพกยาจกานํ อนฺนํ ปานํ วตฺถํ ยานํ มาลาคนฺธวิเลปนํ เสยฺยาวสถปทีเปยฺยํ.\csromanspacer{} เอวํ โข ภิกฺขเว ปุคฺคโล ปเทสวสฺสี โหติ.}

\prose{กถญฺจ ภิกฺขเว ปุคฺคโล สพฺพตฺถาภิวสฺสี โหติ? อิธ ภิกฺขเว เอกจฺโจ ปุคฺคโล สพฺเพสํว เทติ สมณพฺราหฺมณกปณทฺธิกวนิพฺพกยาจกานํ อนฺนํ ปานํ วตฺถํ ยานํ มาลาคนฺธวิเลปนํ เสยฺยาวสถปทีเปยฺยํ.\csromanspacer{} เอวํ โข ภิกฺขเว ปุคฺคโล สพฺพตฺถาภิวสฺสี โหติ.\csromanspacer{} อิเม โข ภิกฺขเว ตโย ปุคฺคลา สนฺโต สํวิชฺชมานา โลกสฺมินฺติ.\csromanspacer{} เอตมตฺถํ ภควา อโวจ, ตตฺเถตํ อิติ วุจฺจติ\csromandash{}}

\csromangathameasure{“น สมเณ น พฺราหฺมเณ, น กปณทฺธิกวนิพฺพเก. \\ ลทฺธาน สํวิภาเชติ, อนฺนํ ปานญฺจ โภชนํ. \\ ตํ เว ‘อวุฏฺฐิกสโม’ติ, อาหุ นํ ปุริสาธมํ. \\ เอกจฺจานํ น ททาติ, เอกจฺจานํ ปเวจฺฉติ. \\ ตํ เว ‘ปเทสวสฺสี’ติ, อาหุ เมธาวิโน ชนา. \\ สุภิกฺขวาโจ ปุริโส, สพฺพภูตานุกมฺปโก. \\ อาโมทมาโน ปกิเรติ, ‘เทถ เทถา’ติ ภาสติ. \\ ยถาปิ เมโฆ ถนยิตฺวา, คชฺชยิตฺวา ปวสฺสติ. \\ ถลํ นินฺนญฺจ ปูเรติ, อภิสนฺทนฺโตว วารินา.}
\csromangathagroup{\csromangathastanza{“น สมเณ น พฺราหฺมเณ, น กปณทฺธิกวนิพฺพเก. \\ ลทฺธาน สํวิภาเชติ, อนฺนํ ปานญฺจ โภชนํ.}\csromangathastanza{ตํ เว ‘อวุฏฺฐิกสโม’ติ, อาหุ นํ ปุริสาธมํ. \\ เอกจฺจานํ น ททาติ, เอกจฺจานํ ปเวจฺฉติ.}\csromangathastanza{ตํ เว ‘ปเทสวสฺสี’ติ, อาหุ เมธาวิโน ชนา. \\ สุภิกฺขวาโจ ปุริโส, สพฺพภูตานุกมฺปโก.}\csromangathastanza{อาโมทมาโน ปกิเรติ, ‘เทถ เทถา’ติ ภาสติ. \\ ยถาปิ เมโฆ ถนยิตฺวา, คชฺชยิตฺวา ปวสฺสติ. \\ ถลํ นินฺนญฺจ ปูเรติ, อภิสนฺทนฺโตว\footnote{อภิสนฺเทนฺโตว (?)} วารินา.}}

\csromanheader{2}{3. ติกนิปาต}
\csromangathameasure{เอวเมวํ อิเธกจฺโจ, ปุคฺคโล โหติ ตาทิโส. \\ ธมฺเมน สํหริตฺวาน, อุฏฺฐานาธิคตํ ธนํ. \\ ตปฺเปติ อนฺนปาเนน, สมฺมา ปตฺเต วนิพฺพเก”ติ.}
\csromangathagroup{\csromangathastanza{\csromanfolio{241}เอวเมวํ อิเธกจฺโจ, ปุคฺคโล โหติ ตาทิโส. \\ ธมฺเมน สํหริตฺวาน, อุฏฺฐานาธิคตํ ธนํ. \\ ตปฺเปติ อนฺนปาเนน, สมฺมา ปตฺเต วนิพฺพเก”ติ.}}

\prose{อยมฺปิ อตฺโถ วุตฺโต ภควตา อิติ เม สุตนฺติ.\csromanspacer{}\csromanspacer{}ฉฏฺฐํ.}

\csromanheader{3}{7. สุขปตฺถนาสุตฺต}
\csromanheader{2}{76. วุตฺตํ เหตํ ภควตา วุตฺตมรหตาติ เม สุตํ\csromandash{}}
\prose{ตีณิมานิ ภิกฺขเว สุขานิ ปตฺถยมาโน สีลํ รกฺเขยฺย ปณฺฑิโต.\csromanspacer{} กตมานิ ตีณิ? “ปสํสา เม อาคจฺฉตู”ติ\footnote{อาคจฺฉนฺตูติ (สฺยา)} สีลํ รกฺเขยฺย ปณฺฑิโต, “โภคา เม อุปฺปชฺชนฺตู”ติ สีลํ รกฺเขยฺย ปณฺฑิโต, “กายสฺส เภทา ปรํ มรณา สุคติํ สคฺคํ โลกํ อุปปชฺชิสฺสามี”ติ สีลํ รกฺเขยฺย ปณฺฑิโต.\csromanspacer{} อิมานิ โข ภิกฺขเว ตีณิ สุขานิ ปตฺถยมาโน สีลํ รกฺเขยฺย ปณฺฑิโตติ.\csromanspacer{} เอตมตฺถํ ภควา อโวจ, ตตฺเถตํ อิติ วุจฺจติ\csromandash{}}

\csromanmatikaanchor{mtk.174}
\csromantocmarkref{subsubsection}{9. ตณฺหาสุตฺต}{mtk.174}
\bookmark[dest={mtk.174},level=3]{9. ตณฺหาสุตฺต}
\csromangathameasure{“สีลํ รกฺเขยฺย เมธาวี, ปตฺถยาโน ตโย สุเข. \\ ปสํสํ วิตฺตลาภญฺจ, เปจฺจ สคฺเค ปโมทนํ. \\ อกโรนฺโตปิ เจ ปาปํ, กโรนฺตมุปเสวติ. \\ สงฺกิโย โหติ ปาปสฺมิํ, อวณฺโณ จสฺส รูหติ. \\ *ยาทิสํ กุรุเต มิตฺตํ, ยาทิสํ จูปเสวติ. \\ ส เว ตาทิสโก โหติ, สหวาโส หิ ตาทิโส. \\ *เสวมาโน เสวมานํ, สมฺผุฏฺโฐ สมฺผุสํ ปรํ. \\ สโร ทิทฺโธ กลาปํว, อลิตฺตมุปลิมฺปติ. \\ อุปเลปภยา ธีโร, เนว ปาปสขา สิยา. \\ *ปูติมจฺฉํ กุสคฺเคน, โย นโร อุปนยฺหติ. \\ กุสาปิ ปูติ วายนฺติ, เอวํ พาลูปเสวนา. \\ *ตครญฺจ ปลาเสน, โย นโร อุปนยฺหติ. \\ ปตฺตาปิ สุรภิ วายนฺติ, เอวํ ธีรูปเสวนา.}
\csromangathagroup{\csromangathastanza{“สีลํ รกฺเขยฺย เมธาวี, ปตฺถยาโน ตโย สุเข. \\ ปสํสํ วิตฺตลาภญฺจ, เปจฺจ สคฺเค ปโมทนํ.}\csromangathastanza{อกโรนฺโตปิ เจ ปาปํ, กโรนฺตมุปเสวติ. \\ สงฺกิโย โหติ ปาปสฺมิํ, อวณฺโณ จสฺส รูหติ.}\csromangathastanza{\csromansymbolfootnote{*}{ขุ 5. 329 ชาตเกปิ.}ยาทิสํ กุรุเต มิตฺตํ, ยาทิสํ จูปเสวติ. \\ ส เว ตาทิสโก โหติ, สหวาโส หิ\footnote{สหวาโสปิ (สี, ก)} ตาทิโส.}\csromangathastanza{\csromansymbolmark{*}เสวมาโน เสวมานํ, สมฺผุฏฺโฐ สมฺผุสํ ปรํ. \\ สโร ทิทฺโธ กลาปํว, อลิตฺตมุปลิมฺปติ.}\csromangathastanza{อุปเลปภยา\footnote{อุปลิมฺปภยา (ก)} ธีโร, เนว ปาปสขา สิยา. \\ \csromansymbolmark{*}ปูติมจฺฉํ กุสคฺเคน, โย นโร อุปนยฺหติ.}\csromangathastanza{กุสาปิ ปูติ วายนฺติ, เอวํ พาลูปเสวนา. \\ \csromansymbolmark{*}ตครญฺจ ปลาเสน, โย นโร อุปนยฺหติ. \\ ปตฺตาปิ สุรภิ วายนฺติ, เอวํ ธีรูปเสวนา.}}

\gambhira{อิติวุตฺตกปาฬิ}
\csromangathameasure{ตสฺมา ปตฺตปุฏสฺเสว, ญตฺวา สมฺปากมตฺตโน. \\ อสนฺเต นุปเสเวยฺย, สนฺเต เสเวยฺย ปณฺฑิโต. \\ อสนฺโต นิรยํ เนนฺติ, สนฺโต ปาเปนฺติ สุคฺคตินฺ”ติ.}
\csromangathagroup{\csromangathastanza{\csromanfolio{242}ตสฺมา ปตฺตปุฏสฺเสว\footnote{ปลาสปุฏสฺเสว (อิ, ก)}, ญตฺวา สมฺปากมตฺตโน. \\ อสนฺเต นุปเสเวยฺย, สนฺเต เสเวยฺย ปณฺฑิโต. \\ อสนฺโต นิรยํ เนนฺติ, สนฺโต ปาเปนฺติ สุคฺคตินฺ”ติ.}}

\csromanmatikaanchor{mtk.175}
\csromantocmarkref{subsubsection}{10. มารเธยฺยสุตฺต}{mtk.175}
\bookmark[dest={mtk.175},level=3]{10. มารเธยฺยสุตฺต}
\prose{อยมฺปิ อตฺโถ วุตฺโต ภควตา อิติ เม สุตนฺติ.\csromanspacer{}\csromanspacer{}สตฺตมํ.}

\csromanheader{3}{8. ภิทุรสุตฺต}
\csromanheader{2}{77. วุตฺตํ เหตํ ภควตา วุตฺตมรหตาติ เม สุตํ\csromandash{}}
\prose{ภิทุรายํ\footnote{ภินฺทนฺตายํ (สฺยา, อิ, ก)} ภิกฺขเว กาโย, วิญฺญาณํ วิราคธมฺมํ, สพฺเพ อุปธี อนิจฺจา ทุกฺขา วิปริณามธมฺมาติ.\csromanspacer{} เอตมตฺถํ ภควา อโวจ, ตตฺเถตํ อิติ วุจฺจติ\csromandash{}}

\csromangathameasure{“กายญฺจ ภิทุรํ ญตฺวา, วิญฺญาณญฺจ วิราคุนํ. \\ อุปธีสุ ภยํ ทิสฺวา, ชาติมรณมจฺจคา. \\ สมฺปตฺวา ปรมํ สนฺติํ, กาลํ กงฺขติ ภาวิตตฺโต”ติ.}
\csromangathagroup{\csromangathastanza{“กายญฺจ ภิทุรํ\footnote{ภินฺทนฺตํ (สฺยา, อิ, ก)} ญตฺวา, วิญฺญาณญฺจ วิราคุนํ\footnote{วิราคิกํ (ก-สี), ปภงฺคุณํ (สฺยา)}. \\ อุปธีสุ ภยํ ทิสฺวา, ชาติมรณมจฺจคา. \\ สมฺปตฺวา ปรมํ สนฺติํ, กาลํ กงฺขติ ภาวิตตฺโต”ติ.}}

\prose{อยมฺปิ อตฺโถ วุตฺโต ภควตา อิติ เม สุตนฺติ.\csromanspacer{}\csromanspacer{}อฏฺฐมํ.}

\csromanheader{3}{9. ธาตุโสสํสนฺทนสุตฺต}
\csromanheader{2}{78. วุตฺต เหตํ ภควตา วุตฺตมรหตาติ เม สุตํ\csromandash{}}
\csromanmatikaanchor{mtk.176}
\csromanmatikaanchor{mtk.177}
\csromantocmarknopage{subsection}{2. ทุติยวคฺค}{mtk.176}
\bookmark[dest={mtk.176},level=2]{2. ทุติยวคฺค}
\csromantocmarkref{subsubsection}{1. ปุญฺญกิริยวตฺถุสุตฺต}{mtk.177}
\bookmark[dest={mtk.177},level=3]{1. ปุญฺญกิริยวตฺถุสุตฺต}
\prose{\csromansymbolfootnote{*}{สํ 1. 369 ปิฏฺเฐ.}ธาตุโส ภิกฺขเว สตฺตา สตฺเตหิ สทฺธิํ สํสนฺทนฺติ สเมนฺติ, หีนาธิมุตฺติกา สตฺตา หีนาธิมุตฺติเกหิ สตฺเตหิ สทฺธิํ สํสนฺทนฺติ สเมนฺติ, กลฺยาณาธิมุตฺติกา สตฺตา กลฺยาณาธิมุตฺติเกหิ สตฺเตหิ สทฺธิํ สํสนฺทนฺติ สเมนฺติ.}

\prose{อตีตมฺปิ ภิกฺขเว อทฺธานํ ธาตุโสว สตฺตา สตฺเตหิ สทฺธิํ สํสนฺทิํสุ สมิํสุ, หีนาธิมุตฺติกา สตฺตา หีนาธิมุตฺติเกหิ สตฺเตหิ สทฺธิํ สํสนฺทิํสุ สมิํสุ, กลฺยาณาธิมุตฺติกา สตฺตา กลฺยาณาธิมุตฺติเกหิ สตฺเตหิ สทฺธิํ สํสนฺทิํสุ สมิํสุ.}

\csromanheader{2}{3. ติกนิปาต}
\prose{\csromanfolio{243}อนาคตมฺปิ ภิกฺขเว อทฺธานํ ธาตุโสว สตฺตา สตฺเตหิ สทฺธิํ สํสนฺทิสฺสนฺติ สเมสฺสนฺติ, หีนาธิมุตฺติกา สตฺตา หีนาธิมุตฺติเกหิ สตฺเตหิ สทฺธิํ สํสนฺทิสฺสนฺติ สเมสฺสนฺติ, กลฺยาณาธิมุตฺติกา สตฺตา กลฺยาณาธิมุตฺติเกหิ สตฺเตหิ สทฺธิํ สํสนฺทิสฺสนฺติ สเมสฺสนฺติ.}

\prose{เอตรหิปิ ภิกฺขเว ปจฺจุปฺปนฺนํ อทฺธานํ ธาตุโสว สตฺตา สตฺเตหิ สทฺธิํ สํสนฺทนฺติ สเมนฺติ, หีนาธิมุตฺติกา สตฺตา หีนาธิมุตฺติเกหิ สตฺเตหิ สทฺธิํ สํสนฺทนฺติ สเมนฺติ, กลฺยาณาธิมุตฺติกา สตฺตา กลฺยาณาธิมุตฺติเกหิ สตฺเตหิ สทฺธิํ สํสนฺทนฺติ สเมนฺตีติ.\csromanspacer{} เอตมตฺถํ ภควา อโวจ, ตตฺเถตํ อิติ วุจฺจติ\csromandash{}}

\csromanmatikaanchor{mtk.178}
\csromantocmarkref{subsubsection}{2. จกฺขุสุตฺต}{mtk.178}
\bookmark[dest={mtk.178},level=3]{2. จกฺขุสุตฺต}
\csromangathameasure{“สํสคฺคา วนโถ ชาโต, อสํสคฺเคน ฉิชฺชติ. \\ ปริตฺตํ ทารุมารุยฺห, ยถา สีเท มหณฺณเว. \\ เอวํ กุสีตมาคมฺม, สาธุชีวีปิ สีทติ. \\ ตสฺมา ตํ ปริวชฺเชยฺย, กุสีตํ หีนวีริยํ. \\ ปวิวิตฺเตหิ อริเยหิ, ปหิตตฺเตหิ ฌายิภิ. \\ นิจฺจํ อารทฺธวิริเยหิ, ปณฺฑิเตหิ สหาวเส”ติ.}
\csromangathagroup{\csromangathastanza{“สํสคฺคา วนโถ ชาโต, อสํสคฺเคน ฉิชฺชติ. \\ ปริตฺตํ ทารุมารุยฺห, ยถา สีเท มหณฺณเว.}\csromangathastanza{เอวํ กุสีตมาคมฺม, สาธุชีวีปิ สีทติ. \\ ตสฺมา ตํ ปริวชฺเชยฺย, กุสีตํ หีนวีริยํ.}\csromangathastanza{ปวิวิตฺเตหิ อริเยหิ, ปหิตตฺเตหิ ฌายิภิ. \\ นิจฺจํ อารทฺธวิริเยหิ, ปณฺฑิเตหิ สหาวเส”ติ.}}

\prose{อยมฺปิ อตฺโถ วุตฺโต ภควตา อิติ เม สุตนฺติ.\csromanspacer{}\csromanspacer{}นวมํ.}

\csromanheader{3}{10. ปริหานสุตฺต}
\csromanheader{2}{79. วุตฺตํ เหตํ ภควตา วุตฺตมรหตาติ เม สุตํ\csromandash{}}
\prose{ตโยเม ภิกฺขเว ธมฺมา เสขสฺส ภิกฺขุโน ปริหานาย สํวตฺตนฺติ.\csromanspacer{} กตเม ตโย? อิธ ภิกฺขเว เสโข ภิกฺขุ กมฺมาราโม โหติ กมฺมรโต กมฺมารามตมนุยุตฺโต, ภสฺสาราโม โหติ ภสฺสรโต ภสฺสารามตมนุยุตฺโต, นิทฺทาราโม โหติ นิทฺทารโต นิทฺทารามตมนุยุตฺโต.\csromanspacer{} อิเม โข ภิกฺขเว ตโย ธมฺมา เสขสฺส ภิกฺขุโน ปริหานาย สํวตฺตนฺติ.}

\prose{ตโยเม ภิกฺขเว ธมฺมา เสขสฺส ภิกฺขุโน อปริหานาย สํวตฺตนฺติ.\csromanspacer{} กตเม ตโย? อิธ ภิกฺขเว เสโข ภิกฺขุ น กมฺมาราโม โหติ น กมฺมรโต น กมฺมารามตมนุยุตฺโต, น ภสฺสาราโม โหติ น ภสฺสรโต น ภสฺสารามตมนุยุตฺโต, น นิทฺทาราโม โหติ น}

\csromanmatikaanchor{mtk.179}
\csromantocmarkref{subsubsection}{3. อินฺทฺริยสุตฺต}{mtk.179}
\bookmark[dest={mtk.179},level=3]{3. อินฺทฺริยสุตฺต}
\gambhira{อิติวุตฺตกปาฬิ}
\prose{\csromanfolio{244}นิทฺทารโต น นิทฺทารามตมนุยุตฺโต.\csromanspacer{} อิเม โข ภิกฺขเว ตโย ธมฺมา เสขสฺส ภิกฺขุโน อปริหานาย สํวตฺตนฺตีติ.\csromanspacer{} เอตมตฺถํ ภควา อโวจ, ตตฺเถตํ อิติ วุจฺจติ\csromandash{}}

\csromangathameasure{“กมฺมาราโม ภสฺสาราโม, นิทฺทาราโม จ อุทฺธโต. \\ อภพฺโพ ตาทิโส ภิกฺขุ, ผุฏฺฐุํ สมฺโพธิมุตฺตมํ. \\ ตสฺมา หิ อปฺปกิจฺจ’สฺส, อปฺปมิทฺโธ อนุทฺธโต. \\ ภพฺโพ โส ตาทิโส ภิกฺขุ, ผุฏฺฐุํ สมฺโพธิมุตฺตมนฺ”ติ.}
\csromangathagroup{\csromangathastanza{“กมฺมาราโม ภสฺสาราโม\footnote{ภสฺสรโต (สพฺพตฺถ)}, นิทฺทาราโม จ อุทฺธโต. \\ อภพฺโพ ตาทิโส ภิกฺขุ, ผุฏฺฐุํ สมฺโพธิมุตฺตมํ.}\csromangathastanza{ตสฺมา หิ อปฺปกิจฺจ’สฺส, อปฺปมิทฺโธ อนุทฺธโต. \\ ภพฺโพ โส ตาทิโส ภิกฺขุ, ผุฏฺฐุํ สมฺโพธิมุตฺตมนฺ”ติ.}}

\prose{อยมฺปิ อตฺโถ วุตฺโต ภควตา อิติ เม สุตนฺติ.\csromanspacer{}\csromanspacer{}ทสมํ.\csromanspacer{} ตติโย วคฺโค.}

\titlehead{\textbf{ตสฺสุทฺทานํ}}
\csromangathameasure{เทฺว ทิฏฺฐี นิสฺสรณํ รูปํ, ปุตฺโต อวุฏฺฐิเกน จ. \\ สุขา จ ภิทุโร ธาตุ, ปริหาเนน เต ทสาติ.}
\csromangathagroup{\csromangathastanza{เทฺว ทิฏฺฐี นิสฺสรณํ รูปํ, ปุตฺโต อวุฏฺฐิเกน จ. \\ สุขา จ ภิทุโร\footnote{ภินฺทนา (สพฺพตฺถ)} ธาตุ, ปริหาเนน เต ทสาติ.}}

\csromanheader{2}{4. จตุตฺถวคฺค}
\csromanheader{3}{1. วิตกฺกสุตฺต}
\csromanmatikaanchor{mtk.180}
\csromantocmarkref{subsubsection}{4. อทฺธาสุตฺต}{mtk.180}
\bookmark[dest={mtk.180},level=3]{4. อทฺธาสุตฺต}
\csromanheader{2}{80. วุตฺต เหตํ ภควตา วุตฺตมรหตาติ เม สุตํ\csromandash{}}
\prose{ตโยเม ภิกฺขเว อกุสลวิตกฺกา.\csromanspacer{} กตเม ตโย? อนวญฺญตฺติปฏิสํยุตฺโต วิตกฺโก, ลาภสกฺการสิโลกปฏิสํยุตฺโต วิตกฺโก, ปรานุทฺทยตาปฏิสํยุตฺโต วิตกฺโก.\csromanspacer{} อิเม โข ภิกฺขเว ตโย อกุสลวิตกฺกาติ.\csromanspacer{} เอตมตฺถํ ภควา อโวจ, ตตฺเถตํ อิติ วุจฺจติ\csromandash{}}

\csromangathameasure{“อนวญฺญตฺติสํยุตฺโต, ลาภสกฺการคารโว. \\ สหนนฺที อมจฺเจหิ, อารา สํโยชนกฺขยา. \\ โย จ ปุตฺตปสุํ หิตฺวา, วิวาเห สํหรานิ จ. \\ ภพฺโพ โส ตาทิโส ภิกฺขุ, ผุฏฺฐุํ สมฺโพธิมุตฺตมนฺ”ติ.}
\csromangathagroup{\csromangathastanza{“อนวญฺญตฺติสํยุตฺโต, ลาภสกฺการคารโว. \\ สหนนฺที อมจฺเจหิ, อารา สํโยชนกฺขยา.}\csromangathastanza{โย จ ปุตฺตปสุํ หิตฺวา, วิวาเห สํหรานิ\footnote{สงฺคหานิ (ก-สี, สฺยา, อิ)} จ. \\ ภพฺโพ โส ตาทิโส ภิกฺขุ, ผุฏฺฐุํ สมฺโพธิมุตฺตมนฺ”ติ.}}

\prose{อยมฺปิ อตฺโถ วุตฺโต ภควตา อิติ เม สุตนฺติ.\csromanspacer{}\csromanspacer{}ปฐมํ.}

\csromanheader{3}{3. ติกนิปาต \textbf{2. สกฺการสุตฺต}}
\csromanheader{2}{81. วุตฺตํ เหตํ ภควตา วุตฺตมรหตาติ เม สุตํ\csromandash{}}
\csromanmatikaanchor{mtk.181}
\csromantocmarkref{subsubsection}{5. ทุจฺจริตสุตฺต}{mtk.181}
\bookmark[dest={mtk.181},level=3]{5. ทุจฺจริตสุตฺต}
\prose{\csromanfolio{245}ทิฏฺฐา มยา ภิกฺขเว สตฺตา สกฺกาเรน อภิภูตา ปริยาทินฺนจิตฺตา กายสฺส เภทา ปรํ มรณา อปายํ ทุคฺคติํ วินิปาตํ นิรยํ อุปปนฺนา.}

\prose{ทิฏฺฐา มยา ภิกฺขเว สตฺตา อสกฺกาเรน อภิภูตา ปริยาทินฺนจิตฺตา กายสฺส เภทา ปรํ มรณา อปายํ ทุคฺคติํ วินิปาตํ นิรยํ อุปปนฺนา.}

\prose{ทิฏฺฐา มยา ภิกฺขเว สตฺตา สกฺกาเรน จ อสกฺกาเรน จ ตทุภเยน อภิภูตา ปริยาทินฺนจิตฺตา กายสฺส เภทา ปรํ มรณา อปายํ ทุคฺคติํ วินิปาตํ นิรยํ อุปปนฺนา.}

\prose{ตํ โข ปนาหํ ภิกฺขเว นาญฺญสฺส สมณสฺส วา พฺราหฺมณสฺส วา สุตฺวา วทามิ. ( )\footnote{(ทิฏฺฐา มยา ภิกฺขเว สตฺตา สกฺกาเรน อภิภูตา ฯเปฯ อสกฺกาเรน อภิภูตา ฯเปฯ สกฺกาเรน จ อสกฺกาเรน จ ตทุภเยน อภิภูตา ปริยาทินฺนจิตฺตา กายสฺส เภทา ปรํ มรณา อปายํ ทุคฺคติํ วินิปาตํ นิรยํ อุปปนฺนา.) (สฺยา) ปุริมวคฺเค (235 ปิฏฺเฐ) มิจฺฉาทิฏฺฐิกสมฺมาทิฏฺฐิกสุตฺเตหิ ปน สเมติ, อนฺวยพฺยติเรกวากฺยานํ ปน อนนฺตริตตา ปาสํสตรา.} อปิ จ ภิกฺขเว ยเทว เม สามํ ญาตํ สามํ ทิฏฺฐํ สามํ วิทิตํ, ตเมวาหํ วทามิ.}

\prose{ทิฏฺฐา มยา ภิกฺขเว สตฺตา สกฺกาเรน อภิภูตา ปริยาทินฺนจิตฺตา กายสฺส เภทา ปรํ มรณา อปายํ ทุคฺคติํ วินิปาตํ นิรยํ อุปปนฺนา.}

\prose{ทิฏฺฐา มยา ภิกฺขเว สตฺตา อสกฺกาเรน อภิภูตา ปริยาทินฺนจิตฺตา กายสฺส เภทา ปรํ มรณา อปายํ ทุคฺคติํ วินิปาตํ นิรยํ อุปปนฺนา.}

\prose{ทิฏฺฐา มยา ภิกฺขเว สตฺตา สกฺกาเรน จ อสกฺกาเรน จ ตทุภเยน อภิภูตา ปริยาทินฺนจิตฺตา กายสฺส เภทา ปรํ มรณา อปายํ ทุคฺคติํ วินิปาตํ นิรยํ อุปปนฺนาติ.\csromanspacer{} เอตมตฺถํ ภควา อโวจ, ตตฺเถตํ อิติ วุจฺจติ\csromandash{}}

\csromanmatikaanchor{mtk.182}
\csromantocmarkref{subsubsection}{6. สุจริตสุตฺต}{mtk.182}
\bookmark[dest={mtk.182},level=3]{6. สุจริตสุตฺต}
\csromangathameasure{“ยสฺส สกฺกริยมานสฺส, อสกฺกาเรน จูภยํ. \\ สมาธิ น วิกมฺปติ, อปฺปมาทวิหาริโน.}
\csromangathagroup{\csromangathastanza{“ยสฺส สกฺกริยมานสฺส, อสกฺกาเรน จูภยํ. \\ สมาธิ น วิกมฺปติ, อปฺปมาทวิหาริโน\footnote{อปฺปมาณวิหาริโน (สี-ฏฺฐ)}.}}

\gambhira{อิติวุตฺตกปาฬิ}
\csromangathameasure{ตํ ฌายินํ สาตติกํ, สุขุมทิฏฺฐิวิปสฺสกํ. \\ อุปาทานกฺขยารามํ, อาหุ สปฺปุริโส อิตี”ติ.}
\csromangathagroup{\csromangathastanza{\csromanfolio{246}ตํ ฌายินํ สาตติกํ, สุขุมทิฏฺฐิวิปสฺสกํ. \\ อุปาทานกฺขยารามํ, อาหุ สปฺปุริโส อิตี”ติ.}}

\prose{อยมฺปิ อตฺโถ วุตฺโต ภควตา อิติ เม สุตนฺติ.\csromanspacer{}\csromanspacer{}ทุติยํ.}

\csromanheader{3}{3. เทวสทฺทสุตฺต}
\csromanheader{2}{82. วุตฺตํ เหตํ ภควตา วุตฺตมรหตาติ เม สุตํ\csromandash{}}
\csromanmatikaanchor{mtk.183}
\csromantocmarkref{subsubsection}{7. โสเจยฺยสุตฺต}{mtk.183}
\bookmark[dest={mtk.183},level=3]{7. โสเจยฺยสุตฺต}
\prose{ตโยเม ภิกฺขเว เทเวสุ เทวสทฺทา นิจฺฉรนฺติ สมยา สมยํ อุปาทาย.\csromanspacer{} กตเม ตโย? ยสฺมิํ ภิกฺขเว สมเย อริยสาวโก เกสมสฺสุํ โอหาเรตฺวา กาสายานิ วตฺถานิ อจฺฉาเทตฺวา อคารสฺมา อนคาริยํ ปพฺพชฺชาย เจเตติ, ตสฺมิํ สมเย\footnote{ตสฺมิํ ภิกฺขเว สมเย (อิ, ก)} เทเวสุ เทวสทฺโท นิจฺฉรติ “เอโส อริยสาวโก มาเรน สทฺธิํ สงฺคามาย เจเตตี”ติ.\csromanspacer{} อยํ ภิกฺขเว ปฐโม เทเวสุ เทวสทฺโท นิจฺฉรติ สมยา สมยํ อุปาทาย.}

\prose{ปุน จปรํ ภิกฺขเว ยสฺมิํ สมเย อริยสาวโก สตฺตนฺนํ โพธิปกฺขิยานํ ธมฺมานํ ภาวนานุโยคมนุยุตฺโต วิหรติ, ตสฺมิํ สมเย เทเวสุ เทวสทฺโท นิจฺฉรติ “เอโส อริยสาวโก มาเรน สทฺธิํ สงฺคาเมตี”ติ.\csromanspacer{} อยํ ภิกฺขเว ทุติโย เทเวสุ เทวสทฺโท นิจฺฉรติ สมยา สมยํ อุปาทาย.}

\prose{ปุน จปรํ ภิกฺขเว ยสฺมิํ สมเย อริยสาวโก อาสวานํ ขยา อนาสวํ เจโตวิมุตฺติํ ปญฺญาวิมุตฺติํ ทิฏฺเฐว ธมฺเม สยํ อภิญฺญา สจฺฉิกตฺวา อุปสมฺปชฺช วิหรติ, ตสฺมิํ สมเย เทเวสุ เทวสทฺโท นิจฺฉรติ “เอโส อริยสาวโก วิชิตสงฺคาโม ตเมว สงฺคามสีสํ อภิวิชิย อชฺฌาวสตี”ติ.\csromanspacer{} อยํ ภิกฺขเว ตติโย เทเวสุ เทวสทฺโท นิจฺฉรติ สมยา สมยํ อุปาทาย.\csromanspacer{} อิเม โข ภิกฺขเว ตโย เทเวสุ เทวสทฺทา นิจฺฉรนฺติ สมยา สมยํ อุปาทายาติ.\csromanspacer{} เอตมตฺถํ ภควา อโวจ, ตตฺเถตํ อิติ วุจฺจติ\csromandash{}}

\csromangathameasure{“ทิสฺวา วิชิตสงฺคามํ, สมฺมาสมฺพุทฺธสาวกํ. \\ เทวตาปิ นมสฺสนฺติ, มหนฺตํ วีตสารทํ.}
\csromangathagroup{\csromangathastanza{“ทิสฺวา วิชิตสงฺคามํ, สมฺมาสมฺพุทฺธสาวกํ. \\ เทวตาปิ นมสฺสนฺติ, มหนฺตํ วีตสารทํ.}}

\csromanheader{2}{3. ติกนิปาต}
\csromanmatikaanchor{mtk.184}
\csromantocmarkref{subsubsection}{8. โมเนยฺยสุตฺต}{mtk.184}
\bookmark[dest={mtk.184},level=3]{8. โมเนยฺยสุตฺต}
\csromantocmarkref{subsubsection}{9-10. ราคสุตฺตานิ}{mtk.185}
\bookmark[dest={mtk.185},level=3]{9-10. ราคสุตฺตานิ}
\csromangathameasure{นโม เต ปุริสาชญฺญ, โย ตฺวํ ทุชฺชยมชฺฌภู. \\ เชตฺวาน มจฺจุโน เสนํ, วิโมกฺเขน อนาวรํ. \\ อิติ เหตํ นมสฺสนฺติ, เทวตา ปตฺตมานสํ. \\ ตญฺหิ ตสฺส น ปสฺสนฺติ, เยน มจฺจุวสํ วเช”ติ.}
\csromangathagroup{\csromangathastanza{\csromanfolio{247}นโม เต ปุริสาชญฺญ, โย ตฺวํ ทุชฺชยมชฺฌภู. \\ เชตฺวาน มจฺจุโน เสนํ, วิโมกฺเขน อนาวรํ.}\csromangathastanza{อิติ เหตํ นมสฺสนฺติ, เทวตา ปตฺตมานสํ. \\ ตญฺหิ ตสฺส น ปสฺสนฺติ, เยน มจฺจุวสํ วเช”ติ.}}

\prose{อยมฺปิ อตฺโถ วุตฺโต ภควตา อิติ เม สุตนฺติ.\csromanspacer{}\csromanspacer{}ตติยํ.}

\csromanheader{3}{4. ปญฺจปุพฺพนิมิตฺตสุตฺต}
\csromanheader{2}{83. วุตฺตํ เหตํ ภควตา วุตฺตมรหตาติ เม สุตํ\csromandash{}}
\prose{ยทา ภิกฺขเว เทโว เทวกายา จวนธมฺโม โหติ, ปญฺจสฺส ปุพฺพนิมิตฺตานิ ปาตุภวนฺติ “มาลา มิลายนฺติ, วตฺถานิ กิลิสฺสนฺติ, กจฺเฉหิ เสทา มุจฺจนฺติ, กาเย ทุพฺพณฺณิยํ โอกฺกมติ, สเก เทโว เทวาสเน นาภิรมตี”ติ.\csromanspacer{} ตเมนํ ภิกฺขเว เทวา “จวนธมฺโม อยํ เทวปุตฺโต”ติ อิติ วิทิตฺวา ตีหิ วาจาหิ อนุโมเทนฺติ\footnote{อนุโมทนฺติ (สี, สฺยา, อิ)} “อิโต โภ สุคติํ คจฺฉ, สุคติํ คนฺตฺวา สุลทฺธลาภํ ลภ, สุลทฺธลาภํ ลภิตฺวา สุปฺปติฏฺฐิโต ภวาหี”ติ.}

\prose{เอวํ วุตฺเต อญฺญตโร ภิกฺขุ ภควนฺตํ เอตทโวจ “กินฺนุ โข ภนฺเต เทวานํ สุคติคมนสงฺขาตํ, กิญฺจ ภนฺเต เทวานํ สุลทฺธลาภสงฺขาตํ, กิํ ปน ภนฺเต เทวานํ สุปฺปติฏฺฐิตสงฺขาตนฺ”ติ.}

\prose{มนุสฺสตฺตํ โข ภิกฺขุ\footnote{ภิกฺขเว (สฺยา, อิ)} เทวานํ สุคติคมนสงฺขาตํ.\csromanspacer{} ยํ มนุสฺสภูโต สมาโน ตถาคตปฺปเวทิเต ธมฺมวินเย สทฺธํ ปฏิลภติ, อิทํ โข ภิกฺขุ2 เทวานํ สุลทฺธลาภสงฺขาตํ.\csromanspacer{} สา โข ปนสฺส สทฺธา นิวิฏฺฐา โหติ มูลชาตา ปติฏฺฐิตา ทฬฺหา อสํหาริยา สมเณน วา พฺราหฺมเณน วา เทเวน วา มาเรน วา พฺรหฺมุนา วา เกนจิ วา โลกสฺมิํ.\csromanspacer{} อิทํ โข ภิกฺขุ2 เทวานํ สุปฺปติฏฺฐิตสงฺขาตนฺติ.\csromanspacer{} เอตมตฺถํ ภควา อโวจ, ตตฺเถตํ อิติ วุจฺจติ\csromandash{}}

\csromangathameasure{“ยทา เทโว เทวกายา, จวติ อายุสงฺขยา. \\ ตโย สทฺทา นิจฺฉรนฺติ, เทวานํ อนุโมทตํ.}
\csromangathagroup{\csromangathastanza{“ยทา เทโว เทวกายา, จวติ อายุสงฺขยา. \\ ตโย สทฺทา นิจฺฉรนฺติ, เทวานํ อนุโมทตํ.}}

\gambhira{อิติวุตฺตกปาฬิ}
\csromangathameasure{อิโต โภ สุคติํ คจฺฉ, มนุสฺสานํ สหพฺยตํ. \\ มนุสฺสภูโต สทฺธมฺเม, ลภ สทฺธํ อนุตฺตรํ. \\ สา เต สทฺธา นิวิฏฺฐ’สฺส, มูลชาตา ปติฏฺฐิตา. \\ ยาวชีวํ อสํหีรา, สทฺธมฺเม สุปฺปเวทิเต. \\ กายทุจฺจริตํ หิตฺวา, วจีทุจฺจริตานิ จ. \\ มโนทุจฺจริตํ หิตฺวา, ยญฺจญฺญํ โทสสญฺหิตํ. \\ กาเยน กุสลํ กตฺวา, วาจาย กุสลํ พหุํ. \\ มนสา กุสลํ กตฺวา, อปฺปมาณํ นิรูปธิํ. \\ ตโต โอปธิกํ ปุญฺญํ, กตฺวา ทาเนน ตํ พหุํ. \\ อญฺเญปิ มจฺเจ สทฺธมฺเม, พฺรหฺมจริเย นิเวสย. \\ อิมาย อนุกมฺปาย, เทวา เทวํ ยทา วิทู. \\ จวนฺตํ อนุโมเทนฺติ, เอหิ เทว ปุนปฺปุนนฺ”ติ.}
\csromangathagroup{\csromangathastanza{\csromanfolio{248}อิโต โภ สุคติํ คจฺฉ, มนุสฺสานํ สหพฺยตํ. \\ มนุสฺสภูโต สทฺธมฺเม, ลภ สทฺธํ อนุตฺตรํ.}\csromangathastanza{สา เต สทฺธา นิวิฏฺฐ’สฺส, มูลชาตา ปติฏฺฐิตา. \\ ยาวชีวํ อสํหีรา, สทฺธมฺเม สุปฺปเวทิเต.}\csromangathastanza{กายทุจฺจริตํ หิตฺวา, วจีทุจฺจริตานิ จ. \\ มโนทุจฺจริตํ หิตฺวา, ยญฺจญฺญํ โทสสญฺหิตํ.}\csromangathastanza{กาเยน กุสลํ กตฺวา, วาจาย กุสลํ พหุํ. \\ มนสา กุสลํ กตฺวา, อปฺปมาณํ นิรูปธิํ.}\csromangathastanza{ตโต โอปธิกํ ปุญฺญํ, กตฺวา ทาเนน ตํ พหุํ. \\ อญฺเญปิ มจฺเจ สทฺธมฺเม, พฺรหฺมจริเย นิเวสย\footnote{นิเปสเย (สี, สฺยา)}.}\csromangathastanza{อิมาย อนุกมฺปาย, เทวา เทวํ ยทา วิทู. \\ จวนฺตํ อนุโมเทนฺติ, เอหิ เทว ปุนปฺปุนนฺ”ติ.}}

\prose{อยมฺปิ อตฺโถ วุตฺโต ภควตา อิติ เม สุตนฺติ.\csromanspacer{}\csromanspacer{}จตุตฺถํ.}

\csromanheader{3}{5. พหุชนหิตสุตฺต}
\csromanheader{2}{84. วุตฺตํ เหตํ ภควตา วุตฺตมรหตาติ เม สุตํ\csromandash{}}
\prose{ตโยเม ปุคฺคลา โลเก อุปฺปชฺชมานา อุปฺปชฺชนฺติ พหุชนหิตาย พหุชนสุขาย โลกานุกมฺปาย อตฺถาย หิตาย สุขาย เทวมนุสฺสานํ.\csromanspacer{} กตเม ตโย? อิธ ภิกฺขเว ตถาคโต โลเก อุปฺปชฺชติ อรหํ สมฺมาสมฺพุทฺโธ วิชฺชาจรณสมฺปนฺโน สุคโต โลกวิทู อนุตฺตโร ปุริสทมฺมสารถิ สตฺถา เทวมนุสฺสานํ พุทฺโธ ภควา, โส ธมฺมํ เทเสติ อาทิกลฺยาณํ มชฺเฌกลฺยาณํ ปริโยสานกลฺยาณํ สาตฺถํ สพฺยญฺชนํ เกวลปริปุณฺณํ ปริสุทฺธํ พฺรหฺมจริยํ ปกาเสติ.\csromanspacer{} อยํ ภิกฺขเว ปฐโม ปุคฺคโล โลเก อุปฺปชฺชมาโน อุปฺปชฺชติ พหุชนหิตาย พหุชนสุขาย โลกานุกมฺปาย อตฺถาย หิตาย สุขาย เทวมนุสฺสานํ.}

\prose{ปุน จปรํ ภิกฺขเว ตสฺเสว สตฺถุ\footnote{สตฺถุโน (สฺยา)} สาวโก อรหํ โหติ ขีณาสโว วุสิตวา กตกรณีโย โอหิตภาโร อนุปฺปตฺตสทตฺโถ}

\vspace{\tipitakaparskip}
\csromancenter{ติกนิปาต ปริกฺขีณภวสํโยชโน สมฺมทญฺญา วิมุตฺโต, โส ธมฺมํ เทเสติ อาทิกลฺยาณํ มชฺเฌกลฺยาณํ ปริโยสานกลฺยาณํ สาตฺถํ สพฺยญฺชนํ เกวลปริปุณฺณํ ปริสุทฺธํ พฺรหฺมจริยํ ปกาเสติ.\csromanspacer{} อยํ ภิกฺขเว ทุติโย ปุคฺคโล โลเก อุปฺปชฺชมาโน อุปฺปชฺชติ พหุชนหิตาย พหุชนสุขาย โลกานุกมฺปาย อตฺถาย หิตาย สุขาย เทวมนุสฺสานํ.}
\prose{\csromanfolio{249}ปุน จปรํ ภิกฺขเว ตสฺเสว สตฺถุ สาวโก เสโข โหติ ปาฏิปโท พหุสฺสุโต สีลวตูปปนฺโน, โสปิ\footnote{โส (?)} ธมฺมํ เทเสติ อาทิกลฺยาณํ มชฺเฌกลฺยาณํ ปริโยสานกลฺยาณํ สาตฺถํ สพฺยญฺชนํ เกวลปริปุณฺณํ ปริสุทฺธํ พฺรหฺมจริยํ ปกาเสติ.\csromanspacer{} อยํ ภิกฺขเว ตติโย ปุคฺคโล โลเก อุปฺปชฺชมาโน อุปฺปชฺชติ พหุชนหิตาย พหุชนสุขาย โลกานุกมฺปาย อตฺถาย หิตาย สุขาย เทวมนุสฺสานํ.\csromanspacer{} อิเม โข ภิกฺขเว ตโย ปุคฺคลา โลเก อุปฺปชฺชมานา อุปฺปชฺชนฺติ พหุชนหิตาย พหุชนสุขาย โลกานุกมฺปาย อตฺถาย หิตาย สุขาย เทวมนุสฺสานนฺติ.\csromanspacer{} เอตมตฺถํ ภควา อโวจ, ตตฺเถตํ อิติ วุจฺจติ\csromandash{}}

\csromangathameasure{“สตฺถา หิ โลเก ปฐโม มเหสิ, \\ ตสฺสนฺวโย สาวโก ภาวิตตฺโต. \\ อถาปโร ปาฏิปโทปิ เสโข, \\ พหุสฺสุโต สีลวตูปปนฺโน. \\ เอเต ตโย เทวมนุสฺสเสฏฺฐา, \\ ปภงฺกรา ธมฺมมุทีรยนฺตา. \\ อปาปุรนฺติ อมตสฺส ทฺวารํ, \\ โยคา ปโมเจนฺติ พหุชฺชนํ เต. \\ เย สตฺถวาเหน อนุตฺตเรน, \\ สุเทสิตํ มคฺคมนุกฺกมนฺติ. \\ อิเธว ทุกฺขสฺส กโรนฺติ อนฺตํ, \\ เย อปฺปมตฺตา สุคตสฺส สาสเน”ติ.}
\csromangathagroup{\csromangathastanza{“สตฺถา หิ โลเก ปฐโม มเหสิ, \\ ตสฺสนฺวโย สาวโก ภาวิตตฺโต. \\ อถาปโร ปาฏิปโทปิ เสโข, \\ พหุสฺสุโต สีลวตูปปนฺโน.}\csromangathastanza{เอเต ตโย เทวมนุสฺสเสฏฺฐา, \\ ปภงฺกรา ธมฺมมุทีรยนฺตา. \\ อปาปุรนฺติ\footnote{อปาปุเรนฺติ (ก)} อมตสฺส ทฺวารํ, \\ โยคา ปโมเจนฺติ\footnote{อปาปุเรนฺติ (ก)} พหุชฺชนํ เต.}\csromangathastanza{เย สตฺถวาเหน อนุตฺตเรน, \\ สุเทสิตํ มคฺคมนุกฺกมนฺติ\footnote{มคฺคมนุคฺคมนฺติ (สี, ก)}. \\ อิเธว ทุกฺขสฺส กโรนฺติ อนฺตํ, \\ เย อปฺปมตฺตา สุคตสฺส สาสเน”ติ.}}

\csromanmatikaanchor{mtk.186}
\csromanmatikaanchor{mtk.187}
\csromantocmarknopage{subsection}{3. ตติยวคฺค}{mtk.186}
\bookmark[dest={mtk.186},level=2]{3. ตติยวคฺค}
\csromantocmarkref{subsubsection}{1. มิจฺฉาทิฏฺฐิกสุตฺต}{mtk.187}
\bookmark[dest={mtk.187},level=3]{1. มิจฺฉาทิฏฺฐิกสุตฺต}
\prose{อยมฺปิ อตฺโถ วุตฺโต ภควตา อิติ เม สุตนฺติ.\csromanspacer{}\csromanspacer{}ปญฺจมํ.}

\gambhira{อิติวุตฺตกปาฬิ}
\csromanheader{3}{6. อสุภานุปสฺสีสุตฺต}
\csromanheader{2}{85. วุตฺตํ เหตํ ภควตา วุตฺตมรหตาติ เม สุตํ\csromandash{}}
\prose{\csromanfolio{250}อสุภานุปสฺสี ภิกฺขเว กายสฺมิํ วิหรถ, อานาปานสฺสติ จ โว อชฺฌตฺตํ ปริมุขํ สูปฏฺฐิตา โหตุ, สพฺพสงฺขาเรสุ อนิจฺจานุปสฺสิโน วิหรถ.\csromanspacer{} อสุภานุปสฺสีนํ ภิกฺขเว กายสฺมิํ วิหรตํ โย สุภาย ธาตุยา ราคานุสโย, โส ปหียติ\footnote{ปหิยฺยติ (ก)}, อานาปานสฺสติยา อชฺฌตฺตํ ปริมุขํ สูปฏฺฐิติตาย เย พาหิรา วิตกฺกาสยา วิฆาตปกฺขิกา, เต น โหนฺติ.\csromanspacer{} สพฺพสงฺขาเรสุ อนิจฺจานุปสฺสีนํ วิหรตํ ยา อวิชฺชา, สา ปหียติ.\csromanspacer{} ยา วิชฺชา, สา อุปฺปชฺชตีติ.\csromanspacer{} เอตมตฺถํ ภควา อโวจ, ตตฺเถตํ อิติ วุจฺจติ\csromandash{}}

\csromangathameasure{“อสุภานุปสฺสี กายสฺมิํ, อานาปาเน ปฏิสฺสโต. \\ สพฺพสงฺขารสมถํ, ปสฺสํ อาตาปิ สพฺพทา. \\ ส เว สมฺมทฺทโส ภิกฺขุ, ยโต ตตฺถ วิมุจฺจติ. \\ อภิญฺญาโวสิโต สนฺโต, ส เว โยคาติโค มุนี”ติ.}
\csromangathagroup{\csromangathastanza{“อสุภานุปสฺสี กายสฺมิํ, อานาปาเน ปฏิสฺสโต. \\ สพฺพสงฺขารสมถํ, ปสฺสํ อาตาปิ สพฺพทา.}\csromangathastanza{ส เว สมฺมทฺทโส ภิกฺขุ, ยโต ตตฺถ วิมุจฺจติ. \\ อภิญฺญาโวสิโต สนฺโต, ส เว โยคาติโค มุนี”ติ.}}

\prose{อยมฺปิ อตฺโถ วุตฺโต ภควตา อิติ เม สุตนฺติ.\csromanspacer{}\csromanspacer{}ฉฏฺฐํ.}

\csromanheader{3}{7. ธมฺมานุธมฺมปฏิปนฺนสุตฺต}
\csromanmatikaanchor{mtk.188}
\csromantocmarkref{subsubsection}{2. สมฺมาทิฏฺฐิกสุตฺต}{mtk.188}
\bookmark[dest={mtk.188},level=3]{2. สมฺมาทิฏฺฐิกสุตฺต}
\csromanheader{2}{86. วุตฺตํ เหตํ ภควตา วุตฺตมรหตาติ เม สุตํ\csromandash{}}
\prose{ธมฺมานุธมฺมปฏิปนฺนสฺส ภิกฺขุโน อยมนุธมฺโม โหติ เวยฺยากรณาย ธมฺมานุธมฺมปฏิปนฺโนยนฺติ, ภาสมาโน ธมฺมญฺเญว ภาสติ โน อธมฺมํ, วิตกฺกยมาโน วา ธมฺมวิตกฺกญฺเญว วิตกฺเกติ โน อธมฺมวิตกฺกํ, ตทุภยํ วา ปน อภินิเวชฺเชตฺวา อุเปกฺขโก วิหรติ สโต สมฺปชาโนติ.\csromanspacer{} เอตมตฺถํ ภควา อโวจ, ตตฺเถตํ อิติ วุจฺจติ\csromandash{}}

\csromangathameasure{“ธมฺมาราโม ธมฺมรโต, ธมฺมํ อนุวิจินฺตยํ. \\ ธมฺมํ อนุสฺสรํ ภิกฺขุ, สทฺธมฺมา น ปริหายติ. \\ จรํ วา ยทิ วา ติฏฺฐํ, นิสินฺโน อุท วา สยํ. \\ อชฺฌตฺตํ สมยํ จิตฺตํ, สนฺติเมวาธิคจฺฉตี”ติ.}
\csromangathagroup{\csromangathastanza{“ธมฺมาราโม ธมฺมรโต, ธมฺมํ อนุวิจินฺตยํ. \\ ธมฺมํ อนุสฺสรํ ภิกฺขุ, สทฺธมฺมา น ปริหายติ.}\csromangathastanza{จรํ วา ยทิ วา ติฏฺฐํ, นิสินฺโน อุท วา สยํ. \\ อชฺฌตฺตํ สมยํ จิตฺตํ, สนฺติเมวาธิคจฺฉตี”ติ.}}

\prose{อยมฺปิ อตฺโถ วุตฺโต ภควตา อิติ เม สุตนฺติ.\csromanspacer{}\csromanspacer{}สตฺตมํ.}

\csromanheader{3}{3. ติกนิปาต \textbf{8. อนฺธกรณสุตฺต}}
\csromanheader{2}{87. วุตฺตํ เหตํ ภควตา วุตฺตมรหตาติ เม สุตํ\csromandash{}}
\prose{\csromanfolio{251}ตโยเม ภิกฺขเว อกุสลวิตกฺกา อนฺธกรณา อจกฺขุกรณา อญฺญาณกรณา ปญฺญานิโรธิกา วิฆาตปกฺขิกา อนิพฺพานสํวตฺตนิกา.\csromanspacer{} กตเม ตโย? กามวิตกฺโก ภิกฺขเว อนฺธกรโณ อจกฺขุกรโณ อญฺญาณกรโณ ปญฺญานิโรธิโก วิฆาตปกฺขิโก อนิพฺพานสํวตฺตนิโก, พฺยาปาทวิตกฺโก ภิกฺขเว อนฺธกรโณ อจกฺขุกรโณ อญฺญาณกรโณ ปญฺญานิโรธิโก วิฆาตปกฺขิโก อนิพฺพานสํวตฺตนิโก, วิหิํสาวิตกฺโก ภิกฺขเว อนฺธกรโณ อจกฺขุกรโณ อญฺญาณกรโณ ปญฺญานิโรธิโก วิฆาตปกฺขิโก อนิพฺพานสํวตฺตนิโก.\csromanspacer{} อิเม โข ภิกฺขเว ตโย อกุสลวิตกฺกา อนฺธกรณา อจกฺขุกรณา อญฺญาณกรณา ปญฺญานิโรธิกา วิฆาตปกฺขิกา อนิพฺพานสํวตฺตนิกา.}

\prose{ตโยเม ภิกฺขเว กุสลวิตกฺกา อนนฺธกรณา จกฺขุกรณา ญาณกรณา ปญฺญาวุทฺธิกา อวิฆาตปกฺขิกา นิพฺพานสํวตฺตนิกา.\csromanspacer{} กตเม ตโย? เนกฺขมฺมวิตกฺโก ภิกฺขเว อนนฺธกรโณ จกฺขุกรโณ ญาณกรโณ ปญฺญาวุทฺธิโก อวิฆาตปกฺขิโก นิพฺพานสํวตฺตนิโก, อพฺยาปาทวิตกฺโก ภิกฺขเว อนนฺธกรโณ จกฺขุกรโณ ญาณกรโณ ปญฺญาวุทฺธิโก อวิฆาตปกฺขิโก นิพฺพานสํวตฺตนิโก, อวิหิํสาวิตกฺโก ภิกฺขเว อนนฺธกรโณ จกฺขุกรโณ ญาณกรโณ ปญฺญาวุทฺธิโก อวิฆาตปกฺขิโก นิพฺพานสํวตฺตนิโก.\csromanspacer{} อิเม โข ภิกฺขเว ตโย กุสลวิตกฺกา อนนฺธกรณา จกฺขุกรณา ญาณกรณา ปญฺญาวุทฺธิกา อวิฆาตปกฺขิกา นิพฺพานสํวตฺตนิกาติ.\csromanspacer{} เอตมตฺถํ ภควา อโวจ, ตตฺเถตํ อิติ วุจฺจติ\csromandash{}}

\csromanmatikaanchor{mtk.189}
\csromantocmarkref{subsubsection}{3. นิสฺสรณิยสุตฺต}{mtk.189}
\bookmark[dest={mtk.189},level=3]{3. นิสฺสรณิยสุตฺต}
\prose{“ตโย วิตกฺเก กุสเล วิตกฺกเย}

\prose{ตโย ปน อกุสเล นิรากเร.}

\csromangathameasure{ส เว วิตกฺกานิ วิจาริตานิ, \\ สเมติ วุฏฺฐีว รชํ สมูหตํ. \\ ส เว วิตกฺกูปสเมน เจตสา, \\ อิเธว โส สนฺติปทํ สมชฺฌคา”ติ.}
\csromangathagroup{\csromangathastanza{ส เว วิตกฺกานิ วิจาริตานิ, \\ สเมติ วุฏฺฐีว รชํ สมูหตํ. \\ ส เว วิตกฺกูปสเมน เจตสา, \\ อิเธว โส สนฺติปทํ สมชฺฌคา”ติ.}}

\prose{อยมฺปิ อตฺโถ วุตฺโต ภควตา อิติ เม สุตนฺติ.\csromanspacer{}\csromanspacer{}อฏฺฐมํ.}

\gambhira{อิติวุตฺตกปาฬิ}
\csromanheader{3}{9. อนฺตรามลสุตฺต}
\csromanheader{2}{88. วุตฺตํ เหตํ ภควตา วุตฺตมรหตาติ เม สุตํ\csromandash{}}
\csromanmatikaanchor{mtk.190}
\csromantocmarkref{subsubsection}{4. สนฺตตรสุตฺต}{mtk.190}
\bookmark[dest={mtk.190},level=3]{4. สนฺตตรสุตฺต}
\prose{\csromanfolio{252}ตโยเม ภิกฺขเว อนฺตรามลา อนฺตรา-อมิตฺตา อนฺตราสปตฺตา อนฺตราวธกา อนฺตราปจฺจตฺถิกา.\csromanspacer{} กตเม ตโย? โลโภ ภิกฺขเว อนฺตรามโล อนฺตรา-อมิตฺโต อนฺตราสปตฺโต อนฺตราวธโก อนฺตราปจฺจตฺถิโก, โทโส ภิกฺขเว อนฺตรามโล อนฺตรา-อมิตฺโต อนฺตราสปตฺโต อนฺตราวธโก อนฺตราปจฺจตฺถิโก, โมโห ภิกฺขเว อนฺตรามโล อนฺตรา-อมิตฺโต อนฺตราสปตฺโต อนฺตราวธโก อนฺตราปจฺจตฺถิโก.\csromanspacer{} อิเม โข ภิกฺขเว ตโย อนฺตรามลา อนฺตรา-อมิตฺตา อนฺตราสปตฺตา อนฺตราวธกา อนฺตราปจฺจตฺถิกาติ.\csromanspacer{} เอตมตฺถํ ภควา อโวจ, ตตฺเถตํ อิติ วุจฺจติ\csromandash{}}

\csromanmatikaanchor{mtk.191}
\csromantocmarkref{subsubsection}{5. ปุตฺตสุตฺต}{mtk.191}
\bookmark[dest={mtk.191},level=3]{5. ปุตฺตสุตฺต}
\csromangathameasure{*“อนตฺถชนโน โลโภ, โลโภ จิตฺตปฺปโกปโน. \\ ภยมนฺตรโต ชาตํ, ตํ ชโน นาวพุชฺฌติ. \\ ลุทฺโธ อตฺถํ น ชานาติ, ลุทฺโธ ธมฺมํ น ปสฺสติ. \\ อนฺธตมํ ตทา โหติ, ยํ โลโภ สหเต นรํ. \\ โย จ โลภํ ปหนฺตฺวาน, โลภเนยฺเย น ลุพฺภติ. \\ โลโภ ปหียเต ตมฺหา, อุทพินฺทูว โปกฺขรา. \\ อนตฺถชนโน โทโส, โทโส จิตฺตปฺปโกปโน. \\ ภยมนฺตรโต ชาตํ, ตํ ชโน นาวพุชฺฌติ. \\ ทุฏฺโฐ อตฺถํ น ชานาติ, ทุฏฺโฐ ธมฺมํ น ปสฺสติ. \\ อนฺธตมํ ตทา โหติ, ยํ โทโส สหเต นรํ. \\ โย จ โทสํ ปหนฺตฺวาน, โทสเนยฺเย น ทุสฺสติ. \\ โทโส ปหียเต ตมฺหา, ตาลปกฺกํว พนฺธนา. \\ อนตฺถชนโน โมโห, โมโห จิตฺตปฺปโกปโน. \\ ภยมนฺตรโต ชาตํ, ตํ ชโน นาวพุชฺฌติ. \\ มูโฬฺห อตฺถํ น ชานาติ, มูโฬฺห ธมฺมํ น ปสฺสติ. \\ อนฺธตมํ ตทา โหติ, ยํ โมโห สหเต นรํ.}
\csromangathagroup{\csromangathastanza{\csromansymbolfootnote{*}{อํ 2. 471 ปิฏฺเฐปิ.}“อนตฺถชนโน โลโภ, โลโภ จิตฺตปฺปโกปโน. \\ ภยมนฺตรโต ชาตํ, ตํ ชโน นาวพุชฺฌติ.}\csromangathastanza{ลุทฺโธ อตฺถํ น ชานาติ, ลุทฺโธ ธมฺมํ น ปสฺสติ. \\ อนฺธตมํ\footnote{อนฺธํ ตมํ (สี)} ตทา โหติ, ยํ โลโภ สหเต นรํ.}\csromangathastanza{โย จ โลภํ ปหนฺตฺวาน, โลภเนยฺเย น ลุพฺภติ. \\ โลโภ ปหียเต ตมฺหา, อุทพินฺทูว โปกฺขรา.}\csromangathastanza{อนตฺถชนโน โทโส, โทโส จิตฺตปฺปโกปโน. \\ ภยมนฺตรโต ชาตํ, ตํ ชโน นาวพุชฺฌติ.}\csromangathastanza{ทุฏฺโฐ อตฺถํ น ชานาติ, ทุฏฺโฐ ธมฺมํ น ปสฺสติ. \\ อนฺธตมํ ตทา โหติ, ยํ โทโส สหเต นรํ.}\csromangathastanza{โย จ โทสํ ปหนฺตฺวาน, โทสเนยฺเย น ทุสฺสติ. \\ โทโส ปหียเต ตมฺหา, ตาลปกฺกํว พนฺธนา.}\csromangathastanza{อนตฺถชนโน โมโห, โมโห จิตฺตปฺปโกปโน. \\ ภยมนฺตรโต ชาตํ, ตํ ชโน นาวพุชฺฌติ.}\csromangathastanza{มูโฬฺห อตฺถํ น ชานาติ, มูโฬฺห ธมฺมํ น ปสฺสติ. \\ อนฺธตมํ ตทา โหติ, ยํ โมโห สหเต นรํ.}}

\csromanheader{2}{3. ติกนิปาต}
\csromangathameasure{โย จ โมหํ ปหนฺตฺวาน, โมหเนยฺเย น มุยฺหติ. \\ โมหํ วิหนฺติ โส สพฺพํ, อาทิจฺโจวุทยํ ตมนฺ”ติ.}
\csromangathagroup{\csromangathastanza{\csromanfolio{253}โย จ โมหํ ปหนฺตฺวาน, โมหเนยฺเย น มุยฺหติ. \\ โมหํ วิหนฺติ โส สพฺพํ, อาทิจฺโจวุทยํ ตมนฺ”ติ.}}

\prose{อยมฺปิ อตฺโถ วุตฺโต ภควตา อิติ เม สุตนฺติ.\csromanspacer{}\csromanspacer{}นวมํ.}

\csromanheader{3}{10. เทวทตฺตสุตฺต}
\csromanheader{2}{89. วุตฺตํ เหตํ ภควตา วุตฺตมรหตาติ เม สุตํ\csromandash{}}
\prose{\csromansymbolfootnote{*}{วิ 4. 366 ปิฏฺเฐปิ.}ตีหิ ภิกฺขเว อสทฺธมฺเมหิ อภิภูโต ปริยาทินฺนจิตฺโต เทวทตฺโต อาปายิโก เนรยิโก กปฺปฏฺโฐ อเตกิจฺโฉ.\csromanspacer{} กตเมหิ ตีหิ? ปาปิจฺฉตาย ภิกฺขเว อภิภูโต ปริยาทินฺนจิตฺโต เทวทตฺโต อาปายิโก เนรยิโก กปฺปฏฺโฐ อเตกิจฺโฉ, ปาปมิตฺตตาย ภิกฺขเว อภิภูโต ปริยาทินฺนจิตฺโต เทวทตฺโต อาปายิโก เนรยิโก กปฺปฏฺโฐ อเตกิจฺโฉ, สติ โข ปน อุตฺตริกรณีเย\footnote{อุตฺตริํ กรณีเย (สฺยา)} โอรมตฺตเกน วิเสสาธิคเมน\footnote{วิเสสาธิคเมน จ (สฺยา, อิ)} อนฺตรา โวสานํ อาปาทิ.\csromanspacer{} อิเมหิ โข ภิกฺขเว ตีหิ อสทฺธมฺเมหิ อภิภูโต ปริยาทินฺนจิตฺโต เทวทตฺโต อาปายิโก เนรยิโก กปฺปฏฺโฐ อเตกิจฺโฉติ.\csromanspacer{} เอตมตฺถํ ภควา อโวจ, ตตฺเถตํ อิติ วุจฺจติ\csromandash{}}

\csromanmatikaanchor{mtk.192}
\csromantocmarkref{subsubsection}{6. อวุฏฺฐิกสุตฺต}{mtk.192}
\bookmark[dest={mtk.192},level=3]{6. อวุฏฺฐิกสุตฺต}
\csromangathameasure{“มา ชาตุ โกจิ โลกสฺมิํ, ปาปิจฺโฉ อุทปชฺชถ. \\ ตทมินาปิ ชานาถ, ปาปิจฺฉานํ ยถา คติ. \\ “ปณฺฑิโต”ติ สมญฺญาโต, “ภาวิตตฺโต”ติ สมฺมโต. \\ ชลํว ยสสา อฏฺฐา, “เทวทตฺโต”ติ วิสฺสุโต. \\ โส สมานมนุจิณฺโณ, อาสชฺช นํ ตถาคตํ. \\ อวีจินิรยํ ปตฺโต, จตุทฺวารํ ภยานกํ. \\ อทุฏฺฐสฺส หิ โย ทุพฺเภ, ปาปกมฺมํ อกุพฺพโต. \\ ตเมว ปาปํ ผุสติ, ทุฏฺฐจิตฺตํ อนาทรํ. \\ สมุทฺทํ วิสกุมฺเภน, โย มญฺเญยฺย ปทูสิตุํ. \\ น โส เตน ปทูเสยฺย, เภสฺมา หิ อุทธิ มหา.}
\csromangathagroup{\csromangathastanza{“มา ชาตุ โกจิ โลกสฺมิํ, ปาปิจฺโฉ อุทปชฺชถ. \\ ตทมินาปิ ชานาถ, ปาปิจฺฉานํ ยถา คติ.}\csromangathastanza{“ปณฺฑิโต”ติ สมญฺญาโต, “ภาวิตตฺโต”ติ สมฺมโต. \\ ชลํว ยสสา อฏฺฐา, “เทวทตฺโต”ติ วิสฺสุโต\footnote{เม สุตํ (ปาฬิยํ)}.}\csromangathastanza{โส สมานมนุจิณฺโณ\footnote{ปมาทมนุจิณฺโณ (ก-สี, สฺยา, อิ), ปมาณมนุจิณฺโณ (ก)}, อาสชฺช นํ ตถาคตํ. \\ อวีจินิรยํ ปตฺโต, จตุทฺวารํ ภยานกํ.}\csromangathastanza{อทุฏฺฐสฺส หิ โย ทุพฺเภ, ปาปกมฺมํ อกุพฺพโต. \\ ตเมว ปาปํ ผุสติ\footnote{ผุสฺเสติ (สฺยา)}, ทุฏฺฐจิตฺตํ อนาทรํ.}\csromangathastanza{สมุทฺทํ วิสกุมฺเภน, โย มญฺเญยฺย ปทูสิตุํ. \\ น โส เตน ปทูเสยฺย, เภสฺมา หิ อุทธิ มหา.}}

\gambhira{อิติวุตฺตกปาฬิ}
\csromangathameasure{เอวเมวํ ตถาคตํ, โย วาเทน วิหิํสติ. \\ สมฺมคฺคตํ สนฺตจิตฺตํ, วาโท ตมฺหิ น รูหติ. \\ ตาทิสํ มิตฺตํ กุพฺเพถ, ตญฺจ เสเวยฺย ปณฺฑิโต. \\ ยสฺส มคฺคานุโค ภิกฺขุ, ขยํ ทุกฺขสฺส ปาปุเณ”ติ.}
\csromangathagroup{\csromangathastanza{\csromanfolio{254}เอวเมวํ\footnote{เอวเมตํ (สฺยา)} ตถาคตํ, โย วาเทน วิหิํสติ. \\ สมฺมคฺคตํ\footnote{เอวเมตํ (สฺยา)} สนฺตจิตฺตํ, วาโท ตมฺหิ น รูหติ.}\csromangathastanza{ตาทิสํ มิตฺตํ กุพฺเพถ, ตญฺจ เสเวยฺย ปณฺฑิโต. \\ ยสฺส มคฺคานุโค ภิกฺขุ, ขยํ ทุกฺขสฺส ปาปุเณ”ติ.}}

\prose{อยมฺปิ อตฺโถ วุตฺโต ภควตา อิติ เม สุตนฺติ.\csromanspacer{}\csromanspacer{}ทสมํ.\csromanspacer{} จตุตฺโถ วคฺโค.}

\titlehead{\textbf{ตสฺสุทฺทานํ}}
\csromangathameasure{วิตกฺกาสกฺการสทฺท, จวนโลเก อสุภํ. \\ ธมฺม-อนฺธการมลํ, เทวทตฺเตน เต ทสาติ.}
\csromangathagroup{\csromangathastanza{วิตกฺกาสกฺการสทฺท, จวนโลเก อสุภํ. \\ ธมฺม-อนฺธการมลํ, เทวทตฺเตน เต ทสาติ.}}

\csromanheader{2}{5. ปญฺจมวคฺค}
\csromanheader{3}{1. อคฺคปฺปสาทสุตฺต}
\csromanheader{2}{90. วุตฺตํ เหตํ ภควตา วุตฺตมรหตาติ เม สุตํ\csromandash{}}
\prose{\csromansymbolfootnote{*}{อํ 1. 343; ขุ 10. 162 ปิฏฺเฐสุปิ.}ตโยเม ภิกฺขเว อคฺคปฺปสาทา.\csromanspacer{} กตเม ตโย? ยาวตา ภิกฺขเว สตฺตา อปทา วา ทฺวิปทา วา จตุปฺปทา วา พหุปฺปทา\footnote{พหุปทา (ก)} วา รูปิโน วา อรูปิโน วา สญฺญิโน วา อสญฺญิโน วา เนวสญฺญินาสญฺญิโน วา, ตถาคโต เตสํ อคฺคมกฺขายติ อรหํ สมฺมาสมฺพุทฺโธ.\csromanspacer{} เย ภิกฺขเว พุทฺเธ ปสนฺนา, อคฺเค เต ปสนฺนา.\csromanspacer{} อคฺเค โข ปน ปสนฺนานํ อคฺโค วิปาโก โหติ.}

\prose{ยาวตา ภิกฺขเว ธมฺมา สงฺขตา วา อสงฺขตา วา, วิราโค เตสํ อคมกฺขายติ ยทิทํ มทนิมฺมทโน ปิปาสวินโย อาลยสมุคฺฆาโต วฏฺฏุปจฺเฉโท ตณฺหกฺขโย วิราโค นิโรโธ นิพฺพานํ.\csromanspacer{} เย ภิกฺขเว วิราเค ธมฺเม ปสนฺนา, อคฺเค เต ปสนฺนา.\csromanspacer{} อคฺเค โข ปน ปสนฺนานํ อคฺโค วิปาโก โหติ.}

\csromanmatikaanchor{mtk.193}
\csromantocmarkref{subsubsection}{7. สุขปตฺถนาสุตฺต}{mtk.193}
\bookmark[dest={mtk.193},level=3]{7. สุขปตฺถนาสุตฺต}
\prose{ยาวตา ภิกฺขเว สํฆา วา คณา วา, ตถาคตสาวกสํโฆ เตสํ อคฺคมกฺขายติ ยทิทํ จตฺตาริ ปุริสยุคานิ อฏฺฐ ปุริสปุคฺคลา,}

\vspace{\tipitakaparskip}
\csromancenter{ติกนิปาต เอส ภควตา สาวกสํโฆ อาหุเนยฺโย ปาหุเนยฺโย ทกฺขิเณยฺโย อญฺชลิกรณีโย อนุตฺตรํ ปุญฺญกฺเขตฺตํ โลกสฺส.\csromanspacer{} เย ภิกฺขเว สํเฆ ปสนฺนา, อคฺเค เต ปสนฺนา.\csromanspacer{} อคฺเค โข ปน ปสนฺนานํ อคฺโค วิปาโก โหติ.\csromanspacer{} อิเม โข ภิกฺขเว ตโย อคฺคปฺปสาทาติ.\csromanspacer{} เอตมตฺถํ ภควา อโวจ, ตตฺเถตํ อิติ วุจฺจติ\csromandash{}}
\vspace{0.5\baselineskip}
\csromangathameasure{“อคฺคโต เว ปสนฺนานํ, อคฺคํ ธมฺมํ วิชานตํ. \\ อคฺเค พุทฺเธ ปสนฺนานํ, ทกฺขิเณยฺเย อนุตฺตเร. \\ อคฺเค ธมฺเม ปสนฺนานํ, วิราคูปสเม สุเข. \\ อคฺเค สํเฆ ปสนฺนานํ, ปุญฺญกฺเขตฺเต อนุตฺตเร. \\ อคฺคสฺมิํ ทานํ ททตํ, อคฺคํ ปุญฺญํ ปวฑฺฒติ. \\ อคฺคํ อายุ จ วณฺโณ จ, ยโส กิตฺติ สุขํ พลํ. \\ อคฺคสฺส ทาตา เมธาวี, อคฺคธมฺมสมาหิโต. \\ เทวภูโต มนุสฺโส วา, อคฺคปฺปตฺโต ปโมทตี”ติ.}
\csromangathagroup{\csromangathastanza{\csromanfolio{255}“อคฺคโต เว ปสนฺนานํ, อคฺคํ ธมฺมํ วิชานตํ. \\ อคฺเค พุทฺเธ ปสนฺนานํ, ทกฺขิเณยฺเย อนุตฺตเร.}\csromangathastanza{อคฺเค ธมฺเม ปสนฺนานํ, วิราคูปสเม สุเข. \\ อคฺเค สํเฆ ปสนฺนานํ, ปุญฺญกฺเขตฺเต อนุตฺตเร.}\csromangathastanza{อคฺคสฺมิํ ทานํ ททตํ, อคฺคํ ปุญฺญํ ปวฑฺฒติ. \\ อคฺคํ อายุ จ วณฺโณ จ, ยโส กิตฺติ สุขํ พลํ.}\csromangathastanza{อคฺคสฺส ทาตา เมธาวี, อคฺคธมฺมสมาหิโต. \\ เทวภูโต มนุสฺโส วา, อคฺคปฺปตฺโต ปโมทตี”ติ.}}

\prose{อยมฺปิ อตฺโถ วุตฺโต ภควตา อิติ เม สุตนฺติ.\csromanspacer{}\csromanspacer{}ปฐมํ.}

\csromanheader{3}{2. ชีวิกสุตฺต}
\csromanheader{2}{91. วุตฺตํ เหตํ ภควตา วุตฺตมรหตาติ เม สุตํ\csromandash{}}
\prose{อนฺตมิทํ ภิกฺขเว ชีวิกานํ ยทิทํ ปิณฺโฑลฺยํ, อภิสาโปยํ\footnote{อภิสาปายํ (สี), อภิลาปายํ (สฺยา, อิ), อภิสปายํ (ก)} ภิกฺขเว โลกสฺมิํ “ปิณฺโฑโล วิจรสิ ปตฺตปาณี”ติ.\csromanspacer{} ตญฺจ โข เอตํ ภิกฺขเว กุลปุตฺตา อุเปนฺติ อตฺถวสิกา อตฺถวสํ ปฏิจฺจ, เนว ราชาภินีตา น โจราภินีตา น อิณฏฺฏา น ภยฏฺฏา น อาชีวิกาปกตา, อปิ จ โข “โอติณฺณมฺหา ชาติยา ชราย มรเณน โสเกหิ ปริเทเวหิ ทุกฺเขหิ โทมนสฺเสหิ อุปายาเสหิ ทุกฺโขติณฺณา ทุกฺขปเรตา, อปฺเปว นาม อิมสฺส เกวลสฺส ทุกฺขกฺขนฺธสฺส อนฺตกิริยา ปญฺญาเยถา”ติ.\csromanspacer{} เอวํ ปพฺพชิโต จายํ ภิกฺขเว กุลปุตฺโต.\csromanspacer{} โส จ โหติ อภิชฺฌาลุ กาเมสุ ติพฺพสาราโค พฺยาปนฺนจิตฺโต ปทุฏฺฐมนสงฺกปฺโป มุฏฺฐสฺสติ อสมฺปชาโน อสมาหิโต วิพฺภนฺตจิตฺโต ปากตินฺทฺริโย.\csromanspacer{} เสยฺยถาปิ ภิกฺขเว ฉวาลาตํ อุภโตปทิตฺตํ มชฺเฌ คูถคตํ เนว คาเม กฏฺฐตฺถํ}

\gambhira{อิติวุตฺตกปาฬิ}
\prose{\csromanfolio{256}ผรติ น อรญฺเญ, ตถูปมาหํ ภิกฺขเว อิมํ ปุคฺคลํ วทามิ “คิหิโภคา ปริหีโน, สามญฺญตฺถญฺจ น ปริปูเรตี”ติ.\csromanspacer{} เอตมตฺถํ ภควา อโวจ, ตตฺเถตํ อิติ วุจฺจติ\csromandash{}}

\csromanmatikaanchor{mtk.194}
\csromantocmarkref{subsubsection}{8. ภิทุรสุตฺต}{mtk.194}
\bookmark[dest={mtk.194},level=3]{8. ภิทุรสุตฺต}
\csromangathameasure{“คิหิโภคา ปริหีโน, สามญฺญตฺถญฺจ ทุพฺภโค. \\ ปริธํสมาโน ปกิเรติ, ฉวาลาตํว นสฺสติ. \\ *กาสาวกณฺฐา พหโว, ปาปธมฺมา อสญฺญตา. \\ ปาปา ปาเปหิ กมฺเมหิ, นิรยํ เต อุปปชฺชเร. \\ *เสยฺโย อโยคุโฬ ภุตฺโต, ตตฺโต อคฺคิสิขูปโม. \\ ยญฺเจ ภุญฺเชยฺย ทุสฺสีโล, รฏฺฐปิณฺฑมสญฺญโต”ติ.}
\csromangathagroup{\csromangathastanza{“คิหิโภคา ปริหีโน, สามญฺญตฺถญฺจ ทุพฺภโค. \\ ปริธํสมาโน ปกิเรติ, ฉวาลาตํว นสฺสติ.}\csromangathastanza{\csromansymbolfootnote{*}{วิ 1. 116; เหฏฺฐา 224 ปิฏฺเฐสุปิ.}กาสาวกณฺฐา พหโว, ปาปธมฺมา อสญฺญตา. \\ ปาปา ปาเปหิ กมฺเมหิ, นิรยํ เต อุปปชฺชเร.}\csromangathastanza{\csromansymbolmark{*}เสยฺโย อโยคุโฬ ภุตฺโต, ตตฺโต อคฺคิสิขูปโม. \\ ยญฺเจ ภุญฺเชยฺย ทุสฺสีโล, รฏฺฐปิณฺฑมสญฺญโต”ติ.}}

\prose{อยมฺปิ อตฺโถ วุตฺโต ภควตา อิติ เม สุตนฺติ.\csromanspacer{}\csromanspacer{}ทุติยํ.}

\csromanheader{3}{3. สงฺฆาฏิกณฺณสุตฺต}
\csromanmatikaanchor{mtk.195}
\csromantocmarkref{subsubsection}{9. ธาตุโสสํสนฺทนสุตฺต}{mtk.195}
\bookmark[dest={mtk.195},level=3]{9. ธาตุโสสํสนฺทนสุตฺต}
\csromanheader{2}{92. วุตฺตํ เหตํ ภควตา วุตฺตมรหตาติ เม สุตํ\csromandash{}}
\prose{สงฺฆาฏิกณฺเณ เจปิ ภิกฺขเว ภิกฺขุ คเหตฺวา ปิฏฺฐิโต ปิฏฺฐิโต อนุพนฺโธ อสฺส ปาเท ปาทํ นิกฺขิปนฺโต, โส จ โหติ อภิชฺฌาลุ กาเมสุ ติพฺพสาราโค พฺยาปนฺนจิตฺโต ปทุฏฺฐมนสงฺกปฺโป มุฏฺฐสฺสติ อสมฺปชาโน อสมาหิโต วิพฺภนฺตจิตฺโต ปากตินฺทฺริโย, อถ โข โส อารกาว มยฺหํ, อหญฺจ ตสฺส.\csromanspacer{} ตํ กิสฺส เหตุ? ธมฺมํ หิ โส ภิกฺขเว ภิกฺขุ น ปสฺสติ, ธมฺมํ อปสฺสนฺโต น มํ ปสฺสติ\footnote{มํ น ปสฺสติ (สฺยา)}.}

\prose{โยชนสเต เจปิ โส ภิกฺขเว ภิกฺขุ วิหเรยฺย, โส จ โหติ อนภิชฺฌาลุ กาเมสุ น ติพฺพสาโรคา อพฺยาปนฺนจิตฺโต อปทุฏฺฐมนสงฺกปฺโป อุปฏฺฐิตสฺสติ สมฺปชาโน สมาหิโต เอกคฺคจิตฺโต สํวุตินฺทฺริโย, อถ โข โส สนฺติเกว มยฺหํ, อหญฺจ ตสฺส.\csromanspacer{} ตํ กิสฺส เหตุ? ธมฺมํ หิ โส ภิกฺขเว ภิกฺขุ ปสฺสติ, ธมฺมํ ปสฺสนฺโต มํ ปสฺสตีติ.\csromanspacer{} เอตมตฺถํ ภควา อโวจ, ตตฺเถตํ อิติ วุจฺจติ\csromandash{}}

\csromangathameasure{“อนุพนฺโธปิ เจ อสฺส, มหิจฺโฉ จ วิฆาตวา. \\ เอชานุโค อเนชสฺส, นิพฺพุตสฺส อนิพฺพุโต. \\ คิทฺโธ โส วีตเคธสฺส, ปสฺส ยาวญฺจ อารกา.}
\csromangathagroup{\csromangathastanza{“อนุพนฺโธปิ เจ อสฺส, มหิจฺโฉ จ วิฆาตวา. \\ เอชานุโค อเนชสฺส, นิพฺพุตสฺส อนิพฺพุโต. \\ คิทฺโธ โส วีตเคธสฺส, ปสฺส ยาวญฺจ อารกา.}}

\csromanheader{2}{3. ติกนิปาต}
\csromangathameasure{โย จ ธมฺมมภิญฺญาย, ธมฺมมญฺญาย ปณฺฑิโต. \\ รหโทว นิวาเต จ, อเนโช วูปสมฺมติ. \\ อเนโช โส อเนชสฺส, นิพฺพุตสฺส จ นิพฺพุโต. \\ อคิทฺโธ วีตเคธสฺส, ปสฺส ยาวญฺจ สนฺติเก”ติ.}
\csromangathagroup{\csromangathastanza{\csromanfolio{257}โย จ ธมฺมมภิญฺญาย, ธมฺมมญฺญาย ปณฺฑิโต. \\ รหโทว นิวาเต จ, อเนโช วูปสมฺมติ.}\csromangathastanza{อเนโช โส อเนชสฺส, นิพฺพุตสฺส จ นิพฺพุโต. \\ อคิทฺโธ วีตเคธสฺส, ปสฺส ยาวญฺจ สนฺติเก”ติ.}}

\prose{อยมฺปิ อตฺโถ วุตฺโต ภควตา อิติ เม สุตนฺติ.\csromanspacer{}\csromanspacer{}ตติยํ.}

\csromanheader{3}{4. อคฺคิสุตฺต}
\csromanheader{2}{93. วุตฺตํ เหตํ ภควตา วุตฺตมรหตาติ เม สุตํ\csromandash{}}
\prose{ตโยเม ภิกฺขเว อคฺคี.\csromanspacer{} กตเม ตโย? ราคคฺคิ, โทสคฺคิ, โมหคฺคิ.\csromanspacer{} อิเม โข ภิกฺขเว ตโย อคฺคีติ.\csromanspacer{} เอตมตฺถํ ภควา อโวจ, ตตฺเถตํ อิติ วุจฺจติ\csromandash{}}

\csromanmatikaanchor{mtk.196}
\csromantocmarkref{subsubsection}{10. ปริหานสุตฺต}{mtk.196}
\bookmark[dest={mtk.196},level=3]{10. ปริหานสุตฺต}
\csromangathameasure{“ราคคฺคิ ทหติ มจฺเจ, รตฺเต กาเมสุ มุจฺฉิเต. \\ โทสคฺคิ ปน พฺยาปนฺเน, นเร ปาณาติปาติโน. \\ โมหคฺคิ ปน สมฺมูเฬฺห, อริยธมฺเม อโกวิเท. \\ เอเต อคฺคี อชานนฺตา, สกฺกายาภิรตา ปชา. \\ เต วุฑฺฒยนฺติ นิรยํ, ติรจฺฉานญฺจ โยนิโย. \\ อสุรํ เปตฺติวิสยํ, อมุตฺตา มารพนฺธนา. \\ เย จ รตฺตินฺทิวา ยุตฺตา, สมฺมาสมฺพุทฺธสาสเน. \\ เต นิพฺพาเปนฺติ ราคคฺคิํ, นิจฺจํ อสุภสญฺญิโน. \\ โทสคฺคิํ ปน เมตฺตาย, นิพฺพาเปนฺติ นรุตฺตมา. \\ โมหคฺคิํ ปน ปญฺญาย, ยายํ นิพฺเพธคามินี. \\ เต นิพฺพาเปตฺวา นิปกา, รตฺตินฺทิวมตนฺทิตา, \\ อเสสํ ปรินิพฺพนฺติ, \\ อเสสํ ทุกฺขมจฺจคุํ. \\ อริยทฺทสา เวทคุโน, \\ สมฺมทญฺญาย ปณฺฑิตา. ชาติกฺขยมภิญฺญาย, \\ นาคจฺฉนฺติ ปุนพฺภวนฺ”ติ.}
\csromangathagroup{\csromangathastanza{“ราคคฺคิ ทหติ มจฺเจ, รตฺเต กาเมสุ มุจฺฉิเต. \\ โทสคฺคิ ปน พฺยาปนฺเน, นเร ปาณาติปาติโน.}\csromangathastanza{โมหคฺคิ ปน สมฺมูเฬฺห, อริยธมฺเม อโกวิเท. \\ เอเต อคฺคี อชานนฺตา, สกฺกายาภิรตา ปชา.}\csromangathastanza{เต วุฑฺฒยนฺติ นิรยํ, ติรจฺฉานญฺจ โยนิโย. \\ อสุรํ เปตฺติวิสยํ, อมุตฺตา มารพนฺธนา.}\csromangathastanza{เย จ รตฺตินฺทิวา ยุตฺตา, สมฺมาสมฺพุทฺธสาสเน. \\ เต นิพฺพาเปนฺติ ราคคฺคิํ, นิจฺจํ อสุภสญฺญิโน.}\csromangathastanza{โทสคฺคิํ ปน เมตฺตาย, นิพฺพาเปนฺติ นรุตฺตมา. \\ โมหคฺคิํ ปน ปญฺญาย, ยายํ นิพฺเพธคามินี.}\csromangathastanza{เต นิพฺพาเปตฺวา นิปกา, รตฺตินฺทิวมตนฺทิตา, \\ อเสสํ ปรินิพฺพนฺติ, \\ อเสสํ ทุกฺขมจฺจคุํ. \\ อริยทฺทสา เวทคุโน, \\ สมฺมทญฺญาย ปณฺฑิตา. ชาติกฺขยมภิญฺญาย, \\ นาคจฺฉนฺติ ปุนพฺภวนฺ”ติ.}}

\prose{อยมฺปิ อตฺโถ วุตฺโต ภควตา อิติ เม สุตนฺติ.\csromanspacer{}\csromanspacer{}จตุตฺถํ.}

\gambhira{อิติวุตฺตกปาฬิ}
\csromanheader{3}{5. อุปปริกฺขสุตฺต}
\csromanheader{2}{94. วุตฺตํ เหตํ ภควตา วุตฺตมรหตาติ เม สุตํ\csromandash{}}
\csromanmatikaanchor{mtk.197}
\csromantocmarknopage{subsection}{4. จตุตฺถวคฺค}{mtk.197}
\bookmark[dest={mtk.197},level=2]{4. จตุตฺถวคฺค}
\prose{\csromanfolio{258}ตถา ตถา ภิกฺขเว ภิกฺขุ อุปปริกฺเขยฺย, ยถา ยถา’สฺส\footnote{ยถา ยถา (พหูสุ)} อุปปริกฺขโต พหิทฺธา จสฺส วิญฺญาณํ อวิกฺขิตฺตํ อวิสฏํ, อชฺฌตฺตํ อสณฺฐิตํ, อนุปาทาย น ปริตสฺเสยฺย.\csromanspacer{} พหิทฺธา ภิกฺขเว วิญฺญาเณ อวิกฺขิตฺเต อวิสเฏ สติ, อชฺฌตฺตํ อสณฺฐิเต, อนุปาทาย อปริตสฺสโต อายติํ ชาติชรามรณทุกฺขสมุทยสมฺภโว น โหตีติ.\csromanspacer{} เอตมตฺถํ ภควา อโวจ, ตตฺเถตํ อิติ วุจฺจติ\csromandash{}}

\csromanmatikaanchor{mtk.198}
\csromantocmarkref{subsubsection}{1. วิตกฺกสุตฺต}{mtk.198}
\bookmark[dest={mtk.198},level=3]{1. วิตกฺกสุตฺต}
\csromangathameasure{“สตฺตสงฺคปฺปหีนสฺส, เนตฺติจฺฉินฺนสฺส ภิกฺขุโน. \\ วิกฺขีโณ ชาติสํสาโร, นตฺถิ ตสฺส ปุนพฺภโว”ติ.}
\csromangathagroup{\csromangathastanza{“สตฺตสงฺคปฺปหีนสฺส, เนตฺติจฺฉินฺนสฺส ภิกฺขุโน. \\ วิกฺขีโณ ชาติสํสาโร, นตฺถิ ตสฺส ปุนพฺภโว”ติ.}}

\prose{อยมฺปิ อตฺโถ วุตฺโต ภควตา อิติ เม สุตนฺติ.\csromanspacer{}\csromanspacer{}ปญฺจมํ.}

\csromanheader{3}{6. กามูปปตฺติสุตฺต}
\csromanheader{2}{95. วุตฺตํ เหตํ ภควตา วุตฺตมรหตาติ เม สุตํ\csromandash{}}
\prose{ติสฺโส อิมา ภิกฺขเว กามูปปตฺติโย\footnote{กามุปฺปตฺติโย (สี)}.\csromanspacer{} กตมา ติสฺโส? ปจฺจุปฏฺฐิตกามา, นิมฺมานรติโน, ปรนิมฺมิตวสวตฺติโน.\csromanspacer{} อิมา โข ภิกฺขเว ติสฺโส กามูปปตฺติโยติ.\csromanspacer{} เอตมตฺถํ ภควา อโวจ, ตตฺเถตํ อิติ วุจฺจติ\csromandash{}}

\csromanmatikaanchor{mtk.199}
\csromantocmarkref{subsubsection}{2. สกฺการสุตฺต}{mtk.199}
\bookmark[dest={mtk.199},level=3]{2. สกฺการสุตฺต}
\csromangathameasure{“ปจฺจุปฏฺฐิตกามา จ, เย เทวา วสวตฺติโน. \\ นิมฺมานรติโน เทวา, เย จญฺเญ กามโภคิโน. \\ อิตฺถภาวญฺญถาภาวํ, สํสารํ นาติวตฺตเร. \\ เอตมาทีนวํ ญตฺวา, กามโภเคสุ ปณฺฑิโต. \\ สพฺเพ ปริจฺจเช กาเม, เย ทิพฺพา เย จ มานุสา. \\ ปิยรูปสาตคธิตํ, เฉตฺวา โสตํ ทุรจฺจยํ. \\ อเสสํ ปรินิพฺพนฺติ, อเสสํ ทุกฺขมจฺจคุํ. \\ อริยทฺทสา เวทคุโน, สมฺมทญฺญาย ปณฺฑิตา. \\ ชาติกฺขยมภิญฺญาย, นาคจฺฉนฺติ ปุนพฺภวนฺ”ติ.}
\csromangathagroup{\csromangathastanza{“ปจฺจุปฏฺฐิตกามา จ, เย เทวา วสวตฺติโน. \\ นิมฺมานรติโน เทวา, เย จญฺเญ กามโภคิโน.}\csromangathastanza{อิตฺถภาวญฺญถาภาวํ, สํสารํ นาติวตฺตเร. \\ เอตมาทีนวํ ญตฺวา, กามโภเคสุ ปณฺฑิโต.}\csromangathastanza{สพฺเพ ปริจฺจเช กาเม, เย ทิพฺพา เย จ มานุสา. \\ ปิยรูปสาตคธิตํ, เฉตฺวา โสตํ ทุรจฺจยํ.}\csromangathastanza{อเสสํ ปรินิพฺพนฺติ, อเสสํ ทุกฺขมจฺจคุํ. \\ อริยทฺทสา เวทคุโน, สมฺมทญฺญาย ปณฺฑิตา. \\ ชาติกฺขยมภิญฺญาย, นาคจฺฉนฺติ ปุนพฺภวนฺ”ติ.}}

\prose{อยมฺปิ อตฺโถ วุตฺโต ภควตา อิติ เม สุตนฺติ.\csromanspacer{}\csromanspacer{}ฉฏฺฐํ.}

\csromanheader{3}{3. ติกนิปาต \textbf{7. กามโยคสุตฺต}}
\csromanheader{2}{96. วุตฺตํ เหตํ ภควตา วุตฺตมรหตาติ เม สุตํ\csromandash{}}
\prose{\csromanfolio{259}กามโยคยุตฺโต ภิกฺขเว ภวโยคยุตฺโต อาคามี โหติ อาคนฺตา\footnote{อาคนฺตฺวา (สฺยา, ก)} อิตฺถตฺตํ, กามโยควิสํยุตฺโต ภิกฺขเว ภวโยคยุตฺโต อนาคามี โหติ อนาคนฺตา อิตฺถตฺตํ, กามโยควิสํยุตฺโต ภิกฺขเว ภวโยควิสํยุตฺโต อรหา โหติ ขีณาสโวติ.\csromanspacer{} เอตมตฺถํ ภควา อโวจ, ตตฺเถตํ อิติ วุจฺจติ\csromandash{}}

\csromangathameasure{“กามโยเคน สํยุตฺตา, ภวโยเคน จูภยํ. \\ สตฺตา คจฺฉนฺติ สํสารํ, ชาติ มรณคามิโน. \\ เย จ กาเม ปหนฺตฺวาน, อปฺปตฺตา อาสวกฺขยํ. \\ ภวโยเคน สํยุตฺตา, “อนาคามี”ติ วุจฺจเร. \\ เย จ โข ฉินฺนสํสยา, ขีณมานปุนพฺภวา. \\ เต เว ปารงฺคตา โลเก, เย ปตฺตา อาสวกฺขยนฺ”ติ.}
\csromangathagroup{\csromangathastanza{“กามโยเคน สํยุตฺตา, ภวโยเคน จูภยํ. \\ สตฺตา คจฺฉนฺติ สํสารํ, ชาติ มรณคามิโน.}\csromangathastanza{เย จ กาเม ปหนฺตฺวาน, อปฺปตฺตา อาสวกฺขยํ. \\ ภวโยเคน สํยุตฺตา, “อนาคามี”ติ วุจฺจเร.}\csromangathastanza{เย จ โข ฉินฺนสํสยา, ขีณมานปุนพฺภวา. \\ เต เว ปารงฺคตา โลเก, เย ปตฺตา อาสวกฺขยนฺ”ติ.}}

\prose{อยมฺปิ อตฺโถ วุตฺโต ภควตา อิติ เม สุตนฺติ.\csromanspacer{}\csromanspacer{}สตฺตมํ.\csromanspacer{} ตติยภาณวารํ.}

\csromanheader{3}{8. กลฺยาณสีลสุตฺต}
\csromanheader{2}{97. วุตฺตํ เหตํ ภควตา วุตฺตมรหตาติ เม สุตํ\csromandash{}}
\csromanmatikaanchor{mtk.200}
\csromantocmarkref{subsubsection}{3. เทวสทฺทสุตฺต}{mtk.200}
\bookmark[dest={mtk.200},level=3]{3. เทวสทฺทสุตฺต}
\prose{กลฺยาณสีโล ภิกฺขเว ภิกฺขุ กลฺยาณธมฺโม กลฺยาณปญฺโญ อิมสฺมิํ ธมฺมวินเย “เกวลี วุสิตวา อุตฺตมปุริโส”ติ วุจฺจติ.}

\prose{กถญฺจ ภิกฺขเว ภิกฺขุ กลฺยาณสีโล โหติ? อิธ ภิกฺขเว ภิกฺขุ สีลวา โหติ, ปาติโมกฺขสํวรสํวุโต วิหรติ อาจารโคจรสมฺปนฺโน, อณุมตฺเตสุ วชฺเชสุ ภยทสฺสาวี สมาทาย สิกฺขติ สิกฺขาปเทสุ.\csromanspacer{} เอวํ โข ภิกฺขเว ภิกฺขุ กลฺยาณสีโล โหติ, อิติ กลฺยาณสีโล.}

\gambhira{อิติวุตฺตกปาฬิ}
\prose{\csromanfolio{260}กลฺยาณธมฺโม จ กถํ โหติ? อิธ ภิกฺขเว ภิกฺขุ สตฺตนฺนํ โพธิปกฺขิยานํ ธมฺมานํ ภาวนานุโยคมนุยุตฺโต วิหรติ.\csromanspacer{} เอวํ โข ภิกฺขเว ภิกฺขุ กลฺยาณธมฺโม โหติ, อิติ กลฺยาณสีโล กลฺยาณธมฺโม.}

\prose{กลฺยาณปญฺโญ จ กถํ โหติ? อิธ ภิกฺขเว ภิกฺขุ อาสวานํ ขยา อนาสวํ เจโตวิมุตฺติํ ปญฺญาวิมุตฺติํ ทิฏฺเฐว ธมฺเม สยํ อภิญฺญา สจฺฉิกตฺวา อุปสมฺปชฺช วิหรติ.\csromanspacer{} เอวํ โข ภิกฺขเว ภิกฺขุ กลฺยาณปญฺโญ โหติ.}

\prose{อิติ กลฺยาณสีโล กลฺยาณธมฺโม กลฺยาณปญฺโญ อิมสฺมิํ ธมฺมวินเย “เกวลี วุสิตวา อุตฺตมปุริโส”ติ วุจฺจตีติ.\csromanspacer{} เอตมตฺถํ ภควา อโวจ, ตตฺเถตํ อิติ วุจฺจติ\csromandash{}}

\csromangathameasure{“ยสฺส กาเยน วาจาย, มนสา นตฺถิ ทุกฺกฏํ. \\ ตํ เว “กลฺยาณสีโล”ติ, อาหุ ภิกฺขุํ หิรีมนํ. \\ ยสฺส ธมฺมา สุภาวิตา, สตฺต สมฺโพธิคามิโน. \\ ตํ เว “กลฺยาณธมฺโม”ติ, อาหุ ภิกฺขุํ อนุสฺสทํ. \\ โย ทุกฺขสฺส ปชานาติ, อิเธว ขยมตฺตโน. \\ ตํ เว “กลฺยาณปญฺโญ”ติ, อาหุ ภิกฺขุํ อนาสวํ. \\ เตหิ ธมฺเมหิ สมฺปนฺนํ, อนีฆํ ฉินฺนสํสยํ. \\ อสิตํ สพฺพโลกสฺส, อาหุ สพฺพปหายินนฺ”ติ.}
\csromangathagroup{\csromangathastanza{“ยสฺส กาเยน วาจาย, มนสา นตฺถิ ทุกฺกฏํ. \\ ตํ เว “กลฺยาณสีโล”ติ, อาหุ ภิกฺขุํ หิรีมนํ\footnote{หิรีมตํ (สฺยา, อิ)}.}\csromangathastanza{ยสฺส ธมฺมา สุภาวิตา, สตฺต\footnote{ปตฺต (สพฺพตฺถ)} สมฺโพธิคามิโน. \\ ตํ เว “กลฺยาณธมฺโม”ติ, อาหุ ภิกฺขุํ อนุสฺสทํ.}\csromangathastanza{โย ทุกฺขสฺส ปชานาติ, อิเธว ขยมตฺตโน. \\ ตํ เว “กลฺยาณปญฺโญ”ติ, อาหุ ภิกฺขุํ อนาสวํ.}\csromangathastanza{เตหิ ธมฺเมหิ สมฺปนฺนํ, อนีฆํ ฉินฺนสํสยํ. \\ อสิตํ สพฺพโลกสฺส, อาหุ สพฺพปหายินนฺ”ติ.}}

\csromanmatikaanchor{mtk.201}
\csromantocmarkref{subsubsection}{4. ปญฺจปุพฺพนิมิตฺตสุตฺต}{mtk.201}
\bookmark[dest={mtk.201},level=3]{4. ปญฺจปุพฺพนิมิตฺตสุตฺต}
\prose{อยมฺปิ อตฺโถ วุตฺโต ภควตา อิติ เม สุตนฺติ.\csromanspacer{}\csromanspacer{}อฏฺฐมํ.}

\csromanheader{3}{9. ทานสุตฺต}
\csromanheader{2}{98. วุตฺตํ เหตํ ภควตา วุตฺตมรหตาติ เม สุตํ\csromandash{}}
\prose{เทฺวมานิ ภิกฺขเว ทานานิ\csromandash{}อามิสทานญฺจ, ธมฺมทานญฺจ, เอตทคฺคํ ภิกฺขเว อิเมสํ ทฺวินฺนํ ทานานํ ยทิทํ ธมฺมทานํ.}

\prose{เทฺวเม ภิกฺขเว สํวิภาคา\csromandash{}อามิสสํวิภาโค จ, ธมฺมสํวิภาโค จ, เอตทคฺคํ ภิกฺขเว อิเมสํ ทฺวินฺนํ สํวิภาคานํ ยทิทํ ธมฺมสํวิภาโค.}

\csromanheader{2}{3. ติกนิปาต}
\prose{\csromanfolio{261}เทฺวเม ภิกฺขเว อนุคฺคหา\csromandash{}อามิสานุคฺคโห จ, ธมฺมานุคฺคโห จ, เอตทคฺคํ ภิกฺขเว อิเมสํ ทฺวินฺนํ อนุคฺคหานํ ยทิทํ ธมฺมานุคฺคโหติ.\csromanspacer{} เอตมตฺถํ ภควา อโวจ, ตตฺเถตํ อิติ วุจฺจติ\csromandash{}}

\csromangathameasure{“ยมาหุ ทานํ ปรมํ อนุตฺตรํ, \\ ยํ สํวิภาคํ ภควา อวณฺณยิ. \\ อคฺคมฺหิ เขตฺตมฺหิ ปสนฺนจิตฺโต, \\ วิญฺญู ปชานํ โก น ยเชถ กาเล. \\ เย เจว ภาสนฺติ สุณนฺติ จูภยํ, \\ ปสนฺนจิตฺตา สุคตสฺส สาสเน. \\ เตสํ โส อตฺโถ ปรโม วิสุชฺฌติ, \\ เย อปฺปมตฺตา สุคตสฺส สาสเน”ติ.}
\csromangathagroup{\csromangathastanza{“ยมาหุ ทานํ ปรมํ อนุตฺตรํ, \\ ยํ สํวิภาคํ ภควา อวณฺณยิ\footnote{อวณฺณยี (สี)}. \\ อคฺคมฺหิ เขตฺตมฺหิ ปสนฺนจิตฺโต, \\ วิญฺญู ปชานํ โก น ยเชถ กาเล.}\csromangathastanza{เย เจว ภาสนฺติ สุณนฺติ จูภยํ, \\ ปสนฺนจิตฺตา สุคตสฺส สาสเน. \\ เตสํ โส อตฺโถ ปรโม วิสุชฺฌติ, \\ เย อปฺปมตฺตา สุคตสฺส สาสเน”ติ.}}

\prose{อยมฺปิ อตฺโถ วุตฺโต ภควตา อิติ เม สุตนฺติ.\csromanspacer{}\csromanspacer{}นวมํ.}

\csromanheader{3}{10. เตวิชฺชสุตฺต}
\csromanheader{2}{99. วุตฺตํ เหตํ ภควตา วุตฺตมรหตาติ เม สุตํ\csromandash{}}
\prose{ธมฺเมนาหํ ภิกฺขเว เตวิชฺชํ พฺราหฺมณํ ปญฺญาเปมิ, นาญฺญํ ลปิตลาปนมตฺเตน.}

\prose{กถญฺจาหํ ภิกฺขเว ธมฺเมน เตวิชฺชํ พฺราหฺมณํ ปญฺญาเปมิ, นาญฺญํ ลปิตลาปนมตฺเตน? อิธ ภิกฺขเว ภิกฺขุ อเนกวิหิตํ ปุพฺเพนิวาสํ อนุสฺสรติ.\csromanspacer{} เสยฺยถิทํ, เอกมฺปิ ชาติํ เทฺวปิ ชาติโย ติสฺโสปิ ชาติโย จตสฺโสปิ ชาติโย ปญฺจปิ ชาติโย ทสปิ ชาติโย วีสมฺปิ ชาติโย ติํสมฺปิ ชาติโย จตฺตาลีสมฺปิ ชาติโย ปญฺญาสมฺปิ ชาติโย ชาติสตมฺปิ ชาติสหสฺสมฺปิ ชาติสตสหสฺสมฺปิ อเนเกปิ สํวฏฺฏกปฺเป อเนเกปิ วิวฏฺฏกปฺเป อเนเกปิ สํวฏฺฏวิวฏฺฏกปฺเป “อมุตฺราสิํ เอวํนาโม เอวํโคตฺโต เอวํวณฺโณ เอวมาหาโร เอวํสุขทุกฺขปฺปฏิสํเวที เอวมายุปริยนฺโต, โส ตโต จุโต อมุตฺร อุทปาทิํ, ตตฺราปาสิํ เอวํนาโม เอวํโคตฺโต เอวํวณฺโณ เอวมาหาโร เอวํสุขทุกฺขปฺปฏิสํเวที เอวมายุปริยนฺโต, โส ตโต จุโต}

\csromanmatikaanchor{mtk.202}
\csromantocmarkref{subsubsection}{5. พหุชนหิตสุตฺต}{mtk.202}
\bookmark[dest={mtk.202},level=3]{5. พหุชนหิตสุตฺต}
\gambhira{อิติวุตฺตกปาฬิ}
\prose{\csromanfolio{262}อิธูปปนฺโน”ติ, อิติ สาการํ ส-อุทฺเทสํ อเนกวิหิตํ ปุพฺเพนิวาสํ อนุสฺสรติ.\csromanspacer{} อยมสฺส ปฐมา วิชฺชา อธิคตา โหติ, อวิชฺชา วิหตา, วิชฺชา อุปฺปนฺนา, ตโม วิหโต, อาโลโก อุปฺปนฺโน.\csromanspacer{} ยถา ตํ อปฺปมตฺตสฺส อาตาปิโน ปหิตตฺตสฺส วิหรโต.}

\prose{ปุน จปรํ ภิกฺขเว ภิกฺขุ ทิพฺเพน จกฺขุนา วิสุทฺเธน อติกฺกนฺตมานุสเกน สตฺเต ปสฺสติ จวมาเน อุปปชฺชมาเน หีเน ปณีเต สุวณฺเณ ทุพฺพณฺเณ สุคเต ทุคฺคเต, ยถากมฺมูปเค สตฺเต ปชานาติ “อิเม วต โภนฺโต สตฺตา กายทุจฺจริเตน สมนฺนาคตา วจีทุจฺจริเตน สมนฺนาคตา มโนทุจฺจริเตน สมนฺนาคตา อริยานํ อุปวาทกา มิจฺฉาทิฏฺฐิกา มิจฺฉาทิฏฺฐิกมฺมสมาทานา, เต กายสฺส เภทา ปรํ มรณา อปายํ ทุคฺคติํ วินิปาตํ นิรยํ อุปปนฺนา.\csromanspacer{} อิเม วา ปน โภนฺโต สตฺตา กายสุจริเตน สมนฺนาคตา วจีสุจริเตน สมนฺนาคตา มโนสุจริเตน สมนฺนาคตา อริยานํ อนุปวาทกา สมฺมาทิฏฺฐิกา สมฺมาทิฏฺฐิกมฺมสมาทานา, เต กายสฺส เภทา ปรํ มรณา สุคติํ สคฺคํ โลกํ อุปปนฺนา”ติ, อิติ ทิพฺเพน จกฺขุนา วิสุทฺเธน อติกฺกนฺตมานุสเกน สตฺเต ปสฺสติ จวมาเน อุปปชฺชมาเน หีเน ปณีเต สุวณฺเณ ทุพฺพณฺเณ สุคเต ทุคฺคเต, ยถากมฺมูปเค สตฺเต ปชานาติ.\csromanspacer{} อยมสฺส ทุติยา วิชฺชา อธิคตา โหติ, อวิชฺชา วิหตา, วิชฺชา อุปฺปนฺนา, ตโม วิหโต, อาโลโก อุปฺปนฺโน.\csromanspacer{} ยถา ตํ อปฺปมตฺตสฺส อาตาปิโน ปหิตตฺตสฺส วิหรโต.}

\prose{ปุน จปรํ ภิกฺขเว ภิกฺขุ อาสวานํ ขยา อนาสวํ เจโตวิมุตฺติํ ปญฺญาวิมุตฺติํ ทิฏฺเฐว ธมฺเม สยํ อภิญฺญา สจฺฉิกตฺวา อุปสมฺปชฺช วิหรติ.\csromanspacer{} อยมสฺส ตติยา วิชฺชา อธิคตา โหติ, อวิชฺชา วิหตา, วิชฺชา อุปฺปนฺนา, ตโม วิหโต, อาโลโก อุปฺปนฺโน.\csromanspacer{} ยถา ตํ อปฺปมตฺตสฺส อาตาปิโน ปหิตตฺตสฺส วิหรโต.\csromanspacer{} เอวํ โข อหํ ภิกฺขเว ธมฺเมน เตวิชฺชํ พฺราหฺมณํ ปญฺญาเปมิ, นาญฺญํ ลปิตลาปนมตฺเตนาติ.\csromanspacer{} เอตมตฺถํ ภควา อโวจ, ตตฺเถตํ อิติ วุจฺจติ\csromandash{}}

\csromangathameasure{*“ปุพฺเพนิวาสํ โย’เวติ, สคฺคาปายญฺจ ปสฺสติ. \\ อโถ ชาติกฺขยํ ปตฺโต, อภิญฺญาโวสิโต มุนิ.}
\csromangathagroup{\csromangathastanza{\csromansymbolmark{*}“ปุพฺเพนิวาสํ โย’เวติ\footnote{โยเวทิ (สพฺพตฺถ) เหฏฺฐา 74, อุปริ 378 ปิฏฺเฐสุปิ.}, สคฺคาปายญฺจ ปสฺสติ. \\ อโถ\footnote{โยเวทิ (สพฺพตฺถ) เหฏฺฐา 74, อุปริ 378 ปิฏฺเฐสุปิ.} ชาติกฺขยํ ปตฺโต, อภิญฺญาโวสิโต มุนิ.}}

\csromanheader{2}{3. ติกนิปาต}
\csromangathameasure{เอตาหิ ตีหิ วิชฺชาหิ, เตวิชฺโช โหติ พฺราหฺมโณ. \\ ตมหํ วทามิ เตวิชฺชํ, นาญฺญํ ลปิตลาปนนฺ”ติ.}
\csromangathagroup{\csromangathastanza{\csromanfolio{263}เอตาหิ ตีหิ วิชฺชาหิ, เตวิชฺโช โหติ พฺราหฺมโณ. \\ ตมหํ วทามิ เตวิชฺชํ, นาญฺญํ ลปิตลาปนนฺ”ติ.}}

\prose{อยมฺปิ อตฺโถ วุตฺโต ภควตา อิติ เม สุตนฺติ.\csromanspacer{}\csromanspacer{}ทสมํ.\csromanspacer{} ปญฺจโม วคฺโค.}

\titlehead{\textbf{ตสฺสุทฺทานํ}}
\csromangathameasure{ปสาท ชีวิต สงฺฆาฏิ, อคฺคิ อุปปริกฺขยา. \\ อุปปตฺติ กาม กลฺยาณํ, ทานํ ธมฺเมน เต ทสาติ.}
\csromangathagroup{\csromangathastanza{ปสาท ชีวิต สงฺฆาฏิ, อคฺคิ อุปปริกฺขยา. \\ อุปปตฺติ\footnote{อุปฺปตฺติ (สี)} กาม กลฺยาณํ, ทานํ ธมฺเมน เต ทสาติ.}}

\nitthitam{\textbf{ติกนิปาโต นิฏฺฐิโต.}}
\csromanmatikaanchor{mtk.203}
\csromantocmarkref{subsubsection}{6. อสุภานุปสฺสีสุตฺต}{mtk.203}
\bookmark[dest={mtk.203},level=3]{6. อสุภานุปสฺสีสุตฺต}
\csromanheader{1}{4. จตุกฺกนิปาต}
\csromanheader{2}{1. พฺราหฺมณธมฺมยาคสุตฺต}
\csromanheader{2}{100. วุตฺตํ เหตํ ภควตา วุตฺตมรหตาติ เม สุตํ\csromandash{}}
\prose{\csromanfolio{264}อหมสฺมิ ภิกฺขเว พฺราหฺมโณ ยาจโยโค สทา ปยตปาณิ\footnote{ปยตปาณี (สี, สฺยา)} อนฺติมเทหธโร อนุตฺตโร ภิสกฺโก สลฺลกตฺโต, ตสฺส เม ตุเมฺห ปุตฺตา โอรสา มุขโต ชาตา ธมฺมชา ธมฺมนิมฺมิตา ธมฺมทายาทา โน อามิสทายาทา.}

\prose{เทฺวมานิ ภิกฺขเว ทานานิ\csromandash{}อามิสทานญฺจ, ธมฺมทานญฺจ, เอตทคฺคํ ภิกฺขเว อิเมสํ ทฺวินฺนํ ทานานํ ยทิทํ ธมฺมทานํ.}

\prose{เทฺวเม ภิกฺขเว สํวิภาคา\csromandash{}อามิสสํวิภาโค จ, ธมฺมสํวิภาโค จ, เอตทคฺคํ ภิกฺขเว อิเมสํ ทฺวินฺนํ สํวิภาคานํ ยทิทํ ธมฺมสํวิภาโค}

\csromanmatikaanchor{mtk.204}
\csromantocmarkref{subsubsection}{7. ธมฺมานุธมฺมปฏิปนฺนสุตฺต}{mtk.204}
\bookmark[dest={mtk.204},level=3]{7. ธมฺมานุธมฺมปฏิปนฺนสุตฺต}
\prose{เทฺวเม ภิกฺขเว อนุคฺคหา\csromandash{}อามิสานุคฺคโห จ, ธมฺมานุคฺคโห จ, เอตทคฺคํ ภิกฺขเว อิเมสํ ทฺวินฺนํ อนุคฺคหานํ ยทิทํ ธมฺมานุคฺคโห.}

\prose{เทฺวเม ภิกฺขเว ยาคา\csromandash{}อามิสยาโค จ, ธมฺมยาโค จ, เอตทคฺคํ ภิกฺขเว อิเมสํ ทฺวินฺนํ ยาคานํ ยทิทํ ธมฺมยาโคติ.\csromanspacer{} เอตมตฺถํ ภควา อโวจ, ตตฺเถตํ อิติ วุจฺจติ\csromandash{}}

\csromangathameasure{“โย ธมฺมยาคํ อยชี อมจฺฉรี, \\ ตถาคโต สพฺพภูตานุกมฺปี. \\ ตํ ตาทิสํ เทวมนุสฺสเสฏฺฐํ, \\ สตฺตา นมสฺสนฺติ ภวสฺส ปารคุนฺ”ติ.}
\csromangathagroup{\csromangathastanza{“โย ธมฺมยาคํ อยชี อมจฺฉรี, \\ ตถาคโต สพฺพภูตานุกมฺปี\footnote{สพฺพสตฺตานุกมฺปี (สฺยา) อฏฺฐกถายมฺปิ.}. \\ ตํ ตาทิสํ เทวมนุสฺสเสฏฺฐํ, \\ สตฺตา นมสฺสนฺติ ภวสฺส ปารคุนฺ”ติ.}}

\prose{อยมฺปิ อตฺโถ วุตฺโต ภควตา อิติ เม สุตนฺติ.\csromanspacer{}\csromanspacer{}ปฐมํ.}

\csromanheader{2}{2. สุลภสุตฺต}
\csromanheader{2}{101. วุตฺตํ เหตํ ภควตา วุตฺตมรหตาติ เม สุตํ\csromandash{}}
\csromanmatikaanchor{mtk.205}
\csromantocmarkref{subsubsection}{8. อนฺธกรณสุตฺต}{mtk.205}
\bookmark[dest={mtk.205},level=3]{8. อนฺธกรณสุตฺต}
\prose{จตฺตาริมานิ ภิกฺขเว อปฺปานิ เจว สุลภานิ จ ตานิ จ อนวชฺชานิ.\csromanspacer{} กกมานิ จตฺตาริ? ปํสุกูลํ ภิกฺขเว จีวรานํ อปฺปญฺจ สุลภญฺจ ตญฺจ อนวชฺชํ,}

\vspace{\tipitakaparskip}
\csromancenter{จตุกฺกนิปาต ปิณฺฑิยาโลโป ภิกฺขเว โภชนานํ อปฺปญฺจ สุลภญฺจ ตญฺจ อนวชฺชํ, รุกฺขมูลํ ภิกฺขเว เสนาสนานํ อปฺปญฺจ สุลภญฺจ ตญฺจ อนวชฺชํ, ปูติมุตฺตํ ภิกฺขเว เภสชฺชานํ อปฺปญฺจ สุลภญฺจ ตญฺจ อนวชฺชํ.\csromanspacer{} อิมานิ โข ภิกฺขเว จตฺตาริ อปฺปานิ เจว สุลภานิ จ ตานิ จ อนวชฺชานิ.\csromanspacer{} ยโต โข ภิกฺขเว ภิกฺขุ อปฺเปน จ ตุฏฺโฐ โหติ สุลเภน จ (อนวชฺเชน จ)\footnote{(...) นตฺถิ สี-อิ-ก-โปตฺถเกสุ จ อํ 1. 335 ปิฏฺเฐ จ.}.\csromanspacer{} อิทมสฺสาหํ “อญฺญตรํ สามญฺญงฺคนฺ”ติ วทามีติ.\csromanspacer{} เอตมตฺถํ ภควา อโวจ, ตตฺเถตํ อิติ วุจฺจติ\csromandash{}}
\vspace{0.5\baselineskip}
\csromangathameasure{“อนวชฺเชน ตุฏฺฐสฺส, อปฺเปน สุลเภน จ. \\ น เสนาสนมารพฺภ, จีวรํ ปานโภชนํ. \\ วิฆาโต โหติ จิตฺตสฺส, ทิสา นปฺปฏิหญฺญติ. \\ เย จสฺส ธมฺมา อกฺขาตา, สามญฺญสฺสานุโลมิกา. \\ อธิคฺคหิตา ตุฏฺฐสฺส, อปฺปมตฺตสฺส ภิกฺขุโน”ติ.}
\csromangathagroup{\csromangathastanza{\csromanfolio{265}“อนวชฺเชน ตุฏฺฐสฺส, อปฺเปน สุลเภน จ. \\ น เสนาสนมารพฺภ, จีวรํ ปานโภชนํ.}\csromangathastanza{วิฆาโต โหติ จิตฺตสฺส, ทิสา นปฺปฏิหญฺญติ. \\ เย จสฺส\footnote{เยปสฺส (สฺยา)} ธมฺมา อกฺขาตา, สามญฺญสฺสานุโลมิกา. \\ อธิคฺคหิตา ตุฏฺฐสฺส, อปฺปมตฺตสฺส ภิกฺขุโน”ติ\footnote{เยปสฺส (สฺยา)}.}}

\prose{อยมฺปิ อตฺโถ วุตฺโต ภควตา อิติ เม สุตนฺติ.\csromanspacer{}\csromanspacer{}ทุติยํ.}

\csromanheader{2}{3. อาสวกฺขยสุตฺต}
\csromanheader{2}{102. วุตฺตํ เหตํ ภควตา วุตฺตมรหตาติ เม สุตํ\csromandash{}}
\prose{\csromansymbolfootnote{+}{ม 1. 9; สํ 1. 267; สํ 2. 124; สํ 3. 380 ปิฏฺเฐสุปิ. * อํ 1. 232 ปิฏฺเฐปิ}ชานโตหํ ภิกฺขเว ปสฺสโต อาสวานํ ขยํ วทามิ, โน อชานโต โน อปสฺสโต.\csromanspacer{} กิญฺจ ภิกฺขเว ชานโต กิํ ปสฺสโต อาสวานํ ขโย โหติ? “อิทํ ทุกฺขนฺ”ติ ภิกฺขเว ชานโต ปสฺสโต อาสวานํ ขโย โหติ, “อยํ ทุกฺขสมุทโย”ติ ภิกฺขเว ชานโต ปสฺสโต อาสวานํ ขโย โหติ, “อยํ ทุกฺขนิโรโธ”ติ ภิกฺขเว ชานโต ปสฺสโต อาสวานํ ขโย โหติ, “อยํ ทุกฺขนิโรธคามินี ปฏิปทา”ติ ภิกฺขเว ชานโต ปสฺสโต อาสวานํ ขโย โหติ.\csromanspacer{} เอวํ โข ภิกฺขเว ชานโต เอวํ ปสฺสโต อาสวานํ ขโย โหตีติ.\csromanspacer{} เอตมตฺถํ ภควา อโวจ, ตตฺเถตํ อิติ วุจฺจติ\csromandash{}}

\csromangathameasure{*“เสขสฺส สิกฺขมานสฺส, อุชุมคฺคานุสาริโน. \\ ขยสฺมิํ ปฐมํ ญาณํ, ตโต อญฺญา อนนฺตรา.}
\csromangathagroup{\csromangathastanza{\csromansymbolmark{*}“เสขสฺส สิกฺขมานสฺส, อุชุมคฺคานุสาริโน. \\ ขยสฺมิํ ปฐมํ ญาณํ, ตโต อญฺญา อนนฺตรา.}}

\csromanmatikaanchor{mtk.206}
\csromantocmarkref{subsubsection}{9. อนฺตรามลสุตฺต}{mtk.206}
\bookmark[dest={mtk.206},level=3]{9. อนฺตรามลสุตฺต}
\gambhira{อิติวุตฺตกปาฬิ}
\csromangathameasure{ตโต อญฺญา วิมุตฺตสฺส, วิมุตฺติญาณมุตฺตมํ. \\ อุปฺปชฺชติ ขเย ญาณํ, ขีณา สํโยชนา อิติ. \\ น เตฺววิทํ กุสีเตน, พาเลนมวิชานตา. \\ นิพฺพานํ อธิคนฺตพฺพํ, สพฺพคนฺถปฺปโมจนนฺ”ติ.}
\csromangathagroup{\csromangathastanza{\csromanfolio{266}ตโต อญฺญา วิมุตฺตสฺส, วิมุตฺติญาณมุตฺตมํ. \\ อุปฺปชฺชติ ขเย ญาณํ, ขีณา สํโยชนา อิติ.}\csromangathastanza{น เตฺววิทํ กุสีเตน, พาเลนมวิชานตา. \\ นิพฺพานํ อธิคนฺตพฺพํ, สพฺพคนฺถปฺปโมจนนฺ”ติ.}}

\prose{อยมฺปิ อตฺโถ วุตฺโต ภควตา อิติ เม สุตนฺติ.\csromanspacer{}\csromanspacer{}ตติยํ.}

\csromanheader{2}{4. สมณพฺราหฺมณสุตฺต}
\proseitem{103}{วุตฺตํ เหตํ ภควตา วุตฺตมรหตาติ เม สุตํ\csromandash{}}

\prose{เย หิ เกจิ ภิกฺขเว สมณา วา พฺราหฺมณา วา “อิทํ ทุกฺขนฺ”ติ ยถาภูตํ นปฺปชานนฺติ, “อยํ ทุกฺขสมุทโย”ติ ยถาภูตํ นปฺปชานนฺติ, “อยํ ทุกฺขนิโรโธ”ติ ยถาภูตํ นปฺปชานนฺติ, “อยํ ทุกฺขนิโรธคามินี ปฏิปทา”ติ ยถาภูตํ นปฺปชานนฺติ.\csromanspacer{} น เม เต ภิกฺขเว สมณา วา พฺราหฺมณา วา สมเณสุ วา สมณสมฺมตา, พฺราหฺมเณสุ วา พฺราหฺมณสมฺมตา, น จ ปเนเต อายสฺมนฺโต สามญฺญตฺถํ วา พฺรหฺมญฺญตฺถํ วา ทิฏฺเฐว ธมฺเม สยํ อภิญฺญา สจฺฉิกตฺวา อุปสมฺปชฺช วิหรนฺติ.}

\prose{เย จ โข เกจิ ภิกฺขเว สมณา วา พฺราหฺมณา วา “อิทํ ทุกฺขนฺ”ติ ยถาภูตํ ปชานนฺติ, “อยํ ทุกฺขสมุทโย”ติ ยถาภูตํ ปชานนฺติ, “อยํ ทุกฺขนิโรโธ”ติ ยถาภูตํ ปชานนฺติ, “อยํ ทุกฺขนิโรธคามินี ปฏิปทา”ติ ยถาภูตํ ปชานนฺติ.\csromanspacer{} เต โข เม ภิกฺขเว สมณา วา พฺราหฺมณา วา สมเณสุ เจว สมณสมฺมตา, พฺราหฺมเณสุ จ พฺราหฺมณสมฺมตา, เต จ ปนายสฺมนฺโต สามญฺญตฺถญฺจ พฺรหฺมญฺญตฺถญฺจ ทิฏฺเฐว ธมฺเม สยํ อภิญฺญา สจฺฉิกตฺวา อุปสมฺปชฺช วิหรนฺตีติ.\csromanspacer{} เอตมตฺถํ ภควา อโวจ, ตตฺเถตํ อิติ วุจฺจติ\csromandash{}}

\csromangathameasure{“เย ทุกฺขํ นปฺปชานนฺติ, อโถ ทุกฺขสฺส สมฺภวํ. \\ ยตฺถ จ สพฺพโส ทุกฺขํ, อเสสํ อุปรุชฺฌติ. \\ ตญฺจ มคฺคํ น ชานนฺติ, ทุกฺขูปสมคามินํ. \\ เจโตวิมุตฺติหีนา เต, อโถ ปญฺญาวิมุตฺติยา. \\ อภพฺพา เต อนฺตกิริยาย, เต เว ชาติชรูปคา.}
\csromangathagroup{\csromangathastanza{“เย ทุกฺขํ นปฺปชานนฺติ, อโถ ทุกฺขสฺส สมฺภวํ. \\ ยตฺถ จ สพฺพโส ทุกฺขํ, อเสสํ อุปรุชฺฌติ.}\csromangathastanza{ตญฺจ มคฺคํ น ชานนฺติ, ทุกฺขูปสมคามินํ. \\ เจโตวิมุตฺติหีนา เต, อโถ ปญฺญาวิมุตฺติยา. \\ อภพฺพา เต อนฺตกิริยาย, เต เว ชาติชรูปคา.}}

\csromanheader{2}{4. จตุกฺกนิปาต}
\csromangathameasure{เย จ ทุกฺขํ ปชานนฺติ, อโถ ทุกฺขสฺส สมฺภวํ. \\ ยตฺถ จ สพฺพโส ทุกฺขํ, อเสสํ อุปรุชฺฌติ. \\ ตญฺจ มคฺคํ ปชานนฺติ, ทุกฺขูปสมคามินํ. \\ เจโตวิมุตฺติสมฺปนฺนา, อโถ ปญฺญาวิมุตฺติยา. \\ ภพฺพา เต อนฺตกิริยาย, น เต ชาติชรูปคา”ติ.}
\csromangathagroup{\csromangathastanza{\csromanfolio{267}เย จ ทุกฺขํ ปชานนฺติ, อโถ ทุกฺขสฺส สมฺภวํ. \\ ยตฺถ จ สพฺพโส ทุกฺขํ, อเสสํ อุปรุชฺฌติ.}\csromangathastanza{ตญฺจ มคฺคํ ปชานนฺติ, ทุกฺขูปสมคามินํ. \\ เจโตวิมุตฺติสมฺปนฺนา, อโถ ปญฺญาวิมุตฺติยา. \\ ภพฺพา เต อนฺตกิริยาย, น เต ชาติชรูปคา”ติ.}}

\prose{อยมฺปิ อตฺโถ วุตฺโต ภควตา อิติ เม สุตนฺติ.\csromanspacer{}\csromanspacer{}จตุตฺถํ.}

\csromanmatikaanchor{mtk.207}
\csromantocmarkref{subsubsection}{10. เทวทตฺตสุตฺต}{mtk.207}
\bookmark[dest={mtk.207},level=3]{10. เทวทตฺตสุตฺต}
\csromanheader{2}{5. สีลสมฺปนฺนสุตฺต}
\csromanheader{2}{104. วุตฺตํ เหตํ ภควตา วุตฺตมรหตาติ เม สุตํ\csromandash{}}
\prose{เย เต ภิกฺขเว ภิกฺขู สีลสมฺปนฺนา สมาธิสมฺปนฺนา ปญฺญาสมฺปนฺนา วิมุตฺติสมฺปนฺนา วิมุตฺติญาณทสฺสนสมฺปนฺนา โอวาทกา วิญฺญาปกา สนฺทสฺสกา สมาทปกา สมุตฺเตชกา สมฺปหํสกา อลํสมกฺขาตาโร สทฺธมฺมสฺส.\csromanspacer{} ทสฺสนมฺปหํ ภิกฺขเว เตสํ ภิกฺขูนํ พหูปการํ วทามิ, สวนมฺปหํ ภิกฺขเว เตสํ ภิกฺขูนํ พหูปการํ วทามิ, อุปสงฺกมนมฺปหํ ภิกฺขเว เตสํ ภิกฺขูนํ พหูปการํ วทามิ, ปยิรุปาสนมฺปหํ ภิกฺขเว เตสํ ภิกฺขูนํ พหูปการํ วทามิ, อนุสฺสรณมฺปหํ\footnote{อนุสฺสติมฺปหํ (สฺยา)} ภิกฺขเว เตสํ ภิกฺขูนํ พหูปการํ วทามิ, อนุปพฺพชฺชมฺปหํ ภิกฺขเว เตสํ ภิกฺขูนํ พหูปการํ วทามิ.\csromanspacer{} ตํ กิสฺส เหตุ? ตถารูเป ภิกฺขเว ภิกฺขู เสวโต ภชโต ปยิรุปาสโต อปริปูโรปิ สีลกฺขนฺโธ ภาวนาปาริปูริํ คจฺฉติ, อปริปูโรปิ สมาธิกฺขนฺโธ ภาวนาปาริปูริํ คจฺฉติ, อปริปูโรปิ ปญฺญากฺขนฺโธ ภาวนาปาริปูริํ คจฺฉติ, อปริปูโรปิ วิมุตฺติกฺขนฺโธ ภาวานาปาริปูริํ คจฺฉติ, อปริปูโรปิ วิมุตฺติญาณทสฺสนกฺขนฺโธ ภาวนาปาริปูริํ คจฺฉติ.\csromanspacer{} เอวรูปา จ เต ภิกฺขเว ภิกฺขู “สตฺถาโร”ติปิ วุจฺจนฺติ, “สตฺถวาหา”ติปิ วุจฺจนฺติ, “รณญฺชหา”ติปิ วุจฺจนฺติ, “ตโมนุทา”ติปิ วุจฺจนฺติ, “อาโลกกรา”ติปิ วุจฺจนฺติ, ‘โอภาสกรา”ติปิ วุจฺจนฺติ, “ปชฺโชตกรา”ติปิ วุจฺจนฺติ, “อุกฺกาธารา”ติปิ วุจฺจนฺติ, “ปภงฺกรา”ติปิ วุจฺจนฺติ, “อริยา”ติปิ วุจฺจนฺติ, “จกฺขุมนฺโต”ติปิ วุจฺจนฺตีติ.\csromanspacer{} เอตมตฺถํ ภควา อโวจ, ตตฺเถตํ อิติ วุจฺจติ\csromandash{}}

\gambhira{อิติวุตฺตกปาฬิ}
\csromangathameasure{“ปาโมชฺชกรณํ ฐานํ, เอตํ โหติ วิชานตํ. \\ ยทิทํ ภาวิตตฺตานํ, อริยานํ ธมฺมชีวินํ. \\ เต โชตยนฺติ สทฺธมฺมํ, ภาสยนฺติ ปภงฺกรา. \\ อาโลกกรณา ธีรา, จกฺขุมนฺโต รณญฺชหา. \\ เยสํ เว สาสนํ สุตฺวา, สมฺมทญฺญาย ปณฺฑิตา. \\ ชาติกฺขยมภิญฺญาย, นาคจฺฉนฺติ ปุนพฺภวนฺ”ติ.}
\csromangathagroup{\csromangathastanza{\csromanfolio{268}“ปาโมชฺชกรณํ ฐานํ\footnote{...กรณฐานํ (สี, สฺยา) * เหฏฺฐา 201 ปิฏฺเฐปิ. + อํ 1. 130, 382 ปิฏฺเฐสุปิ.}, เอตํ โหติ วิชานตํ. \\ ยทิทํ ภาวิตตฺตานํ, อริยานํ ธมฺมชีวินํ.}\csromangathastanza{เต โชตยนฺติ สทฺธมฺมํ, ภาสยนฺติ ปภงฺกรา. \\ อาโลกกรณา ธีรา, จกฺขุมนฺโต รณญฺชหา.}\csromangathastanza{เยสํ เว สาสนํ สุตฺวา, สมฺมทญฺญาย ปณฺฑิตา. \\ ชาติกฺขยมภิญฺญาย, นาคจฺฉนฺติ ปุนพฺภวนฺ”ติ.}}

\prose{อยมฺปิ อตฺโถ วุตฺโต ภควตา อิติ เม สุตนฺติ.\csromanspacer{}\csromanspacer{}ปญฺจมํ.}

\csromanheader{2}{6. ตณฺหุปฺปาทสุตฺต}
\csromanheader{2}{105. วุตฺตํ เหตํ ภควตา วุตฺตมรหตาติ เม สุตํ\csromandash{}}
\prose{จตฺตาโรเม ภิกฺขเว ตณฺหุปฺปาทา, ยตฺถ ภิกฺขุโน ตณฺหา อุปฺปชฺชมานา อุปฺปชฺชติ.\csromanspacer{} กตเม จตฺตาโร? จีวรเหตุ วา ภิกฺขเว ภิกฺขุโน ตณฺหา อุปฺปชฺชมานา อุปฺปชฺชติ, ปิณฺฑปาตเหตุ วา ภิกฺขเว ภิกฺขุโน ตณฺหา อุปฺปชฺชมานา อุปฺปชฺชติ, เสนาสนเหตุ วา ภิกฺขเว ภิกฺขุโน ตณฺหา อุปฺปชฺชมานา อุปฺปชฺชติ, อิติภวาภวเหตุ วา ภิกฺขเว ภิกฺขุโน ตณฺหา อุปฺปชฺชมานา อุปฺปชฺชติ.\csromanspacer{} อิเม โข ภิกฺขเว จตฺตาโร ตณฺหุปฺปาทา, ยตฺถ ภิกฺขุโน ตณฺหา อุปฺปชฺชมานา อุปฺปชฺชตีติ.\csromanspacer{} เอตมตฺถํ ภควา อโวจ, ตตฺเถตํ อิติ วุจฺจติ\csromandash{}}

\csromangathameasure{*“ตณฺหาทุติโย ปุริโส, ทีฆมทฺธาน สํสรํ. \\ อิตฺถภาวญฺญถาภาวํ, สํสรํ นาติวตฺตติ. \\ เอตมาทีนวํ ญตฺวา, ตณฺหํ ทุกฺขสฺส สมฺภวํ. \\ วีตตโณฺห อนาทาโน, สโต ภิกฺขุ ปริพฺพเช”ติ.}
\csromangathagroup{\csromangathastanza{\csromansymbolmark{*}“ตณฺหาทุติโย ปุริโส, ทีฆมทฺธาน สํสรํ. \\ อิตฺถภาวญฺญถาภาวํ, สํสรํ นาติวตฺตติ.}\csromangathastanza{เอตมาทีนวํ ญตฺวา, ตณฺหํ ทุกฺขสฺส สมฺภวํ. \\ วีตตโณฺห อนาทาโน, สโต ภิกฺขุ ปริพฺพเช”ติ.}}

\prose{อยมฺปิ อตฺโถ วุตฺโต ภควตา อิติ เม สุตนฺติ.\csromanspacer{}\csromanspacer{}ฉฏฺฐํ.}

\csromanmatikaanchor{mtk.208}
\csromantocmarknopage{subsection}{5. ปญฺจมวคฺค}{mtk.208}
\bookmark[dest={mtk.208},level=2]{5. ปญฺจมวคฺค}
\csromanheader{2}{7. สพฺรหฺมกสุตฺต}
\csromanmatikaanchor{mtk.209}
\csromantocmarkref{subsubsection}{1. อคฺคปฺปสาทสุตฺต}{mtk.209}
\bookmark[dest={mtk.209},level=3]{1. อคฺคปฺปสาทสุตฺต}
\csromanheader{2}{106. วุตฺตํ เหตํ ภควตา วุตฺตมรหตาติ เม สุตํ\csromandash{}}
\prose{\csromansymbolmark{+}สพฺรหฺมกานิ ภิกฺขเว ตานิ กุลานิ, เยสํ ปุตฺตานํ มาตาปิตโร อชฺฌาคาเร ปูชิตา โหนฺติ.\csromanspacer{} สปุพฺพเทวตานิ ภิกฺขเว ตานิ กุลานิ,}

\vspace{\tipitakaparskip}
\csromancenter{จตุกฺกนิปาต เยสํ ปุตฺตานํ มาตาปิตโร อชฺฌาคาเร ปูชิตา โหนฺติ.\csromanspacer{} สปุพฺพาจริยกานิ ภิกฺขเว ตานิ กุลานิ, เยสํ ปุตฺตานํ มาตาปิตโร อชฺฌาคาเร ปูชิตา โหนฺติ.\csromanspacer{} สาหุเนยฺยกานิ ภิกฺขเว ตานิ กุลานิ, เยสํ ปุตฺตานํ มาตาปิตโร อชฺฌาคาเร ปูชิตา โหนฺติ.}
\prose{\csromanfolio{269}“พฺรหฺมา”ติ ภิกฺขเว มาตาปิตูนํ เอตํ อธิวจนํ, “ปุพฺพเทวตา”ติ ภิกฺขเว มาตาปิตูนํ เอตํ อธิวจนํ, “ปุพฺพาจริยา”ติ ภิกฺขเว มาตาปิตูนํ เอตํ อธิวจนํ,”อาหุเนยฺยา”ติ ภิกฺขเว มาตาปิตูนํ เอตํ อธิวจนํ.\csromanspacer{} ตํ กิสฺส เหตุ? พหุการา ภิกฺขเว มาตาปิตโร ปุตฺตานํ อาปาทกา โปสกา อิมสฺส โลกสฺส ทสฺเสตาโรติ.\csromanspacer{} เอตมตฺถํ ภควา อโวจ, ตตฺเถตํ อิติ วุจฺจติ\csromandash{}}

\csromangathameasure{“พฺรหฺมาติ มาตาปิตโร, ปุพฺพาจริยาติ วุจฺจเร. \\ อาหุเนยฺยา จ ปุตฺตานํ, ปชาย อนุกมฺปกา. \\ ตสฺมา หิ เน นมสฺเสยฺย, สกฺกเรยฺย จ ปณฺฑิโต. \\ อนฺเนน อถ ปาเนน, วตฺเถน สยเนน จ. \\ อุจฺฉาทเนน นฺหาปเนน, ปาทานํ โธวเนน จ. \\ ตาย นํ ปาริจริยาย, มาตาปิตูสุ ปณฺฑิตา. \\ อิเธว นํ ปสํสนฺติ, เปจฺจ สคฺเค ปโมทตี”ติ.}
\csromangathagroup{\csromangathastanza{“พฺรหฺมาติ มาตาปิตโร, ปุพฺพาจริยาติ วุจฺจเร. \\ อาหุเนยฺยา จ ปุตฺตานํ, ปชาย อนุกมฺปกา.}\csromangathastanza{ตสฺมา หิ เน นมสฺเสยฺย, สกฺกเรยฺย จ ปณฺฑิโต. \\ อนฺเนน อถ ปาเนน, วตฺเถน สยเนน จ.}\csromangathastanza{อุจฺฉาทเนน นฺหาปเนน\footnote{นหาปเนน (สี)}, ปาทานํ โธวเนน จ. \\ ตาย นํ ปาริจริยาย, มาตาปิตูสุ ปณฺฑิตา. \\ อิเธว นํ ปสํสนฺติ, เปจฺจ สคฺเค ปโมทตี”ติ.}}

\prose{อยมฺปิ อตฺโถ วุตฺโต ภควตา อิติ เม สุตนฺติ.\csromanspacer{}\csromanspacer{}สตฺตมํ.}

\csromanheader{2}{8. พหุการสุตฺต}
\csromanheader{2}{107. วุตฺตํ เหตํ ภควตา วุตฺตมรหตาติ เม สุตํ\csromandash{}}
\prose{พหุการา\footnote{พหูปการา (สี, อิ)} ภิกฺขเว พฺราหฺมณคหปติกา ตุมฺหากํ, เย โว\footnote{เย เต (สพฺพตฺถ)} ปจฺจุปฏฺฐิตา จีวรปิณฺฑปาตเสนาสนคิลานปจฺจยเภสชฺชปริกฺขาเรหิ.\csromanspacer{} ตุเมฺหปิ ภิกฺขเว พหุการา พฺราหฺมณคหปติกานํ, ยํ\footnote{เย (?)} เนสํ ธมฺมํ เทเสถ อาทิกลฺยาณํ มชฺเฌกลฺยาณํ ปริโยสานกลฺยาณํ สาตฺถํ สพฺยญฺชนํ เกวลปริปุณฺณํ ปริสุทฺธํ พฺรหฺมจริยํ ปกาเสถ.\csromanspacer{} เอวมิทํ ภิกฺขเว อญฺญมญฺญํ นิสฺสาย พฺรหฺมจริยํ วุสฺสติ โอฆสฺส นิตฺถรณตฺถาย สมฺมา ทุกฺขสฺส อนฺตกิริยายาติ.\csromanspacer{} เอตมตฺถํ ภควา อโวจ, ตตฺเถตํ อิติ วุจฺจติ\csromandash{}}

\csromanmatikaanchor{mtk.210}
\csromantocmarkref{subsubsection}{2. ชีวิกสุตฺต}{mtk.210}
\bookmark[dest={mtk.210},level=3]{2. ชีวิกสุตฺต}
\gambhira{อิติวุตฺตกปาฬิ}
\csromangathameasure{“สาคารา อนคารา จ, อุโภ อญฺโญญฺญนิสฺสิตา. \\ อาราธยนฺติ สทฺธมฺมํ, โยคกฺเขมํ อนุตฺตรํ. \\ สาคาเรสุ จ จีวรํ, ปจฺจยํ สยนาสนํ. \\ อนคารา ปฏิจฺฉนฺติ, ปริสฺสยวิโนทนํ. \\ สุคตํ ปน นิสฺสาย, คหฏฺฐา ฆรเมสิโน. \\ สทฺทหานา อรหตํ, อริยปญฺญาย ฌายิโน. \\ อิธ ธมฺมํ จริตฺวาน, มคฺคํ สุคติคามินํ. \\ นนฺทิโน เทวโลกสฺมิํ, โมทนฺติ กามกามิโน”ติ.}
\csromangathagroup{\csromangathastanza{\csromanfolio{270}“สาคารา อนคารา จ, อุโภ อญฺโญญฺญนิสฺสิตา. \\ อาราธยนฺติ สทฺธมฺมํ, โยคกฺเขมํ อนุตฺตรํ.}\csromangathastanza{สาคาเรสุ จ จีวรํ, ปจฺจยํ สยนาสนํ. \\ อนคารา ปฏิจฺฉนฺติ, ปริสฺสยวิโนทนํ.}\csromangathastanza{สุคตํ\footnote{ปุคฺคลํ (สี, ก)} ปน นิสฺสาย, คหฏฺฐา ฆรเมสิโน. \\ สทฺทหานา อรหตํ, อริยปญฺญาย ฌายิโน.}\csromangathastanza{อิธ ธมฺมํ จริตฺวาน, มคฺคํ สุคติคามินํ. \\ นนฺทิโน เทวโลกสฺมิํ, โมทนฺติ กามกามิโน”ติ.}}

\prose{อยมฺปิ อตฺโถ วุตฺโต ภควตา อิติ เม สุตนฺติ.\csromanspacer{}\csromanspacer{}อฏฺฐมํ.}

\csromanheader{2}{9. กุหสุตฺต}
\csromanheader{2}{108. วุตฺตํ เหตํ ภควตา วุตฺตมรหตาติ เม สุตํ\csromandash{}}
\prose{เย เกจิ ภิกฺขเว ภิกฺขู กุหา ถทฺธา ลปา สิงฺคี อุนฺนฬา อสมาหิตา, น เม เต ภิกฺขเว ภิกฺขู มามกา, อปคตา จ เต ภิกฺขเว ภิกฺขู อิมสฺมา ธมฺมวินยา, น จ เต\footnote{น จ เต ภิกฺขเว ภิกฺขู (สี, อิ, ก)} อิมสฺมิํ ธมฺมวินเย วุทฺธิํ วิรูฬฺหิํ เวปุลฺลํ อาปชฺชนฺติ.\csromanspacer{} เย จ โข ภิกฺขเว ภิกฺขู นิกฺกุหา นิลฺลปา ธีรา อตฺถทฺธา สุสมาหิตา, เต โข เม ภิกฺขเว ภิกฺขู มามกา, อนปคตา จ เต ภิกฺขเว ภิกฺขู อิมสฺมา ธมฺมวินยา, เต จ อิมสฺมิํ ธมฺมวินเย\footnote{อิมสฺมิํ จ เต ธมฺมวินเย (สฺยา), เต ภิกฺขเว ภิกฺขู อิมสฺมิํ ธมฺมวินเย (ก)} วุทฺธิํ วิรูฬฺหิํ เวปุลฺลํ อาปชฺชนฺตีติ.\csromanspacer{} เอตมตฺถํ ภควา อโวจ, ตตฺเถตํ อิติ วุจฺจติ\csromandash{}}

\csromanmatikaanchor{mtk.211}
\csromantocmarkref{subsubsection}{3. สงฺฆาฏิกณฺณสุตฺต}{mtk.211}
\bookmark[dest={mtk.211},level=3]{3. สงฺฆาฏิกณฺณสุตฺต}
\csromangathameasure{“กุหา ถทฺธา ลปา สิงฺคี, อุนฺนฬา อสมาหิตา. \\ น เต ธมฺเม วิรูหนฺติ, สมฺมาสมฺพุทฺธเทสิเต. \\ นิกฺกุหา นิลฺลปา ธีรา, อตฺถทฺธา สุสมาหิตา. \\ เต เว ธมฺเม วิรูหนฺติ, สมฺมาสมฺพุทฺธเทสิเต”ติ.}
\csromangathagroup{\csromangathastanza{“กุหา ถทฺธา ลปา สิงฺคี, อุนฺนฬา อสมาหิตา. \\ น เต ธมฺเม วิรูหนฺติ, สมฺมาสมฺพุทฺธเทสิเต.}\csromangathastanza{นิกฺกุหา นิลฺลปา ธีรา, อตฺถทฺธา สุสมาหิตา. \\ เต เว ธมฺเม วิรูหนฺติ, สมฺมาสมฺพุทฺธเทสิเต”ติ.}}

\prose{อยมฺปิ อตฺโถ วุตฺโต ภควตา อิติ เม สุตนฺติ.\csromanspacer{}\csromanspacer{}นวมํ.}

\csromanheader{2}{4. จตุกฺกนิปาต \textbf{10. นทีโสตสุตฺต}}
\csromanheader{2}{109. วุตฺตํ เหตํ ภควตา วุตฺตมรหตาติ เม สุตํ\csromandash{}}
\prose{\csromanfolio{271}เสยฺยถาปิ ภิกฺขเว ปุริโส นทิยา โสเตน โอวุเยฺหยฺย ปิยรูปสาตรูเปน, ตเมนํ จกฺขุมา ปุริโส ตีเร ฐิโต ทิสฺวา เอวํ วเทยฺย “กิญฺจาปิ โข ตฺวํ อมฺโภ ปุริส นทิยา โสเตน โอวุยฺหสิ ปิยรูปสาตรูเปน, อตฺถิ เจตฺถ เหฏฺฐา รหโท ส-อูมิ สาวฏฺโฏ สคโห สรกฺขโส, ยํ ตฺวํ อมฺโภ ปุริส รหทํ ปาปุณิตฺวา มรณํ วา นิคจฺฉสิ มรณมตฺตํ วา ทุกฺขนฺ”ติ.\csromanspacer{} อถ โข โส ภิกฺขเว ปุริโส ตสฺส ปุริสสฺส สทฺทํ สุตฺวา หตฺเถหิ จ ปาเทหิ จ ปฏิโสตํ วายเมยฺย.}

\prose{อุปมา โข เม อยํ ภิกฺขเว กตา อตฺถสฺส วิญฺญาปนาย.\csromanspacer{} อยํ เจตฺถ\footnote{อยํ เจเวตฺถ (สฺยา)} อตฺโถ\csromandash{}“นทิยา โสโต”ติ โข ภิกฺขเว ตณฺหาเยตํ อธิวจนํ.}

\prose{“ปิยรูปํ สาตรูปนฺ”ติ โข ภิกฺขเว ฉนฺเนตํ อชฺฌตฺติกานํ อายตนานํ อธิวจนํ.}

\prose{“เหฏฺฐา รหโท”ติ โข ภิกฺขเว ปญฺจนฺนํ โอรมฺภาคิยานํ สํโยชนานํ อธิวจนํ.}

\csromanmatikaanchor{mtk.212}
\csromantocmarkref{subsubsection}{4. อคฺคิสุตฺต}{mtk.212}
\bookmark[dest={mtk.212},level=3]{4. อคฺคิสุตฺต}
\prose{“อูมิภยนฺ”ติ โข\footnote{ส-อูมีติ โข (พหูสุ)} ภิกฺขเว โกธุปายาสสฺเสตํ อธิวจนํ.}

\prose{“อาวฏฺฏนฺ”ติ โข\footnote{สาวฏฺโฏติ โข (พหูสุ)} ภิกฺขเว ปญฺจนฺเนตํ กามคุณานํ}

\prose{อธิวจนํ.}

\prose{“คหรกฺขโส”ติ โข\footnote{สคโห สรกฺขโสติ โข (พหูสุ)} ภิกฺขเว มาตุคามสฺเสตํ อธิวจนํ.}

\prose{“ปฏิโสโต”ติ โข ภิกฺขเว เนกฺขมฺมสฺเสตํ อธิวจนํ.}

\prose{“หตฺเถหิ จ ปาเทหิ จ วายาโม”ติ โข ภิกฺขเว วีริยารมฺภสฺเสตํ อธิวจนํ.}

\prose{“จกฺขุมา ปุริโส ตีเร ฐิโต”ติ โข ภิกฺขเว ตถาคตสฺเสตํ อธิวจนํ อรหโต สมฺมาสมฺพุทฺธสฺสาติ.\csromanspacer{} เอตมตฺถํ ภควา อโวจ, ตตฺเถตํ อิติ วุจฺจติ\csromandash{}}

\gambhira{อิติวุตฺตกปาฬิ}
\csromangathameasure{“สหาปิ ทุกฺเขน ชเหยฺย กาเม, \\ โยคกฺเขมํ อายติํ ปตฺถยาโน. \\ สมฺมปฺปชาโน สุวิมุตฺตจิตฺโต, \\ วิมุตฺติยา ผสฺสเย ตตฺถ ตตฺถ. \\ ส เวทคู วูสิตพฺรหฺมจริโย, \\ โลกนฺตคู ปารคโตติ วุจฺจตี”ติ.}
\csromangathagroup{\csromangathastanza{\csromanfolio{272}“สหาปิ ทุกฺเขน ชเหยฺย กาเม, \\ โยคกฺเขมํ อายติํ ปตฺถยาโน. \\ สมฺมปฺปชาโน สุวิมุตฺตจิตฺโต, \\ วิมุตฺติยา ผสฺสเย ตตฺถ ตตฺถ. \\ ส เวทคู วูสิตพฺรหฺมจริโย, \\ โลกนฺตคู ปารคโตติ วุจฺจตี”ติ.}}

\prose{อยมฺปิ อตฺโถ วุตฺโต ภควตา อิติ เม สุตนฺติ.\csromanspacer{}\csromanspacer{}ทสมํ.}

\csromanheader{2}{11. จรสุตฺต}
\csromanheader{2}{110. วุตฺตํ เหตํ ภควตา วุตฺตมรหตาติ เม สุตํ\csromandash{}}
\csromanmatikaanchor{mtk.213}
\csromantocmarkref{subsubsection}{5. อุปปริกฺขสุตฺต}{mtk.213}
\bookmark[dest={mtk.213},level=3]{5. อุปปริกฺขสุตฺต}
\prose{จรโต เจปิ ภิกฺขเว ภิกฺขุโน อุปฺปชฺชติ กามวิตกฺโก วา พฺยาปาทวิตกฺโก วา วิหิํสาวิตกฺโก วา, ตญฺเจ ภิกฺขเว ภิกฺขุ อธิวาเสติ นปฺปชหติ น วิโนเทติ น พฺยนฺตีกโรติ\footnote{พฺยนฺติกโรติ (สี, อิ), พฺยนฺติํ กโรติ (ก)} น อนภาวํ คเมติ.\csromanspacer{} จรมฺปิ ภิกฺขเว ภิกฺขุ เอวํภูโต อนาตาปี อโนตฺตาปี\footnote{อโนตฺตปฺปี (สพฺพตฺถ) ทุกนิปาเต 214 ปิฏฺเฐ จ; องฺคุตฺตเร 1. 319 ปิฏฺเฐ จ ปสฺสิตพฺพํ.} สตตํ สมิตํ กุสีโต หีนวีริโยติ วุจฺจติ.}

\prose{ฐิตสฺส เจปิ ภิกฺขเว ภิกฺขุโน อุปฺปชฺชติ กามวิตกฺโก วา พฺยาปาทวิตกฺโก วา วิหิํสาวิตกฺโก วา, ตญฺเจ ภิกฺขเว ภิกฺขุ อธิวาเสติ นปฺปชหติ น วิโนเทติ น พฺยนฺตีกโรติ น อนภาวํ คเมติ.\csromanspacer{} ฐิโตปิ ภิกฺขเว ภิกฺขุ เอวํภูโต อนาตาปี อโนตฺตาปี สตตํ สมิตํ กุสีโต หีนวีริโยติ วุจฺจติ.}

\prose{นิสินฺนสฺส เจปิ ภิกฺขเว ภิกฺขุโน อุปฺปชฺชติ กามวิตกฺโก วา พฺยาปาทวิตกฺโก วา วิหิํสาวิตกฺโก วา, ตญฺเจ ภิกฺขเว ภิกฺขุ อธิวาเสติ นปฺปชหติ น วิโนเทติ น พฺยนฺตีกโรติ น อนภาวํ คเมติ.\csromanspacer{} นิสินฺโนปิ ภิกฺขเว ภิกฺขุ เอวํภูโต อนาตาปี อโนตฺตาปี สตตํ สมิตํ กุสีโต หีนวีริโยติ วุจฺจติ.}

\prose{สยานสฺส เจปิ ภิกฺขเว ภิกฺขุโน ชาครสฺส อุปฺปชฺชติ กามวิตกฺโก วา พฺยาปาทวิตกฺโก วา วิหิํสาวิตกฺโก วา, ตญฺเจ ภิกฺขเว ภิกฺขุ}

\vspace{\tipitakaparskip}
\csromancenter{จตุกฺกนิปาต อธิวาเสติ นปฺปชหติ น วิโนเทติ น พฺยนฺตีกโรติ น อนภาวํ คเมติ.\csromanspacer{} สยาโนปิ ภิกฺขเว ภิกฺขุ ชาคโร เอวํภูโต อนาตาปี อโนตฺตาปี สตตํ สมิตํ กุสีโต หีนวีริโยติ วุจฺจติ.}
\csromanmatikaanchor{mtk.214}
\csromantocmarkref{subsubsection}{6. กามูปปตฺติสุตฺต}{mtk.214}
\bookmark[dest={mtk.214},level=3]{6. กามูปปตฺติสุตฺต}
\prose{\csromanfolio{273}จรโต เจปิ ภิกฺขเว ภิกฺขุโน อุปฺปชฺชติ กามวิตกฺโก วา พฺยาปาทวิตกฺโก วา วิหิํสาวิตกฺโก วา, ตญฺเจ ภิกฺขเว ภิกฺขุ นาธิวาเสติ ปชหติ วิโนเทติ พฺยนฺตีกโรติ อนภาวํ คเมติ.\csromanspacer{} จรมฺปิ ภิกฺขเว ภิกฺขุ เอวํภูโต อาตาปี โอตฺตาปี\footnote{โอตฺตปฺปี (สพฺพตฺถ)} สตตํ สธิตํ อารทฺธวีริโย ปหิตตฺโตติ วุจฺจติ.}

\prose{ฐิตสฺส เจปิ ภิกฺขเว ภิกฺขุโน อุปฺปชฺชติ กามวิตกฺโก วา พฺยาปาทวิตกฺโก วา วิหิํสาวิตกฺโก วา, ตญฺเจ ภิกฺขเว ภิกฺขุ นาธิวาเสติ ปชหติ วิโนเทติ พฺยนฺตีกโรติ อนภาวํ คเมติ.\csromanspacer{} ฐิโตปิ ภิกฺขเว ภิกฺขุ เอวํภูโต อาตาปี โอตฺตาปี สตตํ สมิตํ อารทฺธวีริโย ปหิตตฺโตติ วุจฺจติ.}

\prose{นิสินฺนสฺส เจปิ ภิกฺขเว ภิกฺขุโน อุปฺปชฺชติ กามวิตกฺโก วา พฺยาปาทวิตกฺโก วา วิหิํสาวิตกฺโก วา, ตญฺเจ ภิกฺขเว ภิกฺขุ นาธิวาเสติ ปชหติ วิโนเทติ พฺยนฺตีกโรติ อนภาวํ คเมติ.\csromanspacer{} นิสินฺโนปิ ภิกฺขเว ภิกฺขุ เอวํภูโต อาตาปี โอตฺตาปี สตตํ สมิตํ อารทฺธวีริโย ปหิตตฺโตติ วุจฺจติ.}

\prose{สยานสฺส เจปิ ภิกฺขเว ภิกฺขุโน ชาครสฺส อุปฺปชฺชติ กามวิตกฺโก วา พฺยาปาทวิตกฺโก วา วิหิํสาวิตกฺโก วา, ตญฺเจ ภิกฺขเว ภิกฺขุ นาธิวาเสติ ปชหติ วิโนเทติ พฺยนฺตีกโรติ อนภาวํ คเมติ.\csromanspacer{} สยาโนปิ ภิกฺขเว ภิกฺขุ ชาคโร เอวํภูโต อาตาปี โอตฺตาปี สตตํ สมิตํ อารทฺธวีริโย ปหิตตฺโตติ วุจฺจตีติ.\csromanspacer{} เอตมตฺถํ ภควา อโวจ, ตตฺเถตํ อิติ วุจฺจติ\csromandash{}}

\csromangathameasure{“จรํ วา ยทิ วา ติฏฺฐํ, นิสินฺโน อุท วา สยํ. \\ โย วิตกฺกํ วิตกฺเกติ, ปาปกํ เคหนิสฺสิตํ. \\ กุมฺมคฺคํ ปฏิปนฺโน โส, โมหเนยฺเยสุ มุจฺฉิโต. \\ อภพฺโพ ตาทิโส ภิกฺขุ, ผุฏฺฐุํ สมฺโพธิมุตฺตมํ.}
\csromangathagroup{\csromangathastanza{“จรํ วา ยทิ วา ติฏฺฐํ, นิสินฺโน อุท วา สยํ. \\ โย วิตกฺกํ วิตกฺเกติ, ปาปกํ เคหนิสฺสิตํ.}\csromangathastanza{กุมฺมคฺคํ ปฏิปนฺโน\footnote{กุมฺมคฺคปฺปฏิปนฺโน (อํ 1. 321)} โส, โมหเนยฺเยสุ มุจฺฉิโต. \\ อภพฺโพ ตาทิโส ภิกฺขุ, ผุฏฺฐุํ สมฺโพธิมุตฺตมํ.}}

\gambhira{อิติวุตฺตกปาฬิ}
\csromangathameasure{โย จ จรํ วา ติฏฺฐํ วา, นิสินฺโน อุท วา สยํ. \\ วิตกฺกํ สมยิตฺวาน, วิตกฺกูปสเม รโต. \\ ภพฺโพ โส ตาทิโส ภิกฺขุ, ผุฏฺฐุํ สมฺโพธิมุตฺตมนฺ”ติ.}
\csromangathagroup{\csromangathastanza{\csromanfolio{274}โย จ จรํ วา ติฏฺฐํ วา\footnote{โย จรํ วา ยทิ วา ติฏฺฐํ (สฺยา), โย จรํ วาถ ติฏฺฐํ วา (สี, ก)}, นิสินฺโน อุท วา สยํ. \\ วิตกฺกํ สมยิตฺวาน, วิตกฺกูปสเม รโต. \\ ภพฺโพ โส ตาทิโส ภิกฺขุ, ผุฏฺฐุํ สมฺโพธิมุตฺตมนฺ”ติ.}}

\csromanmatikaanchor{mtk.215}
\csromantocmarkref{subsubsection}{7. กามโยคสุตฺต}{mtk.215}
\bookmark[dest={mtk.215},level=3]{7. กามโยคสุตฺต}
\prose{อยมฺปิ อตฺโถ วุตฺโต ภควตา อิติ เม สุตนฺติ.\csromanspacer{}\csromanspacer{}เอกาทสมํ.}

\csromanheader{2}{12. สมฺปนฺนสีลสุตฺต}
\csromanheader{2}{111. วุตฺตํ เหตํ ภควตา วุตฺตมรหตาติ เม สุตํ\csromandash{}}
\prose{สมฺปนฺนสีลา ภิกฺขเว วิหรถ\footnote{โหถ (สฺยา)} สมฺปนฺนปาติโมกฺขา, ปาติโมกฺขสํวรสํวุตา วิหรถ อาจารโคจรสมฺปนฺนา, อณุมตฺเตสุ วชฺเชสุ ภยทสฺสาวิโน สมาทาย สิกฺขถ สิกฺขาปเทสุ.}

\prose{สมฺปนฺนสีลานํ โว ภิกฺขเว วิหรตํ\footnote{ภวตํ (สฺยา)} สมฺปนฺนปาติโมกฺขานํ ปาติโมกฺขสํวรสํวุตานํ วิหรตํ อาจารโคจรสมฺปนฺนานํ อณุมตฺเตสุ วชฺเชสุ ภยทสฺสาวีนํ สมาทาย สิกฺขตํ สิกฺขาปเทสุ, กิมสฺส อุตฺตริ กรณียํ\footnote{กิมสฺส ภิกฺขเว อุตฺตริ กรณียํ (สพฺพตฺถ) 5-5. อภิชฺฌา วิคตา โหติ, พฺยาปาโท วิคโต (อํ 1. 321) ปิฏฺเฐ ปสฺสิตพฺพํ. 6-6. ถินมิทฺธํ อุทฺธจฺจกุกฺกุจฺจํ วิจิกิจฺฉา (อํ 1. 321)}.}

\prose{จรโต เจปิ ภิกฺขเว ภิกฺขุโน 5 อภิชฺฌา พฺยาปาโท วิคโต5 โหติ, 6 ถินมิทฺธํ วิคตํ โหติ, อุทฺธจฺจกุกฺกุจฺจํ วิคตํ โหติ, วิจิกิจฺฉา6 ปหีนา โหติ, อารทฺธํ โหติ วีริยํ อสลฺลีนํ, อุปฏฺฐิตา สติ อสมฺมุฏฺฐา\footnote{อปฺปมุฏฺฐา (สฺยา)}, ปสฺสทฺโธ กาโย อสารทฺโธ, สมาหิตํ จิตฺตํ เอกคฺคํ.\csromanspacer{} จรมฺปิ ภิกฺขเว ภิกฺขุ เอวํภูโต อาตาปี โอตฺตาปี สตตํ สมิตํ อารทฺธวีริโย ปหิตตฺโตติ วุจฺจติ.}

\prose{ฐิตสฺส เจปิ ภิกฺขเว ภิกฺขุโน อภิชฺฌา พฺยาปาโท วิคโต โหติ ฯเปฯ ถินมิทฺธํ.\csromanspacer{} อุทฺธจฺจกุกฺกุจฺจํ.\csromanspacer{} วิจิกิจฺฉา ปหีนา โหติ, อารทฺธํ}

\csromanmatikaanchor{mtk.216}
\csromantocmarkref{subsubsection}{8. กลฺยาณสีลสุตฺต}{mtk.216}
\bookmark[dest={mtk.216},level=3]{8. กลฺยาณสีลสุตฺต}
\vspace{\tipitakaparskip}
\csromancenter{จตุกฺกนิปาต โหติ วีริยํ อสลฺลีนํ, อุปฏฺฐิตา สติ อสมฺมุฏฺฐา, ปสฺสทฺโธ กาโย อสารทฺโธ, สมาหิตํ จิตฺตํ เอกคฺคํ.\csromanspacer{} ฐิโตปิ ภิกฺขเว ภิกฺขุ เอวํภูโต อาตาปี โอตฺตาปี สตตํ สมิตํ อารทฺธวีริโย ปหิตตฺโตติ วุจฺจติ.}
\prose{\csromanfolio{275}นิสินฺนสฺส เจปิ ภิกฺขเว ภิกฺขุโน อภิชฺฌา พฺยาปาโท วิคโต โหติ ฯเปฯ ถินมิทฺธํ.\csromanspacer{} อุทฺธจฺจกุกฺกุจฺจํ.\csromanspacer{} วิจิกิจฺฉา ปหีนา โหติ, อารทฺธํ โหติ วีริยํ อสลฺลีนํ, อุปฏฺฐิตา สติ อสมฺมุฏฺฐา, ปสฺสทฺโธ กาโย อสารทฺโธ, สมาหิตํ จิตฺตํ เอกคฺคํ.\csromanspacer{} นิสินฺโนปิ ภิกฺขเว ภิกฺขุ เอวํภูโต อาตาปี โอตฺตาปี สตตํ สมิตํ อารทฺธวีริโย ปหิตตฺโตติ วุจฺจติ.}

\prose{สยานสฺส เจปิ ภิกฺขเว ภิกฺขุโน ชาครสฺส อภิชฺฌา พฺยาปาโท วิคโต โหติ ฯเปฯ ถินมิทฺธํ.\csromanspacer{} อุทฺธจฺจกุกฺกุจฺจํ.\csromanspacer{} วิจิกิจฺฉา ปหีนา โหติ, อารทฺธํ โหติ วีริยํ อสลฺลีนํ, อุปฏฺฐิตา สติ อสมฺมุฏฺฐา, ปสฺสทฺโธ กาโย อสารทฺโธ, สมาหิตํ จิตฺตํ เอกคฺคํ.\csromanspacer{} สยาโนปิ ภิกฺขเว ภิกฺขุ ชาคโร เอวํภูโต อาตาปี โอตฺตาปี สตตํ สมิตํ อารทฺธวีริโย ปหิตตฺโตติ วุจฺจตีติ.\csromanspacer{} เอตมตฺถํ ภควา อโวจ, ตตฺเถตํ อิติ วุจฺจติ\csromandash{}}

\csromangathameasure{“ยตํ จเร ยตํ ติฏฺเฐ, ยตํ อจฺเฉ ยตํ สเย. \\ ยตํ สมิญฺชเย ภิกฺขุ, ยตเมนํ ปสารเย. \\ อุทฺธํ ติริยํ อปาจีนํ, ยาวตา ชคโต คติ. \\ สมเวกฺขิตา จ ธมฺมานํ, ขนฺธานํ อุทยพฺพยํ. \\ เอวํ วิหาริมาตาปิํ, สนฺตวุตฺติมนุทฺธตํ. \\ เจโตสมถสามีจิํ, สิกฺขมานํ สทา สตํ. \\ สตตํ ปหิตตฺโตติ, อาหุ ภิกฺขุํ ตถาวิธนฺ”ติ.}
\csromangathagroup{\csromangathastanza{“ยตํ จเร ยตํ ติฏฺเฐ, ยตํ อจฺเฉ ยตํ สเย. \\ ยตํ สมิญฺชเย\footnote{สมฺมิญฺชเย (สี, สฺยา)} ภิกฺขุ, ยตเมนํ ปสารเย.}\csromangathastanza{อุทฺธํ ติริยํ อปาจีนํ, ยาวตา ชคโต คติ. \\ สมเวกฺขิตา จ ธมฺมานํ, ขนฺธานํ อุทยพฺพยํ.}\csromangathastanza{เอวํ วิหาริมาตาปิํ, สนฺตวุตฺติมนุทฺธตํ. \\ เจโตสมถสามีจิํ, สิกฺขมานํ สทา สตํ. \\ สตตํ ปหิตตฺโตติ, อาหุ ภิกฺขุํ ตถาวิธนฺ”ติ.}}

\prose{อยมฺปิ อตฺโถ วุตฺโต ภควตา อิติ เม สุตนฺติ.\csromanspacer{}\csromanspacer{}ทฺวาทสมํ.}

\csromanheader{2}{13. โลกสุตฺต}
\csromanheader{2}{112. วุตฺตํ เหตํ ภควตา วุตฺตมรหตาติ เม สุตํ\csromandash{}}
\prose{โลโก ภิกฺขเว ตถาคเตน อภิสมฺพุทฺโธ, โลกสฺมา ตถาคโต วิสํยุตฺโต.\csromanspacer{} โลกสมุทโย ภิกฺขเว ตถาคเตน}

\gambhira{อิติวุตฺตกปาฬิ}
\prose{\csromanfolio{276}อภิสมฺพุทฺโธ, โลกสมุทโย ตถาคตสฺส ปหีโน.\csromanspacer{} โลกนิโรโธ ภิกฺขเว ตถาคเตน อภิสมฺพุทฺโธ, โลกนิโรโธ ตถาคตสฺส สจฺฉิกโต.\csromanspacer{} โลกนิโรธคามินี ปฏิปทา ภิกฺขเว ตถาคเตน อภิสมฺพุทฺธา, โลกนิโรธคามินี ปฏิปทา ตถาคตสฺส ภาวิตา.}

\prose{\csromansymbolfootnote{*}{ที 3. 111 ปิฏฺเฐปิ.}ยํ ภิกฺขเว สเทวกสฺส โลกสฺส สมารกสฺส สพฺรหฺมกสฺส สสฺสมณพฺราหฺมณิยา ปชาย สเทวมนุสฺสาย ทิฏฺฐํ สุตํ มุตํ วิญฺญาตํ ปตฺตํ ปริเยสิตํ อนุวิจริตํ มนสา, ยสฺมา ตํ ตถาคเตน อภิสมฺพุทฺธํ, ตสฺมา “ตถาคโต”ติ วุจฺจติ.}

\csromanmatikaanchor{mtk.217}
\csromantocmarkref{subsubsection}{9. ทานสุตฺต}{mtk.217}
\bookmark[dest={mtk.217},level=3]{9. ทานสุตฺต}
\prose{ยญฺจ ภิกฺขเว รตฺติํ ตถาคโต อนุตฺตรํ สมฺมาสมฺโพธิํ อภิสมฺพุชฺฌติ, ยญฺจ รตฺติํ อนุปาทิเสสาย นิพฺพานธาตุยา ปรินิพฺพายติ, ยํ เอตสฺมิํ อนฺตเร ภาสติ ลปติ นิทฺทิสติ, สพฺพํ ตํ ตเถว โหติ โน อญฺญถา, ตสฺมา “ตถาคโต”ติ วุจฺจติ.}

\prose{ยถาวาที ภิกฺขเว ตถาคโต ตถาการี, ยถาการี ตถาวาที.\csromanspacer{} อิติ ยถาวาที ตถาการี, ยถาการี ตถาวาที, ตสฺมา “ตถาคโต”ติ วุจฺจติ.}

\prose{สเทวเก ภิกฺขเว โลเก สมารเก สพฺรหฺมเก สสฺสมณพฺราหฺมณิยา ปชาย สเทวมนุสฺสาย ตถาคโต อภิภู อนภิภูโต อญฺญทตฺถุทโส วสวตฺตี, ตสฺมา “ตถาคโต”ติ ปุจฺจตีติ.\csromanspacer{} เอตมตฺถํ ภควา อโวจ, ตตฺเถตํ อิติ วุจฺจติ\csromandash{}}

\csromangathameasure{“สพฺพโลกํ อภิญฺญาย, สพฺพโลเก ยถาตถํ. \\ สพฺพโลกวิสํยุตฺโต, สพฺพโลเก อนูปโย. \\ ส เว สพฺพาภิภู ธีโร, สพฺพคนฺถปฺปโมจโน. \\ ผุฏฺฐา’สฺส ปรมา สนฺติ, นิพฺพานํ อกุโตภยํ. \\ เอส ขีณาสโว พุทฺโธ, อนีโฆ ฉินฺนสํสโย. \\ สพฺพกมฺมกฺขยํ ปตฺโต, วิมุตฺโต อุปธิสงฺขเย.}
\csromangathagroup{\csromangathastanza{“สพฺพโลกํ\footnote{สพฺพํ โลกํ (อํ 1. 332)} อภิญฺญาย, สพฺพโลเก ยถาตถํ. \\ สพฺพโลกวิสํยุตฺโต, สพฺพโลเก อนูปโย\footnote{สพฺพํ โลกํ (อํ 1. 332)}.}\csromangathastanza{ส เว\footnote{สพฺเพ (สพฺพตฺถ) อํ 1. 332 ปิฏฺเฐ ปสฺสิตพฺพํ.} สพฺพาภิภู ธีโร, สพฺพคนฺถปฺปโมจโน. \\ ผุฏฺฐา’สฺส ปรมา สนฺติ, นิพฺพานํ อกุโตภยํ.}\csromangathastanza{เอส ขีณาสโว พุทฺโธ, อนีโฆ ฉินฺนสํสโย. \\ สพฺพกมฺมกฺขยํ ปตฺโต, วิมุตฺโต อุปธิสงฺขเย.}}

\csromanheader{2}{4. จตุกฺกนิปาต}
\csromanmatikaanchor{mtk.218}
\csromantocmarkref{subsubsection}{10. เตวิชฺชสุตฺต}{mtk.218}
\bookmark[dest={mtk.218},level=3]{10. เตวิชฺชสุตฺต}
\csromangathameasure{เอส โส ภควา พุทฺโธ, เอส สีโห อนุตฺตโร. \\ สเทวกสฺส โลกสฺส, พฺรหฺมจกฺกํ ปวตฺตยิ. \\ อิติ เทวา มนุสฺสา จ, เย พุทฺธํ สรณํ คตา. \\ สงฺคมฺม ตํ นมสฺสนฺติ, มหนฺตํ วีตสารทํ. \\ ทนฺโต ทมยตํ เสฏฺโฐ, สนฺโต สมยตํ อิสิ. \\ มุตฺโต โมจยตํ อคฺโค, ติณฺโณ ตารยตํ วโร. \\ อิติ เหตํ นมสฺสนฺติ, มหนฺตํ วีตสารทํ. \\ สเทวกสฺมิํ โลกสฺมิํ, นตฺถิ เต ปฏิปุคฺคโล”ติ.}
\csromangathagroup{\csromangathastanza{\csromanfolio{277}เอส โส ภควา พุทฺโธ, เอส สีโห อนุตฺตโร. \\ สเทวกสฺส โลกสฺส, พฺรหฺมจกฺกํ ปวตฺตยิ.}\csromangathastanza{อิติ เทวา มนุสฺสา จ, เย พุทฺธํ สรณํ คตา. \\ สงฺคมฺม ตํ นมสฺสนฺติ, มหนฺตํ วีตสารทํ.}\csromangathastanza{ทนฺโต ทมยตํ เสฏฺโฐ, สนฺโต สมยตํ อิสิ. \\ มุตฺโต โมจยตํ อคฺโค, ติณฺโณ ตารยตํ วโร.}\csromangathastanza{อิติ เหตํ นมสฺสนฺติ, มหนฺตํ วีตสารทํ. \\ สเทวกสฺมิํ โลกสฺมิํ, นตฺถิ เต ปฏิปุคฺคโล”ติ.}}

\prose{อยมฺปิ อตฺโถ วุตฺโต ภควตา อิติ เม สุตนฺติ.\csromanspacer{}\csromanspacer{}เตรสมํ.}

\nitthitam{\textbf{จตุกฺกนิปาโต นิฏฺฐิโต.}}
\titlehead{\textbf{ตสฺสุทฺทานํ}}
\csromangathameasure{พฺราหฺมณสุลภา ชานํ, สมณสีลา ตณฺหา พฺรหฺมา. \\ พหุการา กุหปุริสา, จร สมฺปนฺนโลเกน เตรสาติ.}
\csromangathagroup{\csromangathastanza{พฺราหฺมณสุลภา\footnote{พฺราหฺมณจตฺตาริ (สพฺพตฺถ)} ชานํ, สมณสีลา ตณฺหา พฺรหฺมา. \\ พหุการา กุหปุริสา\footnote{พฺราหฺมณจตฺตาริ (สพฺพตฺถ)}, จร สมฺปนฺนโลเกน เตรสาติ.}}

\csromanheader{2}{\textbf{สุตฺตสงฺคโห}}
\csromangathameasure{สตฺตวิเสกนิปาตํ, ทุกฺกํ พาวีสสุตฺตสงฺคหิตํ. \\ สมปญฺญาสมถติกํ, เตรส จตุกฺกญฺจ อิติ ยมิทํ. \\ ทฺวิทสุตฺตรสุตฺตสเต, สงฺคายิตฺวา สมาทหิํสุ ปุรา. \\ อรหนฺโต จิรฏฺฐิติยา, ตมาหุ นาเมน อิติวุตฺตนฺติ.}
\csromangathagroup{\csromangathastanza{สตฺตวิเสกนิปาตํ, ทุกฺกํ พาวีสสุตฺตสงฺคหิตํ. \\ สมปญฺญาสมถติกํ, เตรส จตุกฺกญฺจ อิติ ยมิทํ.}\csromangathastanza{ทฺวิทสุตฺตรสุตฺตสเต, สงฺคายิตฺวา สมาทหิํสุ ปุรา. \\ อรหนฺโต จิรฏฺฐิติยา, ตมาหุ นาเมน อิติวุตฺตนฺติ.}}

\nitthitam{\textbf{อิติวุตฺตกปาฬิ นิฏฺฐิตา.}}
\csromanheader{2}{\textbf{ขุทฺทกนิกาย}}
\gambhira{\textbf{สุตฺตนิปาตปาฬิ}}
\namakkaram{นโม ตสฺส ภควตา อรหโต สมฺมาสมฺพุทฺธสฺส.}
\chapterhead{1. อุรควคฺค}
\csromanheader{2}{1. อุรคสุตฺต}
\csromanmatikaanchor{mtk.219}
\csromantocmarknopage{section}{4. จตุกฺกนิปาต}{mtk.219}
\bookmark[dest={mtk.219},level=1]{4. จตุกฺกนิปาต}
\proseitem{1}{\csromanfolio{279}โย\footnote{โย เว (สฺยา)} อุปฺปติตํ วิเนติ โกธํ,}

\csromanmatikaanchor{mtk.220}
\csromantocmarkref{subsection}{1. พฺราหฺมณธมฺมยาคสุตฺต}{mtk.220}
\bookmark[dest={mtk.220},level=2]{1. พฺราหฺมณธมฺมยาคสุตฺต}
\prose{วิสฏํ สปฺปวิสํว โอสเธหิ\footnote{โอสเธภิ (ก)}.}

\csromangathameasure{โส ภิกฺขุ ชหาติ โอรปารํ, \\ อุรโค ชิณฺณมิวตฺตจํ ปุราณํ.}
\csromangathagroup{\csromangathastanza{โส ภิกฺขุ ชหาติ โอรปารํ, \\ อุรโค ชิณฺณมิวตฺตจํ\footnote{ชิณฺณมิว ตจํ (สี, สฺยา, กํ, อิ), ชิณฺณมิวา ตจํ (?) 4. สเรรุหํ (ก)} ปุราณํ.}}

\proseitem{2}{โย ราคมุทจฺฉิทา อเสสํ,}

\prose{ภิสปุปฺผํว สโรรุหํ4 วิคยฺห.}

\csromangathameasure{โส ภิกฺขุ ชหาติ โอรปารํ, \\ อุรโค ชิณฺณมิวตฺตจํ ปุราณํ.}
\csromangathagroup{\csromangathastanza{โส ภิกฺขุ ชหาติ โอรปารํ, \\ อุรโค ชิณฺณมิวตฺตจํ ปุราณํ.}}

\proseitem{3}{โย ตณฺหมุทจฺฉิทา อเสสํ,}

\prose{สริตํ สีฆสรํ วิโสสยิตฺวา.}

\csromangathameasure{โส ภิกฺขุ ชหาติ โอรปารํ, \\ อุรโค ชิณฺณมิวตฺตจํ ปุราณํ.}
\csromangathagroup{\csromangathastanza{โส ภิกฺขุ ชหาติ โอรปารํ, \\ อุรโค ชิณฺณมิวตฺตจํ ปุราณํ.}}

\gambhira{สุตฺตนิปาตปาฬิ}
\csromanmatikaanchor{mtk.221}
\csromantocmarkref{subsection}{2. สุลภสุตฺต}{mtk.221}
\bookmark[dest={mtk.221},level=2]{2. สุลภสุตฺต}
\proseitem{4}{\csromanfolio{280}โย มานมุทพฺพธี อเสสํ,}

\prose{นฬเสตุํว สุทุพฺพลํ มโหโฆ.}

\csromangathameasure{โส ภิกฺขุ ชหาติ โอรปารํ, \\ อุรโค ชิณฺณมิวตฺตจํ ปุราณํ.}
\csromangathagroup{\csromangathastanza{โส ภิกฺขุ ชหาติ โอรปารํ, \\ อุรโค ชิณฺณมิวตฺตจํ ปุราณํ.}}

\proseitem{5}{โย นาชฺฌคมา ภเวสุ สารํ,}

\prose{วิจินํ ปุปฺผมิวา\footnote{ปุปฺผมิว (พหูสุ)} อุทุมฺพเรสุ.}

\csromangathameasure{โส ภิกฺขุ ชหาติ โอรปารํ, \\ อุรโค ชิณฺณมิวตฺตจํ ปุราณํ.}
\csromangathagroup{\csromangathastanza{โส ภิกฺขุ ชหาติ โอรปารํ, \\ อุรโค ชิณฺณมิวตฺตจํ ปุราณํ.}}

\proseitem{6}{ยสฺสนฺตรโต น สนฺติ โกปา,}

\csromanmatikaanchor{mtk.222}
\csromantocmarkref{subsection}{3. อาสวกฺขยสุตฺต}{mtk.222}
\bookmark[dest={mtk.222},level=2]{3. อาสวกฺขยสุตฺต}
\prose{อิติภวาภวตญฺจ\footnote{อิติพฺภวาภวตญฺจ (ก)} วีติวตฺโต.}

\csromangathameasure{โส ภิกฺขุ ชหาติ โอรปารํ, \\ อุรโค ชิณฺณมิวตฺตจํ ปุราณํ.}
\csromangathagroup{\csromangathastanza{โส ภิกฺขุ ชหาติ โอรปารํ, \\ อุรโค ชิณฺณมิวตฺตจํ ปุราณํ.}}

\proseitem{7}{ยสฺส วิตกฺกา วิธูปิตา,}

\prose{อชฺฌตฺตํ สุวิกปฺปิตา อเสสา.}

\csromangathameasure{โส ภิกฺขุ ชหาติ โอรปารํ, \\ อุรโค ชิณฺณมิวตฺตจํ ปุราณํ.}
\csromangathagroup{\csromangathastanza{โส ภิกฺขุ ชหาติ โอรปารํ, \\ อุรโค ชิณฺณมิวตฺตจํ ปุราณํ.}}

\proseitem{8}{โย นาจฺจสารี น ปจฺจสารี,}

\prose{สพฺพํ อจฺจคมา อิมํ ปปญฺจํ.}

\csromangathameasure{โส ภิกฺขุ ชหาติ โอรปารํ, \\ อุรโค ชิณฺณมิวตฺตจํ ปุราณํ.}
\csromangathagroup{\csromangathastanza{โส ภิกฺขุ ชหาติ โอรปารํ, \\ อุรโค ชิณฺณมิวตฺตจํ ปุราณํ.}}

\csromanmatikaanchor{mtk.223}
\csromantocmarkref{subsection}{4. สมณพฺราหฺมณสุตฺต}{mtk.223}
\bookmark[dest={mtk.223},level=2]{4. สมณพฺราหฺมณสุตฺต}
\proseitem{9}{โย นาจฺจสารี น ปจฺจสารี,}

\prose{สพฺพํ วิตถมิทนฺติ ญตฺว\footnote{อุตฺวา (สฺยา, อิ, ก)} โลเก.}

\csromangathameasure{โส ภิกฺขุ ชหาติ โอรปารํ, \\ อุรโค ชิณฺณมิวตฺตจํ ปุราณํ.}
\csromangathagroup{\csromangathastanza{โส ภิกฺขุ ชหาติ โอรปารํ, \\ อุรโค ชิณฺณมิวตฺตจํ ปุราณํ.}}

\proseitem{10}{โย นาจฺจสารี น ปจฺจสารี,}

\prose{สพฺพํ วิตถมิทนฺติ วีตโลโภ.}

\csromangathameasure{โส ภิกฺขุ ชหาติ โอรปารํ, \\ อุรโค ชิณฺณมิวตฺตจํ ปุราณํ.}
\csromangathagroup{\csromangathastanza{โส ภิกฺขุ ชหาติ โอรปารํ, \\ อุรโค ชิณฺณมิวตฺตจํ ปุราณํ.}}

\chapterheadpageread{1. อุรควคฺค}
\proseitem{11}{\csromanfolio{281}โย นาจฺจสารี น ปจฺจสารี,}

\prose{สพฺพํ วิตถมิทนฺติ วีตราโค.}

\csromangathameasure{โส ภิกฺขุ ชหาติ โอรปารํ, \\ อุรโค ชิณฺณมิวตฺตจํ ปุราณํ.}
\csromangathagroup{\csromangathastanza{โส ภิกฺขุ ชหาติ โอรปารํ, \\ อุรโค ชิณฺณมิวตฺตจํ ปุราณํ.}}

\csromanmatikaanchor{mtk.224}
\csromantocmarkref{subsection}{5. สีลสมฺปนฺนสุตฺต}{mtk.224}
\bookmark[dest={mtk.224},level=2]{5. สีลสมฺปนฺนสุตฺต}
\proseitem{12}{โย นาจฺจสารี น ปจฺจสารี,}

\prose{สพฺพํ วิตถมิทนฺติ วีตโทโส.}

\csromangathameasure{โส ภิกฺขุ ชหาติ โอรปารํ, \\ อุรโค ชิณฺณมิวตฺตจํ ปุราณํ.}
\csromangathagroup{\csromangathastanza{โส ภิกฺขุ ชหาติ โอรปารํ, \\ อุรโค ชิณฺณมิวตฺตจํ ปุราณํ.}}

\proseitem{13}{โย นาจฺจสารี น ปจฺจสารี,}

\prose{สพฺพํ วิตถมิทนฺติ วีตโมโห.}

\csromangathameasure{โส ภิกฺขุ ชหาติ โอรปารํ, \\ อุรโค ชิณฺณมิวตฺตจํ ปุราณํ.}
\csromangathagroup{\csromangathastanza{โส ภิกฺขุ ชหาติ โอรปารํ, \\ อุรโค ชิณฺณมิวตฺตจํ ปุราณํ.}}

\proseitem{14}{ยสฺสานุสยา น สนฺติ เกจิ,}

\prose{มูลา จ อกุสลา สมูหตาเส.}

\csromanmatikaanchor{mtk.225}
\csromantocmarkref{subsection}{6. ตณฺหุปฺปาทสุตฺต}{mtk.225}
\bookmark[dest={mtk.225},level=2]{6. ตณฺหุปฺปาทสุตฺต}
\csromangathameasure{โส ภิกฺขุ ชหาติ โอรปารํ, \\ อุรโค ชิณฺณมิวตฺตจํ ปุราณํ.}
\csromangathagroup{\csromangathastanza{โส ภิกฺขุ ชหาติ โอรปารํ, \\ อุรโค ชิณฺณมิวตฺตจํ ปุราณํ.}}

\proseitem{15}{ยสฺส ทรถชา น สนฺติ เกจิ,}

\prose{โอรํ อาคมนาย ปจฺจยาเส.}

\csromangathameasure{โส ภิกฺขุ ชหาติ โอรปารํ, \\ อุรโค ชิณฺณมิวตฺตจํ ปุราณํ.}
\csromangathagroup{\csromangathastanza{โส ภิกฺขุ ชหาติ โอรปารํ, \\ อุรโค ชิณฺณมิวตฺตจํ ปุราณํ.}}

\proseitem{16}{ยสฺส วนถชา น สนฺติ เกจิ,}

\prose{วินิพนฺธาย ภวาย เหตุกปฺปา.}

\csromanmatikaanchor{mtk.226}
\csromantocmarkref{subsection}{7. สพฺรหฺมกสุตฺต}{mtk.226}
\bookmark[dest={mtk.226},level=2]{7. สพฺรหฺมกสุตฺต}
\csromangathameasure{โส ภิกฺขุ ชหาติ โอรปารํ, \\ อุรโค ชิณฺณมิวตฺตจํ ปุราณํ.}
\csromangathagroup{\csromangathastanza{โส ภิกฺขุ ชหาติ โอรปารํ, \\ อุรโค ชิณฺณมิวตฺตจํ ปุราณํ.}}

\proseitem{17}{โย นีวรเณ ปหาย ปญฺจ,}

\prose{อนิโฆ ติณฺณกถํกโถ วิสลฺโล.}

\csromangathameasure{โส ภิกฺขุ ชหาติ โอรปารํ, \\ อุรโค ชิณฺณมิวตฺตจํ ปุราณนฺติ. อุรคสุตฺตํ ปฐมํ.}
\csromangathagroup{\csromangathastanza{โส ภิกฺขุ ชหาติ โอรปารํ, \\ อุรโค ชิณฺณมิวตฺตจํ ปุราณนฺติ.\csromanspacer{} อุรคสุตฺตํ ปฐมํ.}}

\gambhira{สุตฺตนิปาตปาฬิ}
\csromanheader{2}{2. ธนิยสุตฺต}
\proseitem{18}{\csromanfolio{282}ปกฺโกทโน ทุทฺธขีโรหมสฺมิ, (อิติ ธนิโย โคโป,)}

\prose{อนุตีเร มหิยา สมานวาโส.}

\prose{ฉนฺนา กุฏิ อาหิโต คินิ,}

\csromanmatikaanchor{mtk.227}
\csromantocmarkref{subsection}{8. พหุการสุตฺต}{mtk.227}
\bookmark[dest={mtk.227},level=2]{8. พหุการสุตฺต}
\csromanheader{2}{อถ เจ ปตฺถยสี ปวสฺส เทว. (1)}
\proseitem{19}{อกฺโกธโน วิคตขิโลหมสฺมิ\footnote{วิคตขีโลหมสฺมิ (สี, อิ)}, (อิติ ภควา,)}

\prose{อนุตีเร มหิเยกรตฺติวาโส.}

\prose{วิวฏา กุฏิ นิพฺพุโต คินิ,}

\csromanheader{2}{อถ เจ ปตฺถยสี ปวสฺส เทว. (2)}
\proseitem{20}{อนฺธกมกสา น วิชฺชเร, (อิติ ธนิโย โคโป,)}

\prose{กจฺเฉ รูฬฺหติเณ จรนฺติ คาโว.}

\prose{วุฏฺฐิมฺปิ สเหยฺยุมาคตํ,}

\csromanheader{2}{อถ เจ ปตฺถยสี ปวสฺส เทว. (3)}
\csromanmatikaanchor{mtk.228}
\csromantocmarkref{subsection}{9. กุหสุตฺต}{mtk.228}
\bookmark[dest={mtk.228},level=2]{9. กุหสุตฺต}
\proseitem{21}{พทฺธาสิ ภิสี สุสงฺขตา, (อิติ ภควา,)}

\prose{ติณฺโณ ปารคโต วิเนยฺย โอฆํ.}

\prose{อตฺโถ ภิสิยา น วิชฺชติ,}

\csromanheader{2}{อถ เจ ปตฺถยสี ปวสฺส เทว. (4)}
\proseitem{22}{โคปี มม อสฺสวา อโลลา, (อิติ ธนิโย โคโป,)}

\prose{ทีฆรตฺตํ\footnote{ทีฆรตฺต (ก)} สํวาสิยา มนาปา.}

\csromanmatikaanchor{mtk.229}
\csromantocmarkref{subsection}{10. นทีโสตสุตฺต}{mtk.229}
\bookmark[dest={mtk.229},level=2]{10. นทีโสตสุตฺต}
\prose{ตสฺสา น สุณามิ กิญฺจิ ปาปํ,}

\csromanheader{2}{อถ เจ ปตฺถยสี ปวสฺส เทว. (5)}
\proseitem{23}{จิตฺตํ มม อสฺสวํ วิมุตฺตํ, (อิติ ภควา,)}

\prose{ทีฆรตฺตํ ปริภาวิตํ สุทนฺตํ.}

\prose{ปาปํ ปน เม น วิชฺชติ,}

\csromanheader{2}{อถ เจ ปตฺถยสี ปวสฺส เทว. (6)}
\chapterheadpageread{1. อุรควคฺค}
\proseitem{24}{\csromanfolio{283}อตฺตเวตนภโตหมสฺมิ, (อิติ ธนิโย โคโป,)}

\prose{ปุตฺตา จ เม สมานิยา อโรคา.}

\prose{เตสํ น สุณามิ กิญฺจิ ปาปํ,}

\csromanheader{2}{อถ เจ ปตฺถยสี ปวสฺส เทว. (7)}
\proseitem{25}{นาหํ ภตโกสฺมิ กสฺสจิ, (อิติ ภควา,)}

\prose{นิพฺพิฏฺเฐน จรามิ สพฺพโลเก.}

\prose{อตฺโถ ภติยา น วิชฺชติ,}

\csromanheader{2}{อถ เจ ปตฺถยสี ปวสฺส เทว. (8)}
\proseitem{26}{อตฺถิ วสา อตฺถิ เธนุปา, (อิติ ธนิโย โคโป,)}

\csromanmatikaanchor{mtk.230}
\csromantocmarkref{subsection}{11. จรสุตฺต}{mtk.230}
\bookmark[dest={mtk.230},level=2]{11. จรสุตฺต}
\prose{โคธรณิโย ปเวณิโยปิ อตฺถิ.}

\prose{อุสโภปิ ควมฺปตีธ อตฺถิ,}

\csromanheader{2}{อถ เจ ปตฺถยสี ปวสฺส เทว. (9)}
\proseitem{27}{นตฺถิ วสา นตฺถิ เธนุปา, (อิติ ภควา,)}

\prose{โคธรณิโย ปเวณิโยปิ นตฺถิ.}

\prose{อุสโภปิ ควมฺปตีธ นตฺถิ,}

\csromanheader{2}{อถ เจ ปตฺถยสี ปวสฺส เทว. (10)}
\proseitem{28}{ขิลา นิขาตา อสมฺปเวธี, (อิติ ธนิโย โคโป,)}

\prose{ทามา มุญฺชมยา นวา สุสณฺฐานา.}

\prose{น หิ สกฺขินฺติ เธนุปาปิ เฉตฺตุํ\footnote{เฉตุํ (ก)}.}

\csromanheader{2}{อถ เจ ปตฺถยสี ปวสฺส เทว. (11)}
\proseitem{29}{อุสโภริว เฉตฺว\footnote{เฉตฺวา (สฺยา, ก)} พนฺธนานิ, (อิติ ภควา,)}

\prose{นาโค ปูติลตํว ทาลยิตฺวา\footnote{ปูติลตํ ปทาลยิตฺวา (สฺยา, ก)}.}

\prose{นาหํ ปุนุเปสฺสํ\footnote{ปุน อุเปสฺสํ (สี, สฺยา, กํ, อิ), ปุนุเปยฺย (ก)} คพฺภเสยฺยํ,}

\csromanheader{2}{อถ เจ ปตฺถยสี ปวสฺส เทว. (12)}
\proseitem{30}{นินฺนญฺจ ถลญฺจ ปูรยนฺโต,}

\csromanmatikaanchor{mtk.231}
\csromantocmarkref{subsection}{12. สมฺปนฺนสีลสุตฺต}{mtk.231}
\bookmark[dest={mtk.231},level=2]{12. สมฺปนฺนสีลสุตฺต}
\prose{มหาเมโฆ ปวสฺสิ ตาวเทว.}

\prose{สุตฺวา เทวสฺส วสฺสโต,}

\csromanheader{2}{อิมมตฺถํ ธนิโย อภาสถ. (13)}
\gambhira{สุตฺตนิปาตปาฬิ}
\csromangathameasure{ลาภา วต โน อนปฺปกา, \\ เย มยํ ภควนฺตํ อทฺทสาม. \\ สรณํ ตํ อุเปม จกฺขุม,}
\csromangathagroup{\csromangathastanza{\csromanfolio{284}ลาภา วต โน อนปฺปกา, \\ เย มยํ ภควนฺตํ อทฺทสาม. \\ สรณํ ตํ อุเปม จกฺขุม,}}

\csromanheader{2}{สตฺถา โน โหติ ตุวํ มหามุนิ. (14)}
\proseitem{32}{โคปี จ อหญฺจ อสฺสวา,}

\prose{พฺรหฺมจริยํ\footnote{พฺรหฺมจริย (ก)} สุคเต จรามเส.}

\prose{ชาติมรณสฺส ปารคู\footnote{ปารคา (สี, สฺยา, กํ, อิ)},}

\csromanheader{2}{ทุกฺขสฺสนฺตกรา ภวามเส. (15)}
\proseitem{33}{\csromansymbolfootnote{*}{ขุ 3. 9; ขุ 8. 229 ปิฏฺฐาทีสุปิ.}นนฺทติ ปุตฺเตหิ ปุตฺติมา, (อิติ มาโร ปาปิมา,)}

\prose{โคมา\footnote{โคมิโก (สี, อิ), โคปิโก (สฺยา, กํ), โคปิโย (ก)} โคหิ ตเถว นนฺทติ.}

\prose{อุปธี หิ นรสฺส นนฺทนา,}

\csromanmatikaanchor{mtk.232}
\csromantocmarkref{subsection}{13. โลกสุตฺต}{mtk.232}
\bookmark[dest={mtk.232},level=2]{13. โลกสุตฺต}
\csromanheader{2}{น หิ โส นนฺทติ โย นิรูปธิ. (16)}
\proseitem{34}{โสจติ ปุตฺเตหิ ปุตฺติมา, (อิติ ภควา)}

\prose{โคมา โคหิ ตเถว โสจติ.}

\csromangathameasure{อุปธี หิ นรสฺส โสจนา, \\ น หิ โส โสจติ โย นิรูปธีติ. (17) ธนิยสุตฺตํ ทุติยํ.}
\csromangathagroup{\csromangathastanza{อุปธี หิ นรสฺส โสจนา, \\ น หิ โส โสจติ โย นิรูปธีติ. (17) ธนิยสุตฺตํ ทุติยํ.}}

\csromanheader{2}{3. ขคฺควิสาณสุตฺต}
\proseitem{35}{\csromansymbolmark{*}สพฺเพสุ ภูเตสุ นิธาย ทณฺฑํ,}

\prose{อวิเหฐยํ อญฺญตรมฺปิ เตสํ.}

\prose{น ปุตฺตมิจฺเฉยฺย กุโต สหายํ,}

\csromanheader{2}{เอโก จเร ขคฺควิสาณกปฺโป. (1)}
\chapterheadpageread{1. อุรควคฺค}
\proseitem{36}{\csromanfolio{285}สํสคฺคชาตสฺส ภวนฺติ สฺเนหา,}

\prose{สฺเนหนฺวยํ ทุกฺขมิทํ ปโหติ.}

\prose{อาทีนวํ สฺเนหชํ เปกฺขมาโน,}

\csromanheader{2}{เอโก จเร ขคฺควิสาณกปฺโป. (2)}
\proseitem{37}{มิตฺเต สุหชฺเช อนุกมฺปมาโน,}

\prose{หาเปติ อตฺถํ ปฏิพทฺธจิตฺโต.}

\prose{เอตํ ภยํ สนฺถเว\footnote{สนฺธเว (ก)} เปกฺขมาโน,}

\csromanheader{2}{เอโก จเร ขคฺควิสาณกปฺโป. (3)}
\proseitem{38}{วํโส วิสาโลว ยถา วิสตฺโต,}

\prose{ปุตฺเตสุ ทาเรสุ จ ยา อเปกฺขา.}

\prose{วํสกฺกฬีโรว\footnote{วํสกฬีโรว (สี), วํสากฬีโรว (สฺยา, กํ, อิ), วํเสกฬีโรว (นิทฺเทส)} อสชฺชมาโน,}

\csromanheader{2}{เอโก จเร ขคฺควิสาณกปฺโป. (4)}
\proseitem{39}{มิโค อรญฺญมฺหิ ยถา อพทฺโธ\footnote{อพนฺโธ (สฺยา, กํ)},}

\prose{เยนิจฺฉกํ คจฺฉติ โคจราย.}

\prose{วิญฺญู นโร เสริตํ เปกฺขมาโน,}

\csromanheader{2}{เอโก จเร ขคฺควิสาณกปฺโป. (5)}
\csromanmatikaanchor{mtk.233}
\csromantocmarknopage{chapter}{สุตฺตนิปาตปาฬิ}{mtk.233}
\bookmark[dest={mtk.233},level=0]{สุตฺตนิปาตปาฬิ}
\proseitem{40}{อามนฺตนา โหติ สหายมชฺเฌ,}

\prose{วาเส ฐาเน คมเน จาริกาย.}

\prose{อนภิชฺฌิตํ เสริตํ เปกฺขมาโน,}

\csromanmatikaanchor{mtk.234}
\csromanmatikaanchor{mtk.235}
\csromantocmarknopage{chapter}{1. อุรควคฺค}{mtk.234}
\bookmark[dest={mtk.234},level=0]{1. อุรควคฺค}
\csromantocmarkref{subsection}{1. อุรคสุตฺต}{mtk.235}
\bookmark[dest={mtk.235},level=2]{1. อุรคสุตฺต}
\csromanheader{2}{เอโก จเร ขคฺควิสาณกปฺโป. (6)}
\proseitem{41}{ขิฑฺฑา รตี โหติ สหายมชฺเฌ,}

\prose{ปุตฺเตสุ จ วิปุลํ โหติ เปมํ.}

\prose{ปิยวิปฺปโยคํ วิชิคุจฺฉมาโน,}

\csromanheader{2}{เอโก จเร ขคฺควิสาณกปฺโป. (7)}
\gambhira{สุตฺตนิปาตปาฬิ}
\proseitem{42}{\csromanfolio{286}จาตุทฺทิโส อปฺปฏิโฆ จ โหติ,}

\prose{สนฺตุสฺสมาโน อิตรีตเรน.}

\prose{ปริสฺสยานํ สหิตา อฉมฺภี,}

\csromanheader{2}{เอโก จเร ขคฺควิสาณกปฺโป. (8)}
\proseitem{43}{ทุสฺสงฺคหา ปพฺพชิตาปิ เอเก,}

\prose{อโถ คหฏฺฐา ฆรมาวสนฺตา.}

\prose{อปฺโปสฺสุกฺโก ปรปุตฺเตสุ หุตฺวา,}

\csromanheader{2}{เอโก จเร ขคฺควิสาณกปฺโป. (9)}
\proseitem{44}{โอโรปยิตฺวา คิหิพฺยญฺชนานิ\footnote{คิหิคฺยญฺชนานิ (สฺยา, กํ, อิ)},}

\prose{สญฺฉินฺนปตฺโต\footnote{สํสีนปตฺโต (สี, อิ)} ยถา โกวิฬาโร.}

\prose{เฉตฺวาน วีโร คิหิพนฺธนานิ,}

\csromanheader{2}{เอโก จเร ขคฺควิสาณกปฺโป. (10)}
\proseitem{45}{\csromansymbolfootnote{*}{เหฏฺฐา 60; วิ 3. 496; ม 3. 192 ปิฏฺเฐสุปิ.}สเจ ลเภถ นิปกํ สหายํ,}

\prose{สทฺธิํจรํ สาธุวิหาริ’ธีรํ.}

\prose{อภิภุยฺย สพฺพานิ ปริสฺสยานิ,}

\csromanheader{2}{จเรยฺย เตน’ตฺตมโน สตีมา. (11)}
\proseitem{46}{\csromansymbolmark{*}โน เจ ลเภถ นิปกํ สหายํ,}

\prose{สทฺธิํจรํ สาธุวิหาริ’ธีรํ.}

\prose{ราชาว รฏฺฐํ วิชิตํ ปหาย,}

\csromanheader{2}{เอโก จเร มาตงฺค’รญฺเญว นาโค. (12)}
\proseitem{47}{อทฺธา ปสํสาม สหายสมฺปทํ,}

\prose{เสฏฺฐา สมา เสวิตพฺพา สหายา.}

\prose{เอเต อลทฺธา อนวชฺชโภชี,}

\csromanheader{2}{เอโก จเร ขคฺควิสาณกปฺโป. (13)}
\proseitem{48}{\csromansymbolfootnote{+}{ขุ 8. 259 ปิฏฺเฐปิ.}ทิสฺวา สุวณฺณสฺส ปภสฺสรานิ,}

\nitthitam{กมฺมารปุตฺเตน สุนิฏฺฐิตานิ.}
\prose{สํฆฏฺฏมานานิ ทุเว ภุชสฺมิํ,}

\csromanheader{2}{เอโก จเร ขคฺควิสาณกปฺโป. (14)}
\chapterheadpageread{1. อุรควคฺค}
\proseitem{49}{\csromanfolio{287}เอวํ ทุตีเยน\footnote{ทุติเยน (สพฺพตฺถ)} สหา มมสฺส,}

\prose{วาจาภิลาโป อภิสชฺชนา วา.}

\prose{เอตํ ภยํ อายติํ เปกฺขมาโน,}

\csromanheader{2}{เอโก จเร ขคฺควิสาณกปฺโป. (15)}
\proseitem{50}{กามา หิ จิตฺรา มธุรา มโนรมา,}

\prose{วิรูปรูเปน มเถนฺติ จิตฺตํ.}

\prose{อาทีนวํ กามคุเณสุ ทิสฺวา,}

\csromanheader{2}{เอโก จเร ขคฺควิสาณกปฺโป. (16)}
\proseitem{51}{¢ตี จ คณฺโฑ จ อุปทฺทโว จ,}

\prose{โรโค จ สลฺลญฺจ ภยญฺจ เมตํ.}

\prose{เอตํ ภยํ กามคุเณสุ ทิสฺวา,}

\csromanheader{2}{เอโก จเร ขคฺควิสาณกปฺโป. (17)}
\proseitem{52}{สีตญฺจ อุณฺหญฺจ ขุทํ ปิปาสํ,}

\prose{วาตาตเป ฑํสสรีสเป\footnote{qอํสสิริํสเป (สี, สฺยา, กํ, อิ)} จ.}

\prose{สพฺพานิเป’ตานิ อภิสมฺภวิตฺวา,}

\csromanheader{2}{เอโก จเร ขคฺควิสาณกปฺโป. (18)}
\proseitem{53}{นาโคว ยูถานิ วิวชฺชยิตฺวา,}

\prose{สญฺชาตขนฺโธ ปทุมี อุฬาโร.}

\prose{ยถาภิรนฺตํ วิหรํ\footnote{วิหเร (สี, อิ, นิทฺเทส)} อรญฺเญ,}

\csromanheader{2}{เอโก จเร ขคฺควิสาณกปฺโป. (19)}
\proseitem{54}{อฏฺฐาน’ตํ สงฺคณิการตสฺส,}

\prose{ยํ ผสฺสเย\footnote{ผุสฺสเย (สฺยา)} สามยิกํ วิมุตฺติํ.}

\prose{อาทิจฺจพนฺธุสฺส วโจ นิสมฺม,}

\csromanheader{2}{เอโก จเร ขคฺควิสาณกปฺโป. (20)}
\proseitem{55}{ทิฏฺฐีวิสูกานิ อุปาติวตฺโต,}

\prose{ปตฺโต นิยามํ ปฏิลทฺธมคฺโค.}

\prose{อุปฺปนฺนญาโณมฺหิ อนญฺญเนยฺโย,}

\csromanheader{2}{เอโก จเร ขคฺควิสาณกปฺโป. (21)}
\gambhira{สุตฺตนิปาตปาฬิ}
\proseitem{56}{\csromanfolio{288}นิลฺโลลุโป นิกฺกุโห นิปฺปิปาโส,}

\prose{นิมฺมกฺโข นิทฺธนฺตกสาวโมโห.}

\prose{นิราสโย\footnote{นิราสาโส (ก)} สพฺพโลเก ภวิตฺวา,}

\csromanheader{2}{เอโก จเร ขคฺควิสาณกปฺโป. (22)}
\proseitem{57}{ปาปํ สหายํ ปริวชฺชเยถ,}

\prose{อนตฺถทสฺสิํ วิสเม นิวิฏฺฐํ.}

\prose{สยํ น เสเว ปสุตํ ปมตฺตํ,}

\csromanheader{2}{เอโก จเร ขคฺควิสาณกปฺโป. (23)}
\csromanmatikaanchor{mtk.236}
\csromantocmarkref{subsection}{2. ธนิยสุตฺต}{mtk.236}
\bookmark[dest={mtk.236},level=2]{2. ธนิยสุตฺต}
\proseitem{58}{พหุสฺสุตํ ธมฺมธรํ ภเชถ,}

\prose{มิตฺตํ อุฬารํ ปฏิภานวนฺตํ.}

\prose{อญฺญาย อตฺถานิ วิเนยฺย กงฺขํ,}

\csromanheader{2}{เอโก จเร ขคฺควิสาณกปฺโป. (24)}
\proseitem{59}{ขิฑฺฑํ รติํ กามสุขญฺจ โลเก,}

\prose{อนลงฺกริตฺวา อนเปกฺขมาโน.}

\prose{วิภูสนฏฺฐานา วิรโต สจฺจวาที,}

\csromanheader{2}{เอโก จเร ขคฺควิสาณกปฺโป. (25)}
\proseitem{60}{ปุตฺตญฺจ ทารํ ปิตรญฺจ มาตรํ,}

\prose{ธนานิ ธญฺญานิ จ พนฺธวานิ\footnote{พนฺธวานิ จ (อิ)}.}

\prose{หิตฺวาน กามานิ ยโถธิกานิ,}

\csromanheader{2}{เอโก จเร ขคฺควิสาณกปฺโป. (26)}
\proseitem{61}{สงฺโค เอโส ปริตฺตเมตฺถ โสขฺยํ,}

\prose{อปฺปสฺสาโท ทุกฺขเมตฺถ ภิยฺโย.}

\prose{คโฬ เอโส อิติ ญตฺวา มตีมา,}

\csromanheader{2}{เอโก จเร ขคฺควิสาณกปฺโป. (27)}
\proseitem{62}{สนฺทาลยิตฺวาน\footnote{ปทาลยิตฺวาน (ก)} สํโยชนานิ,}

\prose{ชาลํว เภตฺวา สลิลมฺพุจารี.}

\prose{อคฺคีว ทฑฺฒํ อนิวตฺตมาโน,}

\csromanheader{2}{เอโก จเร ขคฺควิสาณกปฺโป. (28)}
\chapterheadpageread{1. อุรควคฺค}
\proseitem{63}{\csromanfolio{289}โอกฺขิตฺตจกฺขู น จ ปาทโลโล,}

\prose{คุตฺตินฺทฺริโย รกฺขิตมานสาโน.}

\prose{อนวสฺสุโต อปริฑยฺหมาโน,}

\csromanheader{2}{เอโก จเร ขคฺควิสาณกปฺโป. (29)}
\proseitem{64}{โอหารยิตฺวา คิหิพฺยญฺชนานิ,}

\prose{สญฺฉนฺนปตฺโต\footnote{สญฺฉินฺนปตฺโต (สฺยา, อิ), ปจฺฉินฺนปตฺโต (ก)} ยถา ปาริฉตฺโต.}

\prose{กาสายวตฺโต อภินิกฺขมิตฺวา,}

\csromanheader{2}{เอโก จเร ขคฺควิสาณกปฺโป. (30)}
\proseitem{65}{รเสสุ เคธํ อกรํ อโลโล,}

\prose{อนญฺญโปสี สปทานจารี.}

\prose{กุเล กุเล อปฺปฏิพทฺธจิตฺโต\footnote{อปฺปฏิพนฺธจิตฺโต (ก)},}

\csromanheader{2}{เอโก จเร ขคฺควิสาณกปฺโป. (31)}
\proseitem{66}{ปหาย ปญฺจาวรณานิ เจตโส,}

\prose{อุปกฺกิเลเส พฺยปนุชฺช สพฺเพ.}

\prose{อนิสฺสิโต เฉตฺว\footnote{เฉตฺวา (สฺยา, อิ, ก)} สิเนหโทสํ\footnote{สฺเนหโทสํ (ก)},}

\csromanheader{2}{เอโก จเร ขคฺควิสาณกปฺโป. (32)}
\proseitem{67}{วิปิฏฺฐิกตฺวาน สุขํ ทุขญฺจ,}

\prose{ปุพฺเพว จ โสมนสฺสโทมนสฺสํ.}

\prose{ลทฺธานุเปกฺขํ สมถํ วิสุทฺธํ,}

\csromanheader{2}{เอโก จเร ขคฺควิสาณกปฺโป. (33)}
\proseitem{68}{อารทฺธวิริโย ปรมตฺถปตฺติยา,}

\prose{อลีนจิตฺโต อกุสีตวุตฺติ.}

\prose{ทฬฺหนิกฺกโม ถามพลูปปนฺโน,}

\csromanheader{2}{เอโก จเร ขคฺควิสาณกปฺโป. (34)}
\proseitem{69}{ปฏิสลฺลานํ ฌานมริญฺจมาโน,}

\prose{ธมฺเมสุ นิจฺจํ อนุธมฺมจารี.}

\prose{อาทีนวํ สมฺมสิตา ภเวสุ,}

\csromanheader{2}{เอโก จเร ขคฺควิสาณกปฺโป. (35)}
\gambhira{สุตฺตนิปาตปาฬิ}
\csromangathameasure{ตณฺหกฺขยํ ปตฺถยมปฺปมตฺโต, \\ อเนฬมูโค สุตวา สตีมา. \\ สงฺขาตธมฺโม นิยโต ปธานวา,}
\csromangathagroup{\csromangathastanza{\csromanfolio{290}ตณฺหกฺขยํ ปตฺถยมปฺปมตฺโต, \\ อเนฬมูโค\footnote{อเนลมูโค (สฺยา, อิ, ก)} สุตวา สตีมา. \\ สงฺขาตธมฺโม นิยโต ปธานวา,}}

\csromanheader{2}{เอโก จเร ขคฺควิสาณกปฺโป. (36)}
\proseitem{71}{สีโหว สทฺเทสุ อสนฺตสนฺโต,}

\prose{วาโตว ชาลมฺหิ อสชฺชมาโน.}

\prose{ปทุมํว โตเยน อลิปฺปมาโน\footnote{อลิมฺปมาโน (สี, สฺยา, ก); ขุ 8. 301 ปิฏฺเฐปิ.},}

\csromanheader{2}{เอโก จเร ขคฺควิสาณกปฺโป. (37)}
\proseitem{72}{สีโห ยถา ทาฐพลี ปสยฺห,}

\prose{ราชา มิคานํ อภิภุยฺย จารี.}

\prose{เสเวถ ปนฺตานิ เสนาสนานิ,}

\csromanheader{2}{เอโก จเร ขคฺควิสาณกปฺโป. (38)}
\proseitem{73}{เมตฺตํ อุเปกฺขํ กรุณํ วิมุตฺติํ,}

\prose{อาเสวมาโน มุทิตญฺจ กาเล.}

\prose{สพฺเพน โลเกน อวิรุชฺฌมาโน,}

\csromanheader{2}{เอโก จเร ขคฺควิสาณกปฺโป. (39)}
\proseitem{74}{ราคญฺจ โทสญฺจ ปหาย โมหํ,}

\prose{สนฺทาลยิตฺวาน สํโยชนานิ.}

\prose{อสนฺตสํ ชีวิตสงฺขยมฺหิ,}

\csromanheader{2}{เอโก จเร ขคฺควิสาณกปฺโป. (40)}
\proseitem{75}{ภชนฺติ เสวนฺติ จ การณตฺถา,}

\prose{นิกฺการณา ทุลฺลภา อชฺช มิตฺตา.}

\csromangathameasure{อตฺตฏฺฐปญฺญา อสุจี มนุสฺสา, \\ เอโก จเร ขคฺควิสาณกปฺโป. (41) ขคฺควิสาณสุตฺตํ ตติยํ.}
\csromangathagroup{\csromangathastanza{อตฺตฏฺฐปญฺญา อสุจี มนุสฺสา, \\ เอโก จเร ขคฺควิสาณกปฺโป. (41) ขคฺควิสาณสุตฺตํ ตติยํ.}}

\csromanmatikaanchor{mtk.237}
\csromantocmarkref{subsection}{3. ขคฺควิสาณสุตฺต}{mtk.237}
\bookmark[dest={mtk.237},level=2]{3. ขคฺควิสาณสุตฺต}
\csromanheader{2}{1. อุรควคฺค \textbf{4. กสิภารทฺวาชสุตฺต}}
\prose{\csromanfolio{291}เอวํ เม สุตํ\csromandash{}เอกํ สมยํ ภควา มคเธสุ วิหรติ ทกฺขิณาคิริสฺมิํ\footnote{ทกฺขิณคิริสฺมิํ (ก)} เอกนาฬายํ พฺราหฺมณคาเม.\csromanspacer{} เตน โข ปน สมเยน กสิภารทฺวาชสฺส พฺราหฺมณสฺส ปญฺจมตฺตานิ นงฺคลสตานิ ปยุตฺตานิ โหนฺติ วปฺปกาเล.\csromanspacer{} อถ โข ภควา ปุพฺพณฺหสมยํ นิวาเสตฺวา ปตฺตจีวรมาทาย เยน กสิภารทฺวาชสฺส พฺราหฺมณสฺส กมฺมนฺโต เตนุปสงฺกมิ.\csromanspacer{} เตน โข ปน สมเยน กสิภารทฺวาชสฺส พฺราหฺมณสฺส ปริเวสนา วตฺตติ.\csromanspacer{} อถ โข ภควา เยน ปริเวสนา เตนุปสงฺกมิ, อุปสงฺกมิตฺวา เอกมนฺตํ อฏฺฐาสิ.}

\prose{อทฺทสา โข กสิภารทฺวาโช พฺราหฺมโณ ภควนฺตํ ปิณฺฑาย ฐิตํ, ทิสฺวาน ภควนฺตํ เอตทโวจ “อหํ โข สมณ กสามิ จ วปามิ จ, กสิตฺวา จ วปิตฺวา จ ภุญฺชามิ.\csromanspacer{} ตฺวมฺปิ สมณ กสสฺสุ จ วปสฺสุ จ, กสิตฺวา จ วปิตฺวา จ ภุญฺชสฺสู”ติ.}

\prose{อหมฺปิ โข พฺราหฺมณ กสามิ จ วปามิ จ, กสิตฺวา จ วปิตฺวา จ ภุญฺชามีติ.\csromanspacer{} น โข ปน มยํ\footnote{น โข ปน สมณ (สฺยา)} ปสฺสาม โภโต โคตมสฺส ยุคํ วา นงฺคลํ วา ผาลํ วา ปาจนํ วา พลีพทฺเท\footnote{พลิวทฺเท (สี, อิ)} วา, อถ จ ปน ภวํ โคตโม เอวมาห “อหมฺปิ โข พฺราหฺมณ กสามิ จ วปามิ จ, กสิตฺวา จ วปิตฺวา จ ภุญฺชามี”ติ.}

\prose{อถ โข กสิภารทฺวาโช พฺราหฺมโณ ภควนฺตํ คาถาย อชฺฌภาสิ\csromandash{}}

\proseitem{76}{\csromansymbolfootnote{*}{สํ 1. 175 ปิฏฺเฐปิ.}“กสฺสโก ปฏิชานาสิ, น จ ปสฺสาม เต กสิํ.}

\prose{กสิํ โน ปุจฺฉิโต พฺรูหิ, ยถา ชาเนมุ เต กสิํ. (1)}

\proseitem{77}{สทฺธา พีชํ ตโป วุฏฺฐิ, ปญฺญา เม ยุคนงฺคลํ.}

\prose{หิรี อีสา มโน โยตฺตํ, สติ เม ผาลปาจนํ. (2)}

\proseitem{78}{กายคุตฺโต วจีคุตฺโต, อาหาเร อุทเร ยโต.}

\prose{สจฺจํ กโรมิ นิทฺทานํ, โสรจฺจํ เม ปโมจนํ. (3)}

\proseitem{79}{วีริยํ เม ธุรโธรยฺหํ, โยคกฺเขมาธิวาหนํ.}

\prose{คจฺฉติ อนิวตฺตนฺตํ, ยตฺถ คนฺตฺวา น โสจติ. (4)}

\gambhira{สุตฺตนิปาตปาฬิ}
\proseitem{80}{\csromanfolio{292}เอวเมสา กสี กฏฺฐา, สา โหติ อมตปฺผลา.}

\prose{เอตํ กสิํ กสิตฺวาน, สพฺพทุกฺขา ปมุจฺจตี”ติ. (5)}

\prose{อถ โข กสิภารทฺวาโช พฺราหฺมโณ มหติยา กํสปาติยา ปายสํ\footnote{ปายาสํ (สพฺพตฺถ)} วฑฺเฒตฺวา ภควโต อุปนาเมสิ “ภุญฺชตุ ภวํ โคตโม ปายสํ ตสฺสโก ภวํ, ยํ หิ ภวํ โคตโม อมตปฺผลํ\footnote{อมตปฺผลมฺปิ (สํ 1. 175) * อุปริ 349 ปิฏฺเฐปิ.} กสิํ กสตี”ติ.}

\proseitem{81}{\csromansymbolmark{*}คาถาภิคีตํ เม อโภชเนยฺยํ,}

\prose{สมฺปสฺสตํ พฺราหฺมณ เนส ธมฺโม.}

\csromangathameasure{คาถาภิคีตํ ปนุทนฺติ พุทฺธา, \\ ธมฺเม สตี พฺราหฺมณ วุตฺติเรสา.}
\csromangathagroup{\csromangathastanza{คาถาภิคีตํ ปนุทนฺติ พุทฺธา, \\ ธมฺเม สตี พฺราหฺมณ วุตฺติเรสา.}}

\csromanheader{2}{(6)}
\csromangathameasure{อญฺเญน จ เกวลินํ มเหสิํ, \\ ขีณาสวํ กุกฺกุจฺจวูปสนฺตํ. \\ อนฺเนน ปาเนน อุปฏฺฐหสฺสุ,}
\csromangathagroup{\csromangathastanza{อญฺเญน จ เกวลินํ มเหสิํ, \\ ขีณาสวํ กุกฺกุจฺจวูปสนฺตํ. \\ อนฺเนน ปาเนน อุปฏฺฐหสฺสุ,}}

\csromanheader{2}{เขตฺตํ หิ ตํ ปุญฺญเปกฺขสฺส โหตีติ. (7)}
\prose{อถ กสฺส จาหํ โภ โคตม อิมํ ปายสํ ทมฺมีติ.\csromanspacer{} น ขฺวาหํ ตํ พฺราหฺมณ ปสฺสามิ สเทวเก โลเก สมารเก สพฺรหฺมเก สสฺสมณพฺราหฺมณิยา ปชาย สเทวมนุสฺสาย, ยสฺส โส ปายโส ภุตฺโต สมฺมา ปริณามํ คจฺเฉยฺย, อญฺญตฺร ตถาคตสฺส วา ตถาคตสาวกสฺส วา.\csromanspacer{} เตน หิ ตฺวํ พฺราหฺมณ ตํ ปายสํ อปฺปหริเต วา ฉฑฺเฑหิ, อปฺปาณเก วา อุทเก โอปิลาเปหีติ.}

\prose{อถ โข กสิภารทฺวาโช พฺราหฺมโณ ตํ ปายสํ อปฺปาณเก อุทเก โอปิลาเปสิ.\csromanspacer{} อถ โข โส ปายโส อุทเก ปกฺขิตฺโต จิจฺจิฏายติ จิฏิจิฏายติ สนฺธูปายติ สมฺปธูปายติ\footnote{สนฺธูมายติ สมฺปธูมายติ (สฺยา) วิ 3. 320 ปิฏฺเฐปิ.}.\csromanspacer{} เสยฺยถาปิ นาม ผาโล ทิวสํสนฺตตฺโต\footnote{ทิวสสนฺตตฺโต (สี, สฺยา, กํ, อิ)} อุทเก ปกฺขิตฺโต จิจฺจิฏายติ จิฏิจิฏายติ สนฺธูปายติ สมฺปธูปายติ.\csromanspacer{} เอวเมวํ โส ปายโส อุทเก ปกฺขิตฺโต จิจฺจิฏายติ จิฏิจิฏายติ สนฺธูปายติ สมฺปธูปายติ.}

\prose{อถ โข กสิภารทฺวาโช พฺราหฺมโณ สํวิคฺโค โลมหฏฺฐชาโต เยน ภควา เตนุปสงฺกมิ, อุปสงฺกมิตฺวา ภควโต ปาเทสุ สิรสา}

\vspace{\tipitakaparskip}
\csromancenter{อุรควคฺค นิปติตฺวา ภควนฺตํ เอตทโวจ “อภิกฺกนฺตํ โภ โคตม, อภิกฺกนฺตํ โภ โคตม.\csromanspacer{} เสยฺยถาปิ โภ โคตม นิกฺกุชฺชิตํ วา อุกฺกุชฺเชยฺย, ปฏิจฺฉนฺนํ วา วิวเรยฺย, มูฬฺหสฺส วา มคฺคํ อาจิกฺเขยฺย, อนฺธกาเร วา เตลปชฺโชตํ ธาเรยฺย ‘จกฺขุมนฺโต รูปานิ ทกฺขนฺตี’ติ\footnote{ทกฺขินฺตีติ (สี, สฺยา, กํ, อิ)}, เอวเมวํ โภตา โคตเมน อเนกปริยาเยน ธมฺโม ปกาสิโต, เอสาหํ ภวนฺตํ โคตมํ สรณํ คจฺฉามิ ธมฺมญฺจ ภิกฺขุสํฆญฺจ, ลเภยฺยาหํ โภโต โคตมสฺส สนฺติเก ปพฺพชฺชํ, ลเภยฺยํ อุปสมฺปทนฺ”ติ.}
\prose{\csromanfolio{293}อลตฺถ โข กสิภารทฺวาโช พฺราหฺมโณ ภควโต สนฺติเก ปพฺพชฺชํ, อลตฺถ อุปสมฺปทํ.\csromanspacer{} อจิรูปสมฺปนฺโน โข ปนายสฺมา ภารทฺวาโช เอโก วูปกฏฺโฐ อปฺปมตฺโต อาตาปี ปหิตตฺโต วิหรนฺโต นจิรสฺเสว, ยสฺสตฺถาย กุลปุตฺตา สมฺมเทว อคารสฺมา อนคาริยํ ปพฺพชนฺติ, ตทนุตฺตรํ พฺรหฺมจริยปริโยสานํ ทิฏฺเฐว ธมฺเม สยํ อภิญฺญา สจฺฉิกตฺวา อุปสมฺปชฺช วิหาสิ. “ขีณา ชาติ, วุสิตํ พฺรหฺมจริยํ, กตํ กรณียํ, นาปรํ อิตฺถตฺตายา”ติ อพฺภญฺญาสิ.\csromanspacer{} อญฺญตโร จ\footnote{อญฺญตโร จ โข (สี, อิ), อญฺญตโร โข (สฺยา, กํ, ก)} ปนายสฺมา ภารทฺวาโช อรหตํ อโหสีติ.\csromanspacer{} กสิภารทฺวาชสุตฺตํ จตุตฺถํ.}

\csromanheader{2}{5. จุนฺทสุตฺต}
\proseitem{83}{ปุจฺฉามิ มุนิํ ปหูตปญฺญํ, (อิติ จุนฺโท กมฺมารปุตฺโต,)}

\prose{พุทฺธํ ธมฺมสฺสามิํ วีตตณฺหํ.}

\prose{ทฺวิปทุตฺตมํ\footnote{ทิปทุตฺตมํ (สี, สฺยา, กํ, อิ)} สารถีนํ ปวรํ,}

\csromanheader{2}{กติ โลเก สมณา ตทิงฺฆ พฺรูหิ. (1)}
\proseitem{84}{จตุโร สมณา น ปญฺจม’ตฺถิ, (จุนฺทาติ ภควา,)}

\prose{เต เต อาวิกโรมิ สกฺขิปุฏฺโฐ.}

\prose{มคฺคชิโน มคฺคเทสโก จ,}

\csromanheader{2}{มคฺเค ชีวติ โย จ มคฺคทูสี. (2)}
\gambhira{สุตฺตนิปาตปาฬิ}
\proseitem{85}{\csromanfolio{294}กํ มคฺคชินํ วทนฺติ พุทฺธา, (อิติ จุนฺโท กมฺมารปุตฺโต,)}

\prose{มคฺคกฺขายี กถํ อตุโลฺย โหติ.}

\prose{มคฺเค ชีวติ เม พฺรูหิ ปุฏฺโฐ,}

\csromanheader{2}{อถ เม อาวิกโรหิ มคฺคทูสิํ\footnote{มคฺคทูสี (ก)}. (3)}
\proseitem{86}{โย ติณฺณกถํกโถ วิสลฺโล,}

\prose{นิพฺพานาภิรโต อนานุคิทฺโธ.}

\prose{โลกสฺส สเทวกสฺส เนตา,}

\csromanheader{2}{ตาทิํ มคฺคชินํ วทนฺติ พุทฺธา. (4)}
\proseitem{87}{ปรมํ “ปรมนฺ”ติ โย’ธ ญตฺวา,}

\prose{อกฺขาติ วิภชเต อิเธว ธมฺมํ.}

\prose{ตํ กงฺขฉิทํ มุนิํ อเนชํ,}

\csromanheader{2}{ทุติยํ ภิกฺขุนมาหุ มคฺคเทสิํ. (5)}
\proseitem{88}{โย ธมฺมปเท สุเทสิเต,}

\prose{มคฺเค ชีวติ สญฺญโต สตีมา.}

\prose{อนวชฺชปทานิ เสวมาโน,}

\csromanheader{2}{ตติยํ ภิกฺขุนมาหุ มคฺคชีวิํ. (6)}
\proseitem{89}{ฉทนํ กตฺวาน สุพฺพตานํ,}

\prose{ปกฺขนฺที กุลทูสโก ปคพฺโภ.}

\prose{มายาวี อสญฺญโต ปลาโป,}

\csromanheader{2}{ปติรูเปน จรํ ส มคฺคทูสี. (7)}
\proseitem{90}{เอเต จ ปฏิวิชฺฌิ โย คหฏฺโฐ,}

\prose{สุตวา อริยสาวโก สปญฺโญ.}

\csromangathameasure{สพฺเพ เน’ตาทิสาติ ญตฺวา, \\ อิติ ทิสฺวา น หาเปติ ตสฺส สทฺธา. \\ กถํ หิ ทุฏฺเฐน อสมฺปทุฏฺฐํ, \\ สุทฺธํ อสุทฺเธน สมํ กเรยฺยาติ. (8) จุนฺทสุตฺตํ ปญฺจมํ.}
\csromangathagroup{\csromangathastanza{สพฺเพ เน’ตาทิสาติ\footnote{สพฺเพ เน ตาทิสาติ (สี, สฺยา, อิ)} ญตฺวา, \\ อิติ ทิสฺวา น หาเปติ ตสฺส สทฺธา. \\ กถํ หิ ทุฏฺเฐน อสมฺปทุฏฺฐํ, \\ สุทฺธํ อสุทฺเธน สมํ กเรยฺยาติ. (8) จุนฺทสุตฺตํ ปญฺจมํ.}}

\csromanheader{2}{1. อุรควคฺค \textbf{6. ปราภวสุตฺต}}
\prose{\csromanfolio{295}เอวํ เม สุตํ\csromandash{}เอกํ สมยํ ภควา สาวตฺถิยํ วิหรติ เชตวเน อนาถปิณฺฑิกสฺส อาราเม.\csromanspacer{} อถ โข อญฺญตรา เทวตา อภิกฺกนฺตาย รตฺติยา อภิกฺกนฺตวณฺณา เกวลกปฺปํ เชตวนํ โอภาเสตฺวา เยน ภควา เตนุปสงฺกมิ, อุปสงฺกมิตฺวา ภควนฺตํ อภิวาเทตฺวา เอกมนฺตํ อฏฺฐาสิ, เอกมนฺตํ ฐิตา โข สา เทวตา ภควนฺตํ คาถาย อชฺฌภาสิ\csromandash{}}

\csromanhangingbegin
\hangingheaditem{91}{“ปราภวนฺตํ ปุริสํ, มยํ ปุจฺฉาม โคตม\footnote{โคตมํ (สี, สฺยา)}.}
\hangingbody{ภวนฺตํ2 ปุฏฺฐุมาคมฺม, กิํ ปราภวโต มุขํ. (1)}
\csromanhangingend

\proseitem{92}{สุวิชาโน ภวํ โหติ, สุวิชาโน\footnote{ทุวิชาโน (สฺยา, ก)} ปราภโว.}

\prose{ธมฺมกาโม ภวํ โหติ, ธมฺมเทสฺสี ปราภโว. (2)}

\proseitem{93}{อิติ เหตํ วิชานาม, ปฐโม โส ปราภโว.}

\prose{ทุติยํ ภควา พฺรูหิ, กิํ ปราภวโต มุขํ. (3)}

\proseitem{94}{อสนฺต’สฺส ปิยา โหนฺติ, สนฺเต น กุรุเต ปิยํ.}

\prose{อสตํ ธมฺมํ โรเจติ, ตํ ปราภวโต มุขํ. (4)}

\proseitem{95}{อิติ เหตํ วิชานาม, ทุติโย โส ปราภโว.}

\prose{ตติยํ ภควา พฺรูหิ, กิํ ปราภวโต มุขํ. (5)}

\proseitem{96}{นิทฺทาสีลี สภาสีลี, อนุฏฺฐาตา จ โย นโร.}

\prose{อลโส โกธปญฺญาโณ, ตํ ปราภวโต มุขํ. (6)}

\proseitem{97}{อิติ เหตํ วิชานาม, ตติโย โส ปราภโว.}

\prose{จตุตฺถํ ภควา พฺรูหิ, กิํ ปราภวโต มุขํ. (7)}

\proseitem{98}{โย มาตรํ\footnote{โย มาตรํ วา (สี, สฺยา, กํ, อิ)} ปิตรํ วา, ชิณฺณกํ คตโยพฺพนํ.}

\prose{ปหุ สนฺโต น ภรติ, ตํ ปราภวโต มุขํ. (8)}

\proseitem{99}{อิติ เหตํ วิชานาม, จตุตฺโถ โส ปราภโว.}

\prose{ปญฺจมํ ภควา พฺรูหิ, กิํ ปราภวโต มุขํ. (9)}

\csromangathameasure{โย พฺราหฺมณํ สมณํ วา, อญฺญํ วาปิ วนิพฺพกํ.}
\csromangathagroup{\csromangathastanza{โย พฺราหฺมณํ\footnote{โย พฺราหฺมณํ วา (สี, สฺยา, กํ, อิ)} สมณํ วา, อญฺญํ วาปิ วนิพฺพกํ.}}

\proseitem{100}{มุสาวาเทน วญฺเจติ, ตํ ปราภวโต มุขํ. (10)}

\gambhira{สุตฺตนิปาตปาฬิ}
\proseitem{101}{\csromanfolio{296}อิติ เหตํ วิชานาม, ปญฺจโม โส ปราภโว.}

\prose{ฉฏฺฐมํ ภควา พฺรูหิ, กิํ ปราภวโต มุขํ. (11)}

\proseitem{102}{ปหูตวิตฺโต ปุริโส, สหิรญฺโญ สโภชโน.}

\prose{เอโก ภุญฺชติ สาทูนิ, ตํ ปราภวโต มุขํ. (12)}

\proseitem{103}{อิติ เหตํ วิชานาม, ฉฏฺฐโม โส ปราภโว.}

\prose{สตฺตมํ ภควา พฺรูหิ, กิํ ปราภวโต มุขํ. (13)}

\proseitem{104}{ชาติตฺถทฺโธ ธนตฺถทฺโธ, โคตฺตตฺถทฺโธ จ โย นโร.}

\prose{สญฺญาติํ อติมญฺเญติ, ตํ ปราภวโต มุขํ. (14)}

\proseitem{105}{อิติ เหตํ วิชานาม, สตฺตโม โส ปราภโว.}

\prose{อฏฺฐมํ ภควา พฺรูหิ, กิํ ปราภวโต มุขํ. (15)}

\proseitem{106}{อิตฺถิธุตฺโต สุราธุตฺโต, อกฺขธุตฺโต จ โย นโร.}

\prose{ลทฺธํ ลทฺธํ วินาเสติ, ตํ ปราภวโต มุขํ. (16)}

\proseitem{107}{อิติ เหตํ วิชานาม, อฏฺฐโม โส ปราภโว.}

\prose{นวมํ ภควา พฺรูหิ, กิํ ปราภวโต มุขํ. (17)}

\proseitem{108}{เสหิ ทาเรหิ อสนฺตุฏฺโฐ\footnote{ทาเรหฺยสนฺตุฏฺโฐ (ก)}, เวสิยาสุ ปทุสฺสติ\footnote{ปทิสฺสติ (สี)}.}

\prose{ทุสฺสติ\footnote{ทิสฺสติ (สี, อิ)} ปรทาเรสุ, ตํ ปราภวโต มุขํ. (18)}

\proseitem{109}{อิติ เหตํ วิชานาม, นวโม โส ปราภโว.}

\prose{ทสมํ ภควา พฺรูหิ, กิํ ปราภวโต มุขํ. (19)}

\proseitem{110}{อตีตโยพฺพโน โปโส, อาเนติ ติมฺพรุตฺถนิํ.}

\prose{ตสฺสา อิสฺสา น สุปติ, ตํ ปราภวโต มุขํ. (20)}

\csromanhangingbegin
\hangingheaditem{111}{อิติ เหตํ วิชานาม, ทสโม โส ปราภโว.}
\hangingbody{เอกาทสมํ ภควา พฺรูหิ, กิํ ปราภวโต มุขํ. (21)}
\csromanhangingend

\proseitem{112}{อิตฺถิํ โสณฺฑิํ วิกิรณิํ, ปุริสํ วาปิ ตาทิสํ.}

\prose{อิสฺสริยสฺมิํ ฐเปติ\footnote{ฐาเปติ (สี, อิ), ถเปติ (ก)}, ตํ ปราภวโต มุขํ. (22)}

\chapterheadpageread{1. อุรควคฺค}
\csromanhangingbegin
\hangingheaditem{113}{\csromanfolio{297}อิติ เหตํ วิชานาม, เอกาทสโม โส ปราภโว.}
\hangingbody{ทฺวาทสมํ ภควา พฺรูหิ, กิํ ปราภวโต มุขํ. (23)}
\csromanhangingend

\proseitem{114}{อปฺปโภโค มหาตโณฺห, ขตฺติเย ชายเต กุเล.}

\prose{โส จ รชฺชํ ปตฺถยติ, ตํ ปราภวโต มุขํ. (24)}

\proseitem{115}{เอเต ปราภเว โลเก, ปณฺฑิโต สมเวกฺขิย.}

\csromangathameasure{อริโย ทสฺสนสมฺปนฺโน, ส โลกํ ภชเต สิวนฺ”ติ. (25) ปราภวสุตฺตํ ฉฏฺฐํ.}
\csromangathagroup{\csromangathastanza{อริโย ทสฺสนสมฺปนฺโน, ส โลกํ ภชเต สิวนฺ”ติ. (25) ปราภวสุตฺตํ ฉฏฺฐํ.}}

\csromanheader{2}{7. วสลสุตฺต}
\prose{เอวํ เม สุตํ\csromandash{}เอกํ สมยํ ภควา สาวตฺถิยํ วิหรติ เชตวเน อนาถปิณฺฑิกสฺส อาราเม.\csromanspacer{} อถ โข ภควา ปุพฺพณฺหสมยํ นิวาเสตฺวา ปตฺตจีวรมาทาย สาวตฺถิํ ปิณฺฑาย ปาวิสิ.\csromanspacer{} เตน โข ปน สมเยน อคฺคิกภารทฺวาชสฺส พฺราหฺมณสฺส นิเวสเน อคฺคิ ปชฺชลิโต โหติ, อาหุติ ปคฺคหิตา.\csromanspacer{} อถ โข ภควา สาวตฺถิยํ สปทานํ ปิณฺฑาย จรมาโน เยน อคฺคิกภารทฺวาชสฺส พฺราหฺมณสฺส นิเวสนํ เตนุปสงฺกมิ.}

\prose{อทฺทสา โข อคฺคิกภารทฺวาโช พฺราหฺมโณ ภควนฺตํ ทูรโตว อาคจฺฉนฺตํ, ทิสฺวาน ภควนฺตํ เอตทโวจ “ตเตฺรว\footnote{อเตฺรว (สฺยา, ก)} มุณฺฑก, ตเตฺรว สมณก, ตเตฺรว วสลก ติฏฺฐาหี”ติ.}

\prose{เอวํ วุตฺเต ภควา อคฺคิกภารทฺวาชํ พฺราหฺมณํ เอตทโวจ “ชานาสิ ปน ตฺวํ พฺราหฺมณ วสลํ วา วสลกรเณ วา ธมฺเม”ติ.\csromanspacer{} น ขฺวาหํ โภ โคตม ชานามิ วสลํ วา วสลกรเณ วา ธมฺเม, สาธุ เม ภวํ โคตโม ตถา ธมฺมํ เทเสตุ, ยถาหํ ชาเนยฺยํ วสลํ วา วสลกรเณ วา ธมฺเมติ.\csromanspacer{} เตน หิ พฺราหฺมณ สุณาหิ สาธุกํ มนสิ กโรหิ, ภาสิสฺสามีติ. “เอวํ โภ”ติ โข อคฺคิกภารทฺวาโช พฺราหฺมโณ ภควโต ปจฺจสฺโสสิ.\csromanspacer{} ภควา เอตทโวจ\csromandash{}}

\gambhira{สุตฺตนิปาตปาฬิ}
\proseitem{116}{\csromanfolio{298}“โกธโน อุปนาหี จ, ปาปมกฺขี จ โย นโร.}

\prose{วิปนฺนทิฏฺฐิ มายาวี, ตํ ชญฺญา วสโล อิติ. (1)}

\proseitem{117}{เอกชํ วา ทฺวิชํ\footnote{ทิชํ (อิ)} วาปิ, โย’ธ ปาณํ วิหิํสติ.}

\prose{ยสฺส ปาเณ ทยา นตฺถิ, ตํ ชญฺญา วสโล อิติ. (2)}

\proseitem{118}{โย หนฺติ ปริรุนฺธติ\footnote{อุปรุนฺเธติ (สฺยา), อุปรุนฺธติ (ก)}, คามานิ นิคมานิ จ.}

\prose{นิคฺคาหโก\footnote{นิคฺฆาตโก (?)} สมญฺญาโต, ตํ ชญฺญา วสโล อิติ. (3)}

\proseitem{119}{คาเม วา ยทิ วารญฺเญ, ยํ ปเรสํ มมายิตํ.}

\prose{เถยฺยา อทินฺนมาเทติ\footnote{อทินฺนํ อาทิยติ (สี, อิ)}, ตํ ชญฺญา วสโล อิติ. (4)}

\proseitem{120}{โย หเว อิณมาทาย, จุชฺชมาโน\footnote{ภุญฺชมาโน (สฺยา, ก)} ปลายติ.}

\prose{‘น หิ เต อิณมตฺถี’ติ, ตํ ชญฺญา วสโล อิติ. (5)}

\proseitem{121}{โย เว กิญฺจิกฺขกมฺยตา, ปนฺถสฺมิํ วชนฺตํ ชนํ.}

\prose{หนฺตฺวา กิญฺจิกฺขมาเทติ, ตํ ชญฺญา วสโล อิติ. (6)}

\proseitem{122}{อตฺตเหตุ ปรเหตุ, ธนเหตุ จ\footnote{ธนเหตุ ว (ก)} โย นโร.}

\prose{สกฺขิปุฏฺโฐ มุสา พฺรูติ, ตํ ชญฺญา วสโล อิติ. (7)}

\proseitem{123}{โย ญาตีนํ สขีนํ วา, ทาเรสุ ปฏิทิสฺสติ.}

\prose{สาหสา\footnote{สหสา (สี, สฺยา)} สมฺปิเยน วา, ตํ ชญฺญา วสโล อิติ. (8)}

\proseitem{124}{โย มาตรํ ปิตรํ วา, ชิณฺณกํ คตโยพฺพนํ.}

\prose{ปหุ สนฺโต น ภรติ, ตํ ชญฺญา วสโล อิติ. (9)}

\proseitem{125}{โย มาตรํ ปิตรํ วา, ภาตรํ ภคินิํ สสุํ.}

\prose{หนฺติ โรเสติ วาจาย, ตํ ชญฺญา วสโล อิติ. (10)}

\proseitem{126}{โย อตฺถํ ปุจฺฉิโต สนฺโต, อนตฺถมนุสาสติ.}

\prose{ปฏิจฺฉนฺเนน มนฺเตติ, ตํ ชญฺญา วสโล อิติ. (11)}

\proseitem{127}{โย กตฺวา ปาปกํ กมฺมํ, ‘มา มํ ชญฺญา’ติ อิจฺฉติ\footnote{อภิ 2. 372 ปิฏฺเฐ ปสฺสิตพฺพํ.}.}

\prose{โย ปฏิจฺฉนฺนกมฺมนฺโต, ตํ ชญฺญา วสโล อิติ. (12)}

\chapterheadpageread{1. อุรควคฺค}
\proseitem{128}{\csromanfolio{299}โย เว ปรกุลํ คนฺตฺวา, ภุตฺวาน\footnote{ภุตฺวา จ (สฺยา, ก)} สุจิโภชนํ.}

\prose{อาคตํ นปฺปฏิปูเชติ, ตํ ชญฺญา วสโล อิติ. (13)}

\proseitem{129}{โย พฺราหฺมณํ สมณํ วา, อญฺญํ วาปิ วนิพฺพกํ.}

\prose{มุสาวาเทน วญฺเจติ, ตํ ชญฺญา วสโล อิติ. (14)}

\proseitem{130}{โย พฺราหฺมณํ สมณํ วา, ภตฺตกาเล อุปฏฺฐิเต.}

\prose{โรเสติ วาจา น จ เทติ, ตํ ชญฺญา วสโล อิติ. (15)}

\proseitem{131}{อสตํ โย’ธ ปพฺรูติ, โมเหน ปลิคุณฺฐิโต.}

\prose{กิญฺจิกฺขํ นิชิคีสาโน\footnote{นิชิคิํสาโน (สี, สฺยา, กํ, อิ)}, ตํ ชญฺญา วสโล อิติ. (16)}

\proseitem{132}{โย จตฺตานํ สมุกฺกํเส, ปเร จ มวชานาติ\footnote{มวชานติ (สี, สฺยา, อิ)}.}

\prose{นิหีโน เสน มาเนน, ตํ ชญฺญา วสโล อิติ. (17)}

\proseitem{133}{โรสโก กทริโย จ, ปาปิจฺโฉ มจฺฉรี สโฐ.}

\prose{อหิริโก อโนตฺตปฺปี, ตํ ชญฺญา วสโล อิติ. (18)}

\proseitem{134}{โย พุทฺธํ ปริภาสติ, อถ วา ตสฺส สาวกํ.}

\prose{ปริพฺพาชํ\footnote{ปริพฺพชํ (ก), ปริพฺพาชกํ (สฺยา, กํ)} คหฏฺฐํ วา, ตํ ชญฺญา วสโล อิติ. (19)}

\proseitem{135}{โย เว อนรหํ\footnote{อนรหา (สี, อิ)} สนฺโต, อรหํ ปฏิชานาติ\footnote{ปฏิชานติ (สี, สฺยา, อิ)}.}

\csromangathameasure{โจโร สพฺรหฺมเก โลเก, เอโส โข วสลาธโม.}
\csromangathagroup{\csromangathastanza{โจโร สพฺรหฺมเก โลเก, เอโส โข วสลาธโม.}}

\prose{เอเต โข วสลา วุตฺตา, มยา เยเต ปกาสิตา. (20)}

\proseitem{136}{น ชจฺจา วสโล โหติ, น ชจฺจา โหติ พฺราหฺมโณ.}

\prose{กมฺมุนา\footnote{กมฺมนา (สี, อิ)} วสโล โหติ, กมฺมุนา โหติ พฺราหฺมโณ. (21)}

\proseitem{137}{ตทมินาปิ ชานาถ,}

\prose{ยถาเมทํ\footnote{ยถาเปทํ (ก)} นิทสฺสนํ.}

\prose{จณฺฑาลปุตฺโต โสปาโก\footnote{สปาโก (?)},}

\csromanheader{2}{มาตงฺโค อิติ วิสฺสุโต. (22)}
\gambhira{สุตฺตนิปาตปาฬิ}
\proseitem{138}{\csromanfolio{300}โส ยสํ ปรมํ ปตฺโต\footnote{โส ยสปฺปรมปฺปตฺโต (สฺยา, ก)}, มาตงฺโค ยํ สุทุลฺลภํ.}

\vspace{\tipitakaparskip}
\csromancenter{อาคจฺฉุํ ตสฺสุปฏฺฐานํ, ขตฺติยา พฺราหฺมณา พหู. (23)}
\proseitem{139}{เทวยานํ อภิรุยฺห, วิรชํ โส มหาปถํ.}

\csromangathameasure{กามราคํ วิราเชตฺวา, พฺรหฺมโลกูปโค อหุ.}
\csromangathagroup{\csromangathastanza{กามราคํ วิราเชตฺวา, พฺรหฺมโลกูปโค อหุ.}}

\prose{น นํ ชาติ นิวาเรสิ, พฺรหฺมโลกูปปตฺติยา. (24)}

\csromanmatikaanchor{mtk.238}
\csromantocmarkref{subsection}{4. กสิภารทฺวาชสุตฺต}{mtk.238}
\bookmark[dest={mtk.238},level=2]{4. กสิภารทฺวาชสุตฺต}
\proseitem{140}{อชฺฌายกกุเล ชาตา, พฺราหฺมณา มนฺตพนฺธวา.}

\prose{เต จ ปาเปสุ กมฺเมสุ, อภิณฺหมุปทิสฺสเร. (25)}

\proseitem{141}{ทิฏฺเฐว ธมฺเม คารยฺหา, สมฺปราเย จ ทุคฺคติ.}

\prose{น เน ชาติ นิวาเรติ, ทุคฺคตฺยา\footnote{ทุคฺคจฺจา (สี, สฺยา, กํ, อิ)} ครหาย วา. (26)}

\proseitem{142}{น ชจฺจา วสโล โหติ,}

\prose{น ชจฺจา โหติ พฺราหฺมโณ.}

\csromangathameasure{กมฺมุนา วสโล โหติ,}
\csromangathagroup{\csromangathastanza{กมฺมุนา วสโล โหติ,}}

\csromanheader{2}{กมฺมุนา โหติ พฺราหฺมโณ”ติ. (27)}
\prose{เอวํ วุตฺเต อคฺคิกภารทฺวาโช พฺราหฺมโณ ภควนฺตํ เอตทโวจ “อภิกฺกนฺตํ โภ โคตม ฯเปฯ อุปาสกํ มํ ภวํ โคตโม ธาเรตุ อชฺชตคฺเค ปาณุเปตํ สรณํ คตนฺ”ติ.\csromanspacer{} วสลสุตฺตํ สตฺตมํ.}

\csromanheader{2}{8. เมตฺตสุตฺต}
\proseitem{143}{\csromansymbolfootnote{*}{เหฏฺฐา 10 ปิฏฺเฐปิ.}กรณียมตฺถกุสเลน,}

\prose{ยนฺต สนฺตํ ปทํ อภิสเมจฺจ.}

\csromangathameasure{สกฺโก อุชู จ สุหุชู จ,}
\csromangathagroup{\csromangathastanza{สกฺโก อุชู จ สุหุชู\footnote{สูชู (สี)} จ,}}

\csromanheader{2}{สูวโจ จสฺส มุทุ อนติมานี. (1)}
\csromangathameasure{สนฺตุสฺสโก จ สุภโร จ, \\ อปฺปกิจฺโจ จ สลฺลหุกวุตฺติ. \\ สนฺตินฺทฺริโย จ นิปโก จ,}
\csromangathagroup{\csromangathastanza{สนฺตุสฺสโก จ สุภโร จ, \\ อปฺปกิจฺโจ จ สลฺลหุกวุตฺติ. \\ สนฺตินฺทฺริโย จ นิปโก จ,}}

\csromanheader{2}{อปฺปคพฺโภ กุเลสฺวนนุคิทฺโธ. (2)}
\chapterheadpageread{1. อุรควคฺค}
\proseitem{145}{\csromanfolio{301}น จ ขุทฺทมาจเร กิญฺจิ,}

\prose{เยน วิญฺญู ปเร อุปวเทยฺยุํ.}

\prose{สุขิโน ว เขมิโน โหนฺตุ,}

\csromanheader{2}{สพฺพสตฺตา\footnote{สพฺเพ สตฺตา (สี, สฺยา)} ภวนฺตุ สุขิตตฺตา. (3)}
\proseitem{146}{เย เกจิ ปาณภูต’ตฺถิ,}

\prose{ตสา วา ถาวรา ว’นวเสสา.}

\prose{ทีฆา วา เย ว มหนฺตา\footnote{มหนฺต (?)},}

\csromanheader{2}{มชฺฌิมา รสฺสกา อณุกถูลา. (4)}
\proseitem{147}{ทิฏฺฐา วา เย ว อทิฏฺฐา\footnote{อทิฏฺฐ (?)},}

\prose{เย ว\footnote{เย จ (สี, สฺยา, กํ, อิ)} ทูเร วสนฺติ อวิทูเร.}

\prose{ภูตา ว สมฺภเวสี ว\footnote{ภูตา วา สมฺภเวสี วา (สฺยา, กํ, อิ, ก)},}

\csromanheader{2}{สพฺพสตฺตา ภวนฺตุ สุขิตตฺตา. (5)}
\proseitem{148}{น ปโร ปรํ นิกุพฺเพถ,}

\prose{นาติมญฺเญถ กตฺถจิ น กญฺจิ\footnote{นํ กญฺจิ (สี, อิ), นํ กิญฺจิ (สฺยา), น กิญฺจิ (ก)}.}

\csromanmatikaanchor{mtk.239}
\csromantocmarkref{subsection}{5. จุนฺทสุตฺต}{mtk.239}
\bookmark[dest={mtk.239},level=2]{5. จุนฺทสุตฺต}
\prose{พฺยาโรสนา ปฏิฆสญฺญ,}

\csromanheader{2}{นาญฺญมญฺญสฺส ทุกฺขมิจฺเฉยฺย. (6)}
\proseitem{149}{มาตา ยถา นิยํปุตฺต,}

\prose{มายุสา เอกปุตฺตมนุรกฺเข.}

\prose{เอวมฺปิ สพฺพภูเตสุ,}

\csromanheader{2}{มานสํ ภาวเย อปริมาณํ. (7)}
\proseitem{150}{เมตฺตญฺจ สพฺพโลกสฺมิ,}

\prose{มานสํ ภาวเย อปริมาณํ.}

\prose{อุทฺธํ อโธ จ ติริยญฺจ,}

\csromanheader{2}{อสมฺพาธํ อเวร’มสปตฺตํ. (8)}
\gambhira{สุตฺตนิปาตปาฬิ}
\proseitem{151}{\csromanfolio{302}ติฏฺฐํ จรํ นิสินฺโน ว\footnote{วา (สี, สฺยา, กํ, อิ)},}

\prose{สยาโน ยาวตา’สฺส วิตมิทฺโธ\footnote{วิคตมิทฺโธ (พหูสุ)}.}

\prose{เอตํ สติํ อธิฏฺเฐยฺย,}

\csromanheader{2}{พฺรหฺมเมตํ วิหารมิธมาหุ. (9)}
\proseitem{152}{ทิฏฺฐิญฺจ อนุปคฺคมฺม,}

\prose{สีลวา ทสฺสเนน สมฺปนฺโน.}

\csromangathameasure{กาเมสุ วินย เคธํ, \\ น หิ ชาตุคฺคพฺภเสยฺย ปุน เรตีติ. (10) เมตฺตสุตฺตํ อฏฺฐมํ.}
\csromangathagroup{\csromangathastanza{กาเมสุ วินย\footnote{วิเนยฺย (สี, สฺยา, อิ)} เคธํ, \\ น หิ ชาตุคฺคพฺภเสยฺย ปุน เรตีติ. (10) เมตฺตสุตฺตํ อฏฺฐมํ.}}

\csromanheader{2}{9. เหมวตสุตฺต}
\proseitem{153}{อชฺช ปนฺนรโส อุโปสโถ, (อิติ สาตาคิโร ยกฺโข,)}

\prose{ทิพฺพา\footnote{ทิพฺยา (สี, สฺยา, กํ, อิ)} รตฺติ อุปฏฺฐิตา.}

\prose{อโนมนามํ สตฺถารํ,}

\csromanheader{2}{หนฺท ปสฺสาม โคตมํ. (1)}
\proseitem{154}{กจฺจิ มโน สุปณิหิโต, (อิติ เหมวโต ยกฺโข,)}

\prose{สพฺพภูเตสุ ตาทิโน.}

\prose{กจฺจิ อิฏฺเฐ อนิฏฺเฐ จ,}

\csromanheader{2}{สงฺกปฺป’สฺส วสีกตา. (2)}
\proseitem{155}{มโน จสฺส สุปณิหิโต, (อิติ สาตาคิโร ยกฺโข,)}

\prose{สพฺพภูเตสุ ตาทิโน.}

\prose{อโถ อิฏฺเฐ อนิฏฺเฐ จ,}

\csromanheader{2}{สงฺกปฺป’สฺส วสีกตา. (3)}
\proseitem{156}{กจฺจิ อทินฺนํ นาทิยติ, (อิติ เหมวโต ยกฺโข,)}

\prose{กจฺจิ ปาเณสุ สญฺญโต.}

\prose{กจฺจิ อารา ปมาทมฺหา,}

\csromanheader{2}{กจฺจิ ฌานํ น ริญฺจติ. (4)}
\chapterheadpageread{1. อุรควคฺค}
\csromanmatikaanchor{mtk.240}
\csromantocmarkref{subsection}{6. ปราภวสุตฺต}{mtk.240}
\bookmark[dest={mtk.240},level=2]{6. ปราภวสุตฺต}
\proseitem{157}{\csromanfolio{303}น โส อทินฺนํ อาทิยติ, (อิติ สาตาคิโร ยกฺโข,)}

\prose{อโถ ปาเณสุ สญฺญโต.}

\prose{อโถ อารา ปมาทมฺหา,}

\csromanheader{2}{พุทฺโธ ฌานํ น ริญฺจติ. (5)}
\proseitem{158}{กจฺจิ มุสา น ภณติ, (อิติ เหมวโต ยกฺโข,)}

\prose{กจฺจิ น ขีณพฺยปฺปโถ.}

\prose{กจฺจิ เวภูติกํ นาห,}

\csromanheader{2}{กจฺจิ สมฺผํ น ภาสติ. (6)}
\proseitem{159}{มุสา จ โส น ภณติ, (อิติ สาตาคิโร ยกฺโข,)}

\prose{อโถ น ขีณพฺยปฺปโถ.}

\prose{อโถ เวภูติยํ นาห,}

\csromanheader{2}{มนฺตา อตฺถํ จ\footnote{อตฺถํโส (สี, อิ, ก)} ภาสติ. (7)}
\proseitem{160}{กจฺจิ น รชฺชติ กาเมสุ, (อิติ เหมวโต ยกฺโข,)}

\prose{กจฺจิ จิตฺตํ อนาวิลํ.}

\prose{กจฺจิ โมหํ อติกฺกนฺโต,}

\csromanheader{2}{กจฺจิ ธมฺเมสุ จกฺขุมา. (8)}
\proseitem{161}{น โส รชฺชติ กาเมสุ, (อิติ สาตาคิโร ยกฺโข,)}

\prose{อโถ จิตฺตํ อนาวิลํ.}

\prose{สพฺพโมหํ อติกฺกนฺโต,}

\csromanheader{2}{พุทฺโธ ธมฺเมสุ จกฺขุมา. (9)}
\proseitem{162}{กจฺจิ วิชฺชาย สมฺปนฺโน, (อิติ เหมวโต ยกฺโข,)}

\prose{กจฺจิ สํสุทฺธจารโณ.}

\prose{กจฺจิ’สฺส อาสวา ขีณา,}

\csromanheader{2}{กจฺจิ นตฺถิ ปุนพฺภโว. (10)}
\proseitem{163}{วิชฺชาย เจว สมฺปนฺโน, (อิติ สาตาคิโร ยกฺโข,)}

\prose{อโถ สํสุทฺธจารโณ.}

\prose{สพฺพ’สฺส อาสวา ขีณา,}

\csromanheader{2}{นตฺถิ ตสฺส ปุนพฺภโว. (11)}
\gambhira{สุตฺตนิปาตปาฬิ}
\proseitem{164}{\csromanfolio{304}สมฺปนฺนํ มุนิโน จิตฺตํ, กมฺมุนา พฺยปฺปเถน จ.}

\prose{วิชฺชาจรณสมฺปนฺนํ, ธมฺมโต นํ ปสํสสิ. (12)}

\proseitem{165}{สมฺปนฺนํ มุนิโน จิตฺตํ, กมฺมุนา พฺยปฺปเถน จ.}

\prose{วิชฺชาจรณสมฺปนฺนํ, ธมฺมโต อนุโมทสิ. (13)}

\proseitem{166}{สมฺปนฺนํ มุนิโน จิตฺตํ, กมฺมุนา พฺยปฺปเถน จ.}

\prose{วิชฺชาจรณสมฺปนฺนํ, หนฺท ปสฺสาม โคตมํ. (14)}

\proseitem{167}{เอณิชงฺฆํ กิสํ วีรํ\footnote{ธีรํ (สฺยา)}, อปฺปาหารํ อโลลุปํ.}

\prose{มุนิํ วนสฺมิํ ฌายนฺตํ, เอหิ ปสฺสาม โคตมํ. (15)}

\proseitem{168}{สีหํเวกจรํ นาคํ, กาเมสุ อนเปกฺขินํ.}

\prose{อุปสงฺกมฺม ปุจฺฉาม, มจฺจุปาสปฺปโมจนํ. (16)}

\csromanhangingbegin
\hangingheaditem{169}{อกฺขาตารํ ปวตฺตารํ, สพฺพธมฺมาน ปารคุํ.}
\hangingbody{พุทฺธํ เวรภยาตีตํ, มยํ ปุจฺฉาม โคตมํ. (17)}
\csromanhangingend

\proseitem{170}{กิสฺมิํ โลโก สมุปฺปนฺโน, (อิติ เหมวโต ยกฺโข,)}

\prose{กิสฺมิํ กุพฺพติ สนฺถวํ\footnote{สนฺธวํ (ก)}.}

\prose{กิสฺส โลโก อุปาทาย,}

\csromanheader{2}{กิสฺมิํ โลโก วิหญฺญติ. (18)}
\proseitem{171}{ฉสุ\footnote{ฉสฺสุ (สี, อิ)} โลโก สมุปฺปนฺโน, (เหวมตาติ ภควา,)}

\prose{ฉสุ กุพฺพติ สนฺถวํ.}

\prose{ฉนฺนเมว อุปาทาย,}

\csromanheader{2}{ฉสุ โลโก วิหญฺญติ. (19)}
\proseitem{172}{กตมํ ตํ อุปาทานํ, ยตฺถ โลโก วิหญฺญติ.}

\prose{นิยฺยานํ ปุจฺฉิโต พฺรูหิ, กถํ ทุกฺขา ปมุจฺจติ\footnote{ปมุญฺจติ (สฺยา)}. (20)}

\proseitem{173}{ปญฺจ กามคุณา โลเก, มโนฉฏฺฐา ปเวทิตา.}

\csromanmatikaanchor{mtk.241}
\csromantocmarkref{subsection}{7. วสลสุตฺต}{mtk.241}
\bookmark[dest={mtk.241},level=2]{7. วสลสุตฺต}
\prose{เอตฺถ ฉนฺทํ วิราเชตฺวา, เอวํ ทุกฺขา ปมุจฺจติ. (21)}

\proseitem{174}{เอตํ โลกสฺส นิยฺยานํ, อกฺขาตํ โว ยถาตถํ.}

\prose{เอตํ โว อหมกฺขามิ, เอวํ ทุกฺขา ปมุจฺจติ. (22)}

\chapterheadpageread{1. อุรควคฺค}
\proseitem{175}{\csromanfolio{305}โก สู’ธ ตรติ โอฆํ, โก’ธ ตรติ อณฺณวํ.}

\prose{อปฺปติฏฺเฐ อนาลมฺเพ, โก คมฺภีเร น สีทติ. (23)}

\proseitem{176}{สพฺพทา สีลสมฺปนฺโน, ปญฺญวา สุสมาหิโต.}

\prose{อชฺฌตฺตจินฺตี\footnote{อชฺฌตฺตสญฺญี (สฺยา, กํ, ก)} สติมา, โอฆํ ตรติ ทุตฺตรํ. (24)}

\proseitem{177}{วิรโต กามสญฺญาย, สพฺพสํโยชนาติโค.}

\prose{นนฺทีภวปริกฺขีโณ, โส คมฺภีเร น สีทติ. (25)}

\proseitem{178}{คมฺภีรปญฺญํ นิปุณตฺถทสฺสิํ,}

\prose{อกิญฺจนํ กามภเว อสตฺตํ.}

\prose{ตํ ปสฺสถ สพฺพธิ วิปฺปมุตฺตํ,}

\csromanheader{2}{ทิพฺเพ ปเถ กมมานํ มเหสิํ. (26)}
\proseitem{179}{อโนมนามํ นิปุณตฺถทสฺสิํ,}

\prose{ปญฺญาททํ กามาลเย อสตฺตํ.}

\prose{ตํ ปสฺสถ สพฺพวิทุํ สุเมธํ,}

\csromanheader{2}{อริเย ปเถ กมมานํ มเหสิํ. (27)}
\proseitem{180}{สุทิฏฺฐํ วต โน อชฺช, สุปฺปภาตํ สุหุฏฺฐิตํ.}

\prose{ยํ อทฺทสาม สมฺพุทฺธํ, โอฆติณฺณมนาสวํ. (28)}

\proseitem{181}{อิเม ทสสตา ยกฺขา, อิทฺธิมนฺโต ยสสฺสิโน.}

\prose{สพฺเพ ตํ สรณํ ยนฺติ, ตฺวํ โน สตฺถา อนุตฺตโร. (29)}

\proseitem{182}{เต มยํ วิจริสฺสาม,}

\prose{คามา คามํ นคา นคํ.}

\csromangathameasure{นมสฺสมานา สมฺพุทฺธํ, \\ ธมฺมสฺส จ สุธมฺมตนฺติ. (30) เหมวตสุตฺตํ นวมํ.}
\csromangathagroup{\csromangathastanza{นมสฺสมานา สมฺพุทฺธํ, \\ ธมฺมสฺส จ สุธมฺมตนฺติ. (30) เหมวตสุตฺตํ นวมํ.}}

\csromanheader{2}{10. อาฬวกสุตฺต}
\prose{\csromansymbolfootnote{*}{สํ 1. 216 ปิฏฺเฐปิ.}เอวํ เม สุตํ\csromandash{}เอกํ สมยํ ภควา อาฬวิยํ วิหรติ อาฬวกสฺส ยกฺขสฺส ภวเน.\csromanspacer{} อถ โข อาฬวโก ยกฺโข เยน ภควา เตนุปสงฺกมิ, อุปสงฺกมิตฺวา ภควนฺตํ เอตทโวจ “นิกฺขม สมณา”ติ.}

\gambhira{สุตฺตนิปาตปาฬิ}
\prose{\csromanfolio{306}“สาธาวุโส”ติ ภควา นิกฺขมิ. “ปวิส สมณา”ติ. “สาธาวุโส”ติ ภควา ปาวิสิ.}

\prose{ทุติยมฺปิ โข ฯเปฯ.\csromanspacer{} ตติยมฺปิ โข อาฬวโก ยกฺโข ภควนฺตํ เอตทโวจ “นิกฺขม สมณา”ติ. “สาธาวุโส”ติ ภควา นิกฺขมิ. “ปวิส สมณา”ติ. “สาธาวุโส”ติ ภควา ปาวิสิ.}

\prose{จตุตฺถมฺปิ โข อาฬวโก ยกฺโข ภควนฺตํ เอตทโวจ “นิกฺขม สมณา”ติ.\csromanspacer{} น ขฺวาหํ ตํ อาวุโส นิกฺขมิสฺสามิ, ยํ เต กรณียํ ตํ กโรหีติ.}

\prose{ปญฺหํ ตํ สมณ ปุจฺฉิสฺสามิ.\csromanspacer{} สเจ เม น พฺยากริสฺสสิ, จิตฺตํ วา เต ขิปิสฺสามิ, หทยํ วา เต ผาเลสฺสามิ, ปาเทสุ วา คเหตฺวา ปารคงฺคาย ขิปิสฺสามีติ.}

\prose{น ขฺวาหํ ตํ อาวุโส ปสฺสามิ สเทวเก โลเก สมารเก สพฺรหฺมเก สสฺสมณพฺราหฺมณิยา ปชาย สเทวมนุสฺสาย, โย เม จิตฺตํ วา ขิเปยฺย, หทยํ วา ผาเลยฺย, ปาเทสุ วา คเหตฺวา ปารคงฺคาย ขิเปยฺย, อปิ จ ตฺวํ อาวุโส ปุจฺฉ, ยทากงฺขสีติ.\csromanspacer{} อถ โข อาฬวโก ยกฺโข ภควนฺตํ คาถาย อชฺฌภาสิ\csromandash{}}

\proseitem{183}{“กิํ สู’ธ วิตฺตํ ปุริสสฺส เสฏฺฐํ,}

\prose{กิํ สุ\footnote{กิํ สู (สี)} สุจิณฺณํ สุขมาวหาติ.}

\prose{กิํ สุ หเว สาทุตรํ รสานํ,}

\csromanheader{2}{กถํชีวิํชีวิตมาหุ เสฏฺฐํ. (1)}
\proseitem{184}{สทฺธีธ วิตฺตํ ปุริสสฺส เสฏฺฐํ,}

\prose{ธมฺโม สุจิณฺโณ สุขมาวหาติ.}

\prose{สจฺจํ หเว สาทุตรํ รสานํ,}

\csromanheader{2}{ปญฺญาชีวิํชีวิตมาหุ เสฏฺฐํ. (2)}
\proseitem{185}{กถํ สุ ตรติ โอฆํ, กถํ สุ ตรติ อณฺณวํ.}

\prose{กถํ สุ ทุกฺขมจฺเจติ, กถํ สุ ปริสุชฺฌติ. (3)}

\proseitem{186}{สทฺธาย ตรติ โอฆํ, อปฺปมาเทน อณฺณวํ.}

\prose{วีริเยน\footnote{วิริเยน (สี, สฺยา, กํ, อิ)} ทุกฺขมจฺเจติ, ปญฺญาย ปริสุชฺฌติ. (4)}

\chapterheadpageread{1. อุรควคฺค}
\proseitem{187}{\csromanfolio{307}กถํ สุ ลภเต ปญฺญํ, กถํ สุ วินฺทเต ธนํ.}

\csromangathameasure{กถํ สุ กิตฺติํ ปปฺโปติ, กถํ มิตฺตานิ คนฺถติ.}
\csromangathagroup{\csromangathastanza{กถํ สุ กิตฺติํ ปปฺโปติ, กถํ มิตฺตานิ คนฺถติ.}}

\prose{อสฺมา โลกา ปรํ โลกํ, กถํ เปจฺจ น โสจติ. (5)}

\proseitem{188}{สทฺทหาโน อรหตํ, ธมฺมํ นิพฺพานปตฺติยา.}

\prose{สุสฺสูสํ\footnote{สุสฺสูสา (สี, อิ)} ลภเต ปญฺญํ, อปฺปมตฺโต วิจกฺขโณ. (6)}

\proseitem{189}{ปติรูปการี ธุรวา, อุฏฺฐาตา วินฺทเต ธนํ.}

\prose{สจฺเจน กิตฺติํ ปปฺโปติ, ททํ มิตฺตานิ คนฺถติ. (7)}

\proseitem{190}{ยสฺเสเต จตุโร ธมฺมา, สทฺธสฺส ฆรเมสิโน.}

\prose{สจฺจํ ธมฺโม\footnote{ทโม (?)} ธิติ จาโค, ส เว เปจฺจ น โสจติ. (8)}

\proseitem{191}{อิงฺฆ อญฺเญปิ ปุจฺฉสฺสุ, ปุถู สมณพฺราหฺมเณ.}

\prose{ยทิ สจฺจา ทมา จาคา, ขนฺตฺยา ภิยฺโย’ธ วิชฺชติ. (9)}

\proseitem{192}{กถํ นุ ทานิ ปุจฺเฉยฺยํ, ปุถู สมณพฺราหฺมเณ.}

\prose{โยหํ\footnote{โสหํ (สี, อิ)} อชฺช ปชานามิ, โย อตฺโถ สมฺปรายิโก. (10)}

\proseitem{193}{อตฺถาย วต เม พุทฺโธ, วาสายา’ฬวิมาคมา.}

\prose{โยหํ3 อชฺช ปชานามิ, ยตฺถ ทินฺนํ มหปฺผลํ. (11)}

\csromanhangingbegin
\hangingheaditem{194}{โส อหํ วิจริสฺสามิ, คามา คามํ ปุรา ปุรํ.}
\hangingbody{นมสฺสมาโน สมฺพุทฺธํ, ธมฺมสฺส จ สุธมฺมตนฺ”ติ.}
\csromanhangingend

\prose{(12) อาฬวกสุตฺตํ ทสมํ.}

\csromanheader{2}{11. วิชยสุตฺต}
\proseitem{195}{จรํ วา ยทิ วา ติฏฺฐํ, นิสินฺโน อุท วา สยํ.}

\prose{สมิญฺเชติ ปสาเรติ, เอสา กายสฺส อิญฺชนา. (1)}

\proseitem{196}{อฏฺฐินหารุสํยุตฺโต\footnote{อฏฺฐินฺหารูหิ สํยุตฺโต (สฺยา, ก)}, ตจมํสาวเลปโน.}

\prose{ฉวิยา กาโย ปฏิจฺฉนฺโน, ยถาภูตํ น ทิสฺสติ. (2)}

\csromanmatikaanchor{mtk.242}
\csromantocmarkref{subsection}{8. เมตฺตสุตฺต}{mtk.242}
\bookmark[dest={mtk.242},level=2]{8. เมตฺตสุตฺต}
\gambhira{สุตฺตนิปาตปาฬิ}
\proseitem{197}{\csromanfolio{308}อนฺตปูโร อุทรปูโร, ยกนเปฬสฺส\footnote{ยกเปฬสฺส (สี, สฺยา)} วตฺถิโน.}

\prose{หทยสฺส ปปฺผาสสฺส, วกฺกสฺส ปิหกสฺส จ. (3)}

\proseitem{198}{สิงฺฆาณิกาย เขฬสฺส, เสทสฺส จ เมทสฺส จ.}

\prose{โลหิตสฺส ลสิกาย, ปิตฺตสฺส จ วสาย จ. (4)}

\proseitem{199}{อถสฺส นวหิ โสเตหิ, อสุจี สวติ สพฺพทา.}

\prose{อกฺขิมฺหา อกฺขิคูถโก, กณฺณมฺหา กณฺณคูถโก. (5)}

\proseitem{200}{สิงฺฆาณิกา จ นาสโต, มุเขน วมเต’กทา.}

\prose{ปิตฺตํ เสมฺหญฺจ วมติ, กายมฺหา เสทชลฺลิกา. (6)}

\proseitem{201}{อถสฺส สุสิรํ สีสํ, มตฺถลุงฺคสฺส ปูริตํ.}

\prose{สุภโต นํ มญฺญติ พาโล, อวิชฺชาย ปุรกฺขโต. (7)}

\proseitem{202}{ยทา จ โส มโต เสติ, อุทฺธุมาโต วินีลโก.}

\prose{อปวิทฺโธ สุสานสฺมิํ, อนเปกฺขา โหนฺติ ญาตโย. (8)}

\proseitem{203}{ขาทนฺติ นํ สุวานา\footnote{สุปาณา (อิ)} จ, สิงฺคาลา3จ วกา กิมี.}

\prose{กากา คิชฺฌา จ ขาทนฺติ, เย จญฺเญ สนฺติ ปาณิโน. (9)}

\proseitem{204}{สุตฺวาน พุทฺธวจนํ, ภิกฺขุ ปญฺญาณวา อิธ.}

\prose{โส โข นํ ปริชานาติ, ยถาภูตญฺหิ ปสฺสติ. (10)}

\proseitem{205}{ยถา อิทํ ตถา เอตํ, ยถา เอตํ ตถา อิทํ.}

\prose{อชฺฌตฺตญฺจ พหิทฺธา จ, กาเย ฉนฺทํ วิราชเย. (11)}

\proseitem{206}{ฉนฺทราควิรตฺโต โส, ภิกฺขุ ปญฺญาณวา อิธ.}

\prose{อชฺฌคา อมตํ สนฺติํ, นิพฺพานํ ปทมจฺจุตํ. (12)}

\proseitem{207}{ทฺวิปาทโกยํ\footnote{ทิปาทโกยํ (สี, สฺยา, กํ, อิ)} อสุจิ, ทุคฺคนฺโธ ปริหีรติ\footnote{ปริหารติ (ก)}.}

\prose{นานากุณปปริปูโร, วิสฺสวนฺโต ตโต ตโต. (13)}

\proseitem{208}{เอตาทิเสน กาเยน, โย มญฺเญ อุณฺณเมตเว\footnote{อุนฺนเมตเว (?)}.}

\csromangathameasure{ปรํ วา อวชาเนยฺย, กิมญฺญตฺร อทสฺสนาติ. (14) วิชยสุตฺตํ เอกาทสมํ.}
\csromangathagroup{\csromangathastanza{ปรํ วา อวชาเนยฺย, กิมญฺญตฺร อทสฺสนาติ. (14) วิชยสุตฺตํ เอกาทสมํ.}}

\csromanheader{2}{1. อุรควคฺค \textbf{12. มุนิสุตฺต}}
\proseitem{209}{\csromanfolio{309}สนฺถวาโต\footnote{สนฺธวโต (ก)} ภยํ ชาตํ,}

\prose{นิเกตา ชายเต รโช.}

\prose{อนิเกตมสนฺถวํ,}

\csromanheader{2}{เอตํ เว มุนิทสฺสนํ. (1)}
\proseitem{210}{โย ชาตมุจฺฉิชฺช น โรปเยยฺย,}

\prose{ชายนฺตมสฺส นานุปฺปเวจฺเฉ.}

\prose{ตมาหุ เอกํ มุนินํ จรนฺตํ,}

\csromanheader{2}{อทฺทกฺขิ โส สนฺติปทํ มเหสิ. (2)}
\proseitem{211}{สงฺขาย วตฺถูนิ ปมาย\footnote{ปหาย (ก-สี, ก), สมาย (ก) ป + มี + ตฺวา = ปมาย, ยถา นิสฺสายาติปทํ.} พีชํ,}

\prose{สิเนหมสฺส นานุปฺปเวจฺเฉ.}

\prose{ส เว มุนี ชาติขยนฺตทสฺสี,}

\csromanheader{2}{ตกฺกํ ปหาย น อุเปติ สงฺขํ. (3)}
\proseitem{212}{อญฺญาย สพฺพานิ นิเวสนานิ,}

\prose{อนิกามยํ อญฺญตรมฺปิ เตสํ.}

\prose{ส เว มุนี วีตเคโธ อคิทฺโธ.}

\csromanmatikaanchor{mtk.243}
\csromantocmarkref{subsection}{9. เหมวตสุตฺต}{mtk.243}
\bookmark[dest={mtk.243},level=2]{9. เหมวตสุตฺต}
\csromanheader{2}{นายูหตี ปารคโต หิ โหติ. (4)}
\proseitem{213}{สพฺพาภิภุํ สพฺพวิทุํ สุเมธํ,}

\prose{สพฺเพสุ ธมฺเมสุ อนูปลิตฺตํ.}

\prose{สพฺพญฺชหํ ตณฺหกฺขเย วิมุตฺตํ,}

\csromanheader{2}{ตํ วาปิ ธีรา มุนิ\footnote{มุนิํ (สี, อิ)} เวทยนฺติ. (5)}
\proseitem{214}{ปญฺญาพลํ สีลวตูปปนฺนํ,}

\prose{สมาหิตํ ฌานรตํ สตีมํ.}

\prose{สงฺคา ปมุตฺตํ อขิลํ อนาสวํ,}

\csromanheader{2}{ตํ วาปิ ธีรา มุนิ เวทยนฺติ. (6)}
\gambhira{สุตฺตนิปาตปาฬิ}
\proseitem{215}{\csromanfolio{310}เอกํ จรนฺตํ มุนิมปฺปมตฺตํ,}

\prose{นินฺทาปสํสาสุ อเวธมานํ.}

\csromangathameasure{สีหํว สทฺเทสุ อสนฺตสนฺตํ, \\ วาตํว ชาลมฺหิ อสชฺชมานํ. \\ ปทฺธํว โตเยน อลิปฺปมานํ. \\ เนตารมญฺเญส’มนญฺญเนยฺยํ.}
\csromangathagroup{\csromangathastanza{สีหํว สทฺเทสุ อสนฺตสนฺตํ, \\ วาตํว ชาลมฺหิ อสชฺชมานํ. \\ ปทฺธํว\footnote{ปทุมํว (สี, สฺยา, อิ)} โตเยน อลิปฺปมานํ\footnote{อลิมฺปมานํ (สฺยา, ก)}. \\ เนตารมญฺเญส’มนญฺญเนยฺยํ.}}

\csromanheader{2}{ตํ วาปิ ธีรา มุนิ เวทยนฺติ. (7)}
\proseitem{216}{โย โอคหเณ ถมฺโภริวาภิชายติ,}

\prose{ยสฺมิํ ปเร วาจา’ปริยนฺตํ\footnote{วาจํ ปริยนฺตํ (ก)} วทนฺติ.}

\prose{ตํ วีตราคํ สุสมาหิตินฺทฺริยํ,}

\csromanheader{2}{ตํ วาปิ ธีรา มุนิ เวทยนฺติ. (8)}
\proseitem{217}{โย เว ฐิตตฺโต ตสรํว อุชฺชุ,}

\prose{ชิคุจฺฉติ กมฺเมหิ ปาปเกหิ.}

\prose{วีมํสมาโน วิสมํ สมญฺจ,}

\csromanheader{2}{ตํ วาปิ ธีรา มุนิ เวทยนฺติ. (9)}
\proseitem{218}{โย สญฺญตตฺโต น กโรติ ปาปํ,}

\prose{ทหโร มชฺฌิโม จ มุนิ\footnote{ทหโร จ มชฺโฌ จ มุนี (สี, สฺยา, กํ, อิ)} ยตตฺโต.}

\prose{อโรสเนยฺโย น โส โรเสติ\footnote{น โรเสติ (สฺยา)} กญฺจิ.}

\csromanheader{2}{ตํ วาปิ ธีรา มุนิ เวทยนฺติ. (10)}
\proseitem{219}{ยทคฺคโต มชฺฌโต เสสโต วา,}

\prose{ปิณฺฑํ ลเภถ ปรทตฺตูปชีวี.}

\prose{นาลํ ถุตุํ โนปิ นิปจฺจวาที,}

\csromanheader{2}{ตํ วาปิ ธีรา มุนิ เวทยนฺติ. (11)}
\proseitem{220}{มุนิํ จรนฺตํ วิรตํ เมถุนสฺมา,}

\prose{โย โยพฺพเน โน’ปนิพชฺฌเต กฺวจิ.}

\prose{มทปฺปมาทา วิรตํ วิปฺปมุตฺตํ.}

\csromanheader{2}{ตํ วาปิ ธีรา มุนิ เวทยนฺติ. (12)}
\chapterheadpageread{1. อุรควคฺค}
\proseitem{221}{\csromanfolio{311}อญฺญาย โลกํ ปรมตฺถทสฺสิํ.}

\prose{โอฆํ สมุทฺทํ อติตริย ตาทิํ.}

\prose{ตํ ฉินฺนคนฺถํ อสิตํ อนาสวํ.}

\csromanheader{2}{ตํ วาปิ ธีรา มุนิ เวทยนฺติ. (13)}
\proseitem{222}{อสมา อุโภ ทูรวิหารวุตฺติโน,}

\csromangathameasure{คิหี ทารโปสี, อมโม จ สุพฺพโต.}
\csromangathagroup{\csromangathastanza{คิหี\footnote{คิหิ (ก)} ทารโปสี, อมโม จ สุพฺพโต.}}

\prose{ปรปาณโรธาย คิหี อสญฺญโต,}

\csromanheader{2}{นิจฺจํ มุนี รกฺขติ ปาณิเน\footnote{ปาณิโน (สี)} ยโต. (14)}
\proseitem{223}{สิขี ยถา นีลคีโว\footnote{นีลคิโว (สฺยา)} วิหงฺคโม,}

\prose{หํสสฺส โนเปติ ชวํ กุทาจนํ.}

\csromangathameasure{เอวํ คิหี นานุกโรติ ภิกฺขุโน, \\ มุนิโน วิวิตฺตสฺส วนมฺหิ ฌายโตติ. (15) มุนิสุตฺตํ ทฺวาทสมํ.}
\csromangathagroup{\csromangathastanza{เอวํ คิหี นานุกโรติ ภิกฺขุโน, \\ มุนิโน วิวิตฺตสฺส วนมฺหิ ฌายโตติ. (15) มุนิสุตฺตํ ทฺวาทสมํ.}}

\vspace{\tipitakaparskip}
\csromancenter{\textbf{อุรควคฺโค ปฐโม.}}
\titlehead{\textbf{ตสฺสุทฺทานํ}}
\csromangathameasure{อุรโค ธนิโย เจว, วิสาณํ จ ตถา กสิ. \\ จุนฺโท ปราภโว เจว, วสโล เมตฺตภาวนา. \\ สาตาคิโร อาฬวโก, วิชโย จ ตถา มุนิ. \\ ทฺวาทเสตานิ สุตฺตานิ, อุรควคฺโคติ วุจฺจตีติ.}
\csromangathagroup{\csromangathastanza{อุรโค ธนิโย เจว, วิสาณํ จ ตถา กสิ. \\ จุนฺโท ปราภโว เจว, วสโล เมตฺตภาวนา.}\csromangathastanza{สาตาคิโร อาฬวโก, วิชโย จ ตถา มุนิ. \\ ทฺวาทเสตานิ สุตฺตานิ, อุรควคฺโคติ วุจฺจตีติ.}}

\chapterheadpageread{2. จูฬวคฺค}
\csromanheader{2}{1. รตนสุตฺต}
\proseitem{224}{\csromanfolio{312}\csromansymbolfootnote{*}{เหฏฺฐา 4 ปิฏฺเฐปิ.}ยานีธ ภูตานิ สมาคตานิ,}

\prose{ภุมฺมานิ\footnote{ภูมานิ (ก)} วา ยานิ ว อนฺตลิกฺเข.}

\csromangathameasure{สพฺเพว ภูตา สุมนา ภวนฺตุ,}
\csromangathagroup{\csromangathastanza{สพฺเพว ภูตา สุมนา ภวนฺตุ,}}

\csromanheader{2}{อโถปิ สกฺกจฺจ สุณนฺตุ ภาสิตํ. (1)}
\csromangathameasure{ตสฺมา หิ ภูตา นิสาเมถ สพฺเพ, \\ เมตฺตํ กโรถ มานุสิยา ปชาย. \\ ทิวา จ รตฺโต จ หรนฺติ เย พลิํ,}
\csromangathagroup{\csromangathastanza{ตสฺมา หิ ภูตา นิสาเมถ สพฺเพ, \\ เมตฺตํ กโรถ มานุสิยา ปชาย. \\ ทิวา จ รตฺโต จ หรนฺติ เย พลิํ,}}

\csromanheader{2}{ตสฺมา หิ เน รกฺขถ อปฺปมตฺตา. (2)}
\csromangathameasure{ยํ กิญฺจิ วิตฺตํ อิธ วา หุรํ วา, \\ สคฺเคสุ วา ยํ รตนํ ปณีตํ. \\ น โน สมํ อตฺถิ ตถาคเตน, \\ อิทมฺปิ พุทฺเธ รตนํ ปณีตํ.}
\csromangathagroup{\csromangathastanza{ยํ กิญฺจิ วิตฺตํ อิธ วา หุรํ วา, \\ สคฺเคสุ วา ยํ รตนํ ปณีตํ. \\ น โน สมํ อตฺถิ ตถาคเตน, \\ อิทมฺปิ พุทฺเธ รตนํ ปณีตํ.}}

\csromanheader{2}{(3)}
\csromangathameasure{ขยํ วิราคํ อมตํ ปณีตํ, \\ ยทชฺฌคา สกฺยมุนี สมาหิโต. \\ น เตน ธมฺเมน สมตฺถิ กิญฺจิ, \\ อิทมฺปิ ธมฺเม รตนํ ปณีตํ.}
\csromangathagroup{\csromangathastanza{ขยํ วิราคํ อมตํ ปณีตํ, \\ ยทชฺฌคา สกฺยมุนี สมาหิโต. \\ น เตน ธมฺเมน สมตฺถิ กิญฺจิ, \\ อิทมฺปิ ธมฺเม รตนํ ปณีตํ.}}

\csromanheader{2}{(4)}
\csromangathameasure{ยํ พุทฺธเสฏฺโฐ ปริวณฺณยี สุจิํ, \\ สมาธิมานนฺตริกญฺญมาหุ. \\ สมาธินา เตน สโม น วิชฺชติ, \\ อิทมฺปิ ธมฺเม รตนํ ปณีตํ,}
\csromangathagroup{\csromangathastanza{ยํ พุทฺธเสฏฺโฐ ปริวณฺณยี สุจิํ, \\ สมาธิมานนฺตริกญฺญมาหุ. \\ สมาธินา เตน สโม น วิชฺชติ, \\ อิทมฺปิ ธมฺเม รตนํ ปณีตํ,}}

\csromanheader{2}{(5)}
\chapterheadpageread{2. จูฬวคฺค}
\proseitem{229}{\csromanfolio{313}เย ปุคฺคลา อฏฺฐ สตํ ปสตฺถา,}

\prose{จตฺตาริ เอตานิ ยุคานิ โหนฺติ.}

\csromangathameasure{เต ทกฺขิเณยฺยา สุคตสฺส สาวกา, \\ เอเตสุ ทินฺนานิ มหปฺผลานิ.}
\csromangathagroup{\csromangathastanza{เต ทกฺขิเณยฺยา สุคตสฺส สาวกา, \\ เอเตสุ ทินฺนานิ มหปฺผลานิ.}}

\csromanheader{2}{(6)}
\proseitem{230}{เย สุปฺปยุตฺตา มนสา ทเฬฺหน,}

\prose{นิกฺกามิโน โคตมสาสนมฺหิ.}

\prose{เต ปตฺติปตฺตา อมตํ วิคยฺห.}

\prose{ลทฺธา มุธา นิพฺพุติํ\footnote{นิพฺพุติ (ก)} ภุญฺชมานา.}

\csromanheader{2}{(7)}
\proseitem{231}{ยถินฺทขีโล ปถวิสฺสิโต\footnote{ปฐวิสฺสิโต (ก-สี), ปฐวิํ สิโต (ก-สี, สฺยา, กํ, อิ)} สิยา,}

\prose{จตุพฺภิ วาเตหิ อสมฺปกมฺปิโย.}

\csromangathameasure{ตถูปมํ สปฺปุริสํ วทามิ, \\ โย อริยสจฺจานิ อเวจฺจ ปสฺสติ.}
\csromangathagroup{\csromangathastanza{ตถูปมํ สปฺปุริสํ วทามิ, \\ โย อริยสจฺจานิ อเวจฺจ ปสฺสติ.}}

\csromanheader{2}{(8)}
\proseitem{232}{เย อริยสจฺจานิ วิภาวยนฺติ,}

\prose{คมฺภีรปญฺเญน สุเทสิตานิ.}

\csromangathameasure{กิญฺจาปิ เต โหนฺติ ภุสํ ปมตฺตา, \\ น เต ภวํ อฏฺฐมมาทิยนฺติ.}
\csromangathagroup{\csromangathastanza{กิญฺจาปิ เต โหนฺติ ภุสํ ปมตฺตา, \\ น เต ภวํ อฏฺฐมมาทิยนฺติ.}}

\csromanheader{2}{(9)}
\proseitem{233}{สหาว’สฺส ทสฺสนสมฺปทาย\footnote{สหาวสทฺทสฺสนสมฺปทาย (ก)},}

\prose{ตยสฺสุ ธมฺมา ชหิตา ภวนฺติ.}

\prose{สกฺกายทิฏฺฐี วิจิกิจฺฉิตญฺจ,}

\csromanheader{2}{สีลพฺพตํ วาปิ ยทตฺถิ กิญฺจิ. (10)}
\gambhira{สุตฺตนิปาตปาฬิ}
\proseitem{234}{\csromanfolio{314}จตูหปาเยหิ จ วิปฺปมุตฺโต,}

\prose{ฉจฺจาภิฐานานิ\footnote{ฉ จาภิฐานานิ (สี, สฺยา)} อภพฺพ กาตุํ\footnote{อภพฺโพ กาตุํ (สี)},}

\csromanheader{2}{(11)}
\proseitem{235}{กิญฺจาปิ โส กมฺม\footnote{กมฺมํ (สี, สฺยา, กํ, อิ)} กโรติ ปาปกํ,}

\prose{กาเยน วาจา อุท เจตสา วา.}

\csromangathameasure{อภพฺพ โส ตสฺส ปฏิจฺฉทาย, \\ อภพฺพตา ทิฏฺฐปทสฺส วุตฺตา.}
\csromangathagroup{\csromangathastanza{อภพฺพ\footnote{อภพฺโพ (พหูสุ)} โส ตสฺส ปฏิจฺฉทาย\footnote{ปฏิจฺฉาทาย (สี)}, \\ อภพฺพตา ทิฏฺฐปทสฺส วุตฺตา.}}

\csromanheader{2}{(12)}
\proseitem{236}{วนปฺปคุมฺเพ ยถ\footnote{ยถา (สี, สฺยา)} ผุสฺสิตคฺเค,}

\csromanmatikaanchor{mtk.244}
\csromantocmarkref{subsection}{10. อาฬวกสุตฺต}{mtk.244}
\bookmark[dest={mtk.244},level=2]{10. อาฬวกสุตฺต}
\prose{คิมฺหานมาเส ปฐมสฺมิํ\footnote{ปฐมสฺมิ (?)} คิเมฺห.}

\csromangathameasure{ตถูปมํ ธมฺมวรํ อเทสยิ, \\ นิพฺพานคามิํ ปรมํ หิตาย. \\ อิทมฺปิ พุทฺเธ รตนํ ปณีตํ,}
\csromangathagroup{\csromangathastanza{ตถูปมํ ธมฺมวรํ อเทสยิ\footnote{อเทสยี (สี)}, \\ นิพฺพานคามิํ ปรมํ หิตาย. \\ อิทมฺปิ พุทฺเธ รตนํ ปณีตํ,}}

\csromanheader{2}{(13)}
\proseitem{237}{วโร วรญฺญู วรโท วราหโร,}

\prose{อนุตฺตโร ธมฺมวรํ อเทสยิ.}

\prose{อิทมฺปิ พุทฺเธ รตนํ ปณีตํ,}

\csromanheader{2}{(14)}
\proseitem{238}{ขีณํ ปุราณํ นว นตฺถิสมฺภวํ,}

\prose{วิรตฺตจิตฺตายติเก ภวสฺมิํ.}

\csromangathameasure{เต ขีณพีชา อวิรูฬฺหิฉนฺทา, \\ นิพฺพนฺติ ธีรา ยถายํ ปทีโป.}
\csromangathagroup{\csromangathastanza{เต ขีณพีชา อวิรูฬฺหิฉนฺทา, \\ นิพฺพนฺติ ธีรา ยถายํ\footnote{ยถยํ (ก)} ปทีโป.}}

\csromanheader{2}{(15)}
\chapterheadpageread{2. จูฬวคฺค}
\proseitem{239}{\csromanfolio{315}ยานีธ ภูตานิ สมาคตานิ,}

\csromanheader{2}{พุทฺธํ นมสฺสาม สุวตฺถิ โหตุ. (16)}
\proseitem{240}{ยานีธ ภูตานิ สมาคตานิ,}

\csromanheader{2}{ธมฺมํ นมสฺสาม สุวตฺถิ โหตุ. (17)}
\proseitem{241}{ยานีธ ภูตานิ สมาคตานิ,}

\prose{สํฆํ นมสฺสาม สุวตฺถิ โหตุ. (18) รตนสุตฺตํ ปฐมํ.}

\csromanheader{2}{2. อามคนฺธสุตฺต}
\proseitem{242}{สามากจิงฺคูลกจีนกานิ จ,}

\prose{ปตฺตปฺผลํ มูลผลํ ควิปฺผลํ.}

\prose{ธมฺเมน ลทฺธํ สตมสฺนมานา\footnote{สตมสมานา (สี, อิ), สตมสฺสมานา (สฺยา, กํ)},}

\csromanheader{2}{น กามกามา อลิกํ ภณนฺติ. (1)}
\nitthitam{243. ยทสฺนมาโน สุกตํ สุนิฏฺฐิตํ,}
\prose{ปเรหิ ทินฺนํ ปยตํ ปณีตํ.}

\prose{สาลีนมนฺนํ ปริภุญฺชมาโน,}

\csromanheader{2}{โส ภุญฺชสี กสฺสป อามคนฺธํ. (2)}
\gambhira{สุตฺตนิปาตปาฬิ}
\proseitem{244}{\csromanfolio{316}“น อามคนฺโธ มม กปฺปตี”ติ,}

\prose{อิจฺเจว ตฺวํ ภาสสิ พฺรหฺมพนฺธุ.}

\csromangathameasure{สาลีนมนฺนํ ปริภุญฺชมาโน, \\ สกุนฺตมํเสหิ สุสงฺขเตหิ. \\ ปุจฺฉามิ ตํ กสฺสป เอตมตฺถํ,}
\csromangathagroup{\csromangathastanza{สาลีนมนฺนํ ปริภุญฺชมาโน, \\ สกุนฺตมํเสหิ สุสงฺขเตหิ. \\ ปุจฺฉามิ ตํ กสฺสป เอตมตฺถํ,}}

\csromanheader{2}{กถํปกาโร ตว อามคนฺโธ. (3)}
\proseitem{245}{ปาณาติปาโต วธเฉทพนฺธนํ,}

\prose{เถยฺยํ มุสาวาโท นิกติวญฺจนานิ จ.}

\prose{อชฺเฌนกุตฺตํ\footnote{อชฺเฌน กุชฺชํ (สี, อิ)} ปรทารเสวนา,}

\csromanheader{2}{เอสามคนฺโธ น หิ มํสโภชนํ. (4)}
\proseitem{246}{เย อิธ กาเมสุ อสญฺญตา ชนา,}

\prose{รเสสุ คิทฺธา อสุจิภาวมสฺสิตา\footnote{อสุจีกมิสฺสิตา (สี, สฺยา, กํ, อิ)}.}

\csromanmatikaanchor{mtk.245}
\csromantocmarkref{subsection}{11. วิชยสุตฺต}{mtk.245}
\bookmark[dest={mtk.245},level=2]{11. วิชยสุตฺต}
\prose{นตฺถิกทิฏฺฐี วิสมา ทุรนฺนยา,}

\csromanheader{2}{เอสามคนฺโธ น หิ มํสโภชนํ. (5)}
\proseitem{247}{เย ลูขสา ทารุณา ปิฏฺฐิมํสิกา\footnote{เย ลูขรสา ทารุณา ปรปิฏฺฐิมํสิกา (ก)},}

\prose{มิตฺตทฺทุโน นิกฺกรุณา’ติมานิโน.}

\prose{อทานสีลา น จ เทนฺติ กสฺสจิ,}

\csromanheader{2}{เอสามคนฺโธ น หิ มํสโภชนํ. (6)}
\proseitem{248}{โกโธ มโท ถมฺโภ ปจฺจุปฏฺฐาปนา\footnote{ปจฺจุฏฺฐาปนา จ (สี, สฺยา), ปจฺจุฏฺฐาปนา (อิ)},}

\prose{มายา อุสูยา ภสฺสสมุสฺสโย จ.}

\prose{มานาติมาโน จ อสพฺภิ สนฺถโว,}

\csromanheader{2}{เอสามคนฺโธ น หิ มํสโภชนํ. (7)}
\proseitem{249}{เย ปาปสีลา อิณฆาตสูจกา,}

\prose{โวหารกูฏา อิธ ปาฏิรูปิกา\footnote{ปาติรูปิกา (?)}.}

\prose{นราธมา เยธ กโรนฺติ กิพฺพิสํ,}

\csromanheader{2}{เอสามคนฺโธ น หิ มํสโภชนํ. (8)}
\chapterheadpageread{2. จูฬวคฺค}
\proseitem{250}{\csromanfolio{317}เย อิธ ปาเณสุ อสญฺญตา ชนา,}

\prose{ปเรส’มาทาย วิเหสมุยฺยุตา.}

\prose{ทุสฺสีลลุทฺทา ผรุสา อนาทรา,}

\csromanheader{2}{เอสามคนฺโธ น หิ มํสโภชนํ. (9)}
\proseitem{251}{เอเตสุ คิทฺธา วิรุทฺธาติปาติโน,}

\prose{นิจฺจุยฺยุตา เปจฺจ ตมํ วชนฺติ เย.}

\prose{ปตนฺติ สตฺตา นิรยํ อวํสิรา,}

\csromanheader{2}{เอสามคนฺโธ น หิ มํสโภชนํ. (10)}
\proseitem{252}{น มจฺฉมํสาน’มนาสกตฺตํ\footnote{น มจฺฉมํสํ น อนาสกตฺตํ (สี, ฏฺฐ - มูลปาโฐ), น มจฺฉมํสานานาสกตฺตํ (สฺยา, ก)},}

\prose{น นคฺคิยํ น มุณฺฑิยํ ชฏาชลฺลํ.}

\csromangathameasure{ขราชินานิ นาคฺคิหุตฺตสฺสุปเสวนา, \\ เย วาปิ โลเก อมรา พหู ตปา. \\ มนฺตาหุตี ยญฺญมุตูปเสวนา,}
\csromangathagroup{\csromangathastanza{ขราชินานิ นาคฺคิหุตฺตสฺสุปเสวนา, \\ เย วาปิ โลเก อมรา พหู ตปา. \\ มนฺตาหุตี ยญฺญมุตูปเสวนา,}}

\csromanheader{2}{โสเธนฺติ มจฺจํ อวิติณฺณกงฺขํ. (11)}
\proseitem{253}{โสเตสุ\footnote{โย เตสุ (ก)} คุตฺโต วิทิตินฺทฺริโย จเร,}

\prose{ธมฺเม ฐิโต อชฺชวมทฺทเว รโต.}

\prose{สงฺคาติโค สพฺพทุกฺขปฺปหีโน,}

\csromanmatikaanchor{mtk.246}
\csromantocmarkref{subsection}{12. มุนิสุตฺต}{mtk.246}
\bookmark[dest={mtk.246},level=2]{12. มุนิสุตฺต}
\csromanheader{2}{น ลิปฺปติ\footnote{น ลิมฺปติ (สฺยา, กํ, ก)} ทิฏฺฐสุเตสุ ธีโร. (12)}
\proseitem{254}{อิจฺเจตมตฺถํ ภควา ปุนปฺปุนํ,}

\prose{อกฺขาสิ นํ\footnote{ตํ (สี, อิ)} เวทยิ มนฺตปารคู.}

\prose{จิตฺราหิ คาถาหิ มุนี ปกาสยิ,}

\csromanheader{2}{นิรามคนฺโธ อสิโต ทุรนฺนโย. (13)}
\proseitem{255}{สุตฺวาน พุทฺธสฺส สุภาสิตํ ปทํ,}

\prose{นิรามคนฺธํ สพฺพทุกฺขปฺปนูทนํ.}

\csromangathameasure{นีจมโน วนฺทิ ตถาคตสฺส, \\ ตตฺเถว ปพฺพชฺชมโรจยิตฺถาติ. (14) อามคนฺธสุตฺตํ ทุติยํ.}
\csromangathagroup{\csromangathastanza{นีจมโน วนฺทิ ตถาคตสฺส, \\ ตตฺเถว ปพฺพชฺชมโรจยิตฺถาติ. (14) อามคนฺธสุตฺตํ ทุติยํ.}}

\gambhira{สุตฺตนิปาตปาฬิ}
\csromanheader{2}{3. หิริสุตฺต}
\proseitem{256}{\csromanfolio{318}หิริํ ตรนฺตํ วิชิคุจฺฉมานํ,}

\prose{ตวาหมสฺมิ\footnote{สขาหมสฺมิ (สี, สฺยา, กํ, อิ)} อิติ ภาสมานํ.}

\prose{สยฺหานิ กมฺมานิ อนาทิยนฺตํ,}

\csromanheader{2}{เนโส มมนฺติ อิติ นํ วิชญฺญา. (1)}
\proseitem{257}{อนนฺวยํ\footnote{อตฺถนฺวยํ (ก)} ปิยํ วาจํ, โย มิตฺเตสุ ปกุพฺพติ.}

\prose{อกโรนฺตํ ภาสมานํ, ปริชานนฺติ ปณฺฑิตา. (2)}

\proseitem{258}{น โส มิตฺโต โย สทา อปฺปมตฺโต,}

\prose{เภทาสงฺกี รนฺธเมวานุปสฺสี.}

\prose{ยสฺมิญฺจ เสติ อุรสีว ปุตฺโต,}

\csromanheader{2}{ส เว มิตฺโต โย ปเรหิ อเภชฺโช. (3)}
\proseitem{259}{ปามุชฺชกรณํ ฐานํ, ปสํสาวหนํ สุขํ.}

\prose{ผลานิสํโส ภาเวติ, วหนฺโต โปริสํ ธุรํ. (4)}

\proseitem{260}{ปวิเวกรสํ ปิตฺวา, รสํ อุปสมสฺส จ.}

\csromangathameasure{นิทฺทโร โหติ นิปฺปาโป, ธมฺมปีติรสํ ปิวนฺติ. (5) หิริสุตฺตํ ตติยํ.}
\csromangathagroup{\csromangathastanza{นิทฺทโร โหติ นิปฺปาโป, ธมฺมปีติรสํ ปิวนฺติ. (5) หิริสุตฺตํ ตติยํ.}}

\csromanheader{2}{4. มงฺคลสุตฺต}
\prose{\csromansymbolfootnote{*}{เหฏฺฐา 3 ปิฏฺเฐปิ.}เอวํ เม สุตํ\csromandash{}เอกํ สมยํ ภควา สาวตฺถิยํ วิหรติ เชตวเน อนาถปิณฺฑิกสฺส อาราเม.\csromanspacer{} อถ โข อญฺญตรา เทวตา อภิกฺกนฺตาย รตฺติยา อภิกฺกนฺตวณฺณา เกวลกปฺปํ เชตวนํ โอภาเสตฺวา เยน ภควา เตนุปสงฺกมิ, อุปสงฺกมิตฺวา ภควนฺตํ อภิวาเทตฺวา เอกมนฺตํ อฏฺฐาสิ, เอกมนฺตํ ฐิตา โข สา เทวตา ภควนฺตํ คาถาย อชฺฌภาสิ\csromandash{}}

\proseitem{261}{“พหู เทวา มนุสฺสา จ, มงฺคลานิ อจินฺตยุํ.}

\prose{อากงฺขมานา โสตฺถานํ, พฺรูหิ มงฺคลมุตฺตมํ. (1)}

\chapterheadpageread{2. จูฬวคฺค}
\proseitem{262}{\csromanfolio{319}อเสวนา จ พาลานํ, ปณฺฑิตานญฺจ เสวนา.}

\prose{ปูชา จ ปูชเนยฺยานํ\footnote{ปูชนียานํ (สี, สฺยา, กํ, อิ)}, เอตํ มงฺคลมุตฺตมํ. (2)}

\proseitem{263}{ปติรูปเทสวาโส จ, ปุพฺเพ จ กตปุญฺญตา.}

\prose{อตฺตสมฺมาปณิธิ\footnote{อตฺตสมฺมาปณีธี (กตฺถจิ)} จ, เอตํ มงฺคลมุตฺตมํ. (3)}

\proseitem{264}{พาหุสจฺจญฺจ สิปฺปญฺจ, วินโย จ สุสิกฺขิโต.}

\prose{สุภาสิตา จ ยา วาจา, เอตํ มงฺคลมุตฺตมํ. (4)}

\proseitem{265}{มาตาปิตุ อุปฏฺฐานํ, ปุตฺตทารสฺส สงฺคโห.}

\prose{อนากุลา จ กมฺมนฺตา, เอตํ มงฺคลมุตฺตมํ. (5)}

\proseitem{266}{ทานญฺจ ธมฺมจริยา จ, ญาตกานญฺจ สงฺคโห.}

\prose{อนวชฺชานิ กมฺมานิ, เอตํ มงฺคลมุตฺตมํ. (6)}

\proseitem{267}{อารตี วิรตี ปาปา, มชฺชปานา จ สํยโม.}

\prose{อปฺปมาโท จ ธมฺเมสุ, เอตํ มงฺคลมุตฺตมํ. (7)}

\proseitem{268}{คารโว จ นิวาโต จ, สนฺตุฏฺฐิ จ กตญฺญุตา.}

\prose{กาเลน ธมฺมสฺสวนํ\footnote{ธมฺมสฺสวณํ (กตฺถจิ), ธมฺมสวนํ (สี, ก)}, เอตํ มงฺคลมุตฺตมํ. (8)}

\proseitem{269}{ขนฺตี จ โสวจสฺสตา, สมณานญฺจ ทสฺสนํ.}

\prose{กาเลน ธมฺมสากจฺฉา, เอตํ มงฺคลมุตฺตมํ. (9)}

\proseitem{270}{ตโป จ พฺรหฺมจริยญฺจ, อริยสจฺจาน ทสฺสนํ.}

\prose{นิพฺพานสจฺฉิกิริยา จ, เอตํ มงฺคลมุตฺตมํ. (10)}

\proseitem{271}{ผุฏฺฐสฺส โลกธมฺเมหิ, จิตฺตํ ยสฺส น กมฺปติ.}

\prose{อโสกํ วิรชํ เขมํ, เอตํ มงฺคลมุตฺตมํ. (11)}

\proseitem{272}{เอตาทิสานิ กตฺวาน,}

\prose{สพฺพตฺถ มปราชิตา.}

\csromangathameasure{สพฺพตฺถ โสตฺถิํ คจฺฉนฺติ, \\ ตํ เตสํ มงฺคลมุตฺตมนฺ”ติ. (12) มงฺคลสุตฺตํ จตุตฺถํ.}
\csromangathagroup{\csromangathastanza{สพฺพตฺถ โสตฺถิํ คจฺฉนฺติ, \\ ตํ เตสํ มงฺคลมุตฺตมนฺ”ติ. (12) มงฺคลสุตฺตํ จตุตฺถํ.}}

\gambhira{สุตฺตนิปาตปาฬิ}
\csromanheader{2}{5. สูจิโลมสุตฺต}
\prose{\csromanfolio{320}\csromansymbolfootnote{+}{สํ 1. 208 ปิฏฺเฐปิ.}เอวํ เม สุตํ\csromandash{}เอกํ สมยํ ภควา คยายํ วิหรติ ฏงฺกิตมญฺเจ สูจิโลมสฺส ยกฺขสฺส ภวเน.\csromanspacer{} เตน โข ปน สมเยน ขโร จ ยกฺโข, สูจิโลโม จ ยกฺโข ภควโต อวิทูเร อติกฺกมนฺติ.\csromanspacer{} อถ โข ขโร ยกฺโข สูจิโลมํ ยกฺขํ เอตทโวจ “เอโส สมโณ”ติ.\csromanspacer{} เนโส สมโณ, สมณโก เอโส, ยาวาหํ ชานามิ\footnote{ยาว ชานามิ (สี, อิ)} “ยทิ วา โส สมโณ\footnote{ยทิ วา สมโณ (สฺยา)}, ยทิ วา โส สมณโก”ติ\footnote{ยทิ วา สมณโกติ (สี, สฺยา, อิ)}.}

\prose{อถ โข สูจิโลโม ยกฺโข เยน ภควา เตนุปสงฺกมิ, อุปสงฺกมิตฺวา ภควโต กายํ อุปนาเมสิ.\csromanspacer{} อถ โข ภควา กายํ อปนาเมสิ.\csromanspacer{} อถ โข สูจิโลโม ยกฺโข ภควนฺตํ เอตทโวจ “ภายสิ มํ สมณา”ติ.\csromanspacer{} น ขฺวาหํ ตํ อาวุโส ภายามิ, อปิ จ เต สมฺผสฺโส ปาปโกติ.}

\prose{ปญฺหํ ตํ สมณ ปุจฺฉิสฺสามิ, สเจ เม น พฺยากริสฺสสิ, จิตฺตํ วา เต ขิปิสฺสามิ, หทยํ วา เต ผาเลสฺสามิ, ปาเทสุ วา คเหตฺวา ปารคงฺคาย ขิปิสฺสามีติ.}

\prose{น ขฺวาหํ ตํ อาวุโส ปสฺสามิ สเทวเก โลเก สมารเก สพฺรหฺมเก สสฺสมณพฺราหฺมณิยา ปชาย สเทวมนุสฺสาย โย เม จิตฺตํ วา ขิเปยฺย, หทยํ วา ผาเลยฺย, ปาเทสุ วา คเหตฺวา ปารคงฺคาย ขิเปยฺย, อปิ จ ตฺวํ อาวุโส ปุจฺฉ ยทากงฺขสีติ.\csromanspacer{} อถ โข สูจิโลโม ยกฺโข ภควนฺตํ คาถาย อชฺฌภาสิ\csromandash{}}

\proseitem{273}{\csromansymbolfootnote{*}{ขุ 8. 252; ขุ 10.126 ปิฏฺเฐสุปิ}“ราโค จ โทโส จ กุโตนิทานา,}

\prose{อรตี รตี โลมหํโส กุโตชา.}

\csromangathameasure{กุโต สมุฏฺฐาย มโนวิตกฺกา,}
\csromangathagroup{\csromangathastanza{กุโต สมุฏฺฐาย มโนวิตกฺกา,}}

\csromanheader{2}{กุมารกา ธงฺกมิโวสฺสชนฺติ. (1)}
\proseitem{274}{ราโค จ โทโส จ อิโตนิทานา,}

\prose{อรตี รตี โลมหํโส อิโตชา.}

\csromangathameasure{อิโต สมุฏฺฐาย มโนวิตกฺกา,}
\csromangathagroup{\csromangathastanza{อิโต สมุฏฺฐาย มโนวิตกฺกา,}}

\csromanheader{2}{กุมารกา ธงฺกมิโวสฺสชนฺติ. (2)}
\chapterheadpageread{2. จูฬวคฺค}
\proseitem{274}{\csromanfolio{321}สฺเนหชา อตฺตสมฺภูตา, นิโคฺรธสฺเสว ขนฺธชา.}

\prose{ปุถู วิสตฺตา กาเมสุ, มาลุวาว วิตตาวเน. (3)}

\proseitem{275}{เย นํ ปชานนฺติ ยโตนิทานํ,}

\prose{เต นํ วิโนเทนฺติ สุโณหิ ยกฺข.}

\csromanmatikaanchor{mtk.247}
\csromanmatikaanchor{mtk.248}
\csromantocmarknopage{chapter}{2. จูฬวคฺค}{mtk.247}
\bookmark[dest={mtk.247},level=0]{2. จูฬวคฺค}
\csromantocmarkref{subsection}{1. รตนสุตฺต}{mtk.248}
\bookmark[dest={mtk.248},level=2]{1. รตนสุตฺต}
\csromangathameasure{เต ทุตฺตรํ โอฆมิมํ ตรนฺติ, \\ อติณฺณปุพฺพํ อปุนพฺภวายา”ติ. (4) สูจิโลมสุตฺตํ ปญฺจมํ.}
\csromangathagroup{\csromangathastanza{เต ทุตฺตรํ โอฆมิมํ ตรนฺติ, \\ อติณฺณปุพฺพํ อปุนพฺภวายา”ติ. (4) สูจิโลมสุตฺตํ ปญฺจมํ.}}

\csromanheader{2}{6. กปิลสุตฺต (ธมฺมจริยสุตฺต)}
\proseitem{276}{ธมฺมจริยํ พฺรหฺมจริยํ, เอตทาหุ วสุตฺตมํ.}

\prose{ปพฺพชิโตปิ เจ โหติ, อคารา อนคาริยํ. (1)}

\proseitem{277}{โส เจ มุขรชาติโก, วิเหสาภิรโต มโค.}

\prose{ชีวิตํ ตสฺส ปาปิโย, รชํ วฑฺเฒติ อตฺตโน. (2)}

\proseitem{278}{กลหาภิรโต ภิกฺขุ, โมหธมฺเมน อาวุโต.}

\prose{อกฺขาตมฺปิ น ชานาติ, ธมฺมํ พุทฺเธน เทสิตํ. (3)}

\proseitem{279}{วิเหสํ ภาวิตตฺตานํ, อวิชฺชาย ปุรกฺขโต.}

\prose{สํกิเลสํ น ชานาติ, มคฺคํ นิรยคามินํ. (4)}

\proseitem{280}{วินิปาตํ สมาปนฺโน, คพฺภา คพฺภํ ตมา ตมํ.}

\prose{ส เว ตาทิสโก ภิกฺขุ, เปจฺจ ทุกฺขํ นิคจฺฉติ. (5)}

\proseitem{281}{คูถกูโป ยถา อสฺส,}

\prose{สมฺปุณฺโณ คณวสฺสิโก.}

\prose{โย จ เอวรูโป อสฺส,}

\csromanheader{2}{ทุพฺพิโสโธ หิ สางฺคโณ. (6)}
\proseitem{282}{ยํ เอวรูปํ ชานาถ, ภิกฺขโว เคหนิสฺสิตํ.}

\prose{ปาปิจฺฉํ ปาปสงฺกปฺปํ, ปาป-อาจารโคจรํ. (7)}

\gambhira{สุตฺตนิปาตปาฬิ}
\csromanmatikaanchor{mtk.249}
\csromanmatikaanchor{mtk.254}
\csromantocmarkref{subsection}{2. อามคนฺธสุตฺต}{mtk.249}
\bookmark[dest={mtk.249},level=2]{2. อามคนฺธสุตฺต}
\proseitem{283}{\csromanfolio{322}สพฺเพ สมคฺคา หุตฺวาน, อภินิพฺพชฺชิยาถ\footnote{อภินิพฺพชฺชยาถ (สี, อิ, อํ 3. 18)} นํ.}

\prose{การณฺฑวํ\footnote{การณฺฑํว (สฺยา, ก) อํ 3. 17, 18 ปิฏฺเฐสุ.} นิทฺธมถ, กสมฺพุํ อปกสฺสถ\footnote{อวกสฺสถ (สี, สฺยา, ก)}. (8)}

\proseitem{284}{ตโต ปลาเป\footnote{ปลาเส (ก)} วาเหถ, อสฺสมเณ สมณมานิเน.}

\prose{นิทฺธมิตฺวาน ปาปิจฺเฉ, ปาป-อาจารโคจเร. (9)}

\proseitem{285}{สุทฺธา สุทฺเธหิ สํวาสํ, กปฺปยโวฺห ปติสฺสตา.}

\csromangathameasure{ตโต สมคฺคา นิปกา, ทุกฺขสฺสนฺตํ กริสฺสถาติ. (10) ธมฺมจริยสุตฺตํ ฉฏฺฐํ.}
\csromangathagroup{\csromangathastanza{ตโต สมคฺคา นิปกา, ทุกฺขสฺสนฺตํ กริสฺสถาติ. (10) ธมฺมจริยสุตฺตํ\footnote{กปิลสุตฺตํ (ฏฺฐ)} ฉฏฺฐํ.}}

\csromanheader{2}{7. พฺราหฺมณธมฺมิกสุตฺต}
\prose{เอวํ เม สุตํ\csromandash{}เอกํ สมยํ ภควา สาวตฺถิยํ วิหรติ เชตวเน อนาถปิณฺฑิกสฺส อาราเม.\csromanspacer{} อถ โข สมฺพหุลา โกสลกา พฺราหฺมณมหาสาลา ชิณฺณา วุฑฺฒา มหลฺลกา อทฺธคตา วโย-อนุปฺปตฺตา เยน ภควา เตนุปสงฺกมิํสุ, อุปสงฺกมิตฺวา ภควตา สทฺธิํ สมฺโมทิํสุ, สมฺโมทนียํ กถํ สารณียํ วีติสาเรตฺวา เอกมนฺตํ นิสีทิํสุ, เอกมนฺตํ นิสินฺนา โข เต พฺราหฺมณมหาสาลา ภควนฺตํ เอตทโวจุํ “สนฺทิสฺสนฺติ นุ โข โภ โคตม เอตรหิ พฺราหฺมณา โปราณานํ พฺราหฺมณานํ พฺราหฺมณธมฺเม”ติ.\csromanspacer{} น โข พฺราหฺมณา สนฺทิสฺสนฺติ เอตรหิ พฺราหฺมณา โปราณานํ พฺราหฺมณานํ พฺราหฺมณธมฺเมติ.\csromanspacer{} สาธุ โน ภวํ โคตโม โปราณานํ พฺราหฺมณานํ พฺราหฺมณธมฺมํ ภาสตุ, สเจ โภโต โคตมสฺส อครูติ.\csromanspacer{} เตน หิ พฺราหฺมณา สุณาถ สาธุกํ มนสิ กโรถ, ภาสิสฺสามีติ. “เอวํ โภ”ติ โข เต พฺราหฺมณมหาสาลา ภควโต ปจฺจสฺโสสุํ.\csromanspacer{} ภควา เอตทโวจ\csromandash{}}

\proseitem{286}{“อิสโย ปุพฺพกา อาสุํ, สญฺญตตฺตา ตปสฺสิโน.}

\prose{ปญฺจ กามคุเณ หิตฺวา, อตฺตทตฺถมจาริสุํ. (1)}

\proseitem{287}{น ปสู พฺราหฺมณานา’สุํ, น หิรญฺญํ น ธานิยํ.}

\prose{สชฺฌายธนธญฺญาสุํ, พฺรหฺมํ นิธิมปาลยุํ. (2)}

\chapterheadpageread{2. จูฬวคฺค}
\proseitem{288}{\csromanfolio{323}ยํ เนสํ ปกตํ อาสิ, ทฺวารภตฺตํ อุปฏฺฐิตํ.}

\prose{สทฺธาปกตเมสานํ, ทาตเว ตทมญฺญิสุํ. (3)}

\proseitem{289}{นานารตฺเตหิ วตฺเถหิ, สยเนหาวสเถหิ จ.}

\prose{ผีตา ชนปทา รฏฺฐา, เต นมสฺสิํสุ พฺราหฺมเณ. (4)}

\proseitem{290}{อวชฺฌา พฺราหฺมณา อาสุํ, อเชยฺยา ธมฺมรกฺขิตา.}

\prose{น เน โกจิ นิวาเรสิ, กุลทฺวาเรสุ สพฺพโส. (5)}

\proseitem{291}{อฏฺฐจตฺตาลีสํ วสฺสานิ, (โกมาร) พฺรหฺมจริยํ จริํสุ เต.}

\prose{วิชฺชาจรณปริเยฏฺฐิํ, อจรุํ พฺราหฺมณา ปุเร. (6)}

\proseitem{292}{น พฺราหฺมณา อญฺญมคมุํ, นปิ ภริยํ กิณิํสุ เต.}

\prose{สมฺปิเยเนว สํวาสํ, สงฺคนฺตฺวา สมโรจยุํ. (7)}

\csromanhangingbegin
\hangingheaditem{293}{อญฺญตฺร ตมฺหา สมยา, อุตุเวรมณิํ ปติ.}
\hangingbody{อนฺตรา เมถุนํ ธมฺมํ, นาสฺสุ คจฺฉนฺติ พฺราหฺมณา. (8)}
\csromanhangingend

\proseitem{294}{พฺรหฺมจริยญฺจ สีลญฺจ, อชฺชวํ มทฺทวํ ตปํ.}

\prose{โสรจฺจํ อวิหิํสญฺจ, ขนฺติญฺจาปิ อวณฺณยุํ. (9)}

\proseitem{295}{โย เนสํ ปรโม อาสิ, พฺรหฺมา ทฬฺหปรกฺกโม.}

\prose{ส วาปิ เมถุนํ ธมฺมํ, สุปินนฺเตปิ นาคมา. (10)}

\proseitem{296}{ตสฺส วตฺตมนุสิกฺขนฺตา, อิเธเก วิญฺญุชาติกา.}

\prose{พฺรหฺมจริยญฺจ สีลญฺจ, ขนฺติญฺจาปิ อวณฺณยุํ. (11)}

\proseitem{297}{ตณฺฑุลํ สยนํ วตฺถํ, สปฺปิเตลญฺจ ยาจิย.}

\prose{ธมฺเมน สโมธาเนตฺวา, ตโต ยญฺญมกปฺปยุํ. (12)}

\proseitem{298}{อุปฏฺฐิตสฺมิํ ยญฺญสฺมิํ, นาสฺสุ คาโว หนิํสุ เต.}

\csromangathameasure{ยถา มาตา ปิตา ภาตา, อญฺเญ วาปิ จ ญาตกา.}
\csromangathagroup{\csromangathastanza{ยถา มาตา ปิตา ภาตา, อญฺเญ วาปิ จ ญาตกา.}}

\prose{คาโว โน ปรมา มิตฺตา, ยาสุ ชายนฺติ โอสธา. (13)}

\proseitem{299}{อนฺนทา พลทา เจตา, วณฺณทา สุขทา ตถา\footnote{สุขทา จ ตา (ก)}.}

\prose{เอตมตฺถวสํ ญตฺวา, นาสฺสุ คาโว หนิํสุ เต. (14)}

\gambhira{สุตฺตนิปาตปาฬิ}
\proseitem{300}{\csromanfolio{324}สุขุมาลา มหากายา, วณฺณวนฺโต ยสสฺสิโน.}

\csromangathameasure{พฺราหฺมณา เสหิ ธมฺเมหิ, กิจฺจากิจฺเจสุ อุสฺสุกา.}
\csromangathagroup{\csromangathastanza{พฺราหฺมณา เสหิ ธมฺเมหิ, กิจฺจากิจฺเจสุ อุสฺสุกา.}}

\prose{ยาว โลเก อวตฺติํสุ, สุขเมธิตฺถ’ยํ ปชา. (15)}

\proseitem{301}{เตสํ อาสิ วิปลฺลาโส, ทิสฺวาน อณุโต อณุํ.}

\prose{ราชิโน จ วิยาการํ, นาริโย สมลงฺกตา. (16)}

\proseitem{302}{รเถ จาชญฺญสํยุตฺเต, สุกเต จิตฺตสิพฺพเน.}

\prose{นิเวสเน นิเวเส จ, วิภตฺเต ภาคโส มิเต. (17)}

\proseitem{303}{โคมณฺฑลปริพฺยูฬฺหํ, นารีวรคณายุตํ.}

\prose{อุฬารํ มานุสํ โภคํ, อภิชฺฌายิํสุ พฺราหฺมณา. (18)}

\proseitem{304}{เต ตตฺถ มนฺเต คนฺเถตฺวา, โอกฺกากํ ตทุปาคมุํ.}

\csromangathameasure{ปหูตธนธญฺโญสิ, ยชสฺสุ พหุ เต วิตฺตํ.}
\csromangathagroup{\csromangathastanza{ปหูตธนธญฺโญสิ, ยชสฺสุ พหุ เต วิตฺตํ.}}

\csromanheader{2}{ยชสฺสุ พหุ เต ธนํ. (19)}
\proseitem{305}{ตโต จ ราชา สญฺญตฺโต, พฺราหฺมเณหิ รเถสโภ.}

\csromangathameasure{อสฺสเมธํ ปุริสเมธํ, \\ สมฺมาปาสํ วาชเปยฺยํ นิรคฺคฬํ.}
\csromangathagroup{\csromangathastanza{อสฺสเมธํ ปุริสเมธํ, \\ สมฺมาปาสํ วาชเปยฺยํ นิรคฺคฬํ.}}

\prose{เอเต ยาเค ยชิตฺวาน, พฺราหฺมณาน’มทา ธนํ. (20)}

\proseitem{306}{คาโว สยนญฺจ วตฺถญฺจ, นาริโย สมลงฺกตา.}

\prose{รเถ จาชญฺญสํยุตฺเต, สุกเต จิตฺตสิพฺพเน. (21)}

\proseitem{307}{นิเวสนานิ รมฺมานิ, สุวิภตฺตานิ ภาคโส.}

\prose{นานาธญฺญสฺส ปูเรตฺวา, พฺราหฺมณาน’มทา ธนํ. (22)}

\csromanhangingbegin
\hangingheaditem{308}{เต จ ตตฺถ ธนํ ลทฺธา, สนฺนิธิํ สมโรจยุํ.}
\hangingbody{เตสํ อิจฺฉาวติณฺณานํ, ภิยฺโย ตณฺหา ปวฑฺฒถ.}
\hangingbody{เต ตตฺถ มนฺเต คนฺเถตฺวา, โอกฺกากํ ปุน มุปาคมุํ. (23)}
\csromanhangingend

\proseitem{309}{ยถา อาโป จ ปถวี จ, หิรญฺญํ ธนธานิยํ.}

\csromangathameasure{เอวํ คาโว มนุสฺสานํ, ปริกฺขาโร โส หิ ปาณินํ.}
\csromangathagroup{\csromangathastanza{เอวํ คาโว มนุสฺสานํ, ปริกฺขาโร โส หิ ปาณินํ.}}

\prose{ยชสฺสุ พหุ เต วิตฺตํ, ยชสฺสุ พหุ เต ธนํ. (24)}

\proseitem{310}{ตโต จ ราชา สญฺญตฺโต, พฺราหฺมเณหิ รเถสโภ.}

\prose{เนกา สตสหสฺสิโย, คาโว ยญฺเญ อฆาตยิ. (25)}

\chapterheadpageread{2. จูฬวคฺค}
\proseitem{311}{\csromanfolio{325}น ปาทา น วิสาเณน, นาสฺสุ หิํสนฺติ เกนจิ.}

\csromangathameasure{คาโว เอฬกสมานา, โสรตา กุมฺภทูหนา.}
\csromangathagroup{\csromangathastanza{คาโว เอฬกสมานา, โสรตา กุมฺภทูหนา.}}

\prose{ตา วิสาเณ คเหตฺวาน, ราชา สตฺเถน ฆาตยิ. (26)}

\proseitem{312}{ตโต เทวา ปิตโร จ\footnote{ตโต จ เทวา ปิตโร (สี, สฺยา)}, อินฺโท อสุรรกฺขสา.}

\prose{อธมฺโม อิติ ปกฺกนฺทุํ, ยํ สตฺถํ นิปตี คเว. (27)}

\proseitem{313}{ตโย โรคา ปุเร อาสุํ, อิจฺฉา อนสนํ ชรา.}

\prose{ปสูนญฺจ สมารมฺภา, อฏฺฐานวุติมาคมุํ. (28)}

\proseitem{314}{เอโส อธมฺโม ทณฺฑานํ, โอกฺกนฺโต ปุราโณ อหุ.}

\prose{อทูสิกาโย หญฺญนฺติ, ธมฺมา ธํสนฺติ\footnote{ธํเสนฺติ (สี, อิ)} ยาชกา. (29)}

\proseitem{315}{เอวเมโส อณุธมฺโม, โปราโณ วิญฺญุครหิโต.}

\prose{ยตฺถ เอทิสกํ ปสฺสติ, ยาชกํ ครหตี\footnote{ครหี (ก)} ชโน. (30)}

\proseitem{316}{เอวํ ธมฺเม วิยาปนฺเน, วิภินฺนา สุทฺทเวสฺสิกา.}

\prose{ปุถู วิภินฺนา ขตฺติยา, ปติํ ภริยา’วมญฺญถ. (31)}

\proseitem{317}{ขตฺติยา พฺราหฺมพนฺธู จ, เย จญฺเญ โคตฺตรกฺขิตา.}

\prose{ชาติวาทํ นิรากตฺวา\footnote{นิรํกตฺวา (พหูสุ) ยถา อนิรากตชฺฌาโนติ. ขุ 2. 179, 282 ปิฏฺเฐสุปิ.}, กามานํ วสมนฺวคุนฺ”ติ. (32)}

\prose{เอวํ วุตฺเต เต พฺราหฺมณมหาสาลา ภควนฺตํ เอตทโวจุํ “อภิกฺกนฺตํ โภ โคตม ฯเปฯ อุปาสเก โน ภวํ โคตโม ธาเรตุ อชฺชตคฺเค ปาณุเปเต สรณํ คเต”ติ.\csromanspacer{} พฺราหฺมณธมฺมิกสุตฺตํ สตฺตมํ.}

\csromanheader{2}{8. ธมฺม (นาวา) สุตฺต}
\proseitem{318}{ยสฺมา หิ ธมฺมํ ปุริโส วิชญฺญา,}

\prose{อินฺทํว นํ เทวตา ปูชเยยฺย.}

\prose{โส ปูชิโต ตสฺมิ ปสนฺนจิตฺโต,}

\csromanheader{2}{พหุสฺสุโต ปาตุกโรติ ธมฺมํ. (1)}
\gambhira{สุตฺตนิปาตปาฬิ}
\proseitem{319}{\csromanfolio{326}ตทฏฺฐิกตฺวาน นิสมฺม ธีโร,}

\prose{ธมฺมานุธมฺมํ ปฏิปชฺชมาโน.}

\prose{วิญฺญู วิภาวี นิปุโณ จ โหติ,}

\csromanheader{2}{โย ตาทิสํ ภชติ อปฺปมตฺโต. (2)}
\proseitem{320}{ขุทฺทญฺจ พาลํ อุปเสวมาโน,}

\prose{อนาคตตฺถญฺจ อุสูยกญฺจ.}

\prose{อิเธว ธมฺมํ อวิภาวยิตฺวา,}

\csromanheader{2}{อวิติณฺณกงฺโข มรณํ อุเปติ. (3)}
\proseitem{321}{ยถา นโร อาปค’โมตริตฺวา,}

\prose{มโหทกํ สลิลํ สีฆโสตํ.}

\prose{โส วุยฺหมาโน อนุโสตคามี,}

\csromanheader{2}{กิํ โส ปเร สกฺขติ ตารเยตุํ. (4)}
\proseitem{322}{ตเถว ธมฺมํ อวิภาวยิตฺวา,}

\prose{พหุสฺสุตานํ อนิสามย’ตฺถํ.}

\prose{สยํ อชานํ อวิติณฺณกงฺโข,}

\csromanheader{2}{กิํ โส ปเร สกฺขติ นิชฺฌเปตุํ. (5)}
\proseitem{323}{ยถาปิ นาวํ ทฬฺหมารุหิตฺวา,}

\prose{ผิเยน\footnote{ปิเยน (สี, สฺยา)} ริตฺเตน สมงฺคิภูโต.}

\prose{โส ตารเย ตตฺถ พหูปิ อญฺเญ,}

\csromanheader{2}{ตตฺรูปยญฺญู กุสโล มุตีมา\footnote{มตีมา (สฺยา, ก)}. (6)}
\proseitem{324}{เอวมฺปิ โย เวทคุ ภาวิตตฺโต,}

\prose{พหุสฺสุโต โหติ อเวธธมฺโม.}

\prose{โส โข ปเร นิชฺฌปเย ปชานํ,}

\csromanheader{2}{โสตาวธานูปนิสูปปนฺเน. (7)}
\proseitem{325}{ตสฺมา หเว สปฺปุริสํ ภเชถ,}

\prose{เมธาวินญฺเจว พหุสฺสุตญฺจ.}

\csromangathameasure{อญฺญาย อตฺถํ ปฏิปชฺชมาโน, \\ วิญฺญาตธมฺโม ส สุขํ ลเภถาติ. (8) นาวาสุตฺตํ อฏฺฐมํ.}
\csromangathagroup{\csromangathastanza{อญฺญาย อตฺถํ ปฏิปชฺชมาโน, \\ วิญฺญาตธมฺโม ส สุขํ\footnote{โส สุขํ (สี)} ลเภถาติ. (8) นาวาสุตฺตํ อฏฺฐมํ.}}

\csromanheader{2}{2. จูฬวคฺค \textbf{9. กิํสีลสุตฺต}}
\csromanmatikaanchor{mtk.250}
\csromantocmarkref{subsection}{3. หิริสุตฺต}{mtk.250}
\bookmark[dest={mtk.250},level=2]{3. หิริสุตฺต}
\proseitem{326}{\csromanfolio{327}กิํสีโล กิํสมาจาโร, กานิ กมฺมานิ พฺรูหยํ.}

\prose{นโร สมฺมา นิวิฏฺฐสฺส, อุตฺตมตฺถญฺจ ปาปุเณ. (1)}

\proseitem{327}{วุฑฺฒาปจายี อนุสูยโก สิยา,}

\prose{กาลญฺญู\footnote{กาลญฺญุ (สี, สฺยา)} จสฺส ครูนํ\footnote{ครุนํ (สี)} ทสฺสนาย.}

\prose{ธมฺมิํ กถํ เอรยิตํ ขณญฺญู,}

\csromanheader{2}{สุเณยฺย สกฺกจฺจ สุภาสิตานิ. (2)}
\proseitem{328}{กาเลน คจฺเฉ ครูนํ สกาสํ,}

\prose{ถมฺภํ นิรํกตฺวา\footnote{นิรากตฺวา (?) นิ + อา + กร + ตฺวา.} นิวาตวุตฺติ.}

\prose{อตฺถํ ธมฺมํ สํยมํ พฺรหฺมจริยํ,}

\csromanheader{2}{อนุสฺสเร เจว สมาจเร จ. (3)}
\proseitem{329}{ธมฺมาราโม ธมฺมรโต,}

\prose{ธมฺเม ฐิโต ธมฺมวินิจฺฉยญฺญู.}

\prose{เนวาจเร ธมฺมสนฺโทสวาทํ,}

\csromanheader{2}{ตจฺเฉหิ นีเยถ สุภาสิเตหิ. (4)}
\proseitem{330}{หสฺสํ ชปฺปํ ปริเทวํ ปโทสํ,}

\csromanmatikaanchor{mtk.251}
\csromantocmarkref{subsection}{4. มงฺคลสุตฺต}{mtk.251}
\bookmark[dest={mtk.251},level=2]{4. มงฺคลสุตฺต}
\prose{มายากตํ กุหนํ คิทฺธิ มานํ.}

\prose{สารมฺภํ กกฺกสํ กสาวญฺจ มุจฺฉํ\footnote{สารมฺภ กกฺกสฺส กสาว มุจฺฉํ (สฺยา, อิ)},}

\csromanheader{2}{หิตฺวา จเร วีตมโท ฐิตตฺโต. (5)}
\proseitem{331}{วิญฺญาตสารานิ สุภาสิตานิ,}

\prose{สุตญฺจ วิญฺญาตสมาธิสารํ.}

\prose{น ตสฺส ปญฺญา จ สุตญฺจ วฑฺฒติ,}

\csromanheader{2}{โย สาหโส โหติ นโร ปมตฺโต. (6)}
\gambhira{สุตฺตนิปาตปาฬิ}
\proseitem{332}{\csromanfolio{328}ธมฺเม จ เย อริยปเวทิเต รตา,}

\prose{อนุตฺตรา เต วจสา มนสา กมฺมุนา จ.}

\csromangathameasure{เต สนฺติโสรจฺจสมาธิสณฺฐิตา, \\ สุตสฺส ปญฺญาย จ สารมชฺฌคูติ. (7) กิํสีลสุตฺตํ นวมํ.}
\csromangathagroup{\csromangathastanza{เต สนฺติโสรจฺจสมาธิสณฺฐิตา, \\ สุตสฺส ปญฺญาย จ สารมชฺฌคูติ. (7) กิํสีลสุตฺตํ นวมํ.}}

\csromanheader{2}{10. อุฏฺฐานสุตฺต}
\proseitem{333}{อุฏฺฐหถ นิสีทถ, โก อตฺโถ สุปิเตน โว.}

\prose{อาตุรานญฺหิ กา นิทฺทา, สลฺลวิทฺธาน รุปฺปตํ. (1)}

\proseitem{334}{อุฏฺฐหถ นิสีทถ, ทฬฺหํ สิกฺขถ สนฺติยา.}

\prose{มา โว ปมตฺเต วิญฺญาย,}

\csromanheader{2}{มจฺจุราชา อโมหยิตฺถ วสานุเค. (2)}
\proseitem{335}{ยาย เทวา มนุสฺสา จ, สิตา ติฏฺฐนฺติ อตฺถิกา.}

\csromangathameasure{ตรเถตํ วิสตฺติกํ, ขโณ โว มา อุปจฺจคา.}
\csromangathagroup{\csromangathastanza{ตรเถตํ วิสตฺติกํ, ขโณ โว\footnote{ขโณ เว (อิ, ก)} มา อุปจฺจคา.}}

\prose{ขณาตีตา หิ โสจนฺติ, นิรยมฺหิ สมปฺปิตา. (3)}

\proseitem{336}{ปมาโท รโช ปมาโท, ปมาทานุปติโต รโช.}

\csromangathameasure{อปฺปมาเทน วิชฺชาย, อพฺพเห สลฺลมตฺตโนติ. (4) อุฏฺฐานสุตฺตํ ทสมํ.}
\csromangathagroup{\csromangathastanza{อปฺปมาเทน วิชฺชาย, อพฺพเห\footnote{อพฺพูเฬฺห (สฺยา, อิ), อพฺพุเห (ก-ฏฺฐ)} สลฺลมตฺตโนติ. (4) อุฏฺฐานสุตฺตํ ทสมํ.}}

\csromanheader{2}{11. ราหุลสุตฺต}
\proseitem{337}{กจฺจิ อภิณฺหสํวาสา, นาวชานาสิ ปณฺฑิตํ.}

\prose{อุกฺกาธาโร\footnote{โอกฺกาธาโร (สฺยา, ก)} มนุสฺสานํ, กจฺจิ อปจิโต ตยา\footnote{ตว (สี-ฏฺฐ)}. (1)}

\proseitem{338}{นาหํ อภิณฺหสํวาสา, อวชานามิ ปณฺฑิตํ.}

\prose{อุกฺกาธาโร มนุสฺสานํ, นิจฺจํ อปจิโต มยา. (2)}

\chapterheadpageread{2. จูฬวคฺค}
\csromanmatikaanchor{mtk.252}
\csromanmatikaanchor{mtk.259}
\csromantocmarkref{subsection}{5. สูจิโลมสุตฺต}{mtk.252}
\bookmark[dest={mtk.252},level=2]{5. สูจิโลมสุตฺต}
\proseitem{339}{\csromanfolio{329}ปญฺจ กามคุเณ หิตฺวา, ปิยรูเป มโนรเม.}

\prose{สทฺธาย ฆรา นิกฺขมฺม, ทุกฺขสฺสนฺตกโร ภว. (3)}

\proseitem{340}{มิตฺเต ภชสฺสุ กลฺยาเณ, ปนฺตญฺจ สยนาสนํ.}

\prose{วิวิตฺตํ อปฺปนิคฺโฆสํ, มตฺตญฺญู โหหิ โภชเน. (4)}

\proseitem{341}{จีวเร ปิณฺฑปาเต จ, ปจฺจเย สยนาสเน.}

\prose{เอเตสุ ตณฺหํ มากาสิ, มา โลกํ ปุนราคมิ. (5)}

\proseitem{342}{สํวุโต ปาติโมกฺขสฺมิํ, อินฺทฺริเยสุ จ ปญฺจสุ.}

\prose{สติ กายคตา ตฺยตฺถุ, นิพฺพิทาพหุโล ภว. (6)}

\proseitem{343}{นิมิตฺตํ ปริวชฺเชหิ, สุภํ ราคูปสญฺหิตํ.}

\prose{อสุภาย จิตฺตํ ภาเวหิ, เอกคฺคํ สุสมาหิตํ. (7)}

\proseitem{344}{อนิมิตฺตญฺจ ภาเวหิ, มานานุสย’มุชฺชห.}

\prose{ตโต มานาภิสมยา, อุปสนฺโต จริสฺสสีติ. (8)}

\prose{อิตฺถํ สุทํ ภควา อายสฺมนฺตํ ราหุลํ อิมาหิ คาถาหิ อภิณฺหํ โอวทตีติ.\csromanspacer{} ราหุลสุตฺตํ เอกาทสมํ.}

\csromanheader{2}{12. นิโคฺรธกปฺป (วงฺคีส) สุตฺต}
\prose{เอวํ เม สุตํ\csromandash{}เอกํ สมยํ ภควา อาฬวิยํ วิหรติ อคฺคาฬเว เจติเย.\csromanspacer{} เตน โข ปน สมเยน อายสฺมโต วงฺคีสสฺส อุปชฺฌาโย นิโคฺรธกปฺโป นาม เถโร อคฺคาฬเว เจติเย อจิรปรินิพฺพุโต โหติ.\csromanspacer{} อถ โข อายสฺมโต วงฺคีสสฺส รโหคตสฺส ปฏิสลฺลีนสฺส เอวํ เจตโส ปริวิตกฺโก อุทปาทิ “ปรินิพฺพุโต นุ โข เม อุปชฺฌาโย, อุทาหุ โน ปรินิพฺพุโต”ติ.\csromanspacer{} อถ โข อายสฺมา วงฺคีโส สายนฺหสมยํ ปฏิสลฺลานา วุฏฺฐิโต เยน ภควา เตนุปสงฺกมิ, อุปสงฺกมิตฺวา ภควนฺตํ อภิวาเทตฺวา เอกมนฺตํ นิสีทิ, เอกมนฺตํ นิสินฺโน โข อายสฺมา วงฺคีโส ภควนฺตํ เอตทโวจ “อิธ มยฺหํ ภนฺเต รโหคตสฺส ปฏิสลฺลีนสฺส เอวํ เจตโส ปริวิตกฺโก อุทปาทิ}

\gambhira{สุตฺตนิปาตปาฬิ}
\prose{\csromanfolio{330}‘ปรินิพฺพุโต นุ โข เม อุปชฺฌาโย, อุทาหุ โน ปรินิพฺพุโต’ติ”.\csromanspacer{} อถ โข อายสฺมา วงฺคีโส อุฏฺฐายาสนา เอกํสํ จีวรํ กตฺวา เยน ภควา เตนญฺชลิํ ปณาเมตฺวา ภควนฺตํ คาถาย อชฺฌภาสิ\csromandash{}}

\proseitem{345}{“ปุจฺฉามิ\footnote{ปุจฺฉาม (สี, สฺยา, อิ)} สตฺถารมโนมปญฺญํ,}

\prose{ทิฏฺเฐว ธมฺเม โย วิจิกิจฺฉานํ เฉตฺตา.}

\prose{อคฺคาฬเว กาลมกาสิ ภิกฺขุ,}

\csromanheader{2}{ญาโต ยสสฺสี อภินิพฺพุตตฺโต. (1)}
\proseitem{346}{นิโคฺรธกปฺโป อิติ ตสฺส นามํ,}

\csromanmatikaanchor{mtk.253}
\csromantocmarkref{subsection}{6. กปิลสุตฺต (ธมฺมจริยสุตฺต)}{mtk.253}
\bookmark[dest={mtk.253},level=2]{6. กปิลสุตฺต (ธมฺมจริยสุตฺต)}
\csromantocmarkref{subsection}{7. พฺรหฺมณธมฺมิกสุตฺต}{mtk.254}
\bookmark[dest={mtk.254},level=2]{7. พฺรหฺมณธมฺมิกสุตฺต}
\prose{ตยา กตํ ภควา พฺราหฺมณสฺส.}

\prose{โส ตํ นมสฺสํ อจริ มุตฺยเปกฺโข,}

\csromanheader{2}{อารทฺธวิริโย ทฬฺหธมฺมทสฺสี. (2)}
\proseitem{347}{ตํ สาวกํ สกฺย\footnote{สกฺก (สี, สฺยา, อิ)} มยมฺปิ สพฺเพ,}

\prose{อญฺญาตุมิจฺฉาม สมนฺตจกฺขุ.}

\prose{สมวฏฺฐิตา โน สวนาย โสตา,}

\csromanheader{2}{ตุวํ โน สตฺถา ตฺวมนุตฺตโรสิ. (3)}
\proseitem{348}{ฉินฺเทว โน วิจิกิจฺฉํ พฺรูหิ เมตํ,}

\prose{ปรินิพฺพุตํ เวทย ภูริปญฺญ.}

\prose{มชฺเฌว\footnote{มชฺเฌ จ (สฺยา, ก)} โน ภาส สมนฺตจกฺขุ,}

\csromanheader{2}{สกฺโกว เทวาน สหสฺสเนตฺโต. (4)}
\proseitem{349}{เย เกจิ คนฺถา อิธ โมหมคฺคา,}

\prose{อญฺญาณปกฺขา วิจิกิจฺฉฐานา.}

\prose{ตถาคตํ ปตฺวา น เต ภวนฺติ,}

\csromanheader{2}{จกฺขุํ หิ เอตํ ปรมํ นรานํ. (5)}
\proseitem{350}{โน เจ หิ ชาตุ ปุริโส กิเลเส,}

\prose{วาโต ยถา อพฺภฆนํ วิหาเน.}

\prose{ตโมวสฺส นิวุโต สพฺพโลโก,}

\csromanheader{2}{น โชติมนฺโตปิ นรา ตเปยฺยุํ. (6)}
\chapterheadpageread{2. จูฬวคฺค}
\proseitem{351}{\csromanfolio{331}ธีรา จ ปชฺโชตกรา ภวนฺติ,}

\prose{ตํ ตํ อหํ วีร\footnote{ธีร (สี, สฺยา)} ตเถว มญฺเญ.}

\prose{วิปสฺสินํ ชานมุปาคมุมฺหา\footnote{ชานมุปคมมฺหา (สี, สฺยา)},}

\csromanheader{2}{ปริสาสุ โน อาวิกโรหิ กปฺปํ. (7)}
\proseitem{352}{ขิปฺปํ คิรํ เอรย วคฺคุ วคฺคุํ,}

\prose{หํโสว ปคฺคยฺห สณิกํ\footnote{สณิํ (สฺยา, อิ)} นิกูช.}

\prose{พินฺทุสฺสเรน สุวิกปฺปิเตน,}

\csromanheader{2}{สพฺเพว เต อุชฺชุคตา สุโณม. (8)}
\proseitem{353}{ปหีนชาติมรณํ อเสสํ,}

\prose{นิคฺคยฺห โธนํ\footnote{โธตํ (สี)} วเทสฺสามิ ธมฺมํ.}

\prose{น กามกาโร หิ ปุถุชฺชนานํ,}

\csromanheader{2}{สงฺเขยฺยกาโร จ\footnote{สงฺเขยฺยกาโรว (ก)} ตถาคตานํ. (9)}
\proseitem{354}{สมฺปนฺนเวยฺยากรณํ ตเวทํ,}

\prose{สมุชฺชุปญฺญสฺส\footnote{สมุชฺชปญฺญสฺส (สฺยา, ก)} สมุคฺคหีตํ.}

\prose{อยมญฺชลี ปจฺฉิโม สุปฺปณามิโต,}

\csromanheader{2}{มา โมหยี ชาน’มโนมปญฺญ. (10)}
\proseitem{355}{ปโรวรํ\footnote{วราวรํ (กตฺถจิ)} อริยธมฺมํ วิทิตฺวา,}

\prose{มา โมหยี ชาน’มโนมวีร.}

\prose{วาริํยถา ฆมฺมนิ ฆมฺมตตฺโต,}

\csromanheader{2}{วาจาภิกงฺขามิ สุตํ ปวสฺส\footnote{สุตสฺส วสฺส (สฺยา)}. (11)}
\proseitem{356}{ยทตฺถิกํ\footnote{ยทตฺถิยํ (อิ), ยทตฺถิตํ (ก)} พฺรหฺมจริยํ อจรี,}

\prose{กปฺปายโน กจฺจิสฺส ตํ อโมฆํ.}

\prose{นิพฺพายิ โส อาทุ ส-อุปาทิเสโส,}

\csromanheader{2}{ยถา วิมุตฺโต อหุ ตํ สุโณม. (12)}
\gambhira{สุตฺตนิปาตปาฬิ}
\proseitem{357}{\csromanfolio{332}อจฺเฉจฺฉิ\footnote{อเฉชฺชิ (ก)} ตณฺหํ อิธ นามรูเป, (อิติ ภควา,)}

\prose{กณฺหสฺส\footnote{ตณฺหาย (ก)} โสตํ ทีฆรตฺตานุสยิตํ.}

\prose{อตาริ ชาติํ มรณํ อเสสํ,}

\csromanheader{2}{อิจฺจพฺรวี ภควา ปญฺจเสฏฺโฐ. (13)}
\csromanhangingbegin
\hangingheaditem{358}{เอส สุตฺวา ปสีทามิ, วโจ เต อิสิสตฺตม.}
\hangingbody{อโมฆํ กิร เม ปุฏฺฐํ, น มํ วญฺเจสิ พฺราหฺมโณ. (14)}
\csromanhangingend

\proseitem{359}{ยถา วาที ตถา การี, อหุ พุทฺธสฺส สาวโก.}

\prose{อจฺฉิทา มจฺจุโน ชาลํ, ตตํ มายาวิโน ทฬฺหํ. (15)}

\proseitem{360}{อทฺทสา ภควา อาทิํ, อุปาทานสฺส กปฺปิโย.}

\csromangathameasure{อจฺจคา วต กปฺปายโน, มจฺจุเธยฺยํ สุทุตฺตรนฺ”ติ. (16) นิโคฺรธกปฺปสุตฺตํ ทฺวาทสมํ.}
\csromangathagroup{\csromangathastanza{อจฺจคา วต กปฺปายโน, มจฺจุเธยฺยํ สุทุตฺตรนฺ”ติ. (16) นิโคฺรธกปฺปสุตฺตํ ทฺวาทสมํ.}}

\csromanheader{2}{13. สมฺมาปริพฺพาชนียสุตฺต}
\proseitem{361}{ปุจฺฉามิ มุนิํ ปหูตปญฺญํ,}

\prose{ติณฺณํ ปารงฺคตํ ปรินิพฺพุตํ ฐิตตฺตํ.}

\prose{นิกฺขมฺม ฆรา ปนุชฺช กาเม,}

\csromanheader{2}{กถํ ภิกฺขุ สมฺมา โส โลเก ปริพฺพเชยฺย. (1)}
\proseitem{362}{ยสฺส มงฺคลา สมูหตา, (อิติ ภควา,)}

\prose{อุปฺปาตา สุปินา จ ลกฺขณา จ.}

\prose{โส มงฺคลโทสวิปฺปหีโน,}

\csromanheader{2}{สมฺมา โส โลเก ปริพฺพเชยฺย. (2)}
\proseitem{363}{ราคํ วินเยถ มานุเสสุ,}

\prose{ทิพฺเพสุ กาเมสุ จาปิ ภิกฺขุ.}

\prose{อติกฺกมฺม ภวํ สเมจฺจ ธมฺมํ,}

\csromanheader{2}{สมฺมา โส โลเก ปริพฺพเชยฺย. (3)}
\chapterheadpageread{2. จูฬวคฺค}
\proseitem{364}{\csromanfolio{333}วิปิฏฺฐิกตฺวาน เปสุณานิ,}

\prose{โกธํ กทรียํ ชเหยฺย ภิกฺขุ.}

\prose{อนุโรธวิโรธวิปฺปหีโน,}

\csromanheader{2}{สมฺมา โส โลเก ปริพฺพเชยฺย. (4)}
\proseitem{365}{หิตฺวาน ปิยญฺจ อปฺปิยญฺจ,}

\prose{อนุปาทาย อนิสฺสิโต กุหิญฺจิ.}

\prose{สํโยชนิเยหิ วิปฺปมุตฺโต,}

\csromanheader{2}{สมฺมา โส โลเก ปริพฺพเชยฺย. (5)}
\proseitem{366}{น โส อุปธีสุ สารเมติ,}

\prose{อาทาเนสุ วิเนยฺย ฉนฺทราคํ.}

\prose{โส อนิสฺสิโต อนญฺญเนยฺโย,}

\csromanheader{2}{สมฺมา โส โลเก ปริพฺพเชยฺย. (6)}
\proseitem{367}{วจสา มนสา จ กมฺมุนา จ,}

\prose{อวิรุทฺโธ สมฺมา วิทิตฺวา ธมฺมํ.}

\prose{นิพฺพานปทาภิปตฺถยาโน,}

\csromanheader{2}{สมฺมา โส โลเก ปริพฺพเชยฺย. (7)}
\proseitem{368}{โย วนฺทติ มนฺติ นุนฺนเมยฺย\footnote{นุณฺณเมยฺย (พหูสุ)},}

\prose{อกฺกุฏฺโฐปิ น สนฺธิเยถ ภิกฺขุ.}

\prose{ลทฺธา ปรโภชนํ น มชฺเช,}

\csromanheader{2}{สมฺมา โส โลเก ปริพฺพเชยฺย. (8)}
\proseitem{369}{โลภญฺจ ภวญฺจ วิปฺปหาย,}

\prose{วิรโต เฉทนพนฺธนา จ\footnote{เฉทนพนฺธนโต (สี, สฺยา)} ภิกฺขุ.}

\prose{โส ติณฺณกถํกโถ วิสลฺโล,}

\csromanheader{2}{สมฺมา โส โลเก ปริพฺพเชยฺย. (9)}
\proseitem{370}{สารุปฺปํ อตฺตโน วิทิตฺวา,}

\prose{โน จ ภิกฺขุ หิํเสยฺย กญฺจิ โลเก.}

\prose{ยถาตถิยํ วิทิตฺวา ธมฺมํ,}

\csromanheader{2}{สมฺมา โส โลเก ปริพฺพเชยฺย. (10)}
\gambhira{สุตฺตนิปาตปาฬิ}
\proseitem{371}{\csromanfolio{334}ยสฺสานุสยา น สนฺติ เกจิ,}

\prose{มูลา จ\footnote{มูลา (สี, สฺยา)} อกุสลา สมูหตาเส.}

\csromanmatikaanchor{mtk.255}
\csromantocmarkref{subsection}{8. ธมฺม (นาวา) สุตฺต}{mtk.255}
\bookmark[dest={mtk.255},level=2]{8. ธมฺม (นาวา) สุตฺต}
\prose{โส นิราโส\footnote{นิราสโย (สี), นิราสโส (สฺยา)} อนาสิสาโน\footnote{อนาสยาโน (สี, อิ), อนาสสาโน (สฺยา)},}

\csromanheader{2}{สมฺมา โส โลเก ปริพฺพเชยฺย. (11)}
\proseitem{372}{อาสวขีโณ ปหีนมาโน,}

\prose{สพฺพํ ราคปถํ อุปาติวตฺโต.}

\prose{ทนฺโต ปรินิพฺพุโต ฐิตตฺโต,}

\csromanheader{2}{สมฺมา โส โลเก ปริพฺพเชยฺย. (12)}
\proseitem{373}{สทฺโธ สุตวา นิยามทสฺสี,}

\prose{วคฺคคเตสุ น วคฺคสาริ ธีโร.}

\prose{โลภํ โทสํ วิเนยฺย ปฏิฆํ,}

\csromanheader{2}{สมฺมา โส โลเก ปริพฺพเชยฺย. (13)}
\proseitem{374}{สํสุทฺธชิโน วิวฏฺฏจฺฉโท,}

\prose{ธมฺเมสุ วสี ปารคู อเนโช.}

\prose{สงฺขารนิโรธญาณกุสโล,}

\csromanheader{2}{สมฺมา โส โลเก ปริพฺพเชยฺย. (14)}
\proseitem{375}{อตีเตสุ อนาคเตสุ จาปิ,}

\prose{กปฺปาตีโต อติจฺจสุทฺธิปญฺโญ.}

\prose{สพฺพายตเนหิ วิปฺปมุตฺโต,}

\csromanheader{2}{สมฺมา โส โลเก ปริพฺพเชยฺย. (15)}
\proseitem{376}{อญฺญาย ปทํ สเมจฺจ ธมฺมํ,}

\prose{วิวฏํ ทิสฺวาน ปหานมาสวานํ.}

\prose{สพฺพุปธีนํ ปริกฺขยาโน\footnote{ปริกฺขยา (อิ)},}

\csromanheader{2}{สมฺมา โส โลเก ปริพฺพเชยฺย. (16)}
\chapterheadpageread{2. จูฬวคฺค}
\proseitem{377}{\csromanfolio{335}อทฺธา หิ ภควา ตเถว เอตํ,}

\prose{โย โส เอวํวิหารี ทนฺโต ภิกฺขุ.}

\csromangathameasure{สพฺพสํโยชนโยควีติวตฺโต, \\ สมฺมา โส โลเก ปริพฺพเชยฺยาติ. (17) สมฺมาปริพฺพาชนียสุตฺตํ เตรสมํ.}
\csromangathagroup{\csromangathastanza{สพฺพสํโยชนโยควีติวตฺโต\footnote{สพฺพสํโยชนิเย จ วีติวตฺโต (สี, สฺยา, อิ)}, \\ สมฺมา โส โลเก ปริพฺพเชยฺยาติ. (17) สมฺมาปริพฺพาชนียสุตฺตํ เตรสมํ.}}

\csromanheader{2}{14. ธมฺมิกสุตฺต}
\prose{เอวํ เม สุตํ\csromandash{}เอกํ สมยํ ภควา สาวตฺถิยํ วิหรติ เชตวเน อนาถปิณฺฑิกสฺส อาราเม.\csromanspacer{} อถ โข ธมฺมิโก อุปาสโก ปญฺจหิ อุปาสกสเตหิ สทฺธิํ เยน ภควา เตนุปสงฺกมิ, อุปสงฺกมิตฺวา ภควนฺตํ อภิวาเทตฺวา เอกมนฺตํ นิสีทิ, เอกมนฺตํ นิสินฺโน โข ธมฺมิโก อุปาสโก ภควนฺตํ คาถาหิ อชฺฌภาสิ\csromandash{}}

\proseitem{378}{“ปุจฺฉามิ ตํ โคตม ภูริปญฺญ,}

\prose{กถํกโร สาวโก สาธุ โหติ.}

\prose{โย วา อคารา อนคารเมติ,}

\csromanheader{2}{อคาริโน วา ปนุปาสกาเส. (1)}
\proseitem{379}{ตุวญฺหิ โลกสฺส สเทวกสฺส,}

\prose{คติํ ปชานาสิ ปรายณญฺจ.}

\csromanmatikaanchor{mtk.256}
\csromantocmarkref{subsection}{9. กิํสีลสุตฺต}{mtk.256}
\bookmark[dest={mtk.256},level=2]{9. กิํสีลสุตฺต}
\prose{น จตฺถิ ตุโลฺย นิปุณตฺถทสฺสี,}

\csromanheader{2}{ตุวญฺหิ พุทฺธํ ปวรํ วทนฺติ. (2)}
\proseitem{380}{สพฺพํ ตุวํ ญาณมเวจฺจ ธมฺมํ,}

\prose{ปกาเสสิ สตฺเต อนุกมฺปมาโน.}

\prose{วิวฏฺฏจฺฉโทสิ สมนฺตจกฺขุ,}

\csromanheader{2}{วิโรจสิ วิมโล สพฺพโลเก. (3)}
\proseitem{381}{อาคญฺฉิ เต สนฺติเก นาคราชา,}

\csromangathameasure{เอราวโณ นาม, “ชิโน”ติ สุตฺวา.}
\csromangathagroup{\csromangathastanza{เอราวโณ นาม, “ชิโน”ติ สุตฺวา.}}

\prose{โสปิ ตยา มนฺตยิตฺวาชฺฌคมา,}

\prose{“สาธู”ติ สุตฺวาน ปตีตรูโป. (4)}

\gambhira{สุตฺตนิปาตปาฬิ}
\proseitem{382}{\csromanfolio{336}ราชาปิ ตํ เวสฺสวโณ กุเวโร,}

\prose{อุเปติ ธมฺมํ ปริปุจฺฉมาโน.}

\prose{ตสฺสาปิ ตฺวํ ปุจฺฉิโต พฺรูสิ ธีร,}

\csromanheader{2}{โส จาปิ สุตฺวาน ปตีตรูโป. (5)}
\proseitem{383}{เย เกจิเม ติตฺถิยา วาทสีลา,}

\prose{อาชีวกา วา ยทิ วา นิคณฺฐา.}

\prose{ปญฺญาย ตํ นาติตรนฺติ สพฺเพ,}

\csromanheader{2}{ฐิโต วชนฺตํ วิย สีฆคามิํ. (6)}
\proseitem{384}{เย เกจิเม พฺราหฺมณา วาทสีลา,}

\prose{วุทฺธา จาปิ พฺราหฺมณา สนฺติ เกจิ.}

\prose{สพฺเพ ตยิ อตฺถพทฺธา ภวนฺติ,}

\csromanheader{2}{เย จาปิ อญฺเญ วาทิโน มญฺญมานา. (7)}
\proseitem{385}{อยญฺหิ ธมฺโม นิปุโณ สุโข จ,}

\prose{โยยํ ตยา ภควา สุปฺปวุตฺโต.}

\prose{ตเมว สพฺเพปิ\footnote{สพฺเพ มยํ (สฺยา)} สุสฺสูสมานา,}

\csromanheader{2}{ตํ โน วท ปุจฺฉิโต พุทฺธเสฏฺฐ. (8)}
\proseitem{386}{สพฺเพปิเม ภิกฺขโว สนฺนิสินฺนา,}

\csromanmatikaanchor{mtk.257}
\csromantocmarkref{subsection}{10. อุฏฺฐานสุตฺต}{mtk.257}
\bookmark[dest={mtk.257},level=2]{10. อุฏฺฐานสุตฺต}
\prose{อุปาสกา จาปิ ตเถว โสตุํ.}

\prose{สุณนฺตุ ธมฺมํ วิมเลนานุพุทฺธํ,}

\csromanheader{2}{สุภาสิตํ วาสวสฺเสว เทวา. (9)}
\proseitem{387}{สุณาถ เม ภิกฺขโว สาวยามิ โว,}

\prose{ธมฺมํ ธุตํ ตญฺจ จราถ สพฺเพ.}

\prose{อิริยาปถํ ปพฺพชิตานุโลมิกํ,}

\csromanheader{2}{เสเวถ นํ อตฺถทโส มุตีมา. (10)}
\proseitem{388}{โน เว วิกาเล วิจเรยฺย ภิกฺขุ,}

\prose{คาเม จ ปิณฺฑาย จเรยฺย กาเล.}

\prose{อกาลจาริํ หิ สชนฺติ สงฺคา,}

\csromanheader{2}{ตสฺมา วิกาเล น จรนฺติ พุทฺธา. (11)}
\csromanmatikaanchor{mtk.258}
\csromantocmarkref{subsection}{11. ราหุลสุตฺต}{mtk.258}
\bookmark[dest={mtk.258},level=2]{11. ราหุลสุตฺต}
\csromantocmarkref{subsection}{12. นิโคฺรธกปฺปสุตฺต}{mtk.259}
\bookmark[dest={mtk.259},level=2]{12. นิโคฺรธกปฺปสุตฺต}
\chapterheadpageread{2. จูฬวคฺค}
\proseitem{389}{\csromanfolio{337}รูปา จ สทฺทา จ รสา จ คนฺธา,}

\prose{ผสฺสา จ เย สมฺมทยนฺติ สตฺเต.}

\prose{เอเตสุ ธมฺเมสุ วิเนยฺย ฉนฺทํ,}

\csromanheader{2}{กาเลน โส ปวิเส ปาตราสํ. (12)}
\proseitem{390}{ปิณฺฑญฺจ ภิกฺขุ สมเยน ลทฺธา,}

\prose{เอโก ปฏิกฺกมฺม รโห นิสีเท.}

\prose{อชฺฌตฺตจินฺตี น มโน พหิทฺธา,}

\csromanheader{2}{นิจฺฉารเย สงฺคหิตตฺตภาโว. (13)}
\proseitem{391}{สเจปิ โส สลฺลเป สาวเกน,}

\prose{อญฺเญน วา เกนจิ ภิกฺขุนา วา.}

\prose{ธมฺมํ ปณีตํ ต’มุทาหเรยฺย,}

\csromanheader{2}{น เปสุณํ โนปิ ปรูปวาทํ. (14)}
\proseitem{392}{วาทญฺหิ เอเก ปฏิเสนิยนฺติ,}

\prose{น เต ปสํสาม ปริตฺตปญฺเญ.}

\prose{ตโต ตโต เน ปสชนฺติ สงฺคา,}

\csromanheader{2}{จิตฺตญฺหิ เต ตตฺถ คเมนฺติ ทูเร. (15)}
\proseitem{393}{ปิณฺฑํ วิหารํ สยนาสนญฺจ,}

\prose{อาปญฺจ สงฺฆาฏิรชูปวาหนํ.}

\prose{สุตฺวาน ธมฺมํ สุคเตน เทสิตํ,}

\csromanheader{2}{สงฺขาย เสเว วรปญฺญสาวโก. (16)}
\proseitem{394}{ตสฺมา หิ ปิณฺเฑ สยนาสเน จ,}

\prose{อาเป จ สงฺฆาฏิรชูปวาหเน.}

\prose{เอเตสุ ธมฺเมสุ อนูปลิตฺโต,}

\csromanheader{2}{ภิกฺขุ ยถา โปกฺขเร วาริพินฺทุ. (17)}
\proseitem{395}{คหฏฺฐวตฺตํ ปน โว วทามิ,}

\prose{ยถากโร สาวโก สาธุ โหติ.}

\prose{น เหส\footnote{น เหโส (สี)} ลพฺภา สปริคฺคเหน,}

\csromanheader{2}{ผสฺเสตุํ โย เกวโล ภิกฺขุธมฺโม. (18)}
\gambhira{สุตฺตนิปาตปาฬิ}
\proseitem{396}{\csromanfolio{338}ปาณํ น หเน\footnote{น หาเน (สี)} น จ ฆาตเยยฺย,}

\prose{น จานุชญฺญา หนตํ ปเรสํ.}

\prose{สพฺเพสุ ภูเตสุ นิธาย ทณฺฑํ,}

\csromanheader{2}{เย ถาวรา เย จ ตสา สนฺติ\footnote{ตสนฺติ (สี, อิ)} โลเก. (19)}
\proseitem{397}{ตโต อทินฺนํ ปริวชฺชเยยฺย,}

\prose{กิญฺจิ กฺวจิ สาวโก พุชฺฌมาโน.}

\prose{น หารเย หรตํ นานุชญฺญา,}

\csromanheader{2}{สพฺพํ อทินฺนํ ปริวชฺชเยยฺย. (20)}
\proseitem{398}{อพฺรหฺมจริยํ ปริวชฺชเยยฺย,}

\prose{องฺคารกาสุํ ชลิตํว วิญฺญู.}

\prose{อสมฺภุณนฺโต ปน พฺรหฺมจริยํ,}

\csromanheader{2}{ปรสฺส ทารํ น อติกฺกเมยฺย. (21)}
\proseitem{399}{สภคฺคโต วา ปริสคฺคโต วา,}

\prose{เอกสฺส เวโก\footnote{เจโก (สี, สฺยา)} น มุสา ภเณยฺย.}

\prose{น ภาณเย ภณตํ นานุชญฺญา,}

\csromanheader{2}{สพฺพํ อภูตํ ปริวชฺชเยยฺย. (22)}
\proseitem{400}{มชฺชญฺจ ปานํ น สมาจเรยฺย,}

\prose{ธมฺมํ อิมํ โรจเย โย คหฏฺโฐ.}

\prose{น ปายเย ปิวตํ นานุชญฺญา,}

\csromanheader{2}{อุมฺมาทนนฺตํ อิติ นํ วิทิตฺวา. (23)}
\proseitem{401}{มทา หิ ปาปานิ กโรนฺติ พาลา,}

\prose{กาเรนฺติ จญฺเญปิ ชเน ปมตฺเต.}

\prose{เอตํ อปุญฺญายตนํ วิวชฺชเย,}

\csromanheader{2}{อุมฺมาทนํ โมหนํ พาลกนฺตํ. (24)}
\proseitem{402}{ปาณํ น หเน น จาทินฺนมาทิเย,}

\prose{มุสา น ภาเส น จ มชฺชโป สิยา.}

\prose{อพฺรหฺมจริยา วิรเมยฺย เมถุนา,}

\csromanheader{2}{รตฺติํ น ภุญฺเชยฺย วิกาลโภชนํ. (25)}
\chapterheadpageread{2. จูฬวคฺค}
\proseitem{403}{\csromanfolio{339}มาลํ น ธาเร น จ คนฺธมาจเร,}

\prose{มญฺเจ ฉมายํ ว สเยถ สนฺถเต.}

\prose{เอตํ หิ อฏฺฐงฺคิกมาหุโปสถํ,}

\csromanheader{2}{พุทฺเธน ทุกฺขนฺตคุนา ปกาสิตํ. (26)}
\proseitem{404}{ตโต จ ปกฺขสฺสุปวสฺสุโปสถํ,}

\prose{จาตุทฺทสิํ ปญฺจทสิญฺจ อฏฺฐมิํ.}

\prose{ปาฏิหาริยปกฺขญฺจ ปสนฺนมานโส,}

\csromanheader{2}{อฏฺฐงฺคุเปตํ สุสมตฺตรูปํ. (27)}
\proseitem{405}{ตโต จ ปาโต อุปวุตฺถุโปสโถ,}

\prose{อนฺเนน ปาเนน จ ภิกฺขุสํฆํ.}

\prose{ปสนฺนจิตฺโต อนุโมทมาโน,}

\csromanheader{2}{ยถารหํ สํวิภเชถ วิญฺญู. (28)}
\proseitem{406}{ธมฺเมน มาตาปิตโร ภเรยฺย,}

\prose{ปโยชเย ธมฺมิกํ โส วณิชฺชํ.}

\csromangathameasure{เอตํ คิหี วตฺตย’มปฺปมตฺโต, \\ สยมฺปเภ นาม อุเปติ เทเว”ติ. (29) ธมฺมิกสุตฺตํ จุทฺทสมํ.}
\csromangathagroup{\csromangathastanza{เอตํ คิหี วตฺตย’มปฺปมตฺโต, \\ สยมฺปเภ นาม อุเปติ เทเว”ติ. (29) ธมฺมิกสุตฺตํ จุทฺทสมํ.}}

\vspace{\tipitakaparskip}
\csromancenter{\textbf{จูฬวคฺโค ทุติโย.}}
\titlehead{\textbf{ตสฺสุทฺทานํ}}
\csromangathameasure{รตนามคนฺโธ หิริ จ, มงฺคลํ สูจิโลเมน. \\ ธมฺมจริยญฺจ พฺราหฺมโณ, นาวา กิํสีลมุฏฺฐานํ. \\ ราหุโล ปุน กปฺโป จ, ปริพฺพาชนิยํ ตถา. \\ ธมฺมิกญฺจ วิทุโน อาหุ, จูฬวคฺคนฺติ จุทฺทสาติ.}
\csromangathagroup{\csromangathastanza{รตนามคนฺโธ หิริ จ, มงฺคลํ สูจิโลเมน. \\ ธมฺมจริยญฺจ พฺราหฺมโณ\footnote{กปิโล พฺราหฺมโณปิ จ (สฺยา, ก)}, นาวา กิํสีลมุฏฺฐานํ.}\csromangathastanza{ราหุโล ปุน กปฺโป จ, ปริพฺพาชนิยํ ตถา. \\ ธมฺมิกญฺจ วิทุโน อาหุ, จูฬวคฺคนฺติ จุทฺทสาติ.}}

\chapterheadpageread{3. มหาวคฺค}
\csromanheader{2}{1. ปพฺพชฺชาสุตฺต}
\proseitem{407}{\csromanfolio{340}ปพฺพชฺชํ กิตฺตยิสฺสามิ, ยถา ปพฺพชิ จกฺขุมา.}

\prose{ยถา วีมํสมาโน โส, ปพฺพชฺชํ สมโรจยิ. (1)}

\csromanmatikaanchor{mtk.260}
\csromantocmarkref{subsection}{13. สมฺมาปริพฺพาชนียสุตฺต}{mtk.260}
\bookmark[dest={mtk.260},level=2]{13. สมฺมาปริพฺพาชนียสุตฺต}
\proseitem{408}{สมฺพาโธยํ ฆราวาโส, รชสฺสายตนํ อิติ.}

\prose{อพฺโภกาโสว ปพฺพชฺชา, อิติ ทิสฺวาน ปพฺพชิ. (2)}

\proseitem{409}{ปพฺพชิตฺวาน กาเยน, ปาปกมฺมํ วิวชฺชยิ.}

\prose{วจีทุจฺจริตํ หิตฺวา, อาชีวํ ปริโสธยิ. (3)}

\proseitem{410}{อคมา ราชคหํ พุทฺโธ, มคธานํ คิริพฺพชํ.}

\prose{ปิณฺฑาย อภิหาเรสิ, อากิณฺณวรลกฺขโณ. (4)}

\proseitem{411}{ตมทฺทสา พิมฺพิสาโร, ปาสาทสฺมิํ ปติฏฺฐิโต.}

\prose{ทิสฺวา ลกฺขณสมฺปนฺนํ, อิมมตฺถํ อภาสถ. (5)}

\proseitem{412}{อิมํ โภนฺโต นิสาเมถ, อภิรูโป พฺรหา สุจิ.}

\prose{จรเณน จ สมฺปนฺโน, ยุคมตฺตญฺจ เปกฺขติ. (6)}

\proseitem{413}{โอกฺขิตฺตจกฺขุ สติมา, นายํ นีจกุลามิว.}

\prose{ราชทูตาภิธาวนฺตุ, กุหิํ ภิกฺขุ คมิสฺสติ. (7)}

\proseitem{414}{เต เปสิตา ราชทูตา, ปิฏฺฐิโต อนุพนฺธิสุํ.}

\prose{กุหิํ คมิสฺสติ ภิกฺขุ, กตฺถ วาโส ภวิสฺสติ. (8)}

\proseitem{415}{สปทานํ จรมาโน, คุตฺตทฺวาโร สุสํวุโต.}

\prose{ขิปฺปํ ปตฺตํ อปูเรสิ, สมฺปชาโน ปฏิสฺสโต. (9)}

\proseitem{416}{ปิณฺฑจารํ จริตฺวาน, นิกฺขมฺม นครา มุนิ.}

\prose{ปณฺฑวํ อภิหาเรสิ, เอตฺถ วาโส ภวิสฺสติ. (10)}

\proseitem{417}{ทิสฺวาน วาสูปคตํ, ตโย\footnote{ตโต (สี, อิ)} ทูตา อุปาวิสุํ.}

\prose{เตสุ เอโกว\footnote{เอโก จ ทูโต (สี, สฺยา, อิ)} อาคนฺตฺวา, ราชิโน ปฏิเวทยิ. (11)}

\chapterheadpageread{3. มหาวคฺค}
\proseitem{418}{\csromanfolio{341}เอส ภิกฺขุ มหาราช, ปณฺฑวสฺส ปุรตฺถโต\footnote{ปุรกฺขโต (สฺยา, ก)}.}

\prose{นิสินฺโน พฺยคฺฆุสโภว, สีโหว คิริคพฺภเร. (12)}

\proseitem{419}{สุตฺวาน ทูตวจนํ, ภทฺทยาเนน ขตฺติโย.}

\prose{ตรมานรูโป นิยฺยาสิ, เยน ปณฺฑวปพฺพโต. (13)}

\proseitem{420}{ส ยานภูมิํ ยายิตฺวา, ยานา โอรุยฺห ขตฺติโย.}

\prose{ปตฺติโก อุปสงฺกมฺม, อาสชฺช นํ อุปาวิสิ. (14)}

\proseitem{421}{นิสชฺช ราชา สมฺโมทิ, กถํ สารณียํ ตโต.}

\prose{กถํ โส วีติสาเรตฺวา, อิมมตฺถํ อภาสถ. (15)}

\proseitem{422}{ยุวา จ ทหโร จาสิ, ปฐมุปฺปตฺติโก\footnote{ปฐมุปฺปตฺติยา (สี), ปฐมุปฺปตฺติโต (สฺยา)} สุสุ.}

\prose{วณฺณาโรเหน สมฺปนฺโน, ชาติมา วิย ขตฺติโย. (16)}

\proseitem{423}{โสภยนฺโต อนีกคฺคํ, นาคสํฆปุรกฺขโต.}

\prose{ททามิ โภเค ภุญฺชสฺสุ, ชาติํ อกฺขาหิ ปุจฺฉิโต. (17)}

\proseitem{424}{อุชุํ ชนปโท ราช, หิมวนฺตสฺส ปสฺสโต.}

\prose{ธนวิริเยน สมฺปนฺโน, โกสเลสุ\footnote{โกสลสฺส (สฺยา, ก)} นิเกติโน. (18)}

\proseitem{425}{อาทิจฺจา\footnote{อาทิจฺโจ (ก)} นาม โคตฺเตน, สากิยา\footnote{สากิโย (ก)} นาม ชาติยา.}

\prose{ตมฺหา กุลา ปพฺพชิโตมฺหิ, น กาเม อภิปตฺถยํ. (19)}

\proseitem{426}{กาเมสฺวาทีนวํ ทิสฺวา, เนกฺขมฺมํ ทฏฺฐุ เขมโต.}

\csromangathameasure{ปธานาย คมิสฺสามิ, เอตฺถ เม รญฺชตี มโนติ. (20) ปพฺพชฺชาสุตฺตํ ปฐมํ.}
\csromangathagroup{\csromangathastanza{ปธานาย คมิสฺสามิ, เอตฺถ เม รญฺชตี มโนติ. (20) ปพฺพชฺชาสุตฺตํ ปฐมํ.}}

\csromanheader{2}{2. ปธานสุตฺต}
\proseitem{427}{ตํ มํ ปธานปหิตตฺตํ, นทิํ เนรญฺชรํ ปติ.}

\prose{วิปรกฺกมฺม ฌายนฺตํ, โยคกฺเขมสฺส ปตฺติยา. (1)}

\proseitem{428}{นมุจี กรุณํ วาจํ, ภาสมาโน อุปาคมิ.}

\prose{กิโส ตฺวมสิ ทุพฺพณฺโณ, สนฺติเก มรณํ ตว. (2)}

\gambhira{สุตฺตนิปาตปาฬิ}
\csromangathameasure{สหสฺสภาโค มรณสฺส, เอกํโส ตว ชีวิตํ.}
\csromangathagroup{\csromangathastanza{\csromanfolio{342}สหสฺสภาโค มรณสฺส, เอกํโส ตว ชีวิตํ.}}

\vspace{\tipitakaparskip}
\csromancenter{ชีว โภ ชีวิตํ เสยฺโย, ชีวํ ปุญฺญานิ กาหสิ. (3)}
\vspace{0.5\baselineskip}
\csromangathameasure{จรโต จ เต พฺรหฺมจริยํ, อคฺคิหุตฺตญฺจ ชูหโต.}
\csromangathagroup{\csromangathastanza{จรโต จ เต พฺรหฺมจริยํ, อคฺคิหุตฺตญฺจ ชูหโต.}}

\vspace{\tipitakaparskip}
\csromancenter{ปหูตํ จียเต ปุญฺญํ, กิํ ปธาเนน กาหสิ. (4)}
\vspace{0.5\baselineskip}
\csromangathameasure{ทุคฺโค มคฺโค ปธานาย, ทุกฺกโร ทุรภิสมฺภโว.}
\csromangathagroup{\csromangathastanza{ทุคฺโค มคฺโค ปธานาย, ทุกฺกโร ทุรภิสมฺภโว.}}

\csromanheader{2}{อิมา คาถา ภณํ มาโร, อฏฺฐา พุทฺธสฺส สนฺติเก. (5)}
\csromangathameasure{ตํ ตถาวาทินํ มารํ, ภควา เอตทพฺรวิ.}
\csromangathagroup{\csromangathastanza{ตํ ตถาวาทินํ มารํ, ภควา เอตทพฺรวิ.}}

\vspace{\tipitakaparskip}
\csromancenter{ปมตฺตพนฺธุ ปาปิม, เยนตฺเถน\footnote{เสนตฺเถน (?), อตฺตโน อตฺเถน (ฏฺฐ - สํวณฺณนา)} อิธาคโต. (6)}
\vspace{0.5\baselineskip}
\csromangathameasure{อณุมตฺโตปิ ปุญฺเญน, อตฺโถ มยฺหํ น วิชฺชติ.}
\csromangathagroup{\csromangathastanza{อณุมตฺโตปิ\footnote{อณุมตฺเตนปิ (สี, สฺยา)} ปุญฺเญน, อตฺโถ มยฺหํ น วิชฺชติ.}}

\proseitem{433}{เยสญฺจ อตฺโถ ปุญฺเญน, เต มาโร วตฺตุมรหติ. (7)}

\csromangathameasure{อตฺถิ สทฺธา ตถา วิริยํ, ปญฺญา จ มม วิชฺชติ.}
\csromangathagroup{\csromangathastanza{อตฺถิ สทฺธา ตถา\footnote{ตโต (สี, อิ), ตโป (สฺยา, ก)} วิริยํ, ปญฺญา จ มม วิชฺชติ.}}

\vspace{\tipitakaparskip}
\csromancenter{เอวํ มํ ปหิตตฺตมฺปิ, กิํ ชีวมนุปุจฺฉสิ. (8)}
\vspace{0.5\baselineskip}
\csromangathameasure{นทีนมปิ โสตานิ, อยํ วาโต วิโสสเย.}
\csromangathagroup{\csromangathastanza{นทีนมปิ โสตานิ, อยํ วาโต วิโสสเย.}}

\vspace{\tipitakaparskip}
\csromancenter{กิญฺจ เม ปหิตตฺตสฺส, โลหิตํ นุปสุสฺสเย. (9)}
\vspace{0.5\baselineskip}
\csromangathameasure{โลหิเต สุสฺสมานมฺหิ, ปิตฺตํ เสมฺหญฺจ สุสฺสติ. \\ มํเสสุ ขียมาเนสุ, ภิยฺโย จิตฺตํ ปสีทติ.}
\csromangathagroup{\csromangathastanza{โลหิเต สุสฺสมานมฺหิ, ปิตฺตํ เสมฺหญฺจ สุสฺสติ. \\ มํเสสุ ขียมาเนสุ, ภิยฺโย จิตฺตํ ปสีทติ.}}

\proseitem{436}{ภิยฺโย สติ จ ปญฺญา จ, สมาธิ มม ติฏฺฐติ. (10)}

\csromangathameasure{ตสฺส เมวํ วิหรโต, ปตฺตสฺสุตฺตมเวทนํ.}
\csromangathagroup{\csromangathastanza{ตสฺส เมวํ วิหรโต, ปตฺตสฺสุตฺตมเวทนํ.}}

\proseitem{437}{กาเมสุ\footnote{กาเม (สี, สฺยา)} นาเปกฺขเต จิตฺตํ, ปสฺส สตฺตสฺส สุทฺธตํ. (11)}

\proseitem{438}{\csromansymbolfootnote{*}{ขุ 7. 73; ขุ 8. 114 ปิฏฺเฐสุปิ.}กามา เต ปฐมา เสนา, ทุติยา อรติ วุจฺจติ.}

\prose{ตติยา ขุปฺปิปาสา เต, จตุตฺถี ตณฺหา ปวุจฺจติ. (12)}

\csromangathameasure{ปญฺจมํ ถินมิทฺธํ เต, ฉฏฺฐา ภีรู ปวุจฺจติ.}
\csromangathagroup{\csromangathastanza{ปญฺจมํ\footnote{ปญฺจมี (สี, อิ)} ถินมิทฺธํ เต, ฉฏฺฐา ภีรู ปวุจฺจติ.}}

\vspace{\tipitakaparskip}
\csromancenter{สตฺตมี วิจิกิจฺฉา เต, มกฺโข ถมฺโภ เต อฏฺฐโม. (13)}
\vspace{0.5\baselineskip}
\csromangathameasure{ลาโภ สิโลโก สกฺกาโร, \\ มิจฺฉาลทฺโธ จ โย ยโส.}
\csromangathagroup{\csromangathastanza{ลาโภ สิโลโก สกฺกาโร, \\ มิจฺฉาลทฺโธ จ โย ยโส.}}

\vspace{\tipitakaparskip}
\csromancenter{โย จตฺตานํ สมุกฺกํเส, ปเร จ อวชานติ. (14)}
\chapterheadpageread{3. มหาวคฺค}
\proseitem{441}{\csromanfolio{343}เอสา นมุจิ เต เสนา, กณฺหสฺสาภิปฺปหารินี.}

\prose{น นํ อสูโร ชินาติ, เชตฺวา จ ลภเต สุขํ. (15)}

\csromanmatikaanchor{mtk.261}
\csromantocmarkref{subsection}{14. ธมฺมิกสุตฺต}{mtk.261}
\bookmark[dest={mtk.261},level=2]{14. ธมฺมิกสุตฺต}
\proseitem{442}{เอส มุญฺชํ ปริหเร, ธิรตฺถุ มม\footnote{อิธ (ก)} ชีวิตํ.}

\prose{สงฺคาเม เม มตํ เสยฺโย, ยญฺเจ ชีเว ปราชิโต. (16)}

\proseitem{443}{ปคาเฬฺหตฺถ น ทิสฺสนฺติ, เอเก สมณพฺราหฺมณา.}

\prose{ตญฺจ มคฺคํ น ชานนฺติ, เยน คจฺฉนฺติ สุพฺพตา. (17)}

\proseitem{444}{สมนฺตา ธชินิํ ทิสฺวา, ยุตฺตํ มารํ สวาหนํ.}

\prose{ยุทฺธาย ปจฺจุคฺคจฺฉามิ, มา มํ ฐานา อจาวยิ. (18)}

\proseitem{445}{ยํ เต ตํ นปฺปสหติ, เสนํ โลโก สเทวโก.}

\prose{ตํ เต ปญฺญาย เภจฺฉามิ\footnote{คจฺฉามิ (สี), เวจฺฉามิ (สฺยา), วชฺฌามิ (ก)}, อามํ ปตฺตํว อสฺมนา\footnote{ปกฺกํว อมฺหนา (ก)}. (19)}

\proseitem{446}{วสีกริตฺวา\footnote{วสิํ กริตฺวา (พหูสุ)} สงฺกปฺปํ, สติญฺจ สูปติฏฺฐิตํ.}

\prose{รฏฺฐา รฏฺฐํ วิจริสฺสํ, สาวเก วินยํ ปุถู. (20)}

\proseitem{447}{เต อปฺปมตฺตา ปหิตตฺตา, มม สาสนการกา.}

\prose{อกามสฺส\footnote{อกามา (ก)} เต คมิสฺสนฺติ, ยตฺถ คนฺตฺวา น โสจเร. (21)}

\proseitem{448}{สตฺต วสฺสานิ ภควนฺตํ, อนุพนฺธิํ ปทาปทํ.}

\prose{โอตารํ นาธิคจฺฉิสฺสํ, สมฺพุทฺธสฺส สตีมโต. (22)}

\proseitem{449}{เมทวณฺณํว ปาสาณํ, วายโส อนุปริยคา.}

\prose{อเปตฺถ มุทุํ\footnote{มุทุ (สี)} วินฺเทม, อปิ อสฺสาทนา สิยา. (23)}

\proseitem{450}{อลทฺธา ตตฺถ อสฺสาทํ, วายเส’ตฺโต อปกฺกมิ.}

\prose{กาโกว เสลมาสชฺช, นิพฺพิชฺชา’เปม โคตมํ. (24)}

\proseitem{451}{ตสฺส โสกปเรตสฺส, วีณา กจฺฉา อภสฺสถ.}

\csromangathameasure{ตโต โส ทุมฺมโน ยกฺโข, ตตฺเถวนฺตรธายถาติ. (25) ปธานสุตฺตํ ทุติยํ.}
\csromangathagroup{\csromangathastanza{ตโต โส ทุมฺมโน ยกฺโข, ตตฺเถวนฺตรธายถาติ. (25) ปธานสุตฺตํ ทุติยํ.}}

\gambhira{สุตฺตนิปาตปาฬิ}
\csromanheader{2}{3. สุภาสิตสุตฺต}
\prose{\csromanfolio{344}เอวํ เม สุตํ\csromandash{}เอกํ สมยํ ภควา สาวตฺถิยํ วิหรติ เชตวเน อนาถปิณฺฑิกสฺส อาราเม.\csromanspacer{} ตตฺร โข ภควา ภิกฺขู อามนฺเตสิ “ภิกฺขโว”ติ. “ภทนฺเต”ติ เต ภิกฺขู ภควโต ปจฺจสฺโสสุํ.\csromanspacer{} ภควา เอตทโวจ\csromandash{}}

\prose{จตูหิ ภิกฺขเว องฺเคหิ สมนฺนาคตา วาจา สุภาสิตา โหติ น ทุพฺภาสิตา อนวชฺชา จ อนนุวชฺชา จ วิญฺญูนํ.\csromanspacer{} กตเมหิ จตูหิ? อิธ ภิกฺขเว ภิกฺขุ สุภาสิตํเยว ภาสติ โน ทุพฺภาสิตํ, ธมฺมํเยว ภาสติ โน อธมฺมํ, ปิยํเยว ภาสติ โน อปฺปิยํ, สจฺจํเยว ภาสติ โน อลิกํ.\csromanspacer{} อิเมหิ โข ภิกฺขเว จตูหิ องฺเคหิ สมนฺนาคตา วาจา สุภาสิตา โหติ โน ทุพฺภาสิตา อนวชฺชา จ อนนุวชฺชา จ วิญฺญูนนฺติ.\csromanspacer{} อิทมโวจ ภควา, อิทํ วตฺวาน สุคโต, อถาปรํ เอตทโวจ สตฺถา\csromandash{}}

\proseitem{452}{\csromansymbolfootnote{*}{สํ 1. 190 ปิฏฺฐาทีสุ.}“สุภาสิตํ อุตฺตมมาหุ สนฺโต,}

\prose{ธมฺมํ ภเณ นาธมฺมํ ตํ ทุติยํ.}

\csromangathameasure{ปิยํ ภเณ นาปฺปิยํ ตํ ตติยํ,}
\csromangathagroup{\csromangathastanza{ปิยํ ภเณ นาปฺปิยํ ตํ ตติยํ,}}

\csromanheader{2}{สจฺจํ ภเณ นาลิกํ ตํ จตุตฺถนฺ”ติ. (1)}
\prose{อถ โข อายสฺมา วงฺคีโส อุฏฺฐายาสนา เอกํสํ จีวรํ กตฺวา เยน ภควา เตนญฺชลิํ ปณาเมตฺวา ภควนฺตํ เอตทโวจ “ปฏิภาติ มํ ภควา, ปฏิภาติ มํ สุคตา”ติ. “ปฏิภาตุ ตํ วงฺคีสา”ติ ภควา อโวจ.\csromanspacer{} อถ โข อายสฺมา วงฺคีโส ภควนฺตํ สมฺมุขา สารุปฺปาหิ คาถาหิ อภิตฺถวิ\csromandash{}}

\proseitem{453}{\csromansymbolmark{*}“ตเมว วาจํ ภาเสยฺย, ยาย’ตฺตานํ น ตาปเย.}

\prose{ปเร จ น วิหิํเสยฺย, สา เว วาจา สุภาสิตา. (2)}

\proseitem{454}{\csromansymbolmark{*}ปิยวาจเมว ภาเสยฺย, ยา วาจา ปฏินนฺทิตา.}

\prose{ยํ อนาทาย ปาปานิ, ปเรสํ ภาสเต ปิยํ. (3)}

\proseitem{455}{\csromansymbolmark{*}สจฺจํ เว อมตา วาจา, เอส ธมฺโม สนนฺตโน.}

\prose{สจฺเจ อตฺเถ จ ธมฺเม จ, อหุ สนฺโต ปติฏฺฐิตา. (4)}

\chapterheadpageread{3. มหาวคฺค}
\csromanmatikaanchor{mtk.262}
\csromanmatikaanchor{mtk.266}
\csromantocmarknopage{chapter}{3. มหาวคฺค}{mtk.262}
\bookmark[dest={mtk.262},level=0]{3. มหาวคฺค}
\proseitem{456}{\csromanfolio{345}ยํ พุทฺโธ ภาสติ วาจํ, เขมํ นิพฺพานปตฺติยา.}

\csromangathameasure{ทุกฺขสฺสนฺตกิริยาย, สา เว วาจานมุตฺตมา”ติ. (5) สุภาสิตสุตฺตํ ตติยํ.}
\csromangathagroup{\csromangathastanza{ทุกฺขสฺสนฺตกิริยาย, สา เว วาจานมุตฺตมา”ติ. (5) สุภาสิตสุตฺตํ ตติยํ.}}

\csromanheader{2}{4. ปูรฬาส (สุนฺทริกภารทฺวาช) สุตฺต}
\prose{เอวํ เม สุตํ\csromandash{}เอกํ สมยํ ภควา โกสเลสุ วิหรติ สุนฺทริกาย นทิยา ตีเร.\csromanspacer{} เตน โข ปน สมเยน สุนฺทริกภารทฺวาโช พฺราหฺมโณ สุนฺทริกาย นทิยา ตีเร อคฺคิํ ชุหติ, อคฺคิหุตฺตํ ปริจรติ.\csromanspacer{} อถ โข สุนฺทริกภารทฺวาโช พฺราหฺมโณ อคฺคิํ ชุหิตฺวา อคฺคิหุตฺตํ ปริจริตฺวา อุฏฺฐายาสนา สมนฺตา จตุทฺทิสา อนุวิโลเกสิ “โก นุ โข อิมํ หพฺยเสสํ ภุญฺเชยฺยา”ติ.\csromanspacer{} อทฺทสา โข สุนฺทริกภารทฺวาโช พฺราหฺมโณ ภควนฺตํ อวิทูเร อญฺญตรสฺมิํ รุกฺขมูเล สสีสํ ปารุตํ นิสินฺนํ, ทิสฺวาน วาเมน หตฺเถน หพฺยเสสํ คเหตฺวา ทกฺขิเณน หตฺเถน กมณฺฑลุํ คเหตฺวา เยน ภควา เตนุปสงฺกมิ.}

\prose{อถ โข ภควา สุนฺทริกภารทฺวาชสฺส พฺราหฺมณสฺส ปทสทฺเทน สีสํ วิวริ.\csromanspacer{} อถ โข สุนฺทริกภารทฺวาโช พฺราหฺมโณ “มุณฺโฑ อยํ ภวํ, มุณฺฑโก อยํ ภวนฺ”ติ ตโตว ปุน นิวตฺติตุกาโม อโหสิ.\csromanspacer{} อถ โข สุนฺทริกภารทฺวาชสฺส พฺราหฺมณสฺส เอตทโหสิ “มุณฺฑาปิ หิ อิเธกจฺเจ พฺราหฺมณา ภวนฺติ, ยํนูนาหํ อุปสงฺกมิตฺวา ชาติํ ปุจฺเฉยฺยนฺ”ติ.\csromanspacer{} อถ โข สุนฺทริกภารทฺวาโช พฺราหฺมโณ เยน ภควา เตนุปสงฺกมิ, อุปสงฺกมิตฺวา ภควนฺตํ เอตทโวจ “กิํชจฺโจ ภวนฺ”ติ.}

\prose{อถ โข ภควา สุนฺทริกภารทฺวาชํ พฺราหฺมณํ คาถาหิ อชฺฌภาสิ\csromandash{}}

\proseitem{457}{น พฺราหฺมโณ โน’มฺหิ น ราชปุตฺโต,}

\prose{น เวสฺสายโน อุท โกจิ โน’มฺหิ.}

\prose{โคตฺตํ ปริญฺญาย ปุถุชฺชนานํ,}

\csromanheader{2}{อกิญฺจโน มนฺต จรามิ โลเก. (1)}
\gambhira{สุตฺตนิปาตปาฬิ}
\proseitem{458}{\csromanfolio{346}สงฺฆาฏิวาสี อคโห\footnote{อคิโห (ก-สี, อิ), อเคโห (กตฺถจิ)} จรามิ,}

\prose{นิวุตฺตเกโส อภินิพฺพุตตฺโต.}

\prose{อลิปฺปมาโน อิธ มาณเวหิ,}

\csromanheader{2}{อกลฺลํ มํ พฺราหฺมณ ปุจฺฉสิ โคตฺตปญฺหํ. (2)}
\proseitem{459}{ปุจฺฉนฺติ เว โภ พฺราหฺมณา พฺราหฺมเณภิ สห “พฺราหฺมโณ โน ภวํ”ติ. (3)}

\csromanhangingbegin
\hangingheaditem{460}{พฺราหฺมโณ หิ เจ ตฺวํ พฺรูสิ, มญฺจ พฺรูสิ อพฺราหฺมณํ.}
\hangingbody{ตํ ตํ สาวิตฺติํ ปุจฺฉามิ, ติปทํ จตุวีสตกฺขรํ. (4)}
\csromanhangingend

\proseitem{461}{กิํ นิสฺสิตา อิสโย มนุชา,}

\prose{ขตฺติยา พฺราหฺมณา\footnote{ปฐมปาทนฺโต.} เทวตานํ,}

\csromanheader{2}{ยญฺญมกปฺปยิํสุ ปุถู อิธ โลเก\footnote{ทุติยปาทนฺโต (สี)}. (5)}
\proseitem{462}{ยทนฺตคู เวทคู ยญฺญกาเล,}

\csromanheader{2}{ยสฺสาหุติํ ลเภ ตสฺสิชฺเฌติ พฺรูมิ. (6)}
\proseitem{463}{อทฺธา หิ ตสฺส หุตมิชฺเฌ, (อิติ พฺราหฺมโณ,)}

\prose{ยํ ตาทิสํ เวทคุมทฺทสาม.}

\prose{ตุมฺหาทิสานญฺหิ อทสฺสเนน,}

\csromanheader{2}{อญฺโญ ชโน ภุญฺชติ ปูรฬาสํ. (7)}
\proseitem{464}{ตสฺมาติห ตฺวํ พฺราหฺมณ อตฺเถน,}

\prose{อตฺถิโก อุปสงฺกมฺม ปุจฺฉ.}

\prose{สนฺตํ วิธูมํ อนีฆํ นิราสํ,}

\csromanheader{2}{อปฺเปวิธ อภิวินฺเท สุเมธํ. (8)}
\proseitem{465}{ยญฺเญ รโตหํ โภ โคตม,}

\prose{ยญฺญํ ยิฏฺฐุกาโม นาหํ ปชานามิ.}

\prose{อนุสาสตุ มํ ภวํ,}

\csromanheader{2}{ยตฺถ หุตํ อิชฺฌเต พฺรูหิ เม’ตํ. (9)}
\prose{เตน หิ ตฺวํ พฺราหฺมณ โอทหสฺสุ โสตํ, ธมฺมํ เต เทเสสฺสามิ\csromandash{}}

\chapterheadpageread{3. มหาวคฺค}
\csromangathameasure{มา ชาติํ ปุจฺฉี จรณญฺจ ปุจฺฉ, \\ กฏฺฐา หเว ชายติ ชาตเวโท. \\ นีจากุลีโนปิ มุนี ธิตีมา,}
\csromangathagroup{\csromangathastanza{\csromanfolio{347}มา ชาติํ ปุจฺฉี จรณญฺจ ปุจฺฉ, \\ กฏฺฐา หเว ชายติ ชาตเวโท. \\ นีจากุลีโนปิ มุนี ธิตีมา,}}

\csromanheader{2}{อาชานิโย โหติ หิรีนิเสโธ. (10)}
\proseitem{467}{\csromansymbolfootnote{*}{สํ 1. 170 ปิฏฺเฐปิ.}สจฺเจน ทนฺโต ทมสา อุเปโต,}

\prose{เวทนฺตคู วูสิตพฺรหฺมจริโย.}

\csromangathameasure{กาเลน ตมฺหิ หพฺยํ ปเวจฺเฉ,}
\csromangathagroup{\csromangathastanza{กาเลน ตมฺหิ หพฺยํ ปเวจฺเฉ,}}

\csromanheader{2}{โย พฺราหฺมโณ ปุญฺญเปกฺโข\footnote{ปุญฺญเปโข (สี, อิ)} ยเชถ. (11)}
\csromangathameasure{เย กาเม หิตฺวา อคหา จรนฺติ, \\ สุสญฺญตตฺตา ตสรํว อุชฺชุํ. \\ กาเลน เตสุ หพฺยํ ปเวจฺเฉ,}
\csromangathagroup{\csromangathastanza{เย กาเม หิตฺวา อคหา จรนฺติ, \\ สุสญฺญตตฺตา ตสรํว อุชฺชุํ. \\ กาเลน เตสุ หพฺยํ ปเวจฺเฉ,}}

\csromanheader{2}{โย พฺราหฺมโณ ปุญฺญเปกฺโข ยเชถ. (12)}
\csromangathameasure{เย วีตราคา สุสมาหิตินฺทฺริยา, \\ จนฺโทว ราหุคฺคหณา ปมุตฺตา. \\ กาเลน เตสุ หพฺยํ ปเวจฺเฉ,}
\csromangathagroup{\csromangathastanza{เย วีตราคา สุสมาหิตินฺทฺริยา, \\ จนฺโทว ราหุคฺคหณา ปมุตฺตา. \\ กาเลน เตสุ หพฺยํ ปเวจฺเฉ,}}

\csromanheader{2}{โย พฺราหฺมโณ ปุญฺญเปกฺโข ยเชถ. (13)}
\csromangathameasure{อสชฺชมานา วิจรนฺติ โลเก, \\ สทา สตา หิตฺวา มมายิตานิ. \\ กาเลน เตสุ หพฺยํ ปเวจฺเฉ,}
\csromangathagroup{\csromangathastanza{อสชฺชมานา วิจรนฺติ โลเก, \\ สทา สตา หิตฺวา มมายิตานิ. \\ กาเลน เตสุ หพฺยํ ปเวจฺเฉ,}}

\csromanheader{2}{โย พฺราหฺมโณ ปุญฺญเปกฺโข ยเชถ. (14)}
\csromangathameasure{โย กาเม หิตฺวา อภิภุยฺยจารี, \\ โย เวทิ ชาตีมรณสฺส อนฺตํ. \\ ปรินิพฺพุโต อุทกรหโทว สีโต,}
\csromangathagroup{\csromangathastanza{โย กาเม หิตฺวา อภิภุยฺยจารี, \\ โย เวทิ ชาตีมรณสฺส อนฺตํ. \\ ปรินิพฺพุโต อุทกรหโทว สีโต,}}

\csromanheader{2}{ตถาคโต อรหติ ปูรฬาสํ. (15)}
\csromangathameasure{สโม สเมหิ วิสเมหิ ทูเร, \\ ตถาคโต โหติ อนนฺตปญฺโญ. \\ อนุปลิตฺโต อิธ วา หุรํ วา,}
\csromangathagroup{\csromangathastanza{สโม สเมหิ วิสเมหิ ทูเร, \\ ตถาคโต โหติ อนนฺตปญฺโญ. \\ อนุปลิตฺโต อิธ วา หุรํ วา,}}

\csromanheader{2}{ตถาคโต อรหติ ปูรฬาสํ. (16)}
\gambhira{สุตฺตนิปาตปาฬิ}
\proseitem{473}{\csromanfolio{348}ยมฺหิ น มายา วสติ น มาโน,}

\prose{โย วีตโลโภ อมโม นิราโส.}

\csromangathameasure{ปณุนฺนโกโธ อภินิพฺพุตตฺโต, \\ โย พฺราหฺมโณ โสกมลํ อหาสิ.}
\csromangathagroup{\csromangathastanza{ปณุนฺนโกโธ อภินิพฺพุตตฺโต, \\ โย พฺราหฺมโณ โสกมลํ อหาสิ.}}

\csromanheader{2}{ตถาคโต อรหติ ปูรฬาสํ. (17)}
\proseitem{474}{นิเวสนํ โย มนโส อหาสิ,}

\prose{ปริคฺคหา ยสฺส น สนฺติ เกจิ.}

\prose{อนุปาทิยาโน อิธ วา หุรํ วา,}

\csromanheader{2}{ตถาคโต อรหติ ปูรฬาสํ. (18)}
\proseitem{475}{สมาหิโต โย อุทตาริ โอฆํ,}

\prose{ธมฺมํ จญฺญาสิ ปรมาย ทิฏฺฐิยา.}

\prose{ขีณาสโว อนฺติมเทหธารี,}

\csromanheader{2}{ตถาคโต อรหติ ปูรฬาสํ. (19)}
\proseitem{476}{ภวาสวา ยสฺส วจี ขรา จ,}

\prose{วิธูปิตา อตฺถคตา น สนฺติ.}

\prose{ส เวทคู สพฺพธิ วิปฺปมุตฺโต,}

\csromanheader{2}{ตถาคโต อรหติ ปูรฬาสํ. (20)}
\proseitem{477}{สงฺคาติโค ยสฺส น สนฺติ สงฺคา,}

\prose{โย มานสตฺเตสุ อมานสตฺโต.}

\prose{ทุกฺขํ ปริญฺญาย สเขตฺตวตฺถุํ,}

\csromanheader{2}{ตถาคโต อรหติ ปูรฬาสํ. (21)}
\proseitem{478}{อาสํ อนิสฺสาย วิเวกทสฺสี,}

\prose{ปรเวทิยํ ทิฏฺฐิมุปาติวตฺโต.}

\prose{อารมฺมณา ยสฺส น สนฺติ เกจิ,}

\csromanheader{2}{ตถาคโต อรหติ ปูรฬาสํ. (22)}
\proseitem{479}{ปโรปรา\footnote{ปโรวรา (สี, อิ)} ยสฺส สเมจฺจ ธมฺมา,}

\prose{วิธูปิตา อตฺถคตา น สนฺติ.}

\prose{สนฺโต อุปาทานขเย วิมุตฺโต,}

\csromanheader{2}{ตถาคโต อรหติ ปูรฬาสํ. (23)}
\chapterheadpageread{3. มหาวคฺค}
\csromangathameasure{สํโยชนํชาติขยนฺตทสฺสี, \\ โย’ปานุทิ ราคปถํ อเสสํ. \\ สุทฺโธ นิโทโส วิมโล อกาโจ,}
\csromangathagroup{\csromangathastanza{\csromanfolio{349}สํโยชนํชาติขยนฺตทสฺสี, \\ โย’ปานุทิ ราคปถํ อเสสํ. \\ สุทฺโธ นิโทโส วิมโล อกาโจ\footnote{อกาโม (สี, สฺยา)},}}

\csromanheader{2}{ตถาคโต อรหติ ปูรฬาสํ. (24)}
\proseitem{481}{โย อตฺตโน อตฺตานํ\footnote{อตฺตนาตฺตานํ (สี, สฺยา)} นานุปสฺสติ,}

\prose{สมาหิโต อุชฺชุคโต ฐิตตฺโต.}

\csromangathameasure{ส เว อเนโช อขิโล อกงฺโข,}
\csromangathagroup{\csromangathastanza{ส เว อเนโช อขิโล อกงฺโข,}}

\csromanheader{2}{ตถาคโต อรหติ ปูรฬาสํ. (25)}
\proseitem{482}{โมหนฺตรา ยสฺส น สนฺติ เกจิ,}

\prose{สพฺเพสุ ธมฺเมสุ จ ญาณทสฺสี.}

\csromangathameasure{สรีรญฺจ อนฺติมํ ธาเรติ, \\ ปตฺโต จ สมฺโพธิมนุตฺตรํ สิวํ. \\ เอตฺตาวตา ยกฺขสฺส สุทฺธิ,}
\csromangathagroup{\csromangathastanza{สรีรญฺจ อนฺติมํ ธาเรติ, \\ ปตฺโต จ สมฺโพธิมนุตฺตรํ สิวํ. \\ เอตฺตาวตา ยกฺขสฺส สุทฺธิ,}}

\csromanmatikaanchor{mtk.263}
\csromantocmarkref{subsection}{1. ปพฺพชฺชาสุตฺต}{mtk.263}
\bookmark[dest={mtk.263},level=2]{1. ปพฺพชฺชาสุตฺต}
\csromanheader{2}{ตถาคโต อรหติ ปูรฬาสํ. (26)}
\proseitem{483}{หุตญฺจ\footnote{หุตฺตญฺจ (สี, ก)} มยฺหํ หุตมตฺถุ สจฺจํ,}

\prose{ยํ ตาทิสํ เวทคุนํ อลตฺถํ.}

\csromangathameasure{พฺรหฺมา หิ สกฺขิ ปฏิคณฺหาตุ เม ภควา,}
\csromangathagroup{\csromangathastanza{พฺรหฺมา หิ สกฺขิ ปฏิคณฺหาตุ เม ภควา,}}

\csromanheader{2}{ภุญฺชตุ เม ภควา ปูรฬาสํ. (27)}
\proseitem{484}{\csromansymbolfootnote{*}{เหฏฺฐา 292; สํ 1. 175 ปิฏฺเฐสุปิ.}คาถาภิคีตํ เม อโภชเนยฺยํ,}

\prose{สมฺปสฺสตํ พฺราหฺมณ เนส ธมฺโม.}

\csromanheader{2}{คาถา ภิคีตํ ปนุทนฺติ พุทฺธา,}
\csromanheader{2}{ธมฺเม สตี พฺราหฺมณ วุตฺติเรสา. (28)}
\proseitem{485}{\csromansymbolmark{*}อญฺเญน จ เกวลินํ มเหสิํ,}

\prose{ขีณาสวํ กุกฺกุจฺจวูปสนฺตํ.}

\csromangathameasure{อนฺเนน ปาเนน อุปฏฺฐหสฺสุ,}
\csromangathagroup{\csromangathastanza{อนฺเนน ปาเนน อุปฏฺฐหสฺสุ,}}

\csromanheader{2}{เขตฺตํ หิ ตํ ปุญฺญเปกฺขสฺส โหติ. (29)}
\gambhira{สุตฺตนิปาตปาฬิ}
\proseitem{486}{\csromanfolio{350}สาธาหํ ภควา ตถา วิชญฺญํ,}

\prose{โย ทกฺขิณํ ภุญฺเชยฺย มาทิสสฺส.}

\prose{ยํ ยญฺญกาเล ปริเยสมาโน,}

\csromanheader{2}{ปปฺปุยฺย ตว สาสนํ. (30)}
\proseitem{487}{สารมฺภา ยสฺส วิคตา, จิตฺตํ ยสฺส อนาวิลํ.}

\prose{วิปฺปมุตฺโต จ กาเมหิ, ถินํ ยสฺส ปนูทิตํ. (31)}

\proseitem{488}{สีมนฺตานํ วิเนตารํ, ชาติมรณโกวิทํ.}

\prose{มุนิํ โมเนยฺยสมฺปนฺนํ, ตาทิสํ ยญฺญมาคตํ. (32)}

\proseitem{489}{ภกุฏิํ\footnote{ภูกุฏิํ (ก-สี), ภากุฏิํ (ก-สี, ม 1. 176)} วินยิตฺวาน, ปญฺชลิกา นมสฺสถ.}

\prose{ปูเชถ อนฺนปาเนน, เอวํ อิชฺฌนฺติ ทกฺขิณา. (33)}

\proseitem{490}{พุทฺโธ ภวํ อรหติ ปูรฬาสํ, ปุญฺญเขตฺตมนุตฺตรํ.}

\prose{อายาโค สพฺพโลกสฺส, โภโต ทินฺนํ มหปฺผลนฺติ. (34)}

\prose{อถ โข สุนฺทริกภารทฺวาโช พฺราหฺมโณ ภควนฺตํ เอตทโวจ “อภิกฺกนฺตํ โภ โคตม, อภิกฺกนฺตํ โภ โคตม.\csromanspacer{} เสยฺยถาปิ โภ โคตม นิกฺกุชฺชิตํ วา อุกฺกุชฺเชยฺย, ปฏิจฺฉนฺนํ วา วิวเรยฺย, มูฬฺหสฺส วา มคฺคํ อาจิกฺเขยฺย, อนฺธกาเร วา เตลปชฺโชตํ ธาเรยฺย ‘จกฺขุมนฺโต รูปานิ ทกฺขนฺตี’ติ.\csromanspacer{} เอวเมวํ โภตา โคตเมน อเนกปริยาเยน ธมฺโม ปกาสิโต, เอสาหํ ภวนฺตํ โคตมํ สรณํ คจฺฉามิ ธมฺมญฺจ ภิกฺขุสํฆญฺจ.\csromanspacer{} ลเภยฺยาหํ โภโต โคตมสฺส สนฺติเก ปพฺพชฺชํ, ลเภยฺยํ อุปสมฺปทนฺ”ติ.\csromanspacer{} อลตฺถ โข สุนฺทริกภารทฺวาโช พฺราหฺมโณ ฯเปฯ อรหตํ อโหสีติ.\csromanspacer{} สุนฺทริกภารทฺวาชสุตฺตํ จตุตฺถํ.}

\csromanheader{2}{5. มาฆสุตฺต}
\prose{เอวํ เม สุตํ\csromandash{}เอกํ สมยํ ภควา ราชคเห วิหรติ คิชฺฌกูเฏ ปพฺพเต.\csromanspacer{} อถ โข มาโฆ มาณโว เยน ภควา เตนุปสงฺกมิ,}

\vspace{\tipitakaparskip}
\csromancenter{มหาวคฺค อุปสงฺกมิตฺวา ภควตา สทฺธิํ สมฺโมทิ, สมฺโมทนียํ กถํ สารณียํ วีติสาเรตฺวา เอกมนฺตํ นิสีทิ, เอกมนฺตํ นิสินฺโน โข มาโฆ มาณโว ภควนฺตํ เอตทโวจ\csromandash{}}
\prose{\csromanfolio{351}อหํ หิ โภ โคตม ทายโก ทานปติ วทญฺญู ยาจโยโค ธมฺเมน โภเค ปริเยสามิ, ธมฺเมน โภเค ปริเยสิตฺวา ธมฺมลทฺเธหิ โภเคหิ ธมฺมาธิคเตหิ เอกสฺสปิ ททามิ ทฺวินฺนมฺปิ ติณฺณมฺปิ จตุนฺนมฺปิ ปญฺจนฺนมฺปิ ฉนฺนมฺปิ สตฺตนฺนมฺปิ อฏฺฐนฺนมฺปิ นวนฺนมฺปิ ทสนฺนมฺปิ ททามิ, วีสายปิ ติํสายปิ จตฺตาลีสายปิ ปญฺญาสายปิ ททามิ, สตสฺสปิ ททามิ ภิยฺโยปิ ททามิ.\csromanspacer{} กจฺจาหํ โภ โคตม เอวํ ททนฺโต เอวํ ยชนฺโต พหุํ ปุญฺญํ ปสวามีติ.}

\prose{ตคฺฆ ตฺวํ มาณว เอวํ ททนฺโต เอวํ ยชนฺโต พหุํ ปุญฺญํ ปสวสิ, โย โข มาณว ทายโก ทานปติ วทญฺญู ยาจโยโค ธมฺเมน โภเค ปริเยสติ, ธมฺเมน โภเค ปริเยสิตฺวา ธมฺมลทฺเธหิ โภเคหิ ธมฺมาธิคเตหิ เอกสฺสปิ ททาติ ฯเปฯ สตสฺสปิ ททาติ ภิยฺโยปิ ททาติ, พหุํ โส ปุญฺญํ ปสวตีติ.\csromanspacer{} อถ โข มาโฆ มาณโว ภควนฺตํ คาถาย อชฺฌภาสิ\csromandash{}}

\proseitem{491}{ปุจฺฉามหํ โคตมํ วทญฺญุํ, (อิติ มาโฆ มาณโว,)}

\csromangathameasure{กาสายวาสิํ อคหํ จรนฺตํ, \\ โย ยาจโยโค ทานปติ คหฏฺโฐ. \\ ปุญฺญตฺถิโก ยชติ ปุญฺญเปกฺโข, \\ ททํ ปเรสํ อิธ อนฺนปานํ.}
\csromangathagroup{\csromangathastanza{กาสายวาสิํ อคหํ\footnote{อคิหํ (สี), อเคหํ (อิ)} จรนฺตํ, \\ โย ยาจโยโค ทานปติ\footnote{อคิหํ (สี), อเคหํ (อิ)} คหฏฺโฐ. \\ ปุญฺญตฺถิโก ยชติ ปุญฺญเปกฺโข\footnote{อคิหํ (สี), อเคหํ (อิ)}, \\ ททํ ปเรสํ อิธ อนฺนปานํ.}}

\csromanheader{2}{กถํ หุตํ ยชมานสฺส สุชฺเฌ. (1)}
\proseitem{492}{โย ยาจโยโค ทานปติ คหฏฺโฐ, (มาฆาติ ภควา,)}

\csromangathameasure{ปุญฺญตฺถิโก ยชติ ปุญฺญเปกฺโข, \\ ททํ ปเรสํ อิธ อนฺนปานํ.}
\csromangathagroup{\csromangathastanza{ปุญฺญตฺถิโก ยชติ ปุญฺญเปกฺโข, \\ ททํ ปเรสํ อิธ อนฺนปานํ.}}

\csromanheader{2}{อาราธเย ทกฺขิเณยฺเยภิ ตาทิ. (2)}
\gambhira{สุตฺตนิปาตปาฬิ}
\proseitem{493}{\csromanfolio{352}โย ยาจโยโค ทานปติ คหฏฺโฐ, (อิติ มาโฆ มาณโว,)}

\csromangathameasure{ปุญฺญตฺถิโก ยชติ ปุญฺญเปกฺโข, \\ ททํ ปเรสํ อิธ อนฺนปานํ.}
\csromangathagroup{\csromangathastanza{ปุญฺญตฺถิโก ยชติ ปุญฺญเปกฺโข, \\ ททํ ปเรสํ อิธ อนฺนปานํ.}}

\csromanheader{2}{อกฺขาหิ เม ภควา ทกฺขิเณยฺเย. (3)}
\csromanmatikaanchor{mtk.264}
\csromantocmarkref{subsection}{2. ปธานสุตฺต}{mtk.264}
\bookmark[dest={mtk.264},level=2]{2. ปธานสุตฺต}
\proseitem{494}{เย เว อสตฺตา\footnote{อลคฺคา (สฺยา)} วิจรนฺติ โลเก,}

\prose{อกิญฺจนา เกวลิโน ยตตฺตา.}

\prose{กาเลน เตสุ หพฺยํ ปเวจฺเฉ,}

\csromanheader{2}{โย พฺราหฺมโณ ปุญฺญเปกฺโข ยเชถ. (4)}
\proseitem{495}{เย สพฺพสํโยชนพนฺธนจฺฉิทา,}

\prose{ทนฺตา วิมุตฺตา อนีฆา นิราสา.}

\prose{กาเลน เตสุ หพฺยํ ปเวจฺเฉ,}

\csromanheader{2}{โย พฺราหฺมโณ ปุญฺญเปกฺโข ยเชถ. (5)}
\proseitem{496}{เย สพฺพสํโยชนวิปฺปมุตฺตา,}

\prose{ทนฺตา วิมุตฺตา อนีฆา นิราสา.}

\prose{กาเลน เตสุ หพฺยํ ปเวจฺเฉ,}

\csromanheader{2}{โย พฺราหฺมโณ ปุญฺญเปกฺโข ยเชถ. (6)}
\proseitem{497}{ราคญฺจ โทสญฺจ ปหาย โมหํ,}

\prose{ขีณาสวา วูสิตพฺรหฺมจริยา.}

\prose{กาเลน เตสุ หพฺยํ ปเวจฺเฉ,}

\csromanheader{2}{โย พฺราหฺมโณ ปุญฺญเปกฺโข ยเชถ. (7)}
\proseitem{498}{เยสุ น มายา วสติ น มาโน,}

\prose{ขีณาสวา วูสิตพฺรหฺมจริยา.}

\prose{กาเลน เตสุ หพฺยํ ปเวจฺเฉ,}

\csromanheader{2}{โย พฺราหฺมโณ ปุญฺญเปกฺโข ยเชถ. (8)}
\chapterheadpageread{3. มหาวคฺค}
\proseitem{499}{\csromanfolio{353}เย วีตโลภา อมมา นิราสา,}

\prose{ขีณาสวา วูสิตพฺรหฺมจริยา.}

\prose{กาเลน เตสุ หพฺยํ ปเวจฺเฉ,}

\csromanheader{2}{โย พฺราหฺมโณ ปุญฺญเปกฺโข ยเชถ. (9)}
\proseitem{500}{เย เว น ตณฺหาสุ อุปาติปนฺนา,}

\prose{วิตเรยฺย โอฆํ อมมา จรนฺติ.}

\prose{กาเลน เตสุ หพฺยํ ปเวจฺเฉ,}

\csromanheader{2}{โย พฺราหฺมโณ ปุญฺญเปกฺโข ยเชถ. (10)}
\proseitem{501}{เยสํ ตณฺหา นตฺถิ กุหิญฺจิ โลเก,}

\prose{ภวาภวาย อิธ วา หุรํ วา.}

\prose{กาเลน เตสุ หพฺยํ ปเวจฺเฉ,}

\csromanheader{2}{โย พฺราหฺมโณ ปุญฺญเปกฺโข ยเชถ. (11)}
\proseitem{502}{เย กาเม หิตฺวา อคหา จรนฺติ,}

\prose{สุสญฺญตตฺตา ตสรํว อุชฺชุํ.}

\prose{กาเลน เตสุ หพฺยํ ปเวจฺเฉ,}

\csromanheader{2}{โย พฺราหฺมโณ ปุญฺญเปกฺโข ยเชถ. (12)}
\proseitem{503}{เย วีตราคา สุสมาหิตินฺทฺริยา,}

\prose{จนฺโทว ราหุคฺคหณา ปมุตฺตา.}

\prose{กาเลน เตสุ หพฺยํ ปเวจฺเฉ,}

\csromanheader{2}{โย พฺราหฺมโณ ปุญฺญเปกฺโข ยเชถ. (13)}
\proseitem{504}{สมิตาวิโน วีตราคา อโกปา,}

\prose{เยสํ คตี นตฺถิธ วิปฺปหาย.}

\prose{กาเลน เตสุ หพฺยํ ปเวจฺเฉ,}

\csromanheader{2}{โย พฺราหฺมโณ ปุญฺญเปกฺโข ยเชถ. (14)}
\proseitem{505}{ชหิตฺวา ชาติมรณํ อเสสํ,}

\prose{กถํกถิํ สพฺพมุปาติวตฺตา.}

\prose{กาเลน เตสุ หพฺยํ ปเวจฺเฉ,}

\csromanheader{2}{โย พฺราหฺมโณ ปุญฺญเปกฺโข ยเชถ. (15)}
\gambhira{สุตฺตนิปาตปาฬิ}
\proseitem{506}{\csromanfolio{354}เย อตฺตทีปา วิจรนฺติ โลเก,}

\prose{อกิญฺจนา สพฺพธิ วิปฺปมุตฺตา.}

\prose{กาเลน เตสุ หพฺยํ ปเวจฺเฉ,}

\csromanheader{2}{โย พฺราหฺมโณ ปุญฺญเปกฺโข ยเชถ. (16)}
\csromanmatikaanchor{mtk.265}
\csromantocmarkref{subsection}{3. สุภาสิตสุตฺต}{mtk.265}
\bookmark[dest={mtk.265},level=2]{3. สุภาสิตสุตฺต}
\csromantocmarkref{subsection}{4. สุนฺทริกภารทฺวาชสุตฺต}{mtk.266}
\bookmark[dest={mtk.266},level=2]{4. สุนฺทริกภารทฺวาชสุตฺต}
\proseitem{507}{เย เหตฺถ ชานนฺติ ยถา ตถา อิทํ,}

\prose{อยมนฺติมา นตฺถิ ปุนพฺภโวติ.}

\prose{กาเลน เตสุ หพฺยํ ปเวจฺเฉ,}

\csromanheader{2}{โย พฺราหฺมโณ ปุญฺญเปกฺโข ยเชถ. (17)}
\proseitem{508}{โย เวทคู ฌานรโต สตีมา,}

\prose{สมฺโพธิปตฺโต สรณํ พหูนํ.}

\prose{กาเลน ตมฺหิ หพฺยํ ปเวจฺเฉ,}

\csromanheader{2}{โย พฺราหฺมโณ ปุญฺญเปกฺโข ยเชถ. (18)}
\proseitem{509}{อทฺธา อโมฆา มม ปุจฺฉนา อหุ,}

\prose{อกฺขาสิ เม ภควา ทกฺขิเณยฺเย.}

\prose{ตฺวเญฺหตฺถ ชานาสิ ยถา ตถา อิทํ,}

\csromanheader{2}{ตถา หิ เต วิทิโต เอส ธมฺโม. (19)}
\proseitem{510}{โย ยาจโยโค ทานปติ คหฏฺโฐ, (อิติ มาโฆ มาณโว,)}

\prose{ปุญฺญตฺถิโก ยชติ ปุญฺญเปกฺโข.}

\prose{ททํ ปเรสํ อิธ อนฺนปานํ,}

\csromanheader{2}{อกฺขาหิ เม ภควา ยญฺญสมฺปทํ. (20)}
\proseitem{511}{ยชสฺสุ ยชมาโน มาฆาติ ภควา,}

\prose{สพฺพตฺถ จ วิปฺปสาเทหิ จิตฺตํ.}

\prose{อารมฺมณํ ยชมานสฺส ยญฺโญ,}

\csromanheader{2}{เอตฺถ ปติฏฺฐาย ชหาติ โทสํ. (21)}
\proseitem{512}{โส วีตราโค ปวิเนยฺย โทสํ,}

\prose{เมตฺตํ จิตฺตํ ภาวยมปฺปมาณํ.}

\prose{รตฺตินฺทิวํ สตตมปฺปมตฺโต,}

\csromanheader{2}{สพฺพา ทิสา ผรติ อปฺปมญฺญํ. (22)}
\chapterheadpageread{3. มหาวคฺค}
\proseitem{513}{\csromanfolio{355}โก สุชฺฌติ มุจฺจติ พชฺฌตี จ,}

\prose{เกนตฺตนา คจฺฉติ\footnote{เกนตฺเถนา คจฺฉติ (ก)} พฺรหฺมโลกํ.}

\csromangathameasure{อชานโต เม มุนิ พฺรูหิ ปุฏฺโฐ, \\ ภควา หิ เม สกฺขิ พฺรหฺมชฺชทิฏฺโฐ. \\ ตุวํ หิ โน พฺรหฺมสโมสิ สจฺจํ,}
\csromangathagroup{\csromangathastanza{อชานโต เม มุนิ พฺรูหิ ปุฏฺโฐ, \\ ภควา หิ เม สกฺขิ พฺรหฺมชฺชทิฏฺโฐ. \\ ตุวํ หิ โน พฺรหฺมสโมสิ สจฺจํ,}}

\csromanheader{2}{กถํ อุปปชฺชติ พฺรหฺมโลกํ ชุติม. (23)}
\proseitem{514}{โย ยชติ ติวิธํ ยญฺญสมฺปทํ, (มาฆาติ ภควา,)}

\prose{อาราธเย ทกฺขิเณยฺเยภิ ตาทิ.}

\prose{เอวํ ยชิตฺวา สมฺมา ยาจโยโค,}

\csromanheader{2}{อุปปชฺชติ พฺรหฺมโลกนฺติ พฺรูมีติ. (24)}
\prose{เอวํ วุตฺเต มาโฆ มาณโว ภควนฺตํ เอตทโวจ “อภิกฺกนฺตํ โภ โคตม ฯเปฯ ‘อชฺชตคฺเค ปาณุเปตํ สรณํ คตนฺ’ติ”.\csromanspacer{} มาฆสุตฺตํ ปญฺจมํ.}

\csromanheader{2}{6. สภิยสุตฺต}
\prose{เอวํ เม สุตํ\csromandash{}เอกํ สมยํ ภควา ราชคเห วิหรติ เวฬุวเน กลนฺทกนิวาเป.\csromanspacer{} เตน โข ปน สมเยน สภิยสฺส ปริพฺพาชกสฺส ปุราณสาโลหิตาย เทวตาย ปญฺหา อุทฺทิฏฺฐา โหนฺติ “โย เต สภิย สมโณ วา พฺราหฺมโณ วา อิเม ปเญฺห ปุฏฺโฐ พฺยากโรติ, ตสฺส สนฺติเก พฺรหฺมจริยํ จเรยฺยาสี”ติ.}

\prose{อถ โข สภิโย ปริพฺพาชโก ตสฺสา เทวตาย สนฺติเก เต ปเญฺห อุคฺคเหตฺวา เย เต สมณพฺราหฺมณา สงฺฆิโน คณิโน คณาจริยา ญาตา ยสสฺสิโน ติตฺถกรา สาธุสมฺมตา พหุชนสฺส.\csromanspacer{} เสยฺยถิทํ, * ปูรโณ กสฺสโป, มกฺขลิโคสาโล, อชิโต เกสกมฺพโล, ปกุโธ\footnote{กกุโธ (สี), ปกุทฺโธ (สฺยา, กํ)} กจฺจาโน, สญฺจโย\footnote{สญฺชโย (สี, สฺยา, กํ, อิ)} เพลฏฺฐปุตฺโต\footnote{เพลฺลฏฺฐิปุตฺโต (สี, อิ), เวฬฏฺฐปุตฺโต (สฺยา)}, นิคณฺโฐ นาฏปุตฺโต\footnote{นาตปุตฺโต (สี, อิ)}, เต อุปสงฺกมิตฺวา เต ปเญฺห ปุจฺฉติ.\csromanspacer{} เต}

\gambhira{สุตฺตนิปาตปาฬิ}
\prose{\csromanfolio{356}สภิเยน ปริพฺพาชเกน ปเญฺห ปุฏฺฐา น สมฺปายนฺติ, อสมฺปายนฺตา โกปญฺจ โทสญฺจ อปฺปจฺจยญฺจ ปาตุกโรนฺติ.\csromanspacer{} อปิ จ สภิยํเยว ปริพฺพาชกํ ปฏิปุจฺฉนฺติ.}

\prose{อถ โข สภิยสฺส ปริพฺพาชกสฺส เอตทโหสิ “เย โข เต โภนฺโต สมณพฺราหฺมณา สงฺฆิโน คณิโน คณาจริยา ญาตา ยสสฺสิโน ติตฺถกรา สาธุสมฺมตา พหุชนสฺส.\csromanspacer{} เสยฺยถิทํ, ปูรโณ กสฺสโป ฯเปฯ นิคณฺโฐ นาฏปุตฺโต, เต มยา ปเญฺห ปุฏฺฐา น สมฺปายนฺติ, อสมฺปายนฺตา โกปญฺจ โทสญฺจ อปฺปจฺจยญฺจ ปาตุกโรนฺติ.\csromanspacer{} อปิ จ มญฺเญเวตฺถ ปฏิปุจฺฉนฺติ.\csromanspacer{} ยํนูนาหํ หีนายาวตฺติตฺวา กาเม ปริภุญฺเชยฺยนฺ”ติ.}

\prose{อถ โข สภิยสฺส ปริพฺพาชกสฺส เอตทโหสิ “อยมฺปิ โข สมโณ โคตโม สงฺฆี เจว คณี จ คณาจริโย จ ญาโต ยสสฺสี ติตฺถกโร สาธุสมฺมโต พหุชนสฺส, ยํนูนาหํ สมณํ โคตมํ อุปสงฺกมิตฺวา อิเม ปเญฺห ปุจฺเฉยฺยนฺ”ติ.}

\prose{อถ โข สภิยสฺส ปริพฺพาชกสฺส เอตทโหสิ “เยปิ โข เต\footnote{เย โข เต (สฺยา), ยํ โข เต (ก)} โภนฺโต สมณพฺราหฺมณา ชิณฺณา วุฑฺฒา มหลฺลกา อทฺธคตา วโย-อนุปฺปตฺตา เถรา รตฺตญฺญู จิรปพฺพชิตา สงฺฆิโน คณิโน คณาจริยา ญาตา ยสสฺสิโน ติตฺถกรา สาธุสมฺมตา พหุชนสฺส.\csromanspacer{} เสยฺยถิทํ, ปูรโณ กสฺสโป ฯเปฯ นิคณฺโฐ นาฏปุตฺโต, เตปิ มยา ปเญฺห ปุฏฺฐา น สมฺปายนฺติ, อสมฺปายนฺตา โกปญฺจ โทสญฺจ อปฺปจฺจยญฺจ ปาตุกโรนฺติ.\csromanspacer{} อปิ จ มญฺเญเวตฺถ ปฏิปุจฺฉนฺติ.\csromanspacer{} กิํ ปน เม สมโณ โคตโม อิเม ปเญฺห ปุฏฺโฐ พฺยากริสฺสติ, สมโณ หิ โคตโม ทหโร เจว ชาติยา, นโว จ ปพฺพชฺชายาติ.}

\prose{อถ โข สภิยสฺส ปริพฺพาชกสฺส เอตทโหสิ “สมโณ โข\footnote{สมโณ โข โคตโม (สฺยา, ก)} “ทหโรติ น อุญฺญาตพฺโพ น ปริโภตพฺโพ.\csromanspacer{} ทหโรปิ เจส สมโณ โคตโม มหิทฺธิโก โหติ มหานุภาโว, ยํนูนาหํ สมณํ โคตมํ อุปสงฺกมิตฺวา อิเม ปเญฺห ปุจฺเฉยฺยนฺ”ติ.}

\prose{อถ โข สภิโย ปริพฺพาชโก เยน ราชคหํ เตน จาริกํ ปกฺกามิ, อนุปุพฺเพน จาริกํ จรมาโน เยน ราชคหํ เวฬุวนํ กลนฺทกนิวาโป, เยน ภควา เตนุปสงฺกมิ, อุปสงฺกมิตฺวา ภควตา}

\vspace{\tipitakaparskip}
\csromancenter{มหาวคฺค สทฺธิํ สมฺโมทิ, สมฺโมทนียํ กถํ สารณียํ วีติสาเรตฺวา เอกมนฺตํ นิสีทิ, เอกมนฺตํ นิสินฺโน โข สภิโย ปริพฺพาชโก ภควนฺตํ คาถาย อชฺฌภาสิ\csromandash{}}
\proseitem{515}{\csromanfolio{357}กงฺขี เวจิกิจฺฉี อาคมํ, (อิติ สภิโย,)}

\prose{ปเญฺห ปุจฺฉิตุํ อภิกงฺขมาโน.}

\prose{เตส’นฺตกโร ภวาหิ\footnote{ภวาหิเม (อิ, ก)} ปเญฺห เม ปุฏฺโฐ,}

\csromanheader{2}{อนุปุพฺพํ อนุธมฺมํ พฺยากโรหิ เม. (1)}
\proseitem{516}{ทูรโต อาคโตสิ สภิย, (อิติ ภควา,)}

\prose{ปเญฺห ปุจฺฉิตุํ อภิกงฺขมาโน.}

\prose{เตส’นฺตกโร ภวามิ\footnote{เตสมนฺตกโรมิ เต (ก)} ปเญฺห เต ปุฏฺโฐ,}

\csromanheader{2}{อนุปุพฺพํ อนุธมฺมํ พฺยากโรมิ เต. (2)}
\proseitem{517}{ปุจฺฉ มํ สภิย ปญฺหํ, ยํกิญฺจิ มนสิจฺฉสิ.}

\prose{ตสฺส ตสฺเสว ปญฺหสฺส, อหํ อนฺตํ กโรมิ เตติ. (3)}

\prose{อถ โข สภิยสฺส ปริพฺพาชกสฺส เอตทโหสิ “อจฺฉริยํ วต โภ, อพฺภุตํ วต โภ, ยํ วตาหํ อญฺเญสุ สมณพฺราหฺมเณสุ โอกาสกมฺมมตฺตมฺปิ\footnote{โอกาสมตฺตมฺปิ (สี, อิ)} นาลตฺถํ.\csromanspacer{} ตํ เม อิทํ สมเณน โคตเมน โอกาสกมฺมํ กตนฺ”ติ อตฺตมโน ปมุทิโต อุทคฺโค ปีติโสมนสฺสชาโต ภควนฺตํ ปญฺหํ อปุจฺฉิ\csromandash{}}

\proseitem{518}{กิํปตฺตินมาหุ ภิกฺขุนํ, (อิติ สภิโย,)}

\prose{โสรตํ เกน กถญฺจ ทนฺตมาหุ.}

\prose{พุทฺโธติ กถํ ปวุจฺจติ,}

\csromanheader{2}{ปุฏฺโฐ เม ภควา พฺยากโรหิ. (4)}
\proseitem{519}{ปชฺเชน กเตน อตฺตนา, (สภิยาติ ภควา,)}

\prose{ปรินิพฺพานคโต วิติณฺณกงฺโข.}

\prose{วิภวญฺจ ภวญฺจ วิปฺปหาย,}

\csromanheader{2}{วุสิตวา ขีณปุนพฺภโว ส ภิกฺขุ. (5)}
\gambhira{สุตฺตนิปาตปาฬิ}
\csromangathameasure{สพฺพตฺถ อุเปกฺขโก สติมา, \\ น โส หิํสติ กญฺจิ สพฺพโลเก. \\ ติณฺโณ สมโณ อนาวิโล,}
\csromangathagroup{\csromangathastanza{\csromanfolio{358}สพฺพตฺถ อุเปกฺขโก สติมา, \\ น โส หิํสติ กญฺจิ สพฺพโลเก. \\ ติณฺโณ สมโณ อนาวิโล,}}

\csromanheader{2}{อุสฺสทา ยสฺส น สนฺติ โสรโต โส. (6)}
\proseitem{521}{\csromansymbolfootnote{+}{ขุ 10. 147 ปิฏฺเฐปิ.}ยสฺสินฺทฺริยานิ ภาวิตานิ,}

\prose{อชฺฌตฺตํ พหิทฺธา จ สพฺพโลเก.}

\csromangathameasure{นิพฺพิชฺฌ อิมํ ปรญฺจ โลกํ,}
\csromangathagroup{\csromangathastanza{นิพฺพิชฺฌ อิมํ ปรญฺจ โลกํ,}}

\csromanheader{2}{กาลํ กงฺขติ ภาวิโต ส ทนฺโต. (7)}
\csromangathameasure{กปฺปานิ วิเจยฺย เกวลานิ, \\ สํสารํ ทุภยํ จุตูปปาตํ. \\ วิคตรชมนงฺคณํ วิสุทฺธํ,}
\csromangathagroup{\csromangathastanza{กปฺปานิ วิเจยฺย เกวลานิ, \\ สํสารํ ทุภยํ จุตูปปาตํ. \\ วิคตรชมนงฺคณํ วิสุทฺธํ,}}

\csromanheader{2}{ปตฺตํ ชาติขยํ ตมาหุ พุทฺธนฺติ. (8)}
\prose{อถ โข สภิโย ปริพฺพาชโก ภควโต ภาสิตํ อภินนฺทิตฺวา อนุโมทิตฺวา อตฺตมโน ปมุทิโต อุทคฺโค ปีติโสมนสฺสชาโต ภควนฺตํ อุตฺตริํ\footnote{อุตฺตริ (ก)} ปญฺหํ อปุจฺฉิ\csromandash{}}

\proseitem{523}{กิํ ปตฺตินมาหุ พฺราหฺมณํ, (อิติ สภิโย,)}

\prose{สมณํ เกน กถญฺจ นฺหาตโกติ.}

\csromangathameasure{นาโคติ กถํ ปวุจฺจติ,}
\csromangathagroup{\csromangathastanza{นาโคติ กถํ ปวุจฺจติ,}}

\csromanheader{2}{ปุฏฺโฐ เม ภควา พฺยากโรหิ. (9)}
\proseitem{524}{\csromansymbolfootnote{*}{ขุ 8. 84 ปิฏฺเฐปิ.}พาหิตฺวา สพฺพปาปกานิ, (สภิยาติ ภควา,)}

\prose{วิมโล สาธุสมาหิโต ฐิตตฺโต.}

\csromangathameasure{สํสารมติจฺจ เกวลี โส,}
\csromangathagroup{\csromangathastanza{สํสารมติจฺจ เกวลี โส,}}

\csromanheader{2}{อสิโต ตาทิ ปวุจฺจเต ส พฺรหฺมา. (10)}
\csromangathameasure{สมิตาวิ ปหาย ปุญฺญปาปํ, \\ วิรโช ญตฺวา อิมํ ปรญฺจ โลกํ. \\ ชาติมรณํ อุปาติวตฺโต,}
\csromangathagroup{\csromangathastanza{สมิตาวิ ปหาย ปุญฺญปาปํ, \\ วิรโช ญตฺวา อิมํ ปรญฺจ โลกํ. \\ ชาติมรณํ อุปาติวตฺโต,}}

\csromanheader{2}{สมโณ ตาทิ ปวุจฺจเต ตถตฺตา. (11)}
\chapterheadpageread{3. มหาวคฺค}
\csromangathameasure{นินฺหาย สพฺพปาปกานิ, \\ อชฺฌตฺตํ พหิทฺธา จ สพฺพโลเก. \\ เทวมนุสฺเสสุ กปฺปิเยสุ,}
\csromangathagroup{\csromangathastanza{\csromanfolio{359}นินฺหาย\footnote{นินหาย (สฺยา)} สพฺพปาปกานิ, \\ อชฺฌตฺตํ พหิทฺธา จ สพฺพโลเก. \\ เทวมนุสฺเสสุ กปฺปิเยสุ,}}

\csromanheader{2}{กปฺปํ เนติ ตมาหุ นฺหาตโกติ. (12)}
\proseitem{527}{\csromansymbolfootnote{*}{ขุ 8. 83 ปิฏฺเฐปิ.}อาคุํ น กโรติ กิญฺจิ โลเก,}

\prose{สพฺพสํโยเค\footnote{สพฺพโยเค (ก)} วิสชฺช พนฺธนานิ.}

\csromangathameasure{สพฺพตฺถ น สชฺชตี วิมุตฺโต,}
\csromangathagroup{\csromangathastanza{สพฺพตฺถ น สชฺชตี วิมุตฺโต,}}

\csromanheader{2}{นาโค ตาทิ ปวุจฺจเต ตถตฺตาติ. (13)}
\csromanheader{2}{อถ โข สภิโย ปริพฺพาชโก ฯเปฯ ภควนฺตํ อุตฺตริํ ปญฺหํ}
\csromanheader{2}{อปุจฺฉิ\csromandash{}}
\proseitem{528}{กํ เขตฺตชินํ วทนฺติ พุทฺธา, (อิติ สภิโย,)}

\prose{กุสลํ เกน กถญฺจ ปณฺฑิโตติ.}

\csromangathameasure{มุนิ นาม กถํ ปวุจฺจติ,}
\csromangathagroup{\csromangathastanza{มุนิ นาม กถํ ปวุจฺจติ,}}

\csromanheader{2}{ปุฏฺโฐ เม ภควา พฺยากโรหิ. (14)}
\proseitem{529}{เขตฺตานิ วิเจยฺย เกวลานิ, (สภิยาติ ภควา,)}

\prose{ทิพฺพํ มานุสกญฺจ พฺรหฺมเขตฺตํ.}

\csromangathameasure{สพฺพเขตฺตมูลพนฺธนา ปมุตฺโต,}
\csromangathagroup{\csromangathastanza{สพฺพเขตฺตมูลพนฺธนา ปมุตฺโต,}}

\csromanheader{2}{เขตฺตชิโน ตาทิ ปวุจฺจเต ตถตฺตา. (15)}
\csromangathameasure{โกสานิ วิเจยฺย เกวลานิ, \\ ทิพฺพํ มานุสกญฺจ พฺรหฺมโกสํ. \\ สพฺพโกสมูลพนฺธนา ปมุตฺโต,}
\csromangathagroup{\csromangathastanza{โกสานิ วิเจยฺย เกวลานิ, \\ ทิพฺพํ มานุสกญฺจ พฺรหฺมโกสํ. \\ สพฺพโกสมูลพนฺธนา ปมุตฺโต,}}

\csromanheader{2}{กุสโล ตาทิ ปวุจฺจเต ตถตฺตา. (16)}
\csromangathameasure{ทุภยานิ วิเจยฺย ปณฺฑรานิ, \\ อชฺฌตฺตํ พหิทฺธา จ สุทฺธิปญฺโญ. \\ กณฺหํ สุกฺกํ อุปาติวตฺโต,}
\csromangathagroup{\csromangathastanza{ทุภยานิ วิเจยฺย ปณฺฑรานิ, \\ อชฺฌตฺตํ พหิทฺธา จ สุทฺธิปญฺโญ. \\ กณฺหํ สุกฺกํ อุปาติวตฺโต,}}

\csromanheader{2}{ปณฺฑิโต ตาทิ ปวุจฺจเต ตถตฺตา. (17)}
\gambhira{สุตฺตนิปาตปาฬิ}
\proseitem{532}{\csromanfolio{360}\csromansymbolfootnote{*}{ขุ 8. 74 ปิฏฺเฐปิ.}อสตญฺจ สตญฺจ ญตฺวา ธมฺมํ,}

\prose{อชฺฌตฺตํ พหิทฺธา จ สพฺพโลเก.}

\csromangathameasure{เทวมนุสฺเสหิ ปูชนีโย,}
\csromangathagroup{\csromangathastanza{เทวมนุสฺเสหิ ปูชนีโย,}}

\csromanheader{2}{สงฺคํ ชาลมติจฺจ โส มุนีติ. (18)}
\prose{อถ โข สภิโย ปริพฺพาชโก ฯเปฯ ภควนฺตํ อุตฺตริํ ปญฺหํ อปุจฺฉิ\csromandash{}}

\proseitem{533}{กิํปตฺตินมาหุ เวทคุํ, (อิติ สภิโย,)}

\prose{อนุวิทิตํ เกน กถญฺจ วิริยวาติ.}

\csromangathameasure{อาชานิโย กินฺติ นาม โหติ,}
\csromangathagroup{\csromangathastanza{อาชานิโย กินฺติ นาม โหติ,}}

\csromanheader{2}{ปุฏฺโฐ เม ภควา พฺยากโรหิ. (19)}
\proseitem{534}{\csromansymbolfootnote{+}{ขุ 8. 64 ปิฏฺเฐปิ.}เวทานิ วิเจยฺย เกวลานิ, (สภิยาติ ภควา,)}

\prose{สมณานํ ยานิธตฺถิ\footnote{ยานิปตฺถิ (สี, สฺยา, อิ)} พฺราหฺมณานํ.}

\csromangathameasure{สพฺพเวทนาสุ วีตราโค,}
\csromangathagroup{\csromangathastanza{สพฺพเวทนาสุ วีตราโค,}}

\csromanheader{2}{สพฺพํ เวทมติจฺจ เวทคู โส. (20)}
\proseitem{535}{อนุวิจฺจ ปปญฺจนามรูปํ,}

\prose{อชฺฌตฺตํ พหิทฺธา จ โรคมูลํ.}

\csromangathameasure{สพฺพโรคมูลพนฺธนา ปมุตฺโต,}
\csromangathagroup{\csromangathastanza{สพฺพโรคมูลพนฺธนา ปมุตฺโต,}}

\csromanheader{2}{อนุวิทิโต ตาทิ ปวุจฺจเต ตถตฺตา. (21)}
\proseitem{536}{\csromansymbolmark{+}วิรโต อิธ สพฺพปาปเกหิ,}

\prose{นิรยทุกฺขํ อติจฺจ วิริยวา โส.}

\csromangathameasure{โส วีริยวา ปธานวา,}
\csromangathagroup{\csromangathastanza{โส วีริยวา ปธานวา,}}

\csromanheader{2}{ธีโร ตาทิ ปวุจฺจเต ตถตฺตา. (22)}
\csromangathameasure{ยสฺสสฺสุ ลุนานิ พนฺธนานิ, \\ อชฺฌตฺตํ พหิทฺธา จ สงฺคมูลํ. \\ สพฺพสงฺคมูลพนฺธนา ปมุตฺโต,}
\csromangathagroup{\csromangathastanza{ยสฺสสฺสุ ลุนานิ พนฺธนานิ, \\ อชฺฌตฺตํ พหิทฺธา จ สงฺคมูลํ. \\ สพฺพสงฺคมูลพนฺธนา ปมุตฺโต,}}

\csromanheader{2}{อาชานิโย ตาทิ ปวุจฺจเต ตถตฺตาติ. (23)}
\chapterheadpageread{3. มหาวคฺค}
\csromanheader{2}{อถ โข สภิโย ปริพฺพาชโก ฯเปฯ ภควนฺตํ อุตฺตริํ ปญฺหํ}
\csromanheader{2}{อปุจฺฉิ\csromandash{}}
\proseitem{538}{\csromanfolio{361}กิํปตฺตินมาหุ โสตฺติยํ, (อิติ สภิโย,)}

\prose{อริยํ เกน กถญฺจ จรณวาติ.}

\csromanmatikaanchor{mtk.267}
\csromantocmarkref{subsection}{5. มาฆสุตฺต}{mtk.267}
\bookmark[dest={mtk.267},level=2]{5. มาฆสุตฺต}
\prose{ปริพฺพาชโก กินฺติ นาม โหติ,}

\csromanheader{2}{ปุฏฺโฐ เม ภควา พฺยากโรหิ. (24)}
\proseitem{539}{สุตฺวา สพฺพธมฺมํ อภิญฺญาย โลเก, (สภิยาติ ภควา,)}

\prose{สาวชฺชานวชฺชํ ยทตฺถิ กิญฺจิ.}

\prose{อภิภุํ อกถํกถิํ วิมุตฺตํ,}

\csromanheader{2}{อนิฆํ สพฺพธิมาหุ โสตฺติโยติ. (25)}
\proseitem{540}{เฉตฺวา อาสวานิ อาลยานิ,}

\prose{วิทฺวา โส น อุเปติ คพฺภเสยฺยํ.}

\prose{สญฺญํ ติวิธํ ปนุชฺช ปงฺกํ,}

\csromanheader{2}{กปฺปํ เนติ ตมาหุ อริโยติ. (26)}
\proseitem{541}{โย อิธ จรเณสุ ปตฺติปตฺโต,}

\prose{กุสโล สพฺพทา อาชานาติ\footnote{อาชานิ (สฺยา)} ธมฺมํ.}

\prose{สพฺพตฺถ น สชฺชติ วิมุตฺตจิตฺโต\footnote{วิมุตฺโต (สี)},}

\csromanheader{2}{ปฏิฆา ยสฺส น สนฺติ จรณวา โส. (27)}
\proseitem{542}{ทุกฺขเวปกฺกํ ยทตฺถิ กมฺมํ,}

\prose{อุทฺธมโธ ติริยํ วาปิ\footnote{ติริยญฺจาปิ (สฺยา)} มชฺเฌ.}

\csromangathameasure{ปริพฺพาชยิตฺวา ปริญฺญจารี, \\ มายํ มานมโถปิ โลภโกธํ. \\ ปริยนฺตมกาสิ นามรูปํ,}
\csromangathagroup{\csromangathastanza{ปริพฺพาชยิตฺวา ปริญฺญจารี, \\ มายํ มานมโถปิ โลภโกธํ. \\ ปริยนฺตมกาสิ นามรูปํ,}}

\csromanheader{2}{ตํ ปริพฺพาชกมาหุ ปตฺติปตฺตนฺติ. (28)}
\prose{อถ โข สภิโย ปริพฺพาชโก ภควโต ภาสิตํ อภินนฺทิตฺวา อนุโมทิตฺวา อตฺตมโน ปมุทิโต อุทคฺโค ปีติโสมนสฺสชาโต อุฏฺฐายาสนา เอกํสํ อุตฺตราสงฺคํ กริตฺวา เยน ภควา เตนญฺชลิํ ปณาเมตฺวา ภควนฺตํ สมฺมุขา สารุปฺปาหิ คาถาหิ อภิตฺถวิ\csromandash{}}

\gambhira{สุตฺตนิปาตปาฬิ}
\proseitem{543}{\csromanfolio{362}ยานิ จ ตีณิ ยานิ จ สฏฺฐิ,}

\prose{สมณปฺปวาทสิตานิ\footnote{สมณปฺปวาทนิสฺสิตานิ (สฺยา, ก)} ภูริปญฺญ.}

\prose{สญฺญกฺขรสญฺญนิสฺสิตานิ,}

\csromanheader{2}{โอสรณานิ วิเนยฺย โอฆตม’คา. (29)}
\proseitem{544}{อนฺตคูสิ ปารคู\footnote{ปารคูสิ (สฺยา, อิ, ก)} ทุกฺขสฺส,}

\prose{อรหาสิ สมฺมาสมฺพุทฺโธ ขีณาสวํ ตํ มญฺเญ.}

\prose{ชุติมา มุติมา ปหูตปญฺโญ,}

\csromanheader{2}{ทุกฺขสฺสนฺตกร อตาเรสิ มํ. (30)}
\proseitem{545}{ยํ เม กงฺขิตมญฺญาสิ,}

\prose{วิจิกิจฺฉา มํ ตารยิ นโม เต.}

\prose{มุนิ โมนปเถสุ ปตฺติปตฺต,}

\csromanheader{2}{อขิล อาทิจฺจพนฺธุ โสรโตสิ. (31)}
\proseitem{546}{ยา เม กงฺขา ปุเร อาสิ, ตํ เม พฺยากาสิ จกฺขุมา.}

\prose{อทฺธา มุนีสิ สมฺพุทฺโธ, นตฺถิ นีวรณา ตว. (32)}

\proseitem{547}{อุปายาสา จ เต สพฺเพ, วิทฺธสฺตา วินฬีกตา.}

\prose{สีติภูโต ทมปฺปตฺโต, ธิติมา สจฺจนิกฺกโม. (33)}

\proseitem{548}{ตสฺส เต นาคนาคสฺส, มหาวีรสฺส ภาสโต.}

\prose{สพฺเพ เทวานุโมทนฺติ, อุโภ นารทปพฺพตา. (34)}

\proseitem{549}{นโม เต ปุริสาชญฺญ, นโม เต ปุริสุตฺตม.}

\prose{สเทวกสฺมิํ โลกสฺมิํ, นตฺถิ เต ปฏิปุคฺคโล. (35)}

\proseitem{550}{ตุวํ พุทฺโธ ตุวํ สตฺถา, ตุวํ มาราภิภู มุนิ.}

\prose{ตุวํ อนุสเย เฉตฺวา, ติณฺโณ ตาเรสิ’มํ ปชํ. (36)}

\proseitem{551}{อุปธี เต สมติกฺกนฺตา, อาสวา เต ปทาลิตา.}

\prose{สีโหสิ อนุปาทาโน, ปหีนภยเภรโว. (37)}

\chapterheadpageread{3. มหาวคฺค}
\proseitem{552}{\csromanfolio{363}ปุณฺฑรีกํ ยถา วคฺคุ, โตเย น อุปลิมฺปติ\footnote{โตเยน น อุปลิปฺปติ (สี), โตเย น อุปลิปฺปติ (อิ), โตเยน น อุปลิมฺปติ (ก)}.}

\csromangathameasure{เอวํ ปุญฺเญ จ ปาเป จ, อุภเย ตฺวํ น ลิมฺปสิ.}
\csromangathagroup{\csromangathastanza{เอวํ ปุญฺเญ จ ปาเป จ, อุภเย ตฺวํ น ลิมฺปสิ.}}

\prose{ปาเท วีร ปสาเรหิ, สภิโย วนฺทติ สตฺถุโนติ. (38)}

\prose{อถ โข สภิโย ปริพฺพาชโก ภควโต ปาเทสุ สิรสา นิปติตฺวา ภควนฺตํ เอตทโวจ “อภิกฺกนฺตํ ภนฺเต ฯเปฯ เอสาหํ ภควนฺตํ สรณํ คจฺฉามิ ธมฺมญฺจ ภิกฺขุสํฆญฺจ, ลเภยฺยาหํ ภนฺเต ภควโต สนฺติเก ปพฺพชฺชํ, ลเภยฺยํ อุปสมฺปทนฺ”ติ.}

\prose{โย โข สภิย อญฺญติตฺถิยปุพฺโพ อิมสฺมิํ ธมฺมวินเย อากงฺขติ ปพฺพชฺชํ, อากงฺขติ อุปสมฺปทํ, โส จตฺตาโร มาเส ปริวสติ, จตุนฺนํ มาสานํ อจฺจเยน อารทฺธจิตฺตา ภิกฺขู ปพฺพาเชนฺติ, อุปสมฺปาเทนฺติ ภิกฺขุภาวาย.\csromanspacer{} อปิ จ เมตฺถ ปุคฺคลเวมตฺตตา วิทิตาติ.}

\prose{สเจ ภนฺเต อญฺญติตฺถิยปุพฺพา อิมสฺมิํ ธมฺมวินเย อากงฺขนฺตา ปพฺพชฺชํ อากงฺขนฺตา อุปสมฺปทํ จตฺตาโร มาเส ปริวสนฺติ, จตุนฺนํ มาสานํ อจฺจเยน อารทฺธจิตฺตา ภิกฺขู ปพฺพาเชนฺติ, อุปสมฺปาเทนฺติ ภิกฺขุภาวาย, อหํ จตฺตาริ วสฺสานิ ปริวสิสฺสามิ, จตุนฺนํ วสฺสานํ อจฺจเยน อารทฺธจิตฺตา ภิกฺขู ปพฺพาเชนฺตุ, อุปสมฺปาเทนฺตุ ภิกฺขุภาวายาติ.\csromanspacer{} อลตฺถ โข สภิโย ปริพฺพาชโก ภควโต สนฺติเก ปพฺพชฺชํ, อลตฺถ อุปสมฺปทํ ฯเปฯ.\csromanspacer{} อญฺญตโร โข ปนายสฺมา สภิโย อรหตํ อโหสีติ.\csromanspacer{} สภิยสุตฺตํ ฉฏฺฐํ.}

\csromanheader{2}{7. เสลสุตฺต}
\prose{\csromansymbolfootnote{*}{ม 2. 347 ปิฏฺเฐปิ.}เอวํ เม สุตํ\csromandash{}เอกํ สมยํ ภควา องฺคุตฺตราเปสุ จาริกํ จรมาโน มหตา ภิกฺขุสํเฆน สทฺธิํ อฑฺฒเตฬเสหิ ภิกฺขุสเตหิ เยน อาปณํ นาม องฺคุตฺตราปานํ นิคโม ตทวสริ.\csromanspacer{} อสฺโสสิ โข เกณิโย ชฏิโล “สมโณ ขลุ โภ โคตโม สกฺยปุตฺโต สกฺยกุลา ปพฺพชิโต องฺคุตฺตราเปสุ จาริกํ จรมาโน มหตา ภิกฺขุสํเฆน สทฺธิํ อฑฺฒเตฬเสหิ ภิกฺขุสเตหิ อาปณํ อนุปฺปตฺโต, ตํ โข ปน ภวนฺตํ}

\gambhira{สุตฺตนิปาตปาฬิ}
\prose{\csromanfolio{364}โคตมํ เอวํ กลฺยาโณ กิตฺติสทฺโท อพฺภุคฺคโต ‘อิติปิ โส ภควา อรหํ สมฺมาสมฺพุทฺโธ วิชฺชาจรณสมฺปนฺโน สุคโต โลกวิทู อนุตฺตโร ปุริสทมฺมสารถิ สตฺถา เทวมนุสฺสานํ พุทฺโธ ภควา’ติ\footnote{ภควา (สฺยา,อิ)}.\csromanspacer{} โส อิมํ โลกํ สเทวกํ สมารกํ สพฺรหฺมกํ สสฺสมณพฺราหฺมณิํ ปชํ สเทวมนุสฺสํ สยํ อภิญฺญา สจฺฉิกตฺวา ปเวเทติ, โส ธมฺมํ เทเสติ อาทิกลฺยาณํ มชฺเฌกลฺยาณํ ปริโยสานกลฺยาณํ สาตฺถํ สพฺยญฺชนํ เกวลปริปุณฺณํ ปริสุทฺธํ พฺรหฺมจริยํ ปกาเสติ, สาธุ โข ปน ตถารูปานํ อรหตํ ทสฺสนํ โหตี”ติ.}

\prose{อถ โข เกณิโย ชฏิโล เยน ภควา เตนุปสงฺกมิ, อุปสงฺกมิตฺวา ภควตา สทฺธิํ สมฺโมทิ, สมฺโมทนียํ กถํ สารณียํ วีติสาเรตฺวา เอกมนฺตํ นิสีทิ, เอกมนฺตํ นิสินฺนํ โข เกณิยํ ชฏิลํ ภควา ธมฺมิยา กถาย สนฺทสฺเสสิ สมาทเปสิ สมุตฺเตเชสิ สมฺปหํเสสิ.\csromanspacer{} อถ โข เกณิโย ชฏิโล ภควตา ธมฺมิยา กถาย สนฺทสฺสิโต สมาทปิโต สมุตฺเตชิโต สมฺปหํสิโต ภควนฺตํ เอตทโวจ “อธิวาเสตุ เม ภวํ โคตโม สฺวาตนาย ภตฺตํ สทฺธิํ ภิกฺขุสํเฆนา”ติ.\csromanspacer{} เอวํ วุตฺเต ภควา เกณิยํ ชฏิลํ เอตทโวจ “มหา โข เกณิย ภิกฺขุสํโฆ อฑฺฒเตฬสานิ ภิกฺขุสตานิ, ตฺวญฺจ พฺราหฺมเณสุ อภิปฺปสนฺโน”ติ.}

\prose{ทุติยมฺปิ โข เกณิโย ชฏิโล ภควนฺตํ เอตทโวจ “กิญฺจาปิ โภ โคตม มหา ภิกฺขุสํโฆ อฑฺฒเตฬสานิ ภิกฺขุสตานิ, อหญฺจ พฺราหฺมเณสุ อภิปฺปสนฺโน, อธิวาเสตุ เม ภวํ โคตโม สฺวาตนาย ภตฺตํ สทฺธิํ ภิกฺขุสํเฆนา”ติ.\csromanspacer{} ทุติยมฺปิ โข ภควา เกณิยํ ชฏิลํ เอตทโวจ “มหา โข เกณิย ภิกฺขุสํโฆ อฑฺฒเตฬสานิ ภิกฺขุสตานิ, ตฺวญฺจ พฺราหฺมเณสุ อภิปฺปสนฺโน”ติ.}

\prose{ตติยมฺปิ โข เกณิโย ชฏิโล ภควนฺตํ เอตทโวจ “กิญฺจาปิ โภ โคตม มหา ภิกฺขุสํโฆ อฑฺฒเตฬสานิ ภิกฺขุสตานิ, อหญฺจ พฺราหฺมเณสุ อภิปฺปสนฺโน, อธิวาเสตุ\footnote{อธิวาเสเตฺวว (สี)} เม ภวํ โคตโม สฺวาตนาย ภตฺตํ สทฺธิํ ภิกฺขุสํเฆนา”ติ.\csromanspacer{} อธิวาเสสิ ภควา ตุณฺหีภาเวน.\csromanspacer{} อถ โข เกณิโย ชฏิโล ภควโต อธิวาสนํ}

\vspace{\tipitakaparskip}
\csromancenter{มหาวคฺค วิทิตฺวา อุฏฺฐายาสนา เยน สโก อสฺสโม เตนุปสงฺกมิ, อุปสงฺกมิตฺวา มิตฺตามจฺเจ ญาติสาโลหิเต อามนฺเตสิ “สุณนฺตุ เม ภวนฺโต มิตฺตามจฺจา ญาติสาโลหิตา, สมโณ เม โคตโม นิมนฺติโต สฺวาตนาย ภตฺตํ สทฺธิํ ภิกฺขุสํเฆน, เยน เม กายเวยฺยาวฏิกํ กเรยฺยาถา”ติ. “เอวํ โภ”ติ โข เกณิยสฺส ชฏิลสฺส มิตฺตามจฺจา ญาติสาโลหิตา เกณิยสฺส ชฏิลสฺส ปฏิสฺสุตฺวา อปฺเปกจฺเจ อุทฺธนานิ ขณนฺติ, อปฺเปกจฺเจ กฏฺฐานิ ผาเลนฺติ, อปฺเปกจฺเจ ภาชนานิ โธวนฺติ, อปฺเปกจฺเจ อุทกมณิกํ ปติฏฺฐาเปนฺติ, อปฺเปกจฺเจ อาสนานิ ปญฺญาเปนฺติ.\csromanspacer{} เกณิโย ปน ชฏิโล สามํเยว มณฺฑลมาฬํ ปฏิยาเทติ.}
\prose{\csromanfolio{365}เตน โข ปน สมเยน เสโล พฺราหฺมโณ อาปเณ ปฏิวสติ, ติณฺณํ เวทานํ ปารคู สนิฆณฺฑุเกฏุภานํ สากฺขรปฺปเภทานํ อิติหาสปญฺจมานํ ปทโก เวยฺยากรโณ โลกายตมหาปุริสลกฺขเณสุ อนวโย, ตีณิ จ มาณวกสตานิ มนฺเต วาเจติ.}

\prose{เตน โข ปน สมเยน เกณิโย ชฏิโล เสเล พฺราหฺมเณ อภิปฺปสนฺโน โหติ.\csromanspacer{} อถ โข เสโล พฺราหฺมโณ ตีหิ มาณวกสเตหิ ปริวุโต ชงฺฆาวิหารํ อนุจงฺกมมาโน อนุวิจรมาโน เยน เกณิยสฺส ชฏิลสฺส อสฺสโม เตนุปสงฺกมิ.\csromanspacer{} อทฺทสา โข เสโล พฺราหฺมโณ เกณิยสฺส ชฏิลสฺส อสฺสเม\footnote{เกณิยสฺสมิเย ชฏิเล (สี, อิ) ม 2. 349 ปิฏฺเฐปิ.} อปฺเปกจฺเจ อุทฺธนานิ ขณนฺเต ฯเปฯ อปฺเปกจฺเจ อาสนานิ ปญฺญาเปนฺเต, เกณิยํ ปน ชฏิลํ สามํเยว มณฺฑลมาฬํ ปฏิยาเทนฺตํ, ทิสฺวาน เกณิยํ ชฏิลํ เอตทโวจ “กิํ นุ โข โภโต เกณิยสฺส อาวาโห วา ภวิสฺสติ, วิวาโห วา ภวิสฺสติ, มหายญฺโญ วา ปจฺจุปฏฺฐิโต, ราชา วา มาคโธ เสนิโย พิมฺพิสาโร นิมนฺติโต สฺวาตนาย สทฺธิํ พลกาเยนา”ติ.}

\prose{น เม โภ เสล อาวาโห วา ภวิสฺสติ วิวาโห วา, นาปิ ราชา มาคโธ เสนิโย พิมฺพิสาโร นิมนฺติโต สฺวาตนาย สทฺธิํ พลกาเยน, อปิ จ โข เม มหายญฺโญ ปจฺจุปฏฺฐิโต, อตฺถิ สมโณ โคตโม สกฺยปุตฺโต สกฺยกุลา ปพฺพชิโต องฺคุตฺตราเปสุ}

\gambhira{สุตฺตนิปาตปาฬิ}
\prose{\csromanfolio{366}จาริกํ จรมาโน มหตา ภิกฺขุสํเฆน สทฺธิํ อฑฺฒเตฬเสหิ ภิกฺขุสเตหิ อาปณํ อนุปฺปตฺโต, ตํ โข ปน ภวนฺตํ โคตมํ ฯเปฯ พุทฺโธ ภควาติ.\csromanspacer{} โส เม นิมนฺติโต สฺวาตนาย ภตฺตํ สทฺธิํ ภิกฺขุสํเฆนา”ติ. “พุทฺโธ”ติ โภ เกณิย วเทสิ,”พุทฺโธ”ติ โภ เสล วทามิ. “พุทฺโธ”ติ โภ เกณิย วเทสิ, “พุทฺโธ”ติ โภ เสล วทามีติ.}

\prose{อถ โข เสลสฺส พฺราหฺมณสฺส เอตทโหสิ “โฆโสปิ โข เอโส ทุลฺลโภ โลกสฺมิํ ยทิทํ พุทฺโธติ, อาคตานิ โข ปนมฺหากํ มนฺเตสุ ทฺวตฺติํสมหาปุริสลกฺขณานิ, เยหิ สมนฺนาคตสฺส มหาปุริสสฺส เทฺวว คติโย ภวนฺติ อนญฺญา\csromandash{}สเจ อคารํ อชฺฌาวสติ ราชา โหติ จกฺกวตฺตี ธมฺมิโก ธมฺมราชา จาตุรนฺโต วิชิตาวี ชนปทตฺถาวริยปฺปตฺโต สตฺตรตนสมนฺนาคโต.\csromanspacer{} ตสฺสิมานิ สตฺต รตนานิ ภวนฺติ.\csromanspacer{} เสยฺยถิทํ, จกฺกรตนํ หตฺถิรตนํ อสฺสรตนํ มณิรตนํ อิตฺถิรตนํ คหปติรตนํ ปริณายกรตนเมว สตฺตมํ.\csromanspacer{} ปโรสหสฺสํ โข ปนสฺส ปุตฺตา ภวนฺติ สูรา วีรงฺครูปา ปรเสนปฺปมทฺทนา, โส อิมํ ปถวิํ สาครปริยนฺตํ อทณฺเฑน อสตฺเถน ธมฺเมน อภิวิชิย อชฺฌาวสติ.\csromanspacer{} สเจ โข ปน อคารสฺมา อนคาริยํ ปพฺพชติ, อรหํ โหติ สมฺมาสมฺพุทฺโธ โลเก วิวฏฺฏจฺฉโท\footnote{วิวตฺตจฺฉทฺโท (สี, อิ)}”.\csromanspacer{} กหํ ปน โภ เกณิย เอตรหิ โส ภวํ โคตโม วิหรติ อรหํ สมฺมาสมฺพุทฺโธติ.}

\prose{เอวํ วุตฺเต เกณิโย ชฏิโล ทกฺขิณํ พาหุํ ปคฺคเหตฺวา เสลํ พฺราหฺมณํ เอตทโวจ “เยเนสา โภ เสล นีลวนราชี”ติ.\csromanspacer{} อถ โข เสโล พฺราหฺมโณ ตีหิ มาณวกสเตหิ สทฺธิํ เยน ภควา เตนุปสงฺกมิ, อถ โข เสโล พฺราหฺมโณ เต มาณวเก อามนฺเตสิ “อปฺปสทฺทา โภนฺโต อาคจฺฉนฺตุ ปเท ปทํ นิกฺขิปนฺตา, ทุราสทา หิ เต ภควนฺโต\footnote{ภวนฺโต (สฺยา, ก)} สีหาว เอกจรา.\csromanspacer{} ยทา จาหํ โภ สมเณน โคตเมน สทฺธิํ มนฺเตยฺยํ, มา เม โภนฺโต อนฺตรนฺตรา กถํ โอปาเตถ, กถาปริโยสานํ เม ภวนฺโต อาคเมนฺตูติ.}

\prose{อถ โข เสโล พฺราหฺมโณ เยน ภควา เตนุปสงฺกมิ, อุปสงฺกมิตฺวา ภควตา สทฺธิํ สมฺโมทิ, สมฺโมทนียํ กถํ สารณียํ วีติสาเรตฺวา เอกมนฺตํ นิสีทิ, เอกมนฺตํ นิสินฺโน โข เสโล พฺราหฺมโณ}

\vspace{\tipitakaparskip}
\csromancenter{มหาวคฺค ภควโต กาเย ทฺวตฺติํสมหาปุริสลกฺขณานิ สมนฺเนสิ\footnote{สมฺมนฺเนสิ (สี, สฺยา)}.\csromanspacer{} อทฺทสา โข เสโล พฺราหฺมโณ ภควโต กาเย ทฺวตฺติํสมหาปุริสลกฺขณานิ เยภุยฺเยน ฐเปตฺวา เทฺว.\csromanspacer{} ทฺวีสุ มหาปุริสลกฺขเณสุ กงฺขติ วิจิกิจฺฉติ นาธิมุจฺจติ น สมฺปสีทติ\csromandash{}โกโสหิเต จ วตฺถคุเยฺห, ปหูตชิวฺหตาย จาติ.}
\prose{\csromanfolio{367}อถ โข ภควโต เอตทโหสิ “ปสฺสติ โข เม อยํ เสโล พฺราหฺมโณ ทฺวตฺติํสมหาปุริสลกฺขณานิ เยภุยฺเยน ฐเปตฺวา เทฺว.\csromanspacer{} ทฺวีสุ มหาปุริสลกฺขเณสุ กงฺขติ วิจิกิจฺฉติ นาธิมุจฺจติ น สมฺปสีทติ\csromandash{} โกโสหิเต จ วตฺถคุเยฺห, ปหูตชิวฺหตาย จา”ติ.\csromanspacer{} อถ โข ภควา ตถารูปํ อิทฺธาภิสงฺขารํ อภิสงฺขาสิ\footnote{อภิสงฺขาเรสิ (สฺยา, ก) ที 1. 101 ปิฏฺเฐปิ.}, ยถา อทฺทส เสโล พฺราหฺมโณ ภควโต โกโสหิตํ วตฺถคุยฺหํ.\csromanspacer{} อถ โข ภควา ชิวฺหํ นิเนเมตฺวา อุโภปิ กณฺณโสตานิ อนุมสิ ปฏิมสิ, อุโภปิ นาสิกโสตานิ อนุมสิ ปฏิมสิ, เกวลมฺปิ นลาฏมณฺฑลํ ชิวฺหาย ฉาเทสิ.}

\prose{อถ โข เสลสฺส พฺราหฺมณสฺส เอตทโหสิ “สมนฺนาคโต โข สมโณ โคตโม ทฺวตฺติํสมหาปุริสลกฺขเณหิ ปริปุณฺเณหิ, โน อปริปุณฺเณหิ.\csromanspacer{} โน จ โข นํ ชานามิ ‘พุทฺโธ วา โน วา’, สุตํ โข ปน เมตํ พฺราหฺมณานํ วุฑฺฒานํ มหลฺลกานํ อาจริยปาจริยานํ ภาสมานานํ ‘เย เต ภวนฺติ อรหนฺโต สมฺมาสมฺพุทฺธา, เต สเก วณฺเณ ภญฺญมาเน อตฺตานํ ปาตุกโรนฺตี’ติ.\csromanspacer{} ยํนูนาหํ สมณํ โคตมํ สมฺมุขา สารุปฺปาหิ คาถาหิ อภิตฺถเวยฺยนฺ”ติ.\csromanspacer{} อถ โข เสโล พฺราหฺมโณ ภควนฺตํ สมฺมุขา สารุปฺปาหิ คาถาหิ อภิตฺถวิ\csromandash{}}

\proseitem{553}{\csromansymbolfootnote{*}{ขุ 2. 330 ปิฏฺเฐปิ.}ปริปุณฺณกาโย สุรุจิ, สุชาโต จารุทสฺสโน.}

\vspace{\tipitakaparskip}
\csromancenter{สุวณฺณวณฺโณสิ ภควา, สุสุกฺกทาโฐสิ วีริยวา. (1)}
\proseitem{554}{นรสฺส หิ สุชาตสฺส, เย ภวนฺติ วิยญฺชนา.}

\prose{สพฺเพ เต ตว กายสฺมิํ, มหาปุริสลกฺขณา. (2)}

\csromangathameasure{ปสนฺนเนตฺโต สุมุโข, พฺรหา อุชุ ปตาปวา.}
\csromangathagroup{\csromangathastanza{ปสนฺนเนตฺโต สุมุโข, พฺรหา อุชุ ปตาปวา.}}

\proseitem{555}{มชฺเฌ สมณสํฆสฺส, อาทิจฺโจว วิโรจสิ. (3)}

\gambhira{สุตฺตนิปาตปาฬิ}
\proseitem{556}{\csromanfolio{368}กลฺยาณทสฺสโน ภิกฺขุ, กญฺจนสนฺนิภตฺตโจ.}

\prose{กิํ เต สมณภาเวน, เอวํ อุตฺตมวณฺณิโน. (4)}

\proseitem{557}{ราชา อรหสิ ภวิตุํ, จกฺกวตฺตี รเถสโภ.}

\prose{จาตุรนฺโต วิชิตาวี, ชมฺพุสณฺฑสฺส\footnote{ชมฺพุมณฺฑสฺส (ก)} อิสฺสโร. (5)}

\proseitem{558}{ขตฺติยา โภคิราชาโน\footnote{โภชราชาโน (สี, สฺยา)}, อนุยนฺตา\footnote{อนุยุตฺตา (สี)} ภวนฺตุ เต.}

\prose{ราชาภิราชา มนุชินฺโท, รชฺชํ กาเรหิ โคตม. (6)}

\proseitem{559}{ราชาหมสฺมิ เสลาติ, (ภควา)}

\prose{ธมฺมราชา อนุตฺตโร.}

\prose{ธมฺเมน จกฺกํ วตฺเตมิ,}

\csromanheader{2}{จกฺกํ อปฺปฏิวตฺติยํ. (7)}
\proseitem{560}{สมฺพุทฺโธ ปฏิชานาสิ, (อิติ เสโล พฺราหฺมโณ,)}

\prose{ธมฺมราชา อนุตฺตโร.}

\prose{ธมฺเมน จกฺกํ วตฺเตมิ,}

\csromanheader{2}{อิติ ภาสสิ โคตม. (8)}
\proseitem{561}{โก นุ เสนาปติ โภโต, สาวโก สตฺถุรนฺวโย.}

\prose{โก เต ตมนุวตฺเตติ, ธมฺมจกฺกํ ปวตฺติตํ. (9)}

\proseitem{562}{มยา ปวตฺติตํ จกฺกํ, (เสลาติ ภควา,)}

\prose{ธมฺมจกฺกํ อนุตฺตรํ.}

\prose{สาริปุตฺโต อนุวตฺเตติ,}

\csromanheader{2}{อนุชาโต ตถาคตํ. (10)}
\proseitem{563}{อภิญฺเญยฺยํ อภิญฺญาตํ, ภาเวตพฺพญฺจ ภาวิตํ.}

\prose{ปหาตพฺพํ ปหีนํ เม, ตสฺมา พุทฺโธสฺมิ พฺราหฺมณ. (11)}

\proseitem{564}{วินยสฺสุ มยิ กงฺขํ, อธิมุจฺจสฺสุ พฺราหฺมณ.}

\prose{ทุลฺลภํ ทสฺสนํ โหติ, สมฺพุทฺธานํ อภิณฺหโส. (12)}

\chapterheadpageread{3. มหาวคฺค}
\proseitem{565}{\csromanfolio{369}เยสํ เว\footnote{เยสํ โว (อิ), ยสฺส เว (สฺยา)} ทุลฺลโภ โลเก,}

\prose{ปาตุภาโว อภิณฺหโส.}

\prose{โสหํ พฺราหฺมณ สมฺพุทฺโธ,}

\csromanheader{2}{สลฺลกตฺโต อนุตฺตโร. (13)}
\proseitem{566}{พฺรหฺมภูโต อติตุโล, มารเสนปฺปมทฺทโน.}

\prose{สพฺพามิตฺเต วสีกตฺวา, โมทามิ อกุโตภโย. (14)}

\proseitem{567}{อิมํ ภวนฺโต นิสาเมถ, ยถา ภาสติ จกฺขุมา.}

\prose{สลฺลกตฺโต มหาวีโร, สีโหว นทตี วเน. (15)}

\csromanmatikaanchor{mtk.268}
\csromantocmarkref{subsection}{6. สภิยสุตฺต}{mtk.268}
\bookmark[dest={mtk.268},level=2]{6. สภิยสุตฺต}
\proseitem{568}{พฺรหฺมภูตํ อติตุลํ, มารเสนปฺปมทฺทนํ.}

\prose{โก ทิสฺวา นปฺปสีเทยฺย, อปิ กณฺหาภิชาติโก. (16)}

\proseitem{569}{โย มํ อิจฺฉติ อเนฺวตุ, โย วา นิจฺฉติ คจฺฉตุ.}

\prose{อิธาหํ ปพฺพชิสฺสามิ, วรปญฺญสฺส สนฺติเก. (17)}

\proseitem{570}{เอวญฺเจ\footnote{เอตญฺเจ (สี, อิ)} รุจฺจติ โภโต, สมฺมาสมฺพุทฺธสาสเน\footnote{สมฺมาสมฺพุทฺธสาสนํ (สี, สฺยา, กํ, อิ)}.}

\prose{มยมฺปิ ปพฺพชิสฺสาม, วรปญฺญสฺส สนฺติเก. (18)}

\proseitem{571}{พฺราหฺมณา ติสตา อิเม, ยาจนฺติ ปญฺชลีกตา.}

\prose{พฺรหฺมจริยํ จริสฺสาม, ภควา ตว สนฺติเก. (19)}

\proseitem{572}{สฺวากฺขาตํ พฺรหฺมจริยํ, (เสลาติ ภควา)}

\prose{สนฺทิฏฺฐิกมกาลิกํ.}

\prose{ยตฺถ อโมฆา ปพฺพชฺชา,}

\csromanheader{2}{อปฺปมตฺตสฺส สิกฺขโตติ. (20)}
\prose{อลตฺถ โข เสโล พฺราหฺมโณ สปริโส ภควโต สนฺติเก ปพฺพชฺชํ, อลตฺถ อุปสมฺปทํ.\csromanspacer{} อถ โข เกณิโย ชฏิโล ตสฺสา รตฺติยา อจฺจเยน สเก อสฺสเม ปณีตํ ขาทนียํ โภชนียํ ปฏิยาทาเปตฺวา ภควโต กาลํ อาโรจาเปสิ “กาโล โภ โคตม นิฏฺฐิตํ ภตฺตนฺ”ติ.\csromanspacer{} อถ โข ภควา ปุพฺพณฺหสมยํ นิวาเสตฺวา ปตฺตจีวรมาทาย เยน เกณิยสฺส ชฏิลสฺส อสฺสโม เตนุปสงฺกมิ, อุปสงฺกมิตฺวา ปญฺญตฺเต อาสเน นิสีทิ สทฺธิํ ภิกฺขุสํเฆน.}

\gambhira{สุตฺตนิปาตปาฬิ}
\prose{\csromanfolio{370}อถ โข เกณิโย ชฏิโล พุทฺธปฺปมุขํ ภิกฺขุสํฆํ ปณีเตน ขาทนีเยน โภชนีเยน สหตฺถา สนฺตปฺเปสิ สมฺปวาเรสิ.\csromanspacer{} อถ โข เกณิโย ชฏิโล ภควนฺตํ ภุตฺตาวิํ โอนีตปตฺตปาณิํ อญฺญตรํ นีจํ อาสนํ คเหตฺวา เอกมนฺตํ นิสีทิ.\csromanspacer{} เอกมนฺตํ นิสินฺนํ โข เกณิยํ ชฏิลํ ภควา อิมาหิ คาถาหิ อนุโมทิ\csromandash{}}

\proseitem{573}{“อคฺคิหุตฺตมุขา ยญฺญา, สาวิตฺตี ฉนฺทโส มุขํ.}

\prose{ราชา มุขํ มนุสฺสานํ, นทีนํ สาคโร มุขํ. (21)}

\csromanhangingbegin
\hangingheaditem{574}{นกฺขตฺตานํ มุขํ จนฺโท, อาทิจฺโจ ตปตํ มุขํ.}
\hangingbody{ปุญฺญํ อากงฺขมานานํ, สํโฆ เว ยชตํ มุขนฺ”ติ. (22)}
\csromanhangingend

\prose{อถ โข ภควา เกณิยํ ชฏิลํ อิมาหิ คาถาหิ อนุโมทิตฺวา อุฏฺฐายาสนา ปกฺกามิ.\csromanspacer{} อถ โข อายสฺมา เสโล สปริโส เอโก วูปกฏฺโฐ อปฺปมตฺโต อาตาปี ปหิตตฺโต วิหรนฺโต นจิรสฺเสว ฯเปฯ.\csromanspacer{} อญฺญตโร โข ปนายสฺมา เสโล สปริโส อรหตํ อโหสิ.}

\prose{อถ โข อายสฺมา เสโล สปริโส เยน ภควา เตนุปสงฺกมิ, อุปสงฺกมิตฺวา เอกํสํ จีวรํ กตฺวา เยน ภควา เตนญฺชลิํ ปณาเมตฺวา ภควนฺตํ คาถาย อชฺฌภาสิ\csromandash{}}

\proseitem{575}{“ยํ ตํ สรณมาคมฺห\footnote{มาคมฺม (สี, สฺยา, ก); ม 2. 354 ปิฏฺเฐปิ.}, อิโต อฏฺฐมิ จกฺขุม.}

\prose{สตฺตรตฺเตน ภควา, ทนฺตมฺห ตว สาสเน. (23)}

\proseitem{576}{ตุวํ พุทฺโธ ตุวํ สตฺถา, ตุวํ มาราภิภู มุนิ.}

\prose{ตุวํ อนุสเย เฉตฺวา, ติณฺโณ ตาเรสิ’มํ ปชํ. (24)}

\proseitem{577}{อุปธี เต สมติกฺกนฺตา, อาสวา เต ปทาลิตา.}

\prose{สีโหสิ\footnote{สีโหว (ม 2. 354 ปิฏฺเฐ.)} อนุปาทาโน, ปหีนภยเภรโว. (25)}

\proseitem{578}{ภิกฺขโว ติสตา อิเม, ติฏฺฐนฺติ ปญฺชลีกตา.}

\csromangathameasure{ปาเท วีร ปสาเรหิ, นาคา วนฺทนฺตุ สตฺถุโน”ติ. (26) เสลสุตฺตํ สตฺตมํ.}
\csromangathagroup{\csromangathastanza{ปาเท วีร ปสาเรหิ, นาคา วนฺทนฺตุ สตฺถุโน”ติ. (26) เสลสุตฺตํ สตฺตมํ.}}

\csromanheader{2}{3. มหาวคฺค \textbf{8. สลฺลสุตฺต}}
\proseitem{579}{\csromanfolio{371}อนิมิตฺตมนญฺญาตํ, มจฺจานํ อิธ ชีวิตํ.}

\prose{กสิรญฺจ ปริตฺตญฺจ, ตญฺจ ทุกฺเขน สํยุตํ. (1)}

\proseitem{580}{น หิ โส อุปกฺกโม อตฺถิ, เยน ชาตา น มิยฺยเร.}

\prose{ชรมฺปิ ปตฺวา มรณํ, เอวํธมฺมา หิ ปาณิโน. (2)}

\proseitem{581}{ผลานมิว ปกฺกานํ, ปาโต ปตนโต\footnote{ปปตโต (สี, อิ, ฏฺฐ)} ภยํ.}

\prose{เอวํ ชาตาน มจฺจานํ, นิจฺจํ มรณโต ภยํ. (3)}

\proseitem{582}{ยถาปิ กุมฺภการสฺส, กตา มตฺติกภาชนา.}

\prose{สพฺเพ เภทนปริยนฺตา\footnote{เภทปริยนฺตา (สฺยา)}, เอวํ มจฺจาน ชีวิตํ. (4)}

\proseitem{583}{ทหรา จ มหนฺตา จ, เย พาลา เย จ ปณฺฑิตา.}

\prose{สพฺเพ มจฺจุวสํ ยนฺติ, สพฺเพ มจฺจุปรายณา. (5)}

\proseitem{584}{เตสํ มจฺจุปเรตานํ, คจฺฉตํ ปรโลกโต.}

\prose{น ปิตา ตายเต ปุตฺตํ, ญาตี วา ปน ญาตเก. (6)}

\proseitem{585}{เปกฺขตํเยว ญาตีนํ, ปสฺส ลาลปตํ ปุถุ.}

\prose{เอกเมโก ว มจฺจานํ, โควชฺโฌ วิย นียติ\footnote{นิยฺยติ (พหูสุ)}. (7)}

\proseitem{586}{เอวมพฺภาหโต โลโก, มจฺจุนา จ ชราย จ.}

\prose{ตสฺมา ธีรา น โสจนฺติ, วิทิตฺวา โลกปริยายํ. (8)}

\proseitem{587}{ยสฺส มคฺคํ น ชานาสิ, อาคตสฺส คตสฺส วา.}

\prose{อุโภ อนฺเต อสมฺปสฺสํ, นิรตฺถํ ปริเทวสิ. (9)}

\csromanhangingbegin
\hangingheaditem{588}{ปริเทวยมาโน เจ, กิญฺจิทตฺถํ อุทพฺพเห.}
\hangingbody{สมฺมูโฬฺห หิํสมตฺตานํ, กยิรา เจ นํ วิจกฺขโณ.  (10)}
\csromanhangingend

\proseitem{589}{น หิ รุณฺเณน โสเกน, สนฺติํ ปปฺโปติ เจตโส.}

\prose{ภิยฺยสฺสุปฺปชฺชเต ทุกฺขํ, สรีรํ จุปหญฺญติ. (11)}

\proseitem{590}{กิโส วิวณฺโณ ภวติ, หิํสมตฺตานมตฺตนา.}

\prose{น เตน เปตา ปาเลนฺติ, นิรตฺถา ปริเทวนา. (12)}

\gambhira{สุตฺตนิปาตปาฬิ}
\proseitem{591}{\csromanfolio{372}โสกมปฺปชหํ ชนฺตุ, ภิยฺโย ทุกฺขํ นิคจฺฉติ.}

\prose{อนุตฺถุนนฺโต กาลงฺกตํ\footnote{กาลกตํ (สี, สฺยา)}, โสกสฺส วสมนฺวคู. (13)}

\proseitem{592}{อญฺเญปิ ปสฺส คมิเน, ยถากมฺมุปเค นเร.}

\prose{มจฺจุโน วสมาคมฺม, ผนฺทนฺเตวิธ ปาณิโน. (14)}

\proseitem{593}{เยน เยน หิ มญฺญนฺติ, ตโต ตํ โหติ อญฺญถา.}

\prose{เอตาทิโส วินาภาโว, ปสฺส โลกสฺส ปริยายํ. (15)}

\proseitem{594}{อปิ วสฺสสตํ ชีเว, ภิยฺโย วา ปน มาณโว.}

\prose{ญาติสํฆา วินา โหติ, ชหาติ อิธ ชีวิตํ. (16)}

\proseitem{595}{ตสฺมา อรหโต สุตฺวา, วิเนยฺย ปริเทวิตํ.}

\prose{เปตํ กาลงฺกตํ ทิสฺวา, เนโส ลพฺภา มยา อิติ. (17)}

\proseitem{596}{ยถา สรณมาทิตฺตํ, วารินา ปรินิพฺพเย\footnote{ปรินิพฺพุโต (สี, ก)}.}

\csromangathameasure{เอวมฺปิ ธีโร สปญฺโญ, ปณฺฑิโต กุสโล นโร.}
\csromangathagroup{\csromangathastanza{เอวมฺปิ ธีโร สปญฺโญ, ปณฺฑิโต กุสโล นโร.}}

\prose{ขิปฺปมุปฺปติตํ โสกํ, วาโต ตูลํว ธํสเย. (18)}

\proseitem{597}{ปริเทวํ ปชปฺปญฺจ, โทมนสฺสญฺจ อตฺตโน.}

\prose{อตฺตโน สุขเมสาโน, อพฺพเห สลฺลมตฺตโน. (19)}

\proseitem{598}{อพฺพูฬฺหสลฺโล อสิโต, สนฺติํ ปปฺปุยฺย เจตโส.}

\csromangathameasure{สพฺพโสกํ อติกฺกนฺโต, อโสโก โหติ นิพฺพุโตติ. สลฺลสุตฺตํ อฏฺฐมํ.}
\csromangathagroup{\csromangathastanza{สพฺพโสกํ อติกฺกนฺโต, อโสโก โหติ นิพฺพุโตติ.\csromanspacer{} สลฺลสุตฺตํ อฏฺฐมํ.}}

\csromanheader{2}{9. วาเสฏฺฐสุตฺต}
\prose{\csromansymbolfootnote{*}{ม 2. 406 ปิฏฺเฐปิ.}เอวํ เม สุตํ\csromandash{}เอกํ สมยํ ภควา อิจฺฉานงฺคเล วิหรติ อิจฺฉานงฺคลวนสณฺเฑ.\csromanspacer{} เตน โข ปน สมเยน สมฺพหุลา อภิญฺญาตา อภิญฺญาตา พฺราหฺมณมหาสาลา อิจฺฉานงฺคเล ปฏิวสนฺติ.\csromanspacer{} เสยฺยถิทํ, จงฺกี พฺราหฺมโณ, ตารุกฺโข พฺราหฺมโณ, โปกฺขรสาติ พฺราหฺมโณ, ชาณุสฺโสณิ\footnote{ชาณุโสณิ (ก)} พฺราหฺมโณ, โตเทยฺโย พฺราหฺมโณ, อญฺเญ จ อภิญฺญาตา อภิญฺญาตา พฺราหฺมณมหาสาลา.\csromanspacer{} อถ โข วาเสฏฺฐภารทฺวาชานํ}

\vspace{\tipitakaparskip}
\csromancenter{มหาวคฺค มาณวานํ ชงฺฆาวิหารํ อนุจงฺกมนฺตานํ อนุวิจรนฺตานํ\footnote{อนุจงฺกมมานานํ อนุวิจรมานานํ (สี, อิ)} อยมนฺตรากถา อุทปาทิ “กถํ โภ พฺราหฺมโณ โหตี”ติ.}
\prose{\csromanfolio{373}ภารทฺวาโช มาณโว เอวมาห “ยโต โข โภ อุภโต สุชาโต โหติ มาติโต จ ปิติโต จ สํสุทฺธคหณิโก ยาว สตฺตมา ปิตามหยุคา อกฺขิตฺโต อนุปกฺกุฏฺโฐ ชาติวาเทน, เอตฺตาวตา โข โภ พฺราหฺมโณ โหตี”ติ.}

\prose{วาเสฏฺโฐ มาณโว เอวมาห “ยโต โข โภ สีลวา จ โหติ วตสมฺปนฺโน\footnote{วตฺตสมฺปนฺโน (สี, สฺยา, ม 2. 407)} จ, เอตฺตาวตา โข โภ พฺราหฺมโณ โหตี”ติ.\csromanspacer{} เนว โข อสกฺขิ ภารทฺวาโช มาณโว วาเสฏฺฐํ มาณวํ สญฺญาเปตุํ, น ปน อสกฺขิ วาเสฏฺโฐ มาณโว ภารทฺวาชํ มาณวํ สญฺญาเปตุํ.}

\prose{อถ โข วาเสฏฺโฐ มาณโว ภารทฺวาชํ มาณวํ อามนฺเตสิ “อยํ โข โภ\footnote{อยํ โภ (สี, สฺยา, ก), อยํ โข (อิ)} ภารทฺวาช สมโณ โคตโม สกฺยปุตฺโต สกฺยกุลา ปพฺพชิโต อิจฺฉานงฺคเล วิหรติ อิจฺฉานงฺคลวนสณฺเฑ, ตํ โข ปน ภวนฺตํ โคตมํ เอวํ กลฺยาโณ กิตฺติสทฺโท อพฺภุคฺคโต ‘อิติปิ ฯเปฯ พุทฺโธ ภควา’ติ.\csromanspacer{} อายาม โภ ภารทฺวาช เยน สมโณ โคตโม เตนุปสงฺกมิสฺสาม, อุปสงฺกมิตฺวา สมณํ โคตมํ เอตมตฺถํ ปุจฺฉิสฺสาม, ยถา โน สมโณ โคตโม พฺยากริสฺสติ, ตถา นํ ธาเรสฺสามา”ติ. “เอวํ โภ”ติ โข ภารทฺวาโช มาณโว วาเสฏฺฐสฺส มาณวสฺส ปจฺจสฺโสสิ.}

\prose{อถ โข วาเสฏฺฐภารทฺวาชา มาณวา เยน ภควา เตนุปสงฺกมิํสุ, อุปสงฺกมิตฺวา ภควตา สทฺธิํ สมฺโมทิํสุ, สมฺโมทนียํ กถํ สารณียํ วีติสาเรตฺวา เอกมนฺตํ นิสีทิํสุ, เอกมนฺตํ นิสินฺโน โข วาเสฏฺโฐ มาณโว ภควนฺตํ คาถาหิ อชฺฌภาสิ\csromandash{}}

\proseitem{599}{อนุญฺญาตปฏิญฺญาตา, เตวิชฺชา มยมสฺมุโภ.}

\prose{อหํ โปกฺขรสาติสฺส, ตารุกฺขสฺสา’ยํ มาณโว. (1)}

\proseitem{600}{เตวิชฺชานํ ยทกฺขาตํ, ตตฺร เกวลิโน’สฺมเส.}

\prose{ปทก’สฺม เวยฺยากรณา, ชปฺเป อาจริยสาทิสา. (2)}

\gambhira{สุตฺตนิปาตปาฬิ}
\proseitem{601}{\csromanfolio{374}เตสํ โน ชาติวาทสฺมิํ, วิวาโท อตฺถิ โคตม.}

\csromangathameasure{ชาติยา พฺราหฺมโณ โหติ, ภารทฺวาโช อิติ ภาสติ.}
\csromangathagroup{\csromangathastanza{ชาติยา พฺราหฺมโณ โหติ, ภารทฺวาโช อิติ ภาสติ.}}

\prose{อหญฺจ กมฺมุนา\footnote{กมฺมนา (สี, อิ) เอวมุปริปิ.} พฺรูมิ, เอวํ ชานาหิ จกฺขุม. (3)}

\csromanhangingbegin
\hangingheaditem{602}{เต น สกฺโกม สญฺญาเปตุํ, อญฺญมญฺญํ มยํ อุโภ.}
\hangingbody{ภวนฺตํ2 ปุฏฺฐุมาคมฺหา, สมฺพุทฺธํ อิติ วิสฺสุตํ. (4)}
\csromanhangingend

\proseitem{603}{จนฺทํ ยถา ขยาตีตํ, เปจฺจ ปญฺชลิกา ชนา.}

\prose{วนฺทมานา นมสฺสนฺติ, เอวํ โลกสฺมิ โคตมํ. (5)}

\proseitem{604}{จกฺขุํ โลเก สมุปฺปนฺนํ, มยํ ปุจฺฉาม โคตมํ.}

\csromangathameasure{ชาติยา พฺราหฺมโณ โหติ, อุทาหุ ภวติ กมฺมุนา.}
\csromangathagroup{\csromangathastanza{ชาติยา พฺราหฺมโณ โหติ, อุทาหุ ภวติ กมฺมุนา.}}

\prose{อชานตํ โน ปพฺรูหิ, ยถา ชาเนมุ พฺราหฺมณํ. (6)}

\proseitem{605}{เตสํ โว อหํ พฺยกฺขิสฺสํ, (วาเสฏฺฐาติ ภควา)}

\prose{อนุปุพฺพํ ยถาตถํ.}

\prose{ชาติวิภงฺคํ ปาณานํ,}

\csromanheader{2}{อญฺญมญฺญา หิ ชาติโย. (7)}
\proseitem{606}{ติณรุกฺเขปิ ชานาถ, น จาปิ ปฏิชานเร.}

\prose{ลิงฺคํ ชาติมยํ เตสํ, อญฺญมญฺญา หิ ชาติโย. (8)}

\proseitem{607}{ตโต กีเฏ ปฏงฺเค จ, ยาว กุนฺถกิปิลฺลิเก.}

\prose{ลิงฺคํ ชาติมยํ เตสํ, อญฺญมญฺญา หิ ชาติโย. (9)}

\proseitem{608}{จตุปฺปเทปิ ชานาถ, ขุทฺทเก จ มหลฺลเก.}

\prose{ลิงฺคํ ชาติมยํ เตสํ, อญฺญมญฺญา หิ ชาติโย. (10)}

\proseitem{609}{ปาทูทเรปิ ชานาถ, อุรเค ทีฆปิฏฺฐิเก.}

\prose{ลิงฺคํ ชาติมยํ เตสํ, อญฺญมญฺญา หิ ชาติโย. (11)}

\proseitem{610}{ตโต มจฺเฉปิ ชานาถ, โอทเก วาริโคจเร.}

\prose{ลิงฺคํ ชาติมยํ เตสํ, อญฺญมญฺญา หิ ชาติโย. (12)}

\proseitem{611}{ตโต ปกฺขีปิ ชานาถ, ปตฺตยาเน วิหงฺคเม.}

\prose{ลิงฺคํ ชาติมยํ เตสํ, อญฺญมญฺญา หิ ชาติโย. (13)}

\chapterheadpageread{3. มหาวคฺค}
\proseitem{612}{\csromanfolio{375}ยถา เอตาสุ ชาตีสุ, ลิงฺคํ ชาติมยํ ปุถุ.}

\prose{เอวํ นตฺถิ มนุสฺเสสุ, ลิงฺคํ ชาติมยํ ปุถุ. (14)}

\proseitem{613}{น เกเสหิ น สีเสน, น กณฺเณหิ น อกฺขิภิ.}

\prose{น มุเขน น นาสาย, น โอฏฺเฐหิ ภมูหิ วา. (15)}

\proseitem{614}{น คีวาย น อํเสหิ, น อุทเรน น ปิฏฺฐิยา.}

\prose{น โสณิยา น อุรสา, น สมฺพาเธ น เมถุเน\footnote{น สมฺพาธา น เมถุนา (สฺยา, ก)}. (16)}

\proseitem{615}{น หตฺเถหิ น ปาเทหิ, นางฺคุลีหิ นเขหิ วา.}

\csromangathameasure{น ชงฺฆาหิ น อูรูหิ, น วณฺเณน สเรน วา.}
\csromangathagroup{\csromangathastanza{น ชงฺฆาหิ น อูรูหิ, น วณฺเณน สเรน วา.}}

\prose{ลิงฺคํ ชาติมยํ เนว, ยถา อญฺญาสุ ชาติสุ. (17)}

\proseitem{616}{ปจฺจตฺตญฺจ สรีเรสุ\footnote{ปจฺจตฺตํ สสรีเรสุ (สี, อิ)}, มนุสฺเสเสฺวตํ น วิชฺชติ.}

\prose{โวการญฺจ มนุสฺเสสุ, สมญฺญาย ปวุจฺจติ. (18)}

\proseitem{617}{โย หิ โกจิ มนุสฺเสสุ, โครกฺขํ อุปชีวติ.}

\prose{เอวํ วาเสฏฺฐ ชานาหิ, กสฺสโก โส น พฺราหฺมโณ. (19)}

\proseitem{618}{โย หิ โกจิ มนุสฺเสสุ, ปุถุสิปฺเปน ชีวติ.}

\prose{เอวํ วาเสฏฺฐ ชานาหิ, สิปฺปิโก โส น พฺราหฺมโณ. (20)}

\proseitem{619}{โย หิ โกจิ มนุสฺเสสุ, โวหารํ อุปชีวติ.}

\prose{เอวํ วาเสฏฺฐ ชานาหิ, วาณิโช โส น พฺราหฺมโณ. (21)}

\proseitem{620}{โย หิ โกจิ มนุสฺเสสุ, ปรเปสฺเสน ชีวติ.}

\prose{เอวํ วาเสฏฺฐ ชานาหิ, เปสฺสิโก\footnote{เปสฺสโก (ก)} โส น พฺราหฺมโณ. (22)}

\proseitem{621}{โย หิ โกจิ มนุสฺเสสุ, อทินฺนํ อุปชีวติ.}

\prose{เอวํ วาเสฏฺฐ ชานาหิ, โจโร เอโส น พฺราหฺมโณ. (23)}

\proseitem{622}{โย หิ โกจิ มนุสฺเสสุ, อิสฺสตฺถํ อุปชีวติ.}

\prose{เอวํ วาเสฏฺฐ ชานาหิ, โยธาชีโว น พฺราหฺมโณ. (24)}

\proseitem{623}{โย หิ โกจิ มนุสฺเสสุ, โปโรหิจฺเจน ชีวติ.}

\prose{เอวํ วาเสฏฺฐ ชานาหิ, ยาชโก โส น พฺราหฺมโณ. (25)}

\gambhira{สุตฺตนิปาตปาฬิ}
\csromangathameasure{โย หิ โกจิ มนุสฺเสสุ, คามํ รฏฺฐญฺจ ภุญฺชติ.}
\csromangathagroup{\csromangathastanza{\csromanfolio{376}โย หิ โกจิ มนุสฺเสสุ, คามํ รฏฺฐญฺจ ภุญฺชติ.}}

\proseitem{624}{เอวํ วาเสฏฺฐ ชานาหิ, ราชา เอโส น พฺราหฺมโณ. (26)}

\proseitem{625}{\csromansymbolfootnote{*}{เหฏฺฐา 70 ปิฏฺเฐปิ.}น จาหํ พฺราหฺมณํ พฺรูมิ, โยนิชํ มตฺติสมฺภวํ.}

\csromangathameasure{โภวาทิ นาม โส โหติ, สเจ โหติ สกิญฺจโน.}
\csromangathagroup{\csromangathastanza{โภวาทิ นาม โส โหติ, สเจ\footnote{ส เว (สี, สฺยา)} โหติ สกิญฺจโน.}}

\vspace{\tipitakaparskip}
\csromancenter{อกิญฺจนํ อนาทานํ, ตมหํ พฺรูมิ พฺราหฺมณํ. (27)}
\proseitem{626}{\csromansymbolmark{*}สพฺพสํโยชนํ เฉตฺวา, โย เว น ปริตสฺสติ.}

\prose{สงฺคาติคํ วิสํยุตฺตํ, ตมหํ พฺรูมิ พฺราหฺมณํ. (28)}

\proseitem{627}{\csromansymbolmark{*}เฉตฺวา นทฺธิํ วรตฺตญฺจ, สนฺทานํ สหนุกฺกมํ.}

\csromangathameasure{อุกฺขิตฺตปลิฆํ พุทฺธํ, ตมหํ พฺรูมิ พฺราหฺมณํ.}
\csromangathagroup{\csromangathastanza{อุกฺขิตฺตปลิฆํ พุทฺธํ, ตมหํ พฺรูมิ พฺราหฺมณํ.}}

\csromanheader{2}{(29)}
\proseitem{628}{\csromansymbolmark{*}อกฺโกสํ วธพนฺธญฺจ, อทุฏฺโฐ โย ติติกฺขติ.}

\vspace{\tipitakaparskip}
\csromancenter{ขนฺตีพลํ พลานีกํ, ตมหํ พฺรูมิ พฺราหฺมณํ. (30)}
\csromancenter{\csromansymbolmark{*}อกฺโกธนํ วตวนฺตํ, สีลวนฺตํ อนุสฺสทํ.}
\csromancenter{ทนฺตํ อนฺติมสารีรํ, ตมหํ พฺรูมิ พฺราหฺมณํ. (31)}
\csromancenter{\csromansymbolmark{*}วาริ โปกฺขรปตฺเตว, อารคฺเคริว สาสโป.}
\csromancenter{โย น ลิมฺปติ กาเมสุ, ตมหํ พฺรูมิ พฺราหฺมณํ. (32)}
\proseitem{631}{\csromansymbolmark{*}โย ทุกฺขสฺส ปชานาติ, อิเธว ขยมตฺตโน.}

\prose{ปนฺนภารํ วิสํยุตฺตํ, ตมหํ พฺรูมิ พฺราหฺมณํ.(33)}

\proseitem{632}{\csromansymbolmark{*}คมฺภีรปญฺญํ เมธาวิํ, มคฺคามคฺคสฺส โกวิทํ.}

\csromanmatikaanchor{mtk.269}
\csromantocmarkref{subsection}{7. เสลสุตฺต}{mtk.269}
\bookmark[dest={mtk.269},level=2]{7. เสลสุตฺต}
\prose{อุตฺตมตฺถมนุปฺปตฺตํ, ตมหํ พฺรูมิ พฺราหฺมณํ. (34)}

\proseitem{633}{\csromansymbolmark{*}อสํสฏฺฐํ คหฏฺเฐหิ, อนาคาเรหิ จูภยํ.}

\vspace{\tipitakaparskip}
\csromancenter{อโนกสาริมปฺปิจฺฉํ, ตมหํ พฺรูมิ พฺราหฺมณํ. (35)}
\csromancenter{\csromansymbolmark{*}นิธาย ทณฺฑํ ภูเตสุ, ตเสสุ ถาวเรสุ จ.}
\csromancenter{โย น หนฺติ น ฆาเตติ, ตมหํ พฺรูมิ พฺราหฺมณํ. (36)}
\csromancenter{\csromansymbolmark{*}อวิรุทฺธํ วิรุทฺเธสุ, อตฺตทณฺเฑสุ นิพฺพุตํ.}
\csromancenter{สาทาเนสุ อนาทานํ, ตมหํ พฺรูมิ พฺราหฺมณํ. (37)}
\chapterheadpageread{3. มหาวคฺค}
\proseitem{636}{\csromanfolio{377}\csromansymbolfootnote{*}{เหฏฺฐา 71, 72 ปิฏฺเฐสุปิ.}ยสฺส ราโค จ โทโส จ, มาโน มกฺโข จ ปาติโต.}

\vspace{\tipitakaparskip}
\csromancenter{สาสโปริว อารคฺคา, ตมหํ พฺรูมิ พฺราหฺมณํ. (38)}
\csromancenter{\csromansymbolmark{*}อกกฺกสํ วิญฺญาปนิํ, คิรํ สจฺจมุทีรเย.}
\csromancenter{ยาย นาภิสเช กญฺจิ, ตมหํ พฺรูมิ พฺราหฺมณํ. (39)}
\csromancenter{\csromansymbolmark{*}โย’ธ ทีฆํ ว รสฺสํ วา, อณุํ ถูลํ สุภาสุภํ.}
\csromancenter{โลเก อทินฺนํ นาทิยติ, ตมหํ พฺรูมิ พฺราหฺมณํ. (40)}
\proseitem{639}{\csromansymbolmark{*}อาสา ยสฺส น วิชฺชนฺติ, อสฺมิํ โลเก ปรมฺหิ จ.}

\prose{นิราสาสํ\footnote{นิราสยํ (สี, สฺยา, อิ), นิราสกํ (?)} วิสํยุตฺตํ, ตมหํ พฺรูมิ พฺราหฺมณํ. (41)}

\proseitem{640}{\csromansymbolmark{*}ยสฺสาลยา น วิชฺชนฺติ, อญฺญาย อกถํกถี.}

\prose{อมโตคธมนุปฺปตฺตํ, ตมหํ พฺรูมิ พฺราหฺมณํ. (42)}

\proseitem{641}{\csromansymbolmark{*}โย’ธ ปุญฺญญฺจ ปาปญฺจ, อุโภ สงฺค’มุปจฺจคา.}

\prose{อโสกํ วิรชํ สุทฺธํ, ตมหํ พฺรูมิ พฺราหฺมณํ. (43)}

\proseitem{642}{\csromansymbolmark{*}จนฺทํว วิมลํ สุทฺธํ, วิปฺปสนฺนมนาวิลํ.}

\prose{นนฺทีภวปริกฺขีณํ, ตมหํ พฺรูมิ พฺราหฺมณํ. (44)}

\proseitem{643}{\csromansymbolmark{*}โย’มํ ปลิปถํ ทุคฺคํ, สํสารํ โมหมจฺจคา.}

\csromangathameasure{ติณฺโณ ปารงฺคโต ฌายี, อเนโช อกถํกถี.}
\csromangathagroup{\csromangathastanza{ติณฺโณ ปารงฺคโต ฌายี, อเนโช อกถํกถี.}}

\vspace{\tipitakaparskip}
\csromancenter{อนุปาทาย นิพฺพุโต, ตมหํ พฺรูมิ พฺราหฺมณํ. (45)}
\proseitem{644}{\csromansymbolmark{*}โย’ธ กาเม ปหนฺตฺวาน, อนาคาโร ปริพฺพเช.}

\prose{กามภวปริกฺขีณํ, ตมหํ พฺรูมิ พฺราหฺมณํ. (46)}

\proseitem{645}{\csromansymbolmark{*}โย’ธ ตณฺหํ ปหนฺตฺวาน, อนาคาโร ปริพฺพเช.}

\prose{ตณฺหาภวปริกฺขีณํ, ตมหํ พฺรูมิ พฺราหฺมณํ. (47)}

\proseitem{646}{\csromansymbolmark{*}หิตฺวา มานุสกํ โยคํ, ทิพฺพํ โยคํ อุปจฺจคา.}

\vspace{\tipitakaparskip}
\csromancenter{สพฺพโยควิสํยุตฺตํ, ตมหํ พฺรูมิ พฺราหฺมณํ. (48)}
\proseitem{647}{\csromansymbolmark{*}หิตฺวา รติญฺจ อรติํ, สีติภูตํ นิรูปธิํ.}

\prose{สพฺพโลกาภิภุํ วีรํ, ตมหํ พฺรูมิ พฺราหฺมณํ. (49)}

\gambhira{สุตฺตนิปาตปาฬิ}
\proseitem{648}{\csromanfolio{378}\csromansymbolfootnote{*}{เหฏฺฐา 73 ปิฏฺเฐปิ.}จุติํ โย เวทิ\footnote{โย’เวติ (?) อิติวุตฺตเก 262 ปิฏฺเฐ อฏฺฐกถาสํวณฺณนา ปสฺสิตพฺพา.} สตฺตานํ, อุปปตฺติญฺจ สพฺพโส.}

\prose{อสตฺตํ สุคตํ พุทฺธํ, ตมหํ พฺรูมิ พฺราหฺมณํ. (50)}

\proseitem{649}{\csromansymbolmark{*}ยสฺส คติํ น ชานนฺติ, เทวา คนฺธพฺพมานุสา.}

\vspace{\tipitakaparskip}
\csromancenter{ขีณาสวํ อรหนฺตํ, ตมหํ พฺรูมิ พฺราหฺมณํ. (51)}
\csromancenter{\csromansymbolmark{*}ยสฺส ปุเร จ ปจฺฉา จ, มชฺเฌ จ นตฺถิ กิญฺจนํ.}
\csromancenter{อกิญฺจนํ อนาทานํ, ตมหํ พฺรูมิ พฺราหฺมณํ. (52)}
\proseitem{651}{\csromansymbolmark{*}อุสภํ ปวรํ วีรํ, มเหสิํ วิชิตาวินํ.}

\prose{อเนชํ นฺหาตกํ พุทฺธํ, ตมหํ พฺรูมิ พฺราหฺมณํ.(53)}

\proseitem{652}{\csromansymbolmark{*}ปุพฺเพนิวาสํ โย เวทิ1, สคฺคาปายญฺจ ปสฺสติ.}

\vspace{\tipitakaparskip}
\csromancenter{อโถ ชาติกฺขยํ ปตฺโต, ตมหํ พฺรูมิ พฺราหฺมณํ. (54)}
\vspace{0.5\baselineskip}
\csromangathameasure{สมญฺญา เหสา โลกสฺมิํ, นามโคตฺตํ ปกปฺปิตํ.}
\csromangathagroup{\csromangathastanza{สมญฺญา เหสา โลกสฺมิํ, นามโคตฺตํ ปกปฺปิตํ.}}

\vspace{\tipitakaparskip}
\csromancenter{สมฺมุจฺจา สมุทาคตํ, ตตฺถ ตตฺถ ปกปฺปิตํ. (55)}
\vspace{0.5\baselineskip}
\csromangathameasure{ทีฆรตฺตมนุสยิตํ, \\ ทิฏฺฐิคตมชานตํ. \\ อชานนฺตา โน ปพฺรุวนฺติ,}
\csromangathagroup{\csromangathastanza{ทีฆรตฺตมนุสยิตํ, \\ ทิฏฺฐิคตมชานตํ. \\ อชานนฺตา โน\footnote{อชานนฺตา เต (ฏฺฐ), ม 2. 412 ปิฏฺเฐปิ ปสฺสิตพฺพํ.} ปพฺรุวนฺติ,}}

\proseitem{654}{“ชาติยา โหติ พฺราหฺมโณ”. (56)}

\csromangathameasure{น ชจฺจา พฺราหฺมโณ โหติ, \\ น ชจฺจา โหติ อพฺราหฺมโณ. \\ กมฺมุนา พฺราหฺมโณ โหติ,}
\csromangathagroup{\csromangathastanza{น ชจฺจา พฺราหฺมโณ โหติ, \\ น ชจฺจา โหติ อพฺราหฺมโณ. \\ กมฺมุนา พฺราหฺมโณ โหติ,}}

\csromanheader{2}{กมฺมุนา โหติ อพฺราหฺมโณ. (57)}
\csromangathameasure{กสฺสโก กมฺมุนา โหติ, \\ สิปฺปิโก โหติ กมฺมุนา. \\ วาณิโช กมฺมุนา โหติ,}
\csromangathagroup{\csromangathastanza{กสฺสโก กมฺมุนา โหติ, \\ สิปฺปิโก โหติ กมฺมุนา. \\ วาณิโช กมฺมุนา โหติ,}}

\csromanheader{2}{เปสฺสิโก โหติ กมฺมุนา. (58)}
\csromangathameasure{โจโรปิ กมฺมุนา โหติ, โยธาชีโวปิ กมฺมุนา.}
\csromangathagroup{\csromangathastanza{โจโรปิ กมฺมุนา โหติ, โยธาชีโวปิ กมฺมุนา.}}

\proseitem{657}{ยาชโก กมฺมุนา โหติ, ราชาปิ โหติ กมฺมุนา. (59)}

\chapterheadpageread{3. มหาวคฺค}
\proseitem{658}{\csromanfolio{379}เอวเมตํ ยถาภูตํ, กมฺมํ ปสฺสนฺติ ปณฺฑิตา.}

\prose{ปฏิจฺจสมุปฺปาททสฺสา, กมฺมวิปากโกวิทา. (60)}

\proseitem{659}{กมฺมุนา วตฺตติ โลโก, กมฺมุนา วตฺตติ ปชา.}

\prose{กมฺมนิพนฺธนา สตฺตา, รถสฺสาณีว ยายโต. (61)}

\proseitem{660}{ตเปน พฺรหฺมจริเยน, สํยเมน ทเมน จ.}

\prose{เอเตน พฺราหฺมโณ โหติ, เอตํ พฺราหฺมณมุตฺตมํ. (62)}

\csromanhangingbegin
\hangingheaditem{661}{ตีหิ วิชฺชาหิ สมฺปนฺโน, สนฺโต ขีณปุนพฺภโว.}
\hangingbody{เอวํ วาเสฏฺฐ ชานาหิ, พฺรหฺมา สกฺโก วิชานตนฺติ.}
\csromanhangingend

\prose{เอวํ วุตฺเต วาเสฏฺฐภารทฺวาชา มาณวา ภควนฺตํ เอตทโวจุํ “อภิกฺกนฺตํ โภ โคตม ฯเปฯ อุปาสเก โน ภวํ โคตโม ธาเรตุ อชฺชตคฺเค ปาณุเปเต\footnote{ปาณุเปตํ (ก)} สรณํ คเต”ติ.\csromanspacer{} วาเสฏฺฐสุตฺตํ นวมํ.}

\csromanheader{2}{10. โกกาลิกสุตฺต}
\prose{เอวํ เม สุตํ\csromandash{}เอกํ สมยํ ภควา สาวตฺถิยํ วิหรติ เชตวเน อนาถปิณฺฑิกสฺส อาราเม.\csromanspacer{} อถ โข โกกาลิโก ภิกฺขุ เยน ภควา เตนุปสงฺกมิ, อุปสงฺกมิตฺวา ภควนฺตํ อภิวาเทตฺวา เอกมนฺตํ นิสีทิ.\csromanspacer{} เอกมนฺตํ นิสินฺโน โข โกกาลิโก ภิกฺขุ ภควนฺตํ เอตทโวจ “ปาปิจฺฉา ภนฺเต สาริปุตฺตโมคฺคลฺลานา ปาปิกานํ อิจฺฉานํ วสํ คตา”ติ.}

\prose{เอวํ วุตฺเต ภควา โกกาลิกํ ภิกฺขุํ เอตทโวจ “มา เหวํ โกกาลิก, มา เหวํ โกกาลิก.\csromanspacer{} ปสาเทหิ โกกาลิก สาริปุตฺตโมคฺคลฺลาเนสุ จิตฺตํ, เปสลา สาริปุตฺตโมคฺคลฺลานา”ติ.}

\prose{ทุติยมฺปิ โข ฯเปฯ.\csromanspacer{} ตติยมฺปิ โข โกกาลิโก ภิกฺขุ ภควนฺตํ เอตทโวจ “กิญฺจาปิ เม ภนฺเต ภควา สทฺธายิโก ปจฺจยิโก.\csromanspacer{} อถ โข ปาปิจฺฉาว สาริปุตฺตโมคฺคลฺลานา ปาปิกานํ อิจฺฉานํ วสํ คตา”ติ.\csromanspacer{} ตติยมฺปิ โข ภควา โกกาลิกํ ภิกฺขุํ เอตทโวจ “มา เหวํ}

\gambhira{สุตฺตนิปาตปาฬิ}
\prose{\csromanfolio{380}โกกาลิก, มา เหวํ โกกาลิก.\csromanspacer{} ปสาเทหิ โกกาลิก สาริปุตฺตโมคฺคลฺลาเนสุ จิตฺตํ, เปสลา สาริปุตฺตโมคฺคลฺลานา”ติ.}

\prose{\csromansymbolfootnote{*}{สํ 1. 152; อํ 3. 394 ปิฏฺเฐสุปิ.}อถ โข โกกาลิโก ภิกฺขุ อุฏฺฐายาสนา ภควนฺตํ อภิวาเทตฺวา ปทกฺขิณํ กตฺวา ปกฺกามิ.\csromanspacer{} อจิรปฺปกฺกนฺตสฺส จ โกกาลิกสฺส ภิกฺขุโน สาสปมตฺตีหิ ปิฬกาหิ สพฺโพ กาโย ผุโฏ\footnote{ผุฏฺโฐ (สฺยา)} อโหสิ, สาสปมตฺติโย หุตฺวา มุคฺคมตฺติโย อเหสุํ, มุคฺคมตฺติโย หุตฺวา กฬายมตฺติโย อเหสุํ, กฬายมตฺติโย หุตฺวา โกลฏฺฐิมตฺติโย อเหสุํ, โกลฏฺฐิมตฺติโย หุตฺวา โกลมตฺติโย อเหสุํ, โกลมตฺติโย หุตฺวา อามลกมตฺติโย อเหสุํ, อามลกมตฺติโย หุตฺวา เพฬุวสลาฏุกมตฺติโย อเหสุํ, เพฬุวสลาฏุกมตฺติโย หุตฺวา พิลฺลมตฺติโย อเหสุํ, พิลฺลมตฺติโย หุตฺวา ปภิชฺชิํสุ, ปุพฺพญฺจ โลหิตญฺจ ปคฺฆริํสุ.\csromanspacer{} อถ โข โกกาลิโก ภิกฺขุ เตเนวาพาเธน กาลมกาสิ, กาลงฺกโต จ โกกาลิโก ภิกฺขุ ปทุมํ นิรยํ อุปปชฺชิ สาริปุตฺตโมคฺคลฺลาเนสุ จิตฺตํ อาฆาเตตฺวา.}

\prose{อถ โข พฺรหฺมา สหมฺปติ อภิกฺกนฺตาย รตฺติยา อภิกฺกนฺตวณฺโณ เกวลกปฺปํ เชตวนํ โอภาเสตฺวา เยน ภควา เตนุปสงฺกมิ, อุปสงฺกมิตฺวา ภควนฺตํ อภิวาเทตฺวา เอกมนฺตํ อฏฺฐาสิ, เอกมนฺตํ ฐิโต โข พฺรหฺมา สหมฺปติ ภควนฺตํ เอตทโวจ “โกกาลิโก ภนฺเต ภิกฺขุ กาลงฺกโต, กาลงฺกโต จ ภนฺเต โกกาลิโก ภิกฺขุ ปทุมํ นิรยํ อุปปนฺโน สาริปุตฺตโมคฺคลฺลาเนสุ จิตฺตํ อาฆาเตตฺวา”ติ.\csromanspacer{} อิทมโวจ พฺรหฺมา สหมฺปติ, อิทํ วตฺวา ภควนฺตํ อภิวาเทตฺวา ปทกฺขิณํ กตฺวา ตตฺเถวนฺตรธายิ.}

\prose{อถ โข ภควา ตสฺสา รตฺติยา อจฺจเยน ภิกฺขู อามนฺเตสิ “อิมํ ภิกฺขเว รตฺติํ พฺรหฺมา สหมฺปติ อภิกฺกนฺตาย รตฺติยา ฯเปฯ.\csromanspacer{} อิทมโวจ ภิกฺขเว พฺรหฺมา สหมฺปติ, อิทํ วตฺวา มํ ปทกฺขิณํ กตฺวา ตตฺเถวนฺตรธายี”ติ.}

\prose{เอวํ วุตฺเต อญฺญตโร ภิกฺขุ ภควนฺตํ เอตทโวจ “กีวทีฆํ นุ โข ภนฺเต ปทุเม นิรเย อายุปฺปมาณนฺ”ติ.\csromanspacer{} ทีฆํ โข ภิกฺขุ ปทุเม นิรเย อายุปฺปมาณํ, ตํ น สุกรํ สงฺขาตุํ “เอตฺตกานิ วสฺสานิ” อิติ วา “เอตฺตกานิ วสฺสสตานิ” อิติ วา “เอตฺตกานิ วสฺสสหสฺสานิ” อิติ วา “เอตฺตกานิ}

\vspace{\tipitakaparskip}
\csromancenter{มหาวคฺค วสฺสสตสหสฺสานิ” อิติ วาติ.\csromanspacer{} สกฺกา ปน ภนฺเต อุปมา\footnote{อุปมํ (สี, สฺยา, ก)} กาตุนฺติ. “สกฺกา ภิกฺขู”ติ ภควา อโวจ\csromandash{}}
\prose{\csromanfolio{381}เสยฺยถาปิ ภิกฺขุ วีสติขาริโก โกสลโก ติลวาโห ตโต ปุริโส วสฺสสตสฺส วสฺสสตสฺส อจฺจเยน เอกเมกํ ติลํ อุทฺธเรยฺย.\csromanspacer{} ขิปฺปตรํ โข โส ภิกฺขุ วีสติขาริโก โกสลโก ติลวาโห อิมินา อุปกฺกเมน ปริกฺขยํ ปริยาทานํ คจฺเฉยฺย, นเตฺวว เอโก อพฺพุโท นิรโย.\csromanspacer{} เสยฺยถาปิ ภิกฺขุ วีสติ อพฺพุทา นิรยา, เอวเมโก นิรพฺพุโท นิรโย.\csromanspacer{} เสยฺยถาปิ ภิกฺขุ วีสติ นิรพฺพุทา นิรยา, เอวเมโก อพโพ นิรโย.\csromanspacer{} เสยฺยถาปิ ภิกฺขุ วีสติ อพพา นิรยา, เอวเมโก อฏโฏ นิรโย.\csromanspacer{} เสยฺยถาปิ ภิกฺขุ วีสติ อฏฏา นิรยา, เอวเมโก อหโห นิรโย.\csromanspacer{} เสยฺยถาปิ ภิกฺขุ วีสติ อหหา นิรยา, เอวเมโก กุมุโท นิรโย.\csromanspacer{} เสยฺยถาปิ ภิกฺขุ วีสติ กุมุทา นิรยา, เอวเมโก โสคนฺธิโก นิรโย.\csromanspacer{} เสยฺยถาปิ ภิกฺขุ วีสติ โสคนฺธิกา นิรยา, เอวเมโก อุปฺปลโก นิรโย.\csromanspacer{} เสยฺยถาปิ ภิกฺขุ วีสติ อุปฺปลกา นิรยา, เอวเมโก ปุณฺฑรีโก นิรโย.\csromanspacer{} เสยฺยถาปิ ภิกฺขุ วีสติ ปุณฺฑรีกา นิรยา, เอวเมโก ปทุโม นิรโย.\csromanspacer{} ปทุมํ โข ปน ภิกฺขุ นิรยํ โกกาลิโก ภิกฺขุ อุปปนฺโน สาริปุตฺตโมคฺคลฺลาเนสุ จิตฺตํ อาฆาเตตฺวาติ, อิทมโวจ ภควา, อิทํ วตฺวาน สุคโต, อถาปรํ เอตทโวจ สตฺถา\csromandash{}}

\proseitem{662}{ปุริสสฺส หิ ชาตสฺส,}

\prose{กุฐารี\footnote{กุธารี (ก)} ชายเต มุเข.}

\prose{ยาย ฉินฺทติ อตฺตานํ,}

\csromanheader{2}{พาโล ทุพฺภาสิตํ ภณํ. (1)}
\proseitem{663}{โย นินฺทิยํ ปสํสติ,}

\prose{ตํ วา นินฺทติ โย ปสํสิโย.}

\prose{วิจินาติ มุเขน โส กลิํ,}

\csromanheader{2}{กลินา เตน สุขํ น วินฺทติ. (2)}
\gambhira{สุตฺตนิปาตปาฬิ}
\csromangathameasure{อปฺปมตฺโต อยํ กลิ, \\ โย อกฺเขสุ ธนปราชโย. \\ สพฺพสฺสาปิ สหาปิ อตฺตนา, \\ อยเมว มหตฺตโร กลิ.}
\csromangathagroup{\csromangathastanza{\csromanfolio{382}อปฺปมตฺโต อยํ กลิ, \\ โย อกฺเขสุ ธนปราชโย. \\ สพฺพสฺสาปิ สหาปิ อตฺตนา, \\ อยเมว มหตฺตโร\footnote{มหนฺตตโร (สี)} กลิ.}}

\csromanheader{2}{โย สุคเตสุ มนํ ปโทสเย. (3)}
\csromangathameasure{สตํ สหสฺสานํ นิรพฺพุทานํ, \\ ฉตฺติํสติ ปญฺจ จ อพฺพุทานิ. \\ ยมริยครหี นิรยํ อุเปติ,}
\csromangathagroup{\csromangathastanza{สตํ สหสฺสานํ นิรพฺพุทานํ, \\ ฉตฺติํสติ ปญฺจ จ อพฺพุทานิ\footnote{อพฺพุทานํ (ก)}. \\ ยมริยครหี นิรยํ อุเปติ,}}

\csromanmatikaanchor{mtk.270}
\csromantocmarkref{subsection}{8. สลฺลสุตฺต}{mtk.270}
\bookmark[dest={mtk.270},level=2]{8. สลฺลสุตฺต}
\csromanheader{2}{วาจํ มนญฺจ ปณิธาย ปาปกํ. (4)}
\proseitem{666}{\csromansymbolfootnote{*}{เหฏฺฐา 57 ปิฏฺเฐปิ.}อภูตวาที นิรยํ อุเปติ,}

\prose{โย วาปิ กตฺวา น กโรมิ’จาห.}

\csromangathameasure{อุโภปิ เต เปจฺจ สมา ภวนฺติ,}
\csromangathagroup{\csromangathastanza{อุโภปิ เต เปจฺจ สมา ภวนฺติ,}}

\csromanheader{2}{นิหีนกมฺมา มนุชา ปรตฺถ. (5)}
\proseitem{667}{\csromansymbolfootnote{+}{เหฏฺฐา 32 ปิฏฺเฐปิ.}โย อปฺปทุฏฺฐสฺส นรสฺส ทุสฺสติ,}

\prose{สุทฺธสฺส โปสสฺส อนงฺคณสฺส.}

\csromangathameasure{ตเมว พาลํ ปจฺเจติ ปาปํ,}
\csromangathagroup{\csromangathastanza{ตเมว พาลํ ปจฺเจติ ปาปํ,}}

\csromanheader{2}{สุขุโม รโช ปฏิวาตํว ขิตฺโต. (6)}
\proseitem{668}{\csromansymbolmark{+}โย โลภคุเณ อนุยุตฺโต,}

\prose{โส วจสา ปริภาสติ อญฺเญ.}

\csromangathameasure{อสทฺโธ กทริโย อวทญฺญู,}
\csromangathagroup{\csromangathastanza{อสทฺโธ กทริโย อวทญฺญู,}}

\csromanheader{2}{มจฺฉริ เปสุณิยํ\footnote{เปสุณิยสฺมิํ (พหูสุ)} อนุยุตฺโต. (7)}
\csromangathameasure{มุขทุคฺค วิภูต อนริย, \\ ภูนหุ ปาปก ทุกฺกฏการิ. \\ ปุริสนฺต กลี อวชาต,}
\csromangathagroup{\csromangathastanza{มุขทุคฺค วิภูต อนริย, \\ ภูนหุ\footnote{ภุนหต (สฺยา, ก)} ปาปก ทุกฺกฏการิ. \\ ปุริสนฺต กลี อวชาต,}}

\csromanheader{2}{มา พหุภาณิธ เนรยิโกสิ. (8)}
\chapterheadpageread{3. มหาวคฺค}
\proseitem{670}{\csromanfolio{383}รชมากิรสี อหิตาย,}

\prose{สนฺเต ครหสิ กิพฺพิสการี.}

\prose{พหูนิ ทุจฺจริตานิ จริตฺวา,}

\csromanheader{2}{คจฺฉสิ โข ปปตํ จิรรตฺตํ. (9)}
\proseitem{671}{น หิ นสฺสติ กสฺสจิ กมฺมํ,}

\prose{เอติ หตํ ลภเตว สุวามิ.}

\prose{ทุกฺขํ มนฺโท ปรโลเก,}

\csromanheader{2}{อตฺตนิ ปสฺสติ กิพฺพิสการี. (10)}
\proseitem{672}{อโยสงฺกุสมาหตฏฺฐานํ,}

\prose{ติณฺหธารมยสูลมุเปติ.}

\prose{อถ ตตฺต-อโยคุฬสนฺนิภํ,}

\csromanheader{2}{โภชนมตฺถิ ตถา ปติรูปํ. (11)}
\proseitem{673}{น หิ วคฺคุ วทนฺติ วทนฺตา,}

\prose{นาภิชวนฺติ น ตาณมุเปนฺติ.}

\prose{องฺคาเร สนฺถเต สยนฺติ\footnote{เสนฺติ (สี, สฺยา, อิ)},}

\csromanheader{2}{คินิสมฺปชฺชลิตํ ปวิสนฺติ. (12)}
\proseitem{674}{ชาเลน จ โอนหิยาน,}

\prose{ตตฺถ หนนฺติ อโยมยกุเฏภิ\footnote{อโยมยกูเฏหิ (สี, สฺยา, อิ)}.}

\prose{อนฺธํว ติมิสมายนฺติ,}

\csromanheader{2}{ตํ วิตตํ หิ ยถา มหิกาโย. (13)}
\proseitem{675}{อถ โลหมยํ ปน กุมฺภิํ,}

\prose{คินิสมฺปชฺชลิตํ ปวิสนฺติ.}

\prose{ปจฺจนฺติ หิ ตาสุ จิรรตฺตํ,}

\csromanheader{2}{อคฺคินิสมาสุ\footnote{คินิสฺสมาสุ (ก)} สมุปฺปิลวาเต. (14)}
\proseitem{676}{อถ ปุพฺพโลหิตมิสฺเส,}

\prose{ตตฺถ กิํ ปจฺจติ กิพฺพิสการี.}

\csromanmatikaanchor{mtk.271}
\csromantocmarkref{subsection}{9. วาเสฏฺฐสุตฺต}{mtk.271}
\bookmark[dest={mtk.271},level=2]{9. วาเสฏฺฐสุตฺต}
\prose{ยํ ยํ ทิสกํ\footnote{ทิสตํ (สี, สฺยา, อิ)} อธิเสติ,}

\csromanheader{2}{ตตฺถ กิลิสฺสติ สมฺผุสมาโน. (15)}
\gambhira{สุตฺตนิปาตปาฬิ}
\proseitem{677}{\csromanfolio{384}ปุฬวาวสเถ สลิลสฺมิํ,}

\prose{ตตฺถ กิํ ปจฺจติ กิพฺพิสการี.}

\prose{คนฺตุํ น หิ ตีรมปตฺถิ,}

\csromanheader{2}{สพฺพสมา หิ สมนฺตกปลฺลา. (16)}
\proseitem{678}{อสิปตฺตวนํ ปน ติณฺหํ,}

\prose{ตํ ปวิสนฺติ สมุจฺฉิทคตฺตา.}

\prose{ชิวฺหํ พลิเสน คเหตฺวา,}

\csromanheader{2}{อารชยารชยา วิหนนฺติ. (17)}
\proseitem{679}{อถ เวตรณิํ ปน ทุคฺคํ,}

\prose{ติณฺหธารขุรธารมุเปนฺติ.}

\prose{ตตฺถ มนฺทา ปปตนฺติ,}

\csromanheader{2}{ปาปกรา ปาปานิ กริตฺวา. (18)}
\proseitem{680}{ขาทนฺติ หิ ตตฺถ รุทนฺเต,}

\prose{สามา สพลา กาโกลคณา จ.}

\prose{โสณา สิงฺคาลา\footnote{สิคาลา (สี, อิ)} ปฏิคิทฺธา\footnote{ปฏิคิชฺฌา (สฺยา, อิ)},}

\csromanheader{2}{กุลลา วายสา จ\footnote{กุลลา จ วายสา (?)} วิตุทนฺติ. (19)}
\proseitem{681}{กิจฺฉา วต’ยํ อิธ วุตฺติ,}

\prose{ยํ ชโน ผุสติ\footnote{ปสฺสติ (สี, สฺยา, อิ)} กิพฺพิสการี.}

\prose{ตสฺมา อิธ ชีวิตเสเส,}

\csromanheader{2}{กิจฺจกโร สิยา นโร น จปฺปมชฺเช. (20)}
\proseitem{682}{เต คณิตา วิทูหิ ติลวาหา,}

\prose{เย ปทุเม นิรเย อุปนีตา.}

\prose{นหุตานิ หิ โกฏิโย ปญฺจ ภวนฺติ,}

\csromanheader{2}{ทฺวาทส โกฏิสตานิ ปุนญฺญา\footnote{ปนญฺเญ (ก)}. (21)}
\chapterheadpageread{3. มหาวคฺค}
\proseitem{683}{\csromanfolio{385}ยาว ทุขา\footnote{ทุกฺขา (สี, สฺยา), ทุกฺข (อิ, ก)} นิรยา อิธ วุตฺตา,}

\prose{ตตฺถปิ ตาว จิรํ วสิตพฺพํ.}

\csromangathameasure{ตสฺมา สุจิเปสลสาธุคุเณสุ, \\ วาจํ มนํ สตตํ 2 ปริรกฺเขติ. (22) โกกาลิกสุตฺตํ ทสมํ.}
\csromangathagroup{\csromangathastanza{ตสฺมา สุจิเปสลสาธุคุเณสุ, \\ วาจํ มนํ สตตํ 2 ปริรกฺเขติ. (22) โกกาลิกสุตฺตํ ทสมํ.}}

\csromanheader{2}{11. นาลกสุตฺต}
\proseitem{684}{อานนฺทชาเต ติทสคเณ ปตีเต,}

\prose{สกฺกญฺจ อินฺทํ สุจิวสเน จ เทเว.}

\prose{ทุสฺสํ คเหตฺวา อติริว โถมยนฺเต,}

\csromanheader{2}{อสิโต อิสิ อทฺทส ทิวาวิหาเร. (1)}
\proseitem{685}{ทิสฺวาน เทเว มุทิตมเน อุทคฺเค,}

\prose{จิตฺติํ กริตฺวาน อิทมโวจ\footnote{กริตฺวา อิทมโวจาสิ (สี)} ตตฺถ.}

\prose{กิํ เทวสํโฆ อติริว กลฺยรูโป,}

\csromanheader{2}{ทุสฺสํ คเหตฺวา รมยถ\footnote{ภมยถ (สี)} กิํ ปฏิจฺจ. (2)}
\proseitem{686}{ยทาปิ อาสี อสุเรหิ สงฺคโม,}

\prose{ชโย สุรานํ อสุรา ปราชิตา.}

\prose{ตทาปิ เนตาทิโส โลมหํสโน,}

\csromanheader{2}{กิมพฺภุตํ ทฏฺฐุ มรู ปโมทิตา. (3)}
\proseitem{687}{เสเฬนฺติ คายนฺติ จ วาทยนฺติ จ,}

\prose{ภุชานิ โผเฏนฺติ\footnote{โปเฐนฺติ (สี, อิ), โปเถนฺติ (ก)} จ นจฺจยนฺติ จ.}

\prose{ปุจฺฉามิ โวหํ เมรุมุทฺธวาสิเน,}

\csromanheader{2}{ธุนาถ เม สํสยํ ขิปฺป มาริสา. (4)}
\gambhira{สุตฺตนิปาตปาฬิ}
\proseitem{688}{\csromanfolio{386}โส โพธิสตฺโต รตนวโร อตุโลฺย,}

\prose{มนุสฺสโลเก หิตสุขตฺถาย\footnote{หิตสุขตาย (สี, สฺยา, อิ)} ชาโต.}

\prose{สกฺยาน คาเม ชนปเท ลุมฺพิเนยฺเย,}

\csromanheader{2}{เตนมฺห ตุฏฺฐา อติริว กลฺยรูปา. (5)}
\proseitem{689}{โส สพฺพสตฺตุตฺตโม อคฺคปุคฺคโล,}

\prose{นราสโภ สพฺพปชานมุตฺตโม.}

\prose{วตฺเตสฺสติ จกฺกมิสิวฺหเย วเน,}

\csromanheader{2}{นทํว สีโห พลวา มิคาภิภู. (6)}
\proseitem{690}{ตํ สทฺทํ สุตฺวา ตุริต’มวสรี โส,}

\prose{สุทฺโธทนสฺส ตท ภวนํ อุปาวิสิ\footnote{อุปาคมิ (สี, อิ)}.}

\prose{นิสชฺช ตตฺถ อิทมโวจาสิ สเกฺย,}

\csromanheader{2}{กุหิํ กุมาโร อหมปิ ทฏฺฐุกาโม. (7)}
\proseitem{691}{ตโต กุมารํ ชลิตมิว สุวณฺณํ,}

\prose{อุกฺกามุเขว สุกุสลสมฺปหฏฺฐํ\footnote{สุกุสเลน สมฺปหฏฺฐํ (ก)}.}

\prose{ททฺทลฺลมานํ\footnote{ททฺทฬฺหมานํ (ก)} สิริยา อโนมวณฺณํ,}

\csromanheader{2}{ทสฺเสสุ ปุตฺตํ อสิตวฺหยสฺส สกฺยา. (8)}
\proseitem{692}{ทิสฺวา กุมารํ สิขิมิว ปชฺชลนฺตํ,}

\prose{ตาราสภํว นภสิคมํ วิสุทฺธํ.}

\prose{สุริยํ ตปนฺตํ สรทริวพฺภมุตฺตํ,}

\csromanheader{2}{อานนฺทชาโต วิปุลมลตฺถ ปีติํ. (9)}
\proseitem{693}{อเนกสาขญฺจ สหสฺสมณฺฑลํ,}

\prose{ฉตฺตํ มรู ธารยุมนฺตลิกฺเข.}

\prose{สุวณฺณทณฺฑา วีติปตนฺติ จามรา,}

\csromanheader{2}{น ทิสฺสเร จามรฉตฺตคาหกา. (10)}
\chapterheadpageread{3. มหาวคฺค}
\proseitem{694}{\csromanfolio{387}ทิสฺวา ชฏี กณฺหสิริวฺหโย อิสิ,}

\prose{สุวณฺณนิกฺขํ วิย ปณฺฑุกมฺพเล.}

\prose{เสตญฺจ ฉตฺตํ ธริยนฺต\footnote{ธาริยนฺต (สฺยา), ธารยนฺตํ (สี, ก)} มุทฺธนิ,}

\csromanheader{2}{อุทคฺคจิตฺโต สุมโน ปฏิคฺคเห. (11)}
\proseitem{695}{ปฏิคฺคเหตฺวา ปน สกฺยปุงฺควํ,}

\prose{ชิคีสโก\footnote{ชิคิํสโต (ก)} ลกฺขณมนฺตปารคู.}

\prose{ปสนฺนจิตฺโต คิรมพฺภุทีรยิ,}

\prose{“อนุตฺตรา’ยํ ทฺวิปทานมุตฺตโม”\footnote{ทิปทานมุตฺตโม (สี, สฺยา, อิ)}. (12)}

\proseitem{696}{อถตฺตโน คมนมนุสฺสรนฺโต,}

\prose{อกลฺยรูโป คฬยติ อสฺสุกานิ.}

\prose{ทิสฺวาน สกฺยา อิสิมโวจุํ รุทนฺตํ,}

\prose{“โน เจ กุมาเร ภวิสฺสติ อนฺตราโย”. (13)}

\proseitem{697}{ทิสฺวาน สเกฺย อิสิมโวจ อกเลฺย,}

\prose{นาหํ กุมาเร อหิตมนุสฺสรามิ.}

\prose{น จาปิมสฺส ภวิสฺสติ อนฺตราโย,}

\csromanheader{2}{น โอรกา’ยํ อธิมานสา\footnote{อธิมนสา (สี, สฺยา)} ภวาถ. (14)}
\proseitem{698}{สมฺโพธิยคฺคํ ผุสิสฺสตา’ยํ กุมาโร,}

\prose{โส ธมฺมจกฺกํ ปรมวิสุทฺธทสฺสี.}

\prose{วตฺเตสฺสตา’ยํ พหุชนหิตานุกมฺปี,}

\csromanheader{2}{วิตฺถาริก’สฺส ภวิสฺสติ พฺรหฺมจริยํ. (15)}
\proseitem{699}{มมญฺจ อายุ น จิรมิธาวเสโส,}

\prose{อถนฺตรา เม ภวิสฺสติ กาลกิริยา.}

\prose{โสหํ น โสสฺสํ\footnote{สุสฺสํ (สี, สฺยา)} อสมธุรสฺส ธมฺมํ,}

\csromanheader{2}{เตนมฺหิ อฏฺโฏ พฺยสนํคโต อฆาวี. (16)}
\gambhira{สุตฺตนิปาตปาฬิ}
\proseitem{700}{\csromanfolio{388}โส สากิยานํ วิปุลํ ชเนตฺวา ปีติํ,}

\prose{อนฺเตปุรมฺหา นิคฺคมา\footnote{นิรคมา (สี, สฺยา), นิคมา (ก-สี), นิรคม (อิ)} พฺรหฺมจารี.}

\prose{โส ภาคิเนยฺยํ สยํ อนุกมฺปมาโน,}

\csromanheader{2}{สมาทเปสิ อสมธุรสฺส ธมฺเม. (17)}
\proseitem{701}{พุทฺโธติ โฆสํ ยท\footnote{ยทิ (สฺยา, ก) 3. สยํ ปริปุจฺฉิยาโน (สี, สฺยา)} ปรโต สุณาสิ,}

\prose{“สมฺโพธิปตฺโต วิวรติ ธมฺมมคฺคํ”.}

\prose{คนฺตฺวาน ตตฺถ สมยํ ปริปุจฺฉมาโน3,}

\csromanheader{2}{จรสฺสุ ตสฺมิํ ภควติ พฺรหฺมจริยํ. (18)}
\proseitem{702}{เตนานุสิฏฺโฐ หิตมเนน ตาทินา,}

\prose{อนาคเต ปรมวิสุทฺธทสฺสินา.}

\prose{โส นาลโก อุปจิตปุญฺญสญฺจโย,}

\csromanheader{2}{ชินํ ปกิกฺขํ\footnote{ปติ + อิกฺขํ = ปติกฺขํ.} ปริวสิ รกฺขิตินฺทฺริโย. (19)}
\proseitem{703}{สุตฺวาน โฆสํ ชินวรจกฺกวตฺตเน,}

\prose{คนฺตฺวาน ทิสฺวา อิสินิสภํ ปสนฺโน.}

\prose{โมเนยฺยเสฏฺฐํ มุนิปวรํ อปุจฺฉิ,}

\nitthitam{สมาคเต อสิตาวฺหยสฺส สาสเนติ. (20) วตฺถุคาถา นิฏฺฐิตา.}
\csromanhangingbegin
\hangingheaditem{704}{อญฺญาตเมตํ วจนํ, อสิตสฺส ยถาตถํ.}
\hangingbody{ตํ ตํ โคตม ปุจฺฉามิ, สพฺพธมฺมาน ปารคุํ. (21)}
\csromanhangingend

\proseitem{705}{อนคาริยุเปตสฺส, ภิกฺขาจริยํ ชิคีสโต.}

\prose{มุนิ ปพฺรูหิ เม ปุฏฺโฐ, โมเนยฺยํ อุตฺตมํ ปทํ. (22)}

\proseitem{706}{โมเนยฺยํ เต อุปญฺญิสฺสํ, (อิติ ภควา,)}

\prose{ทุกฺกรํ ทุรภิสมฺภวํ.}

\prose{หนฺท เต นํ ปวกฺขามิ,}

\csromanheader{2}{สนฺถมฺภสฺสุ ทโฬ ภว. (23)}
\chapterheadpageread{3. มหาวคฺค}
\proseitem{707}{\csromanfolio{389}สมานภาคํ กุพฺเพถ, คาเม อกฺกุฏฺฐวนฺทิตํ.}

\prose{มโนปโทสํ รกฺเขยฺย, สนฺโต อนุณฺณโต จเร. (24)}

\proseitem{708}{อุจฺจาวจา นิจฺฉรนฺติ, ทาเย อคฺคิสิขูปมา.}

\prose{นาริโย มุนิํ ปโลเภนฺติ, ตา สุ ตํ มา ปโลภยุํ. (25)}

\proseitem{709}{วิรโต เมถุนา ธมฺมา, หิตฺวา กาเม ปโรปเร\footnote{ปโรวเร (สี, อิ), วราวเร (สฺยา)}.}

\prose{อวิรุทฺโธ อสารตฺโต, ปาเณสุ ตสถาวเร. (26)}

\proseitem{710}{ยถา อหํ ตถา เอเต, ยถา เอเต ตถา อหํ.}

\prose{อตฺตานํ อุปมํ กตฺวา, น หเนยฺย น ฆาตเย. (27)}

\proseitem{711}{หิตฺวา อิจฺฉญฺจ โลภญฺจ, ยตฺถ สตฺโต ปุถุชฺชโน.}

\prose{จกฺขุมา ปฏิปชฺเชยฺย, ตเรยฺย นรกํ อิมํ. (28)}

\proseitem{712}{อูนูทโร มิตาหาโร, อปฺปิจฺฉ’สฺส อโลลุโป.}

\prose{สทา\footnote{ส เว (อิ)} อิจฺฉาย นิจฺฉาโต, อนิจฺโฉ โหติ นิพฺพุโต. (29)}

\proseitem{713}{ส ปิณฺฑจารํ จริตฺวา, วนนฺตมภิหารเย.}

\prose{อุปฏฺฐิโต รุกฺขมูลสฺมิํ, อาสนูปคโต มุนิ. (30)}

\proseitem{714}{ส ฌานปสุโต ธีโร, วนนฺเต รมิโต สิยา.}

\prose{ฌาเยถ รุกฺขมูลสฺมิํ, อตฺตานมภิโตสยํ. (31)}

\proseitem{715}{ตโต รตฺยา วิวสาเน\footnote{วิวสเน (สี, สฺยา, อิ)}, คามนฺตมภิหารเย.}

\prose{อวฺหานํ นาภินนฺเทยฺย, อภิหารญฺจ คามโต. (32)}

\proseitem{716}{น มุนี คามมาคมฺม, กุเลสุ สหสา จเร.}

\prose{ฆาเสสนํ ฉินฺนกโถ, น วาจํ ปยุตํ ภเณ. (33)}

\proseitem{717}{อลตฺถํ ยทิทํ สาธุ, นาลตฺถํ กุสลํ อิติ.}

\prose{อุภเยนว โส ตาที, รุกฺขํวุ’ปนิวตฺตติ\footnote{รุกฺขํวุ’ปติวตฺตติ (ก), รุกฺขํว อุปาติวตฺตติ (สฺยา)}. (34)}

\proseitem{718}{ส ปตฺตปาณิ วิจรนฺโต, อมูโค มูคสมฺมโต.}

\csromanmatikaanchor{mtk.272}
\csromantocmarkref{subsection}{10. โกกาลิกสุตฺต}{mtk.272}
\bookmark[dest={mtk.272},level=2]{10. โกกาลิกสุตฺต}
\prose{อปฺปํ ทานํ น หีเฬยฺย, ทาตารํ นาวชานิยา. (35)}

\gambhira{สุตฺตนิปาตปาฬิ}
\proseitem{719}{\csromanfolio{390}\csromansymbolfootnote{*}{อภิ 4. 77 ปิฏฺเฐปิ.}อุจฺจาวจา หิ ปฏิปทา, สมเณน ปกาสิตา.}

\vspace{\tipitakaparskip}
\csromancenter{น ปารํ ทิคุณํ ยนฺติ, นยิทํ เอกคุณํ มุตํ. (36)}
\proseitem{720}{ยสฺส จ วิสตา นตฺถิ, ฉินฺนโสตสฺส ภิกฺขุโน.}

\prose{กิจฺจากิจฺจปฺปหีนสฺส, ปริฬาโห น วิชฺชติ. (37)}

\csromangathameasure{โมเนยฺยํ เต อุปญฺญิสฺสํ, ขุรธารูปโม ภเว.}
\csromangathagroup{\csromangathastanza{โมเนยฺยํ เต อุปญฺญิสฺสํ, ขุรธารูปโม ภเว.}}

\proseitem{721}{ชิวฺหาย ตาลุมาหจฺจ, อุทเร สญฺญโต สิยา. (38)}

\csromangathameasure{อลีนจิตฺโต จ สิยา, น จาปิ พหุ จินฺตเย.}
\csromangathagroup{\csromangathastanza{อลีนจิตฺโต จ สิยา, น จาปิ พหุ จินฺตเย.}}

\proseitem{722}{นิรามคนฺโธ อสิโต, พฺรหฺมจริยปรายโณ. (39)}

\csromangathameasure{เอกาสนสฺส สิกฺเขถ, สมณูปาสนสฺส จ. \\ เอกตฺตํ โม นมกฺขาตํ, เอโก เจ อภิรมิสฺสสิ.}
\csromangathagroup{\csromangathastanza{เอกาสนสฺส สิกฺเขถ, สมณูปาสนสฺส จ. \\ เอกตฺตํ โม นมกฺขาตํ, เอโก เจ อภิรมิสฺสสิ.}}

\csromanheader{2}{อถ ภาหิสิ\footnote{ภาสิหิ (สี, สฺยา, อิ)} ทสทิสา. (40)}
\csromangathameasure{สุตฺวา ธีรานํ นิโฆสํ, ฌายีนํ กามจาคินํ.}
\csromangathagroup{\csromangathastanza{สุตฺวา ธีรานํ นิโฆสํ, ฌายีนํ กามจาคินํ.}}

\vspace{\tipitakaparskip}
\csromancenter{ตโต หิริญฺจ สทฺธญฺจ, ภิยฺโย กุพฺเพถ มามโก. (41)}
\vspace{0.5\baselineskip}
\csromangathameasure{ตํ นทีหิ วิชานาถ, โสพฺเภสุ ปทเรสุ จ.}
\csromangathagroup{\csromangathastanza{ตํ นทีหิ วิชานาถ, โสพฺเภสุ ปทเรสุ จ.}}

\vspace{\tipitakaparskip}
\csromancenter{สณนฺตา ยนฺติ กุโสพฺภา\footnote{กุสฺสุพฺภา (สี)}, ตุณฺหี ยนฺติ มโหทธี. (42)}
\vspace{0.5\baselineskip}
\csromangathameasure{ยทูนกํ ตํ สณติ, ยํ ปูรํ สนฺตเมว ตํ.}
\csromangathagroup{\csromangathastanza{ยทูนกํ ตํ สณติ, ยํ ปูรํ สนฺตเมว ตํ.}}

\vspace{\tipitakaparskip}
\csromancenter{อฑฺฒกุพฺภูปโม พาโล, รหโท ปูโรว ปณฺฑิโต. (43)}
\vspace{0.5\baselineskip}
\csromangathameasure{ยํ สมโณ พหุํ ภาสติ, อุเปตํ อตฺถสญฺหิตํ.}
\csromangathagroup{\csromangathastanza{ยํ สมโณ พหุํ ภาสติ, อุเปตํ อตฺถสญฺหิตํ.}}

\vspace{\tipitakaparskip}
\csromancenter{ชานํ โส ธมฺมํ เทเสติ, ชานํ โส พหุ ภาสติ. (44)}
\vspace{0.5\baselineskip}
\csromangathameasure{โย จ ชานํ สํยตตฺโต, ชานํ น พหุ ภาสติ. \\ ส มุนี โมนมรหติ, ส มุนี โมนมชฺฌคาติ. (45) นาลกสุตฺตํ เอกาทสมํ.}
\csromangathagroup{\csromangathastanza{โย จ ชานํ สํยตตฺโต, ชานํ น พหุ ภาสติ. \\ ส มุนี โมนมรหติ, ส มุนี โมนมชฺฌคาติ. (45) นาลกสุตฺตํ เอกาทสมํ.}}

\csromanheader{2}{12. ทฺวยตานุปสฺสนาสุตฺต}
\prose{เอวํ เม สุตํ\csromandash{}เอกํ สมยํ ภควา สาวตฺถิยํ วิหรติ ปุพฺพาราเม มิคารมาตุปาสาเท.\csromanspacer{} เตน โข ปน สมเยน ภควา ตทหุโปสเถ}

\vspace{\tipitakaparskip}
\csromancenter{มหาวคฺค ปนฺนรเส ปุณฺณาย ปุณฺณมาย รตฺติยา ภิกฺขุสํฆปริวุโต อพฺโภกาเส นิสินฺโน โหติ.\csromanspacer{} อถ โข ภควา ตุณฺหีภูตํ ตุณฺหีภูตํ ภิกฺขุสํฆํ อนุวิโลเกตฺวา ภิกฺขู อามนฺเตสิ\csromandash{}}
\prose{\csromanfolio{391}เย เต ภิกฺขเว กุสลา ธมฺมา อริยา นิยฺยานิกา สมฺโพธคามิโน, เตสํ โว ภิกฺขเว กุสลานํ ธมฺมานํ อริยานํ นิยฺยานิกานํ สมฺโพธคามีนํ กา อุปนิสา สวนายาติ อิติ เจ ภิกฺขเว ปุจฺฉิตาโร อสฺสุ, เต เอวมสฺสุ วจนียา\csromandash{}ยาวเทว ทฺวยตานํ ธมฺมานํ ยถาภูตํ ญาณายาติ.\csromanspacer{} กิญฺจ ทฺวยตํ วเทถ?}

\prose{(1) อิทํ ทุกฺขํ, อยํ ทุกฺขสมุทโยติ อยเมกานุปสฺสนา.\csromanspacer{} อยํ ทุกฺขนิโรโธ, อยํ ทุกฺขนิโรธคามินี ปฏิปทาติ อยํ ทุติยานุปสฺสนา.\csromanspacer{} เอวํ สมฺมา ทฺวยตานุปสฺสิโน โข ภิกฺขเว ภิกฺขุโน อปฺปมตฺตสฺส อาตาปิโน ปหิตตฺตสฺส วิหรโต ทฺวินฺนํ ผลานํ อญฺญตรํ ผลํ ปาฏิกงฺขํ\csromandash{}ทิฏฺเฐว ธมฺเม อญฺญา, สติ วา อุปาทิเสเส อนาคามิตาติ.}

\prose{อิทมโวจ ภควา, อิทํ วตฺวาน สุคโต, อถาปรํ เอตทโวจ สตฺถา\csromandash{}}

\proseitem{729}{“เย ทุกฺขํ นปฺปชานนฺติ, อโถ ทุกฺขสฺส สมฺภวํ.}

\csromangathameasure{ยตฺถ จ สพฺพโส ทุกฺขํ, อเสสํ อุปรุชฺฌติ.}
\csromangathagroup{\csromangathastanza{ยตฺถ จ สพฺพโส ทุกฺขํ, อเสสํ อุปรุชฺฌติ.}}

\prose{ตญฺจ มคฺคํ น ชานนฺติ, ทุกฺขูปสมคามินํ. (1)}

\proseitem{730}{เจโตวิมุตฺติหีนา เต, อโถ ปญฺญาวิมุตฺติยา.}

\prose{อภพฺพา เต อนฺตกิริยาย, เต เว ชาติชรูปคา. (2)}

\proseitem{731}{เย จ ทุกฺขํ ปชานนฺติ, อโถ ทุกฺขสฺส สมฺภวํ.}

\csromangathameasure{ยตฺถ จ สพฺพโส ทุกฺขํ, อเสสํ อุปรุชฺฌติ.}
\csromangathagroup{\csromangathastanza{ยตฺถ จ สพฺพโส ทุกฺขํ, อเสสํ อุปรุชฺฌติ.}}

\prose{ตญฺจ มคฺคํ ปชานนฺติ, ทุกฺขูปสมคามินํ. (3)}

\proseitem{732}{เจโตวิมุตฺติสมฺปนฺนา, อโถ ปญฺญาวิมุตฺติยา.}

\prose{ภพฺพา เต อนฺตกิริยาย, น เต ชาติชรูปคา”ติ. (4)}

\prose{(2) สิยา อญฺเญนปิ ปริยาเยน สมฺมา ทฺวยตานุปสฺสนาติ อิติ เจ ภิกฺขเว ปุจฺฉิตาโร อสฺสุ. “สิยา”ติสฺสุ วจนียา.\csromanspacer{} กถญฺจ สิยา? ยํ กิญฺจิ ทุกฺขํ สมฺโภติ, สพฺพํ อุปธิปจฺจยาติ อยเมกานุปสฺสนา.}

\gambhira{สุตฺตนิปาตปาฬิ}
\prose{\csromanfolio{392}อุปธีนํ เตฺวว อเสสวิราคนิโรธา นตฺถิ ทุกฺขสฺส สมฺภโวติ อยํ ทุติยานุปสฺสนา.\csromanspacer{} เอวํ สมฺมา ฯเปฯ อถาปรํ เอตทโวจ สตฺถา\csromandash{}}

\proseitem{733}{“อุปธินิทานา ปภวนฺติ ทุกฺขา,}

\prose{เย เกจิ โลกสฺมิมเนกรูปา.}

\csromangathameasure{โย เว อวิทฺวา อุปธิํ กโรติ, \\ ปุนปฺปุนํ ทุกฺขมุเปติ มนฺโท. \\ ตสฺมา ปชานํ อุปธิํ น กยิรา,}
\csromangathagroup{\csromangathastanza{โย เว อวิทฺวา อุปธิํ กโรติ, \\ ปุนปฺปุนํ ทุกฺขมุเปติ มนฺโท. \\ ตสฺมา ปชานํ อุปธิํ น กยิรา,}}

\csromanheader{2}{ทุกฺขสฺส ชาติปฺปภวานุปสฺสี”ติ. (5)}
\prose{(3) สิยา อญฺเญนปิ ปริยาเยน สมฺมา ทฺวยตานุปสฺสนาติ, อิติ เจ ภิกฺขเว ปุจฺฉิตาโร อสฺสุ. “สิยา”ติสฺสุ วจนียา.\csromanspacer{} กถญฺจ สิยา? ยํ กิญฺจิ ทุกฺขํ สมฺโภติ, สพฺพํ อวิชฺชาปจฺจยาติ อยเมกานุปสฺสนา.\csromanspacer{} อวิชฺชาย เตฺวว อเสสวิราคนิโรธา นตฺถิ ทุกฺขสฺส สมฺภโวติ อยํ ทุติยานุปสฺสนา.\csromanspacer{} เอวํ สมฺมา ฯเปฯ อถาปรํ เอตทโวจ สตฺถา\csromandash{}}

\proseitem{734}{“ชาติมรณสํสารํ, เย วชนฺติ ปุนปฺปุนํ.}

\prose{อิตฺถภาวญฺญถาภาวํ, อวิชฺชาเยว สา คติ. (6)}

\proseitem{735}{อวิชฺชา หายํ มหาโมโห, เยนิทํ สํสิตํ จิรํ.}

\prose{วิชฺชาคตา จ เย สตฺตา, น เต คจฺฉนฺติ\footnote{นาคจฺฉนฺติ (สี, อิ)} ปุนพฺภวนฺ”ติ. (7)}

\prose{(4) สิยา อญฺเญนปิ ฯเปฯ.\csromanspacer{} กถญฺจ สิยา? ยํ กิญฺจิ ทุกฺขํ สมฺโภติ, สพฺพํ สงฺขารปจฺจยาติ อยเมกานุปสฺสนา.\csromanspacer{} สงฺขารานํ เตฺวว อเสสวิราคนิโรธา นตฺถิ ทุกฺขสฺส สมฺภโวติ อยํ ทุติยานุปสฺสนา.\csromanspacer{} เอวํ สมฺมา ฯเปฯ อถาปรํ เอตทโวจ สตฺถา\csromandash{}}

\proseitem{736}{“ยํ กิญฺจิ ทุกฺขํ สมฺโภติ, สพฺพํ สงฺขารปจฺจยา.}

\prose{สงฺขารานํ นิโรเธน, นตฺถิ ทุกฺขสฺส สมฺภโว. (8)}

\proseitem{737}{เอตมาทีนวํ ญตฺวา, ทุกฺขํ สงฺขารปจฺจยา.}

\csromangathameasure{สพฺพสงฺขารสมถา, สญฺญานํ อุปโรธนา.}
\csromangathagroup{\csromangathastanza{สพฺพสงฺขารสมถา, สญฺญานํ อุปโรธนา.}}

\prose{เอวํ ทุกฺขกฺขโย โหติ, เอตํ ญตฺวา ยถาตถํ. (9)}

\csromanhangingbegin
\hangingheaditem{738}{สมฺมทฺทสา เวทคุโน, สมฺมทญฺญาย ปณฺฑิตา.}
\hangingbody{อภิภุยฺย มารสํโยคํ, น คจฺฉนฺติ1 ปุนพฺภวนฺ”ติ.}
\csromanhangingend

\csromanheader{2}{(10)}
\chapterheadpageread{3. มหาวคฺค}
\prose{\csromanfolio{393}(5) สิยา อญฺเญนปิ ฯเปฯ.\csromanspacer{} กถญฺจ สิยา? ยํ กิญฺจิ ทุกฺขํ สมฺโภติ, สพฺพํ วิญฺญาณปจฺจยาติ อยเมกานุปสฺสนา.\csromanspacer{} วิญฺญาณสฺส เตฺวว อเสสวิราคนิโรธา นตฺถิ ทุกฺขสฺส สมฺภโวติ อยํ ทุติยานุปสฺสนา.\csromanspacer{} เอวํ สมฺมา ฯเปฯ อถาปรํ เอตทโวจ สตฺถา\csromandash{}}

\proseitem{739}{“ยํ กิญฺจิ ทุกฺขํ สมฺโภติ, สพฺพํ วิญฺญาณปจฺจยา.}

\prose{วิญฺญาณสฺส นิโรเธน, นตฺถิ ทุกฺขสฺส สมฺภโว. (11)}

\proseitem{740}{เอตมาทีนวํ ญตฺวา, ทุกฺขํ วิญฺญาณปจฺจยา.}

\prose{วิญฺญาณูปสมา ภิกฺขุ, นิจฺฉาโต ปรินิพฺพุโต”ติ. (12)}

\prose{(6) สิยา อญฺเญนปิ ฯเปฯ.\csromanspacer{} กถญฺจ สิยา? ยํ กิญฺจิ ทุกฺขํ สมฺโภติ, สพฺพํ ผสฺสปจฺจยาติ อยเมกานุปสฺสนา.\csromanspacer{} ผสฺสสฺส เตฺวว อเสสวิราคนิโรธา นตฺถิ ทุกฺขสฺส สมฺภโวติ อยํ ทุติยานุปสฺสนา.\csromanspacer{} เอวํ สมฺมา ฯเปฯ อถาปรํ เอตทโวจ สตฺถา\csromandash{}}

\proseitem{741}{“เตสํ ผสฺสปเรตานํ, ภวโสตานุสารินํ.}

\prose{กุมฺมคฺคปฏิปนฺนานํ, อารา สํโยชนกฺขโย. (13)}

\proseitem{742}{เย จ ผสฺสํ ปริญฺญาย, อญฺญายุปสเม\footnote{ปญฺญาย อุปสเม (สฺยา)} รตา.}

\prose{เต เว ผสฺสาภิสมยา, นิจฺฉาตา ปรินิพฺพุตา”ติ. (14)}

\prose{(7) สิยา อญฺเญนปิ ฯเปฯ.\csromanspacer{} กถญฺจ สิยา? ยํ กิญฺจิ ทุกฺขํ สมฺโภติ, สพฺพํ เวทนาปจฺจยาติ อยเมกานุปสฺสนา.\csromanspacer{} เวทนานํ เตฺวว อเสสวิราคนิโรธา นตฺถิ ทุกฺขสฺส สมฺภโวติ อยํ ทุติยานุปสฺสนา.\csromanspacer{} เอวํ สมฺมา ฯเปฯ อถาปรํ เอตทโวจ สตฺถา\csromandash{}}

\proseitem{743}{“สุขํ วา ยทิ วา ทุกฺขํ, อทุกฺขมสุขํ สห.}

\prose{อชฺฌตฺตญฺจ พหิทฺธา จ, ยํ กิญฺจิ อตฺถิ เวทิตํ. (15)}

\proseitem{744}{เอตํ ทุกฺขนฺติ ญตฺวาน, โมสธมฺมํ ปโลกินํ\footnote{ปโลกิตํ (สี)}.}

\csromangathameasure{ผุสฺส ผุสฺส วยํ ปสฺสํ, เอวํ ตตฺถ วิชานติ.}
\csromangathagroup{\csromangathastanza{ผุสฺส ผุสฺส วยํ ปสฺสํ, เอวํ ตตฺถ วิชานติ\footnote{วิรชฺชติ (ก-สี)}.}}

\prose{เวทนานํ ขยา ภิกฺขุ, นิจฺฉาโต ปรินิพฺพุโต”ติ. (16)}

\gambhira{สุตฺตนิปาตปาฬิ}
\prose{\csromanfolio{394}(8) สิยา อญฺเญนปิ ฯเปฯ.\csromanspacer{} กถญฺจ สิยา? ยํ กิญฺจิ ทุกฺขํ สมฺโภติ, สพฺพํ ตณฺหาปจฺจยาติ อยเมกานุปสฺสนา.\csromanspacer{} ตณฺหาย เตฺวว อเสสวิราคนิโรธา นตฺถิ ทุกฺขสฺส สมฺภโวติ อยํ ทุติยานุปสฺสนา.\csromanspacer{} เอวํ สมฺมา ฯเปฯ อถาปรํ เอตทโวจ สตฺถา\csromandash{}}

\csromangathameasure{“ตณฺหาทุติโย ปุริโส, ทีฆมทฺธาน สํสรํ.}
\csromangathagroup{\csromangathastanza{“ตณฺหาทุติโย ปุริโส, ทีฆมทฺธาน สํสรํ.}}

\proseitem{745}{อิตฺถภาวญฺญถาภาวํ, สํสารํ นาติวตฺตติ. (17)}

\csromangathameasure{เอตมาทีนวํ ญตฺวา, ตณฺหํ ทุกฺขสฺส สมฺภวํ.}
\csromangathagroup{\csromangathastanza{เอตมาทีนวํ ญตฺวา, ตณฺหํ\footnote{ตณฺหา (พหูสุ) อิติวุตฺตเก 201 ปิฏฺเฐปิ ปสฺสิตพฺพํ.} ทุกฺขสฺส สมฺภวํ.}}

\proseitem{746}{วีตตโณฺห อนาทาโน, สโต ภิกฺขุ ปริพฺพเช”ติ. (18)}

\prose{(9) สิยา อญฺเญนปิ ฯเปฯ.\csromanspacer{} กถญฺจ สิยา? ยํ กิญฺจิ ทุกฺขํ สมฺโภติ, สพฺพํ อุปาทานปจฺจยาติ อยเมกานุปสฺสนา.\csromanspacer{} อุปาทานานํ\footnote{อุปาทานสฺส (สฺยา, ก)} เตฺวว อเสสวิราคนิโรธา นตฺถิ ทุกฺขสฺส สมฺภโวติ อยํ ทุติยานุปสฺสนา.\csromanspacer{} เอวํ สมฺมา ฯเปฯ อถาปรํ เอตทโวจ สตฺถา\csromandash{}}

\csromangathameasure{“อุปาทานปจฺจยา ภโว, ภูโต ทุกฺขํ นิคจฺฉติ.}
\csromangathagroup{\csromangathastanza{“อุปาทานปจฺจยา ภโว, ภูโต ทุกฺขํ นิคจฺฉติ.}}

\vspace{\tipitakaparskip}
\csromancenter{ชาตสฺส มรณํ โหติ, เอโส ทุกฺขสฺส สมฺภโว. (19)}
\vspace{0.5\baselineskip}
\csromangathameasure{ตสฺมา อุปาทานกฺขยา, สมฺมทญฺญาย ปณฺฑิตา.}
\csromangathagroup{\csromangathastanza{ตสฺมา อุปาทานกฺขยา, สมฺมทญฺญาย ปณฺฑิตา.}}

\vspace{\tipitakaparskip}
\csromancenter{ชาติกฺขยํ อภิญฺญาย, น คจฺฉนฺติ ปุนพฺภวนฺ”ติ. (20)}
\prose{(10) สิยา อญฺเญนปิ ฯเปฯ.\csromanspacer{} กถญฺจ สิยา? ยํ กิญฺจิ ทุกฺขํ สมฺโภติ, สพฺพํ อารมฺภปจฺจยาติ อยเมกานุปสฺสนา.\csromanspacer{} อารมฺภานํ เตฺวว อเสสวิราคนิโรธา นตฺถิ ทุกฺขสฺส สมฺภโวติ อยํ ทุติยานุปสฺสนา.\csromanspacer{} เอวํ สมฺมา ฯเปฯ อถาปรํ เอตทโวจ สตฺถา\csromandash{}}

\csromangathameasure{“ยํ กิญฺจิ ทุกฺขํ สมฺโภติ, สพฺพํ อารมฺภปจฺจยา.}
\csromangathagroup{\csromangathastanza{“ยํ กิญฺจิ ทุกฺขํ สมฺโภติ, สพฺพํ อารมฺภปจฺจยา.}}

\vspace{\tipitakaparskip}
\csromancenter{อารมฺภานํ นิโรเธน, นตฺถิ ทุกฺขสฺส สมฺภโว. (21)}
\vspace{0.5\baselineskip}
\csromangathameasure{เอตมาทีนวํ ญตฺวา, ทุกฺขํ อารมฺภปจฺจยา.}
\csromangathagroup{\csromangathastanza{เอตมาทีนวํ ญตฺวา, ทุกฺขํ อารมฺภปจฺจยา.}}

\vspace{\tipitakaparskip}
\csromancenter{สพฺพารมฺภํ ปฏินิสฺสชฺช, อนารมฺเภ วิมุตฺติโน. (22)}
\csromancenter{\csromansymbolfootnote{*}{ขุ 10.177 ปิฏฺเฐปิ.}อุจฺฉินฺนภวตณฺหสฺส, สนฺตจิตฺตสฺส ภิกฺขุโน.}
\csromancenter{วิกฺขีโณ\footnote{วิติณฺโณ (สี)} ชาติสํสาโร, นตฺถิ ตสฺส ปุนพฺภโว”ติ. (23)}
\chapterheadpageread{3. มหาวคฺค}
\prose{\csromanfolio{395}(11) สิยา อญฺเญนปิ ฯเปฯ.\csromanspacer{} กถญฺจ สิยา? ยํ กิญฺจิ ทุกฺขํ สมฺโภติ, สพฺพํ อาหารปจฺจยาติ อยเมกานุปสฺสนา.\csromanspacer{} อาหารานํ เตฺวว อเสสวิราคนิโรธา นตฺถิ ทุกฺขสฺส สมฺภโวติ อยํ ทุติยานุปสฺสนา.\csromanspacer{} เอวํ สมฺมา ฯเปฯ อถาปรํ เอตทโวจ สตฺถา\csromandash{}}

\proseitem{752}{“ยํ กิญฺจิ ทุกฺขํ สมฺโภติ, สพฺพํ อาหารปจฺจยา.}

\prose{อาหารานํ นิโรเธน, นตฺถิ ทุกฺขสฺส สมฺภโว. (24)}

\proseitem{753}{เอตมาทีนวํ ญตฺวา, ทุกฺขํ อาหารปจฺจยา.}

\prose{สพฺพาหารํ ปริญฺญาย, สพฺพาหารมนิสฺสิโต. (25)}

\csromanmatikaanchor{mtk.273}
\csromantocmarkref{subsection}{11. นาลกสุตฺต}{mtk.273}
\bookmark[dest={mtk.273},level=2]{11. นาลกสุตฺต}
\csromanhangingbegin
\hangingheaditem{754}{อาโรคฺยํ สมฺมทญฺญาย, อาสวานํ ปริกฺขยา.}
\hangingbody{สงฺขาย เสวี ธมฺมฏฺโฐ, สงฺขฺยํ1 โนเปติ เวทคู”ติ. (26)}
\csromanhangingend

\prose{(12) สิยา อญฺเญนปิ ฯเปฯ.\csromanspacer{} กถญฺจ สิยา? ยํ กิญฺจิ ทุกฺขํ สมฺโภติ, สพฺพํ อิญฺชิตปจฺจยาติ อยเมกานุปสฺสนา.\csromanspacer{} อิญฺชิตานํ เตฺวว อเสสวิราคนิโรธา นตฺถิ ทุกฺขสฺส สมฺภโวติ อยํ ทุติยานุปสฺสนา.\csromanspacer{} เอวํ สมฺมา ฯเปฯ อถาปรํ เอตทโวจ สตฺถา\csromandash{}}

\proseitem{755}{“ยํ กิญฺจิ ทุกฺขํ สมฺโภติ, สพฺพํ อิญฺชิตปจฺจยา.}

\prose{อิญฺชิตานํ นิโรเธน, นตฺถิ ทุกฺขสฺส สมฺภโว. (27)}

\proseitem{756}{เอตมาทีนวํ ญตฺวา, ทุกฺขํ อิญฺชิตปจฺจยา.}

\csromangathameasure{ตสฺมา หิ เอชํ โวสชฺช, สงฺขาเร อุปรุนฺธิย.}
\csromangathagroup{\csromangathastanza{ตสฺมา หิ เอชํ โวสชฺช, สงฺขาเร อุปรุนฺธิย.}}

\prose{อเนโช อนุปาทาโน, สโต ภิกฺขุ ปริพฺพเช”ติ. (28)}

\prose{(13) สิยา อญฺเญนปิ ฯเปฯ.\csromanspacer{} กถญฺจ สิยา? นิสฺสิตสฺส จลิตํ โหตีติ อยเมกานุปสฺสนา.\csromanspacer{} อนิสฺสิโต น จลตีติ อยํ ทุติยานุปสฺสนา.\csromanspacer{} เอวํ สมฺมา ฯเปฯ อถาปรํ เอตทโวจ สตฺถา\csromandash{}}

\proseitem{757}{“อนิสฺสิโต น จลติ, นิสฺสิโต จ อุปาทิยํ.}

\prose{อิตฺถภาวญฺญถาภาวํ, สํสารํ นาติวตฺตติ. (29)}

\proseitem{758}{เอตมาทีนวํ ญตฺวา, นิสฺสเยสุ มหพฺภยํ.}

\prose{อนิสฺสิโต อนุปาทาโน, สโต ภิกฺขุ ปริพฺพเช”ติ. (30)}

\gambhira{สุตฺตนิปาตปาฬิ}
\prose{\csromanfolio{396}(14) สิยา อญฺเญนปิ ฯเปฯ.\csromanspacer{} กถญฺจ สิยา? รูเปหิ ภิกฺขเว อรูปา\footnote{อารุปฺปา (สี, อิ)} สนฺตตราติ อยเมกานุปสฺสนา.\csromanspacer{} อรูเปหิ นิโรโธ สนฺตตโรติ อยํ ทุติยานุปสฺสนา.\csromanspacer{} เอวํ สมฺมา ฯเปฯ อถาปรํ เอตทโวจ สตฺถา\csromandash{}}

\proseitem{759}{“เย จ รูปูปคา สตฺตา, เย จ อรูปฏฺฐายิโน\footnote{อารุปฺปวาสิโน (สี, อิ)}.}

\prose{นิโรธํ อปฺปชานนฺตา, อาคนฺตาโร ปุนพฺภวํ. (31)}

\proseitem{760}{เย จ รูเป ปริญฺญาย, อรูเปสุ อสณฺฐิตา\footnote{สุสณฺฐิตา (สี, สฺยา, อิ)}.}

\prose{นิโรเธ เย วิมุจฺจนฺติ, เต ชนา มจฺจุหายิโน”ติ. (32)}

\prose{(15) สิยา อญฺเญนปิ ฯเปฯ.\csromanspacer{} กถญฺจ สิยา? ยํ ภิกฺขเว สเทวกสฺส โลกสฺส สมารกสฺส สพฺรหฺมกสฺส สสฺสมณพฺราหฺมณิยา ปชาย สเทวมนุสฺสาย “อิทํ สจฺจนฺ”ติ อุปนิชฺฌายิตํ, ต’ทมริยานํ “เอตํ มุสา”ติ ยถาภูตํ สมฺมปฺปญฺญาย สุทิฏฺฐํ, อยเมกานุปสฺสนา.\csromanspacer{} ยํ ภิกฺขเว สเทวกสฺส ฯเปฯ สเทวมนุสฺสาย “อิทํ มุสา”ติ อุปนิชฺฌายิตํ, ต’ทมริยานํ “เอตํ สจฺจนฺ”ติ ยถาภูตํ สมฺมปฺปญฺญาย สุทิฏฺฐํ, อยํ ทุติยานุปสฺสนา.\csromanspacer{} เอวํ สมฺมา ฯเปฯ อถาปรํ เอตทโวจ สตฺถา\csromandash{}}

\proseitem{761}{“อนตฺตนิ อตฺตมานิํ\footnote{อตฺตมานี (สฺยา), อตฺตมานํ (อิ, ก)}, ปสฺส โลกํ สเทวกํ.}

\prose{นิวิฏฺฐํ นามรูปสฺมิํ, อิทํ สจฺจนฺติ มญฺญติ. (33)}

\proseitem{762}{เยน เยน หิ มญฺญนฺติ, ตโต ตํ โหติ อญฺญถา.}

\prose{ตญฺหิ ตสฺส มุสา โหติ, โมสธมฺมญฺหิ อิตฺตรํ. (34)}

\proseitem{763}{อโมสธมฺมํ นิพฺพานํ, ตทริยา สจฺจโต วิทู.}

\prose{เต เว สจฺจาภิสมยา, นิจฺฉาตา ปรินิพฺพุตา”ติ. (35)}

\prose{(16) สิยา อญฺเญนปิ ปริยาเยน สมฺมา ทฺวยตานุปสฺสนาติ, อิติ เจ ภิกฺขเว ปุจฺฉิตาโร อสฺสุ. “สิยา”ติสฺสุ วจนียา.\csromanspacer{} กถญฺจ สิยา? ยํ ภิกฺขเว สเทวกสฺส โลกสฺส สมารกสฺส สพฺรหฺมกสฺส สสฺสมณพฺราหฺมณิยา ปชาย สเทวมนุสฺสาย “อิทํ สุขนฺ”ติ อุปนิชฺฌายิตํ, ต’ทมริยานํ “เอตํ ทุกฺขนฺ”ติ ยถาภูตํ สมฺมปฺปญฺญาย สุทิฏฺฐํ, อยเมกานุปสฺสนา.\csromanspacer{} ยํ ภิกฺขเว สเทวกสฺส ฯเปฯ สเทวมนุสฺสาย “อิทํ ทุกฺขนฺ”ติ อุปนิชฺฌายิตํ, ต’ทมริยานํ “เอตํ สุขนฺ”ติ ยถาภูตํ สมฺมปฺปญฺญาย สุทิฏฺฐํ,}

\vspace{\tipitakaparskip}
\csromancenter{มหาวคฺค อยํ ทุติยานุปสฺสนา.\csromanspacer{} เอวํ สมฺมา ทฺวยตานุปสฺสิโน โข ภิกฺขเว ภิกฺขุโน อปฺปมตฺตสฺส อาตาปิโน ปหิตตฺตสฺส วิหรโต ทฺวินฺนํ ผลานํ อญฺญตรํ ผลํ ปาฏิกงฺขํ\csromandash{}ทิฏฺเฐว ธมฺเม อญฺญา, สติ วา อุปาทิเสเส อนาคามิตาติ.\csromanspacer{} อิทมโวจ ภควา, อิทํ วตฺวาน สุคโต, อถาปรํ เอตทโวจ สตฺถา\csromandash{}}
\proseitem{764}{\csromanfolio{397}“รูปา สทฺทา รสา คนฺธา, ผสฺสา ธมฺมา จ เกวลา.}

\prose{อิฏฺฐา กนฺตา มนาปา จ, ยาวต’ตฺถีติ วุจฺจติ. (36)}

\proseitem{765}{สเทวกสฺส โลกสฺส, เอเต โว สุขสมฺมตา.}

\prose{ยตฺถ เจเต นิรุชฺฌนฺติ, ตํ เนสํ ทุกฺขสมฺมตํ. (37)}

\proseitem{766}{สุขนฺติ ทิฏฺฐมริเยหิ, สกฺกายสฺสุ’ปโรธนํ.}

\prose{ปจฺจนีกมิทํ โหติ, สพฺพโลเกน ปสฺสตํ. (38)}

\proseitem{767}{ยํ ปเร สุขโต อาหุ, ตทริยา อาหุ ทุกฺขโต.}

\prose{ยํ ปเร ทุกฺขโต อาหุ, ตทริยา สุขโต วิทู. (39)}

\proseitem{768}{ปสฺส ธมฺมํ ทุราชานํ, สมฺปมูเฬฺหตฺถ’วิทฺทสุ\footnote{สมฺปมูเฬฺหตฺถ อวิทฺทสุ (สี, อิ), สมฺมูเฬฺหตฺถ อวิทฺทสุ (?)}.}

\prose{นิวุตานํ ตโม โหติ, อนฺธกาโร อปสฺสตํ. (40)}

\proseitem{769}{สตญฺจ วิวฏํ โหติ, อาโลโก ปสฺสตามิว.}

\prose{สนฺติเก น วิชานนฺติ, มคา ธมฺมสฺส’โกวิทา. (41)}

\proseitem{770}{ภวราคปเรเตหิ, ภวโสตานุสาริภิ.}

\prose{มารเธยฺยานุปนฺเนหิ, นายํ ธมฺโม สุสมฺพุโธ. (42)}

\proseitem{771}{โก นุ อญฺญตฺร มริเยหิ, ปทํ สมฺพุทฺธุมรหติ.}

\prose{ยํ ปทํ สมฺมทญฺญาย, ปรินิพฺพนฺติ อนาสวา”ติ. (43)}

\prose{อิทมโวจ ภควา, อตฺตมนา เต ภิกฺขู ภควโต ภาสิตํ อภินนฺทุนฺติ.\csromanspacer{} อิมสฺมิํ จ\footnote{อิมสฺมิํ โข (สี)} ปน เวยฺยากรณสฺมิํ ภญฺญมาเน สฏฺฐิมตฺตานํ ภิกฺขูนํ อนุปาทาย อาสเวหิ จิตฺตานิ วิมุจฺจิํสูติ.\csromanspacer{} ทฺวยตานุปสฺสนาสุตฺตํ ทฺวาทสมํ.}

\gambhira{สุตฺตนิปาตปาฬิ}
\titlehead{\textbf{ตสฺสุทฺทานํ}}
\csromangathameasure{สจฺจํ อุปธิ อวิชฺชา จ, สงฺขาเร วิญฺญาณปญฺจมํ. \\ ผสฺสเวทนิยา ตณฺหา, อุปาทานารมฺภ-อาหารา. \\ อิญฺชิตํ จลิตํ รูปํ, สจฺจํ ทุกฺเขน โสฬสาติ. \\ \textbf{มหาวคฺโค ตติโย.}}
\csromangathagroup{\csromangathastanza{\csromanfolio{398}สจฺจํ อุปธิ อวิชฺชา จ, สงฺขาเร วิญฺญาณปญฺจมํ. \\ ผสฺสเวทนิยา ตณฺหา, อุปาทานารมฺภ-อาหารา.}\csromangathastanza{อิญฺชิตํ จลิตํ รูปํ, สจฺจํ ทุกฺเขน โสฬสาติ. \\ \textbf{มหาวคฺโค ตติโย.}}}

\titlehead{\textbf{ตสฺสุทฺทานํ}}
\csromangathameasure{ปพฺพชฺชา จ ปธานญฺจ, สุภาสิตญฺจ สุนฺทริ. \\ มาฆสุตฺตํ สภิโย จ, เสโล สลฺลญฺจ วุจฺจติ. \\ วาเสฏฺโฐ จาปิ โกกาลิ, นาลโก ทฺวยตานุปสฺสนา. \\ ทฺวาทเสตานิ สุตฺตานิ, มหาวคฺโคติ วุจฺจตีติ.}
\csromangathagroup{\csromangathastanza{ปพฺพชฺชา จ ปธานญฺจ, สุภาสิตญฺจ สุนฺทริ. \\ มาฆสุตฺตํ สภิโย จ, เสโล สลฺลญฺจ วุจฺจติ.}\csromangathastanza{วาเสฏฺโฐ จาปิ โกกาลิ, นาลโก ทฺวยตานุปสฺสนา. \\ ทฺวาทเสตานิ สุตฺตานิ, มหาวคฺโคติ วุจฺจตีติ.}}

\chapterheadpageread{4. อฏฺฐกวคฺค}
\csromanheader{2}{1. กามสุตฺต}
\proseitem{772}{\csromanfolio{399}\csromansymbolfootnote{*}{ขุ 7. 1; ขุ 10. 6, 59 ปิฏฺเฐสุปิ.}กามํ กามยมานสฺส, ตสฺส เจ ตํ สมิชฺฌติ.}

\prose{อทฺธา ปีติมโน โหติ, ลทฺธา มจฺโจ ยทิจฺฉติ. (1)}

\proseitem{773}{\csromansymbolmark{*}ตสฺส เจ กามยานสฺส\footnote{กามยมานสฺส (ก)}, ฉนฺทชาตสฺส ชนฺตุโน.}

\prose{เต กามา ปริหายนฺติ, สลฺลวิทฺโธว รุปฺปติ. (2)}

\proseitem{774}{\csromansymbolmark{*}โย กาเม ปริวชฺเชติ, สปฺปสฺเสว ปทา สิโร.}

\prose{โสมํ\footnote{โส อิมํ (สี, อิ)} วิสตฺติกํ โลเก, สโต สมติวตฺตติ. (3)}

\proseitem{775}{\csromansymbolmark{*}เขตฺตํ วตฺถุํ หิรญฺญํ วา, ควสฺสํ\footnote{ควาสฺสํ (สี, สฺยา, อิ)} ทาสโปริสํ.}

\prose{ถิโย พนฺธู ปุถุ กาเม, โย นโร อนุคิชฺฌติ. (4)}

\proseitem{776}{\csromansymbolmark{*}อพลา นํ พลียนฺติ, มทฺทนฺเตนํ ปริสฺสยา.}

\vspace{\tipitakaparskip}
\csromancenter{ตโต นํ ทุกฺขมเนฺวติ, นาวํ ภินฺนมิโวทกํ. (5)}
\proseitem{777}{\csromansymbolmark{*}ตสฺมา ชนฺตุ สทา สโต, กามานิ ปริวชฺชเย.}

\csromangathameasure{เต ปหาย ตเร โอฆํ, นาวํ สิตฺวาว ปารคูติ. (6) กามสุตฺตํ ปฐมํ.}
\csromangathagroup{\csromangathastanza{เต ปหาย ตเร โอฆํ, นาวํ สิตฺวาว\footnote{สิญฺจิตฺวา (สี)} ปารคูติ. (6) กามสุตฺตํ ปฐมํ.}}

\csromanheader{2}{2. คุหฏฺฐกสุตฺต}
\csromangathameasure{สตฺโต คุหายํ พหุนาภิฉนฺโน, \\ ติฏฺฐํ นโร โมหนสฺมิํ ปคาโฬฺห. \\ ทูเร วิเวกา หิ ตถาวิโธ โส,}
\csromangathagroup{\csromangathastanza{สตฺโต คุหายํ พหุนาภิฉนฺโน, \\ ติฏฺฐํ นโร โมหนสฺมิํ ปคาโฬฺห. \\ ทูเร วิเวกา หิ ตถาวิโธ โส,}}

\csromanheader{2}{กามา หิ โลเก น หิ สุปฺปหายา. (1)}
\gambhira{สุตฺตนิปาตปาฬิ}
\csromangathameasure{อิจฺฉานิทานา ภวสาตพทฺธา, \\ เต ทุปฺปมุญฺจา น หิ อญฺญโมกฺขา. \\ ปจฺฉา ปุเร วาปิ อเปกฺขมานา,}
\csromangathagroup{\csromangathastanza{\csromanfolio{400}อิจฺฉานิทานา ภวสาตพทฺธา, \\ เต ทุปฺปมุญฺจา น หิ อญฺญโมกฺขา. \\ ปจฺฉา ปุเร วาปิ อเปกฺขมานา,}}

\csromanheader{2}{อิเม ว กาเม ปุริเม ว ชปฺปํ. (2)}
\proseitem{780}{\csromansymbolfootnote{*}{ขุ 10. 172 ปิฏฺเฐปิ.}กาเมสุ คิทฺธา ปสุตา ปมูฬฺหา,}

\prose{อวทานิยา เต วิสเม นิวิฏฺฐา.}

\csromangathameasure{ทุกฺขูปนีตา ปริเทวยนฺติ,}
\csromangathagroup{\csromangathastanza{ทุกฺขูปนีตา ปริเทวยนฺติ,}}

\csromanheader{2}{กิํสู ภวิสฺสาม อิโต จุตาเส. (3)}
\csromangathameasure{ตสฺมา หิ สิกฺเขถ อิเธว ชนฺตุ, \\ ยํ กิญฺจิ ชญฺญา วิสมนฺติ โลเก. \\ น ตสฺส เหตู วิสมํ จเรยฺย,}
\csromangathagroup{\csromangathastanza{ตสฺมา หิ สิกฺเขถ อิเธว ชนฺตุ, \\ ยํ กิญฺจิ ชญฺญา วิสมนฺติ โลเก. \\ น ตสฺส เหตู วิสมํ จเรยฺย,}}

\csromanheader{2}{อปฺปญฺหิ’ทํ ชีวิตมาหุ ธีรา. (4)}
\csromangathameasure{ปสฺสามิ โลเก ปริผนฺทมานํ, \\ ปชํ อิมํ ตณฺหคตํ ภเวสุ. \\ หีนา นรา มจฺจุมุเข ลปนฺติ,}
\csromangathagroup{\csromangathastanza{ปสฺสามิ โลเก ปริผนฺทมานํ, \\ ปชํ อิมํ ตณฺหคตํ ภเวสุ. \\ หีนา นรา มจฺจุมุเข ลปนฺติ,}}

\csromanheader{2}{อวีตตณฺหาเส ภวาภเวสุ. (5)}
\csromangathameasure{มมายิเต ปสฺสถ ผนฺทมาเน, \\ มจฺเฉว อปฺโปทเก ขีณโสเต. \\ เอตมฺปิ ทิสฺวา อมโม จเรยฺย,}
\csromangathagroup{\csromangathastanza{มมายิเต ปสฺสถ ผนฺทมาเน, \\ มจฺเฉว อปฺโปทเก ขีณโสเต. \\ เอตมฺปิ ทิสฺวา อมโม จเรยฺย,}}

\csromanheader{2}{ภเวสุ อาสตฺติมกุพฺพมาโน. (6)}
\csromangathameasure{อุโภสุ อนฺเตสุ วิเนยฺย ฉนฺทํ, \\ ผสฺสํ ปริญฺญาย อนานุคิทฺโธ. \\ ยทตฺตครหี ตทกุพฺพมาโน,}
\csromangathagroup{\csromangathastanza{อุโภสุ อนฺเตสุ วิเนยฺย ฉนฺทํ, \\ ผสฺสํ ปริญฺญาย อนานุคิทฺโธ. \\ ยทตฺตครหี ตทกุพฺพมาโน,}}

\csromanheader{2}{น ลิปฺปตี\footnote{น ลิมฺปตี (สฺยา, ก) ขุ 7. 42 ปิฏฺเฐปิ.} ทิฏฺฐสุเตสุ ธีโร. (7)}
\csromangathameasure{สญฺญํ ปริญฺญา วิตเรยฺย โอฆํ, \\ ปริคฺคเหสุ มุนิ โนปลิตฺโต. \\ อพฺพูฬฺหสลฺโล จรมปฺปมตฺโต, \\ นาสีสตี โลกมิมํ ปรญฺจาติ. (8) คุหฏฺฐกสุตฺตํ ทุติยํ.}
\csromangathagroup{\csromangathastanza{สญฺญํ ปริญฺญา วิตเรยฺย โอฆํ, \\ ปริคฺคเหสุ มุนิ โนปลิตฺโต. \\ อพฺพูฬฺหสลฺโล จรมปฺปมตฺโต, \\ นาสีสตี\footnote{นาสิํสตี (สี, สฺยา, อิ)} โลกมิมํ ปรญฺจาติ. (8) คุหฏฺฐกสุตฺตํ ทุติยํ.}}

\csromanheader{2}{4. อฏฺฐกวคฺค \textbf{3. ทุฏฺฐฏฺฐกสุตฺต}}
\proseitem{786}{\csromanfolio{401}วทนฺติ เว ทุฏฺฐมนาปิ เอเก,}

\prose{อโถปิ เว สจฺจมนา วทนฺติ.}

\prose{วาทญฺจ ชาตํ มุนิ โน อุเปติ,}

\csromanheader{2}{ตสฺมา มุนี นตฺถิ ขิโล กุหิญฺจิ. (1)}
\proseitem{787}{สกญฺหิ ทิฏฺฐิํ กถมจฺจเยยฺย,}

\prose{ฉนฺทานุนีโต รุจิยา นิวิฏฺโฐ.}

\prose{สยํ สมตฺตานิ ปกุพฺพมาโน,}

\csromanheader{2}{ยถา หิ ชาเนยฺย ตถา วเทยฺย. (2)}
\proseitem{788}{โย อตฺตโน สีลวตานิ ชนฺตุ,}

\prose{อนานุปุฏฺโฐว ปเรส\footnote{ปรสฺส (ก) ขุ 7. 50-1 ปิฏฺเฐสุ ปสฺสิตพฺพํ.} ปาว\footnote{ปาวา (สี, สฺยา, อิ)}.}

\prose{อนริยธมฺมํ กุสลา ตมาหุ,}

\csromanheader{2}{โย อาตุมานํ สยเมว ปาว2. (3)}
\proseitem{789}{สนฺโต จ ภิกฺขุ อภินิพฺพุตตฺโต,}

\prose{อิติ’หนฺติ สีเลสุ อกตฺถมาโน.}

\prose{ตมริยธมฺมํ กุสลา วทนฺติ,}

\csromanheader{2}{ยสฺสุสฺสทา นตฺถิ กุหิญฺจิ โลเก. (4)}
\proseitem{790}{ปกปฺปิตา สงฺขตา ยสฺส ธมฺมา,}

\prose{ปุรกฺขตา\footnote{ปุเรกฺขตา (สี)} สนฺติ อวีวทาตา.}

\prose{ยทตฺตนิ ปสฺสติ อานิสํสํ,}

\csromanheader{2}{ตํ นิสฺสิโต กุปฺปปฏิจฺจสนฺติํ. (5)}
\proseitem{791}{ทิฏฺฐีนิเวสา น หิ สฺวาติวตฺตา,}

\prose{ธมฺเมสุ นิจฺเฉยฺย สมุคฺคหีตํ.}

\prose{ตสฺมา นโร เตสุ นิเวสเนสุ,}

\csromanheader{2}{นิรสฺสตี อาทิยตี จ ธมฺมํ. (6)}
\gambhira{สุตฺตนิปาตปาฬิ}
\csromangathameasure{โธนสฺส หิ นตฺถิ กุหิญฺจิ โลเก, \\ ปกปฺปิตา ทิฏฺฐิ ภวาภเวสุ. \\ มายญฺจ มานญฺจ ปหาย โธโน,}
\csromangathagroup{\csromangathastanza{\csromanfolio{402}โธนสฺส หิ นตฺถิ กุหิญฺจิ โลเก, \\ ปกปฺปิตา ทิฏฺฐิ ภวาภเวสุ. \\ มายญฺจ มานญฺจ ปหาย โธโน,}}

\proseitem{792}{ส เกน คจฺเฉยฺย, อนูปโย โส. (7)}

\csromangathameasure{อุปโย หิ ธมฺเมสุ อุเปติ วาทํ, \\ อนูปยํ เกน กถํ วเทยฺย. \\ อตฺตา นิรตฺตา น หิ ตสฺส อตฺถิ, \\ อโธสิ โส ทิฏฺฐิมิเธว สพฺพนฺติ. (8) ทุฏฺฐฏฺฐกสุตฺตํ ตติยํ.}
\csromangathagroup{\csromangathastanza{อุปโย หิ ธมฺเมสุ อุเปติ วาทํ, \\ อนูปยํ เกน กถํ วเทยฺย. \\ อตฺตา นิรตฺตา\footnote{อตฺตํ นิรตฺตํ (พหูสุ)} น หิ ตสฺส อตฺถิ, \\ อโธสิ โส ทิฏฺฐิมิเธว สพฺพนฺติ. (8) ทุฏฺฐฏฺฐกสุตฺตํ ตติยํ.}}

\csromanheader{2}{4. สุทฺธฏฺฐกสุตฺต}
\proseitem{794}{\csromansymbolfootnote{*}{ขุ 7. 64 ปิฏฺเฐปิ.}ปสฺสามิ สุทฺธํ ปรมํ อโรคํ,}

\prose{ทิฏฺเฐน สํสุทฺธิ นรสฺส โหติ.}

\csromangathameasure{เอวาภิชานํ “ปรมนฺ”ติ ญตฺวา,}
\csromangathagroup{\csromangathastanza{เอวาภิชานํ\footnote{เอตาภิชานํ (สี, อิ)} “ปรมนฺ”ติ ญตฺวา,}}

\csromanheader{2}{สุทฺธานุปสฺสีติ ปจฺเจติ ญาณํ. (1)}
\csromangathameasure{ทิฏฺเฐน เจ สุทฺธิ นรสฺส โหติ, \\ ญาเณน วา โส ปชหาติ ทุกฺขํ. \\ อญฺเญน โส สุชฺฌติ โสปธีโก,}
\csromangathagroup{\csromangathastanza{ทิฏฺเฐน เจ สุทฺธิ นรสฺส โหติ, \\ ญาเณน วา โส ปชหาติ ทุกฺขํ. \\ อญฺเญน โส สุชฺฌติ โสปธีโก,}}

\csromanheader{2}{ทิฏฺฐี หิ นํ ปาว ตถา วทานํ. (2)}
\csromangathameasure{น พฺราหฺมโณ อญฺญโต สุทฺธิมาห, \\ ทิฏฺเฐ สุเต สีลวเต มุเต วา. \\ ปุญฺเญ จ ปาเป จ อนูปลิตฺโต,}
\csromangathagroup{\csromangathastanza{น พฺราหฺมโณ อญฺญโต สุทฺธิมาห, \\ ทิฏฺเฐ สุเต สีลวเต มุเต วา. \\ ปุญฺเญ จ ปาเป จ อนูปลิตฺโต,}}

\csromanheader{2}{อตฺตญฺชโห นยิธ ปกุพฺพมาโน. (3)}
\csromangathameasure{ปุริมํ ปหาย อปรํ สิตาเส, \\ เอชานุคา เต น ตรนฺติ สงฺคํ. \\ เต อุคฺคหายนฺติ นิรสฺสชนฺติ,}
\csromangathagroup{\csromangathastanza{ปุริมํ ปหาย อปรํ สิตาเส, \\ เอชานุคา เต น ตรนฺติ สงฺคํ. \\ เต อุคฺคหายนฺติ นิรสฺสชนฺติ,}}

\csromanheader{2}{กปีว สาขํ ปมุญฺจํ คหายํ\footnote{ปมุขํ คหาย (สฺยา), ปมุญฺจ คหาย (ก)}. (4)}
\chapterheadpageread{4. อฏฺฐกวคฺค}
\proseitem{798}{\csromanfolio{403}สยํ สมาทาย วตานิ ชนฺตุ,}

\prose{อุจฺจาวจํ คจฺฉติ สญฺญสตฺโต.}

\prose{วิทฺวา จ เวเทหิ สเมจฺจ ธมฺมํ,}

\csromanheader{2}{น อุจฺจาวจํ คจฺฉติ ภูริปญฺโญ. (5)}
\proseitem{799}{ส สพฺพธมฺเมสุ วิเสนิภูโต,}

\prose{ยํ กิญฺจิ ทิฏฺฐํ ว สุตํ มุตํ วา.}

\prose{ตเมว ทสฺสิํ วิวฏํ จรนฺตํ,}

\csromanheader{2}{เกนีธ โลกสฺมิ วิกปฺปเยยฺย. (6)}
\proseitem{800}{น กปฺปยนฺติ น ปุเรกฺขโรนฺติ,}

\prose{อจฺจนฺตสุทฺธีติ น เต วทนฺติ.}

\prose{อาทานคนฺถํ คถิตํ วิสชฺช,}

\csromanmatikaanchor{mtk.274}
\csromantocmarkref{subsection}{12. ทฺวยตานุปสฺสนาสุตฺต}{mtk.274}
\bookmark[dest={mtk.274},level=2]{12. ทฺวยตานุปสฺสนาสุตฺต}
\csromanheader{2}{อาสํ น กุพฺพนฺติ กุหิญฺจิ โลเก. (7)}
\proseitem{801}{สีมาติโค พฺราหฺมโณ ตสฺส นตฺถิ,}

\prose{ญตฺวา ว ทิสฺวา ว\footnote{ญตฺวา จ ทิสฺวา จ (ก-สี, ก) ขุ 7. 77 ปิฏฺเฐปิ ปสฺสิตพฺพํ.} สมุคฺคหีตํ.}

\csromangathameasure{น ราคราคี น วิราครตฺโต, \\ ตสฺสีธ นตฺถี ปรมุคฺคหีตนฺติ. (8) สุทฺธฏฺฐกสุตฺตํ จตุตฺถํ.}
\csromangathagroup{\csromangathastanza{น ราคราคี น วิราครตฺโต, \\ ตสฺสีธ นตฺถี ปรมุคฺคหีตนฺติ. (8) สุทฺธฏฺฐกสุตฺตํ จตุตฺถํ.}}

\csromanheader{2}{5. ปรมฏฺฐกสุตฺต}
\proseitem{802}{ปรมนฺติ ทิฏฺฐีสุ ปริพฺพสาโน,}

\prose{ยทุตฺตริ กุรุเต ชนฺตุ โลเก.}

\prose{หีนาติ อญฺเญ ตโต สพฺพมาห,}

\csromanheader{2}{ตสฺมา วิวาทานิ อวีติวตฺโต. (1)}
\proseitem{803}{ยทตฺตนี ปสฺสติ อานิสํสํ,}

\prose{ทิฏฺเฐ สุเต สีลวเต\footnote{สีลพฺพเต (สฺยา)} มุเต วา.}

\prose{ตเทว โส ตตฺถ สมุคฺคหาย,}

\csromanheader{2}{นิหีนโต ปสฺสติ สพฺพมญฺญํ. (2)}
\gambhira{สุตฺตนิปาตปาฬิ}
\proseitem{804}{\csromanfolio{404}ตํ วาปิ คนฺถํ กุสลา วทนฺติ,}

\prose{ยํ นิสฺสิโต ปสฺสติ หีนมญฺญํ.}

\prose{ตสฺมา หิ ทิฏฺฐํ ว สุตํ มุตํ วา,}

\csromanheader{2}{สีลพฺพตํ ภิกฺขุ น นิสฺสเยยฺย. (3)}
\proseitem{805}{ทิฏฺฐิมฺปิ โลกสฺมิํ น กปฺปเยยฺย,}

\prose{ญาเณน วา สีลวเตน วาปิ.}

\prose{สโมติ อตฺตานมนูปเนยฺย,}

\csromanheader{2}{หีโน น มญฺเญถ วิเสสิ วาปิ. (4)}
\proseitem{806}{อตฺตํ ปหาย อนุปาทิยาโน,}

\prose{ญาเณปิ โส นิสฺสยํ โน กโรติ.}

\prose{ส เว วิยตฺเตสุ\footnote{วิยุตฺเตสุ (สี-ฏฺฐ), ทฺวิยตฺเตสุ (ก)} น วคฺคสารี,}

\csromanheader{2}{ทิฏฺฐิมฺปิ\footnote{ทิฏฺฐิมปิ(ก)} โส น ปจฺเจติ กิญฺจิ. (5)}
\proseitem{807}{ยสฺสูภยนฺเต ปณิธีธ นตฺถิ,}

\prose{ภวาภวาย อิธ วา หุรํ วา.}

\prose{นิเวสนา ตสฺส น สนฺติ เกจิ,}

\csromanheader{2}{ธมฺเมสุ นิจฺเฉยฺย สมุคฺคหีตํ. (6)}
\proseitem{808}{ตสฺสีธ ทิฏฺเฐ ว สุเต มุเต วา,}

\prose{ปกปฺปิตา นตฺถิ อณูปิ สญฺญา.}

\prose{ตํ พฺราหฺมณํ ทิฏฺฐิมนาทิยานํ,}

\csromanheader{2}{เกนีธ โลกสฺมิํ วิกปฺปเยยฺย. (7)}
\proseitem{809}{น กปฺปยนฺติ น ปุเรกฺขโรนฺติ,}

\csromangathameasure{ธมฺมาปิ เตสํ น ปฏิจฺฉิตาเส, \\ น พฺราหฺมโณ สีลวเตน เนยฺโย, \\ ปารงฺคโต น ปจฺเจติ ตาทีติ. (8) ปรมฏฺฐกสุตฺตํ ปญฺจมํ.}
\csromangathagroup{\csromangathastanza{ธมฺมาปิ เตสํ น ปฏิจฺฉิตาเส, \\ น พฺราหฺมโณ สีลวเตน เนยฺโย, \\ ปารงฺคโต น ปจฺเจติ ตาทีติ. (8) ปรมฏฺฐกสุตฺตํ ปญฺจมํ.}}

\csromanheader{2}{4. อฏฺฐกวคฺค \textbf{6. ชราสุตฺต}}
\proseitem{810}{\csromanfolio{405}\csromansymbolfootnote{*}{ขุ 7. 90 ปิฏฺเฐปิ.}อปฺปํ วต ชีวิตํ อิทํ,}

\prose{โอรํ วสฺสสตาปิ มิยฺยติ\footnote{มียติ (สี-ฏฺฐ)}.}

\csromangathameasure{โย เจปิ อติจฺจ ชีวติ,}
\csromangathagroup{\csromangathastanza{โย เจปิ อติจฺจ ชีวติ,}}

\csromanheader{2}{อถ โข โส ชรสาปิ มิยฺยติ. (1)}
\csromangathameasure{โสจนฺติ ชนา มมายิเต, \\ น หิ สนฺติ นิจฺจา ปริคฺคหา. \\ วินาภาวสนฺตเมวิทํ,}
\csromangathagroup{\csromangathastanza{โสจนฺติ ชนา มมายิเต, \\ น หิ สนฺติ\footnote{น หิ สนฺตา (สี), น หี สนฺติ (กตฺถจิ)} นิจฺจา ปริคฺคหา. \\ วินาภาวสนฺตเมวิทํ,}}

\csromanheader{2}{อิติ ทิสฺวา นาคาร’มาวเส. (2)}
\csromangathameasure{มรเณนปิ ตํ ปหียติ, \\ ยํ ปุริโส มมิทนฺติ มญฺญติ. \\ เอตมฺปิ วิทิตฺวา ปณฺฑิโต,}
\csromangathagroup{\csromangathastanza{มรเณนปิ ตํ ปหียติ\footnote{ปหิยฺยติ (สี, สฺยา, ก)}, \\ ยํ ปุริโส มมิทนฺติ\footnote{ปหิยฺยติ (สี, สฺยา, ก)} มญฺญติ. \\ เอตมฺปิ วิทิตฺวา\footnote{ปหิยฺยติ (สี, สฺยา, ก)} ปณฺฑิโต,}}

\csromanheader{2}{น มมตฺตาย นเมถ มามโก. (3)}
\proseitem{813}{\csromansymbolfootnote{+}{ขุ 10. 173 ปิฏฺเฐปิ.}สุปิเนน ยถาปิ สงฺคตํ,}

\prose{ปฏิพุทฺโธ ปุริโส น ปสฺสติ.}

\csromangathameasure{เอวมฺปิ ปิยายิตํ ชนํ,}
\csromangathagroup{\csromangathastanza{เอวมฺปิ ปิยายิตํ ชนํ,}}

\csromanheader{2}{เปตํ กาลงฺกตํ น ปสฺสติ. (4)}
\csromangathameasure{ทิฏฺฐาปิ สุตาปิ เต ชนา, \\ เยสํ นามมิทํ ปวุจฺจติ. \\ นามํเยวา’วสิสฺสติ,}
\csromangathagroup{\csromangathastanza{ทิฏฺฐาปิ สุตาปิ เต ชนา, \\ เยสํ นามมิทํ ปวุจฺจติ. \\ นามํเยวา’วสิสฺสติ\footnote{นามเมวา วสิสฺสติ (สี, สฺยา, อิ)},}}

\csromanheader{2}{อกฺเขยฺยํ เปตสฺส ชนฺตุโน. (5)}
\csromangathameasure{โสกปฺปริเทวมจฺฉรํ, \\ น ชหนฺติ คิทฺธา มมายิเต. \\ ตสฺมา มุนโย ปริคฺคหํ,}
\csromangathagroup{\csromangathastanza{โสกปฺปริเทวมจฺฉรํ\footnote{โสกปริเทวมจฺฉรํ (สี, สฺยา, อิ), โสกํ ปริเทวมจฺฉรํ (?)}, \\ น ชหนฺติ คิทฺธา มมายิเต. \\ ตสฺมา มุนโย ปริคฺคหํ,}}

\csromanheader{2}{หิตฺวา อจริํสุ เขมทสฺสิโน. (6)}
\gambhira{สุตฺตนิปาตปาฬิ}
\csromangathameasure{ปติลีนจรสฺส ภิกฺขุโน, \\ ภชมานสฺส วิวิตฺตมาสนํ. \\ สามคฺคิยมาหุ ตสฺส ตํ,}
\csromangathagroup{\csromangathastanza{\csromanfolio{406}ปติลีนจรสฺส ภิกฺขุโน, \\ ภชมานสฺส วิวิตฺตมาสนํ. \\ สามคฺคิยมาหุ ตสฺส ตํ,}}

\csromanheader{2}{โย อตฺตานํ ภวเน น ทสฺสเย. (7)}
\csromangathameasure{สพฺพตฺถ มุนี อนิสฺสิโต, \\ น ปิยํ กุพฺพติ โนปิ อปฺปิยํ. \\ ตสฺมิํ ปริเทวมจฺฉรํ,}
\csromangathagroup{\csromangathastanza{สพฺพตฺถ มุนี อนิสฺสิโต, \\ น ปิยํ กุพฺพติ โนปิ อปฺปิยํ. \\ ตสฺมิํ ปริเทวมจฺฉรํ,}}

\csromanheader{2}{ปณฺเณ วาริ ยถา น ลิมฺปติ\footnote{ลิปฺปติ (สี, อิ)}. (8)}
\csromangathameasure{อุทพินฺทุ ยถาปิ โปกฺขเร, \\ ปทุเม วาริ ยถา น ลิมฺปติ. \\ เอวํ มุนิ โนปลิปฺปติ,}
\csromangathagroup{\csromangathastanza{อุทพินฺทุ ยถาปิ โปกฺขเร, \\ ปทุเม วาริ ยถา น ลิมฺปติ. \\ เอวํ มุนิ โนปลิปฺปติ,}}

\csromanheader{2}{ยทิทํ ทิฏฺฐสุตํ มุเตสุ วา. (9)}
\csromangathameasure{โธโน น หิ เตน มญฺญติ, \\ ยทิทํ ทิฏฺฐสุตํ มุเตสุ วา. \\ นาญฺเญน วิสุทฺธิมิจฺฉติ, \\ น หิ โส รชฺชติ โน วิรชฺชตีติ. (10) ชราสุตฺตํ ฉฏฺฐํ.}
\csromangathagroup{\csromangathastanza{โธโน น หิ เตน มญฺญติ, \\ ยทิทํ ทิฏฺฐสุตํ มุเตสุ วา. \\ นาญฺเญน วิสุทฺธิมิจฺฉติ, \\ น หิ โส รชฺชติ โน วิรชฺชตีติ. (10) ชราสุตฺตํ ฉฏฺฐํ.}}

\csromanheader{2}{7. ติสฺสเมตฺเตยฺยสุตฺต}
\proseitem{820}{\csromansymbolfootnote{*}{ขุ 7. 107 ปิฏฺเฐปิ.}เมถุนมนุยุตฺตสฺส, (อิจฺจายสฺมา ติสฺโส เมตฺเตยฺโย,)}

\prose{วิฆาตํ พฺรูหิ มาริส.}

\csromangathameasure{สุตฺวาน ตว สาสนํ,}
\csromangathagroup{\csromangathastanza{สุตฺวาน ตว สาสนํ,}}

\csromanheader{2}{วิเวเก สิกฺขิสฺสามเส. (1)}
\proseitem{821}{เมถุนมนุยุตฺตสฺส, (เมตฺเตยฺยาติ ภควา,)}

\prose{มุสฺสเต’วาปิ สาสนํ.}

\csromangathameasure{มิจฺฉา จ ปฏิปชฺชติ,}
\csromangathagroup{\csromangathastanza{มิจฺฉา จ ปฏิปชฺชติ,}}

\csromanheader{2}{เอตํ ตสฺมิํ อนาริยํ. (2)}
\chapterheadpageread{4. อฏฺฐกวคฺค}
\csromangathameasure{เอโก ปุพฺเพ จริตฺวาน, เมถุนํ โย นิเสวติ.}
\csromangathagroup{\csromangathastanza{\csromanfolio{407}เอโก ปุพฺเพ จริตฺวาน, เมถุนํ โย นิเสวติ.}}

\vspace{\tipitakaparskip}
\csromancenter{ยานํ ภนฺตํว ตํ โลเก, หีนมาหุ ปุถุชฺชนํ. (3)}
\vspace{0.5\baselineskip}
\csromangathameasure{ยโส กิตฺติ จ ยา ปุพฺเพ, หายเต’วาปิ ตสฺส สา.}
\csromangathagroup{\csromangathastanza{ยโส กิตฺติ จ ยา ปุพฺเพ, หายเต’วาปิ ตสฺส สา.}}

\vspace{\tipitakaparskip}
\csromancenter{เอตมฺปิ ทิสฺวา สิกฺเขถ, เมถุนํ วิปฺปหาตเว. (4)}
\csromancenter{\csromansymbolfootnote{+}{ขุ 10. 173 ปิฏฺเฐปิ.}สงฺกปฺเปหิ ปเรโต โส, กปโณ วิย ฌายติ.}
\csromancenter{สุตฺวา ปเรสํ นิคฺโฆสํ, มงฺกุ โหติ ตถาวิโธ (5)}
\vspace{0.5\baselineskip}
\csromangathameasure{อถ สตฺถานิ กุรุเต, ปรวาเทหิ โจทิโต.}
\csromangathagroup{\csromangathastanza{อถ สตฺถานิ กุรุเต, ปรวาเทหิ โจทิโต.}}

\proseitem{825}{เอส ขฺวสฺส มหาเคโธ, โมสวชฺชํ ปคาหติ. (6)}

\csromangathameasure{ปณฺฑิโตติ สมญฺญาโต, เอกจริยํ อธิฏฺฐิโต.}
\csromangathagroup{\csromangathastanza{ปณฺฑิโตติ สมญฺญาโต, เอกจริยํ อธิฏฺฐิโต.}}

\proseitem{826}{อถาปิ\footnote{ส จาปิ (นิทฺเทเส 119 ปิฏฺเฐ)} เมถุเน ยุตฺโต, มนฺโทว ปริกิสฺสติ\footnote{ปริกิลิสฺสติ (สี)}. (7)}

\csromangathameasure{เอตมาทีนวํ ญตฺวา, มุนิ ปุพฺพาปเร อิธ.}
\csromangathagroup{\csromangathastanza{เอตมาทีนวํ ญตฺวา, มุนิ ปุพฺพาปเร อิธ.}}

\vspace{\tipitakaparskip}
\csromancenter{เอกจริยํ ทฬฺหํ กยิรา, น นิเสเวถ เมถุนํ. (8)}
\vspace{0.5\baselineskip}
\csromangathameasure{วิเวกญฺเญว สิกฺเขถ, เอตํ อริยานมุตฺตมํ.}
\csromangathagroup{\csromangathastanza{วิเวกญฺเญว สิกฺเขถ, เอตํ อริยานมุตฺตมํ.}}

\proseitem{828}{น เตน เสฏฺโฐ มญฺเญถ, ส เว นิพฺพานสนฺติเก. (9)}

\csromangathameasure{ริตฺตสฺส มุนิโน จรโต, กาเมสุ อนเปกฺขิโน. \\ โอฆติณฺณสฺส ปิหยนฺติ, กาเมสุ คธิตา ปชาติ. (10) ติสฺสเมตฺเตยฺยสุตฺตํ สตฺตมํ.}
\csromangathagroup{\csromangathastanza{ริตฺตสฺส มุนิโน จรโต, กาเมสุ อนเปกฺขิโน.}\csromangathastanza{โอฆติณฺณสฺส ปิหยนฺติ, กาเมสุ คธิตา\footnote{คถิตา (สี)} ปชาติ. (10) ติสฺสเมตฺเตยฺยสุตฺตํ สตฺตมํ.}}

\csromanheader{2}{8. ปสูรสุตฺต}
\proseitem{830}{\csromansymbolfootnote{*}{ขุ 7. 124 ปิฏฺเฐปิ.}อิเธว สุทฺธี อิติ วาทยนฺติ\footnote{วาทิยนฺติ (สี, อิ)},}

\prose{นาญฺเญสุ ธมฺเมสุ วิสุทฺธิมาหุ.}

\csromangathameasure{ยํ นิสฺสิตา ตตฺถ สุภํ วทานา,}
\csromangathagroup{\csromangathastanza{ยํ นิสฺสิตา ตตฺถ สุภํ วทานา,}}

\csromanheader{2}{ปจฺเจกสจฺเจสุ ปุถู นิวิฏฺฐา. (1)}
\gambhira{สุตฺตนิปาตปาฬิ}
\proseitem{831}{\csromanfolio{408}เต วาทกามา ปริสํ วิคยฺห,}

\prose{พาลํ ทหนฺตี มิถุ อญฺญมญฺญํ.}

\prose{วทนฺติ เต อญฺญสิตา กโถชฺชํ,}

\csromanheader{2}{ปสํสกามา กุสลา วทานา. (2)}
\proseitem{832}{ยุตฺโต กถายํ ปริสาย มชฺเฌ,}

\prose{ปสํสมิจฺฉํ วินิฆาติ โหติ.}

\prose{อปาหตสฺมิํ ปน มงฺกุ โหติ,}

\csromanheader{2}{นินฺทาย โส กุปฺปติ รนฺธเมสี. (3)}
\proseitem{833}{ยมสฺส วาทํ ปริหีนมาหุ,}

\prose{อปาหตํ ปญฺหวิมํสกาเส.}

\prose{ปริเทวติ โสจติ หีนวาโท,}

\prose{“อุปจฺจคา มนฺ”ติ อนุตฺถุนาติ. (4)}

\proseitem{834}{เอเต วิวาทา สมเณสุ ชาตา,}

\prose{เอเตสุ อุคฺฆาติ นิฆาติ โหติ.}

\prose{เอตมฺปิ ทิสฺวา วิรเม กโถชฺชํ,}

\csromanheader{2}{น ห’ญฺญทตฺถ’ตฺถิ ปสํสลาภา. (5)}
\proseitem{835}{ปสํสิโต วา ปน ตตฺถ โหติ,}

\csromangathameasure{อกฺขาย วาทํ ปริสาย มชฺเฌ, \\ โส หสฺสตี อุนฺนมตี จ เตน,}
\csromangathagroup{\csromangathastanza{อกฺขาย วาทํ ปริสาย มชฺเฌ, \\ โส หสฺสตี อุนฺนมตี\footnote{อุณฺณมตี (พหูสุ) ขุ 7. 131 ปิฏฺเฐปิ ปสฺสิตพฺพํ.} จ เตน,}}

\csromanheader{2}{ปปฺปุยฺย ตมตฺถํ ยถา มโน อหุ. (6)}
\proseitem{836}{ยา อุนฺนตี\footnote{อุนฺนตี (?)} สา’สฺส วิฆาตภูมิ,}

\prose{มานาติมานํ วทเต ปเนโส.}

\prose{เอตมฺปิ ทิสฺวา น วิวาทเยถ,}

\csromanheader{2}{น หิ เตน สุทฺธิํ กุสลา วทนฺติ. (7)}
\proseitem{837}{สูโร ยถา ราชขาทาย ปุฏฺโฐ,}

\prose{อภิคชฺช’เมติ ปฏิสูร’มิจฺฉํ.}

\prose{เยเนว โส, เตน ปเลหิ สูร,}

\csromanheader{2}{ปุพฺเพว นตฺถิ ยทิทํ ยุธาย. (8)}
\chapterheadpageread{4. อฏฺฐกวคฺค}
\csromangathameasure{เย ทิฏฺฐิมุคฺคยฺห วิวาทยนฺติ, \\ “อิทเมว สจฺจนฺ”ติ จ วาทยนฺติ. \\ เต ตฺวํ วทสฺสู น หิ เตธ อตฺถิ,}
\csromangathagroup{\csromangathastanza{\csromanfolio{409}เย ทิฏฺฐิมุคฺคยฺห วิวาทยนฺติ\footnote{วิวาทิยนฺติ (สี, อิ)}, \\ “อิทเมว สจฺจนฺ”ติ จ วาทยนฺติ. \\ เต ตฺวํ วทสฺสู น หิ เตธ อตฺถิ,}}

\csromanheader{2}{วาทมฺหิ ชาเต ปฏิเสนิกตฺตา. (9)}
\csromangathameasure{วิเสนิกตฺวา ปน เย จรนฺติ, \\ ทิฏฺฐีหิ ทิฏฺฐิํ อวิรุชฺฌมานา. \\ เตสุ ตฺวํ กิํ ลเภโถ ปสูร,}
\csromangathagroup{\csromangathastanza{วิเสนิกตฺวา ปน เย จรนฺติ, \\ ทิฏฺฐีหิ ทิฏฺฐิํ อวิรุชฺฌมานา. \\ เตสุ ตฺวํ กิํ ลเภโถ ปสูร,}}

\csromanheader{2}{เยสีธ นตฺถี ปรมุคฺคหีตํ. (10)}
\csromangathameasure{อถ ตฺวํ ปวิตกฺกมาคมา, \\ มนสา ทิฏฺฐิคตานิ จินฺตยนฺโต. \\ โธเนน ยุคํ สมาคมา, \\ น หิ ตฺวํ สกฺขสิ สมฺปยาตเวติ. (11) ปสูรสุตฺตํ อฏฺฐมํ.}
\csromangathagroup{\csromangathastanza{อถ ตฺวํ ปวิตกฺกมาคมา, \\ มนสา ทิฏฺฐิคตานิ จินฺตยนฺโต. \\ โธเนน ยุคํ สมาคมา, \\ น หิ ตฺวํ สกฺขสิ สมฺปยาตเวติ. (11) ปสูรสุตฺตํ อฏฺฐมํ.}}

\csromanheader{2}{9. มาคณฺฑิยสุตฺต}
\proseitem{841}{\csromansymbolfootnote{*}{ขุ 7. 140 ปิฏฺเฐปิ.}ทิสฺวาน ตณฺหํ อรติํ รคญฺจ\footnote{อรติญฺจ ราคํ (สฺยา, ก)},}

\prose{นาโหสิ ฉนฺโท อปิ เมถุนสฺมิํ.}

\csromangathameasure{กิเมวิทํ มุตฺตกรีสปุณฺณํ,}
\csromangathagroup{\csromangathastanza{กิเมวิทํ มุตฺตกรีสปุณฺณํ,}}

\csromanheader{2}{ปาทาปิ นํ สมฺผุสิตุํ น อิจฺเฉ. (1)}
\csromanmatikaanchor{mtk.275}
\csromanmatikaanchor{mtk.276}
\csromantocmarknopage{chapter}{4. อฏฺฐกวคฺค}{mtk.275}
\bookmark[dest={mtk.275},level=0]{4. อฏฺฐกวคฺค}
\csromantocmarkref{subsection}{1. กามสุตฺต}{mtk.276}
\bookmark[dest={mtk.276},level=2]{1. กามสุตฺต}
\csromangathameasure{เอตาทิสํ เจ รตนํ น อิจฺฉสิ, \\ นาริํ นรินฺเทหิ พหูหิ ปตฺถิตํ. \\ ทิฏฺฐิคตํ สีลวตํ นุ ชีวิตํ,}
\csromangathagroup{\csromangathastanza{เอตาทิสํ เจ รตนํ น อิจฺฉสิ, \\ นาริํ นรินฺเทหิ พหูหิ ปตฺถิตํ. \\ ทิฏฺฐิคตํ สีลวตํ นุ ชีวิตํ\footnote{สีลวตานุชีวิตํ (สี, อิ, ก)},}}

\csromanheader{2}{ภวูปปตฺติญฺจ วเทสิ กีทิสํ. (2)}
\vspace{\tipitakaparskip}
\csromancenter{อิทํ วทามีติ น ตสฺส โหติ, (มาคณฺฑิยาติ\footnote{มาคนฺทิยาติ (สี, สฺยา, อิ)} ภควา,)}
\vspace{0.5\baselineskip}
\csromangathameasure{ธมฺเมสุ นิจฺเฉยฺย สมุคฺคหีตํ, \\ ปสฺสญฺจ ทิฏฺฐีสุ อนุคฺคหาย,}
\csromangathagroup{\csromangathastanza{ธมฺเมสุ นิจฺเฉยฺย สมุคฺคหีตํ, \\ ปสฺสญฺจ ทิฏฺฐีสุ อนุคฺคหาย,}}

\csromanheader{2}{อชฺฌตฺตสนฺติํ ปจินํ อทสฺสํ. (3)}
\gambhira{สุตฺตนิปาตปาฬิ}
\proseitem{844}{\csromanfolio{410}วินิจฺฉยา ยานิ ปกปฺปิตานิ, (อิติ มาคณฺฑิโย\footnote{มาคนฺทิโย (สี, สฺยา, อิ)},)}

\prose{เต เว มุนี พฺรูสิ อนุคฺคหาย.}

\prose{อชฺฌตฺตสนฺตีติ ยเมตมตฺถํ,}

\csromanheader{2}{กถํ นุ ธีเรหิ ปเวทิตํ ตํ. (4)}
\proseitem{845}{น ทิฏฺฐิยา น สุติยา น ญาเณน, (มาคณฺฑิยาติ ภควา,)}

\csromangathameasure{สีลพฺพเตนาปิ น สุทฺธิมาห, \\ อทิฏฺฐิยา อสฺสุติยา อญาณา. \\ อสีลตา อพฺพตา โนปิ เตน, \\ เอเต จ นิสฺสชฺช อนุคฺคหาย.}
\csromangathagroup{\csromangathastanza{สีลพฺพเตนาปิ น สุทฺธิมาห, \\ อทิฏฺฐิยา อสฺสุติยา อญาณา. \\ อสีลตา อพฺพตา โนปิ เตน, \\ เอเต จ นิสฺสชฺช อนุคฺคหาย.}}

\csromanheader{2}{สนฺโต อนิสฺสาย ภวํ น ชปฺเป. (5)}
\csromanmatikaanchor{mtk.277}
\csromantocmarkref{subsection}{2. คุหฏฺฐกสุตฺต}{mtk.277}
\bookmark[dest={mtk.277},level=2]{2. คุหฏฺฐกสุตฺต}
\proseitem{846}{โน เจ กิร ทิฏฺฐิยา น สุติยา น ญาเณน, (อิติ มาคณฺฑิโย,)}

\prose{สีลพฺพเตนาปิ น สุทฺธิมาห.}

\csromangathameasure{อทิฏฺฐิยา อสฺสุติยา อญาณา, \\ อสีลตา อพฺพตา โนปิ เตน. \\ มญฺญามหํ โมมุหเมว ธมฺมํ,}
\csromangathagroup{\csromangathastanza{อทิฏฺฐิยา อสฺสุติยา อญาณา, \\ อสีลตา อพฺพตา โนปิ เตน. \\ มญฺญามหํ โมมุหเมว ธมฺมํ,}}

\csromanheader{2}{ทิฏฺฐิยา เอเก ปจฺเจนฺติ สุทฺธิํ. (6)}
\proseitem{847}{ทิฏฺฐญฺจ นิสฺสาย อนุปุจฺฉมาโน, (มาคณฺฑิยาติ ภควา,)}

\prose{สมุคฺคหีเตสุ ปโมหมาคา\footnote{สโมหมาคา (สฺยา, ก)}.}

\prose{อิโต จ นาทฺทกฺขิ อณุมฺปิ สญฺญํ,}

\csromanheader{2}{ตสฺมา ตุวํ โมมุหโต ทหาสิ. (7)}
\proseitem{848}{สโม วิเสสี อุท วา นิหีโน,}

\prose{โย มญฺญตี โส วิวเทถ เตน.}

\prose{ตีสุ วิธาสุ อวิกมฺปมาโน,}

\csromanheader{2}{สโม วิเสสีติ น ตสฺส โหติ. (8)}
\proseitem{849}{สจฺจนฺติ โส พฺราหฺมโณ กิํ วเทยฺย,}

\prose{มุสาติ วา โส วิวเทถ เกน.}

\prose{ยสฺมิํ สมํ วิสมํ วาปิ นตฺถิ,}

\csromanheader{2}{ส เกน วาทํ ปฏิสํยุเชยฺย. (9)}
\chapterheadpageread{4. อฏฺฐกวคฺค}
\proseitem{850}{\csromanfolio{411}\csromansymbolfootnote{*}{สํ 2. 8 ปิฏฺเฐปิ.}โอกํ ปหาย อนิเกตสารี,}

\prose{คาเม อกุพฺพํ มุนิ สนฺถวานิ\footnote{สนฺธวานิ (ก)}.}

\csromanmatikaanchor{mtk.278}
\csromantocmarkref{subsection}{3. ทุฏฺฐฏฺฐกสุตฺต}{mtk.278}
\bookmark[dest={mtk.278},level=2]{3. ทุฏฺฐฏฺฐกสุตฺต}
\csromangathameasure{กาเมหิ ริตฺโต อปุรกฺขราโน,}
\csromangathagroup{\csromangathastanza{กาเมหิ ริตฺโต อปุรกฺขราโน,}}

\csromanheader{2}{กถํ น วิคฺคยฺห ชเนน กยิรา. (10)}
\csromangathameasure{เยหิ วิวิตฺโต วิจเรยฺย โลเก, \\ น ตานิ อุคฺคยฺห วเทยฺย นาโค. \\ ชลมฺพุชํ กณฺฑกวาริชํ ยถา, \\ ชเลน ปงฺเกน จนูปลิตฺตํ. \\ เอวํ มุนี สนฺติวาโท อคิทฺโธ,}
\csromangathagroup{\csromangathastanza{เยหิ วิวิตฺโต วิจเรยฺย โลเก, \\ น ตานิ อุคฺคยฺห วเทยฺย นาโค. \\ ชลมฺพุชํ\footnote{เอลมฺพุชํ (สี, สฺยา)} กณฺฑกวาริชํ ยถา, \\ ชเลน ปงฺเกน จนูปลิตฺตํ.}\csromangathastanza{เอวํ มุนี สนฺติวาโท อคิทฺโธ,}}

\csromanheader{2}{กาเม จ โลเก จ อนูปลิตฺโต. (11)}
\csromangathameasure{น เวทคู ทิฏฺฐิยายโก น มุติยา, \\ ส มานเมติ น หิ ตมฺมโย โส. \\ น กมฺมุนา โนปิ สุเตน เนยฺโย,}
\csromangathagroup{\csromangathastanza{น เวทคู ทิฏฺฐิยายโก\footnote{น เวทคู ทิฏฺฐิยา (ก-สี, สฺยา, อิ) ขุ 7. 157 ปิฏฺเฐปิ.} น มุติยา, \\ ส มานเมติ น หิ ตมฺมโย โส. \\ น กมฺมุนา โนปิ สุเตน เนยฺโย,}}

\csromanheader{2}{อนูปนีโต ส นิเวสเนสุ. (12)}
\csromangathameasure{สญฺญาวิรตฺตสฺส น สนฺติ คนฺถา, \\ ปญฺญาวิมุตฺตสฺส น สนฺติ โมหา. \\ สญฺญญฺจ ทิฏฺฐิญฺจ เย อคฺคเหสุํ, \\ เต ฆฏฺฏยนฺตา วิจรนฺติ โลเกติ. (13) มาคณฺฑิยสุตฺตํ นวมํ.}
\csromangathagroup{\csromangathastanza{สญฺญาวิรตฺตสฺส น สนฺติ คนฺถา, \\ ปญฺญาวิมุตฺตสฺส น สนฺติ โมหา. \\ สญฺญญฺจ ทิฏฺฐิญฺจ เย อคฺคเหสุํ, \\ เต ฆฏฺฏยนฺตา\footnote{ฆฏฺฏมานา (สฺยา, ก)} วิจรนฺติ โลเกติ. (13) มาคณฺฑิยสุตฺตํ นวมํ.}}

\csromanheader{2}{10. ปุราเภทสุตฺต}
\vspace{\tipitakaparskip}
\csromancenter{\csromansymbolfootnote{+}{ขุ 7. 161 ปิฏฺเฐปิ.}กถํทสฺสี กถํสีโล, อุปสนฺโตติ วุจฺจติ.}
\csromancenter{ตํ เม โคตม ปพฺรูหิ, ปุจฺฉิโต อุตฺตมํ นรํ. (1)}
\proseitem{855}{วีตตโณฺห ปุรา เภทา, (อิติ ภควา,)}

\prose{ปุพฺพมนฺต’มนิสฺสิโต.}

\vspace{\tipitakaparskip}
\csromancenter{เวมชฺเฌ นุปสงฺเขยฺโย, ตสฺส นตฺถิ ปุรกฺขตํ. (2)}
\gambhira{สุตฺตนิปาตปาฬิ}
\proseitem{856}{\csromanfolio{412}อกฺโกธโน อสนฺตาสี, อวิกตฺถี อกุกฺกุโจ.}

\prose{มนฺตภาณี\footnote{มนฺตาภาณี (สฺยา, อิ)} อนุทฺธโต, ส เว วาจายโต มุนิ. (3)}

\proseitem{857}{นิราสตฺติ อนาคเต, อตีตํ นานุโสจติ.}

\prose{วิเวกทสฺสี ผสฺเสสุ, ทิฏฺฐีสุ จ น นียติ\footnote{นิยฺยติ (พหูสุ)}. (4)}

\proseitem{858}{ปติลีโน อกุหโก, อปิหาลุ อมจฺฉรี.}

\prose{อปฺปคพฺโภ อเชคุจฺโฉ, เปสุเณยฺเย จ โน ยุโต. (5)}

\proseitem{859}{สาติเยสุ อนสฺสาวี, อติมาเน จ โน ยุโต.}

\prose{สโณฺห จ ปฏิภานวา\footnote{ปฏิภาณวา (สฺยา, อิ) 4. นิสฺสยตา (สี, สฺยา, อิ)}, น สทฺโธ น วิรชฺชติ. (6)}

\proseitem{860}{ลาภกมฺยา น สิกฺขติ, อลาเภ จ น กุปฺปติ.}

\prose{อวิรุทฺโธ จ ตณฺหาย, รเสสุ นานุคิชฺฌติ. (7)}

\proseitem{861}{อุเปกฺขโก สทา สโต, น โลเก มญฺญเต สมํ.}

\prose{น วิเสสี น นีเจยฺโย, ตสฺส โน สนฺติ อุสฺสทา. (8)}

\proseitem{862}{ยสฺส นิสฺสยนา4 นตฺถิ, ญตฺวา ธมฺมํ อนิสฺสิโต.}

\prose{ภวาย วิภวาย วา, ตณฺหา ยสฺส น วิชฺชติ. (9)}

\csromanmatikaanchor{mtk.279}
\csromantocmarkref{subsection}{4. สุทฺธฏฺฐกสุตฺต}{mtk.279}
\bookmark[dest={mtk.279},level=2]{4. สุทฺธฏฺฐกสุตฺต}
\proseitem{863}{ตํ พฺรูมิ อุปสนฺโตติ, กาเมสุ อนเปกฺขินํ.}

\prose{คนฺถา ตสฺส น วิชฺชนฺติ, อตรี โส วิสตฺติกํ. (10)}

\proseitem{864}{น ตสฺส ปุตฺตา ปสโว, เขตฺตํ วตฺถุญฺจ วิชฺชติ.}

\prose{อตฺตา วาปิ นิรตฺตา วา\footnote{อตฺตํ วาปิ นิรตฺตํ วา (พหูสุ)}, น ตสฺมิํ อุปลพฺภติ. (11)}

\proseitem{865}{เยน นํ วชฺชุํ ปุถุชฺชนา, อโถ สมณพฺราหฺมณา.}

\prose{ตํ ตสฺส อปุรกฺขตํ, ตสฺมา วาเทสุ เน’ชติ. (12)}

\proseitem{866}{วีตเคโธ อมจฺฉรี, น อุสฺเสสุ วทเต มุนิ.}

\prose{น สเมสุ น โอเมสุ, กปฺปํ เน’ติ อกปฺปิโย. (13)}

\proseitem{867}{ยสฺส โลเก สกํ นตฺถิ, อสตา จ น โสจติ.}

\csromangathameasure{ธมฺเมสุ จ น คจฺฉติ, ส เว สนฺโตติ วุจฺจตีติ. (14) ปุราเภทสุตฺตํ ทสมํ.}
\csromangathagroup{\csromangathastanza{ธมฺเมสุ จ น คจฺฉติ, ส เว สนฺโตติ วุจฺจตีติ. (14) ปุราเภทสุตฺตํ ทสมํ.}}

\csromanheader{2}{4. อฏฺฐกวคฺค \textbf{11. กลหวิวาทสุตฺต}}
\proseitem{868}{\csromanfolio{413}\csromansymbolfootnote{*}{ขุ 7. 196 ปิฏฺเฐปิ.}กุโตปหูตา กลหา วิวาทา,}

\prose{ปริเทวโสกา สหมจฺฉรา จ.}

\csromangathameasure{มานาติมานา สหเปสุณา จ,}
\csromangathagroup{\csromangathastanza{มานาติมานา สหเปสุณา จ,}}

\csromanheader{2}{กุโตปหูตา เต ตทิงฺฆ พฺรูหิ. (1)}
\csromangathameasure{ปิยปฺปหูตา กลหา วิวาทา, \\ ปริเทวโสกา สหมจฺฉรา จ. \\ มานาติมานา สหเปสุณา จ, \\ มจฺเฉรยุตฺตา กลหา วิวาทา.}
\csromangathagroup{\csromangathastanza{ปิยปฺปหูตา กลหา วิวาทา, \\ ปริเทวโสกา สหมจฺฉรา จ. \\ มานาติมานา สหเปสุณา จ, \\ มจฺเฉรยุตฺตา กลหา วิวาทา.}}

\csromanheader{2}{วิวาทชาเตสุ จ เปสุณานิ. (2)}
\csromangathameasure{ปิยา สุ โลกสฺมิํ กุโตนิทานา, \\ เย จาปิ โลภา วิจรนฺติ โลเก. \\ อาสา จ นิฏฺฐา จ กุโตนิทานา,}
\csromangathagroup{\csromangathastanza{ปิยา สุ\footnote{ปิยานุ (สฺยา), ปิยสฺสุ (ก)} โลกสฺมิํ กุโตนิทานา, \\ เย จาปิ\footnote{ปิยานุ (สฺยา), ปิยสฺสุ (ก)} โลภา วิจรนฺติ โลเก. \\ อาสา จ นิฏฺฐา จ กุโตนิทานา,}}

\csromanheader{2}{เย สมฺปรายาย นรสฺส โหนฺติ. (3)}
\csromangathameasure{ฉนฺทานิทานานิ ปิยานิ โลเก, \\ เย จาปิ โลภา วิจรนฺติ โลเก. \\ อาสา จ นิฏฺฐา จ อิโตนิทานา,}
\csromangathagroup{\csromangathastanza{ฉนฺทานิทานานิ ปิยานิ โลเก, \\ เย จาปิ โลภา วิจรนฺติ โลเก. \\ อาสา จ นิฏฺฐา จ อิโตนิทานา,}}

\csromanheader{2}{เย สมฺปรายาย นรสฺส โหนฺติ. (4)}
\csromangathameasure{ฉนฺโท นุ โลกสฺมิํ กุโตนิทาโน, \\ วินิจฺฉยา จาปิ กุโตปหูตา. \\ โกโธ โมสวชฺชญฺจ กถํกถา จ,}
\csromangathagroup{\csromangathastanza{ฉนฺโท นุ โลกสฺมิํ กุโตนิทาโน, \\ วินิจฺฉยา จาปิ\footnote{วาปิ (สี, สฺยา, อิ)} กุโตปหูตา. \\ โกโธ โมสวชฺชญฺจ กถํกถา จ,}}

\csromanheader{2}{เย วาปิ ธมฺมา สมเณน วุตฺตา. (5)}
\csromangathameasure{สาตํ อสาตนฺติ ยมาหุ โลเก, \\ ตมูปนิสฺสาย ปโหติ ฉนฺโท. \\ รูเปสุ ทิสฺวา วิภวํ ภวญฺจ,}
\csromangathagroup{\csromangathastanza{สาตํ อสาตนฺติ ยมาหุ โลเก, \\ ตมูปนิสฺสาย ปโหติ ฉนฺโท. \\ รูเปสุ ทิสฺวา วิภวํ ภวญฺจ,}}

\csromanheader{2}{วินิจฺฉยํ กุพฺพติ\footnote{กุรุเต (พหูสุ)} ชนฺตุ โลเก. (6)}
\gambhira{สุตฺตนิปาตปาฬิ}
\proseitem{874}{\csromanfolio{414}โกโธ โมสวชฺชญฺจ กถํกถา จ,}

\prose{เอเตปิ ธมฺมา ทฺวยเมว สนฺเต.}

\csromanmatikaanchor{mtk.280}
\csromantocmarkref{subsection}{5. ปรมฏฺฐกสุตฺต}{mtk.280}
\bookmark[dest={mtk.280},level=2]{5. ปรมฏฺฐกสุตฺต}
\prose{กถํกถี ญาณปถาย สิกฺเข,}

\csromanheader{2}{ญตฺวา ปวุตฺตา สมเณน ธมฺมา. (7)}
\proseitem{875}{สาตํ อสาตญฺจ กุโตนิทานา,}

\prose{กิสฺมิํ อสนฺเต น ภวนฺติ เหเต.}

\prose{วิภวํ ภวญฺจาปิ ยเมตมตฺถํ,}

\csromanheader{2}{เอตํ เม ปพฺรูหิ ยโตนิทานํ. (8)}
\proseitem{876}{ผสฺสนิทานํ สาตํ อสาตํ,}

\prose{ผสฺเส อสนฺเต น ภวนฺติ เหเต.}

\prose{วิภวํ ภวญฺจาปิ ยเมตมตฺถํ,}

\csromanheader{2}{เอตํ เต ปพฺรูมิ อิโตนิทานํ. (9)}
\proseitem{877}{ผสฺโส นุ โลกสฺมิ กุโตนิทาโน.}

\prose{ปริคฺคหา จาปิ กุโตปหูตา.}

\prose{กิสฺมิํ อสนฺเต น มมตฺตมตฺถิ,}

\csromanheader{2}{กิสฺมิํ วิภูเต น ผุสนฺติ ผสฺสา. (10)}
\proseitem{878}{นามญฺจ รูปญฺจ ปฏิจฺจ ผสฺโส,}

\prose{อิจฺฉานิทานานิ ปริคฺคหานิ.}

\prose{อิจฺฉาย’สนฺตฺยา น มมตฺตมตฺถิ,}

\csromanheader{2}{รูเป วิภูเต น ผุสนฺติ ผสฺสา. (11)}
\proseitem{879}{กถํ สเมตสฺส วิโภติ รูปํ,}

\prose{สุขํ ทุขญฺจาปิ\footnote{ทุขํ วาปิ (สี, สฺยา)} กถํ วิโภติ.}

\prose{เอตํ เม ปพฺรูหิ ยถา วิโภติ,}

\csromanheader{2}{ตํ ชานิยามาติ\footnote{ชานิสฺสามาติ (สี, ก)} เม มโน อหุ. (12)}
\proseitem{880}{น สญฺญสญฺญี น วิสญฺญสญฺญี,}

\prose{โนปิ อสญฺญี น วิภูตสญฺญี.}

\prose{เอวํ สเมตสฺส วิโภติ รูปํ,}

\csromanheader{2}{สญฺญานิทานา หิ ปปญฺจสงฺขา. (13)}
\chapterheadpageread{4. อฏฺฐกวคฺค}
\csromangathameasure{ยํ ตํ อปุจฺฉิมฺห อกิตฺตยี โน, \\ อญฺญํ ตํ ปุจฺฉาม ตทิงฺฆ พฺรูหิ. \\ เอตฺตาวต’คฺคํ นุ วทนฺติ เหเก, \\ ยกฺขสฺส สุทฺธิํ อิธ ปณฺฑิตาเส.}
\csromangathagroup{\csromangathastanza{\csromanfolio{415}ยํ ตํ อปุจฺฉิมฺห อกิตฺตยี โน, \\ อญฺญํ ตํ ปุจฺฉาม ตทิงฺฆ พฺรูหิ. \\ เอตฺตาวต’คฺคํ นุ\footnote{โน (สี, สฺยา)} วทนฺติ เหเก, \\ ยกฺขสฺส สุทฺธิํ อิธ ปณฺฑิตาเส.}}

\csromanheader{2}{อุทาหุ อญฺญมฺปิ วทนฺติ เอตฺโต. (14)}
\csromangathameasure{เอตฺตาวต’คฺคมฺปิ วทนฺติ เหเก, \\ ยกฺขสฺส สุทฺธิํ อิธ ปณฺฑิตาเส. \\ เตสํ ปเนเก สมยํ วทนฺติ.}
\csromangathagroup{\csromangathastanza{เอตฺตาวต’คฺคมฺปิ วทนฺติ เหเก, \\ ยกฺขสฺส สุทฺธิํ อิธ ปณฺฑิตาเส. \\ เตสํ ปเนเก สมยํ วทนฺติ.}}

\csromanheader{2}{อนุปาทิเสเส กุสลา วทานา. (15)}
\csromangathameasure{เอเต จ ญตฺวา อุปนิสฺสิตาติ, \\ ญตฺวา มุนี นิสฺสเย โส วิมํสี. \\ ญตฺวา วิมุตฺโต น วิวาท’เมติ, \\ ภวาภวาย น สเมติ ธีโรติ. (16) กลหวิวาทสุตฺตํ เอกาทสมํ.}
\csromangathagroup{\csromangathastanza{เอเต จ ญตฺวา อุปนิสฺสิตาติ, \\ ญตฺวา มุนี นิสฺสเย โส วิมํสี. \\ ญตฺวา วิมุตฺโต น วิวาท’เมติ, \\ ภวาภวาย น สเมติ ธีโรติ. (16) กลหวิวาทสุตฺตํ เอกาทสมํ.}}

\csromanheader{2}{12. จูฬพฺยูหสุตฺต\footnote{จูฬวิยูหสุตฺตํ (สี, สฺยา)}}
\proseitem{884}{\csromansymbolfootnote{*}{ขุ 7. 220 ปิฏฺเฐปิ.}สกํสกํทิฏฺฐิปริพฺพสานา,}

\csromanmatikaanchor{mtk.281}
\csromantocmarkref{subsection}{6. ชราสุตฺต}{mtk.281}
\bookmark[dest={mtk.281},level=2]{6. ชราสุตฺต}
\prose{วิคฺคยฺห นานา กุสลา วทนฺติ.}

\csromangathameasure{โย เอวํ ชานาติ ส เวทิ ธมฺมํ,}
\csromangathagroup{\csromangathastanza{โย เอวํ ชานาติ ส เวทิ ธมฺมํ,}}

\csromanheader{2}{อิทํ ปฏิกฺโกส’มเกวลี โส. (1)}
\csromangathameasure{เอวมฺปิ วิคฺคยฺห วิวาทยนฺติ, \\ พาโล ปโร อกฺกุสโลติ จาหุ. \\ สจฺโจ นุ วาโท กตโม อิเมสํ,}
\csromangathagroup{\csromangathastanza{เอวมฺปิ วิคฺคยฺห วิวาทยนฺติ, \\ พาโล ปโร อกฺกุสโลติ\footnote{อกุสโลติ (สี, สฺยา, อิ)} จาหุ. \\ สจฺโจ นุ วาโท กตโม อิเมสํ,}}

\proseitem{885}{สพฺเพว หีเม กุสลา วทานา, (\textbf{2})}

\gambhira{สุตฺตนิปาตปาฬิ}
\proseitem{886}{\csromanfolio{416}ปรสฺส เจ ธมฺมมนานุชานํ,}

\prose{พาโล’มโก\footnote{พาโล มโค (สี, สฺยา, ก)} โหติ นิหีนปญฺโญ.}

\prose{สพฺเพว พาลา สุนิหีนปญฺญา,}

\csromanheader{2}{สพฺเพวิเม ทิฏฺฐิปริพฺพสานา. (3)}
\proseitem{887}{สนฺทิฏฺฐิยา เจว น วีวทาตา.}

\prose{สํสุทฺธปญฺญา กุสลา มุตีมา.}

\prose{น เตสํ โกจิ ปริหีนปญฺโญ\footnote{โกจิปิ นิหีนปญฺโญ (สี, สฺยา, ก)},}

\csromanheader{2}{ทิฏฺฐี หิ เตสมฺปิ ตถา สมตฺตา. (4)}
\proseitem{888}{น วาหเมตํ ตถิยนฺติ\footnote{ตถิวนฺติ (สฺยา, ก)} พฺรูมิ,}

\prose{ยมาหุ พาลา มิถุ อญฺญมญฺญํ.}

\prose{สกํ สกํ ทิฏฺฐิมกํสุ สจฺจํ,}

\csromanheader{2}{ตสฺมา หิ พาโลติ ปรํ ทหนฺติ. (5)}
\proseitem{889}{ยมาหุ สจฺจํ ตถิยนฺติ เอเก,}

\prose{ตมาหุ อญฺเญ\footnote{อญฺเญปิ (สฺยา), อญฺเญ จ (?)} ตุจฺฉํ มุสาติ.}

\prose{เอวมฺปิ วิคยฺห วิวาทยนฺติ,}

\csromanheader{2}{กสฺมา น เอกํ สมณา วทนฺติ. (6)}
\proseitem{890}{เอกํ หิ สจฺจํ น ทุตียมตฺถิ,}

\prose{ยสฺมิํ ปชา โน วิวเท ปชานํ.}

\prose{นานา เต\footnote{นานาโต (ก)} สจฺจานิ สยํ ถุนนฺติ,}

\csromanmatikaanchor{mtk.282}
\csromantocmarkref{subsection}{7. ติสฺสเมตฺเตยฺยสุตฺต}{mtk.282}
\bookmark[dest={mtk.282},level=2]{7. ติสฺสเมตฺเตยฺยสุตฺต}
\csromanheader{2}{ตสฺมา น เอกํ สมณา วทนฺติ. (7)}
\proseitem{891}{กสฺมา นุ สจฺจานิ วทนฺติ นานา,}

\prose{ปวาทิยาเส กุสลา วทานา.}

\prose{สจฺจานิ สุตานิ พหูนิ นานา,}

\csromanheader{2}{อุทาหุ เต ตกฺกมนุสฺสรนฺติ. (8)}
\proseitem{892}{น เหว สจฺจานิ พหูนิ นานา,}

\prose{อญฺญตฺร สญฺญาย นิจฺจานิ โลเก.}

\prose{ตกฺกญฺจ ทิฏฺฐีสุ ปกปฺปยิตฺวา,}

\csromanheader{2}{สจฺจํ มุสาติ ทฺวยธมฺมมาหุ. (9)}
\chapterheadpageread{4. อฏฺฐกวคฺค}
\proseitem{893}{\csromanfolio{417}ทิฏฺเฐ สุเต สีลวเต มุเต วา,}

\prose{เอเต จ นิสฺสาย วิมานทสฺสี.}

\prose{วินิจฺฉเย ฐตฺวา ปหสฺสมาโน,}

\csromanheader{2}{พาโล ปโร อกฺกุสโลติ จาห. (10)}
\proseitem{894}{เยเนว พาโลติ ปรํ ทหาติ,}

\prose{เตนา’ตุมานํ กุสโลติ จาห.}

\prose{สย’มตฺตนา โส กุสลาวทาโน,}

\csromanheader{2}{อญฺญํ วิมาเนติ ตเทว ปาว. (11)}
\proseitem{895}{อติสารทิฏฺฐิยาว โส สมตฺโต,}

\prose{มาเนน มตฺโต ปริปุณฺณมานี.}

\prose{สยเมว สามํ มนสาภิสิตฺโต,}

\csromanheader{2}{ทิฏฺฐี หิ สา ตสฺส ตถา สมตฺตา. (12)}
\proseitem{896}{ปรสฺส เจ หิ วจสา นิหีโน,}

\prose{ตุโม สหา โหติ นิหีนปญฺโญ.}

\prose{อถ เจ สยํ เวทคู โหติ ธีโร,}

\csromanheader{2}{น โกจิ พาโล สมเณสุ อตฺถิ. (13)}
\csromanmatikaanchor{mtk.283}
\csromantocmarkref{subsection}{8. ปสูรสุตฺต}{mtk.283}
\bookmark[dest={mtk.283},level=2]{8. ปสูรสุตฺต}
\proseitem{897}{อญฺญํ อิโต ยา’ภิวทนฺติ ธมฺมํ,}

\prose{อปรทฺธา สุทฺธิ’มเกวลี เต\footnote{สุทฺธิมเกวลีโน (สี)}.}

\prose{เอวมฺปิ ติตฺถฺยา ปุถุโส วทนฺติ,}

\csromanheader{2}{สนฺทิฏฺฐิราเคน หิ เต’ภิรตฺตา\footnote{ตฺยาภิรตฺตา (สฺยา, ก)}. (14)}
\proseitem{898}{อิเธว สุทฺธิํ อิติ วาทยนฺติ,}

\prose{นาญฺเญสุ ธมฺเมสุ วิสุทฺธิมาหุ.}

\prose{เอวมฺปิ ติตฺถฺยา ปุถุโส นิวิฏฺฐา,}

\csromanheader{2}{สกายเน ตตฺถ ทฬฺหํ วทานา. (15)}
\proseitem{899}{สกายเน วาปิ ทฬฺหํ วทาโน,}

\prose{ก’เมตฺถ พาโลติ ปรํ ทเหยฺย.}

\prose{สยํว โส เมธคมาวเหยฺย\footnote{เมธกํ อาวเหยฺย (สี, อิ)},}

\csromanheader{2}{ปรํ วทํ พาลมสุทฺธิธมฺมํ. (16)}
\gambhira{สุตฺตนิปาตปาฬิ}
\csromangathameasure{วินิจฺฉเย ฐตฺวา สยํ ปมาย, \\ อุทฺธํส โลกสฺมิํ วิวาทเมติ. \\ หิตฺวาน สพฺพานิ วินิจฺฉยานิ, \\ น เมธคํ กุพฺพติ ชนฺตุ โลเกติ. (17) จูฬพฺยูหสุตฺตํ ทฺวาทสมํ.}
\csromangathagroup{\csromangathastanza{\csromanfolio{418}วินิจฺฉเย ฐตฺวา สยํ ปมาย, \\ อุทฺธํส\footnote{อุทฺธํ โส (สี, สฺยา, อิ)} โลกสฺมิํ วิวาทเมติ. \\ หิตฺวาน สพฺพานิ วินิจฺฉยานิ, \\ น เมธคํ กุพฺพติ ชนฺตุ โลเกติ. (17) จูฬพฺยูหสุตฺตํ ทฺวาทสมํ.}}

\csromanheader{2}{13. มหาพฺยูหสุตฺต}
\proseitem{901}{\csromansymbolfootnote{*}{ขุ 7. 236 ปิฏฺเฐปิ.}เย เกจิเม ทิฏฺฐิปริพฺพสานา,}

\prose{อิทเมว สจฺจนฺติ วิวาทยนฺติ\footnote{วิวาทิยนฺติ (สี)}.}

\csromangathameasure{สพฺเพว เต นินฺทมนฺวานยนฺติ,}
\csromangathagroup{\csromangathastanza{สพฺเพว เต นินฺทมนฺวานยนฺติ,}}

\csromanheader{2}{อโถ ปสํสมฺปิ ลภนฺติ ตตฺถ. (1)}
\proseitem{902}{อปฺปํ หิ เอตํ น อลํ สมาย,}

\prose{ทุเว วิวาทสฺส ผลานิ พฺรูมิ.}

\csromangathameasure{เอตมฺปิ ทิสฺวา น วิวาทเยถ,}
\csromangathagroup{\csromangathastanza{เอตมฺปิ ทิสฺวา น วิวาทเยถ,}}

\csromanheader{2}{เขมา’ภิปสฺสํ อวิวาทภูมิํ. (2)}
\csromangathameasure{ยา กาจิมา สมฺมุติโย ปุถุชฺชา, \\ สพฺพาว เอตา น อุเปติ วิทฺวา. \\ อนูปโย โส อุปยํ กิเมยฺย,}
\csromangathagroup{\csromangathastanza{ยา กาจิมา สมฺมุติโย ปุถุชฺชา, \\ สพฺพาว เอตา น อุเปติ วิทฺวา. \\ อนูปโย โส อุปยํ กิเมยฺย,}}

\csromanheader{2}{ทิฏฺเฐ สุเต ขนฺติมกุพฺพมาโน. (3)}
\csromangathameasure{สีลุตฺตมา สญฺญเมนาหุ สุทฺธิํ, \\ วตํ สมาทาย อุปฏฺฐิตาเส. \\ อิเธว สิกฺเขม อถสฺส สุทฺธิํ,}
\csromangathagroup{\csromangathastanza{สีลุตฺตมา สญฺญเมนาหุ สุทฺธิํ, \\ วตํ สมาทาย อุปฏฺฐิตาเส. \\ อิเธว สิกฺเขม อถสฺส สุทฺธิํ,}}

\csromanheader{2}{ภวูปนีตา กุสลาวทานา. (4)}
\csromangathameasure{สเจ จุโต สีลวตโต โหติ, \\ ปเวธตี กมฺม วิราธยิตฺวา. \\ ปชปฺปตี ปตฺถยตี จ สุทฺธิํ,}
\csromangathagroup{\csromangathastanza{สเจ จุโต สีลวตโต โหติ, \\ ปเวธตี\footnote{สเวธติ (สี, อิ)} กมฺม วิราธยิตฺวา. \\ ปชปฺปตี ปตฺถยตี จ สุทฺธิํ,}}

\csromanheader{2}{สตฺถาว หีโน ปวสํ ฆรมฺหา. (5)}
\chapterheadpageread{4. อฏฺฐกวคฺค}
\proseitem{906}{\csromanfolio{419}สีลพฺพตํ วาปิ ปหาย สพฺพํ,}

\prose{กมฺมญฺจ สาวชฺชนวชฺชเมตํ.}

\prose{สุทฺธิํ อสุทฺธินฺติ อปตฺถยาโน,}

\csromanheader{2}{วิรโต จเร สนฺติมนุคฺคหาย. (6)}
\proseitem{907}{ตมูปนิสฺสาย ชิคุจฺฉิตํ วา,}

\prose{อถวาปิ ทิฏฺฐํ ว สุตํ มุตํ วา.}

\prose{อุทฺธํสรา สุทฺธิ’มนุตฺถุนนฺติ,}

\csromanheader{2}{อวีตตณฺหาเส ภวาภเวสุ. (7)}
\proseitem{908}{ปตฺถยมานสฺส หิ ชปฺปิตานิ,}

\prose{ปเวธิตํ วาปิ ปกปฺปิเตสุ.}

\csromanmatikaanchor{mtk.284}
\csromantocmarkref{subsection}{9. มาคณฺฑิยสุตฺต}{mtk.284}
\bookmark[dest={mtk.284},level=2]{9. มาคณฺฑิยสุตฺต}
\prose{จุตูปปาโต อิธ ยสฺส นตฺถิ,}

\csromanheader{2}{ส เกน เวเธยฺย กุหิํ ว ชปฺเป\footnote{กุหิญฺจิ ชปฺเป (สี, สฺยา, ก), กุหิํ ปชปฺเป (อิ) นิทฺเทโส ปสฺสิตพฺโพ.}. (8)}
\proseitem{909}{ยมาหุ ธมฺมํ ปรมนฺติ เอเก,}

\prose{ตเมว หีนนฺติ ปนาหุ อญฺเญ.}

\prose{สจฺโจ นุ วาโท กตโม อิเมสํ,}

\csromanheader{2}{สพฺเพว หีเม กุสลาวทานา. (9)}
\proseitem{910}{สกญฺหิ ธมฺมํ ปริปุณฺณมาหุ,}

\prose{อญฺญสฺส ธมฺมํ ปน หีนมาหุ.}

\prose{เอวมฺปิ วิคฺคยฺห วิวาทยนฺติ,}

\csromanheader{2}{สกํ สกํ สมฺมุติมาหุ สจฺจํ. (10)}
\proseitem{911}{ปรสฺส เจ วมฺภยิเตน หีโน,}

\prose{น โกจิ ธมฺเมสุ วิเสสิ อสฺส.}

\prose{ปุถู หิ อญฺญสฺส วทนฺติ ธมฺมํ,}

\csromanheader{2}{นิหีนโต สมฺหิ ทฬฺหํ วทานา. (11)}
\proseitem{912}{สทฺธมฺมปูชาปิ เนสํ ตเถว,}

\prose{ยถา ปสํสนฺติ สกายนานิ.}

\prose{สพฺเพว วาทา\footnote{สพฺเพ ปวาทา (สฺยา)} ตถิยา\footnote{ตถิวา (สพฺพตฺถ)} ภเวยฺยุํ,}

\csromanheader{2}{สุทฺธี หิ เนสํ ปจฺจตฺตเมว. (12)}
\gambhira{สุตฺตนิปาตปาฬิ}
\proseitem{913}{\csromanfolio{420}น พฺราหฺมณสฺส ปรเนยฺยมตฺถิ,}

\prose{ธมฺเมสุ นิจฺเฉยฺย สมุคฺคหีตํ.}

\prose{ตสฺมา วิวาทานิ อุปาติวตฺโต,}

\csromanheader{2}{น หิ เสฏฺฐโต ปสฺสติ ธมฺมมญฺญํ. (13)}
\proseitem{914}{ชานามิ ปสฺสามิ ตเถว เอตํ,}

\prose{ทิฏฺฐิยา เอเก ปจฺเจนฺติ สุทฺธิํ.}

\prose{อทฺทกฺขิ เจ กิญฺหิ ตุมสฺส เตน,}

\csromanheader{2}{อติสิตฺวา อญฺเญน วทนฺติ สุทฺธิํ. (14)}
\proseitem{915}{ปสฺสํ นโร ทกฺขติ\footnote{ทกฺขิติ (สี)} นามรูปํ,}

\prose{ทิสฺวาน วา ญสฺสติ ตานิเมว.}

\prose{กามํ พหุํ ปสฺสตุ อปฺปกํ วา,}

\csromanheader{2}{น หิ เตน สุทฺธิํ กุสลา วทนฺติ. (15)}
\proseitem{916}{นิวิสฺสวาที น หิ สุพฺพินาโย,}

\prose{ปกปฺปิตํ ทิฏฺฐิ ปุรกฺขราโน.}

\prose{ยํ นิสฺสิโต ตตฺถ สุภํ วทาโน,}

\csromanheader{2}{สุทฺธิํวโท ตตฺถ ตถทฺทสา โส. (16)}
\proseitem{917}{น พฺราหฺมโณ กปฺปมุเปติ สงฺขา\footnote{สงฺขํ (สี, สฺยา, อิ)},}

\prose{น ทิฏฺฐิสารี นปิ ญาณพนฺธุ.}

\prose{ญตฺวา จ โส สมฺมุติโย\footnote{สมฺมติโย (สฺยา)} ปุถุชฺชา,}

\csromanheader{2}{อุเปกฺขตี อุคฺคหณนฺติ มญฺเญ. (17)}
\proseitem{918}{วิสฺสชฺช คนฺถานิ มุนีธ โลเก,}

\prose{วิวาทชาเตสุ น วคฺคสารี.}

\prose{สนฺโต อสนฺเตสุ อุเปกฺขโก โส,}

\csromanheader{2}{อนุคฺคโห อุคฺคหณนฺติ มญฺเญ. (18)}
\proseitem{919}{ปุพฺพาสเว หิตฺวา นเว อกุพฺพํ,}

\prose{น ฉนฺทคู โนปิ นิวิสฺสวาที.}

\prose{ส วิปฺปมุตฺโต ทิฏฺฐิคเตหิ ธีโร,}

\csromanheader{2}{น ลิปฺปติ\footnote{น ลิมฺปติ (ก)} โลเก อนตฺตครหี. (19)}
\chapterheadpageread{4. อฏฺฐกวคฺค}
\csromangathameasure{ส สพฺพธมฺเมสุ วิเสนิภูโต, \\ ยํ กิญฺจิ ทิฏฺฐํ ว สุตํ มุตํ วา. \\ ส ปนฺนภาโร มุนิ วิปฺปมุตฺโต, \\ น กปฺปิโย นูปรโต น ปตฺถิโยติ. (20) มหาพฺยูหสุตฺตํ เตรสมํ.}
\csromangathagroup{\csromangathastanza{\csromanfolio{421}ส สพฺพธมฺเมสุ วิเสนิภูโต, \\ ยํ กิญฺจิ ทิฏฺฐํ ว สุตํ มุตํ วา. \\ ส ปนฺนภาโร มุนิ วิปฺปมุตฺโต, \\ น กปฺปิโย นูปรโต น ปตฺถิโยติ. (20) มหาพฺยูหสุตฺตํ เตรสมํ.}}

\csromanheader{2}{14. ตุวฏกสุตฺต}
\csromanmatikaanchor{mtk.285}
\csromantocmarkref{subsection}{10. ปุราเภทสุตฺต}{mtk.285}
\bookmark[dest={mtk.285},level=2]{10. ปุราเภทสุตฺต}
\proseitem{921}{\csromansymbolfootnote{*}{ขุ 7. 263 ปิฏฺเฐปิ.}ปุจฺฉามิ ตํ อาทิจฺจพนฺธุ\footnote{อาทิจฺจพนฺธุํ (สี, สฺยา)},}

\prose{วิเวกํ สนฺติปทญฺจ มเหสิ.}

\csromangathameasure{กถํ ทิสฺวา นิพฺพาติ ภิกฺขุ,}
\csromangathagroup{\csromangathastanza{กถํ ทิสฺวา นิพฺพาติ ภิกฺขุ,}}

\csromanheader{2}{อนุปาทิยาโน โลกสฺมิํ กิญฺจิ. (1)}
\proseitem{922}{มูลํ ปปญฺจสงฺขาย, (อิติ ภควา,)}

\prose{มนฺตา อสฺมีติ สพฺพมุปรุนฺเธ\footnote{สพฺพมุปรุทฺเธ (สฺยา, อิ, ก)}.}

\csromangathameasure{ยา กาจิ ตณฺหา อชฺฌตฺตํ,}
\csromangathagroup{\csromangathastanza{ยา กาจิ ตณฺหา อชฺฌตฺตํ,}}

\csromanheader{2}{ตาสํ วินยา\footnote{วินยาย (?)} สทา สโต สิกฺเข. (2)}
\csromangathameasure{ยํ กิญฺจิ ธมฺมมภิชญฺญา, \\ อชฺฌตฺตํ อถ วาปิ พหิทฺธา. \\ น เตน ถามํ กุพฺเพถ,}
\csromangathagroup{\csromangathastanza{ยํ กิญฺจิ ธมฺมมภิชญฺญา, \\ อชฺฌตฺตํ อถ วาปิ พหิทฺธา. \\ น เตน ถามํ\footnote{มานํ (สี, ก)} กุพฺเพถ,}}

\csromanheader{2}{น หิ สา นิพฺพุติ สตํ วุตฺตา. (3)}
\csromangathameasure{เสยฺโย น เตน มญฺเญยฺย, \\ นีเจยฺโย อถ วาปิ สริกฺโข. \\ ผุฏฺโฐ อเนกรูเปหิ,}
\csromangathagroup{\csromangathastanza{เสยฺโย น เตน มญฺเญยฺย, \\ นีเจยฺโย อถ วาปิ สริกฺโข. \\ ผุฏฺโฐ\footnote{ปุฏฺโฐ (สี, สฺยา, ก)} อเนกรูเปหิ,}}

\csromanheader{2}{นาตุมานํ วิกปฺปยํ ติฏฺเฐ. (4)}
\csromangathameasure{อชฺฌตฺตเมวุ’ปสเม, \\ น อญฺญโต ภิกฺขุ สนฺติเมเสยฺย. \\ อชฺฌตฺตํ อุปสนฺตสฺส,}
\csromangathagroup{\csromangathastanza{อชฺฌตฺตเมวุ’ปสเม, \\ น อญฺญโต ภิกฺขุ สนฺติเมเสยฺย. \\ อชฺฌตฺตํ อุปสนฺตสฺส,}}

\csromanheader{2}{นตฺถิ อตฺตา กุโต นิรตฺตา วา. (5)}
\gambhira{สุตฺตนิปาตปาฬิ}
\proseitem{926}{\csromanfolio{422}มชฺเฌ ยถา สมุทฺทสฺส,}

\prose{อูมิ โน ชายตี ฐิโต โหติ.}

\prose{เอวํ ฐิโต อเนชสฺส,}

\csromanheader{2}{อุสฺสทํ ภิกฺขุ น กเรยฺย กุหิญฺจิ. (6)}
\proseitem{927}{อกิตฺตยี วิวฏจกฺขุ,}

\prose{สกฺขิธมฺมํ ปริสฺสยวินยํ.}

\prose{ปฏิปทํ วเทหิ ภทฺทนฺเต,}

\csromanheader{2}{ปาติโมกฺขํ อถ วาปิ สมาธิํ. (7)}
\proseitem{928}{จกฺขูหิ เนว โลล’สฺส,}

\prose{คามกถาย อาวรเย โสตํ.}

\prose{รเส จ นานุคิชฺเฌยฺย,}

\csromanheader{2}{น จ มมาเยถ กิญฺจิ โลเกสฺมิํ. (8)}
\proseitem{929}{ผสฺเสน ยทา ผุฏฺฐสฺส,}

\prose{ปริเทวํ ภิกฺขุ น กเรยฺย กุหิญฺจิ.}

\prose{ภวญฺจ นาภิชปฺเปยฺย,}

\csromanheader{2}{เภรเวสุ จ น สมฺปเวเธยฺย. (9)}
\csromanmatikaanchor{mtk.286}
\csromantocmarkref{subsection}{11. กลหวิวาทสุตฺต}{mtk.286}
\bookmark[dest={mtk.286},level=2]{11. กลหวิวาทสุตฺต}
\proseitem{930}{อนฺนานมโถ ปานานํ,}

\prose{ขาทนียานํ อโถปิ วตฺถานํ.}

\prose{ลทฺธา น สนฺนิธิํ กยิรา,}

\csromanheader{2}{น จ ปริตฺตเส ตานิ อลภมาโน. (10)}
\proseitem{931}{ฌายี น ปาทโลล’สฺส,}

\prose{วิรเม กุกฺกุจฺจา นปฺปมชฺเชยฺย.}

\prose{อถาสเนสุ สยเนสุ,}

\csromanheader{2}{อปฺปสทฺเทสุ ภิกฺขุ วิหเรยฺย. (11)}
\proseitem{932}{นิทฺทํ น พหุลีกเรยฺย,}

\prose{ชาคริยํ ภเชยฺย อาตาปี.}

\prose{ตนฺทิํ มายํ หสฺสํ ขิฑฺฑํ,}

\csromanheader{2}{เมถุนํ วิปฺปชเห สวิภูสํ. (12)}
\chapterheadpageread{4. อฏฺฐกวคฺค}
\proseitem{933}{\csromanfolio{423}อาถพฺพณํ สุปินํ ลกฺขณํ,}

\prose{โน วิทเห อโถปิ นกฺขตฺตํ.}

\prose{วิรุตญฺจ คพฺภกรณํ,}

\csromanheader{2}{ติกิจฺฉํ มามโก น เสเวยฺย. (13)}
\proseitem{934}{นินฺทาย นปฺปเวเธยฺย,}

\prose{น อุณฺณเมยฺย ปสํสิโต ภิกฺขุ.}

\prose{โลภํ สห มจฺฉริเยน,}

\csromanheader{2}{โกธํ เปสุณิยญฺจ ปนุเทยฺย. (14)}
\proseitem{935}{กยวิกฺกเย น ติฏฺเฐยฺย,}

\prose{อุปวาทํ ภิกฺขุ น กเรยฺย กุหิญฺจิ.}

\prose{คาเม จ นาภิสชฺเชยฺย,}

\csromanheader{2}{ลาภกมฺยา ชนํ น ลปเยยฺย. (15)}
\proseitem{936}{น จ กตฺถิตา สิยา ภิกฺขุ,}

\prose{น จ วาจํ ปยุตฺตํ ภาเสยฺย.}

\prose{ปาคพฺภิยํ น สิกฺเขยฺย,}

\csromanheader{2}{กถํ วิคฺคาหิกํ น กถเยยฺย. (16)}
\proseitem{937}{โมสวชฺเช น นีเยถ,}

\prose{สมฺปชาโน สฐานิ น กยิรา.}

\prose{อถ ชีวิเตน ปญฺญาย,}

\csromanheader{2}{สีลพฺพเตน นาญฺญมติมญฺเญ. (17)}
\proseitem{938}{สุตฺวา รุสิโต พหุํ วาจํ,}

\prose{สมณานํ วา ปุถุชนานํ\footnote{ปุถุวจนานํ (สี, สฺยา, อิ)}.}

\prose{ผรุเสน เน น ปฏิวชฺชา,}

\csromanheader{2}{น หิ สนฺโต ปฏิเสนิ’กโรนฺติ. (18)}
\proseitem{939}{เอตญฺจ ธมฺมมญฺญาย,}

\prose{วิจินํ ภิกฺขุ สทา สโต สิกฺเข.}

\prose{สนฺตีติ นิพฺพุติํ ญตฺวา,}

\csromanheader{2}{สาสเน โคตมสฺส น ปมชฺเชยฺย. (19)}
\gambhira{สุตฺตนิปาตปาฬิ}
\csromangathameasure{อภิภู หิ โส อนภิภูโต, \\ สกฺขิธมฺม’มนีติห’มทสฺสี. \\ ตสฺมา หิ ตสฺส ภควโต สาสเน, \\ อปฺปมตฺโต สทา นมสฺส’มนุสิกฺเขติ. (20) ตุวฏกสุตฺตํ จุทฺทสมํ.}
\csromangathagroup{\csromangathastanza{\csromanfolio{424}อภิภู หิ โส อนภิภูโต, \\ สกฺขิธมฺม’มนีติห’มทสฺสี. \\ ตสฺมา หิ ตสฺส ภควโต สาสเน, \\ อปฺปมตฺโต สทา นมสฺส’มนุสิกฺเขติ. (20) ตุวฏกสุตฺตํ จุทฺทสมํ.}}

\csromanheader{2}{15. อตฺตทณฺฑสุตฺต}
\vspace{\tipitakaparskip}
\csromancenter{\csromansymbolfootnote{*}{ขุ 7. 315 ปิฏฺเฐ.}อตฺตทณฺฑา ภยํ ชาตํ, ชนํ ปสฺสถ เมธคํ.}
\csromancenter{สํเวคํ กิตฺตยิสฺสามิ, ยถา สํวิชิตํ มยา. (1)}
\vspace{0.5\baselineskip}
\csromangathameasure{ผนฺทมานํ ปชํ ทิสฺวา, มจฺเฉ อปฺโปทเก ยถา.}
\csromangathagroup{\csromangathastanza{ผนฺทมานํ ปชํ ทิสฺวา, มจฺเฉ อปฺโปทเก ยถา.}}

\vspace{\tipitakaparskip}
\csromancenter{อญฺญมญฺเญหิ พฺยารุทฺเธ, ทิสฺวา มํ ภยมาวิสิ. (2)}
\vspace{0.5\baselineskip}
\csromangathameasure{สมนฺตมสาโร โลโก, ทิสา สพฺพา สเมริตา.}
\csromangathagroup{\csromangathastanza{สมนฺตมสาโร โลโก, ทิสา สพฺพา สเมริตา.}}

\proseitem{943}{อิจฺฉํ ภวนมตฺตโน, นาทฺทสาสิํ อโนสิตํ. (3)}

\csromanmatikaanchor{mtk.287}
\csromantocmarkref{subsection}{12. จูฬพฺยูหสุตฺต}{mtk.287}
\bookmark[dest={mtk.287},level=2]{12. จูฬพฺยูหสุตฺต}
\csromangathameasure{โอสาเนเตฺวว พฺยารุทฺเธ, ทิสฺวา เม อรตี อหุ.}
\csromangathagroup{\csromangathastanza{โอสาเนเตฺวว พฺยารุทฺเธ, ทิสฺวา เม อรตี อหุ.}}

\proseitem{944}{อเถตฺถ สลฺลมทฺทกฺขิํ, ทุทฺทสํ หทยนิสฺสิตํ. (4)}

\csromangathameasure{เยน สลฺเลน โอติณฺโณ, ทิสา สพฺพา วิธาวติ.}
\csromangathagroup{\csromangathastanza{เยน สลฺเลน โอติณฺโณ, ทิสา สพฺพา วิธาวติ.}}

\proseitem{945}{ตเมว สลฺลมพฺพุยฺห, น ธาวติ น สีทติ. (5)}

\csromangathameasure{ตตฺถ สิกฺขานุคียนฺติ, \\ ยานิ โลเก คธิตานิ, \\ น เตสุ ปสุโต สิยา. \\ นิพฺพิชฺฌ สพฺพโส กาเม,}
\csromangathagroup{\csromangathastanza{ตตฺถ สิกฺขานุคียนฺติ\footnote{สิกฺขานุกิริยนฺติ (ก)}, \\ ยานิ โลเก คธิตานิ, \\ น เตสุ ปสุโต สิยา. \\ นิพฺพิชฺฌ สพฺพโส กาเม,}}

\csromanheader{2}{สิกฺเข นิพฺพานมตฺตโน. (6)}
\csromangathameasure{สจฺโจ สิยา อปฺปคพฺโภ, อมาโย ริตฺตเปสุโณ.}
\csromangathagroup{\csromangathastanza{สจฺโจ สิยา อปฺปคพฺโภ, อมาโย ริตฺตเปสุโณ.}}

\vspace{\tipitakaparskip}
\csromancenter{อกฺโกธโน โลภปาปํ, เววิจฺฉํ วิตเร มุนิ. (7)}
\vspace{0.5\baselineskip}
\csromangathameasure{นิทฺทํ ตนฺทิํ สเห ถีนํ, ปมาเทน น สํวเส.}
\csromangathagroup{\csromangathastanza{นิทฺทํ ตนฺทิํ สเห ถีนํ, ปมาเทน น สํวเส.}}

\proseitem{948}{อติมาเน น ติฏฺเฐยฺย, นิพฺพานมนโส นโร. (8)}

\chapterheadpageread{4. อฏฺฐกวคฺค}
\proseitem{949}{\csromanfolio{425}โมสวชฺเช น นีเยถ, รูเป สฺเนหํ น กุพฺพเย.}

\prose{มานญฺจ ปริชาเนยฺย, สาหสา วิรโต จเร. (9)}

\proseitem{950}{ปุราณํ นาภินนฺเทยฺย, นเว ขนฺติํ น กุพฺพเย.}

\prose{หิยฺยมาเน น โสเจยฺย, อากาสํ น สิโต สิยา. (10)}

\proseitem{951}{เคธํ พฺรูมิ มโหโฆติ, อาชวํ พฺรูมิ ชปฺปนํ.}

\prose{อารมฺมณํ ปกปฺปนํ, กามปงฺโก ทุรจฺจโย. (11)}

\proseitem{952}{สจฺจา อโวกฺกมฺม\footnote{อโวกฺกมํ (นิทฺเทส)} มุนิ, ถเล ติฏฺฐติ พฺราหฺมโณ.}

\prose{สพฺพํ โส\footnote{สพฺพโส (สฺยา, ก)} ปฏินิสฺสชฺช, ส เว สนฺโตติ วุจฺจติ. (12)}

\proseitem{953}{ส เว วิทฺวา ส เวทคู, ญตฺวา ธมฺมํ อนิสฺสิโต.}

\prose{สมฺมา โส โลเก อิริยาโน, น ปิเหตีธ กสฺสจิ. (13)}

\proseitem{954}{โย’ธ กาเม อจฺจตริ, สงฺคํ โลเก ทุรจฺจยํ.}

\prose{น โส โสจติ นาชฺเฌติ, ฉินฺนโสโต อพนฺธโน. (14)}

\proseitem{955}{ยํ ปุพฺเพ ตํ วิโสเสหิ, ปจฺฉา เต มาหุ กิญฺจนํ.}

\prose{มชฺเฌ เจ โน คเหสฺสสิ, อุปสนฺโต จริสฺสสิ. (15)}

\proseitem{956}{สพฺพโส นามรูปสฺมิํ, ยสฺส นตฺถิ มมายิตํ.}

\prose{อสตา จ น โสจติ, ส เว โลเก น ชียติ. (16)}

\proseitem{957}{ยสฺส นตฺถิ ‘อิทํ เม’ติ, ปเรสํ วาปิ กิญฺจนํ.}

\prose{มมตฺตํ โส อสํวินฺทํ, ‘นตฺถิ เม’ติ น โสจติ. (17)}

\proseitem{958}{อนิฏฺฐุรี อนนุคิทฺโธ, อเนโช สพฺพธี สโม.}

\prose{ตมานิสํสํ ปพฺรูมิ, ปุจฺฉิโต อวิกมฺปินํ. (18)}

\proseitem{959}{อเนชสฺส วิชานโต, นตฺถิ กาจิ นิสงฺขติ\footnote{นิสงฺขิติ (สี, อิ)}.}

\prose{วิรโต โส วิยารพฺภา, เขมํ ปสฺสติ สพฺพธิ. (19)}

\proseitem{960}{น สเมสุ น โอเมสุ, น อุสฺเสสุ วทเต มุนิ.}

\csromangathameasure{สนฺโต โส วีตมจฺฉโร, นาเทติ น นิรสฺสตีติ. (20) อตฺตทณฺฑสุตฺตํ ปนฺนรสมํ.}
\csromangathagroup{\csromangathastanza{สนฺโต โส วีตมจฺฉโร, นาเทติ น นิรสฺสตีติ. (20) อตฺตทณฺฑสุตฺตํ ปนฺนรสมํ.}}

\gambhira{สุตฺตนิปาตปาฬิ}
\csromanheader{2}{16. สาริปุตฺตสุตฺต}
\proseitem{961}{\csromanfolio{426}\csromansymbolfootnote{*}{ขุ 7. 352 ปิฏฺเฐปิ.}น เม ทิฏฺโฐ อิโต ปุพฺเพ, (อิจฺจายสฺมา สาริปุตฺโต,)}

\prose{น สุโต อุท กสฺสจิ.}

\csromangathameasure{เอวํ วคฺคุวโท สตฺถา,}
\csromangathagroup{\csromangathastanza{เอวํ วคฺคุวโท สตฺถา,}}

\csromanheader{2}{ตุสิตา คณิมาคโต. (1)}
\csromangathameasure{สเทวกสฺส โลกสฺส, ยถา ทิสฺสติ จกฺขุมา.}
\csromangathagroup{\csromangathastanza{สเทวกสฺส โลกสฺส, ยถา ทิสฺสติ จกฺขุมา.}}

\proseitem{962}{สพฺพํ ตมํ วิโนเทตฺวา, เอโกว รติมชฺฌคา. (2)}

\csromangathameasure{ตํ พุทฺธํ อสิตํ ตาทิํ, อกุหํ คณิมาคตํ.}
\csromangathagroup{\csromangathastanza{ตํ พุทฺธํ อสิตํ ตาทิํ, อกุหํ คณิมาคตํ.}}

\vspace{\tipitakaparskip}
\csromancenter{พหูนมิธ พทฺธานํ, อตฺถิ ปเญฺหน อาคมํ. (3)}
\vspace{0.5\baselineskip}
\csromangathameasure{ภิกฺขุโน วิชิคุจฺฉโต, ภชโต ริตฺตมาสนํ.}
\csromangathagroup{\csromangathastanza{ภิกฺขุโน วิชิคุจฺฉโต, ภชโต ริตฺตมาสนํ.}}

\vspace{\tipitakaparskip}
\csromancenter{รุกฺขมูลํ สุสานํ วา, ปพฺพตานํ คุหาสุ วา. (4)}
\vspace{0.5\baselineskip}
\csromangathameasure{อุจฺจาวเจสุ สยเนสุ, กีวนฺโต ตตฺถ เภรวา.}
\csromangathagroup{\csromangathastanza{อุจฺจาวเจสุ สยเนสุ, กีวนฺโต ตตฺถ เภรวา.}}

\vspace{\tipitakaparskip}
\csromancenter{เยหิ ภิกฺขุ น เวเธยฺย, นิคฺโฆเส สยนาสเน. (5)}
\vspace{0.5\baselineskip}
\csromangathameasure{กตี ปริสฺสยา โลเก, คจฺฉโต อคตํ ทิสํ.}
\csromangathagroup{\csromangathastanza{กตี ปริสฺสยา โลเก, คจฺฉโต อคตํ ทิสํ.}}

\vspace{\tipitakaparskip}
\csromancenter{เย ภิกฺขุ อภิสมฺภเว, ปนฺตมฺหิ สยนาสเน. (6)}
\vspace{0.5\baselineskip}
\csromangathameasure{กฺยาสฺส พฺยปฺปถโย อสฺสุ, กฺยาสฺสสฺสุ อิธ โคจรา.}
\csromangathagroup{\csromangathastanza{กฺยาสฺส พฺยปฺปถโย อสฺสุ, กฺยาสฺสสฺสุ อิธ โคจรา.}}

\proseitem{967}{กานิ สีลพฺพตานาสฺสุ, ปหิตตฺตสฺส ภิกฺขุโน. (7)}

\csromangathameasure{กํ โส สิกฺขํ สมาทาย, เอโกทิ นิปโก สโต.}
\csromangathagroup{\csromangathastanza{กํ โส สิกฺขํ สมาทาย, เอโกทิ นิปโก สโต.}}

\proseitem{968}{กมฺมาโร รชตสฺเสว, นิทฺธเม มลมตฺตโน. (8)}

\proseitem{969}{วิชิคุจฺฉมานสฺส ยทิทํ ผาสุ, (สาริปุตฺตาติ ภควา,)}

\prose{ริตฺตาสนํ สยนํ เสวโต เจ.}

\csromangathameasure{สมฺโพธิกามสฺส ยถานุธมฺมํ,}
\csromangathagroup{\csromangathastanza{สมฺโพธิกามสฺส ยถานุธมฺมํ,}}

\csromanheader{2}{ตํ เต ปวกฺขามิ ยถา ปชานํ. (9)}
\csromangathameasure{ปญฺจนํ ธีโร ภยานํ น ภาเย, \\ ภิกฺขุ สโต สปริยนฺตจารี. \\ qอํสาธิปาตานํ สรีสปานํ,}
\csromangathagroup{\csromangathastanza{ปญฺจนํ ธีโร ภยานํ น ภาเย, \\ ภิกฺขุ สโต สปริยนฺตจารี. \\ qอํสาธิปาตานํ สรีสปานํ,}}

\csromanheader{2}{มนุสฺสผสฺสานํ จตุปฺปทานํ. (10)}
\chapterheadpageread{4. อฏฺฐกวคฺค}
\proseitem{971}{\csromanfolio{427}ปรธมฺมิกานมฺปิ น สนฺตเสยฺย,}

\prose{ทิสฺวาปิ เตสํ พหุเภรวานิ.}

\prose{อถาปรานิ อภิสมฺภเวยฺย,}

\csromanheader{2}{ปริสฺสยานิ กุสลานุ-เอสี. (11)}
\proseitem{972}{อาตงฺกผสฺเสน ขุทาย ผุฏฺโฐ,}

\csromanmatikaanchor{mtk.288}
\csromantocmarkref{subsection}{13. มหาพฺยูหสุตฺต}{mtk.288}
\bookmark[dest={mtk.288},level=2]{13. มหาพฺยูหสุตฺต}
\prose{สีตํ อถุณฺหํ\footnote{อจฺจุณฺหํ (สี, สฺยา), อตุณฺหํ (ก)} อธิวาสเยยฺย.}

\prose{โส เตหิ ผุฏฺโฐ พหุธา อโนโก,}

\csromanheader{2}{วิริยํ ปรกฺกมฺมทฬฺหํ กเรยฺย. (12)}
\proseitem{973}{เถยฺยํ น กาเร\footnote{น กเรยฺย (สี, สฺยา, ก)} น มุสา ภเณยฺย,}

\prose{เมตฺตาย ผสฺเส ตสถาวรานิ.}

\prose{ยทาวิลตฺตํ มนโส วิชญฺญา,}

\csromanheader{2}{กณฺหสฺส ปกฺโขติ วิโนทเยยฺย. (13)}
\proseitem{974}{โกธาติมานสฺส วสํ น คจฺเฉ,}

\prose{มูลมฺปิ เตสํ ปลิขญฺญ ติฏฺเฐ.}

\prose{อถปฺปิยํ วา ปน อปฺปิยํ วา,}

\csromanheader{2}{อทฺธา ภวนฺโต อภิสมฺภเวยฺย. (14)}
\proseitem{975}{ปญฺญํ ปุรกฺขตฺวา กลฺยาณปีติ,}

\prose{วิกฺขมฺภเย ตานิ ปริสฺสยานิ.}

\prose{อรติํ สเหถ สยนมฺหิ ปนฺเต,}

\csromanheader{2}{จตุโร สเหถ ปริเทวธมฺเม. (15)}
\proseitem{976}{กิํสู อสิสฺสามิ กุว วา\footnote{กุธ วา (ก), กุถ วา (นิทฺเทส)} อสิสฺสํ,}

\prose{ทุกฺขํ วต เสตฺถ กฺวชฺช เสสฺสํ.}

\prose{เอเต วิตกฺเก ปริเทวเนยฺเย,}

\csromanheader{2}{วินเยถ เสโข อนิเกตจารี. (16)}
\proseitem{977}{อนฺนญฺจ ลทฺธา วสนญฺจ กาเล,}

\prose{มตฺตํ โส ชญฺญา อิธ โตสนตฺถํ.}

\prose{โส เตสุ คุตฺโต ยตจาริ คาเม,}

\csromanheader{2}{รุสิโตปิ วาจํ ผรุสํ น วชฺชา. (17)}
\gambhira{สุตฺตนิปาตปาฬิ}
\proseitem{978}{\csromanfolio{428}โอกฺขิตฺตจกฺขุ น จ ปาทโลโล,}

\prose{ฌานานุยุตฺโต พหุชาครสฺส.}

\prose{อุเปกฺขมารพฺภ สมาหิตตฺโต,}

\csromanheader{2}{ตกฺกาสยํ กุกฺกุจฺจิยูปฉินฺเท. (18)}
\proseitem{979}{จุทิโต วจีภิ สติมาภินนฺเท,}

\prose{สพฺรหฺมจารีสุ ขิลํ ปภินฺเท.}

\prose{วาจํ ปมุญฺเจ กุสลํ นาติเวลํ,}

\csromanheader{2}{ชนวาทธมฺมาย น เจตเยยฺย. (19)}
\proseitem{980}{อถาปรํ ปญฺจ รชานิ โลเก,}

\prose{เยสํ สตีมา วินยาย สิกฺเข.}

\prose{รูเปสุ สทฺเทสุ อโถ รเสสุ,}

\csromanheader{2}{คนฺเธสุ ผสฺเสสุ สเหถ ราคํ. (20)}
\proseitem{981}{เอเตสุ ธมฺเมสุ วิเนยฺย ฉนฺทํ,}

\prose{ภิกฺขุ สติมา สุวิมุตฺตจิตฺโต.}

\csromangathameasure{กาเลน โส สมฺมา ธมฺมํ ปริวีมํสมาโน, \\ เอโกทิภูโต วิหเน ตมํ โสติ. (21) สาริปุตฺตสุตฺตํ โสฬสมํ.}
\csromangathagroup{\csromangathastanza{กาเลน โส สมฺมา ธมฺมํ ปริวีมํสมาโน, \\ เอโกทิภูโต วิหเน ตมํ โสติ. (21) สาริปุตฺตสุตฺตํ โสฬสมํ.}}

\vspace{\tipitakaparskip}
\csromancenter{\textbf{อฏฺฐกวคฺโค จตุตฺโถ.}}
\titlehead{\textbf{ตสฺสุทฺทานํ}}
\csromangathameasure{กามํ คุหญฺจ ทุฏฺฐา จ, สุทฺธญฺจ ปรมา ชรา. \\ เมตฺเตยฺโย จ ปสูโร จ, มาคณฺฑิ ปุราเภทนํ. \\ กลหํ เทฺว จ พฺยูหานิ, ปุนเทว ตุวฏฺฏกํ. \\ อตฺตทณฺฑวรํ สุตฺตํ, เถรปุฏฺเฐน โสฬส. \\ อิติ เอตานิ สุตฺตานิ, สพฺพานฏฺฐกวคฺคิกาติ.}
\csromangathagroup{\csromangathastanza{กามํ คุหญฺจ ทุฏฺฐา จ, สุทฺธญฺจ ปรมา ชรา. \\ เมตฺเตยฺโย จ ปสูโร จ, มาคณฺฑิ ปุราเภทนํ.}\csromangathastanza{กลหํ เทฺว จ พฺยูหานิ\footnote{วฺยูหานิ (สี)}, ปุนเทว ตุวฏฺฏกํ. \\ อตฺตทณฺฑวรํ สุตฺตํ, เถรปุฏฺเฐน\footnote{วฺยูหานิ (สี)} โสฬส. \\ อิติ เอตานิ สุตฺตานิ, สพฺพานฏฺฐกวคฺคิกาติ.}}

\chapterheadpageread{5. ปารายนวคฺค}
\csromanheader{2}{1. วตฺถุคาถา}
\vspace{\tipitakaparskip}
\csromancenter{\csromansymbolfootnote{*}{ขุ 8. 1 ปิฏฺเฐ จูฬนิทฺเทเสปิ.}โกสลานํ ปุรา รมฺมา, อคมา ทกฺขิณาปถํ.}
\csromancenter{อากิญฺจญฺญํ ปตฺถยาโน, พฺราหฺมโณ มนฺตปารคู. (1)}
\vspace{0.5\baselineskip}
\csromangathameasure{โส อสฺสกสฺส วิสเย, อฬกสฺส สมาสเน.}
\csromangathagroup{\csromangathastanza{\csromanfolio{429}โส อสฺสกสฺส วิสเย, อฬกสฺส\footnote{มุฬกสฺส (สฺยา), มูฬฺหกสฺส (ก), มฬกสฺส (นิทฺเทส)} สมาสเน.}}

\proseitem{983}{วสิ โคธาวรีกูเล, อุญฺเฉน จ ผเลน จ. (2)}

\csromangathameasure{ตสฺเสว อุปนิสฺสาย, คาโม จ วิปุโล อหุ.}
\csromangathagroup{\csromangathastanza{ตสฺเสว อุปนิสฺสาย, คาโม จ วิปุโล อหุ.}}

\proseitem{984}{ตโต ชาเตน อาเยน, มหายญฺญมกปฺปยิ. (3)}

\csromangathameasure{มหายญฺญํ ยชิตฺวาน, ปุน ปาวิสิ อสฺสมํ.}
\csromangathagroup{\csromangathastanza{มหายญฺญํ ยชิตฺวาน, ปุน ปาวิสิ อสฺสมํ.}}

\vspace{\tipitakaparskip}
\csromancenter{ตสฺมิํ ปฏิปวิฏฺฐมฺหิ, อญฺโญ อาคญฺฉิ พฺราหฺมโณ. (4)}
\vspace{0.5\baselineskip}
\csromangathameasure{อุคฺฆฏฺฏปาโท ตสิโต, ปงฺกทนฺโต รชสฺสิโร.}
\csromangathagroup{\csromangathastanza{อุคฺฆฏฺฏปาโท ตสิโต\footnote{ตสฺสิโต (ก)}, ปงฺกทนฺโต รชสฺสิโร.}}

\proseitem{986}{โส จ นํ อุปสงฺกมฺม, สตานิ ปญฺจ ยาจติ. (5)}

\csromangathameasure{ตเมนํ พาวรี ทิสฺวา, อาสเนน นิมนฺตยิ.}
\csromangathagroup{\csromangathastanza{ตเมนํ พาวรี ทิสฺวา, อาสเนน นิมนฺตยิ.}}

\vspace{\tipitakaparskip}
\csromancenter{สุขญฺจ กุสลํ ปุจฺฉิ, อิทํ วจนมพฺรวิ. (6)}
\vspace{0.5\baselineskip}
\csromangathameasure{ยํ โข มม เทยฺยธมฺมํ, สพฺพํ วิสชฺชิตํ มยา.}
\csromangathagroup{\csromangathastanza{ยํ โข มม เทยฺยธมฺมํ, สพฺพํ วิสชฺชิตํ มยา.}}

\proseitem{988}{อนุชานาหิ เม พฺรหฺเม, นตฺถิ ปญฺจสตานิ เม. (7)}

\csromangathameasure{สเจ เม ยาจมานสฺส, ภวํ นานุปทสฺสติ.}
\csromangathagroup{\csromangathastanza{สเจ เม ยาจมานสฺส, ภวํ นานุปทสฺสติ.}}

\vspace{\tipitakaparskip}
\csromancenter{สตฺตเม ทิวเส ตุยฺหํ, มุทฺธา ผลตุ สตฺตธา. (8)}
\vspace{0.5\baselineskip}
\csromangathameasure{อภิสงฺขริตฺวา กุหโก, เภรวํ โส อกิตฺตยิ.}
\csromangathagroup{\csromangathastanza{อภิสงฺขริตฺวา กุหโก, เภรวํ โส อกิตฺตยิ.}}

\proseitem{990}{ตสฺส ตํ วจนํ สุตฺวา, พาวรี ทุกฺขิโต อหุ. (9)}

\proseitem{991}{อุสฺสุสฺสติ อนาหาโร โสกสลฺลสมปฺปิโต.}

\prose{อโถปิ เอวํ จิตฺตสฺส, ฌาเน น รมตี มโน. (10)}

\csromangathameasure{อุตฺรสฺตํ ทุกฺขิตํ ทิสฺวา, เทวตา อตฺถกามินี.}
\csromangathagroup{\csromangathastanza{อุตฺรสฺตํ ทุกฺขิตํ ทิสฺวา, เทวตา อตฺถกามินี.}}

\proseitem{992}{พาวริํ อุปสงฺกมฺม, อิทํ วจนมพฺรวิ. (11)}

\gambhira{สุตฺตนิปาตปาฬิ}
\proseitem{993}{\csromanfolio{430}น โส มุทฺธํ ปชานาติ, กุหโก โส ธนตฺถิโก.}

\prose{มุทฺธนิ มุทฺธปาเต วา, ญาณํ ตสฺส น วิชฺชติ. (12)}

\proseitem{994}{โภตี จรหิ ชานาสิ, ตํ เม อกฺขาหิ ปุจฺฉิตา.}

\prose{มุทฺธํ มุทฺธาธิปาตญฺจ, ตํ สุโณม วโจ ตว. (13)}

\csromanhangingbegin
\hangingheaditem{995}{อหมฺเปตํ น ชานามิ, ญาณเมตฺถ น วิชฺชติ.}
\hangingbody{มุทฺธนิ มุทฺธาธิปาเต จ, ชินานํ เหตฺถ1 ทสฺสนํ. (14)}
\csromanhangingend

\proseitem{996}{อถ โก จรหิ ชานาติ, อสฺมิํ ปถวิมณฺฑเล\footnote{ปุถวิมณฺฑเล (สี, อิ) 3. คนฺตฺวา นมสฺเสมุ (สี, สฺยา, อิ) 4. ทิปทุตฺตมํ (สี, สฺยา, อิ)}.}

\prose{มุทฺธํ มุทฺธาธิปาตญฺจ, ตํ เม อกฺขาหิ เทวเต. (15)}

\csromanmatikaanchor{mtk.289}
\csromantocmarkref{subsection}{14. ตุวฏกสุตฺต}{mtk.289}
\bookmark[dest={mtk.289},level=2]{14. ตุวฏกสุตฺต}
\proseitem{997}{ปุรา กปิลวตฺถุมฺหา, นิกฺขนฺโต โลกนายโก.}

\prose{อปจฺโจ โอกฺกากราชสฺส, สกฺยปุตฺโต ปภงฺกโร. (16)}

\proseitem{998}{โส หิ พฺราหฺมณ สมฺพุทฺโธ, สพฺพธมฺมาน ปารคู.}

\csromangathameasure{สพฺพาภิญฺญาพลปฺปตฺโต, สพฺพธมฺเมสุ จกฺขุมา.}
\csromangathagroup{\csromangathastanza{สพฺพาภิญฺญาพลปฺปตฺโต, สพฺพธมฺเมสุ จกฺขุมา.}}

\prose{สพฺพกมฺมกฺขยํ ปตฺโต, วิมุตฺโต อุปธิกฺขเย. (17)}

\proseitem{999}{พุทฺโธ โส ภควา โลเก, ธมฺมํ เทเสติ จกฺขุมา.}

\prose{ตํ ตฺวํ คนฺตฺวาน ปุจฺฉสฺสุ, โส เต ตํ พฺยากริสฺสติ. (18)}

\proseitem{1000}{สมฺพุทฺโธติ วโจ สุตฺวา, อุทคฺโค พาวรี อหุ.}

\prose{โสก’สฺส ตนุโก อาสิ, ปีติญฺจ วิปุลํ ลภิ. (19)}

\proseitem{1001}{โส พาวรี อตฺตมโน อุทคฺโค,}

\prose{ตํ เทวตํ ปุจฺฉติ เวทชาโต.}

\csromangathameasure{กตมมฺหิ คาเม นิคมมฺหิ วา ปน, \\ กตมมฺหิ วา ชนปเท โลกนาโถ. \\ ยตฺถ คนฺตฺวาน ปสฺเสมุ3,}
\csromangathagroup{\csromangathastanza{กตมมฺหิ คาเม นิคมมฺหิ วา ปน, \\ กตมมฺหิ วา ชนปเท โลกนาโถ. \\ ยตฺถ คนฺตฺวาน ปสฺเสมุ3,}}

\csromanheader{2}{สมฺพุทฺธํ ทฺวิปทุตฺตมํ4. (20)}
\proseitem{1002}{สาวตฺถิยํ โกสลมนฺทิเร ชิโน,}

\prose{ปหูตปญฺโญ วรภูริเมธโส.}

\prose{โส สกฺยปุตฺโต วิธุโร อนาสโว,}

\csromanheader{2}{มุทฺธาธิปาตสฺส วิทู นราสโภ. (21)}
\chapterheadpageread{5. ปารายนวคฺค}
\proseitem{1003}{\csromanfolio{431}ตโต อามนฺตยี สิสฺเส, พฺราหฺมเณ มนฺตปารเค.}

\prose{“เอถ มาณวา อกฺขิสฺสํ, สุณาถ วจนํ มม. (22)}

\proseitem{1004}{ยสฺเสโส ทุลฺลโภ โลเก,}

\prose{ปาตุภาโว อภิณฺหโส.}

\csromangathameasure{สฺวาชฺช โลกมฺหิ อุปฺปนฺโน, \\ สมฺพุทฺโธ อิติ วิสฺสุโต. \\ ขิปฺปํ คนฺตฺวาน สาวตฺถิํ,}
\csromangathagroup{\csromangathastanza{สฺวาชฺช โลกมฺหิ อุปฺปนฺโน, \\ สมฺพุทฺโธ อิติ วิสฺสุโต. \\ ขิปฺปํ คนฺตฺวาน สาวตฺถิํ,}}

\csromanheader{2}{ปสฺสโวฺห ทฺวิปทุตฺตมํ”. (23)}
\proseitem{1005}{กถํ จรหิ ชาเนมุ, ทิสฺวา พุทฺโธติ พฺราหฺมณ.}

\prose{อชานตํ โน ปพฺรูหิ, ยถา ชาเนมุ ตํ มยํ. (24)}

\proseitem{1006}{อาคตานิ หิ มนฺเตสุ, มหาปุริสลกฺขณา.}

\prose{ทฺวตฺติํสานิ จ\footnote{ทฺวตฺติํสา จ (สี, สา, อิ), ทฺวตฺติํส ตานิ (?)} พฺยากฺขาตา, สมตฺตา อนุปุพฺพโส. (25)}

\proseitem{1007}{ยสฺเสเต โหนฺติ คตฺเตสุ, มหาปุริสลกฺขณา.}

\prose{เทฺวเยว ตสฺส คติโย, ตติยา หิ น วิชฺชติ. (26)}

\proseitem{1008}{สเจ อคารํ อาวสติ\footnote{อชฺฌาวสติ (ก)}, วิเชยฺย ปถวิํ อิมํ.}

\prose{อทณฺเฑน อสตฺเถน, ธมฺเมน มนุสาสติ. (27)}

\proseitem{1009}{สเจ จ โส ปพฺพชติ, อคารา อนคาริยํ.}

\prose{วิวฏฺฏจฺฉโท\footnote{วิวตฺตจฺฉทฺโท (สี)} สมฺพุทฺโธ, อรหา ภวติ อนุตฺตโร. (28)}

\proseitem{1010}{ชาติํ โคตฺตญฺจ ลกฺขณํ, มนฺเต สิสฺเส ปุนาปเร.}

\prose{มุทฺธํ มุทฺธาธิปาตญฺจ, มนสาเยว ปุจฺฉถ. (29)}

\proseitem{1011}{อนาวรณทสฺสาวี, ยทิ พุทฺโธ ภวิสฺสติ.}

\prose{มนสา ปุจฺฉิเต ปเญฺห, วาจาย วิสเชสฺสติ. (30)}

\proseitem{1012}{พาวริสฺส วโจ สุตฺวา, สิสฺสา โสฬส พฺราหฺมณา.}

\prose{อชิโต ติสฺส เมตฺเตยฺโย, ปุณฺณโก อถ เมตฺตคู. (31)}

\proseitem{1013}{โธตโก อุปสีโว จ, นนฺโท จ อถ เหมโก.}

\prose{โตเทยฺย-กปฺปา ทุภโย, ชตุกณฺณี จ ปณฺฑิโต. (32)}

\gambhira{สุตฺตนิปาตปาฬิ}
\proseitem{1014}{\csromanfolio{432}ภทฺราวุโธ อุทโย จ, โปสาโล จาปิ พฺราหฺมโณ.}

\prose{โมฆราชา จ เมธาวี, ปิงฺคิโย จ มหา-อิสิ. (33)}

\proseitem{1015}{ปจฺเจกคณิโน สพฺเพ, สพฺพโลกสฺส วิสฺสุตา.}

\prose{ฌายี ฌานรตา ธีรา, ปุพฺพวาสนวาสิตา. (34)}

\proseitem{1016}{พาวริํ อภิวาเทตฺวา, กตฺวา จ นํ ปทกฺขิณํ.}

\prose{ชฏาชินธรา สพฺเพ, ปกฺกามุํ อุตฺตรามุขา. (35)}

\proseitem{1017}{อฬกสฺส ปติฏฺฐานํ, ปุริมาหิสฺสติํ\footnote{ปุริมํ มาหิสฺสติํ (สี, อิ), ปุรํ มาหิสฺสติํ (สฺยา)} ตทา.}

\prose{อุชฺเชนิญฺจาปิ โคนทฺธํ, เวทิสํ วนสวฺหยํ. (36)}

\proseitem{1018}{โกสมฺพิญฺจาปิ สาเกตํ, สาวตฺถิญฺจ ปุรุตฺตมํ.}

\prose{เสตพฺยํ กปิลวตฺถุํ, กุสินารญฺจ มนฺทิรํ. (37)}

\proseitem{1019}{ปาวญฺจ โภคนครํ, เวสาลิํ มาคธํ ปุรํ.}

\prose{ปาสาณกํ เจติยญฺจ, รมณียํ มโนรมํ. (38)}

\proseitem{1020}{ตสิโตวุทกํ สีตํ, มหาลาภํว วาณิโช.}

\prose{ฉายํ ฆมฺมาภิตตฺโตว, ตุริตา ปพฺพตมารุหุํ. (39)}

\proseitem{1021}{ภควา ตมฺหิ สมเย, ภิกฺขุสํฆปุรกฺขโต.}

\prose{ภิกฺขูนํ ธมฺมํ เทเสติ, สีโหว นทตี วเน. (40)}

\proseitem{1022}{อชิโต อทฺทส พุทฺธํ, สตรํสิํว\footnote{วีตรํสิํว (สฺยา), สตรํสีว (ก), ปีตรํสีว (นิทฺเทส)} ภาณุมํ.}

\prose{จนฺทํ ยถา ปนฺนรเส, ปาริปูริํ อุปาคตํ. (41)}

\proseitem{1023}{อถสฺส คตฺเต ทิสฺวาน, ปริปูรญฺจ พฺยญฺชนํ.}

\prose{เอกมนฺตํ ฐิโต หฏฺโฐ, มโนปเญฺห อปุจฺฉถ. (42)}

\proseitem{1024}{อาทิสฺส ชมฺมนํ\footnote{ชปฺปนํ (ก)} พฺรูหิ, โคตฺตํ พฺรูหิ สลกฺขณํ\footnote{พฺรูหิสฺส ลกฺขณํ (นิทฺเทส)}.}

\prose{มนฺเตสุ ปารมิํ พฺรูหิ, กติ วาเจติ พฺราหฺมโณ. (43)}

\proseitem{1025}{วีสํ วสฺสสตํ อายุ, โส จ โคตฺเตน พาวรี.}

\prose{ตีณิสฺส ลกฺขณา คตฺเต, ติณฺณํ เวทาน ปารคู. (44)}

\chapterheadpageread{5. ปารายนวคฺค}
\proseitem{1026}{\csromanfolio{433}ลกฺขเณ อิติหาเส จ, สนิฆณฺฑุสเกฏุเภ.}

\prose{ปญฺจสตานิ วาเจติ, สธมฺเม ปารมิํ คโต. (45)}

\proseitem{1027}{ลกฺขณานํ ปวิจยํ, พาวริสฺส นรุตฺตม.}

\prose{กงฺขจฺฉิท\footnote{ตณฺหจฺฉิท (พหูสุ)} ปกาเสหิ, มา โน กงฺขายิตํ อหุ. (46)}

\proseitem{1028}{มุขํ ชิวฺหาย ฉาเทติ, อุณฺณ’สฺส ภมุกนฺตเร.}

\prose{โกโสหิตํ วตฺถคุยฺหํ, เอวํ ชานาหิ มาณว. (47)}

\proseitem{1029}{ปุจฺฉญฺหิ กิญฺจิ อสุณนฺโต, สุตฺวา ปเญฺห วิยากเต.}

\csromanmatikaanchor{mtk.290}
\csromantocmarkref{subsection}{15. อตฺตทณฺฑสุตฺต}{mtk.290}
\bookmark[dest={mtk.290},level=2]{15. อตฺตทณฺฑสุตฺต}
\prose{วิจินฺเตติ ชโน สพฺโพ, เวทชาโต กตญฺชลี. (48)}

\proseitem{1030}{โก นุ เทโว วา พฺรหฺมา วา, อินฺโท วาปิ สุชมฺปติ.}

\prose{มนสา ปุจฺฉิเต ปเญฺห, กเมตํ ปฏิภาสติ. (49)}

\proseitem{1031}{มุทฺธํ มุทฺธาธิปาตญฺจ, พาวรี ปริปุจฺฉติ.}

\prose{ตํ พฺยากโรหิ ภควา, กงฺขํ วินย โน อิเส. (50)}

\proseitem{1032}{อวิชฺชา มุทฺธาติ ชานาหิ, วิชฺชา มุทฺธาธิปาตินี.}

\prose{สทฺธาสติสมาธีหิ, ฉนฺทวิริเยน สํยุตา. (51)}

\proseitem{1033}{ตโต เวเทน มหตา, สนฺถมฺภิตฺวาน มาณโว.}

\prose{เอกํสํ อชินํ กตฺวา, ปาเทสุ สิรสา ปติ. (52)}

\proseitem{1034}{พาวรี พฺราหฺมโณ โภโต, สห สิสฺเสหิ มาริส.}

\prose{อุทคฺคจิตฺโต สุมโน, ปาเท วนฺทติ จกฺขุม. (53)}

\proseitem{1035}{สุขิโต พาวรี โหตุ, สห สิสฺเสหิ พฺราหฺมโณ.}

\prose{ตฺวญฺจาปิ สุขิโต โหหิ, จิรํ ชีวาหิ มาณว. (54)}

\proseitem{1036}{พาวริสฺส จ ตุยฺหํ วา, สพฺเพสํ สพฺพสํสยํ.}

\prose{กตาวกาสา ปุจฺฉโวฺห, ยํ กิญฺจิ มนสิจฺฉถ. (55)}

\proseitem{1037}{สมฺพุทฺเธน กโตกาโส, นิสีทิตฺวาน ปญฺชลี.}

\nitthitam{อชิโต ปฐมํ ปญฺหํ, ตตฺถ ปุจฺฉิ ตถาคตํ. (56) วตฺถุคาถา นิฏฺฐิตา.}
\gambhira{สุตฺตนิปาตปาฬิ}
\csromanheader{2}{1. อชิตมาณวปุจฺฉา}
\proseitem{1038}{\csromanfolio{434}\csromansymbolfootnote{*}{ขุ 8. 6, 24; ขุ 10. 10, 13, 177 ปิฏฺเฐสุปิ.}เกนสฺสุ นิวุโต โลโก, (อิจฺจายสฺมา อชิโต,)}

\prose{เกนสฺสุ นปฺปกาสติ.}

\vspace{\tipitakaparskip}
\csromancenter{กิสฺสาภิเลปนํ พฺรูสิ, กิํสุ ตสฺส มหพฺภยํ. (1)}
\proseitem{1039}{อวิชฺชาย นิวุโต โลโก, (อชิตาติ ภควา,)}

\prose{เววิจฺฉา ปมาทา นปฺปกาสติ.}

\prose{ชปฺปาภิเลปนํ พฺรูมิ, ทุกฺข’มสฺส มหพฺภยํ. (2)}

\proseitem{1040}{สวนฺติ สพฺพธิ โสตา, (อิจฺจายสฺมา อชิโต,)}

\prose{โสตานํ กิํ นิวารณํ.}

\vspace{\tipitakaparskip}
\csromancenter{โสตานํ สํวรํ พฺรูหิ, เกน โสตา ปิธิยฺยเร\footnote{ปิถิยฺยเร (สี, สฺยา, อิ), ปิถียเร (สี-ฏฺฐ), ปิธียเร (?)}. (3)}
\proseitem{1041}{ยานิ โสตานิ โลกสฺมิํ, (อชิตาติ ภควา,)}

\prose{สติ เตสํ นิวารณํ.}

\vspace{\tipitakaparskip}
\csromancenter{โสตานํ สํวรํ พฺรูมิ, ปญฺญาเยเต ปิธิยฺยเร. (4)}
\proseitem{1042}{ปญฺญา เจว สติ ยญฺจ\footnote{สตี เจว (สี), สตี จ (สฺยา), สตี จาปิ (อิ, นิทฺเทส), สติ จาปิ (นิทฺเทส)}, (อิจฺจายสฺมา อชิโต,)}

\prose{นามรูปญฺจ มาริส.}

\prose{เอตํ เม ปุฏฺโฐ ปพฺรูหิ, กตฺเถตํ อุปรุชฺฌติ. (5)}

\proseitem{1043}{ยเมตํ ปญฺหํ อปุจฺฉิ, อชิต ตํ วทามิ เต.}

\csromangathameasure{ยตฺถ นามญฺจ รูปญฺจ, อเสสํ อุปรุชฺฌติ.}
\csromangathagroup{\csromangathastanza{ยตฺถ นามญฺจ รูปญฺจ, อเสสํ อุปรุชฺฌติ.}}

\prose{วิญฺญาณสฺส นิโรเธน, เอตฺเถตํ อุปรุชฺฌติ. (6)}

\proseitem{1044}{เย จ สงฺขาตธมฺมาเส, เย จ เสขา ปุถู อิธ.}

\vspace{\tipitakaparskip}
\csromancenter{เตสํ เม นิปโก อิริยํ, ปุฏฺโฐ ปพฺรูหิ มาริส. (7)}
\proseitem{1045}{กาเมสุ นาภิคิชฺเฌยฺย, มนสา’นาวิโล สิยา.}

\csromangathameasure{กุสโล สพฺพธมฺมานํ, สโต ภิกฺขุ ปริพฺพเชติ. (8) อชิตมาณวปุจฺฉา ปฐมา.}
\csromangathagroup{\csromangathastanza{กุสโล สพฺพธมฺมานํ, สโต ภิกฺขุ ปริพฺพเชติ. (8) อชิตมาณวปุจฺฉา ปฐมา.}}

\csromanheader{2}{5. ปารายนวคฺค \textbf{2. ติสฺสเมตฺเตยฺยมาณวปุจฺฉา}}
\proseitem{1046}{\csromanfolio{435}\csromansymbolfootnote{*}{ขุ 8. 7, 40 ปิฏฺเฐสุปิ.}โกธ สนฺตุสิโต โลเก, (อิจฺจายสฺมา ติสฺสเมตฺเตยฺโย,)}

\csromanmatikaanchor{mtk.291}
\csromantocmarkref{subsection}{16. สาริปุตฺตสุตฺต}{mtk.291}
\bookmark[dest={mtk.291},level=2]{16. สาริปุตฺตสุตฺต}
\prose{กสฺส โน สนฺติ อิญฺชิตา.}

\csromangathameasure{โก อุภนฺตมภิญฺญาย, มชฺเฌ มนฺตา น ลิปฺปติ.}
\csromangathagroup{\csromangathastanza{โก อุภนฺตมภิญฺญาย, มชฺเฌ มนฺตา น ลิปฺปติ\footnote{ลิมฺปติ (ก)}.}}

\vspace{\tipitakaparskip}
\csromancenter{กํ พฺรูสิ มหาปุริโสติ, โก อิธ สิพฺพินิมจฺจคา. (1)}
\proseitem{1047}{กาเมสุ พฺรหฺมจริยวา, (เมตฺเตยฺยาติ ภควา,)}

\prose{วีตตโณฺห สทา สโต.}

\prose{สงฺขาย นิพฺพุโต ภิกฺขุ, ตสฺส โน สนฺติ อิญฺชิตา. (2)}

\proseitem{1048}{โส อุภนฺตมภิญฺญาย, มชฺเฌ มนฺตา น ลิปฺปติ.}

\csromangathameasure{ตํ พฺรูมิ มหาปุริโสติ, โส อิธ สิพฺพินิมจฺจคาติ. (3) ติสฺสเมตฺเตยฺยมาณวปุจฺฉา ทุติยา.}
\csromangathagroup{\csromangathastanza{ตํ พฺรูมิ มหาปุริโสติ, โส อิธ สิพฺพินิมจฺจคาติ. (3) ติสฺสเมตฺเตยฺยมาณวปุจฺฉา ทุติยา.}}

\csromanheader{2}{3. ปุณฺณกมาณวปุจฺฉา}
\proseitem{1049}{\csromansymbolfootnote{+}{ขุ 8. 7, 46 ปิฏฺเฐสุปิ.}อเนชํ มูลทสฺสาวิํ, (อิจฺจายสฺมา ปุณฺณโก,)}

\prose{อตฺถิ\footnote{อตฺถี (สฺยา)} ปเญฺหน อาคมํ.}

\csromangathameasure{กิํ นิสฺสิตา อิสโย มนุชา, \\ ขตฺติยา พฺราหฺมณา เทวตานํ. \\ ยญฺญมกปฺปยิํสุ ปุถู’ธ โลเก,}
\csromangathagroup{\csromangathastanza{กิํ นิสฺสิตา อิสโย มนุชา, \\ ขตฺติยา พฺราหฺมณา เทวตานํ. \\ ยญฺญมกปฺปยิํสุ ปุถู’ธ โลเก,}}

\csromanheader{2}{ปุจฺฉามิ ตํ ภควา พฺรูหิ เม’ตํ. (1)}
\proseitem{1050}{เย เกจิเม อิสโย มนุชา, (ปุณฺณกาติ ภควา,)}

\prose{ขตฺติยา พฺราหฺมณา เทวตานํ.}

\csromangathameasure{ยญฺญมกปฺปยิํสุ ปุถู’ธ โลเก, \\ อาสีสมานา ปุณฺณก อิตฺถตฺตํ.}
\csromangathagroup{\csromangathastanza{ยญฺญมกปฺปยิํสุ ปุถู’ธ โลเก, \\ อาสีสมานา ปุณฺณก อิตฺถตฺตํ\footnote{อิตฺถภาวํ (สี, สฺยา)}.}}

\csromanheader{2}{ชรํ สิตา ยญฺญมกปฺปยิํสุ. (2)}
\gambhira{สุตฺตนิปาตปาฬิ}
\proseitem{1051}{\csromanfolio{436}เย เกจิเม อิสโย มนุชา, (อิจฺจายสฺมา ปุณฺณโก,)}

\prose{ขตฺติยา พฺราหฺมณา เทวตานํ.}

\csromangathameasure{ยญฺญมกปฺปยิํสุ ปุถู’ธ โลเก, \\ กจฺจิสฺสุ เต ภควา ยญฺญปเถ อปฺปมตฺตา. \\ อตารุํ ชาติญฺจ ชรญฺจ มาริส,}
\csromangathagroup{\csromangathastanza{ยญฺญมกปฺปยิํสุ ปุถู’ธ โลเก, \\ กจฺจิสฺสุ เต ภควา ยญฺญปเถ อปฺปมตฺตา. \\ อตารุํ ชาติญฺจ ชรญฺจ มาริส,}}

\csromanheader{2}{ปุจฺฉามิ ตํ ภควา พฺรูหิ เป’ตํ. (3)}
\proseitem{1052}{อาสีสนฺติ โถมยนฺติ อภิชปฺปนฺติ ชุหนฺติ, (ปุณฺณกาติ ภควา,)}

\prose{กามาภิชปฺปนฺติ ปฏิจฺจ ลาภํ.}

\csromangathameasure{เต ยาชโยคา ภวราครตฺตา,}
\csromangathagroup{\csromangathastanza{เต ยาชโยคา ภวราครตฺตา,}}

\csromanheader{2}{นาตริํสุ ชาติชรนฺติ พฺรูมิ. (4)}
\proseitem{1053}{เต เจ นาตริํสุ ยาชโยคา, (อิจฺจายสฺมา ปุณฺณโก,)}

\prose{ยญฺเญหิ ชาติญฺจ ชรญฺจ มาริส.}

\csromangathameasure{อถ โก จรหิ เทวมนุสฺสโลเก, \\ อตาริ ชาติญฺจ ชรญฺจ มาริส.}
\csromangathagroup{\csromangathastanza{อถ โก จรหิ เทวมนุสฺสโลเก, \\ อตาริ ชาติญฺจ ชรญฺจ มาริส.}}

\csromanheader{2}{ปุจฺฉามิ ตํ ภควา พฺรูหิ เม’ตํ. (5)}
\proseitem{1054}{สงฺขาย โลกสฺมิ ปโรปรานิ\footnote{ปโรวรานิ (สี, สฺยา)}, (ปุณฺณกาติ ภควา,)}

\prose{ยสฺสิญฺชิตํ นตฺถิ กุหิญฺจิ โลเก.}

\csromangathameasure{สนฺโต วิธูโม อนีโฆ นิราโส, \\ อตาริ โส ชาติชรนฺติ พฺรูมีติ. (6) ปุณฺณกมาณวปุจฺฉา ตติยา.}
\csromangathagroup{\csromangathastanza{สนฺโต วิธูโม อนีโฆ นิราโส, \\ อตาริ โส ชาติชรนฺติ พฺรูมีติ. (6) ปุณฺณกมาณวปุจฺฉา ตติยา.}}

\csromanheader{2}{4. เมตฺตคูมาณวปุจฺฉา}
\proseitem{1055}{\csromansymbolfootnote{*}{ขุ 8. 8, 62 ปิฏฺเฐสุปิ.}ปุจฺฉามิ ตํ ภควา พฺรูหิ เม’ตํ, (อิจฺจายสฺมา เมตฺตคู,)}

\prose{มญฺญามิ ตํ เวทคุํ ภาปิตตฺตํ.}

\csromangathameasure{กุโต นุ ทุกฺขา สมุทาคตา อิเม,}
\csromangathagroup{\csromangathastanza{กุโต นุ ทุกฺขา สมุทาคตา อิเม,}}

\csromanheader{2}{เย เกจิ โลกสฺมิ’มเนกรูปา. (1)}
\chapterheadpageread{5. ปารายนวคฺค}
\proseitem{1056}{\csromanfolio{437}ทุกฺขสฺส เว มํ ปภวํ อปุจฺฉสิ, (เมตฺตคูติ ภควา,)}

\prose{ตํ เต ปวกฺขามิ ยถา ปชานํ.}

\prose{อุปธินิทานา ปภวนฺติ ทุกฺขา,}

\csromanheader{2}{เย เกจิ โลกสฺมิ’มเนกรูปา. (2)}
\proseitem{1057}{โย เว อวิทฺวา อุปธิํ กโรติ,}

\prose{ปุนปฺปุนํ ทุกฺข’มุเปติ มนฺโท.}

\prose{ตสฺมา ปชานํ อุปธิํ น กยิรา,}

\csromanheader{2}{ทุกฺขสฺส ชาติปฺปภวานุปสฺสี. (3)}
\proseitem{1058}{ยํ ตํ อปุจฺฉิมฺห อกิตฺตยี โน,}

\prose{อญฺญํ ตํ ปุจฺฉาม\footnote{ปุจฺฉามิ (สี, อิ)} ตทิงฺฆ พฺรูหิ.}

\csromangathameasure{กถํ นุ ธีรา วิตรนฺติ โอฆํ, \\ ชาติํ ชรํ โสกปริทฺทวญฺจ. \\ ตํ เม มุนิ สาธุ วิยากโรหิ,}
\csromangathagroup{\csromangathastanza{กถํ นุ ธีรา วิตรนฺติ โอฆํ, \\ ชาติํ ชรํ โสกปริทฺทวญฺจ. \\ ตํ เม มุนิ สาธุ วิยากโรหิ,}}

\csromanheader{2}{ตถา หิ เต วิทิโต เอส ธมฺโม. (4)}
\proseitem{1059}{กิตฺตยิสฺสามิ เต ธมฺมํ, (เมตฺตคูติ ภควา,)}

\prose{ทิฏฺเฐ ธมฺเม อนีติหํ.}

\prose{ยํ วิทิตฺวา สโต จรํ, ตเร โลเก วิสตฺติกํ. (5)}

\proseitem{1060}{ตญฺจาหํ อภินนฺทามิ, มเหสิ ธมฺมมุตฺตมํ.}

\prose{ยํ วิทิตฺวา สโต จรํ, ตเร โลเก วิสตฺติกํ. (6)}

\proseitem{1061}{ยํ กิญฺจิ สมฺปชานาสิ, (เมตฺตคูติ ภควา,)}

\prose{อุทฺธํ อโธ ติริยญฺจาปิ มชฺเฌ.}

\prose{เอเตสุ นนฺทิญฺจ นิเวสนญฺจ,}

\csromanheader{2}{ปนุชฺช วิญฺญาณํ ภเว น ติฏฺเฐ. (7)}
\proseitem{1062}{เอวํวิหารี สโต อปฺปมตฺโต,}

\prose{ภิกฺขุ จรํ หิตฺวา มมายิตานิ.}

\prose{ชาติํ ชรํ โสกปริทฺทวญฺจ,}

\csromanheader{2}{อิเธว วิทฺวา ปชเหยฺย ทุกฺขํ. (8)}
\gambhira{สุตฺตนิปาตปาฬิ}
\proseitem{1063}{\csromanfolio{438}เอตา’ภินนฺทามิ วโจ มเหสิโน,}

\prose{สุกิตฺติตํ โคตม’นูปธีกํ.}

\csromangathameasure{อทฺธา หิ ภควา ปหาสิ ทุกฺขํ,}
\csromangathagroup{\csromangathastanza{อทฺธา หิ ภควา ปหาสิ ทุกฺขํ,}}

\csromanheader{2}{ตถา หิ เต วิทิโต เอส ธมฺโม. (9)}
\proseitem{1064}{เต จาปิ นูนปฺปชเหยฺยุ ทุกฺขํ,}

\prose{เย ตฺวํ มุนิ อฏฺฐิตํ โอวเทยฺย.}

\csromangathameasure{ตํ ตํ นมสฺสามิ สเมจฺจ นาค,}
\csromangathagroup{\csromangathastanza{ตํ ตํ นมสฺสามิ สเมจฺจ นาค,}}

\csromanheader{2}{อปฺเปว มํ ภควา อฏฺฐิตํ โอวเทยฺย. (10)}
\proseitem{1065}{ยํ พฺราหฺมณํ เวทคุ’มาภิชญฺญา,}

\prose{อกิญฺจนํ กามภเว อสตฺตํ.}

\csromangathameasure{อทฺธา หิ โส โอฆมิมํ อตาริ,}
\csromangathagroup{\csromangathastanza{อทฺธา หิ โส โอฆมิมํ อตาริ,}}

\csromanmatikaanchor{mtk.292}
\csromanmatikaanchor{mtk.293}
\csromantocmarknopage{chapter}{5. ปารายนวคฺค}{mtk.292}
\bookmark[dest={mtk.292},level=0]{5. ปารายนวคฺค}
\csromantocmarkref{subsection}{1. วตฺถุคาถา}{mtk.293}
\bookmark[dest={mtk.293},level=2]{1. วตฺถุคาถา}
\csromanheader{2}{ติณฺโณ จ ปารํ อขิโล อกงฺโข. (11)}
\proseitem{1066}{วิทฺวา จ โย\footnote{โส (สี, สฺยา, อิ)} เวทคู นโร อิธ,}

\prose{ภวาภเว สงฺคมิมํ วิสชฺช.}

\csromangathameasure{โส วีตตโณฺห อนีโฆ นิราโส, \\ อตาริ โส ชาติชรนฺติ พฺรูมีติ. (12) เมตฺตคูมาณวปุจฺฉา จตุตฺถี.}
\csromangathagroup{\csromangathastanza{โส วีตตโณฺห อนีโฆ นิราโส, \\ อตาริ โส ชาติชรนฺติ พฺรูมีติ. (12) เมตฺตคูมาณวปุจฺฉา จตุตฺถี.}}

\csromanheader{2}{5. โธตกมาณวปุจฺฉา}
\proseitem{1067}{\csromansymbolfootnote{*}{ขุ 8, 10, 89 ปิฏฺเฐสุปิ.}ปุจฺฉามิ ตํ ภควา พฺรูหิ เม’ตํ, (อิจฺจายสฺมา โธตโก,)}

\prose{วาจา’ภิกงฺขามิ มเหสิ ตุยฺหํ.}

\vspace{\tipitakaparskip}
\csromancenter{ตว สุตฺวาน นิคฺโฆสํ, สิกฺเข นิพฺพานมตฺตโน. (1)}
\proseitem{1068}{เตนหา’ตปฺปํ กโรหิ, (โธตกาติ ภควา,)}

\prose{อิเธว นิปโก สโต.}

\prose{อิโต สุตฺวาน นิคฺโฆสํ, สิกฺเข นิพฺพานมตฺตโน. (2)}

\chapterheadpageread{5. ปารายนวคฺค}
\proseitem{1069}{\csromanfolio{439}ปสฺสามหํ เทวมนุสฺสโลเก,}

\prose{อกิญฺจนํ พฺราหฺมณมิริยมานํ.}

\csromangathameasure{ตํ ตํ นมสฺสามิ สมนฺตจกฺขุ,}
\csromangathagroup{\csromangathastanza{ตํ ตํ นมสฺสามิ สมนฺตจกฺขุ,}}

\csromanheader{2}{ปมุญฺจ มํ สกฺก กถํกถาหิ. (3)}
\proseitem{1070}{นาหํ สหิสฺสามิ\footnote{สมิสฺสามิ (สฺยา), คมิสฺสามิ (สี), สมีหามิ (อิ)} ปโมจนาย,}

\prose{กถํกถิํ โธตก กญฺจิ โลเก.}

\csromangathameasure{ธมฺมญฺจ เสฏฺฐํ อภิชานมาโน,}
\csromangathagroup{\csromangathastanza{ธมฺมญฺจ เสฏฺฐํ อภิชานมาโน\footnote{อาชานมาโน (สี, สฺยา, อิ)},}}

\csromanheader{2}{เอวํ ตุวํ โอฆมิมํ ตเรสิ. (4)}
\proseitem{1071}{อนุสาส พฺรหฺเม กรุณายมาโน,}

\prose{วิเวกธมฺมํ ยมหํ วิชญฺญํ.}

\csromangathameasure{ยถาหํ อากาโสว อพฺยาปชฺชมาโน,}
\csromangathagroup{\csromangathastanza{ยถาหํ อากาโสว อพฺยาปชฺชมาโน,}}

\csromanheader{2}{อิเธว สนฺโต อสิโต จเรยฺยํ. (5)}
\proseitem{1072}{\csromansymbolfootnote{*}{ขุ 10. 143 ปิฏฺเฐปิ.}กิตฺตยิสฺสามิ เต สนฺติํ, (โธตกาติ ภควา,)}

\prose{ทิฏฺเฐ ธมฺเม อนีติหํ.}

\prose{ยํ วิทิตฺวา สโต จรํ, ตเร โลเก วิสตฺติกํ. (6)}

\proseitem{1073}{ตญฺจาหํ อภินนฺทามิ, มเหสิ สนฺติ’มุตฺตมํ.}

\prose{ยํ วิทิตฺวา สโต จรํ, ตเร โลเก วิสตฺติกํ. (7)}

\proseitem{1074}{ยํ กิญฺจิ สมฺปชานาสิ, (โธตกาติ ภควา,)}

\prose{อุทฺธํ อโธ ติริยญฺจาปิ มชฺเฌ.}

\csromangathameasure{เอตํ วิทิตฺวา สงฺโคติ โลเก, \\ ภวาภวาย มากาสิ ตณฺหนฺติ. (8) โธตกมาณวปุจฺฉา ปญฺจมี.}
\csromangathagroup{\csromangathastanza{เอตํ วิทิตฺวา สงฺโคติ โลเก, \\ ภวาภวาย มากาสิ ตณฺหนฺติ. (8) โธตกมาณวปุจฺฉา ปญฺจมี.}}

\csromanheader{2}{6. อุปสีวมาณวปุจฺฉา}
\proseitem{1075}{\csromansymbolfootnote{+}{ขุ 8. 11, 101 ปิฏฺเฐสุปิ.}เอโก อหํ สกฺก มหนฺตโมฆํ, (อิจฺจายสฺมา อุปสีโว,)}

\prose{อนิสฺสิโต โน วิสหามิ ตาริตุํ.}

\csromangathameasure{อารมฺมณํ พฺรูหิ สมนฺตจกฺขุ,}
\csromangathagroup{\csromangathastanza{อารมฺมณํ พฺรูหิ สมนฺตจกฺขุ,}}

\csromanheader{2}{ยํ นิสฺสิโต โอฆมิมํ ตเรยฺยํ. (1)}
\gambhira{สุตฺตนิปาตปาฬิ}
\proseitem{1076}{\csromanfolio{440}อากิญฺจญฺญํ เปกฺขมาโน สติมา, (อุปสีวาติ ภควา,)}

\prose{นตฺถีติ นิสฺสาย ตรสฺสุ โอฆํ.}

\prose{กาเม ปหาย วิรโต กถาหิ,}

\csromanheader{2}{ตณฺหกฺขยํ นตฺตมหา’ภิปสฺส\footnote{รตฺตมหาภิปสฺส (สฺยา), รตฺตมหํ วิปสฺส (ก) 2. ธิมุตฺโต (ก)}. (2)}
\proseitem{1077}{สพฺเพสุ กาเมสุ โย วีตราโค, (อิจฺจายสฺมา อุปสีโว,)}

\prose{อากิญฺจญฺญํ นิสฺสิโต หิตฺวา มญฺญํ.}

\prose{สญฺญาวิโมกฺเข ปรเม วิมุตฺโต2,}

\csromanheader{2}{ติฏฺเฐ นุ โส ตตฺถ อนานุยายี\footnote{อนานุวายี (สฺยา, ก)}. (3)}
\proseitem{1078}{สพฺเพสุ กาเมสุ โย วีตราโค, (อุปสีวาติ ภควา,)}

\prose{อากิญฺจญฺญํ นิสฺสิโต หิตฺวา มญฺญํ.}

\prose{สญฺญาวิโมกฺเข ปรเม วิมุตฺโต,}

\csromanheader{2}{ติฏฺเฐยฺย โส ตตฺถ อนานุยายี. (4)}
\proseitem{1079}{ติฏฺเฐ เจ โส ตตฺถ อนานุยายี,}

\prose{ปูคมฺปิ วสฺสานํ สมนฺตจกฺขุ.}

\prose{ตตฺเถว โส สีติสิยา วิมุตฺโต,}

\csromanheader{2}{จเวถ วิญฺญาณํ ตถาวิธสฺส. (5)}
\proseitem{1080}{อจฺจี ยถา วาตเวเคน ขิตฺตา\footnote{ขิตฺตํ (สฺยา), ขิตฺโต (อิ)}, (อุปสีวาติ ภควา,)}

\prose{อตฺถํ ปเลติ น อุเปติ สงฺขํ.}

\prose{เอวํ มุนี นามกายา วิมุตฺโต,}

\csromanheader{2}{อตฺถํ ปเลติ น อุเปติ สงฺขํ. (6)}
\proseitem{1081}{อตฺถงฺคโต โส อุท วา โส นตฺถิ,}

\prose{อุทาหุ เว สสฺสติยา อโรโค.}

\prose{ตํ เม มุนี สาธุ วิยากโรหิ,}

\csromanheader{2}{ตถา หิ เต วิทิโต เอส ธมฺโม. (7)}
\chapterheadpageread{5. ปารายนวคฺค}
\proseitem{1082}{\csromanfolio{441}อตฺถงฺคตสฺส น ปมาณ’มตฺถิ, (อุปสีวาติ ภควา,)}

\prose{เยน นํ วชฺชุํ ตํ ตสฺส นตฺถิ.}

\csromangathameasure{สพฺเพสุ ธมฺเมสุ สโมหเตสุ, \\ สมูหตา วาทปถาปิ สพฺเพติ. (8) อุปสีวมาณวปุจฺฉา ฉฏฺฐี.}
\csromangathagroup{\csromangathastanza{สพฺเพสุ ธมฺเมสุ สโมหเตสุ, \\ สมูหตา วาทปถาปิ สพฺเพติ. (8) อุปสีวมาณวปุจฺฉา ฉฏฺฐี.}}

\csromanheader{2}{7. นนฺทมาณวปุจฺฉา}
\proseitem{1083}{\csromansymbolfootnote{*}{ขุ 8. 13, 112 ปิฏฺเฐสุปิ.}“สนฺติ โลเก มุนโย” (อิจฺจายสฺมา นนฺโท,)}

\prose{ชนา วทนฺติ ตยิทํ กถํสุ.}

\csromangathameasure{ญาณูปปนฺนํ โน มุนิํ วทนฺติ,}
\csromangathagroup{\csromangathastanza{ญาณูปปนฺนํ โน มุนิํ\footnote{มุนิ โน (สฺยา, ก)} วทนฺติ,}}

\csromanheader{2}{อุทาหุ เว ชีวิเตนูปปนฺนํ. (1)}
\proseitem{1084}{น ทิฏฺฐิยา น สุติยา น ญาเณน, ( )\footnote{(น สีลพฺพเตน) กตฺถจิ โปตฺถเก.}}

\prose{มุนีธ นนฺท กุสลา วทนฺติ.}

\csromangathameasure{วิเสนิกตฺวา อนีฆา นิราสา,}
\csromangathagroup{\csromangathastanza{วิเสนิกตฺวา อนีฆา นิราสา,}}

\csromanheader{2}{จรนฺติ เย เต มุนโยติ พฺรูมิ. (2)}
\proseitem{1085}{เย เกจิเม สมณพฺราหฺมณาเส, (อิจฺจายสฺมา นนฺโท,)}

\prose{ทิฏฺฐสฺสุเตนาปิ\footnote{ทิฏฺเฐน สุเตนาปิ (สี), ทิฏฺเฐ สุเตนาปิ (สฺยา, อิ, ก)} วทนฺติ สุทฺธิํ.}

\csromangathameasure{สีลพฺพเตนาปิ วทนฺติ สุทฺธิํ, \\ อเนกรูเปน วทนฺติ สุทฺธิํ. \\ กจฺจิสฺสุ เต ภควา ตตฺถ ยตา จรนฺตา, \\ อตารุ ชาติญฺจ ชรญฺจ มาริส.}
\csromangathagroup{\csromangathastanza{สีลพฺพเตนาปิ วทนฺติ สุทฺธิํ, \\ อเนกรูเปน วทนฺติ สุทฺธิํ. \\ กจฺจิสฺสุ เต ภควา ตตฺถ ยตา จรนฺตา, \\ อตารุ ชาติญฺจ ชรญฺจ มาริส.}}

\csromanheader{2}{ปุจฺฉามิ ตํ ภควา พฺรูหิ เม’ตํ. (3)}
\gambhira{สุตฺตนิปาตปาฬิ}
\proseitem{1086}{\csromanfolio{442}เย เกจิเม สมณพฺราหฺมณาเส, (นนฺทาติ ภควา,)}

\prose{ทิฏฺฐสฺสุเตนาปิ วทนฺติ สุทฺธิํ.}

\csromangathameasure{สีลพฺพเตนาปิ วทนฺติ สุทฺธิํ, \\ อเนกรูเปน วทนฺติ สุทฺธิํ. \\ กิญฺจาปิ เต ตตฺถ ยตา จรนฺติ,}
\csromangathagroup{\csromangathastanza{สีลพฺพเตนาปิ วทนฺติ สุทฺธิํ, \\ อเนกรูเปน วทนฺติ สุทฺธิํ. \\ กิญฺจาปิ เต ตตฺถ ยตา จรนฺติ,}}

\csromanheader{2}{นาตริํสุ ชาติชรนฺติ พฺรูมิ. (4)}
\proseitem{1087}{เย เกจิเม สมณพฺราหฺมณาเส, (อิจฺจายสฺมา นนฺโท,)}

\prose{ทิฏฺฐสฺสุเตนาปิ วทนฺติ สุทฺธิํ.}

\csromangathameasure{สีลพฺพเตนาปิ วทนฺติ สุทฺธิํ, \\ อเนกรูเปน วทนฺติ สุทฺธิํ. \\ เต เจ มุนิ พฺรูสิ อโนฆติณฺเณ, \\ อถ โก จรหิ เทวมนุสฺสโลเก. \\ อตาริ ชาติญฺจ ชรญฺจ มาริส,}
\csromangathagroup{\csromangathastanza{สีลพฺพเตนาปิ วทนฺติ สุทฺธิํ, \\ อเนกรูเปน วทนฺติ สุทฺธิํ. \\ เต เจ มุนิ\footnote{สเจ มุนิ (สี)} พฺรูสิ อโนฆติณฺเณ, \\ อถ โก จรหิ เทวมนุสฺสโลเก.}\csromangathastanza{อตาริ ชาติญฺจ ชรญฺจ มาริส,}}

\csromanheader{2}{ปุจฺฉามิ ตํ ภควา พฺรูหิ เม’ตํ. (5)}
\proseitem{1088}{นาหํ สพฺเพ สมณพฺราหฺมณาเส, (นนฺทาติ ภควา,)}

\prose{ชาติชราย นิวุตาติ พฺรูมิ.}

\csromangathameasure{เย สีธ ทิฏฺฐํ ว สุตํ มุตํ วา, \\ สีลพฺพตํ วาปิ ปหาย สพฺพํ. \\ อเนกรูปมฺปิ ปหาย สพฺพํ, \\ ตณฺหํ ปริญฺญาย อนาสวาเส.}
\csromangathagroup{\csromangathastanza{เย สีธ ทิฏฺฐํ ว สุตํ มุตํ วา, \\ สีลพฺพตํ วาปิ ปหาย สพฺพํ. \\ อเนกรูปมฺปิ ปหาย สพฺพํ, \\ ตณฺหํ ปริญฺญาย อนาสวาเส.}}

\csromanheader{2}{เต เว นรา โอฆติณฺณาติ พฺรูมิ. (6)}
\proseitem{1089}{เอตา’ภินนฺทามิ วโจ มเหสิโน,}

\prose{สุกิตฺถิตํ โคตม’นูปธีกํ.}

\csromangathameasure{เย สีธ ทิฏฺฐํ ว สุตํ มุตํ วา, \\ สีลพฺพตํ วาปิ ปหาย สพฺพํ. \\ อเนกรูปมฺปิ ปหาย สพฺพํ, \\ ตณฺหํ ปริญฺญาย อนาสวาเส. \\ อหมฺปิ เต โอฆติณฺณาติ พฺรูมีติ. (7) นนฺทมาณวปุจฺฉา สตฺตมา.}
\csromangathagroup{\csromangathastanza{เย สีธ ทิฏฺฐํ ว สุตํ มุตํ วา, \\ สีลพฺพตํ วาปิ ปหาย สพฺพํ. \\ อเนกรูปมฺปิ ปหาย สพฺพํ, \\ ตณฺหํ ปริญฺญาย อนาสวาเส.}\csromangathastanza{อหมฺปิ เต โอฆติณฺณาติ พฺรูมีติ. (7) นนฺทมาณวปุจฺฉา สตฺตมา.}}

\csromanheader{2}{5. ปารายนวคฺค \textbf{8. เหมกมาณวปุจฺฉา}}
\proseitem{1090}{\csromanfolio{443}\csromansymbolfootnote{*}{ขุ 8. 14, 125 ปิฏฺเฐสุปิ.}เย เม ปุพฺเพ วิยากํสุ, (อิจฺจายสฺมา เหมโก,)}

\prose{หุรํ โคตมสาสนา.}

\csromangathameasure{“อิจฺจาสิ อิติ ภวิสฺสติ”, สพฺพํ ตํ อิติหีติหํ.}
\csromangathagroup{\csromangathastanza{“อิจฺจาสิ อิติ ภวิสฺสติ”, สพฺพํ ตํ อิติหีติหํ.}}

\prose{สพฺพํ ตํ ตกฺกวฑฺฒนํ, นาหํ ตตฺถ อภิรมิํ. (1)}

\proseitem{1091}{ตฺวญฺจ เม ธมฺมมกฺขาหิ, ตณฺหานิคฺฆาตนํ มุนิ.}

\prose{ยํ วิทิตฺวา สโต จรํ, ตเร โลเก วิสตฺติกํ. (2)}

\proseitem{1092}{อิธ ทิฏฺฐิสุตมุตวิญฺญาเตสุ, ปิยรูเปสุ เหมก.}

\prose{ฉนฺทราควิโนทนํ, นิพฺพานปท’มจฺจุตํ. (3)}

\proseitem{1093}{เอตทญฺญาย เย สตา, ทิฏฺฐธมฺมา’ภินิพฺพุตา.}

\csromangathameasure{อุปสนฺตา จ เต สทา, ติณฺณา โลเก วิสตฺติกนฺติ. (4) เหมกมาณวปุจฺฉา อฏฺฐมา.}
\csromangathagroup{\csromangathastanza{อุปสนฺตา จ เต สทา, ติณฺณา โลเก วิสตฺติกนฺติ. (4) เหมกมาณวปุจฺฉา อฏฺฐมา.}}

\csromanheader{2}{9. โตเทยฺยมาณวปุจฺฉา}
\proseitem{1094}{\csromansymbolfootnote{+}{ขุ 8. 14, 130 ปิฏฺเฐสุปิ.}ยสฺมิํ กามา น วสนฺติ, (อิจฺจายสฺมา โตเทยฺโย,)}

\prose{ตณฺหา ยสฺส น วิชฺชติ.}

\csromangathameasure{กถํกถา จ โย ติณฺโณ,}
\csromangathagroup{\csromangathastanza{กถํกถา จ โย ติณฺโณ,}}

\csromanheader{2}{วิโมกฺโข ตสฺส กีทิโส. (1)}
\proseitem{1095}{ยสฺมิํ กามา น วสนฺติ, (โตเทยฺยาติ ภควา,)}

\prose{ตณฺหา ยสฺส น วิชฺชติ.}

\csromangathameasure{กถํกถา จ โย ติณฺโณ,}
\csromangathagroup{\csromangathastanza{กถํกถา จ โย ติณฺโณ,}}

\csromanheader{2}{วิโมกฺโข ตสฺส นา ปโร. (2)}
\proseitem{1096}{นิราสโส โส อุท อาสสาโน,}

\prose{ปญฺญาณวา โส อุท ปญฺญกปฺปี.}

\csromangathameasure{มุนิํ อหํ สกฺก ยถา วิชญฺญํ,}
\csromangathagroup{\csromangathastanza{มุนิํ อหํ สกฺก ยถา วิชญฺญํ,}}

\csromanheader{2}{ตํ เม วิยาจิกฺข สมนฺตจกฺขุ. (3)}
\gambhira{สุตฺตนิปาตปาฬิ}
\proseitem{1097}{\csromanfolio{444}นิราสโส โส น จ อาสสาโน,}

\prose{ปญฺญาณวา โส น จ ปญฺญกปฺปี.}

\csromangathameasure{เอวมฺปิ โตเทยฺย มุนิํ วิชาน, \\ อกิญฺจนํ กามภเว อสตฺตนฺติ. (4) โตเทยฺยมาณวปุจฺฉา นวมา.}
\csromangathagroup{\csromangathastanza{เอวมฺปิ โตเทยฺย มุนิํ วิชาน, \\ อกิญฺจนํ กามภเว อสตฺตนฺติ. (4) โตเทยฺยมาณวปุจฺฉา นวมา.}}

\csromanheader{2}{10. กปฺปมาณวปุจฺฉา}
\proseitem{1098}{\csromansymbolfootnote{*}{ขุ 8. 15, 134 ปิฏฺเฐสุปิ.}มชฺเฌ สรสฺมิํ ติฏฺฐตํ, (อิจฺจายสฺมา กปฺโป,)}

\prose{โอเฆ ชาเต มหพฺภเย.}

\csromanmatikaanchor{mtk.294}
\csromantocmarkref{subsection}{1. อชิตมาณวปุจฺฉา}{mtk.294}
\bookmark[dest={mtk.294},level=2]{1. อชิตมาณวปุจฺฉา}
\csromangathameasure{ชรามจฺจุปเรตานํ, ทีปํ ปพฺรูหิ มาริส.}
\csromangathagroup{\csromangathastanza{ชรามจฺจุปเรตานํ, ทีปํ ปพฺรูหิ มาริส.}}

\vspace{\tipitakaparskip}
\csromancenter{ตฺวญฺจ เม ทีปมกฺขาหิ, ยถายิทํ นาปรํ สิยา. (1)}
\proseitem{1099}{มชฺเฌ สรสฺมิํ ติฏฺฐตํ, (กปฺปาติ ภควา,)}

\prose{โอเฆ ชาเต มหพฺภเย.}

\prose{ชรามจฺจุปเรตานํ, ทีปํ ปพฺรูมิ กปฺป เต. (2)}

\proseitem{1100}{อกิญฺจนํ อนาทานํ, เอตํ ทีปํ อนาปรํ.}

\vspace{\tipitakaparskip}
\csromancenter{นิพฺพานํ อิติ\footnote{นิพฺพานมีติ (สี)} นํ พฺรูมิ, ชรามจฺจุปริกฺขยํ. (3)}
\proseitem{1101}{เอตทญฺญาย เย สตา, ทิฏฺฐธมฺมาภินิพฺพุตา.}

\csromangathameasure{น เต มารวสานุคา, น เต มารสฺส ปทฺธคูติ. (4) กปฺปมาณวปุจฺฉา ทสมา.}
\csromangathagroup{\csromangathastanza{น เต มารวสานุคา, น เต มารสฺส ปทฺธคูติ\footnote{ปฏฺฐคูติ (สฺยา, ก)}. (4) กปฺปมาณวปุจฺฉา ทสมา.}}

\csromanheader{2}{11. ชตุกณฺณิมาณวปุจฺฉา}
\proseitem{1102}{\csromansymbolfootnote{+}{ขุ 8. 15, 141 ปิฏฺเฐสุปิ.}สุตฺวานหํ วีรมกามกามิํ, (อิจฺจายสฺมา ชตุกณฺณิ,)}

\prose{โอฆาติคํ ปุฏฺฐุมกามมาคมํ.}

\csromangathameasure{สนฺติปทํ พฺรูหิ สหชเนตฺต,}
\csromangathagroup{\csromangathastanza{สนฺติปทํ พฺรูหิ สหชเนตฺต,}}

\csromanheader{2}{ยถาตจฺฉํ ภควา พฺรูหิ เม’ตํ. (1)}
\chapterheadpageread{5. ปารายนวคฺค}
\proseitem{1103}{\csromanfolio{445}ภควา หิ กาเม อภิภุยฺย อิริยติ,}

\prose{อาทิจฺโจว ปถวิํ เตชี เตชสา.}

\csromangathameasure{ปริตฺตปญฺญสฺส เม ภูริปญฺญ, \\ อาจิกฺข ธมฺมํ ยมหํ วิชญฺญํ.}
\csromangathagroup{\csromangathastanza{ปริตฺตปญฺญสฺส เม ภูริปญฺญ, \\ อาจิกฺข ธมฺมํ ยมหํ วิชญฺญํ.}}

\csromanheader{2}{ชาติชราย อิธ วิปฺปหานํ. (2)}
\proseitem{1104}{กาเมสุ วินย เคธํ, (ชตุกณฺณีติ ภควา,)}

\prose{เนกฺขมฺมํ ทฏฺฐุ เขมโต.}

\vspace{\tipitakaparskip}
\csromancenter{อุคฺคหีตํ นิรตฺตํ วา, มา เต วิชฺชิตฺถ กิญฺจนํ. (3)}
\proseitem{1105}{ยํ ปุพฺเพ ตํ วิโสเสหิ, ปจฺฉา เต มาหุ กิญฺจนํ.}

\csromanmatikaanchor{mtk.295}
\csromantocmarkref{subsection}{2. ติสฺสเมตฺเตยฺยมาณวปุจฺฉา}{mtk.295}
\bookmark[dest={mtk.295},level=2]{2. ติสฺสเมตฺเตยฺยมาณวปุจฺฉา}
\prose{มชฺเฌ เจ โน คเหสฺสสิ, อุปสนฺโต จริสฺสสิ. (4)}

\proseitem{1106}{สพฺพโส นามรูปสฺมิํ, วีตเคธสฺส พฺราหฺมณ.}

\csromangathameasure{อาสวาสฺส น วิชฺชนฺติ, เยหิ มจฺจุวสํ วเชติ. (5) ชตุกณฺณิมาณวปุจฺฉา เอกาทสมา.}
\csromangathagroup{\csromangathastanza{อาสวาสฺส น วิชฺชนฺติ, เยหิ มจฺจุวสํ วเชติ. (5) ชตุกณฺณิมาณวปุจฺฉา เอกาทสมา.}}

\csromanheader{2}{12. ภทฺราวุธมาณวปุจฺฉา}
\proseitem{1107}{\csromansymbolfootnote{*}{ขุ 8. 17, 148 ปิฏฺเฐสุปิ.}โอกญฺชหํ ตณฺหจฺฉิทํ อเนชํ, (อิจฺจายสฺมา ภทฺราวุโธ,)}

\prose{นนฺทิญฺชหํ โอฆติณฺณํ วิมุตฺตํ.}

\csromangathameasure{กปฺปญฺชหํ อภิยาเจ สุเมธํ,}
\csromangathagroup{\csromangathastanza{กปฺปญฺชหํ อภิยาเจ สุเมธํ,}}

\csromanheader{2}{สุตฺวาน นาคสฺส อปนมิสฺสนฺติ อิโต. (1)}
\proseitem{1108}{นานาชนา ชนปเทหิ สงฺคตา,}

\prose{ตว วีร วากฺยํ อภิกงฺขมานา.}

\csromanmatikaanchor{mtk.296}
\csromantocmarkref{subsection}{3. ปุณฺณกมาณวปุจฺฉา}{mtk.296}
\bookmark[dest={mtk.296},level=2]{3. ปุณฺณกมาณวปุจฺฉา}
\csromangathameasure{เตสํ ตุวํ สาธุ วิยากโรหิ,}
\csromangathagroup{\csromangathastanza{เตสํ ตุวํ สาธุ วิยากโรหิ,}}

\csromanheader{2}{ตถา หิ เต วิทิโต เอส ธมฺโม. (2)}
\proseitem{1109}{อาทานตณฺหํ วินเยถ สพฺพํ, (ภทฺราวุธาติ ภควา,)}

\prose{อุทฺธํ อโธ ติริยญฺจาปิ มชฺเฌ.}

\csromangathameasure{ยํ ยญฺหิ โลกสฺมิมุปาทิยนฺติ,}
\csromangathagroup{\csromangathastanza{ยํ ยญฺหิ โลกสฺมิมุปาทิยนฺติ,}}

\csromanheader{2}{เตเนว มาโร อเนฺวติ ชนฺตุํ. (3)}
\gambhira{สุตฺตนิปาตปาฬิ}
\proseitem{1110}{\csromanfolio{446}ตสฺมา ปชานํ น อุปาทิเยถ,}

\prose{ภิกฺขุ สโต กิญฺจนํ สพฺพโลเก.}

\csromangathameasure{อาทานสตฺเต อิติ เปกฺขมาโน, \\ ปชํ อิมํ มจฺจุเธยฺเย วิสตฺตนฺติ. (4) ภทฺราวุธมาณวปุจฺฉา ทฺวาทสมา.}
\csromangathagroup{\csromangathastanza{อาทานสตฺเต อิติ เปกฺขมาโน, \\ ปชํ อิมํ มจฺจุเธยฺเย วิสตฺตนฺติ. (4) ภทฺราวุธมาณวปุจฺฉา ทฺวาทสมา.}}

\csromanheader{2}{13. อุทยมาณวปุจฺฉา}
\proseitem{1111}{\csromansymbolfootnote{*}{ขุ 8. 17, 154 ปิฏฺเฐสุปิ.}ฌายิํ วิรชมาสีนํ, (อิจฺจายสฺมา อุทโย,)}

\prose{กตกิจฺจํ อนาสวํ.}

\csromangathameasure{ปารคุํ สพฺพธมฺมานํ, อตฺถิ ปเญฺหน อาคมํ.}
\csromangathagroup{\csromangathastanza{ปารคุํ สพฺพธมฺมานํ, อตฺถิ ปเญฺหน อาคมํ.}}

\prose{อญฺญาวิโมกฺขํ ปพฺรูหิ, อวิชฺชาย ปเภทนํ. (1)}

\proseitem{1112}{ปหานํ กามจฺฉนฺทานํ, (อุทยาติ ภควา,)}

\prose{โทมนสฺสาน จูภยํ.}

\vspace{\tipitakaparskip}
\csromancenter{ถินสฺส จ ปนูทนํ, กุกฺกุจฺจานํ นิวารณํ. (2)}
\csromancenter{อุเปกฺขาสติสํสุทฺธํ, ธมฺมตกฺกปุเรชวํ.}
\csromancenter{อญฺญาวิโมกฺขํ ปพฺรูมิ, อวิชฺชาย ปเภทนํ. (3)}
\proseitem{1114}{กิํสุ สํโยชโน โลโก, กิํสุ ตสฺส วิจารณํ.}

\prose{กิสฺสสฺส วิปฺปหาเนน, นิพฺพานํ อิติ วุจฺจติ. (4)}

\proseitem{1115}{นนฺทิสํโยชโน โลโก, วิตกฺกสฺส วิจารณํ.}

\prose{ตณฺหาย วิปฺปหาเนน, นิพฺพานํ อิติ วุจฺจติ. (5)}

\proseitem{1116}{กถํ สตสฺส จรโต, วิญฺญาณํ อุปรุชฺฌติ.}

\csromanmatikaanchor{mtk.297}
\csromantocmarkref{subsection}{4. เมตฺตคูมาณวปุจฺฉา}{mtk.297}
\bookmark[dest={mtk.297},level=2]{4. เมตฺตคูมาณวปุจฺฉา}
\prose{ภควนฺตํ ปุฏฺฐุมาคมฺม, ตํ สุโณม วโจ ตว. (6)}

\proseitem{1117}{อชฺฌตฺตญฺจ พหิทฺธา จ, เวทนํ นาภินนฺทโต.}

\csromangathameasure{เอวํ สตสฺส จรโต, วิญฺญาณํ อุปรุชฺฌตีติ. (7) อุทยมาณวปุจฺฉา เตรสมา.}
\csromangathagroup{\csromangathastanza{เอวํ สตสฺส จรโต, วิญฺญาณํ อุปรุชฺฌตีติ. (7) อุทยมาณวปุจฺฉา เตรสมา.}}

\csromanheader{2}{5. ปารายนวคฺค \textbf{14. โปสาลมาณวปุจฺฉา}}
\proseitem{1118}{\csromanfolio{447}\csromansymbolfootnote{*}{ขุ 8. 18, 163 ปิฏฺเฐสุปิ.}โย อตีตํ อาทิสติ, (อิจฺจายสฺมา โปสาโล,)}

\prose{อเนโช ฉินฺนสํสโย.}

\vspace{\tipitakaparskip}
\csromancenter{ปารคุํ สพฺพธมฺมานํ, อตฺถิ ปเญฺหน อาคมํ. (1)}
\proseitem{1119}{วิภูตรูปสญฺญิสฺส, สพฺพกายปฺปหายิโน.}

\csromangathameasure{อชฺฌตฺตญฺจ พหิทฺธา จ, นตฺถิ กิญฺจีติ ปสฺสโต.}
\csromangathagroup{\csromangathastanza{อชฺฌตฺตญฺจ พหิทฺธา จ, นตฺถิ กิญฺจีติ ปสฺสโต.}}

\vspace{\tipitakaparskip}
\csromancenter{ญาณํ สกฺกา’นุปุจฺฉามิ, กถํ เนยฺโย ตถาวิโธ. (2)}
\proseitem{1120}{วิญฺญาณฏฺฐิติโย สพฺพา, (โปสาลาติ ภควา,)}

\prose{อภิชานํ ตถาคโต.}

\prose{ติฏฺฐนฺตเมนํ ชานาติ, วิมุตฺตํ ตปฺปรายณํ. (3)}

\proseitem{1121}{อากิญฺจญฺญสมฺภวํ ญตฺวา, นนฺที สํโยชนํ อิติ.}

\csromangathameasure{เอวเมตํ อภิญฺญาย, ตโต ตตฺถ วิปสฺสติ. \\ เอตํ ญาณํ ตถํ ตสฺส, พฺราหฺมณสฺส วุสีมโตติ. (4) โปสาลมาณวปุจฺฉา จุทฺทสมา.}
\csromangathagroup{\csromangathastanza{เอวเมตํ อภิญฺญาย, ตโต ตตฺถ วิปสฺสติ. \\ เอตํ\footnote{เอวํ (สฺยา, ก)} ญาณํ ตถํ ตสฺส, พฺราหฺมณสฺส วุสีมโตติ. (4) โปสาลมาณวปุจฺฉา จุทฺทสมา.}}

\csromanheader{2}{15. โมฆราชมาณวปุจฺฉา}
\proseitem{1122}{\csromansymbolfootnote{+}{ขุ 8. 19, 172 ปิฏฺเฐสุปิ.}ทฺวาหํ สกฺกํ อปุจฺฉิสฺสํ, (อิจฺจายสฺมา โมฆราชา,)}

\prose{น เม พฺยากาสิ จกฺขุมา.}

\csromangathameasure{ยาวตติยญฺจ เทวีสิ,}
\csromangathagroup{\csromangathastanza{ยาวตติยญฺจ เทวีสิ,}}

\csromanheader{2}{พฺยากโรตีติ เม สุตํ. (1)}
\proseitem{1123}{อยํ โลโก ปโร โลโก,}

\prose{พฺรหฺมโลโก สเทวโก.}

\csromangathameasure{ทิฏฺฐิํ เต นาภิชานาติ,}
\csromangathagroup{\csromangathastanza{ทิฏฺฐิํ เต นาภิชานาติ,}}

\csromanheader{2}{โคตมสฺส ยสสฺสิโน. (2)}
\proseitem{1124}{เอวํ อภิกฺกนฺตทสฺสาวิํ, อตฺถิ ปเญฺหน อาคมํ.}

\vspace{\tipitakaparskip}
\csromancenter{กถํ โลกํ อเวกฺขนฺตํ, มจฺจุราชา น ปสฺสติ. (3)}
\gambhira{สุตฺตนิปาตปาฬิ}
\proseitem{1125}{\csromanfolio{448}\csromansymbolfootnote{*}{ขุ 10. 7 ปิฏฺเฐ เนตฺติยํปิ.}สุญฺญโต โลกํ อเวกฺขสฺสุ, โมฆราช สทา สโต.}

\csromangathameasure{อตฺตานุทิฏฺฐิํ อูหจฺจ, เอวํ มจฺจุตโร สิยา. \\ เอว โลกํ อเวกฺขนฺตํ, มจฺจุราชา น ปสฺสตีติ. (4) โมฆราชมาณวปุจฺฉา ปนฺนรสมา.}
\csromangathagroup{\csromangathastanza{อตฺตานุทิฏฺฐิํ อูหจฺจ, เอวํ มจฺจุตโร สิยา. \\ เอว โลกํ อเวกฺขนฺตํ, มจฺจุราชา น ปสฺสตีติ. (4) โมฆราชมาณวปุจฺฉา ปนฺนรสมา.}}

\csromanheader{2}{16. ปิงฺคิยมาณวปุจฺฉา}
\proseitem{1126}{\csromansymbolfootnote{+}{ขุ 8. 19, 188 ปิฏฺเฐสุปิ.}ชิณฺโณหมสฺมิ อพโล วีตวณฺโณ, (อิจฺจายสฺมา ปิงฺคิโย,)}

\prose{เนตฺตา น สุทฺธา สวนํ น ผาสุ.}

\csromangathameasure{มา’หํ นสฺสํ โมมุโห อนฺตราว, \\ อาจิกฺข ธมฺมํ ยมหํ วิชญฺญํ.}
\csromangathagroup{\csromangathastanza{มา’หํ นสฺสํ โมมุโห อนฺตราว, \\ อาจิกฺข ธมฺมํ ยมหํ วิชญฺญํ.}}

\csromanheader{2}{ชาติชราย อิธ วิปฺปหานํ. (1)}
\proseitem{1127}{ทิสฺวาน รูเปสุ วิหญฺญมาเน, (ปิงฺคิยาติ ภควา,)}

\prose{รุปฺปนฺติ รูเปสุ ชนา ปมตฺตา.}

\csromangathameasure{ตสฺมา ตุวํ ปิงฺคิย อปฺปมตฺโต,}
\csromangathagroup{\csromangathastanza{ตสฺมา ตุวํ ปิงฺคิย อปฺปมตฺโต,}}

\csromanheader{2}{ชหสฺสุ รูปํ อปุนพฺภวาย. (2)}
\proseitem{1128}{ทิสา จตสฺโส วิทิสา จตสฺโส,}

\prose{อุทฺธํ อโธ ทส ทิสา อิมาโย.}

\csromangathameasure{น ตุยฺหํ อทิฏฺฐํ อสุตํ อมุตํ, \\ อโถ อวิญฺญาตํ กิญฺจน’มตฺถิ โลเก. \\ อาจิกฺข ธมฺมํ ยมหํ วิชญฺญํ,}
\csromangathagroup{\csromangathastanza{น ตุยฺหํ อทิฏฺฐํ อสุตํ อมุตํ\footnote{อสุตํ อมุตํ วา (สี), อสุตามุตํ วา (สฺยา), อสุตํ’มุตํ วา (อิ)}, \\ อโถ อวิญฺญาตํ กิญฺจน’มตฺถิ\footnote{อสุตํ อมุตํ วา (สี), อสุตามุตํ วา (สฺยา), อสุตํ’มุตํ วา (อิ)} โลเก. \\ อาจิกฺข ธมฺมํ ยมหํ วิชญฺญํ,}}

\csromanheader{2}{ชาติชราย อิธ วิปฺปหานํ. (3)}
\proseitem{1129}{ตณฺหาธิปนฺเน มนุเช เปกฺขมาโน, (ปิงฺคิยาติ ภควา,)}

\prose{สนฺตาปชาเต ชรสา ปเรเต.}

\csromangathameasure{ตสฺมา ตุวํ ปิงฺคิย อปฺปมตฺโต, \\ ชหสฺสุ ตณฺหํ อปุนพฺภวายาติ. (4) ปิงฺคิยมาณวปุจฺฉา โสฬสมา.}
\csromangathagroup{\csromangathastanza{ตสฺมา ตุวํ ปิงฺคิย อปฺปมตฺโต, \\ ชหสฺสุ ตณฺหํ อปุนพฺภวายาติ. (4) ปิงฺคิยมาณวปุจฺฉา โสฬสมา.}}

\csromanheader{3}{5. ปารายนวคฺค \textbf{ปารายนตฺถุติคาถา}}
\prose{\csromanfolio{449}\csromansymbolfootnote{*}{ขุ 8. 21, 195 ปิฏฺเฐสุปิ.}อิทมโวจ ภควา มคเธสุ วิหรนฺโต ปาสาณเก เจติเย, ปริจารกโสฬสานํ\footnote{ปริจารกโสฬสนฺนํ (สฺยา, ก)} พฺราหฺมณานํ อชฺฌิฏฺโฐ ปุฏฺโฐ ปุฏฺโฐ ปญฺหํ\footnote{ปเญฺห (สี, อิ)} พฺยากาสิ.\csromanspacer{} เอกเมกสฺส เจปิ ปญฺหสฺส อตฺถมญฺญาย ธมฺมมญฺญาย ธมฺมานุธมฺมํ ปฏิปชฺเชยฺย, คจฺเฉยฺเยว ชรามรณสฺส ปารํ, ปารงฺคมนียา อิเม ธมฺมาติ, ตสฺมา อิมสฺส ธมฺมปริยายสฺส ปารายนนฺเตว\footnote{ปารายณํเตฺวว (สี-ฏฺฐ)} อธิวจนํ.}

\proseitem{1130}{อชิโต ติสฺสเมตฺเตยฺโย, ปุณฺณโก อถ เมตฺตคู.}

\prose{โธตโก อุปสีโว จ, นนฺโท จ อถ เหมโก. (1)}

\csromanmatikaanchor{mtk.298}
\csromantocmarkref{subsection}{5. โธตกมาณวปุจฺฉา}{mtk.298}
\bookmark[dest={mtk.298},level=2]{5. โธตกมาณวปุจฺฉา}
\proseitem{1131}{โตเทยฺย-กปฺปา ทุภโย, ชตุกณฺณี จ ปณฺฑิโต.}

\csromangathameasure{ภทฺราวุโธ อุทโย จ, โปสาโล จาปิ พฺราหฺมโณ.}
\csromangathagroup{\csromangathastanza{ภทฺราวุโธ อุทโย จ, โปสาโล จาปิ พฺราหฺมโณ.}}

\prose{โมฆราชา จ เมธาวี, ปิงฺคิโย จ มหา-อิสิ. (2)}

\proseitem{1132}{เอเต พุทฺธํ อุปาคจฺฉุํ, สมฺปนฺนจรณํ อิสิํ.}

\vspace{\tipitakaparskip}
\csromancenter{ปุจฺฉนฺตา นิปุเณ ปเญฺห, พุทฺธเสฏฺฐํ อุปาคมุํ. (3)}
\csromancenter{เตสํ พุทฺโธ ปพฺยากาสิ, ปเญฺห ปุฏฺโฐ ยถาตถํ.}
\csromancenter{ปญฺหานํ เวยฺยากรเณน, โตเสสิ พฺราหฺมเณ มุนิ. (4)}
\proseitem{1134}{เต โตสิตา จกฺขุมตา, พุทฺเธนาทิจฺจพนฺธุนา.}

\prose{พฺรหฺมจริยมจริํสุ, วรปญฺญสฺส สนฺติเก. (5)}

\proseitem{1135}{เอกเมกสฺส ปญฺหสฺส, ยถา พุทฺเธน เทสิตํ.}

\prose{ตถา โย ปฏิปชฺเชยฺย, คจฺเฉ ปารํ อปารโต. (6)}

\proseitem{1136}{อปารา ปารํ คจฺเฉยฺย, ภาเวนฺโต มคฺคมุตฺตมํ.}

\vspace{\tipitakaparskip}
\csromancenter{มคฺโค โส ปารํ คมนาย, ตสฺมา ปารายนํ อิติ. (7)}
\csromanheader{3}{\textbf{ปารายนานุคีติคาถา}}
\proseitem{1137}{\csromansymbolfootnote{+}{ขุ 8. 202 ปิฏฺเฐปิ.}ปารายนมนุคายิสฺสํ, (อิจฺจายสฺมา ปิงฺคิโย,)}

\csromangathameasure{ยถา’ทฺทกฺขิ ตถา’กฺขาสิ, วิมโล ภูริเมธโส.}
\csromangathagroup{\csromangathastanza{ยถา’ทฺทกฺขิ ตถา’กฺขาสิ, วิมโล ภูริเมธโส.}}

\vspace{\tipitakaparskip}
\csromancenter{นิกฺกาโม นิพฺพโน\footnote{นิพฺพุโต (สฺยา)} นาโค, กิสฺส เหตุ มุสา ภเณ. (1)}
\gambhira{สุตฺตนิปาตปาฬิ}
\proseitem{1138}{\csromanfolio{450}ปหีนมลโมหสฺส, มานมกฺขปฺปหายิโน.}

\prose{หนฺทาหํ กิตฺตยิสฺสามิ, คิรํ วณฺณูปสญฺหิตํ. (2)}

\proseitem{1139}{ตโมนุโท พุทฺโธ สมนฺตจกฺขุ,}

\prose{โลกนฺตคู สพฺพภวาติวตฺโต.}

\prose{อนาสโว สพฺพทุกฺขปฺปหีโน,}

\csromanheader{2}{สจฺจวฺหโย พฺรหฺเม อุปาสิโต เม. (3)}
\proseitem{1140}{ทิโช ยถา กุพฺพนกํ ปหาย,}

\prose{พหุปฺผลํ กานนมาวเสยฺย.}

\prose{เอวมฺปหํ อปฺปทสฺเส ปหาย,}

\csromanheader{2}{มโหทธิํ หํโสริว อชฺฌปตฺโต. (4)}
\csromanmatikaanchor{mtk.299}
\csromantocmarkref{subsection}{6. อุปสีวมาณวปุจฺฉา}{mtk.299}
\bookmark[dest={mtk.299},level=2]{6. อุปสีวมาณวปุจฺฉา}
\proseitem{1141}{เยเม ปุพฺเพ วิยากํสุ,}

\csromangathameasure{หุรํ โคตมสาสนา. “อิจฺจาสิ อิติ ภวิสฺสติ”, \\ สพฺพํ ตํ อิติหิติหํ,}
\csromangathagroup{\csromangathastanza{หุรํ โคตมสาสนา. “อิจฺจาสิ อิติ ภวิสฺสติ”, \\ สพฺพํ ตํ อิติหิติหํ,}}

\csromanheader{2}{สพฺพํ ตํ ตกฺกวฑฺฒนํ. (5)}
\proseitem{1142}{เอโก ตมนุทาสิโน, ชุติมา โส ปภงฺกโร.}

\prose{โคตโม ภูริปญฺญาโณ, โคตโม ภูริเมธโส. (6)}

\proseitem{1143}{โย เม ธมฺมมเทเสสิ, สนฺทิฏฺฐิกมกาลิกํ.}

\prose{ตณฺหกฺขยมนีติกํ, ยสฺส นตฺถิ อุปมา กฺวจิ. (7)}

\proseitem{1144}{กิํนุ ตมฺหา วิปฺปวสสิ, มุหุตฺตมปิ ปิงฺคิย.}

\prose{โคตมา ภูริปญฺญาณา, โคตมา ภูริเมธสา. (8)}

\proseitem{1145}{โย เต ธมฺมมเทเสสิ, สนฺทิฏฺฐิกมกาลิกํ.}

\prose{ตณฺหกฺขยมนีติกํ, ยสฺส นตฺถิ อุปมา กฺวจิ. (9)}

\proseitem{1146}{นาหํ ตมฺหา วิปฺปวสามิ, มุหุตฺตมปิ พฺราหฺมณ.}

\prose{โคตมา ภูริปญฺญาณา, โคตมา ภูริเมธสา. (10)}

\proseitem{1147}{โย เม ธมฺมมเทเสสิ, สนฺทิฏฺฐิกมกาลิกํ.}

\prose{ตณฺหกฺขยมนีติกํ, ยสฺส นตฺถิ อุปมา กฺวจิ. (11)}

\chapterheadpageread{5. ปารายนวคฺค}
\proseitem{1148}{\csromanfolio{451}ปสฺสามิ นํ มนสา จกฺขุนาว,}

\prose{รตฺตินฺทิวํ พฺราหฺมณ อปฺปมตฺโต.}

\prose{นมสฺสมาโน วิวเสมิ รตฺติํ,}

\csromanheader{2}{เตเนว มญฺญามิ อวิปฺปวาสํ. (12)}
\proseitem{1149}{สทฺธา จ ปีติ จ มโน สติ จ,}

\prose{นา’เปนฺติ’เม โคตมสาสนมฺหา.}

\prose{ยํ ยํ ทิสํ วชติ ภูริปญฺโญ,}

\csromanheader{2}{ส เตน เตเนว นโตหมสฺมิ. (13)}
\proseitem{1150}{ชิณฺณสฺส เม ทุพฺพลถามกสฺส,}

\prose{เตเนว กาโย น ปเลติ ตตฺถ.}

\prose{สํกปฺปยนฺตาย\footnote{สํกปฺปยตฺตาย (สี)} วชามิ นิจฺจํ,}

\csromanheader{2}{มโน หิ เม พฺราหฺมณ เตน ยุตฺโต. (14)}
\proseitem{1151}{ปงฺเก สยาโน ปริผนฺทมาโน, ทีปา ทีปํ อุปปฺลวิํ\footnote{อุปลฺลวิํ (สฺยา, นิทฺเทส 23. 220 ปิฏฺเฐสุ.)}.}

\prose{อถทฺทสาสิํ สมฺพุทฺธํ, โอฆติณฺณมนาสวํ. (15)}

\proseitem{1152}{ยถา อหู วกฺกลิ มุตฺตสทฺโธ,}

\prose{ภทฺราวุโธ อาฬวิ โคตโม จ.}

\prose{เอวเมวํ ตฺวมฺปิ ปมุญฺจสฺสุ สทฺธํ,}

\csromanheader{2}{คมิสฺสสิ ตฺวํ ปิงฺคิย มจฺจุเธยฺยสฺส ปารํ\footnote{มจฺจุเธยฺยปารํ (สี)}. (16)}
\csromanmatikaanchor{mtk.300}
\csromantocmarkref{subsection}{7. นนฺทมาณวปุจฺฉา}{mtk.300}
\bookmark[dest={mtk.300},level=2]{7. นนฺทมาณวปุจฺฉา}
\proseitem{1153}{เอส ภิยฺโย ปสีทามิ, สุตฺวาน มุนิโน วโจ.}

\prose{วิวฏฺฏจฺฉโท สมฺพุทฺโธ, อขิโล ปฏิภานวา. (17)}

\proseitem{1154}{อธิเทเว อภิญฺญาย, สพฺพํ เวทิ วโรวรํ\footnote{ปโร วรํ (สี, สฺยา), ปโร ปรํ (นิทฺเทส)}.}

\prose{ปญฺหานนฺตกโร สตฺถา, กงฺขีนํ ปฏิชานตํ. (18)}

\proseitem{1155}{อสํหีรํ อสํกุปฺปํ, ยสฺส นตฺถิ อุปมา กฺวจิ.}

\prose{อทฺธา คมิสฺสามิ น เมตฺถ กงฺขา,}

\csromanheader{2}{เอวํ มํ ธาเรหิ อธิมุตฺตจิตฺตนฺติ. (19)}
\vspace{\tipitakaparskip}
\csromancenter{\textbf{ปารายนวคฺโค ปญฺจโม.}}
\gambhira{สุตฺตนิปาตปาฬิ}
\csromanheader{2}{\textbf{สุตฺตุทฺทานํ}}
\proseitem{1}{\csromanfolio{452}1 อุรโค ธนิโยปิ จ, ขคฺควิสาโณ กสิ จ.}

\csromangathameasure{จุนฺโท ภโว ปุนเทว, วสโล จ กรณียญฺจ. \\ เหมวโต อถ ยกฺโข, วิชยสุตฺตํ มุนิสุตฺตวรนฺติ.}
\csromangathagroup{\csromangathastanza{จุนฺโท ภโว ปุนเทว, วสโล จ กรณียญฺจ. \\ เหมวโต อถ ยกฺโข, วิชยสุตฺตํ มุนิสุตฺตวรนฺติ.}}

\proseitem{2}{ปฐมกฏฺฐวโร วรวคฺโค, ทฺวาทสสุตฺตธโร สุวิภตฺโต.}

\csromangathameasure{เทสิโต จกฺขุมตา วิมเลน, สุยฺยติ วคฺควโร อุรโคติ.}
\csromangathagroup{\csromangathastanza{เทสิโต จกฺขุมตา วิมเลน, สุยฺยติ วคฺควโร อุรโคติ.}}

\proseitem{3}{รตนามคนฺโธ หิริมงฺคลนาโม,}

\prose{สุจิโลมกปิโล จ พฺราหฺมณธมฺโม.}

\csromangathameasure{นาวา กิํสีล-อุฏฺฐหโน จ, \\ ราหุโล จ ปุนปิ วงฺคีโส.}
\csromangathagroup{\csromangathastanza{นาวา\footnote{นาถ (ก)} กิํสีล-อุฏฺฐหโน จ, \\ ราหุโล จ ปุนปิ วงฺคีโส.}}

\proseitem{4}{สมฺมาปริพฺพาชนีโยปิ เจตฺถ,}

\prose{ธมฺมิกสุตฺตวโร สุวิภตฺโต.}

\csromangathameasure{จุทฺทสสุตฺตธโร ทุติยมฺหิ, \\ จูฬกวคฺควโรติ ตมาหุ.}
\csromangathagroup{\csromangathastanza{จุทฺทสสุตฺตธโร ทุติยมฺหิ, \\ จูฬกวคฺควโรติ ตมาหุ.}}

\proseitem{5}{ปพฺพชฺชปธานสุภาสิตนาโม,}

\prose{ปูรฬาโส ปุนเทว มาโฆ จ.}

\csromangathameasure{สภิยํ เกณิยเมว สลฺลนาโม, \\ วาเสฏฺฐวโร กาลิโกปิ จ.}
\csromangathagroup{\csromangathastanza{สภิยํ เกณิยเมว สลฺลนาโม, \\ วาเสฏฺฐวโร กาลิโกปิ จ.}}

\proseitem{6}{นาลกสุตฺตวโร สุวิภตฺโต,}

\prose{ตํ อนุปสฺสี ตถา ปุนเทว.}

\csromangathameasure{ทฺวาทสสุตฺตธโร ตติยมฺหิ, \\ สุยฺยติ วคฺควโร มหานาโม.}
\csromangathagroup{\csromangathastanza{ทฺวาทสสุตฺตธโร ตติยมฺหิ, \\ สุยฺยติ วคฺควโร มหานาโม.}}

\proseitem{7}{กามคุหฏฺฐกทุฏฺฐกนามา,}

\prose{สุทฺธวโร ปรมฏฺฐกนาโม.}

\csromangathameasure{ชรา เมตฺติยวโร สุวิภตฺโต, \\ ปสูรมาคณฺฑิยา ปุราเภโท.}
\csromangathagroup{\csromangathastanza{ชรา เมตฺติยวโร สุวิภตฺโต, \\ ปสูรมาคณฺฑิยา ปุราเภโท.}}

\csromanheader{2}{สุตฺตุทฺทานคาถา}
\csromanhangingbegin
\hangingheaditem{8}{\csromanfolio{453}กลหวิวาโท อุโภ วิยุหา จ,}
\hangingbody{ตุวฏก-อตฺตทณฺฑสาริปุตฺตา.}
\hangingbody{โสฬสสุตฺตธโร จตุตฺถมฺหิ,}
\hangingbody{อฏฺฐกวคฺควโรติ ตมาหุ.}
\csromanhangingend

\csromanhangingbegin
\hangingheaditem{9}{มคเธ ชนปเท รมณีเย,}
\hangingbody{เทสวเร กตปุญฺญนิเวเส.}
\hangingbody{ปาสาณกเจติยวเร สุวิภตฺเต,}
\hangingbody{วสิ ภควา คณเสฏฺโฐ.}
\csromanhangingend

\csromanhangingbegin
\hangingheaditem{10}{อุภยวาสมาคติยมฺหิ\footnote{อุภยํ วา ปุณฺณสมาคตํ ยมฺหิ (สฺยา)},}
\hangingbody{ทฺวาทสโยชนิยา ปริสาย.}
\hangingbody{โสฬสพฺราหฺมณานํ กิร ปุฏฺโฐ,}
\hangingbody{ปุจฺฉาย โสฬสปญฺหกมฺมิยา.}
\hangingbody{นิปฺปกาสยิ ธมฺมมทาสิ.}
\csromanhangingend

\csromanhangingbegin
\hangingheaditem{11}{อตฺถปกาสกพฺยญฺชนปุณฺณํ,}
\hangingbody{ธมฺมมเทเสสิ ปรเขมชนิยํ2.}
\hangingbody{โลกหิตาย ชิโน ทฺวิปทคฺโค,}
\hangingbody{สุตฺตวรํ พหุธมฺมวิจิตฺรํ.}
\hangingbody{สพฺพกิเลสปโมจนเหตุํ,}
\hangingbody{เทสยิ สุตฺตวรํ ทฺวิปทคฺโค.}
\csromanhangingend

\csromanhangingbegin
\hangingheaditem{12}{พฺยญฺชนมตฺถปทํ สมยุตฺตํ\footnote{พฺยญฺชนมตฺถปทสมยุตฺตํ (สฺยา)},}
\hangingbody{อกฺขรสญฺญิต-โอปมคาฬฺหํ.}
\hangingbody{โลกวิจารณญาณปภคฺคํ,}
\hangingbody{เทสยิ สุตฺตวรํ ทฺวิปทคฺโค.}
\csromanhangingend

\csromanhangingbegin
\hangingheaditem{13}{ราคมเล อมลํ วิมลคฺคํ,}
\hangingbody{โทสมเล อมลํ วิมลคฺคํ.}
\hangingbody{โมหมเล อมลํ วิมลคฺคํ,}
\hangingbody{โลกวิจารณญาณปภคฺคํ.}
\hangingbody{เทสยิ สุตฺตวรํ ทฺวิปทคฺโค.}
\csromanhangingend

\gambhira{สุตฺตนิปาตปาฬิ}
\csromanhangingbegin
\hangingheaditem{14}{\csromanfolio{454}เกฺลสมเล อมลํ วิมลคฺคํ,}
\hangingbody{ทุจฺจริตมเล อมลํ วิมลคฺคํ.}
\hangingbody{โลกวิจารณญาณปภคฺคํ,}
\hangingbody{เทสยิ สุตฺตวรํ ทฺวิปทคฺโค.}
\csromanhangingend

\csromanhangingbegin
\hangingheaditem{15}{อาสวพนฺธนโยคากิเลสํ,}
\hangingbody{นีวรณานิ จ ตีณิ มลานิ.}
\hangingbody{ตสฺส กิเลสปโมจนเหตุํ,}
\hangingbody{เทสยิ สุตฺตวรํ ทฺวิปทคฺโค.}
\csromanhangingend

\csromanhangingbegin
\hangingheaditem{16}{นิมฺมลสพฺพกิเลสปนูทํ,}
\hangingbody{ราควิราคมเนชมโสกํ.}
\hangingbody{สนฺตปณีตสุทุทฺทสธมฺมํ,}
\hangingbody{เทสยิ สุตฺตวรํ ทฺวิปทคฺโค.}
\csromanhangingend

\csromanhangingbegin
\hangingheaditem{17}{ราคญฺจ โทสกมภญฺชิตสนฺตํ\footnote{โทสญฺจ ภญฺชิตสนฺตํ (สฺยา)},}
\hangingbody{โยนิจตุคฺคติปญฺจวิญฺญาณํ.}
\hangingbody{ตณฺหารตจฺฉทนตาณลตาปโมกฺขํ2,}
\hangingbody{เทสยิ สุตฺตวรํ ทฺวิปทคฺโค.}
\csromanhangingend

\csromanhangingbegin
\hangingheaditem{18}{คมฺภีรทุทฺทสสณฺหนิปุณํ,}
\hangingbody{ปณฺฑิตเวทนิยํ นิปุณตฺถํ.}
\hangingbody{โลกวิจารณญาณปภคฺคํ,}
\hangingbody{เทสยิ สุตฺตวรํ ทฺวิปทคฺโค.}
\csromanhangingend

\csromanmatikaanchor{mtk.301}
\csromantocmarkref{subsection}{8. เหมกมาณวปุจฺฉา}{mtk.301}
\bookmark[dest={mtk.301},level=2]{8. เหมกมาณวปุจฺฉา}
\csromanhangingbegin
\hangingheaditem{19}{นวงฺคกุสุมมาลคีเวยฺยํ,}
\hangingbody{อินฺทฺริยฌานวิโมกฺขวิภตฺตํ.}
\hangingbody{อฏฺฐงฺคมคฺคธรํ วรยานํ,}
\hangingbody{เทสยิ สุตฺตวรํ ทฺวิปทคฺโค.}
\csromanhangingend

\csromanheader{2}{สุตฺตุทฺทานคาถา}
\csromanhangingbegin
\hangingheaditem{20}{\csromanfolio{455}โสมุปมํ วิมลํ ปริสุทฺธํ,}
\hangingbody{อณฺณวมูปมรตนสุจิตฺตํ.}
\hangingbody{ปุปฺผสมํ รวิมูปมเตชํ,}
\hangingbody{เทสยิ สุตฺตวรํ ทฺวิปทคฺโค.}
\csromanhangingend

\csromanhangingbegin
\hangingheaditem{21}{เขมสิวํ สุขสีตลสนฺตํ,}
\hangingbody{มจฺจุตตาณปรํ ปรมตฺถํ.}
\hangingbody{ตสฺส สุนิพฺพุตทสฺสนเหตุํ,}
\hangingbody{เทสยิ สุตฺตวรํ ทฺวิปทคฺโคติ.}
\csromanhangingend

\gambhira{\textbf{สุตฺตนิปาตปาฬิ สมตฺตา.}}
\orphannote{ขุ 2. 275 ปิฏฺเฐปิ.}

\orphannote{ขุ 10. 140 ปิฏฺเฐปิ.}

\orphannote{ขุ 2. 260; ขุ 10. 140 ปิฏฺเฐสุปิ.}

\orphannote{อุปริ 382 ปิฏฺเฐปิ.}

\orphannote{อุปริ 89 ปิฏฺเฐ; ขุ 10. 29, 108 เนตฺติยมฺปิ.}

\orphannote{ขุ 2. 224, 334 ปิฏฺเฐสุปิ.}

\csromanmatikaanchor{mtk.302}
\csromantocmarkref{subsection}{9. โตเทยฺยมาณวปุจฺฉา}{mtk.302}
\bookmark[dest={mtk.302},level=2]{9. โตเทยฺยมาณวปุจฺฉา}
\orphannote{ขุ 10. 171 ปิฏฺเฐปิ.}

\orphannote{ขุ 10. 108, 173 ปิฏฺเฐสุปิ.}

\orphannote{ขุ 8. 74 ปิฏฺเฐปิ.}

\orphannote{อภิ 4. 433; ขุ 10. 162 ปิฏฺเฐสุปิ.}

\orphannote{ขุ 2. 317; ขุ 10. 7, 53 ปิฏฺเฐสุปิ.}

\orphannote{ขุ 10. 158 ปิฏฺเฐปิ.}

\orphannote{ขุ 10. 32 ปิฏฺเฐปิ.}

\orphannote{ขุ 10. 12 ปิฏฺเฐ เนตฺติยมฺปิ.}

\orphannote{อุปริ 130, 224, 382 ปิฏฺเฐสุปิ.}

\orphannote{วิ 1. 116 ปิฏฺเฐปิ.}

\orphannote{ขุ 10. 30, 109 ปิฏฺเฐสุปิ.}

\orphannote{ขุ 10. 37 ปิฏฺเฐปิ.}

\orphannote{ขุ 10. 31, 185 ปิฏฺเฐสุปิ.}

\orphannote{วิ 3. 12; อภิ 4. 218 ปิฏฺเฐสุปิ.}

\orphannote{อุปริ 376; ม 2. 410 ปิฏฺเฐสุปิ.}

\orphannote{อุปริ 376; ม 2. 410 ปิฏฺเฐสุปิ.}

\orphannote{อุปริ 377; ม 2. 411 ปิฏฺฐปิ.}

\csromanmatikaanchor{mtk.303}
\csromantocmarkref{subsection}{10. กปฺปมาณวปุจฺฉา}{mtk.303}
\bookmark[dest={mtk.303},level=2]{10. กปฺปมาณวปุจฺฉา}
\orphannote{ธมฺมปทสฺส วคฺคสฺสุทฺทานํ ยมกํ ปมาทํ จิตฺตํ, ปุปฺผํ พาลญฺจ ปณฺฑิตํ. รหนฺตํ สหสฺสํ ปาปํ, ทณฺฑํ ชรา อตฺตโลกํ. พุทฺธํ สุขํ ปิยํ โกธํ, มลํ ธมฺมฏฺฐมคฺคญฺจ. ปกิณฺณกํ นิรยํ นาคํ, ตณฺหา ภิกฺขุ จ พฺราหฺมโณ. \textbf{คาถายุทฺทานํ} ยมเก วีสคาถาโย, อปฺปมาทโลกมฺหิ จ. ปิเย ทฺวาทส คาถาโย, จิตฺเต ชรตฺเตกาทส. ปุปฺผพาลสหสฺสมฺหิ, พุทฺธ มคฺค ปกิณฺณเก. โสฬส ปณฺฑิเต โกเธ, นิรเย นาเค จตุทฺทส. อรหนฺเต ทสคฺคาถา, ปาปสุขมฺหิ เตรส. สตฺตรส ทณฺฑธมฺมฏฺเฐ, มลมฺหิ เอกวีสติ. ตณฺหาวคฺเค สตฺตพฺพีส, เตวีส ภิกฺขุวคฺคมฺหิ. พฺราหฺมเณ เอกตาลีส, จตุสฺสตา สเตวีส. (ก)}

\orphannote{วิ 3. 3 ปิฏฺเฐปิ.}

\orphannote{ขุ 10. 130 ปิฏฺเฐปิ.}

\orphannote{วิ 3. 4; อภิ 4. 163 ปิฏฺเฐสุปิ.}

\orphannote{ปฏิพทฺธจิตฺโต (สฺยา), ปฏิพนฺธรูโป (?)}

\orphannote{ทิสฺวา (สี)}

\orphannote{อุปสงฺกมิตฺวา (สพฺพตฺถ) อํ 3. 166 ปิฏฺเฐ ปสฺสิตพฺพํ.}

\orphannote{ขุ 10. 7 ปิฏฺเฐปิ.}

\orphannote{อุปริ 399 ปิฏฺฐาทีสุ.}

\orphannote{จาริยนฺติ (สฺยา)}

\orphannote{เอตมตฺถํ วิทิตฺวา (สี, ก)}

\orphannote{นตฺถิ สีหฬโปตฺถเก.}

\csromanmatikaanchor{mtk.304}
\csromantocmarkref{subsection}{11. ชตุกณฺณิมาณวปุจฺฉา}{mtk.304}
\bookmark[dest={mtk.304},level=2]{11. ชตุกณฺณิมาณวปุจฺฉา}
\orphannote{ขุ 10. 54, 123, 184 ปิฏฺเฐสุปิ.}

\orphannote{ขุ 10. 32, 107, 184 ปิฏฺเฐสุปิ.}

\orphannote{ขุ 10. 55, 180 ปิฏฺเฐสุปิ.}

\orphannote{“เอวํ ภนฺเต”ติ โข ปาฏลิคามิยา อุปาสกา ภควโต ปฏิสฺสุตฺวา (วิ 3. 323; ที 2. 73)}

\orphannote{อํ 1. 232 ปิฏฺเฐปิ.}

\orphannote{มุตีมา (ก) ขุ 8. 282 ปิฏฺเฐปิ.}

\orphannote{สิคาลา (สี, สฺยา, กํ, อิ)}

\orphannote{ที 1. 48 ปิฏฺฐาทีสุ วิตฺถาโร.}

\orphannote{ปกตํ (สฺยา)}

\orphannote{อิมา อุทฺทานคาถาโย สี-อิ-โปตฺถเกสุ น สนฺติ.}

